\documentclass[10pt]{amsart}
\usepackage{amsmath,amssymb,amsthm,mathtools,mathrsfs}
\usepackage[a4paper,hmargin=0.76in,vmargin=0.60in]{geometry}
\usepackage{booktabs}
\usepackage{array}
\usepackage{longtable}
\usepackage{tabularx}
\usepackage{xr-hyper}
\usepackage[hypertexnames=false]{hyperref}
\externaldocument[M:]{sign-radial-intersection-form-idele-class-group}
\hypersetup{hidelinks}
\usepackage[final,expansion=false]{microtype}
\usepackage{url}
\emergencystretch=2.5em

\newtheorem{theorem}{Theorem}[section]
\newtheorem{proposition}[theorem]{Proposition}
\newtheorem{lemma}[theorem]{Lemma}
\newtheorem{corollary}[theorem]{Corollary}
\theoremstyle{definition}
\newtheorem{definition}[theorem]{Definition}
\newtheorem{problem}[theorem]{Problem}
\theoremstyle{remark}
\newtheorem{remark}[theorem]{Remark}

\newcommand{\C}{\mathbb C}
\newcommand{\R}{\mathbb R}
\newcommand{\N}{\mathbb N}
\newcommand{\Z}{\mathbb Z}
\newcommand{\Q}{\mathbb Q}
\newcommand{\Rel}{\operatorname{Re}}
\newcommand{\Iml}{\operatorname{Im}}
\newcommand{\supp}{\operatorname{supp}}
\newcommand{\TV}{\operatorname{TV}}
\newcommand{\Kp}{\mathcal K^{\mathrm{prim}}_{2T}}
\newcommand{\Kw}{\mathcal K_{2T}}
\newcommand{\PT}{\mathcal P_T}
\newcommand{\IT}{I_T}
\newcommand{\lamP}{\lambda^{\mathrm{prim}}_T}
\newcommand{\lamW}{\lambda^{\mathrm W}_T}
\newcommand{\lamA}{\lambda_{\mathrm A}(T)}
\newcommand{\sT}{\sigma_T}
\newcommand{\gh}[1]{\widehat g_{#1}}
\newcommand{\nuh}{\widehat\nu}
\newcommand{\wh}{\widehat w}
\newcommand{\Fh}{\widehat F}
\newcommand{\cF}{c_F}
\newcommand{\Lev}{\mathfrak E}
\newcommand{\Wcl}{\mathcal W_T}
\newcommand{\NT}{\mathcal N_T}
\newcommand{\Aconst}{\mathsf A_T}
\newcommand{\Ainf}{\mathcal A_\infty}
\newcommand{\PiT}{\Pi_T}
\newcommand{\mT}{\mathfrak m_T}
\newcommand{\laminf}{\lambda^{\infty}_{\min}}
\newcommand{\laminfP}{\lambda^{\infty,\mathrm P}_{\min}}
\newcolumntype{Y}{>{\raggedright\arraybackslash}X}

\begin{document}
\title[Column packages for character families]{Column packages for character families: the arithmetic inputs for the sign of the radial intersection form}
\author{David Alejandro Trejo Pizzo}
\date{}
\begin{abstract}
This is the companion to \cite{Main}. That paper reduces the nonnegativity of the radial
intersection form on every window $(-T,T)$ to the existence of character families carrying a
column package with uniform constants \cite[Theorem~\ref{M:thm:alpoge-families}]{Main}, a theorem
of L.~Alpoge. Here we derive that theorem from five arithmetic inputs: a loss bound, a divisor
bound, level counts with almost orthogonality, a windowed bound for Selmer groups, and the
existence of families. The reductions are proved in full; published theorems are quoted with
their statements; the remaining statements of each input are taken from the formal verification
\cite{CharacterFamily}, and the structure of their formal proofs is outlined.
\end{abstract}
\maketitle
\setcounter{tocdepth}{1}
\tableofcontents

\section*{Introduction}

\noindent\textbf{The theorem.}
The main paper \cite{Main} proves that the radial intersection form is nonnegative on every
window $(-T,T)$, assuming the existence of families of characters that carry a column package
(Definition~\ref{M:def:package}) with constants independent of the scale. This existence statement
is Theorem~\ref{M:thm:alpoge-families} of \cite{Main}, and it is due to L.~Alpoge
\cite{CharacterFamily}. Alpoge's proof is a formal development, machine-checked in the archive
\cite{CharacterFamily}; it has not been written out on paper in full. The present annex is a
first written account of it. It derives Theorem~\ref{M:thm:alpoge-families} from five arithmetic
inputs, proves the reductions between them in full, and isolates, for each input, the statements
that are taken from the formal development.

\medskip\noindent\textbf{The five inputs.}
The annex has five sections, one for each input.
\begin{itemize}
\item[A.] \emph{The loss bound} (Section~\ref{i1:sec}): condition (P7) of Definition~\ref{M:def:package},
an upper bound for the excess of the ask term over the gain term.
\item[B.] \emph{The divisor bound} (Section~\ref{i2:sec}): condition (P1), a bound for weighted sums of
local congruence indices by local terms and Selmer lengths.
\item[C.] \emph{Level counts and almost orthogonality} (Section~\ref{i3:sec}): condition (P4) and the
small parameters of (P8) and (P10), through Bessel uniqueness, a Hilbert--Schmidt bound for the
kernel blocks off the diagonal, and counts of automorphic forms.
\item[D.] \emph{The windowed Selmer bound} (Section~\ref{i5:sec}): a bound, uniform over a window of
scales, for the total $l$-adic size of the Selmer groups of Tate twists over the CM fields that occur.
\item[E.] \emph{Existence of families} (Section~\ref{i4:sec}): the construction of hosts, columns and
bundles of automorphic origin, and the proof of Theorem~\ref{M:thm:alpoge-families} from the four
other inputs.
\end{itemize}
The logical order is the following. Section~D supplies the Selmer hypothesis used in Section~B;
Sections~A, B, C and D are the four hypotheses of the existence argument of Section~E, which
proves Theorem~\ref{M:thm:alpoge-families}. Sections~A and C also use definitions and auxiliary
results of Section~E: the notions of families and of closed properties of families
(\S\ref{i4:ssec-closed}), lemmas on closedness and on the ties of the column package, the
short-interval bound for norm primes, and a Siegel--Weil formula. None of these depends on the
main results of Section~E, so there is no circularity. In a diagram, where an arrow means that a
section uses the main result of another:
\[
\begin{array}{ccccc}
\text{D (Selmer)} & \longrightarrow & \text{B (divisor)} & & \\[2pt]
 & \searrow & \downarrow & & \\[2pt]
\text{A (loss)} & \longrightarrow & \text{E (existence)} & \longrightarrow & \text{Theorem~\ref{M:thm:alpoge-families}}\\[2pt]
 & \nearrow & & & \\[2pt]
\text{C (counts)} & & & &
\end{array}
\]
The sections are printed in the order A, B, C, D, E. Each section follows the same plan: an
introduction with the main statement and the organisation of the section; the definitions and
statements; the proofs; an outline of the formally verified steps; and a short summary of its
status.

\medskip\noindent\textbf{Three kinds of statement.}
Every statement of the annex is of one of three kinds.
(1) Statements \emph{proved here}, with a complete written proof; they carry no mark.
(2) \emph{Published theorems}, quoted with their statements and a reference in the header; we check
their hypotheses where they are applied. When we have not been able to confirm a reference
against its source, we add ``(reference to be confirmed)''.
(3) Statements whose header reads ``formally verified, \cite{CharacterFamily}''. They are
proved in the formal development and are not reproved here. Most of them are stated exactly. A
few involve objects of the formal development that we do not define; these are marked ``stated
informally; the exact statement is in \cite{CharacterFamily}''.
A statement may combine kinds, in which case its header says which parts are proved here and
which are formally verified. For each formally verified statement, the subsection ``Outline of
the formally verified steps'' at the end of its section describes the chain of steps of the formal
proof. The outline is not a proof. Each step is classified as follows:
\begin{itemize}
\item[(L)] a published theorem, or a routine specialisation of one; we state it, with what the
formal statement needs beyond the published one;
\item[(N)] a statement that, to our knowledge, is new in the form needed; we state it and indicate
the idea of its proof;
\item[(R)] routine bookkeeping (thresholds, maxima of constants, transport of structure).
\end{itemize}
Combinations such as (L)+(R) are used when a step has parts of both kinds. When we say that we do
not describe the proof of a step, it enters the outline only through its statement.

\medskip\noindent\textbf{The formal development.}
The final theorem of the formal development \cite{CharacterFamily} has no hypotheses. The
published theorems on which its argument rests (local representation theory, automorphic forms
and Galois representations, Galois cohomology, base change, the Siegel--Weil formula, complex
multiplication, and others) are isolated there as separate propositions, and each of them is
proved inside the development. Intermediate formal lemmas take bundles of these statements as
arguments, which is why they read as conditional; the bundles are discharged at the top. When we
say that a formal statement uses a published input, the input is a known theorem re-proved in the
formal development, and the attribution names the known theorem; the published theorems quoted in
this annex are citations for its written proofs. The identification of a formal proposition with
the published theorem it corresponds to is ours and is made from the statements; the validity of
the formal theorem does not depend on it.

\medskip\noindent\textbf{Conventions.}
Sections of this annex are lettered. References with a numbered section, such as
Definition~\ref{M:def:package} or Theorem~\ref{M:thm:alpoge-families}, are to the main paper
\cite{Main}, whose notation we use throughout.
\renewcommand{\thesection}{\Alph{section}}

\section{The loss bound}\label{i1:sec}

This section proves the loss bound, condition (P7) of
Definition~\ref{M:def:package}: for a family of columns, at every large scale
$x$ and at every index whose weight vector has small shares, the ask term
exceeds the gain term by at most $(3|K|+13)\,|\mathrm{TT}|\,m\,b$, where
$|\mathrm{TT}|$ is the number of tops, $m$ the number of lifts,
$b=\exp(x^{\varepsilon})$ and $K$ a constant depending only on the family; the
index loss satisfies the same bound.

\label{i1:ssec-overview}%

The proof has three layers.
\begin{enumerate}
\item An \emph{elementary, real-variable} bound
  (Theorem~\ref{i1:thm-elementary}) for an abstract ask--gain difference,
  prime by prime, after the primes are sorted into four classes
  (\emph{served}, \emph{first kind}, \emph{second kind}, \emph{high}). Its
  only arithmetic inputs are Chebyshev's bound, Mertens' bound and the
  distribution of odd integers in the square classes modulo powers of~$2$.
\item An \emph{instantiation} (Theorem~\ref{i1:thm-setting}): the column data
  of a family point satisfy the abstract hypotheses. The only input that is
  not bookkeeping is the prime-by-prime comparison $\mathrm{LinkG}$
  (Definition~\ref{i1:def-linkG}); the host field enters through elementary
  facts on Frobenius elements in a subfield of $\Q(\zeta_N)$.
\item The \emph{source of $\mathrm{LinkG}$}
  (\S\ref{i1:ssec-linkG-source}). It is reduced to group compatibility of
  the families and the second admissibility condition $\mathrm{Adm2F}$
  (eventual bounds on the number of tops and on the bad primes,
  Definition~\ref{i1:def-groupcompat}). This reduction runs through a
  local statement about lattices in unramified quadratic Hida settings and
  through $37$ admission clauses (Propositions~\ref{i1:prop-linkG},
  \ref{i1:prop-loc} and \ref{i1:prop-adm}). Part of this reduction is
  proved in \S\S\ref{i1:ssec-pf-adm}--\ref{i1:ssec-pf-loc}; the rest is
  formally verified and stated without proof
  (Propositions~\ref{i1:prop-Fm}, \ref{i1:prop-segdeg-rest},
  \ref{i1:prop-closed}, \ref{i1:prop-first}, \ref{i1:prop-gc},
  \ref{i1:prop-linkG-res}); Propositions~\ref{i1:prop-A19c}
  and \ref{i1:prop-A21} are proved in \S\S\ref{i1:ssec-pf-A19c},
  \ref{i1:ssec-pf-A21}.
\end{enumerate}

\emph{What the bound says.} The gain $g$ is a sum over places of explicitly
defined local gains minus taxes. The local gains are \emph{defined} as
truncated, shifted copies of the local ask series (Remark~\ref{i1:rem-gain}).
The loss bound is therefore a \emph{lower} bound for $g$ which holds
essentially by construction. Taken by itself it has no content, since one
could enlarge $g$. The gain is bounded from \emph{above} by the source-slope
condition (P6), and that condition holds because the column lattices are
defined by enlarging fixed lattices by the gains. The real constraint on $g$
is the frame slope inequality (P8) applied to the enlarged lattices, together
with (P1) and the period valuation bound of the column lattices (which is
not the condition $\mathrm{Adm2F}$ of Definition~\ref{i1:def-groupcompat}). In the assembly the present bound is
used only in the form $P_*\le g+\mathrm{LoP}$, $I_O\le\mathrm{LoP}$.

Throughout, $\log$ is the natural logarithm, $\log_p y=\log y/\log p$, sums
over $p$ run over the indicated finite sets of natural numbers, and
$(y)_+=\max(y,0)$.

The objects named in the following theorem are defined in
\S\ref{i1:ssec-family}: the family $\mathrm{Fm}$ and its points
$(D_{x,i},L_{x,i},X_{x,i})$ (Definition~\ref{i1:def-family-data}), the
conditions $\mathrm{ChosenTT}$, $\mathrm{ChosenBadp}$ and $\mathrm{ChW}$
(Definition~\ref{i1:def-chosen}), the comparison $\mathrm{LinkG}$
(Definition~\ref{i1:def-linkG}), and the gain $\mathrm{Ga}$, the index loss
$\mathrm{IOsum}$ and the readings (Definition~\ref{i1:def-losses}).

\begin{theorem}[Loss bound for small-share weights]\label{i1:thm-loss}
Let $\mathrm{Fm}$ be a family at exponent $\varepsilon$ that satisfies
$\mathrm{ChosenTT}$, $\mathrm{ChosenBadp}$ and $\mathrm{LinkG}$, with constant
$C$ in $\mathrm{LinkG}$. Put $K=1+|C|+C_{\mathrm{ord}}$, where
$C_{\mathrm{ord}}=12$ is the absolute constant of Lemma~\ref{i1:lem-orders}.
Then for all large $x$ and all indices $i$ at scale $x$ such that
$\mathrm{ChW}(X_{x,i})$ holds,
\begin{equation}\label{i1:eq-main}
 \mathrm{IOsum}(L_{x,i})\le(3K+13)\,|\mathrm{TT}|\,m\,b,\qquad
 \sum_{t\in\mathrm{TT}}\sum_{j=1}^{m}\mathrm{ASK}(N,a_t,t,w_j)-\mathrm{Ga}(L_{x,i})
 \le(3K+13)\,|\mathrm{TT}|\,m\,b
\end{equation}
for every reading of $L_{x,i}$ and every admissible cutoff $N$.

Moreover (second part) the family satisfies $\mathrm{LossF}(0)$ with constant
$3K+13$: there is $x_0$ such that for every $x\ge x_0$ with
$x\in\mathrm{Win}$ there are an index $i_0$ at scale $x$ and integers
$(\bar a_t,\bar b_t)_{t\in\mathrm{TT}}$ with $\bar a_t$ in the amount window
of $L_{x,i_0}$ and $b'\le\bar b_t\le b_{\rm top}(x)$
(Definition~\ref{i1:def-family-data}(8)), such that~\eqref{i1:eq-main} holds
at every index $i$ at scale $x$ whose amounts and top windows are
$a_t=\bar a_t$, $b_t=\bar b_t$ for all tops $t$. The extra input of this
part is the conclusion $\mathrm{ChosenW}$ of
Proposition~\ref{i1:prop-chosenW}, which follows from $\mathrm{LinkG}$.
Through the ties of the column package built from the family at a large
scale with these weights (Section~\ref{i4:sec}, Lemma~\ref{i4:Llem-L8}), this gives
(P7) with $c_{\rm loss}=3K+13$ and $e_L=0$.
\end{theorem}

In general the window constant is
$1+|C|+|C_b|+C_{\mathrm{ord}}$. On the route used here $C_b=0$, which gives
the value of $K$ above. The constant is not absolute, because it depends on
the constant $C$ of $\mathrm{LinkG}$. The proof gives the sharper bound
$(2\log4\cdot C+13\sigma)|\mathrm{TT}|mb$ with $\sigma=100Kx^{4\varepsilon-1}\to0$.
The formal development asserts only the existence of a constant
$C_{\mathrm{ord}}\ge0$; the value $C_{\mathrm{ord}}=12$ comes from
Lemma~\ref{i1:lem-orders}.

The section is organised as follows. Subsections~\ref{i1:ssec-params}--\ref{i1:ssec-elementary}
introduce the scale parameters and the abstract local data and prove the elementary bound
(Theorem~\ref{i1:thm-elementary}). Subsections~\ref{i1:ssec-family} and~\ref{i1:ssec-inst} define the
column data of a family point and show that they satisfy the abstract hypotheses
(Theorem~\ref{i1:thm-setting}). Subsection~\ref{i1:ssec-linkG-source} defines group compatibility
and states the reduction of $\mathrm{LinkG}$ to local data, and \S\ref{i1:ssec-proof} proves
Theorem~\ref{i1:thm-loss}. The remaining subsections prove the reduction: the admission clauses
(\S\ref{i1:ssec-pf-adm}), the local statement and the statement for families
(\S\ref{i1:ssec-pf-loc}), and four propositions on column lattices and segment modules
(\S\ref{i1:ssec-colpol}). Subsection~\ref{i1:ssec-structure} outlines the formally verified steps,
and the section ends with a summary of its status.

\subsection{Scale parameters and windows}\label{i1:ssec-params}

\begin{definition}[Parameter records]\label{i1:def-par}
A \emph{parameter record} $P$ consists of real numbers
$x,\varepsilon,n,k,N,t,\Lambda,r_G,S_c,\eta,\rho,\vartheta,K_0,C,C_{\mathrm{ord}},Y,b_o$
and natural numbers $e_2,p_0,d$. The derived quantities are
\[
b=e^{x^{\varepsilon}},\ a=xb,\ X'=e^{x^{\varepsilon/3}},\ b'=a/X',\
w''=\frac{b'}{xn},\ X''=\frac{a}{w''}=xnX',\ f=2n,\ \ell_x=\log x,\
\delta_1=x^{-\varepsilon/8},\ \theta=x^{-\varepsilon},
\]
\[
S_k=S_c+k+\eta,\quad Q_1=C\,r_G\,b'\,\ell_x,\quad
Q_2=C\Bigl(\frac{b'}{x}+\eta+S_c+\frac{(w''+t)\ell_x}{\delta_1}\Bigr)+b_o,\quad
J_o=\lfloor x^{\varepsilon/4}\rfloor,\quad I_b=\lfloor x^{\varepsilon/3}\rfloor .
\]
\end{definition}

\begin{definition}[Windows; large $x$]\label{i1:def-windows}
Let $e,K\in\R$. The record $P$ lies in the \emph{windows} $\mathrm W(e,K)$ if
$0<\varepsilon\le1$, $e\ge0$, $K\ge1$ and
\begin{gather*}
x^2\le 2^{e_2}\le n\le x^e,\quad x^2\le k\le x^e,\quad 0\le N,t\le x^e,\quad
x^4+1\le p_0\le x^e,\quad 0\le\Lambda\le x^{2\varepsilon},\quad 0\le r_G\le x^e,\\
0\le S_c,\eta\le e^{x^{\varepsilon/4}},\quad 0\le\rho\le x^{-2},\quad
0\le\vartheta\le x^{-1},\quad x^2\le K_0,\quad x^2\le d,\quad 0\le C\le K,\\
0\le C_{\mathrm{ord}}\le K,\quad 1\le Y\le 2X',\quad 0\le b_o\le e^{x^{\varepsilon}/2}.
\end{gather*}
The real $x$ is \emph{large} for $(\varepsilon,e,K)$ if $x\ge1$ and
\[
100Kx^{4\varepsilon}\le x,\qquad 100Kx^{4e+4}\le e^{x^{\varepsilon/3}},\qquad
100Kx^{4e+4}e^{x^{\varepsilon/3}+x^{\varepsilon/4}}\le e^{x^{\varepsilon}/2}.
\]
Put $U=|T|\,m\,b$ (the \emph{unit}) and $\sigma=100Kx^{4\varepsilon}/x$.
\end{definition}

\begin{lemma}\label{i1:lem-large}
For $\varepsilon\in(0,\frac14)$ and $e,K\in\R$, every sufficiently large $x$ is
large for $(\varepsilon,e,K)$. If $P\in\mathrm W(e,K)$ and $x$ is large, then,
with $z=x^\varepsilon$, $u=x^{\varepsilon/3}$, $E=x^e$, $S=e^{x^{\varepsilon/4}}$:
$x\ge100$, $1\le u\le z$, $\log a=\log x+z$, $\log x\le z$, $e\log x\le z$,
$100KE^4x^4\le X'$, $100KE^4x^4X'S\le e^{z/2}=\sqrt b$, $J_o\le z$, $I_b\le u$,
and $0\le\sigma\le1$.
\end{lemma}

\begin{proof}
Since $1-4\varepsilon>0$, $x^{1-4\varepsilon}\ge100K$ for large $x$. Since
$\log x=o(x^{\varepsilon/3})$, $100Kx^{4e+4}\le e^{x^{\varepsilon/3}}$ for
large $x$, and then $100Kx^{4e+4}e^{u+x^{\varepsilon/4}}\le e^{3u}\le e^{z/2}$
as soon as $z\ge6u$. The listed consequences are these inequalities rewritten
with $x^{4\varepsilon}=z^4$ and $x^{4e+4}=E^4x^4$, together with the
definitions. For example, $x\ge100$ follows from $100Kz^4\le x$ and $Kz^4\ge1$,
and $e\log x\le z$ follows from $(4e+4)\log x\le u$.
\end{proof}

\subsection{Abstract local data}\label{i1:ssec-abstract}

For a finite $I\subset\N$, reals $\alpha\ge0$, $p>0$ and $w:\N\to\R$ put
$\operatorname{ask}(I,\alpha,p,w)=\sum_{i\in I}(\alpha/p^i-w(i))_+$.

\begin{definition}[Local datum]\label{i1:def-LD}
A \emph{local datum} $\mathcal D$ over finite sets $P_s\subset\N$ (primes), $J$
(lifts) and $T$ (tops), and a natural number $d$, consists of:
\begin{itemize}
\item a classification $\mathrm{cl}:P_s\to\{\text{served},\text{first},
 \text{second},\text{high}\}$ and, for each $p$, a set $T_p$ of served tops;
\item finite sets $I^{\rm ex}_{j,t,p},I^d_{j,t}\subset\N$, amounts $\alpha_t$,
 and weights $w_{j,t,p}(i)$, $w^d_{j,t}(i)$;
\item gains $\gamma_j(p)$, taxes $\tau(p)$, local losses $\mathrm{lo}(p)$,
 floors $\mathrm{fl}(p)$, and exponents $L_1(p)$.
\end{itemize}
Put $A_{j,t}(p)=\operatorname{ask}(I^{\rm ex}_{j,t,p},\alpha_t,p,w_{j,t,p})$ and
\[
\mathrm{ASK}_{\mathcal D}=\sum_{j\in J}\sum_{t\in T}\Bigl(\sum_{p\in P_s}2\log p\,A_{j,t}(p)
+\log d\cdot\operatorname{ask}(I^d_{j,t},\alpha_t,d,w^d_{j,t})\Bigr),\quad
\mathrm{Ga}_{\mathcal D}=2\sum_{p\in P_s}\log p\Bigl(\sum_{j\in J}\gamma_j(p)-\tau(p)\Bigr).
\]
A served prime $p$ satisfies the \emph{served comparison} if $T_p\subseteq T$ and
\begin{equation}\label{i1:eq-served}
\sum_{j\in J}\sum_{t\in T_p}\bigl(A_{j,t}(p)-\mathrm{lo}(p)\bigr)-|J|\,\mathrm{fl}(p)
\le\sum_{j\in J}\gamma_j(p)-\tau(p).
\end{equation}
\end{definition}

More generally one may allow a set of places $W$ fibred over $P_s$ with weights. By a
\emph{unit rule} (the weights over $p$ sum to $\log p$, and gains and taxes
are constant on the fibres) such a gain reduces to the form above. In our
instantiation $W=P_s$ and the fibres are singletons, so we use the reduced
form directly.

\begin{definition}[Model losses]\label{i1:def-model}
For a parameter record $P$ and a prime $p$ put
\begin{gather*}
\mathrm{fl}^\star(p)=CfN^2\log_pa+Cf(\log_pa)^2,\qquad
\mathrm{cp}(p)=CS_k\log_p(2a),\qquad
\mathrm{pr}(p)=Q_1\ (p<p_0),\quad Q_2\ (p\ge p_0),\\
\mathrm{or}(p)=C\ell_x\Bigl(\frac{\Lambda a}{pk}+
\begin{cases}a\,p^{-L_1(p)}&p<p_0\\ a/p^2&p\ge p_0\end{cases}\Bigr),\qquad
\mathrm{lo}^\star(p)=\mathrm{cp}(p)+\mathbf 1_{p<X'}\bigl(\mathrm{pr}(p)+\mathrm{or}(p)\bigr).
\end{gather*}
\end{definition}

Let $H_e=(\Z/2^{e+2}\Z)^\times/\{\pm1\}$, a group of order $2^e$, and for odd
$l$ let $h_e(l)\in H_e$ be the class of $l$.

\begin{definition}[Prime-class hypotheses]\label{i1:def-PCLH}
Let $P$ be a parameter record, $\mathcal D$ a local datum over $(P_s,J,T,d(P))$,
$\mathrm{cand},P_R\subset\N$ finite sets, $\mathrm{BAD}(p,t)$ a relation,
$\Sigma:\N\to\R$ and $\mathrm{io}:\N\to\R$. Put $m=|J|$ and
\[
\mathrm{lhs}_W=\sum_{p\in\mathrm{cand}}\frac{\log p}{p}\bigl(2\cdot\mathbf 1_{\theta<\Sigma(p)}+4\Sigma(p)\bigr),\qquad
\mathrm{Bad}=\sum_{p\in P_R}\frac{\log p}{p}\,\#\{t\in T:\mathrm{BAD}(p,t)\}.
\]
The hypotheses $\mathrm{PCLH}$ are the following.
\begin{enumerate}
\item[(H1)] $T\ne\emptyset$; every $p\in P_s$ is an odd prime; $\Sigma\ge0$ on
 $\mathrm{cand}$; $\mathrm{lhs}_W\le x^{-9}$; $\mathrm{Bad}\le|T|x^{-2+\varepsilon}$;
 and $\mathrm{io}(p)\le2\Sigma(p)\cdot2m|T|a/p$ for $p\in\mathrm{cand}$.
\item[(H2)] (\emph{side conditions}) $0\le\alpha_t\le a(1+\vartheta)$; every
 $i\in I^{\rm ex}_{j,t,p}$ is odd with $w_{j,t,p}(i)\ge w''$; every
 $i\in I^d_{j,t}$ is $\ge1$ with $w^d_{j,t}(i)\ge0$; and for served $p$:
 $\mathrm{fl}(p)\le\mathrm{fl}^\star(p)$, $\mathrm{lo}(p)\le\mathrm{lo}^\star(p)$.
 For $p\in P_s$ with $p<p_0$, $K_0\le p^{L_1(p)-1}$.
\item[(H3)] (\emph{classes}) Served primes satisfy $p\le X''$ and the served
 comparison~\eqref{i1:eq-served}. Non-served primes have $\gamma_j(p)=0=\tau(p)$.
 High primes satisfy $p>X''$. Second-kind primes lie in $\mathrm{cand}$ with
 $\theta<\Sigma(p)$.
\item[(H4)] (\emph{first kind}) Every first-kind $p$ satisfies one of two
 conditions. Either $p\le\max(Y,X'')$, $p$ is odd and
 $\operatorname{ord}h_{e_2}(p)<x^{\varepsilon/4}$. Or $p\ge p_0$ and
 $X'(1-\rho)\le p^i\le X'(1+\rho)$ for some $i\ge1$.
\item[(H5)] (\emph{unserved tops are bad}) If $p$ is served and
 $t\in T\setminus T_p$, then $p\in P_R$ and $\mathrm{BAD}(p,t)$.
\end{enumerate}
\end{definition}

\subsection{The elementary loss bound}\label{i1:ssec-elementary}

\begin{lemma}[Chebyshev and Mertens, explicit]\label{i1:lem-cheb}
Let $s$ be a set of primes $\le Y$. Then $\sum_{p\in s}\log p\le Y\log4$, and,
if $Y\ge1$, $\sum_{p\in s}\frac{\log p}p\le\log Y+\log4$.
\end{lemma}

\begin{proof}
Both sums only increase if $s$ is replaced by all primes $\le N=\lfloor
Y\rfloor$. For the first sum we use $\prod_{p\le N}p\le 4^N$. Induct on $N$.
For even $N>2$ the claim follows from the case $N-1$. For $N=2M+1$, every
prime in $(M+1,2M+1]$ divides $\binom{2M+1}{M}\le4^M$. For the second sum, $N!$
is divisible by $p^{\lfloor N/p\rfloor}$, so
$\sum_{p\le N}\lfloor N/p\rfloor\log p\le N\log N$. Since $N/p\le\lfloor
N/p\rfloor+1$, we get $N\sum_{p\le N}\frac{\log p}p\le N\log N+N\log 4$.
\end{proof}

\begin{lemma}\label{i1:lem-tail}
If $p_0\ge3$ and $s$ is a finite set of integers $\ge p_0$, then
$\sum_{n\in s}\frac{\log n}{n^2}\le\frac{2(\log(p_0-1)+1)}{p_0-1}$.
\end{lemma}

\begin{proof}
Put $g(y)=(\log y+1)/y$. For $y\ge3$ one checks
$\log y/y^2\le2(g(y-1)-g(y))$, using $\log y-\log(y-1)\le 1/(y-1)\le\frac12$.
The sum then telescopes.
\end{proof}

\begin{lemma}[Geometric bounds]\label{i1:lem-geom}
Let $\alpha\ge0$ and $w\ge0$ on $I$. (1) If $p\ge2$ and all $i\in I$ are odd,
then $\operatorname{ask}(I,\alpha,p,w)\le\frac43\alpha/p$; if moreover
$\alpha\le\frac32a$, then $2\log p\cdot\operatorname{ask}\le4a\log p/p$. (2) If
$d\ge4$, every $i\in I$ is $\ge1$ and $\alpha\le\frac32a$, then
$\log d\cdot\operatorname{ask}(I,\alpha,d,w)\le2a\log d/d$.
\end{lemma}

\begin{proof}
Since $w\ge0$, each term is at most $\alpha p^{-i}$. Sum the geometric series:
over odd $i$ it is at most $q/(1-q^2)\le\frac43q$ with $q=1/p\le\frac12$, and
over $i\ge1$ it is $1/(d-1)\le\frac43/d$.
\end{proof}

\begin{lemma}[Splitting by classes]\label{i1:lem-split}
Assume (H3). Then
\begin{multline*}
\mathrm{ASK}_{\mathcal D}-\mathrm{Ga}_{\mathcal D}\le\sum_{p\ \text{served}}2\log p\Bigl(
\sum_{j}\sum_{t\in T\setminus T_p}A_{j,t}(p)+|J||T_p|\mathrm{lo}(p)+|J|\mathrm{fl}(p)\Bigr)\\
+\sum_{p\ \text{not served}}2\log p\sum_{j}\sum_{t\in T}A_{j,t}(p)
+\sum_{j,t}\log d\cdot\operatorname{ask}(I^d_{j,t},\alpha_t,d,w^d_{j,t}).
\end{multline*}
\end{lemma}

\begin{proof}
We have $\mathrm{ASK}-\mathrm{Ga}=\sum_p2\log p\,\mathrm{def}(p)+(\text{$d$-terms})$,
where $\mathrm{def}(p)=\sum_{j,t}A_{j,t}(p)-\sum_j\gamma_j(p)+\tau(p)$. For a
served $p$, split $T=(T\setminus T_p)\sqcup T_p$ and apply~\eqref{i1:eq-served}.
For a non-served $p$, $\gamma_j(p)=\tau(p)=0$.
\end{proof}

\begin{lemma}[Bucket bounds]\label{i1:lem-buckets}
Let $s$ be a finite set of primes and $P\in\mathrm W(e,K)$ with $x$ large.
\begin{enumerate}
\item If $s\subseteq[2,X'']$, then
 $\sum_{p\in s}2\log p\,\mathrm{fl}^\star(p)\le2CfX''(N^2\log a+2\log^2a)$ and
 $\sum_{p\in s}2\log p\,\mathrm{cp}(p)\le2CS_k\log(2a)X''$.
\item If every $p\in s$ is $<X'$, then
 $\sum_{p\in s}2\log p\,\mathrm{pr}(p)\le2\log4\,(p_0Q_1+X'Q_2)$. If moreover
 $K_0\le p^{L_1(p)-1}$ for every $p\in s$ with $p<p_0$, then
 \[
 \sum_{p\in s}2\log p\,\mathrm{or}(p)\le2C\ell_xa\Bigl(\bigl(\tfrac{\Lambda}{k}+\tfrac1{K_0}\bigr)(\log X'+\log4)+\tfrac{2(\log(p_0-1)+1)}{p_0-1}\Bigr).
 \]
\item (\emph{high primes}) Suppose $0\le\alpha\le a(1+\vartheta)$, that every
 $p\in s$ exceeds $X''$, and that every $i\in I_p$ is $\ge1$ with
 $w_p(i)\ge w''$. Then
 $\sum_{p\in s}2\log p\operatorname{ask}(I_p,\alpha,p,w_p)\le2\log(2X'')(a\vartheta^2+w''\vartheta)$.
\end{enumerate}
\end{lemma}

\begin{proof}
(1) We have $2\log p\,\mathrm{fl}^\star(p)=2Cf(N^2\log a+\log^2a/\log p)$ and
$\log p>\frac12$. Each term of the second sum equals $2CS_k\log(2a)$. There
are at most $X''$ terms. (2) Split at $p_0$ and apply
Lemma~\ref{i1:lem-cheb}. For $\mathrm{or}$ with $p<p_0$, we have
$ap^{-L_1(p)}\le a/(K_0p)$. Apply Lemma~\ref{i1:lem-cheb} twice and
Lemma~\ref{i1:lem-tail}. (3) Recall $X''w''=a$. For $i\ge2$,
$p^i\ge p^2>2X''$, so $\alpha/p^i\le w''$ and the term vanishes. For $i=1$
the term vanishes unless $p<X''(1+\vartheta)$, and then it is at most
$\alpha/X''-w''\le w''\vartheta$. At most $X''\vartheta+1$ integers lie in
$(X'',X''(1+\vartheta))$.
\end{proof}

\begin{lemma}[Square classes modulo powers of $2$]\label{i1:lem-pm}
Let $0\le j\le e$ and let $l$ be odd. Then $h_e(l)^{2^j}=1$ if and only if
$l\equiv\pm1\pmod{2^{e+2-j}}$.
\end{lemma}

\begin{proof}
We show, by induction on $j$ and for odd $y$ and $m\ge1$, that
$y^{2^j}\equiv\pm1\pmod{2^{m+j}}$ if and only if $y\equiv\pm1\pmod{2^m}$. Put
$Y=y^{2^j}=2k+1$. Then $Y^2+1\equiv2\pmod 4$, so the condition at level $j+1$
reads $2^{m+j+1}\mid Y^2-1=4k(k+1)$. One of $k,k+1$ is odd, so this holds iff
$2^{m+j}\mid Y\mp1$ for one sign, which is the condition at level $j$. Apply
the claim with $m=e+2-j$.
\end{proof}

\begin{lemma}[Primes of small order in $H_e$]\label{i1:lem-orders}
Let $e,J\ge1$, $Y\ge1$, and let $s$ be a set of odd primes $l\le Y$ with
$\operatorname{ord}h_e(l)\le J$. Then
$\sum_{l\in s}\frac{\log l}l\le12J\log(2^eY)^2/2^e$.
\end{lemma}

\begin{proof}
Let $j_0=\min(e,\lfloor\log_2J\rfloor)$ and $M=2^{e+2-j_0}$, so that $2^e\le JM$.
The order of $h_e(l)$ is a power of $2$ that is at most $J$, hence it divides
$2^{j_0}$. By Lemma~\ref{i1:lem-pm}, $l\equiv\pm1\pmod M$. The integers
$l\le Y$ of the form $qM\pm1$ with $q\ge1$ satisfy
$\log l/l\le2\log Y/(qM)$. Summing over $q\le Y$ gives
$\sum_{l\in s}\log l/l\le4\log Y(1+\log Y)/M$. Finally
$\log Y(1+\log Y)\le3\log(2^eY)^2$ because $e\log2\ge\frac12$.
\end{proof}

\begin{lemma}[First-kind primes]\label{i1:lem-first}
Let $e,J\ge1$, $Y\ge1$, $0\le\rho\le\frac12$, $p_0\ge8$ and $X'\ge2$. Let $s$
be a set of odd primes $p$, each satisfying one of two conditions: either
$p\le Y$ and $\operatorname{ord}h_e(p)\le J$; or $p\ge p_0$ and
$X'(1-\rho)\le p^i\le X'(1+\rho)$ for some $i\ge1$. Then
\[
\sum_{p\in s}\frac{\log p}{p}\le12J\frac{\log(2^eY)^2}{2^e}+\lfloor\log X'\rfloor\Bigl(4\rho+\frac1{p_0}\Bigr)\log(2X').
\]
\end{lemma}

\begin{proof}
The primes of the first kind are handled by Lemma~\ref{i1:lem-orders}. Take a
prime of the second kind. Since $p\ge8$ and $p^i\le2X'$, we have
$i\le\log X'$. Monotonicity of $y\mapsto y^i$ gives
$X'^{1/i}(1-\rho)\le p\le X'^{1/i}(1+\rho)$. So $p$ lies in one of
$\lfloor\log X'\rfloor$ bands $[P_i(1-\rho),P_i(1+\rho)]$, where
$P_i=X'^{1/i}\le X'$. Each band contains at most $2\rho P_i+1$ integers
$n\ge p_0$, and each of these has
$\log n/n\le\log(2P_i)\min(2/P_i,1/p_0)$.
\end{proof}

\begin{lemma}[Size bounds]\label{i1:lem-size}
Let $P\in\mathrm W(e,K)$ with $x$ large for $(\varepsilon,e,K)$. Let $m\ge0$,
$T_\#\ge1$ and $U=T_\#mb$. Each of the following eleven quantities is at most
$\sigma U$:
\begin{gather*}
m\,2CfX''(N^2\log a+2\log^2a),\quad 2CS_k\log(2a)mT_\#X'',\quad
2\log4\,mT_\#p_0Q_1,\\
mT_\#\,2C\ell_xa\bigl((\tfrac\Lambda k+\tfrac1{K_0})(\log X'+\log4)+\tfrac{2(\log(p_0-1)+1)}{p_0-1}\bigr),\quad
4amT_\#x^{-2+\varepsilon},\\
4mT_\#a\,C_{\mathrm{ord}}J_o\frac{\log(2^{e_2}\max(Y,X''))^2}{2^{e_2}},\quad
4mT_\#a\,I_b(4\rho+\tfrac1{p_0})\log(2X'),\\
mT_\#2\log(2X'')(a\vartheta^2+w''\vartheta),\quad
mT_\#\tfrac{2a\log d}{d},\quad 2amT_\#x^{-9},\quad amT_\#x^{-9}.
\end{gather*}
Moreover $2\log4\,mT_\#X'Q_2\le(2\log4\cdot C+3\sigma)U$.
\end{lemma}

\begin{proof}
Put $z=x^\varepsilon$, $u=x^{\varepsilon/3}$, $E=x^e$, $S=e^{x^{\varepsilon/4}}$
and $v=b/X'$, so that $b'=xv$ and $w''=v/n$. By Lemma~\ref{i1:lem-large} and
the windows:
\begin{itemize}
\item $\log a\le2z\le x$, $\log X'=u\le z$, and $\log y\le z$ for $0<y\le E$;
\item $100Kz^i/x\le\sigma$ for $i\le4$;
\item $cKE^ix^jX'S\le\sigma b$ and $cKE^ix^jv\le\sigma b$ for $i\le4$, $j\le3$
 and $c\le100$. Indeed $100KE^4x^4X'S\le\sqrt b\le b$ and $100KE^4x^4\le X'$.
\end{itemize}
It then suffices to bound each quantity, with the factor $mT_\#$ removed, by
$\sigma b$. We check each one.
\begin{itemize}
\item Floors: $\le2K(2E)(xEX')(E^2x+2x^2)\le12KE^4x^3X'$.
\item $\mathrm{cp}$: $S_k\le3ES$, so the quantity is $\le6KE^2x^2X'S$.
\item Low $\mathrm{pr}$: $\le3KE^2x^2v$.
\item $\mathrm{or}$: we have $\Lambda/k\le z^2x^{-2}$, $1/K_0\le x^{-2}$ and
 $p_0-1\ge x^4$, so the bracket is $\le10z^3x^{-2}$ and the quantity is
 $\le\frac15b\cdot100Kz^4/x$.
\item Orders: $\log(2^{e_2}\max(Y,X''))\le4z$ and $2^{e_2}\ge x^2$ give
 $\le\frac{64}{100}b\cdot100Kz^3/x$.
\item Bands: $\le4xb\cdot z\cdot5x^{-2}\cdot2z$.
\item High primes: $\log(2X'')\le4z$, $a\vartheta^2\le b/x$ and
 $w''\vartheta\le b/x$.
\item The $d$-term: $\log d/d\le2\log x/x^2$, since $d\ge x^2$.
\item The bad term: $4xb\,x^{-2}z$.
\item The second-kind and io terms: $2ax^{-9}\le b/x$.
\end{itemize}
Finally, expand
$X'Q_2=Cb+CX'(\eta+S_c)+C\frac bn\ell_x\delta_1^{-1}+CX't\ell_x\delta_1^{-1}+X'b_o$.
Four times each of the last four summands is $\le\sigma b$; this uses
$b_o\le\sqrt b$ and $100KE^4x^4X'S\sqrt b\le b$. Hence
$2\log4\,X'Q_2\le(2\log4\cdot C+3\sigma)b$.
\end{proof}

\begin{theorem}[Elementary loss bound]\label{i1:thm-elementary}
Let $\varepsilon\in(0,\frac14)$ and $e,K\in\R$. There is $x_0$ with the
following property. Let $P\in\mathrm W(e,K)$ be a record with
$\varepsilon(P)=\varepsilon$, $C_{\mathrm{ord}}(P)\ge12$ and $x(P)\ge x_0$, and
let the data of Definition~\ref{i1:def-PCLH} satisfy $\mathrm{PCLH}$. Then
\[
\mathrm{ASK}_{\mathcal D}-\mathrm{Ga}_{\mathcal D}\le(2\log4\cdot C+13\sigma)\,|T|\,|J|\,b\le(3C+13)|T||J|b,\qquad
\sum_{p\in\mathrm{cand}}\log p\,\mathrm{io}(p)\le\sigma|T||J|b\le|T||J|b .
\]
\end{theorem}

\begin{proof}
Take $x_0$ such that every $x\ge x_0$ is large (Lemma~\ref{i1:lem-large}). Put
$m=|J|$ and $T_\#=|T|$. Start from Lemma~\ref{i1:lem-split}. Since
$\alpha_t\le a(1+\vartheta)\le\frac32a$, Lemma~\ref{i1:lem-geom} bounds each
$2\log p\,A_{j,t}(p)$ by $4a\log p/p$, and it bounds the $d$-terms by
$mT_\#\cdot2a\log d/d$. The served primes lie in $[2,X'']$, and the model
losses are nonnegative. So (H2) and Lemma~\ref{i1:lem-buckets}(1),(2) bound the
$\mathrm{lo}$ and $\mathrm{fl}$ contributions by the first four quantities of
Lemma~\ref{i1:lem-size} and the $Q_2$-term. The high primes are handled by
Lemma~\ref{i1:lem-buckets}(3).

The remaining contributions are bounded as follows.
\begin{itemize}
\item \emph{Unserved tops at served primes.} By (H5), the sum
 $\sum_p4ma\frac{\log p}{p}|T\setminus T_p|$ is at most
 $4ma\,\mathrm{Bad}\le4amT_\#x^{-2+\varepsilon}$.
\item \emph{Second-kind primes.} For these, $\theta<\Sigma(p)$, so
 $4mT_\#a\frac{\log p}p\le2amT_\#\frac{\log p}p(2\cdot\mathbf 1_{\theta<\Sigma}+4\Sigma)$.
 Their total is at most $2amT_\#\mathrm{lhs}_W\le2amT_\#x^{-9}$.
\item \emph{First-kind primes.} By (H4) and Lemma~\ref{i1:lem-first}, applied
 with $e=e_2$, $J=J_o$ and $Y\mapsto\max(Y,X'')$, their contribution is at
 most the two order and band quantities of Lemma~\ref{i1:lem-size}. The
 windows give $e_2\ge1$, $\rho\le\frac12$, $p_0\ge8$ and
 $\lfloor\log X'\rfloor=I_b$.
\end{itemize}
Lemma~\ref{i1:lem-size} adds the contributions up to
$(2\log4\cdot C+13\sigma)U$. For the io-sum, (H1) gives
$\log p\,\mathrm{io}(p)\le amT_\#\frac{\log p}p\cdot4\Sigma(p)$. The total is
therefore at most $amT_\#\mathrm{lhs}_W\le amT_\#x^{-9}\le\sigma U$. Finally
$2\log4\le3$ and $\sigma\le1$.
\end{proof}

\subsection{The column data of a family point}\label{i1:ssec-family}

We recall the structure of a family point only to the extent that the proof
uses it. The full records are those of the construction of the families
(Section~\ref{i4:sec}). The facts listed in Definition~\ref{i1:def-family-data} are axioms
of a family, valid for $x>x_1(\mathrm{Fm})$.

\begin{definition}[Family data used here]\label{i1:def-family-data}
A \emph{family} $\mathrm{Fm}$ at exponent $\varepsilon\in(0,\frac1{52})$
assigns to every scale $x>x_1(\mathrm{Fm})$ and every index $i$ a column record
$D=D_{x,i}$, size data $L=L_{x,i}$ over $D$ at scale $(x,\varepsilon)$, and an
instance $X=X_{x,i}$. The index set at some scale is non-empty (so a
family has a point, and a column record exists). These carry the following
data.
\begin{enumerate}
\item \emph{Ranks and frame.} There are a finite abelian group $H$ of order
 $n$, the \emph{frame}, together with a surjection $\varpi:H\to H_{e_2}$ that
 is injective on a cyclic subgroup of index $o\le x^{\varepsilon/100}$; ranks
 $k,m,t$ with $t$ odd and $m+3\le k$; the set $\mathrm{TT}\subseteq H$ of
 \emph{tops}, with integer amounts $a\le a_t\le a+a/(kx^{13})$; and the
 quantities $N(m)=k+m(t+1)$, $\Lambda_G$, $S_c=4k\Delta$, $\eta>0$, $r_G$,
 $\rho_B$ entering the parameter record (Definition~\ref{i1:def-parL}).
\item \emph{Size inputs.} With $\vartheta_s=\varepsilon/100$ and
 $\varepsilon_z=\varepsilon/J(\mathrm{Fm})$:
 $x^{2+8\varepsilon\mp\vartheta_s}\gtrless n$,
 $x^{9/2+10\varepsilon\mp\vartheta_s}\gtrless k$,
 $m\le x^{\vartheta_s}\sqrt{x\delta_1}$,
 $\log d\ge x^{\varepsilon_z/5-\varepsilon_z/100}$,
 $|\mathrm{TT}|\le2n\delta_1$, $x^\varepsilon\le\Lambda_G\le x^{\varepsilon+\vartheta_s}$,
 $\log\eta\le x^{\varepsilon/4}$, $\log\Delta\le x^{\varepsilon/4}/2$,
 $1\le t\le x^{C_t}$, $2^{e_2}\mid n$, $\rho_B\le x^{-2}$, and $d$ is prime.
 Here the inequalities written with $\mp$ and $\gtrless$ are two-sided
 bounds, for example $x^{2+8\varepsilon-\vartheta_s}\le n\le x^{2+8\varepsilon+\vartheta_s}$.
\item \emph{Host.} The host consists of primes $q\equiv3\pmod8$ and
 $d\equiv7\pmod{16}$ with $(-d/q)=1$; an integer $Q_S$ with
 $2^{e_2+2}\mid Q_S$, $\gcd(Q_S,qd)=1$ and $(-d/l)=1$ for every odd prime
 $l\mid Q_S$; and the conductor $N=Q_Sqd$. It also consists of a Galois
 number field $E$ with an isomorphism
 $\gamma:\operatorname{Gal}(E/\Q)\to G_D=\{\pm1\}^2\times H$; an embedding
 $E\hookrightarrow\Q(\zeta_N)$; and a homomorphism
 $\mathrm{art}:(\Z/N)^\times\to G_D$. These satisfy:
 \begin{itemize}
 \item $\mathrm{art}(u)$ acts on $E$ as $\zeta_N\mapsto\zeta_N^u$;
 \item the second sign of $\gamma(\sigma)$ is the action of $\sigma$ on
  $\sqrt{-d}\in E$;
 \item $\varpi$ of the $H$-component of $\mathrm{art}(u)$ is the class of
  $u$ in $H_{e_2}$;
 \item $\mathrm{art}(-1)$ has signs $(-1,-1)$ and $H$-component $0$;
 \item $\mathrm{art}(u)\in\{1,\iota_d\}$ when $u\equiv1\pmod{Q_Sq}$, where
  $\iota_d$ has signs $(1,-1)$ and $H$-component $0$.
 \end{itemize}
 A prime $p\nmid N$ is \emph{split} if $(-d/p)=1$. Otherwise it is
 \emph{signal} or \emph{inert} according as the first sign
 $\mathrm{ps}(p)$ of $\mathrm{art}(p)$ is $+1$ or $-1$. Signal and inert
 primes are of \emph{inert type}. The \emph{inert class} is the set of
 inert primes and the \emph{signal class} the set of signal primes; ``$p$ is
 of the inert class'' means that $p$ is inert. Primes dividing $N$ are \emph{ramified}.
 For a prime $p$, $h_P(p)\in H$ denotes the $H$-component of the
 $\gamma$-image of a chosen arithmetic Frobenius above $p$.
\item \emph{Places, weights and fills.} The size data carry the following.
 \begin{itemize}
 \item A finite set of places $w$ over residue characteristics $p(w)$,
  containing all places with $p(w)\le X''$, with weights $l(w)e(w)$ summing to
  $\log p$ over $p$.
 \item For each lift $1\le j\le m$ and each element of
  $\Psi_D=\{g\in G_D:\ \text{second sign of }g=1\}$, a size weight
  $w^{\rm At}_j(g)$. This weight is at least $b$, $b_t+b'/x$ or $w''$
  according as the place is thick, a top over $t$, or thin. Here $b_t\ge b'$
  is the top window.
 \item Local gains $\mathrm{gain}(j,w)$, taxes $\tau(w)$, exponents
  $\mathrm{io}(w)$ and served tops $\mathrm{TT}_w$.
 \end{itemize}
 The \emph{fills} say that these local quantities are the instance quantities
 at $p(w)$: $\mathrm{gain}(j,w)=\mathrm{gain}_P(p(w),j)$,
 $\tau(w)=\tau_P(p(w))$, $\mathrm{io}(w)=\mathrm{io}_P(p(w))$ and
 $\mathrm{TT}_w=\mathrm{TTP}(p(w))$ (Remark~\ref{i1:rem-gain}). They also say
 that the partner weights $w^{\rm part}_{j,t,p}(i)$ are the size weights at
 the partner places.
\item \emph{Instance quantities at a prime $p$.} These are a set
 $\mathrm{VV}(p)$ of \emph{served orbit representatives}, and numbers
 $\Sigma(p)\ge0$, $S_v(p,t_0)$ and $N_v(p,t_0)\le|\mathrm{TT}|$. Put
 $\mathrm{io}_P(p)=2\Sigma(p)\max_{t_0\in\mathrm{VV}(p)}S_v(p,t_0)$, or $0$
 if $\mathrm{VV}(p)=\emptyset$. The prime $p$ is a \emph{candidate},
 $\mathrm{Cand}(p)$, if it is of the inert class, $p<X'$,
 $\operatorname{ord}h_P(p)\ge x^{\varepsilon/4}$, and, if $p\ge p_0(k)$, there
 is no $i\ge1$ with $|p^i-X'|\le\rho_BX'$. Then
 $\mathrm{VV}(p)=\emptyset$ unless $\mathrm{Cand}(p)$ and
 $\Sigma(p)\le x^{-\varepsilon}$. In that case $\mathrm{VV}(p)$ is the set of
 good orbit representatives, and a top is \emph{bad at $p$} if its
 representative is not good. $\mathrm{TTP}(p)=\mathrm{TT}$ if $p$ is signal
 or $p\ge X'$; otherwise $\mathrm{TTP}(p)$ is the set of tops whose
 representative lies in $\mathrm{VV}(p)$.
\item \emph{Truncated ask sums.} For $1\le j\le m$, $t\in\mathrm{TT}$ and $p$
 of inert type,
 \[
 A^L_{j,t}(p)=\sum_{\substack{1\le i\le\log_pa_t+1\\ i\text{ odd}}}\bigl(a_t/p^i-w^{\rm part}_{j,t,p}(i)\bigr)_+ ,
 \]
 and $\mathrm{ASK}(L)=\sum_j\sum_{t}\bigl(\sum_{p\le a_t\ \text{inert type}}2\log p\,A^L_{j,t}(p)+\log d\cdot\operatorname{ask}_d(a_t,w^{\rm At}_j,t)\bigr)$,
 where $\operatorname{ask}_d$ is the $d$-term of
 Definition~\ref{i1:def-groupask}.
\item \emph{Model losses} $\mathrm{lo}_G(p)$ and $\mathrm{fl}_G(p)$ at the
 instance, given by
 \begin{itemize}
 \item $\mathrm{fl}_G=CfN^2\log_pa+Cf(\log_pa)^2$;
 \item $\mathrm{lo}_G=C(S_c+k+\eta)\log_p(2a)+\mathrm{pr}_G+\mathrm{or}_G$.
 \end{itemize}
 Here $\mathrm{pr}_G$ and $\mathrm{or}_G$ vanish unless $p$ is of the inert
 class and $p<X'$. In that case they equal $C r_Gb'\log x$ and
 $C\log x(\Lambda_Ga/(pk)+ap^{-L_1(p)})$ for $p<p_0(k)$, and
 $C(b'/x+\eta+S_c+(w''+t)\log x/\delta_1)+b_o$ and
 $C\log x(\Lambda_Ga/(pk)+a/p^2)$ for $p\ge p_0(k)$. Also
 $L_1(p)-1=\lceil\log_p(k^2x^{78})\rceil$.
\item \emph{Index sets, windows and weights.} The family has an unbounded set
 $\mathrm{Win}\subseteq\R$ of scales, and the set of indices $i$ at scale
 $x$ is nonempty for every $x\in\mathrm{Win}$ (for $x\notin\mathrm{Win}$ it
 may be empty). The set of tops, as a set of elements of the frame, is the
 same for all indices at a given scale. Each point carries integer amounts
 $a_t$ and top windows $b_t$ ($t\in\mathrm{TT}$); the \emph{amount window} of
 $L$ is $[a,\,a+a/(kx^{13})]$, and the family has a function $b_{\rm top}$
 bounding the top windows. A \emph{weight vector} is a pair
 $o=(\bar a_t,\bar b_t)_{t\in\mathrm{TT}}$ of integer families, and the
 index $i$ \emph{carries} $o$ if $a_t=\bar a_t$ and $b_t=\bar b_t$ for all
 tops $t$ of $L_{x,i}$.
\item \emph{Words.} The record has a finite set of \emph{words} $u$; each
 word has a prime $T(u)$, its level prime. These satisfy the \emph{norm
 condition}: $T(u)=\pi\bar\pi$ for some $\pi\in\mathcal O_E$ with
 $\pi\equiv1$ modulo $M\mathcal O_E$, where $M$ is a multiple of $N$ and of
 $4d$. (The words and their level primes are constructed in Section~\ref{i4:sec}.)
\end{enumerate}
\end{definition}

\begin{definition}[$\mathrm{LinkG}$]\label{i1:def-linkG}
For an instance and constants $C$, $b_o$, $\mathrm{LinkG}_{\rm at}(X,C,b_o)$
is the conjunction of five clauses:
\begin{enumerate}
\item \emph{(local lattice clause)} for every word $u$ and every place $w$
 with $p(w)\ne T(u)$:
 \begin{itemize}
 \item[(1a)] $\mathcal N_u\subseteq\mathfrak p_w^{\lfloor e(w)\tau(w)\rfloor}M^{\rm ram}_u$
  after localisation at $w$, that is, for every $y\in\mathcal N_u$ there is
  $s\in\mathcal O\setminus\mathfrak p_w$ with
  $sy\in\mathfrak p_w^{\lfloor e(w)\tau(w)\rfloor}M^{\rm ram}_u$;
 \item[(1b)] $\mathrm{len}_w(u)\le d^*(u)\,(\tau(w)+\mathrm{io}(w))$;
 \item[(1c)] if $\mathrm{VV}(p(w))=\emptyset$ then
  $\tau(w)=\mathrm{io}(w)=\mathrm{len}_w(u)=0$.
 \end{itemize}
 Here $\mathcal O$ is the ring of integers of the field of definition of
 the record, $\mathfrak p_w$ the prime of $w$, $e(w)$ its ramification
 index, and $\mathcal N_u$ (the norm lattice of $u$), $M^{\rm ram}_u$ (the
 ramified lattice of $u$), the local length $\mathrm{len}_w(u)$ and the
 integer $d^*(u)$ are lattice data of the column record which we do not
 define here. \emph{Clause~(1) is therefore stated informally; the exact
 statement is in \cite{CharacterFamily}.} It is not used below.
\item for every $p$ and $t_0\in\mathrm{VV}(p)$:
 $S_v(p,t_0)\le2m\,N_v(p,t_0)\,a/p$;
\item \emph{(prime-by-prime comparison)} for every place $w$ with
 $p=p(w)\le X''$ of inert type:
 \begin{equation}\label{i1:eq-linkG3}
 \sum_{j=1}^m\sum_{t\in\mathrm{TT}_w}\bigl(A^L_{j,t}(p)-\mathrm{lo}_G(p)\bigr)-m\,\mathrm{fl}_G(p)\le\sum_{j=1}^m\mathrm{gain}(j,w)-\tau(w),
 \end{equation}
 together with $0\le\mathrm{gain}(j,w)\le4|\mathrm{TT}|a/p$;
\item \emph{(served orbits)} for every $p$ with $\mathrm{Cand}(p)$ and
 $\Sigma(p)\le x^{-\varepsilon}$, the set of good orbit representatives at
 $p$ is contained in $\mathrm{VV}(p)$ (this holds by the definition of
 $\mathrm{VV}$);
\item \emph{(mean bound)} for every $p$ with $\mathrm{Cand}(p)$ and every
 good orbit representative $t_0$ at $p$:
 \begin{itemize}
 \item[(5a)] the mean of $\kappa^{\rm loc}_{p,t_0}(\bar a,\bar b)$ over the
  pairs $(\bar a,\bar b)$ of integer vectors indexed by
  $\{0,\dots,\operatorname{ord}h_P(p)-1\}$ which vanish outside the positions
  $\mathrm{pos}(p,t_0)$ of the tops in the orbit of $t_0$ and take values in
  the amount window, resp.\ the top window, at those positions, is at most
  $N_v(p,t_0)\,x^{-12}$;
 \item[(5b)] if $p<p_0(k)$, then $\kappa_{p,t_0}(\bar a,\bar b)\le N_v(p,t_0)\,x^{-12}$
  for all weight vectors $(\bar a,\bar b)$.
 \end{itemize}
 Here $\kappa_{p,t_0}(\bar a,\bar b)\ge0$ is the local weight of the orbit
 of $t_0$ at $p$ for the weight vector $(\bar a,\bar b)$ (so that
 $\Sigma(p)=m\sum_{t_0}\kappa_{p,t_0}$, the sum over the good
 representatives, for the weight vector carried by the point),
 $\kappa^{\rm loc}_{p,t_0}$ is its version on the orbit, and
 $\mathrm{pos}(p,t_0)$ is the set of positions of the tops in the orbit; we do
 not define these here. \emph{Clause~(5) is therefore stated informally;
 the exact statement is in \cite{CharacterFamily}.} Below it is used only
 through the consequence that, for every candidate $p$, the mean over the
 weight vectors of $\Sigma(p)$ is at most $m|\mathrm{TT}|x^{-12}$, which
 follows from (5a) by summing over the good representatives
 ($\sum_{t_0}N_v(p,t_0)\le|\mathrm{TT}|$) and multiplying by $m$.
\end{enumerate}
$\mathrm{LinkG}(\mathrm{Fm})$ means: there are $C$, $x_0$ and a function
$\delta$ with $\delta(x)\to0$ such that $\mathrm{LinkG}_{\rm at}(X_{x,i},C,b^{\delta(x)})$
holds for all $x\ge x_0$ and all $i$. Only clauses (2), (3) and (5) are used
below. Clauses (1) and (4) are recorded because the source of $\mathrm{LinkG}$
proves them together with the others. Clauses (2) and (3) are exact.
\end{definition}

\begin{definition}[Small shares, bad primes, top count]\label{i1:def-chosen}
Let $\mathrm{cands}=\{p<X':\mathrm{Cand}(p)\}$ and let $\mathcal P_{\rm lo}$ be
the set of primes $p<X'$ of the inert class with
$\operatorname{ord}h_P(p)\ge x^{\varepsilon/4}$.
\begin{itemize}
\item $\mathrm{ChW}(X)$:
 $\sum_{p\in\mathrm{cands}}\frac{\log p}{p}\bigl(2\cdot\mathbf 1_{x^{-\varepsilon}<\Sigma(p)}+4\Sigma(p)\bigr)\le x^{-9}$.
\item $\mathrm{ChBadp}(X)$:
 $\sum_{p\in\mathcal P_{\rm lo}}\frac{\log p}{p}\,\#\{t:\ t\text{ bad at }p\}\le|\mathrm{TT}|\,x^{-2+\varepsilon}$.
\item $\mathrm{ChTT}(L)$: $\delta_1n/2\le|\mathrm{TT}|\le2\delta_1n$.
\end{itemize}
$\mathrm{ChosenTT}$ and $\mathrm{ChosenBadp}$ mean that $\mathrm{ChTT}$ and
$\mathrm{ChBadp}$ hold at every point of the family with $x$ large.
\end{definition}

\begin{definition}[Losses at a point]\label{i1:def-losses}
$\mathrm{IOsum}(L)=\sum_wl(w)e(w)\mathrm{io}(w)$ and
$\mathrm{Ga}(L)=\sum_w2l(w)e(w)\bigl(\sum_j\mathrm{gain}^\flat(j,w)-\tau(w)\bigr)$,
where $\mathrm{gain}^\flat(j,w)=0$ if $p(w)$ is \emph{dropped}
(Definition~\ref{i1:def-nogain}) and $\mathrm{gain}^\flat=\mathrm{gain}$
otherwise. A \emph{reading} of $L$ presents the same data in additive
coordinates (Definition~\ref{i1:def-groupask}). $\mathrm{ASK}(N,a_t,t,w_j)$
is the ask function of the host arithmetic datum at cutoff $N$, where $N$ is
\emph{admissible} if $a_t/w_j(\rho)\le N$ for all $t$ and all $\rho\in\Phi$.
\end{definition}

\begin{remark}[The instance gains]\label{i1:rem-gain}
The gain $\mathrm{gain}_P(p,j)$ is defined explicitly. It is $0$ if $p>X''$
or if $p$ is of inert type with $h_P(p)=0$. Otherwise it is a sum over orbit
representatives $t_0$ of the shift $h_P(p)$ (or $2h_P(p)$ in the signal
class) acting on $H$:
\begin{itemize}
\item on unserved orbits, the positive part of a shifted ask series
 $\sum_{s}\sum_{e\ \text{odd}}(\mathrm{al}(s)/p^e-\mathrm{capU}_j(s,e))_+$
 minus a penalty $\mathrm{dl}(t_0)$;
\item on served orbits, the maximum of that quantity and
 $\gamma_p\sum_s\mathrm{al}(s)$ minus the non-served part, minus a penalty.
\end{itemize}
Here $\mathrm{al}(s)=a_s-8S_c$, $\gamma_p=p/(p^2-1)$, and $\mathrm{capU}$ is
the fill shifted by an explicit constant. The tax $\tau_P(p)$ is $0$ unless
$\mathrm{VV}(p)\ne\emptyset$. In that case it is an explicit sum of rounded
penalties $m(\lceil\mathrm{tv}(t_0)\rceil+D_p+1)$ over the served orbits.
Inequality~\eqref{i1:eq-linkG3} is therefore a comparison between two
explicitly defined quantities which agree up to explicit errors. The errors
are a shift error of $3\cdot4S_c\log_p(2a)$ per top, the penalties, and the
passage between two forms of the ask series. No lattice or valuation datum
enters~\eqref{i1:eq-linkG3}.
\end{remark}

\medskip\noindent We next record the facts on the host field that the instantiation uses.

\begin{lemma}[Frobenius in the host]\label{i1:lem-host}
Let $p$ be a prime.
\begin{enumerate}
\item If $p\nmid N$, then an arithmetic Frobenius $\sigma$ at a prime above $p$
 exists, and $\gamma(\sigma)=\mathrm{art}(p\bmod N)$. Its signs are
 $(\mathrm{ps}(p),(-d/p))$, and $\varpi$ of its $H$-component is the class of
 $p$ in $H_{e_2}$.
\item $p$ is of inert type if and only if $p\nmid N$ and $(-d/p)=-1$. In that
 case the signs of $\gamma(\sigma)$ are $(1,-1)$ in the signal class and
 $(-1,-1)$ in the inert class, and $p$ is unramified in $E$.
\item The prime $2$ is not of inert type. For every inert-class $p$, $p$ is odd
 and $\operatorname{ord}h_{e_2}(p)=\operatorname{ord}\varpi(h_P(p))$, which
 divides $\operatorname{ord}h_P(p)$.
\item The inertia at $d$ is $\{1,\iota_d\}$. For every $i$ there is exactly
 one $g\in G_D$, in the coset of $\mathrm{Frob}(d)^i$ modulo inertia, with
 $g\notin\Psi_D$.
\end{enumerate}
\end{lemma}

\begin{proof}
Let $l$ be a prime and write $N=l^vN'_l$ with $l\nmid N'_l$. In
$\Z[\zeta_N]$, modulo a prime $Q\mid l$, one has $\zeta^a\equiv\zeta^b$ iff
$a\equiv b\pmod{N'_l}$. Indeed, $(\zeta^{N'_l}-1)^{l^v}\equiv0$, and distinct
$N'_l$-th roots of unity stay distinct modulo $Q$ because $X^{N'_l}-1$ is
separable modulo $l$. Hence $\zeta\mapsto\zeta^u$ is inertial at $Q$ iff
$u\equiv1\pmod{N'_l}$, and it is an arithmetic Frobenius if
$u\equiv l\pmod{N'_l}$. These statements descend to $E$ through the
embedding: lift $\sigma$ to $\Q(\zeta_N)$ and correct the lift by the
decomposition group. This shows that $\sigma$ is inertial (resp.\ a
Frobenius) at $P\mid l$ iff $\gamma(\sigma)=\mathrm{art}(u)$ for some
$u\equiv1$ (resp.\ $u\equiv l$) modulo $N'_l$. For $p\nmid N$ we have
$N'_p=N$, which gives existence and the formula in~(1).

The quadratic Gauss sum $g_d$ satisfies $g_d^2=-d$ and $\tau_u(g_d)=(u/d)g_d$,
since $d\equiv3\pmod 4$. Hence the second sign of $\mathrm{art}(u)$ is
$(u/d)$. By quadratic reciprocity, $(p/d)=(-d/p)$ for odd $p$
(Theorem~\ref{i1:thm-qr}). The claim about $\varpi$ is an axiom of the host.
This proves~(1).

For (2): an odd prime with $(-d/p)=-1$ divides none of $Q_S,q,d$, because
$(-d/l)=1$ for every odd $l\mid Q_S$ and $(-d/q)=1$. The signs follow from~(1).
Unramifiedness follows from the cyclotomic embedding, since the primes
ramified in $\Q(\zeta_N)$ divide $N$.

For (3): $2\mid N$, so $2$ is ramified and not of inert type. The rest
follows from~(1), because a homomorphism divides orders.

For (4): we have $N'_d=Q_Sq$, so the inertia at $d$ is
$\mathrm{art}(\{u\equiv1\bmod Q_Sq\})\subseteq\{1,\iota_d\}$. Equality holds
because $u\equiv1\pmod{Q_Sq}$ with $u\equiv-1\pmod d$ has second sign
$(-1/d)=-1$. The two elements of the coset differ by $\iota_d$, whose second
sign is $-1$, so exactly one of them lies outside $\Psi_D$.
\end{proof}

\begin{theorem}[Quadratic reciprocity; Neukirch, \emph{Algebraic Number Theory}, I (8.6)]\label{i1:thm-qr}
For distinct odd primes $p,\ell$,
$\bigl(\frac p\ell\bigr)\bigl(\frac\ell p\bigr)=(-1)^{\frac{p-1}2\frac{\ell-1}2}$.
Moreover (first supplementary law) $\bigl(\frac{-1}p\bigr)=(-1)^{\frac{p-1}2}$
for every odd prime $p$.
\end{theorem}

\noindent The theorem is applied with $\ell=d\equiv3\pmod4$ and $p$ odd,
$p\nmid N$. These are distinct odd primes, as the theorem requires. The
identity $(p/d)=(-d/p)$ used in Lemma~\ref{i1:lem-host} needs both parts:
$(-d/p)=(-1/p)(d/p)=(-1)^{\frac{p-1}2}(d/p)$ and
$(d/p)(p/d)=(-1)^{\frac{p-1}2}$ since $d\equiv3\pmod4$. (Both the law and
the first supplementary law are stated in Neukirch I (8.6).)

\begin{lemma}[Words are split]\label{i1:lem-Tsplit}
For every word $u$ of the record, the prime $T(u)$ is split. In particular it
is not of inert type.
\end{lemma}

\begin{proof}
By the norm condition of the record, $T=\pi\bar\pi$ with $\pi\equiv1$ modulo
$M\mathcal O_E$, where $M$ is a multiple of $N$ and of $4d$. Then
$T\equiv1\pmod M$, because $\Z\cap M\mathcal O_E=M\Z$. Hence $T\nmid N$,
$T\equiv1\pmod4$ and $T\equiv1\pmod d$, so $(-d/T)=(T/d)=1$.
\end{proof}

\begin{definition}[No gain; dropped primes]\label{i1:def-nogain}
Put $Y_L=X'(1+C_r\rho_B)$ with the family constant $C_r\ge1$. A prime $p$ is
\emph{without gain} if one of the following holds:
\begin{itemize}
\item[(A)] $p<Y_L$, $p$ is of the inert class, and one of
 \begin{itemize}
 \item[(A1)] $\operatorname{ord}h_P(p)<x^{\varepsilon/4}$;
 \item[(A2)] $p\ge p_0(k)$ and $|p^i-X'|\le\rho_BX'$ for some $i\ge1$;
 \item[(A3)] $\mathrm{Cand}(p)$ and $x^{-\varepsilon}<\Sigma(p)$;
 \end{itemize}
\item[(B)] $p>X''$;
\item[(C)] $p$ is not of inert type;
\item[(D)] $p$ is of the inert class and $\operatorname{ord}h_P(p)$ is odd
 and $>1$.
\end{itemize}
A prime is \emph{dropped} if it is without gain and differs from every $T(u)$.
By Lemma~\ref{i1:lem-Tsplit}, a prime of inert type is dropped iff it is
without gain. Every family point satisfies
$\mathrm{gain}_P(p,j)=0=\tau_P(p)$ at dropped $p$.
\end{definition}

\begin{lemma}[Large Frobenius orders are even]\label{i1:lem-even}
If $\operatorname{ord}h_P(p)\ge x^{\varepsilon/4}$, then $\operatorname{ord}h_P(p)$
is even.
\end{lemma}

\begin{proof}
Let $h\in H$ have odd order. Then $\varpi(h)$ has odd order in a $2$-group,
so $\varpi(h)=0$. Hence $\varpi(oh)=0$ with $oh$ in the cyclic subgroup on
which $\varpi$ is injective. So $oh=0$, and the order of $h$ is at most
$o\le x^{\varepsilon/100}<x^{\varepsilon/4}$.
\end{proof}

\subsection{Instantiation at a family point}\label{i1:ssec-inst}

\begin{definition}[The parameter record of a family point]\label{i1:def-parL}
Let $C,b_o\in\R$. The record $\mathrm{par}_L(C,b_o)$ has the following
entries:
\begin{itemize}
\item $x,\varepsilon$ from the scale; $e_2,n,k,t$ from the ranks;
 $N=N(m)$; $d$ the host prime;
\item $p_0=\lceil p_0(k)\rceil$ with $p_0(k)=k^2x^{26}$;
 $\Lambda=\Lambda_G$; $r_G=\lceil8(4+8\varepsilon)/\varepsilon\rceil+1$;
\item $S_c$, $\eta$; $\rho=\rho_B$; $\vartheta=1/(kx^{13})$;
 $K_0=k^2x^{78}$;
\item $C\mapsto\max(C,0)$; $C_{\mathrm{ord}}=12$; $Y=Y_L$; and $b_o$.
\end{itemize}
\end{definition}

\begin{definition}[The local datum of a point]\label{i1:def-DX}
Let $P_s(L)$ be the set of primes of inert type below an explicit bound
$\mathrm{bndP}(L)>\max(a(1+\vartheta),X'')$. The local datum $\mathcal D_X$
over $(P_s(L),\{1,\dots,m\},\mathrm{TT},d)$ has the following entries.
\begin{itemize}
\item Classification: $p$ is \emph{served} if not dropped. A dropped $p$ is
 \emph{high} if $p>X''$, \emph{second} if $\mathrm{Cand}(p)$ and
 $x^{-\varepsilon}<\Sigma(p)$, and \emph{first} otherwise.
\item $T_p=\mathrm{TTP}(p)$; $I^{\rm ex}_{j,t,p}$ is the set of odd
 $i\le\log_pa_t+1$, with $\alpha_t=a_t$ and $w_{j,t,p}=w^{\rm part}_{j,t,p}$.
\item $I^d_{j,t}=\{1,\dots,\lfloor\log_da_t\rfloor+1\}$, with the weights of
 the $d$-term.
\item $\gamma_j(p)=0$ if $p$ is dropped and $\mathrm{gain}_P(p,j)$ otherwise;
 $\tau=\tau_P$.
\item $\mathrm{lo}=\mathrm{lo}_G$, $\mathrm{fl}=\mathrm{fl}_G$, and $L_1$.
\end{itemize}
\end{definition}

\begin{lemma}[Identification]\label{i1:lem-ids}
At a family point with $x>x_1$:
\begin{enumerate}
\item $\mathrm{Ga}_{\mathcal D_X}=\mathrm{Ga}(L)$;
\item $\mathrm{ASK}_{\mathcal D_X}=\mathrm{ASK}(L)$;
\item $\sum_{p\in\mathrm{cands}}\log p\,\mathrm{io}_P(p)=\mathrm{IOsum}(L)$.
\end{enumerate}
\end{lemma}

\begin{proof}
By the fills, the contribution of a place $w$ depends only on $p(w)$. It
vanishes unless $p(w)$ is of inert type and $\le X''$, because primes
$>X''$ are dropped and carry no gain or tax. Group the places by residue
characteristic. The places with $p(w)\le X''$ are all present, and their
weights over $p$ sum to $\log p$; this proves~(1). For~(2), primes
$p>a_t$ contribute nothing, since then $a_t<p^i$ and the weights are
$\ge w''\ge1$; also $a_t<\mathrm{bndP}(L)$. For~(3), use
$\mathrm{io}(w)=\mathrm{io}_P(p(w))$; non-candidates have
$\mathrm{VV}=\emptyset$ and hence $\mathrm{io}_P=0$.
\end{proof}

\begin{lemma}[Remaining hypotheses]\label{i1:lem-remaining}
At a family point with $x>x_1$, assume $\mathrm{LinkG}_{\rm at}(X,C,b_o)$,
$b_o\ge0$ and $p_0(k)\le X'$. Let $P=\mathrm{par}_L(C,b_o)$. Then (H2)--(H5)
hold for $\mathcal D_X$ with $\mathrm{cand}=\mathrm{cands}$,
$P_R=\mathcal P_{\rm lo}$ and $\mathrm{BAD}=$ ``bad at $p$''. Moreover
$\mathrm{io}_P(p)\le2\Sigma(p)\cdot2m|\mathrm{TT}|a/p$ for every $p$.
\end{lemma}

\begin{proof}
\emph{Key observation.} For $p\in P_s(L)$, being without gain implies
$\mathrm{VV}(p)=\emptyset$. Indeed, $\mathrm{VV}(p)\ne\emptyset$ forces
$\mathrm{Cand}(p)$ and $\Sigma(p)\le x^{-\varepsilon}$, and each clause of
Definition~\ref{i1:def-nogain} contradicts this:
\begin{itemize}
\item (A1), (A2) and (A3) contradict $\mathrm{Cand}(p)$ or the bound on
 $\Sigma(p)$;
\item (B) contradicts $p<X'\le X''=xnX'$;
\item (C) contradicts the inert class;
\item (D) contradicts Lemma~\ref{i1:lem-even}.
\end{itemize}
Conversely, let $p<X'$ be of the inert class and not without gain. Then
$p<Y_L$, and the failure of (A1)--(A3) gives $\mathrm{Cand}(p)$ and
$\Sigma(p)\le x^{-\varepsilon}$.

(H2). The amounts satisfy $0<a\le a_t\le a+a/(kx^{13})=a(1+\vartheta)$. The
exponents in $I^{\rm ex}$ are odd. The partner weights are size weights, hence
$\ge\min(b,b'/x,w'')=w''$. The weights of the $d$-term are $0$ or size
weights, hence $\ge0$. The inequality $\mathrm{fl}_G\le\mathrm{fl}^\star$ is
$C\le\max(C,0)$, since $\log_pa\ge0$.

For $\mathrm{lo}_G\le\mathrm{lo}^\star$, compare term by term:
\begin{itemize}
\item the $\mathrm{cp}$ terms agree, by the definition of $S_k$;
\item $\mathrm{pr}_G$ and $\mathrm{or}_G$ are nonzero only when $p<X'$, and
 then they are at most $\mathrm{pr}$ and $\mathrm{or}$; for the case
 $p<p_0(k)$ this uses $p_0(k)\le X'$;
\item $C\le\max(C,0)$ handles the sign of $C$.
\end{itemize}
Finally $K_0=k^2x^{78}=p^{\log_p(k^2x^{78})}\le p^{L_1(p)-1}$.

(H3). Let $p$ be served, i.e.\ not dropped. Then $p$ is not without gain, so
$p\le X''$. Choose a place $w$ over $p$ and apply clause~(3) of
$\mathrm{LinkG}_{\rm at}$. By the fills, \eqref{i1:eq-linkG3} at $w$ is
exactly~\eqref{i1:eq-served} for $\mathcal D_X$. For a non-served (dropped)
$p$, $\gamma_j(p)=0$ by definition. Also $\mathrm{VV}(p)=\emptyset$ by the key
observation, so $\tau(p)=0$. The statements on high and second-kind primes
are read off the classification.

(H4). Let $p$ be of first kind. Then $p$ is dropped, $p\le X''$, and $p$ is
not of second kind, so $p$ satisfies one of the clauses of
Definition~\ref{i1:def-nogain}:
\begin{itemize}
\item (A1): $p<Y_L\le\max(Y,X'')$ and, by Lemma~\ref{i1:lem-host}(3),
 $\operatorname{ord}h_{e_2}(p)\le\operatorname{ord}h_P(p)<x^{\varepsilon/4}$.
\item (A2): this is the band alternative, since $p\ge p_0(k)$ and $p$ is an
 integer, so $p\ge\lceil p_0(k)\rceil$.
\item (D): $\varpi(h_P(p))=0$ (the proof of Lemma~\ref{i1:lem-even}), so
 $\operatorname{ord}h_{e_2}(p)=1$.
\item (A3), (B) and (C) are excluded.
\end{itemize}

(H5). Let $p$ be served and $t\notin T_p$. Since $T_p=\mathrm{TT}$ for signal
$p$ or $p\ge X'$, $p$ is of the inert class with $p<X'$. By the converse
above, $\mathrm{Cand}(p)$ holds, so $p\in\mathcal P_{\rm lo}$ and
$\mathrm{VV}(p)$ is the set of good representatives. Hence $t$ is bad at $p$.

The io bound follows from clause~(2) of $\mathrm{LinkG}_{\rm at}$ and
$N_v\le|\mathrm{TT}|$.
\end{proof}

\begin{theorem}[Loss setting of a family point]\label{i1:thm-setting}
Under the hypotheses of Lemma~\ref{i1:lem-remaining}, assume also that
$\mathrm{TT}\neq\emptyset$, $\mathrm{ChBadp}(X)$ and $\mathrm{ChW}(X)$ hold.
Then $\mathrm{PCLH}$ holds for $P=\mathrm{par}_L(C,b_o)$ and $\mathcal D_X$,
with $\Sigma=\Sigma_X$ and $\mathrm{io}=\mathrm{io}_P$.
\end{theorem}

\begin{proof}
(H1): the elements of $P_s(L)$ are odd primes by
Lemma~\ref{i1:lem-host}(3). We have $\mathrm{lhs}_W\le x^{-9}$ by
$\mathrm{ChW}$, since $\theta=x^{-\varepsilon}$, and
$\mathrm{Bad}\le|\mathrm{TT}|x^{-2+\varepsilon}$ by $\mathrm{ChBadp}$. The io
bound and (H2)--(H5) are Lemma~\ref{i1:lem-remaining}.
\end{proof}

\begin{lemma}[Windows along the family]\label{i1:lem-windows}
Let $C\in\R$ and let $\delta(x)\to0$. Put $e=\max(40,C_t+3)$ and
$K=1+|C|+12$. There is $x_0$ such that for every $x\ge x_0$ and every index
$i$: $\mathrm{par}_{L_{x,i}}(C,b^{\delta(x)})\in\mathrm W(e,K)$ and
$p_0(k)\le X'$.
\end{lemma}

\begin{proof}
We use the elementary fact that $x^A\le\exp(cx^\delta)$ for $x$ large, where
$c,\delta>0$. Take $x$ beyond $x_1$, beyond $3$, $r_G$ and $C_r$, and so
large that the following hold:
\begin{itemize}
\item $x^7\le e^{x^{\varepsilon/4}/2}$, $x^2\le e^{x^{19\varepsilon_z/100}}$
 and $x^{36}\le e^{x^{\varepsilon/3}}$;
\item $\delta(x)<\frac12$.
\end{itemize}
The size inputs give the following:
\begin{itemize}
\item $x^2\le2^{e_2}\le n\le x^3$, from $n\le o2^{e_2}$ and
 $o\le x^{\varepsilon/100}$;
\item $x^4\le k\le x^5$, $m\le x$ and $N\le x^e$;
\item $x^4+1\le2k^2\le p_0(k)\le x^{36}$, so $p_0\le x^{37}\le x^e$ and
 $p_0(k)\le X'$;
\item $0\le\Lambda_G\le x^{2\varepsilon}$;
\item $S_c=4k\Delta\le x^7e^{x^{\varepsilon/4}/2}\le e^{x^{\varepsilon/4}}$;
\item $\eta\le e^{x^{\varepsilon/4}}$;
\item $\vartheta\le x^{-1}$, $K_0\ge x^2$ and $d\ge e^{\log d}\ge x^2$;
\item $1\le Y\le X'(1+x\cdot x^{-2})\le2X'$, using $C_r\le x$;
\item $b_o=e^{\delta(x)x^\varepsilon}\le e^{x^\varepsilon/2}$.
\end{itemize}
The remaining window conditions hold by definition.
\end{proof}

\medskip\noindent It remains to make the readings of a point explicit and to identify the ask series with the series of the abstract datum.

\begin{definition}[Additive coordinates, readings]\label{i1:def-groupask}
Let $M_D=(\Z/2)^2\times H$. Let $t_E:G_D\to M_D$ be the isomorphism that
sends each sign to $0$ or $1$ (swapping the two signs) and is the identity on
$H$. Put $\Phi=t_E(\Psi_D)$. Two maps make the frame explicit:
\begin{itemize}
\item the top places are $\tau_1(h)=((1,1),h)$;
\item $\rho_D(\tau,g)=c\cdot g\cdot\tau$, where $c=\mathrm{art}(-1)$.
\end{itemize}
For a weight $w$ on $G_D$, a prime $p$ and an amount $a$, put
\[
\operatorname{ask}_p(a,w,\tau)=\sum_{\substack{i\ge1\\ \mathrm{Frob}(p)^i\notin\Psi_D}}\bigl(a/p^i-w(\rho_D(\tau,\mathrm{Frob}(p)^i))\bigr)_+,
\]
and define $\operatorname{ask}_d$ in the same way, with $\mathrm{Frob}(d)^i$
replaced by the element $g$ of Lemma~\ref{i1:lem-host}(4).

The host arithmetic datum consists of:
\begin{itemize}
\item the inertia subgroups $I_l=t_E(\mathrm{art}(U_l))$, where
 $U_l=\{u\equiv1\bmod N'_l\}$;
\item the Frobenius classes $F_l=t_E(\mathrm{art}(f_l))$, where
 $f_l\equiv l\pmod{N'_l}$ and $f_l\equiv1\pmod{l^v}$.
\end{itemize}
Its ask function at cutoff $N$, for a top $t$ and a weight $w$ on $M_D$, is
\[
\mathrm{ASK}(N,a,t,w)=\sum_{p\le N\text{ inert}}2\log p\sum_{\substack{1\le i\le N\\ iF_p\notin\Phi}}\bigl(a/p^i-w(\mathrm{partner}(t,iF_p))\bigr)_+
+\log d\sum_{1\le i\le N}\bigl(a/d^i-w(\mathrm{partner}(t,g_i))\bigr)_+ .
\]
Here $\mathrm{partner}(t,y)=t_E(c)+y+((0,0),t)$, and $g_i$ is the unique
$g\notin\Phi$ with $g-iF_d\in I_d$. A \emph{reading} of $L$ is a tuple
$(T,a,w)$ that transports $\mathrm{TT}$, the amounts and the size weights to
$M_D$ along $t_E$. Thus $w_{j-1}=w^{\rm At}_j\circ t_E^{-1}$ on $\Phi$.
\end{definition}

\begin{lemma}[Ask identity]\label{i1:lem-askis}
For $x\ge\max(x_1+1,3)$, every reading of $L_{x,i}$ and every admissible cutoff
$N$ satisfy $\sum_{t\in T}\sum_{j<m}\mathrm{ASK}(N,a_t,t,w_j)=\mathrm{ASK}(L_{x,i})$.
\end{lemma}

\begin{proof}
First we check that the size weights are $\ge1$ on $\Psi_D$. The size
estimates give $1\le\sqrt b\le w''$, $w''\le b'$, $t\eta\le w''/x$ and
$b_t\le(1+\frac1{36})b'$. Together with $x\ge3$, these give $w''\ge1$,
$b'\ge1$ and the bounds $w''+\eta\le2b'$ and $b_t+b'/x+\eta\le2b'$. Each
size weight is at least $b$, $b_t+b'/x$ or $w''$, so it is $\ge1$ on
$\Psi_D$.

Next we use the following facts:
\begin{itemize}
\item the map $t_E$ carries $\rho_D(\tau_1t,g)$ to $\mathrm{partner}(t,t_Eg)$;
\item $\rho_D(\tau_1t,g)\in\Psi_D$ whenever $g\notin\Psi_D$, because the
 second sign of $c$ is $-1$;
\item by Lemma~\ref{i1:lem-host}, $t_E(\mathrm{Frob}(p))=F_p$ for inert-type
 $p$, and the inert primes of the datum are exactly the inert-type primes.
\end{itemize}
Hence each term of the abstract ask function equals the corresponding term of
$\operatorname{ask}_p(a_t,w^{\rm At}_j,\tau_1t)$. Terms with $p^i>N$, or with
$p>N$, vanish on both sides, because $a_t\le Nw$ and the weights are
positive. The $d$-terms match by Lemma~\ref{i1:lem-host}(4).

Finally, for an inert-type $p$ the signs of $\mathrm{Frob}(p)$ have second
entry $-1$, so $\mathrm{Frob}(p)^i\notin\Psi_D$ iff $i$ is odd. For odd $i$
the partner place of $(p,t,i)$ is $\rho_D(\tau_1t,\mathrm{Frob}(p)^i)$:
compare the $H$-components and the signs, using that $\{\pm1\}^2$ has
exponent~$2$. The fills then identify the weights, so
$\operatorname{ask}_p(a_t,w^{\rm At}_j,\tau_1t)=A^L_{j,t}(p)$, which
vanishes for $p>a_t$ because the weights are $\ge1$. Summing over $j$ and $t$
gives the identity.
\end{proof}

\subsection{The comparison $\mathrm{LinkG}$ and its reduction to local data}\label{i1:ssec-linkG-source}

The hypothesis $\mathrm{LinkG}$ is supplied by a reduction to local data. The
reduction runs in three layers: a statement for families
(Proposition~\ref{i1:prop-linkG}), a local statement about lattices
(Proposition~\ref{i1:prop-loc}), and $37$ admission clauses
(Proposition~\ref{i1:prop-adm}). The proofs, and what of them remains formally
verified only, are in \S\S\ref{i1:ssec-pf-adm}--\ref{i1:ssec-pf-linkG}.

\begin{definition}[Group compatibility; $\mathrm{Adm2F}$]\label{i1:def-groupcompat}
A family satisfies \emph{group compatibility} if there is $x_0$ such that for
every $x\ge x_0$ with $x>x_1$ and every index $i$ the column record
$D=D_{x,i}$, with its gathering ties and the lattice comparison of the
family at $(x,i)$, admits a \emph{group compatibility record}, consisting
of five fields:
\begin{enumerate}
\item[(herm)] a hermitian level structure on $D$;
\item[(forms)] for every word $u$, a form datum at $u$;
\item[(cmp)] a comparison tying (herm) and (forms) to the gathering ties
 and to the lattice comparison of the family;
\item[(ring)] for every word $u$, place $w$ with gathering index at
 $(u,w)$, and local index $v$: four ring isomorphisms
 $e_W:W(\mathbb F_{p^{n'_v}})\to W_v$, $e_Q:\Q_{p^{n'_v}}\to\Q_{q,v}$,
 $e_O:O^{\rm rec}_w\to O$, $e_K:K^{\rm rec}_w\to K_f$ (where $p=p(w)$,
 $(W_v,\Q_{q,v},O,K_f)$ are the rings of the local gathering at $(u,w)$, and
 $O^{\rm rec}_w$, $K^{\rm rec}_w$ the local coefficient ring of the record and
 its fraction field), with five compatibilities: $e_Q$ and $e_W$ commute
 with the inclusions $W\subset\Q_q$; $e_K$ and $e_O$ commute with
 $O\subset K_f$; $e_K$ carries the record's map $\Q_{p^{n'_v}}\to K^{\rm rec}_w$
 to the gathering's $\iota_v$; $e_O$ carries the record's map
 $W(\mathbb F_{p^{n'_v}})\to O^{\rm rec}_w$ to the gathering's $j_v$; and $e_O$
 sends the record's uniformiser to the gathering's $w$;
\item[(gat)] for every word $u$ and place $w$ with gathering index: there
 are a Hida setting $(W,O,\iota,j,\sigma,w_0)$ at $(p(w),n')$
 (Definition~\ref{i1:def-hida}), integers $n_V,r,m$ with $m+1=k$, dominant
 base weights $\lambda_{0,v}$ ($v<n_V$), a local identification $I$ of the
 level of (herm) at $p(w)$, a character $\psi$ of the part of that level away
 from $p(w)$ with values in $O^\times$, and $e\in\N$, such that the local
 comparison hypothesis holds for $I$ and the local gathering at $(u,w)$ is
 the gathering at $p(w)$ built from these data.
\end{enumerate}
The hermitian level structure, the form data, the tied comparison, the local
identifications with their comparison hypothesis, and the gathering built
from a Hida setting are objects of the construction of the families
(Section~\ref{i4:sec}) which we do not define here. \emph{This definition is therefore stated
informally; the exact statement is in \cite{CharacterFamily}.} The
predicate is the group compatibility of Section~\ref{i4:sec}.

$\mathrm{Adm2F}$, the \emph{second admissibility condition}, holds if there
is $x_0$ such that for every $x\ge x_0$ with $x>x_1$ and every index $i$:
\[
 \tfrac12\,n\delta_1\le|\mathrm{TT}|\le2n\delta_1
 \quad\text{and}\quad
 \sum_{p\in\mathcal P_{\rm lo}}\frac{\log p}{p}\,\#\{t:\ t\text{ bad at }p\}\le|\mathrm{TT}|\,x^{-2+\varepsilon}
\]
at $(D_{x,i},L_{x,i},X_{x,i})$, with $\mathcal P_{\rm lo}$ as in
Definition~\ref{i1:def-chosen}. In other words, $\mathrm{Adm2F}$ is the
conjunction of $\mathrm{ChosenTT}$ and $\mathrm{ChosenBadp}$; it is a
condition on the sizes of the top sets and on the bad primes, not a period
valuation bound.
\end{definition}

Closedness of a class of families for a property is
Definition~\ref{i4:def-closed}: $\mathcal C$ is closed for $P$ if there is
$\varepsilon_1>0$ such that for every $\varepsilon\in(0,\varepsilon_1]$ there
is a threshold datum $\kappa$ with $P(\mathrm{Fm},C_H)$ for every
$\mathrm{Fm}\in\mathcal C$ of exponent $\varepsilon$ and every $C_H$ at which
$\mathrm{Fm}$ is $\kappa$-large.

\begin{proposition}[$\mathrm{LinkG}$ from group compatibility]\label{i1:prop-linkG}
Let $\mathcal C$ be a class of families. If $\mathcal C$ is closed for group
compatibility and closed for $\mathrm{Adm2F}$, then $\mathcal C$ is closed
for $\mathrm{LinkG}$. Explicitly: there is $\varepsilon_1>0$ such that for
every $\varepsilon\in(0,\varepsilon_1]$ there is a threshold datum $\kappa$
such that $\mathrm{LinkG}(\mathrm{Fm})$ holds for every family
$\mathrm{Fm}\in\mathcal C$ of exponent $\varepsilon$ which is $\kappa$-large
at some $C_H$.
\end{proposition}

\noindent It is proved in \S\ref{i1:ssec-pf-linkG}, by an assembly of closedness
statements, from the local statement below (Proposition~\ref{i1:prop-loc}),
a first theorem on lattices (Proposition~\ref{i1:prop-first}), consequences
of group compatibility (Proposition~\ref{i1:prop-gc}) and a family statement
(Proposition~\ref{i1:prop-linkG-res}); the last three remain formally
verified. The comparison clause~\eqref{i1:eq-linkG3} is a lower bound for the
explicitly defined gains (Remark~\ref{i1:rem-gain}) and needs, besides
elementary bookkeeping, only the orbit-size clause of the good orbits,
proved in \S\ref{i1:ssec-pf-adm} (Lemma~\ref{i1:lem-osize}).

\subsubsection*{Local lattices}

\begin{definition}[Hida settings]\label{i1:def-hida}
A \emph{Hida setting} at $(p,n')$ consists of discrete valuation rings $W,O$
with fraction fields $\Q_q,K_f$, compatible ring maps $\iota:\Q_q\to K_f$,
$j:W\to O$, a ring automorphism $\sigma$ of $\Q_q$ and $w\in O$, such that: $p$
is a prime, $p\ge5$; $n'\ge2$ is even; $p$ is a uniformiser of $W$,
$|W/pW|=p^{n'}$ and $W$ is $p$-adically complete; for every $y\in W$ there is
$z\in W$ with $\sigma(y)=y^p+pz$; $\sigma^{n'}=1$; $j$ is finite; $O$ is
complete; and $w^2=p$. It is \emph{unramified quadratic} if moreover $w$ is
irreducible, $|O/wO|=p^{2n'}$, and every integer is a square in $O$.
\end{definition}

\begin{definition}[Weights, slots and the lattice statement]\label{i1:def-slots}
Fix a Hida setting, $m\ge0$, $k=m+1$, an integer $r\ge1$ and a \emph{dominant}
base weight $\lambda_0=(\lambda_{0,s})_{s<n'}$, $\lambda_{0,s}\in\Z^{k}$
antitone.
\begin{enumerate}
\item Put $\lambda_s=\lambda_{0,s}$ for $s$ even and
 $\lambda_s=(-\lambda_{0,s,k+1-i})_i$ for $s$ odd, $\lambda_{s,\rm last}$ its
 last entry, $\nu_s=\lambda_s-\lambda_{s,\rm last}$, and
 $A_r=\sum_s\sum_{i<r}\lambda_{s,i}$.
\item For a domain $A$ and $\nu\in\N^k$ antitone, the \emph{dual Weyl module}
 $\nabla_A(\nu)\subseteq A[X_{a,c}:a,c<k]$ is the module of polynomials of
 column weight $\nu$ fixed by the column substitutions
 $X_{a,j}\mapsto X_{a,j}+YX_{a,l}$, $l<j$. The row action of a matrix $g$ is
 $g\cdot P=P(X_{a,c}\mapsto\sum_bg_{ab}X_{b,c})$. For invertible $g$ over
 $K_f$ put $\nabla_\lambda(g)=\{F:\det(g)^{\lambda_{\rm last}}g\cdot
 F\in\nabla_O(\lambda-\lambda_{\rm last})\}$, and for an $O$-lattice $\Lambda$
 let $H_0(\lambda,\Lambda)$ be the intersection of the $\nabla_\lambda(g)$
 over all $g$ with column span $\Lambda$.
\item $\mathrm{tw}_s(y)=(\iota\sigma^s(y_i))_i$ for $y\in\Q_q^k$, and
 $\mathrm{ext}_s(L)$ is the $O$-span of $\mathrm{tw}_s(L)$. Put
 $V=\bigotimes_s\nabla_{K_f}(\nu_s)$ (polynomials in variables $X_{s,a,c}$),
 $\mathrm{HL}(L)=V\cap\bigotimes_sH_0(\lambda_s,\mathrm{ext}_s(L))$.
\item An \emph{edge} is a pair $e=(L'',L')$ of $W$-lattices with $L''\subseteq
 L'$, $pL'\subseteq L''$, $|L'/L''|=p^{n'r}$; $\mathrm{src}(e)=L''$,
 $\mathrm{tgt}(e)=L'$. Put
 $\mathrm{Le}_s(e)=\mathrm{ext}_s(L'')+\mathrm{span}_O\{w^{-1}\mathrm{tw}_s(py):y\in L'\}$,
 $\mathrm{HLe}(e)=V\cap\bigotimes_sH_0(\lambda_s,\mathrm{Le}_s(e))$,
 $\mathrm{Re}(e)=w^{-A_r}\mathrm{HLe}(e)$ and
 $\mathrm{Re}'(e)=w^{-A_r}\mathrm{HLe}(e)^\vee$, where
 $N^\vee=\{\varphi\in V^*:\varphi(N)\subseteq O\}$. The pairing is
 $\mathrm{be}(F,G)=G(p^{A_r}F)$.
\item For $t\ge0$: $\mathrm{ball}_t(L)=w^{-t}\mathrm{HL}(L)\cap\bigcap_{e=(L,\cdot)}\mathrm{Re}(e)$
 and $\mathcal N_t(L)=w^{-t}\mathrm{HL}(L)^\vee\cap\bigcap_{e=(\cdot,L)}\mathrm{Re}'(e)$.
\item The \emph{lattice statement} $\mathrm{LP}'(t,D,\kappa,c_\star)$ is: $4(D+1)\le
 c_\star$, and for every $W$-lattice $L$ there are edges $e_i=(L''_i,L)$,
 $i<K$, and $F_i\in\mathrm{ball}_t(L''_i)$ such that (a) $F_i\in
 w^{c_\star}\mathrm{Re}(e')$ for every edge $e'\neq e_i$ with source $L''_i$;
 (b) there are $n\ge(1-\kappa)\dim V$, an $O$-linear $u:O^n\to\mathcal N_t(L)$
 and $M\in O^{n\times K}$ with $\sum_iM_{ji}\,u(y)(p^{A_r}F_i)=p^Dy_j$ for all
 $y\in O^n$, $j<n$. $\mathrm{LP}'_1$ is the same with ``for every lattice''
 replaced by ``for some lattice''.
\end{enumerate}
\end{definition}

The good orbits and their parameters $u(\delta)=\lceil1/\delta\rceil$,
$M_0$, $G'$, $N_1$ are those of Definition~\ref{i1:def-goodorbit} below.
Let $\mathrm{TT}\subseteq\{0,\dots,n'-1\}$, $a_T,b_T:\{0,\dots,n'-1\}\to\Z$
and $c_{pp}\in\N$. A map $\lambda_0:\{0,\dots,n'-1\}\times\{0,\dots,m\}\to\Z$
is a \emph{base weight for $(\mathrm{TT},a_T,b_T,c_{pp})$} if for every
$s$:
\begin{itemize}
\item if $s\in\mathrm{TT}$: $\lambda_0(s,0)=a_T(s)$,
 $\lambda_0(s,m)=-b_T(s)$, and $|\lambda_0(s,i)|\le c_{pp}$ for $0<i<m$;
\item if $s\notin\mathrm{TT}$: $|\lambda_0(s,i)|\le c_{pp}$ for all $i$.
\end{itemize}
It is \emph{dominant} if $i\mapsto\lambda_0(s,i)$ is non-increasing for
every $s$. A \emph{dominant base weight for $(\mathrm{TT},a_T,b_T,c_{pp})$}
is a base weight for $(\mathrm{TT},a_T,b_T,c_{pp})$ which is dominant.

\begin{proposition}[the local statement]\label{i1:prop-loc}
Assume that a column record exists (as it does as soon as a family exists).
There is an absolute constant $C_0>0$ such that for every $C\ge C_0$ the
following holds. Let $p\in\N$, let $(W,O,\iota,j,\sigma,w)$ be an unramified
quadratic Hida setting at $(p,n')$, let
$m,r,L_1,L_2,c_{pp},c_0,D_p,c_\star\in\N$,
$\mathrm{TT}\subseteq\{0,\dots,n'-1\}$, $\delta_e,h,x,a,Y\in\R$, and
$\lambda_0:\{0..n'-1\}\times\{0..m\}\to\Z$, $a_T,b_T:\{0..n'-1\}\to\Z$.
Write $N_1=N_1(r,\delta_e,h,x)$ and $G'=G'(r,\delta_e,h,x)$. Assume:
\begin{enumerate}
\item (orbit hypothesis) $p\ge13$, $r\ge2$; $\mathrm{TT}$ is good at
 $(L_2,r-1,r,\delta_e,h,x)$; $2L_1\le L_2<u(\delta_e)$; $3G'\le n'$;
 $N_1^2\le r|\mathrm{TT}|$;
\item (weight hypothesis) $\lambda_0$ is a dominant base weight for
 $(\mathrm{TT},a_T,b_T,c_{pp})$ and, for $s\in\mathrm{TT}$,
 $a+c_{pp}\le a_T(s)\le2a-c_{pp}$, $0\le b_T(s)$, $b_T(s)+2c_{pp}\le c_0$;
\item (numerical conditions)
 \begin{gather*}
 r\le Y,\qquad CN_1^2\log a\le m+1,\qquad N_1+2r\le m+1-r,\qquad
 256(m+1)^2|\mathrm{TT}|Y^2\le p^{L_1-1},\\
 C|\mathrm{TT}|(m+1)^2\log_pa\le D_p,\qquad 8(D_p+1)\le c_\star,\\
 C(m+1)^3r^2|\mathrm{TT}|\,p^{5G'}\bigl(c_\star+r|\mathrm{TT}|c_0+2rn'c_{pp}\bigr)Y\le a .
 \end{gather*}
\end{enumerate}
Then for every integer
\[
t\ge Cr|\mathrm{TT}|\Bigl(CN_1^2\log a\cdot\frac{a}{p(m+1)}+a\,p^{-L_1}+c_0+(m+1)p^{G'}\Bigr)+2rn'c_{pp}
\]
the lattice statement $\mathrm{LP}'(t,D_p,\kappa,c_\star)$ of
Definition~\ref{i1:def-slots} holds, with $\kappa=C|\mathrm{TT}|r^2(m+1)/Y$.
\end{proposition}

\noindent Proof in \S\ref{i1:ssec-pf-loc}: it is reduced to the admission
clauses (Proposition~\ref{i1:prop-adm}) and to one closed-form implication
(Proposition~\ref{i1:prop-closed}), which remains formally verified.

\begin{proposition}[the admission clauses]\label{i1:prop-adm}
There are constants $C_{La}=C_k=C_0>0$, $C_{rk}>0$, $C_{deg}>0$ such that the
$37$ admission clauses $\mathcal A_1,\dots,\mathcal A_{37}$ listed in
\S\ref{i1:ssec-pf-adm} (Table~\ref{i1:tab-adm}) hold; only $\mathcal A_{27}$
involves the constants.
\end{proposition}

\noindent Proof in \S\ref{i1:ssec-pf-adm}. The orbit clauses, the
arithmetic of the coarse units, the lattice clauses, the basic
$\mathfrak{gl}_n$ clauses and the twisted-valuation clauses ($22$ of the $37$)
are proved there in full, modulo cited classical theorems, and so are
the projection theorems, the path clauses, flat tops and the reduction of
the test side and of the plane count; four precisely stated pieces remain
formally verified (end of \S\ref{i1:ssec-pf-adm}).

\begin{remark}[How the gain is bounded above]\label{i1:rem-upper}
Theorem~\ref{i1:thm-loss} is a lower bound for $g$. The source-slope
condition (P6), $\mu(M_i)+g+F-E_i-\ell_i\le\mu(N_i)$, holds because the column
lattices $N_i$ are defined by enlarging fixed lattices by the gains. The
upper constraint on $g$ is therefore the frame slope inequality (P8) for
these lattices, together with the divisor bound (P1). Its arithmetic content
is a period valuation bound for these lattices; this bound is not the
condition $\mathrm{Adm2F}$ of Definition~\ref{i1:def-groupcompat}, which
concerns the sizes of the top sets and the bad primes. Group compatibility is supplied by the
construction of the families (Section~\ref{i4:sec}), not here.
\end{remark}

\subsection{Proof of Theorem~\ref{i1:thm-loss}}\label{i1:ssec-proof}

\begin{proof}[Proof of the first part]
Let $C$, $x_G$ and $\delta$ be the data of $\mathrm{LinkG}$. Lemma~\ref{i1:lem-windows}
gives $e$, $K=13+|C|$ and $x_0$. Since $\varepsilon<\frac1{52}<\frac14$,
Theorem~\ref{i1:thm-elementary} applies to $(\varepsilon,e,K)$ and gives a
threshold $x_0'$. Take $x$ beyond all these thresholds, beyond those of
$\mathrm{ChosenTT}$, $\mathrm{ChosenBadp}$ and $\mathrm{LinkG}$, and beyond
that of Lemma~\ref{i1:lem-askis}. Let $i$ satisfy $\mathrm{ChW}(X_{x,i})$, put
$b_o=b^{\delta(x)}$ and $P=\mathrm{par}_L(C,b_o)$.

$\mathrm{ChTT}$ gives $\mathrm{TT}\neq\emptyset$, so
Theorem~\ref{i1:thm-setting} yields $\mathrm{PCLH}$. Since
$P\in\mathrm W(e,K)$, Theorem~\ref{i1:thm-elementary} gives
\[
\mathrm{ASK}_{\mathcal D_X}-\mathrm{Ga}_{\mathcal D_X}\le(3\max(C,0)+13)U,\qquad
\textstyle\sum_{p\in\mathrm{cands}}\log p\,\mathrm{io}_P(p)\le U,
\qquad U=|\mathrm{TT}|mb .
\]
Lemma~\ref{i1:lem-ids} converts these into the bounds~\eqref{i1:eq-main}
for $\mathrm{ASK}(L)-\mathrm{Ga}(L)$ and for $\mathrm{IOsum}(L)$, and
Lemma~\ref{i1:lem-askis} replaces $\mathrm{ASK}(L)$ by the reading sums.
Finally $3\max(C,0)+13\le 3K+13$.
\end{proof}

\begin{proposition}[Small-share weights exist]\label{i1:prop-chosenW}
Assume $\mathrm{LinkG}(\mathrm{Fm})$. Then $\mathrm{ChosenW}$ holds: there is
$x_0$ such that for every $x\ge x_0$ with $x\in\mathrm{Win}$ there are an
index $i_0$ at scale $x$ and a weight vector $o=(\bar a_t,\bar b_t)_{t\in\mathrm{TT}}$
of integers with $\bar a_t$ in the amount window of $L_{x,i_0}$ and
$b'\le\bar b_t\le b_{\rm top}(x)$ (Definition~\ref{i1:def-family-data}(8)),
such that $\mathrm{ChW}(X_{x,i})$ holds at every index $i$ at scale $x$
carrying $o$, that is, with $a_t=\bar a_t$ and $b_t=\bar b_t$ for all tops $t$.
\end{proposition}

\begin{proof}
Let $x\in\mathrm{Win}$ be large and take an index $i_0$ at scale $x$ (the
index set is nonempty since $x\in\mathrm{Win}$). Let $\mathcal W$ be the
finite nonempty set of weight vectors $o=(\bar a_t,\bar b_t)$ at $(x,i_0)$
over which the means of clause~(5) of $\mathrm{LinkG}$ are taken (amounts in
the amount window of $L_{x,i_0}$, top windows in the top window of the
family); its elements satisfy the window conditions in the statement.
Average over $\mathcal W$. Termwise,
$2\cdot\mathbf 1_{\theta<\Sigma}+4\Sigma\le(2/\theta+4)\Sigma$. Clause~(5) of
$\mathrm{LinkG}$ bounds the local means of $\kappa$ at each good representative
$t_0$ by $N_v(p,t_0)x^{-12}$; summing over the good representatives (with
$\sum_{t_0}N_v(p,t_0)\le|\mathrm{TT}|$) and multiplying by $m$ bounds the
mean over $\mathcal W$ of $\Sigma_o(p)$ by $m|\mathrm{TT}|x^{-12}$ for every
candidate $p$.
Lemma~\ref{i1:lem-cheb} gives
$\sum_{p<X'}\log p/p\le2x^{\varepsilon/3}$. Hence some $o$ satisfies
\[
\mathrm{lhs}_W\le(2x^\varepsilon+4)\,m|\mathrm{TT}|\,x^{-12}\cdot2x^{\varepsilon/3}\le x^{-9}.
\]
The last inequality uses $m\le x^{1/2+\varepsilon}$,
$|\mathrm{TT}|\le2x^{2+9\varepsilon}$, $\varepsilon<\frac1{52}$ and
$x^{1/4}\ge24$. The quantities $\Sigma_o$ and $\mathrm{cands}$ depend on the
index only through the weight vector it carries.
\end{proof}

\begin{proof}[Proof of the second part]
By Proposition~\ref{i1:prop-chosenW}, $\mathrm{ChosenW}$ holds. For $x\in\mathrm{Win}$
beyond its threshold and that of the first part, take $i_0$ and $o$ from
$\mathrm{ChosenW}$. Every index $i$ at scale $x$ carrying $o$ satisfies
$\mathrm{ChW}(X_{x,i})$, so the first part gives~\eqref{i1:eq-main} at $i$.
This is $\mathrm{LossF}(0)$ with constant $3K+13$.

The passage to (P7) uses the column package built from the family at a large
scale with the chosen weights; it is constructed in Section~\ref{i4:sec}, and its ties
($P_*$, $g$, $I_O$ at a lift $\omega$ are the reading sum, $\mathrm{Ga}$ and
$\mathrm{IOsum}$ of the underlying point $\mathrm{pt}(\omega)$, which carries
the chosen weights; the top count is $r=|\mathrm{TT}|$; the number of lifts
and the scale $(x,\varepsilon,b)$ are those of the point) are those of
Lemma~\ref{i4:Llem-L8}. Substituting them gives (P7) with
$c_{\rm loss}=3K+13$ and $x^{e_L\varepsilon}=1$.
\end{proof}

\subsection{The admission clauses: proof of Proposition~\ref{i1:prop-adm}}\label{i1:ssec-pf-adm}

\begin{table}[ht]
\centering\small
\begin{tabular}{@{}lll@{}}
\toprule
clause & content & here\\
\midrule
$\mathcal A_1$ & positions of tops in a good orbit & Lemma~\ref{i1:lem-A1}\\
$\mathcal A_2$ & sizes of good orbits along the scale & Lemma~\ref{i1:lem-A2}\\
$\mathcal A_3$ & $\mathrm{HL}(\mathrm{src}\,e)\subseteq\mathrm{Re}(e)$, $\mathrm{HL}(\mathrm{tgt}\,e)^\vee\subseteq\mathrm{Re}'(e)$ & Lemma~\ref{i1:lem-LF}\\
$\mathcal A_4$ & perfectness of $\mathrm{be}$ between $\mathrm{Re}(e)$, $\mathrm{Re}'(e)$ & Lemma~\ref{i1:lem-A4}\\
$\mathcal A_5$ & $\mathrm{LP}'_1\Rightarrow\mathrm{LP}'$ & Lemma~\ref{i1:lem-A5}\\
$\mathcal A_6$ & extremal projector & Lemma~\ref{i1:lem-A6}\\
$\mathcal A_7$ & invariant projector & Lemma~\ref{i1:lem-A7}\\
$\mathcal A_8$--$\mathcal A_{10}$ & integral projection theorems & Lemma~\ref{i1:lem-A8}\\
$\mathcal A_{11}$ & projection theorem for nested sets & Lemma~\ref{i1:lem-A11}\\
$\mathcal A_{12}$ & filtered projection criterion & Lemma~\ref{i1:lem-A12}\\
$\mathcal A_{13}$ & valuation lemma & Lemma~\ref{i1:lem-A13}\\
$\mathcal A_{15}$ & path basis & Lemmas~\ref{i1:lem-A15}, \ref{i1:lem-A15I}\\
$\mathcal A_{14}$, $\mathcal A_{16}$--$\mathcal A_{21}$ & Pieri path statements & Lemmas~\ref{i1:lem-A14}--\ref{i1:lem-A21}\\
$\mathcal A_{22}$ & powers of the determinant & Lemma~\ref{i1:lem-A22}\\
$\mathcal A_{23}$ & coarse units & Lemma~\ref{i1:lem-A23}\\
$\mathcal A_{24}$ & factorial units; degree truncation & Lemmas~\ref{i1:lem-A24}, \ref{i1:lem-A24c}\\
$\mathcal A_{25}$, $\mathcal A_{26}$ & derivative lemma; valuations over a DVR & Lemmas~\ref{i1:lem-A25}, \ref{i1:lem-A26}\\
$\mathcal A_{27}$ & plane count with constants & Lemma~\ref{i1:lem-A27}\\
$\mathcal A_{28}$ & Witt identification $W(\mathbb F_{p^{n'}})\cong W$ & Lemma~\ref{i1:lem-A28}\\
$\mathcal A_{29}$ & flat tops & Lemma~\ref{i1:lem-A29}\\
$\mathcal A_{30}$ & elementary divisor basis of an edge & Lemma~\ref{i1:lem-A30}\\
$\mathcal A_{31}$ & bases of $\mathrm{HL}(L)$ & Lemma~\ref{i1:lem-LF}\\
$\mathcal A_{32}$ & test side (duality) & Lemma~\ref{i1:lem-A32}\\
$\mathcal A_{33}$ & nested type bound & Lemma~\ref{i1:lem-A33}\\
$\mathcal A_{34}$ & frames of good orbits & Lemma~\ref{i1:lem-A34}\\
$\mathcal A_{35}$ & renaming and base change of twisted valuations & Lemma~\ref{i1:lem-A35}\\
$\mathcal A_{36}$ & Weyl dimension formula & Lemma~\ref{i1:lem-A36}\\
$\mathcal A_{37}$ & products of sections & Lemma~\ref{i1:lem-A37}\\
\bottomrule
\end{tabular}
\caption{The admission clauses.}\label{i1:tab-adm}
\end{table}

\subsubsection*{Good orbits}

\begin{definition}[Good orbits]\label{i1:def-goodorbit}
Let $n'\ge1$ and $T\subseteq\{0,\dots,n'-1\}$, viewed in $\Z/n'$. Put
$W(i,L)=\{(i+j)\bmod n':0\le j<L\}$ and $\mathrm{fwd}(y,i)=(i-y)\bmod
n'\in[0,n')$. Let $r,u\ge1$, $\ell=(6r+1)u$ and $t=\lfloor n'/\ell\rfloor$. For
$s<t$ and $1\le c\le6r+1$, \emph{cell} $c$ of \emph{tile} $s$ is
$[s\ell+(c-1)u,\,s\ell+cu)$. Tile $s$ is \emph{marked} if $T$ meets each odd
cell in no point and each even cell in exactly one point, which is even.
$T$ is \emph{good} with parameters $(L_2,K,r,u,M_0,G',N_1)$ if
\begin{itemize}
\item[(A)] $|T\cap W(i,L_2)|\le K$ for every $i<n'$;
\item[(Ti)] for every $s_0<t$ some tile $(s_0+j)\bmod t$, $j\le M_0$, is marked;
\item[(In)] $|T\cap W(i,G')|\le N_1$ for every $i<n'$.
\end{itemize}
For $r\ge1$, $0<\delta$, $h\ge0$, $x\ge3$ put $u(\delta)=\lceil1/\delta\rceil$,
$M_0=\lceil(h+3)\log x\cdot10^{6r+1}\rceil$, $G'=(M_0+5)(6r+1)u(\delta)$ and
$N_1=\lceil\max(6eG'\delta,(h+3)\log_2x)\rceil+1$; $T$ is \emph{good at
$(L_2,K,r,\delta,h,x)$} if it is good with $u=u(\delta)$ and these $M_0,G',N_1$.
\end{definition}

\begin{lemma}[Next even top]\label{i1:lem-nexteven}
Let $T$ be good with parameters $(L_2,K,r,u,M_0,G',N_1)$, $(M_0+1)\ell\le n'$
and $(M_0+5)\ell\le G'$. For every $i<n'$ there is $D\in[1,G']$ with
$(i+D)\bmod n'$ an even element of $T$.
\end{lemma}

\begin{proof}
We have $t\ge M_0+1$ and $n'<(t+1)\ell$. Let $s_i<t$ be the tile containing
$i$ ($s_i=t-1$ if $i\ge t\ell$). By (Ti) at $(s_i+1)\bmod t$ there is $j\le
M_0$ with tile $s=(s_i+1+j)\bmod t$ marked; its cell $2$ contains an even
$y\in T$ with $s\ell+u\le y<s\ell+2u$. If $s=s_i+1+j$, then $i<s\ell<y$ and
$D=y-i<(1+j)\ell+2u\le(M_0+5)\ell\le G'$. If $s+t=s_i+1+j$, then $s\le s_i$;
for $i<y$ take $D=y-i<2u$, and for $y\le i$ take $D=y+n'-i<(j+2)\ell+2u\le G'$.
\end{proof}

\begin{lemma}[Orbit positions and size]\label{i1:lem-osize}
Under the hypotheses of Lemma~\ref{i1:lem-nexteven}:
\begin{enumerate}
\item $n'\le|T|G'$ and $|T|G'\le N_1(n'+G')$;
\item for every $y\in T$ the least $d\ge1$ with $(y+d)\bmod n'$ an even element
 of $T$ satisfies $d\le G'$ and $|T\cap\{(y+j)\bmod n':0<j<d\}|\le N_1$;
\item if $K=r-1$, $L_1\le L_2<u$ and $3G'\le n'$, then for every $i\in T$
 fewer than $r$ even $y\in T$ have $1\le\mathrm{fwd}(y,i)\le L_1$.
\end{enumerate}
Moreover, with $\Lambda(r,h)=2(6r+1)((h+3)10^{6r+1}+6)$ and
$C_e=20\Lambda(r,h)$, for $0<\delta\le1$ and $x\ge3$:
$G'\le C_e\log x/\delta$, $N_1\le C_e\log x$, $N_1\le20G'\delta$, and every $T$
good at $(L_2,K,r,\delta,h,x)$ in $\Z/n'$ with $n'\ge3G'$ satisfies
$n'\delta/(C_e\log x)\le|T|\le40n'\delta$.
\end{lemma}

\begin{proof}
(1) The map $T\times[1,G']\to\Z/n'$, $(y,D)\mapsto y-D$, is onto by
Lemma~\ref{i1:lem-nexteven}. Grouping $T$ by $\lfloor y/G'\rfloor$ and applying
(In) to each group gives $|T|\le(\lfloor n'/G'\rfloor+1)N_1$. (2) The least
$d$ is at most the $D$ of Lemma~\ref{i1:lem-nexteven}, and the set lies in
$W((y+1)\bmod n',G')$. (3) Such $y$ lie in $W((i-L_1)\bmod n',L_2)$, which
contains at most $K=r-1$ points of $T$ by (A).

For the bounds put $\ell_x=\log x\ge1$, $R=6r+1$, $P=10^{-R}$. From $1\le
u\delta\le2$ and $M_0+5\le((h+3)/P+6)\ell_x$ we get $G'\delta=(M_0+5)Ru\delta\le
\Lambda\ell_x$ and $G'\delta\ge M_0+5\ge5$. As $\log_2x\le2\ell_x$ and
$G'\delta\ge M_0\ge(h+3)\ell_x$, the maximum in $N_1$ is $6eG'\delta$, so
$N_1<6eG'\delta+2\le16.8\,G'\delta$. This gives the three bounds. If
$n'\ge3G'$, then (1) applies; $n'\delta\le|T|G'\delta\le|T|C_e\ell_x$, and
$|T|G'\le16.8G'\delta\cdot\frac43n'$ gives $|T|\le22.4n'\delta$.
\end{proof}

\begin{lemma}[Clause $\mathcal A_2$]\label{i1:lem-A2}
For every $0<\varepsilon<1/52$ there are $C,x_0$ such that for $x\ge x_0$, with
$h=2+8\varepsilon$, $r=\lceil8(h+2)/\varepsilon\rceil+1$ and
$\delta=x^{-\varepsilon/8}$: $G'\le C\log x/\delta$, $N_1\le C\log x$, and for
every $n'\ge x^{\varepsilon/4}$ we have $3G'\le n'$ and every
$T\subseteq\Z/n'$ good at $(L_2,r-1,r,\delta,h,x)$ satisfies
$n'\delta/(C\log x)\le|T|\le40n'\delta$.
\end{lemma}

\begin{proof}
Take $C=C_e(r,h)$, $A=48\Lambda(r,h)/\varepsilon$ and
$x_0=\max(3,A^{16/\varepsilon})$. Put $y=x^{\varepsilon/16}\ge A$, so
$\delta=y^{-2}$. By the proof of Lemma~\ref{i1:lem-osize},
$G'\le\Lambda\log x\cdot y^2$, and $\log x\le16y/\varepsilon$ gives
$3\Lambda\log x\le Ay\le y^2$; hence $3G'\le y^4=x^{\varepsilon/4}\le n'$. The
rest is Lemma~\ref{i1:lem-osize}.
\end{proof}

\begin{lemma}[Clause $\mathcal A_1$]\label{i1:lem-A1}
Let $r\ge1$, $2L_1\le L_2<u$, $3G'\le n'$, and let $T$ be good at
$(L_2,r-1,r,\delta,h,x)$. Then:
\begin{enumerate}
\item every marked tile $s$ meets $T$ in $3r$ even points $f_0<\dots<f_{3r-1}$
 with $f_i+u<f_{i+1}$, $s\ell+u\le f_i$, $f_i+u<(s+1)\ell$, and
 $\mathrm{fwd}(y,f_i)>u$, $\mathrm{fwd}(f_i,y)>u$ for every other $y\in T$;
\item if $s,s'$ are cyclically consecutive marked tiles with enumerations
 $f,f'$, $a=f_r$, $b=f'_{2r-1}$ and $D_d=\mathrm{fwd}(a,b)$, then $D_d+1\le
 G'$ and $|T\cap\{(a+j)\bmod n':j\le D_d\}|\le N_1$;
\item for every $i\in T$, fewer than $r$ even $y\in T$ have
 $1\le\mathrm{fwd}(y,i)\le L_1$;
\item for every even $y\in T$ there is $d\in[1,G']$ with $(y+d)\bmod n'$ an
 even element of $T$, none for smaller $d'\ge1$, and
 $|T\cap\{(y+j)\bmod n':0<j<d\}|\le N_1$.
\end{enumerate}
\end{lemma}

\begin{proof}
(1) In a marked tile the even cell $2i+2$ contains exactly one point $f_i$,
which is even, and the odd cells are empty; so
$s\ell+(2i+1)u\le f_i<s\ell+(2i+2)u$, which gives the spacing, and every other
point of $T$ is more than $u$ away from $f_i$ in both cyclic directions
because $u\le f_i$ and $f_i+u<n'$. (3) and (4) are
Lemma~\ref{i1:lem-osize}(3),(2). (2) Let $g\ge1$ be the cyclic distance from
$s$ to $s'$; (Ti) at $(s+1)\bmod t$ gives $g\le M_0+1$. With $a$ in tile $s$
and $b$ in tile $s'$, $b-a$ (or $b+n'-a$ if $s'<s$) is less than
$(g+2)\ell\le(M_0+3)\ell<G'\le n'/3$, so it equals $D_d$; the window
$\{(a+j):j\le D_d\}$ lies in $W(a,G')$ and (In) applies.
\end{proof}

\begin{definition}[Frames]\label{i1:def-frames}
A \emph{frame system} on $T\subseteq\Z/n'$ consists of $J\ge1$ and
$\mathrm{fr}:\Z\times\{0,\dots,r-1\}\to\Z$. Its \emph{letters of segment $j$}
are the $y\in[\mathrm{fr}(j-1,0),\mathrm{fr}(j,r-1)]$ with $y\bmod n'\in T$,
and its \emph{L-letters} $\mathrm{LL}(j)$ those with
$\mathrm{fr}(j-1,r-1)<y<\mathrm{fr}(j,0)$. For $L$ and a letter $y$ of segment
$j$, $\mathrm{Gap}_L(j,y)$ says that $y-y'>L$ for the largest letter $y'<y$ of
segment $j$, if any. It is \emph{valid at $(L_1,N_1)$} if for all $j$, $i<r$:
(Fr1) $\mathrm{fr}(j+J,i)=\mathrm{fr}(j,i)+n'$; (Fr2)
$\mathrm{fr}(j,i)<\mathrm{fr}(j,i+1)$; (Fr3)
$\mathrm{fr}(j,r-1)<\mathrm{fr}(j+1,0)$; (Fr4) $\mathrm{fr}(j,i)$ is even with
$\mathrm{fr}(j,i)\bmod n'\in T$; (Fr5) no $y$ strictly between
$\mathrm{fr}(j,i)$ and $\mathrm{fr}(j,i+1)$ has $y\bmod n'\in T$; (Fr6)
$|\mathrm{LL}(j)|\ge r$, each of the first $r$ and of the last $r$ L-letters
$y$ is even with $\mathrm{Gap}_{L_1}(j,y)$, and
$\mathrm{Gap}_{L_1}(j,\mathrm{fr}(j,0))$; (Fr7) segment $j$ has at most $N_1$
letters; (Fr8) for every letter $y>\mathrm{fr}(j-1,r-1)$ of segment $j$, the
set $E$ of even letters of segment $j$ below $y$ has at least $r$ elements
and $y-o>L_1$ for its $r$-th largest element $o$.
\end{definition}

\begin{lemma}[Clause $\mathcal A_{34}$]\label{i1:lem-A34}
Let $n'$ be even, $r\ge2$, $L_1\le L_2<u$, $G'=(M_0+5)(6r+1)u$, $3G'\le n'$,
and let $T$ be good with parameters $(L_2,r-1,r,u,M_0,G',N_1)$. Then there is a
frame system on $T$, valid at $(L_1,N_1)$, such that each frame
$(\mathrm{fr}(j,i))_{i<r}$ reduces modulo $n'$ to the points in cells
$2r+2,2r+4,\dots,4r$ of a marked tile, and every marked tile carries in its
cell $2r+2$ the first point of some frame.
\end{lemma}

\begin{proof}
Call $b\in\Z$ a \emph{marked base} if the intervals
$I_c(b)=[b+(c-1)u,b+cu)$, $1\le c\le6r+1$, contain no lift of a point of $T$
for $c$ odd and exactly one, $\tau(b,c)$, which is even, for $c$ even. Then:
(a) $\tau(b,c)$ is even and reduces into $T$; (b) $\tau(b,c)<\tau(b,c')$ for
even $c<c'$; (c) every lift of a point of $T$ in $[b,b+\ell)$ is some
$\tau(b,c)$; (d) a lift $y<\tau(b,c)$ of a point of $T$ satisfies
$y+u<\tau(b,c)$, since the cell before $I_c(b)$ is empty; (e)
$\tau(b+n',c)=\tau(b,c)+n'$ when both are marked bases.

Let $s_0<\dots<s_{J-1}$ be the marked tiles ($J\ge1$ by (Ti)) and, for
$j=qJ+a$ with $0\le a<J$, $B(j)=qn'+s_a\ell$. Each $B(j)$ is a marked base
(translation by $qn'$ preserves parity as $n'$ is even), $B(j+J)=B(j)+n'$,
$B(j)+\ell\le B(j+1)$, and $B(j+1)-B(j)<(M_0+2)\ell$ by (Ti) as in
Lemma~\ref{i1:lem-A1}(2). Put $\mathrm{fr}(j,i)=\tau(B(j),2r+2+2i)$.

(Fr1) is (e); (Fr2) and (Fr4) are (b) and (a);
(Fr3): $\mathrm{fr}(j,r-1)<B(j)+\ell\le B(j+1)<\mathrm{fr}(j+1,0)$. (Fr5):
such a $y$ lies in $[B(j),B(j)+\ell)$, hence is some $\tau(B(j),c)$ with $c$
even strictly between $2r+2+2i$ and $2r+4+2i$, impossible. (Fr6): the $r$
points $\tau(B(j-1),4r+2+2k)$ and the $r$ points $\tau(B(j),2+2k)$, $k<r$,
are L-letters, so $|\mathrm{LL}(j)|\ge r$. By (c) and (b), the lifts of
points of $T$ in $(\mathrm{fr}(j-1,r-1),B(j-1)+\ell)$ are exactly the $r$
points $\tau(B(j-1),c)$, $c=4r+2,\dots,6r$ even, and those in
$[B(j),\mathrm{fr}(j,0))$ are exactly the $r$ points $\tau(B(j),c)$,
$c=2,\dots,2r$ even; every other L-letter lies in $[B(j-1)+\ell,B(j))$
(recall $B(j-1)+\ell\le B(j)$). Hence the first $r$ L-letters are the
$\tau(B(j-1),4r+2+2k)$ and the last $r$ are the $\tau(B(j),2+2k)$, $k<r$;
they are even by (a), and (d) gives a gap $>u>L_1$ to the previous letter; the same applies to
$\mathrm{fr}(j,0)=\tau(B(j),2r+2)$. (Fr7): the letters of segment $j$ lie in
an interval of length
$\mathrm{fr}(j,r-1)-\mathrm{fr}(j-1,0)<(M_0+3)\ell\le G'\le n'$, which reduces
injectively into a window $W(\cdot,G')$; apply (In). (Fr8): the $r$ frame
points of segment $j-1$ are even letters below $y$, so $|E|\ge r$; if
$y-o\le L_1$, the $r$ largest elements of $E$ would reduce to $r$ points of
$T$ in a window of length $L_2$, against (A) with $K=r-1$. The tile clauses
hold by construction.
\end{proof}

\subsubsection*{Coarse units}

For a prime $p$ write $v_p$ for the $p$-adic valuation on $\Q$. For
$f:\N\to\Q$ let $\Delta f(n)=f(n+1)-f(n)$ and $\mathrm{mah}(f,j)=\Delta^jf(0)$
(the Mahler coefficients); Newton's formula gives
$f(n)=\sum_{j\le n}\mathrm{mah}(f,j)\binom nj$. Say that $f$ \emph{decays at
$(\theta,s)$} if every $\mathrm{mah}(f,j)$ is $0$ or has $v_p\ge\theta j+s$.

\begin{lemma}[Binomial expansion]\label{i1:lem-binexp}
Let $P=p^{t}$ and let $a(\rho,M,l)$ be the coefficient of $X^M$ in
$(1+X)^\rho((1+X)^P-1)^l$. Then $\binom{\rho+Pe}M=\sum_{l\le M}a(\rho,M,l)\binom el$
for all $e\in\N$; $p^{\max(0,l-\lfloor M/P\rfloor)}\mid a(\rho,M,l)$; and
$\prod_{y=1}^\rho(1+Pe/y)=\binom{\rho+Pe}{\rho}$.
\end{lemma}

\begin{proof}
Expand $(1+X)^{\rho+Pe}=(1+X)^\rho(((1+X)^P-1)+1)^e$ binomially and compare
coefficients ($a(\rho,M,l)=0$ for $l>M$ since $X\mid(1+X)^P-1$). For the
divisibility write $(1+X)^P-1=X^P+ph$ with $h\in\Z[X]$ (as $p\mid\binom Pi$
for $0<i<P$) and show by induction on $l$, using
$f(X^P+ph)^{l+1}=X^Pf(X^P+ph)^l+p\,(fh)(X^P+ph)^l$, that $p^{\max(0,l-\lfloor
j/P\rfloor)}$ divides the $X^j$-coefficient of $f(X^P+ph)^l$ for every
$f\in\Z[X]$. The last identity is induction on $\rho$.
\end{proof}

\begin{lemma}[Clause $\mathcal A_{23}$]\label{i1:lem-A23}
Let $p\ge5$, $t\ge0$, $P=p^{t+1}$,
$\mathrm{cU}(P,n)=(Pn)!/((P!)^nn!)$ and
$\mathrm{dU}(P,n,\rho)=\prod_{y=1}^{\rho}(1+Pn/y)$. Then:
\begin{enumerate}
\item $\mathrm{cU}(P,n)=\prod_{i<n}\prod_{1\le y<P}(1+Pi/y)$;
\item with $H_P(s)=\sum_{1\le y<P}y^{-s}$ and $\mathrm{ps}_s(n)=\sum_{i<n}i^s$,
 the $p$-adic logarithm is
 $\log_p\mathrm{cU}(P,n)=\sum_{s\ge0}(-1)^sP^{s+1}H_P(s+1)\mathrm{ps}_{s+1}(n)/(s+1)$;
\item[(2')] for every $s\ge0$ there is $f\in\Q[X]$ of degree exactly $s+1$ with
 $f(n)=\mathrm{ps}_s(n)$ for all $n\in\N$;
\item $v_p(P^sH_P(s))\ge s$, and
 $v_p\bigl(P^sH_P(s)/s\bigr)-1/(p-1)\ge(s+1)/4$ for $s\ge1$;
\item $\mathrm{cU}(P,\cdot)^{\pm1}$ and $\mathrm{dU}(P,\cdot,\rho)^{\pm1}$
 ($\rho<P$) decay at $(1/4,0)$;
\item if $f_1,\dots,f_K$ decay at $(1/4,0)$, then for every $c\in\N$ there are
 $b_J\in\Q$ ($J\in\N^K$), $b_J=0$ unless $\sum J_k\le4c$, $v_p(b_J)\ge0$, such
 that $\prod_kf_k(n_k)-\sum_Jb_J\prod_k\binom{n_k}{J_k}$ is $0$ or has
 $v_p\ge c$ for all $n\in\N^K$.
\end{enumerate}
\end{lemma}

\begin{proof}
(1) $\mathrm{cU}(P,n+1)/\mathrm{cU}(P,n)=\prod_{y=1}^{P}(Pn+y)/(P!(n+1))
=\prod_{1\le y<P}(1+Pn/y)$.

(2) For $|u|_p<1$ put $L(u)=\sum_{k\ge0}(-1)^ku^{k+1}/(k+1)$. The limit
$Q_j(x)=(x^{p^j}-1)/p^j\to L(x-1)$ ($j\to\infty$, $|x-1|_p<1$) follows from the
binomial theorem and $\binom{p^j-1}{k}\to(-1)^k$ $p$-adically; the identity
$Q_j(xy)=p^jQ_j(x)Q_j(y)+Q_j(x)+Q_j(y)$ then gives additivity of $\log_p=L(\cdot-1)$.
Since $v_p(y)\le t$ for $y<P$, each $u=Pi/y$ has $|u|_p<1$; apply (1),
additivity, and exchange the finite sum with the series.

(2') Write $X^s=\sum_{j\le s}c_j\binom Xj$ with $c_j\in\Q$ and $c_s=s!\ne0$.
Since $\sum_{i<n}\binom ij=\binom n{j+1}$, the polynomial
$f=\sum_{j\le s}c_j\binom X{j+1}$ satisfies $f(n)=\mathrm{ps}_s(n)$, and it
has degree $s+1$ with leading coefficient $s!/(s+1)!=1/(s+1)$.

(3) Each term $P^s/y^s$ has $v_p\ge s(t+1)-st=s$, and all terms are
positive. With $v=v_p(s)$, $4v+2\le3\cdot5^v\le3s$ and $1/(p-1)\le1/4$, so
$v_p(P^sH_P(s)/s)-1/(p-1)\ge s-v-\frac14\ge\frac{s+1}4$.

(5) Take $b_J=\prod_k\mathrm{mah}(f_k,J_k)$ for $\sum J_k\le4c$; by Newton's
formula the difference is the sum of the terms with $\sum J_k>4c$, each $0$
or of valuation $>c$.

(4) We use three facts. (M2) If, for $j\le N$, $\mathrm{mah}(f,j)$ and
$\mathrm{mah}(g,j)$ are $0$ or of valuation $\ge\theta j+s$, resp.\
$\ge\theta j+s'$ ($\theta\ge0$), then $\mathrm{mah}(fg,N)$ is $0$ or of
valuation $\ge\theta N+s+s'$; induction on $N$ with
$\Delta(fg)=\Delta f\,g+f\,\Delta g+\Delta f\,\Delta g$. (M3) If $f$ decays at
$(\theta,s)$, $\theta\ge0$, every $f(n)$ is $0$ or has $v_p\ge s$; and if all
values are $0$ or of valuation $\ge s$, so are all Mahler coefficients.
(M4) If $u$ decays at $(\theta,\theta)$, $\theta>0$, then so does
$(1+u)^{-1}-1$ (expand $\sum_{i\le N}(-u)^{i+1}$ plus a remainder of
valuation $\ge(N+2)\theta$, using (M2), (M3)); and $g$ with $g(0)=1$,
$g(n+1)=g(n)(1+u(n))$ decays at $(\theta,0)$ (as $\Delta g=gu$ and (M2)).
Now let $u=\mathrm{dU}(P,\cdot,\rho)-1$, $\rho<P$. By
Lemma~\ref{i1:lem-binexp} with $M=\rho$, $\mathrm{mah}(u,j)=a(\rho,\rho,j)$ for
$j\ge1$ and $\mathrm{mah}(u,0)=0$, and $p^j\mid a(\rho,\rho,j)$ since
$\lfloor\rho/P\rfloor=0$; so $u$ decays at $(1/2,1/2)$. Hence
$\mathrm{dU}^{\pm1}$ decay at $(1/2,0)$ by (M4), and so does
$\mathrm{cU}(P,\cdot)^{\pm1}$, the running product with
$u=\mathrm{dU}(P,\cdot,P-1)-1$ (by (1)). Decay at $(1/2,0)$ implies decay at
$(1/4,0)$.
\end{proof}

\begin{lemma}[Clause $\mathcal A_{24}$, first two parts]\label{i1:lem-A24}
Let $p\ge5$ and $P=p^{t+1}$.
(1) For $\rho<P$: $(\rho+Pe)!=\rho!\,e!\,(P!)^e\,\mathrm{cU}(P,e)\,\mathrm{dU}(P,e,\rho)$.
(2) For finitely many $(t_k,\rho_k)$ with $\rho_k<p^{t_k+1}$, the products
$\prod_k\mathrm{cU}(P_k,e_k)\mathrm{dU}(P_k,e_k,\rho_k)$ are $p$-adic units, and
they, their inverses, and the quotients
$\prod_k(\rho_k+P_ke_k)!/\prod_k\rho_k!\,e_k!\,(P_k!)^{e_k}$ are congruent
modulo $p^c$ to integral binomial combinations of degree $\le4c$ in
$(e_k)_k$.
\end{lemma}

\begin{proof}
(1) is Lemma~\ref{i1:lem-binexp}:
$\mathrm{dU}(P,e,\rho)=(Pe+\rho)!/((Pe)!\rho!)$. (2) By the proof of
Lemma~\ref{i1:lem-A23}(4), $\mathrm{dU}-1$ decays at $(1/2,1/2)$, so by (M3)
$\mathrm{dU}(P,n,\rho)$ is a unit, and so is $\mathrm{cU}$ by (1) of that
lemma. The products decay at $(1/4,0)$ by (M2); apply
Lemma~\ref{i1:lem-A23}(5) and part (1).
\end{proof}

\subsubsection*{Lattice clauses}

\begin{lemma}[Clause $\mathcal A_{37}$]\label{i1:lem-A37}
Let $\lambda_0,\mu_0$ be dominant and $\Lambda_s$ ($s<n'$) $O$-lattices. If
$F\in\bigotimes_sH_0(\lambda_s,\Lambda_s)$ and
$G\in\bigotimes_sH_0(\mu_s,\Lambda_s)$, then
$FG\in\bigotimes_sH_0(\lambda_s+\mu_s,\Lambda_s)$.
\end{lemma}

\begin{proof}
Since $(\lambda_0+\mu_0)_s=\lambda_s+\mu_s$ and products of generators are
generators, it suffices to show
$H_0(\lambda,\Lambda)H_0(\mu,\Lambda)\subseteq H_0(\lambda+\mu,\Lambda)$. Let $g$
span $\Lambda$, $F\in\nabla_\lambda(g)$, $G\in\nabla_\mu(g)$. The row action is
a ring map, and a product of polynomials of column weights $\nu,\nu'$ fixed by
the column substitutions has column weight $\nu+\nu'$ and is fixed by them;
so $\det(g)^{\lambda_{\rm last}+\mu_{\rm last}}g\cdot(FG)\in
\nabla_O(\lambda+\mu-(\lambda+\mu)_{\rm last})$.
\end{proof}

\begin{lemma}[Clause $\mathcal A_4$, from $\mathcal A_{30}$]\label{i1:lem-A4}
For every edge $e$: $G\in\mathrm{Re}'(e)$ iff $\mathrm{be}(F,G)\in O$ for all
$F\in\mathrm{Re}(e)$; and $F\in\mathrm{Re}(e)$ iff $\mathrm{be}(F,G)\in O$ for
all $G\in\mathrm{Re}'(e)$.
\end{lemma}

\begin{proof}
Put $A=A_r$, so $p^Aw^{-A}=w^A$. For $G=w^{-A}\psi$, $\psi\in\mathrm{HLe}(e)^\vee$,
and $F=w^{-A}z$, $z\in\mathrm{HLe}(e)$: $\mathrm{be}(F,G)=\psi(z)$. Conversely
if $\mathrm{be}(\cdot,G)$ is integral on $\mathrm{Re}(e)$ then
$(w^AG)(z)=\mathrm{be}(w^{-A}z,G)\in O$, so $G\in\mathrm{Re}'(e)$. For the
second equivalence take the basis $(x_i)$ and integers $\pi_i\le A$ of clause
$\mathcal A_{30}$ (Lemma~\ref{i1:lem-A30}), with
$\mathrm{Re}(e)=\mathrm{span}_O\{w^{-(A-\pi_i)}x_i\}$ and
$\mathrm{Re}'(e)=\mathrm{span}_O\{w^{-(A-\pi_i)}p^{-\pi_i}x_i^*\}$. Pairing $F$
with these generators of $\mathrm{Re}'(e)$ gives $w^{A-\pi_i}x_i^*(F)\in O$,
i.e.\ $F\in\mathrm{Re}(e)$.
\end{proof}


Throughout, a Hida setting (Definition~\ref{i1:def-hida}) and a dominant
$\lambda_0$ are fixed, with $k=m+1$ and the notation of
Definition~\ref{i1:def-slots}. For a matrix $h$ over $\Q_q$ write $h^{(s)}$
for the matrix with entries $\iota\sigma^s(h_{ij})$; so $\mathrm{tw}_s(hy)=
h^{(s)}\mathrm{tw}_s(y)$. Note $w\ne0$ in $K_f$, as $w^2=p\ne0$. For an
$O$-lattice $\Lambda\subseteq K_f^k$ put
$H(\Lambda_\bullet)=\bigotimes_sH_0(\lambda_s,\Lambda_s)$, the $O$-span of the
products $\prod_sP_s$ with $P_s\in H_0(\lambda_s,\Lambda_s)$ a polynomial in
the variables $X_{s,a,c}$; write $H(L)$ for $\Lambda_s=\mathrm{ext}_s(L)$ and
$H(e)$ for $\Lambda_s=\mathrm{Le}_s(e)$, so that $\mathrm{HL}(L)=V\cap H(L)$ and
$\mathrm{HLe}(e)=V\cap H(e)$. For an $O$-submodule $A\subseteq V^*$ put
${}^\vee\!A=\{x\in V:\varphi(x)\in O\ \forall\varphi\in A\}$.

Two classical theorems are used.

\begin{theorem}[Standard basis of the dual Weyl module; J.~A.~Green,
\emph{Polynomial representations of $GL_n$}, LNM 830, Ch.~4; C.~De~Concini,
D.~Eisenbud, C.~Procesi, \emph{Young diagrams and determinantal varieties}, Invent.\ Math.\ 56 (1980), 129--165, \S1 (reference to be confirmed)]\label{i1:thm-stdbasis}
Let $R$ be an integral domain, $k\ge1$ and $\nu\in\N^k$ antitone. Then
$\nabla_R(\nu)$ is a free $R$-module with a basis of \emph{bideterminants}:
products $\prod_{j<J}\Delta_j$, where $\Delta_j=\det(X_{I_j(\alpha),\beta})_{\alpha,\beta<s_j}$
is a minor on the first $s_j\le k$ columns with strictly increasing row
indices $I_j(0)<\dots<I_j(s_j-1)$, and $\nu_t=\#\{j:t<s_j\}$. The basis is
indexed by the semistandard tableaux of shape $\nu$ (row lengths
$\nu_t$; the minors $\Delta_j$ are its columns) with entries $<k$, and does not depend on $R$; in particular it is also a basis of
$\nabla_K(\nu)$ for every field $K\supseteq R$.
\end{theorem}

\noindent It is applied with $R=O$, $R=K_f$, and in the proof of
Proposition~\ref{i1:prop-Fm} with $R=\Z$, $\Q$, $\Q_p$, a field $K$, an
arbitrary domain, and $O$ a local domain that need not be a discrete valuation
ring; the hypotheses are that $R$ is a domain and $\nu$ antitone, which hold.
(Green's account works over an infinite field; for an arbitrary domain, in
particular $\Z$, the reference is De Concini--Eisenbud--Procesi.)
Proposition~\ref{i1:prop-Fm} depends on this theorem throughout.

\begin{theorem}[Elementary divisors; N.~Bourbaki, \emph{Alg\`ebre},
Ch.~VII, \S4, no.~3, Th.~1; S.~Lang, \emph{Algebra}, III, Thm.~7.8]\label{i1:thm-elemdiv}
If $N\subseteq M$ are free modules of rank $k$ over a principal ideal domain
$A$, there are a basis $f_1,\dots,f_k$ of $M$ and $a_1,\dots,a_k\in A$ with
$N=\bigoplus_ia_iAf_i$.
\end{theorem}

\noindent It is applied to $L''\subseteq L'$ over the discrete valuation ring
$W$; both are free of rank $k$ as lattices.

\begin{lemma}[Transport of the base slot]\label{i1:lem-transport}
Let $\lambda\in\Z^k$ be antitone and $g,h\in GL_k(K_f)$.
\begin{enumerate}
\item The modules $\nabla_O(\nu)$ and $\nabla_{K_f}(\nu)$ are stable under the
 row action of matrices over $O$, resp.\ $K_f$.
\item If $\mathrm{colspan}_O(g)=\mathrm{colspan}_O(h)$ then
 $\nabla_\lambda(g)=\nabla_\lambda(h)$; hence
 $H_0(\lambda,\mathrm{colspan}_O(h))=\nabla_\lambda(h)$.
\item If $F\in\nabla_\lambda(h)$ then
 $\det(g)^{-\lambda_{\rm last}}g^{-1}\cdot F\in\nabla_\lambda(gh)$.
\item If $L=\mathrm{colspan}_W(h_0)$ then
 $\mathrm{ext}_s(L)=\mathrm{colspan}_O(h_0^{(s)})$; for $A\in GL_k(\Q_q)$,
 $\mathrm{ext}_s(AL)=A^{(s)}\mathrm{ext}_s(L)$ and
 $\mathrm{Le}_s(Ae)=A^{(s)}\mathrm{Le}_s(e)$, where $Ae=(AL'',AL')$.
\end{enumerate}
\end{lemma}

\begin{proof}
(1) The row substitution preserves column weights and commutes with the column
substitutions $X_{a,j}\mapsto X_{a,j}+YX_{a,l}$ that define $\nabla$.
(2) Each matrix has its columns in the $O$-span of the other's, so $h=gU$,
$g=hU'$ with $U,U'$ over $O$ and $UU'=1$; thus $\det U\in O^\times$. Since the
row action satisfies $(gU)\cdot F=U\cdot(g\cdot F)$, we get
$\det(h)^{\lambda_{\rm last}}h\cdot F=(\det U)^{\lambda_{\rm last}}U\cdot(\det(g)^{\lambda_{\rm last}}g\cdot F)$,
which lies in $\nabla_O(\nu)$ by (1) when $F\in\nabla_\lambda(g)$. By symmetry
the two modules agree, so the intersection defining $H_0$ is over equal
modules. (3) $(gh)\cdot(g^{-1}\cdot F)=h\cdot F$ and
$\det(gh)^{e}\det(g)^{-e}=\det(h)^e$. (4) $\mathrm{tw}_s$ is additive with
$\mathrm{tw}_s(cy)=\iota\sigma^s(c)\mathrm{tw}_s(y)$ and $\iota\sigma^s(W)\subseteq
O$; so the twists of $W$-combinations of the columns of $h_0$ are the
$O$-combinations of the columns of $h_0^{(s)}$. The other identities follow
from $\mathrm{tw}_s(Ay)=A^{(s)}\mathrm{tw}_s(y)$.
\end{proof}

\begin{lemma}[Edge bases]\label{i1:lem-edgebasis}
For every edge $e=(L'',L')$ there is $g\in GL_k(\Q_q)$ with
$L''=\mathrm{colspan}_W(g)$ and $L'=\mathrm{colspan}_W(gD_p)$, where
$D_c=\mathrm{diag}(c^{-1}[a<r]+[a\ge r])_{a<k}$. For every $s$,
$\mathrm{Le}_s(e)=\mathrm{colspan}_O(g^{(s)}D_w)$, an $O$-lattice, and
$\mathrm{ext}_s(L')=\mathrm{colspan}_O(g^{(s)}D_w^2)$.
\end{lemma}

\begin{proof}
By Theorem~\ref{i1:thm-elemdiv}, $L'$ has a basis $(f_c)$ with
$L''=\bigoplus p^{a_c}Wf_c$, $a_c\ge0$. As $pL'\subseteq L''$, $a_c\le1$, and
$|L'/L''|=p^{n'\sum a_c}=p^{n'r}$ gives $\sum a_c=r$. Order the basis so that
$a_c=1$ exactly for $c<r$ and let $g$ have columns $pf_c$ ($c<r$), $f_c$
($c\ge r$). By Lemma~\ref{i1:lem-transport}(4) and $\iota\sigma^s(p^{-1})=p^{-1}=w^{-2}$,
$\mathrm{ext}_s(L'')=\mathrm{colspan}_O(g^{(s)})$,
$\mathrm{ext}_s(L')=\mathrm{colspan}_O(g^{(s)}D_w^2)$, and the vectors
$w^{-1}\mathrm{tw}_s(py)$, $y\in L'$, span
$\mathrm{colspan}_O(g^{(s)}\mathrm{diag}(w^{-1}[a<r]+w[a\ge r]))$. Adding
$\mathrm{colspan}_O(g^{(s)})$ gives $\mathrm{colspan}_O(g^{(s)}D_w)$ since
$w\in O$.
\end{proof}

\begin{lemma}[Row weights and diagonal shifts]\label{i1:lem-rowwt}
(1) Each bideterminant $b$ of Theorem~\ref{i1:thm-stdbasis} is row
homogeneous: every monomial of $b$ has the same number $\mu_a$ of variables
in row $a$, and $\sum_{a<r}\mu_a\le\sum_{t<r}\nu_t$ for every $r$. Hence
$\mathrm{diag}(t)\cdot b=(\prod_at_a^{\mu_a})b$.
(2) Let $\lambda$ be antitone, $c\in O$ nonzero in $K_f$,
$A=\sum_{a<r}\lambda_a$ and $g\in GL_k(K_f)$. Then
$c^AH_0(\lambda,\mathrm{colspan}_O(g))\subseteq H_0(\lambda,\mathrm{colspan}_O(gD_c))$.
More precisely, if $(b_i)$ is the bideterminant basis of
$\nabla_O(\lambda-\lambda_{\rm last})$ with row weights $\mu_i$, then the
elements $X_i=\det(g)^{-\lambda_{\rm last}}g^{-1}\cdot b_i$ form an $O$-basis
of $H_0(\lambda,\mathrm{colspan}_O(g))$ and a $K_f$-basis of
$\nabla_{K_f}(\lambda-\lambda_{\rm last})$, and the elements $c^{\pi_i}X_i$ with
$\pi_i=\sum_{a<r}\mu_{i,a}+\rho\lambda_{\rm last}\le A$, $\rho=\min(r,k)$,
form an $O$-basis of $H_0(\lambda,\mathrm{colspan}_O(gD_c))$.
\end{lemma}

\begin{proof}
(1) Each monomial of a minor with row set $I_j$ uses exactly one variable from
each row of $I_j$. So $\mu_a=\#\{(j,\alpha):I_j(\alpha)=a\}$, and
$\sum_{a<r}\mu_a\le\sum_j\min(r,s_j)=\sum_{t<r}\nu_t$. The row action of
$\mathrm{diag}(t)$ multiplies $X_{a,c}$ by $t_a$.
(2) Put $e=\lambda_{\rm last}$, $\nu=\lambda-e$. By
Theorem~\ref{i1:thm-stdbasis} and Lemma~\ref{i1:lem-transport}(2), $F\in
H_0(\lambda,\mathrm{colspan}_O(g))=\nabla_\lambda(g)$ iff
$\det(g)^eg\cdot F=\sum_io_ib_i$ with $o_i\in O$, i.e.\ the $X_i$ are an
$O$-basis; they are a $K_f$-basis of $\nabla_{K_f}(\nu)$ since the $b_i$ are and
the row action of $g^{-1}$ is invertible. Now $(gD_c)^{-1}\cdot
b_i=g^{-1}\cdot(D_c^{-1}\cdot b_i)=c^{\sum_{a<r}\mu_{i,a}}g^{-1}\cdot b_i$ by (1),
and $\det(gD_c)=\det(g)c^{-\rho}$, so
$\det(gD_c)^{-e}(gD_c)^{-1}\cdot b_i=c^{\pi_i}X_i$; by the first part (for
$gD_c$) these form an $O$-basis of $H_0(\lambda,\mathrm{colspan}_O(gD_c))$.
Finally $\pi_i\le\sum_{t<r}\nu_t+\rho e=A$ by (1), so $c^AX_i=c^{A-\pi_i}(c^{\pi_i}X_i)$.
\end{proof}

\begin{lemma}[The local lattice facts; clauses $\mathcal A_3$ and $\mathcal A_{31}$]\label{i1:lem-LF}
\begin{enumerate}
\item For every lattice $L$, $H(L)\subseteq V$, so $\mathrm{HL}(L)=H(L)$; and
 $\mathrm{HL}(L)$ is the $O$-span of a $K_f$-basis of $V$ (clause $\mathcal
 A_{31}$). Moreover ${}^\vee(\mathrm{HL}(L)^\vee)=\mathrm{HL}(L)$.
\item For every edge $e=(L'',L')$: $w^{A_r}H(L'')\subseteq H(e)$ and
 $w^{A_r}H(e)\subseteq H(L')$.
\item (Clause $\mathcal A_3$.) $\mathrm{HL}(\mathrm{src}\,e)\subseteq\mathrm{Re}(e)$ and
 $\mathrm{HL}(\mathrm{tgt}\,e)^\vee\subseteq\mathrm{Re}'(e)$ for every edge $e$.
\end{enumerate}
\end{lemma}

\begin{proof}
(1) Write $\mathrm{ext}_s(L)=\mathrm{colspan}_O(h_s)$
(Lemma~\ref{i1:lem-transport}(4)). By Lemma~\ref{i1:lem-rowwt}(2),
$H_0(\lambda_s,\mathrm{ext}_s(L))$ has an $O$-basis $(X_{s,i})_i$ which is a
$K_f$-basis of $\nabla_{K_f}(\nu_s)$. Products of bases of polynomial spaces in
disjoint sets of variables form a basis of the tensor product; so the products
$x_i=\prod_sX_{s,i_s}$ are a $K_f$-basis of $V$ whose $O$-span is $H(L)$. If
$x\in{}^\vee(\mathrm{HL}(L)^\vee)$, the dual basis functionals lie in
$\mathrm{HL}(L)^\vee$, so all coordinates of $x$ are in $O$.
(2) Take $g$ as in Lemma~\ref{i1:lem-edgebasis} and $h_s=g^{(s)}$. Apply
Lemma~\ref{i1:lem-rowwt}(2) with $c=w$ to $h_s$ and to $h_sD_w$: for each $s$,
$w^{A_s}H_0(\lambda_s,\mathrm{ext}_s(L''))\subseteq H_0(\lambda_s,\mathrm{Le}_s(e))$
and $w^{A_s}H_0(\lambda_s,\mathrm{Le}_s(e))\subseteq H_0(\lambda_s,\mathrm{ext}_s(L'))$,
$A_s=\sum_{i<r}\lambda_{s,i}$. Multiply over $s$; $\sum_sA_s=A_r$.
(3) For $x\in\mathrm{HL}(L'')$, $w^{A_r}x\in H(e)\cap V=\mathrm{HLe}(e)$, so
$x\in\mathrm{Re}(e)$. For $\varphi\in\mathrm{HL}(L')^\vee$ and
$z\in\mathrm{HLe}(e)$, $w^{A_r}z\in\mathrm{HL}(L')$ by (2), so
$w^{A_r}\varphi\in\mathrm{HLe}(e)^\vee$ and $\varphi\in\mathrm{Re}'(e)$.
\end{proof}

\begin{lemma}[Clause $\mathcal A_{30}$: elementary divisor basis]\label{i1:lem-A30}
For every edge $e$ there are a $K_f$-basis $(x_i)$ of $V$ and integers
$\pi_i\le A_r$ such that, with $\mathrm{ord}_i=A_r-\pi_i$ and $(x_i^*)$ the
dual basis,
\begin{gather*}
\mathrm{span}_O\{x_i\}=\mathrm{HL}(\mathrm{src}\,e),\quad
\mathrm{span}_O\{p^{\pi_i}x_i\}=\mathrm{HL}(\mathrm{tgt}\,e),\quad
\mathrm{span}_O\{w^{\pi_i}x_i\}=\mathrm{HLe}(e),\\
\mathrm{span}_O\{w^{-\mathrm{ord}_i}x_i\}=\mathrm{Re}(e),\quad
\mathrm{span}_O\{w^{-\mathrm{ord}_i}p^{-\pi_i}x_i^*\}=\mathrm{Re}'(e).
\end{gather*}
\end{lemma}

\begin{proof}
Take $g$ as in Lemma~\ref{i1:lem-edgebasis}, $h_s=g^{(s)}$, and the elements
$X_{s,i}$ and exponents $\pi_{s,i}\le A_s$ of Lemma~\ref{i1:lem-rowwt}(2) for
$\lambda_s$ and $h_s$. For $c\in\{1,w,p\}$ the elements $c^{\pi_{s,i}}X_{s,i}$
form an $O$-basis of $H_0(\lambda_s,\mathrm{colspan}_O(h_sD_c))$, and
$\mathrm{colspan}_O(h_sD_c)$ is $\mathrm{ext}_s(L'')$, $\mathrm{Le}_s(e)$,
$\mathrm{ext}_s(L')$ for $c=1,w,p$ (Lemma~\ref{i1:lem-edgebasis}, $D_p=D_w^2$).
Put $x_i=\prod_sX_{s,i_s}$ and $\pi_i=\sum_s\pi_{s,i_s}\le A_r$. As in
Lemma~\ref{i1:lem-LF}(1) this gives the first three identities (all three
modules lie in $V$). Then $\mathrm{Re}(e)=w^{-A_r}\mathrm{span}_O\{w^{\pi_i}x_i\}$;
and the dual of $\mathrm{span}_O\{w^{\pi_i}x_i\}$ is
$\mathrm{span}_O\{w^{-\pi_i}x_i^*\}$, so
$\mathrm{Re}'(e)=\mathrm{span}_O\{w^{-A_r-\pi_i}x_i^*\}$ and
$w^{-A_r-\pi_i}=w^{-\mathrm{ord}_i}p^{-\pi_i}$.
\end{proof}

\begin{lemma}[Clause $\mathcal A_5$]\label{i1:lem-A5}
For all $t,D,\kappa,c_\star$: $\mathrm{LP}'_1(t,D,\kappa,c_\star)$ implies
$\mathrm{LP}'(t,D,\kappa,c_\star)$.
\end{lemma}

\begin{proof}
For $A\in GL_k(\Q_q)$ let $\rho(A)$ be the $K_f$-linear map of $V$ acting on
the factor $s$ by $P\mapsto\det(A^{(s)})^{-\lambda_{s,\rm last}}(A^{(s)})^{-1}\cdot P$;
it preserves $V$ by Lemma~\ref{i1:lem-transport}(1), and
$\rho(A^{-1})\rho(A)=1$. Multiplication by $A$ is a bijection from edges with
source (target) $L$ to edges with source (target) $AL$. By
Lemma~\ref{i1:lem-transport}(2)--(4), $\rho(A)\mathrm{HL}(L)\subseteq\mathrm{HL}(AL)$
and $\rho(A)\mathrm{HLe}(e)\subseteq\mathrm{HLe}(Ae)$; hence
$\rho(A)\mathrm{Re}(e)\subseteq\mathrm{Re}(Ae)$,
$\rho(A)\mathrm{ball}_t(L)\subseteq\mathrm{ball}_t(AL)$, and
$\varphi\mapsto\varphi\circ\rho(A)$ maps $\mathcal N_t(AL)$ into $\mathcal
N_t(L)$. Now let the data $e_i=(L''_i,L)$, $F_i$, $n$, $u$, $M$ be given at a
lattice $L$, and let $L_1=AL$ be any lattice (any two lattices are related
this way). Put $e'_i=Ae_i$ and $F'_i=\rho(A)F_i\in\mathrm{ball}_t(AL''_i)$.
For an edge $e'\ne e'_i$ with source $AL''_i$, $A^{-1}e'\ne e_i$ has source
$L''_i$, so $F_i\in w^{c_\star}\mathrm{Re}(A^{-1}e')$ and $F'_i\in
w^{c_\star}\mathrm{Re}(e')$. Keep $n,M$ and put $u'(y)=u(y)\circ\rho(A^{-1})$;
it takes values in $\mathcal N_t(L_1)$ (apply the dual statement to
$A^{-1}$), and $u'(y)(p^{A_r}F'_i)=u(y)(p^{A_r}F_i)$. So the statement holds at
$L_1$.
\end{proof}

\begin{lemma}[Clause $\mathcal A_{28}$: Witt identification]\label{i1:lem-A28}
There is a ring isomorphism $\varphi:W(\mathbb F_{p^{n'}})\to W$ with
$\sigma\circ\varphi=\varphi\circ F$, where $W(\cdot)$ denotes $p$-typical Witt
vectors and $F$ their Frobenius.
\end{lemma}

\begin{proof}
We use the standard facts on Witt vectors (ghost components
$w_j(x)=\sum_{i\le j}p^ix_i^{p^{j-i}}$; $w_j(Fx)=w_{j+1}(x)$; functoriality
$W(h)$ commuting with $F$; the ghost map is a ring isomorphism when $p$ is
invertible; and, for a perfect field $\kappa$ of characteristic $p$,
$W(\kappa)$ is a domain in which every nonzero $z$ is $p^mz'$ with $z'_0\ne0$
and $pz'$ has $0$-th coordinate $0$; reference to be confirmed).

\emph{Step 1.} By the Hida axiom, $\sigma$ restricts to a ring endomorphism
$\tau$ of $W$ with $p\mid\tau(x)-x^p$.

\emph{Step 2 (Dwork).} For $x\in W$ let $G(x)\in W(\Q_q)$ have ghost components
$\tau^j(x)$; then $G$ is a ring map with $G(x)_0=x$ and $G(\tau x)=FG(x)$. All
coordinates lie in $W$: if $G(x)_i=b_i\in W$ for $i\le j$, then
$p^{j+1}G(x)_{j+1}=\tau^{j+1}(x)-\sum_{i\le j}p^ib_i^{p^{j+1-i}}
=\sum_{i\le j}p^i(\tau(b_i)^{p^{j-i}}-b_i^{p^{j+1-i}})$, and $p\mid\tau(b_i)-b_i^p$
implies $p^{j-i+1}\mid\tau(b_i)^{p^{j-i}}-b_i^{p^{j+1-i}}$. So $G$ factors
through a ring map $s:W\to W(W)$ with $s(x)_0=x$ and $s\circ\tau=F\circ s$.

\emph{Step 3.} Let $\kappa=W/pW$, a field with $p^{n'}$ elements, and
$\psi=W(\mathrm{red})\circ s:W\to W(\kappa)$; $\psi(x)_0=\bar x$. It is
injective: $x=up^m$ with $u$ a unit gives $\psi(x)=\psi(u)p^m\ne0$. It is
surjective: every $z$ is $\psi(a)+pz'$ (lift $z_0$), so iterating, $z=\psi(S_n)+p^nr_n$
with $S_n=\sum_{i<n}a_ip^i$; by completeness $S_n\to S$, and $z-\psi(S)$ is
divisible by every $p^n$, hence $0$ by the structure of $W(\kappa)$.

\emph{Step 4.} With a field isomorphism $\kappa\cong\mathbb F_{p^{n'}}$ put
$\chi=W(\cong)\circ\psi$ and $\varphi=\chi^{-1}$; then $\chi\circ\tau=F\circ\chi$,
i.e.\ $\varphi\circ F=\tau\circ\varphi=\sigma\circ\varphi$.
\end{proof}

\begin{theorem}[Weyl's dimension formula; W.~Fulton, J.~Harris, \emph{Representation Theory}, Thm.~6.3(1) and Prop.~15.55]\label{i1:thm-weyl}
For $\nu\in\Z^k$ antitone, the number of semistandard tableaux of shape
$\nu-\nu_{\rm last}$ (row lengths $\nu_t-\nu_{\rm last}$) with entries $<k$ (equivalently, of
Gelfand--Tsetlin patterns with top row $\nu$) is
$\prod_{0\le i<j<k}(\nu_i-\nu_j+j-i)/(j-i)$.
\end{theorem}

\begin{lemma}[Clause $\mathcal A_{36}$]\label{i1:lem-A36}
$\dim_{K_f}V=\prod_{s<n'}\prod_{0\le i<j\le m}\frac{\lambda_{s,i}-\lambda_{s,j}+j-i}{j-i}$.
\end{lemma}

\begin{proof}
By the proof of Lemma~\ref{i1:lem-LF}(1), $\dim V=\prod_s\dim\nabla_{K_f}(\nu_s)$;
by Theorem~\ref{i1:thm-stdbasis} each factor is the number of semistandard
tableaux; apply Theorem~\ref{i1:thm-weyl} (the formula is invariant under
$\lambda_s\mapsto\lambda_s-c$).
\end{proof}

\subsubsection*{Representations of $\mathfrak{gl}_n$: clauses $\mathcal A_6$, $\mathcal A_7$}

\begin{definition}[$\mathfrak{gl}_n$-representations]\label{i1:def-gln}
Let $K$ be a field of characteristic $0$. A \emph{$\mathfrak{gl}_n$-representation}
on a $K$-vector space $V$ is a family of linear maps $E_{ab}$ ($a,b<n$) with
$[E_{ab},E_{cd}]=\delta_{bc}E_{ad}-\delta_{da}E_{cb}$, such that every vector lies
in a finite-dimensional subspace stable under all $E_{ab}$, and $V$ is the
sum of the weight spaces $V_\mu=\bigcap_a\ker(E_{aa}-\mu_a)$, $\mu\in\Z^n$.
\begin{enumerate}
\item A \emph{highest weight vector} of weight $\lambda$ is a nonzero $v\in
 V_\lambda$ with $E_{ab}v=0$ for $a<b$. $U$ is the monoid of words in the
 $E_{ab}$, $Uv$ the span of the $uv$. The \emph{isotypic part} $V[\lambda]$
 is the span of the $uv$, $v$ a highest weight vector of weight $\lambda$. An
 \emph{isotypic projector} for $\lambda$ is a linear map which is the identity
 on $V[\lambda]$ and $0$ on $V[\mu]$, $\mu\ne\lambda$; the \emph{invariant
 projector} $\mathrm{Inv}$ is the one for $\lambda=0$. A \emph{weight
 projector} for $\mu$ is the identity on $V_\mu$ and $0$ on $V_\nu$, $\nu\ne\mu$.
 $\mathfrak n^-V=\sum_{b<a}E_{ab}V$.
\item The \emph{Kostant generators} are $E^{(k)}_{ab}=E_{ab}^k/k!$ ($a\ne b$)
 and $\binom{E_{aa}}{k}$; $U_\Z$ is the monoid they generate. For a discrete
 valuation ring $O$ with fraction field $K$, a \emph{Kostant lattice} is a
 finitely generated $O$-submodule stable under $U_\Z$.
\item $\ell(\lambda)=\lambda_0-\lambda_{n-1}$; for a subspace $W$, $\ell(W)$ is
 the largest $k$ such that $\mu-\nu=k(e_a-e_b)$ for weights $\mu,\nu$ of $W$
 and some $a\ne b$ ($0$ if none). $N_0(n)=n(n-1)/2$ and
 $\delta(p,n,k)=\lfloor\log_p(k+n-1)\rfloor$.
\item A \emph{normal order} is a list of all pairs $(a,b)$, $a<b<n$, without
 repetition, such that for $a<b<c$ the pair $(a,c)$ lies between $(a,b)$ and
 $(b,c)$. For $\mu\in\Z^n$ with $s_{ab}:=\mu_a-\mu_b+b-a\ge1$ ($a<b$), a normal
 order and $T\in\N$, the \emph{extremal projector} is the ordered product over
 the list of
 $P_{(a,b)}=\sum_{k\le T}(-1)^k\binom{s_{ab}+k}{k}^{-1}E^{(k)}_{ba}E^{(k)}_{ab}$.
\end{enumerate}
\end{definition}

Throughout, $\mathrm{Pos}$ is the set of finite sums of the $e_a-e_b$, $a<b$,
and $\nu\preceq\lambda$ means $\lambda-\nu\in\mathrm{Pos}$. Two facts are
immediate from the relations: $E_{ab}V_\mu\subseteq V_{\mu+e_a-e_b}$, and weight
spaces for distinct weights are independent.

\begin{lemma}[Highest weight vectors]\label{i1:lem-hw}
Let $V$ be a $\mathfrak{gl}_n$-representation.
\begin{enumerate}
\item A nonzero finite-dimensional stable subspace contains a highest weight
 vector.
\item If $0\ne v\in V_\mu$ and $E_{ab}v=0$ for some $a\ne b$, then $\mu_a\ge\mu_b$
 and $E_{ba}^{\mu_a-\mu_b}v\ne0$, $E_{ba}^{\mu_a-\mu_b+1}v=0$. In particular
 highest weight vectors have dominant weights.
\item If $v$ is a highest weight vector of weight $\lambda$, then $Uv$ is the
 span of the $wv$ with $w$ a word in the $E_{cd}$, $d<c$; so $Uv\subseteq
 Kv+\sum_{\nu\preceq\lambda,\nu\ne\lambda}V_\nu$.
\end{enumerate}
\end{lemma}

\begin{proof}
(1) Take a weight of maximal height $\sum_a(n-a)\mu_a$; the raising operators
increase the height. (2) With $e=E_{ab}$, $f=E_{ba}$, $h=E_{aa}-E_{bb}$ one has
$[e,f]=h$, $[h,f]=-2f$, so $v_i=f^iv$ satisfies $hv_i=(\mu_a-\mu_b-2i)v_i$ and
$ev_{i+1}=(i+1)(\mu_a-\mu_b-i)v_i$; the nonzero $v_i$ have distinct
$h$-eigenvalues, so some $v_{j+1}=0\ne v_j$, and then $\mu_a-\mu_b=j$. (3) The
span $Z$ of the $wv$ is stable: for $a<b$ commute $E_{ab}$ past the letters of
$w$ using the relations, by induction on the length of $w$.
\end{proof}

\begin{lemma}[Complete reducibility]\label{i1:lem-cr}
Let $V$ be a finite-dimensional $\mathfrak{gl}_n$-representation. Every stable
subspace $W\subseteq V$ has a stable complement.
\end{lemma}

\begin{proof}
Let $C=\sum_{a,b}E_{ab}E_{ba}$ and $c(\nu)=\sum_a\nu_a^2+\sum_{a<b}(\nu_a-\nu_b)$.
A direct computation with the relations shows that $C$ commutes with every
$E_{cd}$, and that on $V_\mu$, $C=c(\mu)+2\sum_{a<b}E_{ba}E_{ab}$; so $C$ acts by
$c(\nu)$ on $U H_\nu$, where $H_\nu$ is the space of highest weight vectors of
weight $\nu$ (with $0$). If $\lambda,\mu$ are dominant and $0\ne\lambda-\mu\in\mathrm{Pos}$,
then $c(\mu)<c(\lambda)$: with $\beta=\lambda-\mu$,
$c(\lambda)-c(\mu)=\lambda\cdot\beta+\mu\cdot\beta+\sum_{a<b}(\beta_a-\beta_b)$,
the first two terms are $\ge0$ by dominance and the last is $>0$. Since $C$
acts by the scalar $c(\nu)$ on each $UH_\nu$, a polynomial in $C$ separates any
finite sum $\sum z_\nu$, $z_\nu\in UH_\nu$, into the parts with a given value of
$c(\nu)$.

\emph{Step 1: every stable $W$ equals $W_2=\sum_\mu U(W\cap H_\mu)$.} If not,
the image of $W$ in $V/W_2$ contains a highest weight vector (Lemma~\ref{i1:lem-hw}(1)),
of a dominant weight $\mu$; lift it to $x\in W\cap V_\mu$, so $x\notin W_2$
and $E_{ab}x\in W_2$ ($a<b$). Then $(C-c(\mu))x\in W_2$; using the separating
polynomial, replace $x$ by $x'=q(C)x$ with $q(c(\mu))=1$, so that $x'-x\in W_2$
and $(C-c(\mu))x'$ lies in $L=\sum_{c(\nu)=c(\mu)}U(W\cap H_\nu)$, where $C$ acts by
$c(\mu)$. For $a<b$, $z=E_{ab}x'\in W_2\cap V_{\mu+e_a-e_b}$. An element of $L$
of weight $\mu+e_a-e_b$ is $0$: a nonzero component of an element of $U(W\cap H_\nu)$ of weight
$\mu+e_a-e_b$ needs $\mu+e_a-e_b\preceq\nu$ (Lemma~\ref{i1:lem-hw}(3)), so
$0\ne\nu-\mu\in\mathrm{Pos}$ and $c(\nu)>c(\mu)$. Put
$w_1=(C-c(\mu))z=E_{ab}(C-c(\mu))x'$; it lies in $E_{ab}L$, so
$Cw_1=c(\mu)w_1$. A polynomial $r(C)$ with $r(c(\mu))=1$ maps $z\in W_2$ into
$L$, and $r(C)z$ has weight $\mu+e_a-e_b$, so $r(C)z=0$ by what was just said; but
$r(C)z\in z+Kw_1$. Hence $z\in Kw_1$, so $Cz=c(\mu)z$, $w_1=0$ and $z=0$. So $x'$ is a highest weight
vector of weight $\mu$ in $W$, whence $x'\in W_2$, a contradiction.

\emph{Step 2: construction.} For each $\lambda$ choose $Q_\lambda$ with
$V=(H_\lambda\cap W)\oplus Q_\lambda$, and put $W'=\sum_\lambda U(H_\lambda\cap
Q_\lambda)$. By Step 1 for $V$, $V=\sum U H_\mu$, and splitting each $v\in
H_\mu$ gives $V=W+W'$. If $W\cap W'\ne0$ it contains a highest weight vector $w$
of some weight $\lambda$; separating by $C$, $w$ lies in
$\sum_{c(\nu)=c(\lambda)}U(H_\nu\cap Q_\nu)$, and by Lemma~\ref{i1:lem-hw}(3)
and the inequality for $c$, only $\nu=\lambda$ contributes in weight $\lambda$,
and only through $H_\lambda\cap Q_\lambda$ itself. So $w\in H_\lambda\cap
W\cap Q_\lambda=0$, a contradiction.
\end{proof}

\begin{lemma}[Isotypic decomposition]\label{i1:lem-isodec}
For every $\mathfrak{gl}_n$-representation $V$, $V=\bigoplus_\lambda V[\lambda]$.
Hence isotypic projectors $e_\lambda$ exist, are unique, and every vector is
the finite sum $\sum_\lambda e_\lambda x$. A linear map commuting with all
$E_{ab}$ preserves each $V[\lambda]$ and commutes with each $e_\lambda$.
\end{lemma}

\begin{proof}
\emph{Spanning.} Let $x\in P$, $P$ finite-dimensional stable, and $W=P\cap\sum_\mu V[\mu]$.
A stable complement of $W$ in $P$ (Lemma~\ref{i1:lem-cr}) that is nonzero
would contain a highest weight vector, which lies in some $V[\mu]$, hence in
$W$. \emph{Independence.} First: if $w$ is a highest weight vector of weight
$\lambda$ lying in the span of the $uv$ with $v$ highest weight vectors of
weights in a set $\Omega$, then $\lambda\in\Omega$. (Inside a finite-dimensional
stable $P$ containing everything, take a stable complement $C$ of $Uw$; the
components in $Uw$ of the generating $v$ are $0$ or highest weight vectors of
the same weights, and they are not all $0$ since $w\notin C$; so $Uw$ contains a
highest weight vector $a$ of weight $\nu\in\Omega$. If $\nu\ne\lambda$ then
$\nu\prec\lambda$ by Lemma~\ref{i1:lem-hw}(3), and repeating the argument with
a stable complement of $Ua$ gives $\lambda\prec\nu$, absurd.) If
$V[\mu]\cap\sum_{\nu\ne\mu}V[\nu]\ne0$, it is stable and contains a highest
weight vector, whose weight would be both $\mu$ and $\ne\mu$. The last
assertion: a map $A$ commuting with the $E_{ab}$ sends $uv$ to $u(Av)$, and $Av$
is $0$ or a highest weight vector of the same weight; then $e_\mu Ax=Ax_\mu=Ae_\mu x$.
\end{proof}

\begin{lemma}[Clause $\mathcal A_6$: the extremal projector]\label{i1:lem-A6}
Let $V$ be a finite-dimensional $\mathfrak{gl}_n$-representation, $\mu$ with
$\mu_a-\mu_b+b-a\ge1$ for $a<b$, a normal order, $T\ge\dim V$, and $P$ the
extremal projector. For $x\in V_\mu$: $Px\in V_\mu$; if $x\in\mathfrak n^-V$ then
$Px=0$; if $E_{ab}x=0$ for all $a<b$ then $Px=x$.
\end{lemma}


\begin{proof}
(i) Each $E^{(k)}_{ba}E^{(k)}_{ab}$ preserves weight spaces, so $Px\in V_\mu$.
(ii) If $E_{ab}x=0$ then $P_{(a,b)}x=x$; this gives the last claim. (iii)
$P_{(a,b)}y-y\in E_{ba}V$ ($b>a$), so $Py-y\in\mathfrak n^-V$ for all $y$. (iv)
If $E_{a,a+1}y=0$ for all $a$ then $E_{ab}y=0$ for all $a<b$, by
$E_{ab}=[E_{a,b-1},E_{b-1,b}]$. (v) A vector $y\in V_\mu\cap\mathfrak n^-V$
killed by all raising operators is $0$: otherwise it is a highest weight
vector, $V[\mu]$ has a stable complement $X'$ (Lemma~\ref{i1:lem-cr}), and
$\mathfrak n^-V\subseteq (V[\mu]\cap W_\prec)+X'$ where $W_\prec$ is the sum of
the $V_\nu$ with $\nu\preceq\mu$, $\nu\ne\mu$ (the lowering operators map
$V[\mu]\subseteq\sum_{\nu\preceq\mu}V_\nu$ into $W_\prec$); so $y\in W_\prec\cap
V_\mu=0$.

(vi) \emph{$E_{a,a+1}Px=0$ for $x\in V_\mu$.} This is the heart of the proof,
and rests on three identities. Put
$\mathsf P^{ab}_s=\sum_{k\le T}\alpha_{s,k}E_{ba}^kE_{ab}^k$ with
$\alpha_{s,k}=(-1)^k/((s+1)\cdots(s+k)\,k!)$, so $P_{(a,b)}=\mathsf P^{ab}_{s_{ab}}$;
$E_{ab}$ is nilpotent with $E_{ab}^{\dim V}=0$, so $T$ plays no role.
\begin{itemize}
\item[(K)] \emph{Killing identity.} From the $\mathfrak{sl}_2$ relations
 $ef^{k+1}v=f^{k+1}ev+(k+1)(\gamma-k)f^kv$ ($hv=\gamma v$) and the
 telescoping identity
 $\sum_k\alpha_{s,k}(t_k+k(\gamma+k)t_{k-1})=\frac{s-\gamma}{s+1}\sum_k\alpha_{s+1,k}t_k$
 (valid when $t_T=0$, by $(j+1)\alpha_{s,j+1}=-\alpha_{s+1,j}/(s+1)$ and
 $\alpha_{s,j}=(s+j+1)\alpha_{s+1,j}/(s+1)$) one gets, for $hx=cx$,
 $E_{ab}\mathsf P^{ab}_sx=\frac{s-c-1}{s+1}\mathsf P^{ab}_{s+1}E_{ab}x$, and
 dually $\mathsf P^{ab}_s(E_{ba}z)=\frac{s+1-c}{s+1}E_{ba}\mathsf P^{ab}_{s+1}z$.
 For $x\in V_\nu$ and $s=\nu_a-\nu_b+1$ the first factor vanishes:
 $E_{ab}\mathsf P^{ab}_sx=0$.
\item[(YB)] \emph{Yang--Baxter identity.} For distinct $a,b,c$, $x\in V_\nu$,
 $s=\nu_a-\nu_b+\sigma+1$, $t=\nu_b-\nu_c+\tau+1$ ($\sigma,\tau\ge0$), whenever $s,t\ge0$:
 $\mathsf P^{ab}_s\mathsf P^{ac}_{s+t}\mathsf P^{bc}_tx=\mathsf P^{bc}_t\mathsf P^{ac}_{s+t}\mathsf P^{ab}_sx$.
 For $\sigma=\tau=0$: both sides $L,R$ are killed by $E_{ab}$ and $E_{bc}$, by
 (K) and the following \emph{pair identity} and its mirror image:
 \begin{itemize}
 \item[(PI)] for distinct $a,b,c$, $w\in V_\nu$ and integers $s,u\ge0$ with
  $u-(\nu_a-\nu_c)-1=s-(\nu_a-\nu_b)$:
  $E_{bc}\mathsf P^{ab}_s\mathsf P^{ac}_uw=\mathsf P^{ab}_s\mathsf P^{ac}_uE_{bc}w$;
  and (mirror) for $w\in V_\nu$ and integers $t,u\ge0$ with
  $u-(\nu_a-\nu_c)-1=t-(\nu_b-\nu_c)$:
  $E_{ab}\mathsf P^{bc}_t\mathsf P^{ac}_uw=\mathsf P^{bc}_t\mathsf P^{ac}_uE_{ab}w$.
 \end{itemize}
 \emph{Proof of (PI).} The relations give: $E_{bc}$ commutes with $E_{ba}$
 and with $E_{ac}$; $[E_{bc},E_{ab}]=-E_{ac}$, which commutes with $E_{ab}$;
 $[E_{bc},E_{ca}]=E_{ba}$, which commutes with $E_{ca}$. If $[x,y]=z$ and
 $[z,y]=0$ then $xy^{k}=y^{k}x+ky^{k-1}z$ (induction on $k$). Hence
 $[E_{bc},E_{ba}^kE_{ab}^k]=-kE_{ba}^kE_{ab}^{k-1}E_{ac}$ and
 $[E_{bc},E_{ca}^kE_{ac}^k]=kE_{ca}^{k-1}E_{ba}E_{ac}^k=kE_{ba}E_{ca}^{k-1}E_{ac}^{k-1}E_{ac}$;
 summing against $\alpha_{s,k}$ and using $k\alpha_{s,k}=-\alpha_{s+1,k-1}/(s+1)$
 (and $E^T=0$) gives the commutators
 \[E_{bc}\mathsf P^{ab}_s=\mathsf P^{ab}_sE_{bc}+\tfrac1{s+1}E_{ba}\mathsf P^{ab}_{s+1}E_{ac},\qquad
 E_{bc}\mathsf P^{ac}_u=\mathsf P^{ac}_uE_{bc}-\tfrac1{u+1}E_{ba}\mathsf P^{ac}_{u+1}E_{ac}.\]
 Put $y=\mathsf P^{ac}_uw\in V_\nu$. By the first commutator,
 $E_{bc}\mathsf P^{ab}_sy=\mathsf P^{ab}_sE_{bc}y+\frac1{s+1}E_{ba}\mathsf P^{ab}_{s+1}E_{ac}y$.
 By the second, $E_{bc}y=\mathsf P^{ac}_uE_{bc}w-\frac1{u+1}E_{ba}z$ with
 $z=\mathsf P^{ac}_{u+1}E_{ac}w\in V_{\nu+e_a-e_c}$. By (K) for the pair
 $(a,c)$ (eigenvalue $\nu_a-\nu_c$ of $E_{aa}-E_{cc}$), $E_{ac}y=\frac{u-(\nu_a-\nu_c)-1}{u+1}z$. By the
 dual form of (K) for the pair $(a,b)$ applied to $z$ (eigenvalue
 $\nu_a-\nu_b+1$ of $E_{aa}-E_{bb}$),
 $\mathsf P^{ab}_sE_{ba}z=\frac{s-(\nu_a-\nu_b)}{s+1}E_{ba}\mathsf P^{ab}_{s+1}z$.
 Collecting,
 $E_{bc}\mathsf P^{ab}_s\mathsf P^{ac}_uw-\mathsf P^{ab}_s\mathsf P^{ac}_uE_{bc}w
 =\frac{(u-(\nu_a-\nu_c)-1)-(s-(\nu_a-\nu_b))}{(s+1)(u+1)}E_{ba}\mathsf P^{ab}_{s+1}z=0.$
 For the mirror form, the \emph{dual operators} $E^\vee_{xy}=-E_{yx}$ satisfy
 the same relations, $V_\nu$ is the $(-\nu)$-weight space for them, and
 $\mathsf P^{xy}_s(E^\vee)=\mathsf P^{yx}_s(E)$ (the signs $(-1)^k$ cancel in
 $E^{\vee k}_{yx}E^{\vee k}_{xy}$). The first form for $E^\vee$, with
 $(a,b,c)$ replaced by $(c,b,a)$, weight $-\nu$ and parameters $(t,u)$, is
 the mirror form. Hence $L,R$ are also killed by $E_{ac}=[E_{ab},E_{bc}]$:
 indeed $E_{bc}L=\mathsf P^{ab}_s\mathsf P^{ac}_{s+t}E_{bc}\mathsf P^{bc}_tx=0$
 by (PI) applied to $w=\mathsf P^{bc}_tx$ (the hypothesis reads
 $t-(\nu_b-\nu_c)-1=0$) and (K), and $E_{ab}R=0$ likewise by the mirror form.
 Further $L-R\in\mathfrak n^-V$ for the copy of $\mathfrak{gl}_3$ on $a,b,c$
 (indices in the order $a,b,c$) by (iii), telescoping along the three
 factors. So $L=R$ by (v) for $\mathfrak{gl}_3$.
 \begin{itemize}
 \item[(T)] (a) \emph{Passage.} If (YB) holds for $(\sigma,\tau)$ for every
  finite-dimensional $\mathfrak{gl}_n$-representation (every $n$), all distinct
  $a,b,c$, every weight $\nu$ and every $x\in V_\nu$, then it holds for
  $(\sigma+1,\tau)$ for all of these with $s,t\ge0$. (b) \emph{Symmetry.} The same with
  $(\tau,\sigma)$ in place of $(\sigma+1,\tau)$ (again for all of these with $s,t\ge0$).
 \end{itemize}
 \emph{Proof of (T).} (a) Let $x\in V_\nu$, $s=\nu_a-\nu_b+\sigma+2$,
 $t=\nu_b-\nu_c+\tau+1$. Let $V'=V^n=V\otimes K^n$ with
 $(E'_{xy}w)_z=E_{xy}w_z+\delta_{zx}w_y$; it is a finite-dimensional
 representation, and $w_0$ with $x$ in position $a$ and $0$ elsewhere lies in
 $V'_{\nu+e_a}$. With $\nu'=\nu+e_a$ we have $s=\nu'_a-\nu'_b+\sigma+1$ and
 $t=\nu'_b-\nu'_c+\tau+1$, so (YB)$(\sigma,\tau)$ applies to $w_0$ (with
 $T'\ge\dim V'$; the truncation does not matter by nilpotency). Compare
 $a$-components. If $w_j=0$ and $j\notin\{x,y\}$ then
 $(\mathsf P'^{xy}w)_j=0$; for $x,y\ne a$ the $a$-component of
 $\mathsf P'^{xy}w$ is $\mathsf P^{xy}w_a$; and for a pair $(a,y)$ the
 $a$-component of $\mathsf P'^{ay}w$ is $\mathsf P^{ay}w_a$ when $w_y=0$
 (the subspace $\{w_y=0\}$ is stable under $E'_{ay}$ and on it
 $(E'_{ay}w)_a=E_{ay}w_a$, while $(E'_{ya}w)_a=E_{ya}w_a$ for every $w$).
 Moreover, if $w_b=w_c=0$ then $(E'_{bc}w)_b=E_{bc}w_b+w_c=0$,
 $(E'_{bc}w)_c=E_{bc}w_c=0$, and likewise for $E'_{cb}$; so
 $\mathsf P'^{bc}$ preserves $\{w_b=w_c=0\}$.
 Following the factors from the right, starting from $w_0$ with $b$- and
 $c$-components $0$: on the left side $\mathsf P'^{bc}$ keeps both $0$,
 $\mathsf P'^{ac}$ is applied to a vector with $c$-component $0$, and
 $\mathsf P'^{ab}$ to one with $b$-component $0$ (it was $0$ before
 $\mathsf P'^{ac}$, which does not involve $b$); on the right side the same
 holds with the roles of the factors exchanged. So the $a$-components of the
 two sides are $\mathsf P^{ab}_s\mathsf P^{ac}_{s+t}\mathsf P^{bc}_tx$ and
 $\mathsf P^{bc}_t\mathsf P^{ac}_{s+t}\mathsf P^{ab}_sx$, which is
 (YB)$(\sigma+1,\tau)$. (b) Apply (YB)$(\sigma,\tau)$ to $E^\vee$ with
 $(a,b,c)$ replaced by $(c,b,a)$ and weight $-\nu$: the parameter attached to
 $(c,b)$ is $(-\nu)_c-(-\nu)_b+\sigma+1$ and that attached to $(b,a)$ is
 $(-\nu)_b-(-\nu)_a+\tau+1$; renaming $\sigma\leftrightarrow\tau$ and using
 $\mathsf P^{xy}(E^\vee)=\mathsf P^{yx}(E)$, the identity becomes (YB)$(\tau,\sigma)$
 for $E$ read from right to left. \emph{Conclusion:} from
 (YB)$(0,0)$, (a) gives all (YB)$(\sigma,0)$; if (YB)$(\sigma,\tau)$ holds for
 all $\sigma$, then (b), (a), (b) give (YB)$(\sigma,\tau+1)$ for all
 $\sigma$. So (YB) holds for all $\sigma,\tau\ge0$.
\item[(C)] \emph{Connectivity.} Any two normal orders are connected by moves:
 exchanging two adjacent pairs with disjoint entries, or replacing adjacent
 $(a,b),(a,c),(b,c)$ ($a<b<c$) by $(b,c),(a,c),(a,b)$ or conversely. In
 particular every normal order is connected to one beginning with $(a,a+1)$.
 (Proved in full below.)
\end{itemize}
Exchanges do not change $P$ (the factors commute), and braid moves do not
change $Px$ for $x\in V_\mu$ by (YB) with $\sigma=b-a-1$, $\tau=c-b-1$ (then
$s_{ac}=s_{ab}+s_{bc}$; the factors to the right of the moved block send $x$
into $V_\mu$). By (C) we may therefore assume that the list begins with
$(a,a+1)$, so $Px=P_{(a,a+1)}y$ with $y\in V_\mu$ the image of $x$ under the
remaining factors; and $E_{a,a+1}P_{(a,a+1)}y=0$ by (K), since
$s_{a,a+1}=\mu_a-\mu_{a+1}+1$.

\emph{Proof of (C).} We prove the form used: every normal order can be
moved to one beginning with $(a,a+1)$. Call a list of pairs \emph{$j$-normal}
if it contains each pair $(x,y)$ with $j\le x<y<n$ exactly once and no other
pair, and for $j\le x<y<z$ the pairs $(x,y),(x,z),(y,z)$ occur in this order or
in the reverse order. Moves act on contiguous sublists, so moves of a prefix
are moves of the whole list; moves are reversible; and the reverse of a move
is a move (the reverse of a $j$-normal list is $j$-normal).

\emph{Claim.} In a $j$-normal list, if $(b,c),(j,c),(j,e)$ are consecutive with
$j<b$, then $e=b$. Let $p$ be the position of $(b,c)$; by the triple
$(j,b,c)$, $(j,b)$ comes after $(j,c)$, so if $e\ne b$ it is at position $\ge
p+3$. If $e<b$: $(j,e,c)$ puts $(e,c)$ before $(j,c)$, hence before $p$, and
$(e,b,c)$ then puts $(e,b)$ before $(e,c)$; but $(j,e,b)$, with $(j,e)$ at
$p+2$ before $(j,b)$, puts $(e,b)$ after $p+2$. If $b<e<c$: $(j,b,e)$ puts
$(b,e)$ before $(j,e)$, hence before $p$; $(j,e,c)$ puts $(e,c)$ before
$(j,c)$, hence before $p$; but $(b,e,c)$ with $(b,e)$ before $(b,c)$ puts
$(e,c)$ after $p$. If $e>c$: $(b,e)$ is before $p$ as before; $(j,c,e)$ puts
$(c,e)$ after $p+2$; but $(b,c,e)$ with $(b,e)$ before $(b,c)$ puts $(c,e)$
before $(b,e)$. Each case is contradictory.

\emph{Sorting.} A $j$-normal list can be moved, through $j$-normal lists, to
$\ell_F\ell_G$ with all pairs of $\ell_F$ of first entry $j$ and none in
$\ell_G$. Induct on the sum $m$ of the positions of the pairs with first entry
$j$. Otherwise the list is $\ell_1\,(b,c)\,(j,d)\,\ell_F\ell_G$ with $b\ne j$
and $\ell_F,\ell_G$ as described (look at the last place where a pair with
first entry $\ne j$ is followed by one with first entry $j$). The case $d=b$ is
impossible ($(j,c)$ would have no room between $(j,b)$ and $(b,c)$). If $d=c$,
the triple $(j,b,c)$ is reversed, so $(j,b)$ lies in $\ell_F$ and by the Claim
it is the first pair of $\ell_F$; the braid move
$(b,c),(j,c),(j,b)\mapsto(j,b),(j,c),(b,c)$ keeps $j$-normality (only the triple
$(j,b,c)$ is affected, and it is reversed into the other allowed order) and
lowers $m$ by $2$. If $d\ne b,c$, the two pairs have four distinct entries; the
exchange keeps $j$-normality (no triple contains both) and lowers $m$ by $1$.

\emph{Front.} If the list is $j$-normal and $j+1<n$, it can be moved to one
beginning with $(j,j+1)$: sort; the first pair $(j,d)$ of $\ell_F$ has $d=j+1$,
since otherwise $(j,j+1)$ (in $\ell_F$) and $(j+1,d)$ (in $\ell_G$) would both
come after $(j,d)$, violating the triple $(j,j+1,d)$.

\emph{Back.} A $j$-normal list can be moved to $\ell'R$ with $\ell'$
$(j+1)$-normal: sort the reversed list into $\ell_F\ell_G$ and reverse back;
$\ell_G^{\rm rev}$ contains exactly the pairs with first entry $>j$, in an order
satisfying the triple condition.

\emph{Iteration.} Applying Back $a$ times (to prefixes), a normal order moves
to $\ell_aR$ with $\ell_a$ $a$-normal, and Front applied to $\ell_a$ gives a list
beginning with $(a,a+1)$.

\emph{Conclusion.} If $x\in V_\mu\cap\mathfrak n^-V$, then $Px\in V_\mu$ is killed
by all $E_{a,a+1}$, hence by all raising operators (iv), and
$Px=(Px-x)+x\in\mathfrak n^-V$ by (iii); so $Px=0$ by (v).
\end{proof}

\begin{lemma}[Clause $\mathcal A_7$]\label{i1:lem-A7}
Let $O$ be a discrete valuation ring with fraction field $K$ of characteristic
$0$, $p\in O$ a nonzero non-unit with $p$ prime in $\Z$, $V$ a
finite-dimensional $\mathfrak{gl}_n$-representation over $K$ and $M$ a
Kostant lattice with $KM=V$. Then $\mathrm{Inv}=P_0\circ\pi$ for every weight
projector $\pi$ for $0$ and every extremal projector $P_0$ for $\mu=0$ (normal
order, $T\ge\dim V$); and $p^{N_0(n)\delta(p,n,\ell(V))}\mathrm{Inv}(M)\subseteq M$.
\end{lemma}

\begin{proof}
Integers prime to $p$ are units of $O$ (else $1=xp+yc$ would lie in the maximal
ideal). \emph{The identity.} By Lemma~\ref{i1:lem-isodec} it suffices to test on
$V[\mu]$. On $V[0]$: a highest weight vector $v$ of weight $0$ is killed by
all $E_{ab}$ ($a>b$ too, by Lemma~\ref{i1:lem-hw}(2) with $\mu_a=\mu_b=0$), and
the common kernel is stable, so $V[0]$ lies in it; then $\pi x=x$ and $P_0x=x$
(Lemma~\ref{i1:lem-A6}). On $V[\mu]$, $\mu\ne0$: $V[\mu]$ is spanned by weight
vectors $wv$, $w$ a lowering word (Lemma~\ref{i1:lem-hw}(3)); $\pi(wv)=0$ unless
its weight is $0$, in which case $w\ne1$, $wv\in V_0\cap\mathfrak n^-V$ and
$P_0(wv)=0$ by Lemma~\ref{i1:lem-A6}.

\emph{The bound.} (a) Some weight projector for $0$ maps $M$ into $M$: if all
weights have $|\nu_a|\le m$, there are integers $d_k$ with
$\sum_{k\le2m}d_k\binom{z+m}k=[z=0]$ for $|z|\le m$ (a unitriangular system),
and $\binom{E_{aa}+m}{k}=\sum_{i+j=k}\binom mj\binom{E_{aa}}i$; so
$\pi=\prod_a\sum_kd_k\binom{E_{aa}+m}k$ is an integral combination of Kostant
generators. (b) For $y\in M\cap V_0$ and a factor $P_{(a,b)}$ (here
$s_{ab}=b-a\le n-1$), a nonzero term needs $E^{(k)}_{ab}y\ne0$, of weight
$k(e_a-e_b)$, so $k\le\ell(V)$; and $v_p\binom{t+k}{k}\le\lfloor\log_p(t+k)\rfloor\le\delta$
by Legendre's formula (each digit term is $0$ or $1$ and vanishes for $p^i>t+k$).
So $p^\delta P_{(a,b)}y\in M\cap V_0$. (c) Inducting along the $N_0(n)$ factors,
$p^{N_0(n)\delta}P_0\pi x\in M$ for $x\in M$, and $P_0\pi=\mathrm{Inv}$ (the
identity does not depend on the choice of $\pi$).
\end{proof}


\subsubsection*{Twisted valuations: clauses $\mathcal A_{12}$, $\mathcal A_{13}$, $\mathcal A_{22}$, $\mathcal A_{24}$(3), $\mathcal A_{25}$, $\mathcal A_{26}$, $\mathcal A_{35}$}

\begin{definition}[Alphabets and twisted valuations]\label{i1:def-twist}
An \emph{alphabet} is a set $\iota$ of \emph{letters}, each declared a
\emph{vector} or a \emph{covector}. For a commutative ring $R$ and $n\ge1$ put
$R[\iota]_n=R[X_{l,a}:l\in\iota,\ a<n]$. For a monomial $X^m$,
$\mathrm{md}(m)(l)=\sum_am_{l,a}$ is its letter multidegree; $\mathrm{Pol}_d$ is
the span of the monomials with $\mathrm{md}(m)=d$. For $S\subseteq\iota$:
$\deg_S(d)=\sum_{l\in S}d_l$; $\mathrm{wt}_S(m)_a=\sum_{l\in S\text{ vector}}m_{l,a}-\sum_{l\in S\text{ covector}}m_{l,a}$;
$\mathrm{vec}_S(m)$, $\mathrm{cov}_S(m)$ are the parts of $|m|$ carried by
variables $X_{l,a}$ with $l\in S$ a vector, resp.\ a covector; and
$E^S_{ab}=\sum_{l\in S\text{ vector}}X_{l,a}\partial_{l,b}-\sum_{l\in S\text{ covector}}X_{l,b}\partial_{l,a}$.
Over a field of characteristic $0$ the $E^S_{ab}$ form a
$\mathfrak{gl}_n$-representation on $R[\iota]_n$ in which $X^m$ is a weight
vector of weight $\mathrm{wt}_S(m)$.

Let $g$ be the generic $n\times n$ matrix over
$A_n=R[g_{ab}][\det g^{-1}]$, $T$ a Laurent variable, $D_j$ the diagonal matrix
with $T$ in position $j$ and $1$ elsewhere, and $M_j(k)=gD_j^kg^{-1}$. The
\emph{twist} $\mathrm{tw}^{(j)}_S:R[\iota]_n\to A_n[T^{\pm1}][\iota]_n$ is the
ring map with $X_{l,a}\mapsto\sum_bM_j(1)_{ba}X_{l,b}$ for $l\in S$ a vector,
$X_{l,a}\mapsto\sum_bM_j(-1)_{ab}X_{l,b}$ for $l\in S$ a covector, and
$X_{l,a}\mapsto X_{l,a}$ for $l\notin S$. For $Q\in B[T^{\pm1}][\iota]_n$ let
$\mathcal D(Q)$ be the set of exponents of $T$ occurring in its coefficients.
The \emph{twisted valuation} is $v_S(F)=\min\mathcal D(\mathrm{tw}^{(n-1)}_SF)$
and the \emph{top exponent} is $\lambda_{1,S}(F)=\max\mathcal D(\mathrm{tw}^{(0)}_SF)$
(with $v_S(0)=+\infty$, $\lambda_{1,S}(0)=-\infty$). The \emph{divided
derivative} is $\partial^{[\beta]}X^m=\prod_v\binom{m_v}{\beta_v}X^{m-\beta}$
($0$ unless $\beta\le m$).
\end{definition}

\begin{lemma}[Clause $\mathcal A_{12}$]\label{i1:lem-A12}
Let $O,K$ be as in Lemma~\ref{i1:lem-A7}, $\iota$ an alphabet (Definition~\ref{i1:def-twist}),
$n\ge1$, $S_0\subseteq\dots\subseteq S_{N-1}\subseteq\iota$. For $a\ne b$ and
$l\in\iota$ put $\mathrm{src}(l)=(l,b)$, $\mathrm{dst}(l)=(l,a)$,
$\varepsilon_l=1$ if $l$ is a vector, and $\mathrm{src}(l)=(l,a)$,
$\mathrm{dst}(l)=(l,b)$, $\varepsilon_l=-1$ if $l$ is a covector; let
$E^{S,[k]}_{ab}X^m=\sum_f\prod_{l\in S}\varepsilon_l^{f_l}\binom{m_{\mathrm{src}(l)}}{f_l}
X^{m-\sum_lf_l\mathrm{src}(l)+\sum_lf_l\mathrm{dst}(l)}$, the sum over $f:S\to\N$
with $\sum f_l=k$. Then:
(1) $k!E^{S,[k]}_{ab}=(E^S_{ab})^k$;
(2) $\langle E^S_{ab}F,G\rangle=\langle F,E^S_{ba}G\rangle$ for the Fock pairing
$\langle X^m,X^{m'}\rangle=[m=m']\,m!$;
(3) $\ell(\mathrm{Pol}_d)\le\deg_{S_i}(d)$ for $E^{S_i}$;
(4) for isotypic projectors $e_{i,\lambda}$ of $E^{S_i}$: $e_{i,\lambda}$ commutes
with $E^{S_j}_{ab}$ for $i<j$; an endomorphism that is the identity on the joint
components $J_t=\bigcap_iV_{S_i}[t_i]$, $t\in\mathcal T$, and $0$ on the others
equals $\sum_{t\in\mathcal T}\prod_ie_{i,t_i}$; if $e_{i,\lambda}$ does not
vanish on $\mathrm{Pol}_d$ then $\ell(\lambda)\le\deg_{S_i}(d)$; and if
$\ell(t_i)\le\deg_{S_i}(d)$ for all $t\in\mathcal T$ then
$\max_tN_0(n)\delta(p,n,\ell(\mathrm{Pol}_d)+\ell(t_i))\le N_0(n)\log_p(2\deg_{S_i}(d)+n)$.
\end{lemma}

\begin{proof}
(1) $E^S_{ab}=\sum_lD_l$ with $D_l=\varepsilon_lX_{\mathrm{dst}(l)}\partial_{\mathrm{src}(l)}$,
and the term $T_f$ of $f$ satisfies $D_lT_f=(f_l+1)T_{f+e_l}$ (as
$(f+1)\binom M{f+1}=(M-f)\binom Mf$; sources and destinations are disjoint
since $a\ne b$). Summing, $E^S_{ab}E^{S,[k]}_{ab}=(k+1)E^{S,[k+1]}_{ab}$.
(2) $X_u\partial_v$ and $X_v\partial_u$ are Fock adjoint (compare coefficients
of a pair of monomials); sum over the terms. (3) $E^S_{aa}X^m=\mathrm{wt}_S(m)_aX^m$
with $\mathrm{wt}_S(m)_a=\sum_{l\in S\text{ vec}}m_{l,a}-\sum_{l\in S\text{ cov}}m_{l,a}$,
and for $m,m'$ of multidegree $d$, $\mathrm{wt}_S(m)_a-\mathrm{wt}_S(m')_b\le\deg_S(d)$.
(4) For $i<j$, $E^{S_j}_{ab}=E^{S_i}_{ab}+E^{S_j\setminus S_i}_{ab}$ and the
second summand commutes with $E^{S_i}$; so $E^{S_j}_{ab}$ preserves each
$V_{S_i}[\mu]$ and commutes with $e_{i,\lambda}$ (Lemma~\ref{i1:lem-isodec}).
Consequently the $e_{i,\mu}$ commute pairwise and preserve each other's
isotypic parts; by induction on $N$ the $J_t$ span, and on $J_t$ the product
$\prod_ie_{i,t'_i}$ is $[t'=t]$. If $e_{i,\lambda}F\ne0$ with
$F\in\mathrm{Pol}_d$, the multidegree projection commutes with $E^{S_i}$, so
some highest weight vector of weight $\lambda$ has a nonzero component in
$\mathrm{Pol}_d$, of weight $\lambda$; by (3), $\ell(\lambda)\le\deg_{S_i}(d)$.
The numerical bound follows from (3) and $\lfloor\cdot\rfloor\le\cdot$.
\end{proof}

\begin{lemma}[Clause $\mathcal A_{35}$: renaming and base change]\label{i1:lem-A35}
(1) If $f:\iota\to\kappa$ is an injective map of alphabets preserving types and
$F\circ f$ denotes the renamed polynomial, then $v_S(F\circ f)=v_{f^{-1}S}(F)$ and
$\lambda_{1,S}(F\circ f)=\lambda_{1,f^{-1}S}(F)$. (2) If $O$ is a domain with
fraction field $K$, $v_S$ and $\lambda_{1,S}$ of $F\in O[\iota]_n$ do not change
under $O\to K$.
\end{lemma}

\begin{proof}
(1) On variables, the twist of $X_{f(l),a}$ is the renaming of the twist of
$X_{l,a}$ (with $f^{-1}S$), so the twists of $F\circ f$ and $F$ correspond
under an injective renaming of monomials which keeps coefficients. (2) The
twist commutes with the coefficient map $A_n(O)\to A_n(K)$, which is
injective.
\end{proof}

\begin{lemma}[Clause $\mathcal A_{13}$: valuation lemma]\label{i1:lem-A13}
Let $R$ be a field, $S\subseteq\iota$, $F,G\in R[\iota]_n$.
\begin{enumerate}
\item $v_S(FG)=v_SF+v_SG$, $v_S(F+G)\ge\min(v_SF,v_SG)$,
 $\lambda_{1,S}(FG)=\lambda_{1,S}F+\lambda_{1,S}G$,
 $\lambda_{1,S}(F+G)\le\max(\lambda_{1,S}F,\lambda_{1,S}G)$.
\item Let $M_G$ be the $R$-span of the $uG$, $u$ in the monoid generated by the
 operators $E^{S,[k]}_{ab}$ ($a\ne b$; Lemma~\ref{i1:lem-A12}) and
 $B^{S,k}_a:X^m\mapsto\binom{\mathrm{wt}_S(m)_a}{k}X^m$ ($k\in\N$), and say that
 $\mu\in\Z^n$ \emph{occurs} in $M_G$ if $M_G$ contains a nonzero polynomial all of
 whose monomials have weight $\mu$. Then $\theta\le v_S(G)$ iff $\theta\le\mu_{n-1}$
 for every $\mu$ occurring in $M_G$, and $\lambda_{1,S}(G)\le M$ iff $\mu_0\le M$
 for every such $\mu$.
\item If $R$ has characteristic $0$ and $(e_\mu)$ are the isotypic projectors of
 $E^S$, then $\theta\le v_S(G)$ iff $\theta\le\mu_{n-1}$ whenever $e_\mu G\ne0$,
 and $\lambda_{1,S}(G)\le M$ iff $\mu_0\le M$ whenever $e_\mu G\ne0$.
\end{enumerate}
\end{lemma}

\begin{proof}
(1) The twist is a ring map into $A_n[\iota]_n[T^{\pm1}]$, and $A_n[\iota]_n$ is a
domain; so the lowest (highest) $T$-parts of the twists of $F$ and $G$ multiply
to the nonzero lowest (highest) part of the twist of $FG$. The inequalities for
sums are clear.

(2) For a ring map $\phi:R\to B$ and $h,h'\in M_n(B)$ let $\sigma^\phi_{h,h'}$ be
the ring map which is $\phi$ on coefficients, sends $X_{l,c}$ to
$\sum_bh_{bc}X_{l,b}$ ($l\in S$ a vector), $\sum_bh'_{cb}X_{l,b}$ ($l\in S$ a
covector), and fixes the other variables. Then
$\sigma_{h_1,h_1'}\circ\sigma^\phi_{h_2,h_2'}=\sigma^\phi_{h_1h_2,h_2'h_1'}$ and
$\mathrm{tw}^{(j)}=\sigma_{gD_jg^{-1},gD_j^{-1}g^{-1}}$. Two consequences of the
binomial theorem: (T) $\sigma_{\mathrm{diag}(u),\mathrm{diag}(u)^{-1}}X^m=\prod_au_a^{\mathrm{wt}_S(m)_a}X^m$;
(Tr) for the transvection $\tau(t)=1+te_{ab}$,
$\sigma_{\tau(t),\tau(-t)}F=\sum_kt^kE^{S,[k]}_{ab}F$. Moreover (C) a subspace
stable under all $B^{S,k}_a$ contains the weight components of its elements,
since integral combinations of the $w\mapsto\binom wk$ realise every indicator
function $[w=w_0]$ on a finite set of integers.

\emph{Twist exponents.} Let $\Gamma(G)=\sigma_{g^{-1},g}(G)$. By the composition
rule, $\mathrm{tw}^{(j)}G=\sigma_{g,g^{-1}}(\sigma_{D_j,D_j^{-1}}\Gamma(G))$; by (T)
the inner polynomial has coefficient $\Gamma(G)_mT^{\mathrm{wt}_S(m)_j}$ at
$X^m$, and the outer substitution, which is invertible and does not involve
$T$, preserves the set of $T$-exponents. Hence
$\mathcal D(\mathrm{tw}^{(j)}G)=\{\mathrm{wt}_S(m)_j:X^m\text{ in the support of }\Gamma(G)\}$.

It remains to see that the weights of the monomials of $\Gamma(G)$ and the
weights occurring in $M_G$ satisfy the same one-coordinate conditions.
(A) Every weight of a monomial of $\Gamma(G)$ occurs in $M_G$: over the
fraction field $L$ of $A_n$, the $L$-span $W_L$ of $M_G$ is preserved by
$\sigma_{h,h'}$ for transvections (by (Tr)) and diagonal matrices (by (T) and
(C)), hence by $\sigma_{\psi(g^{-1}),\psi(g)}$ (Gaussian elimination); so
$\Gamma(G)\in W_L$, and a monomial in the support of an element of $W_L$ has a
weight occurring in $M_G$, by (C). (B) Conversely, let $\mathcal P$ be the set
of $y$ such that for every ring map $\phi:R\to k_1$ and $g_0g'_0=1$ over $k_1$,
all monomials of $\sigma^\phi_{g'_0,g_0}(y)$ have weights satisfying the
condition. Then $G\in\mathcal P$ (specialise $g\mapsto g_0$ in $\Gamma(G)$), and
$\mathcal P$ is stable under the $E^{S,[k]}_{ab}$ and the $B^{S,k}_a$: apply the
definition over $k_1[T^{\pm1}]$ to $\tau(T)g'_0$, resp.\ $\mathrm{diag}g'_0$,
and read off the coefficients of the powers of $T$ via (Tr), resp.\ (T). So
$M_G\subseteq\mathcal P$, and taking $\phi=\mathrm{id}$, $g_0=1$ gives the
condition for every weight occurring in $M_G$.

(3) In characteristic $0$, $M_G=UG$ (Lemma~\ref{i1:lem-A12}(1); the
$\binom{E_{aa}}k$ are polynomials in $E_{aa}$), and the conditions are monotone:
if they hold for $\mu$ and $\mu_{n-1}\le\nu_{n-1}$, $\nu_0\le\mu_0$, they hold for
$\nu$. If $e_\mu G\ne0$, $UG$ contains a highest weight vector of weight $\mu$
(Lemma~\ref{i1:lem-isodec}), so $\mu$ occurs. Conversely, a weight vector $x\in
UG$ lies in $\sum_{e_\mu G\ne0}V[\mu]$ (the $e_\mu$ commute with $U$), and every
weight $\nu'$ of $V[\mu]$ satisfies $\mu_{n-1}\le\nu'_{n-1}$, $\nu'_0\le\mu_0$
(Lemma~\ref{i1:lem-hw}(3)); by monotonicity and (2) the claim follows.
\end{proof}

\begin{lemma}[Clause $\mathcal A_{26}$]\label{i1:lem-A26}
Let $O$ be a discrete valuation ring, $\iota$ any alphabet, $S\subseteq\iota$.
(1) For $F,G\in O[\iota]_n$ the four relations of Lemma~\ref{i1:lem-A13}(1)
hold. (2) For $l\in S$: if $l$ is a vector, $v_S(X_{l,c})\ge0$ and
$\lambda_{1,S}(X_{l,c})\le1$; if $l$ is a covector, $v_S(X_{l,c})\ge-1$ and
$\lambda_{1,S}(X_{l,c})\le0$. (3) If $y:\{0..n-1\}\to S$ is injective with all
$y_a$ vectors, $v_S(\det(X_{y_a,b}))=1$.
\end{lemma}

\begin{proof}
(1) Pass to a finite alphabet (Lemma~\ref{i1:lem-A35}(1)) and to $K$
(Lemma~\ref{i1:lem-A35}(2)), and apply Lemma~\ref{i1:lem-A13}(1). (2) The entries
of $M_j(\pm1)$ involve only $T^0,T^{\pm1}$. (3) The twist multiplies the matrix
$(X_{y_a,b})$ on the right by $M_{n-1}(1)$, of determinant $T$.
\end{proof}

\begin{lemma}[Clause $\mathcal A_{25}$: derivative lemma]\label{i1:lem-A25}
(1) For a field $R$ and exponent vectors $\alpha,\beta$:
$v_S(F)\le v_S(X^\alpha\partial^{[\beta]}F)+\mathrm{cov}_S(\alpha)+\mathrm{vec}_S(\beta)$ and
$\lambda_{1,S}(X^\alpha\partial^{[\beta]}F)\le\lambda_{1,S}(F)+\mathrm{vec}_S(\alpha)+\mathrm{cov}_S(\beta)$;
in particular both hold with $|\alpha|+|\beta|$.
(2) For a prime $p\ge5$, $\rho_k<p^t$, $M_k\in\N$ and $c\in\N$, the product
$\prod_k\binom{\rho_k+p^te_k}{M_k}$ is congruent modulo $p^c$, for all $e\in\N^K$,
to an integral combination of the $\prod_k\binom{e_k}{l_k}$ with
$\sum_kl_k<\sum_kM_k/p^t+c$ and $M_k\le\rho_k+p^tl_k$.
(3) If $P_k\ge1$ and every monomial $X^m$ of $F$ has $m_k\equiv\rho_k\pmod{P_k}$,
then every monomial of $\partial_kF$ has $m_k\equiv\rho_k-1$ if $\rho_k\ge1$ and
$m_k\equiv P_k-1$ if $\rho_k=0$.
\end{lemma}

\begin{proof}
(3) is immediate from $\partial_kX^m=m_kX^{m-e_k}$. (2) Let
$a_k(l)=a(\rho_k,M_k,l)$ (Lemma~\ref{i1:lem-binexp}); also $a_k(l)=0$ if
$M_k<l$ or $\rho_k+p^tl<M_k$, since $X^l$ divides and the degree is
$\le\rho_k+p^tl$. Expand $\prod_k\binom{\rho_k+p^te_k}{M_k}=\sum_l\prod_ka_k(l_k)\binom{e_k}{l_k}$
and drop the terms with $\sum_k\max(0,l_k-\lfloor M_k/p^t\rfloor)\ge c$, which
are divisible by $p^c$ (Lemma~\ref{i1:lem-binexp}); the others have
$\sum_kl_k<c+\sum_kM_k/p^t$.

(1) For $e=\pm1$ let $\mathcal F^e_t$ be the polynomials over $A_n[T^{\pm1}]$ whose
coefficients involve only $T^k$ with $ek\ge t$; these form a multiplicative
filtration, and $t\le v_S(F)$ iff $\mathrm{tw}^{(n-1)}F\in\mathcal F^1_t$, while
$\lambda_{1,S}(F)\le M$ iff $\mathrm{tw}^{(0)}F\in\mathcal F^{-1}_{-M}$. Since the
entries of $M_j(k)$ involve only $T^0,T^k$, a variable $X_{l,c}$ twists into
$\mathcal F^1_{-[l\in S\text{ cov}]}$ and into $\mathcal F^{-1}_{-[l\in S\text{ vec}]}$;
so multiplication by $X^\alpha$ costs $\mathrm{cov}_S(\alpha)$, resp.\
$\mathrm{vec}_S(\alpha)$. For derivatives, $\partial^{[\beta]}F$ is the
coefficient of $Y^\beta$ in $F(X+Y)$, and the substitution $\Theta$ with
$X_{l,b}\mapsto X_{l,b}+\sum_cM_j(-1)_{cb}Y_{l,c}$ ($l\in S$ vector),
$X_{l,b}+\sum_cM_j(1)_{bc}Y_{l,c}$ ($l\in S$ covector), $X_{l,b}+Y_{l,b}$
(otherwise) satisfies $\Theta\circ\mathrm{tw}^{(j)}=\mathrm{tw}^{(j)}\circ(X\mapsto
X+Y)$ because $M_j(-1)M_j(1)=1$. Comparing $T$-exponents of the $Y$-terms shows
that the coefficient of $Y^\beta$ costs $\mathrm{vec}_S(\beta)$ for $e=1$ and
$\mathrm{cov}_S(\beta)$ for $e=-1$.
\end{proof}

\begin{lemma}[Clause $\mathcal A_{24}$, third part: degree truncation]\label{i1:lem-A24c}
For a field $R$, finitely many exponent vectors $m$ with $|m|\le4c$ and
$q_m\in R$: $v_S(F)\le v_S(\sum_mq_mX^m\partial^{[m]}F)+4c$ and
$\lambda_{1,S}(\sum_mq_mX^m\partial^{[m]}F)\le\lambda_{1,S}(F)+4c$.
\end{lemma}

\begin{proof}
By the proof of Lemma~\ref{i1:lem-A25}(1) with $\alpha=\beta=m$, each twisted term
lies in $\mathcal F^1_{t-|m|}\subseteq\mathcal F^1_{t-4c}$ when $t\le v_S(F)$
(and in $\mathcal F^{-1}_{-M-4c}$ when $M\ge\lambda_{1,S}(F)$); the filtration
steps are additive groups.
\end{proof}

\begin{lemma}[Clause $\mathcal A_{22}$: powers of the determinant]\label{i1:lem-A22}
Let $R$ be a field, $\iota$ finite, $r\ge1$, $S\subseteq\iota$, and
$A:\{0..r-1\}\to\iota$ injective with every $A(i)\in S$ a vector. Put
$N_A=(X_{A(i),j})_{i,j<r}$, $\det_A=\det N_A$, and
$\deg_{\rm oth}(d)=\sum d_y$ over the vectors $y\in S\setminus A(\{0..r-1\})$.
(1) If $G\in\mathrm{Pol}_d\subseteq R[\iota]_r$, $\theta\le v_S(G)$ and
$\deg_{\rm oth}(d)<\theta$, then $\det_A^{\theta-\deg_{\rm oth}(d)}$ divides $G$.
(2) If $G=\det_A^mH$, then $v_{S'}(G)\le v_{S'}(H)+m$ for every $S'\subseteq\iota$.
\end{lemma}

\begin{proof}
Write $w_S(m)=\mathrm{wt}_S(m)_{r-1}$ and $\mathrm{last}=r-1$. Three facts.
(a) \emph{Specialisation.} If a ring map $\Psi$ applies a ring map to the
coefficients and substitutes values for the variables, every $T$-exponent of
$\Psi(Q)$ lies in $\mathcal D(Q)$. (b) $v_{S'}(\det_A)\le1$: specialise $g$ to
$1$ and $X_{y,b}$ to $[y=A(b)]$; the twist of $\det_A$ becomes
$\det\mathrm{diag}(1,\dots,1,T^{[A(\mathrm{last})\in S']})$. (c) \emph{Frame
change.} If $g_0g'_0=1$ over a ring $k$ and $E$ is the substitution
$X_{y,c}\mapsto\sum_a(g'_0)_{ac}X_{y,a}$ ($y\in S$ vector),
$\sum_a(g_0)_{ca}X_{y,a}$ ($y\in S$ covector), identity otherwise, then every
monomial $X^m$ of $E(G)$ has $w_S(m)\ge v_S(G)$: specialising $g\mapsto g_0$
and $X\mapsto E(X)$ in the twist gives $\delta(E(G))$, where $\delta$ multiplies
$X^m$ by $T^{w_S(m)}$; apply (a).

(2) By Lemma~\ref{i1:lem-A13}(1) and (b), $v_{S'}(G)=mv_{S'}(\det_A)+v_{S'}(H)\le
m+v_{S'}(H)$.

(1) Let $c=\mathrm{adj}(N_A)_{\mathrm{last},\mathrm{last}}$. In the factorial
ring $P=R[\iota]_r$, $c\ne0$ and $c$ is coprime to $\det_A$. Indeed, each
$\mathrm{adj}(N_A)_{a,\mathrm{last}}$ is the determinant of $N_A$ with row
$\mathrm{last}$ replaced by $e_a$; clearing column $a$ with that row shows it
lies in the subring generated by the $X_{A(i),b}$ with $i\ne\mathrm{last}$,
$b\ne a$, and evaluating at a permutation matrix shows it is nonzero; the
substitution $X_{A(\mathrm{last}),b}\mapsto\delta_{ba}$ maps $\det_A$ to it and
fixes every polynomial free of the letter $A(\mathrm{last})$. A divisor of a
nonzero polynomial in such a subring lies in it (compare degrees in each
variable). So a common divisor $q$ of $\det_A$ and $c$ is free of the letter
$A(\mathrm{last})$, hence divides every $\mathrm{adj}(N_A)_{a,\mathrm{last}}$,
hence involves no variable of column $a$ for any $a$: it is a constant.
Over $k=P[c^{-1}]$ let
$g_0$ be the identity with last column replaced by
$\mathrm{cc}=(\mathrm{adj}(N_A)_{a,\mathrm{last}})_a$; it is invertible since
$\mathrm{cc}_{\mathrm{last}}=c$. Let $F$ be the $k$-algebra map sending
$X_{y,c'}$ to $\sum_aX_{y,a}(g_0)_{ac'}$ ($y\in S$ vector),
$\sum_a(g_0^{-1})_{c'a}X_{y,a}$ ($y\in S$ covector), and $X_{y,c'}$ otherwise,
evaluated in $k$. Then $F\circ E=\mathrm{id}$ on $P\subseteq k$ for the
substitution $E$ of (c) with $g'_0=g_0^{-1}$, and
$F(X_{A(l),\mathrm{last}})=(N_A\,\mathrm{adj}N_A)_{l,\mathrm{last}}=\delta_{l,\mathrm{last}}\det_A$.
Expand $G=F(E(G))$ over the monomials $X^m$ of $E(G)$. A monomial with
$m(A(l),\mathrm{last})>0$ for some $l\ne\mathrm{last}$ contributes $0$.
Otherwise $\theta\le w_S(m)$ by (c), and termwise comparison gives
$w_S(m)\le m(A(\mathrm{last}),\mathrm{last})+\deg_{\rm oth}(d)$ (here $E$
preserves multidegrees); so the term contains
$\det_A^{m(A(\mathrm{last}),\mathrm{last})}$, divisible by $\det_A^t$,
$t=\theta-\deg_{\rm oth}(d)$. Hence $\det_A^t\mid G$ in $k$, i.e.\
$Gc^M=\det_A^tw$ in $P$ for some $M$ and $w$; as $c$ is coprime to $\det_A$,
$\det_A^t\mid G$.
\end{proof}

\subsubsection*{Pieri paths and the remaining clauses}

\begin{definition}[Pieri paths and path polynomials]\label{i1:def-pieri}
Fix $n\ge1$ and letters $0,\dots,N-1$, each a vector or a covector, with
degrees $B_i\in\N$. For antitone $\alpha,\beta\in\Z^n$, $\beta/\alpha$ is a
\emph{horizontal strip of size $b$} if $\beta_0\ge\alpha_0\ge\beta_1\ge\alpha_1\ge\dots\ge\beta_{n-1}\ge\alpha_{n-1}$
and $\sum(\beta_a-\alpha_a)=b$.
\begin{enumerate}
\item A \emph{Pieri path} is a sequence $\mu_0=0,\mu_1,\dots,\mu_N$ of antitone
 vectors in $\Z^n$ such that $\mu_{i+1}/\mu_i$ (letter $i$ a vector) or
 $\mu_i/\mu_{i+1}$ (letter $i$ a covector) is a horizontal strip of size
 $B_i$. Its \emph{end} is $\mu_N$.
\item For antitone $\mu$ and $c\in\Z$ put $\mathrm{colht}(\mu,c)=\#\{a:c\le\mu_a\}$.
 A \emph{filling} of $\mu$ assigns to each column $c$ a set
 $I(c)\subseteq\{0..n-1\}$ with $|I(c)|=\mathrm{colht}(\mu,c)$ for $c\ge1$ and
 $|I(c)|=n-\mathrm{colht}(\mu,c)$ for $c\le0$. The \emph{canonical filling} has
 $I^{\rm can}(c)$ the first $\mathrm{colht}(\mu,c)$ indices ($c\ge1$), resp.\ the
 last $n-\mathrm{colht}(\mu,c)$ indices ($c\le0$). Put $\hat I(c)=I(c)$ for
 $c\ge1$ and $\hat I(c)=\{0..n-1\}\setminus I(c)$ for $c\le0$, so
 $|\hat I(c)|=\mathrm{colht}(\mu,c)$. The filling is \emph{standard} if for every
 $c$ and every $a<|\hat I(c+1)|$ the $a$-th smallest element of $\hat I(c)$ is
 at most the $a$-th smallest element of $\hat I(c+1)$. $\mathrm{St}(\mu)$ is the
 set of standard fillings of $\mu$.
\item A \emph{form} is a function $K$ from subsets of $\{0..n-1\}$ to a ring.
 With $\mathrm{sg}(J,m)=(-1)^{\#\{j\in J:j>m\}}$, $\mathbf 1(J)=[J=\emptyset]$ and
 $x\in A^n$, put $(K\wedge x)(J)=\sum_{m\in J}\mathrm{sg}(J,m)K(J\setminus m)x_m$
 and $(K\lrcorner x)(J)=\sum_{m\notin J}\mathrm{sg}(J,m)x_mK(J\cup m)$. For a
 Pieri path $\mu$ and a column $c$, the \emph{column word} is $w_c^{(0)}=\mathbf 1$,
 $w_c^{(i+1)}=w_c^{(i)}$ if $\mathrm{colht}(\mu_{i+1},c)=\mathrm{colht}(\mu_i,c)$, and
 otherwise $w_c^{(i+1)}=w_c^{(i)}\wedge X_i$ if letter $i$ is a vector and
 $c\ge1$ or a covector and $c\le0$, and $w_c^{(i+1)}=w_c^{(i)}\lrcorner X_i$ in the
 other two cases; here $X_i=(X_{i,b})_{b<n}$. The \emph{path polynomial} of
 $\mu$ and a filling $I$ of $\mu_N$ is $\Psi_{\mu,I}=\prod_cw^{(N)}_c(I(c))\in\Z[X_{i,b}]$
 (a finite product).
\end{enumerate}
\end{definition}

\begin{lemma}[Clause $\mathcal A_{15}$, reduction]\label{i1:lem-A15}
Assume, for all $n$, letter types and degrees $B$: (I) over every field the
$\Psi_{\mu,I}$ ($\mu$ a Pieri path, $I\in\mathrm{St}(\mu_N)$) are linearly
independent, and (C) the number of such pairs $(\mu,I)$ is
$\prod_{i<N}\binom{n+B_i-1}{B_i}$. Then for every ring $R$ the $\Psi_{\mu,I}$ form
an $R$-basis of $\mathrm{Pol}_B\subseteq R[X_{i,b}]$ (clause $\mathcal A_{15}$,
for $R$ a discrete valuation ring).
\end{lemma}

\begin{proof}
The monomials of letter multidegree $B$ are $\prod_i\binom{n+B_i-1}{B_i}$ in
number and form a basis of $\mathrm{Pol}_B$; each $\Psi_{\mu,I}$ has integer
coefficients and lies in $\mathrm{Pol}_B$ (each letter enters $B_i$ times, by
the strip sizes). By (C) the coefficient matrix over $\Z$ is square; by (I) over
$\mathbb F_q$ its determinant is prime to every prime $q$; so it is $\pm1$ and
the matrix is invertible over every $R$.
\end{proof}


\begin{theorem}[Kostant's commutation formula; B.~Kostant, Groups over $\Z$, Proc.\ Sympos.\ Pure Math.\ 9 (1966) 90--98; J.~E.~Humphreys, \emph{Introduction to Lie Algebras and Representation Theory}, GTM 9, \S26.2, Lemma]\label{i1:thm-kostant}
Let $e,f,h$ be elements of an associative algebra over a field of
characteristic $0$ with $[h,e]=2e$, $[h,f]=-2f$, $[e,f]=h$, and put
$e^{(k)}=e^k/k!$, $f^{(l)}=f^l/l!$. Then for all $k,l\ge0$
\[
 e^{(k)}f^{(l)}=\sum_{j=0}^{\min(k,l)}f^{(l-j)}\binom{h-k-l+2j}{j}e^{(k-j)},
\]
where $\binom{h-c}{j}=(h-c)(h-c-1)\cdots(h-c-j+1)/j!$.
\end{theorem}

\noindent It is applied in a $\mathfrak{gl}_n$-representation (characteristic
$0$) with $e=E_{ab}$, $f=E_{ba}$, $h=E_{aa}-E_{bb}$, $a\ne b$, which satisfy
the relations. On $x$ of weight $\nu$ the vector $E^{(k-j)}_{ab}x$ has
$h$-eigenvalue $\nu_a-\nu_b+2(k-j)$, so the formula becomes
$E^{(k)}_{ab}E^{(l)}_{ba}x=\sum_j\binom{\nu_a-\nu_b+k-l}{j}E^{(l-j)}_{ba}E^{(k-j)}_{ab}x$.

\begin{lemma}[Clauses $\mathcal A_8$, $\mathcal A_9$, $\mathcal A_{10}$]\label{i1:lem-A8}
Let $O,K,p$ be as in Lemma~\ref{i1:lem-A7}, $V$ a finite-dimensional
$\mathfrak{gl}_n$-representation over $K$, $M$ a Kostant lattice with $KM=V$,
and $(e_\lambda)$ the isotypic projectors (Lemma~\ref{i1:lem-isodec}).
\begin{itemize}
\item[$\mathcal A_8$] Let $v$ be a highest weight vector of weight $\lambda$
 with $Uv=V$, $\Delta(v)$ the $O$-span of $U_\Z v$, and $\nabla(\lambda,v)$ the
 set of $x$ with $ux\in Ov+\sum_{\mu\ne\lambda}V_\mu$ for all $u\in U_\Z$. Then
 $p^{N_0(n)\delta(p,n,\ell(\lambda))}\nabla(\lambda,v)\subseteq\Delta(v)$, and if
 $M\cap V_\lambda=Ov$ then $\Delta(v)\subseteq M\subseteq\nabla(\lambda,v)$.
\item[$\mathcal A_9$] For a finite set $\Sigma$ of dominant weights,
 $p^{\max_{\lambda\in\Sigma}N_0(n)\delta(p,n,\ell(V)+\ell(\lambda))}\sum_{\lambda\in\Sigma}e_\lambda(M)\subseteq M$.
\item[$\mathcal A_{10}$] For dominant $\lambda$, a weight projector $\pi$ for
 $\lambda$, $c\in\Z$ and $Y\in M$ with $e_\lambda Y=Y$: if $f(\pi(uY))\in p^cO$
 for all $u\in U_\Z$ and all $O$-linear $f:M\to O$, then
 $Y\in p^{c-N_0(n)\delta(p,n,\ell(\lambda))}M$.
\end{itemize}
\end{lemma}

\begin{proof}
\emph{Preliminaries.} (P1) \emph{The $\ell^\infty$ bound.} Let $L_\infty(V)$ be
the largest $|\nu_a|$ over the weights $\nu$ of $V$ with $\sum_a\nu_a=0$ ($0$ if
none). The proof of Lemma~\ref{i1:lem-A7} gives
$p^{N_0(n)\delta(p,n,L_\infty(V))}\mathrm{Inv}(M)\subseteq M$: in step (b) a
nonzero $E^{(k)}_{ab}y$ has weight $k(e_a-e_b)$ with coordinate sum $0$, so
$k\le L_\infty(V)$, and truncating the extremal projector at $L_\infty(V)$ is
harmless. (P2) If $Uv=V$ for a highest weight vector $v$ of weight $\lambda$,
then $V$ is spanned by the $wv$, $w$ a lowering word, each $wv$ ($w\ne1$) has a
weight of strictly larger $\sum_c c\nu_c$; hence $V_\lambda=Kv$,
$V=Kv\oplus\sum_{\mu\ne\lambda}V_\mu$, and every coordinate of every weight lies
in $[\lambda_{n-1},\lambda_0]$ (use the $\mathfrak{sl}_2$-strings of
Lemma~\ref{i1:lem-hw}(2) along $E_{a0}$, $E_{0a}$, resp.\ $E_{n-1,a}$, $E_{a,n-1}$).
(P3) $M$ and $M\cap H$, for a subspace $H$, are free over $O$ (finitely
generated torsion-free), with bases that are $K$-bases of $V$, resp.\ $H$, when
they span. (P4) Weight projectors can be chosen as integral combinations of
products of the $\binom{E_{aa}}k$ (proof of Lemma~\ref{i1:lem-A7}); they
preserve Kostant lattices.

\emph{$\mathcal A_8$.} Let $\xi$ be the functional with $\xi(v)=1$ and
$\xi(V_\mu)=0$ for $\mu\ne\lambda$; so $x\in\nabla(\lambda,v)$ iff $\xi(ux)\in O$
for all $u\in U_\Z$. (a) $\Delta(v)\subseteq\nabla(\lambda,v)$: the $O$-span
$\Lambda$ of the words in lowering divided powers applied to $v$ lies in
$Ov+\sum_{\mu\ne\lambda}V_\mu$ and is stable under $U_\Z$; for raising divided
powers this is an induction on the height, using Kostant's formula
(Theorem~\ref{i1:thm-kostant}) in the form
$E^{(k)}_{ab}E^{(l)}_{ba}x=\sum_{j}\binom{\nu_a-\nu_b+k-l}{j}E^{(l-j)}_{ba}E^{(k-j)}_{ab}x$
($x$ of weight $\nu$) and the normal-ordering identity
$X^{(k)}Y^{(l)}=\sum_jY^{(l-j)}Z^{(j)}X^{(k-j)}$ when $Z=[X,Y]$ commutes with
$X$ and $Y$. (Proof of the latter: in formal power series in commuting
variables $s,t$, $e^{sX}Ye^{-sX}=Y+sZ$, hence $e^{sX}e^{tY}e^{-sX}=e^{t(Y+sZ)}=e^{tY}e^{stZ}$,
that is, $e^{sX}e^{tY}=e^{tY}e^{stZ}e^{sX}$; compare the coefficients of $s^kt^l$.)
(b) $\mathcal L=\{f\in\mathrm{End}V:f(\nabla(\lambda,v))\subseteq\Delta(v)\}$ is a
Kostant lattice spanning $\mathrm{End}V$ for the action $E_{ab}\cdot f=[E_{ab},f]$:
stability follows from
$(E_{ab}\cdot)^{(k)}f=\sum_m(-1)^{k-m}E^{(m)}_{ab}fE^{(k-m)}_{ab}$ and a similar
formula for the $\binom{E_{aa}\cdot}k$; finite generation from the functionals
$f\mapsto\xi(uf(u'v))$, which are integral on $\mathcal L$ and separate points
(by Lemma~\ref{i1:lem-cr} and $V_\lambda=Kv$); spanning since $\nabla(\lambda,v)$
is finitely generated and $\Delta(v)$ spans. (c) The invariant projector of
$\mathrm{End}V$ is $T(f)=(\dim V)^{-1}\mathrm{tr}(f)\,\mathrm{id}_V$ (invariant
elements commute with $\mathfrak{gl}_n$ and so are scalars by $V_\lambda=Kv$;
other isotypic parts are traceless). (d) $L_\infty(\mathrm{End}V)\le\ell(\lambda)$
by (P2). (e) $\phi=\xi(\cdot)v\in\mathcal L$, so by (P1) for $\mathcal L$,
$p^{D}(\dim V)^{-1}\mathrm{id}_V\in\mathcal L$ with $D=N_0(n)\delta(p,n,\ell(\lambda))$;
multiplying by $\dim V$, $p^D\nabla(\lambda,v)\subseteq\Delta(v)$. If
$M\cap V_\lambda=Ov$, then $\Delta(v)\subseteq M$, and for $x\in M$, $u\in U_\Z$
the $\lambda$-component $\xi(ux)v$ of $ux\in M$ lies in $M\cap V_\lambda=Ov$ by
(P4), so $x\in\nabla(\lambda,v)$.

\emph{$\mathcal A_9$.} Fix $\lambda\in\Sigma$, $x\in M$. If $V$ has no highest
weight vector of weight $\lambda$, $e_\lambda x=0$. Otherwise take one, $v_0$, let
$L=Uv_0$ with an $O$-basis $b_1,\dots,b_d$ of $M\cap L$ (P3),
$H=\mathrm{Hom}_K(L,V)$ with $E_{ab}\cdot g=E_{ab}g-gE_{ab}$, $P$ its invariant
projector, and $\mathcal H=\{g:g(M\cap L)\subseteq M\}$. Then (i) $\mathcal H\cong
M^d$ is a Kostant lattice spanning $H$ (as in (b)); (ii)
$L_\infty(H)\le\ell(V)+\ell(\lambda)$ (by (P2) for $L$ and the definition of
$\ell(V)$); (iii) $d\sum_iP(b_i^*(\cdot)x)(b_i)=e_\lambda x$. For (iii), check it
on each $V[\mu]$: if $\mu=\lambda$, write $y=\varphi(l)$ for an equivariant
$\varphi:L\to V$ (which exists, sending $v_0$ to any highest weight vector of
weight $\lambda$: the graph of $Uv_0\to U w$ is generated by $(v_0,w)$ and its
first projection is injective by Lemma~\ref{i1:lem-cr} and $L_\lambda=Kv_0$);
equivariant maps intertwine invariant projectors, so by (c) for $L$,
$dP(b_i^*(\cdot)y)(b_i)=b_i^*(l)\varphi(b_i)$, summing to $y$. If $\mu\ne\lambda$,
$P(b_i^*(\cdot)y)$ is equivariant with values in $V[\mu]$, so it kills $v_0$ and
hence $L$. Now $b_i^*(\cdot)x\in\mathcal H$, and (P1) with (ii) gives
$p^{N_0(n)\delta(p,n,\ell(V)+\ell(\lambda))}e_\lambda x\in M$; sum over $\Sigma$.

\emph{$\mathcal A_{10}$.} Put $D=N_0(n)\delta(p,n,\ell(\lambda))$. $Y\in
V[\lambda]$. Let $H$ be the space of highest weight vectors of weight $\lambda$
(with $0$) and $h_1,\dots,h_k$ an $O$-basis of $M\cap H$ (P3); then
$Y=\sum_jy_j$ with $y_j\in Uh_j$, and for $u\in U_\Z$, $\pi(uy_j)=a_jh_j$ (by (P2)
for $Uh_j$). The hypothesis and (P3) give $\pi(uY)=\sum_ja_jh_j\in p^cM\cap
H=p^c\sum_jOh_j$, so $a_j\in p^cO$. Hence $p^{-c}y_j\in\nabla(\lambda,h_j)$
computed in $Uh_j$, and by $\mathcal A_8$ for $Uh_j$, $p^{D-c}y_j\in\Delta(h_j)\subseteq
M$. Sum over $j$.
\end{proof}

\begin{lemma}[Gelfand--Tsetlin counts; input (C) of $\mathcal A_{15}$]\label{i1:lem-counts}
For antitone $\mu\in\Z^n$ let $\mathrm{gt}(n,\mu)$ be defined by
$\mathrm{gt}(0,\cdot)=1$ and $\mathrm{gt}(n+1,\mu)=\sum_{\kappa\in\mathrm{box}(\mu)}\mathrm{gt}(n,\kappa)$,
$\mathrm{box}(\mu)=\{\kappa\in\Z^n:\mu_{i+1}\le\kappa_i\le\mu_i\}$.
\begin{enumerate}
\item $|\mathrm{St}(\mu)|=\mathrm{gt}(n,\mu)$.
\item (Pieri) $\sum_\nu\mathrm{gt}(n,\nu)=\binom{n+b-1}{b}\mathrm{gt}(n,\mu)$, the
 sum over $\nu$ with $\nu/\mu$ a horizontal strip of size $b$; likewise over
 $\nu$ with $\mu/\nu$ such a strip.
\item (Input (C)) The number of pairs $(\mu,I)$, $\mu$ a Pieri path and
 $I\in\mathrm{St}(\mu_N)$, is $\prod_{i<N}\binom{n+B_i-1}{B_i}$.
\end{enumerate}
\end{lemma}

\begin{proof}
(1) Write $T(c)=\hat I(c)$; $|T(c)|=h_\mu(c):=\mathrm{colht}(\mu,c)$. Induct on
$n$: removing the largest index $n$ from each $T(c)$ leaves sets of sizes $h(c)$
with $h$ nonincreasing and $h_\mu-1\le h\le h_\mu$; such $h$ are exactly the
$h_\kappa$, $\kappa\in\mathrm{box}(\mu)$ ($\kappa_i=\max\{c\in[\mu_{i+1},\mu_i]:i<h(c)\}$),
and for fixed $\kappa$ removal is a bijection onto the standard fillings of
$\kappa$ (the row condition is unaffected below $n$; the inverse adds $n$ where
$h_\kappa(c)<h_\mu(c)$). (2) Induct on $n$. Expanding $\mathrm{gt}(n+1,\nu)$
gives a sum over pairs $(\nu,l)$, $l\in\mathrm{box}(\nu)$. With
$\mathrm{Lo}_j=\max(\mu_j,l_j)$, $\mathrm{Lo}_n=\mu_n$, $\mathrm{Hi}_i=\min(\mu_i,l_i)$,
the reflection $\kappa_i=\mathrm{Lo}_{i+1}+\mathrm{Hi}_i-\nu_{i+1}$ is a bijection
onto the pairs $(\kappa,l)$ with $\kappa\in\mathrm{box}(\mu)$ and $l/\kappa$ a
strip of size $\le b$ (the interval conditions are exchanged, and
$\sum\mathrm{Lo}+\sum\mathrm{Hi}=|\mu|+|l|$ matches the sizes). By induction and
the hockey-stick identity $\sum_{c\le b}\binom{n+c-1}c=\binom{n+b}b$ the sum is
$\binom{n+b}{b}\mathrm{gt}(n+1,\mu)$. The dual case follows by
$\mu\mapsto(-\mu_{n-1-i})_i$. (3) Let $\mathcal C(i,\mu^0)$ be the set of strip
sequences $\mu_i=\mu^0,\dots,\mu_N$ with a standard filling of $\mu_N$. By (1),
(2) and induction on $N-i$,
$|\mathcal C(i,\mu^0)|=\mathrm{gt}(n,\mu^0)\prod_{l\ge i}\binom{n+B_l-1}{B_l}$; and
$\mathrm{gt}(n,0)=1$.
\end{proof}

\begin{lemma}[Input (I) of $\mathcal A_{15}$]\label{i1:lem-A15I}
For every field $k$, every $n\ge1$ and all letter types and degrees $B$, the
path polynomials $\Psi_{\mu,I}$ ($\mu$ a Pieri path, $I\in\mathrm{St}(\mu_N)$)
are linearly independent over $k$.
\end{lemma}

We use the notation of Definition~\ref{i1:def-pieri}. For a matrix $M$ and a
form $K$ put $(M\cdot K)(J)=\sum_{J'}K(J')\det M_{J',J}$, where $\det M_{J',J}$
is the minor with rows $J'$ and columns $J$ (increasing order), $0$ if
$|J'|\ne|J|$. The \emph{generic frame change} is the following. Let $k(g)$ be
the field of rational functions in the entries of a generic $n\times n$ matrix
$g$, and let $g$ act on $k(g)[X_{i,b}]$ by the algebra map sending the row
$X_i=(X_{i,b})_b$ to $X_ig$ for a vector letter and to $X_i(g^{-1})^T$ for a
covector letter. Grade polynomials by the torus weight, in which $X_{i,b}$ has
weight $e_b$ (vector) or $-e_b$ (covector). For $\mu\in\Z^n$ put
$\mathrm{ht}(\mu)=\sum_{j\le n}\sum_{a<j}\mu_a$; if $\nu\preceq\mu$ and
$\mathrm{ht}(\mu)\le\mathrm{ht}(\nu)$ then $\nu=\mu$. For a filling $I$ of an
antitone $\mu$ let $\mathrm{brk}(\mu,I)\in k(g)$ be the product over the
columns $c$ of the minors of $g$ ($c\ge1$), resp.\ of $(g^{-1})^T$ ($c\le0$),
with rows $I^{\rm can}(c)$ and columns $I(c)$, and for any filling $I'$ let
$\mathrm{iw}(I')_a=\#\{c\ge1:a\in I'(c)\}-\#\{c\le0:a\in I'(c)\}$.

\begin{proof}
The proof rests on three facts, proved below.
\begin{itemize}
\item[(EV)] Every nonzero torus-weight component of $g\cdot\Psi_{\mu,I}$ has a
 weight $\nu\preceq\mu_N$, and the component of weight $\mu_N$ is
 $\mathrm{brk}(\mu_N,I)\,\Psi_{\mu,I^{\rm can}}$.
\item[(BR)] For fixed antitone $\lambda$, the $\mathrm{brk}(\lambda,I)$,
 $I\in\mathrm{St}(\lambda)$, are linearly independent over $k$.
\item[(PP)] For fixed $\lambda$, the $\Psi_{\mu,I^{\rm can}}$ for distinct
 Pieri paths $\mu$ ending at $\lambda$ are linearly independent over $k(g)$.
\end{itemize}
By (BR) and (PP) the products $\mathrm{brk}(\lambda,I)\Psi_{\mu,I^{\rm can}}$
($I$ standard, $\mu$ ending at $\lambda$) are linearly independent over $k$: a
relation gives, by (PP), $\sum_Ia_{I,\mu}\mathrm{brk}(\lambda,I)=0$ for each
$\mu$, and then all $a_{I,\mu}=0$ by (BR). Now let $\sum_xa_x\Psi_x=0$ with some
$a_x\ne0$, and choose $x_0=(\mu^0,I^0)$ with $a_{x_0}\ne0$ and
$\mathrm{ht}(\mu^0_N)$ maximal. Apply $g$ and take the component of weight
$\mu^0_N$. A pair $x=(\mu,I)$ contributes only if $\mu^0_N\preceq\mu_N$ (EV),
and then $\mu_N=\mu^0_N$ by maximality. By the second half of (EV),
$\sum_{x:\mu_N=\mu^0_N}a_x\mathrm{brk}(\mu^0_N,I_x)\Psi_{\mu_x,I^{\rm can}}=0$,
so $a_{x_0}=0$, a contradiction.

\emph{Two identities for forms.} For every commutative ring $A$, form $K$ and
$x,y\in A^n$:
\[\text{(W)}\ M\cdot(K\wedge x)=(M\cdot K)\wedge(xM),\qquad
\text{(Ct)}\ M\cdot(K\lrcorner y)=(M\cdot K)\lrcorner\bigl(y(M^{-1})^T\bigr),\]
for every matrix $M$, resp.\ every invertible $M$. For (W), expand the
$J$-entry of the right side; by Laplace expansion along the added row $a$,
$\sum_{b\in J}\mathrm{sg}(J,b)\det M_{J',J\setminus b}M_{ab}$ is $0$ if $a\in J'$
and $\mathrm{sg}(J'\cup a,a)\det M_{J'\cup a,J}$ otherwise, and the substitution
$J''=J'\cup a$ turns the result into the left side. Since $M\cdot\mathbf 1=\mathbf 1$
and the indicator of $\{a_1<\dots<a_k\}$ is
$\mathbf 1\wedge e_{a_1}\wedge\dots\wedge e_{a_k}$, every form is built from
$\mathbf 1$ by sums, scalars and wedges, and an identity stable under these
steps holds for all forms (\emph{form induction}). For (Ct) one uses also
\[(K\wedge x)\lrcorner y=\Bigl(\sum_ax_ay_a\Bigr)K-(K\lrcorner y)\wedge x,\]
checked at each $J$: the diagonal terms add to $(\sum x_ay_a)K(J)$ and the
cross terms cancel in pairs by
$\mathrm{sg}(J,m)\mathrm{sg}(J\cup m,m')=-\mathrm{sg}(J,m')\mathrm{sg}(J\setminus m',m)$
($m\ne m'$). Then (Ct) follows by form induction: for $\mathbf 1$ both sides
vanish, and for $K\wedge x$ both sides reduce, by (W), the hypothesis and the
displayed identity, to $(\sum_ax_ay_a)M\cdot K-((M\cdot K)\lrcorner
y')\wedge xM$ with $y'=y(M^{-1})^T$, using $\sum_a(xM)_ay'_a=\sum_ax_ay_a$.

\emph{(E1) Column words under $g$.} $g\cdot w_c(J)=\sum_{J'}\det(M_c)_{J',J}\,w_c(J')$
with $M_c=g$ for $c\ge1$ and $M_c=(g^{-1})^T$ for $c\le0$. Indeed for $c\ge1$
the wedging letters are vectors (image $X_ig$) and the contracting letters are
covectors (image $X_i(g^{-1})^T$); for $c\le0$ the roles are exchanged and
$X_ig=X_i(M_c^{-1})^T$. So (W), (Ct) with $M=M_c$ and induction along the
recursion give $g\cdot w_c=M_c\cdot w_c$.

\emph{(E2) Weights.} By induction along the recursion, $w_c(J)$ is
torus-homogeneous of weight $\epsilon_c\mathbf 1_J$ ($\epsilon_c=1$ for $c\ge1$,
$-1$ for $c\le0$): a wedge step multiplies by $X_{i,m}$ of weight
$\epsilon_ce_m$ (the type of the letter matches the sign of $c$), a
contraction step by $X_{i,m}$ of weight $-\epsilon_ce_m$ and enlarges the set
by $m$. Outside a finite range of columns the height is constant, $w_c=\mathbf 1$
and $I(c)=\emptyset$; so $\Psi_{\mu,I'}$ is homogeneous of weight $\mathrm{iw}(I')$
for every filling $I'$ with the column sizes of $\mu_N$.

\emph{(E3) Dominance.} Let $\lambda$ be antitone, $h_c=\mathrm{colht}(\lambda,c)$
and $R_j=\{0..j-1\}$. Then $I^{\rm can}(c)=R_{h_c}$ ($c\ge1$), resp.\ its
complement ($c\le0$), and $\mathrm{iw}(I^{\rm can})=\lambda$. For $I'$ with the
same column sizes, $\sum_{m<j}\mathrm{iw}(I')_m=\sum_c\epsilon_c|R_j\cap I'(c)|$,
and column by column $\epsilon_c|R_j\cap I'(c)|\le\epsilon_c|R_j\cap I^{\rm can}(c)|$
(for $c\ge1$, $|R_j\cap I'(c)|\le\min(j,h_c)$; for $c\le0$,
$|R_j\cap I'(c)|\ge\max(0,j-h_c)$ by inclusion--exclusion), with equality for
$j=n$. So $\mathrm{iw}(I')\preceq\lambda$. If equality holds, then at $j=h_c$
every column difference vanishes, which forces $I'(c)=I^{\rm can}(c)$ for all
$c$.

\emph{Proof of (EV).} By (E1) and multiplicativity,
$g\cdot\Psi_{\mu,I}=\sum_{I'}\bigl(\prod_c\det(M_c)_{I'(c),I(c)}\bigr)\Psi_{\mu,I'}$;
terms with a non-square minor vanish, the others have weight
$\mathrm{iw}(I')\preceq\mu_N$ by (E2), (E3), and weight $\mu_N$ only for
$I'=I^{\rm can}$, with coefficient $\mathrm{brk}(\mu_N,I)$.

\emph{Proof of (BR).} (J) \emph{Jacobi's complementary minors:} for invertible
$M$, $\det(M)\det((M^{-1})^T)_{I,J}=\pm\det M_{I^c,J^c}$. With
$(\star K)(J)=(-1)^{\epsilon(J)}K(J^c)$, $\epsilon(J)=\sum_{j\in J}(n+1+j)$, one
checks $\star(K\wedge y)=(\star K)\lrcorner y$ from the sign rules, and then
$M\cdot\star K=\det(M)\star((M^{-1})^T\cdot K)$ by form induction with (W),
(Ct); evaluating on indicator forms gives (J). Choose $lo\le0\le hi$ with all
$\lambda_m\in[lo,hi]$. By (J) the factor of $\mathrm{brk}(\lambda,I)$ at
$c\le0$ is $\pm\det(g)^{-1}\det g_{R_{h_c},\hat I(c)}$, and at $c\ge1$ it is
$\det g_{R_{h_c},\hat I(c)}$. Hence
$\mathrm{brk}(\lambda,I)=\pm D\prod_{c\in(lo,hi]}\det g_{R_{|\hat I(c)|},\hat I(c)}$
with $D=\det(g)^{-\#\{c\in(lo,hi]:c\le0\}}$ independent of $I$; and $I\mapsto
T_I=(\hat I(c))_{c\in(lo,hi]}$ is injective, with $T_I$ semistandard (columns
read left to right have weakly increasing $a$-th entries) when $I$ is
standard. (A) \emph{Products of initial minors of semistandard column
sequences are linearly independent.} Order the positions $(a,b)$
lexicographically and monomials lexicographically by exponents. For strictly
increasing $r$, $q$, $\det X_{r,q}$ has leading monomial $\prod_iX_{r_i,q_i}$
with coefficient $1$ (for a permutation $\pi\ne1$ with least moved index $a$,
the term of $\pi$ has exponent $0$ at $(r_a,q_a)$ and agrees with the diagonal
before). So the product attached to $T$ has leading monomial
$\prod_t\prod_aX_{a,T_t(a)}$, which determines, for each row $a$, the
multiset of the $T_t(a)$; as $t\mapsto T_t(a)$ is nondecreasing, it determines
$T$. Distinct leading monomials with coefficient $1$ give independence.

\emph{Proof of (PP).} Induction on $N$; for $N=0$ there is one path and
$\Psi=1$. For $N+1$ letters grade by the exponent vector of the last letter's
variables $X_{N,\cdot}$, ordered lexicographically along the positions in
increasing order for a covector and in decreasing order for a vector. Put
$\Delta(\mu)=\prod_mX_{N,m}^{|\mu_{N+1,m}-\mu_{N,m}|}$. \emph{Claim:} the top
graded component of $\Psi_{\mu,I^{\rm can}}$ is
$\pm\Psi_{\mu|_N,I^{\rm can}}\Delta(\mu)$. Column by column: if the height does
not change at the last step, the word does not change; otherwise it changes by
one (horizontal strip), say between $t$ and $t+1$, the canonical sets of
$\mu_{N+1}$ and $\mu_N$ at $c$ differ exactly by $t$, and the last step is a sum
of terms $\pm w(\cdot)X_{N,m}$ in which every admissible $m\ne t$ is smaller
than $t$ in the chosen order. So the top part of the word is
$\pm w^{(N)}_c(I^{\rm can}(\mu_N)(c))X_{N,t}$, and the number of columns with
$t=m$ is $|\mu_{N+1,m}-\mu_{N,m}|$ by the strip condition. Now in a relation
$\sum_\mu a_\mu\Psi_{\mu,I^{\rm can}}=0$ take the largest grade among the
$\Delta(\mu)$ with $a_\mu\ne0$; the paths with that $\Delta$ have the same
$\mu_N$ (it is determined by $\mu_{N+1}=\lambda$, the exponents and the letter
type), distinct truncations, and their top components $\pm\Psi_{\mu|_N,I^{\rm can}}\Delta$
are independent by induction. So those $a_\mu$ vanish, a contradiction.
\end{proof}

\begin{lemma}[Clause $\mathcal A_{14}$]\label{i1:lem-A14}
Let $O$ be a discrete valuation ring with residue field $k_0$, and $E=E^S$ for
$S$ the set of all letters, with divided powers $E^{[k]}_{ab}$ as in
Lemma~\ref{i1:lem-A12}. For every Pieri path $\mu$, $\Psi_{\mu,I^{\rm can}}$ is
a highest weight vector of weight $\mu_N$ over $O$ and over $k_0$: it is
nonzero, all its monomials have weight $\mu_N$, and $E^{[k]}_{ab}\Psi_{\mu,I^{\rm can}}=0$
for $a<b$, $k\ge1$. For fixed $\mu_N$ these vectors are linearly independent
over $k_0$.
\end{lemma}

\begin{proof}
$I^{\rm can}$ is standard ($\hat I^{\rm can}(c)=R_{h_c}$ for every $c$ and
$h_c\ge h_{c+1}$), so non-vanishing and independence follow from the path
basis (Lemmas~\ref{i1:lem-A15}, \ref{i1:lem-A15I}), and the weight is $\mu_N$
by (E2), (E3). For the raising operators work over $\Z$. $E_{ab}$ is a
derivation with $E_{ab}X_{l,m}=[m=b]X_{l,a}$ (vector) and $-[m=a]X_{l,b}$
(covector). Put $(\varepsilon_mK)(J)=[m\in J]\mathrm{sg}(J,m)K(J\setminus m)$,
$(\iota_mK)(J)=[m\notin J]\mathrm{sg}(J,m)K(J\cup m)$, so $K\wedge x=\sum_mx_m\varepsilon_mK$
and $K\lrcorner x=\sum_mx_m\iota_mK$; they satisfy the exterior relations
($\varepsilon$'s and $\iota$'s anticommute among themselves,
$\varepsilon_a\iota_b+\iota_b\varepsilon_a=[a=b]$). Let
$\mathrm{Op}_c=\varepsilon_b\iota_a$ for $c\ge1$ and $-\varepsilon_a\iota_b$ for
$c\le0$. Then $E_{ab}(w^{(j)}_c(J))=(\mathrm{Op}_cw^{(j)}_c)(J)$ for all $j$: at
$j=0$ both sides vanish, and a step $K\mapsto\sum_mx_m\mathrm{op}_mK$ reduces,
by the derivation rule, to $\sum_m(E_{ab}x_m)\mathrm{op}_m=\sum_mx_m[\mathrm{Op}_c,\mathrm{op}_m]$,
which follows in each of the four cases (letter type, sign of $c$) from
$[\varepsilon_b\iota_a,\varepsilon_m]=[m=a]\varepsilon_b$ and
$[\varepsilon_b\iota_a,\iota_m]=-[m=b]\iota_a$. For $a<b$, $\mathrm{Op}_c$
vanishes at $I^{\rm can}(c)$ for every form: for $c\ge1$ it needs $b\in R_{h_c}$
and $a\notin R_{h_c}$; for $c\le0$ it needs $a\notin R_{h_c}$ and $b\in R_{h_c}$.
So $E_{ab}\Psi_{\mu,I^{\rm can}}=0$ over $\Z$; by Lemma~\ref{i1:lem-A12}(1),
$k!E^{[k]}_{ab}\Psi=E_{ab}^k\Psi=0$, so $E^{[k]}_{ab}\Psi=0$ over $\Z$, hence over
every ring.
\end{proof}

\begin{lemma}[Clause $\mathcal A_{16}$]\label{i1:lem-A16}
Let $R$ be a field or a discrete valuation ring (in particular $O$ or its
residue field), and $\theta_i,M_i\in\Z$ ($i<N$). A polynomial
$G\in\mathrm{Pol}_B$ satisfies $\theta_i\le v_{\{j\le i\}}(G)$ and
$\lambda_{1,\{j\le i\}}(G)\le M_i$ for all $i<N$ iff $G$ is in the $R$-span of
the $\Psi_{\mu,I}$ ($I$ standard) with $\theta_i\le(\mu_{i+1})_{n-1}$ and
$(\mu_{i+1})_0\le M_i$ for all $i$.
\end{lemma}

\begin{proof}
Let $W$ be the set of pairs $(\mu,I)$ and $W_{\theta,M}$ the subset in the
statement. We use the path basis (P) over $R$ (Lemmas~\ref{i1:lem-A15},
\ref{i1:lem-counts}, \ref{i1:lem-A15I}) and two further facts:
\begin{itemize}
\item[(Wd)] for $(\mu,I)\in W$ and $i<N$, $(\mu_{i+1})_{n-1}\le v_{\{k\le i\}}(\Psi_{\mu,I})$
 and $\lambda_{1,\{k\le i\}}(\Psi_{\mu,I})\le(\mu_{i+1})_0$ (over $\Z$);
\item[(FL)] over a field $F$, for $i\le N$, a lower set $\pi$ of dominant
 weights for $\preceq$ and $G\in\mathrm{Pol}_B$: $G$ is in the span of the
 $\Psi_{\mu,I}$ with $\mu_i\in\pi$ iff every dominant weight occurring in $M_G$
 (Lemma~\ref{i1:lem-A13}(2), $S=S_i=\{k<i\}$) lies in $\pi$.
\end{itemize}

\emph{Span implies bounds.} By Definition~\ref{i1:def-twist}, $\theta\le v_S(G)$
says that all $T$-exponents of $\mathrm{tw}^{(n-1)}_SG$ are $\ge\theta$, and
$\lambda_{1,S}(G)\le M$ that those of $\mathrm{tw}^{(0)}_SG$ are $\le M$. Exponent
sets of sums and scalar multiples lie in the union, and mapping coefficients
from $\Z$ to $R$ shrinks them. So (Wd) gives the bounds on the span.

\emph{Bounds imply span.} Over a field $F$: write $G=\sum_wc_w\Psi_w$ by (P).
Fix $i<N$ and let $\pi$ be the dominant $\nu$ with $\theta_i\le\nu_{n-1}$ and
$\nu_0\le M_i$, a lower set ($\nu\preceq\nu'$ gives $\nu_0\le\nu'_0$ and
$\nu_{n-1}\ge\nu'_{n-1}$). By Lemma~\ref{i1:lem-A13}(2) every weight occurring in
$M_G$ for $S=\{k<i+1\}$ lies in $\pi$, so by (FL) and independence the support of
$(c_w)$ lies in $\{\mu_{i+1}\in\pi\}$. Over a discrete valuation ring, pass to
the fraction field and use that $O$ injects into it.

\emph{Proof of (Wd).} Use the filtrations $\mathcal F^e_t$ of the proof of
Lemma~\ref{i1:lem-A25}(1) over $A_n$ with $R=\Z$; $S=\{k\le i\}$, $j_0=i+1$.
(1) For a twist index $j$ put $u_b=g_{bj}$, $v_m=(g^{-1})_{jm}$; then
$\sum_mu_mv_m=1$, $M_j(k)=1+(T^k-1)uv^T$, and for $l\in S$,
$\mathrm{tw}^{(j)}X_l=X_l+(T-1)v\sum_bu_bX_{l,b}$ (vector) or
$X_l+(T^{-1}-1)u\sum_bv_bX_{l,b}$ (covector). (2) With $W_y=\cdot\wedge y$,
$C_y=\cdot\lrcorner y$ and $\Lambda(u,v,\theta)=1+(\theta-1)W_vC_u$, the
exterior relations give, for $\theta\theta'=1$,
$\Lambda(K\wedge x)=\Lambda(K)\wedge(x+(\theta-1)v\sum_bu_bx_b)$ and
$\Lambda(K\lrcorner x)=\Lambda(K)\lrcorner(x+(\theta'-1)u\sum_bv_bx_b)$; also
$\Lambda K=K$ if $K$ is supported on $\emptyset$ and $\Lambda K=\theta K$ if $K$ is
supported on the full set. (3) Hence, with $\Lambda_c=\Lambda(u,v,T)$ ($c\ge1$)
and $\Lambda(v,u,T^{-1})$ ($c\le0$), $\mathrm{tw}^{(j)}w^{(k)}_c=\Lambda_cw^{(k)}_c$
for $k\le j_0$: along a path a vector letter only raises heights (wedge for
$c\ge1$, contraction for $c\le0$) and a covector only lowers them, and (1),
(2) match each case. Moreover $w^{(k)}_c$ is supported on sets of the column
size of $\mu_k$ at $c$ (induction). (4) $X$, $u$, $v$ lie in $\mathcal F^e_0$,
$T^k\in\mathcal F^e_{ek}$, and $W_y,C_y$ preserve levels for $y\in\mathcal F^e_0$;
so $\Lambda(u,v,\theta)K\in\mathcal F^e_{\min(a,0)}$ for $K\in\mathcal F^e_0$,
$\theta\in\mathcal F^e_a$. (5) With $x=(\mu_{j_0})_{n-1}$, $y=(\mu_{j_0})_0$,
$\epsilon(z,c)=[c\le z]$ ($c\ge1$) and $-[z<c]$ ($c\le0$), the four cases
(height $n$ or $0$ where $\Lambda_c$ is $\theta$ or $1$, Step (4) otherwise)
give $\mathrm{tw}^{(n-1)}w^{(j_0)}_c\in\mathcal F^1_{\epsilon(x,c)}$ and
$\mathrm{tw}^{(0)}w^{(j_0)}_c\in\mathcal F^{-1}_{-\epsilon(y,c)}$. (6) Later
letters are untwisted and keep the levels. (7) Multiplying over the columns,
$\sum_c\epsilon(x,c)=x$, so $x\le v_S(\Psi_{\mu,I})$ and likewise
$\lambda_{1,S}(\Psi_{\mu,I})\le y$.

\emph{Proof of (FL).} Write $\mathrm{wt}_i=\mathrm{wt}_{S_i}$ and use the
substitutions $\sigma^\phi_{h,h'}$ of the proof of Lemma~\ref{i1:lem-A13}(2)
with $S=S_i$. (1) \emph{Truncation:} $w^{(N)}_c(J)=\sum_{J'}w^{(i)}_c(J')Q_c(J',J)$
with $Q_c$ involving only letters $\ge i$. (2) For $gg'=1$ over any commutative
ring, every monomial $e$ of $\sigma_{g,g'}\Psi_{\mu,I}$ has $\mathrm{wt}_i(e)\preceq\mu_i$:
expand by (1) over fillings $I'$ with the column sizes of $\mu_i$; the first
factor is $\Psi_{\mu|_i,I'}$, on which $\sigma_{g,g'}$ is the full action, and
(EV) (whose proof works over every ring) bounds its weights; the second factor
does not change $\mathrm{wt}_i$. (3) \emph{Span implies the condition:} by the
argument of part (B) of the proof of Lemma~\ref{i1:lem-A13}(2) applied to the
property ``$\nu\preceq\nu'$ for some $\nu'\in\pi$'', (2) shows that every weight
occurring in $M_G$ has this property; dominant ones lie in $\pi$. (4) \emph{Top
case $i=N$:} if all dominant weights of the monomials of $g\cdot G$ (over
$F(g)$) lie in $\pi$, write $G=\sum c_x\Psi_x$; for $x_1$ in the support take
$x_0$ with $\mu^{x_1}_N\preceq\mu^{x_0}_N$ and $\mathrm{ht}$ maximal; as in the
proof of Lemma~\ref{i1:lem-A15I} the weight $\mu^{x_0}_N$ component of $g\cdot G$
is nonzero, so $\mu^{x_0}_N\in\pi$ and $\mu^{x_1}_N\in\pi$. (5) \emph{General
$i$:} write $G=\sum_\beta G_\beta X^\beta$ with $\beta$ monomials in the letters
$\ge i$; the hypothesis passes to each $G_\beta$ (the substitution fixes
$X^\beta$). If every dominant weight occurring in $M_G$ lies in $\pi$, then by
part (A) of the proof of Lemma~\ref{i1:lem-A13}(2) the hypothesis of (4) holds
for each $G_\beta$, so $G$ lies in the span $W'$ of the
$\Psi_{\mu',I'}X^\beta$ ($\mu'$ of length $i$, $\mu'_i\in\pi$). The span $U$ of
the $\Psi_{\mu,I}$ with $\mu_i\in\pi$ lies in $W'$ by (2), (4); by
Lemma~\ref{i1:lem-A15I} and the count $|\mathcal C(i,\mu^0)|$ of
Lemma~\ref{i1:lem-counts}(3), $\dim U$ equals the number of pairs $(\mu',I')$
times $\prod_{l\ge i}\binom{n+B_l-1}{B_l}$, which bounds $\dim W'$. So $U=W'\ni G$.
\end{proof}

\begin{lemma}[Clause $\mathcal A_{17}$: splitting of paths]\label{i1:lem-A17}
Let $r,K\ge1$. For sequences $\nu_i\in\Z^r$, $\kappa_i\in\Z^K$ ($i\le N$), the
concatenation $\mu_i=(\nu_i,\kappa_i)$ is a Pieri path of degrees $B$ in $r+K$
coordinates iff there are $b'_i+b''_i=B_i$ such that $\nu$ is a Pieri path of
degrees $b'$, $\kappa$ one of degrees $b''$, and at each letter $i$ the
\emph{interlacing condition} holds: $\kappa_{i+1,0}\le\nu_{i,r-1}$ for a vector,
$\kappa_{i,0}\le\nu_{i+1,r-1}$ for a covector.
\end{lemma}

\begin{proof}
A strip $(\alpha,\gamma)\to(\beta,\delta)$ of size $B$ splits into strips
$\alpha\to\beta$, $\gamma\to\delta$ of sizes $b'+b''=B$ plus the single linking
inequality $\delta_0\le\alpha_{r-1}$, and $(\alpha,\gamma)$ is antitone iff
both blocks are and $\gamma_0\le\alpha_{r-1}$. This gives ($\Rightarrow$). For
($\Leftarrow$), $\kappa_{i,0}\le\nu_{i,r-1}$ for all $i$ by induction (for a
vector letter from the interlacing and $\nu_{l,r-1}\le\nu_{l+1,r-1}$; for a
covector from $\kappa_{l+1,0}\le\kappa_{l,0}$), so each $\mu_i$ is antitone and
each step is a strip of size $B_i$.
\end{proof}

\noindent The formal clause $\mathcal A_{17}$ assumes in addition $N+r\le K$;
the statement above does not need it. Similarly, Lemmas~\ref{i1:lem-A7},
\ref{i1:lem-A19}, \ref{i1:lem-A36} and \ref{i1:lem-A37} are stated here in
less generality than the corresponding formal clauses, in the generality
in which they are used. (Propositions~\ref{i1:prop-A19c} and
\ref{i1:prop-A21} are, conversely, stated without the completeness
assumption of their formal clauses.)

\begin{lemma}[Clause $\mathcal A_{11}$: nested projections]\label{i1:lem-A11}
Let $O,K,p$ be as in Lemma~\ref{i1:lem-A7}, $\iota$ an alphabet, $n\ge1$,
$S_0\subseteq\dots\subseteq S_{N-1}\subseteq\iota$, and $E^S$ the operators of
Definition~\ref{i1:def-twist} on $K[\iota]_n$. Call $F$ \emph{integral} if its
coefficients lie in $O$ and \emph{Fock-integral} if
$v_p(c_m)+v_p(m!)/2\ge0$ for every coefficient $c_m$. Then: (1) both
conditions are stable under the Kostant generators of every $E^S$; (2)
isotypic projectors $e_{i,\lambda}$ of $E^{S_i}$ commute with $E^{S_j}$
($i\le j$), with $E^{S'}$ for $S'\cap S_i=\emptyset$, and with each other, and
are self-adjoint for the Fock pairing; (3) for $d$ a multidegree and a finite set
$\mathcal T$ of tuples, the operator
$p^{\sum_iN_0(n)\lfloor\log_p(2\deg_{S_i}(d)+n-1)\rfloor}\sum_{t\in\mathcal T}\prod_ie_{i,t_i}$
maps integral, resp.\ Fock-integral, elements of $\mathrm{Pol}_d$ to integral,
resp.\ Fock-integral, elements.
\end{lemma}

\begin{proof}
(1) The letter operators $D^l_{ab}$ ($E^S_{ab}=\sum_{l\in S}D^l_{ab}$) commute,
so $(E^S_{ab})^{(k)}=\sum_f\prod_l(D^l_{ab})^{(f_l)}$, and
$(D^l_{ab})^{(f)}X^m=\pm\binom{\beta}{f}X^{m'}$ where $m'$ moves $f$ units from
a source exponent $\beta$ to a target exponent $\alpha$; $\binom{E^S_{aa}}k$
multiplies $X^m$ by $\binom{\mathrm{wt}_S(m)_a}k$. This gives integrality; for
Fock-integrality use $m!\binom{\alpha+f}f=m'!\binom\beta f$, so
$v_p(m!)\le v_p(m'!)+2v_p\binom\beta f$. (2) Operators of disjoint letter sets
commute; an operator commuting with all $E_{ab}$ preserves isotypic parts, and
an operator preserving all isotypic parts commutes with every isotypic
projector. With $E^{S_j}=E^{S_i}+E^{S_j\setminus S_i}$ this gives the
commutations, as in Lemma~\ref{i1:lem-A12}(4). Self-adjointness: the Fock pairing
is contravariant (Lemma~\ref{i1:lem-A12}(2)); for highest weight vectors $v,h$
of weights $\lambda\ne\mu$, $\langle v,h\rangle=0$ (compare $E_{aa}$-eigenvalues)
and $\langle v,E_{cd}y\rangle=\langle E_{dc}v,y\rangle=0$ for $d<c$, so
$V[\lambda]\perp V[\mu]$ by Lemma~\ref{i1:lem-hw}(3), and isotypic projectors are
self-adjoint. (3) The integral and the Fock-integral polynomials of
$\mathrm{Pol}_d$ are finitely generated $O$-modules spanning $\mathrm{Pol}_d$ and
stable under the Kostant generators by (1). By Lemma~\ref{i1:lem-A12}(3),(4),
$\ell(\mathrm{Pol}_d)+\ell(\lambda)\le2\deg_{S_i}(d)$ when $e_{i,\lambda}$ does not
vanish on $\mathrm{Pol}_d$, so clause $\mathcal A_9$ (Lemma~\ref{i1:lem-A8}) for
the restriction of $E^{S_i}$ to $\mathrm{Pol}_d$ (an isotypic projector restricts
to one) shows that $p^{N_0(n)\lfloor\log_p(2\deg_{S_i}(d)+n-1)\rfloor}e_{i,\lambda}$
preserves both lattices. Each factor preserves $\mathrm{Pol}_d$; multiply the
factors and sum over $\mathcal T$.
\end{proof}

\subsubsection*{Column polynomials and the lifted filtration: clauses $\mathcal A_{18}$--$\mathcal A_{21}$}

\begin{definition}[Plane and tangent parts]\label{i1:def-pltan}
Fix $r,K\ge1$, $k=r+K$, and letters $0..N-1$ with types, $N+r\le K$. The
indices $a<r$ are \emph{plane}, the indices $a\ge r$ \emph{tangent}; for
$J\subseteq\{0..k-1\}$ put $J_p=J\cap[0,r)$ and $J_t=\{b:r+b\in J\}$. The
\emph{plane} and \emph{tangent embeddings} are $\rho^{pl}:X_{l,a}\mapsto X_{l,a}$
($a<r$) and $\rho^{tan}:X_{l,b}\mapsto X_{l,r+b}$ ($b<K$). For Pieri paths $\nu$
in $r$ coordinates and $\kappa$ in $K$ coordinates (the \emph{plane} and
\emph{tangent} parts of their concatenation $\mu$, when it is a Pieri path) and
fillings $I',I''$ put $\Psi^{pl}_{\nu,I'}=\rho^{pl}\Psi_{\nu,I'}$,
$\Psi^{tan}_{\kappa,I''}=\rho^{tan}\Psi_{\kappa,I''}$. For a path $\mu$ in $k$
coordinates and a column $c$ write $h_i=\mathrm{colht}(\mu_i,c)$; letter $i$ is
\emph{used} at $c$ if $h_i\ne h_{i+1}$ and \emph{costly} at $c$ if
($h_i<h_{i+1}$ and $r\le h_i$) or ($h_{i+1}<h_i$ and $h_i\le r$). Given weights
$\omega_0>\dots>\omega_{N-1}>0$, give $X_{l,a}$ the weight $\omega_l$ if ($l$ is a
vector and $a\ge r$) or ($l$ is a covector and $a<r$), and $0$ otherwise; put
$\mathrm{oc}(c)=\sum_{i\text{ costly at }c}\omega_i$, and let the \emph{tangent
column} be $\mathrm{pt}_c(J)=\rho^{pl}(w^\nu_c(J_p))\,\rho^{tan}(w^\kappa_c(J_t))$,
where $w^\nu_c$, $w^\kappa_c$ are the column words of $\nu$ and $\kappa$.
\end{definition}

\begin{lemma}[Clause $\mathcal A_{18}$]\label{i1:lem-A18}
For every Pieri path $\mu$ in $k$ coordinates, every column $c$ and every
$J$, all monomials of the column word $w_c(J)$ of $\mu$ have weight
$\ge\mathrm{oc}(c)$, and its component of weight $\mathrm{oc}(c)$ is
$\mathrm{pt}_c(J)$.
\end{lemma}

\begin{proof}
Say $f$ has \emph{order} $\ge b$ if all its monomials have weight $\ge b$, and
write $[f]_o$ for its weight-$o$ part. For $a<r$ in $J$,
$\#\{j\in J:j>a\}=\#\{a'\in J_p:a'>a\}+|J_t|$, and for $r+b\in J$ the count
only involves $J_t$; so in the expansions of $(K\wedge x)(J)$ and $(K\lrcorner
x)(J)$ a tangent term has the sign computed in $J_t$ and a plane term the sign
computed in $J_p$ times $(-1)^{|J_t|}$. Since $\mu_i$ is antitone, the plane
part has height $\min(h_i,r)$ and the tangent part $(h_i-r)^+$; a rise
$h_i<h_{i+1}=h_i+1$ happens only at a vector letter, a fall only at a
covector. Write $o_j=\sum_{i<j\text{ costly}}\omega_i$, $p^{(j)},t^{(j)}$ for the
stages of $w^\nu_c,w^\kappa_c$ (embedded).

\emph{Case $c\ge1$} ($h_0=0$). By induction on $j$: (A$_j$) for every $u\ge0$
with $u\le\omega_i$ ($i<j$), $w^{(j)}(J)$ has order
$\ge o_j+(\min(h_j,r)-|J_p|)^+u$; (D$_j$) $[w^{(j)}(J)]_{o_j}=p^{(j)}(J_p)t^{(j)}(J_t)$.
In the step $j\to j+1$ apply (A$_j$) with $u=\omega_j+1$ (allowed since
$\omega_i>\omega_j$). There are four cases. A rise with $h_j<r$ (wedge, vector,
not costly): plane terms have weight $0$ and raise $\min(h,r)$ and $|J_p|$
together, tangent terms carry $\omega_j$ and have order $>o_j$; only plane
terms with $J_t=\emptyset$ survive in the weight-$o_j$ part, giving
$(p^{(j)}\wedge X^{pl}_j)(J_p)t^{(j)}(J_t)$. A rise with $h_j\ge r$ (costly):
plane terms have order $\ge o_j+\omega_j+1$ by (A$_j$), tangent terms
contribute $p^{(j)}(J_p)(t^{(j)}\wedge X^{tan}_j)(J_t)$ at weight $o_{j+1}$. A
fall with $h_j>r$ (contraction, covector, not costly): plane terms carry
$\omega_j$ and have order $>o_j$; tangent terms give
$p^{(j)}(J_p)(t^{(j)}\lrcorner X^{tan}_j)(J_t)$. A fall with $h_j\le r$ (costly):
plane terms give $(p^{(j)}\lrcorner X^{pl}_j)(J_p)t^{(j)}(J_t)$ at weight
$o_{j+1}$; a tangent term $X_{j,r+b}w^{(j)}(J\cup\{r+b\})$ is nonzero only if
$|J|+1=h_j$, and then (A$_j$) gives it order $\ge o_j+(h_j-|J_p|)(\omega_j+1)>o_{j+1}$.
In each case the bounds (A$_{j+1}$) follow by the same estimates.

\emph{Case $c\le0$} ($h_0=k$). As $h$ drops by at most one per letter,
$h_j\ge k-N\ge r$, so $p^{(j)}=\mathbf 1$; the analogous induction with
(A$'_j$): order $\ge o_j+|J_p|u$, applies to a rise (contraction, costly; only
tangent terms survive) and a fall (wedge, not costly; plane terms carry
$\omega_j$). At $j=N$, $u=0$ the claims follow.
\end{proof}

\begin{definition}[Joint types]\label{i1:def-jt}
Let $F$ be a field of characteristic $0$, $S$ a set of letters and $e$ an
injection of $\{0..n-1\}$ into $\{0..k-1\}$ (the plane embedding $a\mapsto a$,
$n=r$, or the tangent embedding $b\mapsto r+b$, $n=K$). Let
$E^{S,e}_{ab}=\sum_{l\in S\text{ vector}}X_{l,e(a)}\partial_{l,e(b)}-\sum_{l\in S\text{ covector}}X_{l,e(b)}\partial_{l,e(a)}$,
and $\mathrm{iso}(S,e,\lambda)$ the isotypic part of weight $\lambda$ for this
representation. The \emph{joint type} of a pair $(\nu,\kappa)$ of sequences is
$\mathrm{JT}(\nu,\kappa)=\bigcap_{i\le N}\mathrm{iso}(\{l<i\},\mathrm{pl},\nu_i)\cap\mathrm{iso}(\{l<i\},\mathrm{tan},\kappa_i)$.
\end{definition}

\begin{lemma}[Clause $\mathcal A_{19}$, parts (a), (b)]\label{i1:lem-A19}
Let $O$ be a discrete valuation ring of characteristic $0$, $\mathcal P$ a set
of pairs $(\nu,\kappa)$ whose concatenation is a Pieri path of degrees $B$ in
$k$ coordinates, and $X(\mathcal P)$ the $O$-span of the products
$\Psi^{pl}_{\nu,I'}\Psi^{tan}_{\kappa,I''}$ ($(\nu,\kappa)\in\mathcal P$, $I',I''$
standard). (a) $X(\mathcal P)$ is saturated in $O[X]$. (b) For
$(\nu,\kappa)\in\mathcal P$ and standard $I',I''$, the filling $I(c)=I'(c)\sqcup
(r+I''(c))$ of the concatenation $\mu$ has: every monomial of $\Psi_{\mu,I}$ has
weight $\ge\sum_c\mathrm{oc}(c)$, and the component of that weight is
$\Psi^{pl}_{\nu,I'}\Psi^{tan}_{\kappa,I''}$.
\end{lemma}

\begin{proof}
(b) By Lemma~\ref{i1:lem-A18} each factor $w_c(I(c))$ has order
$\ge\mathrm{oc}(c)$ with lowest part $\mathrm{pt}_c(I(c))$, and lowest parts of
products multiply. Outside their ranges the column words of $\nu,\kappa$ are
$\mathbf 1$ and $I',I''$ are empty, so the product of the $\mathrm{pt}_c$ is
$\Psi^{pl}_{\nu,I'}\Psi^{tan}_{\kappa,I''}$. (a) By Lemma~\ref{i1:lem-A17}
every element of $\mathcal P$ is a pair of Pieri paths of degrees $b'+b''=B$.
The products over all such pairs and standard fillings form a basis of
$\mathrm{Pol}_B$ (in $k$ coordinates): $\mathrm{Pol}_B=\bigoplus_{b'+b''=B}\rho^{pl}\mathrm{Pol}_{b'}\cdot\rho^{tan}\mathrm{Pol}_{b''}$,
each factor has its path basis (Lemmas~\ref{i1:lem-A15}, \ref{i1:lem-A15I}),
and products of independent families in disjoint variables are independent
(compare coefficients of monomials in one set of variables). If $a\ne0$ and
$af\in X(\mathcal P)$, then $f\in\mathrm{Pol}_B$; expanding $f$ in this basis and
comparing with $af$, the coefficients of $f$ off $\mathcal P$ are killed by $a$,
hence vanish.
\end{proof}

\begin{proposition}[Clause $\mathcal A_{19}$(c)]\label{i1:prop-A19c}
In the setting of Lemma~\ref{i1:lem-A19}, with $k_0$ the fraction field of $O$,
let $\mathcal P$ be the set of pairs whose concatenation is a Pieri path of
degrees $B$ and which satisfy each inequality of a finite list, each of one of
the forms $\theta\le\nu_{i,r-1}$, $\kappa_{i,0}\le M$,
$\kappa_{j,0}+m\le\nu_{i,r-1}$. Then
$X(\mathcal P)=\{f\in O[X]:\bar f\in\sum_{q\in\mathcal P}\mathrm{JT}(q)\}$, $\bar f$
the image of $f$ over $k_0$.
\end{proposition}

\noindent The formal clause assumes in addition that $O$ is complete with
residue characteristic $p$; the proof below does not use it.

\noindent The proof is given in \S\ref{i1:ssec-pf-A19c}; it combines a
dimension count for commuting isotypic decompositions with the path bases
and an evaluation criterion for the plane and tangent operator families.

\begin{definition}[Lifted substitution and the $w$-filtration]\label{i1:def-lift}
Let $W=W(\mathbb F_{p^f})$ with Frobenius $\sigma$, $\sigma^{(s)}=\sigma^{s\bmod f}$;
let $O$ be a $W$-algebra which is a discrete valuation ring with uniformiser
$w$, $p\in wO$, and fraction field $k_0$; let the letters carry types and
\emph{positions} $\mathrm{pos}(l)\in\Z$. For $g\in M_k(W)$ with $\det g$ a unit,
$g^{(s)}$ is the image of $\sigma^{(s)}(g)$ in $M_k(k_0)$. Put $\delta(l,b)=1$ if
($l$ is a vector and $b\ge r$) or ($l$ is a covector and $b<r$), else $0$. The
\emph{lifted substitution} $L_g$ is the $k_0$-algebra map
$X_{l,a}\mapsto\sum_bw^{\delta(l,b)}X_{l,b}((g^{(\mathrm{pos}\,l)})^{-1})_{ba}$
(vector), $X_{l,a}\mapsto\sum_bw^{\delta(l,b)}X_{l,b}\,g^{(\mathrm{pos}\,l)}_{ab}$
(covector). $F$ lies in the \emph{$j$-th lifted $w$-filtration} $W^j_g$ if all
coefficients of $L_gF$ lie in $w^jO$. The \emph{rational plane}
$\mathrm{rp}(g)\subseteq\mathbb F_{p^f}^k$ is the span of the rows $i<r$ of the
reduction of $g^{-1}$.
\end{definition}

\begin{lemma}[Clause $\mathcal A_{20}$]\label{i1:lem-A20}
(1) If $\mathrm{rp}(g)=\mathrm{rp}(g')$ then $W^0_g=W^0_{g'}$. (2) If $F$ has
coefficients in $O$, the coefficient of $X^m$ in $L_gF$ lies in
$w^{\sum\delta(l,b)m_{l,b}}O$.
\end{lemma}

\begin{proof}
For an additive $\delta':\N^{\rm vars}\to\N$ the polynomials whose
$m$-coefficient lies in $w^{\delta'(m)}O$ form an $O$-subalgebra containing
$X_v$ ($\delta'(e_v)=0$) and $wX_v$ ($\delta'(e_v)=1$); so an algebra map sending
each variable into it sends integral polynomials into it. (2) The entries of
$(g^{(s)})^{\pm1}$ lie in $O$, so each $L_gX_{l,a}$ lies in the subalgebra for
$\delta'=\sum\delta m$. (1) If $\mathrm{rp}(g)\subseteq\mathrm{rp}(g')$, $b<r\le c$,
then $(g^{-1}g')_{bc}\equiv0\bmod p$ (the functional $x\mapsto\sum_jx_j\bar
g'_{jc}$ kills $\mathrm{rp}(g')$, hence row $b$ of $\bar g^{-1}$), and its image
under $\sigma^{(s)}$ lies in $pO\subseteq wO$. Let $N_l(c,b)$ be the image of
$\sigma^{(\mathrm{pos}\,l)}((g'^{-1}g)_{cb})$ (vector) or of
$\sigma^{(\mathrm{pos}\,l)}((g^{-1}g')_{bc})$ (covector). In all cases there is
$A_l(c,b)\in O$ with $w^{\delta(l,b)}A_l(c,b)=w^{\delta(l,c)}N_l(c,b)$: the only
nontrivial case $\delta(l,b)=1$, $\delta(l,c)=0$ is a plane-tangent block entry,
divisible by $w$ as just shown. The integral substitution
$\Theta:X_{l,b}\mapsto\sum_cX_{l,c}A_l(c,b)$ satisfies $\Theta\circ L_g=L_{g'}$
(a direct computation), so $L_gF$ integral implies $L_{g'}F$ integral; and
symmetrically.
\end{proof}

\begin{definition}[Canonical translate column]\label{i1:def-ctc}
In the notation of Definition~\ref{i1:def-lift}, with $p\ge5$, $f\ge1$,
$N\ge1$, $w^2=p$, $N+r\le K$,
positions monotone with $\mathrm{pos}(i)\le z$, a Pieri path $\mu$ in $k$
coordinates, $g\in GL_k(W)$ and a column $c$ such that
$\mathrm{pos}(i)<\mathrm{pos}(j)$ whenever $i<j$, $i$ is used and $j$ is costly
at $c$: put $\omega_i=p^{z-\mathrm{pos}(i)}$ (so $\mathrm{oc}=\mathrm{oc}(c)$ is
defined), $x_j(b)=\sum_aX_{j,a}g^{(\mathrm{pos}\,j)}_{ab}$ (vector),
$x_j(b)=\sum_aX_{j,a}((g^{(\mathrm{pos}\,j)})^{-1})_{ba}$ (covector). The
\emph{canonical translate column} is $Z_N(z)$, where $Z_0(s)=\mathbf 1$,
$Z_{j+1}(s)=Z_j(s)$ if $j$ is not used, and otherwise
$Z_{j+1}(s)(J)=\bigl(\mathrm{step}_j(x_j,Z_j(\mathrm{pos}\,j))(J)\bigr)^{p^{s-\mathrm{pos}(j)}}$,
$\mathrm{step}_j$ being the column-word step ($\wedge x_j$ or $\lrcorner x_j$) at
letter $j$. $\mathrm{Pow}_z$ is the $O$-algebra map $X_{l,a}\mapsto X_{l,a}^{p^{z-\mathrm{pos}(l)}}$.
\end{definition}

\begin{lemma}[Clause $\mathcal A_{21}$]\label{i1:lem-A21}
For every $J$: (1) $Z_N(z)(J)$ has coefficients in $O$; (2) every monomial of
$Z_N(z)(J)$ has letter degree $p^{z-\mathrm{pos}(i)}$ at each used letter $i$
(and $0$ at the others; this addition is not part of the clause
$\mathcal A_{21}$, but it is proved below together with it); (3) $Z_N(z)(J)\in W^{\mathrm{oc}}_{g'}$ for every
$g'\in GL_k(W)$; (4) $L_g(Z_N(z)(J))-w^{\mathrm{oc}}\mathrm{Pow}_z(\mathrm{pt}_c(J))$
has coefficients in $w^{\mathrm{oc}+1}O$; (5) $\mathrm{oc}\le2p^{z-\mathrm{pos}(0)}$.
\end{lemma}

\begin{proof}[Proof of (1), (2), (5)]
(1) The $x_j(b)$ are integral, and the steps and powers preserve integrality.
(2) Letter-homogeneity is additive under products and multiplicative under
powers; $x_j(b)$ has letter degree $e_j$; by induction (using monotone
positions), $Z_j(s)$ has degree $p^{s-\mathrm{pos}(i)}$ at each used $i<j$, since
$p^{s-\mathrm{pos}(j)}p^{\mathrm{pos}(j)-\mathrm{pos}(i)}=p^{s-\mathrm{pos}(i)}$.
(5) Costly letters are used, so distinct costly letters have distinct exponents
$z-\mathrm{pos}(i)\in[0,E]$, $E=z-\mathrm{pos}(0)$, and
$\mathrm{oc}\le\sum_{e\le E}p^e\le2p^E$.
\end{proof}

\begin{proposition}[$\mathcal A_{21}$(3),(4)]\label{i1:prop-A21}
Parts (3) and (4) of Lemma~\ref{i1:lem-A21} hold.
\end{proposition}

\noindent The formal clause assumes in addition that $O$ is complete with
residue characteristic $p$; the proof below does not use it.

\noindent The proof is given in \S\ref{i1:ssec-pf-A21}: an induction along
the column with a twisting structure comparing $Z_j$ with its
Frobenius-twisted image, and a reduction modulo $w$ of the translate column.

\subsubsection*{Segment modules: clause $\mathcal A_{33}$}

\begin{definition}[Segment modules]\label{i1:def-seg}
Fix a frame system (Definition~\ref{i1:def-frames}) on $\mathrm{TT}$, $r\ge1$, $K$,
$k=r+K$, and type the letters $x\in\Z$ by: $x$ is a vector iff $x$ is even. Let
$\mathcal S_O=O[X_{(x,c)}:x\in\Z,\ c<k]$, with plane and tangent embeddings
$\iota_{pl}$, $\iota_{tan}$ as in Definition~\ref{i1:def-pltan}. A letter $x$ of
segment $j$ is a \emph{$W$-letter} if $x\le\mathrm{fr}(j-1,r-1)$ and \emph{late}
otherwise. For an exponent $m$ put $\mathrm{od}(m,x)=\sum_{b<K}m(x,r+b)$ for $x$
even and $\sum_{a<r}m(x,a)$ for $x$ odd. For $m:\Z\to\Z$ and $\theta,o:\Z\to\R$,
the \emph{segment module} $\mathcal X'_j(m,\theta,o)$ is the intersection of
\begin{itemize}
\item the module of polynomials all of whose monomials satisfy
 $\mathrm{od}(\cdot,x)\le o(x)$ for every late letter $x$ of $j$, and
\item the $O$-span of the products $\iota_{pl}(A)\iota_{tan}(G)$, with $A$ a
 polynomial in the $X_{(x,a)}$, $x$ a letter of $j$, and $G$ one in the
 $X_{(x,b)}$, $x$ a late letter of $j$, for which there is $M:\Z\to\Z$ with,
 for every late letter $x$ of $j$: if $x$ is even, $\lambda_{1,\{y\le x\}}(G)\le
 M(x)$ and $M(x)+m(x)\le v_{\{y<x\}}(A)$; if $x$ is odd,
 $\lambda_{1,\{y<x\}}(G)\le M(x)$ and $M(x)+m(x)\le v_{\{y\le x\}}(A)$; and
 $\theta(x)\le v_{\{y<x\}}(A)$ when the latter is finite.
\end{itemize}
Let $e_0<\dots<e_{N-1}$ be the letters of segment $j$ and $\rho$ the renaming
$X_{(i,c)}\mapsto X_{(e_i,c)}$. A pair $(\nu,\kappa)$ of Pieri paths (in $r$,
resp.\ $K$ coordinates; letter $i$ of the type of $e_i$) of degrees $(b',b'')$
is \emph{admissible} if for every $i<N$: if $e_i$ is a $W$-letter, $b''_i=0$;
otherwise $\theta(e_i)\le\nu_{i,r-1}$, and $\kappa_{i+1,0}+m(e_i)\le\nu_{i,r-1}$,
$b''_i\le o(e_i)$ if $e_i$ is even, $\kappa_{i,0}+m(e_i)\le\nu_{i+1,r-1}$,
$b'_i\le o(e_i)$ if $e_i$ is odd.
\end{definition}

\begin{lemma}[Clause $\mathcal A_{33}$]\label{i1:lem-A33}
Let $O$ be a discrete valuation ring with residue field $k_O$ and fraction
field $k_0$, the frame system valid at $(L_1,N_1)$, $j$ a frame index with
$N+r\le K$, $m\ge0$. Then:
\begin{itemize}
\item[(i)] $\mathcal X'_j(m,\theta,o)$ is $\rho$ of the $O$-span of the
 products $\Psi^{pl}_{\nu,I'}\Psi^{tan}_{\kappa,I''}$ over admissible pairs and
 standard $I',I''$;
\item[(ii)] the renamed products over all pairs of Pieri paths and standard
 fillings are linearly independent over $O$;
\item[(iii)] $\mathcal X'_j(m,\theta,o)$ is saturated;
\item[(iv)] if $k_0$ has characteristic $0$, $f\in\mathcal X'_j(m,\theta,o)$ iff
 $\bar f\in\rho\bigl(\sum_{\text{admissible}}\mathrm{JT}(\nu,\kappa)\bigr)$;
\item[(v)] every product $\iota_{pl}(A)\iota_{tan}(G)$ as in
 Definition~\ref{i1:def-seg} which also satisfies the degree bounds lies in
 $\mathcal X'_j(m,\theta,o)$;
\item[(vi)] over $k_O$, such a product is the reduction of an element of
 $\mathcal X'_j(m,\theta,o)$.
\end{itemize}
\end{lemma}

\begin{proof}
(v) is the definition. (ii) The $\Psi_{\mu,I}$ over all degrees are
independent (Lemmas~\ref{i1:lem-A15}, \ref{i1:lem-A15I}, and disjoint
multidegree supports); products of independent families in disjoint variables
are independent; $\rho$ is injective.

(i) Write $v_i(A)=v_{\{k<i\}}(A)$, $\lambda_i(A)=\lambda_{1,\{k<i\}}(A)$ for
polynomials in the letters $0..N-1$; renaming gives $v_{\{y\le e_i\}}(\rho A)=v_{i+1}(A)$
and $v_{\{y<e_i\}}(\rho A)=v_i(A)$ (Lemma~\ref{i1:lem-A35}). (a) By
Lemma~\ref{i1:lem-A16}, $(\mu_i)_{\rm last}\le v_i(\Psi_{\mu,I})$ and
$\lambda_i(\Psi_{\mu,I})\le(\mu_i)_0$. (b) \emph{Thresholds:} for $A\in\mathrm{Pol}_B$
and finite sets $B_v,B_l\subseteq\Z$, $A$ lies in the span of the $\Psi_{\mu,I}$
such that every $X\in B_v$ with $X\le v_i(A)$ has $X\le(\mu_i)_{\rm last}$ and
every $Y\in B_l$ with $\lambda_i(A)\le Y$ has $(\mu_i)_0\le Y$: apply
Lemma~\ref{i1:lem-A16} with $\theta_{i-1}$ the largest admissible threshold
$\le v_i(A)$ and $M_{i-1}$ the smallest $\ge\lambda_i(A)$. (c) For a product as in
(v) write $A=\rho A_0$, $G=\rho G_0$ and split into multihomogeneous parts
(which only improve $v$ and $\lambda$); applying (b) to $A_{b'}$ with
$B_v=\{M(e_i)+m(e_i)\}\cup\{\lceil\theta(e_i)\rceil\}$ and to $G_{b''}$ with
$B_l=\{M(e_i)\}$, every pair of generators that occurs satisfies the inequalities
of admissibility (with $b''_i=0$ at $W$-letters as $G$ has no such variables).
(d) A renamed product of degrees $(b',b'')$ has $\mathrm{od}(\cdot,e_i)=b''_i$
($e_i$ even), resp.\ $b'_i$ ($e_i$ odd), on every monomial; so the projection
onto the monomials satisfying the degree bounds keeps admissible products and
kills the others. Together: every element of $\mathcal X'_j$ is in the
admissible span. Conversely, for admissible $(\nu,\kappa)$ take
$M(e_i)=\kappa_{i+1,0}$ ($e_i$ even), $\kappa_{i,0}$ ($e_i$ odd); by (a) the
renamed product satisfies the inequalities of (v) and, by (d), the degree
bounds.

(iii) By (i) and (ii), $\mathcal X'_j$ is the span of a subfamily of an
independent family that spans all polynomials in the renamed variables
(path bases in every degree). If $a\ne0$ and $af\in\mathcal X'_j$, $f$ is a
renamed polynomial, and comparing the expansions of $f$ and $af$ shows that the
coefficients of $f$ outside the subfamily vanish.

(iv) \emph{Bihomogeneity:} in characteristic $0$, $\mathrm{JT}(\nu,\kappa)$ consists
of polynomials of plane degree $b'_i$ and tangent degree $b''_i$ at every
letter: the $E^{S,e}_{aa}$-weight of an isotypic part is a sum of signed
degrees, and comparing $S=\{k<i+1\}$ with $\{k<i\}$ isolates the degree at
letter $i$, which is $\sum_a((\nu_{i+1})_a-(\nu_i)_a)=\pm b'_i$. \emph{Comparison
of pair sets:} by Lemma~\ref{i1:lem-A17}, $(\nu,\kappa)$ is admissible iff it lies
in the set $Q_{b'+b''}$ of pairs whose concatenation is a Pieri path and which
satisfy the segment inequalities ($\kappa_{i+1,0}+m(e_i)\le\nu_{i,r-1}$ resp.\
$\kappa_{i,0}+m(e_i)\le\nu_{i+1,r-1}$, and $\lceil\theta(e_i)\rceil\le\nu_{i,r-1}$,
at the late letters), and $(b',b'')$ satisfies the order bounds and $b''_i=0$ at
$W$-letters. The interlacing condition at a late letter follows from the
segment inequality and $m\ge0$; at a $W$-letter $e_i$ all $e_l$, $l\le i$, are
$W$-letters, hence frame points of segment $j-1$ by (Fr5), hence even by
(Fr4), and then $\kappa_{i+1}=0\le\nu_{i,r-1}$. Now project to bidegree
$(b',b'')$: by bihomogeneity the projection fixes the joint types and the
products of that bidegree and kills the others. For $f$ in the admissible
span, Proposition~\ref{i1:prop-A19c} applied to $Q_{b'+b''}$ and the projection
put $\bar f$ in the sum of the admissible joint types; conversely, each
projection of $f$ with an order-admissible bidegree lies in $X(Q_{b'+b''})$ by
Proposition~\ref{i1:prop-A19c} and hence, projecting again, in the admissible
span, and the other projections vanish since $O\to k_0$ is injective.

(vi) By Lemma~\ref{i1:lem-A16} over $k_O$ and the argument of (i), the product
lies in the $k_O$-span of the admissible renamed products; these are
reductions of the same products over $O$, and sums and scalar multiples lift.
\end{proof}

\subsubsection*{Flat tops and the test side: clauses $\mathcal A_{29}$, $\mathcal A_{32}$}

In this paragraph a Hida setting and the notation of Definition~\ref{i1:def-slots}
and of the paragraph before Theorem~\ref{i1:thm-stdbasis} are fixed, $k=m+1$, and
$\varphi:W(\mathbb F_{p^{n'}})\to W$ is the isomorphism of Lemma~\ref{i1:lem-A28}.
For a letter $s<n'$ of type vector or covector let $\theta_s:K_f[x_0..x_m]\to
K_f[X_{a,c}]$ be the algebra map with $\theta_s(x_c)=X_{c,0}$ (vector) or
$\theta_s(x_c)=\mathrm{adj}(X)_{m,c}$ (covector), $X=(X_{a,c})_{a,c\le m}$; put
$A_s(g)=g$ (vector), $(g^{-1})^T$ (covector); for a family $N=(N_s)$ of
matrices let $\mathrm{Lin}_N$ be the substitution $X_{s,b}\mapsto\sum_{b'}(N_s)_{bb'}X_{s,b'}$
on polynomials in variables $X_{s,b}$ ($s<n'$, $b\le m$), and let $\Theta_0$ send
$X_{s,c}$ to $\theta_s(x_c)$ placed in block $s$. For degrees $d$, $\mathrm{Pol}_d$
is as in Definition~\ref{i1:def-twist}. The weight $\lambda$ is \emph{flat} if
$\lambda_s=d_se_0$ for vector $s$ and $\lambda_s=-d_se_m$ for covector $s$. For
$g\in GL_k(W(\mathbb F_{p^{n'}}))$ put $G_s=(\varphi g)^{(s)}$ and
$D_w=\mathrm{diag}(w^{-1}[c<r]+[c\ge r])$; $F$ lies in the \emph{$i$-th lifted
$w$-filtration for $(r,g)$} if all coefficients of $\mathrm{ls}(F)$ lie in $w^iO$,
where $\mathrm{ls}$ sends $X_{s,c}$ to $\sum_b\delta'_b(G_s^{-1})_{bc}X_{s,b}$
(vector; $\delta'_b=w$ for $b\ge r$, $1$ otherwise) and to
$\sum_b\delta_b(G_s)_{cb}X_{s,b}$ (covector; $\delta_b=w$ for $b<r$, $1$
otherwise).

\begin{lemma}[Clause $\mathcal A_{29}$: flat tops]\label{i1:lem-A29}
Let $1\le r\le m$, $\lambda$ flat for degrees $d$, and $h\in GL_k(\Q_q)$. Put
$h_s=h^{(s)}$, $\kappa:K\to K_f$ the coefficient map of the setting, and
$\Theta(F)=\Theta_0(\mathrm{Lin}_{(A_s(h_s^{-1}))_s}\kappa F)$. Then $\Theta$ is an
$O$-isomorphism from $\mathrm{Pol}_d$ onto $V$; $F$ has coefficients in $O$ iff
$\Theta F\in\mathrm{HL}(\mathrm{colspan}_Wh)$; and for every $g$ as above and every
edge $e$ with edge basis $h((\varphi g)^{-1})^T$ (i.e.\ $\mathrm{src}\,e$ and
$\mathrm{tgt}\,e$ are spanned by its columns and by those of its product with
$D_p$), $F$ lies in the $i$-th lifted $w$-filtration for $(r,g)$ iff $\Theta F\in
w^i\mathrm{Re}(e)$.
\end{lemma}

\begin{proof}
Put $e_s=\lambda_{s,\rm last}$, $\nu_s=\lambda_s-e_s$: for a vector $e_s=0$,
$\nu_s=d_se_0$; for a covector $e_s=-d_s$, $\nu_s=d_s(e_0+\dots+e_{m-1})$. Four
facts.

(F1) $\nabla_O(\nu_s)=\theta_s(\{\text{integral forms of degree }d_s\})$, and
likewise over $K_f$. The $\theta_s(x_c)$ are signed initial minors (of size $1$,
resp.\ $m$); initial minors of size $z$ lie in $\nabla(e_0+\dots+e_{z-1})$ and
dual Weyl modules multiply, which gives $\subseteq$. Conversely, in the
bideterminant basis of Theorem~\ref{i1:thm-stdbasis} of shape $\nu_s$, all
minors have size $\le1$ (vector, as $\nu_{s,1}=0$), resp.\ exactly $m$ or $0$
(covector, as $\nu_{s,m}=0$ and $\nu_{s,m-1}=\nu_{s,0}=d_s$), so each basis element
is $\theta_s$ of an integral monomial.

(F2) $\det(g)^{e_s}g\cdot\theta_s(Q)=\theta_s(A_s(g)\cdot Q)$ for $Q$ of degree
$d_s$, where $A_s(g)\cdot Q$ is the linear substitution: for a vector this is
the row action on $X_{c,0}$; for a covector $g\cdot\mathrm{adj}(X)=\mathrm{adj}(X)\mathrm{adj}(g)$ and
$\mathrm{adj}(g)^T=\det(g)(g^{-1})^T$, whose factor $\det(g)^{d_s}$ is cancelled by
$e_s=-d_s$.

(F3) If $M_s=\{\theta_s(N_s\cdot P):P\text{ integral of degree }d_s\}$, then the
$O$-span of the products of elements of the $M_s$ is
$\{\Theta_0(\mathrm{Lin}_NF):F\in\mathrm{Pol}_d\text{ integral}\}$ (split monomials into
block parts).

(F4) $\Theta_0$ is injective on $\mathrm{Pol}_d$: the algebra map $\psi_s$ with
$\psi_{\rm v}(X_{a,t})=[t=0]x_a$ and $\psi_{\rm c}(X)=P$,
$P_{at}=[a=t]x_m-[a=m,t\ne m]x_t$ (lower triangular, so
$\mathrm{adj}(P)_{m,t}=x_m^{m-1}x_t$) gives $\psi_s\theta_s(x_c)=q_sx_c$ with
$q_s\in\{1,x_m^{m-1}\}$, so $\Psi_0\Theta_0F=\prod_s(q_s^{[s]})^{d_s}F$ on
$\mathrm{Pol}_d$, and the factor is a nonzero polynomial.

By (F1), (F2) and the definition of $\nabla_\lambda(g)$, $x\in\nabla_{\lambda_s}(g)$
iff $x=\theta_s(A_s(g^{-1})\cdot P)$ with $P$ integral of degree $d_s$; with
(F3), $\bigotimes_s\nabla_{\lambda_s}(g_s)=\{\Theta_0(\mathrm{Lin}_{(A_s(g_s^{-1}))}F):F\text{ integral}\}$,
and $V=\Theta_0(\mathrm{Pol}_d\otimes K_f)$. So $\Theta$ is onto $V$, and injective
by (F4). More generally, by (F4),
$\Theta_0(\mathrm{Lin}_{(A_s(h_s^{-1}))}F)\in c\bigotimes_s\nabla_{\lambda_s}(M_s)$ iff
$\mathrm{Lin}_{(A_s(h_s^{-1})A_s(M_s))}F\in c\cdot\{\text{integral}\}$. With $M=h$,
$c=1$ and Lemma~\ref{i1:lem-transport}(2),(4)
($H_0(\lambda_s,\mathrm{ext}_s(\mathrm{colspan}_Wh))=\nabla_{\lambda_s}(h_s)$) this is
the integrality statement. For the filtration, with
$M_s=h_s(G_s^{-1})^TD_w$ one checks $\kappa\,\mathrm{ls}(F)=\mathrm{Lin}_BF$ with
$B_s=w^{[s\text{ vector}]}A_s(h_s^{-1})A_s(M_s)$ ($M_s$ is invertible), and the
scalar $w$ on the vector blocks multiplies $\mathrm{Lin}$ on $\mathrm{Pol}_d$ by
$w^{A}$, $A=\sum_{s\text{ vector}}d_s$. So the filtration condition reads
$\Theta F\in w^{i-A}\bigotimes_s\nabla_{\lambda_s}(M_s)$. By
Lemma~\ref{i1:lem-edgebasis} (whose computation uses only the two column-span
equalities of an edge basis) $\mathrm{Le}_s(e)=\mathrm{colspan}_O(M_s)$, so
$\bigotimes_s\nabla_{\lambda_s}(M_s)=\mathrm{HLe}(e)$; and for flat $\lambda$ and
$1\le r\le m$, $A_r=\sum_{s\text{ vector}}d_s=A$. Since
$\mathrm{Re}(e)=w^{-A_r}\mathrm{HLe}(e)$, the condition is $\Theta F\in
w^i\mathrm{Re}(e)$.
\end{proof}

\begin{proposition}[The forms on dual Weyl modules]\label{i1:prop-Fm}
Let $p$ be a prime, $O$ a local domain with fraction field $K$ of
characteristic $0$ in which $p$ is not a unit, $n\le k$ and $\nu\in\N^n$
antitone. There is a $K$-bilinear form $B$ on $K[X_{a,t}:a<k,\ t<n]$ such
that: (a) $B(g\cdot P,Q)=B(P,g^T\cdot Q)$ for the row action and all $g\in
M_k(K)$; (b) if $Q\in\nabla_K(\nu)$ and $B(P,Q)=0$ for all $P\in\nabla_K(\nu)$
then $Q=0$; (c) if $Q\in\nabla_K(\nu)$ and $B(P,Q)\in O$ for all
$P\in\nabla_O(\nu)$ then $Q\in\nabla_O(\nu)$; (d) for every $l$ with ($\nu_t\le
l$ for all $t$) or ($n=k$ and $\nu_s\le\nu_t+l$ for all $s,t$),
$B(P,p^{N_0(k)\delta(p,k,l)}Q)\in O$ for $P,Q\in\nabla_O(\nu)$; (e) if $\nu=0$,
$B(P,Q)\in O$ for $P,Q\in\nabla_O(\nu)$.
\end{proposition}

\noindent Here $K[X_{a,t}]$ carries the row action of
Definition~\ref{i1:def-slots}(2), and $\nabla_K(\nu)$, $\nabla_O(\nu)$ are the
dual Weyl modules on $k\times n$ matrices (defined at the beginning of
\S\ref{i1:ssec-pf-Fm}; for $n=k$ and a domain they are those of
Definition~\ref{i1:def-slots}), with $\nabla_O(\nu)$ identified with its image
in $K[X]$. The form is a multiple of the Fock pairing; the proof, in
\S\ref{i1:ssec-pf-Fm}, uses the standard basis
(Theorem~\ref{i1:thm-stdbasis}), the representation theory of
$\mathfrak{gl}_k$ of \S\ref{i1:ssec-pf-adm} and an integral
extremal-projector argument.

\begin{lemma}[Clause $\mathcal A_{32}$: test side]\label{i1:lem-A32}
For every $r$: the dual $e^*=(\mathrm{tgt}(e)^*,\mathrm{src}(e)^*)$ of an edge is an
edge and $e^{**}=e$, where $L^*=\{x:\sum_ix_iy_i\in W\ \forall y\in L\}$; and there
is a $K_f$-isomorphism $\Psi:V^*\to V$ with, for every lattice $L$ and edge $e$,
$\Psi(\mathrm{HL}(L)^\vee)\subseteq\mathrm{HL}(L^*)\subseteq p^{-\kappa_1}\Psi(\mathrm{HL}(L)^\vee)$
and $\Psi(\mathrm{Re}'(e))\subseteq\mathrm{Re}(e^*)\subseteq p^{-\kappa_1}\Psi(\mathrm{Re}'(e))$,
$\kappa_1=\sum_{s:\ell(\lambda_s)>0}N_0(k)\delta(p,k,\ell(\lambda_s))$.
\end{lemma}

\begin{proof}
\emph{Duals.} $\mathrm{colspan}_W(g)^*=\mathrm{colspan}_W(g^{-T})$ ($g^{-T}=(g^{-1})^T$),
so $L^{**}=L$. For an edge with basis $g$ (Lemma~\ref{i1:lem-edgebasis}) and
$h=g^{-T}D_p^{-1}$, $\mathrm{tgt}(e)^*=\mathrm{colspan}(h)$ and
$\mathrm{src}(e)^*=\mathrm{colspan}(hD_p)$; and every such pair is an edge
($p\,\mathrm{colspan}(hD_p)\subseteq\mathrm{colspan}(h)$ with quotient $(W/pW)^r$).

\emph{Step 1.} $W$ has characteristic $0$ ($p$ is irreducible), $W\to O$ is
injective and local ($O$ is finite over $W$), so $p$ is not a unit in $O$.
Proposition~\ref{i1:prop-Fm} with $n=k$, $\nu=\nu_s$ (antitone, entries
$\le\ell(\lambda_s)$) gives a form $B_s$ with $B_s(P,p^{N_s}Q)\in O$ on
$\nabla_O(\nu_s)$, $N_s=0$ if $\ell(\lambda_s)=0$ (part (e)) and
$N_s=N_0(k)\delta(p,k,\ell(\lambda_s))$ otherwise (part (d)); $\sum_sN_s=\kappa_1$.

\emph{Step 2.} Fix $s$, $V_s=\nabla_{K_f}(\nu_s)$, $e=\lambda_{s,\rm last}$, and
$\beta_s(Q)=B_s(\cdot,Q)$; by (b) and dimensions $\beta_s:V_s\to V_s^*$ is an
isomorphism. By (a), $B_s(\det(M)^eM\cdot P,\det(M^{-T})^eM^{-T}\cdot Q)=B_s(P,Q)$.
Hence (up) if $\beta_s(Q)(P)\in O$ for all $P\in\nabla_{\lambda_s}(M)$, then
$Q\in\nabla_{\lambda_s}(M^{-T})$ (apply (c) to $\det(M^{-T})^eM^{-T}\cdot Q$, testing
against $P_0=\det(M)^eM\cdot P$); and (low) if $Q\in\nabla_{\lambda_s}(M^{-T})$ and
$P\in\nabla_{\lambda_s}(M)$ then $\beta_s(p^{N_s}Q)(P)\in O$ (Step 1).

\emph{Step 3.} Fix bases $b_s$ of the $V_s$, let $x$ be their products and $x'$
the products of the bases $b_s^\beta$ with $\beta_s(b^\beta_{s,j})=$ the $j$-th
coordinate functional; put $\Psi(\phi)=\sum_a\phi(x_a)x'_a$. Coordinates of pure
tensors are products, and a direct comparison shows
$\Psi(\mathrm{coord}^z_j)=z'_j$ for every choice of bases $Z_s$ (with products $z$,
$z'$). For invertible $M_s$, choose $Z_s$ with $O$-span $\nabla_{\lambda_s}(M_s)$
(Lemma~\ref{i1:lem-rowwt}(2)); then $H(M)=\bigotimes_s\nabla_{\lambda_s}(M_s)$ is
the span of the $z_i$, its dual is the span of the $\mathrm{coord}^z_j$, and
$\Psi$ maps it onto the span of the $z'_j$. By (up) each $z'_j\in H(M^{-T})$, and
by (low) $p^{\sum N_s}H(M^{-T})\subseteq\mathrm{span}_O(z')$.

\emph{Conclusion.} For $L=\mathrm{colspan}_W(g)$ take $M_s=g^{(s)}$; then
$\mathrm{HL}(L^*)=H(M^{-T})$ (twisting commutes with $g\mapsto g^{-T}$), which gives
the inclusions for $\mathrm{HL}$. For an edge with basis $g$ take
$M_s=g^{(s)}D_w$ (Lemma~\ref{i1:lem-edgebasis}); $e^*$ has edge basis
$g^{-T}D_p^{-1}$, so $\mathrm{Le}_s(e^*)$ is spanned by the columns of
$(g^{(s)})^{-T}D_p^{-1}D_w=(g^{(s)})^{-T}D_w^{-1}=M_s^{-T}$ (as $p=w^2$), i.e.\
$\mathrm{HLe}(e^*)=H(M^{-T})$; and the factor
$w^{-A_r}$ in $\mathrm{Re}$, $\mathrm{Re}'$ commutes with $\Psi$.
\end{proof}

\subsubsection*{The plane count: clause $\mathcal A_{27}$}

\begin{definition}[The plane count data]\label{i1:def-A27}
Notation of Definition~\ref{i1:def-seg}; an unramified quadratic Hida setting
(Definition~\ref{i1:def-hida}) with ring $O$, $w^2=p$, residue field $k_O$;
$\mathrm{pos}(x)=x\bmod n'$; $\mathrm{dist}(j,x)=\max(0,\mathrm{fr}(j,r-1)-x)$;
$\mathrm{Late}_j$ the late letters of $j$, $\mathrm{LL}_j$ the L-letters; for a
letter $x$ of $j$, $\ell_j(x)$ the increasing list of even letters of $j$ below
$x$, of length $L$, $\mathrm{old}_j(x)=\ell_j(x)[L-r]$ and
$\mathrm{row}_j(x,i)=\ell_j(x)[L-r+i]$ (truncations as in the formal
development: $L-r$ is cut at $0$, an index $\ge L$ gives $0$). Let
$\mathcal R_O=O[x_{t,c}:t\in\Z/n',\ c<k]$.
\begin{enumerate}
\item \emph{Data.} $X=(F,\deg,L_1,N_1,N_s,E_0,m^*,E',a,Y,Q,C_{La},C_k)$, and
 $\mathrm{Hyp}(X)$: $p\ge13$, $r\ge2$, $n'$ even; $F$ valid at $(L_1,N_1)$;
 $a\le\deg t\le2a$ ($t\in\mathrm{TT}$); every segment has at most $N_s$
 letters; $N_s\le N_1$, $N_s\le2|\mathrm{TT}|$, $N_s+r\le K$;
 $p^{\mathrm{dist}(j,x)}\le Q$ for all letters; $Y\ge1$; $C_kN_s^2\log a\le k$;
 $128k^2N_sY^2\le p^{L_1-1}$; $E'_t\le E_0$; $32pYk^2N_sQ^2(E_0+m^*+1)\le a$.
\item \emph{Generators.} $\mathrm{La}=C_{La}N_s^2\log a$;
 $\mathrm{Bo}(j,x)=\deg(\mathrm{pos}\,\mathrm{old}_j(x))/p^{\mathrm{dist}(j,\mathrm{old}_j(x))}$,
 $\mathrm{sg}=\lfloor\mathrm{Bo}/(k^2Y)\rfloor$, $\mathrm{mg}=m^*+\mathrm{sg}$,
 $\mathrm{th}=\mathrm{Bo}/2$, $o(j,x)=\mathrm{La}\,a/(pk\,p^{\mathrm{dist}(j,x)})+c_xap^{-(L_1+\mathrm{dist}(j,x))}$
 ($c_x=10$ for $x$ even, $1$ for $x$ odd), $d_t=(\deg t-E'_t)^+$ on
 $\mathrm{TT}$ and $0$ elsewhere. $\varphi_j:\mathcal S_O\to\mathcal R_O$ sends
 $X_{(x,c)}\mapsto x_{\mathrm{pos}(x),c}^{p^{\mathrm{dist}(j,x)}}$; a \emph{label} is an
 exponent $l$ with $l(\mathrm{pos}(x),c)<p^{\mathrm{dist}(j,x)}$ for all $j$, late
 $x$, $c$. $\mathrm{Gen}$ is the set of
 $\bigl(\prod_{j<J}\varphi_j(y_j)\bigr)x^l$, $y_j\in\mathcal X'_j(\mathrm{mg}(j,\cdot),\mathrm{th}(j,\cdot),o(j,\cdot))$,
 $l$ a label, supported on exponents with row sums $d_t$; $\mathcal
 N=\mathrm{span}_O\mathrm{Gen}$.
\end{enumerate}
The \emph{segment-degree objects} are the following.
\begin{enumerate}\setcounter{enumi}{2}
\item \emph{Discriminant.} $\Delta_l=\prod_{j<J}\varphi_j\iota_{pl}\prod_{x\in\mathrm{Late}_j}\det(X_{(\mathrm{row}_j(x,a),b)})_{a,b<r}^{\mathrm{sg}(j,x)}$.
\item \emph{Coarse data.} $E_t=\sum_{j<J}\sum_{x\in\mathrm{Late}_j}\mathrm{sg}(j,x)\sum_{i<r}[\mathrm{pos}\,\mathrm{row}_j(x,i)=t]p^{\mathrm{dist}(j,\mathrm{row}_j(x,i))}$;
 $a'_t=(d_t-E_t)^+$; $R_j=\lfloor\log_p\lfloor a/(8Yk^2N_sp^{\mathrm{dist}(j,\mathrm{fr}(j-1,0))})\rfloor\rfloor$
 ($\lfloor\log_p0\rfloor=0$); $P_j(x)=p^{\mathrm{dist}(j,x)+R_j}$;
 $Q_j=p^{\max(0,\mathrm{fr}(j+1,r-1)-\mathrm{fr}(j,r-1)+R_{j+1}-R_j)}$; $m_0=kN_s+1$;
 $\varphi^c_j:X_{(x,c)}\mapsto x_{\mathrm{pos}(x),c}^{P_j(x)}$; a \emph{coarse label}
 is an $l$ with $l(\mathrm{pos}(x),c)<P_j(x)$ for all $j$, late $x$, $c$;
 $\mathrm{Res}_j$ is the set of exponents $m$ with $m(\mathrm{fr}(j,i),a)<Q_j$
 ($i,a<r$); $S_b$ the exponents with $\sum_cm(x,c)=b_x$ for all $x$;
 $\mathrm{Bw}_j(y)=\lfloor a'_{\mathrm{pos}(y)}/P_j(y)\rfloor-(2r+K+1)$;
 $\theta^c_j(b)(x)=\frac23(\mathrm{Bw}_j(\mathrm{old}_j(x))-m_0)$ if
 $\mathrm{old}_j(x)$ is a $W$-letter of $j$, $\frac23b(\mathrm{old}_j(x))$ otherwise;
 $o^c_j(x)=\mathrm{La}\,a/(pkP_j(x))+[x\text{ even}]4ap^{-(\mathrm{dist}(j,\mathrm{old}_j(x))+R_j)}$.
\item \emph{Modules.} $M_j(b)=\mathcal X'_j(m_0,\theta^c_j(b),o^c_j)\cap\mathrm{Pol}_{S_b}\cap\mathrm{Pol}_{\mathrm{Res}_j}$;
 $\Lambda\mathrm{Gen}$ is the set of $\bigl(\prod_{j<J}\varphi^c_j(y_j)\bigr)x^l$ with
 $y_j\in M_j(b_j)$ for some $b$, $l$ a coarse label, supported on exponents
 with row sums $a'$.
\item \emph{Index sets.} $\mathrm{LA}=\prod_t\mathrm{la}(t)$, $\mathrm{la}(t)$ the
 $v\in\N^k$ with $\sum v_c\le a'_t$ and, for every $j$ and late $x$ with
 $\mathrm{pos}(x)=t$, $v_c<P_j(x)$ and $\sum_cv_c\equiv a'_t\bmod P_j(x)$. For
 $\lambda\in\mathrm{LA}$, $\mathrm{nz}_\lambda(j,x)=\lfloor(a'_{\mathrm{pos}(x)}-\sum_c\lambda(\mathrm{pos}(x),c))/P_j(x)\rfloor$;
 $\mathrm{EE}(\lambda)$ is the set of $(e_{j,i})_{j<J,i<r}$ with $0\le e_{j,i}<rQ_j$,
 $e_{j,i}\equiv\mathrm{nz}_\lambda(j,\mathrm{fr}(j,i))\bmod Q_j$ (indices $j$ mod $J$).
 The padded vector $b^{\lambda,e}_j$ vanishes off the letters of $j$, is
 $\lfloor(\mathrm{nz}_\lambda(j-1,x)-e_{j-1,i})^+/Q_{j-1}\rfloor$ at $x=\mathrm{fr}(j-1,i)$,
 $e_{j,i}$ at $x=\mathrm{fr}(j,i)$ ($0$ at other late $x\ge\mathrm{fr}(j,0)$), and
 $\mathrm{nz}_\lambda(j,x)$ at late $x<\mathrm{fr}(j,0)$. $\tau_j(b,e)$ replaces the
 values at $\mathrm{fr}(j,i)$ by $e_i$ (and is $0$ at other letters in
 $[\mathrm{fr}(j,0),\mathrm{fr}(j,r-1)]$). $\mathrm{Adm}_j(b)$:
 $a'_{\mathrm{pos}(x)}/P_j(x)-k<b_x\le a'_{\mathrm{pos}(x)}/P_j(x)$ for $x\in\mathrm{LL}_j$.
 $b|_j$ is $b$ on $\mathrm{Late}_j$ and $0$ elsewhere.
\item \emph{Paths.} Let $e_j(0)<\dots<e_j(N_j-1)$ be the letters of $j$;
 $\mathrm{Bd}_j(b)(x)=(\mathrm{Bw}_j(x)-m_0)^+$ at $W$-letters and $b_x$ otherwise;
 $\beta_j(b)=\sum\mathrm{Bd}_j(b)(x)$ over letters $x<\mathrm{fr}(j,0)$ of $j$.
 $\mu$ ($k$ coordinates) is a \emph{path for $b$ with degrees $dd$} if it is a
 Pieri path of degrees $dd$ with $dd_i=\mathrm{Bd}_j(b)(e_j(i))$ when
 $e_j(i)<\mathrm{fr}(j,0)$ and $\mu_{i+1,m}=\beta_j(b)$ for $m\le i'$ when
 $e_j(i)=\mathrm{fr}(j,i')$; it is \emph{good} if for late $e_j(i)$:
 $\theta^c_j(b)(e_j(i))\le\mu^{pl}_{i,r-1}$, and the tangent (even letter),
 resp.\ plane (odd letter), part of the step $i$ has size $\le o^c_j(e_j(i))$.
 $T_j(b)$ is the set of $(\mu,I)$, $\mu$ good, $I$ a standard filling of
 $\mu^{tan}$; $\mathrm{Idx}_j(b)$ the set of $(\mu,I,h)$, $h:\{0..r-1\}^2\to\N$,
 with $\mu$ good of degrees $dd$, $I$ standard for $\mu^{tan}$,
 $h_v+2\beta_j(b)<Q_j$, and $dd_i+\sum_ch(i',c)=b(\mathrm{fr}(j,i'))$ when
 $e_j(i)=\mathrm{fr}(j,i')$. With $D_{w,j}=\det(X_{(\mathrm{fr}(j-1,a),c)})_{a,c<r}$,
 $\mathrm{fam}_j(b,(\mu,I,h))=\iota_{pl}\bigl(D_{w,j}^{m_0}\prod_{i<r}X_{(\mathrm{fr}(j-1,i),0)}^{(b(\mathrm{fr}(j-1,i))-\mathrm{Bw}_j(\mathrm{fr}(j-1,i)))^+}\prod_vX_{(\mathrm{fr}(j,v_1),v_2)}^{h_v}\bigr)\cdot\mathrm{ren}_j(\Psi^{pl}_{\mu^{pl},I^{\rm can}}\Psi^{tan}_{\mu^{tan},I})$,
 $\mathrm{ren}_j$ renaming letter $i$ to $e_j(i)$.
\item \emph{Segment bound.} $\mathrm{sb}'_j(b)=(1-a^{-N_s^2})\prod_{x\in\mathrm{LL}_j}\binom{b_x+k-1}{k-1}\prod_{i<r}\binom{(\mathrm{Bw}_j(\mathrm{fr}(j-1,i))-m_0)^++k-1}{k-1}$,
 $c_j=(Q_j^{r-1}\max(0,1-\frac r{2Y}))^r$, $\mathrm{sb}_j=\mathrm{sb}'_jc_j$.
\end{enumerate}
\end{definition}

\begin{proposition}[The segment-degree part of $\mathcal A_{27}$]\label{i1:prop-segdeg}
Let $C_{La},C_k\ge1088$ and $X$ satisfy $\mathrm{Hyp}(X)$; put
$\mathrm{bound}=(1-|\mathrm{TT}|r^2k/Y)\prod_{t\in\mathrm{TT}}\binom{\deg t+k-1}{k-1}$,
and write $\overline G$ for the $k_O$-span of the reductions of a set $G$. Then:
(I1) $\bar\Delta_l\ne0$; (I2) for every $u\in\Lambda\mathrm{Gen}$ there is
$g\in\mathrm{Gen}$ with $\overline{\Delta_lu}=\bar g$; (I3)
$\sum_{\lambda\in\mathrm{LA}}\sum_{e\in\mathrm{EE}(\lambda)}\prod_{j<J}\mathrm{rk}_OM_j(b^{\lambda,e}_j)\le\dim\overline{\Lambda\mathrm{Gen}}$;
(I4) each $b^{\lambda,e}_j$ vanishes off the letters of $j$, satisfies
$\mathrm{Bw}_j(x)\le b^{\lambda,e}_{j,x}$ at the $W$-letters and
$\mathrm{Adm}_j(b^{\lambda,e}_j)$; (I5) for every family $G_j(b)\ge0$ depending only
on $b|_j$ with $\mathrm{sb}_j(b)\le\sum_eG_j(\tau_j(b,e))$ whenever
$\mathrm{Adm}_j(b)$, for every $\rho\in\N^r$ (the sum over $e$ with $0\le e_i<rQ_j$,
$e_i\equiv\rho_i\bmod Q_j$; the same range in (I9)),
$p^{-n'rK}\mathrm{bound}\le\sum_\lambda\sum_e\prod_jG_j(b^{\lambda,e}_j)$; (I6),
(I7) under the hypotheses of (I4) on $b$, the $\mathrm{fam}_j(b,q)$,
$q\in\mathrm{Idx}_j(b|_j)$, lie in $M_j(b)$ and are $O$-linearly independent;
(I8) if $\mathrm{Adm}_j(b)$, $T_j(b|_j)$ is finite and $\mathrm{sb}'_j(b)\le\#T_j(b|_j)$;
(I9) if moreover $T_j(b|_j)$ is finite, $\#T_j(b|_j)\cdot c_j\le\sum_e\#\mathrm{Idx}_j(\tau_j(b,e)|_j)$
for every $\rho$;
(Deg) every exponent $m$ of every element of $\mathcal N$ has
$\mathrm{vs}(m)\le10|\mathrm{TT}|(\mathrm{La}\,a/(pk)+ap^{-L_1}+kQ)$, where
$\mathrm{vs}(m)=\sum_{t<n'}\bigl([t\text{ even}]\sum_{c\ge r}m(t,c)+[t\text{ odd}]\sum_{c<r}m(t,c)\bigr)$
(the degree in the tangent variables at even $t$ and the plane variables at
odd $t$).
\end{proposition}

\noindent (I1), (I4), (I9) and (Deg) are proved in \S\ref{i1:ssec-pf-segdeg}; the
other items are the following proposition.

\begin{proposition}[the counting part of $\mathcal A_{27}$; formally verified, \cite{CharacterFamily}]\label{i1:prop-segdeg-rest}
Under the hypotheses of Proposition~\ref{i1:prop-segdeg}, the items (I2),
(I3), (I5), (I6), (I7), (I8) of that proposition hold. The formal statements
of (I2), (I3), (I6), (I7), (I8) carry as additional hypotheses the clauses
$\mathcal A_{12}$, $\mathcal A_{13}$, $\mathcal A_{16}$, $\mathcal A_{17}$,
$\mathcal A_{19}$ (including Proposition~\ref{i1:prop-A19c}),
$\mathcal A_{26}$, $\mathcal A_{33}$, $\mathcal A_{34}$, the conclusion of
Lemma~\ref{i1:lem-planes} and the three counts (Theorem~\ref{i1:thm-weyl},
Lemma~\ref{i1:lem-counts}); those of (I2), (I3), (I6), (I7) assume in
addition that $O$ is complete with $p$ in its maximal ideal.
\end{proposition}

\noindent All these additional hypotheses are established or cited in this
section, and the completeness of $O$ holds in the Hida setting; so the
proposition applies wherever it is used. (I5) is proved formally without
them. (I9) is proved in \S\ref{i1:ssec-pf-segdeg}.

\noindent These are the segment computations: (I2) reduces the coarse
generators through Frobenius identities on column words; (I3), (I6), (I7) use
Lemmas~\ref{i1:lem-A15I}, \ref{i1:lem-A16}, \ref{i1:lem-A17}, \ref{i1:lem-A33};
(I5), (I8) are counts of Gelfand--Tsetlin patterns
(Lemma~\ref{i1:lem-counts}) with logarithmic estimates of Weyl dimensions.  Its formally verified steps are outlined in \S\ref{i1:str-segdeg}.

\begin{lemma}[Rational planes]\label{i1:lem-planes}
For a field $\mathbb F$ with $q$ elements and $r\le k$, there are at least
$q^{r(k-r)}$ subspaces of dimension $r$ in $\mathbb F^k$.
\end{lemma}

\begin{proof}
Let $P(n,r)=\prod_{i<r}(q^n-q^i)$, the number of independent $r$-tuples in
$\mathbb F^n$. Each $r$-subspace has $P(r,r)$ bases, so their number $N_r$
satisfies $N_rP(r,r)=P(n,r)$. The Gaussian binomials $\binom nr_q$
($\binom n0_q=1$, $\binom0{r+1}_q=0$, $\binom{n+1}{r+1}_q=\binom nr_q+q^{r+1}\binom n{r+1}_q$)
satisfy $\binom nr_qP(r,r)=P(n,r)$ by induction, so $N_r=\binom nr_q$; and
$q^{r(n-r)}\le\binom nr_q$ by induction, from $q^{(r+1)(d+1)}=q^{r+1}q^{(r+1)d}\le
q^{r+1}\binom n{r+1}_q$.
\end{proof}

\begin{lemma}[Clause $\mathcal A_{27}$, modulo Proposition~\ref{i1:prop-segdeg}]\label{i1:lem-A27}
With $C_0=1088$, $C_{rk}=1$, $C_{deg}=10$: for every $X$ with $\mathrm{Hyp}(X)$ and
$C_{La}=C_k=C_0$, put
$\mathrm{bound}=(1-C_{rk}|\mathrm{TT}|r^2k/Y)\prod_{t\in\mathrm{TT}}\binom{\deg t+k-1}{k-1}$.
Then $\mathcal N$ contains an $O$-direct summand of $\mathcal R_O$ with a
basis of elements of $\mathrm{Gen}$, of rank at least
$p^{-n'rK}\,\mathrm{bound}$ and at least $\mathrm{bound}$ divided by the
number of $r$-planes in $\mathbb F_{p^{n'}}^k$; and every exponent $m$ of an
element of $\mathcal N$ has
$\mathrm{vs}(m)\le C_{deg}|\mathrm{TT}|(\mathrm{La}\,a/(pk)+ap^{-L_1}+kQ)$
($\mathrm{vs}$ as in (Deg)).
\end{lemma}

\begin{proof}
(R1) Put $G_j(b)=\#\mathrm{Idx}_j(b|_j)$. Under the hypotheses of (I4),
$G_j(b)\le\mathrm{rk}_OM_j(b)$: $M_j(b)$ is finitely generated (it lies in the
span of finitely many monomials), and by (I6), (I7) it contains an independent
family indexed by $\mathrm{Idx}_j(b|_j)$. If $\mathrm{Adm}_j(b)$, then
$\mathrm{sb}_j(b)\le\sum_eG_j(\tau_j(b,e))$ by (I8) times $c_j$ and (I9). (R2) By
(I5), (R1), (I4) and (I3), $p^{-n'rK}\mathrm{bound}\le\dim\overline{\Lambda\mathrm{Gen}}$.
(R3) Multiplication by $\bar\Delta_l$ is injective on $k_O[x]$ (a domain; (I1))
and maps $\overline{\Lambda\mathrm{Gen}}$ into $\overline{\mathrm{Gen}}$ (I2), which is
finite-dimensional (supports lie in a finite set $S_d$); so
$p^{-n'rK}\mathrm{bound}\le\dim\overline{\mathrm{Gen}}$. (R4) \emph{Selection:} for a
set $G$ of polynomials supported in a finite set $s$, there is a direct summand
$S\subseteq\mathrm{span}_OG$ of $O[x]$ with a basis of elements of $G$ and
$\mathrm{rk}\,S=\dim\overline G$. Choose $v_\beta\in G$ whose reductions form a basis
of $\overline G$ and let $S$ be their span. If $\sum c_\beta v_\beta=wy$, reducing
gives $w\mid c_\beta$ and $\sum(c_\beta/w)v_\beta=y$; hence the $v_\beta$ are
independent (a nonzero coefficient would be divisible by every power of $w$)
and $S$ is saturated. $\mathrm{Pol}_s/S$ is finitely generated and torsion free,
hence free over the principal ideal domain $O$, so $S$ is a direct summand of
$\mathrm{Pol}_s$, which is a direct summand of $O[x]$. (R5) Apply (R4) to
$G=\mathrm{Gen}$, $s=S_d$: $\mathrm{rk}\,S=\dim\overline{\mathrm{Gen}}\ge p^{-n'rK}\mathrm{bound}$.
By Lemma~\ref{i1:lem-planes} the number of $r$-planes is at least
$p^{n'rK}>0$, which gives the second bound (trivially if $\mathrm{bound}\le0$).
The degree bound is (Deg).
\end{proof}

\noindent With Lemmas~\ref{i1:lem-A15I}--\ref{i1:lem-A27}, all the admission
clauses are proved in this paper except one: the counting part of
$\mathcal A_{27}$, Proposition~\ref{i1:prop-segdeg-rest}.
(Propositions~\ref{i1:prop-A19c}, \ref{i1:prop-A21}, \ref{i1:prop-Fm} and the
items (I1), (I4), (Deg) of Proposition~\ref{i1:prop-segdeg} are proved in
\S\S\ref{i1:ssec-pf-A19c}, \ref{i1:ssec-pf-A21}, \ref{i1:ssec-pf-Fm},
\ref{i1:ssec-pf-segdeg}.) In addition, the formal proof of $\mathcal A_{27}$ takes as inputs
Weyl's dimension formula for standard fillings, the branching count and the
Pieri count; here these are Theorem~\ref{i1:thm-weyl} and
Lemma~\ref{i1:lem-counts}.

\begin{proof}[Proof of Proposition~\ref{i1:prop-adm}]
Clauses $\mathcal A_1$, $\mathcal A_2$, $\mathcal A_{34}$ are
Lemmas~\ref{i1:lem-A1}, \ref{i1:lem-A2}, \ref{i1:lem-A34}; $\mathcal A_{23}$
and the first two parts of $\mathcal A_{24}$ are Lemmas~\ref{i1:lem-A23},
\ref{i1:lem-A24}; $\mathcal A_4$ and $\mathcal A_{37}$ are
Lemmas~\ref{i1:lem-A4}, \ref{i1:lem-A37}; $\mathcal A_3$, $\mathcal A_{31}$ are
Lemma~\ref{i1:lem-LF}; $\mathcal A_5$, $\mathcal A_6$, $\mathcal A_7$, $\mathcal A_{12}$,
$\mathcal A_{13}$, $\mathcal A_{22}$, the third part of $\mathcal A_{24}$,
$\mathcal A_{25}$, $\mathcal A_{26}$, $\mathcal A_{28}$, $\mathcal A_{30}$,
$\mathcal A_{35}$, $\mathcal A_{36}$ are Lemmas~\ref{i1:lem-A5}, \ref{i1:lem-A6},
\ref{i1:lem-A7}, \ref{i1:lem-A12}, \ref{i1:lem-A13}, \ref{i1:lem-A22},
\ref{i1:lem-A24c}, \ref{i1:lem-A25}, \ref{i1:lem-A26}, \ref{i1:lem-A28},
\ref{i1:lem-A30}, \ref{i1:lem-A35}, \ref{i1:lem-A36}; $\mathcal A_{15}$ is
Lemma~\ref{i1:lem-A15} with its inputs (C) (Lemma~\ref{i1:lem-counts}) and (I)
(Lemma~\ref{i1:lem-A15I}); $\mathcal A_8$--$\mathcal A_{10}$ are Lemma~\ref{i1:lem-A8};
$\mathcal A_{11}$, $\mathcal A_{14}$, $\mathcal A_{16}$, $\mathcal A_{17}$,
$\mathcal A_{18}$, $\mathcal A_{20}$ are Lemmas~\ref{i1:lem-A11}, \ref{i1:lem-A14},
\ref{i1:lem-A16}, \ref{i1:lem-A17}, \ref{i1:lem-A18}, \ref{i1:lem-A20};
$\mathcal A_{19}$ is Lemma~\ref{i1:lem-A19} with Proposition~\ref{i1:prop-A19c};
$\mathcal A_{21}$ is Lemma~\ref{i1:lem-A21} with Proposition~\ref{i1:prop-A21};
$\mathcal A_{29}$, $\mathcal A_{32}$, $\mathcal A_{33}$ are Lemmas~\ref{i1:lem-A29},
\ref{i1:lem-A32}, \ref{i1:lem-A33} (the second with Proposition~\ref{i1:prop-Fm});
$\mathcal A_{27}$ is Lemma~\ref{i1:lem-A27} with Proposition~\ref{i1:prop-segdeg},
and its constants are those of the statement.
\end{proof}

\subsection{Proofs of Propositions~\ref{i1:prop-loc} and~\ref{i1:prop-linkG}}\label{i1:ssec-pf-loc}

Write $\mathrm{Adm}(C_{La},C_k,C_{rk},C_{deg})$ for the conjunction of the
$37$ clauses of Table~\ref{i1:tab-adm} and $\mathrm{aux}_3(C)$ for the
statement of Proposition~\ref{i1:prop-loc} at the constant $C$ (everything
after ``for every $C\ge C_0$'').

\begin{proposition}[the closed form; formally verified, \cite{CharacterFamily}]\label{i1:prop-closed}
For all constants, $\mathrm{Adm}(C_{La},C_k,C_{rk},C_{deg})$ implies: there is
$C_0>0$ such that $\mathrm{aux}_3(C)$ holds for every $C\ge C_0$.
\end{proposition}

\noindent Its proof in the formal development is a valuation chain bound: the closed
form is derived from two statements for all constants, a membership statement
($\mathrm{RestMEM}$: Lemma~5a of the chain together with a word bound, ten
operator steps and a no-operator step) and a Gram statement (Proposition~5 of
the chain, from three parts), via a test-side reduction (cleaning, plane-Gram,
cloud and edge bounds). Its valuation core (the chain bound and the link bound
for the column lattices) has a written proof, which we do not reproduce; the
operator steps, the parts of Proposition~5, the cloud, flatness and
test-side inputs are proved in the formal development. The statements
involved are formulated in terms of the signed chain data, level polynomials
and Fock models of the formal development, which are not defined in this
paper; so the proposition is stated without proof. (The explicit minor map
and the extended closed form, which appear among the closed-form statements of
the formal development, are not on this proof path; they enter the proof of
Proposition~\ref{i1:prop-linkG} through the first theorem,
Proposition~\ref{i1:prop-first}.)  Its formally verified steps are outlined in \S\ref{i1:str-closed}.

\begin{proof}[Proof of Proposition~\ref{i1:prop-loc}]
By Proposition~\ref{i1:prop-adm} there are constants with
$\mathrm{Adm}(C_0',C_0',C_{rk},C_{deg})$. Proposition~\ref{i1:prop-closed}
then gives $C_0>0$ with $\mathrm{aux}_3(C)$ for all $C\ge C_0$. Neither
statement depends on a column record, so $C_0$ is absolute.
\end{proof}

\medskip\noindent\label{i1:ssec-pf-linkG}We now turn to Proposition~\ref{i1:prop-linkG}.

The proof in the formal development has three layers: (a) a local lattice
estimate at each word $u$ and place $w$, the \emph{service lattice condition};
(b) a gluing step which replaces the lattice data of the families by zero
chain data, so that every valuation certificate holds trivially, and then
checks the five clauses of $\mathrm{LinkG}_{\rm at}$ one by one; (c) the
transport of two theorems on lattices to a family and an assembly of
closedness statements. We prove (a) from one explicitly stated all-but
estimate, the zero-chain step of (b), and the whole of (c) except the first
theorem itself (Proposition~\ref{i1:prop-first}); what remains is stated
exactly as Propositions~\ref{i1:prop-first}, \ref{i1:prop-gc} and
\ref{i1:prop-linkG-res}.

\begin{definition}[Local data at a place]\label{i1:def-locdatum}
Fix an instance (Definition~\ref{i1:def-family-data}), a word $u$ and a place $w$
with $p=p(w)$. The \emph{local datum} at $(u,w)$ consists of a discrete
valuation ring $O$ with fraction field $K_f$ (normalised valuation $v_O$), a
$K_f$-vector space $F$, two $O$-submodules $M\subseteq M^{\rm na}$ of $F$
(the \emph{lattice} and the \emph{naive lattice}), commuting $O$-linear
operators $y_1,\dots,y_{n_V}$ on $F$, an \emph{orbit labelling}
$t_0\mapsto v(t_0)\in\{1..n_V\}$, and a $K_f$-subspace $V_{\rm ram}$. Write
$d=\mathrm{rk}_OM$, $d^{\rm na}=\mathrm{rk}_OM^{\rm na}$, $M_{\rm ram}=M\cap V_{\rm
ram}$. Besides $\mathrm{VV}(p)$, $\Sigma(p)$, $S_v(p,t_0)$, $\tau_P(p)$ and
$\mathrm{io}_P(p)$ the instance provides integers $\mathrm{ep}(p)\ge0$, $D_p\ge0$,
reals $\mathrm{tv}(t_0)$ and weights $\kappa_v(p,t_0)\ge0$, with
$\tau_P(p)=\mathrm{ep}(p)+\sum_{t_0\in\mathrm{VV}(p)}m(\lceil\mathrm{tv}(t_0)\rceil+D_p+1)$
and $\Sigma(p)=m\sum_{t_0\in\mathrm{VV}(p)}\kappa_v(p,t_0)$ when
$\mathrm{VV}(p)\ne\emptyset$, $\tau_P(p)=\mathrm{io}_P(p)=0$ otherwise, and
$S_v(p,t_0)\in\Z$. For $V'\subseteq\mathrm{VV}(p)$ put
$y(V')=\prod_{t_0\in V'}y_{v(t_0)}^m$, $s(V')=\sum_{t_0\in\mathrm{VV}(p)\setminus V'}S_v(p,t_0)$
and
\[
\mathrm{NP}=\mathrm{span}_O\{y(V')(p^{s(V')}x):V'\subseteq\mathrm{VV}(p),\ x\in M\}\cap p^{\lceil\tau_P(p)\rceil}M\cap V_{\rm ram}.
\]
The \emph{service lattice condition} at $(u,w)$ with parameter $c_\star$ is:
if (i) $p^{\mathrm{ep}(p)}M^{\rm na}\subseteq M\subseteq M^{\rm na}$, (ii) $m\ge1$,
(iii) $v$ is injective on $\mathrm{VV}(p)$, (iv) for each $t_0\in\mathrm{VV}(p)$
the lattice statement $\mathrm{LP}'(\lceil\mathrm{tv}(t_0)\rceil,D_p,\kappa_v(p,t_0),c_\star)$
of Definition~\ref{i1:def-slots} holds for the slot of $v(t_0)$, (v) $M_{\rm
ram}$ is saturated in $M$, and (vi) $\Sigma(p)\le\frac12$, then
$\mathrm{NP}\subseteq p^{\lceil\tau_P(p)\rceil}M_{\rm ram}$ and
$\ell_O(M_{\rm ram}/\mathrm{NP})\le v_O(p)\,d\,(\tau_P(p)+\mathrm{io}_P(p))$.
\end{definition}

\begin{definition}[All-but statements]\label{i1:def-allbut}
For a ring $R$, an $R$-module $X$, $\pi\in R$, submodules $P,M$ and $g,E\in\N$,
$(M,P)$ is \emph{$(g,E)$-all-but} if there is $C\subseteq M$ with $|C|\le g$ and
$\pi^Ex\in P+\mathrm{span}_RC$ for all $x\in M$. For a finite set $V$ of labels,
integers $S'_v,t_v\ge0$, $D_\ast,S,m\ge0$ and reals $\kappa_v\ge0$ put
$N(V)=\mathrm{span}_O\{p^{\sum_{v\in V\setminus B}S'_v}\prod_{v\in B}y_v^m(x):B\subseteq V,\ x\in M\}$,
$N^{\rm na}(V)$ likewise with $M^{\rm na}$, and $\tau(V)=\sum_{v\in V}m(t_v+D_\ast+1)$.
The \emph{all-but estimate} (AB) for these data is: for every $\beta\ge0$,
$(M^{\rm na},N^{\rm na}(V))$ is
$(\lfloor d^{\rm na}m^{\beta+1}e_{\beta+1}(\kappa)\rfloor,\tau(V)+\beta S)$-all-but
for $\pi=p$, $e_k$ the $k$-th elementary symmetric polynomial.
\end{definition}

\begin{lemma}[Lengths from all-but statements]\label{i1:lem-allbut}
Let $R$ be a discrete valuation ring, $X$ an $R$-module, $\pi\in R$.
(1) If $P\subseteq M'\subseteq M$ and $(M,P)$ is $(g,E)$-all-but, so is $(M',P)$.
(2) If $B\subseteq A$, $N_A\subseteq A$, $\pi^eN_A\subseteq N_B\subseteq B$ and
$(A,N_A)$ is $(g,E)$-all-but, then $(B,N_B)$ is $(g,E+e)$-all-but. (3) Let
$M\subseteq X$ be free of rank $d$, $N\subseteq X$, $0\ne a\in R$, $\tau,S,V\in\N$,
$0\le s\le\frac12$; if for every $\beta<V$, $(M,N)$ is $(ds^{\beta+1},\tau+\beta
S)$-all-but for $a$, and $a^{VS}M\subseteq N$, then
$\ell_R(M/(N\cap a^\tau M))\le\ell_R(R/aR)(d\tau+2Ssd)$.
\end{lemma}

\begin{proof}
(1) With $C$ for $M$, let $f:R^C\to X/P$, $\alpha\mapsto\sum\alpha_c[c]$, and
$U=f^{-1}(\text{image of }\pi^EM')$; $U$ is free with a basis $b_1..b_n$, $n\le|C|$
($R$ is a principal ideal domain); choose $m_i\in M'$ with $[\pi^Em_i]=f(b_i)$;
then $C'=\{m_i\}$ works. (2) Apply (1) to $M'=B+N_A\subseteq A$ and $P=N_A$, and
multiply the resulting set by $\pi^e$. (3) Reduce to $a=\varpi$ a uniformiser
(writing $a=u\varpi^e$ multiplies $\tau,S$ by $e=\ell_R(R/aR)$). Let
$Q=M/(N\cap\varpi^\tau M)$ and $A_0=\varpi^\tau Q$; $\ell(Q/A_0)\le\tau d$. The
map $A_0\to X/N$ induced by $M\to X/N$ is injective with image in $E_0$, where
$E_j$ is the image of $\varpi^{\tau+jS}M$; $E_{j+1}\subseteq E_j$,
$\varpi^SE_j\subseteq E_{j+1}$, $E_V=0$, and $E_j$ is generated by at most
$ds^{j+1}$ elements, so $\ell(E_j/E_{j+1})\le Sds^{j+1}$. Summing,
$\ell(A_0)\le Sd\sum_{j<V}s^{j+1}\le2Ssd$ as $s\le\frac12$.
\end{proof}

\begin{proposition}[Service lattice condition from the all-but estimate]\label{i1:prop-SLC}
Suppose that at $(u,w)$ the operators $y_v$ commute and preserve $M$ and
$M^{\rm na}$, $M^{\rm na}$ is finitely generated, and (AB) holds for
$V=v(\mathrm{VV}(p))$, $t_v=\lceil\mathrm{tv}(t_0)\rceil$, $S'_v=S_v(p,t_0)$
($v=v(t_0)$), $D_\ast=D_p$, $S=\max_{t_0}S_v(p,t_0)$, $\kappa_v=\kappa_v(p,t_0)$,
whenever (i)--(vi) of Definition~\ref{i1:def-locdatum} hold. Then the service
lattice condition holds at $(u,w)$ for every $c_\star$.
\end{proposition}

\begin{proof}
If $\mathrm{VV}(p)=\emptyset$, then $\tau_P=\mathrm{io}_P=0$ and
$\mathrm{NP}=M_{\rm ram}$ (take $V'=\emptyset$). Otherwise assume (i)--(vi) and
relabel by $v$ (injective by (iii)), so $e+\tau(V)=\tau_P(p)$ with
$e=\mathrm{ep}(p)$ and $m\sum_V\kappa_v=\Sigma(p)$. $M$ is finitely generated and
torsion free, hence free of rank $d$, and $x\mapsto p^ex$ embeds $M^{\rm na}$ into
$M$, so $d^{\rm na}\le d$. \emph{Step 1.} Since the $y_v$ preserve the
lattices, $N(V)\subseteq M$, $N^{\rm na}(V)\subseteq M^{\rm na}$ and
$p^eN^{\rm na}(V)\subseteq N(V)$; Lemma~\ref{i1:lem-allbut}(2) transfers (AB) to
$(M,N(V))$ with exponent $\tau(V)+\beta S+e$, and since
$k!\,e_k(\kappa)\le(\sum\kappa_v)^k$, the sets have at most $d\Sigma^{\beta+1}$
elements. \emph{Step 2.} $p^{\sum_VS'_v}M\subseteq N(V)$ ($B=\emptyset$) and
$\sum_VS'_v\le|V|S$. \emph{Step 3.} Lemma~\ref{i1:lem-allbut}(3) with $a=p$,
$\tau=\tau_P$, $s=\Sigma(p)\le\frac12$ gives
$\ell_O(M/(N(V)\cap p^{\tau_P}M))\le v_O(p)d(\tau_P+2S\Sigma(p))$, and
$2S\Sigma(p)=\mathrm{io}_P(p)$. \emph{Step 4.} The generators of $N(V)$ are
generators of the first factor of $\mathrm{NP}$, so $N(V)\cap p^{\tau_P}M\cap
V_{\rm ram}\subseteq\mathrm{NP}$; by saturation (v), $(N\cap p^{\tau_P}M)\cap
M_{\rm ram}\subseteq p^{\tau_P}M_{\rm ram}$, which gives the inclusion, and
$M_{\rm ram}/(M_{\rm ram}\cap N(V)\cap p^{\tau_P}M)$ injects into
$M/(N(V)\cap p^{\tau_P}M)$, which gives the length bound.
\end{proof}

\begin{definition}[Chain certificates]\label{i1:def-chain}
\emph{Lattice data} $Y$ at $(u,w)$ are $O$-modules $\mathrm{Int}(n)$ ($n\ge0$)
with $O$-linear \emph{chain maps} $K(n):\mathrm{Int}(n)\to\mathrm{Int}(n+1)$; put
$\mathrm{ch}(0)=\mathrm{id}$, $\mathrm{ch}(n+1)=K(n)\circ\mathrm{ch}(n)$. For a vector
$x$ and $g\in\R$ the \emph{valuation certificate} $\mathrm{vKge}_Y(x,g)$ says
that some irreducible $\varpi\in O$ satisfies
$\mathrm{ch}(m)(x)\in\varpi^{\lceil g\,v_O(p)\rceil}\mathrm{Int}(m)$. $Y^0$ denotes
$Y$ with every chain map replaced by $0$.
\end{definition}

\begin{lemma}[Zero chains]\label{i1:lem-zeroch}
If $m\ge1$, then $\mathrm{vKge}_{Y^0}(x,g)$ holds for all $x$ and $g$.
\end{lemma}

\begin{proof}
$\mathrm{ch}(m)=0$ for $m\ge1$; $O$ has an irreducible element, and $0$ lies in
every submodule.
\end{proof}

\medskip

The remaining inputs of $\mathrm{LinkG}$ are organised around two theorems
on lattices, the \emph{second theorem} (the local statement,
Proposition~\ref{i1:prop-loc}) and the \emph{first theorem}
(Proposition~\ref{i1:prop-first} below), transported to a family through
its constants.

\begin{definition}[Loss constants of a family]\label{i1:def-famconst}
A family $\mathrm{Fm}$ of exponent $\varepsilon$ carries real constants,
independent of the point $(x,i)$: $C_\Lambda$, which is an entry of the list
$c_{\rm mid}(C_H)$ of Definition~\ref{i4:def-closed}; the eleven constants
\[
\rho_G,\ k_{\rm II},\ t_{\rm II},\ D_{\rm II},\ s_{\rm II},\ k'_{\rm II},\
D_{\rm I},\ \kappa_1,\ \kappa_2,\ k_{\rm I},\ C_N,
\]
which are entries of the list $c_{\rm last}$; and two normalised constants
$\Pr_1=1$ and $C_e=1$. For a point at scale $x$ put $h=2+8\varepsilon$,
$r_G=\lceil8(h+2)/\varepsilon\rceil+1$, $\delta_1=x^{-\varepsilon/8}$ and
$N_1^{\rm fam}(x)=N_1(r_G,\delta_1,h,x)$ (Definition~\ref{i1:def-goodorbit}).
The \emph{window log bound} at the point is
$N_1^{\rm fam}(x)\le C_N\log x$. A property of points holds
\emph{eventually} if there is $x_0$ such that it holds at every point $(x,i)$
with $x\ge x_0$.
\end{definition}

\begin{definition}[The first theorem]\label{i1:def-FT}
Let $p,n',m,G',N_1,c_{cl},D_p,c_\star\in\N$, $\mathrm{TT}\subseteq\Z/n'$ (an
element is \emph{even} or \emph{odd} according to its representative in
$\{0,\dots,n'-1\}$) and $Y,Y_b,a,b',b_l,b_u\in\R$; put $k=m+1$ and
$N=|\mathrm{TT}|$.
\begin{enumerate}
\item \emph{Segments.} Let $e_0<\dots<e_{J-1}$ be the even elements of
 $\mathrm{TT}$ and $e_J=e_0$. For $j<J$ let $g_j=(e_{j+1}-e_j)\bmod n'$ if this
 is nonzero and $g_j=n'$ otherwise, and let $s_j$ be the number of odd
 $y\in\mathrm{TT}$ with $0<(y-e_j)\bmod n'<g_j$. For $y\in\mathrm{TT}$ let
 $d(y)$ be the least $d\ge1$ with $(x+d)\bmod n'=y$ for some even
 $x\in\mathrm{TT}$ ($d(y)=0$ if there is none), and $y^-=(y-d(y))\bmod n'$.
\item \emph{Hypotheses.} Put $\Lambda_I=C_0N_1^2\log a$ and
 $\delta_I=C_HN_1\Lambda_I/k+2(b_u-b_l)/b'$. $\mathrm{HypI}(C_0,C_H)$ is the
 conjunction of: $1\le Y$; $k^2Y^2\le p$; $J\ge2$; $g_j\le G'$ and $s_j\le N_1$
 for all $j<J$; $1\le N_1$; $2N_1+2\le m$; $8\log(NY)\le\Lambda_I$; $0<b'$,
 $b_l\le b'\le b_u$; $Y_bN\Lambda_Ib'/k\le b_u-b_l\le b'/2$; $C_H\le Y_b$;
 $2Y_bN\Lambda_I<k$; $kY^2Y_bN\Lambda_I\log a\le a/b'$; $\delta_I<\frac16$;
 $p<a/b'$; $C_HN\Lambda_I^2\le k$; for no odd $d\in\N$ is
 $(a/b')(1-2\delta_I)\le p^d\le(a/b')(1+2\delta_I)$; $c_{cl}\le\Lambda_Ib'/k$;
 $C_HNY^2(N_1^5+k^3)\log_pa\le D_p$; $8(D_p+1)\le c_\star$; and
 $C_Hk^3Np^{4G'}(c_\star+Nc_{cl}+n'c_{cl})Y\le b'$.
\item \emph{Windows and weights.} $I_a=\Z\cap[\lceil a\rceil,\lfloor a+a/(kY)\rfloor]$,
 $I_b=\Z\cap[\lceil b_l\rceil,\lfloor b_u\rfloor]$. For
 $\kappa:\Z^{\mathrm{TT}}\times\Z^{\mathrm{TT}}\to\R$,
 $\mathrm{mean}(\kappa)=\bigl(|I_a|^N|I_b|^N\bigr)^{-1}\sum_{\alpha\in I_a^{\mathrm{TT}}}\sum_{\beta\in I_b^{\mathrm{TT}}}\kappa(\alpha,\beta)$
 ($=0$ if a window is empty). For $a_w,b_w:\Z/n'\to\Z$ and
 $\lambda_0:\{0..n'-1\}\times\{0..m\}\to\Z$, $\mathrm{WeightI}$ says: $\lambda_0$
 is a dominant base weight for $(\mathrm{TT},a_w,b_w,c_{cl})$ (the definition
 before Proposition~\ref{i1:prop-loc}), and $a_w(s)\in I_a$, $b_w(s)\in I_b$
 for $s\in\mathrm{TT}$.
\item \emph{Threshold.} With $Q=p^{G'}$ and
 $c=c_\star+CNk\log_pa+4Nc_{cl}+2n'c_{cl}$,
 \[
 t_I(C_0,C)=\sum_{y\in\mathrm{TT}\ \mathrm{odd}}\min\Bigl(b_w(y),\frac{a_w(y^-)}{p^{d(y)}}\Bigr)
 +CNN_1\Bigl(\frac a{p^2}+\frac{b'}p+\frac{\Lambda_Ia}{pk}+\frac{\Lambda_Ib'}k\Bigr)
 +CN\bigl(Qc+kQ+k\log_pa+c_{cl}\bigr)+2n'c_{cl}.
 \]
\end{enumerate}
$\mathrm{FT}(C_0,C_H,C)$ is the statement: for all data as above with
$\mathrm{HypI}(C_0,C_H)$ there is $\kappa:\Z^{\mathrm{TT}}\times\Z^{\mathrm{TT}}\to[0,\infty)$
with $\mathrm{mean}(\kappa)\le CNN_1/Y$ such that for every unramified
quadratic Hida setting at $(p,n')$, all $a_w,b_w,\lambda_0$ with
$\mathrm{WeightI}$ and every integer $t\ge t_I(C_0,C)$, the lattice statement
$\mathrm{LP}'(t,D_p,\kappa(a_w|_{\mathrm{TT}},b_w|_{\mathrm{TT}}),c_\star)$ of
Definition~\ref{i1:def-slots} holds for $r=1$ and the base weight
$\lambda_0$.
\end{definition}

\begin{definition}[The theorems at a family]\label{i1:def-thmfam}
Let $\mathrm{aux}_3(C)$ be the statement of Proposition~\ref{i1:prop-loc} at
the constant $C$ (\S\ref{i1:ssec-pf-loc}). A family satisfies the
\emph{second theorem} if there is $C$ with $0\le C\le\min(k_{\rm II},t_{\rm
II},C_\Lambda,D_{\rm II},s_{\rm II},k'_{\rm II})$ such that $\mathrm{aux}_3(C)$
holds eventually (it does not depend on the point). It satisfies the
\emph{first theorem} if there are $C'_H,C$ with $0\le C'_H\le k_{\rm I}$,
$C'_H\le D_{\rm I}$, $2C'_H\le\rho_G$, $0\le C$, $2C\le\kappa_1$, $C\le\kappa_2$
and $\mathrm{FT}(C_\Lambda,C'_H,C)$ (at every point; $\mathrm{FT}$ does not
depend on the point).
\end{definition}

\begin{proposition}[the first theorem; formally verified, \cite{CharacterFamily}]\label{i1:prop-first}
Assume that a column record exists (as it does as soon as a family exists).
There are $T_0>0$ and functions $T_2,C:\R\to\R$ such that for every
$C_0\ge T_0$: $T_2(C_0)>0$, $C(C_0)>0$, and $\mathrm{FT}(C_0,C_H,C(C_0))$
holds for every $C_H\ge T_2(C_0)$.
\end{proposition}

\noindent The formal statement is quantified over column records, on which
$\mathrm{FT}$ does not depend; this is why a record is assumed to exist. Its
formal proof has three layers. (a) \emph{Closed admission of the projection
lemma:} for every $C_c\ge0$ there are nested thresholds $T_0$, $T_1(C_0)$,
$T_2(C_0,C_{E1},C_{E2},C_m)$ beyond which ten constants satisfy a list of
clauses on the column record (a correspondence of path pairs, two counting
clauses, a lower-set clause for column families, and twelve tool clauses on
$\mathfrak{gl}_n$-modules). These are derived from the admission clauses of
\S\ref{i1:ssec-pf-adm} and from the \emph{extended closed form} of the column
systems, a closed-form implication for the window and segment ranks whose top
layer is five implications between three interface statements (on columns,
on pre-images and on a truncated core). (b) \emph{The second part of the
projection lemma in closed form:} closed admission implies the conclusion
$\mathrm{FT}$ with thresholds. It is reduced, by an explicit pairing on dual
Weyl modules (Proposition~\ref{i1:prop-Fm}), a Smith-type diagonalisation
(Theorem~\ref{i1:thm-elemdiv}) and a contraction statement, to the
\emph{explicit minor map} statement, a rank statement for the minors of the
twisted column embeddings. (c) \emph{Choice of thresholds:} from (a), (b) at a
fixed record $D_0$ one gets $T_0$ and, for every $C_0\ge T_0$, numbers
$T_2,C>0$ with $\mathrm{FT}(C_0,C_H,C)$ for $C_H\ge T_2$; choosing them for each
$C_0$ gives the functions, and since $\mathrm{FT}$ does not depend on the
record the conclusion holds for every record. The intermediate statements of
(a) and (b) are formulated with about thirty predicates of the formal
development on the segment modules and their ranks which are not defined in
this paper, so they are not stated separately here.  Its formally verified steps are outlined in \S\ref{i1:str-first}.

\begin{lemma}[Closedness of the transported statements]\label{i1:lem-thmfam}
Every class $\mathcal C$ of families is closed
(Definition~\ref{i4:def-closed}) for each of the following properties:
(a) the second theorem; (b) the first theorem; (c) $C_\Lambda>0$ and
$C_e\ge0$; (d) $\Pr_1\ge1$, $C_N\ge1$ and $C_e\ge1$; (e) the window log bound
holds eventually. Properties (a) and (b) are taken with
Propositions~\ref{i1:prop-loc} and~\ref{i1:prop-first}.
\end{lemma}

\begin{proof}
In each case we take $\varepsilon_1=1$ and exhibit, for each
$\varepsilon\in(0,1]$, a threshold datum
$\kappa=(\kappa_{\rm pre},\kappa_t,\kappa_{\rm mid},\kappa_{\rm last})$; let
$\mathrm{Fm}\in\mathcal C$ of exponent $\varepsilon$ be $\kappa$-large at
$C_H$.

(a) Let $C_\ast>0$ be the constant $C_0$ of Proposition~\ref{i1:prop-loc},
so $\mathrm{aux}_3(C_\ast)$ holds: a column record exists, namely any point of
$\mathrm{Fm}$, which exists by Definition~\ref{i1:def-family-data}. Take $\kappa=(0,0,C_\ast,C_\ast)$ (constant
functions). As $C_\Lambda$ is an entry of $c_{\rm mid}(C_H)$, $C_\Lambda\ge
C_\ast$; as $k_{\rm II},t_{\rm II},D_{\rm II},s_{\rm II},k'_{\rm II}$ are
entries of $c_{\rm last}$, they are $\ge C_\ast$. So $C=C_\ast$ witnesses the
second theorem, with any $x_0$.

(b) Let $T_0,T_2,C$ be as in Proposition~\ref{i1:prop-first} (a column record
exists, namely any point of $\mathrm{Fm}$, which exists by
Definition~\ref{i1:def-family-data}). Take $\kappa_{\rm pre}=0$,
$\kappa_t=0$, $\kappa_{\rm mid}\equiv T_0$, and let $\kappa_{\rm last}(l)$ be
the maximum of $0$ and of $\max(2T_2(c),2C(c))$ over the entries $c$ of the
list $l$. Then $C_\Lambda\ge T_0$, and since $C_\Lambda$ is an entry of
$c_{\rm out}\frown c_{\rm mid}(C_H)$, every entry of $c_{\rm last}$ is at
least $M=\max(2T_2(C_\Lambda),2C(C_\Lambda))$. Put $C'_H=T_2(C_\Lambda)>0$ and
$C'=C(C_\Lambda)>0$. Then $C'_H\le2C'_H\le M\le k_{\rm I}$, $C'_H\le D_{\rm
I}$, $2C'_H\le\rho_G$, $2C'\le\kappa_1$ and $C'\le2C'\le\kappa_2$; and
$\mathrm{FT}(C_\Lambda,C'_H,C')$ holds by Proposition~\ref{i1:prop-first}
with $C_0=C_\Lambda\ge T_0$ and $C_H=T_2(C_\Lambda)$.

(c) Take $\kappa=(0,0,1,0)$: $C_\Lambda\ge1>0$, and $C_e=1$.

(d) Take $\kappa=(0,0,0,1)$: $C_N\ge1$ as an entry of $c_{\rm last}$, and
$\Pr_1=C_e=1$.

(e) With $h=2+8\varepsilon$ and $r_G=\lceil8(h+2)/\varepsilon\rceil+1$, which
depend only on $\varepsilon$, let $C_e(r_G,h)=20\Lambda(r_G,h)$ be the constant
of Lemma~\ref{i1:lem-osize}, and take $\kappa=(0,0,0,C_e(r_G,h))$. Then
$C_N\ge C_e(r_G,h)$. For $x\ge x_0=\max(x_1,3)$ we have $0<\delta_1=x^{-\varepsilon/8}\le1$,
so Lemma~\ref{i1:lem-osize} gives $N_1^{\rm fam}(x)\le C_e(r_G,h)\log x\le
C_N\log x$, as $\log x\ge0$. So the bound holds eventually, with $x_0=\max(x_1,3)$.
\end{proof}

\begin{proposition}[consequences of group compatibility; formally verified, \cite{CharacterFamily}]\label{i1:prop-gc}
\emph{Stated informally; the exact statement is in \cite{CharacterFamily}. The three predicates below concern the local gatherings
of Section~\ref{i4:sec} and are not defined here; group compatibility itself is
Definition~\ref{i1:def-groupcompat}.} If $\mathcal C$ is
closed for group compatibility, then it is closed for each of: (i)
eventually, at every word $u$ and place $w$ with gathering index, the local
gathering has trivial stabilisers and finitely many classes; (ii)
eventually, the standard-subring condition holds at the column record; (iii)
the descent of the norm lattices holds for the family.
\end{proposition}

\noindent Each is Lemma~\ref{i4:lem-closed}(c) applied to a pointwise
implication: (i) comes from the fields (herm) and (gat) of the group
compatibility record, (ii) from the field (ring), (iii) from the field
(cmp), each in a few lines of the formal development once the gathering
objects are available.  Its formally verified steps are outlined in \S\ref{i1:str-gc}.

\begin{proposition}[the remaining inputs of $\mathrm{LinkG}$; formally verified, \cite{CharacterFamily}]\label{i1:prop-linkG-res}
\emph{Stated informally; the exact statement is in \cite{CharacterFamily}. The
three predicates (i)--(iii) of Proposition~\ref{i1:prop-gc} and
clauses (1), (5) of Definition~\ref{i1:def-linkG} are not defined here; all other notions are
defined above.} Let $\mathrm{Fm}$ be a family satisfying: the second theorem
and the first theorem (Definition~\ref{i1:def-thmfam}); $C_\Lambda>0$;
$C_e\ge1$; the window log bound eventually; (i) and (ii) of
Proposition~\ref{i1:prop-gc} eventually and (iii); and $\mathrm{Adm2F}$
(Definition~\ref{i1:def-groupcompat}). Then $\mathrm{LinkG}(\mathrm{Fm})$
holds.
\end{proposition}

\begin{proof}[Proof of Proposition~\ref{i1:prop-linkG}]
Let $\mathcal C$ be closed for group compatibility and for $\mathrm{Adm2F}$.
By Lemma~\ref{i1:lem-thmfam}, Proposition~\ref{i1:prop-gc} and the hypothesis,
$\mathcal C$ is closed for each of the properties (a)--(e) of
Lemma~\ref{i1:lem-thmfam}, (i)--(iii) of Proposition~\ref{i1:prop-gc}, and
$\mathrm{Adm2F}$; by Lemma~\ref{i4:lem-closed}(b) it is closed for their
conjunction. That conjunction contains every hypothesis of
Proposition~\ref{i1:prop-linkG-res} ($C_e\ge1$ comes from (d)), so it implies
$\mathrm{LinkG}(\mathrm{Fm})$ for every family; by
Lemma~\ref{i4:lem-closed}(c), $\mathcal C$ is closed for $\mathrm{LinkG}$.
The explicit form of the statement is the definition of closedness.
\end{proof}

\noindent\emph{On the proof of Proposition~\ref{i1:prop-linkG-res}.} The
formal proof first derives, eventually along the family, nine family
admission clauses: the type-two and type-one hypotheses from the second and
first theorems (by comparing their constants with the loss constants), the
share bound of the least serving weights, the exponent clause (from
$C_e\ge1$), the orbit clauses and the size bounds (from the window log bound
and Lemmas~\ref{i1:lem-osize}, \ref{i1:lem-A2}), and $\mathrm{Adm2F}$; then the
orbit lattice bounds (hypothesis (iv) of Definition~\ref{i1:def-locdatum},
from the two theorems, (ii) and the window log bound), the lattice sandwich
(hypothesis (i), from the exponent clause) and the fact that the operators
$y_v$ preserve the lattices. The last step, which we carry out here, glues
these into $\mathrm{LinkG}$. Replacing the lattice data by $Y^0$
(Lemma~\ref{i1:lem-zeroch}) makes every valuation certificate true, so the
unit bound, the service lower valuation and the valuation clause of the
gluing statement hold whatever the gains are. The two lattice clauses of the
gluing statement ($\mathrm{NP}\subseteq p^{\lceil\tau_P\rceil}M_{\rm ram}$ and
the length bound) are the service lattice condition with
$c_\star=8(D_p+1)$, which holds by Proposition~\ref{i1:prop-SLC} once (AB)
holds; (AB) itself is the formal consequence of the trivial stabilisers (i)
and of the preservation of the lattices. The hypotheses (i)--(vi) of the
service lattice condition are the sandwich, $m\ge1$, injectivity of the orbit
labelling on $\mathrm{VV}(p)$, the orbit lattice bounds, the saturation of
$M\cap V_{\rm ram}$ (if $c\ne0$ and $cx\in V_{\rm ram}$ then $x\in V_{\rm ram}$),
and $\Sigma(p)\le x^{-\varepsilon}\le\frac12$ on served primes. Then the
clauses of $\mathrm{LinkG}_{\rm at}$ are checked one by one, with constant
$4|C|+1$ where $C$ is the constant of the family statement of the gains:
clause (1) from the gluing statement (if $\mathrm{VV}(p)=\emptyset$ the
inclusion at exponent $0$ is the unserved local inclusion; otherwise the
length bound with the lattice-length identity, which comes from (iii), after
cancelling $v_O(p)>0$); clause (2) is the service bound; clause (4) holds by
the definition of $\mathrm{VV}$; clause (5) is the mean bound; and clause
(3), the comparison \eqref{i1:eq-linkG3}, is the transport through the fills
($\mathrm{gain}(j,w)=\mathrm{gain}_P(p(w),j)$, $\tau(w)=\tau_P(p(w))$) of the
gain bound, which the formal development proves from the explicit gains of
Remark~\ref{i1:rem-gain} (star inequality on each served orbit, a sum over
orbits, and the passage between the two forms of the ask series), using
only the orbit-size clause of the good orbits (Lemmas~\ref{i1:lem-osize},
\ref{i1:lem-A2}) and elementary estimates.  Its formally verified steps are outlined in \S\ref{i1:str-linkG-res}.

\subsection{Column polynomials, dual Weyl modules and segment degrees}\label{i1:ssec-colpol}

This subsection contains the proofs of the four propositions on column lattices stated in \S\ref{i1:ssec-pf-adm}: Propositions~\ref{i1:prop-A19c}, \ref{i1:prop-A21}, \ref{i1:prop-Fm} and the parts (I1), (I4), (I9), (Deg) of Proposition~\ref{i1:prop-segdeg}.

\subsubsection{Proof of Proposition~\ref{i1:prop-A19c}}\label{i1:ssec-pf-A19c}

Throughout, $F$ is a field of characteristic $0$, $S_i=\{l<i\}$, and for an
injection $e$ of $\{0..m-1\}$ into the columns we use the operators
$E^{S,e}_{ab}$ and the spaces $\mathrm{iso}(S,e,\lambda)$ of
Definition~\ref{i1:def-jt} ($e=\mathrm{id}$ in $m$ columns, $e=\mathrm{pl}$, or
$e=\mathrm{tan}$). For a set of letters $S$, $V^S[\mu]=\mathrm{iso}(S,\mathrm{id},\mu)$.
By (P) we mean the path basis: for every ring $R$ and multidegree $d$, the
$\Psi_{\nu,I}$ ($\nu$ a Pieri path of degrees $d$, $I$ standard) form an
$R$-basis of $\mathrm{Pol}_d$, finitely many (Lemmas~\ref{i1:lem-A15},
\ref{i1:lem-A15I}, \ref{i1:lem-counts}); by (FL) the first-letters lower set
property in the proof of Lemma~\ref{i1:lem-A16}. The ring $O$ has
characteristic $0$, so $k_0$ does. We use two inputs, proved first.
\begin{itemize}
\item[(JC)] Over $F$, if $\lambda=(\lambda_i)_{i\le N}$ is not a Pieri path of
 degrees $d$ (in $n$ columns), then $\mathrm{Pol}_d\cap\bigcap_{i\le N}V^{S_i}[\lambda_i]=0$.
\item[(PP)] (1) The products $v_x=\Psi^{pl}_{\nu,I'}\Psi^{tan}_{\kappa,I''}$,
 $x=((\nu,I'),(\kappa,I''))$ in the set $X_B$ of \emph{pair paths} ($\nu,\kappa$
 Pieri paths in $r$, resp.\ $K$, coordinates of degrees $b',b''$ with
 $b'+b''=B$) with standard fillings, form an $R$-basis of $\mathrm{Pol}_B$ in
 $k$ coordinates over every commutative ring $R$. (2) For families of dominance lower sets $\pi'_i,\pi''_i$
 ($i\le N$) of dominant weights,
 $\{f\in\mathrm{Pol}_B:\bar f\in\sum_{q:\,q_{1,i}\in\pi'_i,\,q_{2,i}\in\pi''_i\,\forall i}\mathrm{JT}(q)\}$
 is the $O$-span of the $v_x$ with $\nu_{x,i}\in\pi'_i$, $\kappa_{x,i}\in\pi''_i$ for all $i$.
\end{itemize}

\emph{Isotypic projectors.} For each $(S,e)$ the $E^{S,e}_{ab}$ form a
$\mathfrak{gl}_m$-representation: each is a derivation, and the relation
$[E_{ab},E_{cd}]=\delta_{bc}E_{ad}-\delta_{da}E_{cb}$ between derivations needs
only to be checked on variables, where it is a direct computation using
injectivity of $e$ (separately for vector and covector letters). Monomials are
weight vectors, of weight $a\mapsto\sum_{l\in S\text{ vector}}m_{l,e(a)}-\sum_{l\in S\text{ covector}}m_{l,e(a)}$,
and the operators preserve multidegrees, so every polynomial lies in a
finite-dimensional stable subspace. By Lemmas~\ref{i1:lem-cr} and
\ref{i1:lem-isodec} there are isotypic projectors $e^{S,e}_\lambda$. We use:
(O1) if a linear map $T$ preserves every isotypic component, it commutes with
every projector ($e_\mu Tx=\sum_\nu e_\mu Te_\nu x=Te_\mu x$); (O2) if $T$
commutes with all $E^{S,e}_{cd}$, it preserves every component (a component is
spanned by the $uv$, $u$ a word, $v$ a highest weight vector, and $Tuv=uTv$
with $Tv$ zero or a highest weight vector of the same weight); (O3) each
$E^{S,e}_{ab}$ preserves each component.

\emph{(Q1) Commutation.} If $S\subseteq S'$, the projectors of $(S,e)$ and
$(S',e)$ commute: $E^{S',e}_{ab}=E^{S,e}_{ab}+E^{S'\setminus S,e}_{ab}$, and the
second summand commutes with all $E^{S,e}_{cd}$ (a commutator of derivations,
zero on every variable because the two operators act on variables of disjoint
sets of letters); so $E^{S',e}_{ab}$ preserves the components of $(S,e)$
(O2, O3), the projectors of $(S,e)$ commute with it (O1), hence preserve the
components of $(S',e)$ (O2) and commute with its projectors (O1). If $e,e'$
have disjoint images the projectors of $(S,e)$ and $(S',e')$ commute, by the
same argument ($E^{S',e'}_{cd}$ and $E^{S,e}_{ab}$ act on disjoint sets of
columns). \emph{(Q2) Degrees.} The projection $p_d$ onto $\mathrm{Pol}_d$
commutes with every $E^{S,e}_{ab}$, hence (O2, O1) with every projector; so the
projectors preserve $\mathrm{Pol}_d$.

\emph{(Q3) Evaluation.} For $y\in F^{N\times k}$ let $T_y$ be the $F$-algebra
map sending $X_{l,e(a)}\mapsto X_{l,a}$ and every other variable $X_{l,c}$ to the
constant $y_{l,c}$. Then $T_y\circ E^{S,e}_{ab}=E^{S,\mathrm{id}}_{ab}\circ T_y$
(both sides satisfy the Leibniz rule along $T_y$ and agree on variables), so
$T_y$ maps $\mathrm{iso}(S,e,\lambda)$ into $\mathrm{iso}(S,\mathrm{id},\lambda)$
(words to words, highest weight vectors to highest weight vectors of the same
weight or $0$). If $x\ne0$, some $T_yx\ne0$: since $F$ is infinite, $x$ has a
nonzero value at some point $y'$, and $T_yx$ evaluated at
$z_{l,a}=y'_{l,e(a)}$ is $x(y')$ when $y=y'$ off the variables of $e$.
Consequently $\mathrm{iso}(S,e,\lambda)=0$ for $\lambda$ not dominant
(Lemma~\ref{i1:lem-hw}(2)), and also for $e=\mathrm{id}$.

\emph{(Q4) Lower sets in $n$ columns.} Let $i\le N$ and $\pi$ a dominance
lower set of dominant weights. Then
$\mathrm{Pol}_d\cap\sum_{\mu\in\pi}V^{S_i}[\mu]=\mathrm{span}_F\{\Psi_{\nu,I}:\nu_i\in\pi\}$
($\nu$ Pieri of degrees $d$, $I$ standard). By (FL) and (P) it suffices to
show, for $G\in\mathrm{Pol}_d$: $e^{S_i}_\mu G\ne0$ implies $\mu\in\pi$ iff every
dominant weight occurring in $M_G$ (Lemma~\ref{i1:lem-A13}(2), $S=S_i$) lies in
$\pi$. By the proof of Lemma~\ref{i1:lem-A13}(3), $M_G=UG$ in characteristic
$0$. ($\Leftarrow$) If $e_\mu G\ne0$, $UG$ contains a highest weight vector of
weight $\mu$, a nonzero polynomial all of whose monomials have weight $\mu$; so
$\mu$ occurs, and is dominant. ($\Rightarrow$) Let $\nu$ dominant occur through
$0\ne x\in UG$ and $m$ a monomial of $x$. As the projectors commute with words,
$e_\mu G=0$ implies $e_\mu x=0$; since $x=\sum_\mu e_\mu x$, $m$ is a monomial of
some $e_\mu x\ne0$, so $e_\mu G\ne0$ and $\mu\in\pi$. By
Lemma~\ref{i1:lem-hw}(3) all monomials of elements of $V^{S_i}[\mu]$ have
weights $\preceq\mu$; so $\nu\preceq\mu$ and $\nu\in\pi$.

\emph{(Q5) Two abstract facts.} Let $V$ be an $F$-vector space, $\kappa$ a
finite nonempty set, and for $i\in\kappa$: a partially ordered set
$\Lambda_i$ with a subset of \emph{dominant} labels, subspaces $W_i(a)$ with
$V=\bigoplus_aW_i(a)$ and projections $e^i_a$, such that all $e^i_a,e^j_b$
commute and $W_i(a)=0$ for $a$ not dominant. Let $\mathrm{Pl}\subseteq V$ be stable
under all $e^i_a$, with a finite basis $(b_x)_{x\in X}$ and labels
$P(x)\in\prod_i\Lambda_i$, all $P_i(x)$ dominant, such that (low) for each $i$
and each set $\pi$ of dominant labels closed downwards among dominant labels,
$\mathrm{Pl}\cap\sum_{a\in\pi}W_i(a)=\mathrm{span}\{b_x:P_i(x)\in\pi\}$. Then
(a) $\dim(\mathrm{Pl}\cap\bigcap_iW_i(\lambda_i))=\#\{x:P(x)=\lambda\}$ for every
$\lambda$; (b) for such downward closed $\pi_i\subseteq\Lambda_i$,
$\mathrm{Pl}\cap\sum_{q:\,q_i\in\pi_i\,\forall i}\bigcap_iW_i(q_i)=\mathrm{span}\{b_x:P_i(x)\in\pi_i\ \forall i\}$.

We use (s): $y\in\sum_{a\in\pi}W_i(a)$ iff $e^i_ay=0$ for $a\notin\pi$; so these
sums are stable under all $e^j_b$; and (int): if $y$ is in the span of
$\{b_x:x\in A_k\}$ for each $k$ of a nonempty family, it is in the span over
$\bigcap A_k$ (unique coordinates). (a) Let
$L_i=\{a\text{ dominant}:a\le\lambda_i\}$, $L_i^-=L_i\setminus\{\lambda_i\}$ (both
downward closed), $A=\mathrm{Pl}\cap\bigcap_i\sum_{L_i}W_i$,
$A^-=\sum_j(A\cap\sum_{L_j^-}W_j)$, $U=\mathrm{Pl}\cap\bigcap_iW_i(\lambda_i)$. By (low)
and (int), $A$ is spanned by the $b_x$ with all $P_i(x)\in L_i$, and $A^-$ by
those among them with $P(x)\ne\lambda$ (if $P_j(x)\ne\lambda_j$ then
$P_j(x)\in L_j^-$). $U\subseteq A$ (if $\lambda_i$ is not dominant,
$W_i(\lambda_i)=0$). $A$ is stable under all $e^j_b$. For $y\in A$, $y-e^k_{\lambda_k}y\in A\cap\sum_{L_k^-}W_k$
by (s). With $E=\prod_ie^i_{\lambda_i}$, telescoping gives $y-Ey\in A^-$ for
$y\in A$; $E$ maps $\mathrm{Pl}$ into $U$, is the identity on $U$ and kills $A^-$
(on the $j$-th summand $e^j_{\lambda_j}$ kills by (s) and commutes with the
other factors). So $A=U\oplus A^-$, and $\dim U=\dim A-\dim A^-=\#\{x:P(x)=\lambda\}$
(as $P(x)=\lambda$ forces $P_i(x)\in L_i$). (b) $\subseteq$: each
$\bigcap_iW_i(q_i)$ with $q_i\in\pi_i$ lies in $\bigcap_i\sum_{\pi_i}W_i$; apply
(low) for each $i$ and (int). $\supseteq$: let $P_i(x)\in\pi_i$ for all $i$; then
$b_x\in\mathrm{Pl}\cap\sum_{\pi_i}W_i$ for each $i$. By induction on a set
$s\subseteq\kappa$: if $y\in\sum_{\pi_i}W_i$ for all $i$ and $y\in W_i(q_i)$ for
$i\notin s$ (some $q$ with $q_i\in\pi_i$), then $y$ is in the joint sum; for
$s\cup\{k\}$ write $y=\sum_{a\in\pi_k}e^k_ay$ (by (s)) and apply the hypothesis
to each $e^k_ay$, which keeps the other memberships by commutation. Take $s=\kappa$.

\emph{Proof of (JC).} Apply (Q5)(a) with $V=F[\iota]_n$, $\kappa=\{0..N\}$,
$\Lambda_i=\Z^n$ with dominance order, dominant meaning antitone,
$W_i=V^{S_i}$ (commuting projectors by (Q1), zero off dominant labels by
(Q3)), $\mathrm{Pl}=\mathrm{Pol}_d$ (stable by (Q2)), basis the $\Psi_{\nu,I}$ of
(P), $P_i(\nu,I)=\nu_i$, and (low) given by (Q4). The dimension is the number of
$(\nu,I)$ with $\nu=\lambda$, which is $0$ if $\lambda$ is not a Pieri path of
degrees $d$.

\emph{Proof of (PP).} (1) Independence of the $v_x$ is the argument
of the proof of Lemma~\ref{i1:lem-A19}(a) (products of independent path
bases in disjoint variables are independent). Each $v_x$ lies in $\mathrm{Pol}_B$ by (P) in $r$ and $K$
columns (degrees at each letter add). A monomial of multidegree $B$ in $k$
columns factors uniquely as a plane monomial of some multidegree $b'$ times a
tangent one of multidegree $b''=B-b'$; expanding each factor by (P) and
renaming shows it is a combination of the $v_x$. (2) First over $F$: for $i\le N$
and $\pi$ as above,
$\mathrm{Pol}_B\cap\sum_{\mu\in\pi}\mathrm{iso}(S_i,\mathrm{pl},\mu)=\mathrm{span}\{v_x:\nu_{x,i}\in\pi\}$,
and likewise for $\mathrm{tan}$ and $\kappa_{x,i}$. Indeed: (a) $x\in\sum_\pi\mathrm{iso}(S_i,\mathrm{pl},\cdot)$
iff $T_yx\in\sum_\pi V^{S_i}[\cdot]$ for all $y$ (forward by (Q3); backward:
$T_ye^{pl}_\mu x=e_\mu T_yx$ by (Q3), so if $\mu\notin\pi$ and $e^{pl}_\mu x\ne0$,
some $y$ gives $0\ne T_ye^{pl}_\mu x=e_\mu T_yx=0$). (b) $T_yv_x=t\Psi_{\nu_x,I'_x}$
with $t\in F$ the value of $\Psi_{\kappa_x,I''_x}$ at the tangent coordinates of
$y$. ($\supseteq$) If $\nu_{x,i}\in\pi$ then $\Psi_{\nu_x,I'_x}\in\sum_\pi V^{S_i}$ by
(Q4) in $r$ columns, so $v_x$ is in the left side by (a), (b). ($\subseteq$) Write
$G=G_1+G_2$ by (1), $G_1$ in the span of the $v_x$ with $\nu_{x,i}\in\pi$,
$G_2$ in the span of the others; then $G_2$ is in the left side. If $G_2\ne0$
choose $y$ with $z=T_yG_2\ne0$; then $z\in\sum_\pi V^{S_i}$ by (a) and $z$ lies in
the span of the $\Psi_{\nu,I}$ (all degrees) with $\nu_i\notin\pi$ by (b). For each
multidegree $b$, $p_bz$ lies in $\mathrm{Pol}_b\cap\sum_\pi V^{S_i}$ (by (Q2)),
which is spanned by the $\Psi_{\nu,I}$ of degree $b$ with $\nu_i\in\pi$ (Q4),
and also in the span of those with $\nu_i\notin\pi$; by (P) $p_bz=0$. So $z=0$,
a contradiction. Now apply (Q5)(b) with $\kappa$ two copies of $\{0..N\}$,
$W=\mathrm{iso}(S_i,\mathrm{pl},\cdot)$ on the first and
$\mathrm{iso}(S_i,\mathrm{tan},\cdot)$ on the second (commuting by (Q1), zero
off dominant labels by (Q3)), $\mathrm{Pl}=\mathrm{Pol}_B$ (Q2), basis $(v_x)$
(finite by (P)), labels $((\nu_{x,i})_i,(\kappa_{x,i})_i)$, and (low) just
proved; since $\bigcap_i\mathrm{iso}(S_i,\mathrm{pl},q_{1,i})\cap\mathrm{iso}(S_i,\mathrm{tan},q_{2,i})=\mathrm{JT}(q)$,
we get (2) over $F$. Over $O$: put $Y=\sum\mathrm{JT}(q)$ over $k_0$ and $A$ the set
of $x$ with $\nu_{x,i}\in\pi'_i$, $\kappa_{x,i}\in\pi''_i$. The $\bar v_x$ are
independent over $k_0$ and $\mathrm{span}(\bar v_x)\cap Y=\mathrm{span}\{\bar v_x:x\in A\}$.
If $f=\sum c_xv_x\in\mathrm{Pol}_B$ has $\bar f\in Y$, then $\bar c_x=0$ for $x\notin A$,
hence $c_x=0$ ($O\subseteq k_0$); the converse is clear.

\begin{proof}[Proof of Proposition~\ref{i1:prop-A19c}]
\emph{Step 1.} (J1) If $(\nu,\kappa)$ is a pair path then
$\mathrm{JT}(\nu,\kappa)\subseteq\mathrm{Pol}_B$: by the total-weight argument
of (J2) below applied to $\mathrm{pl}$ and to $\mathrm{tan}$, every monomial of an
element of $\mathrm{JT}(\nu,\kappa)$ has plane degree $b'_i$ and tangent degree
$b''_i$ at letter $i$. (J2) If $\mathrm{JT}(\nu,\kappa)\ne0$ over $F$,
then $\nu$ and $\kappa$ are Pieri paths. For $\nu$: take $0\ne x$ in the joint
type and $y$ with $x'=T_yx\ne0$ (plane embedding, (Q3)). Then $x'\in\bigcap_iV^{S_i}[\nu_i]$
in $r$ columns. In characteristic $0$ all monomials of an element of
$V^{S}[\lambda]$ have total weight $\sum_a\lambda_a$ (the operator
$\sum_aE^S_{aa}$ commutes with every $E^S_{cd}$, so it acts by
$\sum_a\lambda_a$ on each $uv$, $v$ a highest weight vector of weight
$\lambda$; and it is diagonal on monomials, with eigenvalue the total weight);
comparing $S_{i+1}$ and $S_i$,
every monomial of $x'$ has degree $d_i=\pm(\sum_a(\nu_{i+1})_a-\sum_a(\nu_i)_a)$
at letter $i$ (sign $+$ for a vector). So $x'\in\mathrm{Pol}_d\cap\bigcap_iV^{S_i}[\nu_i]$,
and $\nu$ is a Pieri path of degrees $d$ by (JC). Likewise for $\kappa$.

\emph{Step 2 (downward closure).} Let $x\in\mathcal P$ and $q=(q_1,q_2)$ a pair of
Pieri paths with $q_{1,i}\preceq x_{1,i}$ and $q_{2,i}\preceq x_{2,i}$ for all
$i$. Then $q\in\mathcal P$. Indeed $\lambda\preceq\mu$ gives
$\lambda_0\le\mu_0$, $\mu_{\rm last}\le\lambda_{\rm last}$ and equal coordinate
sums; equal sums give equal degrees at every letter, so $q$ has degrees $B$;
by Lemma~\ref{i1:lem-A17}, applied to $x$ and then to $q$, the concatenation
of $q$ is a Pieri path, the interlacing passing from $x$ to $q$ via
$q_{2,i+1,0}\le x_{2,i+1,0}\le x_{1,i,r-1}\le q_{1,i,r-1}$ for a vector
letter $i$ and $q_{2,i,0}\le x_{2,i,0}\le x_{1,i+1,r-1}\le q_{1,i+1,r-1}$ for a
covector letter $i$ (the middle inequality is the interlacing of $x$,
direction ($\Rightarrow$) of Lemma~\ref{i1:lem-A17}; the outer two are
$\lambda_0\le\mu_0$ and $\mu_{\rm last}\le\lambda_{\rm last}$ for
$\lambda\preceq\mu$); and each inequality of
the list ($\theta\le\nu_{i,r-1}$, $\kappa_{i,0}\le M$,
$\kappa_{j,0}+m\le\nu_{i,r-1}$) passes from $x$ to $q$ by the same two
comparisons.

\emph{Step 3.} For $x\in\mathcal P$ let $\Pi(x)$ be the set of pairs $q$ with
$q_{1,i}\in\mathrm{prin}(x_{1,i})$, $q_{2,i}\in\mathrm{prin}(x_{2,i})$ for all
$i$, where $\mathrm{prin}(\mu)$ is the set of dominant $\lambda\preceq\mu$, a
dominance lower set of dominant weights; $x\in\Pi(x)$ (Pieri paths are
antitone). By Steps 1 (J2) and 2,
$\sum_{q\in\Pi(x)}\mathrm{JT}(q)\subseteq\sum_{q\in\mathcal P}\mathrm{JT}(q)$.
($\subseteq$) A generator $v_x$ of $X(\mathcal P)$ has $\bar v_x\in\sum_{\Pi(x)}\mathrm{JT}$
by (PP)(2) (as $x\in\Pi(x)$), so $\bar v_x\in\sum_{\mathcal P}\mathrm{JT}$; the
condition is $O$-linear. ($\supseteq$) Let $Y'$ be the set of $y$ over $k_0$ with
$ay=\bar f'$ for some $a\ne0$ in $O$ and $f'\in X(\mathcal P)$, a $k_0$-subspace.
For $q\in\mathcal P$ and $y\in\mathrm{JT}(q)$, write $ay=\bar f'$ with $a\ne0$ and
$f'$ over $O$; then $\bar f'\in\mathrm{JT}(q)\subseteq\mathrm{Pol}_B$ by (J1), so
$f'\in\mathrm{Pol}_B$, and $\bar f'\in\sum_{\Pi(q)}\mathrm{JT}$; by (PP)(2), $f'$ is
in the span of the $v_{x'}$ with $x'$ in $\Pi(q)$, which lie in $\mathcal P$ by
Step 2; so $f'\in X(\mathcal P)$ and $y\in Y'$. Hence if
$\bar f\in\sum_{\mathcal P}\mathrm{JT}\subseteq Y'$, then $a\bar f=\bar f'$ with
$f'\in X(\mathcal P)$, so $af=f'$ and $f\in X(\mathcal P)$ by
Lemma~\ref{i1:lem-A19}(a).
\end{proof}

\subsubsection{Proof of Proposition~\ref{i1:prop-A21}}\label{i1:ssec-pf-A21}

Notation of Definitions~\ref{i1:def-lift}, \ref{i1:def-ctc}; $W=W(\mathbb F_{p^f})$,
$w^2=p$ in $O$, $k=r+K$, $h_j=\mathrm{colht}(\mu_j,c)$, and $k_O=O/wO$ is the
residue field. Forms, $\wedge$, $\lrcorner$, $\mathbf 1$ and $M\cdot K$ are as
in the proof of Lemma~\ref{i1:lem-A15I}, with the identities (W), (Ct) proved
there. Besides these we use the following standard facts on $W$ (reference to
locate, as in the proof of Lemma~\ref{i1:lem-A28}): (i) the $0$-th Witt
coordinate $\mathrm{cc}:W\to\mathbb F_{p^f}$ is a ring homomorphism; (ii) the
Frobenius $\sigma$ satisfies $\mathrm{cc}(\sigma x)=\mathrm{cc}(x)^p$; (iii) an
element with $\mathrm{cc}=0$ lies in $pW$ (the residue field being perfect).

\emph{(F1) Frobenius congruence.} For $s\le s'$ and $\beta\in W$,
$\sigma^{(s)}(\beta)^{p^{s'-s}}-\sigma^{(s')}(\beta)\in pW$. With $a=s\bmod f$,
$a'=s'\bmod f$ we have $a+(s'-s)=a'+ft$, $t\ge0$, so the $0$-th coordinate of
the difference is $\mathrm{cc}(\beta)^{p^{a'+ft}}-\mathrm{cc}(\beta)^{p^{a'}}=0$,
since $y^{p^{ft}}=y$ on $\mathbb F_{p^f}$; and $\mathrm{cc}$ is additive modulo
$p$ in the sense that the difference has $\mathrm{cc}$ equal to the difference
of the $\mathrm{cc}$'s (Witt subtraction in coordinate $0$ is subtraction).

\emph{(F2) Minors.} For a $k\times k$ matrix $m$ over a commutative ring and
$I\subseteq\{0..k-1\}$ put $\Lambda(m)_{I,J}=(m\cdot\delta_I)(J)$, $\delta_I$ the
indicator of $I$, so that $m\cdot K=\sum_IK(I)\Lambda(m)_{I,\cdot}=:K\Lambda(m)$.
By (W), (Ct): ($\Lambda$1) $(K\wedge x)\Lambda(m)=(K\Lambda(m))\wedge(xm)$;
($\Lambda$2) if $mm'=1$, $(K\lrcorner y)\Lambda(m)=(K\Lambda(m))\lrcorner(m'y)$
(as $y(m^{-1})^T$ is the column $m'y$ written as a row); ($\Lambda$3)
$\Lambda(1)=1$ and $\Lambda(mm')=\Lambda(m)\Lambda(m')$: writing
$\delta_I=\delta_{I\setminus b}\wedge e_b$, $b=\max I$, and using ($\Lambda$1)
twice, $\delta_I\Lambda(mm')=(\delta_{I\setminus b}\Lambda(mm'))\wedge e_bmm'$,
which by induction on $|I|$ is
$((\delta_{I\setminus b}\Lambda(m))\Lambda(m'))\wedge(e_bm)m'=(\delta_I\Lambda(m))\Lambda(m')$;
in particular $mm'=1$ gives $\Lambda(m)\Lambda(m')=1$; ($\Lambda$4)
$\mathbf 1\Lambda(m)=\mathbf 1$; ($\Lambda$5) $\Lambda(\varphi m)=\varphi\Lambda(m)$
entrywise for every ring map $\varphi$.

\emph{(F3) One Frobenius step.} Let $A$ be a domain, $\varpi\in A$ with
$\varpi^2=p$, $q=p^e$, and $L,L',L_2,L'_2$ matrices over $A$ indexed by subsets,
with $LL'=1$, $L'_2L_2=1$ and $p\mid(L'_{IJ})^q-(L'_2)_{IJ}$ for all $I,J$. If
$\varpi^o\mid(KL)(J)$ for all $J$, then for every $J$: (a)
$\varpi^{oq}\mid(K^qL_2)(J)$, $K^q$ the entrywise power; (b) if
$\varpi^{o+1}\mid(KL)(J)$ then $\varpi^{oq+1}\mid(K^qL_2)(J)$, and
$\varpi^{oq+2}\mid(K^qL_2)(J)$ if moreover $e>0$. Indeed write
$(KL)(J)=\varpi^oZ_J$; then $K(I)=\varpi^o\sum_JZ_JL'_{JI}$. Since
$p\mid\binom qi$ for $0<i<q$, $p\mid(\sum_ia_i)^q-\sum_ia_i^q$ for any finite
sum; hence $K^q=\varpi^{oq}(Z^qL'_2+p\rho)$ for some $\rho$, and
$(K^qL_2)(J)=\varpi^{oq}(Z_J^q+p(\rho L_2)(J))$. This is (a). If
$\varpi\mid Z_J$ (which follows from $\varpi^{o+1}\mid\varpi^oZ_J$, $A$ a domain)
then $\varpi\mid Z_J^q$ and $\varpi^2=p$, giving (b); if $e>0$ then $q\ge2$ and
$\varpi^2\mid Z_J^q$.

\emph{(F4) Twisting structures.} Let $A=O[X]$ and let
$\zeta_j,y_j\in A^k$ ($j<N$). Write $\mathrm{step}_j(x,K)$ for the column-word
step at letter $j$: $K\wedge x$ if ($h_j<h_{j+1}$, $c\ge1$) or ($h_{j+1}<h_j$,
$c<1$); $K\lrcorner x$ if ($h_j<h_{j+1}$, $c<1$) or ($h_{j+1}<h_j$, $c\ge1$);
$K$ if $h_j=h_{j+1}$. A \emph{twisting structure} is a family of maps $T_s$
($s\in\Z$) on forms over $A$ with: (Tw1)
$T_{\mathrm{pos}(j)}(\mathrm{step}_j(\zeta_j,K))=\mathrm{step}_j(y_j,T_{\mathrm{pos}(j)}K)$;
(Tw2) $T_s(\mathbf 1)(J)=0$ for $J\ne\emptyset$; (Tw3) if $s\le s'$,
$q=p^{s'-s}$ and $w^o\mid T_sK(J)$ for all $J$, then $w^{oq}\mid T_{s'}(K^q)(J)$
for all $J$, and for each $J$ with $w^{o+1}\mid T_sK(J)$ also
$w^{oq+1}\mid T_{s'}(K^q)(J)$, and $w^{oq+2}\mid T_{s'}(K^q)(J)$ if $s<s'$;
(Tw4) if $w^o\mid T_sK(J)$ for all $J$ then $w^o\mid K(J)$ for all $J$.
\emph{Construction.} Let $M,M'$ be matrices over $W$ with $MM'=M'M=1$,
$M^{(s)}$ the image of $\sigma^{(s)}(M)$ in $O$, and suppose
$\zeta_jM^{(\mathrm{pos}\,j)}=y_j$ when $\mathrm{step}_j$ is a wedge and
$M'^{(\mathrm{pos}\,j)}\zeta_j=y_j$ when it is a contraction. Then
$T_sK=K\Lambda(M^{(s)})$ is a twisting structure. (Tw1) is ($\Lambda$1),
resp.\ ($\Lambda$2) with $M^{(s)}M'^{(s)}=1$; (Tw2) is ($\Lambda$4); (Tw3) is
(F3) in the domain $A$ with $\varpi=w$, $L=\Lambda(M^{(s)})$,
$L'=\Lambda(M'^{(s)})$, $L_2=\Lambda(M^{(s')})$, $L'_2=\Lambda(M'^{(s')})$
(inverse pairs by ($\Lambda$3)), the congruence coming from ($\Lambda$5) and
(F1) applied to the entries of $\Lambda(M')$; (Tw4): $K=(K\Lambda(M^{(s)}))\Lambda(M'^{(s)})$.

\emph{(F5) Column conditions.} Since $\mu$ is a Pieri path in $k$ coordinates,
$N+r\le K$: (R) $h_j<h_{j+1}$ implies $j$ is a vector and $h_{j+1}=h_j+1$; (Fa)
$h_{j+1}<h_j$ implies $j$ is a covector and $h_j=h_{j+1}+1$; (S) $c\ge1$
implies $h_0=0$; (D) $c<1$ implies $h_j>r$ for all $j\le N$. Indeed if
$\nu\subseteq\mu'$ is a horizontal strip then
$\mathrm{colht}(\nu,c)\le\mathrm{colht}(\mu',c)\le\mathrm{colht}(\nu,c)+1$
($a\mapsto a-1$ injects the indices $a\ge1$ with $c\le\mu'_a$ into those with
$c\le\nu_{a-1}$); a vector step is a strip $\mu_j\subseteq\mu_{j+1}$, a
covector step the reverse. As $\mu_0=0$, $h_0=0$ for $c\ge1$ and $h_0=k$ for
$c\le0$, and then $h_j\ge k-N\ge2r>r$.

\emph{(F6) Bookkeeping.} Put $y_j(b)=w^{\delta(j,b)}X_{j,b}$ and
$o_j(s)=\sum_{i<j\text{ costly}}p^{s-\mathrm{pos}(i)}$, so $o_N(z)=\mathrm{oc}$. Let
$\mathrm{Nm}(h,J)$ mean: if $c\ge1$, $|J|=h$, $\{0..r-1\}\subseteq J$ when $r\le h$,
and $J\subseteq\{0..r-1\}$ when $h\le r$; if $c<1$, $J\cap\{0..r-1\}=\emptyset$.
For a used letter $j$, a subset $J$ and an index $m$, write $J'=J\setminus m$
($m\in J$, wedge step) or $J'=J\cup m$ ($m\notin J$, contraction step).
(Nom) If ($\delta(j,m)=1$ iff $j$ is costly) and $\mathrm{Nm}(h_j,J')$, then
$\mathrm{Nm}(h_{j+1},J)$. (Van) If $j$ is costly and $\delta(j,m)=0$, then
$\neg\mathrm{Nm}(h_j,J')$. Both are case checks using (R), (Fa), (D) and the
definition of costly. For instance, a wedge with $c\ge1$ is a vector rise,
costly iff $r\le h_j$, so $\delta(j,m)=1$ iff $m\ge r$; in (Nom), if $h_j<r$
then $J'$ and $m$ lie in $\{0..r-1\}$, and if $r\le h_{j+1}$ then either
$\{0..r-1\}\subseteq J'$ ($r\le h_j$) or $J$ is an $r$-subset of
$\{0..r-1\}$ ($h_j=r-1$); in (Van), $m<r\le h_j$ and $\mathrm{Nm}$ would force
$m\in J'$. A wedge with $c<1$ is a covector fall with $h_j>r$ (D), not costly,
so $m\ge r$ and $J$ avoids $\{0..r-1\}$. A contraction with $c<1$ is a vector
rise, costly; (Nom) is clear from $J\subseteq J'$, and in (Van) $m<r$
contradicts $\mathrm{Nm}$. A contraction with $c\ge1$ is a covector fall,
costly iff $h_j\le r$, so $\delta(j,m)=1$ iff $m<r$; (Nom) follows from
$J\subseteq J'$ and $|J|=h_{j+1}$ (if $h_{j+1}\le r<h_j$ then $J=\{0..r-1\}$),
and in (Van) $m\ge r$ contradicts $J'\subseteq\{0..r-1\}$. (Rec) Let
$Z^\zeta_j(s)$ be the canonical translate column built with the letters
$\zeta_j$ in place of $x_j$. If $\mathrm{pos}(i)\le s\le s'$ for all used
$i<j$, then $Z^\zeta_j(s')=Z^\zeta_j(s)^{p^{s'-s}}$ entrywise and
$o_j(s')=p^{s'-s}o_j(s)$ (induction on $j$; costly letters are used).

\emph{(F7) Order claim.} Given a twisting structure, for $j\le N$, every $s$
with $\mathrm{pos}(i)\le s$ for used $i<j$, and every $J$: (i)
$w^{o_j(s)}\mid T_s(Z^\zeta_j(s))(J)$; (ii) if $\neg\mathrm{Nm}(h_j,J)$ then
$w^{o_j(s)+1}\mid T_s(Z^\zeta_j(s))(J)$. Induction on $j$. For $j=0$:
$o_0=0$, $\mathrm{Nm}(h_0,\emptyset)$ holds (by (S) if $c\ge1$), so (ii) concerns
$J\ne\emptyset$ and (Tw2) applies. An unused $j$ changes nothing. Let $j$ be
used, $s_0=\mathrm{pos}(j)\le s$, so $\mathrm{pos}(i)\le s_0$ for $i<j$. By (Tw1)
$V=T_{s_0}(\mathrm{step}_j(\zeta_j,Z^\zeta_j(s_0)))=\mathrm{step}_j(y_j,T_{s_0}Z^\zeta_j(s_0))$,
whose $J$-entry is a signed sum of terms $y_j(m)\,T_{s_0}Z^\zeta_j(s_0)(J')$.
\emph{Claim:} $w^{o_{j+1}(s_0)}\mid V(J)$, and $w^{o_{j+1}(s_0)+1}\mid V(J)$ if
$\neg\mathrm{Nm}(h_{j+1},J)$. If $j$ is costly, the hypothesis on positions gives
$\mathrm{pos}(i)\le s_0-1$ for used $i<j$; the induction hypothesis at $s_0-1$,
(Tw3) with $q=p$ and (Rec) give $w^{o_j(s_0)}\mid T_{s_0}Z^\zeta_j(s_0)(J')$
always and $w^{o_j(s_0)+2}$ when $\neg\mathrm{Nm}(h_j,J')$. A term with
$\delta(j,m)=1$ carries an extra $w$; a term with $\delta(j,m)=0$ has
$\neg\mathrm{Nm}(h_j,J')$ by (Van). So all terms are divisible by
$w^{o_j(s_0)+1}=w^{o_{j+1}(s_0)}$; if $\neg\mathrm{Nm}(h_{j+1},J)$, a term with
$\delta(j,m)=1$ has $\neg\mathrm{Nm}(h_j,J')$ by (Nom), hence is divisible by
$w^{o_j(s_0)+3}$, and all terms by $w^{o_{j+1}(s_0)+1}$. If $j$ is not
costly, the induction hypothesis at $s_0$ gives $w^{o_j(s_0)}$ for all terms,
and if $\neg\mathrm{Nm}(h_{j+1},J)$, terms with $\delta(j,m)=1$ carry an extra
$w$ and terms with $\delta(j,m)=0$ have $\neg\mathrm{Nm}(h_j,J')$ by (Nom). Finally
$Z^\zeta_{j+1}(s)=\mathrm{step}_j(\zeta_j,Z^\zeta_j(s_0))^{p^{s-s_0}}$, and (Tw3)
from $s_0$ to $s$ with $o_{j+1}(s)=p^{s-s_0}o_{j+1}(s_0)$ gives (i), (ii).

\emph{(F8) Substitution.} For $g'\in GL_k(W)$ and $s=\mathrm{pos}(l)$ put
$\zeta_l=y_l(g'^{-1}g)^{(s)}$ (vector $l$) and $\zeta_l=(g^{-1}g')^{(s)}y_l$
(covector $l$). Then $L_{g'}(x_l(b))=\zeta_l(b)$, by expanding the definitions
and $(g^{(s)})^{-1}=(g^{-1})^{(s)}$. As $L_{g'}$ is a ring map,
$L_{g'}(Z_N(z)(J))=Z^\zeta_N(z)(J)$.

\emph{Proof of (3).} Let $M=g^{-1}g'$, $M'=g'^{-1}g$. Then $\zeta_lM^{(s)}=y_l$
for a vector and $M'^{(s)}\zeta_l=y_l$ for a covector. If $c\ge1$, wedge steps
are at vectors and contraction steps at covectors ((R), (Fa)), and the
construction of (F4) applies to $(M,M')$; if $c<1$ the roles are exchanged and
it applies to $(M'^T,M^T)$. The order claim (i) at $j=N$, $s=z$ and (Tw4) give
$w^{\mathrm{oc}}\mid Z^\zeta_N(z)(J)=L_{g'}(Z_N(z)(J))$ in $O[X]$, i.e.\
$Z_N(z)(J)\in W^{\mathrm{oc}}_{g'}$.

\emph{(F9) Nominal column.} For $j<N$ let $n_j(b)=X_{j,b}$ if ($\delta(j,b)=1$
iff $j$ is costly) and $n_j(b)=0$ otherwise; $S_0=\mathbf 1$,
$S_{j+1}=\mathrm{step}_j(n_j,S_j)$. (Supp) If $S_j(J)\ne0$ then
$\mathrm{Nm}(h_j,J)$, by induction with (Nom). (Vanish) If $j$ is costly and
$\delta(j,m)=0$ then $S_j(J')=0$, by (Supp) and (Van). (Split) $S_N(J)=\mathrm{pt}_c(J)$.
More precisely, with $P_j,Q_j$ the stages of the column words of $\nu,\kappa$
(embedded by $\rho^{pl},\rho^{tan}$), $S_j(J)=P_j(J_p)Q_j(J_t)$. By the first
paragraph of the proof of Lemma~\ref{i1:lem-A18}: for $K(J)=P(J_p)Q(J_t)$,
(a) if $x$ vanishes on plane indices, $(K\wedge x)(J)=P(J_p)(Q\wedge x^{tan})(J_t)$
and likewise for $\lrcorner$; (b) if $x$ vanishes on tangent indices and
$Q(J'')=0$ for $J''\ne\emptyset$, $(K\wedge x)(J)=(P\wedge x^{pl})(J_p)Q(J_t)$ and
likewise for $\lrcorner$. The plane and tangent heights are $\min(h_j,r)$ and
$(h_j-r)^+$, and for $c\ge1$, $Q_j$ vanishes off $\emptyset$ when $h_j\le r$.
Going through the six cases (rise or fall; $c\ge1$ with $h_j$ below or above
$r$, or $c<1$, where $h_j>r$ by (D)), in each case exactly one of the plane
and tangent words moves, $n_j$ is $X_{j,\cdot}$ exactly on its indices
($\delta=1$ iff costly), and (a) or (b) applies.

\emph{(F10) Reduction.} Let $\pi_s:\Z[X]\to k_O[X]$, $X_{i,b}\mapsto X_{i,b}^{p^{s-\mathrm{pos}(i)}}$;
if the variables of $F$ have $\mathrm{pos}(i)\le s\le s'$ then
$\pi_s(F)^{p^{s'-s}}=\pi_{s'}(F)$ (Frobenius of $k_O[X]$, which has
characteristic $p$ as $p=w^2$). Let $Z^y_j(s)$ be built with the letters
$y_j$. \emph{Claim:} if $\mathrm{pos}(i)\le s$ for used $i<j$, then
$Z^y_j(s)(J)=w^{o_j(s)}Y$ with $Y\in O[X]$ and $\bar Y=\pi_s(S_j(J))$. Induction
on $j$; unused steps change nothing. Let $j$ be used, $s_0=\mathrm{pos}(j)$. If
$j$ is costly, by induction at $s_0-1$ and (Rec),
$Z^y_j(s_0)(J')=w^{o_j(s_0)}Y_{J'}^p$ with
$\overline{Y_{J'}^p}=\pi_{s_0}(S_j(J'))$; each term $y_j(m)Y_{J'}^p$ is $wu$
with $\bar u=\pi_{s_0}(n_j(m))\pi_{s_0}(S_j(J'))$: if $\delta(j,m)=1$ take
$u=X_{j,m}Y^p_{J'}$; if $\delta(j,m)=0$, (Vanish) gives $\bar Y_{J'}=0$, so
$Y_{J'}=wT$ and $u=X_{j,m}w^{p-1}T^p$, $\bar u=0=\pi_{s_0}(n_j(m))\cdots$
(as $n_j(m)=0$). Summing, $\mathrm{step}_j(y_j,Z^y_j(s_0))(J)=w^{o_j(s_0)+1}Y'$
with $\bar Y'=\pi_{s_0}(S_{j+1}(J))$. If $j$ is not costly, the reduction of
$y_j$ is $\pi_{s_0}n_j$ ($y_j(b)=wX_{j,b}$ exactly when $n_j(b)=0$), and the
claim at $s_0$ follows directly. Raising to the power $p^{s-s_0}$ and using
(Rec) and the Frobenius remark gives the claim at $s$.

\emph{Proof of (4).} By (F8) with $g'=g$, $\zeta_l=y_l$, so
$L_g(Z_N(z)(J))=Z^y_N(z)(J)=w^{\mathrm{oc}}Y$ with
$\bar Y=\pi_z(S_N(J))=\pi_z(\mathrm{pt}_c(J))$ by (F10) and (F9). With
$Q=\mathrm{Pow}_z(\mathrm{pt}_c(J))\in O[X]$, $\bar Q=\pi_z(\mathrm{pt}_c(J))$, so
$Y-Q=wT$ with $T$ integral, and
$L_g(Z_N(z)(J))-w^{\mathrm{oc}}\mathrm{Pow}_z(\mathrm{pt}_c(J))=w^{\mathrm{oc}+1}T$.
$\square$


\subsubsection{Proof of Proposition~\ref{i1:prop-Fm}}\label{i1:ssec-pf-Fm}

The form will be a rescaled Fock pairing. Parts (a), (b), (c), (e) follow
from three properties of the Fock pairing (adjointness, pairing with the
extremal vector, coefficients as words) and from the triangularity of the
standard basis. Part (d) needs a \emph{sandwich}
$p^{N}\nabla_O(\nu)\subseteq\Delta_O(\nu)\subseteq\nabla_O(\nu)$ between the
dual Weyl module and the Weyl module, which is the main work: it is proved
over $\Z_p$ by an extremal-projector argument for $\mathfrak{gl}_k$, then
descended to $\Z$ and lifted to $O$.

\emph{Notation.} Throughout $k\ge1$, $n\le k$, and $\nu\in\N^n$ is
antitone; rows are indexed by $a,b,c<k$, columns by $s,t,j,l<n$. For a
commutative ring $R$ put $R[X]=R[X_{a,t}:a<k,\ t<n]$; for an exponent $m$ let
$\deg_a(m)=\sum_tm_{a,t}$ (its \emph{row degrees}), and let $\bar\nu\in\N^k$ be
$\nu$ followed by zeros. $\nabla_R(\nu)\subseteq R[X]$ is the module of
polynomials homogeneous of degree $\nu_t$ in the variables of column $t$ for
every $t$ and fixed by the column substitutions $X_{a,j}\mapsto
X_{a,j}+YX_{a,l}$ ($l<j<n$, all $a$; $Y$ an indeterminate). For $n=k$ and a
domain $R$ this is Definition~\ref{i1:def-slots}(2); in general it is that
module for $\bar\nu$ ($k$ columns), since the variables of the columns $\ge n$
do not occur and the substitutions with $j\ge n$ fix such polynomials. As
$O\to K$ is injective, $\nabla_O(\nu)$ is identified with its image in $K[X]$,
as in the statement. The \emph{row action} of $g\in M_k(R)$ is the
$R$-algebra map $g\cdot X_{a,t}=\sum_bg_{ab}X_{b,t}$. For $a,b<k$ and $q\ge0$,
$E^{(q)}_{a,b}f$ is the coefficient of $T^q$ in $f(X_{b,t}\mapsto
X_{b,t}+TX_{a,t})$. For a word $w=((a_1,b_1,q_1),\dots,(a_z,b_z,q_z))$ put
$E_w=E^{(q_1)}_{a_1,b_1}\circ\dots\circ E^{(q_z)}_{a_z,b_z}$ and
$w^T=((b_z,a_z,q_z),\dots,(b_1,a_1,q_1))$, so $(w^T)^T=w$. Let $D_s$ be the
leading principal $s\times s$ minor of $(X_{a,t})$, $e_\nu=\prod_{s<n}D_{s+1}^{\nu_s-\nu_{s+1}}$
($\nu_n=0$; the \emph{extremal vector}), and $\Delta_R(\nu)$ the $R$-span of
the $E_we_\nu$ (the \emph{Weyl module}). $I$ is the $k\times n$ matrix with
entries $1$ at $(a,a)$ and $0$ elsewhere, and $m_\nu=\sum_{s<n}\nu_s\epsilon_{(s,s)}$.
The \emph{Fock pairing} on $R[X]$ is
$\langle f,g\rangle_R=\sum_mf_mg_m\,m!$, $m!=\prod_{a,t}m_{a,t}!$. The
operators $E^{(q)}_{a,b}$, the row action, $e_\nu$ and evaluation at $I$ are
defined over $\Z$ and commute with every ring map applied to coefficients.

\emph{(N) Facts on $\nabla$.} Let $R$ be a commutative ring.
(N1) $E^{(q)}_{a,b}\nabla_R(\nu)\subseteq\nabla_R(\nu)$: the substitution
$X_{b,\cdot}\mapsto X_{b,\cdot}+TX_{a,\cdot}$ replaces a variable by a sum of
variables of the same column, so it preserves column degrees coefficientwise
in $T$, and it commutes with every column substitution (both composites send
$X_{c,s}$ to $X_{c,s}+[c=b]TX_{a,s}+[s=j]Y(X_{c,l}+[c=b]TX_{a,l})$); compare
coefficients of $T^q$.
(N2) $e_\nu\in\nabla_R(\nu)$: $D_s$ has degree $1$ in the columns $t<s$ and
$0$ in the others, and a column substitution either does not involve $D_s$
($j\ge s$) or adds a multiple of one of its columns to another; products
multiply column degrees additively and the substitutions are ring maps, and
$\sum_{t\le s<n}(\nu_s-\nu_{s+1})=\nu_t$.
(N3) $\Delta_R(\nu)\subseteq\nabla_R(\nu)$, by (N1) and (N2), induction on
the length of the word.
(N4) A ring map $\varphi:R\to R'$ maps $\nabla_R(\nu)$ into $\nabla_{R'}(\nu)$
(supports only shrink and $\varphi$ commutes with the substitutions); for
$\Z\to\Q$ an integral polynomial lies in $\nabla_\Z(\nu)$ iff its image lies in
$\nabla_\Q(\nu)$ (same support; injectivity of $\Z[X][Y]\to\Q[X][Y]$).
(N5) $E^{(q)}_{a,b}e_\nu=0$ for $a<b$, $q\ge1$: the substitution
$X_{b,\cdot}\mapsto X_{b,\cdot}+TX_{a,\cdot}$ fixes every $D_{s+1}$ (it adds a
multiple of the row $a<b$ inside the minor, or does not touch it).
(N6) \emph{Diagonal evaluation.} If $g$ is column homogeneous of degree
$\nu$ and $D$ is a $k\times n$ matrix with $D_{a,t}=0$ for $a\ne t$, then
$g(D)=g_{m_\nu}\prod_sD_{s,s}^{\nu_s}$ (a monomial not supported on the
$X_{s,s}$ vanishes at $D$; one supported there is $m_\nu$ by column
homogeneity). As every $D_{s+1}(I)=1$, $e_\nu(I)=(e_\nu)_{m_\nu}=1$; $e_\nu$ has
row degrees $\bar\nu$ ($D_{s+1}$ has degree $1$ in each row $a\le s$).
(N7) \emph{Highest weight line.} Let $L$ be an infinite field and $f\in
\nabla_L(\nu)$ with $E^{(q)}_{a,b}f=0$ for all $a<b$, $q\ge1$. Then $f=f(I)e_\nu$.
Put $g=f-f(I)e_\nu$. It lies in $\nabla_L(\nu)$; it is invariant under the
substitutions $X_{b,\cdot}\mapsto X_{b,\cdot}+tX_{a,\cdot}$ ($a<b$, $t\in L$;
hypothesis and (N5)); and $g(D)=0$ for every $D$ as in (N6), since
$g_{m_\nu}=f(I)-f(I)(e_\nu)_{m_\nu}=0$. Let $x\in L^{k\times n}$ with
$Q(x)\ne0$, $Q=\prod_{s<n}D_{s+1}$. Row operations (adding a multiple of row
$a$ to a row $b>a$) and column operations (a multiple of column $l$ to a column
$j>l$) change neither $g(x)$ nor any leading minor; by elimination
(the pivot $x_{s,s}$ at step $s$ is the ratio of two consecutive leading
minors, so nonzero) $x$ is moved to a matrix vanishing off the diagonal, so
$g(x)=0$. Thus $gQ$ vanishes on $L^{k\times n}$, so $gQ=0$ ($L$ infinite), and
$g=0$ as $Q(I)=1$.

\emph{(B) The standard basis.} Order exponents column-lexicographically:
positions $(a,t)$ are ordered by $t$, then by $a$, and $m\succ m'$ if $m>m'$
at the first position where they differ; this is a monomial order and a
well-order; $\mathrm{lead}(f)$ is the largest exponent of $f\ne0$. By
Theorem~\ref{i1:thm-stdbasis} (for $\bar\nu$; as $\bar\nu_t=0$ for $t\ge n$,
all minors have size $s_j\le n$), for every domain $R$ the bideterminants
$b_\mathsf T=\prod_j\det(X_{I_j(\alpha),\beta})_{\alpha,\beta<s_j}$, $\mathsf T$
running over the semistandard tableaux of shape $\nu$ with entries $<k$
(columns $I_j$, strictly increasing), form an $R$-basis of $\nabla_R(\nu)$;
they have integer coefficients and do not depend on $R$. Put
$e(\mathsf T)=\sum_j\sum_{\alpha<s_j}\epsilon_{(I_j(\alpha),\alpha)}$.
(B1) $\mathrm{lead}(b_\mathsf T)=e(\mathsf T)$ with coefficient $1$;
$e(\mathsf T)$ is lower triangular ($e_{a,t}=0$ for $a<t$); and
$\mathsf T\mapsto e(\mathsf T)$ is injective. For one minor with strictly
increasing $I$, the term of a permutation $\sigma\ne1$ differs from the
diagonal term first in the first column $\beta$ with $\sigma\beta\ne\beta$,
where $\sigma\beta>\beta$, so its entry sits in the later row
$I(\sigma\beta)>I(\beta)$; so the diagonal term leads, with coefficient $1$,
and leading exponents add under products. $I_j(\alpha)\ge\alpha$ gives lower
triangularity. Finally $e(\mathsf T)_{a,\alpha}$ is the number of entries $a$ in
row $\alpha$ of $\mathsf T$, and a row of a semistandard tableau is weakly
increasing, so it is determined by this multiset.
(B2) \emph{Row-degree bounds.} Let $R$ be a domain, $f\in\nabla_R(\nu)$ and
$m$ an exponent of $f$. Then $\sum_{c\le a}\deg_c(m)\le\sum_{c\le a}\bar\nu_c$
for every $a$; $\deg_a(m)\le\nu_0$ if $n>0$; and if $n=k$,
$\deg_a(m)\ge\nu_{n-1}$. The exponents of $f$ are exponents of basis elements
$b_\mathsf T$; a monomial of the $j$-th minor uses the $s_j$ distinct rows
$I_j(\alpha)$ once each. So it has at most $\min(s_j,a+1)$ variables in rows
$\le a$, and $\sum_j\min(s_j,a+1)=\sum_{c\le a}\#\{j:s_j>c\}=\sum_{c\le a}\bar\nu_c$;
row $a$ occurs at most $\#\{j:s_j>0\}=\nu_0$ times; and if $s_j=n=k$, $I_j$ is
a bijection, so row $a$ occurs at least $\#\{j:s_j=n\}=\nu_{n-1}$ times.
(B3) For every domain $R$, $\nabla_R(\nu)$ is spanned by bideterminants of
shape $\nu$, each of which lies in $\nabla_\Z(\nu)$ (it is the image of an
integral basis element).

\emph{(P1) The pairing.} $\langle\cdot,\cdot\rangle_R$ is $R$-bilinear and
symmetric, and $\langle\varphi f,\varphi g\rangle=\varphi\langle f,g\rangle$ for
a ring map $\varphi$ (the support of $\varphi f$ lies in that of $f$). It is
diagonal in the monomials; in particular $\langle f,g\rangle=\langle
f,g'\rangle$ if $g'$ is the part of $g$ on the exponents of $f$.

\emph{(P2) Row action.} $\langle g\cdot P,Q\rangle_R=\langle P,g^T\cdot
Q\rangle_R$ for $g\in M_k(R)$. First $\langle fX_v,h\rangle=\langle
f,\partial_vh\rangle$ (compare coefficients: $(m+\epsilon_v)!=(m_v+1)m!$).
Second, $g\cdot h$ and $h$ have the same constant term. Third, the chain rule
$\partial_{(b,t)}(g\cdot h)=\sum_ag_{ab}\,g\cdot(\partial_{(a,t)}h)$ (induction
on $h$, Leibniz rule). Now induct on $P$: constants and sums are clear; for
$P=P'X_{(c,t)}$, $g\cdot P=\sum_bg_{cb}(g\cdot P')X_{(b,t)}$, so by the first
remark and induction $\langle g\cdot P,Q\rangle=\sum_bg_{cb}\langle P',g^T\cdot
\partial_{(b,t)}Q\rangle$, while $\langle P'X_{(c,t)},g^T\cdot Q\rangle=\langle
P',\partial_{(c,t)}(g^T\cdot Q)\rangle$, which is the same sum by the chain rule
for $g^T$.

\emph{(P3) Words.} $\langle E_wf,g\rangle_R=\langle f,E_{w^T}g\rangle_R$.
Apply (P2) over $R[T]$ to $g_T=1+Te_{ba}$ ($e_{ba}$ the matrix unit). Then
$g_T\cdot f=\sum_qT^qE^{(q)}_{a,b}f$ and $g_T^T\cdot h=\sum_qT^qE^{(q)}_{b,a}h$,
so by bilinearity and (P1)
$\sum_qT^q\langle E^{(q)}_{a,b}f,h\rangle_R=\sum_qT^q\langle f,E^{(q)}_{b,a}h\rangle_R$
in $R[T]$; compare coefficients. For words induct on the length, using
$E_{w^T}E^{(q)}_{b,a}=E_{((a,b,q)\cdot w)^T}$.

\emph{(P4) The norm.} $N_\nu=\langle e_\nu,e_\nu\rangle_\Z=\sum_m(e_\nu)_m^2m!$
is a positive integer ($e_\nu\ne0$ by (N6)), and $c:=\langle
e_\nu,e_\nu\rangle_K=N_\nu\cdot1_K\ne0$ as $K$ has characteristic $0$.

\emph{(P5) Pairing with the extremal vector.} For $F\in\nabla_K(\nu)$,
$\langle e_\nu,F\rangle_K=F(I)\,c$. Let $G$ be the part of $F$ of row degrees
$\bar\nu$. Column substitutions keep each variable in its row, so they
commute with taking this part, and $G\in\nabla_K(\nu)$. For $a<b$ and $q\ge1$,
$E^{(q)}_{a,b}G\in\nabla_K(\nu)$ (N1) has row degrees $\bar\nu+q(\epsilon_a-\epsilon_b)$,
whose sum over the rows $\le a$ exceeds $\sum_{c\le a}\bar\nu_c$ by $q$; by (B2)
it is $0$. By (N7) ($K$ is infinite), $G=G(I)e_\nu$. By (P1) and (N6),
$\langle e_\nu,F\rangle=\langle e_\nu,G\rangle=G(I)c$, and
$G(I)=G_{m_\nu}=F_{m_\nu}=F(I)$ by (N6), as $m_\nu$ has row degrees $\bar\nu$.

\emph{(P6) Coefficients as words.} Let $R$ be a commutative ring, $F\in R[X]$
column homogeneous of degree $\nu$, and $e$ a lower triangular exponent with
column sums $\nu$. List the pairs $(a,s)$ with $s<a$ as $l_1,\dots,l_M$, with
$s$ decreasing and, for fixed $s$, $a$ increasing, and let $w(e)$ be the word
whose $i$-th letter is $(s,a,e_{a,s})$ for $l_i=(a,s)$. Then
$F_e=(E_{w(e)}F)(I)$. With indeterminates $Y_l$ and the $R[Y]$-algebra maps
$\tau_l:X_{a,\cdot}\mapsto X_{a,\cdot}+Y_lX_{s,\cdot}$ ($l=(a,s)$), put
$\Phi=\tau_{l_1}\circ\dots\circ\tau_{l_M}$. (a) Since $\tau_lP=\sum_qY_l^qE^{(q)}_{s,a}P$
for $P\in R[X]$, the $\tau_l$ are $R[Y]$-linear and the $Y_l$ are distinct,
the coefficient of $Y^d$ in $\Phi(F)$ is $E_{w_d}F$ ($w_d$ the word with
exponents $d$). (b) $\Phi(F)|_{X=I}=F(\ell)$ with $\ell_{a,t}=\delta_{at}+[t<a]Y_{(a,t)}$:
by induction over initial segments of the list, since a pair $(a',s')$
appended after pairs $(a'',s'')$ with $s''\ge s'$ meets no earlier pair with
row $s'$, the new term is $Y_{(a',s')}$ at column $s'$ only. (c) In $F(\ell)$
only lower triangular exponents $m$ survive, contributing $F_mY^{m^{\rm off}}$
($m^{\rm off}$ the off-diagonal part), and $m^{\rm off}=e^{\rm off}$ forces
$m=e$ since the diagonal entries are fixed by the column sums. Compare the
coefficients of $Y^{e^{\rm off}}$.

\emph{(P7) Lower triangular coefficients.} Let $Q\in\nabla_K(\nu)$. (a) If
$Q_e\in O$ for every lower triangular $e$, then $Q\in\nabla_O(\nu)$. (b) If
$Q_e=0$ for every lower triangular $e$, then $Q=0$. Write
$Q=\sum_\mathsf Tc_\mathsf Tb_\mathsf T$ (B). For (a), suppose some
$c_\mathsf T\notin O$ and choose such a $\mathsf T_0$ with $e(\mathsf T_0)$
maximal. By (B1), $(b_\mathsf T)_{e(\mathsf T_0)}=0$ unless $e(\mathsf
T)\succeq e(\mathsf T_0)$, equals $1$ for $\mathsf T=\mathsf T_0$, and is an
integer always; and $c_\mathsf T\in O$ when $e(\mathsf T)\succ e(\mathsf
T_0)$. So $Q_{e(\mathsf T_0)}\in c_{\mathsf T_0}+O$, while $Q_{e(\mathsf
T_0)}\in O$ ($e(\mathsf T_0)$ is lower triangular): a contradiction. So all
$c_\mathsf T\in O$ and $Q\in\nabla_O(\nu)$. (b) is the same argument with
``$\notin O$'' replaced by ``$\ne0$''.

\emph{(P8) The sandwich.} Let $N'=N_0(k)\delta(p,k,l)=\binom k2\lfloor
\log_p(l+k-1)\rfloor$ with $l$ as in (d). Then $\Delta_O(\nu)\subseteq
\nabla_O(\nu)$ (N3) and $p^{N'}\nabla_O(\nu)\subseteq\Delta_O(\nu)$. Every
integer $c'$ prime to $p$ is a unit of $O$ (otherwise $c'$ and $p$ lie in the
maximal ideal, but $up+vc'=1$ for some $u,v\in\Z$), so the second inclusion is
the case $A=O$ of the following statement, proved below in steps (S0)--(S4):
\begin{quote}
(Sw) \emph{If $A$ is a domain in which every integer prime to $p$ is a unit,
and $l$ is as in (d), then $p^{N'}\nabla_A(\nu)\subseteq\Delta_A(\nu)$.}
\end{quote}

\emph{(M1)--(M6): the form, given (P8).} (M1) For $Q\in\nabla_K(\nu)$ and a
word $w$: $\langle E_we_\nu,Q\rangle_K=(E_{w^T}Q)(I)\,c$, by (P3), (N1) and
(P5). (M2) For $Q\in\nabla_K(\nu)$ and $e$ lower triangular, either $Q_e=0$
or $\langle E_{w(e)^T}e_\nu,Q\rangle_K=Q_e\,c$: if the column sums of $e$ are
not $\nu$, $Q_e=0$; otherwise apply (M1) to $w(e)^T$ and (P6). (M3) If
$P\in\Delta_O(\nu)$ and $Q\in\nabla_O(\nu)$, then $\langle P,Q\rangle_K\in Oc$:
by bilinearity it suffices to take $P=E_we_\nu$, and by (M1) the pairing is
$c$ times $(E_{w^T}Q)(I)$, the image of an element of $O$ (N4). (M4) If
$Q\in\nabla_K(\nu)$ and $\langle P,Q\rangle_K\in Oc$ for all $P\in\Delta_O(\nu)$,
then $Q\in\nabla_O(\nu)$: by (M2) every lower triangular coefficient of $Q$ is
$0$ or equals $z$ with $zc=\langle E_{w(e)^T}e_\nu,Q\rangle_K\in Oc$; apply
(P7)(a). (M5) If $Q\in\nabla_K(\nu)$ and $\langle P,Q\rangle_K=0$ for all
$P\in\nabla_K(\nu)$, then $Q=0$: as in (M4), using $E_{w(e)^T}e_\nu\in
\nabla_K(\nu)$ (N3) and (P7)(b). (M6) Put $B(P,Q)=c^{-1}\langle P,Q\rangle_K$.
Part (a) is (P2) over $K$; (b) is (M5); (c): $\nabla_O(\nu)\supseteq
\Delta_O(\nu)$ (N3), so the hypothesis gives $\langle P,Q\rangle_K\in Oc$ for
$P\in\Delta_O(\nu)$, and (M4) applies; (d): for $P,Q\in\nabla_O(\nu)$,
$p^{N'}Q\in\Delta_O(\nu)$ by (P8), so $B(P,p^{N'}Q)=c^{-1}\langle
p^{N'}Q,P\rangle_K\in O$ by symmetry and (M3); (e): if $\nu=0$ then $e_\nu=1$
and $\nabla_O(0)$ consists of constants, so $\nabla_O(0)=\Delta_O(0)$ and (M3)
applies. This proves Proposition~\ref{i1:prop-Fm} once (Sw) is proved.

\medskip\noindent\emph{Proof of (Sw).} Write $N'$ as above.

\emph{(S0) Over $\Z$.} (a) If $g\in\Z[X]$, $0\ne c'\in\Z$ and $c'g\in
\nabla_\Z(\nu)$, then $g\in\nabla_\Z(\nu)$ (same support; $\Z[X][Y]$ is
torsion free). (b) $\Delta_\Q(\nu)=\nabla_\Q(\nu)$. Let
$M=\Delta_\Q(\nu)\subseteq\nabla_\Q(\nu)$ and $N''$ its orthogonal in
$\nabla_\Q(\nu)$ for the Fock pairing (positive definite over $\Q$).
$\nabla_\Q(\nu)$ is finite dimensional (exponents with column sums $\nu$), $M$
is stable under all $E^{(q)}_{a,b}$, hence so is $N''$ by (P3) and (N1). A
subspace stable under $E^{(1)}_{a,a}$ (which multiplies a monomial by its
degree in row $a$) contains the row-degree components of its elements
(Lagrange interpolation in the eigenvalues). If $N''\ne0$, choose $0\ne v\in
N''$ with all exponents of the same row degrees and with
$\sum_aa\deg_a$ minimal; for $a<b$, $q\ge1$, $E^{(q)}_{a,b}v\in N''$ has
smaller height, so it is $0$, and $v=v(I)e_\nu$ by (N7), $v(I)\ne0$; then
$e_\nu\in M\cap N''$ and $\langle e_\nu,e_\nu\rangle=0$, a contradiction. So
$N''=0$ and $M=\nabla_\Q(\nu)$ (orthogonal decomposition with respect to the
positive definite form). (c) There is $0\ne N_1\in\Z$ with
$N_1\nabla_\Z(\nu)\subseteq\Delta_\Z(\nu)$: $\nabla_\Z(\nu)$ has a finite basis
(B); by (b) each basis element is a $\Q$-combination of the $E_we_\nu$, and
clearing denominators (operators commute with $\Z\to\Q$, and $\Z[X]\to\Q[X]$ is
injective) gives $c_\mathsf T\ne0$ with $c_\mathsf Tb_\mathsf T\in
\Delta_\Z(\nu)$; take $N_1=\prod c_\mathsf T$.

\emph{(S1) Spread.} For every domain $R$ and $l$ as in (d), any two exponents
$m,m'$ of elements of $\nabla_R(\nu)$ satisfy $\deg_a(m)\le\deg_a(m')+l$ for
all $a$. In case $\nu_t\le l$ for all $t$: $\deg_a(m)\le\nu_0\le l$ (B2). In
case $n=k$ and $\nu_s\le\nu_t+l$ for all $s,t$: $\deg_a(m)\le\nu_0\le
\nu_{n-1}+l\le\deg_a(m')+l$ (B2).

\emph{(S2) Over $\Z_p$: $p^{N'}\nabla_{\Z_p}(\nu)\subseteq\Delta_{\Z_p}(\nu)$.}
On $V=\Q_p[X]$ let $E_{ab}=\sum_{t<n}X_{a,t}\partial/\partial X_{b,t}$. These
satisfy $[E_{ab},E_{cd}]=\delta_{bc}E_{ad}-\delta_{da}E_{cb}$, $E_{cc}$ multiplies a
monomial by its degree in row $c$, and every polynomial lies in the finite
dimensional stable space of its column degrees; so $V$ is a
$\mathfrak{gl}_k$-representation (Definition~\ref{i1:def-gln}) whose weight
spaces are the row-degree spaces; let $\pi_\mu$ be the projection onto $V_\mu$.
Moreover (RG) $E^{(q)}_{a,b}=E_{ab}^q/q!$ for $a\ne b$ (both are
$\Q_p$-linear and satisfy the same Leibniz rule
$\Theta_q(fg)=\sum_{i+j=q}\Theta_if\,\Theta_jg$ with the same values on constants and
variables: $X_{c,t}\mapsto X_{c,t}$, $[c=b]X_{a,t}$, $0$ for $q=0,1,\ge2$),
and $E^{(1)}_{a,a}=E_{aa}$. A stable subspace contains the weight components
of its elements (as above). Write $\preceq$, $\mathrm{Pos}$, $U$, $\langle
S\rangle=US$ as in Definition~\ref{i1:def-gln}, and
$c(\mu)=\sum_a\mu_a^2+\sum_{a<b}(\mu_a-\mu_b)$, $C=\sum_{a,b}E_{ab}E_{ba}$.
From the proofs of Lemmas~\ref{i1:lem-hw} and \ref{i1:lem-cr} we use: (H1)
a nonzero finite-dimensional stable subspace contains a highest weight
vector; (H2) highest weight vectors have dominant weights; (H3) if $v$ is a
highest weight vector of weight $\lambda$, every weight of $Uv$ is
$\preceq\lambda$ and $\pi_\lambda Uv=\Q_pv$; (H4) $C$ commutes with all $E_{ab}$
and acts by $c(\lambda)$ on $Uv$; (H5) $c(\mu)<c(\lambda)$ if
$\lambda,\mu$ are dominant and $0\ne\lambda-\mu\in\mathrm{Pos}$; (H6) for a
finite-dimensional representation $W'$, $W'=\sum_\mu U(W'\cap H_\mu)$, $H_\mu$
the highest weight vectors of weight $\mu$ together with $0$ (Step~1 of the
proof of Lemma~\ref{i1:lem-cr}). From these: (H7) if $W'$ is finite
dimensional, $v\in W'$ a highest weight vector of weight $\lambda$ and $0\ne
x\in Uv$ a highest weight vector of weight $\mu$, then $\mu=\lambda$ and $x\in
\Q_pv$: $\mu\preceq\lambda$ by (H3), and if $\mu\ne\lambda$ then
$c(\lambda)x=Cx=c(\mu)x$ by (H4) (applied to $v$ and to $x$) with
$c(\mu)<c(\lambda)$ by (H2), (H5); so $\mu=\lambda$ and $x=\pi_\lambda x\in\Q_pv$.

Now let $W=\nabla_{\Q_p}(\nu)$: it is finite dimensional and stable ((N1),
(RG)), and by (S1) (over $\Q_p$) any two weights $\mu,\mu'$ of $W$ satisfy
$|\mu_c-\mu'_c|\le l$. Let $M=\Z_p[X]\subseteq V$: an additive group stable
under $\Z_p$, under every $\pi_\mu$ and under every $E_{ab}^q/q!=E^{(q)}_{a,b}$
($a\ne b$). (h1) $e_\nu\in W$ is a highest weight vector of weight $\bar\nu$
((N2), (N6), (N5)). (h2) Every highest weight vector $v\in W$ is a multiple of
$e_\nu$: by (RG) $E^{(q)}_{a,b}v=0$ for $a<b$, $q\ge1$, and (N7). (h3) Hence
$W=Ue_\nu$ by (H6). Let $\mathrm{Dv}$ be the smallest additive subgroup of
$V$ containing $\Z_pe_\nu$ and stable under $\Z_p$ and all $E_{ab}^q/q!$
($a\ne b$); it is contained in the image of $\Delta_{\Z_p}(\nu)$, which has
these stability properties ((RG): $E^{(q)}_{a,b}E_we_\nu=E_{(a,b,q)w}e_\nu$).
(h4) Every highest weight vector $h\in W\cap M$ lies in $\Z_pe_\nu$: $h=c'e_\nu$
by (h2) and $c'=h(I)\in\Z_p$ by (N6).

(WI) \emph{For every $x\in W\cap M$, $p^{N'}x\in\mathrm{Dv}$.} Granting (WI):
for $x\in\nabla_{\Z_p}(\nu)$ with image $X'\in W\cap M$ (N4), $p^{N'}X'$ is the
image of some $y\in\Delta_{\Z_p}(\nu)$, and $y=p^{N'}x$ since
$\Z_p[X]\to\Q_p[X]$ is injective. This is (S2). The proof of (WI) is given
after (S4).

\emph{(S3) Descent: for $f\in\nabla_\Z(\nu)$ there is $c'\in\Z$, $p\nmid c'$,
with $c'p^{N'}f\in\Delta_\Z(\nu)$.} Write $N_1=p^ic'$ with $p\nmid c'$ (S0)(c).
By (S2), $p^{N'}f=Y'$ in $\Z_p[X]$ with $Y'\in\Delta_{\Z_p}(\nu)$. Every element
of $\Delta_{\Z_p}(\nu)$ is $d+p^iz$ with $d\in\Delta_\Z(\nu)$ and $z\in\Z_p[X]$
(true for generators; preserved by sums and $\Z_p$-multiples, writing
$a'=a_0+p^ia''$ with $a_0\in\N$). So $g=p^{N'}f-d\in\Z[X]$ has all
coefficients divisible by $p^i$ in $\Z_p$, hence in $\Z$: $g=p^ig_1$ with
$g_1\in\Z[X]$; $g\in\nabla_\Z(\nu)$ (N3), so $g_1\in\nabla_\Z(\nu)$ (S0)(a), and
$c'p^{N'}f=c'd+N_1g_1\in\Delta_\Z(\nu)$.

\emph{(S4) Ascent.} By (B3), $\nabla_A(\nu)$ is spanned by bideterminants $b$
which are images of elements $b_\Z\in\nabla_\Z(\nu)$. By (S3), $c'p^{N'}b_\Z\in
\Delta_\Z(\nu)$ with $p\nmid c'$; mapping to $A[X]$ (words and $e_\nu$ commute
with base change) gives $c'p^{N'}b\in\Delta_A(\nu)$, and $c'$ is a unit of $A$.
This proves (Sw), given (WI).

\medskip\noindent\emph{Proof of (WI).} Write $L=(0,1,\dots,k-1)$ and
$N'=\binom k2\lfloor\log_p(l+k-1)\rfloor$.

\emph{The operators.} Let $Y$ be a vector space over a field of characteristic
$0$ with operators $E_{ab}$ ($a,b$ in an index set) satisfying the relations
$[E_{ab},E_{cd}]=\delta_{bc}E_{ad}-\delta_{da}E_{cb}$. For a finite set $S$ of
indices, $z\notin S$ and $q\ge0$ put $\Omega^{(0)}_{S,z}=1$,
$\Omega^{(q+1)}_{S,z}=\sum_{i\in S}E_{zi}\Omega^{(q)}_{S,z}E_{iz}$,
$c_{m,q}=(-1)^qm!/((q+m)!\,q!)$, $P_{S,z}=\sum_{q=0}^lc_{\#S,q}\Omega^{(q)}_{S,z}$,
and for a list $(z_1,\dots,z_j)$ of distinct indices
$\mathrm{Inv}(z_1,\dots,z_j)=P_{\{z_2,\dots,z_j\},z_1}\mathrm{Inv}(z_2,\dots,z_j)$,
$\mathrm{Inv}()=1$. Let $\mathcal R=\sum_{a\ne b}\mathrm{im}(E_{ab})$.

(J1) \emph{Naturality.} A linear map intertwining the $E_{ab}$ intertwines the
$\Omega$, the $P$ and $\mathrm{Inv}$ (induction).

(J2) \emph{Invariance.} Let $N\subseteq Y$ be stable, with (ev): if $x\in N$,
$E_{zz}x=jx$, $j\in\Z$, $|j|>l$, then $x=0$. If $u\in N$ and $E_{aa}u=0$ for the
indices $a$ of a list $\mathcal L$, then $\mathrm{Inv}(\mathcal L)u\in N$, it is
killed by all $E_{ab}$, $a,b\in\mathcal L$, and $\mathrm{Inv}(\mathcal L)u-u\in
\mathcal R$. Fix $S$, $z\notin S$, $m=\#S$, $\Omega^{(q)}=\Omega^{(q)}_{S,z}$.
(a) For $a,b\in S$, $E_{ab}$ commutes with $\Omega^{(q)}$ (the commutators with
$E_{zi}$ and $E_{iz}$ produce $-E_{zb}\Omega E_{az}+E_{zb}\Omega E_{az}$);
so do $E_{zz}$ and $E_{cc}$ for $c\notin S\cup\{z\}$; the $E_{iz}$ ($i\in S$)
commute with each other, and so do the $E_{zi}$. (b) For $j\in S$ put
$C_j=\sum_{i\in S}E_{iz}E_{ji}-E_{jz}E_{zz}+(m+1)E_{jz}$. From the relations,
$C_jE_{kz}=E_{kz}C_j+2E_{kz}E_{jz}$ ($k\in S$) and
$E_{jz}\Omega^{(q+1)}=\Omega^{(q+1)}E_{jz}+\sum_iE_{zi}[E_{jz},\Omega^{(q)}]E_{iz}+\Omega^{(q)}C_j$,
whence by induction
$E_{jz}\Omega^{(q+1)}=\Omega^{(q+1)}E_{jz}+(q+1)\Omega^{(q)}(C_j+qE_{jz})$.
(c) If $E_{ab}u=0$ ($a,b\in S$) and $E_{zz}u=0$, then $C_ju=(m+1)E_{jz}u$, so
$E_{jz}\Omega^{(q+1)}u=\Omega^{(q+1)}E_{jz}u+(q+1)(q+1+m)\Omega^{(q)}E_{jz}u$;
as $c_{m,0}=1$ and $c_{m,q}+(q+1)(q+1+m)c_{m,q+1}=0$, the sum telescopes:
$E_{jz}P_{S,z}u=c_{m,l}\Omega^{(l)}E_{jz}u$. (d) If $x\in N$, $E_{zz}x=-jx$
($j\ge0$) and $l<j+q$, then $\Omega^{(q)}x=0$ (induction on $q$; (ev) for
$q=0$; $E_{iz}$ lowers the $E_{zz}$-eigenvalue by $1$). (e) If $i\ne z$, $y\in
N$ and $E_{iz}y=E_{ii}y=E_{zz}y=0$, then $E_{zi}y=0$: $y_j=E_{zi}^jy$ has
$E_{zz}$-eigenvalue $j$ and $E_{iz}y_{j+1}=-j(j+1)y_j$; by (ev) $y_j=0$ for
$j>l$, and descending, $y_1=0$. (f) Let $u$ be as in (c) and $w'=P_{S,z}u\in
N$. Then $E_{ab}w'=0$ for $a,b\in S\cup\{z\}$: for $a,b\in S$ and $a=b=z$ by
(a); for $E_{jz}$ by (c) and (d) ($E_{jz}u$ has eigenvalue $-1$ and $l<1+l$);
for $E_{zb}$ by (e). And $w'-u=\sum_{q\ge1}c_{m,q}\sum_iE_{zi}(\cdots)\in
\mathcal R$. (g) Induct on the list: for $\mathcal L=(z,\mathcal L')$,
$v=\mathrm{Inv}(\mathcal L')u$ satisfies the claims for $\mathcal L'$ and
$E_{zz}v=0$ (by (a), $E_{zz}$ commutes with $\mathrm{Inv}(\mathcal L')$); apply
(f) with $S=\mathcal L'$.

(J3) \emph{Integrality.} Over $\Q_p$, let $M'\subseteq Y$ be an additive
subgroup stable under $\Z_p$ and under all $E_{ab}^q/q!$ ($a\ne b$), and
$\mathcal L$ a list of $j$ distinct indices. Then $p^{\binom j2\lfloor\log_p(l+j-1)\rfloor}
\mathrm{Inv}(\mathcal L)M'\subseteq M'$. (a) The maps $\ell_i(X)=E_{zi}XE_{iz}$
commute ((J2)(a)), so
$\Omega^{(q)}_{S\cup\{i\},z}=\sum_{j'\le q}\binom q{j'}E_{zi}^{j'}\Omega^{(q-j')}_{S,z}E_{iz}^{j'}$.
(b) $v_p\binom{n_1}{j'}\le\lfloor\log_pn_1\rfloor$ for $0<n_1$, $j'\le n_1$
(Legendre's formula, as in the proof of Lemma~\ref{i1:lem-A7}). (c) For $z\notin S$, $q+c''\le N''$ and $u\in
M'$: $p^{\#S\lfloor\log_pN''\rfloor}\frac{c''!}{(q+c'')!\,q!}\Omega^{(q)}_{S,z}u\in
M'$. By induction on $S$ with (a): for the $j'$-th term put
$u'=E_{iz}^{j'}u/j'!\in M'$, apply the induction hypothesis with $(c''+j',q-j')$,
apply $E_{zi}^{j'}/j'!$, and multiply by $p^{\lfloor\log_pN''\rfloor}/\binom{c''+j'}{j'}\in
\Z_p$ (by (b)); the product of the coefficients is exactly
$\frac{c''!}{(q+c'')!q!}\binom q{j'}$. (d) Hence
$p^{\#S\lfloor\log_p(l+\#S)\rfloor}P_{S,z}M'\subseteq M'$ ($c''=\#S$,
$N''=l+\#S$). (e) Induct on $\mathcal L$, using $\binom j2=\binom{j-1}2+(j-1)$.

(J4) \emph{Hom spaces.} For spaces $V_1,V_2$ with operators satisfying the
relations, $E^{\rm h}_{ab}\varphi=E_{ab}\varphi-\varphi E_{ab}$ on
$\mathrm{Hom}(V_1,V_2)$ satisfies them, and
$(E^{\rm h}_{ab})^q\varphi/q!=\sum_{i+j=q}(-1)^j(E_{ab}^i/i!)\varphi(E_{ab}^j/j!)$.
So if $L_1\subseteq V_1$ and $M_2\subseteq V_2$ are stable under the divided
powers and $M_2$ under $\Z_p$, then $\{\varphi:\varphi(L_1)\subseteq M_2\}$ is
stable under $\Z_p$ and the divided powers. If $V_1=V_2$ is finite
dimensional, the trace vanishes on $\mathcal R$.

(J5) \emph{Schur trace.} Let $Z=Ue_\nu=W$, $d=\dim Z\ge1$, and $\varphi\in
\mathrm{End}(Z)$ with $E^{\rm h}_{cc}\varphi=0$ for all $c$. Then
$\mathrm{Inv}(L)\varphi=(\mathrm{tr}\varphi/d)\mathrm{id}$. Condition (ev) holds
on $\mathrm{End}(Z)$: if $E^{\rm h}_{zz}\varphi=j\varphi$ with $|j|>l$ and
$\varphi(x)\ne0$ for a weight vector $x$ of weight $\mu$, then $\varphi(x)$ has
$E_{zz}$-eigenvalue $\mu_z+j$, so a weight of $W$ differs from $\mu$ by more
than $l$ in coordinate $z$, which (S1) excludes. By (J2), (J4),
$\psi=\mathrm{Inv}(L)\varphi$ commutes with all $E_{ab}$ and $\psi-\varphi\in
\mathcal R$. Then $\psi e_\nu$ has weight $\bar\nu$, so $\psi e_\nu=c'e_\nu$ by
(H3); $\ker(\psi-c')$ is stable and contains $e_\nu$, so $\psi=c'\mathrm{id}$;
and $c'd=\mathrm{tr}\,\psi=\mathrm{tr}\,\varphi$ (J4).

(J6) \emph{Conclusion.} Let $\tau=\pi_{\bar\nu}|_W\in\mathrm{End}(W)$; as $W=Ue_\nu$, $\tau$ is a
projection onto $\Q_pe_\nu$ (H3), of trace $1$, and $E^{\rm h}_{cc}\tau=0$. By
(J5), $\mathrm{Inv}(L)\tau=d^{-1}\mathrm{id}_W$. Let $\iota:W\to V$ be the
inclusion and regard $\tau$ as an element of $\mathrm{Hom}(W,V)$; by (J1)
(postcomposition with $\iota$ intertwines the $E^{\rm h}$),
$\mathrm{Inv}(L)(\iota\tau)=d^{-1}\iota$. Now $W\cap M$ is stable under the
divided powers, and $\tau(W\cap M)\subseteq\mathrm{Dv}$: for $y\in W\cap M$,
$\pi_{\bar\nu}y\in M$ is $0$ or a highest weight vector (it is a multiple of
$e_\nu$), so lies in $\Z_pe_\nu$ by (h4). By (J4) the set
$\{\varphi\in\mathrm{Hom}(W,V):\varphi(W\cap M)\subseteq\mathrm{Dv}\}$ satisfies
the hypotheses of (J3), so $p^{N'}d^{-1}\iota$ maps $W\cap M$ into $\mathrm{Dv}$.
Multiplying by the integer $d$ gives (WI). \hfill$\square$

\noindent\emph{Remark.} The standard basis is taken from
Theorem~\ref{i1:thm-stdbasis}. The argument for (WI) needs no hypothesis on
lattice bases of subspaces of $W$.

\subsubsection{Proof of Proposition~\ref{i1:prop-segdeg}: the parts (I1), (I4), (I9), (Deg)}\label{i1:ssec-pf-segdeg}

Notation of Definition~\ref{i1:def-A27}; assume $\mathrm{Hyp}(X)$, in
particular that the frame system is valid. Write $\epsilon(j)=\mathrm{fr}(j,r-1)$.

\emph{(F) Frame facts.} (a) $i\mapsto\mathrm{fr}(j,i)$ is strictly increasing
(Fr2), so $\mathrm{fr}(j,0)\le\mathrm{fr}(j,i)\le\epsilon(j)$; with (Fr3),
$\epsilon$ is strictly increasing and $\mathrm{fr}(j,i)<\mathrm{fr}(j',i')$ for
$j<j'$. The frame points $\mathrm{fr}(j-1,i)$ are letters of $j$ (Fr4).
(b) A $W$-letter $x$ of $j$ lies in $[\mathrm{fr}(j-1,0),\epsilon(j-1)]$ and
reduces into $\mathrm{TT}$, so by (Fr5) it is a frame point
$\mathrm{fr}(j-1,i')$, hence even (Fr4). (c) The late letters of $j$ are the
$x\in(\epsilon(j-1),\epsilon(j)]$ with $\mathrm{pos}(x)\in\mathrm{TT}$. (d) Every
$x\in\Z$ has a $j$ with $\epsilon(j-1)<x\le\epsilon(j)$ ($\epsilon$ is strictly
increasing and integer valued); hence for every $t\in\mathrm{TT}$ there are
$j$ and a late letter $x$ of $j$ with $\mathrm{pos}(x)=t$. (e) $\mathrm{pos}$ is
injective on every interval $(c,c+n']$. (f) $\sum_{j<J}|\mathrm{Late}_j|\le
|\mathrm{TT}|$: by (c) the $\mathrm{Late}_j$, $j<J$, are disjoint with union the
$x\in(\epsilon(-1),\epsilon(J-1)]$ reducing into $\mathrm{TT}$, and
$\epsilon(J-1)=\epsilon(-1)+n'$ (Fr1), so (e) applies. (g) As $n'$ is even,
$x$ and $\mathrm{pos}(x)\in\{0,\dots,n'-1\}$ have the same parity.

\emph{Proof of (Deg).} For an exponent $m$ on $\Z/n'\times\{0..k-1\}$ let
$\mathrm{vs}(m)=\sum_{t}\bigl([t\text{ even}]\sum_{c\ge r}m(t,c)+[t\text{ odd}]\sum_{c<r}m(t,c)\bigr)$
($t$ read in $\{0,\dots,n'-1\}$); it is additive, and
$\mathrm{vs}(m)=\sum_{t,c}m(t,c)\omega(t,c)$ with $\omega\in\{0,1\}$.
(a) $Q,N_s,a\ge1$: $\epsilon(0)$ is a letter of segment $0$, so
$1\le p^{\mathrm{dist}(0,\epsilon(0))}\le Q$ and $N_s\ge1$; then
$32pYk^2N_sQ^2(E_0+m^*+1)\le a$ gives $a\ge1$. So $\mathrm{La}\ge0$ and
$B_1:=\mathrm{La}\,a/(pk)+10ap^{-L_1}\ge0$. (b) \emph{One segment.} Let $y\in
\mathcal X'_j(\mathrm{mg}(j,\cdot),\mathrm{th}(j,\cdot),o(j,\cdot))$ and $m'$ an
exponent of $\varphi_j(y)$. By Definition~\ref{i1:def-seg}, $y$ involves only
variables $X_{(x,c)}$ with $x$ a letter of $j$ and, if $c\ge r$, $x$ late; and
$m'=\sum_vm(v)p^{\mathrm{dist}(j,v_1)}\epsilon_{(\mathrm{pos}(v_1),v_2)}$ for an
exponent $m$ of $y$. A variable $X_{(x,c)}$ with $x$ a $W$-letter has $c<r$
and $x$ even ((F)(b),(g)), so $\omega=0$ on its image. Grouping the other
terms by the late letter $x=v_1$, $\mathrm{vs}(m')\le\sum_{x\in\mathrm{Late}_j}
p^{\mathrm{dist}(j,x)}\mathrm{od}(m,x)$ ((F)(g) identifies $\omega$ with the
coordinates counted by $\mathrm{od}$), and $\mathrm{od}(m,x)\le o(j,x)$ by the
first condition of Definition~\ref{i1:def-seg}, while
$p^{\mathrm{dist}(j,x)}o(j,x)=\mathrm{La}\,a/(pk)+c_xap^{-L_1}\le B_1$. So
$\mathrm{vs}(m')\le|\mathrm{Late}_j|B_1$. (c) \emph{Generators.} Let
$P=\bigl(\prod_{j<J}\varphi_j(y_j)\bigr)x^l\in\mathrm{Gen}$ and $m$ an exponent of
$P$. Then $m=l+\sum_jm'_j$ with $m'_j$ an exponent of $\varphi_j(y_j)$, so by (b)
and (F)(f) $\mathrm{vs}(m-l)\le|\mathrm{TT}|B_1$. For $t\notin\mathrm{TT}$,
$\sum_cm(t,c)=d_t=0$, so $l(t,\cdot)=0$; for $t\in\mathrm{TT}$ pick by (F)(d) a
late $x$ with $\mathrm{pos}(x)=t$; then $l(t,c)<p^{\mathrm{dist}(j,x)}\le Q$.
Hence $\mathrm{vs}(l)\le|\mathrm{TT}|kQ$ and
$\mathrm{vs}(m)\le|\mathrm{TT}|(\mathrm{La}\,a/(pk)+10ap^{-L_1}+kQ)\le10|\mathrm{TT}|(\mathrm{La}\,a/(pk)+ap^{-L_1}+kQ)$.
(d) Exponents of sums and multiples of elements of $\mathrm{Gen}$ are
exponents of the summands, so the bound holds on $\mathcal N$.

\emph{Proof of (I1).} (a) For $x\in\mathrm{Late}_j$, $a\mapsto
\mathrm{pos}(\mathrm{row}_j(x,a))$ is injective on $\{0..r-1\}$. Let $S$ be the
set of even letters of $j$ below $x$, $\ell=\ell_j(x)$, $L=|S|$. The $r$
points $\mathrm{fr}(j-1,i)$ are distinct even letters of $j$ below $x$, so
$L\ge r$ and $\mathrm{row}_j(x,a)=\ell[L-r+a]$. Put $y_1=\ell[L-r]$. Every
$\ell[i]$, $i\ge L-r$, is $<y_1+n'$: otherwise $z=\ell[i]\ge y_1+n'$, and the
$r$ translates $\mathrm{fr}(0,i')+qn'$ ($=\mathrm{fr}(qJ,i')$, (Fr1)) lying in
$[y_1,y_1+n')$ are distinct even letters of $j$ (they lie in
$[\mathrm{fr}(j-1,0),\epsilon(j)]$; (Fr4)) below $z<x$; with $z$ this gives
$r+1$ elements of $S$ that are $\ge y_1$, while there are at most $r$. So the
$\mathrm{row}_j(x,a)$ are distinct points of $(y_1-1,y_1-1+n']$, on which
$\mathrm{pos}$ is injective (F)(e). (b) Reduction modulo $w$ commutes with
$\varphi_j$, $\iota_{pl}$, products and determinants, and $k_O[x]$ is a
domain, so it suffices that each reduced
$\varphi_j(\iota_{pl}\det(X_{(\mathrm{row}_j(x,a),b)})_{a,b<r})$ is nonzero.
Evaluate at the point $v$ with $v(t,c)=1$ if $t=\mathrm{pos}(\mathrm{row}_j(x,c))$
and $c<r$, and $0$ otherwise: by (a) the entry $(a,b)$ becomes
$v(\mathrm{pos}(\mathrm{row}_j(x,a)),b)^{p^{\mathrm{dist}}}=\delta_{ab}$, so the
value is $\det(1)=1$.

\emph{Numerical facts.} Put $c_0=8Yk^2N_s$ and
$T_j=a/(c_0p^{\mathrm{dist}(j,\mathrm{fr}(j-1,0))})$. (n1) $T_j\ge1$: as
$\mathrm{fr}(j-1,0)$ is a letter of $j$, $p^{\mathrm{dist}}\le Q\le Q^2$, and
$4p(E_0+m^*+1)c_0Q^2\le a$. (n2) $p^{R_j}\le T_j<p^{R_j+1}$ (from
$\lfloor T_j\rfloor\ge1$). (n3)
$R_j+\mathrm{fr}(j,0)-\mathrm{fr}(j-1,0)\le\epsilon(j+1)-\epsilon(j)+R_{j+1}$: by
(n2) at $j$ and $j+1$, $p^{R_j+\mathrm{dist}(j,\mathrm{fr}(j-1,0))}\le a/c_0<
p^{R_{j+1}+1+\mathrm{dist}(j+1,\mathrm{fr}(j,0))}$, and the distances are
$\epsilon(j)-\mathrm{fr}(j-1,0)$, $\epsilon(j+1)-\mathrm{fr}(j,0)$. (n4) For
$v\le\epsilon(j)$, $P_j(v)Q_j=P_{j+1}(v)$: by (n3) and
$\mathrm{fr}(j,0)>\mathrm{fr}(j-1,0)$ the maximum in $Q_j$ is the positive
number $\epsilon(j+1)-\epsilon(j)+R_{j+1}-R_j$, and
$\epsilon(j)-v+R_j+(\epsilon(j+1)-\epsilon(j)+R_{j+1}-R_j)=\mathrm{dist}(j+1,v)+R_{j+1}$.
(n5) By (Fr1), $\mathrm{dist}$, $R_j$, $Q_j$ are $J$-periodic in $j$ (with
the shift by $n'$ in $x$), so $e_{j,i}<rQ_j$ for $e\in\mathrm{EE}(\lambda)$ and
all $j\in\Z$ (indices mod $J$). (n6) For $\lambda\in\mathrm{LA}$ and
$x\in\mathrm{Late}_j$, $\sum_c\lambda(\mathrm{pos}(x),c)<kP_j(x)$ (each term is
$<P_j(x)$). (n7) For integers $n_2,s,e'\ge0$, $P,Q'\ge1$ with $s<kP$, $e'<rQ'$:
$\lfloor n_2/(PQ')\rfloor-(k+r+1)\le\lfloor(\lfloor(n_2-s)^+/P\rfloor-e')^+/Q'\rfloor$.
With $A,U,W'$ the three floors, $APQ'\le n_2<(U+1+k)P$ gives $AQ'\le U+k$, and
$U-e'<(W'+1)Q'$ with $e'<rQ'$ gives $AQ'<(W'+1+r+k)Q'$.

\emph{Proof of (I4).} Fix $\lambda$, $e$, $j$ and $b=b^{\lambda,e}_j$; it
vanishes off the letters of $j$ by definition. (a) Let $x$ be a $W$-letter of
$j$; by (F)(b) $x=\mathrm{fr}(j-1,i)$ for a unique $i$, a late letter of $j-1$,
and $b_x=\lfloor(\mathrm{nz}_\lambda(j-1,x)-e_{j-1,i})^+/Q_{j-1}\rfloor$. By (n4)
at $j-1$, $\mathrm{Bw}_j(x)=\lfloor a'_{\mathrm{pos}(x)}/(P_{j-1}(x)Q_{j-1})\rfloor
-(k+r+1)$ ($2r+K+1=k+r+1$). Apply (n7) with $n_2=a'_{\mathrm{pos}(x)}$,
$s=\sum_c\lambda(\mathrm{pos}(x),c)$, $P=P_{j-1}(x)$, $Q'=Q_{j-1}$, $e'=e_{j-1,i}$
((n6), (n5)): $\mathrm{Bw}_j(x)\le b_x$. (b) Let $x\in\mathrm{LL}_j$; it is late
with $x<\mathrm{fr}(j,0)$, so $b_x=\mathrm{nz}_\lambda(j,x)$. By the definition of
$\mathrm{la}(\mathrm{pos}(x))$, $s=\sum_c\lambda(\mathrm{pos}(x),c)\le a'_{\mathrm{pos}(x)}$
and $s\equiv a'_{\mathrm{pos}(x)}\pmod{P_j(x)}$, so
$b_x=(a'_{\mathrm{pos}(x)}-s)/P_j(x)$, and $0\le s<kP_j(x)$ gives
$\mathrm{Adm}_j(b)$.


\emph{Proof of (I9).} Assume $\mathrm{Hyp}(X)$ and $C_{La},C_k\ge1088$; fix
$j$ and $b$ with $\mathrm{Adm}_j(b)$. Write $\beta=\beta_j(b|_j)$, $Q=Q_j$,
$D_1=\mathrm{dist}(j,\mathrm{fr}(j,0))$, and recall that the letters of $j$
are $e_j(0)<\dots<e_j(N_j-1)$.

(a) \emph{The last letters.} A letter $y\ge\mathrm{fr}(j,0)$ of $j$ is a frame
point $\mathrm{fr}(j,i)$: otherwise, by induction on $i$, $\mathrm{fr}(j,i)<y$
for all $i<r$ (if $\mathrm{fr}(j,i)<y<\mathrm{fr}(j,i+1)$, then (Fr5) is
violated, as $y\bmod n'\in\mathrm{TT}$), contradicting
$y\le\epsilon(j)$. As the $\mathrm{fr}(j,i)$ are letters of $j$ ((F)(a),
(Fr4)), the letters $\ge\mathrm{fr}(j,0)$ are exactly the $r$ frame points;
so $N_j\ge r$, $e_j(i)<\mathrm{fr}(j,0)$ iff $i<N_j-r$, and
$e_j(N_j-r+t)=\mathrm{fr}(j,t)$, an even letter, i.e.\ a vector.

(b) \emph{Oldest letters.} Let $x\in\mathrm{Late}_j$. The $r$ points
$\mathrm{fr}(j-1,i)$ are even letters of $j$ below $x$, so $L\ge r$ and
$\mathrm{old}_j(x)=\ell[L-r]$ is a letter of $j$ ($\ell=\ell_j(x)$). If
$\mathrm{old}_j(x)$ is not a $W$-letter and $\mathrm{old}_j(x)\ge\mathrm{fr}(j,0)$,
then $\ell[L-r],\dots,\ell[L-1]$ and $x$ are $r+1$ distinct letters
$\ge\mathrm{fr}(j,0)$, against (a). Hence, if $\mathrm{old}_j(x)$ is not a
$W$-letter, then $\mathrm{fr}(j-1,r-1)<\mathrm{old}_j(x)<\mathrm{fr}(j,0)$,
that is, $\mathrm{old}_j(x)\in\mathrm{LL}_j$.

(c) \emph{A lower bound on paths.} Let $\mu$ be a Pieri path in $k$
coordinates of degrees $dd$, and $d^-(s)=\sum\{dd_l:\ l<s,\ l\text{ a
covector}\}$. Then $\mu_{s,c}\ge-d^-(s)$ for all $s,c$, by induction on $s$: a
vector step does not decrease any coordinate ($\mu_{s+1}/\mu_s$ is a
horizontal strip), and at a covector step
$\mu_{s,c}-\mu_{s+1,c}\le\sum_{c'}(\mu_{s,c'}-\mu_{s+1,c'})=dd_s$.

(d) \emph{Degrees at frame points.} Let $\mu$ be a path for $b|_j$ with
degrees $dd$, and $e_j(i)\ge\mathrm{fr}(j,0)$. Then $dd_i\le2\beta$. Indeed by
(a), $e_j(i)=\mathrm{fr}(j,t)$ with $i=N_j-r+t$, a vector; with
$f=\mu_{i+1}$, $g=\mu_i$ we have $f_{c+1}\le g_c$, so
$dd_i=\sum_c(f_c-g_c)=f_0-g_{k-1}+\sum_{c<k-1}(f_{c+1}-g_c)\le f_0-g_{k-1}$.
Now $f_0=\beta$ by the head condition of a path ($c=0\le t$), and
$g_{k-1}\ge-d^-(i)$ by (c). Finally $d^-(i)\le\beta$: a covector letter $l$
has $l<N_j-r$ by (a), so $e_j(l)<\mathrm{fr}(j,0)$ and
$dd_l=\mathrm{Bd}_j(b|_j)(e_j(l))\ge0$; summing over all $l<N_j-r$ gives
$\beta$.

(e) \emph{Residue counts.} Let $r\ge2$, $Q\ge1$, $\beta\ge0$ with
$2\beta<Q$, $T=Q-2\beta$, and $c,\rho\in\Z$. Then
$\#\{w\in[0,T)^r:\ c+\sum_iw_i\equiv\rho\pmod Q\}\ge Q^{r-1}(1-2r\beta/Q)$.
Let $\mathrm{cnt}_s(c)=\#\{u\in[0,T)^s:\sum u\equiv c\}$. Splitting off a
coordinate, $\mathrm{cnt}_{s+1}(c)=\sum_{a<T}\mathrm{cnt}_s(c-a)$, and
$\mathrm{cnt}_1(c)=[c\bmod Q<T]$. So $\mathrm{cnt}_2(c)$ is the number of
$a<T$ with $(c-a)\bmod Q<T$; the other $a$ give distinct residues in
$[T,Q)$ (as $a<T\le Q$), at most $Q-T$ of them, so
$\mathrm{cnt}_2(c)\ge T-(Q-T)$, and by induction
$\mathrm{cnt}_{s+2}(c)\ge T^s(T-(Q-T))$. Write $r=s+2$ and $b''=2\beta$, so
$Q^{r-1}(1-2r\beta/Q)=Q^s(Q-(s+2)b'')$. If $2b''\le Q$, then
$Q^s(Q-(s+2)b'')\le(Q-b'')^s(Q-2b'')=T^s(T-(Q-T))$ by induction on $s$, the
step being $Q^{s+1}(Q-(s+3)b'')=(Q-b'')Q^s(Q-(s+2)b'')-(s+2)Q^sb''^2$; if
$2b''>Q$ the left side is negative. Apply this to $\mathrm{cnt}_r(\rho-c)$.

(f) \emph{The input $(\beta)$: $4Y\beta\le Q$.} (f1) For a letter $x$ of $j$
with $x<\mathrm{fr}(j,0)$: $0\le\mathrm{Bd}_j(b|_j)(x)\le a'_{\mathrm{pos}(x)}/P_j(x)\le2a/P_j(x)$.
At a $W$-letter, $\mathrm{Bd}_j=(\mathrm{Bw}_j(x)-m_0)^+$ and $\mathrm{Bw}_j(x)\le
a'_{\mathrm{pos}(x)}/P_j(x)$; otherwise $x\in\mathrm{LL}_j$ and
$\mathrm{Bd}_j(b|_j)(x)=b_x\le a'_{\mathrm{pos}(x)}/P_j(x)$ by
$\mathrm{Adm}_j(b)$; and $a'_t\le d_t\le\deg t\le2a$, using $E_t\ge0$ and
$E'_t\ge0$ ($E'_t\in\N$ in the data). (f2) For $p\ge2$ and a
finite set $s\subseteq\N$ of integers $\ge m_1$,
$\sum_{D\in s}p^{-D}\le2p^{-m_1}$ (geometric series). (f3) $L_1\ge1$:
otherwise $128k^2N_sY^2\le p^{L_1-1}<1$, while the left side is $\ge1$.
(f4) Let $s$ be the set of letters $x<\mathrm{fr}(j,0)$ of $j$. By
$\mathrm{Gap}_{L_1}(j,\mathrm{fr}(j,0))$ (Fr6), $\mathrm{fr}(j,0)-\max s>L_1$,
so for $x\in s$: $\mathrm{dist}(j,x)=\epsilon(j)-x>\epsilon(j)-\mathrm{fr}(j,0)+L_1=D_1+L_1$.
(f5) By (f1), $P_j(x)=p^{R_j+\mathrm{dist}(j,x)}$, injectivity of
$x\mapsto\mathrm{dist}(j,x)$ on $s$, (f4) and (f2) with $m_1=D_1+L_1+1$:
$\beta\le2a\,p^{-R_j}\cdot2p^{-(D_1+L_1+1)}$. (f6) Let
$d=\epsilon(j+1)-\epsilon(j)\ge0$ ((F)(a)). Then
$\mathrm{dist}(j+1,\mathrm{fr}(j,0))=D_1+d$, so (n2) at $j+1$ gives
$a<p^{R_{j+1}+1}\,8Yk^2N_s\,p^{D_1+d}$; and by (n3),
$Q=p^{d+R_{j+1}-R_j}$ (the exponent is $\ge0$). (f7) Hence
\[
\frac\beta Q\le\frac{4a}{p^{D_1+L_1+1+d+R_{j+1}}}<\frac{32Yk^2N_s}{p^{L_1}}\le\frac1{4pY}\le\frac1{4Y},
\]
the third inequality being $128k^2N_sY^2\le p^{L_1-1}$ (with (f3)).

(g) \emph{Assembly.} Let $T_j(b|_j)$ be finite, $\rho\in\N^r$ and
$c_*=Q^{r-1}(1-\frac r{2Y})$; recall $c_j=(Q^{r-1}\max(0,1-\frac r{2Y}))^r$. If
$1-\frac r{2Y}\le0$, then $c_j=0$ and there is nothing to prove. Otherwise
$c_j=c_*^r$; by (f) and $Y\ge1$, $4\beta\le Q$, so $T'=Q-2\beta\ge1$ and
$2\beta<Q$. For $\tau=(\mu,I)\in T_j(b|_j)$ fix degrees $dd_\tau$ of $\mu$ and
put $\delta_\tau(i')=dd_\tau(N_j-r+i')$, which is $\le2\beta$ by (d). Let $A$ be
the set of pairs $(\tau,h)$ with $\tau\in T_j(b|_j)$ and
$h\in[0,T')^{r\times r}$ such that
$\varphi(\tau,h)_{i'}=\delta_\tau(i')+\sum_ch(i',c)\equiv\rho_{i'}\pmod Q$
for all $i'<r$.
(g1) $\varphi(A)\subseteq E=\prod_{i'<r}\{0\le e<rQ:\ e\equiv\rho_{i'}\pmod Q\}$,
as $\varphi_{i'}\le2\beta+r(T'-1)<rT'+2r\beta=rQ$.
(g2) For $e\in E$, the map $((\mu,I),h)\mapsto(\mu,I,h)$ injects
$\varphi^{-1}(e)$ into $\mathrm{Idx}_j(\tau_j(b,e)|_j)$. Indeed $b|_j$ and
$\tau_j(b,e)|_j$ agree below $\mathrm{fr}(j,0)$, and the value of
$\mathrm{Bd}_j$ at $W$-letters does not depend on $b$; the path conditions
depend on $b$ only through $\mathrm{Bd}_j(b)$ below $\mathrm{fr}(j,0)$ and
$\beta_j(b)$, and the good conditions only through $\theta^c_j(b)(x)$, $x$
late, which by (b) uses $b$ only at a point of $\mathrm{LL}_j$ (and
$o^c_j$ does not depend on $b$). So $\mu$ is a good path for
$\tau_j(b,e)|_j$, with the same degrees $dd_\tau$ and the same $\beta$.
Moreover $h_v+2\beta<T'+2\beta=Q$, and for $e_j(i)=\mathrm{fr}(j,i')$ (so
$i=N_j-r+i'$ by (a)) we have
$dd_{\tau,i}+\sum_ch(i',c)=e_{i'}=\tau_j(b,e)|_j(\mathrm{fr}(j,i'))$, as
$\mathrm{fr}(j,i')$ is late.
(g3) For fixed $\tau$, every $h$ whose rows lie in
$W_{i'}=\{w\in[0,T')^r:\ \delta_\tau(i')+\sum w\equiv\rho_{i'}\pmod Q\}$ gives
an element of $A$; by (e) and $2r\beta/Q\le r/(2Y)$ (which is $4Y\beta\le Q$),
$\#W_{i'}\ge c_*$. So at least $c_*^r$ elements of $A$ lie over $\tau$.
(g4) $\mathrm{Idx}_j(\tau_j(b,e)|_j)$ is finite: $(\mu,I,h)\mapsto((\mu,I),h)$
maps it injectively into $T_j(b|_j)\times[0,Q)^{r\times r}$ by the comparison
of (g2). Counting $A$ by $\tau$ and by the value of $\varphi$, using (g1)--(g3):
$\#T_j(b|_j)\,c_*^r\le\#A=\sum_{e\in E}\#\varphi^{-1}(e)\le\sum_{e\in E}\#\mathrm{Idx}_j(\tau_j(b,e)|_j)$,
which is (I9), the sum over $e\in E$ being the range of (I5). This proves (I9).

The items (I2), (I3), (I5)--(I8) are not proved here; they form
Proposition~\ref{i1:prop-segdeg-rest}. Their proofs are long: they consist of the reduction of the
coarse generators through column-word identities (I2), the injectivity and
degree bookkeeping of the coarse family (I3), the counting with binomial and
logarithmic Weyl-dimension estimates (I5), (I8), and the path-basis
membership and independence (I6), (I7), which rest on
Lemmas~\ref{i1:lem-A15I}, \ref{i1:lem-A16}, \ref{i1:lem-A17},
\ref{i1:lem-A33} and \ref{i1:lem-counts}. Of these, (I5) is the shortest,
but it rests on a further package of
bounds for the coarse top counts $E_t$, $a'_t$ which we do not reproduce,
so we leave it stated.


\subsection{Outline of the formally verified steps}\label{i1:ssec-structure}

We outline the formal proofs of the five statements of this section that are formally verified, Propositions~\ref{i1:prop-first},
\ref{i1:prop-closed}, \ref{i1:prop-gc}, \ref{i1:prop-linkG-res} and \ref{i1:prop-segdeg-rest}. The outline is not a proof. Each
step is classified as (L), (N) or (R) as explained in the introduction to this annex; when we say that we do not describe the proof of
a step, the step enters only through its statement. Throughout, $O$ is a discrete valuation ring of characteristic $0$ with
uniformiser $w$ and residue characteristic $p$; $H^0(\mu)$ is the dual Weyl module (module of global sections) of highest weight
$\mu$ of $GL_n$ over $O$, realised as a lattice of polynomial functions, and $V(\mu)$ is the Weyl module.

\medskip\noindent\label{i1:str-first}\emph{The first theorem (Proposition~\ref{i1:prop-first}).}
The three layers (a), (b), (c) described after the statement form the top of the formal proof. Layer (c) is (R): it is the choice
of thresholds at a fixed record, written in the remark after the statement. The inputs of (a) and (b) are a list of \emph{tool
statements} on $\mathfrak{gl}_n$-modules over $O$ and a list of \emph{admission} and \emph{closed-form} statements on the column
systems. The tool statements are the following.
\begin{enumerate}
\item[(T1)] (L) \emph{Weyl's dimension formula} for $GL_{m+1}$, in the product form $\dim V(\lambda)=\prod_{i<j}(\lambda_i-\lambda_j+j-i)/(j-i)$,
for each component of a dominant multi-weight; this is Theorem~\ref{i1:thm-weyl}.
\item[(T2)] (L) \emph{Surjectivity of multiplication}: for a lattice $\Lambda$ and antitone weights $\lambda,\mu$,
$H^0(\lambda+\mu)$ is the $O$-span of the products $FG$ with $F\in H^0(\lambda)$, $G\in H^0(\mu)$. The multiplication
$H^0(\lambda)\otimes H^0(\mu)\to H^0(\lambda+\mu)$ is surjective over $\Z$ and over every field (projective normality of flag
varieties; see J.~C.~Jantzen, \emph{Representations of algebraic groups}, 2nd ed., II.14.20, and the work of Ramanathan
and Mehta--Ramanathan cited there (reference to be confirmed)), hence over $O$.
\item[(T3)] (L) \emph{No zero divisors modulo $w$}: if $FG\in wH^0(\lambda+\mu)$, then $F\in wH^0(\lambda)$ or $G\in
wH^0(\mu)$. In substance this is published: $H^0(\lambda)\otimes_Ok=H^0_k(\lambda)$ by Kempf vanishing and base change (Jantzen II.4.5,
II.8.8 (reference to be confirmed)), and the sections over the residue field $k$ form part of the multi-homogeneous coordinate ring
of the flag variety over $k$, which is a domain.
\item[(T4)] (L) \emph{Standard monomial theory over $\Z$}: the $\Z$-span of the bracket polynomials
(products of minors indexed by semistandard tableaux) of shape $\mu$ equals the intersection of their $\Q$-span with the
$\nabla$-lattice. This is Theorem~\ref{i1:thm-stdbasis} (De Concini--Eisenbud--Procesi), together with the characteristic-free
Schur functors of Akin--Buchsbaum--Weyman (Adv.\ Math.\ 44 (1982) 207--278). The formal statement is the
equality of two $\Z$-lattices, that is, the statement that the bideterminants form a $\Z$-basis of the $\Z$-form; the match with
the published wording is (R).
\item[(T5)] (L) \emph{Weyl duality over $O$}: the lattice generated by a highest weight vector under the divided powers (Kostant's
$\Z$-form, the Weyl module) is the dual lattice of the dual Weyl module under the contravariant pairing (Jantzen II.2.13,
II.8; Kostant, Theorem~\ref{i1:thm-kostant}).
\item[(T6)] (L) \emph{Witt vectors}: an absolutely unramified complete discrete valuation ring with residue field $\mathbb F_{p^{n'}}$ and
a Frobenius lift is isomorphic to $W(\mathbb F_{p^{n'}})$ compatibly with Frobenius (Serre, \emph{Local Fields}, II \S\S5--6).
\item[(T7)] (L) \emph{The extremal projector} of $\mathfrak{gl}_n$ (Asherova--Smirnov--Tolstoy, Theoret.\ Math.\ Phys.\ 8 (1971);
Zhelobenko, \emph{Representations of reductive Lie algebras}, 1994) enters through one admission clause of
the covariant degrees; we do not describe its proof here.
\item[(T8)] (N) The remaining tool clauses: two flatness statements (for the top components and for the Hilbert lattices), a basis of the Weyl
modules adapted to heights with an inclusion of the dual lattices, a perfectness statement
for an edge pairing, a spanning statement for Hilbert lattices, the test side (a duality $\Psi$ between $V(\lambda)$ and its dual with
$e\mapsto e^\vee$ an involution on edges), and the passage $\mathrm{LP}'_1\Rightarrow\mathrm{LP}'$. They are statements about lattices in
Weyl modules attached to the Hida setting; we do not describe their proofs here.
\end{enumerate}
The admission and closed-form statements are the following.
\begin{enumerate}
\item[(A)] (R), given this section: the admission clauses (Proposition~\ref{i1:prop-adm}) feed the closed admission of layer (a)
directly.
\item[(E)] (N) \emph{The extended closed form} of the column systems: a closed-form implication for the window and segment ranks.
Its top layer consists of five implications between three interface statements (on columns, on pre-images and on a truncated core);
we do not describe the proofs below that layer.
\item[(M)] (N) \emph{The explicit minor map}: a rank statement for the minors of the twisted column embeddings. Layer (b) reduces the
second part of the projection lemma to it through the forms on dual Weyl modules (Proposition~\ref{i1:prop-Fm}, proved here), the
Smith normal form (Theorem~\ref{i1:thm-elemdiv}, (L)) and a contraction statement ((N); we do not describe its proof here).
\end{enumerate}
Thus (T1)--(T6) are published theorems (proved in the formal development), (T7) is published in substance,
and (T8), (E), (M) are new. By number of statements this is about half; by length of the formal proofs the published part is much
smaller, since the column-system statements (E), (M) form the bulk.

\medskip\noindent\label{i1:str-closed}\emph{The closed form (Proposition~\ref{i1:prop-closed}).}
The formal proof derives the closed form from a single \emph{valuation chain} statement: for every prime $p$, every discrete
valuation ring $O$ of characteristic $0$, every valid frame of even period, signed chain data tied to it, and polynomials $A_i$
supported on the letters of the segment $i$, lower bounds $t(j,u)$ for the valuations of the $A_i$ along the chain imply the same lower
bounds for the digits of the level polynomials of the plane substitution. The chain statement is proved from five statements:
(V1) the valuation of an evaluation, (V2) the valuation along a filtration, (V3) nonnegativity, (V4) a monomial phase identity, and
(V5) the composition of phases. Statements (V1)--(V3) are (R): they are properties of valuations of polynomials over a discrete valuation
ring of the type of Gauss's lemma, published in substance. Statements (V4), (V5) are (N) identities for the signed chain data; we do
not describe their proofs here. The derivation of the closed form from the chain statement, via the membership statement
$\mathrm{RestMEM}$ and the Gram statement (Proposition~5 of the chain) with the test-side reduction, is (N); see the remark after the
proposition. The published inputs are therefore elementary valuation theory, and the content is new.

\medskip\noindent\label{i1:str-gc}\emph{Consequences of group compatibility (Proposition~\ref{i1:prop-gc}).}
Each of (i)--(iii) is obtained, eventually along the family, from one field of the group compatibility record by a pointwise
implication and Lemma~\ref{i4:lem-closed}(c).
For (i), \emph{trivial stabilisers}, which is (L) and (R), one uses the fields (herm) and (gat). The stabiliser of a level datum is a subgroup of a totally definite
unitary group preserving a lattice, hence finite, and its elements are congruent to the identity at places above the two distinct
primes $g_0,g_1$; such elements are trivial. This is Minkowski's lemma in the form: a finite-order element of $GL_n(\mathcal O)$
congruent to $1$ modulo ideals of two distinct residue characteristics is $1$ (H.~Minkowski, J.\ Reine Angew.\ Math.\ 101 (1887);
J.-P.~Serre, \emph{Rigidit\'e du foncteur de Jacobi d'\'echelon $n\ge3$}, S\'em.\ Cartan 1960/61 (reference to be confirmed)). The formal
proof uses the total definiteness of the Gram matrix and the congruences at the two places; we describe only its top level. The finiteness
of the classes is (R).
Item (ii), the \emph{standard subring}, is (R) and follows from the field (ring); item (iii), the \emph{descent of the norm lattices},
is (R) and follows from the field (cmp) through a valuation-ratio statement.
The proposition is thus routine given the record, with one classical lemma. It is stated informally because the three
predicates are those of Section~\ref{i4:sec}.

\medskip\noindent\label{i1:str-linkG-res}\emph{The remaining inputs of $\mathrm{LinkG}$ (Proposition~\ref{i1:prop-linkG-res}).}
The formal proof is the chain described after the statement:
(G1) (R) the type-two and type-one bounds, from the second and first theorems at the family, by comparing their constants with
the loss constants, and a family $\kappa$-bound, also (R);
(G2) (R) the exponent clause and the sandwich, from $C_e\ge1$: the Weyl sandwich of the local model (which uses
Proposition~\ref{i1:prop-Fm}) and an exponent bound;
(G3) (N) the orbit lattice bounds (hypothesis (iv) of Definition~\ref{i1:def-locdatum}): a linear-programming input from the two
theorems, (ii) and the window log bound, with a base-change statement for the local model;
(G4) (R) the order and size bounds, from the window log bound and Lemmas~\ref{i1:lem-osize}, \ref{i1:lem-A2}, and the fact that the
operators $y_v$ preserve the lattices;
(G5) (R), given Proposition~\ref{i1:prop-SLC} and Lemma~\ref{i1:lem-zeroch}: the gluing into $\mathrm{LinkG}$, written after the statement;
(G6) (N) the \emph{gain bound}: the star inequality on each served orbit, a sum over orbits, and the passage between the two forms of the
ask series, by elementary estimates whose proof we do not describe here.
No published theorem beyond those already used enters; the content is the bookkeeping of the family constants and the two
new estimates (G3), (G6). The proposition is stated informally because clauses (1), (5) of Definition~\ref{i1:def-linkG} and the predicates of
Proposition~\ref{i1:prop-gc} are not defined here.

\medskip\noindent\label{i1:str-segdeg}\emph{The counting part of $\mathcal A_{27}$ (Proposition~\ref{i1:prop-segdeg-rest}).}
(I2) (L) and (N): \emph{reduction of the coarse generators}. Over a ring of characteristic $p$, $F(x)^{p^e}=F^{(p^e)}(x^{p^e})$ (the
Frobenius endomorphism; for polynomials over $\mathbb Z/p$ this is the identity $f(x)^p=f(x^p)$, which is (L)). The formal proof applies it to
the determinants of the segment frames reduced modulo $w$, factorises the coarse labels along the
period of the frame ((R)) and reduces column words in four steps ((N)); the
lifted identities give, for $u\in\Lambda\mathrm{Gen}$, an element $g\in\mathrm{Gen}$ with $\overline{\Delta_lu}=\bar g$.
(I3) (N) \emph{Injectivity and degrees of the coarse family}: the map sending $\lambda$ to its monomial labels is injective and degree
controlled (an image formula, injectivity and a degree bound); summing ranks gives (I3).
(I6), (I7) (N), given Lemmas~\ref{i1:lem-A15I}, \ref{i1:lem-A16}, \ref{i1:lem-A17}, \ref{i1:lem-A33}: the family elements indexed by
good paths lie in $M_j(b)$ and are linearly independent (the path basis). The background, (L), is the theory of Gelfand--Tsetlin bases and their integral
forms; the independence used is that of Lemma~\ref{i1:lem-A15I}.
(I5), (I8) (L) and (N): \emph{counts}. These combine Gelfand--Tsetlin pattern counts (Lemma~\ref{i1:lem-counts}, that is, the branching rule and
the Pieri rule, (L)) with logarithmic estimates of Weyl dimensions (Theorem~\ref{i1:thm-weyl}, (L)) and binomial estimates ((R)); (I5)
uses in addition lower bounds for the coarse top counts ((N)) and a threshold power bound
deduced from them and from the oldest-letter lemma (proved in \S\ref{i1:ssec-pf-segdeg}).
The published inputs are the Frobenius identity, the branching and Pieri counts and Weyl's dimension formula, all
already cited in this section; the rest, about four fifths, is new combinatorics of the segment modules.

\medskip\noindent\textbf{Status.}
\emph{Proved here:} Theorem~\ref{i1:thm-loss} (both parts) from $\mathrm{LinkG}$, together with the elementary loss bound
(Theorem~\ref{i1:thm-elementary}), the instantiation (Theorem~\ref{i1:thm-setting}) and the small-share weights
(Proposition~\ref{i1:prop-chosenW}); the reduction of $\mathrm{LinkG}$ to group compatibility and $\mathrm{Adm2F}$
(Proposition~\ref{i1:prop-linkG}) modulo the formally verified statements below; the local statement
(Proposition~\ref{i1:prop-loc}) from the admission clauses and the closed form; all $37$ admission clauses
(Proposition~\ref{i1:prop-adm}) except the counting part of $\mathcal A_{27}$; and Propositions~\ref{i1:prop-A19c}, \ref{i1:prop-A21},
\ref{i1:prop-Fm} and the parts (I1), (I4), (I9), (Deg) of Proposition~\ref{i1:prop-segdeg}.
\emph{Published theorems, quoted:} quadratic reciprocity (Theorem~\ref{i1:thm-qr}), the standard basis of the dual Weyl module
(Theorem~\ref{i1:thm-stdbasis}), elementary divisors (Theorem~\ref{i1:thm-elemdiv}), Weyl's dimension formula
(Theorem~\ref{i1:thm-weyl}), Kostant's commutation formula (Theorem~\ref{i1:thm-kostant}) and standard facts on Witt vectors.
\emph{Formally verified, \cite{CharacterFamily}:} Propositions~\ref{i1:prop-first}, \ref{i1:prop-closed}, \ref{i1:prop-gc},
\ref{i1:prop-linkG-res} and \ref{i1:prop-segdeg-rest}. Group compatibility and $\mathrm{Adm2F}$ for the families, and the ties of the
column package, are supplied by Section~\ref{i4:sec}.


\section{The divisor bound}\label{i2:sec}

This section proves the divisor bound, condition (P1) of
Definition~\ref{M:def:package}. The proof has three layers.
\begin{enumerate}
\item An elementary reduction (\S\ref{i2:ssec-red}): the divisor bound follows from
  a \emph{local index bound} $\mathrm{LinkN}$, a \emph{Selmer hypothesis} $\mathrm{HypV}$
  and a rank equality. The Selmer hypothesis is used twice: once to make
  $\mathrm{LinkN}$ applicable, and once to bound the Selmer sum it produces.
\item The proof of $\mathrm{LinkN}$ (\S\S\ref{i2:ssec-perprime}--\ref{i2:ssec-cuts}). A
  weighted sum of local congruence indices is split by residue characteristic; at each
  prime $l$ the local index is bounded by local terms plus the $l$-part of a Selmer
  length (an \emph{admissible table}); the local terms are summed over $l$ by
  elementary prime-sum estimates. A \emph{cut} of the Hecke data reduces the general
  case to one where these tables exist, at a cost of $\log|G|$.
\item The local inputs (\S\ref{i2:ssec-local}): at the generator primes the local
  term comes from a crystalline lifting theorem with explicit exponents; at all other
  primes from a global lattice-index (Ribet-type) argument. Here the published
  $p$-adic Hodge theory enters, and here most of the steps are taken from the
  formal verification \cite{CharacterFamily} rather than proved.
\end{enumerate}
The main result is Theorem~\ref{i2:thm-div}, stated in \S\ref{i2:ssec-data} after the data it
concerns; Proposition~\ref{i2:prop-P1} turns it into condition (P1).
Subsection~\ref{i2:ssec-num} records the numerical estimate for the period constant
$\log\Omega$, and \S\ref{i2:ssec-P1} the passage to condition (P1). The proofs of the
propositions stated in \S\ref{i2:ssec-local} occupy \S\S\ref{i2:ssec-pf-local}--\ref{i2:ssec-pf-gen}, and
\S\ref{i2:ssec-proof-P1} proves Proposition~\ref{i2:prop-P1}. Subsection~\ref{i2:ssec-struct} outlines the
formally verified steps, and the section ends with a summary of its status.
Throughout, $\log$ is the natural logarithm, and for a real sequence $(a_l)_{l\in\N}$ we
write $\sum'_la_l$ for its sum if it is summable and $0$ otherwise. This convention is the
reason for the summability premise in Definition~\ref{i2:def-LinkN}.

\subsection{Data and the main statement}\label{i2:ssec-data}

We recall only the parts of the data that enter this section.

\begin{definition}[Data, size data, families]\label{i2:def-data}
A \emph{datum} $D$ consists, among other things, of:
\begin{itemize}
\item ranks $n,m,f,k,t,r\in\N$ with $f>0$, $t$ odd, $r=k+m-1$, $n_V=r+t$, and
  $k,m$ even, $2\le m$, $m+3\le k$, $r+5\le t$, $3m<t$, $m\le f\le k$, $f=2n$ (in particular
  $m\le t$ and $r\le t$);
\item a \emph{host}: a CM field $E=E_D$, Galois and abelian over $\Q$, of degree $2f$; primes
  $q\equiv3\pmod 8$ and $d\equiv7\pmod{16}$; an even integer $Q_S$ prime to $qd$ (the
  \emph{small} primes are those dividing $Q_S$); two \emph{generator primes} $g_0,g_1$; a
  finite set $T$ of tops; the root discriminant $\mathrm{rd}_E$ and $v=\mathrm{rd}_E/2$;
\item a totally complex Galois coefficient field $L_D$ and a finite set of its places $w$,
  each with residue characteristic $p(w)$, ramification index $e_w$, weight
  $\lambda_w=\log\mathrm N(w)/[L_D:\Q]$ and a twist label $\mu_w$; a place is
  \emph{twisted} if $\mu_w\ne\mu_0$ (the base label);
\item for each lift $j\in\{1,\dots,m\}$ and place $w$ a finite set $\mathrm{locs}(j,w)$ of
  \emph{local indices} $i$, namely the maximal ideals of the local Hecke algebra
  $\mathbb T_{j,w}$ (the $O_{j,w}$-subalgebra of $\mathrm{End}(V_{j,w})$ generated by the
  Hecke operators); each $i$ carries a residue characteristic $l_i=p(w)$, a finite set of
  eigensystems $X_i$, a subset $Z_i$ and a subset $Z^{\rm good}_i\subseteq Z_i$ of
  \emph{good} eigensystems with multiplicities $\mathrm{mult}_i(\pi)$, and a normalised
  congruence length $v^{cc}_i(\pi)\in[0,\infty]$ at $\pi\in Z^{\rm good}_i$; its local index
  $\eta^{\rm loc}_i\in\Q_{\ge0}$ is the normalised length of the congruence module of
  $Z^{\rm good}_i$ (Lemma~\ref{i2:lem-length-le}); and
  $\eta_j(w)=\sum_{i\in\mathrm{locs}(j,w)}\eta^{\rm loc}_i$;
\item ranks $\mathrm{rk}_j(w)\in\N$, and a finite \emph{level group}
  $G=\prod_{k=0,1}\mathrm{GL}_r(\mathbb F_{g_k})^{f}$;
\item \emph{Selmer length data} $\ell_D(l,e)\in\Q_{\ge0}$ and, for every prime $p$, $p$-adic
  Galois data $\mathrm{ld}_D(p)$ (Definition~\ref{i2:def-LData}) such that, for $0<e<t$,
  if $\ell(\mathrm{Sha}(\mathrm{ld}_D(p),e))=N<\infty$ then $\ell_D(p,e)=N$.
  \hfill\refstepcounter{equation}(\theequation)\label{i2:eq-sha}
\end{itemize}
A \emph{size datum} $L$ for $D$ carries a scale $(x,\varepsilon)$ with $x>1$,
$0<\varepsilon<1/52$; we put
\[
\mathsf L=x^\varepsilon,\quad b=e^{\mathsf L},\quad a=xb,\quad B_0=f\,m\,b\,x^\varepsilon,\quad
N_c=\frac{|T_2|\,k\,b\,x^\varepsilon}{n\cdot nx^{-\varepsilon}}
\]
($T_2$ a finite set of the host). It also carries a finite \emph{support} of places, the
integers $d_{Z_0}(j)=\dim_\C Z_j(\mu_0)$, further real parameters used in
Definition~\ref{i2:def-budgets}, and the real number
$\Lambda=\Rel\bigl(L'(1,\chi_d)/L(1,\chi_d)\bigr)$, $\chi_d$ the Legendre character
modulo $d$.

A \emph{family} $\mathcal F$ with exponent $\varepsilon$ consists of a threshold $x_1>1$,
index sets $I(x)$, and data $D_{x,i}$ with size data $L_{x,i}$ ($i\in I(x)$) of exponent
$\varepsilon$ and with $x(L_{x,i})=x$ for $x>x_1$.
\end{definition}

\begin{definition}[$l$-adic Galois data and Tate--Shafarevich groups]\label{i2:def-LData}
Let $F$ be a number field with algebraic closure $\Omega$ and $l$ a prime. An
\emph{$l$-datum} consists of the set $S$ of places of $F$ above $l$, a place of $\Omega$
above each $v\in S$, and a continuous character $\chi:G_S\to\Z_l^\times$ inflating to the
$l$-adic cyclotomic character, $G_S$ being the quotient of $\mathrm{Gal}(\Omega/F)$ by the
closed normal subgroup generated by the inertia groups outside $S$. For $e\in\Z$,
\[
\mathrm{Sha}(\mathrm{ld},e)=\ker\Bigl(H^1\bigl(G_S,\Q_l/\Z_l(\chi^{-e})\bigr)\to\prod_{v\in S}H^1\bigl(G_v,\Q_l/\Z_l(\chi^{-e})\bigr)\Bigr),
\]
and $\ell(\mathrm{Sha}(\mathrm{ld},e))\in\N\cup\{\infty\}$ is its length as a $\Z_l$-module.
\end{definition}

\begin{definition}[Selmer terms, weighted index sums]\label{i2:def-vv}
For a datum $D$ and $l\in\N$ put $\mathrm{VV}_D(l)=\sum_{e=1}^{t-1}\ell_D(l,e)$ and
$\mathrm{vv}_D(l)=\mathrm{VV}_D(l)\log l$ if $l$ is prime, $\mathrm{vv}_D(l)=0$ otherwise.
$D$ is \emph{Sha-finite} if $\ell(\mathrm{Sha}(\mathrm{ld}_D(p),e))<\infty$ for all primes $p$
and $0<e<t$. For a size datum $L$ and $j\in\N$,
\[
\eta\Sigma_L(j)=\sum_{\substack{w\in\mathrm{supp}\,L\\ \mu_w\ne\mu_0}}\lambda_we_w\,\eta_j(w).
\]
\end{definition}

\begin{definition}[Local index bound]\label{i2:def-LinkN}
For a size datum $L$ of $D$, $j\in\N$ and $C\in\R$, $\mathrm{LinkN}_{\rm at}(L,j,C)$ is the
statement: if $(\mathrm{vv}_D(l))_l$ is summable, then
\[
\eta\Sigma_L(j)\le d_{Z_0}(j)\Bigl(C\bigl(fmb\,\mathsf L^2+(\log x)^2N_c\bigr)+\textstyle\sum'_l\mathrm{vv}_D(l)\Bigr).
\]
$\mathrm{LinkN}(\mathcal F)$: there are $C_N,x_0$ such that $\mathrm{LinkN}_{\rm at}(L_{x,i},j,C_N)$
holds for all $x\ge x_0$, all $i\in I(x)$ with $D_{x,i}$ Sha-finite, and all lifts $j$.
$\mathrm{SmallLevelRankEq}(\mathcal F)$: eventually $d_{Z_0}(j)=D_0$ for all $i$ and lifts $j$,
where $D_0=D_0(D_{x,i})\in\N$ is a fixed rank attached to the datum.
\end{definition}

\begin{definition}[Selmer hypothesis]\label{i2:def-HypV}
For a number field $E$, $t\in\N$ and $B\in\R$, $\mathrm{VV}_\le(E,t,B)$ means: for every
algebraic closure $\Omega$ of $E$ and every choice of $l$-data $\mathrm{ld}(l)$ for all primes
$l$, (1) $\ell(\mathrm{Sha}(\mathrm{ld}(p),e))<\infty$ for all $p$ and $0<e<t$, and (2) for every
finite set $S$ of primes,
$\sum_{l\in S}\log l\sum_{e=1}^{t-1}\ell(\mathrm{Sha}(\mathrm{ld}(l),e))\le B$.
$\mathrm{HypV}(\mathcal F)$: eventually $\mathrm{VV}_\le(E_{D_{x,i}},t,\sqrt{b})$ for all $i\in I(x)$.
\end{definition}

The hypothesis $\mathrm{HypV}$ is supplied for the families actually used by the windowed
Selmer bound of Section~\ref{i5:sec}; it is used here only as a hypothesis.

\begin{theorem}[divisor bound]\label{i2:thm-div}
Let $\mathcal F$ be a family with exponent $\varepsilon>0$ satisfying
$\mathrm{HypV}(\mathcal F)$, $\mathrm{LinkN}(\mathcal F)$ and
$\mathrm{SmallLevelRankEq}(\mathcal F)$. Then there is $x_0$ such that for all
$x\ge x_0$, $i\in I(x)$ and lifts $j$,
\[
 \eta\Sigma_{L_{x,i}}(j)\le D_0\,(|C_N|+1)\,x^{\varepsilon}\,(B_0+N_c).
\]
\end{theorem}

\subsection{Reduction to the local index bound}\label{i2:ssec-red}

\begin{lemma}[Selmer record]\label{i2:lem-record}
If $\mathrm{VV}_\le(E_D,t,B)$, then $D$ is Sha-finite, $(\mathrm{vv}_D(l))_l$ is summable,
$\mathrm{vv}_D\ge0$, and $\sum'_l\mathrm{vv}_D(l)\le B$.
\end{lemma}

\begin{proof}
Apply $\mathrm{VV}_\le$ to the closure and the $l$-data of $D$. Clause (1) is
Sha-finiteness; by it and \eqref{i2:eq-sha},
$\mathrm{VV}_D(p)=\sum_{e<t}\ell(\mathrm{Sha}(\mathrm{ld}_D(p),e))$ for every prime $p$. Hence
$\mathrm{vv}_D\ge0$, and by clause (2) every finite partial sum of $\mathrm{vv}_D$ (non-primes
contribute $0$) is at most $B$. A nonnegative sequence with partial sums bounded by $B$ is
summable with sum $\le B$.
\end{proof}

\begin{lemma}[Pointwise divisor bound]\label{i2:lem-pointwise}
Let $L$ be a size datum of $D$, $j\in\{1,\dots,m\}$, $C\in\R$, and assume
$\mathrm{LinkN}_{\rm at}(L,j,C)$, $\mathrm{vv}_D$ summable with $\sum'_l\mathrm{vv}_D(l)\le b$,
$d_{Z_0}(j)=D_0$ and $(\log x)^2\le x^\varepsilon$. Then
$\eta\Sigma_L(j)\le D_0(|C|+1)x^\varepsilon(B_0+N_c)$.
\end{lemma}

\begin{proof}
Put $y=x^\varepsilon\ge1$; then $\mathsf L=y$, $fmby^2=yB_0$, $N_c\ge0$ and $(\log x)^2N_c\le yN_c$.
The bracket $yB_0+(\log x)^2N_c$ is nonnegative, so
$C(yB_0+(\log x)^2N_c)\le|C|\,y(B_0+N_c)$. Since $fmy\ge1$ we have $b\le B_0\le y(B_0+N_c)$,
so $\sum'\mathrm{vv}_D\le y(B_0+N_c)$. Insert both bounds in $\mathrm{LinkN}_{\rm at}$ and use
$D_0\ge0$.
\end{proof}

\begin{proof}[Proof of Theorem~\ref{i2:thm-div}]
Since $\log x=o(x^{\varepsilon/2})$, we have $(\log x)^2\le x^\varepsilon$ for large $x$. Let
$x$ exceed this threshold, $x_1+1$, and the thresholds of the three hypotheses. By
$\mathrm{HypV}$ and Lemma~\ref{i2:lem-record}, $D_{x,i}$ is Sha-finite, $\mathrm{vv}$ is
summable, and $\sum'\mathrm{vv}\le\sqrt b\le b$ (as $b\ge1$). Sha-finiteness makes
$\mathrm{LinkN}_{\rm at}(L_{x,i},j,C_N)$ available, the rank equality gives $d_{Z_0}(j)=D_0$,
and Lemma~\ref{i2:lem-pointwise} concludes.
\end{proof}

\begin{remark}\label{i2:rem-empty}
The Selmer hypothesis is used twice: to discharge the Sha-finiteness premise and the
summability guard of $\mathrm{LinkN}$, and to bound the resulting Selmer sum. A member may
have empty sets $\mathrm{locs}(j,w)$: this happens exactly when $V_{j,w}=0$ (the Hecke
algebra contains the identity, so it is the zero ring only if $V_{j,w}=0$, and a nonzero ring
has a maximal ideal), and then $\eta_j(w)=0$ is the true value. The bound is only ever used
as an upper bound, so no lower bound on the local indices is needed.
\end{remark}

\subsection{The per-prime bound and its summation}\label{i2:ssec-perprime}

For a prime $l$ let $\mathrm{tw}_L(l)$ be the set of twisted $w\in\mathrm{supp}\,L$ with
$p(w)=l$, and $\eta_L(j,l)=\sum_{w\in\mathrm{tw}_L(l)}\lambda_we_w\eta_j(w)$; thus
$\eta\Sigma_L(j)=\sum_l\eta_L(j,l)$, a finite sum.

\begin{theorem}[Neukirch, \emph{Algebraic Number Theory} (1999), I (8.2)]\label{i2:thm-fund}
Let $K$ be a number field and $l$ a prime with $lO_K=\prod_{i}\mathfrak P_i^{e_i}$,
$f_i=[O_K/\mathfrak P_i:\mathbb F_l]$. Then $\sum_ie_if_i=[K:\Q]$.
\end{theorem}
\noindent It applies to $L_D$ and to $E$, which are number fields.

\begin{lemma}[Weights over $l$]\label{i2:lem-meanweight}
For every prime $l$, $\sum_{w\in\mathrm{tw}_L(l)}\lambda_we_w\le\log l$. For every finite
$N\ne0$, $\sum_{w\ \rm twisted}\lambda_we_w\,v_{p(w)}(N)\le\log N$.
For every prime $l$, $\sum_{\lambda\mid l}n_\lambda\le2f$ and $\sum_{\lambda\mid l}(n_\lambda+1)\le4f$,
the sums over the places $\lambda$ of $E$ above $l$ with local degrees $n_\lambda=e_\lambda f_\lambda$.
\end{lemma}

\begin{proof}
For $w\mid l$, $\mathrm N(w)=l^{f_w}$, so $\lambda_we_w=\log l\cdot e_wf_w/[L_D:\Q]$; the
places of $\mathrm{tw}_L(l)$ are distinct primes above $l$, so Theorem~\ref{i2:thm-fund}
gives the first claim. The second follows by grouping by $l=p(w)$ and
$\log N=\sum_lv_l(N)\log l$. The third is Theorem~\ref{i2:thm-fund} for $E$, $[E:\Q]=2f$, with
$n_\lambda+1\le2n_\lambda$.
\end{proof}

\medskip

The local information at $l$ is packaged in a \emph{table}. We first define the per-prime
budgets.

\begin{definition}[Budgets]\label{i2:def-budgets}
The size datum supplies reals $b\ge1$, $\mathsf L\ge0$, $S_c>0$, $\eta_{\rm sz}>0$, $b'>0$,
$w''\ge0$, thresholds $l_A<l^*$, and the host supplies $L_w>0$, a \emph{class} for each prime
(split, inert -- comprising two subclasses -- or one of the special classes for
$l\in\mathrm{small}$, $l\in\{q,d\}$), and
$\mathrm{ord}_H(l)=\min\{o>0:l^o\equiv\pm1\bmod 2^{e+2}\}$. For a constant $C$ put
$h_B=b+S_c+\eta_{\rm sz}+t+2$, $e_P(l)=C(1+h_B/l+\log h_B/\log l)$, $\eta_3=\eta_{\rm sz}+t+2$,
$w_B=\max\bigl(w''+\eta_{\rm sz},\sup_{t'\in T}b_{t'}+b'/x+\eta_{\rm sz}\bigr)$, $w_P=w_B+(t+1)/2$,
$G_P=w_P+S_c+\eta_3$, $j_A(l)=\max(1,\lceil\log l_A/\log(l-1)\rceil)$, and
\begin{align*}
\mathrm{rowBr}(l)&=h_B/(l-1)^{j_A(l)+1}+(G_P+1)(j_A(l)+1)+\mathsf L,\\
\mathrm{flatB}(l)&=(w_P+\eta_3)/l+S_c+\mathsf L+\eta_3I_{\rm band}(l),
\end{align*}
where $I_{\rm band}(l)=1$ iff some $s\in[w'',w_B]$ has
$|ls-b|\le C(l\eta_3+lS_c+b/l+s+l\mathsf L)$, and $0$ otherwise. Further
\[
\mathrm{lenB}(l)=\begin{cases}(2fm+f)e_P(l),& l\ \text{split},\\
4fCr\,\mathrm{rowBr}(l),& l\ \text{inert},\ l<l_A,\ 4L_w<\mathrm{ord}_H(l),\\
4fCr\bigl(h_B/l^2+\mathrm{flatB}(l)\bigr)+2fCm\,e_P(l),&\text{otherwise,}\end{cases}
\]
$\mathrm{bbB}=Cfnt^2n_V\mathsf L^2$, and, testing the cases in order,
\[
\lambda^B(l)=\begin{cases}Cfmb\log x+Cx^C\mathsf L^2,& l\in\mathrm{small}\ \text{or}\ l=3,\\
2Cfmb,& l\in\{q,d\},\\ 0,& l>l^*,\ l\notin\mathrm{exc},\\ \mathrm{bbB}+\mathrm{lenB}(l),&\text{otherwise;}\end{cases}
\qquad
\lambda^G(l)=\begin{cases}Cx^C\eta_{\rm sz},& l\in\{g_0,g_1\},\\ \lambda^B(l),&\text{otherwise.}\end{cases}
\]
Here $\mathrm{exc}$ is a finite set of exceptional primes, part of the table below. The
\emph{base exceptional set} of $L$ is
\[
\mathrm{Exc}(L)=\bigl\{l\ \text{prime}:\ l>l^*,\ l\mid n\ \text{for some integer}\ n\in[W-\eta_{\rm sz}-2s_+-3,\ W+S_c+2\eta_{\rm sz}+2s_++4]\ \text{and some}\ W\in\mathcal W\bigr\},
\]
where $s_+=\lfloor2b/l^*\rfloor$ and $\mathcal W=\{b,\ w''\}\cup\{b_{t'}+\lfloor b'/x\rfloor:t'\in T\}$ is the set of
base sizes ($w''$, $b_{t'}$, $b'$ as above).
\end{definition}

For $C>0$ and $l\ge2$ all of $h_B,e_P,\mathrm{flatB},\mathrm{rowBr},\mathrm{lenB},\mathrm{bbB},\lambda^B$ are
nonnegative.

\begin{definition}[Admissible tables]\label{i2:def-AdmT}
A \emph{table} for $L$ consists of reals $s_i(\pi)$, local lengths $\mathrm{len}_i(\lambda)\ge0$ at
places $\lambda$ of $E$, reals $\mathrm{BB}_i$ (for twisted local indices $i$ and
$\pi\in Z^{\rm good}_i$), a band predicate $\mathrm{band}_i(\lambda)$, the set $\mathrm{exc}$, and
reals $\mathrm{gc}_j(l),\mathrm{gs}_j(l)$ at the generator primes. It is \emph{admissible with
constant $C$}, $\mathrm{AdmT}(L,\mathrm{tb},C)$, if the following hold, for twisted $i$ with
$l=l_i$ (and $l\notin\{g_0,g_1\}$ unless said otherwise), $\pi\in Z^{\rm good}_i$ and lifts $j$:
\begin{enumerate}
\item[(adm1)] $\eta^{\rm loc}_i\le\sum_{\pi\in Z^{\rm good}_i}\mathrm{mult}_i(\pi)\bigl(v^{cc}_i(\pi)+v_l(|G|)\bigr)$;
\item[(adm2)] $\sum_{i\in\mathrm{locs}(j,w)}\sum_\pi\mathrm{mult}_i(\pi)\le\mathrm{rk}_j(w)$ at twisted $w$ with $p(w)\notin\{g_0,g_1\}$;
\item[(adm3)] $\mathrm{rk}_j(w)=d_{Z_0}(j)$ at twisted $w$;
\item[(ASM)] $v^{cc}_i(\pi)\le s_i(\pi)+\sum_{\lambda\mid l}\mathrm{len}_i(\lambda)+\mathrm{VV}_D(l)+v_l(h_E)+\mathrm{BB}_i+l^{-2}$,
  $h_E$ the class number of $E$;
\item[(adm5)] if $l\notin\mathrm{small}\cup\{q,d\}$: $\mathrm{BB}_i\le\mathrm{bbB}$;
\item[(adm6)] $s_i(\pi)+l^{-2}\le\nu_l+2l^{-2}$ (any twisted $i$), with $\nu_l$ as in Lemma~\ref{i2:lem-nu};
  and $\mathrm{gs}_j(l)+l^{-2}\le2l^{-2}$ for $l\in\{g_0,g_1\}$;
\item[(A)] if $l\ge5$ is split and ($l\le l^*$ or $l\in\mathrm{exc}$):
  $2\sum_{\lambda\mid l}\mathrm{len}_i(\lambda)\le\sum_{\lambda\mid l}(2n_\lambda m+1)e_P(l)$;
\item[(B)] if $l$ is inert, $l<l_A$, $4L_w<\mathrm{ord}_H(l)$: $\mathrm{len}_i(\lambda)\le C(n_\lambda+1)r\,\mathrm{rowBr}(l)$;
\item[(Med)] if $l$ is inert, $l\ge l_A$, and ($l\le l^*$ or $l\in\mathrm{exc}$); or (SmI) $l$ inert,
  $l<l_A$, $\mathrm{ord}_H(l)\le4L_w$:
  $\mathrm{len}_i(\lambda)\le C(n_\lambda+1)r\bigl(h_B/l^2+\mathrm{flatCD}_i(\lambda,l)\bigr)+Cn_\lambda m\,e_P(l)$,
  where $\mathrm{flatCD}_i(\lambda,l)$ is $\mathrm{flatB}(l)$ with $I_{\rm band}(l)$ replaced by $[\mathrm{band}_i(\lambda)]$;
\item[(Band)] if $\mathrm{band}_i(\lambda)$ for some $\lambda\mid l$, then $I_{\rm band}(l)=1$;
\item[(Van)] if $l>l^*$, $l\notin\mathrm{exc}$: $\mathrm{len}_i\equiv0$ above $l$ and $\mathrm{BB}_i=0$;
\item[(Sm)] if $l\in\mathrm{small}$ or $l=3$: $\sum_{\lambda\mid l}\mathrm{len}_i(\lambda)+\mathrm{BB}_i\le Cfmb\log x+Cx^C\mathsf L^2$;
\item[(Ram)] if $l\in\{q,d\}$: $\sum_{\lambda\mid l}\mathrm{len}_i(\lambda)+\mathrm{BB}_i\le2Cfmb$;
\item[(gen)] for $l\in\{g_0,g_1\}$:
  $\eta_L(j,l)\le\log l\cdot d_{Z_0}(j)\bigl(\mathrm{gc}_j(l)+\mathrm{gs}_j(l)+\mathrm{VV}_D(l)+v_l(h_E)+l^{-2}\bigr)$
  and $\mathrm{gc}_j(l)\le Cx^C\eta_{\rm sz}$;
\item[(Par)] parameter bounds: $0<C_{Lw}\le C$ and $0<C_{YL}\le C$ for two auxiliary positive
  constants $C_{Lw},C_{YL}$ of the parameter choice that accompanies the size datum (so $C>0$); $k,n,t\le x^C$; $x^{-\varepsilon}\le\delta_2$,
  where $\delta_2=|T_2|/n\in(0,1]$ is the density of the thick type set; $\log h_E\le Cx^C$; and
  $\log|G|\le Cx^C\mathsf L$. Nothing else (see Remark~\ref{i2:rem-par}; the size facts used with (Par)
  in Proposition~\ref{i2:prop-point}, such as $Q_S\le n^2$, $|T|\le2n$, $\log q,\log d\le\mathsf L$ and the formulas for
  $l^*,l_A,S_c$, come from the host axioms and the parameter choice accompanying the size datum, not from (Par));
\item[(exc)] $|\mathrm{exc}|\le(|T|+2)\bigl(S_c+3\eta_{\rm sz}+t+8b/l^*+9\bigr)$ and $l\le3b$ for
  $l\in\mathrm{exc}$, $|T|$ the number of tops. (Under (Par) and $\mathrm{Large}$ this gives
  $|\mathrm{exc}|\le388U^8$, $U=x^{C+2}e^{x^{\varepsilon/4}}$; see the proof of Proposition~\ref{i2:prop-point}.)
\end{enumerate}
$D$ is \emph{admitted with constant $C$} if Sha-finiteness of $D$ implies the existence of an
admissible table with constant $C$.
\end{definition}

\begin{remark}[What (Par) does and does not contain]\label{i2:rem-par}
The bound ``the ranks and the
denominators of the size data are at most $10U^7$'', $U=x^{C+2}e^{x^{\varepsilon/4}}$, and the size
relations among $l^*,l_A,b',S_c$ are not clauses of (Par): they are not
hypotheses but consequences of (Par), of $\mathrm{Large}(x,\varepsilon,C)$
(Definition~\ref{i2:def-Large}) and of the definitions of the thresholds,
\[
l^*=\frac{mb}{\mathrm{Den}},\qquad l_A=\frac{k\mathsf L}{m\delta_2},\qquad
\mathrm{Den}=nt^2(N+t)\mathsf L^2+kn^2(S_c+t+\mathsf L),\quad N=k+m(t+1),\quad \mathsf L=x^\varepsilon .
\]
Precisely: from $k,n,t\le x^C\le U$, $m+3\le k$, $r=k+m-1$, $f=2n$ and $n_V=r+t$ one gets
$n,k,t,m\le U$, $f,r\le2U$, $N+t\le4U^2$, $n_V\le4U^2$; with $\mathsf L\le x\le U$ and $S_c\le4U$
(from the definition of $S_c$ in the size datum and $k\le x^C$) this gives
$\mathrm{Den}\le U\cdot U^2\cdot4U^2\cdot U^2+U\cdot U^2\cdot6U\le10U^7$, since $U\ge1$.
The relations used in the proof of Proposition~\ref{i2:prop-point} ($\log l^*\le\mathsf L$,
$l^*\le b/\mathsf L$, $b'<l^*$, $l_A<l^*$, $r/l_A\le2m/\mathsf L$, and the upper bound for $l_A$,
which uses $x^{-\varepsilon}\le\delta_2$) are derived in the same way from these formulas, (Par) and
$\mathrm{Large}$. So the only quantity named ``$10U^7$'' is the bound $\mathrm{Den}\le10U^7$, a lemma,
and ``ranks $\le10U^7$'' should be read as the weaker statement implied by the bounds just listed.
\end{remark}

Condition (ASM) is the ``local index plus Selmer'' structure: at $l$ the congruence length
is bounded by local terms plus the $l$-part $\mathrm{VV}_D(l)$ of the Selmer length.

\begin{lemma}[Local length budget]\label{i2:lem-lamB}
Under $\mathrm{AdmT}(L,\mathrm{tb},C)$, for twisted $i$ with $l=l_i\notin\{g_0,g_1\}$,
$\sum_{\lambda\mid l}\mathrm{len}_i(\lambda)+\mathrm{BB}_i\le\lambda^B(l)$.
\end{lemma}

\begin{proof}
Follow the cases of $\lambda^B$. Small $l$ or $l=3$: (Sm). Otherwise $l\ge5$ (as
$2\in\mathrm{small}$). $l\in\{q,d\}$: (Ram). $l>l^*$, $l\notin\mathrm{exc}$: (Van). In the last
case $\mathrm{BB}_i\le\mathrm{bbB}$ by (adm5), and $l$ is split or inert. Split: by (A) and
Lemma~\ref{i2:lem-meanweight}, $2\sum\mathrm{len}_i\le(2m+1)e_P(l)\sum n_\lambda\le(2m+1)e_P(l)2f$.
Inert with $l<l_A$, $4L_w<\mathrm{ord}_H(l)$: (B) and $\sum(n_\lambda+1)\le4f$. Otherwise
$\mathrm{flatCD}_i\le\mathrm{flatB}$ (by (Band) when the band predicate holds, and because a
nonnegative term is dropped otherwise), and (Med) or (SmI) summed over $\lambda$ gives
$\mathrm{lenB}(l)$.
\end{proof}

\begin{lemma}[Per-prime bound]\label{i2:lem-etaAt}
Let $j$ be a lift, $l\notin\{g_0,g_1\}$, $U\ge0$, assume (adm1)--(adm3), and
$v^{cc}_i(\pi)+v_l(|G|)\le U$ for twisted $i$ with $l_i=l$ and $\pi\in Z^{\rm good}_i$. Then
$\eta_L(j,l)\le\log l\cdot d_{Z_0}(j)\,U$.
\end{lemma}

\begin{proof}
For $w\in\mathrm{tw}_L(l)$, every $i\in\mathrm{locs}(j,w)$ is twisted with $l_i=l$; by (adm1)
and (adm2), (adm3),
$\eta_j(w)\le U\sum_{i}\sum_\pi\mathrm{mult}_i(\pi)\le U\,\mathrm{rk}_j(w)=U\,d_{Z_0}(j)$. Multiply
by $\lambda_we_w\ge0$, sum, and apply Lemma~\ref{i2:lem-meanweight}.
\end{proof}

\begin{lemma}\label{i2:lem-logl2}
For every finite set $T$ of integers $\ge2$, $\sum_{l\in T}\log l/l^2\le2$.
\end{lemma}

\begin{proof}
Let $\phi(u)=(\log u+2)/u$. For $u,M\ge1$ one has
$\log(u+M)/(u+M)^2\le(\phi(u)-\phi(u+M))/M$: with $v=u+M$ and
$X=u\log u+M-u\log v\ge0$ (from $\log(v/u)\le M/u$), the claim is equivalent to
$Mu\log v\le v(v\log u-u\log v+2M)$, and the difference of the two sides is
$(u+2M)X+Mu+M^2\log u\ge0$. Telescoping with $u=M=1$ gives
$\sum_{2\le l\le J+1}\log l/l^2\le\phi(1)=2$.
\end{proof}

\begin{lemma}[Cyclotomic order terms]\label{i2:lem-nu}
For a prime $l$ let $\nu_l$ be the least $k\ge0$ such that some $g\in\mathrm{Gal}(\Omega_D/E)$
has $l^{k+1}\nmid\chi_l(g)^{tr}-1$ ($\chi_l$ the $l$-adic cyclotomic character; $\nu_l=0$ if no
such $k$ exists). Then $\sum_{l\in\mathcal S}\nu_l\log l\le\nu(D):=\log Q_S+\log(tr)$ for
every finite set $\mathcal S$ of primes.
\end{lemma}

\begin{proof}
Put $n=tr$ (odd, as $t$ is odd and $r=k+m-1$ with $k,m$ even) and $N=Q_Sqd$. The host
axioms give an embedding $E\subseteq\Q(\zeta_N)$ and a map
$\mathrm{art}:(\Z/N)^\times\to\mathrm{Gal}(E/\Q)$ through which $\zeta_N\mapsto\zeta_N^v$ acts on
$E$, with $\mathrm{art}(v)^2=1$ if $v\equiv1\bmod Q_Sd$, $\mathrm{art}(v)\in\{1,\iota\}$ if
$v\equiv1\bmod Q_Sq$ ($\iota$ complex conjugation), and $q,d>4tr+1$.

(a) If $u\in\Z_l^\times$ is approximated to every order by values $\chi_l(g)$,
$g\in\mathrm{Gal}(\Omega/E)$, and $l^{s+1}\nmid u^n-1$, then $\nu_l\le s$.
(b) If $\mathrm{art}(v)=1$ and $u\equiv v\bmod l^{v_l(N)}$, then $u$ is so approximated: by the
Chinese remainder theorem pick $a$ modulo $\mathrm{lcm}(N,l^k)$ with $a\equiv v\bmod N$,
$a\equiv u\bmod l^k$; by irreducibility of cyclotomic polynomials there is $\sigma$ with
$\sigma\zeta=\zeta^a$; it fixes $E$ and $\chi_l(\sigma)\equiv u\bmod l^k$.
(c) $l\notin\{q,d\}$: take $v=1$. If $l\mid Q_S$, $s=v_l(Q_S)$, $u=1+l^s$: by the
lifting-the-exponent lemma $v_l(u^n-1)=s+v_l(n)$ ($v_2(n)=0$), so
$\nu_l\le v_l(Q_S)+v_l(n)$. If $l\nmid Q_S$, $u=-1$ gives $u^n-1=-2$, a unit at $l$ (as
$2\mid Q_S$), so $\nu_l=0$.
(d) $l\in\{q,d\}$: let $u$ lift a generator of $\mathbb F_l^\times$; $l-1\nmid2n$ because
$l>4n+1$, so $l\nmid u^{2n}-1$. Choose $v\equiv1\bmod N/l$, $v\equiv u\bmod l$. Then
$\mathrm{art}(v^2)=1$ (for $l=q$ from $\mathrm{art}(v)^2=1$; for $l=d$ from
$\mathrm{art}(v)\in\{1,\iota\}$), and (a), (b) with $(v^2,u^2)$ give $\nu_l=0$.
(e) Summing $\nu_l\le v_l(Q_S)+v_l(n)$ against $\log l$ gives the claim.
\end{proof}

\begin{theorem}[Summed local bound]\label{i2:thm-summed}
Assume $\mathrm{AdmT}(L,\mathrm{tb},C)$, $\mathrm{vv}_D$ summable, and let $j$ be a lift and
$\mathcal S=\{p(w):w\in\mathrm{supp}\,L\ \text{twisted}\}$. Then
\[
\eta\Sigma_L(j)\le d_{Z_0}(j)\Bigl(6+\nu(D)+\log h_E+\log|G|+\sum_{l\in\mathcal S}\log l\,\lambda^G(l)\Bigr)+d_{Z_0}(j)\textstyle\sum'_l\mathrm{vv}_D(l).
\]
\end{theorem}

\begin{proof}
Put $U(l)=\nu_l+2/l^2+\lambda^G(l)+\mathrm{VV}_D(l)+v_l(h_E)+v_l(|G|)$. We claim
$\eta_L(j,l)\le\log l\cdot d_{Z_0}(j)U(l)$ for $l\in\mathcal S$. At $l\in\{g_0,g_1\}$ this is
(gen) with (adm6). Otherwise, by (ASM), (adm6) and Lemma~\ref{i2:lem-lamB},
\[
v^{cc}_i(\pi)\le\bigl(\nu_l+2l^{-2}\bigr)+\lambda^B(l)+\mathrm{VV}_D(l)+v_l(h_E),
\]
and Lemma~\ref{i2:lem-etaAt} applies with $U=U(l)$. Summing over $l\in\mathcal S$: the
$\nu_l$ and $2/l^2$ parts give at most $\nu(D)+4\le\nu(D)+6$ (Lemmas~\ref{i2:lem-nu},
\ref{i2:lem-logl2}; the constant $6$ of the statement is the one recorded in the formalization); the $\mathrm{VV}$ part is a finite partial sum of the nonnegative summable
sequence $\mathrm{vv}_D$; and $\sum_lv_l(N)\log l\le\log N$ for $N=h_E,|G|$.
\end{proof}

\begin{definition}[Large scales]\label{i2:def-Large}
$\mathrm{Large}(x,\varepsilon,C)$ means: $x\ge100$; $8x\le e^{x^{\varepsilon/3}}$; $4\le e^{x^\varepsilon}$;
$x^{C+2}e^{x^{\varepsilon/3}}\le e^{x^\varepsilon}$; $48C\le x^{1/2}$;
$(C+3)^2(1+\log x)^2\le x^\varepsilon$; $x^{20C+40}e^{20x^{\varepsilon/4}}\le e^{x^{\varepsilon/3}}$;
$10^8(C+1)^3\le x$. Each clause holds for all large $x$ (each reduces to $K\log x\le x^r/2$
eventually, for fixed $K$ and $r>0$, or to a power of $x$ being large), hence so does
$\mathrm{Large}$.
\end{definition}

\begin{proposition}[Point bound]\label{i2:prop-point}
Let $C>0$, $\mathrm{Large}(x,\varepsilon,C)$ and $\mathrm{AdmT}(L,\mathrm{tb},C)$. For every finite
set $\mathcal S$ of primes,
\[
6+\nu(D)+\log h_E+\log|G|+\sum_{l\in\mathcal S}\log l\,\lambda^G(l)\le2000(C+1)^3fmb\,\mathsf L^2.
\]
\end{proposition}

\begin{proof}
Put $X=fmb\mathsf L^2$ and $U=x^{C+2}e^{x^{\varepsilon/4}}$. By $\mathrm{Large}$,
$\mathsf L,x^C,\eta_{\rm sz}\le U$ and
$4\cdot10^6(C+1)^2U^{18}\le x^{20C+40}e^{20x^{\varepsilon/4}}\le e^{x^{\varepsilon/3}}\le b\le X$.
Split $\log l\,\lambda^G(l)$ according to the case of $\lambda^G$ into five sub-sums.
\emph{Generator primes:} $\log g_k\le(C+1)U$ and $Cx^C\eta_{\rm sz}\le(C+1)U^2$, so the sum is
$\le2(C+1)^2U^3$.
\emph{Small primes and $3$:} the small primes divide $Q_S\le n^2\le x^{2C}$, so
$\sum\log l\le w:=2C\log x+2$, and $w\le\mathsf L$, $w\log x\le\mathsf L^2$ by the sixth clause of
$\mathrm{Large}$; with $x^C\mathsf L\le b$ the sum is $\le2CX$.
\emph{$q,d$:} $\log q,\log d\le\mathsf L$ (the conductor of $E$ is at most $b$), so the sum is
$\le4CX$.
\emph{$5\le l\le l^*$:} using the consequences of (Par) and $\mathrm{Large}$ recorded in Remark~\ref{i2:rem-par} ($\log l^*\le\mathsf L$, $l^*\mathsf L\le b$, $h_B\le2b$,
$r/l_A\le2m/\mathsf L$, $(3+\log l_A)^2\le\mathsf L$, $rb'l_A\log x\le mb$) and the elementary
estimates $\sum_{l\le Y}\log l\le Y\log Y$, $\sum_{l\le Y}\log l/l\le(1+\log Y)^2$,
$\sum_{2\le l\le Y}\log l/(l-1)\le2(1+\log Y)^2$, $\sum_{l\ge A}\log l/l^2\le4(\log A+2)/A$, one
bounds the contributions of $\mathrm{bbB}$, of the $e_P$ terms, of $\mathrm{rowBr}$, of the
$h_B/l^2$ terms (split into $l\ge l_A$ and the primes of small order modulo a power of $2$),
and of $\mathrm{flatB}$, by $CX$, $11C(3+2C)X$, $80CX$, $(64+32)CX$ and
$4C(102+4(2C+1)^2)X$; their total is $(64C^3+86C^2+634C)X\le1000(C+1)^3X$.
\emph{Exceptional primes $>l^*$:} first, $|\mathrm{exc}|\le388U^8$: by (exc), with $|T|\le2n\le2U$ (parameter
choice), $S_c\le4U$, $\eta_{\rm sz}\le U$, $t\le U$ and $8b/l^*=8\mathrm{Den}/m\le80U^7$ (Remark~\ref{i2:rem-par}), one has
$|T|+2\le4U$ and $S_c+3\eta_{\rm sz}+t+8b/l^*+9\le97U^7$. By (exc), $l\le3b$; each term
$\log l\,(\mathrm{bbB}+\mathrm{lenB}(l))$ is $\le3000(C+1)^2U^{10}$, and there are $\le388U^8$ of them: total
$\le1164000(C+1)^2U^{18}$.
\emph{Global terms:} by (Par), $6+\nu(D)+\log h_E+\log|G|\le20(C+1)^2U^{18}$.
Since all the $U^{18}$-terms together are $\le b\le X$, the left side is at most
$(1+12C+1000(C+1)^3)X\le2000(C+1)^3X$.
\end{proof}

\begin{theorem}[Admissible families satisfy $\mathrm{LinkN}$]\label{i2:thm-export}
Let $\mathcal F$ be a family such that, for large $x$, every $D_{x,i}$ is admitted with constant
$C_a$. Then $\mathrm{LinkN}(\mathcal F)$ holds with $C_N=2000(|C_a|+1)^3$.
\end{theorem}

\begin{proof}
For large $x$, $\mathrm{Large}(x,\varepsilon,C_a)$ holds. If $D_{x,i}$ is Sha-finite and
$\mathrm{vv}$ summable, pick an admissible table; (Par) gives $C_a>0$.
Theorem~\ref{i2:thm-summed} and Proposition~\ref{i2:prop-point} give the bound with
$2000(C_a+1)^3fmb\mathsf L^2\le C_N\bigl(fmb\mathsf L^2+(\log x)^2N_c\bigr)$.
\end{proof}

\subsection{Local conclusions and cuts}\label{i2:ssec-cuts}

Admissible tables are produced from \emph{local conclusions}, which are stated in the
form in which the local theory delivers them.

\begin{definition}[Local conclusions]\label{i2:def-Concl}
\emph{Hub data} $X$ for $D$ consist of local lengths $s^X_i(\pi)\ge0$,
$\mathrm{len}^0_i(\pi,\lambda)$, error terms $\mathrm{err}_i(\pi)$, a band predicate and a finite
set $\mathrm{exc}$. For twisted $i$, $l=l_i$, $\pi\in Z^{\rm good}_i$ and $\beta\in\R$ write
\[
\mathrm{Bd}_i(\pi,\beta):\qquad v^{cc}_i(\pi)\le s^X_i(\pi)+\beta+\mathrm{VV}_D(l)+v_l(h_E)+l^{-2}.
\]
Let $\mathrm{RangeAD}(l)$ mean: $l\notin\mathrm{small}\cup\{3,q,d\}$ and not ($l>l^*$ and
$l\notin\mathrm{exc}$). $\mathrm{Concl}(L,X,C)$ is the conjunction of: $C\ge0$; (adm1);
(adm2) at all twisted $w$ (``$\mathrm{MultLe}$''); (adm3); (kill)
$s^X_i(\pi)\le\nu_{l}+l^{-2}$; (num) conditions (A), (B), (Med), (SmI), (Band) for the
table with $\mathrm{len}=\sup_\pi\mathrm{len}^0_i(\pi,\cdot)$; (mid) for
$l\notin\{g_0,g_1\}$ with $\mathrm{RangeAD}(l)$:
$\mathrm{Bd}_i(\pi,\sum_{\lambda\mid l}\mathrm{len}^0_i(\pi,\lambda)+\mathrm{err}_i(\pi))$ and
(bb) $\mathrm{err}_i(\pi)\le\mathrm{bbB}$; (large) for $l>l^*$, $l\notin\{g_0,g_1\}\cup\mathrm{exc}$:
$\mathrm{Bd}_i(\pi,0)$; (exc) as in Definition~\ref{i2:def-AdmT}; (two) a table $\mathrm{tf}$ satisfying (Sm) with
$\mathrm{Bd}_i(\pi,\sum_{\lambda\mid l}\mathrm{len}^{\rm tf}_i(\lambda)+\mathrm{BB}^{\rm tf}_i)$ for
$l\in\mathrm{small}\cup\{3\}$; (ram) for $l\in\{q,d\}$: $\mathrm{Bd}_i(\pi,\mathrm{midG}(l))$ with
some $\mathrm{midG}(l)\le2Cfmb$; (typ) for $l\in\{g_0,g_1\}$:
$\eta_L(j,l)\le\log l\cdot d_{Z_0}(j)\bigl(Cx^C\eta_{\rm sz}+2l^{-2}+\mathrm{VV}_D(l)+v_l(h_E)\bigr)$;
(Par); (small) $3,q,d$ and all small primes are $\le l^*$, and $3$ is not inert.
$\mathrm{ConclF}(\mathcal F)$: there are $C,x_0$ such that for $x\ge x_0$ and Sha-finite
$D_{x,i}$ some hub data satisfy $\mathrm{Concl}(L_{x,i},X,C)$.
\end{definition}

\begin{lemma}[Tables from local conclusions]\label{i2:lem-adm-of-concl}
$\mathrm{Concl}(L,X,C)$ implies $\mathrm{AdmT}(L,\mathrm{tb},C)$ for the table with
$s_i=s^X_i$, band and $\mathrm{exc}$ from $X$, $\mathrm{gc}_j(l)=Cx^C\eta_{\rm sz}$,
$\mathrm{gs}_j(l)=l^{-2}$, and (cases in order, $\sup\emptyset=0$)
\[
\mathrm{len}_i(\lambda)=\begin{cases}0,&l\in\mathrm{small},\\ \beta^{\rm tf}_i/\#\{\lambda'\mid3\},&l=3,\\
\sup_\pi\mathrm{len}^0_i(\pi,\lambda),&\mathrm{RangeAD}(l),\\0&\text{else,}\end{cases}\qquad
\mathrm{BB}_i=\begin{cases}\beta^{\rm tf}_i,&l\in\mathrm{small},\\ \mathrm{midG}(l),&l\in\{q,d\},\\
\sup_\pi\mathrm{err}_i(\pi),&\mathrm{RangeAD}(l),\\0&\text{else,}\end{cases}
\]
where $\beta^{\rm tf}_i=\sum_{\lambda\mid l}\mathrm{len}^{\rm tf}_i(\lambda)+\mathrm{BB}^{\rm tf}_i$.
\end{lemma}

\begin{proof}
(adm1)--(adm3) and (Par) are clauses of $\mathrm{Concl}$. For (ASM) it suffices to find $\beta$
with $\mathrm{Bd}_i(\pi,\beta)$ and $\beta\le\sum_\lambda\mathrm{len}_i(\lambda)+\mathrm{BB}_i$:
take $\beta=\beta^{\rm tf}_i$ for $l$ small or $l=3$ (clause (two); at $l=3$ the constant
$\beta/\#\{\lambda'\mid3\}$ summed over the places above $3$ is $\beta$, and $\mathrm{BB}_i=0$
because $3\notin\{q,d\}$), $\beta=\mathrm{midG}(l)$ for $l\in\{q,d\}$, the (mid) value for
$\mathrm{RangeAD}(l)$ (each term is at most its supremum), and $\beta=0$ otherwise (clause
(large)). (adm5) follows from (bb); (adm6) from (kill); (gen) is (typ); (Ram) and (Sm) are
the defining bounds; (Van) holds because, by (small), $l>l^*$ with $l\notin\mathrm{exc}$ forces
the last case. Conditions (A)--(Band) concern only split or inert $l$, never $l=3$; on
$\mathrm{RangeAD}$ rows they are (num), and elsewhere the length is $0$ and the right sides
are nonnegative.
\end{proof}

\begin{corollary}\label{i2:cor-concl-link}
$\mathrm{ConclF}(\mathcal F)$ implies $\mathrm{LinkN}(\mathcal F)$.
\end{corollary}

\begin{proof}
Lemma~\ref{i2:lem-adm-of-concl} and Theorem~\ref{i2:thm-export}.
\end{proof}

\medskip

The local conclusions are not available for the datum itself but only after discarding the
eigensystems that are not good. This is done by a cut.

\begin{definition}[Cuts]\label{i2:def-cut}
A \emph{cut} of $D$ keeps some local indices $i$ and, for kept $i$, a nonempty set
$X^c_i\subseteq X_i$. The cut datum $D^c$ has as local indices the kept $i$, with the Hecke
algebra restricted to $X^c_i$ acting on $M_i/M^{\rm out}_i$, where $M^{\rm out}_i$ is the
intersection of the kernels of the Hecke operators vanishing on $X_i\setminus X^c_i$,
$M^{\rm out}_i=\bigcap\{\ker f:f\in T_i,\ f|_{X_i\setminus X^c_i}=0\}$; everything else
(ranks, host, Selmer data, scale, $d_{Z_0}$) is unchanged. The cut is \emph{good} if $i$ is
kept iff $Z^{\rm good}_i\ne\emptyset$, $Z^{\rm good}_i\subseteq X^c_i$,
$Z_i\cap X^c_i\subseteq Z^{\rm good}_i$, and for kept $i$ there is an $O_i$-linear
endomorphism $\alpha_i$ of $M_i$ commuting with the Hecke algebra and acting on the
eigenspace $E_x$ by $|G|$ for $x\in Z^{\rm good}_i$ and by $0$ for $x\notin X^c_i$.
\end{definition}

\begin{lemma}[Congruence modules under a cut]\label{i2:lem-length-le}
Let $O$ be a principal ideal domain with fraction field $K$, $M$ a finitely generated
$O$-module, $X$ a finite set, $T\subseteq O^X$ a subalgebra separating points, with an $O$-algebra
homomorphism $T\to\mathrm{End}_O(M)$ (all that is used is that $f\in T$ acts on $E_x$ below by $f(x)$, which is the
definition of $E_x$, and that the zero function acts by $0$).
Put $V=K\otimes M$, $L$ the image of $M$, $E_x=\bigcap_{f\in T}\ker(f-f(x))\subseteq V$,
$\mathrm{mult}(x)=\dim E_x$, and for $Y\subseteq X$, $E_Y=\sum_{x\in Y}E_x$,
$M^{\rm up}(Y)=E_Y\cap(L+E_{X\setminus Y})$, $M^{\rm low}(Y)=L\cap E_Y$. Let $S\subseteq X$, form
$M^c=M/M^{\rm out}$ with $M^{\rm out}=\bigcap\{\ker f:f\in T,\ f|_{X\setminus S}=0\}$ (as in
Definition~\ref{i2:def-cut}), $Z\subseteq X$, $Y'=Z\cap S$. If $0\ne A\in O$ and
$\alpha\in\mathrm{End}_O(M)$ acts by $A$ on $E_x$ ($x\in Z$) and by $0$ on $E_x$ ($x\notin S$),
then
\[
\ell\bigl(M^{\rm up}(Z)/M^{\rm low}(Z)\bigr)\le\ell\bigl(M^{c,\rm up}(Y')/M^{c,\rm low}(Y')\bigr)+\Bigl(\sum_{x\in Z}\mathrm{mult}(x)\Bigr)\ell(O/A),
\]
and over a discrete valuation ring the left side is finite.
\end{lemma}

\begin{proof}
\emph{Idempotent-like elements.} Since $T$ separates points, products
$\prod_{y\ne x}(f^{(y)}-f^{(y)}(y))$ give $f_x\in T$ with $f_x(x)\ne0=f_x(y)$ for $y\ne x$.
With $d=\prod_xf_x(x)\ne0$ and $g_x=(\prod_{y\ne x}f_y(y))f_x$ we get $\sum_xg_x=d$,
$g_xV\subseteq E_x$, and $g_x E_y=0$ for $y\ne x$.
\emph{The projection.} Let $q:V\to V^c=K\otimes M^c$. Then $q(E_x)\subseteq E^c_x$ for
$x\in S$, $q(E_x)=0$ for $x\notin S$ ($f_x$ vanishes on $S$, so acts by $0$ on $M^c$, but by
$f_x(x)\ne0$ on $q(E_x)$), and $\ker q=E_{X\setminus S}$ (if $q(v)=0$ then $g_xv=0$ for
$x\in S$, so $dv=\sum_{x\notin S}g_xv$). Hence $q$ induces
$\Phi:U/U_0\to U'/U'_0$ between the two congruence modules, and
$\ell(U/U_0)\le\ell(\ker\Phi)+\ell(U'/U'_0)$.
\emph{$A$ kills $\ker\Phi$.} If $u\in U$ and $q(u)=q(1\otimes m_0)$, then
$u-1\otimes m_0\in E_{X\setminus S}$ is killed by $\alpha$, so $Au=\alpha u=1\otimes\alpha m_0\in L\cap E_Z=U_0$.
\emph{Rank.} For $u=l'+e\in U$ ($l'\in L$, $e\in E_{X\setminus Z}$),
$du=\sum_{x\in Z}g_xl'\in L$; so $U$ is finitely generated and torsion free, and the
preimage of $\ker\Phi$ is free of rank $\le\dim E_Z\le\sum_{x\in Z}\mathrm{mult}(x)$. A quotient
of $(O/A)^r$ has length $\le r\ell(O/A)$. Finiteness: $dU\subseteq U_0$.
\end{proof}

\begin{lemma}[Cut comparison]\label{i2:lem-etaJ}
For a good cut, a lift $j$ and a place $w$,
$\eta_j(w)\le\eta^c_j(w)+\bigl(\sum_{i\in\mathrm{locs}(j,w)}\sum_{\pi\in Z^{\rm good}_i}\mathrm{mult}_i(\pi)\bigr)v_{p(w)}(|G|)$.
\end{lemma}

\begin{proof}
Non-kept $i$ have $Z^{\rm good}_i=\emptyset$ and $\eta^{\rm loc}_i=0$. For kept $i$ apply
Lemma~\ref{i2:lem-length-le} with $S=X^c_i$, $Z=Z^{\rm good}_i$, $A=|G|$ and $\alpha=\alpha_i$,
and divide by $\ell(O_i/l_i)$ (the normalisation of $\eta^{\rm loc}$):
$\ell(O_i/|G|)/\ell(O_i/l_i)=v_{l_i}(|G|)$ because $|G|$ is an integer and $O_i$ is a discrete
valuation ring finite over $\Z_{l_i}$.
\end{proof}

\begin{lemma}[Size of the level group and class number]\label{i2:lem-logG}
Let $\mathcal F$ be a family and $C_0\ge1$ such that for all large $x$ and all $i$: $k,n,t\le x^{C_0}$,
$\log q,\log d\le x^{C_0}$, $Q_S\le n^2$, $\log g_0,\log g_1\le x$, $x^{-\varepsilon}\le\delta_2$, and the auxiliary
constants of (Par) satisfy $0<C_{Lw},C_{YL}\le C_0$ (these are the size conditions imposed where families are
constructed). Then there is $C_1\ge C_0$
such that for large $x$, $\log|G|\le C_1x^{C_1}\mathsf L$ and $\log h_E\le C_1x^{C_1}$; more
precisely (Par) holds with constant $C_1$.
\end{lemma}

\begin{proof}
\emph{Class number.} Every ideal class of a number field $K$ of degree $n$ contains an
integral ideal of norm at most $M_K\sqrt{|d_K|}$ with Minkowski's constant $M_K\le1$; the
number of integral ideals of norm $m$ is at most the divisor function $d_n(m)$, and
$\sum_{m\le Y}d_n(m)\le Y(1+\log Y)^{n-1}$. Hence
$\log h_E\le f\log\mathrm{rd}_E+2f\log(1+f\log\mathrm{rd}_E)$, and
$\log\mathrm{rd}_E\le2\log n+\tfrac12(\log q+\log d)$ (conductor--discriminant formula, $Q_S\le n^2$, and the
fact, from the host axioms, that $q$ and $d$ have ramification index $2$ in $E$). \emph{Level group.} $|\mathrm{GL}_r(\mathbb F_g)|\le g^{r^2}$, so
$\log|G|\le2fr^2\max_k\log g_k$, polynomial in $x$. The remaining clauses of (Par) are the size
conditions of the family, and enlarging the constant preserves them.
\end{proof}

\begin{theorem}[Transfer from a cut family]\label{i2:thm-transfer}
Let $\mathcal F^c$ be obtained from $\mathcal F$ by good cuts, and assume $\mathrm{MultLe}$
and (adm3) for $\mathcal F$ at twisted places for $x>x_1$. Then
$\mathrm{LinkN}(\mathcal F^c)\Rightarrow\mathrm{LinkN}(\mathcal F)$, with constant $C_N+C_1$.
\end{theorem}

\begin{proof}
By Lemma~\ref{i2:lem-etaJ}, $\mathrm{MultLe}$ and (adm3),
$\eta_j(w)\le\eta^c_j(w)+d_{Z_0}(j)v_{p(w)}(|G|)$ at twisted $w$; weighting and
Lemma~\ref{i2:lem-meanweight} give
$\eta\Sigma_L(j)\le\eta\Sigma^c_L(j)+d_{Z_0}(j)\log|G|$. The cut does not change the data in
$\mathrm{LinkN}_{\rm at}$, and $\log|G|\le C_1x^{C_1}\mathsf L\le C_1b\mathsf L^2\le C_1fmb\mathsf L^2$
for large $x$ (as $x^{C_1}\le e^{x^\varepsilon}$ eventually).
\end{proof}

\begin{lemma}[Multiplicity at most rank]\label{i2:lem-multle}
Assume: (i) for each lift $j$, the finite-dimensional space $M_j$ is the internal direct sum of
isotypic subspaces $M_j(\sigma)$; (ii) $\sum_{\pi\in Z^{\rm good}_i}\mathrm{mult}_i(\pi)=\dim\sum_{\sigma\in R(i)}M_{j}(\sigma)$
when $Z^{\rm good}_i\ne\emptyset$, $R(i)$ the good row types at $i$; (iii) distinct twisted local
indices at the same $(j,w)$ have disjoint $R(i)$; (iv) there is a linear $\theta_P:M^t\to M_j$
with $Z_j(\mu_w)=\mathrm{im}\,\theta_P^*$, and $\theta_P\theta_P^*$ acts on each nonzero
$M_j(\sigma)$ by a nonzero scalar. Then $\mathrm{rk}_j(w)=\dim M_j$ and $\mathrm{MultLe}$ holds.
\end{lemma}

\begin{proof}
By (iv), if $\theta_P^*v=0$ then each component $\lambda_\sigma v_\sigma$ of
$\theta_P\theta_P^*v$ vanishes, so $\theta_P^*$ is injective and $\mathrm{rk}_j(w)=\dim M_j$. By (ii)
and (i), $\sum_\pi\mathrm{mult}_i(\pi)\le\sum_{\sigma\in R(i)}\dim M_j(\sigma)$; by (iii) the sum over
$i\in\mathrm{locs}(j,w)$ is at most $\dim M_j$.
\end{proof}

\subsection{The local inputs}\label{i2:ssec-local}

We now discharge $\mathrm{ConclF}$ for the cut families. The published results used are the
following.

\begin{theorem}[Tate local duality; Neukirch--Schmidt--Wingberg, \emph{Cohomology of Number Fields}, 2nd ed., Thm.~7.2.6]\label{i2:thm-tate}
Let $K/\Q_l$ be finite, $M$ a finite discrete $G_K$-module and $M'=\mathrm{Hom}(M,\mu)$. For
$0\le i\le2$ the cup product $H^i(K,M)\times H^{2-i}(K,M')\to H^2(K,\mu)=\Q/\Z$ is a perfect pairing of
finite groups.
\end{theorem}
\noindent The coefficient modules to which it is applied are finite
torsion quotients of lattices over finite extensions of $\Q_l$.

\begin{theorem}[Tate's local Euler characteristic; ibid., Thm.~7.3.1]\label{i2:thm-euler}
For $M$ as above, $\#H^0(K,M)\,\#H^2(K,M)/\#H^1(K,M)=\|\#M\|_K$.
\end{theorem}
\noindent The hypotheses hold as for Theorem~\ref{i2:thm-tate}. That the groups
$H^i(K,M)$ are finite and vanish for $i\ge3$ is a separate standard statement (NSW, \emph{op.\ cit.}, Thm.~7.1.8(i),(iii)), used together with this one.

\begin{theorem}[Bloch--Kato, Grothendieck Festschrift I, Progr.\ Math.\ 86 (1990), Prop.~3.8 and Cor.~3.8.4]\label{i2:thm-bk}
For a de Rham representation $V$ of $G_K$, $H^1_f(K,V)$ and $H^1_f(K,V^*(1))$ are exact
annihilators under the local Tate pairing, and
$\dim H^1_f(K,V)=\dim H^0(K,V)+\dim D_{\rm dR}(V)/\mathrm{Fil}^0D_{\rm dR}(V)$. Integral form: if
$T\subset V$ and $T'\subset V^*(1)$ are dual $G_K$-stable lattices and $H^1_f(K,V/T)$ denotes the image of
$H^1_f(K,V)$, then a class in $H^1(K,V/T)$ lies in $H^1_f(K,V/T)$ if and only if its cup product with every class of
$H^1_f(K,T')$ (the preimage of $H^1_f(K,V^*(1))$) vanishes in $H^2(K,(V/T)\otimes T')\to H^2(K,(\Q_l/\Z_l)(1))$.
\end{theorem}
\noindent The integral form is stated in \S3 of \emph{loc.\ cit.}\ (reference to be confirmed).
The hypotheses hold: potentially crystalline representations are de Rham; the coefficients are passed from $\Q_l$ to a
finite extension $L$ by restriction of scalars.

\begin{theorem}[Chebotarev; Neukirch, \emph{Algebraic Number Theory} (1999), VII (13.4)]\label{i2:thm-cheb}
Let $L/K$ be a finite Galois extension of number fields with group $G$, $\sigma\in G$ and $\langle\sigma\rangle$ its
conjugacy class. The set of primes $\mathfrak p$ of $K$ unramified in $L$ whose Frobenius class is $\langle\sigma\rangle$
has Dirichlet density $\#\langle\sigma\rangle/\#G$.
\end{theorem}
\noindent It is used in Proposition~\ref{i2:prop-cut}(3) and Proposition~\ref{i2:prop-ptsel}(PF),
in a form (degree-one primes, inertia in a given open subgroup) derived from (13.4) and the fact that degree-one primes
have density one (Neukirch VII \S13, remark after (13.1)).

\begin{theorem}[Global Artin reciprocity; Neukirch, \emph{Algebraic Number Theory}, Ch.~VI (5.5), (5.6), (6.1), (6.9); Tate, in
Cassels--Fr\"ohlich (1967), Ch.~VII]\label{i2:thm-artin}
For a number field $K$ there is a continuous surjective homomorphism (the global Artin map) from the idele class group
$C_K$ onto $\mathrm{Gal}(K^{\rm ab}/K)$ whose restriction to $K_v^\times$ has dense image in the decomposition group at
$v$ and is the local Artin map there. In particular the Galois group of the maximal abelian extension of $K$
unramified at all finite primes and split at all real places (the small Hilbert class field, Neukirch VI (6.9)) is
isomorphic to the ideal class group of $K$. If $K$ is totally imaginary there are no real places, and this is the
maximal everywhere unramified abelian quotient of $G_K$.
\end{theorem}
\noindent With real places the maximal abelian extension
unramified at the finite primes alone has Galois group the narrow class group (Neukirch VI (6.8)); this case does not
occur here: the theorem is applied only to the CM field $E$, which is totally imaginary. Used in Propositions~\ref{i2:prop-colconds} (trivial block),
\ref{i2:prop-cec} and \ref{i2:prop-large}, and through $\mathrm{Hypotheses}_2$.

\begin{theorem}[Bhatt--Scholze, \emph{Prismatic $F$-crystals and crystalline Galois representations}, Cambridge J.\ Math.\ 11 (2023), Thm.~1.2]\label{i2:thm-bs}
For $K/\Q_p$ finite (more generally $K$ complete discretely valued of mixed characteristic with
perfect residue field), the etale realisation is an equivalence between prismatic $F$-crystals
on the absolute prismatic site of $O_K$ and $G_K$-stable $\Z_p$-lattices in crystalline
representations of $G_K$.
\end{theorem}
\noindent It is applied with $K$ a finite (tamely ramified) extension of $\Q_l$.

\begin{theorem}[Fontaine, Ast\'erisque 223 (1994), Exp.~II, III]\label{i2:thm-fontaine}
Crystalline (resp.\ de Rham) representations are stable under subrepresentations, quotients,
tensor products and duals, with $D_{\rm cris}$ compatible with these operations;
$B_{\rm cris}^{\varphi=1}\cap\mathrm{Fil}^0B_{\rm dR}=\Q_p$ (Exp.~II Thm.~5.3.7(iii));
$\dim D_B(V)\le\dim V$ for a $(\Q_p,G_K)$-regular ring $B$ (Exp.~III Prop.~1.4.2(i)); and for $K$ unramified and
integers $(w_\tau)$ there is a crystalline character with Hodge--Tate weight $w_\tau$ at $\tau$
(reference to be confirmed).
\end{theorem}
\noindent Closure under subquotients, tensor products and duals is Exp.~III Prop.~1.5.2 with \S5.1.

\begin{theorem}[Herr, Bull.\ Soc.\ Math.\ France 126 (1998), 563--600, Thm.~2.1]\label{i2:thm-herr}
For $p\ne2$ and a $p$-torsion representation $V$ of $G_K$ (a finite $\Z_p$-module with continuous action) with
etale $(\varphi,\Gamma_K)$-module $D$, the complex
$0\to D\xrightarrow{(\varphi-1,\gamma-1)}D\oplus D\xrightarrow{(\gamma-1)\oplus(1-\varphi)}D\to0$
computes $H^\bullet(G_K,V)$ functorially in $V$.
\end{theorem}
\noindent It is applied with $l\ge5$. We use it for $\Z_p$-representations of finite type, by the following passage to the limit.

\begin{corollary}\label{i2:cor-herr}
The conclusion of Theorem~\ref{i2:thm-herr} holds for every $\Z_p$-representation $V$ of $G_K$ of finite type, with
continuous cohomology $H^\bullet(G_K,V)$.
\end{corollary}

\begin{proof}
Put $V_n=V/p^nV$. These are $p$-torsion representations, and $V=\varprojlim_nV_n$ since $V$ is finitely generated over
$\Z_p$. Write $D(-)=(\mathcal O_{\widehat{\mathcal E^{\rm ur}}}\otimes_{\Z_p}-)^{H_K}$. As $V$ is finitely presented and
$\mathcal O_{\widehat{\mathcal E^{\rm ur}}}$ is $p$-adically complete, $\mathcal O_{\widehat{\mathcal E^{\rm ur}}}\otimes V=
\varprojlim_n\mathcal O_{\widehat{\mathcal E^{\rm ur}}}\otimes V_n$, and taking $H_K$-invariants commutes with inverse limits,
so $D(V)=\varprojlim_nD(V_n)$. The maps $D(V_{n+1})\to D(V_n)$ are surjective, because $D$ is exact on $p$-torsion
representations (Fontaine's equivalence, on which Theorem~\ref{i2:thm-herr} rests). Let $C_n$ be the complex of
Theorem~\ref{i2:thm-herr} for $D(V_n)$ and $C$ the one for $D(V)$. Then $C=\varprojlim_nC_n$ termwise, with surjective
transition maps, so there are exact sequences
$0\to{\varprojlim_n}^1H^{i-1}(C_n)\to H^i(C)\to\varprojlim_nH^i(C_n)\to0$ (Weibel, \emph{An Introduction to Homological
Algebra}, Thm.~3.5.8). By Theorem~\ref{i2:thm-herr}, functorially in $n$, $H^i(C_n)\cong H^i(G_K,V_n)$, a finite group
(NSW, Thm.~7.1.8(iii)). An inverse system of finite groups satisfies the Mittag-Leffler condition, so the
$\varprojlim^1$ term vanishes and $H^i(C)\cong\varprojlim_nH^i(G_K,V_n)$. On the Galois side NSW, Thm.~2.7.5 gives
$0\to{\varprojlim_n}^1H^{i-1}(G_K,V_n)\to H^i(G_K,V)\to\varprojlim_nH^i(G_K,V_n)\to0$, and the first term vanishes for
the same reason. Composing, $H^i(C)\cong H^i(G_K,V)$. Each step is natural in $V$, so the isomorphism is functorial.
\end{proof}

\begin{theorem}[Wach; Berger, Compositio Math.\ 140 (2004), Prop.~II.1.1, Prop.~III.4.2, Thm.~III.4.4]\label{i2:thm-wach}
For $K/\Q_p$ unramified and $V$ crystalline with Hodge--Tate weights in $[-h,0]$, every
$G_K$-stable lattice $T$ has a Wach module $N(T)$ of height $\le h$ with
$\mathbf A_K\otimes N(T)=D(T)$, $T\mapsto N(T)$ is a bijection between the $G_K$-stable lattices of $V$ and the
Wach lattices of $N(V)$, and
$N(V)/\pi N(V)\cong D_{\rm cris}(V)$.
\end{theorem}
\noindent It is used only at
large primes, where $l$ is unramified in $E$.

In the formalization each of these statements (in the form used) is itself a theorem of the
development; we list them because they are the mathematical content behind the steps below. We write
$\mathrm{Hypotheses}_2$ for the conjunction of: Fontaine's comparison for $B_{\rm cris}$
(Theorem~\ref{i2:thm-fontaine}), global Artin reciprocity (Theorem~\ref{i2:thm-artin}), the local Bloch--Kato statements
(Theorem~\ref{i2:thm-bk}), local Tate duality (Theorem~\ref{i2:thm-tate}), the local Euler characteristic and local
finiteness (Theorem~\ref{i2:thm-euler} and the remark after it), and the Bhatt--Scholze comparison
(Theorem~\ref{i2:thm-bs}). Chebotarev's theorem (Theorem~\ref{i2:thm-cheb}) is used in addition, in the cut and in the
construction of point families.

\medskip\noindent We first construct good cuts.

Let $G=\prod_{k=0,1}\prod_v\mathrm{GL}_r(\mathbb F_{g_k})$ be the level group of the rows, so that
$\C[G]=\bigotimes_{(k,v)}\C[\mathrm{GL}_r(\mathbb F_{g_k})]$. For a monic $f$ of degree $r$ over
$\mathbb F_g$ let $e_{\rm ser}(f)\in\C[\mathrm{GL}_r(\mathbb F_g)]$ be the sum of the primitive central
idempotents whose Deligne--Lusztig series is labelled by the characteristic polynomial $f$; for
distinct $f$ these are orthogonal central idempotents. For a tuple $p=(p_{kv})$ put
$e(p)=\bigotimes_{(k,v)}e_{\rm ser}(p_{kv})$, and let
\[
e_0=\sum_{p\ \mathrm{good}}e(p)\in Z(\C[G]),
\]
the sum over the good tuples (those allowed by the good row types of the datum). Distinct tuples
give orthogonal central idempotents, so $e_0$ is a central idempotent. Each local index $i$ has a
complex space $W_i$ of \emph{type vectors}, on which $\C[G]$ acts; for an automorphic form $P$ of
the row of $i$ its type space $W_i(P)\subseteq W_i$ is nonzero, and a \emph{type vector over
$x\in X_i$} is an element of $W_i(P)$ for some $P$ with eigensystem $x$. Write $\psi(z)$ for the
action of $\iota(z)$ on $W_i$, where $\iota:O_i\hookrightarrow\C$ is the fixed embedding.

\begin{definition}[Central-action tie]\label{i2:def-tie}
$D$ satisfies the \emph{central-action tie} if for every local index $i$:
\begin{enumerate}
\item[(T1)] $O_i$ contains a primitive $|G|$-th root of unity;
\item[(T2)] there is an $O_i$-algebra homomorphism $z_O:Z(O_i[G])\to\mathrm{End}_{O_i}(M_i)$
  whose image commutes with the Hecke algebra $T_i$;
\item[(T3)] (series localisation) to each $x\in Z_i$ is attached a tuple $p(x)=(p_{kv}(x))$, with
  $x\in Z^{\rm good}_i$ iff $p(x)$ is good, such that $e_{\rm ser}(f)$, in the factor $(k,v)$, acts on
  every type vector over $x$ by $1$ if $f=p_{kv}(x)$ and by $0$ otherwise;
\item[(T4)] (eigen-span) if $z\in Z(O_i[G])$, $c\in O_i$ and $\psi(z)F=\iota(c)F$ for every type
  vector $F$ over $x$, then $z_O(z)$ acts on $E_x$ by $c$.
\end{enumerate}
\end{definition}

In (T3) the tuple $p(x)$ is read off from the parameter list of any form with eigensystem $x$; that
it does not depend on the form is a clause of the datum (two forms with the same eigensystem have
parameter lists that are permutations of each other, and goodness and the labels are invariant
under permutation). (T3) is not itself a clause of the datum. The datum contains the weaker
\emph{series localisation} clause $\mathrm{SerLoc}$: for a parameter $\sigma'$ and a vector $F$ of its parameter space, if
$e_{\rm ser}(f)$ in the factor $(k,v)$ acts nonzero on $F$, then $f$ is the label of $\sigma'$ at $(k,v)$. (T3) is derived
from $\mathrm{SerLoc}$ together with clauses of the Galois preconditions (the identification of $e_{\rm ser}$ with the
series projectors, the coefficient and label ties, and the clauses (BC) and (SM1) of the datum; to avoid a clash with the
clauses (adm1)--(adm6) of Definition~\ref{i2:def-AdmT}, we call (BC) the base-change clause of the datum (every form has an
admissible parameter list, see (Ad) before Lemma~\ref{i2:lem-T3Z}), (LGC$'$) its local--global compatibility clause and
(SM1) its clause of strong multiplicity one, which includes density of split places; (LGC$'$) and (SM1) are named, not
defined, here); this
derivation, restricted to points of $Z_i$, is Lemma~\ref{i2:lem-T3Z}, proved in \S\ref{i2:ssec-proof-cec}. Likewise (T1)--(T4) are the part of the
tie used by Proposition~\ref{i2:prop-cut}; the full tie required by Theorem~\ref{i2:thm-linkN} is
Definition~\ref{i2:def-tie4}. At the members of the main-chain families, (T2) and (T4) are not assumed but
constructed from the member structure; this is done in \S\ref{i2:ssec-proof-cut}.

\begin{proposition}[Good cuts with a commuting endomorphism]\label{i2:prop-cut}
Let $D$ satisfy the Galois preconditions and the central-action tie. Keep $i$ iff
$Z^{\rm good}_i\ne\emptyset$, and for kept $i$ let $X^c_i$ be the set of $x\in X_i$ over which
there is a type vector $F$ with $e_0F\ne0$. Then:
\begin{enumerate}
\item there is a central $z\in Z(O_i[G])$ with $\iota(z)=|G|e_0$;
\item this is a good cut (Definition~\ref{i2:def-cut}), with $\alpha_i=z_O(z)$;
\item (formally verified, \cite{CharacterFamily}) if moreover $D$ satisfies the tie $\mathrm{Tie4}$ of
  Definition~\ref{i2:def-tie4} and $x$ is a norm choice for which $(D,x)$ satisfies the Galois preconditions, then, under
  Chebotarev's theorem, local finiteness, local Tate duality and the local Euler characteristic, the cut datum $D^c$ with
  the restricted norm choice satisfies the norm-choice facts and $\mathrm{Tie4}$.
\end{enumerate}
\end{proposition}

\begin{proof}
(1) \emph{Integrality.} Let $\Gamma$ be a finite group of order $n$ and $e=\sum_he_hh$ a central
idempotent of $\C[\Gamma]$. Left multiplication $L_g$ on $\C[\Gamma]$ has trace $n$ if $g=1$ and $0$
otherwise, so $n\,e_h=\mathrm{tr}(L_{h^{-1}e})$. Since $e$ is central and idempotent,
$T=L_{h^{-1}e}$ satisfies $T^{k}=L_{h^{-k}e}$, hence $T^{n+1}=T$ as $h^n=1$; its minimal
polynomial divides the separable polynomial $X(X^n-1)$, so its eigenvalues are $0$ or $n$-th roots
of unity and $n\,e_h=\mathrm{tr}(T)\in\Z[\zeta_n]$. Apply this to $\Gamma=G$, $e=e_0$: $|G|e_0$ has
coefficients in $\Z[\zeta_{|G|}]\subseteq\iota(O_i)$ by (T1). Lifting the coefficients gives
$z\in O_i[G]$ with $\iota(z)=|G|e_0$; $z$ is central because $\iota:O_i[G]\to\C[G]$ is injective
and $\iota(z)$ is central.

(2) \emph{Action of $e_0$ on type vectors.} Let $x\in Z_i$ and $F$ a type vector over $x$. By (T3),
$e(p)F=\prod_{(k,v)}[p_{kv}=p_{kv}(x)]\,F$, which is $F$ if $p=p(x)$ and $0$ otherwise. Summing over
good $p$: $e_0F=F$ if $p(x)$ is good, i.e.\ if $x\in Z^{\rm good}_i$, and $e_0F=0$ otherwise.

\emph{Goodness.} The kept indices are those with $Z^{\rm good}_i\ne\emptyset$ by construction. If
$x\in Z^{\rm good}_i$, choose a form with eigensystem $x$ and a nonzero $F$ in its type space;
then $e_0F=F\ne0$, so $x\in X^c_i$. Thus $Z^{\rm good}_i\subseteq X^c_i$, and $X^c_i\ne\emptyset$
for kept $i$. If $x\in Z_i\cap X^c_i$, some $F$ over $x$ has $e_0F\ne0$, so $x\in Z^{\rm good}_i$ by
the previous paragraph. The endomorphism $\alpha_i=z_O(z)$ is $O_i$-linear and commutes with $T_i$
by (T2). For $x\in Z^{\rm good}_i$ and every type vector $F$ over $x$,
$\psi(z)F=|G|e_0F=\iota(|G|)F$, so $\alpha_i=|G|$ on $E_x$ by (T4). For $x\notin X^c_i$ we have
$e_0F=0$, hence $\psi(z)F=0$, for all $F$ over $x$, so $\alpha_i=0$ on $E_x$ by (T4).

(3) This part is not proved here: the Galois preconditions and the norm-choice facts are predicates with many clauses
that are not defined in this paper. The structure of its formal proof is described in \S\ref{i2:sst-cut3}; apart from
restriction of clauses to the smaller eigensystem sets, its content is clause (G5) for the cut datum (from local Tate
duality and the Euler characteristic) and the duties clause (from Chebotarev's theorem).
\end{proof}

\begin{remark}\label{i2:rem-cut-extra}
The formal notion of goodness has one further clause, used only by the norm-choice facts: for
$x\in X^c_i$ and any form with eigensystem $x$, a condition at the split places above the
generator primes. It is proved by the same mechanism: a type vector $F$ over $x$ with $e_0F\ne0$
lies in the type space of a good tuple, which gives the condition for one form over $x$, and the
permutation clause of the datum transports it to every form over $x$.
\end{remark}

\medskip\noindent We next describe the four local steps at a generator prime.

Let $i$ be twisted at a generator prime $l=g_k\ne2$; by the degree step $l$ has $2f$ places
$w$ in $E$, all with $E_w=\Q_l$. Let $O$ be the coefficient ring at $w$, $G=G_{\Q_l}$,
$N=n_V+1$.

\begin{definition}[Local family conditions]\label{i2:def-locfam}
For a family $(M_x)_{x}$ of continuous representations of $G$ over $L_w$:
(Pot$_n$) there are a finite Galois $K'/\Q_l$, a Frobenius $\mathrm{Fr}$ and a
$\mathrm{Fr}$-stable set $\Sigma\ni1$ of at most $n$ characters of the group
$I(K'/\Q_l)=I_l\cdot G_{K'}\subseteq G$ (the closed subgroup generated by the inertia group $I_l$ of $G$ and by
$G_{K'}$; its image in $\mathrm{Gal}(K'/\Q_l)$ is the inertia group of $K'/\Q_l$), trivial on $G_{K'}$, of order prime to $l$, such that every $M_x$ is crystalline over $K'$ and
inertia acts on $D_{\rm pcris}(M_x)$ through $\Sigma$; (Jump$_{n;a,b}$) some $U\subset\Z$ with
$0\in U\subseteq[-a,b]$, $|U|\le n+1$, contains all Hodge--Tate jumps of all $M_x$; a
\emph{subfamily} is a representation $W$ with equivariant maps $p_j:W\to M_{x_j}$ of trivial
common kernel; (Lift$_e$) for subfamilies $W_1,W_2$, stable lattices $T_\nu\subset W_\nu$ and
every surjective equivariant $F:T_1\to T_2$, every finite (crystalline) cocycle $c$ of $T_2$
satisfies $l^ec\in F_*Z^1_f(T_1)+B^1(T_2)$. Here a cocycle is \emph{finite} if its class lies in
$H^1_f$, i.e.\ it becomes a coboundary in $V\otimes B_{\rm cris}$.
\end{definition}

\begin{lemma}[Jump ranges]\label{i2:lem-jumps}
Let $\eta,t,s,S\in\N$ with $t\ge1$ odd and $\eta+t+1<s$, $c_t=(t-1)/2$, and let $U_w$, $|U_w|\le n_V$,
contain all jumps of the twisted family at $w$, each $a\in U_w$ satisfying $0\le a<t$, or
$|a-c_t|\le\eta+1$, or $0<\pm(a-c_t)\le s+S-c_t$ with a fixed sign $\pm$ (the sign $\mathrm{sg}(w)$ of $w$). If the sign is $+$,
the family satisfies (Jump$_{N;\eta+t+1,s+S}$) with $U=U_w\cup\{0\}$; if it is $-$, the dual family
$(r_x^\vee(1))$ does, with $U=\{0\}\cup\{-a-1:a\in U_w\}$.
\end{lemma}

\begin{proof}
$|U|\le n_V+1$ and $0\in U$. For the sign $+$, let $u\in U_w$. In the three cases: $u\ge0$; or
$u\ge c_t-\eta-1\ge-\eta-1\ge-(\eta+t+1)$; or $u>c_t\ge0$. And $u\le\max(t-1,c_t+\eta+1,s+S)\le s+S$, as
$\eta+t+1<s$. For the sign $-$, the jumps of $r_x^\vee(1)$ are the $-a-1$, $a$ a jump of $r_x$, so we need
$-(s+S)-1\le a\le\eta+t$ for $a\in U_w$: either $0\le a<t$; or $|a-c_t|\le\eta+1$, whence
$a\le c_t+\eta+1\le\eta+t$ and $a\ge c_t-\eta-1\ge-\eta-1\ge-(s+S)-1$; or $0<c_t-a\le s+S-c_t$, whence $a<c_t\le t$ and
$a\ge2c_t-(s+S)=t-1-(s+S)\ge-(s+S)-1$.
\end{proof}

\begin{lemma}[Frobenius periods]\label{i2:lem-period}
Let $e=l^a-1$, and for $c\in\Z/e$ let $\mathrm{per}(c)$ be the least $m\ge1$ with $l^mc=c$. If
$\Sigma'\subseteq\Z/e$ has $l\Sigma'\subseteq\Sigma'$, $0\in\Sigma'$, $|\Sigma'|\le N$, $N\ge2$, and
$r(\Sigma')=\max_{b,b'\in\Sigma'}\mathrm{per}(b-b')$, then $1\le r(\Sigma')$ and $4r(\Sigma')\le N^2$.
\end{lemma}

\begin{proof}
Multiplication by $l$ is a permutation of $\Z/e$ (as $l^a=1$), and orbits of elements of
$\Sigma'$ lie in $\Sigma'$, of size $\mathrm{per}$. Let $(x,y)$ attain the maximum;
$\mathrm{per}(x-y)\mid\mathrm{per}(x)\mathrm{per}(y)$. If $y$ is in the orbit of $x\ne0$, then
$\mathrm{per}(x-y)\le\mathrm{per}(x)\le N-1$ (the orbit misses $0$), and $4(N-1)\le N^2$. Otherwise the
two orbits are disjoint and $4\mathrm{per}(x)\mathrm{per}(y)\le(\mathrm{per}(x)+\mathrm{per}(y))^2\le N^2$.
\end{proof}

\begin{lemma}[Change of coefficients]\label{i2:lem-coef}
Let $L\subseteq L'$ be finite extensions of $\Q_l$, $T_1,T_2$ stable $O_L$-lattices, $F:T_1\to T_2$
surjective and equivariant. If $F\otimes O_{L'}$ has the lifting property with exponent $e$, so
has $F$. For every $a\ge1$ there is such $L'\supseteq L$ containing $\Q_{l^a}$, hence a primitive
$(l^a-1)$-th root of unity.
\end{lemma}

\begin{proof}
$O_{L'}$ is free over $O_L$; writing $1=\sum\beta_ib_i$ in a basis, the ideal $(\beta_i)$ is
the unit ideal (its generator $\delta$ has $\delta^{-1}\in O_{L'}\cap L=O_L$), so there are
$\alpha_i\in O_L$ with $\sum\alpha_i\beta_i=1$ and $\varrho=\sum\alpha_ib_i^*$ is an $L$-linear
retraction $L'\to L$ with $\varrho(O_{L'})\subseteq O_L$. Equivariant maps carry finite cocycles to
finite cocycles. If $c$ is a finite cocycle of $T_2$, then
$l^e(c\otimes1)=F'\circ d+\partial y$, and applying $\mathrm{id}\otimes\varrho$ gives
$l^ec=F\circ((\mathrm{id}\otimes\varrho)d)+\partial((\mathrm{id}\otimes\varrho)y)$. For the second claim
take $L'=L\Q_{l^a}$, which contains the Teichm\"uller lifts of $\mathbb F_{l^a}^\times$.
\end{proof}

\begin{proposition}[local lifting theorem]\label{i2:prop-lift}
Let $C\ge40$, $l$ a prime, $K'/\Q_l$ finite Galois, $L/\Q_l$ finite, $T_1,T_2$ $G$-stable
$O_L$-lattices in representations $V_1,V_2$ of $G=G_{\Q_l}$ that are crystalline over $K'$, and
$F:T_1\to T_2$ surjective and $G$-equivariant. Let $\mathrm{Fr}$ be a Frobenius and $\Sigma$ a
$\mathrm{Fr}$-stable set of characters of $I(K'/\Q_l)$ of order prime to $l$, with $1\in\Sigma$,
$|\Sigma|\le N$, $N\ge2$, such that inertia acts on $D_{\rm pcris}(V_\nu)$ through $\Sigma$
($\nu=1,2$). Let $U\subset\Z$ with $0\in U\subseteq[-\eta,H]$, $|U|\le N+1$, $1\le\eta<H$,
containing the Hodge--Tate jumps of $V_1,V_2$, and assume $C(HN^2+N^3)<l$. Let $r'(\Sigma)$ be
the maximum over $a,b\in\Sigma$ of the length of the $\mathrm{Fr}$-orbit of $ab^{-1}$. Then
$4r'(\Sigma)\le N^2$, and $F$ has the lifting property (Lift$_e$) with
$e=C\bigl(Nr'(\Sigma)+\eta(N+r'(\Sigma))\bigr)$.
\end{proposition}

\noindent The proof is in \S\ref{i2:ssec-proof-lift}. It reduces, by a change of coefficients and
Proposition~\ref{i2:prop-goodlat}, to lattices crystalline over the tame field $K_\varpi$, and proves the lifting
there (Proposition~\ref{i2:prop-liftgood}) from a sequence of Breuil--Kisin steps. The lifting of maps
(Proposition~\ref{i2:bk-c}, from Theorem~\ref{i2:thm-bs}) and the $\zeta$-bound (Proposition~\ref{i2:bk-zeta}) are proved
there; the remaining steps (Propositions~\ref{i2:bk-b}, \ref{i2:bk-ex}, \ref{i2:bk-cl}, \ref{i2:bk-mon},
\ref{i2:bk-tri}, \ref{i2:bk-int}) are stated exactly and cited as formally verified.

\begin{lemma}[Lifting for the twisted family]\label{i2:lem-steplift}
Let $(M_x)$ satisfy (Pot$_N$) and (Jump$_{N;\eta',H}$) with $1\le\eta'<H$, $N\ge2$ and
$C_T(HN^2+N^3)<l$, $C_T\ge40$. Then the family has (Lift$_{e}$) with
$e=C_T\bigl(N\lfloor N^2/4\rfloor+\eta'(N+\lfloor N^2/4\rfloor)\bigr)$. The same holds for the dual
family $(M_x^\vee(1))$ when it satisfies the same conditions, and (Pot$_n$) passes to the dual family with
$\Sigma^{-1}$.
\end{lemma}

\begin{proof}
A subfamily $W$ embeds into a finite sum of members, so (Theorem~\ref{i2:thm-fontaine}) it is
crystalline over $K'$, inertia acts on $D_{\rm pcris}(W)$ through $\Sigma$, and its jumps lie in
$U$. Proposition~\ref{i2:prop-lift} gives $r'\le\lfloor N^2/4\rfloor$ and the lifting property
with exponent $C_T(Nr'+\eta'(N+r'))\le e$; the lifting property is monotone in the exponent.
For the dual family use $\Sigma^{-1}$, which contains $1$, has the same size and is
$\mathrm{Fr}$-stable.
\end{proof}

We fix the objects used at a generator prime $l=g_k$. A \emph{point family} $r=(r_x)_{x\in X_i}$ (see
\S\ref{i2:ssec-proof-hub}) consists of continuous $r_x:G_E\to\mathrm{GL}_{n_V}(O_i)$ whose Frobenius characteristic
polynomials at unramified places are given by the Hecke eigenvalues of $x$ (the \emph{good-row} property). The
\emph{twisted family} at $w\mid l$ is $(r_x\otimes\chi^{-1}|_{G_{E_w}})_{x}$, $\chi$ the first character of the datum
at $i$, and the \emph{dual family} $(r_x^\vee(1))_x$ consists of the Tate duals of the twisted members (as in
Lemma~\ref{i2:lem-jumps} we keep writing $r_x$ for the twisted member).
A point $\pi$ has a \emph{semisimple decomposition} (for a given $\rho_\pi:G_E\to\mathrm{GL}_r(O_i)$) if
$r_\pi\otimes K_i\cong\rho_\pi\otimes K_i\oplus\bigoplus_{c=1}^t\chi_c$ with the characters $\chi_c$ of the datum. A
set $X_s\subseteq X_i$ is a \emph{type locus} (for $(i,k)$) if there are squarefree $s_v\in\mathbb F_{g_k}[X]$, one for each
place $v$ of the reflex field above $g_k$, none vanishing at the label $c_v^{\pm1}$ of $\chi$ at $v$, such that every
$x\in X_s$ is the eigensystem of a form whose type at each $v$ has Deligne--Lusztig label $s_v$. A representation $M$
\emph{embeds into a twisted member} if $M\otimes K_i$ is isomorphic to a subrepresentation of some $r_x\otimes\chi^{-1}$,
$x\in X_s$. The \emph{cut cocycles} of $M$ at $w$ are defined at the start of the proof of (b) below.

\begin{proposition}[potential crystallinity and the local count]\label{i2:prop-potcount}
Let $i$ be a local index with $l_i=l=g_k$ and $w\mid l$ with $E_w=\Q_l$.
(a) Assume the following clauses of the datum: (a1) the type-locus local--global compatibility (for a point family with
the good-row property, $x$ in a type locus and $w\mid l$, $r_x$ is de Rham at $w$ with monodromy zero and inertia
polynomial $\prod_{y\in M_x}(X-\omega_y)$ as recalled in the proof); (a2) Fontaine's comparison at the places above $l$;
(a3) the coefficient, character and Hodge--Tate ties at $i$; (a4) the inertia-polynomial tie (for $g$ in inertia a
power of the inertia polynomial at $g$ kills $D_{\rm pcris}$); (a5) stability of potential crystallinity under twists by
potentially crystalline characters, with the inertial characters multiplied. Let $\pi\in Z_i$ have a semisimple
decomposition and let $X_s\ni\pi$ be a type locus. Then the twisted family $(r_x)_{x\in X_s}$ at $w$ satisfies
(Pot$_{n_V+1}$).
(b) (direct case) If the twisted family $(r_x)_{x\in X_s}$ at $w$ satisfies (Pot$_n$) and (Lift$_e$), and $M$ is a rank-$k$
representation that embeds into a twisted member, then the span of the cut cocycles of $M$ at $w$ has length at
most $2ke\,\ell_O(O/l)$ modulo $H^1_f$.
(c) (dual case; formally verified, \cite{CharacterFamily}) If $r_x\otimes K_i$ is a semisimple representation of $G_E$ for
every $x\in X_s$, the twisted family at $w$ satisfies (Pot$_n$), the dual family satisfies (Lift$_e$), and $M$ is a rank-$k$
representation that embeds into a twisted member, then the span of the cut cocycles of $M$ at $w$ has length at most
$2ke\,\ell_O(O/l)$ modulo $H^1_f$.
\end{proposition}
\noindent The proofs of (a) and (b) are in \S\ref{i2:ssec-proof-potcount}; (b) rests on the pushout bound,
Proposition~\ref{i2:prop-pushout}. Clause (a1) is not a primitive clause of the datum but a formally verified consequence
of the depth-zero admissibility of types and of the clauses (BC), (LGC$'$), (SM1) of the datum; (a2)--(a5) are
clauses of the datum. The argument of (b) does not transfer directly to (c): if $X=Y\otimes L$ is a subrepresentation
of a sum of twisted members, its Tate dual is a quotient, not a subrepresentation, of a sum of dual members. The formal
proof of (c) uses the global semisimplicity of the members to make $X$ and $U$ direct summands, so that their Tate duals
are subrepresentations of sums of dual members, and then applies Proposition~\ref{i2:prop-pushout} in the second case of
Lemma~\ref{i2:lem-absorb}; its structure is in \S\ref{i2:sst-potc}. The semisimplicity hypothesis is a clause of the
good family (goodness of the $r_x$) wherever (c) is applied.

\medskip\noindent At the twisted indices we use a global lattice-index bound, for which we need the following data.

\emph{Local admissibility data.} Fix twisted $i$, $l=l_i$, a point family $(r,\rho)$ and $\pi\in Z_i$. A
\emph{decomposition} of the places of $E$ above $l$ is a partition into rows and columns; \emph{admissibility data} for it
are row lengths $\ell_w\ge0$ (rows $w$), column lift exponents $e'_w\in\N$ (columns $w$) and a block number $c_B\in\N$. The
\emph{local admissibility conditions} $\mathrm{Adm}_i(\pi;\ell,e',c_B)$ are the conjunction of twenty clauses of the
formalization: the Galois-family conditions (G1)--(G6) at $i$ and the tame-generator conditions; the local Bloch--Kato
and $H^2$-duality clauses at the places above $l$; the block condition with constant $c_B$; the row-length condition
(the local condition at a row $w$ has colength $\le\ell_w$); the column lifting condition with exponents $e'_w$;
continuity and semisimplicity of the family; the semisimple decomposition at $\pi$; the cyclotomic inertia-image,
local-coefficient and cyclotomic-completion clauses; Fontaine's comparison at the places above $l$; and the
unramified class-field bound (the module of everywhere unramified homomorphisms $G_E\to K_i/O_i$ has normalised length
$\le v_l(h_E)$). This list is the formal predicate used at inert $l$ and at $q,d$; the closed formal theorem behind
Proposition~\ref{i2:prop-colconds} uses nineteen clauses of it (\S\ref{i2:sst-colconds}), and at split $l$ the formal
proof of Proposition~\ref{i2:prop-locadm}(1) uses a different predicate with 38 clauses (\S\ref{i2:sst-locadm}). We do
not define these clauses individually; statements that take them as hypotheses
(Propositions~\ref{i2:prop-selmer}, \ref{i2:prop-colconds}, \ref{i2:prop-locadm}) are therefore stated informally as far as
these clauses are concerned; the exact statements are in \cite{CharacterFamily}.

\begin{proposition}[Selmer bound for the congruence module]\label{i2:prop-selmer}
Let $i$ be twisted with $l=l_i\notin\{g_0,g_1\}$, $\pi\in Z_i$, a decomposition of the places above $l$ with data
$(\ell,e',c_B)$ satisfying $\mathrm{Adm}_i(\pi;\ell,e',c_B)$, the Selmer group of the $\rho_\pi$-block of finite length,
and $D$ Sha-finite. Then, with $e_O=\ell(O/l)$, $j_w$ the number of column generators at $w$, $h^0_w$ the length of
the $\rho_\pi$-invariants of the decomposition group at a column $w$, $H^0=\max_wh^0_w$, $\deg_w$ the local degree,
$v_\delta$ as in Proposition~\ref{i2:prop-homterms} and $\zeta_v$ the roots of unity of the datum at $v\in Q_i$,
\[
\begin{aligned}
v^{cc}_i(\pi)\le{}&s_i(\pi)+\sum_{e=1}^{t-1}\ell_D(l,e)+v_l(h_E)+\frac1{e_O}\Bigl(\sum_{\rm rows}\ell_w
+\sum_{\rm cols}\bigl(e_Oe'_w+j_w(e_Oe'_w+v_\delta)+2v_\delta+h^0_w\bigr)\Bigr)\\
&+tc_B\sum_{\rm rows}(\deg_w+2)+\frac{4n_V(|\mathrm{col}|+1)v_\delta+t|\mathrm{col}|H^0}{e_O}
+t\sum_{v\in Q_i}\frac{\ell(O/(\zeta_v-1))}{e_O}.
\end{aligned}
\]
\end{proposition}
\noindent This proposition does not apply at the generator primes; there the corresponding inequality is
Proposition~\ref{i2:prop-stepglobal}.
\noindent The proof is in \S\ref{i2:ssec-proof-selmer}: the determinant (lattice) bound is proved there, and the bound of the resulting $\mathrm{Hom}$-terms (Proposition~\ref{i2:prop-homterms}) is proved there from an algebraic bound for each block and the formally verified arithmetic input Proposition~\ref{i2:prop-colconds}.

\medskip

At $l=g_k$ the local index is bounded through a separate cut for each good point, and the global inequality is not
Proposition~\ref{i2:prop-selmer}. For $X_s\subseteq X_i$ and $\pi\in X_s$ let
\[
v^{\rm cut}_i(X_s,\pi)=\ell\bigl(O_i/\{f(\pi):f\in T_i,\ f|_{X_s\setminus Z_i}=0\}\bigr)/\ell(O_i/l)
\]
be the normalised congruence length of $\pi$ relative to the points of $X_s$ outside $Z_i$. For $v\in Q_i$ write
$v(\zeta_v-1)=\ell(O_i/(\zeta_v-1))/\ell(O_i/l)$, and let $v_\delta$ be the normalised valuation of the element $\delta_O$
of Lemma~\ref{i2:lem-vdelta} at $i$.

\begin{proposition}[global inequality at a generator prime]\label{i2:prop-stepglobal}
Assume global Artin reciprocity. Let $i$ be a local index, $l=l_i$, $(r,\rho)$ a point family with the good-row property,
$\pi\in Z_i$, $X_s\subseteq X_i$ with $\pi\in X_s$, and $e_\rho,e_\chi\in\N$. Write $\chi_1,\dots,\chi_t$ for the characters
of the datum at $i$ and $\tilde r_x=r_x\otimes\chi_1^{-1}$ for the twisted members. Assume the global admissibility
clauses at $(i,\pi)$:
\begin{itemize}
\item[(G4)] for $x\in X_i$ and every place $v\nmid l$ of $E$ with $v\notin S_i\cup Q_i$, $r_x$ is trivial on the inertia
group at $v$;
\item[(G6)] if $x\in X_i$ and $r_x\otimes K_i$ has a nonzero vector on which $G_E$ acts by $\chi_1$, then $x\in Z_i$;
\item the Galois-family conditions (G1), (G5) at $i$, the tame-generator conditions at $S_i\cup Q_i$ (which include
$\zeta_v\ne1$ for $v\in Q_i$), residual absolute irreducibility of $\rho_\pi$, local-coefficient compatibility, the
semisimple decomposition of $\pi$, and the existence of idempotent data with
good congruence objects at $\pi$ (the data of Proposition~\ref{i2:prop-latbound});
\end{itemize}
and at every $w\mid l$, writing $W_0=\rho_\pi\otimes\chi_1^{-1}$ and $W_c=\chi_c\chi_1^{-1}$ ($1\le c\le t$) for the blocks of
the twisted member $\tilde r_\pi$: the Hodge--Tate character condition $\mathrm{HT}_w$ (for $2\le c\le t$, $H^1_f$ of
$W_c$ at $w$ consists of coboundaries; for $c=1$, $H^1_f$ of the trivial block $W_1$ at $w$ consists of classes
trivial on inertia), and the bounds by $e_\rho$ (resp.\ $e_\chi$) for the span of the cut cocycles, relative to $X_s$, of
$W_0$ (resp.\ of each $W_c$, $1\le c\le t$) modulo $H^1_f$, in the normalisation of
Proposition~\ref{i2:prop-potcount}(b). If the Selmer group of $W_0$ is finite and
$\ell(\mathrm{Sha}(\mathrm{ld}_D(l),e))<\infty$ for $0<e<t$, then
\[
v^{\rm cut}_i(X_s,\pi)\le\#\{w\mid l\}\,(e_\rho+te_\chi)+s_i(\pi)+\sum_{0<e<t}\ell_D(l,e)+v_l(h_E)+4n_Vv_\delta
+t\sum_{v\in Q_i}v(\zeta_v-1).
\]
\end{proposition}
\noindent The proof is in \S\ref{i2:ssec-proof-stepglobal}: a determinant bound for the cut congruence ideal by the
lengths of $t+1$ modules of global cut cocycles, one per block, followed by a bound for each block (local conditions above
$l$, coboundaries, peeling of the places of $Q_i$, and a global term: the Selmer group, the class group by Artin
reciprocity, or $\mathrm{Sha}$). The arithmetic inputs form Proposition~\ref{i2:prop-cutcoc}, proved there; the assembly is
proved. Of the global admissibility clauses, the proof uses (G4), (G5),
(G6), the tame-generator conditions, the semisimple decomposition and the idempotent data; (G1), residual irreducibility
and local-coefficient compatibility are hypotheses of the formal statement that its proof does not use. The clauses (G1),
(G5), the tame-generator conditions, residual irreducibility and local-coefficient compatibility are named, not defined,
here. (A variant of the formal statement has in
addition the tame/wild conditions at $S_i\cup Q_i$ as a hypothesis, with the same proof.)

\begin{proposition}[paired-character condition]\label{i2:prop-pcc}
Let $D$ satisfy $\mathrm{Hypotheses}_2$, the norm-choice facts and $\mathrm{Tie4}$ (Definition~\ref{i2:def-tie4}), and let
$C_T\ge40$ be an integer and $C\ge8C_T$. Let $i$ be a local index with $l_i=g_k$, $k'\ne k$, $(r,\rho)$ a point family with the good-row
property such that every $\pi\in Z_i$ has a semisimple decomposition, and let $\eta,s,S\in\N$, a sign function
$\mathrm{sg}$ on the places above $l$ and finite sets $U_w\subset\Z$ satisfy:
$\eta+t+1<s$; $\log g_{k'}\le2\log g_k$; $C_T\bigl((s+S)(n_V+1)^2+(n_V+1)^3\bigr)<g_k$; ($\mathrm{SignOK}$) conjugate
distinct places above $l$ have opposite signs; ($\mathrm{HostJ}$) for every $w\mid l$, $|U_w|\le n_V$, $U_w$ contains all
Hodge--Tate jumps at $w$ of all twisted members $\tilde r_x$, and every $a\in U_w$ satisfies $a\ne-1,t$ and one of $0\le a<t$,
$|a-c_t|\le\eta+1$, $0<\pm(a-c_t)\le s+S-c_t$ with $\pm=\mathrm{sg}(w)$; the Selmer groups of the blocks $W_0=\rho_\pi\otimes\chi_1^{-1}$,
$\pi\in Z^{\rm good}_i$, are finite; and $\ell(\mathrm{Sha}(\mathrm{ld}_D(l),e))<\infty$ for $0<e<t$. Then
\[
\eta^{\rm loc}_i\le\sum_{\pi\in Z^{\rm good}_i}\mathrm{mult}_i(\pi)\Bigl(4fn_Ve_T+2fr^2+s_i(\pi)+\sum_{0<e<t}\ell_D(l,e)+v_l(h_E)
+4n_Vv_\delta+t\sum_{v\in Q_i}v(\zeta_v-1)\Bigr),
\]
$e_T=Cn_V^2(n_V+\eta+t)$.
\end{proposition}
\noindent The proof is in \S\ref{i2:ssec-proof-pcc}. It cuts at each good point separately to a type locus (Proposition~\ref{i2:prop-percut}, proved, with
Lemma~\ref{i2:lem-bcperm} proved modulo the formally verified Lemma~\ref{i2:lem-isorep}; the set $X^c_i$ of Proposition~\ref{i2:prop-cut} is not a type locus),
bounds $\eta^{\rm loc}_i$ by the per-point cut congruence lengths (Lemmas~\ref{i2:lem-famcore} and \ref{i2:lem-percutfam},
proved), bounds each of these by Proposition~\ref{i2:prop-stepglobal} with the local exponents of
Proposition~\ref{i2:prop-potcount} and Lemmas~\ref{i2:lem-jumps}, \ref{i2:lem-steplift} (the supply of their hypotheses from
the norm-choice facts and $\mathrm{Tie4}$ is the formally verified Proposition~\ref{i2:prop-pcclocal}), and ends with an
elementary computation (Lemma~\ref{i2:lem-vG}).

\begin{proposition}[generator primes]\label{i2:prop-typ}
Let $C\ge0$, $C_T\in\N$, and assume: the conclusion of Proposition~\ref{i2:prop-pcc} with constants $(C,C_T)$ for every
twisted $i$ at $l\in\{g_0,g_1\}$, for data $(r,\rho,\eta,s,S,\mathrm{sg},U)$ whose hypotheses are supplied by the hub data
(this is part of the formally verified setup, Proposition~\ref{i2:prop-hub}); $\mathrm{KillAt}$ (below); $\mathrm{MultLe}$
and (adm3); and (Par) with constant $C_1$. Then, with $C'=2000(C+1)+5C_1$, for every lift $j$ and $l\in\{g_0,g_1\}$,
\[
\eta_L(j,l)\le\log l\cdot d_{Z_0}(j)\bigl(C'x^{C'}\eta_{\rm sz}+2l^{-2}+\mathrm{VV}_D(l)+v_l(h_E)\bigr).
\]
\end{proposition}

\begin{proof}
\emph{Local bound.} For twisted $i$ at $l$: if $Z^{\rm good}_i=\emptyset$ then $\eta^{\rm loc}_i=0$; otherwise the paired-character
bound gives $\eta^{\rm loc}_i\le\sum_\pi\mathrm{mult}_i(\pi)\,\mathrm{rhs}_i(\pi)$, $\mathrm{rhs}_i(\pi)$ the bracket of
Proposition~\ref{i2:prop-pcc}.
\emph{Size of $\mathrm{rhs}$.} By (Par) the polynomial terms $4fn_Ve_T+2fr^2+4n_Vv_\delta$ are at most $C'x^{C'}\eta_{\rm sz}$
(elementary comparison of monomials, $x\ge1$); $\mathrm{KillAt}$ gives $t\sum_{Q_i}v(\zeta_v-1)\le l^{-2}$ and
$s_i(\pi)\le\nu_l+l^{-2}$, and $\nu_{g_k}=0$; $\sum_e\ell_D(l,e)=\mathrm{VV}_D(l)$. So
$\mathrm{rhs}_i(\pi)\le U_l:=C'x^{C'}\eta_{\rm sz}+2l^{-2}+\mathrm{VV}_D(l)+v_l(h_E)$.
\emph{Summation} as in Lemma~\ref{i2:lem-etaAt}: $\eta_j(w)\le U_l\sum_i\sum_\pi\mathrm{mult}_i(\pi)\le U_l\,\mathrm{rk}_j(w)
=U_l\,d_{Z_0}(j)$ by $\mathrm{MultLe}$ and (adm3), and Lemma~\ref{i2:lem-meanweight} gives the factor $\log l$.
\end{proof}

\begin{proposition}[the remaining primes]\label{i2:prop-other}
Under the setup of the family (Theorems~\ref{i2:thm-tate}--\ref{i2:thm-wach}, the norm-choice
facts, $\mathrm{KillAt}$, Sha-finiteness), for twisted $i$ with $l=l_i\notin\{g_0,g_1\}$ and
$\pi\in Z^{\rm good}_i$:
\begin{enumerate}
\item (medium primes) if $\mathrm{RangeAD}(l)$:
$\mathrm{Bd}_i(\pi,\sum_{\lambda\mid l}\mathrm{len}^0_i(\pi,\lambda)+\mathrm{err}_i(\pi))$ with
$\mathrm{err}_i(\pi)\le7200fnt^2n_V\mathsf L^2$, and the lengths satisfy (A), (B), (Med), (SmI);
\item ($q$ and $d$) if $l\in\{q,d\}$: $\mathrm{Bd}_i(\pi,\mathrm{midG}(l))$ with
$\mathrm{midG}(l)=Cf\bigl(r(b/l+S_c+t+\eta_{\rm sz}+\log b)+mb\bigr)+C(frt)^C\log b\le2Cfmb$ eventually;
\item (small primes and $3$) if $l\in\mathrm{small}\cup\{3\}$:
$\mathrm{Bd}_i(\pi,Cfmb\log x+Cx^C(\log b)^2)$;
\item (large primes) if $l>l^*$, $l\notin\mathrm{exc}$, and $\mathrm{exc}=\mathrm{Exc}(L)$, under the further
hypotheses of Proposition~\ref{i2:prop-large}: $\mathrm{Bd}_i(\pi,0)$.
\end{enumerate}
\end{proposition}
\noindent The proof of items (1)--(3) is in \S\ref{i2:ssec-proof-other}: Proposition~\ref{i2:prop-selmer} for a
suitable decomposition of the places above $l$, with the local admissibility data of the formally verified
Proposition~\ref{i2:prop-locadm}, followed by elementary comparisons. Item (4) is the formally verified
Proposition~\ref{i2:prop-large}.

\medskip

For a twisted local index $i$ at $l$, $\mathrm{KillAt}$ is the pair of inequalities
\[
s_i(\pi)\le\nu_l+l^{-2}\quad(\pi\in Z^{\rm good}_i),\qquad
t\sum_{v\in Q_i}\frac{\ell(O/(\zeta_v-1))}{\ell(O/l)}\le l^{-2},
\]
together with finiteness of the Selmer group of the $\rho_\pi$-block. Hub data at $i$ carry a
\emph{twist datum} at $w=w_i$: Bloch--Kato lengths $\mathrm{BK}(\sigma)\in[0,\infty]$ indexed by
signatures $\sigma$, tame values $\mathrm{tv}(P)\ge0$ at the auxiliary places $P$ of the twist, and a
base value $\nu$, equal to $\nu_l$ since the twist datum has residue characteristic $l$. Two
conditions on it are used:
(SC) \emph{signature choice}: for each $\pi\in Z^{\rm good}_i$ there is a good signature $\sigma$
with $s_i(\pi)\le\mathrm{BK}(\sigma)$ and ($\mathrm{BK}(\sigma)<\infty\Rightarrow$ the Selmer group of the
$\rho_\pi$-block is finite); and
$\sum_{v\in Q_i}\ell(O/(\zeta_v-1))/\ell(O/l)\le\sum_P\mathrm{tv}(P)$;
(GT) \emph{good twist}: there are $\epsilon_T(\sigma)\ge0$ with $\sum_\sigma\epsilon_T(\sigma)\le l^{-2}$ and
$\mathrm{BK}(\sigma)\le\nu+\epsilon_T(\sigma)$ for good $\sigma$, and $t\sum_P\mathrm{tv}(P)\le l^{-2}$.

\begin{lemma}[The twist kills the Selmer groups]\label{i2:lem-killat}
(SC) and (GT) imply $\mathrm{KillAt}$.
\end{lemma}

\begin{proof}
For $\pi\in Z^{\rm good}_i$ take $\sigma$ from (SC). By (GT), $\mathrm{BK}(\sigma)\le\nu_l+\epsilon_T(\sigma)<\infty$,
so the Selmer group is finite and
$s_i(\pi)\le\nu_l+\epsilon_T(\sigma)\le\nu_l+\sum_{\sigma'}\epsilon_T(\sigma')\le\nu_l+l^{-2}$. Also
$t\sum_{v\in Q_i}\ell(O/(\zeta_v-1))/\ell(O/l)\le t\sum_P\mathrm{tv}(P)\le l^{-2}$.
\end{proof}

\begin{proposition}[existence of hub data]\label{i2:prop-hub}
Assume $\mathrm{Hypotheses}_2$. Let $\mathcal F$ be a family with $\mathrm{TwGood}(\mathcal F)$ and the family facts
$\mathrm{FactsF}(\mathcal F,C,C_T)$ (for $x>x_1$: the norm-choice facts, the commuting-endomorphism condition of
Proposition~\ref{i2:prop-cec} and the conclusion of Proposition~\ref{i2:prop-pcc} with constants $(C,C_T)$), and let
$C_2\in\R$. There is $x_0$ such that for $x\ge x_0$, $x>x_1$, every $i$ with $D=D_{x,i}$ Sha-finite and every finite set
$\mathrm{exc}$ there are: a Wach functor, a Fontaine functor and a comparison between them for $D$, and hub data $X$ for
$D$ with excluded set $\mathrm{exc}$, such that the jump parameters of $L_{x,i}$ are admissible and $X$ satisfies:
(facts) the norm-choice facts of $D$; (canon) the canonicity conditions with row constants
$(C_2,320,25,23,32,32)$; (par) admissibility of the parameter choice, and (Par); (big) $\mathsf L\ge8$; (split) the split
bound of the host is $\ge60$; (kill) $\mathrm{KillAt}$ at every twisted local index; (sha) $D$ is Sha-finite.
\end{proposition}
\noindent The canonicity predicate (it fixes each entry of $X$ to an explicit value
in terms of the size datum, the point family and the twist datum; see the proof), the norm-choice facts, $\mathrm{TwGood}$
and the admissibility of the jump parameters are not defined in this paper, so this statement is stated informally; the exact
statement is in \cite{CharacterFamily}. What is used later is: (kill), (par), (big), (split), and the explicit values of the
entries of $X$ listed in the proof. The proof is in \S\ref{i2:ssec-proof-hub}: the hub data are written down explicitly
from a point family, and $\mathrm{KillAt}$ is assembled from the conditions (SC), (GT) of a Selmer-signature selection
(Lemma~\ref{i2:lem-killat}); these inputs and the size thresholds form the smaller formally verified
Proposition~\ref{i2:prop-ptsel}.

\begin{proposition}[commuting endomorphism and clause (adm1)]\label{i2:prop-cec}
(i) Assume $\mathrm{Hypotheses}_2$ (in particular global Artin reciprocity). If $D$ satisfies the norm-choice facts and
$\mathrm{Tie4}$, then $D$ satisfies the \emph{commuting-endomorphism condition}: for every twisted $i$ with
$l_i\notin\{g_0,g_1\}$ there is an $O_i$-linear endomorphism $\alpha$ of $M_i$ commuting with $T_i$, acting on $E_x$ by
$|G|$ for $x\in Z^{\rm good}_i$ and by $0$ for $x\in Z_i\setminus Z^{\rm good}_i$.
(ii) The commuting-endomorphism condition implies (adm1) at every twisted $i$ with $l_i\notin\{g_0,g_1\}$.
\end{proposition}
\noindent Here $E_x$ is the rational Hecke eigenspace of $x$ (Lemma~\ref{i2:lem-length-le}) and, in (ii), $v^{cc}_i(\pi)=
\ell(O_i/\mathfrak c_\pi)/\ell(O_i/l_i)$ with $\mathfrak c_\pi=\{f(\pi):f\in T_i,\ f|_{X_i\setminus Z_i}=0\}$. The proof is in
\S\ref{i2:ssec-proof-cec}. Part (ii) is proved completely (Lemma~\ref{i2:lem-famcore}). Part (i) is the argument of
Proposition~\ref{i2:prop-cut}(1),(2) with $Z_i$ in place of $X^c_i$; its input on the action of the
idempotent $e_0$ on type vectors over points of $Z_i$ is Lemma~\ref{i2:lem-T3Z}, proved from explicit clauses of the datum (it is the
derivation of (T3) from $\mathrm{SerLoc}$ mentioned after Definition~\ref{i2:def-tie}). In the main chain the
commuting-endomorphism condition for the cut family is part of the family facts, which are produced for the cut data by
Proposition~\ref{i2:prop-cut}(3) and part (i).

\begin{proposition}[local conclusions along the family]\label{i2:prop-conclF}
Assume $\mathrm{Hypotheses}_2$, $C\ge0$, $C_T\in\N$, $\mathrm{TwGood}(\mathcal F)$ and the family facts
$\mathrm{FactsF}(\mathcal F,C,C_T)$ of Proposition~\ref{i2:prop-hub}. Then $\mathrm{ConclF}(\mathcal F)$.
\end{proposition}

\begin{proof}
Let $C_s$ be the maximum of $1$, $C_T$, the constants of Propositions~\ref{i2:prop-typ} and \ref{i2:prop-other} and of
Lemma~\ref{i2:lem-logG}. For large $x$ and Sha-finite $D=D_{x,i}$, take $\mathrm{exc}=\mathrm{Exc}(L_{x,i})$, the base
exceptional set (Definition~\ref{i2:def-budgets}), and choose the hub data of Proposition~\ref{i2:prop-hub} with this
excluded set; they satisfy $\mathrm{KillAt}$. Then: (adm1) is Proposition~\ref{i2:prop-cec}(ii), the
commuting-endomorphism condition being part of the family facts; (adm2) is Lemma~\ref{i2:lem-multle}; (adm3) a rank
clause; (kill) is $\mathrm{KillAt}$; (num), (mid), (bb), (ram) are Proposition~\ref{i2:prop-other}(1), (2); (two) is
Proposition~\ref{i2:prop-other}(3) (in the formalization the bound at $2$, $3$ and the small primes is a separate,
formally verified input along the family); (large) is Proposition~\ref{i2:prop-large}, whose hypothesis
$\mathrm{CongBound}$ holds for $x$ beyond a threshold by a formally verified congruence bound at large primes (from
$\mathrm{Hypotheses}_2$ and the norm-choice facts), and whose hypotheses $3\le l^*$ and $\mathrm{exc}=\mathrm{Exc}(L)$ hold by
(small) and by the choice of $\mathrm{exc}$; (typ) is Proposition~\ref{i2:prop-typ}; (Par) is Lemma~\ref{i2:lem-logG}; (exc)
holds for $\mathrm{Exc}(L)$ by a formally verified count of the exceptional primes along the family; (small) is a size
bound of the family. Every clause is monotone in the constant.
\end{proof}

\begin{definition}[The tie $\mathrm{Tie4}$]\label{i2:def-tie4}
A datum $D$ satisfies $\mathrm{Tie4}$ if: ($\chi$-tie) the characters $\chi_c$ of the datum are tied to the labels of the
types at the generator primes; (DZ, EigSpan) there is a pending datum $Q$ for $D$ (an action
$z_O:Z(O_i[G])\to\mathrm{End}_{O_i}(M_i)$ for each $i$, with the coefficient data of (T1)) satisfying the depth-zero
admissibility of types and the eigen-span clause (T4) of Definition~\ref{i2:def-tie} for $z_O$; (InePoly) the
inertia-polynomial tie (clause (a4) of Proposition~\ref{i2:prop-potcount}); and (SerLoc) the series localisation clause
recalled after Definition~\ref{i2:def-tie}. The clauses ($\chi$-tie) and (DZ) are named, not defined, here.
\end{definition}

\begin{theorem}[$\mathrm{LinkN}$]\label{i2:thm-linkN}
Assume $\mathrm{Hypotheses}_2$ and Chebotarev's theorem (Theorem~\ref{i2:thm-cheb}). Let $\mathcal F$ be a family with
$\mathrm{TwGood}(\mathcal F)$ such that for all $x>x_1$ and $i$ the member $D_{x,i}$ satisfies the Galois preconditions (for
its norm choice) and $\mathrm{Tie4}$. Then $\mathrm{LinkN}(\mathcal F)$.
\end{theorem}
\noindent The notions $\mathrm{TwGood}$ (the twist data of the members are good and conjugate along the
family), the Galois preconditions (eighteen clauses: coefficient, residual, largeness, admissibility and joint conditions on
the eigensystems and their Galois representations, and the norm-choice tie) and the clauses ($\chi$-tie), (DZ) of
$\mathrm{Tie4}$ are not defined in this paper, so the statement is stated informally as far as these terms are concerned;
the exact statement is in \cite{CharacterFamily}. In the
main chain the hypotheses are discharged where the theorem is applied; in particular the parts (T2), (T4) of the eigen-span
clause are constructed at members in \S\ref{i2:ssec-proof-cut} (Proposition~\ref{i2:prop-T2T4}).

\begin{proof}
By Proposition~\ref{i2:prop-pcc} and Proposition~\ref{i2:prop-cec}(i) there are $C\ge0$, $C_T\in\N$ such that every datum
satisfying the norm-choice facts and $\mathrm{Tie4}$ satisfies the commuting-endomorphism condition and the paired-character
condition with $(C,C_T)$; one may take $(C,C_T)=(320,40)$. For $x>x_1$ cut each member by Proposition~\ref{i2:prop-cut}: the cut is good, and the cut datum
satisfies (i) the norm-choice facts and (iii) $\mathrm{Tie4}$ (Proposition~\ref{i2:prop-cut}(3), formally verified, using
Chebotarev's theorem, local finiteness, local Tate duality and the Euler characteristic). (ii) $\mathrm{TwGood}$ passes to the
cut family because the cut does not change the host, the ranks, the places or the twist data (formally verified). Hence the
cut family $\mathcal F^c$ satisfies $\mathrm{TwGood}$ and the family facts $\mathrm{FactsF}(\mathcal F^c,C,C_T)$, and
Proposition~\ref{i2:prop-conclF} gives $\mathrm{ConclF}(\mathcal F^c)$, hence $\mathrm{LinkN}(\mathcal F^c)$
(Corollary~\ref{i2:cor-concl-link}). Theorem~\ref{i2:thm-transfer} transfers it to $\mathcal F$; its hypotheses are goodness of
the cuts, (adm3) (a rank clause of the Galois preconditions), and $\mathrm{MultLe}$ (Lemma~\ref{i2:lem-multle}, whose
hypotheses are clauses of the Galois preconditions).
\end{proof}

\subsection{The numerical divisor estimate}\label{i2:ssec-num}

Put $\log\Omega=\gamma+\log(\pi/v)-\Lambda$, $\Lambda=\Rel L'/L(1,\chi_d)$, $\gamma$ Euler's
constant and $v=\mathrm{rd}_E/2$.

\begin{lemma}\label{i2:lem-logOm}
(1) $\log\Omega=\Rel\bigl(\gamma+\log(2\pi/\mathrm{rd}_E)-L'/L(1,\chi)\bigr)$ for a character
$\chi$ modulo $d$ with $\chi(n)=\bigl(\tfrac{-d}{n}\bigr)$ for odd $n$.
(2) If, for $\psi\in\{\chi_d,\psi_q,\chi_d\psi_q\}$ and every primitive character inducing $\psi$
modulo $8qd$, $|L'/L(1,\chi)|\le4\log x$ (as provided by the analytic pair condition of the host,
Section~\ref{i5:sec}), then $|\Lambda|\le4\log x$; hence $|\log\Omega-\gamma-\log(\pi/v)|\le x^\varepsilon$ for
large $x$.
\end{lemma}

\begin{proof}
(1) Take $\chi=\chi_d$; for $d\equiv3\pmod4$ and odd $n$, quadratic reciprocity gives
$\bigl(\tfrac{-d}{n}\bigr)=\bigl(\tfrac nd\bigr)$; and $2\pi/\mathrm{rd}_E=\pi/v$.
(2) $\chi_d$ is nontrivial, hence primitive of conductor $d$, and induces the Kronecker
character modulo $8qd$ (same values on $a$ prime to $8qd$, by the reciprocity just used). So
$|\Lambda|\le|L'/L(1,\chi_d)|\le4\log x$, and $4\log x\le x^\varepsilon$ eventually.
\end{proof}

\subsection{From the divisor bound to condition (P1)}\label{i2:ssec-P1}

The divisor bound is consumed through \emph{division kernel packages}: data attached to each
member which tie the quantities of the column package to those of the member. Write $\mathrm{DIVinN}(\mathcal F)$ for the
conclusion of Theorem~\ref{i2:thm-div} (there are $C,x_0$ with $\eta\Sigma_{L_{x,i}}(j)\le D_0\,C\,x^\varepsilon(B_0+N_c)$ for
$x\ge x_0$, all $i$ and lifts $j$), and $\mathrm{LinkEta}(\mathcal F)$ for the statement that, for $x$ large, the local
indices of every member cover its twisted places.

\begin{definition}[division kernel constant bounds]\label{i2:def-DKCB}
Fix division kernel data $\mathrm{Rk}$ at the members beyond some $x_0$, and $C_{\rm abs},C_H\in\R$.
$\mathrm{DKCB}(\mathcal F,\mathrm{Rk},C_{\rm abs},C_H)$ is the conjunction of ten clauses:
(pkg) if $\mathrm{DIVinN}(\mathcal F)$, there are $C,x_0$ such that for $x\ge x_0$ every member carries a division kernel
package for its data with constants $(C,C_{\rm abs})$;
(lq2) eventually, the second admissibility clause of the local quadratic parts (named (lq2) to avoid a clash with
(adm2) of Definition~\ref{i2:def-AdmT} and with the datum clauses);
(cut) eventually, the cut-up condition for every used weight;
(c3) for some $C$, eventually the third complement bound with constant $C$;
(hh) if $\mathrm{DIVinN}(\mathcal F)$, the frame-number bound with exponent $C_H$;
(c2), (c5) the second complement and the complement-count conditions along the family;
(bu), (rr) eventually, rationality of the Bessel data and the full $RR$-condition;
(la) eventually, $L_a=\log(a\cdot\mathrm{rd}_E)\ge0$.
\end{definition}
\noindent The division kernel package is described informally in Definition~\ref{i2:def-divpkg}(c) (for the column data given by $\mathrm{Rk}$);
that definition names, without defining, the clauses of the indefinite package and the input record of the divisor
inequality. The
predicates in (lq2), (cut), (c3), (hh), (c2), (c5), (bu), (rr) belong to the division data of another part of the paper
and are named, not defined, here; to that extent this definition is stated informally, and the exact definition is in
\cite{CharacterFamily}. The
constants $C_0\ge28$, $c_e$ and the inequalities $50C_0+10\le C$, $60c_e\le C$, $C_{\rm abs}\ge1$ are hypotheses of the
constructor of a package, not clauses of $\mathrm{DKCB}$.

\begin{proposition}[from the divisor bound to condition (P1)]\label{i2:prop-P1}
Assume $\mathrm{DKCB}(\mathcal F,\mathrm{Rk},C_{\rm abs},C_H)$, $\mathrm{LinkEta}(\mathcal F)$ and $\mathrm{DIVinN}(\mathcal F)$.
Then there are $C,x_0$ such that for all $x\ge x_0$, $i\in I(x)$ the nine clauses of the division link hold at
$(L_{x,i},C,C_{\rm abs},C_H)$; the seventh of them is: for every word $u$,
\[
\mu^\vee+h(u)-F\le\mathrm{MAIN}(u)+C\,|T|\,a\,L_a+C\,m\,(B_0+N_c)\,x^{C_{\rm abs}\varepsilon},
\]
with $\mathrm{MAIN}(u)=\sum_{t\in T}\bigl(m(\phi(a_t)-a_t\log\Omega)-a_t\log T_u\bigr)$, $\phi(y)=y\log y-y$, and
$\mu^\vee$, $h(u)$, $F$ the dual height, the word heights and the frame number of the division data.
\end{proposition}
\noindent The proof is in \S\ref{i2:ssec-proof-P1}: an assembly of the clauses of $\mathrm{DKCB}$, proved, from the
Proposition~\ref{i2:prop-wordfacts} (the divisor inequality of one package, and the reduction of
clauses one, six and seven of the division link to it; its part (i) is proved from Lemmas~\ref{i2:lem-divclauses} and
\ref{i2:lem-tietransfer}). $\mathrm{LinkEta}$ is a hypothesis of the formal statement that is
not used in its proof. The division package is described informally in Definition~\ref{i2:def-divpkg}, which names some of its clauses without defining
them, and the remaining predicates of $\mathrm{DKCB}$ are named; so the statement is stated informally as far as
those are concerned (the exact statement is in \cite{CharacterFamily}); the displayed clause is exact.

Proposition~\ref{i2:prop-P1} is condition (P1) of
Definition~\ref{M:def:package}, with $c_{\rm div}=C$ and divisor exponent $e_D=C_{\rm abs}$ ($e_D=2$
in the top-level instance), and the displayed inequality is its content. The chain is: Section~\ref{i5:sec} gives $\mathrm{HypV}$;
Theorem~\ref{i2:thm-linkN} gives $\mathrm{LinkN}$; Theorem~\ref{i2:thm-div} gives the divisor
bound; Proposition~\ref{i2:prop-P1} gives (P1).

\subsection{The local count and the Selmer bound for the congruence module}\label{i2:ssec-pf-local}

\subsubsection{Proof of Proposition~\ref{i2:prop-potcount}}\label{i2:ssec-proof-potcount}

Throughout, $O$ is the ring of integers of a finite extension $L$ of $\Q_l$ with uniformiser
$\varpi$, $\ell$ is $O$-length, $K/\Q_l$ is finite and $G=G_K$. For a $G$-stable lattice $T$ in a
representation $U$, $Z^1_{\rm div}(U/T)$ denotes the locally constant cocycles $G\to U/T$, and
$H^1_f(U/T)$ the image of the finite cocycles of $U$ plus the coboundaries of $U/T$. A \emph{Tate
dual} of $U$ is a representation $U'$ with a perfect pairing $P:U\times U'\to L$,
$P(gu,gu')=\chi_{\rm cyc}(g)P(u,u')$; the dual lattice is $T'=\{u':P(T,u')\subseteq O\}$.

\begin{definition}[Pushout span]\label{i2:def-push}
For representations $X,U$ and a lattice $T\subset U$, $\mathrm{Push}(X,U,T)$ is the $O$-span of the
cochains $g\mapsto\varphi(y_g)$, where $Y$ is a $G$-stable lattice in $X$, $\varphi:Y\to U/T$ is
$G$-equivariant and $y$ is a finite cocycle of $Y$.
\end{definition}

The local inputs are the following, all published: (i) rational orthogonality of
Theorem~\ref{i2:thm-bk}: a cocycle of a de Rham $U$ is finite iff its cup product with every
finite cocycle of $U'$ is a coboundary of $L(1)$; (ii) its integral form for divisible quotients
(Bloch--Kato, \emph{loc.\ cit.}, \S3): $c\in Z^1_{\rm div}(U/T)$ lies in $H^1_f(U/T)$ iff its cup
product with every finite cocycle of $T'$ is a coboundary of $(L/O)(1)$; (iii)
$H^2(G,(L/O)(1))\cong L/O$ (Theorem~\ref{i2:thm-tate}, passed to the limit over $\Z/l^r(1)$); (iv)
$H^1(G,T)$ is finitely generated over $O$ (Theorem~\ref{i2:thm-euler}, by $l$-adic completeness);
(v) Tate duals of de Rham representations are de Rham (Theorem~\ref{i2:thm-fontaine}). For
Fontaine's theory we use, from Theorem~\ref{i2:thm-fontaine}: (F1) in $0\to A_0\to A_1\to A_2\to0$
with $A_1$ crystalline, $A_0$ and $A_2$ are crystalline, and a split extension of crystalline
representations is crystalline; (F2) twists by crystalline characters and (F3) Tate duals of
crystalline representations are crystalline; (F4) crystalline implies de Rham.

\begin{lemma}[Length through a pairing]\label{i2:lem-pair}
Let $\mathrm{pr}:A\times H\to L/O$ be $O$-bilinear, $M\subseteq H$, $S\subseteq A$ submodules,
$\alpha\in O$, and $\mathrm{ann}(M)=\{s:\mathrm{pr}(s,M)=0\}$. If $\alpha S\subseteq\mathrm{ann}(M)$, then
$\ell(S/(S\cap\mathrm{ann}M))\le\ell(M/\alpha M)$.
\end{lemma}

\begin{proof}
$s\mapsto\mathrm{pr}(s,\cdot)|_M$ embeds $S/(S\cap\mathrm{ann}M)$ into $\mathrm{Hom}_O(M/\alpha M,L/O)$, since
$\mathrm{pr}(s,\alpha m)=\mathrm{pr}(\alpha s,m)=0$. For $Q$ of finite length,
$\ell\,\mathrm{Hom}_O(Q,L/O)=\ell(Q)$ by induction on $\ell(Q)$ ($L/O$ is injective and
$\mathrm{Hom}(O/\varpi,L/O)\cong O/\varpi$).
\end{proof}

\begin{lemma}[Length of $M/l^eM$]\label{i2:lem-quot}
Let $M$ be finitely generated over $O$, $N\subseteq(L/O)^d$ and $\delta:N\to M$ $O$-linear with
$M_{\rm tors}\subseteq\delta(N)$. Then $\ell(M/l^eM)\le(\mathrm{rk}_OM+d)\,e\,\ell(O/l)$.
\end{lemma}

\begin{proof}
$M\cong O^{\rm rk}\oplus Q$ with $Q$ torsion, and $\ell(Q/l^eQ)\le\mu(Q)e\ell(O/l)$, $\mu$ the minimal
number of generators. A submodule of $(L/O)^d$ is $(L/O)^s\oplus\bigoplus_{j\le d-s}O/\varpi^{n_j}$;
its divisible part maps to $0$ in the finite $Q$, so $\mu(Q)\le d$.
\end{proof}

\begin{lemma}[Torsion of $H^1_f$]\label{i2:lem-tors}
For a lattice $T\subset U$, the torsion of $H^1_f(T)$ lies in the image of
$\delta:N\to H^1_f(T)$, $v\mapsto[g\mapsto gv-v]$, where
$N=\{v\in U:gv-v\in T\ \forall g\}/T\subseteq U/T\cong(L/O)^{\mathrm{rk}T}$.
\end{lemma}

\begin{proof}
Coboundaries of $U$ are finite. If $\varpi^kc=(g\mapsto gt-t)$, $t\in T$, then $v=\varpi^{-k}t$
represents an element of $N$ with $\delta(v)=[c]$.
\end{proof}

\begin{lemma}[Absorption]\label{i2:lem-absorb}
Let $U$ be de Rham with Tate dual $U'$ and dual lattices $T,T'$; $X$ a representation with Tate
dual $X'$; $Y\subset X$ a lattice, $\varphi:Y\to U/T$ equivariant, $y$ a finite cocycle of $Y$.
If the family generated by $X,U$ has $(\mathrm{Lift}_e)$, or $X$ is de Rham and the family generated
by $X',U'$ has $(\mathrm{Lift}_e)$, then $l^e(\varphi\circ y)\in H^1_f(U/T)$.
\end{lemma}

\begin{proof}
Lift $\varphi$ to an $L$-linear $\psi:X\to U$; then $g\psi(v)-\psi(gv)\in T$ for $v\in Y$.
\emph{First case.} $\Lambda=\{(x,u):x\in Y,\ u-\psi x\in T\}$ is a $G$-stable lattice in $X\oplus U$
mapping onto $Y$. By $(\mathrm{Lift}_e)$, $l^ey=\mathrm{pr}_1\tilde y+\partial t$ with $\tilde y$ a
finite cocycle of $\Lambda$, $t\in Y$; modulo $T$,
$l^e\varphi(y_g)\equiv\mathrm{pr}_2\tilde y_g+g\varphi(t)-\varphi(t)$, and $\mathrm{pr}_2\tilde y$ is a
finite cocycle of $U$.
\emph{Second case.} With $\psi^t:U'\to X'$ the transpose, $\Lambda'=\{(f,x')\in U'\oplus X':f\in T',\
P_X(v,x')-P_U(\psi v,f)\in O\ \forall v\in Y\}$ is a $G$-stable lattice (by the cyclotomic
equivariance of the pairings and $g\psi v-\psi gv\in T$) mapping onto $T'$. By (ii) and (v) it suffices
that for each finite cocycle $d$ of $T'$ the cup product of $d$ with $l^e(\varphi\circ y)$ is a
coboundary of $(L/O)(1)$. By $(\mathrm{Lift}_e)$, $l^ed=\mathrm{pr}_1\kappa+\partial b$ with $\kappa$
finite in $\Lambda'$. The cup product of $\mathrm{pr}_2\kappa$ (finite in the de Rham $X'$) with $y$ is
a coboundary of $L(1)$ by (i); and it differs from the cup product of $\mathrm{pr}_1\kappa$ with
$\varphi\circ y$ (computed with the lift $\psi\circ y$) by a cochain with values in $O$, by the
definitions of $T'$ and $\Lambda'$.
\end{proof}

\begin{proposition}[Pushout bound]\label{i2:prop-pushout}
Under (i)--(v), with $X,U,T,T'$ and the lifting hypothesis as in Lemma~\ref{i2:lem-absorb},
\[
\ell\bigl((\mathrm{Push}(X,U,T)+H^1_f(U/T))/H^1_f(U/T)\bigr)\le e\bigl(\mathrm{rk}_OH^1_f(T')+\mathrm{rk}_OT\bigr)\ell(O/l).
\]
\end{proposition}

\begin{proof}
Put $S=\mathrm{Push}(X,U,T)\subseteq Z^1_{\rm div}(U/T)$ ($T$ is open in $U$) and $H=H^1_f(U/T)$; then
$l^eS\subseteq H$ by Lemma~\ref{i2:lem-absorb}. The cup product into $H^2(G,(L/O)(1))\cong L/O$ (iii)
is a pairing $Z^1_{\rm div}(U/T)\times H^1(G,T')\to L/O$, and by (ii) the annihilator of
$M=H^1_f(T')$ is $H\cap Z^1_{\rm div}(U/T)$. Lemma~\ref{i2:lem-pair} gives
$\ell((S+H)/H)\le\ell(M/l^eM)$. By (iv) $M$ is finitely generated, by Lemma~\ref{i2:lem-tors} its
torsion comes from a submodule of $(L/O)^{\mathrm{rk}T'}$ with $\mathrm{rk}T'=\mathrm{rk}T$, and
Lemma~\ref{i2:lem-quot} concludes.
\end{proof}

\begin{proof}[Proof of Proposition~\ref{i2:prop-potcount}(b)]
A cut cocycle of $M$ at $w$ has the form $g\mapsto\phi(\gamma_gE-E)$ for a datum
$\mathcal D=(x_1,\dots,x_a;E;Y;\phi)$: points of $X_s$ (repetitions allowed), a vector $E$ of $\bigoplus_jr_{x_j}$ (members of
the twisted family), a
finitely generated $G$-stable $O$-submodule $Y$ containing all $\gamma E-E$, and an equivariant
$\phi:Y\to(L/O)^k$ (with the action of $M$). Concatenating data shows that the sets $P(\mathcal D)$
of cut cocycles form a directed system, and the length of a directed union of spans modulo a fixed
submodule is the supremum of the lengths; so it suffices to bound one $P(\mathcal D)$. Let $U=M$ (a
representation of $G=G_{E_w}$, $E_w=\Q_l$), $T_M\cong O^k$ its lattice, and $X=Y\otimes L$.
\emph{(1) $U$ is de Rham:} by (Pot$_n$) the member $r_{x_0}\supseteq M$ is crystalline over $K'$, so
$U|_{G_{K'}}$ is (F1), hence de Rham (F4) and $U$ is de Rham.
\emph{(2) Rank:} take $U'=\chi_{\rm cyc}\otimes{}^tU(\cdot^{-1})$ with the dot product and $T'=O^k$;
$U'$ is de Rham by (F3), (F4). By the dimension formula of Theorem~\ref{i2:thm-bk},
$\dim H^1_f(U')=\dim H^0(U')+\dim D_{\rm dR}(U')/\mathrm{Fil}^0\le\dim D_{\rm dR}(U')=k$, since
$H^0(U')\otimes K\hookrightarrow\mathrm{Fil}^0D_{\rm dR}(U')$; and
$\mathrm{rk}_OH^1_f(T')=\dim_LH^1_f(U')$ because $H^1_f(T')\otimes L\to H^1_f(U')$ is an isomorphism
(every continuous cochain has bounded image, so a multiple of any finite cocycle, or of a
coboundary vector, is integral).
\emph{(3)} $X'=\chi_{\rm cyc}\otimes{}^tX(\cdot^{-1})$ is a Tate dual of $X$.
\emph{(4) Lifting:} a representation with jointly injective maps to copies of $X$ and $U$ is, via the
coordinate projections of $Y\subseteq\bigoplus_jr_{x_j}$ and the embedding $M\hookrightarrow r_{x_0}$ (which
has a linear left inverse), a subfamily of the twisted family; so $(\mathrm{Lift}_e)$ applies.
\emph{(5) Inclusion:} $y_g=\gamma_gE-E$ is a cocycle of $Y$ that becomes a coboundary in
$A=\bigoplus_jr_{x_j}$, crystalline over $K'$. Such a cocycle is finite: if $E\notin X$, the
extension $X\oplus L$ with $g(w,c)=(gw+cy_g,c)$ embeds into $A$, hence is crystalline over $K'$
(F1); counting dimensions of $D_{\rm cris}$ in $0\to X\to X\oplus L\to L\to0$ gives a lift
$(z_1,1)$ of $1$, whose invariance says $y_g\otimes1=g(-z_1)-(-z_1)$ in $X\otimes B_{\rm cris}$ for
$g\in G_{K'}$; and a cocycle of $G_K$ vanishing on the normal subgroup $G_{K'}$ of finite index $n$
is a coboundary (average over $G_K/G_{K'}$ and divide by $n$). So $P(\mathcal D)\subseteq\mathrm{Push}(X,U,T_M)$.
Proposition~\ref{i2:prop-pushout} gives the bound $e(k+k)\ell(O/l)=2ke\,\ell(O/l)$. (This argument uses the
first case of Lemma~\ref{i2:lem-absorb} only; it does not give part (c), see the remark after
Proposition~\ref{i2:prop-potcount}.)
\end{proof}

\begin{proof}[Proof of Proposition~\ref{i2:prop-potcount}(a)]
The hypotheses (a1)--(a5) supply: the type-locus local--global compatibility (for $x\in X_s$, $r_x$ is de Rham at
$w$ with monodromy zero, hence potentially crystalline, with inertia polynomial
$\prod_{y\in M_x}(X-\omega_y)$, $M_x$ the multiset of roots of $S\cdot(X-u_x)^t$, inverted at the
conjugate place, where $S\in\mathbb F_g[X]$, $S(0)\ne0$, $\deg S\le r$, is common to the type locus,
$u_x\in\{c,c^{-1}\}$ and $c$ is the label of $\chi$ at $w$); the inertia-polynomial tie (for $g$ in
inertia a power of the inertia polynomial evaluated at $g$ kills $D_{\rm pcris}$); and stability of
potential crystallinity under twists by potentially crystalline characters, with the inertial
characters multiplied.
\emph{(1) Labels.} Put $Y=\mathrm{roots}(S)\cup\{c,c^{-1}\}$ (inverted at the conjugate place); then
$|Y|\le r+2\le n_V+1$, and every $y\in Y$ satisfies $y^{N_0}=1$, $N_0$ a multiple of
$|\mathrm{GL}_r(\mathbb F_g)|$ (a root of a polynomial of degree $\le r$ lies in some $\mathbb F_{g^b}$,
$b\le r$). The character $\chi$ is de Rham, so has monodromy zero, with inertia polynomial
$X-\omega_u$, $u=c^{\pm1}$.
\emph{(2) Common field.} Let $K'$ be the Galois closure of the compositum of fields over which the
finitely many $r_x$ and $\chi$ become crystalline and over whose inertia the $\omega_y$ are trivial.
\emph{(3) Inertial characters.} Transport each $\omega_y$ by the Teichm\"uller lift to a character
$\Psi_y$ of the inertia of $K_w$ with values in $\mu_{N_0}(L)$, multiplicative in $y$; extend it to
$I(K'/\Q_l)=\Gamma\cdot G_{K'}$ (Definition~\ref{i2:def-locfam}; $\Gamma$ the inertia group of $K_w=\Q_l$) by
$\mathrm{ch}_y(su)=\Psi_y(s)$ (well defined because $\Psi_y$ is trivial on
$\Gamma\cap G_{K'}$, multiplicative because $G_{K'}$ is normal); the orders divide that of an
element of $\overline{\mathbb F}_g^\times$, hence are prime to $l$. All $\mathrm{ch}_y$ are powers of one
character $c_\delta$; choose $\gamma$ with $c_\delta(\gamma)$ generating $c_\delta(I)$. Then the values
$\mathrm{ch}_y(\gamma)$ are distinct and every $g\in I$ is $\gamma^mu$ with all $\mathrm{ch}_y(u)=1$. Since
$G_{K'}$ acts trivially on $D_{\rm pcris}$, the separable $X^n-1$ kills it under every $g$; combined
with the inertia-polynomial tie this gives that $u$ acts trivially and $D_{\rm pcris}$ is the sum of
the $\gamma$-eigenspaces for the values $\mathrm{ch}_y(\gamma)$, i.e.\ inertia acts through
$\Sigma_0=\mathrm{ch}(Y)$, $|\Sigma_0|\le n_V+1$; and on $D_{\rm pcris}(\chi)$ through $c_0=\mathrm{ch}(u)\in\Sigma_0$.
\emph{(4) Twist.} The members $r_x\otimes\chi^{-1}$ are crystalline over $K'$ with inertia through
$\Sigma_0c_0^{-1}\ni1$. As $E_w=\Q_l$, a Frobenius $\mathrm{Fr}\in G_{\Q_l}$ exists. Let $\Sigma$ consist of
$1$ and the characters of $\Sigma_0c_0^{-1}$ with a nonzero eigenspace in some member; inertia acts
through $\Sigma$, and $\Sigma$ is $\mathrm{Fr}$-stable, since $\mathrm{Fr}$ carries an $a$-eigenvector to a
nonzero $a(\mathrm{Fr}^{-1}\cdot\mathrm{Fr})$-eigenvector of $D_{\rm pcris}$, which must lie in a summand of
$\Sigma_0c_0^{-1}$ (eigenspaces of distinct characters are independent). This is (Pot$_{n_V+1}$).
\end{proof}

\subsubsection{Proof of Proposition~\ref{i2:prop-selmer}}\label{i2:ssec-proof-selmer}

The proof has two halves. The first, a determinant bound for the congruence ideal by
$\mathrm{Hom}$-lengths into the diagonal blocks of a Ribet-type lattice, is proved here. The second,
the bound of each $\mathrm{Hom}$-length by a Selmer length plus local terms, is
Proposition~\ref{i2:prop-homterms}; its algebraic part is proved here, and its arithmetic
input is the smaller formally verified Proposition~\ref{i2:prop-colconds}.

\begin{lemma}[Block-ring decomposition: determinant form]\label{i2:lem-blockring}
Let $O$ be a discrete valuation ring with fraction field $K$, valuation $v$, residue field $k$,
$V_O=K/O$. Let $T\subseteq O^X$, $T'\subseteq O^{X'}$ be subalgebras, $Z\subseteq X$, $p:X'\to X$,
$b_1,\dots,b_q\in O^{X'}$, $J=\{c\in T:c|_{X\setminus Z}=0\}$, $v_{cc}(x)=\inf_{c\in J}v(c(x))$, and
$R'$ a left-Noetherian $O$-algebra of functions $X'\to M_d(K)$. Assume:
(H1) every $f'\in T'$ is $\sum_m(a_m\circ p)b_m$ with $a_m\in T$;
(H2) for $x\notin Z$, $\sum_ma_mb_m(x')=0$ for all $x'\in p^{-1}(x)$ forces $a=0$;
(H3) for every $\Theta\in M_g(R')$ the function $x'\mapsto\det\Theta(x')$ lies in $T'$;
(H4) $\rho_\pi=\mathrm{ev}_{\pi'}:R'\to M_d(O)$ is block upper triangular with surjective diagonal
blocks $\rho_\kappa:R'\to M_{d_\kappa}(O)$;
(H5) $N$ is an $R'$-module with $g_0$ generators, and for every $x'$ with $p(x')\notin Z$ there is a
nonzero $R'$-map $N\to K^d$ (action through $\mathrm{ev}_{x'}$);
(H6) $|k|>g_0\sum_\kappa d_\kappa$.
Then $v_{cc}(p(\pi'))\le\sum_\kappa\ell\,\mathrm{Hom}_{R'}(N,V_O^{d_\kappa})$.
\end{lemma}

\begin{proof}
\emph{Presentation.} $N=\mathrm{coker}(P)$ with $P\in M_{m\times g_0}(R')$, and
$\mathrm{Hom}_{R'}(N,M)=\ker(P:M^{g_0}\to M^m)$.
\emph{One matrix.} For $A\in M_{M\times G}(O)$, by elementary divisors $AO^G=\bigoplus_{q\le n}a_qOU_{f(q)}$
with $U$ invertible; if $C'\in M_{G\times M}(O)$ makes the minor $((C'U)_{e(p),f(q)})$ a unit, then
$\ker(C'A)=\ker A$ on $V_O^G$; such $C'$ exist.
\emph{Generic left factor.} Choose $r_{\kappa,a,b}\in R'$ with $\rho_\kappa(r_{\kappa,a,b})=E_{ab}$ (H4),
and let $C(\tau)=\sum_\jmath\tau_\jmath c_\jmath\in M_{g_0\times m}(R')$ with $c_\jmath$ the matrix units carrying
these elements. For each $\kappa$ the residue of the relevant minor of $\rho_\kappa(C(\tau))U_\kappa$ is a
polynomial $Q_\kappa$ of degree $\le g_0d_\kappa$, nonzero (it is a unit at the $\tau$ realising
$C'_\kappa$). By (H6), $\prod_\kappa Q_\kappa$ has a nonzero value at a point of $k$; lifting it gives
$C$ with $\ker\rho_\kappa(CP)=\ker P_\kappa$ on $V_O^{g_0d_\kappa}$ for all $\kappa$, so
$\ell\ker\rho_\kappa(CP)=\ell\,\mathrm{Hom}_{R'}(N,V_O^{d_\kappa})$.
\emph{Determinant.} By (H3), $f=\det(CP)\in T'$; by (H5), $f(x')=0$ when $p(x')\notin Z$. By (H1),
$f=\sum_m(a_m\circ p)b_m$, and (H2) gives $a_m\in J$; hence $v(f(\pi'))\ge v_{cc}(p(\pi'))$.
\emph{Blocks.} By (H4), $f(\pi')=\prod_\kappa\det\rho_\kappa(CP)$, and $v(\det A)=\ell\ker(A$ on
$V_O^n)$ for square $A$ (Smith form). Summing gives the claim.
\end{proof}

\begin{lemma}[Large residue fields]\label{i2:lem-bigres}
For a discrete valuation ring $O$ and $M\in\N$ there is a discrete valuation ring $O'\supseteq O$,
finite free over $O$, with $v_{O'}|_O=v_O$ and residue field of more than $M$ elements.
\end{lemma}

\begin{proof}
For finite $k$ take $k(\alpha)$ with more than $M$ elements, $\bar g$ the minimal polynomial of
$\alpha$, $g$ a monic lift and $O'=O[X]/(g)$. It is free, a domain ($g$ irreducible), integral over
$O$, and its unique maximal ideal is $\varpi O'$ (the kernel of reduction to $k(\alpha)$), so it is a
discrete valuation ring with uniformiser $\varpi$.
\end{proof}

\begin{proposition}[Lattice bound]\label{i2:prop-latbound}
Let $X$ be finite, $T\subseteq O^X$, $r:G\to M_d(O^X)$ a family with block decomposition into a block
$0$ of size $r$ and blocks $1,\dots,t$ of size one, $\pi\in X$, and $e_B,e_1,s,g_0$ idempotent data
satisfying the primitive-congruence hypotheses, absolute irreducibility of the residual block $0$ at
$\pi$, and the distinct-eigenvalue hypothesis (the block polynomial
$P_x=\det(Xe_B(x)-X_0(x)+1-e_B(x))$, $X_0=e_Br(g_0)e_B+s$, has $t$ distinct roots for every $x$).
Let $R'$ be the associated ring of functions, $Z\subseteq X$, and $N$ an $R'$-module with $g$
generators and a nonzero map to the fibre at every point over $X\setminus Z$. Then
$v_{cc}(\pi)\le\sum_{c=0}^t\ell\,\mathrm{Hom}_{R'}(N,V_O^{d_c})$.
\end{proposition}

\begin{proof}
\emph{Base change.} Take $O'$ from Lemma~\ref{i2:lem-bigres} with $M=gd$. All hypotheses are
polynomial identities or statements about roots, preserved by the injective $O\to O'$; the
congruence valuation at $\pi$ and each $\mathrm{Hom}$-length are multiplied by the same ramification
index $1$ (Smith form). So it suffices to work over $O'$, where (H6) holds.
\emph{Ring data.} Points $X'$ are pairs $(x,\alpha)$ with $\alpha$ an ordering of the roots of $P_x$,
which lie in $O$ since $O$ is integrally closed; $p(x,\alpha)=x$; $T'$ is generated by $T\circ p$ and
the root functions; $R'$ is the $T'$-span of the values of $r$; $\pi'=(\pi,(\chi_1(g_0)\varepsilon(g_0)^{-i})_i)$.
(H1), (H2): $T'$ is the image of the universal splitting algebra $A$ of $P$ over $T$, free of rank
$t!$ (built by adjoining one root at a time to $T[X]/(P)$); at a fixed $x$, $A\otimes O\otimes K$ is the
product of copies of $K$ indexed by the $t!$ orderings of the distinct roots, so evaluation on the
fibre is injective. (H3): the coefficients of $\det(1-u\,r(g))$ lie in $T'$, and an elimination
argument in $R'[u]/(u^{N+1})$ (peel off the lowest-degree term $x$ via $(1-x)^{-1}=\sum_{k\le N}x^k$)
shows that every coefficient of $\det(1-\sum_jc_ju^{d_j}r(\gamma_j))$ lies in $T'$; hence so do the
coefficients of generic determinants $\det(\sum_iY_ia_i)$, $a_i\in R'$, and, by Schur complements over
$T'[[u]]$, $\det\Theta$ for block matrices $\Theta$ over $R'$. (H4): evaluation at $\pi'$ is block
upper triangular; the blocks of size one are onto via scalars, and block $0$ is onto by Burnside's
theorem for the absolutely irreducible residual block and Nakayama's lemma. $R'$ is left Noetherian,
being a submodule of the finite free $M_d(O)^{X'}$ ($X'$ is finite). (H5) is the hypothesis on $N$.
Lemma~\ref{i2:lem-blockring} concludes.
\end{proof}

\begin{proposition}[the $\mathrm{Hom}$-terms]\label{i2:prop-homterms}
In the setting of Proposition~\ref{i2:prop-selmer}, let $N$ be the module spanned by the column
cocycles (Definition~\ref{i2:def-colmod}), $M_0=\rho_\pi$, $M_1$ the trivial character, $M_c=\varepsilon^{1-c}$ ($2\le c\le t$), and
$e_O=\ell(O/l)$. Then $N$ satisfies the hypothesis of Proposition~\ref{i2:prop-latbound}, and
\[
\begin{aligned}
\sum_{c=0}^t\ell\,\mathrm{Hom}(N,M_c\otimes K/O)\le{}&e_O\Bigl(s_i(\pi)+\sum_{e=1}^{t-1}\ell_D(l,e)+v_l(h_E)\Bigr)+\sum_{\rm rows}\ell_w\\
&+\sum_{\rm cols}\bigl(e_Oe'_w+j_w(e_Oe'_w+v_\delta)+2v_\delta+h^0_w\bigr)\\
&+te_Oc_B\sum_{\rm rows}(\deg_w+2)+4(N_0+t)(|\mathrm{col}|+1)v_\delta\\
&+t|\mathrm{col}|H^0+t\sum_{v\in Q_i}v(\zeta_v-1),
\end{aligned}
\]
with $N_0=r$ the rank of $M_0$, $\ell_w$ the row lengths, $e'_w$ the lift exponents, $j_w$ the numbers of column generators,
$h^0_w$ the lengths of the $\rho_\pi$-invariants, $c_B$ the cyclotomic block number and
$v_\delta=e_O\sum_{0<k<t}(s_0+v_l(k))$, where $\varepsilon(g_0)=1+l^{s_0}y$.
\end{proposition}
\noindent The proof occupies the rest of this subsection: an algebraic bound for each block
(Theorem~\ref{i2:thm-homabs}, proved), the arithmetic properties of the column module
(Proposition~\ref{i2:prop-colconds}, formally verified), and the computation of $v_\delta$
(Lemma~\ref{i2:lem-vdelta}).

\medskip

Let $O$ be a discrete valuation ring with fraction field $K$, uniformiser $\varpi$, normalised
valuation $v$, and $V_O=K/O$; $l\in O$ is a nonunit, $e_O=\ell(O/l)=v(l)$. Let $A$ be an
$O$-algebra (in the application, the ring $R'$ of Proposition~\ref{i2:prop-latbound}), $T'$ a
commutative ring of scalars of $A$ (central), $G$ a group, $r:G\to A^\times$ and
$\chi:G\to T'^\times$ homomorphisms, and put $u(g)=\chi(g)^{-1}r(g)$; $g\mapsto u(g)$ is again
multiplicative. Fix $E,e_1\in A$ and $\delta\in T'$ with $E^2=\delta E$, $e_1^2=e_1$, and a finite
$s_E\subseteq G$ and $t_E:G\to T'$ with $E=\sum_{g\in s_E}t_E(g)r(g)$. Put
\[
y_g=(u(g)-1)E\in A,\qquad Y=\sum_gAy_g .
\]
Then $y_{gh}=y_g+u(g)y_h$.

\begin{definition}[Column module]\label{i2:def-colmod}
Let $\iota$ be a finite set (the columns), and for $l\in\iota$ let $D_l\subseteq G$ be a subgroup,
$j_l\in\N$, $K_{l,a}:D_l\to A$ ($a<j_l$) and $k^0_l:D_l\to AE$. The \emph{column module}
$N=N(D,j,K,k^0)$ is the left $A$-module generated by $Y$ (as a submodule) and by symbols $\nu_l$,
$\sigma_{l,a}$ ($a<j_l$), $\sigma^0_l$, subject to
\[
E\nu_l=0,\qquad(1-e_1)\sigma_{l,a}=0,\qquad(\delta-E)\sigma^0_l=0,\]
\[y_g+(u(g)-1)\nu_l-\delta\sum_aK_{l,a}(g)\sigma_{l,a}-\delta k^0_l(g)\sigma^0_l=0\quad(g\in D_l).
\]
For a left $A$-module $M$ on which $u(G)$ acts, the map
$\Psi:\mathrm{Hom}_A(N,M)\to Z^1(G,M)$, $\Psi(\varphi)(g)=\varphi(y_g)$, takes values in
cocycles for the action $g\cdot m=u(g)m$, by $y_{gh}=y_g+u(g)y_h$.
\end{definition}

The last relation says that on $D_l$ the cocycle $g\mapsto y_g$ equals the coboundary of $-\nu_l$
plus $\delta$ times a combination of the column generators; these are the column cocycles of
Proposition~\ref{i2:prop-homterms}. In the application $M=V_O^{d}$ is a diagonal block $M_c$,
with $A$ acting through $\rho_c$.

\begin{lemma}[Kernel and coboundaries]\label{i2:lem-psiker}
Let $M=V_O^d$ with an $A$-action commuting with $O$, on which $\delta$ acts as $\delta_O\in O$.
\begin{enumerate}
\item $\ker\Psi=\{\varphi:\varphi|_Y=0\}$.
\item If $E$ acts as $0$ on $M$ and $\tau=\sum_{g\in s_E}t_E(g)\chi(g)$ acts as $\delta_O$, then
$\Psi(\varphi)=dv$ (that is, $\varphi(y_g)=u(g)v-v$ for all $g$) implies $\delta_O^2v=0$.
\end{enumerate}
\end{lemma}

\begin{proof}
(1) $Y$ is spanned by the $y_g$. (2) Put $y_1=\sum_{g\in s_E}t_E(g)\chi(g)y_g\in Y$. Since
$t_E(g)\chi(g)u(g)=t_E(g)r(g)$, $y_1=(E-\tau)E$, and
$\varphi(y_1)=\sum t_E(g)\chi(g)(u(g)v-v)=Ev-\tau v=-\delta_Ov$. As $\tau$ is central and $E^2=\delta E$,
$Ey_1=E(E-\tau)E=\delta(E-\tau)E=\delta y_1$. Hence
$0=E\varphi(y_1)=\varphi(Ey_1)=\delta_O\varphi(y_1)=-\delta_O^2v$.
\end{proof}

For a subgroup $D\subseteq G$ write $h^0_D(M)=\ell(M^{D})$, invariants for the $u$-action.

\begin{lemma}[Length of $\ker\Psi$]\label{i2:lem-colker}
Let $M=V_O^d$ be as in Lemma~\ref{i2:lem-psiker}, with $\delta$ acting as $\delta_O\ne0$, $v_\delta=v(\delta_O)$.
\begin{enumerate}
\item[(a)] (Block $\rho_\pi$.) Suppose $E$ acts as $0$, $e_1M\cong V_O$, and the \emph{column lifting
condition} holds: for $l\in\iota$ and $s\in(e_1M)^{j_l}$, if the cochain $g\mapsto\sum_aK_{l,a}(g)s_a$ on
$D_l$ is a coboundary, then $l^{e'_l}s_a=0$ for all $a$. Then
$\ell(\ker\Psi)\le\sum_l\bigl(j_l(e_Oe'_l+v_\delta)+h^0_{D_l}(M)\bigr)+|\iota|\,d\,v_\delta$.
\item[(b)] (Blocks $c\ge2$.) If $e_1$ and $E$ act as $0$, then
$\ell(\ker\Psi)\le\sum_lh^0_{D_l}(M)+|\iota|\,d\,v_\delta$.
\item[(c)] (Trivial block.) If $u(G)$ acts trivially, $e_1$ acts as $0$, $E$ acts as $\delta_O$, and for each $l$
some $g\in D_l$ has $k^0_l(g)$ acting as $l^{e'_l}\delta_Oc$ with $c\in O^\times$, then
$\ell(\ker\Psi)\le\sum_ld\,(e_Oe'_l+2v_\delta)+|\iota|\,d\,v_\delta$.
\end{enumerate}
\end{lemma}

\begin{proof}
An element of $\mathrm{Hom}_A(N,M)$ is determined by its values on $Y$ and on the generators
$\nu_l,\sigma_{l,a},\sigma^0_l$; on $\ker\Psi$ by the latter alone. We use
$\ell(V_O^d[\beta])=d\,v(\beta)$ for $0\ne\beta\in O$, and $v(l^{e'}\beta)=e_Oe'+v(\beta)$.
(a) For $\varphi\in\ker\Psi$: $\varphi(\sigma_{l,a})\in e_1M$ by the second relation;
$\delta_O\varphi(\sigma^0_l)=E\varphi(\sigma^0_l)=0$ by the third; and since $k^0_l(g)\in AE$ acts as $0$,
the fourth gives $(u(g)-1)\varphi(\nu_l)=\sum_aK_{l,a}(g)\,\delta_O\varphi(\sigma_{l,a})$ on $D_l$. By the column
lifting condition applied to $s_a=\delta_O\varphi(\sigma_{l,a})$, each $\varphi(\sigma_{l,a})$ lies in
$(e_1M)[l^{e'_l}\delta_O]\cong V_O[l^{e'_l}\delta_O]$. The map
$\varphi\mapsto(\varphi(\sigma_{l,a}),\varphi(\sigma^0_l))_{l,a}$ thus lands in a module of length
$\sum_l(j_l(e_Oe'_l+v_\delta)+dv_\delta)$. On its kernel the displayed relation gives
$\varphi(\nu_l)\in M^{D_l}$, and $\varphi\mapsto(\varphi(\nu_l))_l$ is injective there. Adding the two
lengths gives (a). (b) Now $\varphi(\sigma_{l,a})=0$, $\varphi(\sigma^0_l)\in M[\delta_O]$, and the fourth
relation gives $(u(g)-1)\varphi(\nu_l)=0$ on $D_l$; so $\varphi\mapsto(\varphi(\nu_l),\varphi(\sigma^0_l))_l$
is injective into $\bigoplus_l(M^{D_l}\oplus M[\delta_O])$. (c) Now $\varphi(\sigma_{l,a})=0$,
$\delta_O\varphi(\nu_l)=E\varphi(\nu_l)=0$, $u(g)-1$ acts as $0$, and the fourth relation at the given
$g$ reads $0=\delta_O\,l^{e'_l}\delta_Oc\,\varphi(\sigma^0_l)$; so $\varphi(\sigma^0_l)\in M[l^{e'_l}\delta_O^2]$ and
$\varphi\mapsto(\varphi(\sigma^0_l),\varphi(\nu_l))_l$ is injective into a module of length
$\sum_l(d(e_Oe'_l+2v_\delta)+dv_\delta)$.
\end{proof}

\begin{theorem}[Bound for one block]\label{i2:thm-homabs}
Let $M=V_O^d$ be a block as in Lemma~\ref{i2:lem-colker}, $H=\mathrm{Hom}_A(N,M)$, and let
$\mathrm{rows}$, $\mathrm{qs}$ and $\iota$ be finite sets of places $w$, each with a subgroup
$D_w\subseteq G$ (for $w\in\iota$ the column subgroup), and $\mathrm{res}_w:H^1(G,M)\to H^1(D_w,M)$ the
restrictions. Let $Q\subseteq H^1(G,M)$ and $F_w\subseteq B_w\subseteq H^1(D_w,M)$ be $O$-submodules,
$F_w=B_w$ for $w\in\iota$, such that the class of every $\Psi(\varphi)$ lies in $Q$ and its restriction to
$D_w$ lies in $B_w$ ($w\in\mathrm{rows}\cup\mathrm{qs}\cup\iota$). Suppose
\[
\ell\Bigl(Q\cap\bigcap_{w}\mathrm{res}_w^{-1}F_w\Bigr)\le s,\qquad \ell(B_w/F_w)\le\lambda_w .
\]
Then $\ell(H)\le\ell(\ker\Psi)+\ell(\Psi(H)\cap B^1)+s+\sum_w\lambda_w$. In particular:
\begin{enumerate}
\item[(a)] for the block $\rho_\pi$ (rank $d=N_0$, hypotheses of Lemma~\ref{i2:lem-colker}(a) and of
Lemma~\ref{i2:lem-psiker}(2)),
$\ell(H)\le\sum_w\lambda_w+\sum_{l\in\iota}\bigl(j_l(e_Oe'_l+v_\delta)+h^0_{D_l}\bigr)+s+4N_0(|\iota|+1)v_\delta$;
\item[(b)] for a block $c\ge2$ ($d=1$, hypotheses of Lemma~\ref{i2:lem-colker}(b) and
Lemma~\ref{i2:lem-psiker}(2), $h^0_{D_l}\le H^0$),
$\ell(H)\le s+\sum_w\lambda_w+4(|\iota|+1)v_\delta+|\iota|H^0$;
\item[(c)] for the trivial block ($d=1$, hypotheses of Lemma~\ref{i2:lem-colker}(c)),
$\ell(H)\le s+\sum_w\lambda_w+\sum_{l\in\iota}(e_Oe'_l+2v_\delta)+4(|\iota|+1)v_\delta$.
\end{enumerate}
\end{theorem}

\begin{proof}
Let $\mathrm{cl}:Z^1\to H^1$ be the class map. Additivity of length along
$\ker\Psi\subseteq H$ and $\Psi(H)\cap\ker\mathrm{cl}\subseteq\Psi(H)$ gives
$\ell(H)=\ell(\ker\Psi)+\ell(\Psi(H)\cap B^1)+\ell(\mathrm{cl}\,\Psi(H))$, and
$\mathrm{cl}\,\Psi(H)\subseteq S:=Q\cap\bigcap_w\mathrm{res}_w^{-1}B_w$. The map
$S\to\bigoplus_wB_w/F_w$ has kernel $Q\cap\bigcap_w\mathrm{res}_w^{-1}F_w$, so $\ell(S)\le s+\sum_w\lambda_w$.
For the coboundary term: if $\Psi(\varphi)=dv$ then $v\in M[\delta_O^2]$ by Lemma~\ref{i2:lem-psiker}(2), so
$\Psi(H)\cap B^1\subseteq d(M[\delta_O^2])$ has length $\le2dv_\delta$; for the trivial block $B^1=0$.
Now (a): $\ell(\ker\Psi)\le\sum_l(\dots)+|\iota|N_0v_\delta$ by Lemma~\ref{i2:lem-colker}(a), and
$|\iota|N_0v_\delta+2N_0v_\delta\le4N_0(|\iota|+1)v_\delta$. (b) and (c) likewise, with
$|\iota|v_\delta+2v_\delta\le4(|\iota|+1)v_\delta$ and $|\iota|v_\delta\le4(|\iota|+1)v_\delta$.
\end{proof}

\begin{lemma}[The valuation $v_\delta$]\label{i2:lem-vdelta}
Let $l$ be prime, $O$ finite over $\Z_l$, $\varepsilon\in O^\times$ with $\varepsilon-1=l^{s_0}y$, $y\in\Z_l^\times$ (the formal statement allows $y\in O^\times$; the
case $y\in\Z_l^\times$ suffices for the cyclotomic character),
$s_0\ge1$, and $s_0\ge2$ if $l=2$. For $t\ge1$, $u_1\in O^\times$ and
$\delta_O=\prod_{0<k<t}(u_1-u_1\varepsilon^{-k})$:
$v(\delta_O)=e_O\sum_{0<k<t}(s_0+v_l(k))\le e_O(ts_0+v_l(t!))$.
\end{lemma}

\begin{proof}
$v(u_1-u_1\varepsilon^{-k})=v(\varepsilon^k-1)$, and $\varepsilon\in\Z_l^\times$; by the lifting-the-exponent lemma
(valid under the stated conditions on $s_0$), $v_l(\varepsilon^k-1)=s_0+v_l(k)$, and $v=e_Ov_l$ on $\Z_l$.
Sum over $k$; $\sum_{0<k<t}v_l(k)\le v_l(t!)$.
\end{proof}

\begin{proposition}[arithmetic of the column module; formally verified, \cite{CharacterFamily}]\label{i2:prop-colconds}
In the setting of Proposition~\ref{i2:prop-selmer} (local admissibility conditions at $(i,\pi)$, $l\ne g_0,g_1$,
$\pi\in Z_i$, Selmer finiteness of the $\rho_\pi$-block, Sha-finiteness), the ring data of
Proposition~\ref{i2:prop-latbound} carry a column module $N$ (Definition~\ref{i2:def-colmod}, with
$G=G_E$, $\chi$ the character $\chi_1$, $D_w$ the decomposition groups, and
$\delta_O$ as in Lemma~\ref{i2:lem-vdelta} with $\varepsilon$ the cyclotomic character and $g_0$ the chosen element)
such that:
\begin{enumerate}
\item $N$ is finitely generated and, for every point $x'$ over $X_i\setminus Z_i$, has a nonzero
$A$-linear map to the fibre at $x'$;
\item (block $\rho_\pi$) the hypotheses of Lemma~\ref{i2:lem-colker}(a) hold with the lift exponents $e'_w$,
and those of Theorem~\ref{i2:thm-homabs} hold with $\mathrm{qs}=\emptyset$, $\lambda_w=\ell_w$ (the row
lengths) and $s=e_Os_i(\pi)$;
\item (blocks $c\ge2$) the hypotheses of Lemma~\ref{i2:lem-colker}(b) and of Theorem~\ref{i2:thm-homabs} hold
with $\mathrm{qs}=Q_i$, $\lambda_w=e_Oc_B(\deg_w+2)$ at rows, $\lambda_v=v(\zeta_v-1)$ at $v\in Q_i$,
$h^0_{D_l}\le H^0$, and $s=e_O\ell_D(l,c-1)$;
\item (trivial block) the hypotheses of Lemma~\ref{i2:lem-colker}(c) and of Theorem~\ref{i2:thm-homabs} hold
with the same $\lambda_w$ and $s=e_Ov_l(h_E)$.
\end{enumerate}
\end{proposition}
\noindent\emph{What the proof involves} (the structure is in \S\ref{i2:sst-colconds}). Item (1) is a rank count in the fibres (the column
conditions bound the span of the column cocycles in each fibre, and the image of $E$ has rank one
with no invariant vector). The column lifting conditions in (2) and the element $k^0$ in (4) come from
the local lifting theorem at the column places with exponents $e'_w$. The row bounds in (2) come from
Bloch--Kato orthogonality (Theorem~\ref{i2:thm-bk}) and Tate duality (Theorem~\ref{i2:thm-tate}): the
local condition at a row is the annihilator of the image of $H^1_f$ of the dual, and its colength is
bounded by $\ell_w$. In (3), (4) the row bounds are the statement that $l^{c_B}$ times a cut class is
finite; at $v\in Q_i$ the inertia restriction is killed by $\zeta_v-1$. The global bounds are the
definition of $s_i(\pi)$; for $c\ge2$, classes unramified outside $l$ and locally trivial above $l$,
i.e.\ $\mathrm{Sha}^1(E,\Q_l/\Z_l(1-c))$ after change of coefficients, of normalised length
$\ell_D(l,c-1)$ by \eqref{i2:eq-sha}; for the trivial block, everywhere unramified homomorphisms
$G_E\to V_O$, which factor through the class group by global Artin reciprocity. The row conditions, the column lifting, the unramified classes and the
$\delta$-clause are separate formal proofs.
The proposition is stated informally; the exact statement is in \cite{CharacterFamily}. The column subgroups $D_l$,
the element $g_0$ and the column data $j_l,K_{l,a},k^0_l$ are defined exactly in \S\ref{i2:sst-colconds}. Not defined
here: the local admissibility conditions $\mathrm{Adm}_i(\pi;\ell,e',c_B)$ (named before Proposition~\ref{i2:prop-selmer};
nineteen of its clauses are used), including the block clause that gives $c_B$ its meaning. Global Artin reciprocity is
used in the form of Theorem~\ref{i2:thm-artin}.

\begin{proof}[Proof of Proposition~\ref{i2:prop-homterms}]
Item (1) of Proposition~\ref{i2:prop-colconds} is the hypothesis of Proposition~\ref{i2:prop-latbound}.
Apply Theorem~\ref{i2:thm-homabs}(a) to $M_0$, (b) to each $M_c$, $2\le c\le t$, and (c) to $M_1$, with the
data of Proposition~\ref{i2:prop-colconds}(2)--(4), and add. The global terms give
$e_O(s_i(\pi)+\sum_{e=1}^{t-1}\ell_D(l,e)+v_l(h_E))$. The row terms give $\sum_{\rm rows}\ell_w$ from $M_0$ and
$e_Oc_B\sum_{\rm rows}(\deg_w+2)$ from each of the $t$ blocks $M_1,\dots,M_t$, and the places of $Q_i$ give
$\sum_{v\in Q_i}v(\zeta_v-1)$ from each of these $t$ blocks. The column terms give
$j_w(e_Oe'_w+v_\delta)+h^0_w$ from $M_0$ and $e_Oe'_w+2v_\delta$ from $M_1$; and $(t-1)|\mathrm{col}|H^0\le t|\mathrm{col}|H^0$
from the blocks $c\ge2$. The coboundary terms give $4N_0(|\mathrm{col}|+1)v_\delta+t\cdot4(|\mathrm{col}|+1)v_\delta$.
Finally $v_\delta$ is computed by Lemma~\ref{i2:lem-vdelta}.
\end{proof}

\begin{proof}[Proof of Proposition~\ref{i2:prop-selmer}]
By Proposition~\ref{i2:prop-homterms}, $N$ satisfies the hypothesis of
Proposition~\ref{i2:prop-latbound}, which gives $e_Ov^{cc}_i(\pi)=\ell(O/\mathfrak c_\pi)\le\sum_c\ell\,\mathrm{Hom}(N,M_c\otimes K/O)$.
Proposition~\ref{i2:prop-homterms} bounds the right side; dividing by $e_O$ gives the statement,
the local terms being all terms other than $s_i(\pi)$, the Sha terms and $v_l(h_E)$.
\end{proof}

\subsection{The central action at members and the hub data}\label{i2:ssec-pf-hubgroup}

\subsubsection{The central-action tie at members}\label{i2:ssec-proof-cut}

We construct the central action of (T2) and prove (T4) at a member of a main-chain family. (T1) is
a clause of the coefficient rings of the member, and (T3) is the series localisation clause of the
datum together with the permutation clause recalled after Definition~\ref{i2:def-tie}. Fix $i$, put
$K_i=\mathrm{Frac}(O_i)$, $V=K_i\otimes_{O_i}M_i$, and for $o\in X_i\cup\{\emptyset\}$ let $W_i(o)$ be
the sum of the type spaces of the forms with eigensystem $o$ ($o=\emptyset$: eigensystem not in
$X_i$).

\emph{Member structure.} A member provides:
(M1) an $L$-structure on type vectors (the \emph{integral type vectors} $y$, with images
$j(y)\in W_i$, $j$ of an $L$-basis being a $\C$-basis of $W_i$), and an $L\to K_i$-semilinear
projection $y\mapsto\mathrm{pr}(y)\in V$ whose image spans $V$;
(M2) for each Hecke operator $T$ an $L$-rational operator $\tilde T$ on integral type vectors with
$\mathrm{pr}(\tilde Ty)=T\,\mathrm{pr}(y)$, whose $\C$-linear extension acts on each type space by a
scalar;
(M3) a form $P$ has eigensystem $x$ iff every $\tilde T$ acts on $W_i(P)$ by $\iota(T(x))$;
(M4) $\dim_\C W_i(x)=\dim_{K_i}E_x$;
(M5) $W_i$ is the sum of all type spaces;
(M6) an injective realization of $W_i$ in functions on the finite adelic group, carrying $W_i(P)$
into the isotypic component of the automorphic representation of $P$, through which central
$z\in\C[G]$ acts by $\sum_gz_gR(g)$, $R$ the right translation;
(M7) a finite set $\mathrm{Gen}$ of integral type vectors whose projections $\mathrm{pr}(y)$,
$y\in\mathrm{Gen}$, span the image $1\otimes M_i$ of $M_i$ in $V$ over $O_i$, such that, with
$\mathrm{Lat}$ the $O_L$-span of $\mathrm{Gen}$ ($O_L$ the ring of integers of $L$, mapped into $O_i$),
every central $c\in\Z[G]$ and $x\in\mathrm{Gen}$ admit $y\in\mathrm{Lat}$ with $j(y)=\psi(c)j(x)$.
(In the formalization the last property is derived from a lattice model: integral type vectors
are vectors of a space on which $G$ acts by an $L$-rational representation $\varrho$ preserving a
lattice $\Lambda$, $\mathrm{Lat}$ corresponds to the vectors $u$ with $Eu\in\Lambda$, $E$ the level
projector, and central $c$ acts by an operator $A_c$ with $EA_c=(\sum_tc_t\varrho(t))E$; then
$E(A_cu)=\sum_tc_t\varrho(t)(Eu)\in\Lambda$ because the $c_t$ are integers.)
Moreover the sums $W_i(o)$, $o\in X_i\cup\{\emptyset\}$, are independent, and the Hecke algebra $T_i\subseteq O_i^{X_i}$ separates the points of $X_i$.

\begin{lemma}[Reducedness]\label{i2:lem-reduced}
$V=\bigoplus_{x\in X_i}E_x$.
\end{lemma}

\begin{proof}
As in the first step of the proof of Lemma~\ref{i2:lem-length-le} there are $g_x\in T_i$ and
$0\ne d\in O_i$ with $\sum_xg_x=d$, $g_x(y)=0$ for $y\ne x$, $g_x(x)\ne0$ and $g_xV\subseteq E_x$.
So $dv=\sum_xg_xv\in\sum_xE_x$ for all $v$. If $\sum_xu_x=0$ with $u_x\in E_x$, applying $g_y$ gives
$g_y(y)u_y=0$, so $u_y=0$.
\end{proof}

\begin{lemma}[Comparison section]\label{i2:lem-section}
There is an injective additive map $\Phi:V\to W_i$, semilinear along $\iota$, with
$\Phi(E_x)\subseteq W_i(x)$ for all $x$ and $j(y)-\Phi(\mathrm{pr}\,y)\in W_i(\emptyset)$ for every
integral type vector $y$.
\end{lemma}

\begin{proof}
Choose a $K_i$-basis $(b_k)_{k\in\mathcal K}$ of $V$ adapted to Lemma~\ref{i2:lem-reduced},
$b_k\in E_{x_k}$, with coordinate functionals $c_k$; then $c_k(Tv)=T(x_k)c_k(v)$. By (M1) there is a
$\C$-linear $\Psi_k:W_i\to\C$ with $\Psi_k(j(y))=\iota c_k(\mathrm{pr}\,y)$ (define it on the basis $j(b)$;
both sides are semilinear along $L\to\C$). By (M2), $\Psi_k\circ\tilde T=\iota(T(x_k))\Psi_k$, so if
$\Psi_k$ is nonzero on $W_i(P)$ then $\tilde T$ acts on $W_i(P)$ by $\iota(T(x_k))$ for all $T$, and $P$
has eigensystem $x_k$ by (M3). Thus $p=(\Psi_k)_k:W_i\to\C^{\mathcal K}$ vanishes on $W_i(\emptyset)$
and its $k$-th coordinate vanishes on $W_i(y)$ for $y\ne x_k$. Its image contains
$(\iota c_k(v))_k$ for every $v\in V$ (true for $v=\mathrm{pr}\,y$, and these $v$ span), so $p$ is onto.
With $G_W=\sum_xW_i(x)$, (M5) gives $W_i=G_W+W_i(\emptyset)$, so $p|_{G_W}$ is onto, and
$\dim G_W\le\sum_x\dim E_x=|\mathcal K|$ by (M4) and Lemma~\ref{i2:lem-reduced}. Hence $p|_{G_W}$ is
bijective and $\ker p=W_i(\emptyset)$. Put $\Phi(u)=(p|_{G_W})^{-1}((\iota c_k(u))_k)$: additive,
$\iota$-semilinear, injective. For $u\in E_x$ write $\Phi(u)=a+c$, $a\in W_i(x)$,
$c\in\sum_{y\ne x}W_i(y)$; the coordinates of $p(c)$ with $x_k=x$ vanish by the support property, and
those with $x_k\ne x$ vanish because $c_k(u)=0$ and $p(a)_k=0$; so $p(c)=0$, $c=0$. Finally
$p(j(y))=p(\Phi(\mathrm{pr}\,y))$, so $j(y)-\Phi(\mathrm{pr}\,y)\in\ker p=W_i(\emptyset)$.
\end{proof}

\begin{lemma}[Two facts about the central action on type vectors]\label{i2:lem-ab}
(a) For central $z\in\C[G]$, $\psi(z)$ preserves every type space $W_i(P)$.
(b) For central $c\in\Z[G]$, $\psi(c)$ preserves $J(M_i)$, where $J(m)=\Phi(1\otimes m)$.
\end{lemma}

\begin{proof}
(a) By (M6), the realization of $\psi(z)F$ is $\sum_gz_gR(g)$ applied to that of $F$; right
translations preserve isotypic components, and the realization is injective.
(b) \emph{Intertwining.} Let $c\in\C[G]$ be central and $x,y$ integral type vectors with
$j(y)=\psi(c)j(x)$. Then $\Phi(\mathrm{pr}\,y)=\psi(c)\Phi(\mathrm{pr}\,x)$: write
$j(x)=\Phi(\mathrm{pr}\,x)+n_x$ and $j(y)=\Phi(\mathrm{pr}\,y)+n_y$ with $n_x,n_y\in W_i(\emptyset)$
(Lemma~\ref{i2:lem-section}). By (a), $\psi(c)$ preserves $W_i(\emptyset)$ and
$G_W=\sum_{x'}W_i(x')$, which contains the image of $\Phi$ (Lemmas~\ref{i2:lem-reduced},
\ref{i2:lem-section}). So $\Phi(\mathrm{pr}\,y)-\psi(c)\Phi(\mathrm{pr}\,x)=\psi(c)n_x-n_y$ lies in
$G_W\cap W_i(\emptyset)=0$.
\emph{Stability.} Let $c\in\Z[G]$ be central and $m\in M_i$. By (M7) write
$1\otimes m=\sum_ko_k\,\mathrm{pr}(x_k)$ with $o_k\in O_i$, $x_k\in\mathrm{Gen}$, and choose
$y_k\in\mathrm{Lat}$ with $j(y_k)=\psi(c)j(x_k)$. Since $\mathrm{pr}$ is semilinear and
$\mathrm{pr}(\mathrm{Gen})\subseteq1\otimes M_i$, $\mathrm{pr}(y_k)=1\otimes m_k$ for some $m_k\in M_i$.
Then, by semilinearity of $\Phi$ and the intertwining,
$\psi(c)J(m)=\sum_k\iota(o_k)\Phi(\mathrm{pr}\,y_k)=J\bigl(\sum_ko_km_k\bigr)$.
\end{proof}

\begin{proposition}[(T2) and (T4) at members]\label{i2:prop-T2T4}
$J:M_i\to W_i$ is injective, and $z_O(z)=J^{-1}\psi(z)J$ defines an $O_i$-algebra homomorphism
$Z(O_i[G])\to\mathrm{End}_{O_i}(M_i)$ satisfying (T2) and (T4).
\end{proposition}

\begin{proof}
$J$ is injective because $M_i\to V$ is (the lattice is torsion free) and $\Phi$ is. The centre
$Z(O_i[G])$ is the $O_i$-span of the class sums, which lie in $Z(\Z[G])$; by
Lemma~\ref{i2:lem-ab}(b) and semilinearity ($\psi(oz)=\iota(o)\psi(z)$, $J(om)=\iota(o)J(m)$),
$\psi(z)J(M_i)\subseteq J(M_i)$, so $z_O(z)$ is well defined; it is an $O_i$-algebra homomorphism
because $\psi$ is. By $\iota$-semilinearity of $\Phi$ and $\C$-linearity of $\psi(z)$,
\[
\Phi\circ(1\otimes z_O(z))=\psi(z)\circ\Phi\qquad\text{on }V.
\]
\emph{(T2).} For $u\in E_x$, $\Phi(z_O(z)u)=\psi(z)\Phi(u)\in W_i(x)$ by Lemma~\ref{i2:lem-ab}(a).
Writing $z_O(z)u=\sum_yu_y$ ($u_y\in E_y$, Lemma~\ref{i2:lem-reduced}), $\sum_y\Phi(u_y)\in W_i(x)$
with $\Phi(u_y)\in W_i(y)$; the type spaces of forms with distinct eigensystems are independent, so
$\Phi(u_y)=0$, hence $u_y=0$, for $y\ne x$. Thus $z_O(z)$ preserves each $E_x$, on which each Hecke
operator is a scalar; so $z_O(z)$ commutes with $T_i$ on $V$, hence on $M_i$.
\emph{(T4).} If $\psi(z)$ acts by $\iota(c)$ on every type vector over $x$, then for $u\in E_x$,
$\Phi(z_O(z)u)=\psi(z)\Phi(u)=\iota(c)\Phi(u)=\Phi(cu)$ since $\Phi(u)\in W_i(x)$; injectivity gives
$z_O(z)u=cu$.
\end{proof}

\noindent\emph{This proof is not the formal one} (the formal proof of (T4) goes through a different route, see below);
it is a proof from the member structure (M1)--(M7), which we have not matched field by field with the formal member
structure. The formalization defines $z_O$ by the same transport (the map $J$ and stability of
$J(M_i)$ are proved there exactly as in Lemma~\ref{i2:lem-ab}(b), from the generating set and the
integrality clause); its proof of (T4) uses in addition the rationality of central elements of
$L[G]$ on integral type vectors, which the argument above does not need. In particular no
identification of a lattice of type vectors with $J(M_i)$ is assumed: stability of $J(M_i)$ is a
consequence of (M7) and fact (a).


\subsubsection{Proof of Proposition~\ref{i2:prop-hub}}\label{i2:ssec-proof-hub}

A \emph{point family} for a datum $D$ is a choice, for every local index $i$ and point $x\in X_i$,
of a continuous representation $r_x:G_E\to\mathrm{GL}_{n_V}(O_i)$ attached to $x$ (its
characteristic polynomials of Frobenius at the unramified places are given by the Hecke
eigenvalues of $x$), and, for $\pi\in Z_i$, of a representation $\rho_\pi:G_E\to\mathrm{GL}_r(O_i)$
which is the $\rho$-block of $r_\pi$, i.e.\ $r_\pi\otimes K$ has semisimplification
$\rho_\pi\oplus\bigoplus_{c=1}^t\chi_c$, together with the good-row property of each $r$ at $i$. For
the twist datum $I$ at $w=w_i$ we use the notation of the definition of (SC), (GT) before
Lemma~\ref{i2:lem-killat}.

\begin{proposition}[point family, Selmer signatures, thresholds; formally verified, \cite{CharacterFamily}]\label{i2:prop-ptsel}
Assume $\mathrm{Hypotheses}_2$. Let $\mathcal F$ be a family with the family facts and $\mathrm{TwGood}$.
\begin{enumerate}
\item[(TH)] There are $x_a,x_b,x_c$ such that: for $x\ge x_a$ the parameter choice of $D_{x,i}$ is
admissible and its parameters satisfy (Par); for $x\ge x_b$ the special primes are small,
$\mathsf L\ge8$ and the split bound is $\ge60$; for $x\ge x_c$, $\mathrm{TwGood}$ provides twist
conjugacy and data $(N_0,S)$ for which the twist datum $(\mathrm{twist}_E,N_0,S)$ satisfies (GT). For every $x>x_1$ the
twist datum of $D_{x,i}$ is tied to its norm choice and its auxiliary places are the places of $Q_i$ (the norm-choice tie,
a field of the norm-choice facts).
\item[(PF)] (Chebotarev and cuspidality.) Every member $D=D_{x,i}$, $x>x_1$, satisfying the family facts has a point
family $(r,\rho)$ in the above sense except for the good-row property, which is supplied separately by the family facts;
the representations $r_x$ are themselves supplied by the family facts, and the content is the construction of $\rho$.
\item[(SEL)] (Selmer-signature selection.) If moreover the twist datum is tied to the norm choice, and the triple,
good-row, residual absolute irreducibility and same-twist conditions of the family facts hold, then
for every point family $(r,\rho)$, every twisted $i$ and every $\pi\in Z^{\rm good}_i$ there is a good
signature $\sigma$ with $s_i(\pi)\le\mathrm{BK}(\sigma)$ and ($\mathrm{BK}(\sigma)<\infty\Rightarrow$ the Selmer group
of the $\rho_\pi$-block is finite).
\item[(CR)] (Conductor roots.) Under the same tie, for $v\in Q_i$,
$\ell(O/(\zeta_v-1))/\ell(O/l)\le\mathrm{tv}(v)$.
\end{enumerate}
Moreover, for every member the Bhatt--Scholze comparison (Theorem~\ref{i2:thm-bs}) provides the Wach and
Fontaine functors of $D$ and a comparison between them, used at the large primes.
\end{proposition}
\noindent\emph{What the proof involves} (structure in \S\ref{i2:sst-ptsel}). (TH) collects size statements along the
family. (PF) extends a trace identity from Frobenius elements to $G_E$ by Chebotarev's theorem (in a degree-one form proved
in the formalization) and continuity, and splits off $\rho$ by linear algebra. (SEL) uses a theorem of Carayol type (equal
traces and absolutely irreducible residual representation give conjugate lattices) and the invariance of Bloch--Kato
Selmer lengths under conjugation and flat change of coefficients. (CR) is an elementary tame-inertia argument.
The proposition is stated informally; the exact statement is in \cite{CharacterFamily}. The admissibility of the
parameter choice and the tame values are defined in \S\ref{i2:sst-ptsel}. Not defined here: the family facts,
$\mathrm{TwGood}$ and twist conjugacy, the split bound, the twist datum $(\mathrm{twist}_E,N_0,S)$, its Bloch--Kato lengths
$\mathrm{BK}(\sigma)$, good signatures, and the conditions of the family facts named in (SEL); \S\ref{i2:sst-ptsel} explains
why they do not fit in a few pages.

\begin{proof}[Proof of Proposition~\ref{i2:prop-hub}]
Let $x_0=\max(x_a,x_b,x_c)$ (Proposition~\ref{i2:prop-ptsel}(TH)). Fix $x\ge x_0$, $x>x_1$, $i$ with
$D=D_{x,i}$ Sha-finite, and $\mathrm{exc}$. The family facts give the norm-choice facts of $D$; let
$(N_0,S)$ be the data of (TH) for $D$, and $(r,\rho)$ a point family (PF).

\emph{The canonical hub data.} Put $h=\lfloor2b\rfloor+1$ and define $X$ by: the point family
$(r,\rho)$; the lengths $\mathrm{len}^0_i(\pi,\lambda)$ and row lengths given by the explicit row
budgets of the size datum (at inert $l$ the budgets $\mathrm{rowBr}$ and $\mathrm{flatB}$-type
expressions of Definition~\ref{i2:def-budgets}, at the other primes the split-type budget with the
exponent $e_P$); the error terms $\mathrm{err}_i(\pi)$ given by the explicit local error expressions
(\S\ref{i2:ssec-proof-other}); the cyclotomic block numbers $c_B(i)=\mathrm{cB}(l_i,\lfloor2b\rfloor,t,4nn_V)$;
the column constants
$c_0(i)=5+\lfloor\log_l(2h)\rfloor+4nn_V\bigl(4+2\lfloor\log_l(2h)\rfloor\bigr)$; column lift exponents $e_S(l)$ as in
Proposition~\ref{i2:prop-locadm}(1);
$\eta=\lfloor\eta_{\rm sz}\rfloor$; the given $\mathrm{exc}$; the band predicate determined by the
budgets; and the twist datum $(\mathrm{twist}_E,N_0,S)$. The canonicity conditions (with row constants
$(C_2,320,25,23,32,32)$, which enter only through the row budgets) require exactly
that each of these entries has the stated value (with the row decomposition at the split primes
chosen as in the hypothesis of the declaration of embedding types, and all places as rows at
inert primes), so they hold by construction. The parameter bounds hold by (TH), and
$\eta=\lfloor\eta_{\rm sz}\rfloor\le\eta_{\rm sz}\le2\eta_{\rm sz}$ when $\eta_{\rm sz}\ge0$.

\emph{(GT)} holds for the twist datum $(\mathrm{twist}_E,N_0,S)$ by (TH).

\emph{(SC).} By (TH) the twist datum is tied to the norm choice, so (SEL) gives, for twisted $i$ and
$\pi\in Z^{\rm good}_i$, the good signature required by the first half of (SC). For the second
half, the auxiliary places of the twist at $w_i$ are the places of $Q_i$, and (CR) gives
$\sum_{v\in Q_i}\ell(O/(\zeta_v-1))/\ell(O/l)\le\sum_{v\in Q_i}\mathrm{tv}(v)=\sum_P\mathrm{tv}(P)$.

$\mathsf L\ge8$ and the split bound $\ge60$ hold by (TH); the Wach and Fontaine functors and their comparison are
those of Proposition~\ref{i2:prop-ptsel}; the admissibility of the jump parameters and the norm-choice facts are part of
(TH) and of the family facts. Finally $\mathrm{KillAt}$ follows from (SC) and (GT) by Lemma~\ref{i2:lem-killat}. Since the
canonicity predicate is described, not defined, in this paper, the step ``(canon) holds by construction'' is taken from
the formal verification \cite{CharacterFamily}.
\end{proof}

\subsubsection{Proof of Proposition~\ref{i2:prop-other}}\label{i2:ssec-proof-other}

Items (1)--(3) are Proposition~\ref{i2:prop-selmer} (with Proposition~\ref{i2:prop-homterms})
applied to a decomposition of the places of $E$ above $l$ into rows and columns, followed by
elementary comparisons. Item (4) is the formally verified Proposition~\ref{i2:prop-large} below.
Throughout, $i$ is twisted with $l=l_i\notin\{g_0,g_1\}$, $\pi\in Z^{\rm good}_i$, $D$ is Sha-finite,
and the hub data are those of Proposition~\ref{i2:prop-hub}, so $\mathrm{KillAt}$ holds; write
$\bar v_\delta=v_\delta/e_O$ and $\deg_w$ for the local degree of a place $w\mid l$. For an index
set $W$ of places above $l$ we use
\begin{equation}\label{i2:eq-degsum}
\#W\le\sum_{w\in W}\deg_w\le2f ,
\end{equation}
the second inequality being the fundamental identity for $E$ (degree $2f$ over $\Q$ at the
relevant level, Lemma~\ref{i2:lem-meanweight}). In Proposition~\ref{i2:prop-selmer} the local
terms are, after division by $e_O$ (Proposition~\ref{i2:prop-homterms}),
\begin{equation}\label{i2:eq-localterms}
\begin{split}&\frac1{e_O}\Bigl(\sum_{\rm rows}\ell_w+\sum_{\rm cols}\bigl(e_Oe'_w+j_w(e_Oe'_w+v_\delta)+2v_\delta+h^0_w\bigr)\Bigr)\\
&\quad+tc_B\sum_{\rm rows}(\deg_w+2)+4n_V(|\mathrm{col}|+1)\bar v_\delta+\frac{t|\mathrm{col}|H^0}{e_O}+\varepsilon_Q,\end{split}
\end{equation}
with $N_0+t=r+t=n_V$ and $\varepsilon_Q=t\sum_{v\in Q_i}\ell(O/(\zeta_v-1))/\ell(O/l)\le l^{-2}$ by $\mathrm{KillAt}$.

\begin{proposition}[local admissibility data; formally verified, \cite{CharacterFamily}]\label{i2:prop-locadm}
For $x$ beyond a threshold, with the hub data of Proposition~\ref{i2:prop-hub}, the local
admissibility conditions of Proposition~\ref{i2:prop-selmer} hold at $(i,\pi)$ for the following
decompositions, with the following bounds.
\begin{enumerate}
\item ($\mathrm{RangeAD}(l)$.) The decomposition of the hub data (all places rows if $l$ is inert;
the declared rows and columns if $l$ is split), with column lift exponents $e_S(l)=\lceil23(2+2\lfloor h_B\rfloor/l+\log_l(2\lfloor h_B\rfloor+1))\rceil$
at the columns, whose contribution is included in $\mathrm{len}^0_i(\pi,\cdot)$. The row and column terms of
\eqref{i2:eq-localterms} are at most $\sum_{\lambda\mid l}\mathrm{len}^0_i(\pi,\lambda)+\Sigma^{\rm col}$, where
$\Sigma^{\rm col}=\sum_{u\in{\rm col}}\bigl(\deg_um\bar v_\delta+2\bar v_\delta+h^0_u+r\,c_0(u)\bigr)$ (and $0$ at inert $l$),
$h^0_u$ the length of the $\rho_\pi$-invariants of the decomposition group at $u$, and $H^0/e_O\le h_0$. The hub
constants satisfy, with $c=120$: $c_B\le c\,tnn_V\mathsf L$, $h_0\le ct$, and at split $l$, for $u$ a column,
$c_0(u)\le c\,nn_V\mathsf L$ and $h^0_u\le c\,rnn_V\mathsf L$. Moreover the table with
$\mathrm{len}=\sup_\pi\mathrm{len}^0_i(\pi,\cdot)$ and the band predicate of the hub data satisfies (A), (B), (Med), (SmI)
and (Band) at $l$, for every constant $C\ge\max(3C_2+1,22365)$ (from the canonicity of the hub data).
\item ($l\in\{q,d\}$.) The decomposition with no columns, with row lengths
$\ell_w/e_O\le K\bigl((\deg_w+1)rX'+3\deg_wmY'\bigr)+1$, $K=160$, $X'=4b/l+2S_c+2\log_l(2b)+1+\eta_{\rm sz}+t$,
$Y'=2b+1$, and block constant $c_B\le60\,tf(r+t)\log b$.
\item ($l\in\mathrm{small}\cup\{3\}$.) A split decomposition, with column lift exponents $e'$, column constant $c_0$
and block constant $c_B$ satisfying, with $c=1000$, $h_B\ge b$ the hub size, $s_l$ the local level and
$\Lambda_l$ the abelian size at $l$: $e'\le c\,h_Bs_l$, $c_0\le cfn_V\Lambda_l$, $h^0_u\le c\,rfn_V\Lambda_l$,
$c_B\le c\,tfn_V\Lambda_l$, $h_0\le c\Lambda_l$; for $l=2$, $s_0\ge2$.
\end{enumerate}
\end{proposition}
\noindent\emph{What the proof involves} (structure in \S\ref{i2:sst-locadm}). The row lengths are the colengths of the
local Bloch--Kato conditions, bounded through the local lifting theorem (constant $K=160$) and the local Euler
characteristic; the block constants come from the cyclotomic exponents at the level of the twist; the invariant bounds are
estimates for $H^0$ of the local representations. At inert $l$ and at $q,d$ the admissibility predicate is the twenty-clause
one; at split $l$ it is a different formal predicate with 38 clauses. The proposition is stated informally; the exact
statement is in \cite{CharacterFamily}. Not defined here are the local admissibility conditions, the canonical hub data and their row
constants, the hub size $h_B$, the local level $s_l$ and the abelian size $\Lambda_l$ (explicit quantities of the size
datum at $l$), the declared rows and columns at split $l$, and the constants $c_0(u)$, $h^0_u$, $h_0$ beyond the bounds
stated. Of these, all but the two admissibility predicates, $s_l$ and $\Lambda_l$ can be defined exactly in under two pages
each; see \S\ref{i2:sst-locadm}.

\begin{proof}[Proof of Proposition~\ref{i2:prop-other}, item (1)]
By Proposition~\ref{i2:prop-locadm}(1), Proposition~\ref{i2:prop-selmer} applies, and by
\eqref{i2:eq-localterms} and $\varepsilon_Q\le l^{-2}$ it gives $\mathrm{Bd}_i(\pi,\sum_\lambda\mathrm{len}^0_i(\pi,\lambda)+\mathrm{err}_i(\pi))$ with
\[
\mathrm{err}_i(\pi)=\Sigma^{\rm col}+t\,c_B\sum_{\rm rows}(\deg_w+2)+4n_V(|\mathrm{col}|+1)\bar v_\delta+t|\mathrm{col}|h_0 .
\]
\emph{The bound (bb).} Since $l\notin\mathrm{small}$, $s_0=1$, so by Lemma~\ref{i2:lem-vdelta}
$\bar v_\delta\le t+v_l(t!)\le2t$. Also $f,t,n,n_V\ge1$, $m\le t$, $r\le t$, and $\mathsf L\ge1$. By
\eqref{i2:eq-degsum}, $\sum_{\rm rows}(\deg_w+2)\le2f+2\cdot2f=6f$ and $|\mathrm{col}|\le2f$. With $c=120$:
\[
\begin{aligned}
t\,c_B\textstyle\sum_{\rm rows}(\deg_w+2)&\le6\,cfnt^2n_V\mathsf L,\\
4n_V(|\mathrm{col}|+1)\bar v_\delta&\le4n_V\cdot3f\cdot2t\le24\,cfnt^2n_V\mathsf L,\\
t|\mathrm{col}|h_0&\le t\cdot2f\cdot ct\le2\,cfnt^2n_V\mathsf L,\\
\Sigma^{\rm col}&\le2f\cdot t\cdot2t+2f\bigl(4t+2\,ctnn_V\mathsf L\bigr)\le16\,cfnt^2n_V\mathsf L,
\end{aligned}
\]
using $\sum_u\deg_u\le2f$, $|\mathrm{col}|\le2f$, $r\,c_0(u)+h^0_u\le2\,crnn_V\mathsf L\le2\,ctnn_V\mathsf L$.
Hence $\mathrm{err}_i(\pi)\le48\,cfnt^2n_V\mathsf L\le60\cdot120\,fnt^2n_V\mathsf L^2=7200\,fnt^2n_V\mathsf L^2$.
At inert $l$ the column terms are absent and the same estimate holds. The lengths satisfy (A), (B),
(Med), (SmI) and (Band) by the last clause of Proposition~\ref{i2:prop-locadm}(1). Concerning $\Sigma^{\rm col}$: with $j_u=\deg_um$ column generators, the column term of \eqref{i2:eq-localterms} at $u$ is
$\deg_um\bar v_\delta+2\bar v_\delta+h^0_u/e_O$ plus the lift-exponent terms $(1+\deg_um)e_S(l)$, which
Proposition~\ref{i2:prop-locadm}(1) includes in $\mathrm{len}^0_i(\pi,\cdot)$; the bound carries in addition the
nonnegative term $r\,c_0(u)$, which only weakens it.
\end{proof}

\begin{proof}[Proof of Proposition~\ref{i2:prop-other}, item (2)]
Let $l\in\{q,d\}$; then $l\ne g_0,g_1$, $\nu_l=0$, and $l>4tr+1\ge9$. By
Proposition~\ref{i2:prop-locadm}(2) and Proposition~\ref{i2:prop-selmer} with no columns,
\[
v^{cc}_i(\pi)\le\sum_w\ell_w/e_O+s_i(\pi)+\mathrm{VV}_D(l)+v_l(h_E)+tc_B\sum_w(\deg_w+2)+4n_V\bar v_\delta+\varepsilon_Q .
\]
Since $l\in\{q,d\}$ has ramification index $2$ in $E$ (host axioms), $\deg_w\ge2$ for every $w\mid l$, so
$\#\{w\mid l\}\le f$; with \eqref{i2:eq-degsum}, $\sum_w(\deg_w+1)\le3f$, $\sum_w(\deg_w+2)\le4f$ and
$\sum_w\deg_w\le2f$, so
$\sum_w\ell_w/e_O\le K(3frX'+6fmY')+f$ and $tc_B\sum_w(\deg_w+2)\le240\,t^2f^2(r+t)\log b$; as $l>t$,
$v_l(t!)=0$ and $\bar v_\delta\le t$ (Lemma~\ref{i2:lem-vdelta} with $s_0=1$), so $4n_V\bar v_\delta\le4(r+t)t$; and
$\varepsilon_Q\le l^{-2}$. Now compare with
$Cf\bigl(r(b/l+S_c+t+\eta_{\rm sz}+\log b)+mb\bigr)+Q(f,r,t)\log b$, where $C\ge13K$ and $Q$ is a
polynomial with natural coefficients and $Q(f,r,t)\ge260f^2t^2(r+t)$. Since $l\ge5$, $\log b>1$ and
$\log_l(2b)\le(\log2+\log b)/\log5\le\tfrac32\log b$, we get $3X'\le12b/l+6S_c+3(\eta_{\rm sz}+t)+12\log b$, so
$3KfrX'\le13Kfr(b/l+S_c+t+\eta_{\rm sz}+\log b)$; since $b\ge6$ ($\mathsf L\ge8$), $6KfmY'=6Kfm(2b+1)\le13Kfmb$; and
$f+240t^2f^2(r+t)\log b+4(r+t)t\le260f^2t^2(r+t)\log b$ as $f,t\ge1$, $\log b>1$. Hence the local terms are at most
$Cf\bigl(r(b/l+S_c+t+\eta_{\rm sz}+\log b)+mb\bigr)+Q(f,r,t)\log b$. Finally, a polynomial with natural coefficients
satisfies $Q(f,r,t)\le(\sum\text{coefficients})(frt)^{\deg Q}\le C(frt)^C$ for $f,r,t\ge1$ once $C$ is at least the sum of
the coefficients and the degree of $Q$; so the local terms are at most $\mathrm{midG}(l)$, which is
$\mathrm{Bd}_i(\pi,\mathrm{midG}(l))$. Lastly $\mathrm{midG}(l)\le2Cfmb$ for large $x$: $r/l<1\le m$, and the remaining terms are bounded
by the size relations of the family (Remark~\ref{i2:rem-par}) by a quantity of size
$x^{O(1)}e^{x^{\varepsilon/4}}\mathsf L$, while $b=e^{\mathsf L}=e^{x^\varepsilon}$; this is the clause (small) of the family.
\end{proof}

\begin{proof}[Proof of Proposition~\ref{i2:prop-other}, item (3)]
Let $l=2$, or $l=3\notin\mathrm{small}$, or $l\in\mathrm{small}\setminus\{2\}$; in each case $l\ne g_0,g_1$, and for
$l=2$ the condition $s_0\ge2$ of Lemma~\ref{i2:lem-vdelta} holds. By Proposition~\ref{i2:prop-locadm}(3) and
Proposition~\ref{i2:prop-selmer} for the split decomposition, $v^{cc}_i(\pi)\le s_i(\pi)+\mathrm{VV}_D(l)+v_l(h_E)$ plus
\[\begin{gathered}
\sum_{u\in{\rm col}}\bigl(e'+\deg_um(e'+\bar v_\delta)+2\bar v_\delta+h^0_u+\deg_ume'+rc_0\bigr)\\+t\sum_{\rm rows}(\deg_w+2)c_B
+4n_V(|\mathrm{col}|+1)\bar v_\delta+t|\mathrm{col}|h_0+\varepsilon_Q .
\end{gathered}\]
With the bounds of Proposition~\ref{i2:prop-locadm}(3), \eqref{i2:eq-degsum}, $s_0+1\le\Lambda_l$ (so
$\bar v_\delta\le t\Lambda_l+v_l(t!)$) and $\varepsilon_Q\le l^{-2}$, collecting the monomials gives
$\mathrm{Bd}_i(\pi,\mathrm{cost}_l)$ with $\mathrm{cost}_l=6000\,fmh_Bs_l+30000\,f^2t^2n_V\Lambda_l$ (the constants
of the formalization). The column term above contains, besides the column length
$e'+\deg_um(e'+\bar v_\delta)+2\bar v_\delta+h^0_u$ of \eqref{i2:eq-localterms} (with $j_u=\deg_um$), the nonnegative terms
$\deg_ume'$ and $rc_0$; these come from the formal split-type bound and only weaken the inequality. Along the family, with the admissible-parameter constant $C_0$, the local level
and abelian size are bounded by the parameters ($N$ and $Q_S$ have the same $l$-adic valuation at these
$l$), and with $\log x\ge1$ one gets $\mathrm{cost}_l\le Cfmb\log x+Cx^C\mathsf L^2$ for
$C=18AC_0+2B(3C_0+9)+5C_0+1$, $(A,B)=(6000,30000)$. This is item (3) with the maximum of these constants over the
three cases.
\end{proof}

\begin{proposition}[large primes; formally verified, \cite{CharacterFamily}]\label{i2:prop-large}
Let $D$ be a datum with size datum $L$, $(r,\rho)$ a point family for $D$, and $\mathrm{exc}$ a finite set of primes. Assume:
\begin{enumerate}
\item[(CB)] $\mathrm{CongBound}(D,L)$: for every local index $i$ with $l_i>l^*$, $l_i\notin\{g_0,g_1\}\cup\mathrm{Exc}(L)$,
  satisfying the weight tie and the Hecke--Galois condition at $i$, every point family $r'$ with the good-row property,
  every $\pi\in Z_i$ and every rank-$r$ part $\rho'$ of $r'_\pi$,
  $v^{cc}_i(\pi)\le s_i(\pi)+\sum_{0<e<t}\ell(\mathrm{Sha}(\mathrm{ld}_D(l_i),e))+v_{l_i}(h_E)+t\sum_{v\in Q_i}v(\zeta_v-1)$
  (in $[0,\infty]$, $s_i(\pi)$ computed with $\rho'$);
\item the weight parametrisation clause of $D$; $\mathrm{KillAt}$ for $\rho$; $D$ Sha-finite;
\item $3\le l^*$, and $\mathrm{exc}=\mathrm{Exc}(L)$ (Definition~\ref{i2:def-budgets}).
\end{enumerate}
Then for every twisted $i$ with $l=l_i>l^*$, $l\notin\{g_0,g_1\}\cup\mathrm{exc}$, and every $\pi\in Z^{\rm good}_i$:
$\mathrm{Bd}_i(\pi,0)$.
\end{proposition}
\noindent The weight tie and the Hecke--Galois condition in (CB), and the weight parametrisation clause, are named, not
defined, here; to that extent the statement is stated informally, and the exact statement is in \cite{CharacterFamily}. The hypothesis (CB) is itself a formally
verified consequence of $\mathrm{Hypotheses}_2$ and the norm-choice facts along a family, for $x$ beyond a threshold; its
proof is the large-prime argument sketched below. The hypothesis $\mathrm{exc}=\mathrm{Exc}(L)$ cannot be dropped: for a
smaller excluded set the statement would claim $\mathrm{Bd}_i(\pi,0)$ at primes where it is not proved.
\noindent\emph{What the proof involves} (structure in \S\ref{i2:sst-large}). Given (CB), the proposition is the
identification of the right side of (CB) with that of $\mathrm{Bd}_i(\pi,0)$, by Sha-finiteness, \eqref{i2:eq-sha} and
$\mathrm{KillAt}$, after deriving the weight tie from the weight parametrisation clause. The content is in (CB): at large
unramified $l$ outside $\mathrm{Exc}(L)$ there are no local terms at all: the cut cocycles of the $\rho_\pi$-block are
finite, those of the blocks $\varepsilon^{1-c}$ are locally trivial, and those of the trivial block are unramified. This is
of Fontaine--Laffaille type; the formalization proves it through a torsion comparison for Wach lattices
(Theorems~\ref{i2:thm-wach}, \ref{i2:thm-herr} with Corollary~\ref{i2:cor-herr}, \ref{i2:thm-bs}); we do not give a written proof of this
comparison, and we do not replace it by a published statement.

\subsection{The local lifting theorem: proof of Proposition~\ref{i2:prop-lift}}\label{i2:ssec-proof-lift}

The proof has three layers: a reduction to lattices that are crystalline over one tame field
(Proposition~\ref{i2:prop-goodlat}, proved); the lifting theorem for such lattices
(Proposition~\ref{i2:prop-liftgood}), proved from a sequence of Breuil--Kisin steps; and the steps
themselves. Of these, the lifting of maps (Proposition~\ref{i2:bk-c}) and the $\zeta$-bound
(Proposition~\ref{i2:bk-zeta}, the first half of the final step) are proved here; the others are
stated exactly and cited as formally verified. Throughout, $l$ is prime, $G=G_{\Q_l}$,
$I_l\subseteq G$ is inertia, and all lattices are $G$-stable $O_L$-lattices.

\subsubsection*{Good lattices}

\begin{definition}[Tame data and good lattices]\label{i2:def-good}
Let $a\ge1$, $e=l^a-1$, and fix $\varpi=(\varpi_n)_{n\ge0}$ in $\overline\Q_l$ with $\varpi_0^e=-l$ and
$\varpi_{n+1}^l=\varpi_n$; $K_\varpi=\Q_{l^a}(\varpi_0)$ is tamely ramified of index $e$. Let $L\supseteq\Q_{l^a}$, so $L$ contains
$\mu_e$. For $c\in\Z/e$ let $\mathrm{per}(c)$ be the least $m\ge1$ with $l^mc=c$, and for
$\Sigma'\subseteq\Z/e$ put $r(\Sigma')=\max_{c,c'\in\Sigma'}\mathrm{per}(c-c')$. For $U\subset\Z$ finite, a lattice $T$
in $V$ is \emph{good for $(\Sigma',U)$} if $V$ is crystalline over $K_\varpi$; for $g\in I_l$ and
$\zeta=g(\varpi_0)/\varpi_0\in\mu_e$, $\prod_{c\in\Sigma'}(g-\zeta^c)$ kills $D_{\rm cris}(V|_{K_\varpi})$; and the Hodge--Tate
jumps of $V$ lie in $U$.
\end{definition}

\begin{proposition}[Good lattices over a tame field]\label{i2:prop-goodlat}
Let $a=(l-1)!$, $e=l^a-1$, $L\supseteq\Q_{l^a}$, $K'/\Q_l$ finite Galois, $\mathrm{Fr}$ a Frobenius and $\Sigma$ a
$\mathrm{Fr}$-stable set of characters of $I(K'/\Q_l)$ with values in $L^\times$, $1\in\Sigma$, $|\Sigma|\le N<l$, each trivial on
$I_l\cap G_{K'}$ and of order prime to $l$. There is $\Sigma'\subseteq\Z/e$ with $l\Sigma'\subseteq\Sigma'$, $0\in\Sigma'$,
$|\Sigma'|\le N$ and $r(\Sigma')=r'(\Sigma)$, such that every lattice in a representation crystalline over $K'$, on
whose $D_{\rm pcris}$ inertia acts through $\Sigma$, is crystalline over $K_\varpi$ and satisfies the tame-type condition
of Definition~\ref{i2:def-good} for $\Sigma'$.
\end{proposition}

\begin{proof}
Let $\omega(g)=g(\varpi_0)/\varpi_0$ on $I_l$. We use four facts.
(T1) $\omega(I_l)=\mu_e(L)$, and every character $\psi$ of $I_l$ with open kernel and order prime to $l$ satisfies
$\psi(\Phi g\Phi^{-1})=\psi(g)^l$ for a lift $\Phi$ of Frobenius, and $\psi=\omega^c$ for a unique $c\in\Z/e$ if
$\psi^e=1$. (This is the structure of the tame quotient: $\psi$ factors through
$\mathrm{Gal}(\Q_l^{\rm ur}(\rho)/\Q_l^{\rm ur})\cong\mu_M$ for a root $\rho^M=-l$, $l\nmid M$, by the facts that
inertia elements fixing such a $\rho$ act on a finite extension through an $l$-group, and that $X^M+l$ is
irreducible over $\Q_l^{\rm ur}$; Frobenius acts on $\mu_M$ by the $l$-th power.)
(T2) There is a lift $\Phi$ of the arithmetic Frobenius with $\Phi^{-1}\mathrm{Fr}\in I=I(K'/\Q_l)$, and every
$g\in I$ is $su$ with $s\in I_l$, $u\in G_{K'}$ (Galois theory for $K'\cap\Q_l^{\rm ur}$).
(T3) $D_{\rm pcris}(V)=D_{\rm cris}(V|_{K'})$, by definition.
(T4) (Galois descent of crystallinity.) Let $K'/\Q_l$ be finite Galois with $V|_{K'}$ crystalline, and $F/\Q_l$ finite. If every
$g\in I_l\cap G_F$ fixes $D_{\rm cris}(V|_{K'})$, then $V|_F$ is crystalline; and a polynomial in $g\in I_l$ killing
$D_{\rm cris}(V|_{K'})$ kills $D_{\rm cris}(V|_F)$. The first assertion is Fontaine's criterion: a potentially crystalline
$V$ is crystalline over $F$ if and only if $I_F=I_l\cap G_F$ acts trivially on $D_{\rm pcris}(V)=D_{\rm cris}(V|_{K'})$
(Fontaine, Ast\'erisque 223, Exp.~VIII, with Exp.~III \S1 for the regularity of $B_{\rm cris}$; and Fontaine--Ouyang (reference to be confirmed)). For the second, let $K''=K'F$; as $V|_{K'}$ is crystalline,
$D_{\rm cris}(V|_{K''})=K''_0\otimes_{K'_0}D_{\rm cris}(V|_{K'})$ (same references), and
$D_{\rm cris}(V|_F)=(B_{\rm cris}\otimes V)^{G_F}\subseteq(B_{\rm cris}\otimes V)^{G_{K''}}=D_{\rm cris}(V|_{K''})$. An element $g\in I_l$
fixes $K''_0\subseteq\Q_l^{\rm ur}$, so it acts $K''_0$-linearly on $D_{\rm cris}(V|_{K''})$, and a polynomial in $g$ killing
$D_{\rm cris}(V|_{K'})$ kills its $K''_0$-span, hence $D_{\rm cris}(V|_F)$. (The formal proof, \S\ref{i2:sst-bk} end, follows the
same route through Speiser's lemma.)

\emph{Exponents.} For $\chi\in\Sigma$, $\psi=\chi|_{I_l}$ has open kernel and order prime to $l$; by (T1), (T2) and
$\mathrm{Fr}$-stability, for each $\chi$ there is $b\in\Sigma$ with $\chi=b^l$ on $I_l$. Iterating inside the finite set
$\Sigma$ gives $\chi^{l^k}=\chi$ with $1\le k\le|\Sigma|<l$, so $k\mid a$ and $\chi^e=1$ on $I_l$. Hence $\chi=\omega^{c(\chi)}$
on $I_l$; put $\Sigma'=c(\Sigma)$. Then $0=c(1)\in\Sigma'$, $|\Sigma'|\le N$, and $l\Sigma'\subseteq\Sigma'$.
\emph{Periods.} For $\chi_1,\chi_2\in\Sigma$ the $\mathrm{Fr}$-orbit length of $\chi_1\chi_2^{-1}$ equals
$\mathrm{per}(c(\chi_1)-c(\chi_2))$: on $I_l$ conjugation by $\mathrm{Fr}$ raises to the $l$-th power, on $G_{K'}$ both
characters are trivial, and (T2) covers $I$. So $r(\Sigma')=r'(\Sigma)$.
\emph{Lattices.} Inertia acts on $D_{\rm cris}(V|_{K'})$ through $\Sigma$. For $g\in I_l\cap G_{K_\varpi}$, $\omega(g)=1$, so
$\chi(g)=1$ for all $\chi\in\Sigma$ and $g$ fixes $D_{\rm cris}(V|_{K'})$; by (T4) with $F=K_\varpi$, $V|_{K_\varpi}$ is crystalline. For
$g\in I_l$, every $\chi(g)=\omega(g)^{c(\chi)}$ is a root of $\prod_{c\in\Sigma'}(X-\omega(g)^c)$, which therefore kills
$D_{\rm cris}(V|_{K'})$, hence $D_{\rm cris}(V|_{K_\varpi})$ by (T4).
\end{proof}

\subsubsection*{Breuil--Kisin modules}

Fix $l\ne2$, $a$, $e$, $\varpi$, $L\supseteq\Q_{l^a}$ as above, $O=O_L$.
\emph{Rings.} $\mathfrak S_L=O[[u]]$; $E(u)=u^e+l$; $\varphi$ is the $L$-linear map $u\mapsto u^l$ on $L[[u]]$;
$S$ is the ring of $\sum a_nu^n$ with $|a_n|\,|\lfloor n/e\rfloor!|\le1$ and $\to0$ (the $l$-adically completed
divided-power envelope of $(E)$); $\mathrm{Fil}^jS=\bigcap_n(\mathrm{Fil}^j_0+l^nS)$ with $\mathrm{Fil}^j_0$ the ideal of $S$
generated by the $E^i/i!$, $i\ge j$; $R$ is the ring of series converging on $|u|\le|\varpi_0|$.
\emph{$\varphi$-modules.} $\mathrm{Vec}=(L[[u]]^n)^{\Z/a}$. A $\varphi$-module is a family
$A_\tau\in\mathrm{GL}_n(\mathfrak S_L[1/E])$, with $\varphi(v)_\tau=A_\tau\varphi(v_{\tau-1})$. For $B\subseteq L[[u]]$,
$M_B=(B^n)^{\Z/a}$ and $\mathcal L_B$ is the product of the $B$-spans of the columns of the $A_\tau$;
$M=M_{\mathfrak S_L}$, $\mathcal L=\mathcal L_{\mathfrak S_L}$, $M\!M=\mathcal L_S$, $M\!M_l=\mathcal L_{S[1/l]}$,
$M^{\rm an}=M_R$, $\mathrm{Pol}=\mathrm{Vec}/M^{\rm an}$. Morphisms are families of matrices over $\mathfrak S_L$
commuting with $\varphi$.
\emph{Operators.} Kisin's section $\xi:D=(L^n)^{\Z/a}\to\mathrm{Vec}$ is the unique $\varphi$-equivariant map reducing to
the identity mod $u$ ($\varphi_D(d)_\tau=A_\tau(0)d_{\tau-1}$); $\nabla=\xi\circ\frac d{du}\circ\xi^{-1}$, $N_\nabla=u\nabla$,
\[
\zeta(y)=-\sum_{j\ge1}(l\varphi)^j(u\nabla y),\qquad\mathrm{Mon}(y)=\overline{E\,u\nabla y-E\,\zeta(y)}\in\mathrm{Pol};
\]
$\zeta(y)=l\varphi(\zeta(y)-u\nabla y)$.
\emph{Heights, jumps, descent.} The height condition for $(\eta,H)$ is $E^HM\subseteq\mathcal L$, $E^\eta\mathcal L\subseteq M$;
the jumps are the $i$ with $\mathrm{Fil}^{i+1}\ne\mathrm{Fil}^i$, $\mathrm{Fil}^i=(\mathcal L_{\mathfrak S_L[1/l]}\cap E^iM_{\mathfrak S_L[1/l]})+E\mathcal L_{\mathfrak S_L[1/l]}$.
$\Gamma=\mu_e\rtimes\Z/a=\mathrm{Gal}(K_\varpi/\Q_l)$ acts on $L[[u]]^{\Z/a}$; a tame descent datum is a semilinear
representation $\Delta$ of $\Gamma$ on $\mathrm{Vec}$ preserving $M$ and commuting with $\varphi$. Put
$M^{\rm inv}=M\cap\mathrm{Vec}^\Gamma$, $Y=\{y\in\mathrm{Vec}^\Gamma:E^\eta y\in M\}$, $Y'=u^eY$,
$Z=\{\bar y:y\in E^{-\eta}M\}\cap\mathrm{Pol}^\Gamma$, $Z_l$ its $L$-span analogue, $\mathrm{Lc}$ the endomorphism of
$\mathrm{Pol}$ induced by $EN_\nabla$. The characters $\chi_\tau\subseteq\Z/e$ are those through which $\mu_e$ acts on
$D$ in component $\tau$; the datum has type $(\Sigma',U)$ if $\chi_\tau\subseteq\Sigma'$ and the jumps lie in $U$.
\emph{Breuil--Kisin datum of height $(\eta,H)$}: a $\varphi$-module with tame descent datum, the height condition,
$\nabla(M\!M)\subseteq M\!M$ and $EN_\nabla(M)\subseteq M^{\rm an}$.
\emph{Galois realisation.} With $W_O=O\otimes W(C^\flat)$, $A_O=O\otimes A_{\rm inf}$, $B_O=O\otimes B_{\rm cris}$ and
$u\mapsto[\varpi^\flat]$, a \emph{Galois datum} for a lattice $T$ is a $\varphi$-compatible isomorphism
$j:W_O^n\to W_O\otimes T$, a compatible $j_B:B_O^n\to B_O\otimes T$, compatibility of $\Delta$ with $G$ (each $g$
covering $\gamma\in\Gamma$ acts on $j(M_{\mathfrak S_L[1/E]})$ as $\Delta(\gamma)$), $G$-stability of the $A_O$-lattice, and
horizontality ($G_{K_\varpi}$ fixes $j_B(\xi(d))$ when $\xi(d)\in M\!M_l$). A \emph{pack} for $T$ is a Breuil--Kisin datum
with a Galois datum which over $G_{K_\varpi}$ is the evaluation of a prismatic $F$-crystal on the Breuil--Kisin
prism of $\varpi_0$.
\emph{Classes.} For $y\in Y$, a potential is $b=j_B(\sigma)-b_0$ with $\varphi(m_0)-m_0=y$ in $W_O^n$,
$\sigma\in M\!M_l\cap\mathrm{Vec}^\Gamma$, $u\nabla\sigma=\zeta(y)$, and $b_0$ the crystalline value of $\varphi(m_0)$; a
cocycle $\kappa:G\to T[1/l]$ is a class of $y$ if $1\otimes\kappa(g)=(g-1)b$ for a potential $b$.
\emph{Constants.} $r=r(\Sigma')$, $e_W=2(N-1)r$, $e_L=(N+1)\eta$, $c_1=\eta(r+1)$, $c_2=c_1+e_W$,
$k=c_2+\eta+e_L+e_W$.

\medskip

The statements below agree with the formal ones, the formal statements carrying in addition a fixed
standard normalisation of the local Galois data (a choice of a place above $l$ and of $A_{\rm inf}$-coordinates), the
standing assumption $l\ne2$, and the published inputs (Bhatt--Scholze, Fontaine) as explicit arguments; these inputs are themselves isolated as separate
propositions and proved inside the formal development. The notion of
pack is described in words above; a pack includes that $j$ is the \'etale comparison of the prismatic $F$-crystal
evaluated on the Breuil--Kisin prism, which is what Proposition~\ref{i2:bk-c} uses.

In this subsection the Bhatt--Scholze theorem (Theorem~\ref{i2:thm-bs}) and Fontaine's results
(Theorem~\ref{i2:thm-fontaine}) are assumed, as in the formalization, where both are theorems of the
development. The \emph{numerical conditions} are: $N\ge2$, $l\Sigma'\subseteq\Sigma'\ni0$, $|\Sigma'|\le N$,
$0\in U\subseteq[-\eta,H]$, $|U|\le N+1$, $1\le\eta<H$, $40(HN^2+N^3)<l$. They imply
$8H+1<l$, $10H+2\le l$, $14H<l$, $(r+2)\eta<l+1$, $1\le r$ and $4r\le N^2$ (Lemma~\ref{i2:lem-period}).

\begin{proposition}[packs; formally verified, \cite{CharacterFamily}]\label{i2:bk-b}
Every lattice good for $(\Sigma',U)$, $U\subseteq[-\eta,H]$, has a pack whose Breuil--Kisin datum has height
$(\eta,H)$ and type $(\Sigma',U)$.
\end{proposition}
\noindent\emph{Structure} (\S\ref{i2:sst-bk}). The module is the value of the Bhatt--Scholze crystal on the
Breuil--Kisin prism; the connection, the $B_{\rm cris}$-comparison, the tame descent datum and the condition on $EN_\nabla$
are derived from the crystal. Without descent data, the existence of the module is Kisin's theorem
(Theorem~\ref{i2:thm-kisin}); the version with tame descent data and the integral conditions on $\nabla$ and $EN_\nabla$
is not known to us in the literature in this form.

\begin{proposition}[lifting of maps]\label{i2:bk-c}
Let $T_1,T_2$ have packs $P_1,P_2$ and let $f:T_1\to T_2$ be $G$-equivariant. There is a unique morphism of
Breuil--Kisin data $F:P_1\to P_2$ (commuting with $\varphi$ and with the descent data) with
$j_2\circ F=(1\otimes f)\circ j_1$.
\end{proposition}

\begin{proof}
Uniqueness: $j_2$ and $\mathfrak S_L^{\Z/a}\to W_O$ are injective. Existence as a morphism of $\varphi$-modules: by
Theorem~\ref{i2:thm-bs} (full faithfulness, applied to $f$ restricted to $G_{K_\varpi}$ and viewed $\Z_l$-linearly) there is a
morphism $u:\mathcal E_1\to\mathcal E_2$ of the prismatic $F$-crystals of the packs realising $f$. Since each pack is
the evaluation of its $F$-crystal on the Breuil--Kisin prism, $(1\otimes f)(j_1(m))$ lies in $j_2(M_2)$ for
$m\in M_1$; applying this to the standard basis vectors gives a matrix $Y$ over $\mathfrak S_L$ with
$j_2(Yv)=(1\otimes f)j_1(v)$, and $\varphi$-compatibility of $j_1,j_2,f$ with injectivity of $j_2$ gives
$Y\Phi_1=\Phi_2\varphi(Y)$ for the Frobenius matrices. Descent data: each $\gamma\in\Gamma$ is covered by some $g\in G$
(restriction $G\to\Gamma$ is onto). For a basis vector $e_i$ of $M_1$, writing $w$ for the image in $W_O^n$,
$w(F\Delta_1(\gamma)e_i)=F\,g\,w(e_i)=g\,w(Fe_i)=w(\Delta_2(\gamma)Fe_i)$ by the Galois compatibility of
$\Delta_1$, the $G$-equivariance of $F$ on $W_O$-points (from that of $f$ and the $j_\nu$) and that of $\Delta_2$. As
$w$ is injective and both sides are semilinear, $F\Delta_1(\gamma)=\Delta_2(\gamma)F$. Finally $F$ is $O$-linear,
although the theorem was applied $\Z_l$-linearly: for $o\in O$, $j_2(F(om))=(1\otimes f)(j_1(om))=o\,(1\otimes f)(j_1(m))
=j_2(oF(m))$, since $1\otimes f$ and $j_1,j_2$ are $O$-linear, and $j_2$ is injective.
\end{proof}
\noindent\emph{The hypotheses of Theorem~\ref{i2:thm-bs}.} $K_\varpi$ is finite over $\Q_l$ (perfect
residue field), and $T_\nu$, as a $\Z_l$-module, is finite free with $T_\nu[1/l]$ crystalline over $K_\varpi$
(goodness). Only full faithfulness is used. Its specialisation to the Breuil--Kisin prism is Kisin's full
faithfulness (Kisin 2006, Prop.~2.1.12), which is not needed separately.

\begin{proposition}[exactness; formally verified, \cite{CharacterFamily}]\label{i2:bk-ex}
Let $0\to T_0\xrightarrow fT_1\xrightarrow gT_2\to0$ be exact, with packs and lifted morphisms $F,G$
(Proposition~\ref{i2:bk-c}). Then $F$ is injective on $M_0$, $\ker(G|_{M_1})=F(M_0)$, and $M_2/G(M_1)$ has finite length
over $\mathfrak S_L$.
\end{proposition}
\noindent\emph{Structure} (\S\ref{i2:sst-bk}). The formal statement holds for any $\varphi$-modules with identifications and
compatible morphisms (no packs, descent data or height needed). The sequence splits over $W_O$; Cauchy--Binet gives
quasi-inverses $Y'F=f\cdot1$, $GY=g\cdot1$ over $\mathfrak S_L$ with $f,g$ having a unit coefficient, from which saturation of
$F(M_0)$ follows; finiteness of the cokernel uses in addition the $\varphi$-structure. We do not know this statement in the literature in this
form.

\begin{proposition}[the class map; formally verified, \cite{CharacterFamily}]\label{i2:bk-cl}
For a pack of height $(\eta,H)$ for $T$: (i) every $y\in Y$ has a potential; (ii) two classes of $y$ differ by a
coboundary $g\mapsto(g-1)x$, $x\in T$; (iii) if $\kappa$ is a class of $y$, $l^k\kappa$ is a class of $l^ky$; (iv) a class of
$y$ is a class of $y+\varphi(m)-m$ for $m\in M^{\rm inv}$. (v) (Naturality.) If $F:P_1\to P_2$ lifts $f$ as in
Proposition~\ref{i2:bk-c} and $\kappa$ is a class of $y$, then $f\circ\kappa$ is a class of $F(y)$.
\end{proposition}
\noindent\emph{Structure} (\S\ref{i2:sst-bk}). The formal statement holds for any Breuil--Kisin datum with a Galois
datum. Artin--Schreier--Witt surjectivity of $\varphi-1$ on $W_O\otimes T$ gives $m_0$; $\sigma$ comes from the extension of
the module by the trivial module along $y$; (ii)--(iv) use uniqueness and additivity of crystalline values; (v) uses
naturality of $j_B$. A class of $y$ lies in Bloch--Kato's $H^1_f$ by definition; no published statement with this
potential was identified.

\begin{proposition}[monodromy-kernel representatives; formally verified, \cite{CharacterFamily}]\label{i2:bk-mon}
Under the numerical conditions, let $T$ be good for $(\Sigma',U)$ with a pack of height $(\eta,H)$ and $c$ a finite
cocycle of $T$. There is $y\in Y$ with $\mathrm{Mon}(y)=0$ such that $l^{e_W}c$ is a class of $y$.
\end{proposition}
\noindent\emph{Structure} (\S\ref{i2:sst-bk}). The cocycle defines an extension $0\to T\to T_c\to O\to0$ of good
lattices; Propositions~\ref{i2:bk-b}, \ref{i2:bk-c}, \ref{i2:bk-ex} give an exact triple of packs, whose cokernel on
invariants is killed by $l^{e_W}$ by a Lagrange interpolation argument on the $u$-adic degrees in each isotypic part;
lifting the fixed vector of the trivial pack and taking $\varphi(x)-x$ gives $y$. No published statement with the exponent
$e_W$ was identified.

\begin{proposition}[lifting through exact triples; formally verified, \cite{CharacterFamily}]\label{i2:bk-tri}
Under the numerical conditions, let $P_0\xrightarrow FP_1\xrightarrow GP_2$ be an exact triple of Breuil--Kisin data of
height $(\eta,H)$, $P_1$ of type $(\Sigma',U)$, the characters of $P_2$ in $\Sigma'$. If $y_2'\in Y_2'$, $m_2\in M_2^{\rm inv}$ and
$\mathrm{Mon}(y_2'+\varphi m_2)=0$, there are $y'\in Y_1'$, $m\in M_1^{\rm inv}$ with $\mathrm{Mon}(y'+\varphi m)=0$ and
$G(y'+\varphi m)=l^k(y_2'+\varphi m_2)$.
\end{proposition}
\noindent\emph{Structure} (\S\ref{i2:sst-bk}). Lift with loss $l^{e_W}$; the monodromy defect lies in $Z_1$ after
$l^{c_1}$, pulls back to $Z_0$ after $l^{c_2}$, and is corrected by an element of $Y_0'$ using the semisimplicity of
$\mathrm{Lc}$ with eigenvalues $nel$, $n$ in a shifted jump set, at cost $l^{\eta+e_L}$. No published statement
with these exponents was identified.

\begin{proposition}[the $\zeta$-bound]\label{i2:bk-zeta}
Let $P$ be a Breuil--Kisin datum of height $(\eta,H)$ and type $(\Sigma',U)$, under the numerical conditions. If
$y'\in Y'$, $m\in M^{\rm inv}$ and $\mathrm{Mon}(y'+\varphi m)=0$, then $l^{c_1}\zeta(y'+\varphi m)\in M\!M$.
\end{proposition}
\noindent The proof is at the end of this subsection.

\begin{proposition}[integral classes; formally verified, \cite{CharacterFamily}]\label{i2:bk-int}
Let $P$ be a Breuil--Kisin datum of height $(\eta,H)$ with a Galois datum for $T$ (for instance a pack), $\eta<H$,
$8H+1<l$. If $y\in Y$, $\mathrm{Mon}(y)=0$ and
$l^c\zeta(y)\in M\!M$, there is a finite cocycle $\kappa_1:G\to T$ which is a class of $l^cy$.
\end{proposition}
\noindent\emph{Structure} (\S\ref{i2:sst-bk}). Write $g=hk$ with $h\in G_{K_\varpi}$ and $k$ in the Galois group of the
Kummer tower; $(k-1)b$ is integral, and $\mathrm{Mon}(y)=0$ with $l^c\zeta(y)\in M\!M$ make $l^c(h-1)b$ integral, using
Fontaine's identification of $\varphi$-fixed vectors in $\mathrm{Fil}^0B_{\rm cris}$
(Theorem~\ref{i2:thm-fontaine}) and a divided-power estimate in $A_{\rm cris}$ valid when $2(H+2\eta)<l-1$, which follows from
$\eta<H$ and $8H+1<l$. No published statement with these bounds was identified.

\begin{remark}[Literature]\label{i2:rem-bklit}
The method is that of Breuil--Kisin modules with tame descent data over $K_\varpi$ (Kisin; Savitt;
Caraiani--Emerton--Gee--Savitt), with the module obtained from Bhatt--Scholze's theorem. Steps that are published
theorems in the needed form: full faithfulness of the realisation (Theorem~\ref{i2:thm-bs}), used in
Propositions~\ref{i2:bk-b}, \ref{i2:bk-c}; Fontaine's $B_{\rm cris}^{\varphi=1}\cap\mathrm{Fil}^0=\Q_l$ and the comparison
isomorphisms (Theorem~\ref{i2:thm-fontaine}), used in Propositions~\ref{i2:bk-b}, \ref{i2:bk-int}; Bloch--Kato's
description of $H^1_f$ by $B_{\rm cris}$-coboundaries, behind Propositions~\ref{i2:bk-cl}, \ref{i2:bk-mon}. What is new, as
far as we could determine (\S\ref{i2:sst-bk}): the integral forms of the connection conditions with tame descent data,
the exactness of lifted triples up to finite length, the cokernel bound with exponent $e_W$ (a Lagrange interpolation over
the $u$-adic degrees of the isotypic parts, with $e_W$ controlled by the Frobenius periods $r$), the orbit bound with
exponent $c_1=\eta(r+1)$, the semisimplicity of $\mathrm{Lc}$ with eigenvalues $nel$ and the resulting exponent $e_L$, and the
$A_{\rm cris}$ estimate giving integral classes for $2(H+2\eta)<l-1$. All bounds are uniform in $l$ under the numerical
conditions. Fontaine--Laffaille theory does not apply here because $K_\varpi$ is ramified of index $l^a-1$.
\end{remark}

\subsubsection*{Lifting for good lattices}

\begin{proposition}[Lifting for good lattices]\label{i2:prop-liftgood}
Under the numerical conditions, with $L\supseteq\Q_{l^a}$, let $T_1,T_2$ be good for $(\Sigma',U)$ and $F:T_1\to T_2$
surjective and $G$-equivariant. Then $F$ has the lifting property with exponent $7(Nr+\eta(N+r))$.
\end{proposition}

\begin{proof}
$l>40HN^2\ge160H$, so $l\ne2$ and $8H+1<l$.
(a) $T_0=\ker F$ is good for $(\Sigma',U)$ and $0\to T_0\to T_1\to T_2\to0$ is exact: $T_0$ is saturated; a
subrepresentation of a crystalline one is crystalline (Theorem~\ref{i2:thm-fontaine}); the inclusion
$V_0\to V_1$ is injective on $D_{\rm cris}$ and equivariant, so $\prod_c(g-\zeta^c)$ kills $D_{\rm cris}(V_0|_{K_\varpi})$; and the
jumps of $V_0$ are jumps of $V_1$.
(b), (c) Packs $P_0,P_1,P_2$ of height $(\eta,H)$ and type $(\Sigma',U)$ (Proposition~\ref{i2:bk-b}); lifts of
$T_0\to T_1$ and $F$ (Proposition~\ref{i2:bk-c}); the triple is exact (Proposition~\ref{i2:bk-ex}).
(e) Given a finite cocycle $c$ of $T_2$, Proposition~\ref{i2:bk-mon} gives $y_2\in Y_2$ with $\mathrm{Mon}(y_2)=0$
and $l^{e_W}c$ a class of $y_2$.
(f) $l^\eta Y_2\subseteq Y_2'+M_2^{\rm inv}$: since $E-l=u^e$, $l^\eta=E^\eta-u^e\sum_{i<\eta}l^iE^{\eta-1-i}$; for $y\in Y$,
$s=\sum_{i<\eta}l^iE^{\eta-1-i}y\in Y$ and $E^\eta y\in M^{\rm inv}$, so $l^\eta y=-u^es+E^\eta y$. Write
$l^\eta y_2=y_2'+m_2$. Then $y_2'+\varphi(m_2)=l^\eta y_2+(\varphi(m_2)-m_2)$ has $\mathrm{Mon}=0$ ($\mathrm{Mon}$ is additive, vanishes
on $\varphi(m)-m$ since $\zeta(\varphi m-m)=u\nabla\varphi m$ and $EN_\nabla(M)\subseteq M^{\rm an}$), and by
Proposition~\ref{i2:bk-cl}(iii),(iv) its class is $l^{\eta+e_W}c$.
(g) Proposition~\ref{i2:bk-tri} gives $y'\in Y_1'$, $m\in M_1^{\rm inv}$ with $\mathrm{Mon}(y'+\varphi m)=0$ and
$G(y'+\varphi m)=l^k(y_2'+\varphi m_2)$.
(h) By Proposition~\ref{i2:bk-zeta}, $l^{c_1}\zeta(y'+\varphi m)\in M\!M$; by Proposition~\ref{i2:bk-int} (with $y=y'+\varphi m\in Y_1$)
there is a finite integral cocycle $\kappa_1$ of $T_1$ which is a class of $l^{c_1}(y'+\varphi m)$.
(i) By naturality (Proposition~\ref{i2:bk-cl}(v)), $F\circ\kappa_1$ is a class of
$l^{c_1}G(y'+\varphi m)=l^{c_1+k}(y_2'+\varphi m_2)$, which by (f) and Proposition~\ref{i2:bk-cl}(iii) has the class
$l^{c_1+k+\eta+e_W}c$; by Proposition~\ref{i2:bk-cl}(ii) the two differ by a coboundary $\beta$:
$l^{c_1+k+\eta+e_W}c=F\circ\kappa_1-\beta$.
(j) $c_1+k+\eta+e_W=\eta(2r+N+5)+6(N-1)r\le7(Nr+\eta(N+r))$, since $6(N-1)r\le7Nr$ and $N+5\le7N+5r$. Multiplying
(i) by the remaining power of $l$ gives $l^{7(Nr+\eta(N+r))}c\in F_*Z^1_f(T_1)+B^1(T_2)$.
\end{proof}

\begin{proof}[Proof of Proposition~\ref{i2:prop-lift}]
(1) \emph{Coefficients.} Put $a=(l-1)!$. By Lemma~\ref{i2:lem-coef} there is $L'\supseteq L\Q_{l^a}$, and it suffices to
prove the lifting property for $F\otimes O_{L'}$ with the same exponent. The hypotheses pass to the extension:
for a basis of $L'/L$, the crystalline and de Rham periods, the $\Sigma$-eigenspaces and the filtration steps of
$L'\otimes_LV$ are the $L'$-spans of those of $V$, of the same dimension (coordinates in the basis); characters
of $\Sigma$ viewed in $L'^\times$ keep $1\in\Sigma$, $\mathrm{Fr}$-stability, orbit lengths and $|\Sigma|$; and
surjectivity is witnessed by an $O_L$-matrix $X$ with $AX=1$.
(2) \emph{Good lattices.} $N\le N^3<l$ and the characters of $\Sigma$ have order prime to $l$, so
Proposition~\ref{i2:prop-goodlat} gives $\Sigma'$ with $r(\Sigma')=r'(\Sigma)$ for which $T_1,T_2$ are good for $(\Sigma',U)$
(the jump condition is a hypothesis).
(3) \emph{Lifting.} The numerical conditions hold, $40(HN^2+N^3)\le C(HN^2+N^3)<l$, and
Proposition~\ref{i2:prop-liftgood} gives the exponent $7(Nr+\eta(N+r))$, $r=r'(\Sigma)$.
(4) Lemma~\ref{i2:lem-period} gives $4r'(\Sigma)\le N^2$. As $C\ge7$ and the lifting property is monotone in the
exponent, $F$ has the lifting property with exponent $e$.
\end{proof}

\subsubsection*{Proof of the $\zeta$-bound}

Write $x=l$ when it is used as a scalar, $w(n)=|\lfloor n/e\rfloor!|$, and $\mathrm{Fil}M\!M=\mathrm{Fil}^{10H}S\cdot M\!M$.

\begin{lemma}[Elementary facts]\label{i2:lem-zfacts}
Let $P$ be a $\varphi$-module.
(F1) There is $\gamma_0\in S$ with $\gamma_0\varphi(E)=l$.
(F2) $\varphi(S)\subseteq S$.
(F3) If $E^kv\in M_S$ then $l^k\varphi(v)\in M\!M$.
(F4) If $E^\eta\mathcal L\subseteq M$, then $E^\eta M\!M\subseteq M_S$ and $l^\eta\varphi(M\!M)\subseteq M\!M$.
(F5) For $w\in M\!M$ and $n\ge0$, $u^{el^{n+1}}w=l^{l^n}w'$ with $w'\in M\!M$.
(F6) If $v\in\mathrm{Vec}$ and for every $k$ there is $s\in M\!M$ such that for every $m$ there is $g\in M\!M$ with
$v-(s+l^kg)\in u^m\mathrm{Vec}$, then $v\in M\!M$.
(F7) If $x\in M^{\rm an}$ and $E^kx\in M$ for some $k$, then $x\in M$.
\end{lemma}

\begin{proof}
(F1) $\gamma_0=\sum_k(-l^{-1})^ku^{elk}$: $\varphi(E)=u^{el}+l$, and $|l|^{-k}w(elk)=|l|^{-k}|(lk)!|\le|k!|\to0$ since $l^kk!\mid(lk)!$.
(F2) $\lfloor n/e\rfloor!$ divides $\lfloor ln/e\rfloor!$. (F3) $\varphi(E^kv)\in\mathcal L_S$ by (F2), and
$l^k\varphi(v)=\gamma_0^k\varphi(E^kv)$. (F4) The columns of the $A_\tau$ are the $\varphi(b_i)\in\mathcal L$, so
$E^\eta M\!M\subseteq M_S$; apply (F3). (F5) $l^{-l^n}u^{el^{n+1}}\in S$ because $|(l^{n+1})!|\le|l|^{l^n}$.
(F6) Apply $A_\tau^{-1}$ (power series entries) and compare coefficients: each coordinate $f$ of
$A_\tau^{-1}v_\tau$ has $|f_n|w(n)\le1$, and $|f_n|w(n)<|l|^k$ for $n$ large, for every $k$; so $f\in S$.
(F7) It suffices that $h\in R$, $Eh\in\mathfrak S_L$ imply $h\in\mathfrak S_L$. Write $(Eh)_{n+e}=h_n+lh_{n+e}$. Then
$h_n=\sum_{j\ge0}(Eh)_{n+e+je}(-l)^j$: the difference $g$ satisfies $g_n=(-l)^jg_{n+je}$, so
$|g_n|\rho^n=|g_{n+je}|\rho^{n+je}\to0$ ($\rho=|l|^{1/e}$, convergence on the disc), i.e.\ $g=0$. Hence $|h_n|\le1$.
\end{proof}

\begin{lemma}[Integrality computation]\label{i2:lem-zint}
Let $P$ be a Breuil--Kisin datum of height $(\eta,H)$, $\eta<H$, $10H+2\le l$.
(0) Every $m\in M\!M$ is $p+f$ with $p\in\mathcal L$, $f\in\mathrm{Fil}M\!M$.
(1) For $f\in\mathrm{Fil}M\!M$, $E^{-2H}f\in M^{\rm an}$.
(2) For $y\in E^{-\eta}M$, $E^{2H}EN_\nabla(y)=p+f$ with $p\in\mathcal L$, $f\in\mathrm{Fil}M\!M$.
\end{lemma}

\begin{proof}
Two facts on $S$: (A) for $j<l$, every $f\in S$ is $g+\phi$ with $g\in\mathfrak S_L$, $\phi\in\mathrm{Fil}^jS$; (B)
$\mathrm{Fil}^jS\subseteq E^jR$. For (A), expand $u^{qe+r}=u^r(E-l)^q=\sum_k\binom qkk!(-l)^{q-k}u^rE^k/k!$; the coefficients are
bounded by $|a_{qe+r}|w(qe+r)$ since $|l^{q-k}|\le|(q-k)!|$, so collecting the terms with $k<j$ (where $|k!|=1$) gives
$g\in\mathfrak S_L$, and the rest lies in $\mathrm{Fil}^j_0+l^nS$ for every $n$ after truncating the tail of $f$. For (B),
$\mathrm{Fil}^j_0\subseteq E^jR$ directly, and $E^jR$ is closed in the norm $\sup|f_n|\rho^n$ (by induction on $j$, using the
division by $E$ of (F7), which is continuous); elements of $\mathrm{Fil}^jS$ are limits of elements of $\mathrm{Fil}^j_0$ in this
norm since $|l|^n|h_k|\rho^k\le|l|^n|h_k|w(k)$.
(0) Write $m_\tau=A_\tau c_\tau$, $c_\tau\in S^n$, and apply (A) with $j=10H<l$ to each coordinate. (1) By (B)
$f$ is a sum of terms $E^{10H}q\,m$, $q\in R$, $m\in M\!M$, and $E^{-2H}E^{10H}qm=(E^{8H-\eta}q)(E^\eta m)\in M_R$ by (F4).
(2) $E^{H+\eta}y\in E^HM\subseteq\mathcal L\subseteq M\!M$, so $N_\nabla(E^{H+\eta}y)\in M\!M$ ($\nabla$ preserves $M\!M$). Apply (0) to
$E^{H-\eta}EN_\nabla(E^{H+\eta}y)$ and the Leibniz rule $E^kEN_\nabla(y)=EN_\nabla(E^ky)-k\,eu^eE^ky$ with $k=H+\eta$; the
correction $E^{H-\eta}(H+\eta)eu^eE^{H+\eta}y$ lies in $\mathcal L$.
\end{proof}

\begin{lemma}[Fine bound]\label{i2:lem-zfine}
Let $P$ be a Breuil--Kisin datum of height $(\eta,H)$, $1\le\eta<H$, $14H<l$. For $y'\in Y'$,
$\zeta(y')=l^{l-14H}p+l^{-3H}f$ with $p\in\mathcal L$, $f\in\mathrm{Fil}M\!M$.
\end{lemma}

\begin{proof}
Write $y'=u^ey_0$, $E^\eta y_0\in M$, $N=u\nabla$.
(a) For $j\le l$ and $w\in M\!M$, $u^{el}w=l^{l-j}s+f$ with $s\in M\!M$, $f\in\mathrm{Fil}^jS\cdot M\!M$: expand
$u^{el}=(E-l)^l$ and split the terms with $i<j$ and $i\ge j$.
(b) For $n\ge0$ there is $w\in M\!M$ with $l^{H+2\eta+n(\eta-1)}(l\varphi)^{n+1}(Ny')=u^{el^{n+1}}w$. For $n=0$:
$l\varphi(Ny')=lu^{el}(e\varphi(y_0)+\varphi(Ny_0))$, and $l^{H+2\eta+1}\varphi(y_0)$, $l^{H+2\eta+1}\varphi(Ny_0)\in M\!M$ by (F3), since
$E^\eta y_0\in M$ and $E^{H+2\eta+1}Ny_0\in M_S$ (from $E^{H+\eta}y_0\in\mathcal L\subseteq M\!M$, the Leibniz rule and (F4)). The step
multiplies by $l^{\eta-1}l\varphi$ and uses (F4).
(c) $\zeta(y')+l\varphi(Ny')=l^{l-14H}r$ with $r\in M\!M$: the terms $(l\varphi)^{n+2}(Ny')$ are, by (b) and (F5),
$l^{-B}l^{l^{n+1}}w'$ with $B=H+2\eta+(n+1)(\eta-1)$, and $B+(l-14H)+n\le l^{n+1}$; (F6) concludes.
Now apply (a) with $j=10H$ to the $w_1$ of (b) for $n=0$, and Lemma~\ref{i2:lem-zint}(0) to the elements of $M\!M$
obtained; with $A=l-14H$, $B'=3H-2\eta$, $C'=H+2\eta$, $D'=2H-2\eta$ (all $\ge0$),
$\zeta(y')=l^A(p_2-l^{B'}p_1)+l^{-3H}\bigl(l^{l-11H}f_2-l^{D'}(l^{l-10H}f_1+f)\bigr)$, using $l-10H=A+B'+C'$ and $l-11H=A+D'+C'$.
\end{proof}

\begin{proof}[Proof of Proposition~\ref{i2:bk-zeta}]
\emph{Step 1: $l^{c_1}\zeta(\varphi m)\in M\!M$ for $m\in M^{\rm inv}$.} Let $w=u\nabla\varphi(m)\in M\!M$; $w\equiv0\bmod u$ and $w$
is $\mu_e$-invariant. If $w\in u^j\mathrm{Vec}$ and $-j\notin\chi_\tau$ for all $\tau$, then $w\in u^{j+1}\mathrm{Vec}$
($\mu_e$ acts on the reduction of $u^{-j}w$ through the character $-j$). Hence $w\in u^s\mathrm{Vec}$ with $s=e$ or
$-s\in\chi_\tau$; in the latter case $\rho=\mathrm{per}(-s)\le r$ ($0\in\Sigma'$) gives $e\mid(l^\rho-1)s$; in both cases
$e\le sl^r$. Then $\varphi(w)=u^{sl}x'$ with $l^\eta x'\in M\!M$ (coefficientwise, as $|\lfloor lq/e\rfloor!|\le|\lfloor(q+s)/e\rfloor!|$,
and (F4)), and $(l\varphi)^{n+1}w=l^{n+1}u^{sl^{n+1}}\varphi^n(x')$ with $l^{(n+1)\eta}\varphi^n(x')\in M\!M$. For $n\le r$ the factor
$l^{\eta(r+1)}$ suffices; for $n=r+q$ ($q\ge1$) write $u^{sl^{r+q+1}}=l^{l^q}\cdot u^{sl^{r+q+1}-el^{q+1}}\cdot(u^{el^{q+1}}/l^{l^q})$ with the
last factor in $S$, and $\eta(r+1)+(r+q+1)+l^q\ge(r+q+1)\eta+q$ ($q\eta\le ql\le l^q$). So
$l^{\eta(r+1)}(l\varphi)^{n+1}w=l^{\max(n-r,0)}g_n$, $g_n\in M\!M$, and (F6) gives $l^{c_1}\zeta(\varphi m)\in M\!M$.

\emph{Step 2.} Let $y=y'+\varphi m$ and $w=E\zeta(y)-EN_\nabla y$; $\mathrm{Mon}(y)=0$ means $w\in M^{\rm an}$. From
$\zeta(y)=l\varphi(\zeta(y)-u\nabla y)$ and (F1), $\zeta(y)=\gamma_0\varphi(w)$. By Step~1 and Lemma~\ref{i2:lem-zint}(0),
$l^{c_1}\zeta(\varphi m)=p'+f'$; by Lemma~\ref{i2:lem-zint}(2) (as $y\in E^{-\eta}M$), $E^{2H}EN_\nabla y=p''+f''$; by
Lemma~\ref{i2:lem-zfine}, $\zeta(y')=l^{l-14H}p_1+l^{-3H}f_1$. Combining,
$l^{c_1}E^{2H}w=p+l^{-3H}f$ with $p\in\mathcal L$, $f\in\mathrm{Fil}M\!M$ and $E^{-2H}f\in M^{\rm an}$ (Lemma~\ref{i2:lem-zint}(1)).
Put $X_0=E^{-2H}p$ and $Q=E^{-2H}l^{-3H}f$, so $l^{c_1}w=X_0+Q$, $Q\in M^{\rm an}$, hence $X_0\in M^{\rm an}$; and
$E^{2H+\eta}X_0=E^\eta p\in M$, so $X_0\in M$ by (F7) and $\varphi(X_0)\in\mathcal L\subseteq M\!M$. For $Q$: since $10H<l$,
$\varphi(\mathrm{Fil}^{10H}S)\subseteq l^{10H}S$ (Legendre: $v_l(i!)\le i-10H$ for $i\ge10H$, and $\varphi(E)=l\,d$ with $d\in S$), and
$l^\eta\varphi(M\!M)\subseteq M\!M$ (F4); so $l^\eta\varphi(f)=l^{10H}v'$ with $v'\in M\!M$, and
$\varphi(Q)=\varphi(E)^{-2H}l^{-3H}l^{10H-\eta}v'=\gamma_0^{2H}l^{5H-\eta}v'\in M\!M$. Therefore
$l^{c_1}\zeta(y)=\gamma_0(\varphi(X_0)+\varphi(Q))\in M\!M$.
\end{proof}

\subsection{Commuting endomorphisms and the generator primes}\label{i2:ssec-pf-gen}

\subsubsection{An index bound for families of commuting endomorphisms}\label{i2:ssec-famcore}

\begin{lemma}[Index bound from a family of endomorphisms]\label{i2:lem-famcore}
Let $O$ be a discrete valuation ring with valuation $v$ and fraction field $K$, and let $M$, $X$, $T\subseteq O^X$
(separating points), $V=K\otimes M$, $L$, $E_x$, $\mathrm{mult}(x)$, $E_Y$, $M^{\rm up}(Y)$, $M^{\rm low}(Y)$ be as in
Lemma~\ref{i2:lem-length-le}. Let $Y\subseteq X$ and let $\mathcal E$ be a set of $O$-linear endomorphisms of $M$ commuting
with $T$ such that every $\beta\in\mathcal E$ acts by $0$ on $E_x$ for $x\notin Y$. Suppose that for every $\pi\in Y$ there
are $\alpha_\pi\in\mathcal E$ and $0\ne a_\pi\in O$ with $\alpha_\pi=a_\pi$ on $E_\pi$. Then
\[
\ell\bigl(M^{\rm up}(Y)/M^{\rm low}(Y)\bigr)\le\sum_{\pi\in Y}\mathrm{mult}(\pi)\,v(a_\pi).
\]
\end{lemma}

\begin{proof}
\emph{Step 0: decomposition.} As in the proof of Lemma~\ref{i2:lem-length-le} there are $g_x\in T$ and $0\ne d\in O$ with
$\sum_xg_x=d$, $g_xV\subseteq E_x$ and $g_x=0$ on $E_y$ for $y\ne x$; since $g_x$ acts on $E_x$ by the scalar $g_x(x)=d$,
this gives $V=\bigoplus_xE_x$. An endomorphism $\beta$ commuting with $T$ preserves each $E_x$ (if $fe=f(x)e$ for all
$f\in T$, then $f\beta e=\beta fe=f(x)\beta e$); in particular every $\beta\in\mathcal E$ maps $V$ into $E_Y$, and every
sum of eigenspaces is $\mathcal E$-stable.

\emph{Step 1: reduction to $L$.} Let $P=d^{-1}\sum_{x\in Y}g_x$, the projection onto $E_Y$ along $E_{X\setminus Y}$. For
$m\in L$, $Pm\in E_Y$ and $Pm=m-(m-Pm)\in L+E_{X\setminus Y}$, so $m\mapsto Pm$ maps $L$ to $M^{\rm up}(Y)$. It is
surjective: if $u=m+e\in E_Y$ with $m\in L$, $e\in E_{X\setminus Y}$, then $u=Pu=Pm$. Elements of $L\cap E_Y$ map to
themselves, hence into $M^{\rm low}(Y)$, and elements of $L\cap E_{X\setminus Y}$ map to $0$. Therefore
\[
\ell\bigl(M^{\rm up}(Y)/M^{\rm low}(Y)\bigr)\le\ell\bigl(L/(L\cap E_Y+L\cap E_{X\setminus Y})\bigr).
\]

\emph{Step 2: the induction.} For a finitely generated $\mathcal E$-stable $O$-submodule $N\subseteq V$ and an
$\mathcal E$-stable $K$-subspace $U\subseteq V$ put $D(N,U)=\mathcal EN+(U\cap N)$, where $\mathcal EN$ is the $O$-span
of the $\beta n$ ($\beta\in\mathcal E$, $n\in N$). We claim: if $S\subseteq Y$ and $N\subseteq E_S+U$, then
$\ell(N/D(N,U))\le\sum_{x\in S}\mathrm{mult}(x)v(a_x)$. For $S=\emptyset$, $N\subseteq U$ and $D(N,U)=N$. Let
$S=S'\cup\{x\}$ with $x\notin S'$, put $U_1=E_{S'}+U$ (an $\mathcal E$-stable subspace), $\beta=\alpha_x$, $a=a_x$, and
$N_1=N\cap U_1$, which is finitely generated, $\mathcal E$-stable, and contained in $E_{S'}+U$.
(a) $aN\subseteq\mathcal EN+N_1$: write $n=e+u_1$ with $e\in E_x$, $u_1\in U_1$; then $\beta n=ae+\beta u_1$, so
$an-\beta n=au_1-\beta u_1$ lies in $U_1$ and in $N$.
(b) $\ell(N/(aN+N_1))\le\mathrm{mult}(x)v(a)$: $F=N/N_1$ embeds into $(E_x+U_1)/U_1$, a $K$-space of dimension
$\le\dim E_x$; being finitely generated and torsion free, $F$ is free of rank $\rho\le\mathrm{mult}(x)$, and
$N/(aN+N_1)\cong F/aF$ has length $\rho\,v(a)$.
(c) Put $A=D(N,U)$ and $B=\mathcal EN+N_1$. Then $A\subseteq B=A+N_1$, the kernel of $N/A\to N/B$ is the image of $N_1$,
and $D(N_1,U)=\mathcal EN_1+U\cap N_1\subseteq A$; hence
$\ell(N/A)\le\ell(N/B)+\ell(N_1/D(N_1,U))$. By (a), $\ell(N/B)\le\ell(N/(aN+N_1))$, which (b) bounds, and the induction
hypothesis for $(S',N_1)$ bounds the last term.

\emph{Step 3.} Apply Step 2 with $S=Y$, $U=E_{X\setminus Y}$, $N=L$ ($L$ is $\mathcal E$-stable since $\mathcal E$ consists of
endomorphisms of $M$). As $\mathcal EL\subseteq L\cap E_Y$ (Step 0), $D(L,U)\subseteq L\cap E_Y+L\cap E_{X\setminus Y}$, so
$\ell(L/(L\cap E_Y+L\cap E_{X\setminus Y}))\le\ell(L/D(L,U))$. Combine with Step 1.
\end{proof}

\subsubsection{Proof of Proposition~\ref{i2:prop-cec}}\label{i2:ssec-proof-cec}

\noindent We recall the part of the member structure that the next lemma uses. For each form $P$
of the row of $i$ there is a set $\mathrm{Adm}(P)$ of \emph{admissible parameter lists} (finite lists of parts; each part
carries a multiplicity $a\ge0$, a cuspidal parameter and a dimension). Write $\tau$ and $\tau'$ for twisting every part of a
list by $\eta_T$ and by $\eta_T^{-1}$, $\mu_i$ for the part attached to the central character of $i$, and
$m_i=\tau'(\mu_i)$ (the part with $a=1$, parameter $\eta_T^{-1}\mu_i$ and dimension $t$). The clauses used are:
\begin{enumerate}
\item[(Ad)] (BC) every form has an admissible list; every part of an admissible list has $a\ge1$; and $\mathrm{Adm}(P)$
  is stable under permutation of lists;
\item[(Tw)] (from the base tie tie0 and (SM1)) $\tau'(\tau(p))=p$ for every part $p$ with $a\ge1$;
\item[(Perm)] if $P_0,P_1$ have the same eigensystem, every list in $\mathrm{Adm}(P_0)$ is a permutation of every list in
  $\mathrm{Adm}(P_1)$;
\item[(Z)] $x\in Z_i$ iff some form $P_0$ with eigensystem $x$ has an admissible list containing $m_i$;
\item[(Zg)] $x\in Z^{\rm good}_i$ iff $x\in Z_i$ and there are a list $\Sigma$ with $\Sigma$ good and a form $P$ with
  eigensystem $x$ such that $\tau'(\Sigma\cdot\mu_i)\in\mathrm{Adm}(P)$ (we write $\Sigma\cdot\mu$ for $\Sigma$ followed by
  $\mu$, and say that $\Sigma$ is \emph{a parameter of $x$});
\item[(Par)] (pending datum) if $\Sigma$ is a parameter of $x$, there is a parameter $\sigma'$ whose local components
  are those of $\Sigma$ at every place; $\sigma'$ is good iff $\Sigma$ is good;
\item[(Mt)] the parameter space of $\sigma'$ is the sum of the type spaces $W_i(P)$ over the forms $P$ for which some
  $\Sigma$ with the local components of $\sigma'$ has $\tau'(\Sigma\cdot\mu_i)\in\mathrm{Adm}(P)$;
\item[(Lab)] the label of $\sigma'$ at $(k,v)$ is a monic polynomial $\mathrm{lab}_{kv}(\sigma')$ of degree $r$; $\sigma'$
  is good iff the tuple $(\mathrm{lab}_{kv}(\sigma'))_{kv}$ is good; for each $(k,v)$ the $e_{\rm ser}(f)$ ($f$ monic of
  degree $r$) are orthogonal central idempotents with sum $1$;
\item[(SerLoc)] if $F$ is in the parameter space of $\sigma'$ and $e_{\rm ser}(f)$ in the factor $(k,v)$ does not kill $F$,
  then $f=\mathrm{lab}_{kv}(\sigma')$.
\end{enumerate}
In the formalization (Z), (Zg), (Mt), (Lab) are fields of the datum, (Par) is a field of the pending datum, (Ad) is the
datum clause (BC) together with a lemma producing admissible lists, (Perm) is a clause of the datum (see also
Lemma~\ref{i2:lem-bcperm}), (Tw) is a lemma (from tie0 and (SM1): the two parts have the same local components almost
everywhere by tie0, so are equal by the uniqueness clause of (SM1)), and the equivalence of goodness of $\sigma'$ and
$\Sigma$ is a formally verified lemma.

\begin{lemma}[action of $e_0$ on type vectors over $Z_i$]\label{i2:lem-T3Z}
Assume (Ad), (Tw), (Perm), (Z), (Zg), (Par), (Mt), (Lab) and (SerLoc). Let $x\in Z_i$, $P$ a form with eigensystem $x$ and
$F\in W_i(P)$. Then $e_0F=F$ if $x\in Z^{\rm good}_i$, and $e_0F=0$ if $x\notin Z^{\rm good}_i$.
\end{lemma}

\begin{proof}
\emph{Step 1: one parameter.} Let $F$ lie in the parameter space of $\sigma'$ and put $p_{kv}=\mathrm{lab}_{kv}(\sigma')$.
By (SerLoc), $e_{\rm ser}(f)F=0$ in the factor $(k,v)$ for $f\ne p_{kv}$; as the $e_{\rm ser}(f)$ sum to $1$,
$e_{\rm ser}(p_{kv})F=F$. The factors commute (they act through distinct tensor factors and are central), so for a
tuple $q$, $e(q)F=F$ if $q=p$ and $e(q)F=0$ otherwise. Summing over good $q$: $e_0F=F$ if $p$ is good and $e_0F=0$ if
not, that is (Lab), according as $\sigma'$ is good or not.

\emph{Step 2: a parameter of $x$ through $P$.} If $\Sigma$ is a parameter of $x$ through some form $P_2$ with eigensystem
$x$, then $\tau'(\Sigma\cdot\mu_i)\in\mathrm{Adm}(P)$: by (Ad) $P$ has an admissible list $L$, by (Perm)
$\tau'(\Sigma\cdot\mu_i)$ is a permutation of $L$, and (Ad) gives the claim. By (Par) take $\sigma'$ with the local
components of $\Sigma$; by (Mt), $F\in W_i(P)$ lies in the parameter space of $\sigma'$, and $\sigma'$ is good iff
$\Sigma$ is good.

\emph{Case $x\in Z^{\rm good}_i$.} (Zg) gives a good parameter $\Sigma$ of $x$; Step 2 and Step 1 give $e_0F=F$.

\emph{Case $x\in Z_i\setminus Z^{\rm good}_i$.} By (Z) there are $P_0$ over $x$ and $L_0\in\mathrm{Adm}(P_0)$ with
$m_i\in L_0$. Let $L\in\mathrm{Adm}(P)$ (Ad); by (Perm) $L$ is a permutation of $L_0$, so $m_i\in L$, say
$L=L_1\,m_i\,L_2$. Then $L_1L_2m_i$ is a permutation of $L$, hence admissible for $P$. Put $\Sigma=\tau(L_1L_2)$. Since
$\tau'$ acts part by part, $\tau'(\Sigma\cdot\mu_i)=\tau'(\tau(L_1L_2))\,m_i=L_1L_2m_i$ by (Tw) (all parts of $L$ have
$a\ge1$). So $\Sigma$ is a parameter of $x$ through $P$; take $\sigma'$ as in Step 2. If $\sigma'$ were good, $\Sigma$
would be good and (Zg) would give $x\in Z^{\rm good}_i$. So $\sigma'$ is not good and Step 1 gives $e_0F=0$.
\end{proof}
\noindent This is the consequence of (T3) used in Proposition~\ref{i2:prop-cut}(2), restricted to $Z_i$; the proof above is complete (in the formalization Step 1 and the second case are separate
lemmas). The clauses named in the
hypothesis are clauses of the norm-choice facts and of the member.

\begin{proof}[Proof of Proposition~\ref{i2:prop-cec}]
(i) Let $i$ be twisted with $l_i\notin\{g_0,g_1\}$. By $\mathrm{Tie4}$ there is a pending datum: coefficient data with (T1),
an $O_i$-algebra map $z_O:Z(O_i[G])\to\mathrm{End}_{O_i}(M_i)$ commuting with $T_i$ (T2), satisfying the eigen-span clause
(T4). By the integrality argument of Proposition~\ref{i2:prop-cut}(1), which uses only (T1), there is a central
$z\in Z(O_i[G])$ with $\iota(z)=|G|e_0$. Put $\alpha=z_O(z)$; it is $O_i$-linear and commutes with $T_i$. Let $x\in Z_i$ and
$F$ a type vector over $x$. By Lemma~\ref{i2:lem-T3Z} (its hypotheses are clauses of the norm-choice facts and
$\mathrm{SerLoc}$ is a clause of $\mathrm{Tie4}$), $\psi(z)F=|G|e_0F$ equals $\iota(|G|)F$ if $x\in Z^{\rm good}_i$ and $0$
otherwise. By (T4), $\alpha$ acts on $E_x$ by $|G|$ for $x\in Z^{\rm good}_i$ and by $0$ for $x\in Z_i\setminus Z^{\rm good}_i$.
Global Artin reciprocity is not used in this argument (the formal proof obtains the central element through a package
that also contains steps using it).

(ii) Let $i$ be twisted with $l=l_i\notin\{g_0,g_1\}$, $O=O_i$, and $\alpha$ as in the commuting-endomorphism condition. Let
$\mathcal E=\{\alpha\circ f:f\in T_i,\ f|_{X_i\setminus Z_i}=0\}$ and $Y=Z^{\rm good}_i$. Each $\alpha f$ commutes with $T_i$
($T_i$ is commutative). It acts by $0$ on $E_x$ for $x\notin Y$: if $x\notin Z_i$, $f$ acts on $E_x$ by $f(x)=0$; if
$x\in Z_i\setminus Z^{\rm good}_i$, $f$ preserves $E_x$ and $\alpha$ kills it. Let $\pi\in Y$. The ideal $\mathfrak c_\pi$
is nonzero: the function $f_\pi$ of the proof of Lemma~\ref{i2:lem-length-le} vanishes at every $y\ne\pi$, in particular on
$X_i\setminus Z_i$, and $f_\pi(\pi)\ne0$. As $O$ is a discrete valuation ring, $\mathfrak c_\pi=f(\pi)O$ for some $f$ vanishing
on $X_i\setminus Z_i$, and $\alpha f$ acts on $E_\pi$ by $a_\pi=|G|f(\pi)\ne0$, with
$v(a_\pi)=\ell(O/|G|)+\ell(O/\mathfrak c_\pi)$. Lemma~\ref{i2:lem-famcore} gives
$\ell(M^{\rm up}(Y)/M^{\rm low}(Y))\le\sum_{\pi\in Y}\mathrm{mult}_i(\pi)\bigl(\ell(O/\mathfrak c_\pi)+\ell(O/|G|)\bigr)$. Dividing
by $\ell(O/l)$, and using $\ell(O/|G|)/\ell(O/l)=v_l(|G|)$ (as in Lemma~\ref{i2:lem-etaJ}), gives (adm1) at $i$.
\end{proof}

\subsubsection{Proof of Proposition~\ref{i2:prop-stepglobal}}\label{i2:ssec-proof-stepglobal}

Fix the data of Proposition~\ref{i2:prop-stepglobal}; $O=O_i$, $K=K_i$, $e_O=\ell(O/l)$, $n_l=\#\{w\mid l\}$. The twisted
member $\tilde r_\pi$ has, by the semisimple decomposition, the constituents $W_0=\rho_\pi\otimes\chi_1^{-1}$ (rank $r$; the
$\rho_\pi$-block) and $W_c=\chi_c\chi_1^{-1}$, $1\le c\le t$ (the character blocks; $W_1$ is trivial).

\emph{Cut cocycle modules.} Fix idempotent data with good congruence objects at $\pi$, i.e.\ the ring data of
Proposition~\ref{i2:prop-latbound}: $p:X'\to X_i$, a point $\pi'$ over $\pi$, the ring $R'$ of functions $X'\to M_{n_V}(O)$
spanned over $T'$ by the values of $r$, the diagonal block maps $\psi_0:R'\to M_r(O)$ and $\psi_c:R'\to O$ of the evaluation
at $\pi'$, an element $E\in R'$, and $\delta_O\in O$, $\delta_O\ne0$, with $\ell(O/\delta_O)=e_Ov_\delta$ (a clause of the data;
$\delta_O=\prod_{1<c\le t}(\chi_1(g_0)-\chi_c(g_0))\ne0$ by the cyclotomic clause). For $\gamma\in G_E$ let $y_\gamma=\chi_1(\gamma)^{-1}r(\gamma)E-E$ and let $Y$ be the $R'$-span of the $y_\gamma$
(compare Definition~\ref{i2:def-colmod}). For a block $c\in\{0,\dots,t\}$ of rank $k_c$ let $C_c=C_c(X_s)$ be the
$O$-span of the cochains $\gamma\mapsto\varphi(y_\gamma)\in(K/O)^{k_c}$, where $\varphi:A\to(K/O)^{k_c}$ is $O$-linear on the
module $A$ of all functions $X'\to M_{n_V}(O)$, $\varphi(ay)=\psi_c(a)\varphi(y)$ for $a\in R'$ and $y\in Y$, and $\varphi(y)=0$
for every $y\in Y$ vanishing at all $x'$ with $p(x')\in X_s$.
Write $Z^1$, $B^1$ for continuous cocycles and coboundaries of $G_E$ with values in $W_c\otimes K/O$, and, for a set $C$
of cochains and a place $v$, $\mathrm{res}_{I_v}C$ for the $O$-module of restrictions to the inertia group at $v$. Consider
the conditions: (NC) for $x\in X_s\setminus Z_i$, $\tilde r_x\otimes K$ has no nonzero $G_E$-invariant vector;
(U) for $x\in X_s$ and $v\nmid l$, $v\notin S_i\cup Q_i$, $\tilde r_x$ is trivial on $I_v$; ($\mathrm Q_v$)
$(r_x(\gamma)-1)(r_x(\gamma)-\chi_1(\gamma))=0$ for $x\in X_s$ and $\gamma\in I_v$; ($\mathrm P_v$) the polynomial relation
of (G5) at $v$ (a polynomial over $T_i$ whose value at $\zeta_v$, evaluated at $\pi$, is a unit, kills
$r_x(\gamma)-\chi_1(\gamma)$ after evaluation at the tame generator).

We write $\mathrm{Sel}_\pi\subseteq H^1(G_E,W_0\otimes K/O)$ for the classes whose restriction at every $w\mid l$ lies in the
image of the Bloch--Kato finite part of the lattice, and whose restriction to $H^1(I_v,W_0\otimes K/O)$ vanishes for every
$v\nmid l$, $v\notin S_i\cup Q_i$; a cochain is \emph{in $H^1_f$ at $w$} if its restriction to the
decomposition group at $w$ lies in that finite part (equivalently, the localised cochain lies in the span of the finite
cocycles and the coboundaries; this is the form in which (c) is formally verified, and the equivalence is a formally
verified lemma). Two facts are used in (b): $a_\gamma=\chi_1(\gamma)^{-1}r(\gamma)\in R'$ (definition of $R'$), and
$\psi_c(a_\gamma)$ is the matrix of $\gamma$ on $W_c$ (from the clauses $p(\pi')=\pi$ and the block form of $r_\pi$ and of
$a(\pi')$, $a\in R'$).

\begin{proposition}[cut-cocycle bounds]\label{i2:prop-cutcoc}
In this setting:
\begin{enumerate}
\item[(a)] (determinant) under (NC), $\ell(O/\mathfrak c^{\rm cut}_\pi)\le\sum_{c=0}^t\ell(C_c)$, where
$\mathfrak c^{\rm cut}_\pi=\{f(\pi):f\in T_i,\ f|_{X_s\setminus Z_i}=0\}$;
\item[(b)] $C_c\subseteq Z^1(G_E,W_c\otimes K/O)$;
\item[(c)] (above $l$) every element of $C_c$ restricts at each $w\mid l$ into the $O_w$-span of the cut cocycles of $W_c$
at $w$ relative to $X_s$; consequently, if that span has length $\le e\,\ell(O_w/l)$ modulo $H^1_f$ at every $w\mid l$, there
is $C'\subseteq C_c$ whose elements are in $H^1_f$ at every $w\mid l$, with $\ell(C_c)\le n_le\,e_O+\ell(C')$;
\item[(d)] under (U), $\mathrm{res}_{I_v}C_c=0$ for $v\nmid l$, $v\notin S_i\cup Q_i$;
\item[(e)] (coboundaries) $\ell(C_0\cap B^1)\le2r\,\ell(O/\delta_O)$ and $\ell(C_c\cap B^1)\le2\ell(O/\delta_O)$, $c\ge1$;
\item[(f)] (Selmer) if $C\subseteq Z^1(G_E,W_0\otimes K/O)$ is in $H^1_f$ at every $w\mid l$ and
$\mathrm{res}_{I_v}C=0$ for $v\nmid l$, $v\notin S_i\cup Q_i$, then $\ell(C)\le\ell(C\cap B^1)+\ell(\mathrm{Sel}_\pi)$,
$\mathrm{Sel}_\pi$ the Selmer group of the $\rho_\pi$-block (so $\ell(\mathrm{Sel}_\pi)=e_Os_i(\pi)$);
\item[(g)] (inertia at $S_i\cup Q_i$) for $c\ge1$: if $v\in S_i\cup Q_i$ satisfies ($\mathrm Q_v$), then
$\ell(\mathrm{res}_{I_v}C_c)\le\ell(O/(\zeta_v-1))$; if $v\in S_i$ satisfies ($\mathrm P_v$), then $\mathrm{res}_{I_v}C_c=0$;
\item[(h)] ($\mathrm{Sha}$) for $2\le c\le t$: if $C\subseteq Z^1(G_E,W_c\otimes K/O)$ consists of cocycles that are
coboundaries at every $w\mid l$ and restrict to $0$ on $I_v$ for every $v\nmid l$, then
$\ell(C)\le\ell(C\cap B^1)+e_O\,\ell(\mathrm{Sha}(\mathrm{ld}_D(l),c-1))$.
\end{enumerate}
\end{proposition}
\noindent All items are proved below. Item (a) is a determinant argument for the cut ideal (Lemma~\ref{i2:lem-cutdet}),
with the modules $C_c$ in place of the $\mathrm{Hom}$-modules of Proposition~\ref{i2:prop-latbound}; the first assertion of
(c) constructs, from a global $\varphi$, local cut-cocycle data at $w$; (h) inflates the cocycles to $G_S$ and compares with
the Tate--Shafarevich group of the cyclotomic twist. Item (h) uses the clause $\chi_c=\chi_1\varepsilon^{1-c}$ of the data
($\varepsilon$ the $l$-adic cyclotomic character), and (a) uses the clause (E5) below.

\begin{proof}[Proof of (b), (d), (f)]
(b) A finite sum of cocycles is a cocycle, so it suffices to treat $c(\gamma)=\varphi(y_\gamma)$. Since $r$ is a
homomorphism and $\chi_1$ a character, $a_{\gamma\delta}=a_\gamma a_\delta$, hence
$y_{\gamma\delta}=a_\gamma a_\delta E-E=a_\gamma y_\delta+y_\gamma$ with $a_\gamma\in R'$ and $y_\delta\in Y$. Therefore
$c(\gamma\delta)=c(\gamma)+\psi_c(a_\gamma)c(\delta)=c(\gamma)+\gamma\cdot c(\delta)$. Continuity: $A$ is free of finite rank
over $O$, so $\varphi(A)$ is a finitely generated submodule of $(K/O)^{k_c}$, hence finite, and $\varphi$ kills $l^NA$ for some
$N$. So $c(\gamma)$ depends only on $r(\gamma)\bmod l^N$ and $\chi_1(\gamma)\bmod l^N$, which are locally constant in $\gamma$.

(d) Let $v\nmid l$, $v\notin S_i\cup Q_i$, and $\gamma\in I_v$. By (U), $a_\gamma$ acts trivially at every $x'$ with
$p(x')\in X_s$ (it is the twisted member there), so $y_\gamma(x')=(a_\gamma(x')-1)E(x')=0$ at these $x'$, and $\varphi(y_\gamma)=0$ by
the defining property of $\varphi$. Thus every generator, hence every element, of $C_c$ vanishes on $I_v$.

(f) Let $h:C\to H^1(G_E,W_0\otimes K/O)$ send a cocycle to its class. Its kernel is $C\cap B^1$. By hypothesis the class of
each $c\in C$ is in the finite part at every $w\mid l$, and its restriction to $I_v$ ($v\nmid l$, $v\notin S_i\cup Q_i$) is the
class of the zero cocycle; so $h(C)\subseteq\mathrm{Sel}_\pi$. Hence $\ell(C)=\ell(C\cap B^1)+\ell(h(C))\le\ell(C\cap B^1)+
\ell(\mathrm{Sel}_\pi)$.
\end{proof}

\begin{proof}[Proof of the second assertion of (c)]
Write $\lambda_w$ for the map sending a cochain $z$ of $G_E$ to its localisation at $w$ (restriction to the decomposition
group, transported to $O_w$-coefficients); it is semilinear along the ring isomorphism $O\cong O_w$ of the local data. A semilinear injection along a surjective ring map
does not decrease length, so lengths of images are bounded by lengths over $O$, and $\ell(O_w/l)\le\ell(O/l)=e_O$. For $w\mid l$ let $L_w$ be the $O_w$-span of the
cut cocycles at $w$ and $H_w$ the finite part. For any $O$-submodule $C_0\subseteq C_c$, the map
$C_0\to L_w/(L_w\cap H_w)$, $z\mapsto\lambda_w(z)$, is defined by the first assertion and has kernel
$C_0\cap\lambda_w^{-1}(H_w)$, so $\ell(C_0)\le\ell(C_0\cap\lambda_w^{-1}(H_w))+e\,e_O$. Enumerate the $n_l$ places $w\mid l$
and intersect successively: $C^{(0)}=C_c$, $C^{(m)}=C^{(m-1)}\cap\lambda_{w_m}^{-1}(H_{w_m})$. Then
$\ell(C_c)\le n_le\,e_O+\ell(C')$ with $C'=C^{(n_l)}$, whose elements are in $H^1_f$ at every $w\mid l$.
\end{proof}

\noindent For (e) and (g) we use the following clauses of the good congruence objects, all part of the ring data of
Proposition~\ref{i2:prop-latbound}: (E1) $E\in R'$ ($E=\prod_{1<c\le t}(X_0-r_ce_B)$ with $X_0,e_B$ in the image of $R$ and the
$r_c\in T'$); (E2) $E(x')^2=\delta(x')E(x')$ for all $x'$, where $\delta=\prod_{1<c\le t}(r_1-r_c)\in T'$ and $\delta(\pi')=\delta_O$;
(E3) in the block decomposition at $\pi'$, $E(\pi')$ is $0$ on $W_0$ and on $W_c$ for $c\ge2$, and is $\delta_O$ on $W_1$;
(E4) for $a\in R'$, $a(\pi')$ is block diagonal with blocks $\psi_0(a),\psi_1(a),\dots,\psi_t(a)$; hence
$\psi_c(f\cdot1)=f(\pi')$ for $f\in T'$, $\psi_c(E)=0$ for $c\ne1$, and $\psi_c(r(\gamma))=\chi_c(\gamma)$, $\psi_0(r(\gamma))=\rho_\pi(\gamma)$
(as $p(\pi')=\pi$ and $r_\pi$ is block diagonal with blocks $\rho_\pi,\chi_1,\dots,\chi_t$). Write $M_c(\gamma)$ for the matrix of $\gamma$ on
$W_c$, so $M_c(\gamma)=\chi_1(\gamma)^{-1}\psi_c(r(\gamma))$.

\begin{proof}[Proof of (e)]
For $c=1$ the block $W_1$ is trivial, so $B^1=0$ and there is nothing to prove. Let $c\ne1$, of rank $k_c$ ($k_0=r$,
$k_c=1$ for $c\ge2$). We show: if $w\in(K/O)^{k_c}$ and the coboundary $d w:\gamma\mapsto M_c(\gamma)w-w$ lies in $C_c$, then
$\delta_O^2w=0$. Granting this, $C_c\cap B^1$ is contained in the image under $d$ of the $\delta_O^2$-torsion of
$(K/O)^{k_c}$, which is $\cong(O/\delta_O^2)^{k_c}$, of length $2k_c\ell(O/\delta_O)$; this gives both bounds.

By (E1) write $E=\sum_jf_j\,r(g_j)$ (finite sum, $f_j\in T'$, $g_j\in G_E$). \emph{Step 1.} For every $z\in C_c$,
\[
\delta_O\sum_j\chi_1(g_j)f_j(\pi')\,z(g_j)=0.
\]
The left side is $O$-linear in $z$, so take $z=\varphi(y_\bullet)$ with $\varphi$ as in the definition of $C_c$. Put
$s_j=\chi_1(g_j)f_j\in T'$ and $u=\sum_js_jy_{g_j}\in Y$. Pointwise, $s_jy_{g_j}=f_jr(g_j)E-s_jE$, so by (E2)
$u=E^2-(\sum_js_j)E=(\delta-\sum_js_j)E$, and therefore $Eu=\delta u$ (the factor $\delta-\sum s_j$ is scalar-valued). Applying
$\varphi$ and its $R'$-equivariance ($E,\delta\in R'$, $u\in Y$): $\psi_c(E)\varphi(u)=\psi_c(\delta)\varphi(u)$, i.e.\
$0=\delta_O\varphi(u)$ by (E4). Finally $\varphi(u)=\sum_j\psi_c(s_j)\varphi(y_{g_j})=\sum_j\chi_1(g_j)f_j(\pi')z(g_j)$.

\emph{Step 2.} Evaluating $E=\sum_jf_jr(g_j)$ at $\pi'$ and comparing blocks (E3), (E4):
$\sum_jf_j(\pi')\psi_c(r(g_j))=0$ for $c\ne1$, and $\sum_jf_j(\pi')\chi_1(g_j)=\delta_O$. Hence
$\sum_j\chi_1(g_j)f_j(\pi')M_c(g_j)=\sum_jf_j(\pi')\psi_c(r(g_j))=0$.

\emph{Step 3.} Take $z=dw$. Then $\sum_j\chi_1(g_j)f_j(\pi')(M_c(g_j)w-w)=0-\delta_Ow$ by Step 2, and Step 1 gives
$-\delta_O^2w=0$.
\end{proof}

\begin{proof}[Proof of (g)]
Let $v\in S_i\cup Q_i$; then $v\nmid l$ (a clause of the data). We use: (i) $\chi_c=\chi_1$ on $I_v$ for $1\le c\le t$ (by the clause
$\chi_c=\chi_1\varepsilon^{1-c}$ of the data, the characters $\chi_c\chi_1^{-1}$ are powers of the cyclotomic character, which is
unramified outside $l$), so $W_c$ is trivial on $I_v$ and, by (b), every $z\in C_c$ restricts to a continuous homomorphism $I_v\to K/O$;
(ii) the datum fixes $\gamma_v\in I_v$ such that $I_v$ is contained in the closure of the subgroup generated by $\gamma_v$
and the wild inertia group $P_v$, and $\zeta_v=\chi_1(\gamma_v)$. A continuous homomorphism $h:I_v\to K/O$ kills $P_v$: $P_v$ is
pro-$p$ with $p\ne l$ the residue characteristic of $v$, $h(P_v)$ is finite (compact image in a discrete group) and
$l$-primary, and a finite quotient of a pro-$p$ group is a $p$-group. Hence $h$ is determined by $h(\gamma_v)$: its kernel
is a closed subgroup containing $P_v$ and, if $h(\gamma_v)=0$, also $\gamma_v$, hence all of $I_v$.

\emph{Case ($\mathrm Q_v$).} Let $\gamma\in I_v$ and $z=\varphi(y_\bullet)$ a generator. The element
$r(\gamma)y_\gamma-y_\gamma$ lies in $Y$ ($r(\gamma)\in R'$), and its value at $x'$ is
$\chi_1(\gamma)^{-1}(r_x(\gamma)-1)(r_x(\gamma)-\chi_1(\gamma))E(x')$, $x=p(x')$, which vanishes when $x\in X_s$ by ($\mathrm Q_v$). So
$\varphi$ kills it, i.e.\ $\psi_c(r(\gamma))z(\gamma)=z(\gamma)$, i.e.\ $\chi_c(\gamma)z(\gamma)=z(\gamma)$ by (E4); this is linear in $z$,
so it holds for all $z\in C_c$. With $\gamma=\gamma_v$ and (i): $(\zeta_v-1)z(\gamma_v)=0$. Thus $z\mapsto z(\gamma_v)$ maps
$\mathrm{res}_{I_v}C_c$ injectively (by (ii)) into the $(\zeta_v-1)$-torsion of $K/O$, whose length is at most
$\ell(O/(\zeta_v-1))$ (it is $\cong O/(\zeta_v-1)$ if $\zeta_v\ne1$, and the bound is infinite otherwise).

\emph{Case ($\mathrm P_v$), $v\in S_i$.} Let $P$ be the polynomial of ($\mathrm P_v$): coefficients in $T_i$, and for $x\in X_s$,
$\gamma\in I_v$, $P_x(r_x(\gamma_v))\,(r_x(\gamma)-\chi_1(\gamma))=0$, where $P_x$ has the coefficients evaluated at $x$, and
$P_\pi(\zeta_v)\in O^\times$. Let $a\in R'$ be $x'\mapsto P_{p(x')}(r_{p(x')}(\gamma_v))=\sum_kc_k(p(x'))r(\gamma_v^k)(x')$ (the
$c_k\circ p$ lie in $T'$). For $\gamma\in I_v$, $ay_\gamma\in Y$ and its value at $x'$ is
$\chi_1(\gamma)^{-1}P_x(r_x(\gamma_v))(r_x(\gamma)-\chi_1(\gamma))E(x')=0$ for $p(x')\in X_s$. So
$0=\varphi(ay_\gamma)=\psi_c(a)\varphi(y_\gamma)$. By (E4) and (i), $\psi_c(a)=\sum_kc_k(\pi)\chi_c(\gamma_v)^k=P_\pi(\zeta_v)$, a unit;
so $\varphi(y_\gamma)=0$. Every generator, hence every element, of $C_c$ vanishes on $I_v$.
\end{proof}

\noindent For (a) we use, besides (E1)--(E4), the following facts.
\begin{itemize}
\item From the ring data of Proposition~\ref{i2:prop-latbound} (see its proof):
\begin{itemize}
\item[(H1)] every $f\in T'$ is $\sum_m(a_m\circ p)b_m$ with $a_m\in T_i$;
\item[(H2)] at every $x\in X_i$, the values $a_m(x)$ are determined by the values of $f$ on $p^{-1}(x)$;
\item[(H3)] for block matrices $\Theta_1,\dots,\Theta_n$ over $R'$, the coefficients of $x'\mapsto\det(\sum_kX_k\Theta_k(x'))$
lie in $T'$. This is the argument for (H3) in that proof; it consists of polynomial identities, so it runs with coefficients
in $T'[X_1,\dots,X_n]$.
\end{itemize}
\item From the same ring data:
\begin{itemize}
\item the block maps $\psi_c:R'\to M_{k_c}(O)$ are surjective;
\item $R'$ is left Noetherian.
\end{itemize}
\item (E5) $E(x')\ne0$ for every $x'$. This is a clause of the good congruence objects; there $E(x')$ has rank one.
\end{itemize}

\begin{lemma}[determinant lemma]\label{i2:lem-cutdet}
Let $O$ be a discrete valuation ring with fraction field $K$, $A$ an $O$-algebra, and $J$ a finite set. For $j\in J$ let
$\psi_j:A\to M_{d_j}(O)$ be surjective $O$-algebra maps, let $g\ge1$, let $\mathrm{Rel}\subseteq A^g$ be a left $A$-submodule,
and let $I\subseteq O$ be an ideal. For $\theta=(\theta_1,\dots,\theta_g)\in\mathrm{Rel}^g$, $\theta_p=(\theta_{p,q})_q$, write
$\Psi_j(\theta)=[\psi_j(\theta_{p,q})]_{p,q}\in M_{gd_j}(O)$. Assume that for all $n$ and all $\theta_1,\dots,\theta_n\in\mathrm{Rel}^g$
every coefficient of $\prod_j\det\bigl(\sum_kX_k\Psi_j(\theta_k)\bigr)\in O[X_1,\dots,X_n]$ lies in $I$. Then
\[
\ell(O/I)\le\sum_j\ell(\mathrm{Sol}_j),\qquad
\mathrm{Sol}_j=\Bigl\{w\in\bigl((K/O)^{d_j}\bigr)^g:\ \sum_q\psi_j(\theta_q)w_q=0\ \text{for all }\theta\in\mathrm{Rel}\Bigr\}.
\]
\end{lemma}

\begin{proof}
\emph{One matrix per $j$.} Fix $j$ and $d=d_j$. Let $N_j\subseteq O^{gd}$ be the set of first rows of the $d\times gd$ matrices
$[\psi_j(\theta_q)]_q$, $\theta\in\mathrm{Rel}$, whose other rows vanish. We use that $\psi_j$ is onto: choose $e_{ab}\in A$
with $\psi_j(e_{ab})=E_{ab}$.
\begin{itemize}
\item $N_j$ is an $O$-module: sums and $O$-multiples of such $\theta$ are again such.
\item Every row $a$ of every relation matrix lies in $N_j$: replace $\theta$ by $e_{1a}\theta\in\mathrm{Rel}$.
\item Every $v\in N_j$ occurs as row $a$, with all other rows zero: replace $\theta$ by $e_{a1}\theta$.
\end{itemize}
As $N_j$ is a submodule of $O^{gd}$, it is generated by $gd$ elements (allowing zeros), $v_{p,a}$ with $p\le g$, $a\le d$.
Let $\theta^{(p,a)}\in\mathrm{Rel}$ have $v_{p,a}$ as row $a$ and zeros elsewhere, and $\Theta_{j,p}=\sum_a\theta^{(p,a)}$. Then
$\Theta_j=(\Theta_{j,p})_p\in\mathrm{Rel}^g$ and $M_j:=\Psi_j(\Theta_j)$ is the square matrix with rows $v_{p,a}$.

\emph{$v(\det M_j)\le\ell(\mathrm{Sol}_j)$.} Every relation row is an $O$-combination of rows of $M_j$, so the kernel of $M_j$ on
$(K/O)^{gd}$ lies in $\mathrm{Sol}_j$.
\begin{itemize}
\item If $\det M_j=0$, a nonzero $w\in K^{gd}$ with $M_jw=0$ gives the submodule $\{cw\bmod O^{gd}:c\in K\}$ of infinite
length, so $\ell(\mathrm{Sol}_j)=\infty$.
\item Otherwise, by the Smith normal form, the kernel has length $v(\det M_j)$.
\end{itemize}

\emph{Content.} Enumerate $J=\{j_1,\dots,j_n\}$, put $\theta_k=\Theta_{j_k}$ and $F_j(X)=\det(\sum_kX_k\Psi_j(\theta_k))$. Then
$F_{j_k}$ takes the value $\det M_{j_k}$ at the $k$-th unit vector. We may assume every $\det M_j\ne0$, so $F_j\ne0$; write
$F_j=c_jG_j$ with $c_j\in O$ and $G_j$ primitive. By Gauss's lemma (the reduction of $\prod_jG_j$ modulo the maximal ideal
is a product of nonzero polynomials over a field), $\prod_jG_j$ has a unit coefficient. The hypothesis then gives
$\prod_jc_j\in I$. Since $c_j$ divides every value of $F_j$ on $O^n$,
\[
\ell(O/I)\le\sum_jv(c_j)\le\sum_jv(\det M_j)\le\sum_j\ell(\mathrm{Sol}_j).\qedhere
\]
\end{proof}

\begin{proof}[Proof of (a)]
\emph{Generators and relations.} The module $A$ of all functions $X'\to M_{n_V}(O)$ is finite over $O$, so $Y\subseteq A$ is
a finitely generated module over the left Noetherian ring $R'$. Hence $Y=\sum_{q\le g}R'y_{\gamma_q}$ for some
$\gamma_1,\dots,\gamma_g$. Let $\mathrm{Rel}$ be the set of $\theta\in R'^g$ for which $\sum_q\theta_qy_{\gamma_q}$ vanishes at every
$x'$ with $p(x')\in X_s$; it is a left $R'$-submodule. Apply Lemma~\ref{i2:lem-cutdet} with $A=R'$, $J=\{0,\dots,t\}$,
$\psi_c$ the block maps (surjective), and $I=\mathfrak c^{\rm cut}_\pi$, an ideal of $O$.

\emph{The hypothesis of the lemma.} Let $\theta_1,\dots,\theta_n\in\mathrm{Rel}^g$. Let $\Theta_k(x')$ be the block matrix
$[\theta_{k,p,q}(x')]_{p,q}$, and $f_m(x')$ the coefficients of $\det(\sum_kX_k\Theta_k(x'))$. By (H3), $f_m\in T'$.
\begin{itemize}
\item[(i)] Let $p(x')=x\in X_s\setminus Z_i$. If every $y_\gamma(x')$ vanished, every column of $E(x')$ would be invariant under
$\tilde r_x$, and (NC) would give $E(x')=0$, contradicting (E5). So there are $\gamma_*$ and $\eta\in K^{n_V}$ with
$y_{\gamma_*}(x')\eta\ne0$. Writing $y_{\gamma_*}=\sum_qt_qy_{\gamma_q}$ ($t_q\in R'$), the vector $v=(y_{\gamma_q}(x')\eta)_q$ is nonzero.
\item[(ii)] Each $\theta_{k,p}\in\mathrm{Rel}$ gives $\sum_q\theta_{k,p,q}(x')y_{\gamma_q}(x')=0$. So $(\sum_kX_k\Theta_k(x'))v=0$ over
$K(X)$, the determinant vanishes identically, and $f_m(x')=0$.
\item[(iii)] Thus $f_m$ vanishes on $p^{-1}(X_s\setminus Z_i)$. Write $f_m=\sum_\nu(a_\nu\circ p)b_\nu$ by (H1). By (H2), each $a_\nu$
vanishes on $X_s\setminus Z_i$, so $a_\nu(\pi)\in\mathfrak c^{\rm cut}_\pi$ and $f_m(\pi')\in\mathfrak c^{\rm cut}_\pi$.
\item[(iv)] By (E4), after a fixed permutation of coordinates, $\sum_kX_k\Theta_k(\pi')$ is block diagonal with blocks
$\sum_kX_k\Psi_c(\theta_k)$. So the $f_m(\pi')$ are the coefficients of $\prod_c\det(\sum_kX_k\Psi_c(\theta_k))$.
\end{itemize}
Lemma~\ref{i2:lem-cutdet} gives $\ell(O/\mathfrak c^{\rm cut}_\pi)\le\sum_c\ell(\mathrm{Sol}_c)$.

\emph{$\ell(\mathrm{Sol}_c)\le\ell(C_c)$.} Let $w\in\mathrm{Sol}_c$. On $R'^g$ put $F(\theta)=\sum_q\theta_qy_{\gamma_q}\in Y$ and
$G(\theta)=\sum_q\psi_c(\theta_q)w_q$.
\begin{itemize}
\item $\ker F\subseteq\mathrm{Rel}\subseteq\ker G$. So $\varphi_0(F(\theta))=G(\theta)$ is a well-defined $O$-linear map on $Y=F(R'^g)$.
\item $\varphi_0$ kills the elements of $Y$ vanishing over $X_s$ (their preimages lie in $\mathrm{Rel}$).
\item $\varphi_0(ay)=\psi_c(a)\varphi_0(y)$ for $a\in R'$.
\item $(K/O)^{k_c}$ is an injective $O$-module (divisible over a principal ideal domain). So $\varphi_0$ extends to an
$O$-linear $\varphi:A\to(K/O)^{k_c}$, which is of the kind allowed in the definition of $C_c$.
\end{itemize}
Then $z=(\gamma\mapsto\varphi(y_\gamma))\in C_c$ and $z(\gamma_q)=w_q$. So the $O$-linear map $C_c\to((K/O)^{k_c})^g$,
$z\mapsto(z(\gamma_q))_q$, has image containing $\mathrm{Sol}_c$, and $\ell(\mathrm{Sol}_c)\le\ell(C_c)$.
\end{proof}

\begin{proof}[Proof of the first assertion of (c)]
Let $x'_1,\dots,x'_m$ enumerate $p^{-1}(X_s)$, a finite set. Define $\Lambda:A\to\bigoplus_{j\le m,\,s\le n_V}O^{n_V}$ by
$\Lambda(z)=(z(x'_j)\epsilon_s)_{j,s}$, with $\epsilon_s$ the standard basis. Let $G_E$ act on the target through the twisted
members $\tilde r_{p(x'_j)}$ (one copy for each $(j,s)$). Then $\Lambda(a_\gamma z)=\gamma\cdot\Lambda(z)$, because
$a_\gamma(x')=\tilde r_{p(x')}(\gamma)$.

Put $E'=\Lambda(E)$, so that $\Lambda(y_\gamma)=\gamma E'-E'$. Let $Y_0$ be the $O$-span of the $y_\gamma$. It is stable under every
$a_\gamma$, since $a_\gamma y_\delta=y_{\gamma\delta}-y_\gamma$, and $Y_0\subseteq Y$. So $Y_1=\Lambda(Y_0)$ is a finitely generated, $G_E$-stable
$O$-module containing all $\gamma E'-E'$.

Let $\varphi$ be as in the definition of $C_c$, and define $\varphi'(\Lambda(y))=\varphi(y)$ for $y\in Y_0$.
\begin{itemize}
\item It is well defined: $\Lambda(y)=0$ means that $y$ vanishes over $X_s$.
\item It is $G_E$-equivariant: $\varphi'(\gamma\Lambda(y))=\varphi(a_\gamma y)=\psi_c(a_\gamma)\varphi(y)=M_c(\gamma)\varphi'(\Lambda(y))$, by (E4).
\end{itemize}
Thus $(p(x'_j)_{j,s};E';Y_1;\varphi')$ is a cut-cocycle datum for $W_c$ relative to $X_s$ (with repeated points), and its
cocycle is $\gamma\mapsto\varphi(y_\gamma)$. Restricting to the decomposition group at $w$ and transporting along $O\cong O_w$ gives
a cut cocycle at $w$. Every element of $C_c$ is an $O$-combination of such cocycles, and $\lambda_w$ is semilinear, so its
localisation lies in the $O_w$-span of the cut cocycles.
\end{proof}

\begin{proof}[Proof of (h)]
We use the clause $\chi_c=\chi_1\varepsilon^{1-c}$ of the data, with $\varepsilon$ the $l$-adic cyclotomic character of $G_E$; so $W_c=\varepsilon^{1-c}$.

\emph{Coordinates.} A $\Z_l$-basis of $O$ gives a $\Z_l$-linear isomorphism $p:K/O\to(\Q_l/\Z_l)^d$, $d=[O:\Z_l]=e_Of_O$, and $p$
commutes with the action of $G_E$ through the scalars $\varepsilon^{1-c}(\gamma)\in\Z_l^\times$. For $z\in C$, each coordinate $z_j=p_j\circ z$
is a continuous cocycle $G_E\to\Q_l/\Z_l(\varepsilon^{1-c})$, trivial on $I_v$ for $v\nmid l$ and a coboundary on the decomposition
groups at the places above $l$.

\emph{Inflation.} Let $N\subseteq G_E$ be the closed normal subgroup generated by the inertia groups at the places not above
$l$. Then $\varepsilon$ is trivial on $N$, since $\ker\varepsilon$ is closed, normal and contains these inertia groups. Let
$H=\{n\in N:z_j(n)=0\}$.
\begin{itemize}
\item $H$ is a subgroup, since $z_j(nn')=z_j(n)+z_j(n')$ on $N$.
\item $H$ is closed, since $z_j$ is continuous with discrete target.
\item $H$ is normal in $G_E$: for $n\in H$ and $g\in G_E$, using $z_j(g^{-1})=-g^{-1}z_j(g)$ and $n=1$ on the module,
$z_j(gng^{-1})=z_j(g)+gz_j(n)+gnz_j(g^{-1})=z_j(g)-z_j(g)=0$.
\item $H$ contains the inertia groups outside $l$.
\end{itemize}
Hence $H=N$, and $z_j(an)=z_j(a)+az_j(n)=z_j(a)$. So $z_j$ is the inflation of a continuous cocycle $\bar z_j$ of $G_S$ with values
in $\Q_l/\Z_l(\chi^{-(c-1)})$, $\chi$ the character of $\mathrm{ld}_D(l)$ (Definition~\ref{i2:def-LData}).

\emph{Local triviality.} If $z_j(d)=dm-m$ for $d$ in a decomposition group $D$ and
$g\in G_E$, then for $d'=gdg^{-1}$ we get $z_j(d')=d'm'-m'$ with $m'=gm-z_j(g)$. This is a direct computation from the
cocycle relation. So $z_j$ is a coboundary on the decomposition group of every place of $\Omega$ above $l$, in particular at
the places chosen in $\mathrm{ld}_D(l)$. Hence $[\bar z_j]\in\mathrm{Sha}(\mathrm{ld}_D(l),c-1)$.

\emph{Count.} Let $\Psi:C\to\mathrm{Sha}(\mathrm{ld}_D(l),c-1)^d$, $z\mapsto([\bar z_j])_j$; it is $\Z_l$-linear. If $\Psi(z)=0$, then each
$\bar z_j$ is a coboundary of some $m_j$, so $z$ is the coboundary of $p^{-1}((m_j)_j)$. Thus $\ker\Psi\subseteq C\cap B^1$, and
$\ell_{\Z_l}(C/(C\cap B^1))\le d\,\ell(\mathrm{Sha}(\mathrm{ld}_D(l),c-1))$. For $O$-modules, $\ell_{\Z_l}=f_O\ell_O$. Dividing by
$f_O$ gives $\ell(C)\le\ell(C\cap B^1)+e_O\ell(\mathrm{Sha}(\mathrm{ld}_D(l),c-1))$; this is trivially true if a length is infinite.
\end{proof}

\begin{lemma}[unramified homomorphisms]\label{i2:lem-unrhom}
Assume global Artin reciprocity (Theorem~\ref{i2:thm-artin}). If $C\subseteq Z^1(G_E,K/O)$ (trivial action) is an
$O$-submodule with $\mathrm{res}_{I_v}C=0$ for every finite place $v$, then $\ell(C)\le e_Ov_l(h_E)$.
\end{lemma}

\begin{proof}
Cocycles for the trivial action are continuous homomorphisms $G_E\to K/O$; those trivial on all inertia groups (the inertia groups above $v$ are conjugate and $K/O$ is
abelian, so one of them suffices) factor
through the Galois group of the maximal unramified abelian extension ($E$ is totally complex, so there are no
conditions at infinity), which is isomorphic to the class group $A=\mathrm{Cl}_E$ by Theorem~\ref{i2:thm-artin}. Write
$A\cong\prod_j\Z/n_j$. Then $\mathrm{Hom}(A,K/O)\cong\bigoplus_j(K/O)[n_j]\cong\bigoplus_jO/n_jO$ has length
$\sum_jv(n_j)=e_O\sum_jv_l(n_j)=e_Ov_l(h_E)$, and $C$ is a submodule of it.
\end{proof}

\begin{proof}[Proof of Proposition~\ref{i2:prop-stepglobal}]
Put $d=\ell(O/\delta_O)=e_Ov_\delta$ and $\mathcal Z=\sum_{v\in Q_i}\ell(O/(\zeta_v-1))$, finite since $\zeta_v\ne1$.

\emph{The conditions.} (NC): if $u\ne0$ is $G_E$-invariant for $\tilde r_x$, then $r_xu=\chi_1u$, so $x\in Z_i$ by (G6).
(U): for $\gamma\in I_v$, $r_x(\gamma)=1$ by (G4); also $r_\pi(\gamma)=1$, and $\chi_1(\gamma)$ is an eigenvalue of
$r_\pi(\gamma)$ by the semisimple decomposition, so $\chi_1(\gamma)=1$ and $\tilde r_x(\gamma)=1$. At $v\in S_i$, (G5) and the
tame-generator conditions give either ($\mathrm Q_v$) with $\zeta_v-1$ a unit, or ($\mathrm P_v$); at $v\in Q_i$ they give
($\mathrm Q_v$) and $\zeta_v\ne1$. The cut congruence length is $v^{\rm cut}_i(X_s,\pi)=\ell(O/\mathfrak c^{\rm cut}_\pi)/e_O$, so
by Proposition~\ref{i2:prop-cutcoc}(a) it suffices to bound $\sum_c\ell(C_c)$.

\emph{Block $c=0$.} By (c) with $e=e_\rho$ there is $C'\subseteq C_0$ in $H^1_f$ above $l$ with
$\ell(C_0)\le n_le_\rho e_O+\ell(C')$. By (b), (d) and (U), (f) applies to $C'$, and $\ell(C'\cap B^1)\le\ell(C_0\cap B^1)\le2rd$
by (e). So $\ell(C_0)\le n_le_\rho e_O+2rd+e_Os_i(\pi)$.

\emph{Peeling $Q_i$.} For $c\ge1$, (c) with $e=e_\chi$ gives $C'\subseteq C_c$ in $H^1_f$ above $l$ with
$\ell(C_c)\le n_le_\chi e_O+\ell(C')$. Let $C''=\{z\in C':\mathrm{res}_{I_v}z=0\ \text{for all }v\in Q_i\}$. Then $C'/C''$ embeds
into $\bigoplus_{v\in Q_i}\mathrm{res}_{I_v}C'$, and $\mathrm{res}_{I_v}C'\subseteq\mathrm{res}_{I_v}C_c$ has length
$\le\ell(O/(\zeta_v-1))$ by (g) and ($\mathrm Q_v$). So $\ell(C')\le\ell(C'')+\mathcal Z$. The elements of $C''$ restrict to $0$ on
$I_v$ for every $v\nmid l$: for $v\in Q_i$ by construction; for $v\in S_i$ by (g) (in the case ($\mathrm Q_v$) the bound is
$\ell(O/(\zeta_v-1))=0$); for the other $v$ by (d).

\emph{Block $c=1$.} At $w\mid l$ the elements of $C''$ are in $H^1_f$ of the trivial block, hence trivial on inertia by
$\mathrm{HT}_w$. So $C''$ restricts to $0$ on every inertia group, and Lemma~\ref{i2:lem-unrhom} gives
$\ell(C'')\le e_Ov_l(h_E)$. Hence $\ell(C_1)\le n_le_\chi e_O+e_Ov_l(h_E)+\mathcal Z$.

\emph{Blocks $2\le c\le t$.} By $\mathrm{HT}_w$ the elements of $C''$ are coboundaries at every $w\mid l$; by (b), (h) and (e),
$\ell(C'')\le2d+e_O\ell(\mathrm{Sha}(\mathrm{ld}_D(l),c-1))$. Hence $\ell(C_c)\le n_le_\chi e_O+2d+e_O\ell(\mathrm{Sha}(\mathrm{ld}_D(l),c-1))+\mathcal Z$.

\emph{Sum.} Adding, and using $\ell(\mathrm{Sha}(\mathrm{ld}_D(l),e))=\ell_D(l,e)$ for $0<e<t$ (finite by hypothesis,
\eqref{i2:eq-sha}),
\[
\ell(O/\mathfrak c^{\rm cut}_\pi)\le n_l(e_\rho+te_\chi)e_O+e_Os_i(\pi)+e_O\Bigl(v_l(h_E)+\sum_{0<e<t}\ell_D(l,e)\Bigr)+2(r+t-1)d+t\mathcal Z .
\]
Divide by $e_O$: $d/e_O=v_\delta$, $2(r+t-1)\le4n_V$ as $n_V=r+t$, and $\mathcal Z/e_O=\sum_{v\in Q_i}v(\zeta_v-1)$.
\end{proof}

\subsubsection{Proof of Proposition~\ref{i2:prop-pcc}}\label{i2:ssec-proof-pcc}

Throughout, $D$, $i$, $k\ne k'$, $(r,\rho)$, $\eta,s,S,\mathrm{sg},U$ are as in Proposition~\ref{i2:prop-pcc}, $l=g_k$,
$O=O_i$, $e_O=\ell(O/l)$, $N=n_V+1$ and $e_{\max}=C_T\bigl(N\lfloor N^2/4\rfloor+(\eta+t+1)(N+\lfloor N^2/4\rfloor)\bigr)$.

\begin{proposition}[per-point cut]\label{i2:prop-percut}
Let a pending datum for $D$ be given with (T2) and the eigen-span clause (T4), and assume (SerLoc), (BC), the
local--global compatibility clause (LGC$'$), the base tie tie0, (SM1), the member clauses (Z), (Zg), (Par), (Mt), (Lab)
recalled before Lemma~\ref{i2:lem-T3Z}, and the central-element clause (CE): for every $i$
and every tuple $s=(s_{kv})$ there is a central $z\in Z(O_i[G])$ with $\iota(z)=|G|\,e(s)$. Let $l_i=g_k$. For every
$\pi\in Z^{\rm good}_i$ there are a type locus $X_s(\pi)$ for $(i,k)$, an $O$-linear endomorphism $\beta_\pi$ of $M_i$
commuting with $T_i$, and $0\ne b_\pi\in O$, such that: $\pi\in X_s(\pi)$, $\pi\in Z_i$,
$Z_i\cap X_s(\pi)\subseteq Z^{\rm good}_i$; $\beta_\pi=b_\pi$ on $E_\pi$; $\beta_\pi=0$ on $E_x$ for $x\notin X_s(\pi)$; and
$\ell(O/b_\pi)/e_O\le v_{g_k}(|G|)$.
\end{proposition}
\noindent The proof below is complete modulo Lemma~\ref{i2:lem-isorep}; formally, the conclusion is the conjunction of the cut property and the type-locus property. (CE) holds under (T1) by the integrality argument of
Proposition~\ref{i2:prop-cut}(1) applied to the central idempotent $e(s)$. In Proposition~\ref{i2:prop-pcc} the hypotheses
are supplied by the norm-choice facts and $\mathrm{Tie4}$ (whose pending datum carries (T2) and (T4)). The proof below uses one smaller formally
verified statement:

\begin{lemma}[permutation clause from the datum]\label{i2:lem-bcperm}
Assume tie0, (LGC$'$) and (SM1). Let $P_0,P_1$ be forms of the row of $i$ that both have an eigensystem, and the same
one $x$. Then every list in $\mathrm{Adm}(P_0)$ is a permutation of every list in $\mathrm{Adm}(P_1)$. In particular
clause (Perm) before Lemma~\ref{i2:lem-T3Z} holds.
\end{lemma}

\noindent We fix the notation of the proof. Places are finite places $w$ of $E$; a place is \emph{split} if it is split
over $F^+$. For a form $P$ and a place $w$, $P_w$ is the local component of $P$ at $w$ (a local parameter at $w$); for a
list $\Sigma$ of parts, $\Pi(\Sigma)_w$ is the local component at $w$ of the isobaric sum of the Speh representations
attached to the parts of $\Sigma$, and for a local parameter $a$ and $c\in\R$, $a(c)$ is its twist by $|\cdot|^c$. Put
$c_V=(1-n_V)/2$. The clauses of the member structure and of the definition of $\mathrm{Adm}(P)$ used are:
\begin{enumerate}
\item[(A1)] every part of a list $\Sigma\in\mathrm{Adm}(P)$ has $a\ge1$, $d\ge1$ and conjugate self-dual cuspidal
  parameter, and $\Pi(\Sigma)_w=P_w$ at every split $w$;
\item[(A2)] for a list $\Sigma$ of parts, $\Pi(\Sigma)_w$ is unramified for all but finitely many $w$ (each cuspidal
  parameter is unramified almost everywhere and the Speh construction preserves unramifiedness);
\item[(A3)] if $P$ has eigensystem $x$ and $w$ is split, not above $l=l_i$, not in $S_i\cup Q_i$, then the Frobenius
  polynomial of $P_w(c_V)$ is the polynomial $\prod_\lambda(X-\lambda)$ attached to $x$ at $w$; if moreover $w$ is not
  above $g_0$, $g_1$ or the auxiliary prime of the host ($p_{\rm de}$, or $p_{\rm de}'$ when $l=p_{\rm de}$), then
  $P_w$ is unramified.
\end{enumerate}
\begin{enumerate}
\item[(A4)] (Frobenius clause) if $w\nmid l$, $\rho$ is a continuous representation of $G_E$ and the Weil--Deligne
  representation $\mathrm{WD}(\rho|_{G_{E_w}})$ is unramified, then the characteristic polynomial of $\rho(\mathrm{Frob}_w)$
  is the Frobenius polynomial of $\mathrm{WD}(\rho|_{G_{E_w}})$.
\end{enumerate}
Local parameters are abstract in the automorphic datum; by clauses of the datum, an unramified local parameter is
determined by the multiset of its Frobenius eigenvalues, unramifiedness is preserved by twisting, and twisting is additive
with twist by $0$ the identity, so twisting commutes with the eigenvalue description. From (SM1) we use two of its clauses: \emph{(SM1-dens)} for every finite set $S'$ of places, the Frobenius
elements at split places $w\notin S'$ (for all places of $\bar E$ above $w$) are dense in $G_E$; and \emph{(SM1-iso)}
if two lists $\Sigma,\Sigma'$ whose parts all have $a\ge1$, $d\ge1$ and conjugate self-dual parameter satisfy
$\Pi(\Sigma)_w=\Pi(\Sigma')_w$ for all but finitely many $w$, then $\Sigma'$ is a permutation of $\Sigma$ (strong
multiplicity one for isobaric sums of discrete conjugate self-dual representations). From tie0 and (LGC$'$) we use
one formally verified consequence:

\begin{lemma}[Galois representation of an admissible list; formally verified, \cite{CharacterFamily}]\label{i2:lem-isorep}
Assume tie0 and (LGC$'$). For every form $P$ and $\Sigma\in\mathrm{Adm}(P)$ there is a continuous semisimple
$\rho:G_E\to\mathrm{GL}_{n_V}(\bar\Q_l)$ such that for every $w\nmid l$ the Weil--Deligne representation of
$\rho|_{G_{E_w}}$ is related to $\Pi(\Sigma)_w(c_V)$ by the weak compatibility relation of (LGC$'$); this relation
implies equality when $\Pi(\Sigma)_w(c_V)$ is unramified. Moreover, for $w\mid l$ such that $\Pi(\Sigma)_w$ satisfies the tame and
wild conditions of (LGC$'$), $\rho|_{G_{E_w}}$ is de Rham, its Weil--Deligne representation is related to $\Pi(\Sigma)_w(c_V)$ by
the relation of (LGC$'$) at $l$, and its Hodge--Tate weights are those of $\Pi(\Sigma)_w$ prescribed by (LGC$'$).
\end{lemma}
\noindent The second part is used only in (a1) of Proposition~\ref{i2:prop-pcclocal}. Its proof (structure in
\S\ref{i2:sst-isorep}; no Chebotarev or Brauer--Nesbitt argument is involved) takes for $\rho$ the direct sum, over the parts $(a,\pi,d)$ of $\Sigma$ ($\pi$ a cuspidal parameter on
$\mathrm{GL}_a$) and over $0\le k<d$, of a representation attached by (LGC$'$) to the twist of $\pi$ by
$|\cdot|^{e_k/2}$, $e_k=d-2k-n_V$ (when $e_k+a-1$ is odd the half-integral twist is first made algebraic with the
character ties of tie0), and checks the relation summand by summand; the regularity of the infinity type of $P$ is used
to make the twists algebraic.

\begin{lemma}[split base change]\label{i2:lem-splitbc}
Assume tie0, (LGC$'$), (SM1-dens) and (SM1-iso). Let $P,P'$ be forms with $P_w=P'_w$ for all but finitely many split
$w$, and $\Sigma\in\mathrm{Adm}(P)$, $\Sigma'\in\mathrm{Adm}(P')$. Then $\Sigma'$ is a permutation of $\Sigma$.
\end{lemma}

\begin{proof}
Let $\rho,\rho'$ be given by Lemma~\ref{i2:lem-isorep} for $(P,\Sigma)$ and $(P',\Sigma')$. Let $S'$ be the finite set of
places $w$ that lie above $l$, or at which $\Pi(\Sigma)_w$ or $\Pi(\Sigma')_w$ is ramified (A2), or which are split
with $P_w\ne P'_w$. For $w\notin S'$, $\Pi(\Sigma)_w(c_V)$ and $\Pi(\Sigma')_w(c_V)$ are unramified, so by
Lemma~\ref{i2:lem-isorep} $\mathrm{WD}(\rho|_{G_{E_w}})=\Pi(\Sigma)_w(c_V)$ and similarly for $\rho'$; in particular
$\rho,\rho'$ are unramified at $w$ and, by (A4), the characteristic polynomial of $\rho(\mathrm{Frob}_w)$ is the Frobenius
polynomial of $\Pi(\Sigma)_w(c_V)$. If $w\notin S'$ is split, (A1) gives $\Pi(\Sigma)_w=P_w=P'_w=\Pi(\Sigma')_w$, so
$\rho(\sigma)$ and $\rho'(\sigma)$ have the same characteristic polynomial for every Frobenius element $\sigma$ at $w$.
By (SM1-dens) these $\sigma$ are dense in $G_E$; the coefficients of the characteristic polynomial are continuous
functions of $\rho(\sigma)$, $\rho'(\sigma)$, and $\bar\Q_l$ is Hausdorff, so $\det(X-\rho(\sigma))=\det(X-\rho'(\sigma))$
for all $\sigma\in G_E$. Now let $w\notin S'$ be arbitrary (split or not) and $\sigma$ a Frobenius element at $w$. The
Frobenius eigenvalues of $\Pi(\Sigma)_w(c_V)$ and of $\Pi(\Sigma')_w(c_V)$ are the roots of the same polynomial, so these
unramified parameters coincide, and untwisting by $-c_V$ gives $\Pi(\Sigma)_w=\Pi(\Sigma')_w$ for all $w\notin S'$. The parts
of $\Sigma,\Sigma'$ satisfy the hypotheses of (SM1-iso) by (A1), so $\Sigma'$ is a permutation of $\Sigma$.
\end{proof}

\begin{proof}[Proof of Lemma~\ref{i2:lem-bcperm}]
(SM1-dens) and (SM1-iso) are clauses of (SM1), so Lemma~\ref{i2:lem-splitbc} applies; it suffices to show
$P_{0,w}=P_{1,w}$ for all split $w$ outside the finite set $S''$ of places above $l$, $g_0$, $g_1$ and the auxiliary prime,
and of $S_i\cup Q_i$. For such $w$ both $P_{0,w}(c_V)$ and $P_{1,w}(c_V)$ are unramified by (A3), and their Frobenius
polynomials both equal the polynomial attached to $x$ at $w$ (A3), so they have the same Frobenius eigenvalues and
are equal; untwisting gives $P_{0,w}=P_{1,w}$.
\end{proof}

\begin{proof}[Proof of Proposition~\ref{i2:prop-percut}]
We use the notation and clauses (Ad)--(SerLoc) introduced before Lemma~\ref{i2:lem-T3Z}; (Ad) is (BC), (Tw) follows
from tie0 and (SM1), and (Perm) is Lemma~\ref{i2:lem-bcperm}. By (Zg), $\pi\in Z_i$ and there are a good $\Sigma$ and a form
$P_0$ over $\pi$ with $\tau'(\Sigma\cdot\mu_i)\in\mathrm{Adm}(P_0)$; by (Par) take $\sigma$ with the local components of $\Sigma$,
and let $s_{kv}=\mathrm{lab}_{kv}(\sigma)$. By (CE) take $z$ central with $\iota(z)=|G|e(s)$, and put $\beta=z_O(z)$ (T2),
$b=|G|\in O$, and
\[
X_s=\{x\in X_i:\ \beta\ne0\ \text{on}\ E_x\}.
\]
$\beta$ is $O$-linear and commutes with $T_i$ by (T2); $b\ne0$ as $|G|\ne0$ in $O\subset\C$; and $\ell(O/|G|)/e_O=v_{g_k}(|G|)$
since $l=g_k$ (as in Lemma~\ref{i2:lem-etaJ}). By definition $\beta=0$ on $E_x$ for $x\notin X_s$.

\emph{$\beta=b$ on $E_\pi$.} Let $P$ be any form over $\pi$ and $F\in W_i(P)$. As in Step 2 of the proof of
Lemma~\ref{i2:lem-T3Z}, $\tau'(\Sigma\cdot\mu_i)\in\mathrm{Adm}(P)$, so by (Mt) $F$ is in the parameter space of $\sigma$, and
by Step 1 of that proof $e_{\rm ser}(s_{kv})F=F$ for every $(k,v)$, so $e(s)F=F$ and $\psi(z)F=|G|F=\iota(b)F$. By (T4),
$\beta=b$ on $E_\pi$.

\emph{$\pi\in X_s$.} $E_\pi\ne0$: its dimension is the multiplicity of $\pi$, which equals the dimension of the sum of
the type spaces of the forms over $\pi$ (a clause of the member), and these are nonzero (clauses of the member); there is at
least one such form, $P_0$ of (Zg). As $\beta=b\ne0$ on $E_\pi$,
$\pi\in X_s$.

\emph{A witness over each $x\in X_s$.} If $x\in X_s$, there are a form $P$ over $x$ and $F\in W_i(P)$ with
$e_{\rm ser}(s_{kv})F\ne0$ (in the factor $(k,v)$) for every $(k,v)$. Otherwise for every such $P,F$ some factor
$e_{\rm ser}(s_{kv})$ kills $F$, hence so does the product $e(s)$ (the factors commute), so $\psi(z)F=0$ for all $F$
over $x$ and (T4) with $c=0$ gives $\beta=0$ on $E_x$, a contradiction.

\emph{$Z_i\cap X_s\subseteq Z^{\rm good}_i$.} Let $x\in Z_i\cap X_s$ with witness $P,F$. The second case of the proof of
Lemma~\ref{i2:lem-T3Z} (which uses only (Ad), (Tw), (Perm), (Z)) gives a parameter $\Sigma_1$ of $x$ through a form over
$x$; by Step 2 of that proof and (Par), $F$ lies in the parameter space of some $\sigma_1$ with the local components of
$\Sigma_1$. By (SerLoc) and $e_{\rm ser}(s_{kv})F\ne0$, $s_{kv}=\mathrm{lab}_{kv}(\sigma_1)$ for all $(k,v)$. The labels of
$\sigma$ and $\sigma_1$ are read off from the local components of $\Sigma$ and $\Sigma_1$ at the places $(k,v)$, and
goodness of a list depends only on these labels; so $\Sigma_1$ is good, and (Zg) gives $x\in Z^{\rm good}_i$.

\emph{$X_s$ is a type locus for $(i,k)$.} Take $s_v=s_{kv}$. Since $\Sigma$ is good, each $s_v$ is squarefree and has
neither $c_v$ nor $c_v^{-1}$ as a root (goodness of a list contains these two conditions at every place above $g_k$). For $x\in X_s$ take the witness $P,F$. By definition $e_{\rm ser}(s_v)$ is the sum of the
primitive central idempotents $e_\theta$ of the irreducible representations $\theta$ of $\mathrm{GL}_r(\mathbb F_{g_k})$
with label $s_v$; as $e_{\rm ser}(s_v)F\ne0$, some $e_\theta$ with label $s_v$ does not kill $F$, i.e.\ $\theta$ belongs to the
type of $P$ at $v$. So $x$ is the eigensystem of a form whose type at each $v$ has label $s_v$.
Take $X_s(\pi)=X_s$, $\beta_\pi=\beta$, $b_\pi=b$.
\end{proof}

\begin{lemma}[from a per-point cut to an endomorphism]\label{i2:lem-percutfam}
In the situation of Proposition~\ref{i2:prop-percut}, for $\pi\in Z^{\rm good}_i$ there are an $O$-endomorphism $\alpha_\pi$ of
$M_i$ commuting with $T_i$ and $0\ne a_\pi\in O$ with $\alpha_\pi=a_\pi$ on $E_\pi$, $\alpha_\pi=0$ on $E_x$ for
$x\notin Z^{\rm good}_i$, and $\ell(O/a_\pi)/e_O\le v^{\rm cut}_i(X_s(\pi),\pi)+v_{g_k}(|G|)$.
\end{lemma}

\begin{proof}
Write $X_s=X_s(\pi)$. The ideal $\mathfrak c^{\rm cut}_\pi=\{f(\pi):f\in T_i,\ f|_{X_s\setminus Z_i}=0\}$ is nonzero (the function
$f_\pi$ of Lemma~\ref{i2:lem-length-le} vanishes off $\pi$, and $\pi\in Z_i$), so it equals $g(\pi)O$ for some $g\in T_i$ vanishing
on $X_s\setminus Z_i$ with $g(\pi)\ne0$. Put $\alpha_\pi=g\circ\beta_\pi$ and $a_\pi=g(\pi)b_\pi$. It commutes with $T_i$. On $E_\pi$
it is $g(\pi)b_\pi$. On $E_x$ with $x\notin Z^{\rm good}_i$: if $x\notin X_s$, $\beta_\pi$ kills $E_x$; if $x\in X_s$, then
$x\notin Z_i$ (as $Z_i\cap X_s\subseteq Z^{\rm good}_i$), so $g(x)=0$, and $g$ kills $\beta_\pi E_x\subseteq E_x$. Finally
$\ell(O/a_\pi)=\ell(O/\mathfrak c^{\rm cut}_\pi)+\ell(O/b_\pi)$.
\end{proof}

\begin{proposition}[local inputs at a generator prime; partly formally verified, \cite{CharacterFamily}]\label{i2:prop-pcclocal}
Under the hypotheses of Proposition~\ref{i2:prop-pcc}, let $\pi\in Z^{\rm good}_i$ and let $X_s\ni\pi$ be a type locus for
$(i,k)$. Then: $l$ has exactly $2f$ places $w$ in $E$, all with $E_w=\Q_l$; the global admissibility clauses of
Proposition~\ref{i2:prop-stepglobal} hold at $(i,\pi)$; and at every $w\mid l$ the condition $\mathrm{HT}_w$ holds, as do the
clauses (a1)--(a5) of Proposition~\ref{i2:prop-potcount}(a), the hypotheses of Proposition~\ref{i2:prop-potcount}(b),(c) other
than (Pot) and (Lift), and continuity of the twisted and dual members.
\end{proposition}
\noindent\emph{What is formally verified, item by item.} The formal statement only collects the items; it is the assembly written out in the proof of Proposition~\ref{i2:prop-pcc} below. The items
are of two kinds. (i) Items obtained from the norm-choice facts and $\mathrm{Tie4}$, some read off as clauses and the others by separately
formally verified statements (listed below; their structure is in \S\ref{i2:sst-pcclocal}): (a2)--(a5) of
Proposition~\ref{i2:prop-potcount}(a), the hypotheses of Proposition~\ref{i2:prop-potcount}(b),(c) other than (Pot), (Lift)
and the block embedding, and the global clauses (G1), (G4), (G5), (G6), the tame-generator conditions, the coefficient
compatibility and the good congruence objects of Proposition~\ref{i2:prop-stepglobal}. The separately verified
statements among these are: the tame-generator conditions at every local index (from the local behaviour of the first
local character at wild inertia away from $l$); (G4) at $i$ from the Galois-family conclusion of the datum; (G5) and (G6)
at $i$ from the corresponding Galois-family conclusions of the datum, given the good-point hypothesis at $i$ (itself
derived from $\pi\in Z^{\rm good}_i$); local-coefficient compatibility from the branch data and the degree clause at the
places above $l$; the dual-jump, potential de Rham and tensor-of-potentially-crystalline inputs of
Proposition~\ref{i2:prop-potcount}(b),(c), together with the branch tie and the inertia equality, from the local tie
records of the datum; the good congruence objects from the type-prime congruence lemmas; and the crystalline period
bound with constant $C_T\ge40$, the local pushout proposition, the twisted-base clause (for some Hodge--Tate weight data
and Fontaine data $HW,F$) and the Breuil--Kisin admissibility, from the published
theorems as formalised. (ii) Six steps. We state them exactly; (Deg), (HT), (Cont), (Sub) are proved below, and (Res), (a1)
are separate formally verified statements \cite{CharacterFamily}.
\begin{enumerate}
\item[(Deg)] If $l_i=g_k$, then $l$ has exactly $2f$ places $w$ in $E$, and $E_w=\Q_l$ for each of them.
\item[(HT)] Under the clause $\chi_c=\chi_1\varepsilon^{1-c}$ ($1\le c\le t$, $\varepsilon$ the $l$-adic cyclotomic character of $G_E$) of the
data, for every $w\mid l$ with $E_w=\Q_l$ the condition $\mathrm{HT}_w$ of Proposition~\ref{i2:prop-stepglobal} holds, with
$H^1_f$ of $W_c\otimes K/O$ taken as in \S\ref{i2:ssec-proof-potcount} (images of finite cocycles of $W_c\otimes K$ plus
coboundaries). (The formal proof uses instead the twisted-base clause of the local data, the inertia equality and
$l\ne2$.)
\item[(Cont)] If $(r_x)$ is a good family (each $r_x$ continuous with the Frobenius clauses), then for every $x$ the
twisted member $\tilde r_x=\chi_1^{-1}r_x$ and the dual member $r_x^\vee:g\mapsto\chi_1(g)\epsilon(g)\,{}^tr_x(g^{-1})$
($\epsilon=\chi_1\chi_2^{-1}$) are continuous.
\item[(Sub)] If $(r_x)$ is good, $\pi$ has a semisimple decomposition (a $K$-linear isomorphism
$e:K^{n_V}\to K^r\times K^t$ with $e\,r_\pi(g)=(\rho_\pi(g)\oplus\mathrm{diag}(\chi_1(g),\dots,\chi_t(g)))\,e$ for all $g$), and
$\pi\in X_s$, then each block $W_0=\rho_\pi\otimes\chi_1^{-1}$ and $W_c=\chi_c\chi_1^{-1}$ ($1\le c\le t$) is continuous and
admits an injective $K$-linear $G_E$-map into $\tilde r_x\otimes K$ for some $x\in X_s$.
\item[(Res)] Under the residual absolute irreducibility clause of the datum (the Galois-family conditions are hypotheses
of the formal statement that its proof does not use), if
$(r_x)$ is good, $\pi\in Z^{\rm good}_i$ and $\pi$ has a semisimple decomposition, then the reductions modulo the maximal
ideal of the matrices $\rho_\pi(g)$, $g\in G_E$, span $M_r(k_O)$ (the residual span clause). Its proof shows that the
reduction of $\rho_\pi$ is continuous (open kernel) and has, at the split places outside $l$, $S_i$, $Q_i$, Frobenius
characteristic polynomials equal to the reductions of the Hecke polynomials of $\pi$ (from the characteristic polynomial
of the decomposition and the Frobenius clauses of the good family), so that the residual irreducibility clause applies.
\item[(a1)] Under the DZ clause of the pending datum, tie0, (BC), (LGC$'$) and (SM1), the type-locus local--global
compatibility (clause (a1) of Proposition~\ref{i2:prop-potcount}(a)) holds.
\end{enumerate}
So the statement above is exact as the conjunction of (i) and (ii); what remains formally verified are the separately
verified statements of (i) and the steps (Res), (a1); their structure is described in \S\ref{i2:sst-pcclocal}.

\begin{proof}[Proof of (HT)]
Let $w\mid l$, so $G_{E_w}=G_{\Q_l}$. The restriction of $\varepsilon$ to $G_{\Q_l}$ is the cyclotomic character of $\Q_l$: both are
defined by the action on $l$-power roots of unity. So $W_c\otimes K\cong K(1-c)$ as a representation of $G_{\Q_l}$, with lattice
$W_c\otimes O$; write $U=K(n)$, $n=1-c$, $T=O(n)$. Restriction of scalars to $\Q_l$ multiplies all dimensions by $[K:\Q_l]$, and
$U$ is crystalline, hence de Rham.

\emph{$2\le c\le t$, so $n\le-1$.} We have $H^0(\Q_l,U)=0$, since $\varepsilon^n$ is nontrivial. Also
$\mathrm{Fil}^0D_{\rm dR}(U)=D_{\rm dR}(U)$, since $D_{\rm dR}(\Q_l(n))$ is spanned by $t^{-n}\otimes1$ and $-n\ge0$. By the dimension
formula of Theorem~\ref{i2:thm-bk}, $H^1_f(\Q_l,U)=0$. So every finite cocycle of $U$ is a coboundary $g\mapsto gu-u$, and its image
in $U/T$ is a coboundary of $U/T$. Hence $H^1_f$ of $U/T$ consists of coboundaries.

\emph{$c=1$, so $U=K$ with trivial action.} The same formula gives $\dim_{\Q_l}H^1_f(\Q_l,K)=[K:\Q_l]=\dim H^1_{\rm ur}(\Q_l,K)$. We also
have $H^1_{\rm ur}\subseteq H^1_f$: an unramified homomorphism $h$ becomes $g\mapsto h(\mathrm{Fr})(g-1)b$ in
$\widehat{\Q_l^{\rm ur}}\subseteq B_{\rm cris}$, where $(\mathrm{Fr}-1)b=1$ is solvable in $W(\bar{\mathbb F}_l)$. So
$H^1_f(\Q_l,K)=H^1_{\rm ur}(\Q_l,K)$; this is Bloch--Kato's Example~3.9 (reference to be confirmed). Coboundaries of the
trivial module vanish, so a finite cocycle of $K$ is a homomorphism trivial on inertia, and so is its image in $K/O$.
\end{proof}

\begin{proof}[Proof of (Deg), (Cont), (Sub)]
(Deg) By a clause of the datum, $g_k$ splits completely in $E$: the number of places above $g_k$ equals $[E:\Q]$; and
$[E:\Q]=4\,\mathrm{ord}=2f$. As $E/\Q$ is Galois, all places above $l$ have the same local degree $[E_w:\Q_l]$ and
$\#\{w\mid l\}\cdot[E_w:\Q_l]=[E:\Q]$; hence $[E_w:\Q_l]=1$. The local field of the branch datum at $w$ is $E_w$, by the
density and valuation clauses of the branch datum, so the formal conclusion (its degree over $\Q_l$ is $1$) follows.

(Cont) The characters $\chi_c:G_E\to O^\times$ are continuous (clause of the data), inversion on $G_E$ and transposition are
continuous, and $r_x$ is continuous (goodness). $\tilde r_x(g)=\chi_1(g^{-1})r_x(g)$ and
$r_x^\vee(g)=\chi_1(g)\chi_1(g)\chi_2(g^{-1})\,{}^tr_x(g^{-1})$ are products of continuous functions.

(Sub) Take $x=\pi\in X_s$. \emph{Continuity of $\rho_\pi$.} For indices $p,q$, the entry $\rho_\pi(g)_{pq}$ is
$L(r_\pi(g)u)$ for the $K$-linear form $L=\mathrm{pr}_p\circ\mathrm{pr}_1\circ e$ and $u=e^{-1}(\epsilon_q,0)$; expanding,
$\rho_\pi(g)_{pq}=\sum_{a,b}\lambda_{ab}r_\pi(g)_{ab}$ with fixed $\lambda_{ab}\in K$. Choose $0\ne d\in O$ with
$d\lambda_{ab}\in O$; then $g\mapsto d\,\rho_\pi(g)_{pq}$ is continuous into $O$, and so is $g\mapsto\rho_\pi(g)_{pq}$, because
multiplication by $d=u\varpi^m$ is a homeomorphism of $O$ onto $\varpi^mO$ for the $\varpi$-adic topology (the topology of $O$ is the $\varpi$-adic one, a clause of the data). Hence $W_0$ and the
$W_c$ (products of continuous characters with $\rho_\pi$ or $1$) are continuous. \emph{Embeddings.} Let
$j_0=e^{-1}\circ(\mathrm{id},0):K^r\to K^{n_V}$ and, for $1\le c\le t$, $j_c=e^{-1}\circ(0,\epsilon_c):K\to K^{n_V}$; both are
injective. By the decomposition, $r_\pi(g)j_0=j_0\rho_\pi(g)$ and $r_\pi(g)j_c=\chi_c(g)j_c$; multiplying by $\chi_1(g)^{-1}$
gives $\tilde r_\pi(g)j_0=j_0W_0(g)$ and $\tilde r_\pi(g)j_c=j_cW_c(g)$.
\end{proof}

\begin{lemma}[$l$-part of the level group]\label{i2:lem-vG}
If $\log g_{k'}\le2\log g_k$, then $v_{g_k}(|G|)\le2fr^2$.
\end{lemma}

\begin{proof}
Write $g=g_k$, $g'=g_{k'}$; they are distinct primes and $g'\le g^2$. $|G|=|\mathrm{GL}_r(\mathbb F_g)|^f|\mathrm{GL}_r(\mathbb F_{g'})|^f$
and $|\mathrm{GL}_r(\mathbb F_q)|=q^{r(r-1)/2}\prod_{a=1}^r(q^a-1)$. Since $g\nmid g^a-1$,
$v_g(|\mathrm{GL}_r(\mathbb F_g)|)=r(r-1)/2$. For $1\le a\le r$, $g^{v_g(g'^a-1)}\le g'^a-1<g^{2a}$, so $v_g(g'^a-1)<2a$ and
$v_g(|\mathrm{GL}_r(\mathbb F_{g'})|)\le\sum_{a=1}^r2a=r(r+1)$. Hence $v_g(|G|)\le f\bigl(r(r-1)/2+r(r+1)\bigr)\le2fr^2$.
\end{proof}

\begin{proof}[Proof of Proposition~\ref{i2:prop-pcc}]
\emph{Per-point bound.} Let $\pi\in Z^{\rm good}_i$ and $X_s=X_s(\pi)$ (Proposition~\ref{i2:prop-percut}). Let $w\mid l$. By
Proposition~\ref{i2:prop-pcclocal} and Proposition~\ref{i2:prop-potcount}(a) the twisted family $(r_x)_{x\in X_s}$ satisfies
(Pot$_N$) at $w$. If $\mathrm{sg}(w)=+$, Lemma~\ref{i2:lem-jumps} (hypothesis (HostJ)) gives (Jump$_{N;\eta+t+1,s+S}$) for the
twisted family, and Lemma~\ref{i2:lem-steplift} (with $\eta'=\eta+t+1<H=s+S$ and $C_T(HN^2+N^3)<g_k$) gives
(Lift$_{e_{\max}}$); if $\mathrm{sg}(w)=-$, the same holds for the dual family, which satisfies (Pot$_N$) by
Lemma~\ref{i2:lem-steplift}. By the semisimple decomposition of $\pi$, the $\rho_\pi$-block (rank $r$) and each character block
(rank $1$) embed into the twisted member at $\pi\in X_s$, so Proposition~\ref{i2:prop-potcount}(b) (sign $+$) or (c) (sign $-$)
bounds the spans of their cut cocycles modulo $H^1_f$ by $e_\rho=2re_{\max}$ and $e_\chi=2e_{\max}$. With $\mathrm{HT}_w$ and the
global clauses (Proposition~\ref{i2:prop-pcclocal}), the finiteness hypotheses, $\pi\in Z_i$ and global Artin reciprocity
(part of $\mathrm{Hypotheses}_2$), Proposition~\ref{i2:prop-stepglobal} and $\#\{w\mid l\}=2f$ give
\[
v^{\rm cut}_i(X_s,\pi)\le2f(2r+2t)e_{\max}+R_i(\pi),\qquad R_i(\pi)=s_i(\pi)+\sum_{0<e<t}\ell_D(l,e)+v_l(h_E)+4n_Vv_\delta+t\sum_{v\in Q_i}v(\zeta_v-1).
\]
\emph{Size of $e_{\max}$.} Put $n=n_V$ and $r'=\lfloor N^2/4\rfloor$. Then $N\le2n$, $Nr'\le N^3/4\le2n^3$,
$N+r'\le2n+n^2\le3n^2$ and $\eta+t+1\le n+\eta+t$, so $e_{\max}\le C_T\bigl(2n^3+3n^2(n+\eta+t)\bigr)\le5C_Tn^2(n+\eta+t)\le e_T$, since
$C=8C_T$. As $r+t=n_V$, $2f(2r+2t)e_{\max}\le4fn_Ve_T$.

\emph{Index bound.} Let $\alpha_\pi,a_\pi$ be as in Lemma~\ref{i2:lem-percutfam} and $\mathcal E=\{\alpha_\pi:\pi\in Z^{\rm good}_i\}$.
Lemma~\ref{i2:lem-famcore} with $Y=Z^{\rm good}_i$, divided by $e_O$, gives
$\eta^{\rm loc}_i\le\sum_\pi\mathrm{mult}_i(\pi)\bigl(v^{\rm cut}_i(X_s(\pi),\pi)+v_{g_k}(|G|)\bigr)$. Insert the per-point bound and
Lemma~\ref{i2:lem-vG}.
\end{proof}

\subsection{Proof of Proposition~\ref{i2:prop-P1}}\label{i2:ssec-proof-P1}

Number the clauses of the division link at $(L,C,C_{\rm abs},C_H)$ as in the formalization: (1) the one-map clause (the
image of each word lattice under its bilinear map lies in the congruence lattice); (2) the frame bound
$s_{\rm Frame}+\mu_{\max}\le x^{C_H}b$; (3) vanishing above the support; (4) the Bessel identity (rationality of the Bessel
data and the thirteenth admissibility clause); (5) the full $RR$-condition; (6) for each word $u$ with nonzero bilinear map,
its height is at most $h(u)$; (7) the displayed inequality; (8) the degree--weight bound with constant $C$; (9) the
comparison of $h(u)+\sum_ta_t\log T_u$ between words with constant $C$.

\noindent\emph{Hermitian lattices and the division data.} Let $K$ be the coefficient field of the division data (a totally complex number
field; $n_K$ is the number of its infinite places $v$, so $n_K=[K:\Q]/2$). A \emph{Hermitian lattice} $A$ is a finite-dimensional $K$-space
$V_A$ with an $O_K$-lattice $L_A$ and, at each infinite place $v$, a positive definite Hermitian majorant
$\mathrm{maj}_{A,v}$ on $V_A\otimes_vC$ (given by a basis and frames); $L_A^\vee=\{f\in V_A^*:f(L_A)\subseteq O_K\}$ is its dual
lattice. For $d\ge1$, an $O_K$-lattice $N\subseteq K^d$, frames $X=(X_v)$ in $M_d(\C)$ and a $K$-isomorphism $\phi:K^d\to V_A$,
the \emph{coordinate tie} $\mathrm{Tie}(A;N,X;\phi)$ means $L_A=\phi(N)$ and
$\mathrm{maj}_{A,v}(\phi x,\phi y)=\mathrm{maj}_{X,v}(x,y)$ for all $v,x,y$, where $\mathrm{maj}_{X,v}$ is the coordinate majorant
with frame $X_v$, $\mathrm{maj}_{X,v}(x,y)=(X_v^{-1}x)^*(X_v^{-1}y)$ (frames are invertible matrices; in the formalization the
inverse of a singular matrix is $0$, and invertibility is then a consequence of a tie, see
Lemma~\ref{i2:lem-tietransfer}). For a $K$-linear $B:V\to W^*$ between spaces with majorants, $\|B\|_v$ is its operator norm at $v$
(majorant on $V$, dual majorant on $W^*$) and $\mathrm{ht}(B)=\frac1{n_K}\sum_v\log\|B\|_v^2$. The division data of a member
(defined in the divisor-bound part of the paper) consist of Hermitian lattices $\mathrm{Nu}(u)$ (one per word $u$) and
$\mathrm{RR}$, $K$-linear maps $B_k(u):V_{\mathrm{Nu}(u)}\to V_{\mathrm{RR}}^*$, the word heights $h(u)$, the dual height
$\mu^\vee$ of $\mathrm{RR}$ (the height of $L_{\mathrm{RR}}^\vee$ with the dual majorants), the congruence lattice
$\mathrm{cc}=L_{\mathrm{RR}}^\vee$, and a goodness clause (the frames of $\mathrm{Nu}(u)$ and $\mathrm{RR}$ are invertible at every
$v$). The height of a Hermitian lattice is defined from the covolume of the lattice for the majorants, so it depends only
on the isometry class.

\begin{definition}[division clauses and division package; stated informally]\label{i2:def-divpkg}
(a) A \emph{column core} $\mathcal P$ is the core of an indefinite package of the column data of $D$, with two constants
$C_p,C_{Hp}$ and a convex face of the column data; it carries a base side and a column side subject to the clauses of the
divisor-bound part of the paper (positivity, frame, level and passage clauses among others), which are named and not
defined here. Of the column side we use: a finite set of members $w$, integers $a_w\ge1$, frames $X^c_w=(X^c_{w,v})$ in
$M_{a_w}(\C)$ and $K$-linear maps $B_w:K^{a_w}\to(K^n)^*$; a lattice $\mathrm{RR}_c\subseteq K^n$ with frames $X^E$; numbers
$\kappa_v>0$ and $\mu_{CC}$, the height of $\mathrm{RR}_c^\vee$ with the dual majorants of the frames $X^E$; and its frame number $F_o(\mathcal P)$ (the weighted sum of the local frame exponents of the
core). Write $\overline{\log\kappa}=\frac1{n_K}\sum_v\log\kappa_v$. For an $O_K$-lattice $N\subseteq K^{a}$ and a function $\sigma$
on the finite places, $N(\sigma)=\mathfrak a_\sigma N$ where $\mathfrak a_\sigma$ is the fractional ideal with exponents $-\lfloor\sigma(w)\rfloor$
(that is, $\mathfrak a_\sigma=\prod_w\mathfrak p_w^{-\lfloor\sigma(w)\rfloor}$).

(b) Given a column core $\mathcal P$, a member $w$, numbers $h_R$ and $\mathrm{rhs}$, a lattice $N\subseteq K^{a_w}$ and $\sigma$,
the \emph{division clauses} are: (D1) $B_w(\phi)\in\mathrm{RR}_c^\vee$ for every $\phi\in N(\sigma)$; (D2) there are $c_v>0$
with $\|B_w\|_v^2\le c_v$ (coordinate majorants with frames $X^c_w$, dual frames of $X^E$) and
$\frac1{n_K}\sum_v\log c_v=h_R-\overline{\log\kappa}$; (D3) $\mu_{CC}-\overline{\log\kappa}+h_R-2F_o(\mathcal P)\le\mathrm{rhs}$;
(D4) the height of $N(\sigma)$ equals that of $N$ plus $2\sum_w\ell_w\lfloor\sigma(w)\rfloor$.

(c) A \emph{division package with constants $(C,C_{\rm abs})$} for a member $D$ with size data $L$ consists of: a column
core $\mathcal P$ (the core of an indefinite package, as in (a), whose member set and rows are tied to the column data of
$D$; these ties also give its well-formedness); for every word $u$ a member $w_u$, an isomorphism
$e_u:K^{a_{w_u}}\to V_{\mathrm{Nu}(u)}$ and a lattice $N^c_{w_u}$ with $\mathrm{Tie}(\mathrm{Nu}(u);N^c_{w_u},X^c_{w_u};e_u)$; an
isomorphism $e_{\rm RR}:K^n\to V_{\mathrm{RR}}$ with $\mathrm{Tie}(\mathrm{RR};\mathrm{RR}_c,X^E;e_{\rm RR})$; the compatibility
$B_k(u)(e_ux)(e_{\rm RR}y)=B_{w_u}(x)(y)$; the tiedness of the column branch data (a clause of the column data); and for
every word $u$: record numbers with $h_R(u)$, $\mathrm{rhs}(u)$, a lattice $N(u)\subseteq K^{a_{w_u}}$, the shift
$\sigma_u$ (the integer-valued function $w\mapsto\lfloor\mathrm{ur}(w,\mathrm{sg}_L(u,w))\rfloor$ given by the size data,
transported to the places of $K$), the input record of the divisor inequality for
$(\mathcal P,h_R(u),\mathrm{rhs}(u),w_u,N(u),\sigma_u)$ (a column chain, height and norm bounds, local bases and error
bounds; named and not defined here), the maps of the column branch data into the ambient spaces of $D$, and the ties
\[
N^c_{w_u}=N(u)(\sigma_u),\quad h_R(u)-\overline{\log\kappa}=h(u),\quad 2F_o(\mathcal P)=F,\quad
\mathrm{rhs}(u)\le\mathrm{MAIN}(u)+C|T|aL_a+Cm(B_0+N_c)x^{C_{\rm abs}\varepsilon}.
\]
\end{definition}

\begin{proposition}[the divisor inequality of a package]\label{i2:prop-wordfacts}
(i) Let $x>x_1$ and $i\in I(x)$, and suppose the member $D=D_{x,i}$ carries a division package with constants
$(C,C_{\rm abs})$ (Definition~\ref{i2:def-divpkg}), satisfies (lq2) and the cut-up condition for every used weight; let
the family field fixing the modified level $N'$ hold at $D$ (it does for $x>x_1$). Then for every word $u$:
(1) $B_k(u)(L_{\mathrm{Nu}(u)})\subseteq\mathrm{cc}$; (6) if $B_k(u)\ne0$ then $\mathrm{ht}(B_k(u))\le h(u)$; and (7)
$\mu^\vee+h(u)-F\le\mathrm{MAIN}(u)+C|T|aL_a+Cm(B_0+N_c)x^{C_{\rm abs}\varepsilon}$.
(ii) Clause (3) holds at every member with $x>x_1$ (formally verified, \cite{CharacterFamily}; a case check of definitions,
\S\ref{i2:sst-div}), and the word heights of every such member are given by the explicit formula $h(u)=h^{\rm expl}(u)$
used in (hh) and (c3) (this is a field of the family, so it needs no proof).
\end{proposition}

\noindent Since Definition~\ref{i2:def-divpkg} is stated informally, so is part (i); the exact statement is in
\cite{CharacterFamily}. Part (i) is exact for the formal package, whose unnamed clauses are those named in the definition. Part (i) rests on the formally verified Lemma~\ref{i2:lem-divclauses} and on Lemma~\ref{i2:lem-tietransfer}, proved
below.

\begin{lemma}[division clauses of a package; formally verified, \cite{CharacterFamily}]\label{i2:lem-divclauses}
Under the hypotheses of Proposition~\ref{i2:prop-wordfacts}(i), for every word $u$ the division clauses (D1)--(D4) hold
for $(\mathcal P,w_u,h_R(u),\mathrm{rhs}(u),N(u),\sigma_u)$.
\end{lemma}
\noindent This is a statement about the formal package, stated informally; the exact statement is in
\cite{CharacterFamily}. It uses the clauses of the indefinite
package and the input record that Definition~\ref{i2:def-divpkg} names without defining.
\noindent Its proof first derives the \emph{value-two input} (every pairing value of $\phi\in N(u)$ with a member vector at
$w$ is $0$ or has $w$-valuation $\ge\sigma_u(w)$) from the tiedness of the branch data, (lq2), the cut-up condition and
the modified-level field. Then (D1) is a local valuation argument, (D4) a covolume computation, (D2) a height chain through
the pairings of the column core, and (D3) the slope (divisor) inequality of the base side carried to the column side by the
passage clause; the method is the Arakelov geometry of Hermitian lattices. Structure in \S\ref{i2:sst-div}.

\begin{lemma}[transport along coordinate ties]\label{i2:lem-tietransfer}
In a division package of a member whose division data are good: $\det X^c_{w_u,v}$ and $\det X^E_v$ are nonzero at every
$v$; $\mathrm{ht}(B_k(u))$ (for the majorants of $\mathrm{Nu}(u)$ and the dual majorants of $\mathrm{RR}$) equals
$\mathrm{ht}(B_{w_u})$ (coordinate majorants); and $\mu_{CC}=\mu^\vee$.
\end{lemma}

\begin{proof}
\emph{Frames.} With frames defined as invertible matrices, the first assertion is part of the definition. With the formal
convention, suppose $X_v$ is singular. Then $\mathrm{maj}_{X,v}$ is identically $0$. But it equals the pull-back of the positive
definite $\mathrm{maj}_{A,v}$ along the isomorphism $\phi$, which is a contradiction.

\emph{Heights of maps.} For $B:V\to W^*$ and majorants on $V,W$, the operator norm for the dual majorant is
$\|B\|_v=\sup_{x,y\ne0}|B(x)(y)|/(\|x\|_v\|y\|_v)$. The ties make $e_u:K^{a_{w_u}}\to V_{\mathrm{Nu}(u)}$ and
$e_{\rm RR}:K^n\to V_{\mathrm{RR}}$ isometries at every $v$, and $B_k(u)(e_ux)(e_{\rm RR}y)=B_{w_u}(x)(y)$. Hence
$\|B_k(u)\|_v=\|B_{w_u}\|_v$ for every $v$, and the heights agree.

\emph{Dual heights.} The transpose of $e_{\rm RR}^{-1}$ maps $(K^n)^*$ onto $V_{\mathrm{RR}}^*$, carries $\mathrm{RR}_c^\vee$ onto
$L_{\mathrm{RR}}^\vee$ (because $L_{\mathrm{RR}}=e_{\rm RR}(\mathrm{RR}_c)$), and is an isometry for the dual majorants (because $e_{\rm RR}$ is
one). Isometric Hermitian lattices have the same height, so $\mu_{CC}=\mu^\vee$.
\end{proof}

\begin{proof}[Proof of Proposition~\ref{i2:prop-wordfacts}(i)]
Fix $u$ and write $w=w_u$, $e=e_u$. By Lemma~\ref{i2:lem-divclauses}, (D1)--(D4) hold.

(1) Let $f\in B_k(u)(L_{\mathrm{Nu}(u)})$. By the lattice part of the tie and $N^c_w=N(u)(\sigma_u)$, $f=B_k(u)(e\phi)$ with
$\phi\in N(u)(\sigma_u)$. For $r\in L_{\mathrm{RR}}=e_{\rm RR}(\mathrm{RR}_c)$, $r=e_{\rm RR}y$ with $y\in\mathrm{RR}_c$, and
$f(r)=B_w(\phi)(y)\in O_K$ by (D1). So $f\in L_{\mathrm{RR}}^\vee=\mathrm{cc}$.

(6) If $B_k(u)\ne0$ then $B_w\ne0$: otherwise $B_k(u)(e x)(e_{\rm RR}y)=0$ for all $x,y$, and $e,e_{\rm RR}$ are onto. A nonzero
map has positive operator norm at every $v$ for positive definite majorants; the coordinate majorants are positive
definite since the frames are invertible (Lemma~\ref{i2:lem-tietransfer}). Hence, with $c_v$ from (D2),
$\log\|B_w\|_v^2\le\log c_v$, and $\mathrm{ht}(B_w)\le\frac1{n_K}\sum_v\log c_v=h_R(u)-\overline{\log\kappa}=h(u)$ by the height
tie. By Lemma~\ref{i2:lem-tietransfer}, $\mathrm{ht}(B_k(u))=\mathrm{ht}(B_w)\le h(u)$.

(7) By Lemma~\ref{i2:lem-tietransfer} and the ties, $\mu^\vee+h(u)-F=\mu_{CC}-\overline{\log\kappa}+h_R(u)-2F_o(\mathcal P)$,
which is $\le\mathrm{rhs}(u)$ by (D3), and $\mathrm{rhs}(u)$ is bounded by the last tie of the package.
\end{proof}

\begin{proof}[Proof of Proposition~\ref{i2:prop-P1}]
Assume $\mathrm{DIVinN}(\mathcal F)$. Take from (pkg) constants $C_1,x_{\rm p}$ (packages with constants $(C_1,C_{\rm abs})$ for
$x\ge x_{\rm p}$), from (c3) constants $C_3,x_3$, from (c2) a constant $C_2$, and thresholds for (lq2), (cut), (hh), (c5),
(bu), (rr), (la). Let $x_0$ be the maximum of all these thresholds, of the threshold of the division kernel data and of
$x_1+1$, and $C=\max(C_1,C_2,C_3,0)$. Let $x\ge x_0$ and $i\in I(x)$. Clauses (1), (6), (7) hold with $C_1$ by
Proposition~\ref{i2:prop-wordfacts}(i); (1) and (6) do not involve $C$, and the right side of (7) is nondecreasing in $C$ since
$|T|aL_a\ge0$ ($a\ge1$ and $L_a\ge0$ by (la)) and $m(B_0+N_c)x^{C_{\rm abs}\varepsilon}\ge0$ (since $B_0\ge0$ and
$N_c\ge0$, clauses of the package). Clause (2) is (hh) with
$h=h^{\rm expl}$ (Proposition~\ref{i2:prop-wordfacts}(ii)); (3) is Proposition~\ref{i2:prop-wordfacts}(ii); (4) is (bu) and
(c5); (5) is (rr); (8) is (c2), monotone in its constant since $B_0\ge0$; (9) is (c3), monotone in its constant since $B_0\ge0$,
$L_a\ge0$ and $a\ge1$, with $h=h^{\rm expl}$.
$\mathrm{LinkEta}(\mathcal F)$ is not used.
\end{proof}

\subsection{Outline of the formally verified steps}\label{i2:ssec-struct}

We outline, for every statement of this section that is formally verified, how its formal proof is organised. Each proof
is split into steps, and each step is classified as (L), (N) or (R) as explained in the introduction to this annex. These
are outlines, not proofs. Every (N) step, and every (L) step whose published form differs from the one used (the
difference is stated each time), rests on \cite{CharacterFamily}. When we say that we do not describe the proof of a step,
the step enters only through its statement.

\medskip\noindent\label{i2:sst-bk}\emph{The Breuil--Kisin steps.}

\emph{Method.} The lattices of Proposition~\ref{i2:prop-liftgood} are crystalline over the tamely ramified field
$K_\varpi=\Q_{l^a}(\varpi_0)$, $\varpi_0^e=-l$, $e=l^a-1$, with $\mathrm{Gal}(K_\varpi/\Q_l)=\Gamma=\mu_e\rtimes\Z/a$. The
method is that of Breuil--Kisin modules with tame descent data over $K_\varpi$. This is the setting of the descent data of
Caraiani--Emerton--Gee--Savitt, \emph{Moduli stacks of two-dimensional Galois representations}, arXiv:1908.07019 (its
section on extensions of rank one Breuil--Kisin modules with descent data; reference to be confirmed). The Breuil--Kisin module is not produced by Kisin's construction. It
is the value of the Bhatt--Scholze prismatic $F$-crystal (Theorem~\ref{i2:thm-bs}) on the Breuil--Kisin prism
$(W(\mathbb F_{l^a})[[u]],(E))$, $u\mapsto[\varpi^\flat]$. Every further structure is then derived from the crystal:
\begin{itemize}
\item the connection;
\item the $B_{\rm cris}$-comparison;
\item the descent datum;
\item the condition on $EN_\nabla$.
\end{itemize}
For comparison, the classical existence and full faithfulness theorem is the following.

\begin{theorem}[Kisin 2006, Lemma~2.1.15 and Prop.~2.1.12]\label{i2:thm-kisin}
Let $K/\Q_p$ be finite, $p>2$, $\pi$ a uniformiser with Eisenstein polynomial $E$ over $W(k)$, $\mathfrak S=W(k)[[u]]$, and
$K_\infty=K(\pi^{1/p^\infty})$. For every $G_K$-stable $\Z_p$-lattice $T$ in a crystalline representation with Hodge--Tate
weights in $[0,h]$ (in Kisin's normalisation) there is a finite free $\mathfrak S$-module $\mathfrak M$ with a Frobenius
whose linearisation has cokernel killed by $E^h$, such that $T|_{G_{K_\infty}}\cong T_{\mathfrak S}(\mathfrak M)$. The
functor $\mathfrak M\mapsto T_{\mathfrak S}(\mathfrak M)$ from such modules to $G_{K_\infty}$-lattices is fully faithful.
\end{theorem}
\noindent Reference: Kisin, \emph{Crystalline representations and $F$-crystals}, in Algebraic Geometry and Number
Theory, Progr.\ Math.\ 253 (2006) (the statement, in particular the $E^h$ bound, is a reference to be confirmed). It is
not used: the formal proofs take Theorem~\ref{i2:thm-bs} instead.

\noindent\emph{Proposition~\ref{i2:bk-b} (packs).} The formal statement agrees with the one above. It also has the
normalisation of the local data, $l\ne2$, $L\supseteq\Q_{l^a}$, and Theorems~\ref{i2:thm-bs} and \ref{i2:thm-fontaine} as
hypotheses.
\begin{enumerate}
\item[(1)] (L) \emph{Module and \'etale comparison.} Theorem~\ref{i2:thm-bs} over $K_\varpi$ gives a prismatic $F$-crystal
$\mathcal E$ whose \'etale realisation is $T|_{G_{K_\varpi}}$. Its value on the Breuil--Kisin prism is finite projective
over the local ring $W(\mathbb F_{l^a})[[u]]$, hence free. A basis gives the matrices $A_\tau$ and the isomorphism $j$. The
$O$-linear structure comes from full faithfulness applied to multiplication by elements of $O$ (R). This step replaces
the existence part of Theorem~\ref{i2:thm-kisin}. Beyond the published statement it uses only the passage to
$O$-coefficients and freeness, both routine.
\item[(2)] (N) \emph{Integral connection.} Evaluate $\mathcal E$ on a crystalline (divided-power) prism with two maps from
$\mathfrak S$. This gives a stratification matrix $\epsilon\equiv1$ modulo the diagonal ideal, compatible with $\varphi$. Its
$n$-th divided coefficient is the matrix of $\nabla^n$ on the $\varphi$-basis, so $\nabla(M\!M)\subseteq M\!M$ with
$S$-coefficients. Moreover $g\in G_{K_\varpi}$ acts on $A_{\rm cris}\otimes M$ by the divided Taylor series
$\sum_k\mu_g^{[k]}\nabla^k$. This step uses $l\ne2$. Rationally, $\nabla$ is Kisin's connection $N_\nabla$ (Kisin 2006,
\S1.2). We do not know a reference for the integral statement with the completed divided-power envelope $S$.
Related work: Gao, \emph{Breuil--Kisin modules and integral $p$-adic Hodge theory} (J.\ Eur.\ Math.\ Soc.\ 25 (2023), 3979--4032), and
Du--Liu, \emph{A prismatic approach to $(\varphi,\hat G)$-modules and $F$-crystals} (J.\ Eur.\ Math.\ Soc.\ 28 (2026), 2581--2636).
\item[(3)] (L) \emph{$B_{\rm cris}$-comparison.} An isomorphism $j_B:B_O^n\to B_O\otimes T$ compatible with $j$ and $\varphi$, and
crystalline values, unique and additive. Both come from the height bounds
$\mu^\beta(A_O\otimes T)\subseteq A_O\text{-span of }j(M)\subseteq\mu^{-\alpha}(A_O\otimes T)$, using that $\mu=[\varepsilon]-1$ is a unit of
$B_{\rm cris}$ (Fontaine, Theorem~\ref{i2:thm-fontaine}). Rationally this is $T\otimes B_{\rm cris}=\mathfrak M\otimes
B_{\rm cris}$ (Kisin 2006, \S2.1).
\item[(4)] (R) \emph{Horizontality.} $G_{K_\varpi}$ fixes $j_B(\xi(d))$ when $\xi(d)\in M\!M_l$. This follows from the Taylor
formula of (2), since $\nabla\xi=0$.
\item[(5)] (R) \emph{Stability.} The $A_O$-lattice is $G_{K_\varpi}$-stable. The Galois group of the Kummer tower
$K_\infty=\Q_l((-l)^{1/l^\infty})$ fixes $j(M[1/E])$ pointwise.
\item[(6)] (L, with (N) form) \emph{Conjugates preserve $M$.} An element $g\in G_{\Q_l}$ over $\gamma\in\Gamma$ lifts to an
automorphism of the prism: $u\mapsto\zeta u$, Frobenius on $W(\mathbb F_{l^a})$. The pull-back of $\mathcal E$ along it
realises the same lattice. Full faithfulness (Theorem~\ref{i2:thm-bs}) gives a map to $\mathcal E$, so $g(j(M))=j(M)$.
\item[(7)] (N in form, L in content) \emph{Descent datum.} $G_{\Q_l}=G_{K_\varpi}\cdot G_{K_\infty}$, and the covering map
$G_{K_\infty}\to\Gamma$ is a homomorphism. Then $\Delta(\gamma)=j^{-1}\circ k\circ j$ on the basis, extended semilinearly,
is a tame descent datum compatible with $G_{\Q_l}$. This is the descent datum of Caraiani--Emerton--Gee--Savitt in the same tame setting,
derived here from full faithfulness rather than from Kisin's functor.
\item[(8)] (N) \emph{The condition $EN_\nabla(M)\subseteq M^{\rm an}$.} First, $E^nEN_\nabla(m)\in M$ for some $n$, by algebra.
The remaining pole at $E$ is removed in the $E$-adic completion, identified with $B_{\rm dR}^+$ by $u\mapsto\varpi_0$. This
uses $\mathfrak M\otimes B_{\rm dR}^+=T\otimes B_{\rm dR}^+$ and the fact that the monodromy image lies in the
$B_{\rm dR}^+$-lattice (from Theorem~\ref{i2:thm-fontaine} and horizontality). On $|u|\le|\varpi_0|$ the only zero of Kisin's
$\lambda=\prod_n\varphi^n(E/E(0))$ is that of $E$. So this is an integral form, on the closed disc, of Kisin's condition that
$\lambda N_\nabla$ preserves $\mathcal M$ (Kisin 2006, \S1.2--1.3). We do not know a reference for the integral form.
\item[(9)] (L) \emph{Jumps and height.} The jumps of the filtration $\mathcal L\cap E^iM$ lie among the Hodge--Tate jumps of
$V$ (comparison with the $B_{\rm dR}$-filtration; Kisin 2006, \S1.2, 2.1, and Theorem~\ref{i2:thm-fontaine}). Jumps in
$[-\eta,H]$ are then equivalent to $E^HM\subseteq\mathcal L$ and $E^\eta\mathcal L\subseteq M$, a definitional
reformulation. There is no loss here.
\item[(10)] (L/R) \emph{Characters.} Kisin's section identifies $D=M/uM$ with $D_{\rm cris}(V|_{K_\varpi})$ inside
$B_{\rm cris}\otimes V$. Inertia elements over $\zeta\in\mu_e$ act on $j_B(\xi(d))$ through $\Delta$. So the polynomial of
Definition~\ref{i2:def-good} kills $D$, and the characters lie in $\Sigma'$.
\item[(11)] (R) Assembly.
\end{enumerate}
\emph{New:} (2), (6)--(8) in the stated integral form with descent data. No explicit constant and no condition on $l$
other than $l\ne2$ occurs in this step.

\noindent\emph{Proposition~\ref{i2:bk-ex} (exactness).} The formal statement is more general than the one above. It holds
for any three $\varphi$-modules with identifications $j_\nu$ and any morphisms $F,G$ of $\varphi$-modules with
$j_1F=(1\otimes f)j_0$, $j_2G=(1\otimes g)j_1$. It needs no packs, no descent data, no height and no $l\ne2$. It uses no
published input.
\begin{enumerate}
\item[(1)] (R) \emph{Splitting.} The exact sequence of lattices splits over $O$. Over $W_O$ this gives matrices $Z,Y^W$ with
$G^WZ=1$, $Y^WF^W=1$, and $\ker G^W=\mathrm{im}\,F^W$.
\item[(2)] (N, short) \emph{Unit minors.} By Cauchy--Binet, $1=\det(G^WZ)$ is a $W_O$-combination of images of maximal
minors of $G_\tau$. Hence some minor $g_\tau\in\mathfrak S_L$ has a coefficient of norm $1$. The adjugate gives
$G_\tau Y_\tau=g_\tau\cdot1$ over $\mathfrak S_L$; likewise $Y'_\tau F_\tau=f_\tau\cdot1$.
\item[(3)] (R) Injectivity of $F$: $f_\tau v_\tau=Y'_\tau F_\tau v_\tau$.
\item[(4)] (R) $GF=0$, since it holds on $W_O$-points and $\mathfrak S_L^{\Z/a}\to W_O$ is injective.
\item[(5)] (N, short) $\ker G\subseteq F(M_0)$, from a $W_O$-preimage and the saturation of $F(M_0)$, which follows from
the left quasi-inverse.
\item[(6)] (N) \emph{Finite length of $M_2/G(M_1)$.} This uses $G_\tau Y_\tau=g_\tau\cdot1$ with $g_\tau$ having a unit coefficient,
and also the $\varphi$-structure. A unit coefficient alone does not suffice: $\mathfrak S/(u-l)\cong O$ has infinite
length. We do not describe the proof of this step here.
\end{enumerate}
\emph{Literature.} Kisin 2006, Lemma~2.1.2 concerns exactness of the realisation in the opposite direction. The
statement above (exactness of the lifted triple up to a cokernel of finite length) is not known to us in the literature.

\noindent\emph{Proposition~\ref{i2:bk-cl} (the class map).} The formal statement is more general: it holds for any
Breuil--Kisin datum of height $(\eta,H)$ with a Galois datum, not necessarily a pack. It uses no published input. The
Artin--Schreier--Witt statement below is proved in the development.
\begin{enumerate}
\item[(1)] (L) \emph{Artin--Schreier--Witt.} $\varphi-1$ is onto $W(C^\flat)\otimes T$, with fixed points $T$. This gives
$m_0$. The statement is standard; reference to be confirmed.
\item[(2)] (N) \emph{The section $\sigma$.} Extend the $\varphi$-module by a trivial rank-one module along $y$
($\varphi(x)=x+y$), with $\Delta$ fixing $x$. Put $t=\varphi(x)-\xi^+(\varphi(x)\bmod u)$. Then $t$ is the image of a
$\Gamma$-invariant $\sigma\in M\!M_l$ with $u\nabla\sigma=\zeta(y)$. Moreover $\varphi\sigma-\sigma-\varphi(y)=\xi(d)$ with $d$
$\Gamma$-fixed.
\item[(3)] (N/L) \emph{Crystalline value $b_0$.} $\mu^{H+\eta}\varphi(m_0)$ lies in the $A_O$-span of the $\varphi(e_i)$; as $\mu$ is a unit of
$B_{\rm cris}$, $b_0=\mu^{-(H+\eta)}j_B(\cdot)$. The exponent $H+\eta$ is internal and does not appear in the statement.
\item[(4)] (R) \emph{(ii).} Let $m_0,m_0'$ be two choices. Then $m_0-m_0'$ is $\varphi$-fixed, so it equals $1\otimes x$ with
$x\in T$. Also $\sigma-\sigma'$ has $u\nabla=0$, so it lies in $\xi(D)$. Write $g=hk$ with $h\in G_{K_\varpi}$ and $k$ in
the Galois group of the Kummer tower. Then $k$ acts on $j_B$ of $\Gamma$-invariant elements through $\Delta$, hence
trivially, and $h$ fixes them by horizontality. So two classes differ by $g\mapsto(g-1)x$.
\item[(5)] (R) \emph{(iii), (iv).} For (iii), scale all data by $l^k$. For (iv), $\varphi(m)-m$ has potential $0$, and potentials are
additive.
\item[(6)] (R and one lemma) \emph{(v).} $(Fm_0,F\sigma,\Phi b_0)$ is a potential of $F(y)$, since $F$ commutes with $\varphi$,
$\nabla$ and $\zeta$. The lemma used is the naturality of $j_B$. It is proved on $1\otimes x$
through crystalline values and extended $B$-linearly.
\end{enumerate}
\emph{Relation to Bloch--Kato.} If $\kappa$ is a class of $y$, then $1\otimes\kappa(g)=(g-1)b$ with
$b\in B_O\otimes T$. So $\kappa$ lies in $\ker\bigl(H^1(V)\to H^1(B_{\rm cris}\otimes V)\bigr)=H^1_f(V)$, which is Bloch--Kato's
definition (reference to be confirmed). The statement (i) gives only a potential. That the coboundary of the potential takes
values in $T[1/l]$ is supplied by Proposition~\ref{i2:bk-int}, under $\mathrm{Mon}(y)=0$.

\noindent\emph{Proposition~\ref{i2:bk-mon} (monodromy-kernel representatives).} The formal statement agrees with the
one above. A finite cocycle is a continuous $T$-valued cocycle whose image in $V$ is of the form $g\mapsto(g-1)x$,
$x\in B_{\rm cris}\otimes V$, over all of $G_{\Q_l}$.
\begin{enumerate}
\item[(1)] (L, with a routine (N) remark) \emph{Extension lattice.} Let $T_c=T\oplus O$ with
$h(t,x)=(\rho(h)t+xc(h),x)$. It is crystalline over $K_\varpi$ because $c$ is finite: classes in $H^1_f$ are crystalline
extensions (Bloch--Kato \S3). Its jumps lie in $U$ because $0\in U$. Its type is $\Sigma'$: $c$ is finite over all
of $G_{\Q_l}$, so $1\in D(O)$ has a $G_{\Q_l}$-fixed lift in $D_{\rm cris}(V_c|_{K_\varpi})$. Hence $\prod_{c\in\Sigma'}(g-\zeta^c)$
still kills $D_{\rm cris}(V_c|_{K_\varpi})$, using $0\in\Sigma'$.
\item[(2)] (cited formal) Packs for $T_c$ and $O$ (Proposition~\ref{i2:bk-b}), the lifts $F:P_T\to P_c$, $G:P_c\to P_O$
(Proposition~\ref{i2:bk-c}), and exactness (Proposition~\ref{i2:bk-ex}).
\item[(3)] (N) \emph{Cokernel bound.} Let $F:P_1\to P_2$ be a morphism of Breuil--Kisin data with characters in $\Sigma'$
and cokernel on $M$ of finite length. Then $l^{e_W}$ kills the cokernels of $F$ on $\mathcal L$, $M$, $E^\eta M$,
$M^{\rm inv}$, $Y$ and $Y'$. The idea:
\begin{itemize}
\item Finite length gives $u^{eA'}x=F(y)+l^bz$ with $A'=H+\eta+b$ and $b=e_W+1$.
\item Project to a $\mu_e$-isotypic part. By the descent datum, the $u$-adic degrees occurring there lie in classes
$c-s$ modulo $e$, $s\in\Sigma'$.
\item Modulo high degree, $N_\nabla=u\,d/du$ acts on $u^i\mathcal L/u^{i+1}\mathcal L$ by $i$. Products $P_1(N_\nabla)$,
$P_2(N_\nabla)$ over the degrees below $eA'$ and in $[eA',2eA')$ push isotypic vectors to higher order.
\item \emph{Lagrange step:} $l^{e_W}=\alpha P_1+\beta P_2$ with $\alpha,\beta\in\Z_l[X]$. This follows from a digit lemma: if
$0<c'<e$, $e\mid c'(l^{\rho'}-1)$ and $|k|<l$, then $v_l(c'+ke)\le\rho'$. Take $\rho'=\mathrm{per}(s-s')\le r$. Then for a
fixed root the sum of $v_l(i-j)$ over the other roots is at most $(|\Sigma'|-1)r\le(N-1)r$.
\item Iterate in $l$.
\end{itemize}
The Lagrange step needs only $(N-1)r$. The factor $2$ in $e_W=2(N-1)r$ is not used there.
\item[(4)] (L/R) The trivial pack has a $\varphi$-fixed $u\in M_O^{\rm inv}$ realising $1\in O$.
\item[(5)] (R) By (3) on $M^{\rm inv}$ there is $x_2\in M_c^{\rm inv}$ with $G(x_2)=l^{e_W}u$. Then $G(\varphi x_2-x_2)=0$. Exactness
gives $y\in Y_T$ with $F(y)=\varphi x_2-x_2$, and $\mathrm{Mon}(y)=0$ because $\mathrm{Mon}$ vanishes on $\varphi(m)-m$ and $F$
is injective on $\mathrm{Pol}$.
\item[(6)] (R) A horizontal section $\sigma$ with $F(\sigma)=\varphi x_2-\xi(\varphi x_2\bmod u)$, $\Gamma$-invariant by injectivity
of $F$.
\item[(7)] (N, bookkeeping with $B_{\rm cris}$) \emph{The class.} Solve $\varphi m_0-m_0=y$ with $F(m_0)=x_2-l^{e_W}v$, where
$v=(0,1)\in T_c$. Then the coboundary of the potential $j_B(\sigma)-b_0$ maps under $F$ to $l^{e_W}(g-1)v=l^{e_W}c(g)$.
This uses the naturality of the $B_{\rm cris}$-comparison and the injectivity of $1\otimes f$.
\end{enumerate}
\emph{Literature.} Fontaine--Laffaille would give loss $0$ over an unramified base with weights in $[0,l-2]$.
Here the base $K_\varpi$ has ramification index $e=l^a-1\ge l$, so it does not apply; the loss $l^{e_W}$ is the price of
the ramification. The closest published analogues are cokernel bounds for Kisin modules of exact sequences (Liu, Ann.\
Sci.\ ENS 2007; Caruso--Liu; reference to be confirmed). We do not describe here the reduction of $N_\nabla$ modulo high degree,
the passage of (3) to all six modules, or the proofs of the crystallinity and jump statements of (1).

\noindent\emph{Proposition~\ref{i2:bk-tri} (lifting through exact triples).} The formal statement agrees with the one
above; it also assumes $L$ complete. It is pure linear algebra on Breuil--Kisin data with tame descent; no Galois input.
The type of $P_0$ is not assumed: its jumps lie among those of $P_1$.
\begin{enumerate}
\item[(1)] (N, step (3) of Proposition~\ref{i2:bk-mon}) Lift $y_2'$, $m_2$ to $y_1\in Y_1'$, $m_1\in M_1^{\rm inv}$ with loss
$l^{e_W}$.
\item[(2)] (N) \emph{Defect.} $\omega=\mathrm{Mon}(y_1+\varphi m_1)$ satisfies $l^{c_1}\omega\in Z_1$, $c_1=\eta(r+1)$. The main input:
if every character has Frobenius period $\le r$ and $(r+2)\eta<l+1$, then $l^{\eta(r+1)}\zeta(\varphi m)\in M\!M$ for $m\in M^{\rm inv}$.
Inertia invariance imposes $u$-adic order conditions on each isotypic piece. These improve along Frobenius orbits of
length $\le r$, at cost $l^\eta$ per application of $\varphi$ (the height). It also uses the fine $\zeta$-bound and the
integrality computation (Lemmas~\ref{i2:lem-zint}, \ref{i2:lem-zfine}). Here $r$ enters through $\varphi$-orbits. We do not
describe the proof of the orbit bound here.
\item[(3)] (N/R) \emph{Pull-back.} $G(\omega)=0$. Exactness on $Z$, with cokernel killed by $l^{e_W}$, gives $w_0\in Z_0$ with
$F(w_0)=l^{c_2}\omega$, $c_2=c_1+e_W$.
\item[(4)] (N) \emph{Semisimplicity.} On $Z_l$, $\mathrm{Lc}$ is semisimple with eigenvalues in $\{nel:n\in U+[1,\eta]\}$, and
$\omega\in\mathrm{Lc}(Z_l)$. With $R(X)=\prod_{n\ne0}(X-nel)$, $R$ kills the range of $\mathrm{Lc}$, and $R(0)w=-\mathrm{Lc}(R_1(\mathrm{Lc})w)$
there. This gives $w'\in Z_{0,l}$ with $F(\mathrm{Lc}\,w')=\omega$.
\item[(5)] (N) \emph{Integrality at cost $l^{\eta+e_L}$.} Each nonzero eigenvalue $nel$ has $v_l=1$, since $|n|<l$ and
$e\equiv-1\pmod l$. So $v_l(R(0))\le|U|\eta\le e_L=(N+1)\eta$. The extra $\eta$ comes from $l^\eta Y\subseteq Y'+M^{\rm inv}$. This
gives $y\in Y_0'$ with $\mathrm{Lc}(\bar y)=l^{c_2+\eta+e_L}\mathrm{Lc}\,w'$.
\item[(6)] (N) \emph{Realisation.} There is $y_0\in Y_0'$ with $\mathrm{Mon}(y_0)=\mathrm{Lc}(\bar y)$. This uses $l$-adic completeness of
$Y'$, $14H+(N+2)\eta<l$ and $10H+2\le l$.
\item[(7)] (R) $y'=l^ny_1-F(y_0)$ and $m=l^nm_1$, $n=c_2+\eta+e_L$; then $\mathrm{Mon}=0$ and $G(y'+\varphi m)=l^k(y_2'+\varphi m_2)$.
\end{enumerate}
\emph{Literature.} The semisimplicity in (4) is an integral, tame-descent analogue of the action of $N_\nabla$ on
$D\otimes\mathcal O$ in Kisin 2006, \S1.2--1.3. We do not know a published statement with these exponents. Tame
descent data on Breuil--Kisin modules: Savitt (Duke Math.\ J.\ 128 (2005), 141--197); Caraiani--Emerton--Gee--Savitt (arXiv:1908.07019).

\noindent\emph{Proposition~\ref{i2:bk-int} (integral classes).} The formal statement needs only a Breuil--Kisin datum of height $(\eta,H)$ with a Galois
datum: no type, no numerical conditions other than $\eta<H$ and $8H+1<l$, and no pack; the statement above is given in
this form. It assumes
Theorem~\ref{i2:thm-fontaine}, which is proved in the development.
\begin{enumerate}
\item[(1)] (cited formal) $y$ has a potential $b$ (Proposition~\ref{i2:bk-cl}(i)).
\item[(2)] (L) $G_{\Q_l}=G_{K_\varpi}\cdot G_{K_\infty}$ with $K_\infty$ the Kummer tower of $-l$. This holds because $K_\varpi$ and
$K_\infty$ are linearly disjoint and $K_\infty$ covers $\Gamma$ (Kummer theory and tame ramification; standard).
\item[(3)] (N) \emph{Kummer part.} For $k\in G_{K_\infty}$, $(k-1)b\in1\otimes T$, because $k$ acts on the image of $M$ through
$\Delta$. This is the analogue of the fact, used by Kisin and Liu, that $G_{K_\infty}$ fixes $\mathfrak M$.
\item[(4)] (L) \emph{Qualitative integrality on $G_{K_\varpi}$.} For $h\in G_{K_\varpi}$, $(h-1)b$ is $\varphi$-fixed. It lies in
$\mathrm{Fil}^0B_{\rm dR}\otimes T$ because $\mathrm{Mon}(y)=0$: this is exactly the condition that the horizontal section
gives a $\mathrm{Fil}^0$ period. Fontaine's $B_{\rm cris}^{\varphi=1}\cap\mathrm{Fil}^0B_{\rm dR}=\Q_l$
(Theorem~\ref{i2:thm-fontaine}) puts it in $\Q_l\otimes T$, i.e.\ $l^n(h-1)b\in T$ for some $n$.
\item[(5)] (N) \emph{Explicit exponent.} Let $\mu=[\varepsilon]-1$.
\begin{itemize}
\item[(a)] $l^c(h-1)j_B(\sigma)$ lies in the $A_{\rm cris}$-image, from $l^c\zeta(y)\in M\!M$.
\item[(b)] $\mu^{H+\eta}$ times the $A_{\rm cris}$-lattice lies in $A_{\rm cris}\otimes T$.
\item[(c)] $\mu^{H+2\eta}b_0$ lies in the $A_{\rm inf}$-lattice.
\item[(d)] \emph{Valuation lemma:} if $l^n\mid\mu^jx$ in $A_{\rm cris}$ with $0\ne x\in\Z_l$, then $(l-1)n\le(l-1)v_l(x)+2j$. With
$j=H+2\eta$ and $2j<l-1$, an element $u$ with $l^nu\in T$ and $\mu^jl^cu\in A_{\rm cris}\otimes T$ has $l^cu\in T$.
\end{itemize}
The inequality $2(H+2\eta)<l-1$ follows from $\eta<H$ and $8H+1<l$. This is the only use of $8H+1<l$ besides $l\ne2$.
\item[(6)] (R) $X(g)=l^c(g-1)b\in T$ combines over $g=hk$. It is a finite cocycle because
$B_O\otimes T\cong B_{\rm cris}\otimes V$, and it is a class of $l^cy$ by Proposition~\ref{i2:bk-cl}(iii).
\end{enumerate}
\emph{Literature.} The condition $8H+1<l$ is not a Fontaine--Laffaille range; it is used only through $2(H+2\eta)<l-1$.
Bloch--Kato (\S4) use Fontaine--Laffaille to compute $H^1_f$ of lattices over an unramified base with weights in
$[0,l-2]$. That does not apply over $K_\varpi$. The $h=hk$ decomposition corresponds to Liu's
$(\varphi,\hat G)$-modules (Liu, \emph{A note on lattices in semi-stable representations}, Math.\ Ann.\ 2010; reference to
be confirmed). The valuation lemma (d) is a divided-power estimate of Breuil's type (reference to be confirmed). We do not
describe here the proofs of (5)(a), (5)(d) and (3).

\noindent\emph{Fact (T4) of Proposition~\ref{i2:prop-goodlat}.} Fact (T4) is proved in the text from Fontaine's theory. The formal proof
follows the standard route, and every step is (L):
\begin{enumerate}
\item[(1)] restrict to $K''=K'F$;
\item[(2)] $D_{\rm cris}(V|_{K''})$ is contained in the $K''_0$-span of $D_{\rm cris}(V|_{K'})$, by regularity of $B_{\rm cris}$;
\item[(3)] inertia fixes $\Q_l^{\rm ur}$;
\item[(4)] an element of $G_F$ fixing $K''_0$ is a product of an element of $I_l\cap G_F$ and an element of $G_{K''}$, which
uses that $K'$ is Galois;
\item[(5)] Speiser's lemma descends the semilinear action (reference to be confirmed);
\item[(6)] a dimension count.
\end{enumerate}
The formal statement assumes $K'/\Q_l$ Galois; so does (T4).

\medskip\noindent\label{i2:sst-large}\emph{Large primes.}

The proposition is (R) given its hypothesis (CB): instantiate (CB) and identify terms.
\begin{itemize}
\item The weight tie comes from the weight parametrisation clause, and the Hecke--Galois condition from
$\pi\in Z^{\rm good}_i$.
\item Sha-finiteness and \eqref{i2:eq-sha} turn the Sha lengths into $\ell_D(l,e)$.
\item $\mathrm{KillAt}$ gives finiteness of $s_i(\pi)$ and $t\sum_{v\in Q_i}v(\zeta_v-1)\le l^{-2}$.
\item Pass from $[0,\infty]$ to $\R$.
\end{itemize}
The derivation of the weight tie is a formal statement about clauses not defined here. The content is the proof of (CB)
along a family, for $x$ beyond a threshold. It uses Theorems~\ref{i2:thm-tate}, \ref{i2:thm-euler}, \ref{i2:thm-artin},
\ref{i2:thm-bs}, \ref{i2:thm-fontaine} as arguments of the formal statement; the development proves each of them as a
separate proposition. The Bloch--Kato and local-finiteness components of
$\mathrm{Hypotheses}_2$ are passed to it but not used.
\begin{enumerate}
\item[(1)] (N, elementary) \emph{Position data off $\mathrm{Exc}(L)$.} For $x\ge x_0$ and prime $l>l^*$: $l\ge5$, $t+1<l$,
$l\ne q,d$, and $l\notin\mathrm{small}$. If moreover $l\notin\mathrm{Exc}(L)$, then for every $n$, sign function and type
function there are position data with the following properties:
\begin{itemize}
\item all weights lie in $[1,h-1]$;
\item $\eta+4s+6<l$;
\item the separation is $>s+2$;
\item no integer in the windows of Definition~\ref{i2:def-budgets} is divisible by $l$.
\end{itemize}
The last property is where $\mathrm{Exc}(L)$ enters.
\item[(2)] (R) $l$ is unramified in $E$, and the cyclotomic inertia image is as required.
\item[(3)] (R/L) \emph{Global image and ramification conditions} for $r$ and the rank-$r$ part $\rho$:
\begin{itemize}
\item continuity and semisimplicity;
\item the conditions at $S$ and $Q$: a quadratic relation with $\zeta_v-1$ a unit, or a polynomial condition;
\item a relation at $Q$ with $\zeta_v\ne1$.
\end{itemize}
Each is a short verification.
\item[(4)] (N) \emph{Idempotents.} An orthogonal idempotent decomposition with unit root differences. The cyclotomic inertia
generator separates $1,\varepsilon^{-1},\dots,\varepsilon^{-(t-1)}$ modulo $l$, since $t+1<l$ and $l$ is unramified.
\item[(5)] (L) Wach and Fontaine functors from Theorem~\ref{i2:thm-bs}, as in Proposition~\ref{i2:prop-ptsel}.
\item[(6)] (N; the large part) \emph{Torsion comparison.} Let $K/\Q_l$ be unramified, $l\ge5$, with position hypotheses. Let
$\rho$ be crystalline with prescribed Hodge--Tate multiplicities, and $\chi_c=\chi_1\varepsilon^{-(c-1)}$ crystalline with one jump.
Then, for every finite $O$-module $X$, a $\chi_1$-twisted cocycle is \emph{admissible} if and only if:
\begin{itemize}
\item for the $\rho$-block, it is finite plus a coboundary;
\item for $c\ge2$, it is a coboundary;
\item for $c=1$, it is unramified.
\end{itemize}
Here admissible means that its extension is a subquotient of lattices in sums of crystalline representations with the
prescribed weights and residual characteristic polynomial. The proof assembles 24 clauses:
\begin{itemize}
\item Wach lattices (Theorem~\ref{i2:thm-wach}, here derived from Theorem~\ref{i2:thm-bs});
\item Herr's complex (Theorem~\ref{i2:thm-herr} and Corollary~\ref{i2:cor-herr}) with the field-of-norms comparison;
\item Theorems~\ref{i2:thm-tate}, \ref{i2:thm-euler};
\item the Bloch--Kato dimension formula over an unramified base (proved in the development);
\item crystalline characters with given jumps (Theorem~\ref{i2:thm-fontaine}, Exp.~III \S5.1);
\item clauses on Sen lattices and $A_{\rm inf}$.
\end{itemize}
We do not give a written proof of this step. \emph{Literature:} Fontaine--Laffaille together with Bloch--Kato \S4 gives a statement of the
same kind: over an unramified base, in the Fontaine--Laffaille range, torsion crystalline classes are described by
Fontaine--Laffaille modules. We have not identified ``admissible'' with the Fontaine--Laffaille category,
so we do not substitute a citation.
\item[(7)] (L) Crystalline characters with prescribed jumps over unramified $K$ (Lubin--Tate products;
Theorem~\ref{i2:thm-fontaine}).
\item[(8)] (N) \emph{Local conditions above $l$.} At each $w\mid l$:
\begin{itemize}
\item cut cocycles of the $\rho_\pi$-block are finite;
\item those of $\chi_c$, $c\ge2$, are coboundaries;
\item those of $\chi_1$ are unramified.
\end{itemize}
The steps:
\begin{itemize}
\item extend coefficients;
\item cut cocycles are torsion;
\item lift to $T\otimes O_{L'}/\lambda^a$;
\item check admissibility;
\item apply (6);
\item descend coefficients.
\end{itemize}
\item[(9)] (N + L) \emph{Global count.} The determinant bound $\ell(O/\mathfrak c_\pi)\le\ell(C_\rho)+\sum_c\ell(C_c)$
(the argument of Proposition~\ref{i2:prop-cutcoc}(a)). Then, block by block:
\begin{itemize}
\item $\rho$-block: local terms vanish by (8), coboundary terms vanish since $\delta$ is a unit, and the Selmer bound
gives $e_Os_i(\pi)$;
\item character blocks: peeling $Q_i$ costs $\sum\ell(O/(\zeta_v-1))$;
\item trivial block: Artin reciprocity gives $v_l(h_E)$ (Lemma~\ref{i2:lem-unrhom});
\item blocks $c\ge2$: the $\mathrm{Sha}$ bound (Proposition~\ref{i2:prop-cutcoc}(h)).
\end{itemize}
This is the argument of \S\ref{i2:ssec-proof-stepglobal} with all local terms $0$.
\item[(10)] (R) Pass to $[0,\infty]$. \quad (11) (R) Thresholds.
\end{enumerate}
\emph{New:} (1), (4), (6), (8) in the present form, uniformly for all $l$ outside $\mathrm{Exc}(L)$ above $l^*$.
We do not describe here the proofs of the 24 clauses of (6) or the proofs in (3), (4).

\medskip\noindent\label{i2:sst-locadm}\emph{Local admissibility data.}

Proposition~\ref{i2:prop-locadm} is not one formal statement. Each item has its own source.
\begin{enumerate}
\item[(1)] \emph{Item (1).}
\begin{itemize}
\item At inert $l$, the admissibility conditions come from the word conditions of the hub data and the Wach coordinates.
\item At split $l$ they come from a declared decomposition. This uses Chebotarev in the degree-one form, local finiteness
and Theorem~\ref{i2:thm-fontaine}.
\item The invariant bounds with $c=120$ come from the setup.
\item The table conditions (A), (B), (Med), (SmI), (Band) come from the canonicity of the hub data. Their row constants are
$(C_2,320,25,23,32,32)$, and they hold for every constant $C\ge\max(3C_2+1,22365)$. (Band) uses in addition (Par).
\end{itemize}
The first three are (N), and we do not describe their proofs here; the table conditions are (R): elementary comparisons.
\item[(2)] \emph{Item (2).} The inert-type admissibility with all places as rows, $e'=0$, and
$c_B\le60\,tf(r+t)\log b$. This is (N); we do not describe its proof here. The constant $K=160$ is the constant of the first
conclusion of the element-denominator clause.
\item[(3)] \emph{Item (3).} Constants $1000$, $6000$, $30000$; we do not describe the proof here.
\end{enumerate}
The statement cannot be made exact within a reasonable length, for the following reasons.
\begin{itemize}
\item At inert $l$ (and at $q$, $d$) the admissibility predicate has exactly twenty clauses. About half are one line each:
(G1), (G3), (G4), (G6), continuity, semisimplicity, the cyclotomic clauses, the unramified class-field bound. Three are a
paragraph each: (G2), (G5), the tame-generator conditions. The block clause, the row-length clause, the lift clause, the
Bloch--Kato and $H^2$ clauses, Fontaine's comparison and the local-coefficient compatibility need about three to four pages
together.
\item At split $l$ the formal predicate is a different one, with 38 clauses. It feeds a split-type bound rather than
Proposition~\ref{i2:prop-selmer}. We do not describe its 18 additional clauses here.
\end{itemize}
These two predicates are named and not defined. The following notions could each be defined exactly in under two pages
(in the formal development they are canonicity data with nine fields, word data, and row and column choices):
\begin{itemize}
\item the canonical hub data, their row constants $(C_2,320,25,23,32,32)$ and the row budgets;
\item the joint parameter conditions (word data with base, box and band clauses; about two pages);
\item the declared rows and columns (half a page);
\item $h^0_u$ (normalised length of the $\rho_\pi$-invariants of the decomposition group at $u$) and $h_0$ (the same for
the twisted lines $\varepsilon^{-e}$, $0<e<t$).
\end{itemize}
The hub size is $h_B=b+S_c+\eta_{\rm sz}+t+2$ (Definition~\ref{i2:def-budgets}). We do not describe here the local level $s_l$
or the abelian size $\Lambda_l$.

\medskip\noindent\label{i2:sst-ptsel}\emph{Point families, Selmer signatures and thresholds.}

The formal statement is an assembly. Only the Bhatt--Scholze component of $\mathrm{Hypotheses}_2$ is used. Chebotarev's
theorem, in the degree-one form, is proved in the development (see (PF1)), so it is not a hypothesis. There is no class
field existence theorem, no Grunwald--Wang, no Greenberg--Wiles and no explicit Chebotarev bound in this proposition. The
twisting characters with prescribed local behaviour are part of the hypothesis $\mathrm{TwGood}$, not constructed here.

\emph{(TH).}
\begin{itemize}
\item[(TH1)] (R, elementary asymptotics) For $x\ge x_0(\varepsilon,C,A)$:
\begin{itemize}
\item $100\le w''(n)$ and $100(\eta_{\rm sz}+S_c+t+2)\le w''(n)$;
\item $w''\le b'$ and $b_T+b'/x\le\frac32b'$;
\item $C(L_w+1)b'\le b$.
\end{itemize}
These follow from $b/b'=e^{x^{\varepsilon/3}}/x$ against polynomial growth of the other parameters.
\item[(TH2)] (R) This implies admissibility of the parameter choice and (Par), by linear arithmetic.
\item[(TH3)] (R) Special primes are small, $\mathsf L\ge8$, the split bound is $\ge60$, and $C\le b$ along the family.
\item[(TH4)] (R) The threshold $x_c$ is the witness in $\mathrm{TwGood}$, which gives (GT) with $\varepsilon'=l^{-2}$. The tie of the
twist datum to the norm choice, and the identification of its auxiliary places with $Q_i$, are \emph{not} part of
$\mathrm{TwGood}$. They are a field of the norm-choice facts (the norm-choice tie) and hold for all $x>x_1$; (TH) is stated in this form.
\end{itemize}

\emph{(PF).}
\begin{itemize}
\item[(PF1)] (L) \emph{Chebotarev, degree-one form.} Let $U\subseteq G_F$ be open normal, $g\in G_F$, and $A$ finite. Then
there is a degree-one prime $v\notin A$, unramified over $\Q$, with inertia in $U$ and Frobenius in $gU$. The proof
reduces to cyclotomic fields: take a prime $m$ with $[L:F]\mid m-1$ and cross with $F(\zeta_m)$. It then uses infinitely
many degree-one primes in a class of $(\Z/m)^\times$, by $L(1,\chi\circ N)\ne0$. This is a form of
Theorem~\ref{i2:thm-cheb}. For the degree-one, inertia-in-$U$ form: Lenstra--Stevenhagen, \emph{Chebotarev and
his density theorem}, Math.\ Intelligencer 18 (1996), no.~2, 26--37 (background; reference to be confirmed for the
degree-one form as stated).
\item[(PF2)] (L, a clause of the family facts) Cuspidal, conjugate self-dual, regular algebraic $\pi$ have $l$-adic Galois
representations (reference to be confirmed; Harris--Taylor, Shin, Chenevier--Harris, BLGGT, Caraiani). The representations
$r_x$ themselves are supplied by the family facts.
\item[(PF3)] (N, routine) \emph{Trace identity.} $\mathrm{tr}\,r=T+\sum_c\bar\chi_c$ on Frobenii at almost all split places, by
matching Hecke eigenvalues across the isobaric parts. It then holds on $G_E$ by (PF1) and continuity.
\item[(PF4)] (L, linear algebra) Semisimplicity and the trace identity give eigenvectors for the $\chi_c$ and a
continuous complement $\rho$ over $O$.
\item[(PF5)] (R) Assembly by choice.
\end{itemize}
The resulting point family does not contain the good-row property, which is derived separately from the family facts.

\emph{(SEL).} The formal statement has, besides the norm-choice tie, the hypotheses: the triple condition, the good-row
property, residual absolute irreducibility and the same-twist condition. All four are supplied by the family facts.
\begin{itemize}
\item[(SEL1)] (R) A row, parameters and a good triple $s$.
\item[(SEL2)] (L) $\mathrm{tr}\,\rho_\pi=\mathrm{tr}(\text{twist of }\rho_s)$ on $G_E$, by (PF1) and continuity.
\item[(SEL3)] (L) \emph{Carayol--Serre.} Two representations over a complete discrete valuation ring with equal traces and
absolutely irreducible residual representation are conjugate. The proof: $O[G]\to M_a(O)$ is onto, then Nakayama.
Reference: Carayol, \emph{Formes modulaires et repr\'esentations galoisiennes \`a valeurs dans un anneau local complet}
Contemp.\ Math.\ 165 (1994), 213--237, Thm.~1; see also Serre, \emph{Repr\'esentations lin\'eaires sur des anneaux locaux, d'apr\`es Carayol} (1993), Th\'eor\`eme~1 (for any local ring $A$).
\item[(SEL4)] (L/R) Conjugate lattices have equal Selmer lengths. Bloch--Kato finite conditions are compatible with
injective coefficient extension and with conjugating the $l$-adic embedding (Theorem~\ref{i2:thm-bk}; flat base change of
$H^1_f$: reference to be confirmed).
\item[(SEL5)] (R) Flat base change to a common discrete valuation ring.
\item[(SEL6)] (N, short) If the twist is trivial on the decomposition groups above $l$, then the Selmer length is at most
$\mathrm{BK}(\sigma)$, by an explicit lift of cocycles comparing local conditions.
\item[(SEL7)] (R) Packaging.
\end{itemize}

\emph{(CR).}
\begin{itemize}
\item[(CR1)] (R) $\zeta_v$ is the value at $\gamma_v$ of the first twist factor.
\item[(CR2)] (L) The twist has finite $l$-power order on $I_v$, $v\nmid l$, so it kills wild inertia. Its image is generated by
the image of $\gamma_v$, since tame inertia is procyclic (Serre, \emph{Local Fields}, Ch.~IV, \S2, Prop.~7, Cor.~1, at each finite level, then the inverse limit).
\item[(CR3)] (R) Normalised valuations are independent of the discrete valuation ring inducing the same place.
\item[(CR4)] (N, elementary) $(u-1)\mid(u^m-1)$, so $v(\psi(\gamma_v)-1)\le v(\psi(\sigma)-1)$ whenever $\psi(\sigma)\ne1$.
\end{itemize}

\emph{Wach and Fontaine.}
\begin{itemize}
\item[(W1)] (L) Theorem~\ref{i2:thm-bs}.
\item[(W2)] (R) Coefficients: $O_L$ to $\Z_l$.
\item[(W3)] (N) Wach lattices at all layers, by evaluating the $F$-crystal on a constructed prism. This replaces
Theorem~\ref{i2:thm-wach}, which is not used.
\item[(W4)] (R) The functors.
\end{itemize}

The statement can be made exact only in part. Exact definitions of a few lines each exist for the following notions:
\begin{itemize}
\item the admissible parameter choice: $\eta_{\rm sz}+t+S_c<w''(n)$, $3w''(n)\le b$, $b_T+b'/x\le b$, $t\le\eta_{\rm sz}$;
\item (Par);
\item the scale ($w''(n)=b'/(xn)$, $b'=a/X'$, $X'=e^{x^{\varepsilon/3}}$);
\item small primes;
\item the point family;
\item the tame value $\mathrm{tv}(v)=\inf\{v(\psi(\sigma)-1)/v(l):\sigma\in I_v,\ \psi(\sigma)\ne1\}$;
\item the clauses of (GT) used.
\end{itemize}
The following notions cannot be defined within a reasonable length, and are therefore named and not defined:
\begin{itemize}
\item $\mathrm{TwGood}$, whose good-twist clause uses the theta and Hecke-module data of the twist (some 300 fields);
\item the family facts (core norm facts, norm-choice facts, extra clauses);
\item the canonicity of hub data;
\item the jump admissibility;
\item the good signature (which needs the Hecke module of the twist datum);
\item the hypotheses of (SEL) listed above.
\end{itemize}
We do not describe here the proofs of the asymptotic lemmas, of flat base change of Selmer groups, of the comparison of
Bloch--Kato local conditions, or of $\mathrm{KillAt}$ from (SC), (GT), nor the definitions of the local Bloch--Kato condition
at $\mathrm{Spl}\cup Q$ and of $\nu$.

\medskip\noindent\label{i2:sst-colconds}\emph{The column module.}

There is no single formal statement equal to Proposition~\ref{i2:prop-colconds}. The proposition is the conjunction of the
roughly twenty inputs of the formal Selmer inequality, discharged by one closed theorem from $\mathrm{Adm}_i$. The formal
version works with the normalised family, conjugated to standard block form. The objects named in the statement have the
following exact meanings.
\begin{itemize}
\item \emph{Decomposition:} a partition of the places of $E$ above $l$ into rows and columns.
\item \emph{Column subgroups:} $D_l$ is the decomposition group $D(A_w)\subseteq G_E$ of a fixed place $A_w$ of
$\Omega$ above the column place $w$.
\item $g_0$: an element of inertia at a place above $l$ with $\varepsilon(g_0)=1+l^{s_0}u_0$, $u_0\in\Z_l^\times$.
\item $j_l=\dim_{L_w}H^1_f(G_{E_w},W_0\otimes L_w)$, $W_0=\chi_1^{-1}\rho_\pi$.
\item $K_{l,a}$: $j_l$ continuous cocycles $D_l\to O^r$ (block $0$, action through $\psi_0$). Their images in $K^r$ are
finite, and their classes form a basis, modulo torsion, of the classes of $O^r$-valued cocycles that become finite in
$K^r$. The formal construction uses $O^r$-valued cocycles and a tensor map into $\mathrm{Hom}(K/O,M)$; the passage to the
$A$-valued form of Definition~\ref{i2:def-colmod} is not described here.
\item $k^0_l=\delta_Ou_l$, with $u_l:D_l\to O$ a continuous homomorphism, trivial on inertia, taking the value $1$ at some
element (a Frobenius coordinate). There is no factor $l^{e'}$ in the formal input.
\end{itemize}
The item that stays undefined is $\mathrm{Adm}_i$. The closed theorem uses exactly nineteen of its clauses:
\begin{itemize}
\item (G1)--(G6) and the tame-generator conditions;
\item the Bloch--Kato and Bloch--Kato-annihilator clauses;
\item continuity and semisimplicity of the family;
\item Fontaine's comparison above $l$;
\item local finiteness of $H^1$ above $l$;
\item local-coefficient compatibility;
\item existence of good congruence objects;
\item the unramified input;
\item the row-length clause, the lift clause and the block clause $c_B$.
\end{itemize}
The block clause can be defined in about a page: at each row and good coefficient extension, $l^{c_B}$ times a cut
cocycle with the jumps of $\tilde r_\pi$ is finite, together with a comparison for the lines $\varepsilon^{-e}$. The row-length
and lift clauses take about half a page each; we do not describe the remaining ones. The formal block clause bounds the
relevant length by $e_Oc_B(\deg_w+1)$, while the value $e_Oc_B(\deg_w+2)$ used in item (3) appears in the formal Selmer
inequality; the weaker bound used here is consistent with both.

\emph{Steps.}
\begin{itemize}
\item[(C1)] (R) Item (1): no fixed vector off $Z_i$ (from (G6)), plus a rank count of the column span in each fibre:
$\mathrm{rank}\le j_l+\dim\ker$.
\item[(C2)] (L/N) Row bounds: Bloch--Kato orthogonality and Tate duality (Theorems~\ref{i2:thm-bk}, \ref{i2:thm-tate}),
with $\ell_w$ from the row-length clause.
\item[(C3)] (N) Column lifting: the local lifting theorem (Proposition~\ref{i2:prop-lift}) at the column places.
\item[(C4)] (R) Global terms: $s_i(\pi)$ by definition; $\mathrm{Sha}$ via \eqref{i2:eq-sha}; the trivial block via the unramified
input and Lemma~\ref{i2:lem-unrhom}.
\item[(C5)] (N) Block bounds from the block clause.
\end{itemize}
We do not describe these steps below the hypotheses they use.

\medskip\noindent\label{i2:sst-cut3}\emph{The cut datum.}

The formal statement of Proposition~\ref{i2:prop-cut}(3) assumes the full tie $\mathrm{Tie4}$ of $D$, with an explicit norm
choice $x$ satisfying the Galois preconditions. It concludes for the cut datum with the restricted norm choice. Part (3) is stated in this form. Chebotarev (degree-one form) is proved in the development; local
finiteness, local Tate duality and the local Euler characteristic enter as arguments of the formal
statement and are proved as separate propositions elsewhere in the development.
\begin{enumerate}
\item[(1)] (proved here) A good cut exists: parts (1), (2) and Remark~\ref{i2:rem-cut-extra}.
\item[(2)] (R) Galois preconditions transfer field by field. The cut only shrinks $X_i$, keeping points, eigensystems and
Galois representations. One clause (the first conclusion) is re-derived; we do not describe this here.
\item[(3)] (R) The cut forces local--global compatibility at the split places above $g_0,g_1$ for every parameter with a
point, from the extra clause of Remark~\ref{i2:rem-cut-extra}.
\item[(4)] (L + N) \emph{(G5) for the cut datum.} Local--global compatibility from (3), then the sufficient condition for the
Galois clause. The local inputs enter only here. Let $w\in S_i$ and $\bar\rho$ be residual with
$\mathrm{charpoly}(\bar\rho(\gamma_w))$ not vanishing at $\bar\chi_1(\gamma_w)$. The twisted residual representation has:
\begin{itemize}
\item $H^0=0$, by the non-root condition;
\item $H^2=0$, by Theorem~\ref{i2:thm-tate};
\item $H^1=0$, by Theorem~\ref{i2:thm-euler}, since $w\nmid l$.
\end{itemize}
\item[(5)] (L) (G6) for the cut datum; we do not describe the proof here.
\item[(6)] (R) Jump clauses, from (Par).
\item[(7)] (L) The duties clause uses Chebotarev (Theorem~\ref{i2:thm-cheb}, degree-one form); we do not describe the proof here.
\item[(8)] (R) \emph{$\mathrm{Tie4}$ transfers.} $z_O(z)$ descends to $M_i/M^{\rm out}_i$ since it commutes with $T_i$. The
eigenspace clause restricts. The other components are unchanged.
\end{enumerate}
\emph{New:} only (4) has content beyond restriction.

\medskip\noindent\label{i2:sst-pcclocal}\emph{Local inputs at a generator prime: the remaining items.}

\emph{(a1).} The formal statement agrees with the one above. For each point $x$, form over $x$, place $v$ of the reflex
field, and $\tau$ in the type at $v$ whose label is squarefree and avoids $c^{\pm1}$, there is $u\in\{c,c^{-1}\}$ with the
following properties at $w$ and at the conjugate place: $r_x$ is de Rham, has monodromy $0$, and has inertia polynomial
$\prod_{y\in M}(X-\omega_y)$ (respectively with inverted roots).
\begin{enumerate}
\item[(1)] (L, clause DZ of the pending datum) The local parameter at $w$ has monodromy $0$ and residual label
$\mathrm{lab}(\tau)(X-u)^t$.
\item[(2)] (N, short) The local parameter satisfies the tame and wild conditions and the root-multiplicity condition. The
reason: $\mathrm{lab}(\tau)$ is squarefree, so its only repeated root is $u\in\mathbb F_g^\times$, which satisfies
$u^{g-1}=1$.
\item[(3)] (L, clause (BC)) A base-change list $\Sigma$ for the form.
\item[(4)] (cited formal) Lemma~\ref{i2:lem-isorep}, including its clauses at $w\mid l$: a semisimple $\rho$ attached to
$\Sigma$.
\item[(5)] (L) $r_x\otimes\bar\Q_l\cong\rho$. Characteristic polynomials of Frobenius agree at split places outside a
finite set. Split Frobenii are dense (Theorem~\ref{i2:thm-cheb}). Continuity then gives equality of characteristic
polynomials, and Brauer--Nesbitt gives the isomorphism. For Brauer--Nesbitt see Curtis--Reiner (30.16) (reference to be confirmed). Both are clauses of (SM1).
\item[(6)] (L, tie0) Conjugate representations have the same de Rham property, Weil--Deligne representation and
Hodge--Tate weights.
\item[(7)] (R) Monodromy $0$ transfers.
\item[(8)] (N, short) \emph{Inertia polynomial.} The pending datum gives $\prod_{y\in M}(X-\omega_y(\sigma))$. Each $\omega_y$ is
multiplicative on inertia, hence trivial on commutators. So the parameter is abelian on inertia, and the relation of
(LGC$'$) gives equality of inertia polynomials.
\end{enumerate}
Literature for the input (LGC$'$) at $w\mid l$: BLGGT 2014, Thm.~2.1.1(4), covers only $\pi_w$ with an Iwahori-fixed
vector; in general: Barnet-Lamb--Gee--Geraghty--Taylor, \emph{Local-global compatibility for $l=p$}, I (Ann.\ Fac.\ Sci.\
Toulouse Math.\ 21 (2012), 57--92) and II (Ann.\ Sci.\ ENS 47 (2014), 165--179), up to semisimplification, and Caraiani,
Algebra Number Theory 8 (2014), 1597--1646, for the Frobenius-semisimplified statement with monodromy, which the full Weil--Deligne relation needs. The use
in (8) compares characteristic polynomials on inertia only, which is insensitive to semisimplification, so for that use
BLGGT $l=p$ I, II suffice.

\emph{The separately verified items of kind (i).} Each is a short derivation from clauses of the datum that are named, not
defined, here:
\begin{itemize}
\item (G1), (G4): (R), read off.
\item (G5), (G6): (R), from the datum's Galois-family conclusions and the good-point hypothesis.
\item Local-coefficient compatibility: (R), from the branch data.
\item The tame-generator conditions: we do not describe their proof here.
\item The dual-jump, potential de Rham and tensor inputs, the branch tie and the inertia equality: from the local tie
records; we do not describe their proofs here.
\item The good congruence objects: (R).
\item The crystalline period bound, the pushout proposition, the twisted-base clause and Breuil--Kisin admissibility:
clauses of the datum. A derivation of the twisted-base clause from Theorems~\ref{i2:thm-bk}, \ref{i2:thm-tate},
\ref{i2:thm-euler} exists in the development; we do not describe it here.
\end{itemize}

\medskip\noindent\label{i2:sst-isorep}\emph{The Galois representation of an admissible list.}

The proof of Lemma~\ref{i2:lem-isorep} is a construction. No Chebotarev, Brauer--Nesbitt or strong multiplicity one is
used; those enter Lemma~\ref{i2:lem-splitbc} and (a1) instead.
\begin{enumerate}
\item[(1)] (L, clause of (LGC$'$)) For cuspidal, conjugate self-dual, regular algebraic $\pi$ on $\mathrm{GL}_a$ there is $\rho_\pi$
with the weak relation at $w\nmid l$. At $w\mid l$, under the tame and wild conditions, $\rho_\pi$ is de Rham with the
relation of (LGC$'$) and Hodge--Tate weights $-\inf(\pi)+(a-1)/2$. Literature: BLGGT 2014, Thm.~2.1.1 (at $w\mid l$
only for Iwahori-spherical $\pi_w$), with BLGGT $l=p$ I, II and Caraiani 2014 at the other $w\mid l$, as above.
\item[(2)] (R/L, tie0) The properties are stable under direct sums (de Rham, Weil--Deligne and Hodge--Tate data are
additive; Theorem~\ref{i2:thm-fontaine}).
\item[(3)] (R/L, tie0) Stability under twist by the Galois character of an algebraic Hecke character (Hodge--Tate weights
shift).
\item[(4)] (N, short) \emph{Half twists by $|\cdot|^{e/2}$.}
\begin{itemize}
\item If $e+a-1=2k$, take $\rho_\pi$ twisted by the character of $\eta_T\eta_T^{-1}$ with shift $k$.
\item If $e+a-1=2k+1$, then $\pi\otimes\eta_T$ is regular algebraic ($\eta_T$ has half-integral infinity type, a tie0 clause).
Take $\rho_{\pi\otimes\eta_T}$ twisted by the character of $\eta_T^{-1}$ with shift $k+\tfrac12$.
\end{itemize}
\item[(5)] (R) For a Speh part $(a,\pi,d)$, sum over $k<d$. The infinity weights are distinct and lie in $\Z+(n_V-1)/2$
(regularity of $P$).
\item[(6)] (R) Sum over the parts; the rank is $n_V$.
\item[(7)] (R) Repackaging.
\end{enumerate}
We do not describe here the proofs of the clauses on $\eta_T$ and of the regularity of the infinity weights.

\medskip\noindent\label{i2:sst-potc}\emph{The dual case of the local count.}

Proposition~\ref{i2:prop-potcount}(c) is not obtained by applying (b) to $T^*(1)$: the cut cocycles are still those of
$M$, and only the lifting hypothesis moves to the dual family. In the formal statement the pushout bound is a clause of
the datum; here it is Proposition~\ref{i2:prop-pushout}, proved. The formal statement also uses the semisimplicity of the
members over $K$, included in the statement.
\begin{enumerate}
\item[(1)] (R) Reduction to one cut-cocycle datum, as in (b).
\item[(2)] (N, elementary) \emph{Complements.} The $G_E$-stable $K$-span $W$ of $Y$ in $\bigoplus_b\tilde r_{x_b}\otimes K$ has a
$G_E$-stable complement $C$. This follows from semisimplicity of each member, by induction over blocks. Here
semisimplicity is global, of $G_E$, not local.
\item[(3)] (N, elementary) \emph{Dual realisation.} Pair $\bigoplus\tilde r_{x_b}$ with $\bigoplus r^\vee_{x_b}$ by the dot
product. Then $\mathrm{pr}(gu,gu')=\chi_{\rm cyc}(g)\mathrm{pr}(u,u')$. The annihilator $X'=\mathrm{ann}(C)$ is $G$-stable, and the pairing
$X\times X'\to L$ is perfect, since $W\oplus C$ is everything. So $X'$ is a Tate dual of $X$ inside a sum of dual members.
\item[(4)] (L) $X$ is de Rham: it is a subrepresentation of a sum crystalline over $K'$ (Theorem~\ref{i2:thm-fontaine}), and
potentially de Rham implies de Rham.
\item[(5)] (N + L) The same for $U=M$. The rank bound
$\mathrm{rk}\,H^1_f(T')=\dim H^1_f(U')\le\dim U'=k$ comes from Theorem~\ref{i2:thm-bk} with $E_w=\Q_l$. The complementation of
the local matrix representation, the de Rham property of $U'$ and the existence of the dual lattice enter through their
statements; we do not describe their proofs here.
\item[(6)] (R) $X'$, $U'$ are subfamilies of the dual family, so (Lift$_e$) applies to them.
\item[(7)] (L) Cut cocycles are pushouts of finite cocycles: step (5) of the proof of (b).
\item[(8)] (proved) Proposition~\ref{i2:prop-pushout}, second case of Lemma~\ref{i2:lem-absorb}, gives $\le e(k+k)\ell(O/l)$.
\end{enumerate}
So the proof of (c) is the proof of (b) with step (4) of that proof replaced by (2)--(6) above. The only inputs not proved
here are the three sub-lemmas of (5).

\medskip\noindent\label{i2:sst-div}\emph{The division clauses.}

\emph{Method of Lemma~\ref{i2:lem-divclauses}.} The method is the Arakelov geometry of Hermitian $O_K$-lattices
(slopes and covolumes), not a lattice-point count.
\begin{enumerate}
\item[(V)] (R) \emph{Value-two input.} For every finite $w$, $\phi\in N(u)$ and member vector $F$, the pairing value is $0$
or has $w$-valuation $\ge\sigma_u(w)$. This uses four clauses: the modified-level field gives $e_u(N(u))=\mathrm{Nu}$-lattice
after cancelling the invertible ideal $\mathfrak a_{\sigma}$; tiedness of the branch data writes the value as an
$O_K$-combination; the cut-up condition; and (lq2). The property is closed under sums and $O_K$-multiples.
\item[(D1)] (N, elementary) Write $\phi\in\mathfrak a_\sigma N$ as $\sum a_i\phi_i$ with $v_w(a_i)\ge-\lfloor\sigma(w)\rfloor$, and write $x\in\mathrm{RR}_c$
locally through a member basis; the latter uses the gluing clause of $\mathrm{RR}_c$, a local--global principle for
lattices over Dedekind domains (reference to be confirmed). By (V), each term has $v_w\ge0$, so the value is integral at every
$w$.
\item[(D4)] (L/R) $\mathrm{covol}(\mathfrak aN)=N(\mathfrak a)^d\mathrm{covol}(N)$, so the height changes by
$\frac2{[K:\Q]}\sum_w\lfloor\sigma(w)\rfloor\log N\mathfrak p_w$. Standard; reference to be confirmed (e.g.\ Bost; Gaudron).
\item[(D2)] (N) \emph{Height chain.} $B_w$ factors through the chain of pairings of the column core. Per place,
\[
\log\|B_w\|_v^2\le\log p_v^2+\sum_j\log k_{j,v}^2-\log\kappa_v-\min_p\sum_j\log\lambda_{j,p,v},
\]
using the quotient metric and the action on Hecke parts. Averaging, with the clause bounding
$\sum_j\log\lambda_{j,p,v}$, gives the mean bound $h_R-\overline{\log\kappa}$. We do not describe the proofs of these statements here.
\item[(D3)] (N; the divisor inequality proper) A chain of inequalities:
\begin{itemize}
\item the slope bound of the base side;
\item the passage clause carrying it to the column side, with $\mathrm{de}\le\frac13$;
\item the explicit numerical record of the package.
\end{itemize}
This is the longest part of the formal proof; we do not describe the proofs of the slope bound or of the numerical record here.
\end{enumerate}
(V), (D1) and (D4) are elementary once the clauses they use are defined. Those clauses are among the about seventy
fields of the package and its input record, which are named and not defined here; so the lemma is taken from the formal verification.
The goodness clause of the division data is invertibility of the frames (nonzero complex determinant), not unimodularity.

\emph{Proposition~\ref{i2:prop-wordfacts}(ii).} The second assertion, $h(u)=h^{\rm expl}(u)$, is a field of the family; it needs no
proof. The first assertion (clause (3)) is a short case check. Suppose $p>X''(n)$, $p\ne T_u$,
$p\notin\{g_0,g_1\}$, $p\ne d$ and $p\nmid D\Gamma_u$. Then:
\begin{itemize}
\item the gain vanishes, by the support field, or the filling clause (the $p$-part of the gain is $0$ above $X''(n)$), or
the non-inertness field;
\item the error term $E$ vanishes, because every branch of its defining case formula is excluded and
$v_p(D\Gamma_u)=0$.
\end{itemize}
It is (R), but it requires the definitions of the gain, of $E$ and of four family fields, which are not given here.


\medskip\noindent\textbf{Status.}
\emph{Proved here:} the reduction of the divisor bound to $\mathrm{LinkN}$ (Theorem~\ref{i2:thm-div}), the summation over
primes, the cut transfer, parts (1), (2) of Proposition~\ref{i2:prop-cut} (with the central action at members in
\S\ref{i2:ssec-proof-cut}), Proposition~\ref{i2:prop-potcount}(a) relative to its hypotheses (a1)--(a5) and (b), the
lattice bound and the $\mathrm{Hom}$-term algebra, the generator-prime summation (Proposition~\ref{i2:prop-typ}), the
comparisons of Proposition~\ref{i2:prop-other}(1)--(3), the step from the twist conditions to $\mathrm{KillAt}$
(Lemma~\ref{i2:lem-killat}), the cut-cocycle bounds (Proposition~\ref{i2:prop-cutcoc}), the per-point cut
(Proposition~\ref{i2:prop-percut}) and Lemma~\ref{i2:lem-T3Z} (both modulo the formally verified statements below), condition
$\mathrm{HT}_w$, the transport along ties (Lemma~\ref{i2:lem-tietransfer}), the assembly of
Theorem~\ref{i2:thm-linkN}, and the passage to condition (P1) (Proposition~\ref{i2:prop-P1}) from the divisor inequality
of a package.
\emph{Published theorems, quoted:} Theorems~\ref{i2:thm-fund}--\ref{i2:thm-artin}, \ref{i2:thm-bs}, \ref{i2:thm-fontaine},
\ref{i2:thm-herr} and \ref{i2:thm-wach}. We also use the member structure (M1)--(M7), the hypotheses of
Proposition~\ref{i2:prop-potcount}(a), and the derivation of (T3) from the series localisation.
\emph{Formally verified, \cite{CharacterFamily}:} Propositions~\ref{i2:prop-cut}(3), \ref{i2:prop-potcount}(c),
\ref{i2:prop-pcclocal} (steps (Res), (a1) and the separately verified items of kind (i)), \ref{i2:prop-colconds},
\ref{i2:prop-ptsel}, \ref{i2:prop-locadm}, \ref{i2:prop-large}, the first assertion of Proposition~\ref{i2:prop-wordfacts}(ii),
Lemmas~\ref{i2:lem-isorep} and \ref{i2:lem-divclauses}, and the Breuil--Kisin steps (Propositions~\ref{i2:bk-b}, \ref{i2:bk-ex},
\ref{i2:bk-cl}, \ref{i2:bk-mon}, \ref{i2:bk-tri}, \ref{i2:bk-int}).
Some clauses are named but not defined here: the local admissibility conditions, the Galois preconditions, the norm-choice
facts, $\mathrm{TwGood}$, the canonicity of hub data, the clauses ($\chi$-tie) and (DZ) of $\mathrm{Tie4}$, the datum clauses
(LGC$'$) and (SM1), and the division data of $\mathrm{DKCB}$. Statements that use them are stated informally, as indicated at
each of them; their exact forms are in \cite{CharacterFamily}. The Selmer hypothesis $\mathrm{HypV}$ is supplied by
Section~\ref{i5:sec}.




\section{Level counts and almost orthogonality}\label{i3:sec}

This section proves, or reduces to clearly stated propositions, the third
arithmetic input of Theorem~\ref{M:thm:alpoge-families}: the almost-orthogonality
of the kernel blocks (Theorem S), the irreducibility defect, and the level
counts.  Statements marked \emph{formally verified} are proved in
\cite{CharacterFamily}; we state them and indicate what their proofs
involve, but do not reproduce them.  The main result is
Theorem~\ref{i3:thm-main}, stated in \S\ref{i3:ssec-setting} together with the
notions it uses.

The section is organised as follows.  Subsections~\ref{i3:ssec-setting}--\ref{i3:ssec-count}
contain the definitions and the statements: Bessel uniqueness under a spectral gap
and its archimedean source (\S\S\ref{i3:ssec-bessel}, \ref{i3:ssec-gap}), the off-diagonal
bound (\S\ref{i3:ssec-S}), its passage to almost orthogonality (\S\ref{i3:ssec-linkS}),
irreducibility (\S\ref{i3:ssec-irr}) and the level counts (\S\ref{i3:ssec-count}).
Subsections~\ref{i3:ssec-grp-count}--\ref{i3:ssec-grp-St} contain the proofs, and
\S\ref{i3:ssec-structure} outlines the formally verified steps.  The section ends with a
summary of its status.

\subsection{Setting and main theorem}\label{i3:ssec-setting}

A \emph{family} is a collection $\mathcal F=(\mathcal D(x,i),\mathcal L(x,i))$
indexed by a real parameter $x>1$ and a finite set $I_x$.  Each
$\mathcal D=\mathcal D(x,i)$ is a column datum: a CM extension $E/F$ with
$f=[F:\Q]$, ranks $k,m$ and $r=k+m-1$, a finite set of \emph{words} $u$, a
prime level $T_u$ and a positive integer $d_u$ for each word, an expected
dimension $D_*>0$, a total dimension $D_{\rm all}$ and a row space of dimension
$D_{\rm rows}=(1-\delta)D_{\rm all}$; this defines the irreducibility defect
$\delta=\delta(\mathcal D)$.  The record $\mathcal L(x,i)$ carries the scale
\[
 L=x^{\varepsilon},\qquad b=e^{L},
\]
where $0<\varepsilon<1/52$ is the exponent of the family, and an auxiliary
exponent $\varepsilon_z=\varepsilon/J$, where $J\ge1$ is a constant of the
family (so $0<\varepsilon_z\le\varepsilon$ is the same for all members).  The datum has a kernel operator
whose $(u,u')$ block we denote $K(u,u')$; $\|\cdot\|_{HS}$ is the
Hilbert--Schmidt norm.  We put
\[
 \Gamma=\Gamma(\mathcal L)=\exp(C_\Gamma m^5fL),
\]
with $C_\Gamma$ a constant of the family: strict positivity $C_\Gamma>0$ is
part of the definition of the record $\mathcal L(x,i)$, and all records of a
family carry the same value of $C_\Gamma$, so $C_\Gamma>0$ is a fixed positive
constant of $\mathcal F$.  A statement holds
\emph{eventually} for $\mathcal F$ if there is $x_0$ such that it holds for
all $x\ge x_0$ and all $i\in I_x$.

\emph{Closed statements.}  A family carries finitely many real constants, the
same for all its members.  A fixed finite list of them are the \emph{closure
constants}; among these are $C_\Gamma$, the constant $C_{k,\rm low}$ and the
six constants whose minimum is $C_l^{\mathcal F}$
(\S\ref{i3:ssec-pf-slots}).  The others are the \emph{outer constants}
(among them $C_{YL}$, $C_{Tu}$, $J$).  A \emph{threshold system} $\mathcal K$
assigns to each closure constant a lower bound which is a function of the outer
constants (and, for the later ones in the list, of the earlier closure
constants); $\mathcal F$ is \emph{$\mathcal K$-large} if each closure constant
of $\mathcal F$ is at least its bound.  For a class $\mathcal C$ of families
and a property $P$ of families, we say that \emph{$P$ is closed on $\mathcal
C$} if there is $\varepsilon_1>0$ such that for every $0<\varepsilon\le
\varepsilon_1$ there is a threshold system $\mathcal K$ such that $P(\mathcal
F)$ holds for every $\mathcal F\in\mathcal C$ of exponent $\varepsilon$ which
is $\mathcal K$-large.  Taking maxima of thresholds and the minimum of the
$\varepsilon_1$ shows that finitely many closed properties are closed
simultaneously.  When a closed statement below lists \emph{Large conditions}
(explicit inequalities such as $1000\le C_l^{\mathcal F}$), the threshold
system may be, and is, chosen so that every $\mathcal K$-large family satisfies
them; further thresholds that are not explicit here are said to be ``of
\cite{CharacterFamily}''.  The \emph{own constants} of $\mathcal F$ are
$C'=C_{k,\rm low}$ and $C_l=C_l^{\mathcal F}$; ``the standing hypotheses at
$(C_0,C_{k,\rm low},C_l^{\mathcal F})$'' (\S\ref{i3:ssec-pf-slots}) are
those at the own constants of the member, with window constant $C_0$.

\begin{definition}[Almost orthogonality]\label{i3:def-linkS}
Let $C_0\ge1$ and $\eta_0,\eta_1,e_0\in\R$.  The datum $\mathcal D$ is
\emph{$(C_0,\Gamma,\eta_0,\eta_1,e_0)$-almost orthogonal} if the following
implication holds: if $C_0^{-1}\le d_u\Gamma^2/D_*\le C_0$ for every word $u$
(the \emph{window}), then
\begin{enumerate}
\item[(i)] $\|K(u,u')\|_{HS}^2\le(1+\eta_0)\,d_ud_{u'}/D_*$ for all $u\ne u'$;
\item[(ii)] $\|K(u,u)-\mathrm{id}\|_{HS}^2\le e_0\,d_u$ for all $u$;
\item[(iii)] $D_{\rm all}\le(1+\eta_1)D_*$.
\end{enumerate}
The family $\mathcal F$ satisfies $\mathrm{LinkS}(C_0)$ if eventually there
are $\eta_0,\eta_1\in\R$, $e_0\ge0$ with
$\eta_0,\eta_1,\sqrt{e_0}\le\exp(-x^{\varepsilon/8})$ such that
$\mathcal D(x,i)$ is $(C_0,\Gamma(\mathcal L(x,i)),\eta_0,\eta_1,e_0)$-almost
orthogonal.  It satisfies $\mathrm{LinkIRR}$ if eventually
$0\le\delta(\mathcal D(x,i))\le\exp(-x^{\varepsilon_z/7})$.
\end{definition}

\begin{definition}[Level counts]\label{i3:def-count}
Let $v=|d_E|^{1/(2f)}/2$ (the host parameter), $g_0,g_1$ the two auxiliary
primes of $\mathcal D$, $w_j(\tau)$ ($1\le j\le m$, $\tau$ an infinite place
of $F$) its archimedean weights, $a=xb$, $L_a=\log(2av)$.  For $C\in\R$ the
datum satisfies $\mathrm{Count}(C)$ if for every word $u$
\begin{enumerate}
\item $\displaystyle\sum_{j=1}^m\frac1f\sum_\tau\log\frac{w_j(\tau)v}{\pi}
 +m\log(g_0g_1)-C\Bigl(L_a+\frac{m^2L_a}{k}\Bigr)\le\log T_u$;
\item $\log T_u\le CmL$;
\item $|\log T_u-\log T_{u'}|\le C\log k$ for every word $u'$.
\end{enumerate}
The family satisfies $\mathrm{CountF}$ if there is $C\ge0$ with eventually
$\mathrm{Count}(C)$.
\end{definition}

\begin{theorem}[Level counts and almost orthogonality]\label{i3:thm-main}
Let $\mathcal C$ be the working class of
Theorem~\ref{M:thm:alpoge-families}, with its spectral selection.  (The unitary
Siegel--Weil formula (SW) of Definition~\ref{i3:S-def-SW}, which the proofs use, holds by
Proposition~\ref{i3:prop-SW}; it is not a hypothesis.)
\begin{enumerate}
\item[(a)] (Off-diagonal bound.)  If every selected member is locally tied and
the rank row and the loop row hold with constant $3$, then the conclusion of
Proposition~\ref{i3:prop-S-closed} is closed on $\mathcal C$: eventually, for
all words $u\ne u'$, if the standing hypotheses hold at
$(3,C_{k,\rm low},C_l^{\mathcal F})$ then
$\|K(u,u')\|_{HS}^2\le(1+b^{-6})\,d_ud_{u'}/D_*$.
\item[(b)] (Almost orthogonality.)  Under the slot hypotheses
(Sel), (Rk), (Lp), (Ty), (Mn), (Inst) of Theorem~\ref{i3:thm-linkS-slots},
the property ``$\mathrm{TuSize}(\mathcal F)$ implies $\mathrm{LinkS}(3)$'' is
closed on $\mathcal C$.
\item[(c)] (Irreducibility.)  Assume also Green's parametrisation (G)
(Definition~\ref{i3:def-green}), that every member of $\mathcal C$ satisfies
$\mathrm{TuSize}$, and the automorphic tie hypothesis (AT) of
Proposition~\ref{i3:prop-irr-row}.  Then $\mathrm{LinkIRR}$ is closed on
$\mathcal C$.
\item[(d)] (Level counts.)  Every family of $\mathcal C$ whose members
eventually have the volume tie (VT) of Proposition~\ref{i3:prop-massratio}
satisfies the torus-pair and variation bounds of
Proposition~\ref{i3:prop-part5}, and hence $\mathrm{CountF}$.  (VT) holds
eventually on every family of a class on which the real-place second-layer
ties are closed (Definition~\ref{i4:def-hyps}(vii)).
\end{enumerate}
\end{theorem}

\noindent Part (d) is condition (P4) of Definition~\ref{M:def:package}.  Parts
(b) and (c) control, through the frame data, the small parameters $\eta$, $e$,
$\delta$ in (P8) and (P10).  The proof occupies the rest of the section:
(a) in \S\ref{i3:ssec-bessel}--\S\ref{i3:ssec-S}, (b) in
\S\ref{i3:ssec-linkS}, (c) in \S\ref{i3:ssec-irr}, (d) in
\S\ref{i3:ssec-count}.  The statement (SW) is Proposition~\ref{i3:prop-SW}; the hypotheses (G), the slot hypotheses,
(AT) and the volume tie (VT) are stated where they are used; Proposition~\ref{i3:prop-slots-SW}
discharges (Sel), (Rk), (Lp) from (SW) for classes built from a raw selection
with the bundle condition.

\subsection{Bessel uniqueness under a spectral gap}\label{i3:ssec-bessel}

The archimedean input of Theorem S is a multiplicity-one statement for Bessel
models of discrete series which are far from the walls.  Fix integers
$k\ge1$, $r\ge0$, $n=2r+k+1$, a Hermitian matrix $H=H_{k,r}$ of size $n$ and
$U=\{g\in GL_n(\C):g^*Hg=H\}$.  The basis is split as
$e_0,\dots,e_{r-1}$, a \emph{middle block} of size $k+1$, and
$e^*_{r-1},\dots,e^*_0$, with $H$ pairing $e_a$ with $e^*_a$, so that $H$ has
signature $(k+r,r+1)$.  For $z\in\C^n$ supported on the middle block with
$T(z)=-\Re\,{}^t\bar zHz>0$ let
\[
 U_z=\{h\in U:\ he_a=e_a,\ he^*_a=e^*_a\ (a<r),\ hz=z\},
\]
the pointwise fixer of the $e_a$, $e^*_a$ and $z$.  It acts on the orthogonal
complement of $z$ in the middle block, which is definite of dimension $k$, so
$U_z$ is a copy of the compact group $U(k)$.  (The stabiliser of $z$ alone is
non-compact for $r\ge1$ and is not used.)  Let $N\subset U$ be the standard
unipotent subgroup of the Bessel datum.  For
$c\in(\C^\times)^r$ put
\[
 \psi_{z,c}(g)=\exp\Bigl(i\,\Re\sum_{a<r}c_a\ell_a(g)\Bigr),\qquad
 \ell_a(g)=g_{a,a+1}\ (a+1<r),\quad \ell_{r-1}(g)=(gz)_{r-1}.
\]

\begin{definition}\label{i3:def-gap}
A vector $\lambda\in\R^n$ is a \emph{Harish-Chandra parameter} if it is
strictly decreasing and $\lambda_a-\frac{n-1}2\in\Z$ for all $a$.  For a sign
pattern $\epsilon:\{1,\dots,n\}\to\{0,1\}$ and $h\in\R$, $\lambda$ has
\emph{gap at least $h$} if $|\lambda_a-\lambda_b|\ge h$ whenever
$\epsilon_a\ne\epsilon_b$.
\end{definition}

If $\epsilon$ records the two blocks of the signature, the pairs with
$\epsilon_a\ne\epsilon_b$ are the non-compact roots, so the gap is a uniform
lower bound for $|\langle\lambda,\alpha^\vee\rangle|$ over non-compact roots.

\begin{definition}\label{i3:def-bessel}
Let $\sigma$ be a unitary representation of $U$ on $H_\sigma$, with smooth
vectors $H_\sigma^\infty$, and $\pi$ a representation of $U_z$ on a
finite-dimensional inner product space $W$.  A \emph{smooth Bessel map} for
$(z,\psi,\sigma,\pi)$ is a $\C$-linear map $B:H_\sigma^\infty\to W$ with
$B(\sigma(hn)v)=\psi(n)\pi(h)B(v)$ for $h\in U_z$, $n\in N$, which is
continuous in the sense that $\|B(v)\|\le C\sum_j\|d\sigma(X_j)v\|$ for some
$C$ and finitely many elements $X_j$ of the enveloping algebra of
$\mathrm{Lie}(U)$.
\end{definition}

\medskip\noindent\emph{Discrete series with a Harish-Chandra parameter.}

Let $\epsilon:\{1,\dots,n\}\to\{0,1\}$ be a sign pattern with $p=\#\epsilon^{-1}(1)$, $q=n-p$, $q_0=\min(p,q)$, and
$J_\epsilon=\mathrm{diag}(\epsilon'_1,\dots,\epsilon'_n)$, $\epsilon'_a=1$ if $\epsilon_a=1$ and $-1$ otherwise.  Put
$U_\epsilon=\{g:g^*J_\epsilon g=J_\epsilon\}$, $\mathfrak g_\epsilon$ its Lie algebra, $K=U_\epsilon\cap U(n)\cong U(p)\times U(q)$ and
$\mathfrak p=\{X\in\mathfrak g_\epsilon:X^*=X\}$.  Choosing a unitary $P$ with $P^*J_\epsilon P=H$ (where $\epsilon$ records the
signature of $H$), $g\mapsto PgP^{-1}$ is an isomorphism $U\to U_\epsilon$ preserving $\|\cdot\|_{HS}$; we transport
representations along it without comment.  For a unitary representation $\sigma$ of $U_\epsilon$ and a matrix unit $e_{ab}$
write $E_{ab}=d\sigma(e_{ab})$ for the complexified derivative on smooth vectors; the $E_{ab}$ satisfy the relations of
$\mathfrak{gl}_n$ and $\langle E_{ab}x,y\rangle=s_{ab}\langle x,E_{ba}y\rangle$ with $s_{ab}=\epsilon'_a\epsilon'_b$.  A pair $(a,b)$ is
\emph{compact} if $\epsilon_a=\epsilon_b$.  Let $\rho_c(a)$ (resp.\ $\rho_n(a)$) be half the number of $b>a$ minus half
the number of $b<a$ with $(a,b)$ compact (resp.\ non-compact), $\rho=\rho_c+\rho_n$, and for a Harish-Chandra parameter
$\lambda$ put
\[
 \Lambda_\lambda=\lambda+\rho_n-\rho_c,\qquad c(\lambda)=|\lambda|^2-|\rho|^2 .
\]
A \emph{$K$-type of highest weight $\mu$} of $\sigma$ is a finite-dimensional irreducible $K$-stable subspace containing a
vector $v\ne0$ with $E_{aa}v=\mu_av$ for all $a$, all of whose weights are dominated by $\mu$.

\begin{definition}\label{i3:def-ds}
A \emph{discrete series with parameter $(\epsilon,\lambda)$} is a nonzero irreducible unitary representation $\sigma$ of
$U_\epsilon$ (or of $U$, through $P$) all of whose matrix coefficients are square integrable, which has a $K$-type of highest
weight $\Lambda_\lambda$, and such that $|\Lambda_\lambda+2\rho_c|\le|\mu+2\rho_c|$ for every $K$-type of highest weight
$\mu$ of $\sigma$ (Vogan minimality of $\Lambda_\lambda$).
\end{definition}

We use three standard results of representation theory.

\begin{theorem}[Harish-Chandra, \emph{Discrete series for semisimple Lie groups II}, Acta Math.\ 116 (1966) 1--111]\label{i3:thm-hc}
Let $G$ be a connected semisimple Lie group with finite centre, $K$ a maximal compact subgroup.  If
$\mathrm{rank}\,G=\mathrm{rank}\,K$, then $G$ has square-integrable irreducible unitary representations, and to every
regular parameter $\lambda$ (non-singular, with $\lambda-\rho$ integral for the compact Cartan subgroup) there corresponds a
discrete series representation with infinitesimal character $\lambda$, whose formal degree is a positive constant multiple
of $\prod_{\alpha>0}|\langle\lambda,\alpha\rangle|$.
\end{theorem}

\noindent What the theorem delivers is existence, the parametrisation by infinitesimal character, and the shape of the formal
degree up to a constant depending on the Haar measure.  It does not apply verbatim to $U(p,q)$, whose centre $U(1)$ is not
finite; the extension to reductive groups of Harish-Chandra's class is due to Harish-Chandra (\emph{Harmonic analysis on real reductive groups} I, J.\ Funct.\ Anal.\ 19 (1975) 104--204; III, Ann.\ of Math.\ (2) 104 (1976) 117--201; reference to be confirmed).  The
characterisation of the discrete series by a Vogan-minimal lowest $K$-type $\Lambda_\lambda=\lambda+\rho_n-\rho_c$, which is the
form of Definition~\ref{i3:def-ds}, is not in this paper; it is due to Schmid (Ann.\ of Math.\ (2) 102 (1975) 535--564), Hecht--Schmid (Invent.\ Math.\ 31 (1975) 129--154) and Vogan (Ann.\ of Math.\ (2) 109 (1979) 1--60) (reference to be confirmed).
We use the theorem only as background, to explain why Definition~\ref{i3:def-ds} describes the classical discrete series; the
arguments below construct an explicit model (Proposition~\ref{i3:prop-model}) and do not depend on it.

\begin{theorem}[Schur orthogonality; Godement, \emph{Sur les relations d'orthogonalit\'e de V.~Bargmann}, C.~R.\ Acad.\ Sci.\ Paris 225 (1947), 521--523 and 657--659]\label{i3:thm-godement}
Let $G$ be a unimodular locally compact group with Haar measure $\mu$ and $\sigma$ an irreducible unitary representation
of $G$ on $\mathcal H$ having one nonzero square-integrable matrix coefficient.  Then all matrix coefficients are square
integrable and there is $d_\sigma>0$ (the formal degree) with
\[
 \int_G\langle v_1,\sigma(g)u_1\rangle\overline{\langle v_2,\sigma(g)u_2\rangle}\,d\mu(g)=d_\sigma^{-1}\langle u_2,u_1\rangle\langle v_1,v_2\rangle
 \qquad(u_i,v_i\in\mathcal H).
\]
\end{theorem}

\noindent\emph{Convention.}  In this section all inner products are conjugate-linear in the first variable and linear in the
second, as in the formal statements; the display above is Godement's relation written in this convention.  With it, every
coefficient written below as $\overline{\langle w,\sigma(g)w\rangle}$ is the kernel of the projection onto $\C w$:
$d_\sigma\int_G\overline{\langle w,\sigma(g)w\rangle}\,\sigma(g)v\,d\mu(g)=\langle w,v\rangle w$ for a unit vector $w$, by the
display with $(u_1,v_1,u_2,v_2)=(v,u,w,w)$ paired against an arbitrary $u$.

\noindent $U$ is reductive, hence unimodular, and a discrete series has square-integrable coefficients by definition;
the compact group $U_z\cong U(k)$ with an irreducible finite-dimensional $\pi$ also satisfies the hypotheses.

\begin{theorem}[Theorem of the highest weight; Knapp, \emph{Lie Groups Beyond an Introduction}, 2nd ed., Thm.~5.110]\label{i3:thm-knapp}
For a compact connected Lie group, every dominant analytically integral weight is the highest weight of an irreducible
finite-dimensional representation, unique up to equivalence.
\end{theorem}

\noindent It is applied to $K\cong U(p)\times U(q)$, compact and connected, with the weight $\Lambda_\lambda$ restricted to
the two blocks; that these restrictions are dominant and integral is checked in Lemma~\ref{i3:lem-ktype}.

\begin{proposition}\label{i3:prop-bessel}
Let $\lambda$ be a Harish-Chandra parameter of gap at least $h$ for $\epsilon$, with $h\ge120n^3$.  Let $\sigma$ be a
discrete series representation of $U$ with parameter $(\epsilon,\lambda)$, $z$ supported on the middle block with $T(z)>0$,
$c\in(\C^\times)^r$, and $\pi$ a continuous irreducible unitary representation of $U_z$ on a finite-dimensional space $W$.  Then any two smooth Bessel maps for
$(z,\psi_{z,c},\sigma,\pi)$ are linearly dependent.
\end{proposition}

\noindent The proof is in \S\ref{i3:ssec-pf-bessel}.  It rests on Theorems~\ref{i3:thm-godement} and \ref{i3:thm-knapp}
on the boundedness of Bessel maps on orbits (Proposition~\ref{i3:prop-orbbd}, proved in \S\ref{i3:ssec-pf-orbbd}), and on two
narrower statements which are formally verified: the representation-theoretic properties of the model
(Proposition~\ref{i3:prop-model}) and the analytic and Gelfand--Kazhdan inputs for the Bessel operator
(Proposition~\ref{i3:prop-GKin}).  The published
multiplicity-one theorems for Bessel models (Jiang--Sun--Zhu, GAFA 20 (2010) 690--709; Sun--Zhu, Ann.\ of Math.\ (2) 175 (2012) 23--44) need no gap but are not used.


\subsection{The gap is supplied by the archimedean weights}\label{i3:ssec-gap}

The datum has, at each infinite place $\tau$, integers $v_a(\tau)$
($0\le a<k-1$) and archimedean weights $b_j(\tau)$ ($1\le j\le m$).  Write
$r=k+m-1=2r_S+k+1$ with $k$ even.  Put, for $0\le a<r$,
\[
 \lambda^{SH}_a(\tau)=\begin{cases}v_a(\tau),&a<k-1,\\
 -b_{a-k+2}(\tau),&a\ge k-1,\end{cases}\qquad
 \epsilon^{SH}_a=\begin{cases}1,&a<k-1,\\
 1\text{ iff }a-k+1\text{ is even},&a\ge k-1.\end{cases}
\]
The \emph{standing hypotheses} at constants $(C',C_l,\dots)$ are a list of
explicit inequalities among $k,m,f,r,L,b$, the auxiliary primes and the
archimedean data.  Those used below are
\begin{equation}\label{i3:eq-standing}
\begin{gathered}
 8C_\Gamma\le C',\qquad C'm^5fL\le k,\qquad m^4\le f\le k,\qquad m\ge2,\quad k\ge m+3,\\
 C_l\,r^{19}f^5L\le b_{j+1}(\tau)-b_j(\tau)\quad(1\le j<m),\qquad
 C_lkm^4+rm^{10}<\lambda^{SH}_i(\tau)-\lambda^{SH}_{i+1}(\tau)\quad(i+1<r).
\end{gathered}
\end{equation}
In the families the weights are chosen with spacing
$h_0\le b_{j+1}-b_j\le h_0+1$, $h_0=X\,r^{19}f^5L$, so the node-gap clause
holds as soon as $C_l\le X$; producing such weights is part of Section~\ref{i4:sec}.

\begin{lemma}\label{i3:lem-gap}
Under \eqref{i3:eq-standing}, for every $\tau$, $\lambda^{SH}(\tau)$ is a
Harish-Chandra parameter with gap at least $h_0=C_lr^{19}f^5L$ for
$\epsilon^{SH}$.  If moreover $C_l\ge1000$, $L\ge\frac12$, $f\ge2$, and
$\lambda'\in\R^r$ satisfies $\sum_i|\lambda'_i-\lambda^{SH}_i(\tau)|\le C_nkm^3$
with $C_n\le C_l$, then $\lambda'$ has gap at least $h=h_0/2\ge120r^3$.
\end{lemma}

\begin{proof}
The last clause of \eqref{i3:eq-standing} has positive left side, so
$\lambda^{SH}$ is strictly decreasing; its entries are integers and
$\frac{r-1}2=r_S+k/2\in\Z$.  Let $a<b$ with
$\epsilon^{SH}_a\ne\epsilon^{SH}_b$.  Since $\epsilon^{SH}=1$ on
$\{0,\dots,k-1\}$, $b\ge k$; put $j=b-k+1$, so $1\le j<m$.  Monotonicity gives
$\lambda^{SH}_a-\lambda^{SH}_b\ge\lambda^{SH}_{b-1}-\lambda^{SH}_b
=b_{j+1}(\tau)-b_j(\tau)\ge h_0$.

For $\lambda'$: $km^3\le r^4$ and $f^5L\ge16$ give $4C_nkm^3\le h_0/2$, and by
the triangle inequality
$|\lambda'_a-\lambda'_b|\ge|\lambda^{SH}_a-\lambda^{SH}_b|-|\lambda'_a-\lambda^{SH}_a|-|\lambda'_b-\lambda^{SH}_b|
\ge|\lambda^{SH}_a-\lambda^{SH}_b|-\sum_i|\lambda'_i-\lambda^{SH}_i|\ge h_0/2$.
Finally $h_0/2\ge8C_lr^{19}\ge120r^3$.
\end{proof}

The proof in \cite{CharacterFamily} also checks two further numerical
conditions on $h$ needed in the cell decomposition; each has left side
$\le30r^5$ while $h/(8\sqrt m)\ge30r^5$, so they hold by the same estimate.
Thus Proposition~\ref{i3:prop-bessel} applies at every infinite place of every
datum satisfying the standing hypotheses with $1000\le C_l$, to every
Harish-Chandra parameter within $C_nkm^3$ of $\lambda^{SH}$ when $C_n\le C_l$.
(The conditions $L\ge\frac12$ and $f\ge2$ of the lemma hold for every member of
a family, since $L=x^\varepsilon>1$ and $f\ge m^4\ge16$ by
\eqref{i3:eq-standing}; the lemma does not show
that $\lambda'$ is a Harish-Chandra parameter, which has to be supplied
separately.)

\subsection{The off-diagonal bound}\label{i3:ssec-S}

For a local Bessel datum $A$ at an
infinite place we write $\sigma=\sigma_A$ for its discrete series, $v_0=v_0(A)\ne0$ for its lowest $K$-type vector,
$B,B'$ for its two smooth Bessel maps (for $z$ and $z'$), $c_N>0$ for its normalising constant, and
$\Phi(g)=B(\sigma(g)v_0)$, $\Phi'(g)=B'(\sigma(g)v_0)$.  For a cell datum $(\gamma,\nu)$ and two functions
$F:U\to W$, $F':U\to W'$ let $\mathcal J_{\nu,\gamma}[F,F']$ be the cell integral of \cite{CharacterFamily} (an integral over
the cell of $\gamma$ of a sesquilinear expression in $F,F'$; we only use that it is sesquilinear in $(F,F')$), so that the
local Bessel integral is $J_A(\nu,\gamma)=c_N^{-1}\mathcal J_{\nu,\gamma}[\Phi,\Phi']$.  For $K$-finite $a,b$ and
$\xi\in W$, $\xi'\in W'$, $\mathrm{Orb}_{\gamma,\nu}(a,b,\xi,\xi')$ is the orbital integral over $(\gamma,\nu)$ of the
matrix coefficient $g\mapsto\langle\sigma(g)a,b\rangle$ twisted by $\xi,\xi'$.  Let $w_0=v_0/\|v_0\|$ and let $\mathcal
D_0$ be the span of the vectors $d\sigma(X_1)\cdots d\sigma(X_j)w_0$ ($X_i\in\mathrm{Lie}(U)$).

\begin{proposition}\label{i3:prop-amp}
Assume the standing hypotheses with $1000\le C_l$ and $C_n\le C_l$.  Let $(\tau,w,w')$ be admissible local data at an
infinite place, $\lambda'$ within $C_nkm^3$ of $\lambda^{SH}(\tau)$, $A$ a good local Bessel datum for $\lambda'$ (part of
which is that $\lambda'$ is a Harish-Chandra parameter) with Bessel maps $B,B'$, $a,b$
$K$-finite vectors and $\xi,\xi'$ with $\langle Ba,\xi\rangle\langle\xi',B'b\rangle\ne0$.  Then there is $\kappa\ne0$ such
that for every cell datum $(\gamma,\nu)$
\[
 \mathrm{Orb}_{\gamma,\nu}(a,b,\xi,\xi')=\kappa\,J_A(\nu,\gamma),\qquad \kappa=d_\sigma\,s\,K_1^{-1},
\]
where $d_\sigma>0$ is the formal degree (Theorem~\ref{i3:thm-godement}), $s=\langle la,\xi\rangle\langle\xi',l'b\rangle$
and $K_1=c_N^{-1}(c_0\|v_0\|)\overline{(c_0'\|v_0\|)}$, with $l,l',c_0,c_0'$ as in the proof.
\end{proposition}

\noindent The proof (\S\ref{i3:ssec-pf-amp}) is short given Proposition~\ref{i3:prop-bessel} and the following
narrower statement on the local data of the family, which is formally verified.

\begin{proposition}[formally verified, \cite{CharacterFamily}]\label{i3:prop-amp-in}
In the situation of Proposition~\ref{i3:prop-amp} (without the non-vanishing hypothesis on $a,b$):
\begin{enumerate}
\item[(A1)] (Bessel pair) there are smooth Bessel maps $l$ for $(z,\psi_{z,c},\sigma,\pi)$ and $l'$ for
 $(z',\psi_{z',c},\sigma,\pi')$ such that the Bessel operator of $\sigma$ (Definition~\ref{i3:def-beop}) factors as
 $\beta(x,y)=l(x)\otimes\overline{l(y)}$, and similarly for $l'$, for $x\in\mathcal D_0$ and $y$ smooth;
\item[(A2)] (non-vanishing) there are $x,y\in\mathcal D_0$ with $\beta(x,y)\ne0$, and $x',y'\in\mathcal D_0$ with
 $\beta'(x',y')\ne0$;
\item[(A3)] (amplitude identity) for every cell datum $(\gamma,\nu)$, all $K$-finite $v,v'$ and all $\xi,\xi'$,
\[
 \langle lv,\xi\rangle\langle\xi',l'v'\rangle\cdot\mathcal J_{\nu,\gamma}\bigl[l(\sigma(\cdot)w_0),\,l'(\sigma(\cdot)w_0)\bigr]
 =d_\sigma^{-1}\,\mathrm{Orb}_{\gamma,\nu}(v,v',\xi,\xi').
\]
\end{enumerate}
\end{proposition}

\noindent\emph{What the proof involves.}  (A1) is the construction of the Bessel functional from the multiplicity-one map of
Lemma~\ref{i3:lem-thmM} at the two vectors $z,z'$, and the smoothness of $\mathcal D_0$; (A2) is a non-vanishing computation
on the span of derivatives of the lowest vector; (A3) unfolds the orbital integral over the cell with Schur orthogonality,
the inversion invariance of the relative stabiliser measure and the absolute convergence of the cell integrals (which uses
the second numerical condition on $h$ of \S\ref{i3:ssec-gap}).  In the formalization these are the inputs to the amplitude
identity and its convergence, which is a long verification.  Its formally verified steps are outlined in \S\ref{i3:str-amp}.


\begin{definition}[Unitary Siegel--Weil formula in the convergent range]\label{i3:S-def-SW}
Let $F$ be a totally real number field, $E/F$ a totally imaginary quadratic extension with nontrivial automorphism
$x\mapsto\bar x$, $\mathbb A=\mathbb A_F$, $F_\infty=F\otimes\R$, and for a finite place $v$ let $F_v$ be the completion
and $\mathcal O_v$ its ring of integers.  Let $M\ge1$, let $J\in GL_M(E)$ be hermitian ($\bar J^t=J$), put
$h(x,y)=\bar x^{\,t}Jy$ for $x,y\in E^M$ (conjugate-linear in the first variable) and let $U=U_J$ be the unitary group,
$U(R)=\{g\in GL_M(R\otimes_FE):\bar g^{\,t}Jg=J\}$ for $F$-algebras $R$.  For $n\ge0$ let $V_n=(E^M)^n$ (tuples
$X=(X_1,\dots,X_n)$ of column vectors), with $U$ acting by $gX=(gX_1,\dots,gX_n)$ and with the Gram map
$\mathrm{gr}(X)=(h(X_a,X_b))_{a,b}\in\mathrm{Herm}_n$.  Fix a basis $\beta$ of $E/F$ and let $L_v\subset V_n(F_v)$ be
the set of tuples whose $\beta$-coordinates lie in $\mathcal O_v$.  A \emph{test function} is
$\Phi=\Phi_\infty\otimes\bigotimes_v\Phi_v$ on $V_n(\mathbb A)$ with $\Phi_\infty$ smooth and compactly supported on
$V_n(F_\infty)$, each $\Phi_v$ locally constant and compactly supported, and $\Phi_v=\mathbf 1_{L_v}$ for almost all
$v$.  Fix a nontrivial character $\psi$ of $\mathbb A/F$; for $R=F_\infty$ or $F_v$ give $V_n(R)$ and
$\mathrm{Herm}_n(R)$ the Haar measures that are self-dual for $\psi_R$ and the pairings
$\langle X,Y\rangle=\sum_a\mathrm{tr}_{E/F}h(X_a,Y_a)$ and $\langle S,S'\rangle=\frac12\mathrm{tr}_{E/F}\mathrm{tr}(SS')$.
A \emph{local density} of $\Phi_R$ is a function $g_R$ on the open set $\mathrm{Herm}_n^\times(R)$ of nondegenerate
matrices, continuous there, with $\int_{\mathrm{gr}^{-1}(K)}\Phi_R\,dX=\int_Kg_R(S)\,dS$ for every compact
$K\subset\mathrm{Herm}_n^\times(R)$.  For nondegenerate $T\in\mathrm{Herm}_n(F)$ and $h\in U(\mathbb A)$ put
$\theta_T(h)=\sum_{X\in V_n(F),\ \mathrm{gr}(X)=T}\Phi(h^{-1}X)$.  We say that \emph{the unitary Siegel--Weil formula holds
in the convergent range} (statement (SW)) if, for all such $F,E,M,J$:
\begin{enumerate}
\item[(i)] $U(F)$, which is discrete in $U(\mathbb A)$, has a fundamental domain for every Haar measure $\mu$ on
$U(\mathbb A)$, and every fundamental domain $\mathcal F$ has $\mu(\mathcal F)<\infty$ (a \emph{fundamental domain} is a
null-measurable set whose $U(F)$-translates cover $U(\mathbb A)$ up to a $\mu$-null set and are pairwise disjoint up to
$\mu$-null sets);
\item[(ii)] for every $n$ with $2n<M$, every nondegenerate $T\in\mathrm{Herm}_n(F)$, every test function $\Phi$, every
Haar measure $\mu$, fundamental domain $\mathcal F$ and character $\psi$ as above, there are local densities
$g_\infty$ of $\Phi_\infty$ and $g_v$ of $\Phi_v$ such that $\prod_vg_v(T)$ converges (unconditionally, over the
finite places) and
\[
\frac1{\mu(\mathcal F)}\int_{\mathcal F}\theta_T(h)\,d\mu(h)=g_\infty(T)\prod_vg_v(T).
\]
\end{enumerate}
\end{definition}

\noindent The left side of (ii) is the mean of $\theta_T$ over $U(F)\backslash U(\mathbb A)$, so it does not depend
on $\mu$; the right side does not depend on $\psi$ or $\beta$ (by the product formula).  (SW) is a theorem: it consists
of two published results, which we now state.

\begin{theorem}[Borel--Harish-Chandra; Borel and Harish-Chandra, Ann.\ of Math.\ 75 (1962) 485--535; reference to be
confirmed]\label{i3:thm-BHC}
Let $G$ be a connected reductive group over a number field $F$ without nontrivial $F$-rational characters.  Then $G(F)$
is a discrete subgroup of $G(\mathbb A_F)$ of finite covolume.
\end{theorem}

\begin{theorem}[Siegel--Weil formula, $T$-th coefficient; Weil, \emph{Sur la formule de Siegel dans la th\'eorie des
groupes classiques}, Acta Math.\ 113 (1965) 1--87; for unitary groups A.~Ichino, Math.\ Z.\ 247 (2004) 241--277; reference to be confirmed for the $T$-th coefficient form]\label{i3:thm-SWT}
In the notation of Definition~\ref{i3:S-def-SW}, let $2n<M$, $T\in\mathrm{Herm}_n(F)$ nondegenerate and $\Phi$ a
Schwartz--Bruhat function on $V_n(\mathbb A)$ which is a pure tensor.  Then the theta integral converges absolutely,
the local densities $g_\infty(T)$, $g_v(T)$ exist (they are the values at $T$ of the local Siegel integrals
$W_{T,v}(\Phi_v)=\int_{\mathrm{Herm}_n(F_v)}\int_{V_n(F_v)}\Phi_v(X)\psi_v(\langle S,\mathrm{gr}(X)-T\rangle)\,dX\,dS$, with the
self-dual measures), $\prod_vg_v(T)$ converges absolutely, and the mean of $\theta_T$ over $U(F)\backslash U(\mathbb A)$
equals $g_\infty(T)\prod_vg_v(T)$.
\end{theorem}

\noindent Theorem~\ref{i3:thm-SWT} is the $T$-th Fourier coefficient, at the identity, of the Siegel--Weil identity
$E(g,s_0,\Phi)=I(g,\Phi)$ in Weil's convergent range, Theorem~\ref{i4:thm-sw} of the existence section: the $T$-th
coefficient of the theta integral $I$ at $g=1$ is the mean of $\theta_T$; the $T$-th coefficient of the Siegel
Eisenstein series is, for nondegenerate $T$ and in the range of absolute convergence ($2n<M$), the Euler product of
the local Siegel integrals $W_{T,v}$ (with the self-dual measures, for which $\mathrm{Herm}_n(F)\backslash
\mathrm{Herm}_n(\mathbb A)$ has volume one); and $W_{T,v}(\Phi_v)$ is, by Fourier inversion on $\mathrm{Herm}_n(F_v)$,
the density at $T$ of the push-forward of $\Phi_v\,dX$ under $\mathrm{gr}$.  The normalisation (mean, self-dual
measures, constant $1$) is Weil's Tamagawa normalisation; the exact place in Weil's paper where it is stated is to be
confirmed against the source.

\begin{proposition}[(SW) holds]\label{i3:prop-SW}
The unitary Siegel--Weil formula of Definition~\ref{i3:S-def-SW} holds.
\end{proposition}

\begin{proof}
(i) $U(F)$ is a discrete subgroup of $U(\mathbb A)$ because $F$ is discrete in $\mathbb A$.  The group $U$ is
connected reductive, and its rational characters are trivial: any character factors through $\det:U\to U_1$, the norm
one torus of $E/F$, which is anisotropic over $F$.  Theorem~\ref{i3:thm-BHC} gives finite covolume.  A countable
discrete subgroup of a second countable locally compact group has a Borel fundamental domain, and any two measurable
fundamental domains have the same Haar measure (the covolume), which is finite.

(ii) Let $2n<M$.  A test function in the sense of Definition~\ref{i3:S-def-SW} is a Schwartz--Bruhat pure tensor, so
Theorem~\ref{i3:thm-SWT} applies and gives local densities at $T$, the convergence of $\prod_vg_v(T)$, and the identity.
The densities of Theorem~\ref{i3:thm-SWT} are continuous on $\mathrm{Herm}_n^\times$ and satisfy the defining property of
a local density (for $n=0$ both sides are $\Phi(0)$).  The hypothesis $2n<M$ is the range of Theorem~\ref{i3:thm-SWT}:
it implies $n=0$ or $n+2\le M$, and, for every Witt index $r\le M/2$ of $J$, $n+r<M$ and $M>2n$, the
hypotheses of Theorem~\ref{i4:thm-sw}.
\end{proof}

\noindent\emph{The formal proof.}  In \cite{CharacterFamily} (SW) is proved, not assumed, and the proof does not pass
through Eisenstein series; the published statements above are cited so that the reader need not rely on it.  Its
structure: (L) part (i) is proved by reduction theory (Siegel sets for $U_J$: a torus embedding, the modulus of the
torus, properness of unipotent coordinates and cocompactness of the unipotent quotients), which is the proof of
Theorem~\ref{i3:thm-BHC} specialised to $U_J$.  For (ii) one proves the stronger statement for all $n$ with $n=0$ or
$n+2\le M$, for the measure $I_T=\mu(\mathcal F)^{-1}\int_{\mathcal F}\sum_{\mathrm{gr}(X)=T}\delta_{h^{-1}X}\,d\mu(h)$
on $V_n(\mathbb A)$, in the form ``$I_T$ is the restricted product of the local fibre measures'', by induction on $n$:
(R) $n=0$ is trivial; (N) \emph{rank one} ($n=1$, $M\ge3$): $I_T$ is locally finite, $U(\mathbb A)$-invariant and
carried by the fibre $\mathrm{gr}^{-1}(T)$, which is a single $U(F_v)$-orbit at every place (Witt), so $I_T=c\cdot$
(product of fibre measures) for a constant $c\ge0$; then $c=1$ is forced by dilation: for integers $a\ge1$ one compares
the $I_T$-mass of the tuples $X$ with $aX\in\prod_vL_v$ against the product of the local masses, shows that their
difference is bounded independently of $a$, while the product grows like $a^{[F:\Q](2M-2)}$ along a sequence of $a$;
(R) the formula is invariant under change of basis of $J$ and of $T$; (L) \emph{induction step}: write
$T\simeq t\oplus T_1$; if $J$ represents $t$, then $J\simeq t\oplus J_1$ (Witt), and unfolding over the first vector
reduces the formula for $(J,T)$ to the rank-one formula for $(J,t)$ and the formula for $(J_1,T_1)$, the local
densities factoring accordingly; if $J$ does not represent $t$ globally, both sides vanish (the local side at a place
where $t$ is not represented; we do not describe the details of this case).  This is Siegel's inductive argument; the
formal proof is self-contained.  In the main line of \cite{CharacterFamily}, the input (SW) of the off-diagonal theorem
and of the row statements below is supplied by this proved statement.

\begin{proposition}\label{i3:prop-S-closed}
Assume (SW).  Let $\mathcal C$ be a class of families with a spectral selection
(Definition~\ref{i3:S-def-selection}) such that every selected member is locally tied, and such that the rank
row and the loop row (Definition~\ref{i3:S-def-rows}) hold with constant $3$.  Then the following property of
$\mathcal F$ is closed on $\mathcal C$ (\S\ref{i3:ssec-setting}): eventually, at every member and for all words
$u\ne u'$, if the standing hypotheses hold at $(3,C_{k,\rm low},C_l^{\mathcal F})$, then
\[
 \|K_0(u,u')\|^2\le\bigl(1+b^{-6}\bigr)\frac{d_ud_{u'}}{D_*}\qquad\text{and}\qquad
 \|K(u,u')\|_{HS}^2\le\bigl(1+b^{-6}\bigr)\frac{d_ud_{u'}}{D_*},
\]
where $K_0(u,u')$ is the term operator of \S\ref{i3:S-sss-abstract}.  The Large conditions include
$29\le C_{k,\rm low}$ and $16\le C_l^{\mathcal F}$ (which give the gap hypothesis of
\S\ref{i3:S-sss-components}), $29\le C_l^{\mathcal F}$ (admissibility of the selection),
$1000+c_{\rm orb}\le C_l^{\mathcal F}$ with $c_{\rm orb}=|c|\ge0$ for the absolute
constant $c$ of the convergence of the cell integrals (which gives $1000\le C_l^{\mathcal F}$ for
Lemma~\ref{i3:lem-gap}), the thresholds $C_{k,\rm low}\ge C'_0$, $C_l^{\mathcal F}\ge C_{l,0}$ of
Proposition~\ref{i3:S-prop-NG}, the thresholds $40$, $40$, $1$ of Lemma~\ref{i3:S-lem-rankrow}, and further thresholds of
\cite{CharacterFamily}.
\end{proposition}

\noindent The first inequality is the formally verified off-diagonal theorem; the second follows from it by the
Hilbert--Schmidt identification, Theorem~\ref{i3:S-thm-HS}.  In the standing hypotheses the window constant is
$C_0=3$ (clause (N10)).

\begin{remark}\label{i3:S-rem-change}
Note the form of the statement: it does not assert, for each family, the existence of constants $C',C_l$; what is proved is the closed form
above, at the family's own constants and for $\varepsilon\le\varepsilon_1$.  In \cite{CharacterFamily} the off-diagonal theorem takes three global inputs: (SW), the
finite dimensionality (FD) of the spaces of automorphic vectors of a fixed archimedean type
(Lemma~\ref{i3:S-lem-FD}), and the \emph{non-generic row} (NG) (Proposition~\ref{i3:S-prop-NG}), which is a bound,
not a non-degeneracy statement.  Both (FD) and (NG) are proved in \cite{CharacterFamily}, the second from
(SW); and (SW) itself holds (Proposition~\ref{i3:prop-SW}).  Local tying,
the rank row and the loop row are hypotheses on the class.  The formal statement asserts the bound on
$\|K_0(u,u')\|^2$; the bound on $\|K(u,u')\|_{HS}^2$ follows by Theorem~\ref{i3:S-thm-HS}.
\end{remark}


\noindent The proof is given in \S\ref{i3:S-ssec-proof}.

\emph{Slot hypotheses.}  For a class $\mathcal C$ with a spectral selection we use the following hypotheses,
named as in \cite{CharacterFamily}:
(Sel) every selected member is locally tied;
(Rk) the rank row holds with constant $3$;
(Lp) the loop row holds with constant $3$ (Definition~\ref{i3:S-def-rows});
(Ty), (Mn), (Inst) as follows (as formulated below).  Below, ``tied'' means: $\mathcal D$ satisfies the ties, $(X,Y)$ is a
concretely tied package over $\mathcal D$, and $X(u,u)$ satisfies (AC) where stated; $(C_g,c_b,C_w,c_T,C_T,C_0)$ are the outer
constants, with $c_b,c_T>0$.
\begin{itemize}
\item[(Ty)] For all outer constants there are $C_H,C_h,C_p\ge0$ such that for every $C_s\ge C_p+4$ there is $C^{(0)}$ such that for
every $C\ge C^{(0)}$ there are $(C'_0,C_{l,0})$ such that for $C'\ge C'_0$, $C_l\ge C_{l,0}$, every datum $\mathcal D$ with size data
satisfying the $K_2$-hypotheses at $(C',C_l)$, every tied package and every word $u$: for $1\le j',j''\le r_S$, labels
$\delta',\delta''$ and every finite set $S\subset\mathrm{all}(\mathrm{deg})$ (Definition~\ref{i3:def-terms}; so
$\mathrm{row}(t_1)=\mathrm{row}(t_2)$) of pairs $(t_1,t_2)$ of degenerate terms of $X(u,u)$ with $t_1,t_2$ relevant, $j(t_1)=j'$, $j(t_2)=j''$,
$\mathrm{lab}(t_1)=\delta'$, $\mathrm{lab}(t_2)=\delta''$,
\[
\sum_{(t_1,t_2)\in S}\|\mathrm{vec}_{t_1}\|\,\|\mathrm{vec}_{t_2}\|\le d_\pi(u)\,m(u)\,\kappa(j',\delta')\,\kappa(j'',\delta''),
\]
where $d_\pi(u)=\dim V_\pi(u)$ and $m(u)\ge0$ (the mass of the rows of $u$) are the dimension and mass factors
of the package, and $\kappa(j,\delta)=k_{2a}(j,\delta)\,k_{2e}(j,\delta)$ with
\[
\begin{aligned}
k_{2a}(j,\delta)&=\Bigl(\frac Ck\Bigr)^{j(k-1-j)f}\,\Delta(j,\delta)^{-\frac{k-1-j}2}\,H_b(C)^{Cr_S^4}\,S(H_1,j,\delta)^{-C_sk},\\
k_{2e}(j,\delta)&=\Delta(j,\delta)^{\frac k2}\,H_b(C_h)^{C_hr_S^4}\,S(H_1,j,\delta)^{C_pk},
\end{aligned}
\qquad H_b(c)=b^{cm^2f},\quad H_1=b^{C_Hm}.
\]
\emph{Inner objects.}  For a term $t$ of $X(u,u)$ with flag pair $(F_u,F_t)$: $t$ is \emph{relevant} if the pair is in
relative position called ``relevant'' (the isotropic parts of the two flags meet the defining lattice relevantly); $j(t)=r_S-\dim_E
I(t)$, where $I(t)$ is the intersection of the isotropic parts of the two flags; $\mathrm{lab}(t)=(\delta_l)_{l\ge1}$ with
$\delta_l\in E$ the $l$-th block minor of the flag pair for $l\le j(t)$ and $\delta_l=0$ for $l>j(t)$.  For a rank $j$ and a
label $\delta$: $\Delta(j,\delta)$ is the local norm index of the top minor $\delta_j$, and
$S(H,j,\delta)=\prod_\tau\max\bigl(1,\max_{l}\|\delta_l\|_\tau/(2H)^{2r_S}\bigr)\ge1$, the inner maximum over $1\le l\le j$ with
$\delta_l\ne0$ and $\|\cdot\|_\tau$ the normalised archimedean size at $\tau$.  The only
properties of $\Delta$ and $S$ used in the proofs below are $S\ge1$ and the bounds (Del), (Cos) of
Proposition~\ref{i3:prop-kppdiag-in}.
\item[(Mn)] For all outer constants and every $A$ there is $C_B\ge0$ and $(C'_0,C_{l,0})$ such that for $C'\ge C'_0$, $C_l\ge C_{l,0}$,
every datum satisfying the counting hypotheses at $(u,u)$ with constants $(C',C_l,\ldots,A)$, every tied package with $X(u,u)$
satisfying (AC): $X(u,u)$ satisfies (RM$_{\gamma_B,d}$) with $\gamma_B=\exp(C_Bm^5f\log b)$ and $d=d_u^2/D_*$.
\item[(Inst)] At every selected member $\mathcal D=\mathcal D(x,i)$ ($x$ beyond the selection threshold): (i) \emph{volume--signature
independence}: for any two hermitian spaces $W,W_0$ over $E$ satisfying the signature condition of the datum, of the same rank, of constant signatures $(P,Q)$, $(P_0,Q_0)$, with
the volume measures of the datum and fundamental domains $\mathcal F,\mathcal F_0$ of $U_W(F)$, $U_{W_0}(F)$, one has
$\mu(\mathcal F)=\mu_0(\mathcal F_0)$; (ii) the \emph{class items} at $n=\mathrm{rank}(Y_0)$: for every definite hermitian space $W$ of rank $n$
(signature $(n,0)$ everywhere) satisfying the signature condition, the mass of the standard lattice $L_0$ equals the genus mass $m_I(n)$ of the datum; there are
finitely many $g_1,\dots,g_m$ with $U_W(\mathbb A)=\bigsqcup_iU_W(F)g_iK_{L_0}$ (disjoint); and for every such decomposition
$\sum_i|\mathrm{Stab}(g_i)|^{-1}=\mathrm{mass}(L_0)$.
\end{itemize}
The objects $\mathrm{vec}_t$ and the package data $V_\pi(u)$, $m(u)$ are part of the construction of the term
system in \cite{CharacterFamily} (Proposition~\ref{i3:prop-terms}); (Ty), (Mn), (Inst) enter only as hypotheses of
Propositions~\ref{i3:prop-kpp-diag}, \ref{i3:prop-diag-rows} and \ref{i3:prop-rowtie}.

\begin{proposition}\label{i3:prop-slots-SW}
Assume (SW).  Let $\mathcal A$ be a class of families with a raw selection satisfying the bundle condition of
\cite{CharacterFamily} (a choice, for each member and each large $x$, of an automorphic realisation of the
column bundle of $\mathcal D(x,i)$), and let $\mathcal D_3(\mathcal A)$ be the class of members of $\mathcal A$
equipped with a design at constant $3$, with the spectral selection induced by the raw selection.  Then on
$\mathcal D_3(\mathcal A)$ the hypotheses (Sel), (Rk) and (Lp) hold, and so does the frame row of
Lemma~\ref{i3:S-lem-rankrow}.  In particular the conclusion of Proposition~\ref{i3:prop-S-closed} is closed on
$\mathcal D_3(\mathcal A)$.
\end{proposition}

\begin{proof}
This is the composition of three formally verified
inputs and Proposition~\ref{i3:prop-S-closed}.  (Sel): the spectral selection induced by a raw selection with the bundle
condition comes with a spectral condition at every selected member, one clause of which is local tying.  (Rk) and the frame row: both are derived from (SW).  (Lp): the loop row is
derived from the finite dimensionality (FD) of Lemma~\ref{i3:S-lem-FD} and the volume ratio
identity supplied by (SW).  With (SW), (Sel), (Rk) and
(Lp) available, Proposition~\ref{i3:prop-S-closed} applies to $\mathcal D_3(\mathcal A)$.  Membership in
$\mathcal D_3(\mathcal A)$ includes, besides the design, a list of size and level clauses (among them $C_{YL}\le9$); the (Lp) step uses the size clause on $T_u$, together with Bessel multiplicity one, the
discrete-series coefficient bound and Wu decay, all formally verified.
\end{proof}

\noindent The three inputs are formally verified \cite{CharacterFamily}; the
hypotheses (Ty), (Mn), (Inst) are not discharged by this proposition.  The hypothesis (SW) holds by
Proposition~\ref{i3:prop-SW}, so the proposition applies unconditionally with respect to it.

\begin{proposition}\label{i3:prop-conc1}
Assume (SW), and let $\mathcal C$ be a class with a spectral selection satisfying (Sel), (Rk), (Lp), (Ty), (Mn),
(Inst).  The following property is closed on $\mathcal C$: eventually, at every member satisfying the standing
hypotheses at $(3,C_{k,\rm low},C_l^{\mathcal F})$, there is an
operator $K''$ (the Kloosterman part, Definition~\ref{i3:def-terms}) with
\begin{enumerate}
\item $K''(u,u')=0$ for $u\ne u'$, and $\|K''(u,u)\|_{HS}^2\le d_ue^{-kf\log k}$;
\item $\|K(u,u')-K''(u,u')\|_{HS}^2\le(1+b^{-6})\,d_ud_{u'}/D_*$ for $u\ne u'$;
\item $\|K(u,u)-K''(u,u)-\mathrm{id}\|_{HS}^2\le\Gamma\,d_u^2/D_*$;
\end{enumerate}
and the count bound $|D_{\rm all}-D_*|\le e^{-c_7rf\log r}D_*$ holds with
$c_7=\frac1{16}$.  The Large conditions are those of Propositions~\ref{i3:prop-S-closed},
\ref{i3:prop-terms}, \ref{i3:prop-kpp-off}, \ref{i3:prop-kpp-diag}, \ref{i3:prop-diag-rows} and
$2\le C_l^{\mathcal F}$, $2C_0(A)\le C_l^{\mathcal F}$ (with $A=2(C_{YL}+1)$ and $C_0(A)$ of
Proposition~\ref{i3:prop-counting}), $8C_\Gamma\le C_{k,\rm low}$, and $C_\Gamma\ge\max(C_B,C_4)+1$ with $C_B,C_4$ as in
Proposition~\ref{i3:prop-diag-rows} (used in Lemma~\ref{i3:lem-diagnum}).
\end{proposition}

\noindent The proof is given in \S\ref{i3:ssec-pf-conc1}.  Item 2 is
Proposition~\ref{i3:prop-S-closed} (whose hypotheses (SW), (Sel), (Rk), (Lp) are assumed) together with
item 1.  Item 1, item 3 and
the count bound are reduced, by Hilbert-space arguments and explicit numerics
proved there, to four formally verified statements about the term expansion
of the kernel (Propositions~\ref{i3:prop-terms}, \ref{i3:prop-kpp-off},
\ref{i3:prop-kpp-diag}, \ref{i3:prop-diag-rows}) and to the counting theorem
(Proposition~\ref{i3:prop-counting}) with the row tie
(Proposition~\ref{i3:prop-rowtie}).


\subsection{From Theorem S to almost orthogonality}\label{i3:ssec-linkS}

This step is elementary.

\begin{lemma}\label{i3:lem-eF}
Assume \eqref{i3:eq-standing}, $C_0\ge1$, $c_7>0$, $k\ge5$, and write
$L=y^8$ with $y\ge1$.  Suppose $(2C_0+2)e^{-C_\Gamma y^8}\le e^{-2y}$ and
$y\le8c_7C_\Gamma y^8$.  Put $e_F=2C_0/\Gamma+2e^{-kf\log k}$.  Then
$e_F\ge0$, $\sqrt{e_F}\le e^{-y}$, $b^{-6}\le e^{-y}$ and
$e^{-c_7rf\log r}\le e^{-y}$.
\end{lemma}

\begin{proof}
Let $\Lambda=C_\Gamma m^5fL=\log\Gamma\ge C_\Gamma y^8$.  From \eqref{i3:eq-standing},
$8\Lambda\le C'm^5fL\le k$, hence $\Lambda\le kf\log k\le rf\log r$.  So
$e_F\le(2C_0+2)e^{-\Lambda}\le(2C_0+2)e^{-C_\Gamma y^8}\le e^{-2y}$.  Next
$b^{-6}=e^{-6y^8}\le e^{-y}$, and $c_7rf\log r\ge c_7k\ge8c_7\Lambda\ge
8c_7C_\Gamma y^8\ge y$.
\end{proof}

\begin{lemma}\label{i3:lem-linkS-concl}
If the conclusions of Proposition~\ref{i3:prop-conc1} hold and the window
$C_0^{-1}\le d_u\Gamma^2/D_*\le C_0$ holds for all $u$, then (i)--(iii) of
Definition~\ref{i3:def-linkS} hold with $\eta_0=b^{-6}$, $e_0=e_F$,
$\eta_1=e^{-c_7rf\log r}$.
\end{lemma}

\begin{proof}
(i) is item 2 with $K''(u,u')=0$; (iii) is the upper half of the count bound.
For (ii) write $K(u,u)-\mathrm{id}=(K-K''-\mathrm{id})(u,u)+K''(u,u)$ and use
$\|A+B\|_{HS}^2\le2\|A\|_{HS}^2+2\|B\|_{HS}^2$.  The window gives
$d_u\le C_0D_*\Gamma^{-2}$, so $\Gamma d_u^2/D_*\le(C_0/\Gamma)d_u$; with
item 1 the total is $\le e_Fd_u$.
\end{proof}

\begin{theorem}\label{i3:thm-linkS-standing}
Let $C_0\ge1$, $c_7>0$, and suppose that eventually (a) the window implies the
standing hypotheses, and (b) the standing hypotheses imply the conclusions of
Proposition~\ref{i3:prop-conc1} with constant $c_7$.  Then $\mathcal F$
satisfies $\mathrm{LinkS}(C_0)$.
\end{theorem}

\begin{proof}
Put $y=x^{\varepsilon/8}$, so $L=y^8$.  The constant $C_\Gamma$ is the same
for all members and strictly positive (\S\ref{i3:ssec-setting}); hence, as
$c_7>0$, for $y$ large
$\log(2C_0+2)+2y\le C_\Gamma y^8$ and $1\le8c_7C_\Gamma y^7$, which are the hypotheses of Lemma~\ref{i3:lem-eF}.  Take $x$ beyond all thresholds.  If the window holds, (a), (b)
and Lemmas~\ref{i3:lem-eF}, \ref{i3:lem-linkS-concl} give almost orthogonality
with $\eta_0,\eta_1,\sqrt{e_0}\le e^{-y}=\exp(-x^{\varepsilon/8})$.  If the
window fails, take $\eta_0=\eta_1=e_0=0$: the implication of
Definition~\ref{i3:def-linkS} holds vacuously.
\end{proof}

\begin{remark}\label{i3:rem-window}
When the window $1/3\le d_u\Gamma^2/D_*\le3$ fails, $\mathrm{LinkS}(3)$
carries no information.  The window must therefore be supplied by the
consumer of this input (the frame construction of Section~\ref{i4:sec}).
\end{remark}

\begin{theorem}\label{i3:thm-linkS-slots}
Assume (SW), and let $\mathcal C$ be a class (the working class) with a spectral selection satisfying the slot
hypotheses (Sel), (Rk), (Lp), (Ty), (Mn), (Inst) of \S\ref{i3:ssec-S}.  Write $\mathrm{TuSize}(\mathcal F)$
for: eventually $b^{m/2}\le T_u$ and $\log T_u\le C_{Tu}mL$ for all words $u$.  Then the property
``$\mathrm{TuSize}(\mathcal F)$ implies that $\mathcal F$ satisfies $\mathrm{LinkS}(3)$'' is closed on
$\mathcal C$, with the Large conditions of Proposition~\ref{i3:prop-conc1}.
\end{theorem}

\noindent The proof is given in \S\ref{i3:ssec-pf-slots}.  It is
Theorem~\ref{i3:thm-linkS-standing} with $C_0=3$: hypothesis (b) is
Proposition~\ref{i3:prop-conc1}, and hypothesis (a) is an elementary but long
comparison of sizes (Proposition~\ref{i3:prop-standing-family}), which uses
only $\mathrm{TuSize}$, the size and parameter clauses carried by the members
of the working class, and the Large condition $8C_\Gamma\le C_{k,\rm low}$.
In \cite{CharacterFamily} the theorem is stated with the slots as explicit
hypotheses; besides those listed it takes the decay of the archimedean
coefficients under the gap, the finite dimensionality (FD), the non-generic row
and the off-diagonal bound, which are proved there (the last two from (SW)),
and the frame row and the count bound, which are closed under the
hypotheses above (Lemma~\ref{i3:S-lem-rankrow} and
Proposition~\ref{i3:prop-conc1}).  For classes built from a raw selection with
the bundle condition, (Sel), (Rk), (Lp) follow from (SW)
(Proposition~\ref{i3:prop-slots-SW}); (Ty), (Mn), (Inst) remain hypotheses.

\subsection{Irreducibility}\label{i3:ssec-irr}

\begin{definition}\label{i3:def-allrows}
An \emph{all-rows datum} $A$ for $\mathcal D$ consists of labellings
$\mathrm{lab}_k$ ($k=0,1$) of the $D_{\rm all}$ rows by local types at $g_k$,
an injection $\mathrm{inc}$ of the $D_{\rm rows}$ rows into all rows, weights,
and for each cylinder set $S$ of label pairs a dimension $\dim_A(S)$ and an
expected value $D^*_A(S)=D_*\cdot\mathrm{share}_k(S)$ for cylinders over
$S\subset\mathrm{Tup}_k$, with $\mathrm{share}_k\ge0$.  It is \emph{faithful}
if $\dim_A(S)$ is the number of rows with label in $S$ and a row has a good
label pair if and only if it lies in the image of $\mathrm{inc}$.
$\mathrm{CountAllRows}(A,\varepsilon_C)$ means $0\le\varepsilon_C<1$ and
$|\dim_A(S)-D^*_A(S)|\le\varepsilon_CD^*_A(S)$ for all cylinders $S$, including
$S$ = everything (with $D^*_A=D_*>0$).  $\mathrm{WeightAllRows}(A)$ means
$\mathrm{share}_k(S)\le W_g(f,m,r,g_k,q,d)$ for every set $S$ of bad types,
where
\[
 W_g(f,m,N,g,q,d)=\frac{3f}{g}+\frac{23f\log g}{\sqrt g}
 +2fm\Bigl(\frac1{4\min(q,d)}+\frac{1+\log N}{g-1}\Bigr)
 +\frac N2\Bigl(\frac78\Bigr)^f+e^{-f/(5\log N)}.
\]
\end{definition}

\begin{lemma}\label{i3:lem-share}
Let $0\le\varepsilon<1$, $D_*>0$, $0\le P_j\le W_j$ ($j=1,2$),
$|D_{\rm all}-D_*|\le\varepsilon D_*$, $|D_j-P_jD_*|\le\varepsilon P_jD_*$ and
$\mathrm{bad}\le D_1+D_2$.  Then
$\mathrm{bad}/D_{\rm all}\le\frac{1+\varepsilon}{1-\varepsilon}(W_1+W_2)$.
\end{lemma}

\begin{proof}
$D_{\rm all}\ge(1-\varepsilon)D_*>0$ and $D_j\le(1+\varepsilon)W_jD_*$.
\end{proof}

\begin{proposition}\label{i3:prop-delta}
Let $A$ be a faithful all-rows datum for $\mathcal D$ with
$\mathrm{CountAllRows}(A,\varepsilon_C)$ and $\mathrm{WeightAllRows}(A)$, and
suppose $\frac{1+\varepsilon_C}{1-\varepsilon_C}(W_g(g_0)+W_g(g_1))\le E$.
Then $0\le\delta(\mathcal D)\le E$.
\end{proposition}

\begin{proof}
Let $\mathrm{Bad}_k$ be the rows whose $k$-th label is bad and $S_k$ its set of
labels.  By faithfulness $\dim_A(\mathrm{cyl}_kS_k)=\#\mathrm{Bad}_k$ and the
rows with two good labels are exactly the image of $\mathrm{inc}$.  Hence
$0\le D_{\rm all}-D_{\rm rows}\le\#\mathrm{Bad}_0+\#\mathrm{Bad}_1$, and
Lemma~\ref{i3:lem-share} with $P_k=\mathrm{share}_k(S_k)\le W_g(g_k)$ gives
$\delta=(D_{\rm all}-D_{\rm rows})/D_{\rm all}\le E$.
\end{proof}

\begin{lemma}\label{i3:lem-wg}
Let $0<\varepsilon\le1$ and $C_s\in\R$.  For all large $x$: if
$g\ge e^{8x^\varepsilon}$, $x^2\le f\le x^{C_s}$, $m\le x^{C_s}$,
$100\le N\le x^{C_s}$ and $\min(q,d)\ge e^{x^{\varepsilon/6}}$, then
$W_g(f,m,N,g,q,d)\le x^{3C_s+3}e^{-x^{\varepsilon/6}}$.
\end{lemma}

\begin{proof}
Let $P=x^{C_s}\ge1$, $u=x^{\varepsilon/6}$, $\eta=e^{-u}$, and take
$x\ge\max(10,100C_s^2+1)$.  Then: $3f/g\le3P\eta$; with $s=g^{1/4}\ge e^{2x^\varepsilon}$,
$23f\log g/\sqrt g\le92f/s\le92P\eta$; $1/(4\min(q,d))\le\eta$ and, since
$1+\log N\le N\le P$ and $\eta(g-1)\ge\frac12$, $(1+\log N)/(g-1)\le2P\eta$,
so the third term is $\le2P^2(\eta+2P\eta)$; $(7/8)^f\le e^{-f/8}\le\eta$ as
$f\ge x^2\ge8u$, so the fourth term is $\le P\eta$; and
$5u\log N\le5x\cdot C_s\cdot2\sqrt x\le x^2\le f$, so the last term is
$\le\eta$.  Summing, $W_g\le103P^3\eta\le x^3P^3\eta$.
\end{proof}

\begin{theorem}\label{i3:thm-linkIRR}
Let $c>0$, $C_s\in\R$.  Suppose that eventually there is a faithful all-rows
datum $A$ for $\mathcal D(x,i)$ with $\mathrm{CountAllRows}(A,\varepsilon_C)$,
$\varepsilon_C=\exp(-crf\log r)$, $\mathrm{WeightAllRows}(A)$, and
$g_j\ge e^{8x^{\varepsilon_z}}$, $x^2\le f\le x^{C_s}$, $m\le x^{C_s}$,
$100\le r\le x^{C_s}$, $\min(q,d)\ge e^{x^{\varepsilon_z/6}}$.  Then
$\mathcal F$ satisfies $\mathrm{LinkIRR}$.
\end{theorem}

\begin{proof}
For $x\ge\max(1,1/c)$, $crf\log r\ge cx^2\ge1$, so $\varepsilon_C\le e^{-1}\le\frac12$
and $\frac{1+\varepsilon_C}{1-\varepsilon_C}\le3$.  By Lemma~\ref{i3:lem-wg},
$W_g(g_j)\le x^{3C_s+3}e^{-x^{\varepsilon_z/6}}$, and since
$x^{\varepsilon_z/6}-x^{\varepsilon_z/7}$ grows faster than any multiple of
$\log x$, for large $x$ we have
$x^{3C_s+3}e^{-x^{\varepsilon_z/6}}\le\frac1{20}e^{-x^{\varepsilon_z/7}}$.
Hence $3(W_g(g_0)+W_g(g_1))\le e^{-x^{\varepsilon_z/7}}$, and
Proposition~\ref{i3:prop-delta} applies with $E=e^{-x^{\varepsilon_z/7}}$.
\end{proof}

\begin{proposition}\label{i3:prop-irr-row}
Assume (SW) and Green's parametrisation (G) (Definition~\ref{i3:def-green}).
Let $\mathcal C$ be a class (the working class) all of whose members satisfy
$\mathrm{TuSize}$, and assume the \emph{automorphic tie hypothesis} (AT): the
following property is closed on $\mathcal C$: eventually every member
$\mathcal D(x,i)$ carries an automorphic realisation $X$ (a Hida setting),
a good operator $G$ on it and an isometry $\iota$ of the row space of
$\mathcal D(x,i)$ onto that of $X$, satisfying the automorphic tie, the
row-dictionary tie and the label tie of \cite{CharacterFamily}.  Then the
following property is closed on $\mathcal C$: eventually there is a faithful
all-rows datum satisfying the hypotheses of Theorem~\ref{i3:thm-linkIRR} with
the constant $c=c_{Ic}=\frac1{16}$ and $C_s=6$.  The Large conditions include
$2\le C_l^{\mathcal F}$, $2C_0(A)\le C_l^{\mathcal F}$ with $A=2(C_{YL}+1)$ and
$C_0(A)$ of Proposition~\ref{i3:prop-counting}, and $8C_\Gamma\le C_{k,\rm low}$.
\end{proposition}

\noindent The proof is given in \S\ref{i3:ssec-pf-irr}.  It is complete
except for formally verified inputs, stated exactly there: the counting
theorem for automorphic rows (Proposition~\ref{i3:prop-counting}), the
construction of the all-rows datum from a tied automorphic realisation
(Proposition~\ref{i3:prop-allrows-built}), the local estimates on random
squarefree polynomials (Proposition~\ref{i3:prop-W})
and the admissibility of the local weights (Proposition~\ref{i3:prop-admLq},
which uses (G)).  The automorphic realisation itself, with its three ties, is
the hypothesis (AT); it is not constructed here.

\subsection{Level counts}\label{i3:ssec-count}

For each word $u$ let $d^*_u>0$ and define $X_u>0$ by
$X_u^{fk}=D_*/d^*_u$.  Let $\Gamma_u\ge1$ be the same for all words.  The
levels are selected in the \emph{level window}
\begin{equation}\label{i3:eq-window}
 T_u^*\le T_u\le\Bigl(1+\frac1{fk}\Bigr)T^*_u,\qquad T^*_u=X_u\Gamma_u^{-2/(fk)}
\end{equation}
(that the level $T_u$ is a prime in this window, of the shape described in clause (A) of
\S\ref{i3:ssec-pf-slots}, is a clause of the datum; the construction of the members in the existence
section fulfils it by Lemma~\ref{i3:lem-levelprime} below, which shows that such a prime exists).

The existence of the prime rests on a prime number theorem in short intervals for the ray classes of an imaginary
quadratic field.  For a prime $d\equiv3\bmod4$ and an integer $N\ge1$ say that a prime $p$ is
\emph{$(d,N)$-principal} if there are $A,B\in\Z$ with $A^2+dB^2=4p$, $N\mid B$ and $2N\mid A-2-B$; equivalently, with
$K=\Q(\sqrt{-d})$, $p=\pi\bar\pi$ for some $\pi=\frac{A+B\sqrt{-d}}2\in\mathcal O_K$ with
$\pi\equiv1\bmod N\mathcal O_K$ (since $\mathcal O_K=\Z[\frac{1+\sqrt{-d}}2]$ and
$\pi-1=\frac{A-2-B}2+B\frac{1+\sqrt{-d}}2$).

\begin{theorem}[Primes in short intervals in ray classes of $\Q(\sqrt{-d})$]\label{i3:thm-shortprimes}
There are constants $A,c,c'>0$ such that for every prime $d\equiv3\bmod4$ with $d\ge100$, every integer $N\ge3$ and
all reals $X\ge1$ and $h$ with $\log(4dN^2)\le c\sqrt{\log X}$ and $e^{-c\sqrt{\log X}}\le h\le1$, the number of
$(d,N)$-principal primes $p$ with $\lfloor X\rfloor<p\le\lfloor X(1+h)\rfloor$ is at least
$c'hX/((4dN^2)^A\log X)$.
\end{theorem}

\noindent This is Proposition~\ref{i4:thm-shortnorm} of the existence section, word for word (with $M=N$, $A_0=A$,
$B_0=B$), where it is proved; it is a prime ideal theorem in short intervals for the trivial class of the ray class group of
conductor $N$ of $K$, uniform in the conductor up to $\log(4dN^2)\ll\sqrt{\log X}$ with a polynomial loss $(4dN^2)^{-A}$.

\begin{lemma}[A prime in the level window]\label{i3:lem-levelprime}
Let $0<\varepsilon<1/52$ and $c_s>0$.  There is $x_0$ such that the following holds for all $x\ge x_0$, with
$L=x^\varepsilon$ and $b=e^L$.  Let $m,f,k\ge1$ be integers with $x^{1/3}\le m$ and $fk\le x^{12}$; let $d$ be a
prime with $d\equiv3\bmod4$ and $d\ge100$, $N\ge1$ an integer with $d\mid N$ and $\log N\le3x^{\varepsilon/5}(\log
x)^2$, and $g_0,g_1\ge1$ integers with $g_0,g_1\le b^9$; let $T^*>0$ with $\log T^*\ge c_smL$.  Then there is a
$(d,N(g_0g_1)^2)$-principal prime $T$ with
\[
 T^*<T\le\Bigl(1+\frac1{fk}\Bigr)T^*.
\]
\end{lemma}

\begin{proof}
Let $A,c,c'$ be the constants of Theorem~\ref{i3:thm-shortprimes} and apply it with $X=T^*$, $h=1/(fk)$ and
$N'=N(g_0g_1)^2$; $N'\ge d\ge100$.  Since $\log T^*\ge c_smL\ge c_sx^{1/3+\varepsilon}$, we have $T^*>1$ and
$\sqrt{\log T^*}\ge\sqrt{c_s}\,x^{1/6+\varepsilon/2}$.  On the other side $\log d\le\log N$ and $\log g_i\le9L$, so
\[
 \log(4dN'^2)\le\log4+3\log N+4\log g_0+4\log g_1\le\log4+9x^{\varepsilon/5}(\log x)^2+72x^{\varepsilon}.
\]
As $\varepsilon<1/52$, each of the three terms is at most $\frac{c\sqrt{c_s}}3x^{1/6+\varepsilon/2}$ once $x$ is large
in terms of $\varepsilon,c,c_s$; so $\log(4dN'^2)\le c\sqrt{\log T^*}$.  Next $h\le1$ as $fk\ge1$, and
$h\ge x^{-12}\ge\exp(-c\sqrt{c_s}x^{1/6})\ge e^{-c\sqrt{\log T^*}}$ for $x$ large.  The theorem gives at least
$c'hT^*/((4dN'^2)^A\log T^*)>0$, hence at least one, $(d,N')$-principal prime $T$ with $\lfloor T^*\rfloor<T\le
\lfloor T^*(1+h)\rfloor$; then $T^*<T\le(1+\frac1{fk})T^*$.
\end{proof}

\noindent In the construction of the members, the lemma is applied to the host's prime $d$, the modulus
$N=\mathrm{QS}\cdot qd$ of clause (A) and the admissible primes $g_0,g_1$ (for which the construction gives $g_i\le2(kb)^8\le b^9$, as
$k\le x^5$ and $2x^{41}\le b$ for large $x$), with $T^*=T^*_u$; the size clauses on $m$, $fk$ and on $\log N$ and the
lower bound $\log T^*_u\ge c_smL$ are verified there.  The element $\pi\in\mathcal O_K\subset\mathcal O_E$ with
$T=\pi\bar\pi$ and $\pi\equiv1\bmod\mathrm{QS}\cdot qd(g_0g_1)^2$ is the one required by (A).

\begin{lemma}\label{i3:lem-power}
Let $f,k\ge1$ with $2f\le2^k$, $T\ge2$, $X>0$, $\Gamma\ge1$, $D_*,d^*>0$ with
$X^{fk}=D_*/d^*$, and $Y\le T\le Y(1+\frac1{fk})$ with $Y=X\Gamma^{-2/(fk)}$.
Then $\frac12D_*\Gamma^{-2}\le d^*(T^k-1)^f\le eD_*\Gamma^{-2}$,
$\log T\le\log X+\log2$ and $\log X-2\log\Gamma/(fk)\le\log T$.
\end{lemma}

\begin{proof}
$d^*Y^{fk}=D_*\Gamma^{-2}$.  Upper: $(T^k-1)^f\le T^{fk}\le(1+\frac1{fk})^{fk}Y^{fk}\le eY^{fk}$.
Lower: $T^k\ge2^k\ge2f$, so by Bernoulli $(1-T^{-k})^f\ge1-fT^{-k}\ge\frac12$,
and $(T^k-1)^f\ge\frac12T^{fk}\ge\frac12Y^{fk}$.  Finally $Y\le X$ and
$T\le2Y$ give the upper bound for $\log T$, and $\log T\ge\log Y$ the lower.
\end{proof}

Lemma~\ref{i3:lem-power} says that the dimension $d^*_u(T_u^k-1)^f$ of the
level-$T_u$ component is within a factor $2e$ of $D_*\Gamma^{-2}$; this is how
the level window produces the dimension window of Remark~\ref{i3:rem-window}.

For a member put
\[
 \mathrm{low}_S=\sum_{j=1}^m\frac1f\sum_\tau\log\frac{w_j(\tau)v}{\pi}
 +m\log(g_0g_1)-11L .
\]

\begin{proposition}\label{i3:prop-V}
For the working class there is a constant $C_Y$ such that, for every family of
the class whose members eventually have the volume tie (VT) of
Proposition~\ref{i3:prop-massratio}, eventually, at every member and for all
words $u,u'$:
\begin{enumerate}
\item[(V1)] $\mathrm{low}_S\le\log X_u\le(15+3C_Y)mL$;
\item[(V3)] $|\log X_u-\log X_{u'}|\le1+\log k+\frac2k$;
\item[(V4)] $\log\Gamma_u\le fk$.
\end{enumerate}
\end{proposition}

\begin{proposition}\label{i3:prop-part5}
For the working class there is a constant $C$ such that, for every family of
the class whose members eventually have the volume tie (VT) of
Proposition~\ref{i3:prop-massratio}, eventually, at every member and for every
word $u$ and $u'$:
\begin{enumerate}
\item[(a)] (torus-pair bound) for every $\Gamma\ge1$ with
$X_u\Gamma^{-2/(fk)}\le T_u\le X_u\Gamma^{-2/(fk)}(1+\frac1{fk})$,
\[
 \log T_u\le CmL,\qquad
 \sum_{j=1}^m\frac1f\sum_\tau\log\frac{w_j(\tau)v}{\pi}+m\log(g_0g_1)
 -C\Bigl(L+\frac{m^2L}k\Bigr)-\frac{2\log\Gamma}{fk}\le\log T_u;
\]
\item[(b)] (variation) $|\log X_u-\log X_{u'}|\le C\log k$;
\item[(c)] $\log\Gamma_u\le fk$.
\end{enumerate}
\end{proposition}

\begin{proof}
Let $C=\max(11,16+3C_Y)$ with $C_Y$ from Proposition~\ref{i3:prop-V}, and fix
a member beyond its threshold.  Recall $k\ge m+3\ge5$, so $\log k\ge1$,
$L\ge1$, $mL\ge1$, and $f\le k\le2^{k-1}$.

(a) Let $\Gamma\ge1$ and $u$ satisfy the window.  $T_u\ge2$ is prime and
$2f\le2^k$, so Lemma~\ref{i3:lem-power} with $X=X_u$ gives
$\log T_u\le\log X_u+\log2\le(15+3C_Y)mL+1\le CmL$ by (V1), and
$\log X_u-2\log\Gamma/(fk)\le\log T_u$.  Since $C\ge11$,
$C(L+m^2L/k)\ge11L$, so the left side of the lower bound in (a) is at most
$\mathrm{low}_S-2\log\Gamma/(fk)\le\log X_u-2\log\Gamma/(fk)\le\log T_u$
by (V1).

(b) $2/k\le1\le\log k$, so by (V3)
$|\log X_u-\log X_{u'}|\le3\log k\le C\log k$.

(c) is (V4).
\end{proof}

\noindent Proposition~\ref{i3:prop-V} is proved in \S\ref{i3:ssec-pf-V}, from explicit formulas for the two degrees $D_*$, $d^*_u$ as products of volumes, local indices, $\sigma$-degrees and Weyl dimensions.  Its inputs are the inductive Siegel mass formula and the tie of the base volume to the mass (Proposition~\ref{i3:prop-massratio}), the formal degree formula in the gap range (Proposition~\ref{i3:thm-HC-fdeg}, formally verified; Remark~\ref{i3:rem-HC-fdeg} relates it to Harish-Chandra's formula) and the data bound of the class (Lemma~\ref{i3:lem-dataB}).  The splitting of a rank-one sublattice in the mass formula uses the following classical fact.

\begin{theorem}[Landherr, Abh.\ Math.\ Sem.\ Hamburg 11 (1936) 245--248]\label{i3:thm-landherr}
Non-degenerate hermitian forms over a CM field $E/F$ are isometric if and only
if they have the same rank, the same determinant modulo $N_{E/F}(E^\times)$,
and the same signature at every infinite place.
\end{theorem}

\noindent Applied to $\mathrm{diag}(a,b)$ and $\mathrm{diag}(n,abn)$ with
$a,b,n\in F$ totally positive (both positive definite everywhere, with
determinants differing by $n\bar n$), it shows that every totally positive $n$
is represented by every totally positive definite hermitian space of rank
$\ge2$.

\begin{proposition}\label{i3:prop-countF}
Let $\mathcal F$ be a family of the working class.  If the conclusions of
Proposition~\ref{i3:prop-part5} hold eventually with
constant $C$, then $\mathcal F$ satisfies $\mathrm{Count}(\max(C,0)+2)$
eventually; in particular $\mathrm{CountF}$.
\end{proposition}

\begin{proof}
Let $C^+=\max(C,0)$.  We have $k\ge m+3\ge3$, so $\log k\ge1$; $L\ge1$; and
$L\le L_a$ since $e^L\le xe^L\cdot2v$ ($x>1$, $2v\ge2\pi$).  Apply (a) with
$\Gamma=\Gamma_u$, whose hypothesis is the level window \eqref{i3:eq-window}.

Item 2: $\log T_u\le CmL\le(C^++2)mL$.

Item 1: $C(L+m^2L/k)\le C^+(L_a+m^2L_a/k)$, and by (c)
$2\log\Gamma_u/(fk)\le2\le2(L_a+m^2L_a/k)$; substituting in (a) gives item 1
with constant $C^++2$.

Item 3: from \eqref{i3:eq-window},
$\log T^*_u\le\log T_u\le\log T^*_u+\log(1+\frac1{fk})\le\log T^*_u+1$ and
$\log T^*_u=\log X_u-\frac2{fk}\log\Gamma_u$.  Since $\Gamma_u=\Gamma_{u'}$,
$|\log T_u-\log T_{u'}|\le|\log X_u-\log X_{u'}|+1\le C\log k+1\le(C^++2)\log k$
by (b).
\end{proof}

This completes the proof of Theorem~\ref{i3:thm-main}: (a) is
Proposition~\ref{i3:prop-S-closed}; (b) is Theorem~\ref{i3:thm-linkS-slots},
built on Theorem~\ref{i3:thm-linkS-standing}; (c) is
Theorem~\ref{i3:thm-linkIRR} with Proposition~\ref{i3:prop-irr-row} (within a
$\mathcal K$-large family, the conclusion of
Proposition~\ref{i3:prop-irr-row} gives the hypotheses of
Theorem~\ref{i3:thm-linkIRR}, so $\mathrm{LinkIRR}$ is closed); (d) is
Proposition~\ref{i3:prop-countF} with Proposition~\ref{i3:prop-part5}, the
latter proved from Proposition~\ref{i3:prop-V}.

\subsection{Almost orthogonality and the level counts}\label{i3:ssec-grp-count}

We prove Theorem~\ref{i3:thm-linkS-slots} and Proposition~\ref{i3:prop-V}, and state the counting theorem on which the level counts rest.

\subsubsection{Proof of Theorem~\ref{i3:thm-linkS-slots}}\label{i3:ssec-pf-slots}

\medskip\noindent\emph{The records of the working class.}
Each member of the working class carries, besides the datum, a \emph{size
record}, a \emph{parameter choice} and the record $\mathcal L$; their
constants $C_\Gamma$, $C_{YL}$, $C_{Tu}$, $C_{k,\rm low}$ and
$C_l^{\mathcal F}=\min(C_{fL},C_w,C_\Delta,C_{\rm count},C_{h0r},C_{\rm gap})$
are positive and the same for all members.  Write $\theta=\varepsilon/100$,
$X_\theta=x^\theta$ (so $L=X_\theta^{100}$), $f=2n_F$, $r_S=m/2-1$,
$N=k+m(t+1)$ with a further rank $t\ge r+5$, $a=xb$,
$X'=\exp(x^{\varepsilon/3})$, $b'=a/X'$, $w''=b'/(xn_F)$, $\mathrm{rd}_E$ the
root discriminant and $v=\mathrm{rd}_E/2$.  At an infinite place $\tau$ the
datum has rationals $v_i(\tau)$ ($0\le i<k-1$), node weights $b_j(\tau)$ and
$e_j(\tau)$ ($1\le j\le m$); each word $u$ has a column parameter
$u_l(u,\tau)$ ($1\le l\le k$), and $\Delta>0$ is the spacing of the words,
$S_c=4k\Delta$.  Some places are \emph{top} places, with integers $a_T,b_T$.
We use the following clauses (all part of the definition of the class):
\begin{itemize}
\item[(L)] $a\le a_T\le a+a/(kx^{13})$, $b'\le b_T$ and
$b_1-b_T>\frac{k-1}2+3$ and $u_1=a_T+\frac{k-1}2$ at top places; $0\le v_i\le S_c$; $S_c\le u_1-v_0$ at
non-top places; $|b_j-e_{j'}|\le\eta$ (node width $\eta>0$) and $b_j\le b^2$;
$e_j\ge w''$, $\ge b$, $\ge b_T+b'/x$ at thin, thick, top places; the auxiliary
integer of the host satisfies $\mathrm{QS}\le4n_Fx^\theta$.
\item[(P)] $C_{k,\rm low}m^5fLx^\varepsilon\le k$; $\log f\le\sqrt{r_S}$;
$kS_c\le w''$; $t\eta\le w''/x$; at non-top places $u_1-v_0\le2S_c$ and
$u_k=v_{k-2}-\frac12$; $b_m\le2b_1$; $b_1\le2b_T$ at top places;
$(NfL)^{C_{\rm count}}\le\Delta$; $\eta/m\le b_{j+1}-b_j$; with
$\mathrm{gap}=C_{\rm gap}km^4+Nm^{10}$, $\mathrm{gap}<v_i-v_{i+1}$,
$\mathrm{gap}<v_{k-2}+b_1$, $\mathrm{gap}<b_{j+1}-b_j$; there is $Y$ with
$(kb)^8\le Y\le g_i\le2Y$ and $\log Y\le C_{YL}L$.
\item[(S)] $x^{2+8\varepsilon-\theta}\le n_F\le x^{2+8\varepsilon+\theta}$;
$x^{-\theta}\sqrt{x\delta_1}\le m\le x^\theta\sqrt{x\delta_1}$ with
$\delta_1=x^{-\varepsilon/8}$; $k\le x^{9/2+10\varepsilon+\theta}$;
$X_\theta m^5fLx^\varepsilon\le k$, $X_\theta m^3L\log\mathrm{rd}_E\le k(\log k)^2$,
$X_\theta L\le f\log k$; $\log q,\log d\le x^{\varepsilon/5+\theta}$ and
$\mathrm{rd}_E\le\frac{\mathrm{QS}}2\sqrt{qd}$ for the auxiliary primes $q,d$
of the host, with $x^{\varepsilon_z/6}\le\log q$, $x^{\varepsilon_z/5-\varepsilon_z/100}\le\log q$ and $x^{\varepsilon_z/5-\varepsilon_z/100}\le\log d$; $q\ne d$ are odd primes with $\max(q,d)\ge16$ (a clause of the host, not checked here); $X_\theta r^{19}f^5L\le\eta/m$; $X_\theta mk\Delta\le w''$,
$X_\theta m\log v\le k$; $\sqrt b\le w''$.
\item[(A)] eventually $\mathrm{rd}_E\ge2$, $d_E=d_F^2$, every $T_u$, $g_0$,
$g_1$ splits completely in $E$, and $T_u=\pi\bar\pi$ with
$\pi\equiv1\bmod \mathrm{QS}\cdot qd(g_0g_1)^2\mathcal O_E$; and the host is
\emph{admissible}: each $g_i$ is prime, split in $E$ and in the ray class field
of the host, and $g_i\equiv1\bmod\mathrm{QS}\cdot qd$.  Since $16\mid\mathrm{QS}$
(a theorem of \cite{CharacterFamily}), this gives $16qd\mid g_i-1$.
\end{itemize}
The inequality $8C_\Gamma\le C_{k,\rm low}$ is not a clause of the class: it
is a Large condition of the closed statements in which it is used
(\S\ref{i3:ssec-setting}).
The words satisfy $u_{l+1}(u,\tau)=v_{l-1}(\tau)-\frac12-\Delta\ell_l(u)$
($1\le l\le k-2$) with letters $\ell_l(u)\in\{0,1\}$ determining $u$.

The \emph{standing hypotheses} at constants $(C_0,C',C_l)$ are the clauses
(1)--(25) and the numerical clauses (N1)--(N14) verified below; those used in
this section are collected in \eqref{i3:eq-standing}.

\begin{proposition}\label{i3:prop-standing-family}
Let $\mathcal F$ be a family of the working class with
$8C_\Gamma\le C_{k,\rm low}$, and assume $\mathrm{TuSize}(\mathcal F)$.  For
every $C_0\ge1$, eventually, every member satisfying the window
$C_0^{-1}\le d_u\Gamma^2/D_*\le C_0$ satisfies the standing hypotheses at
$(C_0,C_{k,\rm low},C_l^{\mathcal F})$.  The size clauses (1)--(25) hold
eventually at every member, without the window.
\end{proposition}

\begin{proof}
Take $x$ so large that $X_\theta\ge\max(2C_l,2,5/\theta)$, beyond the
thresholds of $\mathrm{TuSize}$ and of (A).  Only (N10) below uses the window.  \emph{Scale facts.}
$L>1$, $b\ge2$; $\log y\le y^\theta/\theta$ gives $\log x\le X_\theta/\theta$;
from $X_\theta\ge2$, $5/\theta\le X_\theta$: $5X_\theta/\theta\le X_\theta^2$,
$X_\theta^3+1\le L/2$, $L\ge4$; by (S), $\log k\le5\log x$ (as
$\varepsilon<1/52$), hence $\log k\le L$ and $\log x\le L$, i.e.\ $k,x\le b$;
and $0<w''\le b'\le a\le b^2$, $w''\le b$.

\emph{Size clauses.} (1) $8C_\Gamma\le C'$ is the hypothesis $8C_\Gamma\le C_{k,\rm low}$.  (2) $C'm^5fL\le k$ by (P)
as $x^\varepsilon\ge1$.  (3) $m^4\le x^{4\theta+2-\varepsilon/4}\le n_F\le f$ by
(S), since $5\theta\le8\varepsilon+\varepsilon/4$.  (4) $f\le k$ and (5)
$\log f\le\sqrt{r_S}$ are rank and (P) clauses.  (6)--(7) $C_lL\le f\log k$
and $C_lm^2L\le k$ from $C_l\le X_\theta$ and (S).  (8) $e_j\ge w''$ at every
place (by (L) and $b_T\ge b'\ge w''$, $b\ge w''$), $\eta\le t\eta\le
w''/x\le w''/2$, so $b_j\ge w''/2\ge C_lmk\Delta$ by (S).  (9)--(10)
$C_lm\log v\le k$ and $C_lm^3L\log\mathrm{rd}_E\le k(\log k)^2$ by (S).
(11) By (S) and (L), $\log\mathrm{rd}_E\le\log2+3\log x+X_\theta^{21}
\le1+3X_\theta/\theta+X_\theta^{21}\le X_\theta^{23}$, so
$C_l\log\mathrm{rd}_E\le X_\theta^{24}\le L$.  (12)--(15)
$(kb)^8\le Y\le g_i\le2Y\le2b^{C_{YL}}\le b^{C_{YL}+1}$.  (16) $b_j\le b^2$ by
(L); and $b_j\ge w''/2\ge\sqrt b/2\ge e^{L/2-1}\ge\max(k^{C_l},b^{1/4})$, since
$C_l\log k\le X_\theta^3\le L/2-1$ and $L\ge4$.  (17) $b_m\le2b_1$.  (18)
$u_1\le b^3$: at top places $u_1=a_T+\frac{k-1}2\le2a+b\le b^3$; elsewhere
$u_1\le3S_c\le kS_c\le w''\le b^3$.  (19) $v\ge1$.  (20)
$(kfL)^{C_l}\le(NfL)^{C_{\rm count}}\le\Delta$ as $1\le kfL\le NfL$ and
$C_l\le C_{\rm count}$.  (21) Node gaps:
$C_lr^{19}f^5L\le X_\theta r^{19}f^5L\le\eta/m\le b_{j+1}-b_j$.  (22) Row gaps:
$C_lkm^4+rm^{10}\le\mathrm{gap}$, and $\lambda^{SH}_i-\lambda^{SH}_{i+1}$ is
$v_i-v_{i+1}$, $v_{k-2}+b_1$ or $b_{j+1}-b_j$, each $>\mathrm{gap}$ by (P).
(23) $b^{m/2}\le T_u\le b^{C_{Tu}m}$ is $\mathrm{TuSize}$.  (24)
$C_l\le X_\theta\le m$ by (S), as $2\theta+\varepsilon/16\le\frac12$.  (25)
$C_l>0$.

\emph{Numerical clauses.} (N1) $0\le v_{l-1}\le4k\Delta$.  (N2) Each column
entry $u_{l+1}(u,\tau)$ is $v_{l-1}-\frac12$ for all $\tau$ or
$v_{l-1}-\frac12-\Delta$ for all $\tau$.  (N3) $4k\Delta\le u_1-v_0$: at non-top
places by (L); at top places $u_1-v_0\ge a-S_c\ge(k-1)S_c\ge S_c$ since $a\ge
w''\ge kS_c$.  (N4)--(N5) the non-top clauses of (P).  (N6) At top places
$b_1-(b_T+\frac{k-1}2)\ge3$ and $b_1/2\le b_T$.  (N7) The interlacing of the
column parameters with $v$ and $-b_1$ is part of the definition of the word
data.  (N8) Distinct words have distinct column parameters: equal parameters
give equal letters by (N2) as $\Delta>0$.  (N9) $T_u$ splits completely, and
$T_u\equiv1\bmod(g_0g_1)^2$: writing $\pi=1+N_0\alpha$ we get
$T_u-1=N_0(\alpha+\bar\alpha+N_0\alpha\bar\alpha)$, so $N_0\mid T_u-1$ in
$\mathcal O_E$; the quotient $(T_u-1)/N_0$ is rational and integral over $\Z$, hence in $\Z$.  (N10) $\Gamma_u=\Gamma>1$ and the
window gives $\Gamma^{-2}D_*/C_0\le d_u\le\Gamma^{-2}D_*C_0$.  (N11)--(N13)
splitting of $g_0,g_1$ and $d_E=d_F^2$ are (A).  (N14) $T_u\ne g_0,g_1$, since
$T_u\equiv1$ and $2\le g_i<(g_0g_1)^2$.
\end{proof}

\begin{proof}[Proof of Theorem~\ref{i3:thm-linkS-slots}]
Let $\varepsilon_1$ and, for $\varepsilon\le\varepsilon_1$, the threshold
system $\mathcal K$ be those of Proposition~\ref{i3:prop-conc1}; they include
$8C_\Gamma\le C_{k,\rm low}$.  Let $\mathcal F\in\mathcal C$ be
$\mathcal K$-large with $\mathrm{TuSize}(\mathcal F)$, and apply
Theorem~\ref{i3:thm-linkS-standing} to $\mathcal F$ with $C_0=3$.  Hypothesis
(a) is Proposition~\ref{i3:prop-standing-family}.  Hypothesis (b) is the
conclusion of Proposition~\ref{i3:prop-conc1} for $\mathcal F$, whose count
constant $c_7=\frac1{16}>0$ comes with it.
\end{proof}

\noindent This proof rests on Proposition~\ref{i3:prop-conc1} and its
hypotheses, and on nothing else.


\subsubsection{Proof of Proposition~\ref{i3:prop-V}}\label{i3:ssec-pf-V}

Throughout this subsection indices start at $0$, and $\tau$ runs over the
$f$ infinite places of $F$.

\medskip\noindent\emph{The two degrees.}
The ranks are even, $m\ge2$, $k\ge m+3$, $m\le f\le k$, $r=k+m-1$,
$r_S=\frac m2-1$.  At each place $\tau$ the datum has rationals $v_i(\tau)$
($0\le i<k-1$), node weights $b_j(\tau)$ ($0\le j<m$), archimedean weights
$w_j(\tau)$, and for each word $u$ a \emph{column parameter}
$u_l(u,\tau)\in\Z+\frac12$ ($0\le l<k$), with
(H1) $v_i<u_i$ and $u_{i+1}<v_i$ ($i<k-1$); (H2) $-b_0<u_{k-1}$;
(H3) $b_{j+1}-b_j\ge1$; (H4) $u_0(u,\tau)$, $u_{k-1}(u,\tau)$ do not depend on
$u$ (clauses of the word data, part of the formally verified data bound of
Lemma~\ref{i3:lem-dataB}).  The row parameter $\lambda(\tau)\in\R^r$ and the row signs $\epsilon$ are
$\lambda_a=v_a$ ($a<k-1$), $\lambda_a=-b_{a-k+1}$ ($a\ge k-1$), and
$\epsilon_a=+$ iff $a<k-1$ or $a-k+1$ is even; there are
$p=k-1+\frac m2$ plus signs and $q=\frac m2$ minus signs.  (This is
$\lambda^{SH}$, $\epsilon^{SH}$ of \S\ref{i3:ssec-gap}, re-indexed.)

For $y\in\R^N$ let $\mathrm{vand}(y)=\prod_{a<c}(y_a-y_c)$,
$\Phi(N)=\prod_{i<N}i!$ and $\dim V(y)=\mathrm{vand}(y)/\Phi(N)$ (the Weyl
dimension when $y$ is a strictly decreasing Harish-Chandra parameter).  For a
prime $g$ put $\pi_n(g)=\prod_{i=1}^n(1-g^{-i})$.  The column and row
\emph{local indices} at $g$ are the indices in $GL_k(\mathbb F_g)$ of the upper
unitriangular group, and in $GL_r(\mathbb F_g)$ of the group $R_{k,m}(g)$ of
upper unitriangular matrices whose entries above the diagonal in the last
column of the middle block vanish inside that block, for the partition
$r=r_S+(k+1)+r_S$ of the indices (the entries in the positions $(i,j)$,
$i<j$, with $i,j$ not both in the middle block are free); $I_\Pi$
and $I_\Sigma$ are the products over $g\in\{g_0,g_1\}$ of their $f$-th powers.
The two degrees of the datum are
\[
 D_*=\mathrm{vol}(Y_0)\,I_\Sigma\prod_\tau d_\sigma(\tau),\qquad
 d^*_u=\mathrm{mass}_I(k)\,I_\Pi\prod_\tau\dim V(u(u,\tau)),
\]
where $Y_0$ is the base hermitian space (signature $(k+r_S,r_S+1)$, rank $r$),
$\mathrm{vol}(Y_0)>0$ its volume, $\mathrm{mass}_I(M)>0$ the mass of the genus
of the standard unimodular hermitian lattice of rank $M$, and $d_\sigma(\tau)>0$
the \emph{$\sigma$-degree} at $\tau$: the formal degree, for the Haar measure
fixed by the datum, of a discrete series of $U(p,q)$ with Harish-Chandra
parameter $(\epsilon,\lambda(\tau))$.  Then $\log X_u=\frac1{fk}\log(D_*/d^*_u)$.

\medskip\noindent\emph{Three inputs.}

\begin{proposition}[formally verified, \cite{CharacterFamily}]\label{i3:prop-massratio}
Let $E/F$ be a CM field unramified at all finite primes, $f=[F:\Q]$, $v=|d_E|^{1/(2f)}/2$, and for $M\ge1$ let
$\mathrm{mass}_I(M)$ be the mass of the standard lattice $O_E^M$ with the form $\sum_ix_i\bar y_i$ (the sum of
$|\mathrm{Aut}|^{-1}$ over the classes of its genus).  For $M\ge3$,
\[
\frac{\mathrm{mass}_I(M)}{\mathrm{mass}_I(M-1)}=R(E,M):=
\Bigl(\frac{(M-1)!}{\pi^M}\Bigr)^f\frac{v^{fM}}{(2v)^{f/2}}L(E,M),
\]
with $L(E,M)=\prod_{\mathfrak p}(1-\eta(\mathfrak p)^MN\mathfrak p^{-M})^{-1}>0$,
$\eta$ the quadratic character of $E/F$.  Moreover, if a member of the working class has the \emph{volume tie}
\begin{itemize}
\item[(VT)] for every non-degenerate hermitian space $W$ over $E$ satisfying the signature condition of the datum, with
constant signature $(P,Q)$, and every fundamental domain $\mathcal F_W$ for $U_W(F)$ in $U_W(\mathbb A)$, the measure fixed by the
datum (compact-dual normalisation at infinity) gives $\mu(\mathcal F_W)=\mathrm{vol}(W)$, the volume recorded by the datum,
\end{itemize}
then it satisfies $\mathrm{vol}(Y_0)=\mathrm{mass}_I(r)$.
\end{proposition}

\noindent The first statement is the Siegel mass ratio, the second is the mass ratio for the datum, and
the third is that the base volume equals the mass under (VT) (see the paragraph after this one).  In the formal
statement $L(E,M)$ is written as a product over the primes $\mathfrak P$ of $E$, with factor $(1-(-1)^MN\mathfrak P^{-M/2})^{-1}$
if $\bar{\mathfrak P}=\mathfrak P$ and $(1-N\mathfrak P^{-M})^{-1/2}$ otherwise; grouping the primes of $E$ above each $\mathfrak p$
(inert: $N\mathfrak P=N\mathfrak p^2$, $\eta(\mathfrak p)=-1$; split: two conjugate primes of norm $N\mathfrak p$, $\eta(\mathfrak p)=1$;
no ramified primes) gives the product above.  The first statement is the inductive form of the Siegel mass formula,
proved by splitting off a rank-one sublattice spanned by a vector of prescribed totally positive norm and factoring the
local densities; the splitting uses that every binary form $\mathrm{diag}(a,b)$ with $a,b$ totally positive represents every
totally positive $n\in F$, which is the consequence of
Theorem~\ref{i3:thm-landherr} recorded after it.  It is a ratio of two instances of the exact mass formula for unimodular
hermitian lattices (Gan--Hanke--Yu, Duke Math.\ J.\ 107 (2001) 103--133; reference to be confirmed); since we have not matched the normalisations, we take it from the formal verification.

\emph{The second statement and where (VT) comes from.}  Under (VT) two formally verified, unconditional facts give
$\mathrm{vol}(Y_0)=\mathrm{mass}_I(r)$: (1) the \emph{volume ratio}: for $W$ of rank $\ge3$ and $W'\subset W$ the orthogonal complement
of a vector of norm $\pm1$, $\mu(\mathcal F_W)/\mu(\mathcal F_{W'})=R(E,\mathrm{rank}\,W)$ (a consequence of the rank-one Siegel--Weil
computation); (2) the \emph{rank-two genus mass}: the mass of the genus of a totally positive definite space of rank $2$ is
$\mathrm{mass}_I(2)$.  From (VT) and (2), $\mathrm{vol}(W)=\mathrm{mass}_I(2)$ for definite $W$ of rank $2$.  The proof is then a two-stage
induction: first on the rank of totally definite spaces, from rank $2$ upwards, each step splitting off a vector of norm $1$
and using (1); then, for the indefinite signature $(k+r_S,r_S+1)$ of $Y_0$, on the number of negative signs, each step
splitting off a vector of norm $-1$ and using (1) again.  This gives $\mathrm{vol}(W)=\mathrm{mass}_I(\mathrm{rank}\,W)$ along the way, in particular
$\mathrm{vol}(Y_0)=\mathrm{mass}_I(r)$.  Condition (VT) itself is not proved here: it is the volume clause of the
archimedean (real-place) tie of the datum, and holds eventually at the members of every family on which the real-place
second-layer ties are closed, which is field (vii) of Definition~\ref{i4:def-hyps}.  Statements that use the second
assertion (Propositions~\ref{i3:prop-V} and \ref{i3:prop-part5}, Theorem~\ref{i3:thm-main}(d)) therefore assume (VT).  Its formally verified steps are outlined in \S\ref{i3:str-mass}.

\emph{The Haar measures.}  A \emph{compact-dual normalisation} is a family of Haar measures $\mu_{P,Q}$ on the groups
$U(P,Q)$ such that $\mu_{N,0}$ is the probability Haar measure of the compact group $U(N)$ and, for all $P,Q$, the densities
of $\mu_{P,Q}$ and of $\mu_{P+Q,0}$ at the identity, with respect to one and the same Lebesgue measure on the Lie algebras
transported through $X\mapsto X+iY$ coordinates, coincide.  The datum fixes such a family, and its archimedean scale tie says that
the Haar measure of the datum on $U_\epsilon\cong U(p,q)$ is $\mu_{p,q}$ transported by the coordinate permutation of
$\epsilon$.  For a discrete series $\sigma$ and $w\ne0$ the formal degree is
$d(\sigma)=\|w\|^4/\int|\langle w,\sigma(g)w\rangle|^2\,d\mu(g)$ (independent of $w$ by Schur orthogonality,
Theorem~\ref{i3:thm-godement}).

\begin{proposition}[formally verified, \cite{CharacterFamily}]\label{i3:thm-HC-fdeg}
Let $N\ge2$, $\epsilon$ a sign pattern with $p\ge1$ plus signs and $q\ge1$ minus signs, $\mu$ the Haar measure on
$U_\epsilon$ given by a compact-dual normalisation (as above), and $\lambda$ a Harish-Chandra parameter with gap
$h>32Npq+2$ for $\epsilon$.  Then a discrete series with parameter $(\epsilon,\lambda)$ (Definition~\ref{i3:def-ds}) exists,
and every discrete series $\sigma$ with parameter $(\epsilon,\lambda)$ and every $w\ne0$ satisfy
\[
 \frac{\|w\|^4}{\int_{U_\epsilon}|\langle w,\sigma(g)w\rangle|^2\,d\mu(g)}=\prod_{a<c}\frac{|\lambda_a-\lambda_c|}{c-a}=\dim V(\lambda).
\]
\end{proposition}

\noindent Existence needs only the gap $>\max(2N,8N\min(p,q))$, which is implied by $h>32Npq+2$.  The
last equality is Weyl's dimension formula for the irreducible representation $V(\lambda)$ of $U(N)$ with highest weight
$\lambda-\rho$.

\begin{remark}[Comparison with Harish-Chandra's formula]\label{i3:rem-HC-fdeg}
For $G=U(p,q)$ with the diagonal torus as compact Cartan subgroup, the positive roots are $e_a-e_c$ ($a<c$) and
$\langle\rho,e_a-e_c\rangle=c-a$.  Harish-Chandra's formal degree formula (Harish-Chandra, \emph{Discrete series~II}, Acta
Math.\ 116 (1966), stated there for connected semisimple groups with finite centre, with the formal degree determined
up to a positive constant depending on the Haar measure; its extension to reductive groups of Harish-Chandra's class
such as $U(p,q)$: reference to be confirmed), in the normalisation in which the Haar measure of $G$ is compatible with that of
the compact dual $U(p+q)$ of total mass $1$, gives $d(\pi_\lambda)=\dim V(\lambda)$ (this value of the constant in the
compact-dual normalisation: e.g.\ Atiyah--Schmid, Invent.\ Math.\ 42 (1977) 1--62, erratum 54 (1979) 189--192, reference to be confirmed; it can be checked directly for $SL_2(\R)$ and is not used in any proof)
for the discrete series $\pi_\lambda$ of Harish-Chandra parameter $\lambda$, for every regular $\lambda$.  Combined with the
characterisation of $\pi_\lambda$ by its lowest $K$-type $\lambda+\rho_n-\rho_c$ (Schmid; Hecht--Schmid; Vogan), which
identifies the representations of Definition~\ref{i3:def-ds} with $\pi_\lambda$, this would give
Proposition~\ref{i3:thm-HC-fdeg} for all gaps.  So the formal statement is a special case (gap $>32Npq+2$) of a
classical formula.  We nevertheless take it from the formal verification and do not cite it: we have not matched the exact
statements of the lowest-$K$-type characterisation and of the measure normalisation in the sources with the definition of discrete series and the Haar normalisation used here; the formal proof does not use them (it constructs a model of the discrete
series and computes the formal degree directly).  Its formally verified steps are outlined in \S\ref{i3:str-HCfdeg}.
\end{remark}

\noindent In particular $d_\sigma(\tau)$ does not depend on the choice of the discrete series with parameter
$(\epsilon,\lambda(\tau))$.  Hence:

\begin{lemma}\label{i3:lem-sigdeg}
Assume that $\lambda(\tau)$ is a Harish-Chandra parameter (Definition~\ref{i3:def-gap}: strictly decreasing, with
$\lambda_a(\tau)-\frac{r-1}2\in\Z$) and that $d_\sigma(\tau)$ is the formal degree for the compact-dual normalisation of the Haar
measure on $U_\epsilon$; both are records of the datum (its Harish-Chandra data and its archimedean scale tie, with the tie of
$d_\sigma$ to the formal degree).  If $\lambda(\tau)$ has gap at least $h_*=8192\,q(pq)^3\log r\cdot9$ for
$\epsilon$, then $d_\sigma(\tau)=\dim V(\lambda(\tau))$ and
$d_\sigma(\tau)\le d_\tau(\lambda)|\lambda|^{pq}$, where $|\lambda|=\lambda_0-\lambda_{r-1}$ and
\[
d_\tau(\lambda)=\frac{\prod_{a<c,\ \epsilon_a=\epsilon_c}(\lambda_a-\lambda_c+n(a,c))}{\Phi(p)\Phi(q)},\qquad
n(a,c)=\#\{a<e<c:\epsilon_e\ne\epsilon_a\}.
\]
\end{lemma}

\begin{proof}
$h_*>32rpq+2$ (as $r\le2pq$), so Proposition~\ref{i3:thm-HC-fdeg} (with
$N=r$; here $p,q\ge1$) gives the equality.  Split $\mathrm{vand}(\lambda)$ into same-sign and opposite-sign
pairs: each same-sign factor is $\le\lambda_a-\lambda_c+n(a,c)$, each of the $pq$
opposite-sign factors is $\le|\lambda|$; and $\Phi(p)\Phi(q)\le\Phi(p+q)$.
\end{proof}

\begin{lemma}[formally verified, \cite{CharacterFamily}]\label{i3:lem-dataB}
Eventually, every member of the working class satisfies, with
$C_0=C_Y=\max(5,C_{YL}+1)$, $\Delta$ the spacing of its words and $w$ its
archimedean weights, the following \emph{data bound}:
(B1) $\pi\le v$, $\log(2v)\le L$; (B2) $f,k\le x^{C_0}$, $m^2L\le k$,
$\log(2k)\le L$, $44\le m\le\min(k,x)$; (B3) $(kb)^8\le g_i$ and
$\log g_i\le C_0L$; (B4) $\Delta\ge1$, $0\le v_i\le4k\Delta$,
$v_i-v_{i+1}=4\Delta$; (B5) for $1\le l\le k-2$,
$u_l\in\{v_{l-1}-\frac12,\ v_{l-1}-\frac12-\Delta\}$; (B6) at top places
$u_0-u_{k-1}\le2a$, $4k\Delta\le u_0-v_0$ and $u_{k-1}\le-4k\Delta$; at other
places $4k\Delta\le u_0-v_0\le8k\Delta$ and $u_{k-1}=v_{k-2}-\frac12$; (B7)
$b_j<w_j\le(1+\frac1x)b_j$, $w_j<b_{j+1}$, $4k\Delta\le b_0/x$, $w_{m-1}\le2b$;
(B8) $b_{j+1}-b_j\ge r^{19}f^5L/C_0$.
\end{lemma}

\noindent\emph{What the proof involves.}  Most clauses are clauses of the
records (L), (P), (S) of \S\ref{i3:ssec-pf-slots} or follow from them by the
scale facts proved there, but two do not: (B1) $\pi\le v$ uses a lower bound
for the root discriminant carried by the host field
($\mathrm{rd}_E\ge2^{e+1-2^{-e}}\sqrt d$ with $e\ge2$, $d\ge7$), and the
equality $v_i-v_{i+1}=4\Delta$ in (B4) is a structural clause of the word data
(the records above only give $\mathrm{gap}<v_i-v_{i+1}$).  The interlacing
conditions (H1)--(H4) are likewise clauses of the word data.  For the rest:
(B1) and $\log(2k)\le L$ from (S) via $2v=\mathrm{rd}_E\le
\frac{\mathrm{QS}}2\sqrt{qd}$ and the bounds on $q,d$; $44\le m\le x$ from
$m\ge x^{1/4}$ and $m\le x^\theta\sqrt x$; (B3) from $(kb)^8\le Y\le g_i\le2Y$,
$\log Y\le C_{YL}L$; (B4)--(B6) are the word clauses and (N1)--(N6); (B7) is
(L) with $\eta\le w''/x$; (B8) is the node gap clause (21).  The comparisons
are of the same kind as in the proof of Proposition~\ref{i3:prop-standing-family},
and are carried out in \cite{CharacterFamily}; we do not reproduce them.

\medskip\noindent\emph{Elementary facts.}

\begin{lemma}\label{i3:lem-V-elem}
Under the data bound: $b_j\ge x$, $w_j\ge1$, $w$ is increasing and
$w_j\le2b$; $u(u,\tau)$ is strictly decreasing with consecutive differences
$\ge1$; $\lambda(\tau)$ is strictly decreasing with all entries in
$[-b_{m-1},v_0]$, and has gap at least $h_0=r^{19}f^5L/C_0$ for $\epsilon$.
\end{lemma}

\begin{proof}
$4k\Delta\ge1$ and $4k\Delta\le b_0/x$ give $b_0\ge x$; (B7) and induction give
the rest of the first sentence.  (H1) interlaces $u$ with $v$; entries lie in
$\Z+\frac12$.  For the gap: if $\epsilon_a\ne\epsilon_c$, $a<c$, then
$c\ge k$ (as $\epsilon_{k-1}=+$); with $j=c-k+1$, either $a<k-1$ (then $j$ is odd) and
$\lambda_a-\lambda_c=v_a+b_j\ge b_j-b_0\ge h_0$, or $a=k-1+i$ and
$\lambda_a-\lambda_c=b_j-b_i\ge(j-i)h_0$, by (B8).
\end{proof}

\begin{lemma}\label{i3:lem-V-index}
For a prime $g$, the column index is $g^{k(k+1)/2}\pi_k(g)$ and the row index
is $g^{k(k+1)/2+m(2k+m-1)/2}\pi_r(g)$; and
$\frac12\le1-2g^{-(k+1)}\le\pi_r(g)/\pi_k(g)=\prod_{j=k+1}^r(1-g^{-j})\le1$.
Hence
$I_\Sigma/I_\Pi=(g_0g_1)^{fm(2k+m-1)/2}\prod_{i}\prod_{j=k+1}^r(1-g_i^{-j})^f$.
\end{lemma}

\begin{proof}
$|GL_n(\mathbb F_g)|=g^{n^2}\pi_n(g)$.  The unitriangular group has order
$g^{k(k-1)/2}$.  $R_{k,m}(g)$ consists of the matrices agreeing with the
identity off the free positions, namely the pairs $i<j$ not both in the middle
block ($r_S(2k+1+2r_S)$ of them) and the pairs inside the middle block avoiding
its last column ($k(k-1)/2$); subtracting from $r^2$ and using $m=2r_S+2$ gives
the exponent.  For the ratio, with $y=g^{-1}\le\frac12$, induction on $n$ gives
$\prod_{j=k+1}^{k+n}(1-y^j)\ge1-2y^{k+1}+2y^{k+n+1}$: the step multiplies by
$1-s$, $s=y^{k+n+1}$, and $(A+2s)(1-s)\ge A+2ys$ for $A=1-2y^{k+1}$, $y\le\frac12$.
\end{proof}

\begin{lemma}[The identity for $\log X_u$]\label{i3:lem-V-identity}
For every word $u$, with $\ell_v=\log\frac v\pi$,
\[
\begin{aligned}
fk\log X_u={}&f\Bigl((m-1)k+\tfrac{m(m-1)}2\Bigr)\ell_v-f\tfrac{m-1}2\log(2v)
+\sum_{M=k+1}^r\log L(E,M)+f\tfrac{m(2k+m-1)}2\log(g_0g_1)\\
&+f\sum_{i=0,1}\sum_{j=k+1}^r\log(1-g_i^{-j})
+\sum_\tau\log\frac{d_\sigma(\tau)\Phi(r)}{\mathrm{vand}(u(u,\tau))}.
\end{aligned}
\]
\end{lemma}

\begin{proof}
By Proposition~\ref{i3:prop-massratio} and (VT), $\mathrm{vol}(Y_0)=\mathrm{mass}_I(r)
=\mathrm{mass}_I(k)\prod_{M=k+1}^rR(E,M)$, so
$D_*/d^*_u=\prod_MR(E,M)\cdot(I_\Sigma/I_\Pi)\cdot\prod_\tau
d_\sigma(\tau)/\dim V(u(u,\tau))$, all factors positive.  Take logarithms; use
Lemma~\ref{i3:lem-V-index}, $\sum_{M=k+1}^{r}\log(M-1)!=\log\Phi(r)-\log\Phi(k)$,
$\sum_{M=k+1}^rM=(m-1)k+\frac{m(m-1)}2$, and
$\dim V(u)=\mathrm{vand}(u)/\Phi(k)$: the terms $\pm f(\log\Phi(r)-\log\Phi(k))$
cancel.
\end{proof}

\begin{lemma}\label{i3:lem-V-Lvalue}
For $M\ge2$: $|\log L(E,M)|\le f\log\zeta(M)\le f$.
\end{lemma}

\begin{proof}
Each Euler factor at $\mathfrak p$ lies between $(1+N\mathfrak p^{-M})^{-1}$ and
$(1-N\mathfrak p^{-M})^{-1}$, and above each rational prime $p$ there are at most
$f$ primes $\mathfrak p$, each with $N\mathfrak p\ge p$; so
$|\log L(E,M)|\le f\sum_p|\log(1-p^{-M})|=f\log\zeta(M)$.  Finally
$\zeta(M)\le\zeta(2)<2<e$.
\end{proof}

\medskip\noindent\emph{Bounds at one place.}

\begin{lemma}[Lower bound at a place]\label{i3:lem-V-placelow}
Assume the data bound, $2(\log8+1+\log x)\le L$ and that $\lambda(\tau)$ has gap
$\ge h_*$.  Then for every word $u$,
\[
\log\frac{d_\sigma(\tau)\Phi(r)}{\mathrm{vand}(u(u,\tau))}\ge
(k-1)\sum_{j<m}\log b_j(\tau)-5kL.
\]
\end{lemma}

\begin{proof}
By Lemma~\ref{i3:lem-sigdeg}, $d_\sigma\Phi(r)=\mathrm{vand}(\lambda)$.
\emph{Row parameter:} adding the entries $-b_0,\dots,-b_{m-1}$ one at a time
multiplies $\mathrm{vand}$ by $\prod_i(v_i+b_{j})\prod_{j'<j}(b_j-b_{j'})\ge
b_j^{k-1}$, so $\mathrm{vand}(\lambda)\ge\mathrm{vand}(v)\prod_jb_j^{k-1}$.
\emph{Column parameter:} all differences of $u$ are at most $D=8a$ (by (B6),
using $4k\Delta\le b_0\le w_{m-1}\le2b$ at non-top places).  Write
$\mathrm{vand}(u)$ as the product of the factors with first index $0$ ($k-1$
factors), with last index $k-1$ ($k-2$ further factors), and the rest; each of
the former is $\le D$, and by interlacing $0<u_i-u_j\le v_{i-1}-v_j$ for
$1\le i<j\le k-2$, so the rest is $\le\mathrm{vand}(v)$ (the omitted
consecutive differences of $v$ are $\ge1$).  Hence
$\mathrm{vand}(u)\le D^{2k}\mathrm{vand}(v)$ and
$\log\frac{\mathrm{vand}(\lambda)}{\mathrm{vand}(u)}\ge(k-1)\sum_j\log b_j
-2k(\log8+\log x+L)$.  The hypothesis on $L$ bounds the last term by $5kL$.
\end{proof}

\begin{lemma}[Upper bound at a place]\label{i3:lem-V-placeup}
Under the hypotheses of Lemma~\ref{i3:lem-V-placelow}, for every word $u$,
$\log d_\sigma(\tau)-\log\dim V(u(u,\tau))\le km(5\log(2k)+5L+12)$.
\end{lemma}

\begin{proof}
By Lemma~\ref{i3:lem-sigdeg}, $d_\sigma\le d_\tau(\lambda)|\lambda|^{pq}$, and
$|\lambda|=v_0+b_{m-1}\le2b_{m-1}\le4e^L$.  In the numerator of $d_\tau$ the
pairs $a<c<k-1$ give $\mathrm{vand}(v)$ (there $n(a,c)=0$); every other row
contributes at most $m$ factors, each $\le\mathrm{Bd}=v_0+b_{m-1}+2k\le8ke^L$.
So $d_\tau\le\mathrm{vand}(v)\mathrm{Bd}^{mr}/(\Phi(p)\Phi(q))$.  Interlacing
gives $\mathrm{vand}(v)\le\mathrm{vand}(u)$, and $\Phi(k)\le\Phi(p)$.  Hence
$\log d_\sigma-\log\dim V(u)\le mr\log\mathrm{Bd}+pq\log(4e^L)$, and
$mr\le2km$, $pq\le km$, $\log\mathrm{Bd}\le\log(2k)+L+2$ give the claim.
\end{proof}

\medskip\noindent\emph{Proof of (V1).}

Fix $x$ beyond the thresholds of Lemma~\ref{i3:lem-dataB} and so large that
$L\ge9$, $\log x\le L/4$, $L\ge4(\log8+1)$ and
$L\ge9\cdot8192\,C_0$.  By Lemma~\ref{i3:lem-V-elem},
$\lambda(\tau)$ has gap $\ge h_0=r^{19}f^5L/C_0\ge9\cdot8192\,r^9\ge h_*$
(since $q(pq)^3\log r\le r^8$), so Lemmas~\ref{i3:lem-V-placelow},
\ref{i3:lem-V-placeup} apply at every place.

\emph{Lower bound.}  Write $S_W=\sum_\tau\sum_{j<m}\log w_j(\tau)$.  By (B7),
$\log w_j-\log b_j\le\frac1x$ and $\log w_j\le\log(2b)\le2L$, so
$k\log w_j-(k-1)\log b_j\le\frac kx+2L$.  With Lemma~\ref{i3:lem-V-placelow},
summing over $j<m$ and using $m\le x$, $m\le k$,
$\sum_\tau\log\frac{d_\sigma\Phi(r)}{\mathrm{vand}(u)}-kS_W\ge-f(5kL+k+2kL)$.
The $L$-values contribute $\ge-f(m-1)$ (Lemma~\ref{i3:lem-V-Lvalue}), the local
factors $\ge-2f$ (Lemma~\ref{i3:lem-V-index}).  Since
$\log\frac{w_jv}\pi=\log w_j+\ell_v$,
$fk\,\mathrm{low}_S=kS_W+fkm\ell_v+fkm\log(g_0g_1)-11fkL$.  Subtracting this
from Lemma~\ref{i3:lem-V-identity}, and using $0\le\ell_v\le L$,
$\log(2v)\le L$, $\log(g_0g_1)\ge0$:
\[
fk(\log X_u-\mathrm{low}_S)\ge f\Bigl(11kL-kL-\tfrac{m-1}2L-(m-1)-2-1-5kL-k-2kL\Bigr)
\ge f\Bigl(3kL-\tfrac m2L-m-k-3\Bigr)\ge0,
\]
as $m\le k$ and $L\ge9$.  So $\mathrm{low}_S\le\log X_u$.

\emph{Upper bound.}  In Lemma~\ref{i3:lem-V-identity}: the first term is
$\le f((m-1)k+\frac{m(m-1)}2)L$; the second is $\le0$; the $L$-values give
$\le fm$; $\log(g_0g_1)\le2C_YL$ by (B3); the local factors are $\le0$; and by
Lemma~\ref{i3:lem-V-placeup} and $\log\Phi(r)-\log\Phi(k)\le(m-1)r\log r\le2kmL$,
each place gives $\le km(10L+12)+2kmL$.  Using $m-1\le k$, $m\le kL$, and
$12\le\frac43L$, the total is $\le fk(15+3C_Y)mL$.\qed

\medskip\noindent\emph{Proof of (V3).}

\begin{lemma}\label{i3:lem-V-sep}
Let $\Delta>0$, $y,y'\in\R^k$ with: for $0<c<k$,
$y_0-y_c,\ y'_0-y'_c\ge4k\Delta$ and $|(y_0-y_c)-(y'_0-y'_c)|\le\Delta$; for
$1\le a<c<k$, $y_a-y_c,\ y'_a-y'_c\ge3(c-a)\Delta$ and
$|(y_a-y_c)-(y'_a-y'_c)|\le2\Delta$.  Then
$|\log\mathrm{vand}(y)-\log\mathrm{vand}(y')|\le\frac14+\frac23k(1+\log k)$.
\end{lemma}

\begin{proof}
For $X,X'\ge c_0>0$, $|\log X-\log X'|\le|X-X'|/c_0$.  The pairs with $a=0$
contribute at most $k\cdot\frac{\Delta}{4k\Delta}=\frac14$; a pair $a<c$ with
$a\ge1$ contributes $\le\frac23\frac1{c-a}$, and
$\sum_{c>a}\frac1{c-a}\le1+\log k$.
\end{proof}

For two words $u,u'$ and a place $\tau$, the hypotheses of
Lemma~\ref{i3:lem-V-sep} hold for $y=u(u,\tau)$, $y'=u(u',\tau)$: by (B4), (B5)
each $u_l$ ($1\le l\le k-2$) lies in the interval
$[v_0-4\Delta(l-1)-\frac12-\Delta,\ v_0-4\Delta(l-1)-\frac12]$ for both words,
$u_0,u_{k-1}$ are common (H4), and $u_0-v_0\ge4k\Delta$, $u_{k-1}\le
v_0-4\Delta(k-2)-\frac12$ by (B6).  Now
$d^*_u=K_0\prod_\tau\mathrm{vand}(u(u,\tau))$ with $K_0$ independent of $u$, so
\[
|\log X_u-\log X_{u'}|\le\frac1{fk}\sum_\tau|\log\mathrm{vand}(u(u,\tau))-\log\mathrm{vand}(u(u',\tau))|
\le\frac1{4k}+\frac23(1+\log k)\le1+\log k+\frac2k.\qed
\]

\medskip\noindent\emph{Proof of (V4).}
The records of the working class carry $\Gamma_u=\Gamma=\exp(C_\Gamma m^5fL)$
for every word and the size clause $8C_\Gamma m^5fL\le k$.  Hence
$\log\Gamma_u\le k/8\le fk$.\qed

\noindent This proves Proposition~\ref{i3:prop-V}, with $C_Y=\max(5,C_{YL}+1)$,
modulo Propositions~\ref{i3:prop-massratio} and \ref{i3:thm-HC-fdeg} and
Lemma~\ref{i3:lem-dataB}, all formally verified.

\subsubsection{The counting theorem}\label{i3:ssec-counting}

Both propositions use the same count of automorphic forms.  An
\emph{automorphic counting problem} $I$ consists of a CM extension $E/F$ with
$f=[F:\Q]$ and root discriminant $\mathrm{rd}$, a hermitian space of
dimension $n$ over $E$ with the same signature $(p,q)$ at every infinite place
$\tau$ of $F$, discrete series representations $\sigma_\tau$ of the real
unitary groups with Harish-Chandra parameters $\lambda_\tau$ and formal degrees
$d_\tau$, a fundamental domain $\mathcal F_I$ of Haar measure $\mu(\mathcal F_I)$,
a level $g$, an auxiliary prime $g'$, and an admissible Hecke idempotent $e$
of level $g$ (locally constant, idempotent, hermitian, of positive type,
supported in the stabiliser $K$ of a self-dual lattice, bi-invariant under the
principal congruence subgroup $K(g)$, and supported on elements unipotent
modulo $g'$).  Its \emph{count} $D_{\rm count}(I)$ is the dimension of the
space of intertwiners from $\bigotimes_\tau\sigma_\tau$ into the
$e$-part of the discrete spectrum, and its \emph{expected value} is
\[
 D^*(I)=\Bigl(\prod_\tau d_\tau\Bigr)\,\mu(\mathcal F_I)\,\mathrm{Re}\,e(1).
\]
For $n_0\in\N$, $A,b,h\in\R$ the problem satisfies $\mathrm{CH}(n_0,A,b,h)$ if:
$n_0\le n$; $q^2\le n$; $2\le b$;
$\log\log\mathrm{rd}\le n$; $|\lambda_{\tau,a}|\le b^2$;
$1\le g\le b^A$; $g'$ is prime, $4n<g'$ and $g'$ splits completely in $E$;
each $\lambda_\tau$ is a Harish-Chandra parameter (Definition~\ref{i3:def-gap})
with gap at least $h$ (between entries of different signature);
$\log f\le\sqrt n$.  (Only constant signatures are covered by the formal
statement below.)

\begin{proposition}[formally verified, \cite{CharacterFamily}]\label{i3:prop-counting}
Assume the unitary Siegel--Weil formula in the convergent range.  For every
$A\in\R$ there are $n_0(A)\in\N$ and $C_0(A)\ge0$ such that every automorphic
counting problem whose idempotent is $e=\overline{\chi}\mathbf 1_{K'}/\mu_f(K')$ for an open compact $K'\subset K$ and a
unitary character $\chi$ of $K'$ (the level conditions verified in the proof of Proposition~\ref{i3:prop-rowtie}),
which satisfies the instance tie (it is the instance of a datum with the volume tie, (Inst)(i), (Inst)(ii) and the
archimedean scale tie), with $q\ge1$, and satisfying $\mathrm{CH}(n_0(A),A,b,h)$ with
$h\ge C_0(A)\,n^{19}f^5\log(b\,\mathrm{rd})$ and
$\log\log b\le\frac1{16}nf\log n$ has
\[
 |D_{\rm count}(I)-D^*(I)|\le e^{-\frac1{16}nf\log n}\,D^*(I).
\]
\end{proposition}

\noindent\emph{What remains.}  This is the trace-formula count of the
discrete series of a definite or indefinite unitary group with large
Harish-Chandra gap: the kernel $\sum_\gamma\phi_e(x^{-1}\gamma x)$ built from
matrix coefficients of lowest $K$-type vectors is integrated over
$\mathcal F_I$, the identity contributes $D^*(I)$, and the remaining
contributions are bounded by the decay of these coefficients under the gap,
volume bounds for the hermitian lattices of small covolume (derived from the
Siegel--Weil formula), and the rarity of lattices with short sublattices.  It
occupies the volume-counting part of \cite{CharacterFamily}, a long verification;
no published statement with these explicit constants is known to us
(reference to be confirmed).  Its formally verified steps are outlined in \S\ref{i3:str-count}.

\subsection{The irreducibility defect}\label{i3:ssec-grp-irr}

We prove Proposition~\ref{i3:prop-irr-row}.

\subsubsection{Proof of Proposition~\ref{i3:prop-irr-row}}\label{i3:ssec-pf-irr}

\medskip\noindent\emph{The local weights.}
Fix $k\in\{0,1\}$ and the prime $g=g_k$.  (In this subsection and in
Definition~\ref{i3:def-allrows} the index $k\in\{0,1\}$ of $g_k$,
$\mathrm{lab}_k$, $\mathrm{share}_k$, $\mathrm{cyl}_k$ is unrelated to the rank
$k$ of the datum, which appears here only through $r=k+m-1$.)  The datum fixes a set $\Pi_k$ of $f$
places of $E$ above $g$, a set $J$ with $|J|=m$ and characters
$\chi_k:\Pi_k\times J\to\mathbb F_g^\times$.  A \emph{label} is a monic
polynomial of degree $r$ over $\mathbb F_g$ with non-zero constant term; for
$c\in\mathbb F_g^r$ let $P_c=X^r+\sum_{i<r}c_iX^i$.  For a nonzero
polynomial $P$ with monic irreducible factors $\Phi_j$ (with multiplicity) put
$\mathrm{sq}(P)=\prod_j(g^{\deg\Phi_j}-1)^{-1}$.  For a nontrivial additive
character $\psi$ of $\mathbb F_g$ the datum has two elements
$e_U(\psi),e_J(\psi)$ of the group algebra of $GL_r(\mathbb F_g)$, each
defining a weight $\mathrm{wt}(e,\cdot)$ on monic polynomials of degree $r$,
vanishing off labels.

\begin{definition}[Good tuples]\label{i3:def-good}
A tuple $(P_v)_{v\in\Pi_k}$ of monic polynomials of degree $r$ over
$K=\mathbb F_g$ is \emph{good} if (G1) each $P_v(0)\ne0$; (G2) each $P_v$ is
squarefree; (G3) for every prime $l\notin B(g,r)$, where
$B(g,r)=\{g\}\cup\{l:\exists d\le r,\ g^d<l^{4v_l(g^d-1)}\}$, the prime-to-$l$
part is injective on the roots of each $P_v$ in $\bar K$; (G4) for every prime
$l\ne g$, all $v,j$ and every root $\theta$ of $P_v$, the prime-to-$l$ parts of
$\theta$ and of $\chi_k(v,j)^{\pm1}$ differ; (G5) for every $a$ with
$1\le a\le r/2$ some $P_v$ has no divisor of degree $a$; (G6) some $P_v$ has an
irreducible factor of prime degree $p$ with $r<2p$ and $p+3\le r$.  A label
tuple is \emph{bad} if it is not good.
\end{definition}

\begin{proposition}\label{i3:prop-W}
Let $K$ be a finite field with $|K|=g\ge2^{12}$, $N\ge100$, $q\ne d$ odd primes
with $\max(q,d)\ge16$ and $16qd\mid g-1$, $\iota,J$ finite sets,
$\chi:\iota\times J\to K^\times$, and $P\ge0$ a function on $K^N$ with $\sum_cP(c)=1$ and
$P(c)=\mathrm{sq}(P_c)$ whenever $P_c$ is squarefree with $P_c(0)\ne0$.  Then
the $P^{\otimes\iota}$-mass of the tuples $(c_v)_{v\in\iota}$ such that $(P_{c_v})_v$ is not good for $\chi$ is at
most
\[
 \frac{2|\iota|}g+\frac{23|\iota|\log g}{\sqrt g}
 +2|\iota||J|\Bigl(\frac1{4\min(q,d)}+\frac{1+\log N}{g-1}\Bigr)
 +\frac N2\Bigl(\frac78\Bigr)^{|\iota|}+e^{-|\iota|/(5\log N)}.
\]
\end{proposition}

\noindent This is exactly the formal statement: there the
divisibility is written as $16\mid g-1$, $q\mid g-1$, $d\mid g-1$, which for distinct odd primes $q,d$ is equivalent to
$16qd\mid g-1$; goodness has six clauses, (G1)--(G6) of Definition~\ref{i3:def-good} (the set
$B(g,r)$ consists of primes: a prime $l$ with $g^d<l^{4v_l(g^d-1)}$ for some $1\le d\le r$ divides $g^d-1$, so the bound
$l\le g^r$ in $B(g,r)$ is automatic).  The abstract statement is itself formally verified, and it is the one
used in the proof of Proposition~\ref{i3:prop-irr-row}, where the goodness of a datum tuple
(goodness for $\chi$ and for $\chi^{-1}$) is the conjunction of (G1)--(G6), with (G4) for
$\chi^{\pm1}$.  The proof (\S\ref{i3:ssec-pf-W}) is a union bound over five local estimates.

\begin{definition}[Green's parametrisation]\label{i3:def-green}
For a prime $g$ fix an injective homomorphism $\theta:\bar{\mathbb F}_g^\times\to\C^\times$
and generators $\gamma_d$ of $\mathbb F_{g^d}^\times$ compatible under norms; for a root
$x$ of a monic irreducible $P$ of degree $d$ with $P(0)\ne0$ put
$\theta_x(\gamma_d^n)=\theta(x)^n$.  Hypothesis (G) is: for every such $g$,
$\theta$, $(\gamma_d)$ and every $d\ge1$ there is a bijection $P\mapsto\chi_P$
from the monic irreducible polynomials of degree $d$ over $\mathbb F_g$ with
non-zero constant term onto the irreducible cuspidal characters of
$GL_d(\mathbb F_g)$, with $\chi_P(1)=\prod_{i=1}^{d-1}(g^i-1)$ and
$\chi_P(y)=(-1)^{d-1}\sum_{i<d}\theta_x(z^{g^i})$ whenever $y$ has irreducible
characteristic polynomial with root $z$ and $x$ is a root of $P$.  This is the
cuspidal part of Green's parametrisation of the irreducible characters of
$GL_d(\mathbb F_g)$ (J.~A.~Green, Trans.\ AMS 80 (1955), 402--447; reference to be confirmed); it enters as a hypothesis.
\end{definition}

\emph{The weights.}  The datum fixes, for $j\in\{0,1\}$, idempotents $\varepsilon_P$ of the group algebra
$\C[GL_r(\mathbb F_{g_j})]$ indexed by monic polynomials $P$ of degree $r$ (the series idempotents), with $\varepsilon_P=0$
unless $P$ is a label and $\sum_{P\text{ label}}\varepsilon_P=1$.  Under the normalised-types compatibility, $\varepsilon_P$ is the
sum of the primitive central idempotents of the irreducible characters with label $P$; in particular it is central.  For $a\in\C[GL_r(\mathbb F_{g_j})]$ with $\mathrm{Re}\,a(1)\ne0$
put $\mathrm{wt}(a,P)=\mathrm{Re}\,(a\varepsilon_P)(1)/\mathrm{Re}\,a(1)$, where $x(1)$ is the coefficient of the identity.
For a nontrivial additive character $\psi$ of $\mathbb F_{g_j}$ let $U\subset GL_r(\mathbb F_{g_j})$ be the upper
unitriangular group, $e_U(\psi)=|U|^{-1}\sum_{u\in U}\psi(\sum_{i<r-1}u_{i,i+1})^{-1}u$ (the Gelfand--Graev idempotent),
and $e_J(\psi)=|U_J|^{-1}\sum_{u\in U_J}\psi(\phi_J(u))^{-1}u$, where $U_J=\{u\in U:u_{a,s+k}=0$ for $s\le a<s+k\}$ and
$\phi_J$ is the sum of the superdiagonal entries along the chain of $U_J$, with $s=m/2-1$ and $k$ the rank
(the idempotents $e_U$, $e_J$ and the map $\phi_J$).  Both have identity coefficient
$|U|^{-1}$, resp.\ $|U_J|^{-1}$, which is $>0$.

\begin{proposition}\label{i3:prop-admLq}
Assume (G).  Let $\mathcal D$ be the datum of a selected member of the working class (it carries the
normalised-types compatibility of the twist selection), $j\in\{0,1\}$ and $\psi$ a nontrivial additive character of
$\mathbb F_{g_j}$.  Then:
\begin{enumerate}
\item[(i)] $\mathrm{wt}(e_U(\psi),P)\ge0$ and $\mathrm{wt}(e_J(\psi),P)\ge0$ for every $P$, both vanish unless $P$ is a label,
 and both have total mass $1$: $\sum_{c\in\mathbb F_{g_j}^r}\mathrm{wt}(e,P_c)=1$ for $e=e_U(\psi),e_J(\psi)$;
\item[(ii)] $\mathrm{wt}(e_U(\psi),P)=\mathrm{sq}(P)$ for every squarefree label $P$;
\item[(iii)] $(1-\frac1{g_j})\,\mathrm{wt}(e_U(\psi),P)\le\mathrm{wt}(e_J(\psi),P)$ for every label $P$;
\item[(iv)] $\mathrm{wt}(e_J(\psi),P)\le(1+\frac1{g_j})^r\,\mathrm{wt}(e_U(\psi),P)$ for every squarefree label $P$.
\end{enumerate}
\end{proposition}

\noindent Item (ii) is stated for squarefree labels, which is all that
Proposition~\ref{i3:prop-W} needs.

\noindent Write $G=GL_r(\mathbb F_{g_j})$, $M=r(r-1)/2$, and for $a\in\C[G]$ and a character $\chi$ of $G$ put
$\mathrm{tr}(a|\chi)=\sum_{x\in G}a(x)\chi(x)$, the trace of $a$ acting in a representation with character $\chi$.  The rank of the datum
satisfies $r=k+2s+1$ with $s=m/2-1$ and $k\ge1$.

\begin{proposition}[(J), (Mir) proved; (Sq) proved from (G) and the cited (HC), (GG)]\label{i3:prop-admLq-in}
Let $g$ be a prime, $\psi$ a nontrivial additive character of $\mathbb F_g$, $r=k+2s+1$ with $k\ge1$, and $G=GL_r(\mathbb F_g)$.
\begin{enumerate}
\item[(J)] $e_J(\psi)$ is an idempotent of $\C[G]$ with $e_J(\psi)(1)=g^k\,g^{-M}$.
\item[(Mir)] For every character $\chi$ of $G$, $(g^k-g^{k-1})\,\mathrm{Re}\,\mathrm{tr}(e_U(\psi)|\chi)\le\mathrm{Re}\,\mathrm{tr}(e_J(\psi)|\chi)$.
\item[(Sq)] Assume (G).  In the situation of Proposition~\ref{i3:prop-admLq}, (ii) and (iv) hold.
\end{enumerate}
\end{proposition}

\noindent The proof of Proposition~\ref{i3:prop-admLq-in} is given in \S\ref{i3:ssec-pf-admLq-in}: (J) and (Mir) are
proved completely (the latter from the identities (Ex), (Der) of Proposition~\ref{i3:prop-mirder}, proved there); (Sq) is
proved from these, the identity (BZ) (proved there from a cited theorem of Zelevinsky), published facts (HC), (GG) on the
characters of $GL_r(\mathbb F_g)$, and the comparison of the labels of the datum with Green's labels (Lemma~\ref{i3:lem-labcomp},
proved there).

\begin{proof}[Proof of Proposition~\ref{i3:prop-admLq}]
\emph{Vanishing and total mass.}  If $P$ is not a label then $\varepsilon_P=0$; and
$\sum_P\mathrm{wt}(e,P)=\mathrm{Re}\,(e\sum_P\varepsilon_P)(1)/\mathrm{Re}\,e(1)=1$.  Since $c\mapsto P_c$ is a bijection from
$\mathbb F_{g_j}^r$ onto the monic polynomials of degree $r$, the sum over $c$ is the sum over $P$.

\emph{Three facts on $\C[G]$.}  For an irreducible character $\chi$ let $z_\chi=\frac{\chi(1)}{|G|}\sum_{x}\chi(x^{-1})x$ be its
primitive central idempotent.  (F1) Every central idempotent of $\C[G]$ is the sum of the $z_\chi$ over a set of irreducible
characters: the $z_\chi$ are a basis of the centre consisting of orthogonal idempotents, and $\varepsilon=\sum c_\chi z_\chi$ with
$\varepsilon^2=\varepsilon$ forces $c_\chi^2=c_\chi$.  By the normalised-types compatibility, $\varepsilon_P=\sum_{\chi\in I_P}z_\chi$ with $I_P$ the set of irreducible characters with label $P$.  (F2) For
$a\in\C[G]$, $(az_\chi)(1)=\sum_xa(x)z_\chi(x^{-1})=\frac{\chi(1)}{|G|}\mathrm{tr}(a|\chi)$.  (F3) If $a$ is idempotent and $\chi$ is the
character of a representation $\rho$, then $\rho(a)$ is an idempotent operator, so $\mathrm{tr}(a|\chi)=\mathrm{tr}\,\rho(a)$ is its
rank, a nonnegative integer.  Combining,
\[
\mathrm{Re}\,(a\varepsilon_P)(1)=\sum_{\chi\in I_P}\frac{\chi(1)}{|G|}\,\mathrm{Re}\,\mathrm{tr}(a|\chi),\qquad \frac{\chi(1)}{|G|}>0,\tag{$*$}
\]
and for idempotent $a$ every summand is $\ge0$.

\emph{The two idempotents.}  On the upper unitriangular group $U$ the superdiagonal entries add under multiplication:
$(uu')_{i,i+1}=u_{i,i+1}+u'_{i,i+1}$, since $u_{i,j}=0$ for $j<i$ and $u_{ii}=1$.  So $\lambda(u)=\psi(\sum_iu_{i,i+1})$ is a homomorphism
$U\to\C^\times$ and $e_U(\psi)=|U|^{-1}\sum_u\lambda(u)^{-1}u$ satisfies
$e_U^2=|U|^{-2}\sum_{u,u'}\lambda(uu')^{-1}uu'=|U|^{-1}\sum_{y}\lambda(y)^{-1}y=e_U$.  Its identity coefficient is $|U|^{-1}=g^{-M}$.
By Proposition~\ref{i3:prop-admLq-in}(J), $e_J(\psi)$ is idempotent with identity coefficient $g^kg^{-M}$.  Both identity coefficients are
real and positive.

\emph{Nonnegativity in (i).}  For $e=e_U(\psi)$ or $e_J(\psi)$, $\mathrm{wt}(e,P)=\mathrm{Re}(e\varepsilon_P)(1)/\mathrm{Re}\,e(1)$ is a quotient of
a nonnegative number, by $(*)$ and (F3), by a positive one.

\emph{(iii).}  By $(*)$ for $e_U$ and $e_J$, Proposition~\ref{i3:prop-admLq-in}(Mir) applied to each $\chi\in I_P$, and
$\chi(1)/|G|>0$,
\[
(g^k-g^{k-1})\,\mathrm{Re}\,(e_U\varepsilon_P)(1)\le\mathrm{Re}\,(e_J\varepsilon_P)(1).
\]
Divide by $g^kg^{-M}=\mathrm{Re}\,e_J(1)$ and use $(g^k-g^{k-1})/(g^kg^{-M})=(1-\frac1g)/g^{-M}=(1-\frac1g)/\mathrm{Re}\,e_U(1)$: this is
$(1-\frac1g)\mathrm{wt}(e_U,P)\le\mathrm{wt}(e_J,P)$, for every $P$ (in particular every label), with $g=g_j$.

\emph{(ii) and (iv)} are Proposition~\ref{i3:prop-admLq-in}(Sq).
\end{proof}

\begin{lemma}[Comparison transfer]\label{i3:lem-transfer}
Let $\iota$ be a finite set, $T_v$ ($v\in\iota$) finite sets, $g\ge1$, and
$w^U_v,w^J_v:T_v\to\R$ with $w^U_v\ge0$,
$\sum_tw^U_v(t)=\sum_tw^J_v(t)=1$ and $(1-\frac1g)w^U_v\le w^J_v$.  Then for
every $S\subset\prod_vT_v$,
$\sum_{x\in S}\prod_vw^J_v(x_v)\le\sum_{x\in S}\prod_vw^U_v(x_v)+|\iota|/g$.
\end{lemma}

\begin{proof}
Let $\alpha=1-\frac1g$, $n=|\iota|$ and $\Sigma^U_S=\sum_S\prod w^U_v$.  Both
product weights have total mass $1$, and pointwise
$\prod_vw^J_v\ge\alpha^n\prod_vw^U_v$.  Hence
$\sum_S\prod w^J_v=1-\sum_{S^c}\prod w^J_v\le1-\alpha^n(1-\Sigma^U_S)\le
(1-\alpha^n)+\Sigma^U_S$, and $1-\alpha^n\le n/g$ by Bernoulli's inequality.
\end{proof}

\medskip\noindent\emph{The all-rows datum.}
Let $\psi^{(k)}(t)=e^{2\pi it/g_k}$, a nontrivial additive character of
$\mathbb F_{g_k}$.  For the cylinder sets of Definition~\ref{i3:def-allrows}
we write $\mathrm{cyl}_kS$ ($S$ a set of $\Pi_k$-tuples).

\begin{proposition}[formally verified, \cite{CharacterFamily}]\label{i3:prop-allrows-built}
Assume the unitary Siegel--Weil formula in the convergent range and that the
row spaces are finite dimensional.  Let $\mathcal D$ be the datum of a member
of the working class, and suppose that $\mathcal D$ carries an automorphic
realisation $X$, a good operator and a row isometry $\iota$ satisfying the
three ties of the hypothesis (AT) of Proposition~\ref{i3:prop-irr-row}.  Then,
for $x$ beyond a threshold depending only on the constants of the family,
$\mathcal D$ carries an all-rows datum $A$ such that:
\begin{enumerate}
\item[(a)] $A$ is faithful; all entries of every row label are labels;
\item[(b)] for $k=0,1$, $\mathrm{share}_k(S)=\sum_{t\in S}\prod_{v\in\Pi_k}
 \mathrm{wt}(e_J(\psi^{(k)}),t_v)$ for every set $S$ of $\Pi_k$-tuples;
\item[(c)] for $k=0,1$ and every $S$, $\dim_A(\mathrm{cyl}_kS)=D_{\rm count}(I)$
 and $D^*_A(\mathrm{cyl}_kS)=D^*(I)$ for an automorphic counting problem
 $I=I_{k,S}$ of rank $n=r$, degree $f$, root discriminant
 $\mathrm{rd}_E$, level $g_0g_1$, auxiliary prime $g_{1-k}$, whose
 Harish-Chandra parameters are the $\lambda^{SH}(\tau)$ of
 \S\ref{i3:ssec-gap} with signatures $\epsilon^{SH}$; for $S$ the set of all
 tuples, $\dim_A=D_{\rm all}$ and $D^*_A=D_*$.
\end{enumerate}
\end{proposition}

\noindent\emph{What remains.}  The realisation of the rows of $\mathcal D$
as automorphic forms on a unitary group (a Hida setting), with its three ties,
is the hypothesis; it is not constructed here.  Given it, the Hecke idempotent
is cut by the cylinder idempotents $e*\varepsilon_{E_S}$, where $E_S$ is built
from the series idempotents of $GL_r(\mathbb F_{g_k})$ at the $2f$ places
above $g_0,g_1$ through the reduction map $K\to\prod GL_r(\mathbb F_{g_k})$.
One proves that these are admissible idempotents (surjectivity of the
reduction, kernel $K(g)$, unipotence modulo $g_{1-k}$, idempotence and
positivity through the finite Hecke algebra $\C[K/K(g)]$), that the
corresponding dimensions are additive over label tuples, and that the
labelled row basis spans exactly the good part, which yields the faithful
labelling.  The proof is long; it uses the
discrete-series decay for integrability.  Its formally verified steps are outlined in \S\ref{i3:str-allrows}.

\begin{lemma}[Counting hypotheses]\label{i3:lem-countclauses}
Assume the size clauses (1)--(25) of the standing hypotheses at $(C',C_l)$
(the numerical clauses (N1)--(N14), in particular the window, are not used),
put $C_g=C_{YL}+1$, assume the record clauses (L) and (A) of
\S\ref{i3:ssec-pf-slots}, and let $n_0\in\N$,
$A,C,c\in\R$ with $0\le C$, $0<c$, $2C_g\le A$, $n_0\le k$, $2\le C_l$,
$2C\le C_l$, and $\log L\le ck$.  Then every counting problem $I$ as in
Proposition~\ref{i3:prop-allrows-built}(c), and every counting problem with
the same rank, degree, root discriminant, level, auxiliary prime $g_0$ or
$g_1$ and parameters $\lambda^{SH}$, satisfies $\mathrm{CH}(n_0,A,b,h)$ with
$h=Cr^{19}f^5\log(b\,\mathrm{rd}_E)$, and $\log\log b\le c\,rf\log r$.
\end{lemma}

\begin{proof}
We use the size clauses listed in the proof of
Proposition~\ref{i3:prop-standing-family}.  The signature $\epsilon^{SH}$ does
not depend on $\tau$, so the problem has the same signature at every infinite
place, and each $\lambda^{SH}(\tau)$ is a Harish-Chandra parameter by
Lemma~\ref{i3:lem-gap} (whose hypotheses \eqref{i3:eq-standing} are among the
size clauses).  First $k\ge5$, $k\le r\le2k$,
$m\ge2$, $L>1$, $\log b=L$; from $C_lm^2L\le k$ with $C_l\ge2$:
$m^2\le k$ and $C_lL\le k$; from clause (11), $\log\mathrm{rd}_E\le L/C_l$.
Then: $n_0\le k\le r$; the signature has $q\le r_S+1\le m$
(the indices $a$ with $\epsilon^{SH}_a=0$ are $a=k-1+2i+1$, $0\le i\le r_S$),
so $q^2\le m^2\le r$; $\log\log\mathrm{rd}_E\le\log\mathrm{rd}_E\le L\le
k\le r$; $|\lambda^{SH}_a|\le b^2$, since an entry is $v_a\in[0,4k\Delta]$ (by (L)) or
$-b_j$ with $0\le b_j\le b^2$, and $4k\Delta\le C_lmk\Delta\le b_1\le b^2$;
$g_0g_1\le b^{2C_g}\le b^A$; for $g\in\{g_0,g_1\}$,
$g\ge(kb)^8\ge k^8\ge5^7k>4r$, and $g$ splits completely by (A).  The gap:
$\lambda^{SH}(\tau)$ has gap at least $C_lr^{19}f^5L$ by
Lemma~\ref{i3:lem-gap}, and $Cr^{19}f^5\log(b\,\mathrm{rd}_E)=
Cr^{19}f^5(L+\log\mathrm{rd}_E)\le2Cr^{19}f^5L\le C_lr^{19}f^5L$.  Further
$\log f\le\sqrt{r_S}\le\sqrt r$ (clause (5)), and
$\log\log b=\log L\le ck\le c\,rf\log r$.
\end{proof}

\begin{proof}[Proof of Proposition~\ref{i3:prop-irr-row}]
\emph{Constants.}  Let $C_g=C_{YL}+1$, $A=2C_g$, and $n_0=n_0(A)$,
$C=C_0(A)$ from Proposition~\ref{i3:prop-counting}; put
$c=c_{Ic}=\frac1{16}$.  Take $\varepsilon_1$ and, for
$\varepsilon\le\varepsilon_1$, a threshold system containing the Large
conditions $2\le C_l^{\mathcal F}$, $2C\le C_l^{\mathcal F}$,
$8C_\Gamma\le C_{k,\rm low}$ and those of the hypothesis (AT), and let
$\mathcal F\in\mathcal C$ be large for it.  By hypothesis $\mathcal F$
satisfies $\mathrm{TuSize}$, and eventually each member carries a realisation
as in (AT), so Proposition~\ref{i3:prop-allrows-built} applies to it.

\emph{Eventually the counting hypotheses hold.}  Put $y=x^{\varepsilon/8}$, so
$L=y^8$.  Choose $Y_1$ with $n_0y\le8C_\Gamma y^8$ and $8y\le\frac1{16}\cdot
8C_\Gamma y^8$ for $y\ge Y_1$ (possible as $C_\Gamma>0$: take
$Y_1=\max(1,n_0/(8C_\Gamma),16/C_\Gamma)$: for $y\ge Y_1$ one has $y^7\ge y$,
so $8C_\Gamma y^8\ge n_0y$ and $\frac12C_\Gamma y^8\ge8y$).  By
Proposition~\ref{i3:prop-standing-family}, which applies since $\mathcal F$
satisfies $\mathrm{TuSize}$ and $8C_\Gamma\le C_{k,\rm low}$, eventually
every member satisfies the size clauses (1)--(25) at
$(C_{k,\rm low},C_l^{\mathcal F})$, without the window (clause (23) is
$\mathrm{TuSize}$).  From
$8C_\Gamma\le C_{k,\rm low}$ and $C_{k,\rm low}m^5fL\le k$ we get
$8C_\Gamma y^8\le k$, hence for $y\ge Y_1$: $n_0\le k$ and
$\log L=8\log y\le8y\le\frac1{16}k$.  Lemma~\ref{i3:lem-countclauses} applies
with $c=\frac1{16}$.

\emph{CountAllRows.}  Let $A$ be the datum of
Proposition~\ref{i3:prop-allrows-built}.  For every $k$ and $S$, the problem
$I_{k,S}$ satisfies the hypotheses of Proposition~\ref{i3:prop-counting}, so
$|\dim_A(\mathrm{cyl}_kS)-D^*_A(\mathrm{cyl}_kS)|\le
e^{-\frac1{16}rf\log r}D^*_A(\mathrm{cyl}_kS)$.  Since $r\ge2$, $f\ge1$,
$\varepsilon_C=e^{-\frac1{16}rf\log r}<1$; and $D_*>0$.  This is
$\mathrm{CountAllRows}(A,\varepsilon_C)$ with $c_{Ic}=\frac1{16}$.

\emph{WeightAllRows.}  Fix $k$, $g=g_k$, and put
$w^J_v(c)=\mathrm{wt}(e_J(\psi^{(k)}),P_c)$ for $v\in\Pi_k$ and
$P(c)=\mathrm{wt}(e_U(\psi^{(k)}),P_c)$.  By Proposition~\ref{i3:prop-admLq}
these are probability weights on $\mathbb F_g^r$, vanishing off labels, with
$(1-\frac1g)P\le w^J_v$ (on labels by (iii), and off labels both vanish) and $P(c)=\mathrm{sq}(P_c)$ on squarefree labels.  Let $S$ be a set
of bad tuples.  By Proposition~\ref{i3:prop-allrows-built}(b) and since
$c\mapsto P_c$ is injective, $\mathrm{share}_k(S)$ is at most the
$\prod_vw^J_v$-mass of the non-good coefficient tuples, which by
Lemma~\ref{i3:lem-transfer} is at most the $P^{\otimes\Pi_k}$-mass of the
non-good tuples plus $f/g$.  The host primes $q,d$ of the member are distinct
odd primes with $\max(q,d)\ge16$ (clause (S)) and $16qd\mid g-1$ (host
admissibility in (A): $g\equiv1\bmod\mathrm{QS}\cdot qd$ with
$16\mid\mathrm{QS}$); $g\ge(kb)^8\ge2^{12}$ and $r\ge100$ eventually.
Proposition~\ref{i3:prop-W} with
$K=\mathbb F_g$, $N=r$, $|\iota|=f$, $|J|=m$ gives $\mathrm{share}_k(S)\le W_g(f,m,r,g,q,d)$ (the $f/g$ turns $2f/g$ into
$3f/g$).  This is $\mathrm{WeightAllRows}(A)$.

\emph{Size conditions.}  Take $C_s=6$ and $x$ large.  By (S) and (P) of
\S\ref{i3:ssec-pf-slots}: $g_j\ge(kb)^8\ge b^8=e^{8x^\varepsilon}\ge
e^{8x^{\varepsilon_z}}$ (as $\varepsilon_z\le\varepsilon$);
$f=2n_F\ge2x^{2+8\varepsilon-\theta}\ge x^2$ and
$f\le2x^{2+8\varepsilon+\theta}\le x^6$; $m\le x^\theta\sqrt x\le x^6$;
$r\ge k\ge f\ge x^2\ge100$ and $r\le2k\le2x^{9/2+10\varepsilon+\theta}\le x^6$
(as $\varepsilon<\frac1{52}$).  Finally $\min(q,d)\ge e^{x^{\varepsilon_z/6}}$
because the size record of the class contains
$\log q,\log d\ge x^{\varepsilon_z/5-\varepsilon_z/100}$, and
$\varepsilon_z/5-\varepsilon_z/100>\varepsilon_z/6$.  All hypotheses of
Theorem~\ref{i3:thm-linkIRR} hold.
\end{proof}

\subsection{Term systems}\label{i3:ssec-grp-terms}

We prove Proposition~\ref{i3:prop-conc1}.

\subsubsection{Proof of Proposition~\ref{i3:prop-conc1}}\label{i3:ssec-pf-conc1}

\medskip\noindent\emph{Term systems.}

\begin{definition}[Term systems]\label{i3:def-terms}
A \emph{term system} consists of a set $\tau$ of terms, a set $\mathrm{Rw}$ of
rows, a finite set $I$, complex inner product spaces $H$, $H_1$ ($H_1$
complete), maps $\mathrm{op}:\tau\to H$, $\mathrm{vec}:\tau\to H_1$,
$\mathrm{row}:\tau\to\mathrm{Rw}$, $\iota:\mathrm{Rw}\to I$, predicates
$\mathrm{isId}$, $\mathrm{gener}$ on $\tau$, a function $G:\tau\to\R$, a matrix
$\Gamma^\sharp\in\C^{I\times I}$, a number $\Gamma_1>0$, a corank
$s_0:\tau\times\tau\to\N$ and a relation $\mathrm{scal}$ on pairs.  Put
$\mathrm{id}=\{\mathrm{isId}\}$, $\mathrm{deg}=\{G=0,\neg\mathrm{isId}\}$,
$\mathrm{gen}=\{G\ne0,\mathrm{gener}\}$, $\mathrm{ng}=\{G\ne0,\neg\mathrm{gener}\}$,
$\mathrm{nz}=\{G\ne0\}$, $\mathrm{part}(S)=\sum_{t\in S}\mathrm{op}_t$,
$K_0=\mathrm{part}(\mathrm{id})+\mathrm{part}(\mathrm{gen})+\mathrm{part}(\mathrm{ng})$,
$v(S,i)=\sum_{t\in S,\,\iota\mathrm{row}(t)=i}\mathrm{vec}_t$,
$\Phi(S)=\Gamma_1\sum_i\|v(S,i)\|^2$.  For $S\subset\tau$ let
$\mathrm{all}(S)$ be the pairs $(t,t')\in S^2$ with equal rows, and inside it
$\mathrm{big}(S)=\{s_0\ge2\}$, $\mathrm{one}(S)=\{s_0=1\}$,
$\mathrm{nsc}(S)=\{s_0=1,\neg\mathrm{scal}\}$,
$\mathrm{sc}(S)=\{s_0=1,\mathrm{scal}\}$.  For a set $Q$ of pairs,
$\Sigma(Q)=\Gamma_1\sum_{Q}\mathrm{Re}\langle\mathrm{vec}_t,\mathrm{vec}_{t'}\rangle$
and $|\Sigma|(Q)=\Gamma_1\sum_Q|\langle\mathrm{vec}_t,\mathrm{vec}_{t'}\rangle|$.
Conditions: (AC) $\sum\|\mathrm{op}_t\|$, $\sum\|\mathrm{vec}_t\|<\infty$;
(ID) $\mathrm{isId}\Rightarrow G=0$; (NE$_S$)
$\|\mathrm{part}(S)\|^2=\mathrm{Re}\sum_{i,j}\Gamma^\sharp_{ij}\langle
v(S,i),v(S,j)\rangle$; (GG$_P$) $P\Gamma_1\mathbf 1-\Gamma^\sharp$ is positive
semidefinite; (RO) $\langle\mathrm{vec}_t,\mathrm{vec}_{t'}\rangle=0$ if the
rows differ but have the same image under $\iota$; (CP) $s_0(t,t')\ge1$ if
$t,t'$ have the same row and $G(t)\ne0\ne G(t')$.
\end{definition}

For a datum $\mathcal D$ and words $u,u'$, the kernel block is expanded over
the double cosets of the unitary group (the terms); $\mathrm{op}_t$ is the
orbital operator of the term, $\mathrm{vec}_t$ its vector in the row model, and
$G(t)$ the Gram determinant of its flag pair.  The \emph{Kloosterman part}
$K''(u,u')$ is the image of $\mathrm{part}(\mathrm{deg})$, the sum of the
degenerate terms, under the identification of the next proposition.

\emph{The thresholds.}  Each of the propositions below is a closed statement whose threshold system contains, besides
those of the hypotheses, one explicit threshold on the own constants $(C_{k,\rm low},C_l^{\mathcal F})$, built from the
constants that the corresponding formally verified input provides at the outer constants
$\mathbf c=(C_{YL}+1,\frac14,3,\frac12,C_{Tu},C_0)$ with $C_0=3$:
\begin{itemize}
\item \emph{off-diagonal Kloosterman threshold} the off-diagonal Kloosterman input gives, at $\mathbf c$, constants $(C'_0,C_{l,0})$ such that its conclusion holds for $C'\ge C'_0$, $C_l\ge C_{l,0}$;
the threshold is $\max(3,C'_0,C_{l,0})$, imposed on both $C_{k,\rm low}$ and $C_l^{\mathcal F}$;
\item \emph{diagonal Kloosterman threshold} the same
construction from the diagonal Kloosterman input (whose hypotheses are those of (Ty) with $X(u,u)$ satisfying (AC));
\item \emph{diagonal-row threshold} the diagonal-row input gives at $(\mathbf c,A=1)$
constants $C_B,C_4\ge0$ and, for every $C_G\ge\max(C_B,C_4)+1$, constants $(C'_0,C_{l,0})$; if $C_\Gamma\ge\max(C_B,C_4)+1$ the
threshold is $\max(3,C'_0,C_{l,0})$ at $C_G=C_\Gamma$, otherwise $3$ (and the Large condition $C_\Gamma\ge\max(C_B,C_4)+1$ is
imposed separately);
\item the numerical thresholds $29\le C_{k,\rm low}$, $29\le C_l^{\mathcal F}$ $40\le C_l^{\mathcal F}$
(Lemma~\ref{i3:S-lem-rankrow}), and $1000+c_{\rm orb}\le C_l^{\mathcal F}$, with $c_{\rm orb}:=|c|$ for the constant $c$
provided by the orbital-integral convergence.
\end{itemize}
In the first three items the formal threshold is $\max(3,c,c')$ for witnesses $c,c'$ chosen by the formal existence
statements and otherwise unspecified; where the text below states a smaller explicit threshold, it is the written proof
that shows the smaller value suffices.
``Closed (with thresholds of \cite{CharacterFamily})'' below means closed with these thresholds added.

\begin{proposition}[formally verified, \cite{CharacterFamily}]\label{i3:prop-terms}
Let $\mathcal C$ be the working class, with a spectral selection satisfying
(Sel) and (Rk).  The following is closed on $\mathcal C$ (with thresholds of
\cite{CharacterFamily}, including the Large conditions $29\le C_{k,\rm low}$,
$29\le C_l^{\mathcal F}$, $40\le C_l^{\mathcal F}$ of the admissibility of the
selection, under which eventually every selected package is concretely tied): eventually, at every member satisfying the standing
hypotheses at $(3,C_{k,\rm low},C_l^{\mathcal F})$, for all words
$u,u'$ there are a term system $X(u,u')$ and a linear isometry $j$ from $H$
onto the Hilbert--Schmidt operators $C_{u'}\to C_u$ with
$j(\sum_t\mathrm{op}_t)=K(u,u')$ and $j(\mathrm{part}(\mathrm{deg}))=K''(u,u')$,
satisfying (AC), (ID), (RO), (CP), (NE$_S$) for
$S\in\{\mathrm{deg},\mathrm{gen},\mathrm{ng}\}$ and (GG$_{P_g}$) with
$P_g=(1+3/g_0)^f(1+3/g_1)^f$; for $u=u'$ moreover
$j(\mathrm{part}(\mathrm{id}))=\mathrm{id}$, and for $u\ne u'$,
$\mathrm{part}(\mathrm{id})=0$.
\end{proposition}

\noindent\emph{What remains.}  The Poincar\'e-series model of the column
spaces, the unfolding of the kernel into orbital terms with absolute
convergence (from the discrete-series decay), the Parseval identity (NE) with
the finite-field Gram matrix at the auxiliary primes and its bound by $P_g$
(a Gelfand--Graev computation on the mirabolic subgroup), and orthogonality
of distinct rows.  These are the structural components of Theorem~S, shared
with Proposition~\ref{i3:prop-S-closed}; this is a long verification.  Its formally verified steps are outlined in \S\ref{i3:str-terms}.

\begin{proposition}\label{i3:prop-kpp-off}
In the situation of Proposition~\ref{i3:prop-terms}, the following is closed
on $\mathcal C$ with the further threshold $3$ (that is, $C_{k,\rm low}\ge3$ and $C_l^{\mathcal F}\ge3$) beyond
those of Proposition~\ref{i3:prop-terms}, in particular those of the admissibility of the selection: eventually,
at every member satisfying the standing hypotheses at $(3,C_{k,\rm low},C_l^{\mathcal F})$, if $u\ne u'$ then
$\mathrm{op}_t=0$ for every degenerate term $t$ of $X(u,u')$.
\end{proposition}

\noindent (In the formal statement the threshold is $\max(3,C'_0,C_{l,0})$ for witnesses $(C'_0,C_{l,0})$ of the
formal input that are not specified; the proof given here shows that $3$ suffices.  The formal statement used in the
main line takes the concrete ties of the selected packages as a hypothesis, which is why (Rk) and the admissibility
thresholds are assumed.)  The proof
(\S\ref{i3:ssec-pf-kppoff}) reduces the vanishing to a Schur-lemma argument, proved there, and to statements on the orbital operators
(Propositions~\ref{i3:prop-VGeo} and \ref{i3:prop-cov}), proved there from the formally verified inputs of
Propositions~\ref{i3:prop-kppoff-in} and \ref{i3:prop-cov-in}.

\begin{proposition}\label{i3:prop-kpp-diag}
In the situation of Proposition~\ref{i3:prop-terms}, assume the type bound (Ty).  The following is closed on
$\mathcal C$ (with the Kloosterman-diagonal threshold, defined in the proof): for every word $u$ the degenerate pairs of
$X(u,u)$ satisfy
$P_g\,|\Sigma|(\mathrm{all}(\mathrm{deg}))\le d_u\,e^{-kf\log k}$, the family
$(t,t')\mapsto\|\mathrm{vec}_t\|\|\mathrm{vec}_{t'}\|$ being summable on $\mathrm{all}(\mathrm{deg})$.
\end{proposition}

\noindent (The bound on the degenerate pairs uses
only (Ty) and the ties of the members, which are a record of every family of the class.)  The proof is in
\S\ref{i3:ssec-pf-kppdiag}: the degenerate pairs are grouped by rank and label, the type bound (Ty) bounds each
group, the labels of a given rank are summed with a coset count, and a numerical comparison closes the estimate.

\begin{proposition}\label{i3:prop-diag-rows}
Assume (SW), and let the working class $\mathcal C$ have a spectral selection
satisfying (Sel), (Rk), (Lp) and (Mn).  Let $C_B\ge0$ be the constant of (Mn) at the outer constants
$\mathbf c$ and $A=1$, and $C_4\ge0$ the constant $C_n$ of Proposition~\ref{i3:S-prop-NG} at $\mathbf c$.  The
following is closed on $\mathcal C$, with the Large conditions $2<C_l^{\mathcal F}$, those of the admissibility of
the selection ($29\le C_{k,\rm low}$, $29\le C_l^{\mathcal F}$, $40\le C_l^{\mathcal F}$),
$C_{k,\rm low}\ge\max(3,C'_0)$ and $C_l^{\mathcal F}\ge\max(3,C_{l,0})$, where $(C'_0,C_{l,0})$ is the
coordinatewise maximum of the thresholds of (Mn) at $(\mathbf c,1)$ and of Proposition~\ref{i3:S-prop-NG} at
$\mathbf c$: eventually, at every member satisfying the standing hypotheses at $(3,C_{k,\rm low},C_l^{\mathcal F})$,
for every word $u$ the term system $X(u,u)$ satisfies, with $d=d_u^2/D_*$, $\varepsilon_L=e^{-\frac13kf\log k}$,
$m=2r_S+2$, $L=\log b$, $\Gamma_B=e^{C_Bm^5fL}$ and $\varepsilon'=2^{-(k-m-3)}e^{C_4m^5fL}$:
\[
 |\Sigma|(\mathrm{big}(\mathrm{nz}))\le\varepsilon_L d,\quad
 |\Sigma|(\mathrm{nsc}(\mathrm{gen}))\le\varepsilon_L d,\quad
 \Sigma(\mathrm{sc}(\mathrm{gen}))\le\Gamma_Bd,\quad
 \Sigma(\mathrm{one}(\mathrm{ng}))\le\varepsilon'd .
\]
The further Large condition $C_\Gamma\ge\max(C_B,C_4)+1$, which the formal statement also imposes, is used only in Lemma~\ref{i3:lem-diagnum}.
\end{proposition}

\begin{proof}
The four estimates are, respectively, the two loop rows (RB$_{\varepsilon_L,d}$), (RS$_{\varepsilon_L,d}$), the main row
(RM$_{\Gamma_B,d}$) and the non-generic row (RN$_{\varepsilon',d}$) of Definition~\ref{i3:S-def-terms} for $X(u,u)$.

\emph{Constants and thresholds.}  The loop row (Lp) is closed for every function $a_F$ of the family constants
(Definition~\ref{i3:S-def-rows}); take $a_F\equiv1$.  The selection is admissible (closed statement): exactly as in the proof of
Proposition~\ref{i3:prop-S-closed}, the frame row (Lemma~\ref{i3:S-lem-rankrow}(i), from (Rk)), (FD)
(Lemma~\ref{i3:S-lem-FD}) and the decay (C12) show that eventually every selected package is concretely tied, the
member being locally tied by (Sel).  Take $\varepsilon_1$ at most those of (Lp) and of the
admissibility, and for $\varepsilon\le\varepsilon_1$ a threshold system containing theirs and the Large conditions of the
statement.  Let $\mathcal F$ be large, and $x$ beyond the thresholds of the ``eventually'' clauses of (Lp), of the
admissibility and of the selection.  Write $C'=C_{k,\rm low}$, $C_l=C_l^{\mathcal F}$; so $C'\ge3\ge1$, $C_l>2$.

\emph{Counting hypotheses.}  If the member satisfies the standing hypotheses at $(3,C',C_l)$ (so $C_0=3$), then, because
$2<C_l$ and $1\le C'$, it satisfies the counting hypotheses at $(u,u)$ with constants $(C',C_l,\mathbf c)$ and $A=1$,
for every word $u$, by Lemma~\ref{i3:lem-cnthyp} (\S\ref{i3:ssec-pf-cnthyp}), which also defines the counting hypotheses.  They
include the clause $\mathrm{Stand}(u,u)$ used by Proposition~\ref{i3:S-prop-NG}.

\emph{The four rows.}  Let $(X,Y)$ be the selected package; it is concretely tied, the member satisfies the ties
(a record of every family of the class), and $X(u,u)$ satisfies (AC) by
Proposition~\ref{i3:prop-terms}.  (i), (ii): the loop row with $a_F\equiv1$, applied with $u'=u$, gives
(RB$_{\varepsilon_L,d}$) and (RS$_{\varepsilon_L,d}$) with $d=d_ud_u/D_*$.  (iii): (Mn) at $(\mathbf c,1)$, applicable
since $C'\ge C'_0$, $C_l\ge C_{l,0}$, gives (RM$_{\gamma_B,d}$) with $\gamma_B=\exp(C_Bm^5f\log b)=\Gamma_B$.
(iv): Proposition~\ref{i3:S-prop-NG} (which assumes (SW)) at $(u,u)$ gives (RN$_{\varepsilon'',d}$) with
$\varepsilon''=2^{-(k-2r_S-5)}e^{C_4(2r_S+2)^5f\log b}$; since $m=2r_S+2$, $k-2r_S-5=k-m-3$ and $\log b=L$, this is
$\varepsilon'$.  In all four, $d_u$ is the dimension of the column space $C_u$ of the package and $D_*$ its
expected value; the concrete tie identifies them with $d_u$ and $D_*$ of the datum.
\end{proof}

\medskip\noindent\emph{Hilbert-space lemmas.}

\begin{lemma}\label{i3:lem-pairs}
Let $X$ be a term system with (AC) and (RO).
\begin{enumerate}
\item $\Phi(S)=\Sigma(\mathrm{all}(S))$ for every $S$.
\item Under (CP), $\Phi(\mathrm{gen})=\Sigma(\mathrm{big}(\mathrm{gen}))
 +\Sigma(\mathrm{nsc}(\mathrm{gen}))+\Sigma(\mathrm{sc}(\mathrm{gen}))$ and
 $\Phi(\mathrm{ng})=\Sigma(\mathrm{big}(\mathrm{ng}))+\Sigma(\mathrm{one}(\mathrm{ng}))$.
\item $|\Sigma|$ is monotone in $Q$, and $|\Sigma(Q)|\le|\Sigma|(Q)$.
\item Under (NE$_S$) and (GG$_P$), $\|\mathrm{part}(S)\|^2\le P\,\Phi(S)$.
\end{enumerate}
\end{lemma}

\begin{proof}
(1) For $h:\tau\to H_1$ absolutely summable,
$\|\sum_th(t)\|^2=\sum_{(t,t')}\mathrm{Re}\langle h(t),h(t')\rangle$ (the double
family is absolutely summable and Fubini applies).  Apply this to
$h=\mathbf 1[t\in S,\iota\mathrm{row}(t)=i]\mathrm{vec}_t$ and sum over $i$:
$\Phi(S)$ is $\Gamma_1$ times the sum over pairs in $S^2$ whose rows have the
same image under $\iota$; by (RO) only pairs with equal rows contribute.
(2) By (CP), $s_0\ge1$ on $\mathrm{all}(\mathrm{gen})$ and
$\mathrm{all}(\mathrm{ng})$; split by $s_0\ge2$, $s_0=1$ and $\mathrm{scal}$.
(3) Sums of nonnegative summable families are monotone; $|\mathrm{Re}\,z|\le|z|$.
(4) If $c\mathbf 1-\Gamma^\sharp=\sum_lv_lv_l^*$ is positive semidefinite, then
$c\sum_i\|a_i\|^2-\mathrm{Re}\sum_{i,j}\Gamma^\sharp_{ij}\langle a_i,a_j\rangle
=\sum_l\|\sum_i\overline{v_l(i)}a_i\|^2\ge0$; take $c=P\Gamma_1$,
$a_i=v(S,i)$.
\end{proof}

\begin{lemma}[Diagonal pair bound]\label{i3:lem-diagabstract}
Let $X$ be a term system with (AC), (ID), (RO), (CP), (NE$_{\rm gen}$),
(NE$_{\rm ng}$), (GG$_P$), $P,d\ge0$, and suppose
$|\Sigma|(\mathrm{big}(\mathrm{nz}))\le\varepsilon d$,
$|\Sigma|(\mathrm{nsc}(\mathrm{gen}))\le\varepsilon d$,
$\Sigma(\mathrm{sc}(\mathrm{gen}))\le A'd$,
$\Sigma(\mathrm{one}(\mathrm{ng}))\le\varepsilon'd$.  Then
\[
 \|K_0-\mathrm{part}(\mathrm{id})\|^2\le
 \Bigl(\sqrt{P(A'+2\varepsilon)}+\sqrt{P(\varepsilon+\varepsilon')}\Bigr)^2d.
\]
\end{lemma}

\begin{proof}
$K_0-\mathrm{part}(\mathrm{id})=\mathrm{part}(\mathrm{gen})+\mathrm{part}(\mathrm{ng})$.
Since $\mathrm{big}(\mathrm{gen}),\mathrm{big}(\mathrm{ng})\subset
\mathrm{big}(\mathrm{nz})$, Lemma~\ref{i3:lem-pairs}(2),(3) give
$\Phi(\mathrm{gen})\le(A'+2\varepsilon)d$ and
$\Phi(\mathrm{ng})\le(\varepsilon+\varepsilon')d$, and (4) gives
$\|\mathrm{part}(\mathrm{gen})\|^2\le P(A'+2\varepsilon)d$,
$\|\mathrm{part}(\mathrm{ng})\|^2\le P(\varepsilon+\varepsilon')d$.  The
triangle inequality concludes.
\end{proof}

\begin{lemma}[Diagonal numerics]\label{i3:lem-diagnum}
Assume the standing hypotheses, and let $C_B,C_4$ be the constants of
Proposition~\ref{i3:prop-diag-rows} with $C_\Gamma\ge\max(C_B,C_4)+1$.  Then
$\bigl(\sqrt{P_g(\Gamma_B+2\varepsilon_L)}+\sqrt{P_g(\varepsilon_L+\varepsilon')}\bigr)^2\le\Gamma$.
\end{lemma}

\begin{proof}
Let $M=\max(C_B,C_4)\,m^5fL\ge0$.  By the standing hypotheses
$(kb)^8\le g_j$ and $f\le k$, so $f/g_j\le b^{-7}$ and
$P_g\le e^{3f/g_0+3f/g_1}\le e^{6b^{-7}}\le e$.  Also $\varepsilon_L\le1$,
$\Gamma_B\le e^M$, and $\varepsilon'\le e^M$ since $k-m-3\ge0$.  Using
$(\sqrt a+\sqrt c)^2\le2(a+c)$, the left side is at most
$2P_g(\Gamma_B+3\varepsilon_L+\varepsilon')\le2e(2e^M+3)\le10e\cdot e^M<e^4e^M$.
Since $C_\Gamma\ge\max(C_B,C_4)+1$ and $m^5fL\ge512L\ge4$ (as $m\ge2$,
$f\ge m^4\ge16$ and $L=x^\varepsilon>1$),
$\Gamma=e^{C_\Gamma m^5fL}\ge e^Me^{m^5fL}\ge e^4e^M$.
\end{proof}

\medskip\noindent\emph{The count bound.}

\begin{proposition}\label{i3:prop-rowtie}
Let the working class have a spectral selection satisfying (Sel) and (Inst), such that eventually the
instance of each selected $X$ is the instance of the member $\mathcal D$ and $\mathcal D$ has the volume tie (VT) of
Proposition~\ref{i3:prop-massratio} (both hold for the spectral selection induced by a raw selection, as in
Proposition~\ref{i3:prop-slots-SW}: the first is a clause of its spectral condition, the second a record of its twist
selection).  Then eventually every selected member $\mathcal D$, with selected package $(X,Y)$, has
$D_{\rm all}=D_{\rm count}(I)$ and $D_*=D^*(I)$ for the automorphic counting problem $I=I(X)$ of the package: the
hermitian space of $\mathcal D$ (rank $r$, signature $(k+r_S,r_S+1)$ at every infinite place), degree $f$, root
discriminant $\mathrm{rd}_E$, level $g_0g_1$, auxiliary prime $g_0$, the idempotent $e=\overline{\chi_2}\mathbf 1_{K_2}/\mu_f(K_2)$ of
the package, and Harish-Chandra parameters $\lambda^{SH}(\tau)$ with signatures $\epsilon^{SH}$.  Moreover $I$ satisfies the
instance tie (it is the instance of a datum with the volume tie, (Inst)(i), (Inst)(ii) and the
archimedean scale tie) and has $q=r_S+1\ge1$.
\end{proposition}

\begin{proof}
Take $x$ beyond the selection threshold, so that the selected package $(X,Y)$ exists, with its spectral ties.
\emph{The two equalities.}  Among the spectral ties of the package is the automorphic tie, which states
precisely that the row count and the expected value of $\mathcal D$ are the count $D_{\rm count}$ and the expected value
$D^*$ of the automorphic instance of $X$ with the idempotent $e$; so
$D_{\rm all}=D_{\rm count}(I)$ and $D_*=D^*(I)$ by definition.
\emph{The idempotent.}  $K_2$ is a subgroup of the stabiliser $K$ of the standard lattice containing $K(g)$ (Lemma~\ref{i3:lem-K2group}), hence open
($K(g)$ is open, Step~1 of \S\ref{i3:S-ssec-kppvol}), hence closed (an open subgroup is closed), hence compact as a closed subset of
the compact $K$ (Proposition~\ref{i3:S-prop-bridge}(K)); and $\chi_2$ is a unitary
character of $K_2$ (Proposition~\ref{i3:prop-chi2}); then
$e*e=e$, $e^*=e$ and $e$ is of positive type, by orthogonality of characters on $K_2$:
$\int_{K_2}\overline{\chi_2(y)}\,\overline{\chi_2(y^{-1}x)}\,dy=\mu_f(K_2)\overline{\chi_2(x)}$ for $x\in K_2$.  Every element of $K_2$ is
unipotent modulo $g_0$, i.e.\ $(x-1)^n\widehat L\subset g_0\widehat L$: this is Proposition~\ref{i3:prop-unip}, which
uses only the integrality data of the package (frame vectors in the lattice, integral Gram matrix, integral weight
idempotents) and $g_0\mid g$, $n\ge2r_S+k+1$.  The prime $g_0$ splits completely in $E$ by the host clause (A).
(These are the level-character idempotent conditions on $(K_2,\chi_2,e)$ that the formal row tie
establishes; the remaining conditions on $e$ in \S\ref{i3:ssec-counting} are clauses of the Hida setting
of the package.)
\emph{Signature and instance tie.}  The instance of $X$ has at every infinite place $p=k+r_S$ plus signs among
$n=k+2r_S+1$, so $q=r_S+1\ge1$.  The instance tie is
obtained from four inputs: the instance of $X$ is the instance of $\mathcal D$ (hypothesis), which also gives $q\ge1$
as just computed; the volume tie (hypothesis); the archimedean scale tie, a record of the datum; and (Inst), which
supplies volume--signature independence and the class items.
\end{proof}

\begin{proof}[Proof of Proposition~\ref{i3:prop-conc1}]
The hypotheses of Propositions~\ref{i3:prop-terms}, \ref{i3:prop-kpp-off},
\ref{i3:prop-kpp-diag}, \ref{i3:prop-diag-rows}, \ref{i3:prop-S-closed},
\ref{i3:prop-rowtie} and \ref{i3:prop-counting} are among those of
Proposition~\ref{i3:prop-conc1}.  Combine their $\varepsilon_1$ and threshold
systems (\S\ref{i3:ssec-setting}) with the Large conditions listed in the
statement, let $\mathcal F$ be large, take $x$ beyond the thresholds of the
``eventually'' clauses of these propositions for $\mathcal F$, and let a member
satisfy the standing hypotheses at $(3,C_{k,\rm low},C_l^{\mathcal F})$.

\emph{Item 1.}  If $u\ne u'$, every degenerate term vanishes
(Proposition~\ref{i3:prop-kpp-off}), so $\mathrm{part}(\mathrm{deg})=0$ and
$K''(u,u')=0$.  For $u=u'$: by Lemma~\ref{i3:lem-pairs}(1),(3),(4) with
$S=\mathrm{deg}$, and since $j$ is isometric,
$\|K''(u,u)\|_{HS}^2=\|\mathrm{part}(\mathrm{deg})\|^2\le P_g\Phi(\mathrm{deg})
=P_g\Sigma(\mathrm{all}(\mathrm{deg}))\le P_g|\Sigma|(\mathrm{all}(\mathrm{deg}))
\le d_ue^{-kf\log k}$ by Proposition~\ref{i3:prop-kpp-diag}.

\emph{Item 2.}  For $u\ne u'$, $K-K''=K$ by item 1, and
Proposition~\ref{i3:prop-S-closed} (second inequality) gives the bound.

\emph{Item 3.}  By (ID), (AC) and completeness, the four sets
$\mathrm{id},\mathrm{deg},\mathrm{gen},\mathrm{ng}$ partition $\tau$, so
$\sum_t\mathrm{op}_t=K_0+\mathrm{part}(\mathrm{deg})$.  Applying $j$:
$K(u,u)-K''(u,u)-\mathrm{id}=j(K_0-\mathrm{part}(\mathrm{id}))$.
Lemma~\ref{i3:lem-diagabstract} with $P=P_g$, $A'=\Gamma_B$, $\varepsilon=\varepsilon_L$ and the estimates
of Proposition~\ref{i3:prop-diag-rows}, then Lemma~\ref{i3:lem-diagnum} (applicable by the Large condition
$C_\Gamma\ge\max(C_B,C_4)+1$ of the statement), give
$\|K(u,u)-K''(u,u)-\mathrm{id}\|_{HS}^2\le\Gamma d_u^2/D_*$.

\emph{Count bound.}  Let $C_g=C_{YL}+1$, $A=2C_g$, $n_0=n_0(A)$, $C=C_0(A)$ as
in Proposition~\ref{i3:prop-counting}.  The Large conditions give $2\le C_l$,
$2C\le C_l$, $8C_\Gamma\le C_{k,\rm low}$, and, as in the proof of
Proposition~\ref{i3:prop-irr-row}, for $y=x^{\varepsilon/8}$
beyond $Y_1=\max(1,n_0/(8C_\Gamma),16/C_\Gamma)$ one has $n_0\le k$ and
$\log L\le\frac1{16}k$; the size clauses (1)--(25) used by
Lemma~\ref{i3:lem-countclauses} are part of the standing hypotheses assumed
here, and (SW) is assumed for Proposition~\ref{i3:prop-counting}.
Lemma~\ref{i3:lem-countclauses} shows that the problem $I$ of
Proposition~\ref{i3:prop-rowtie} satisfies the hypotheses of
Proposition~\ref{i3:prop-counting}, whence
$|D_{\rm all}-D_*|\le e^{-\frac1{16}rf\log r}D_*$, i.e.\ $c_7=\frac1{16}$.
\end{proof}

\begin{remark}\label{i3:rem-c1-status}
After this reduction, Theorem~\ref{i3:thm-linkS-slots} rests on the formally verified
Propositions~\ref{i3:prop-S-closed} (through its own leaves), \ref{i3:prop-counting}, \ref{i3:prop-terms},
\ref{i3:prop-kppoff-in} (item (NF) and the integrability in (FD); (NR), (Levi) and the existence of $\Omega_t$ are proved),
\ref{i3:prop-levtr} (to which item (b) of Proposition~\ref{i3:prop-St-rest} is reduced; its items (a), (c) are proved, (c) from (NF);
(A) and (S) are proved), \ref{i3:prop-kppdiag-in} (items (J), (Del), (Cos),
(Dim); (Irr) and (Num) are proved) and \ref{i3:S-prop-NG}.  Propositions~\ref{i3:prop-chi2}, \ref{i3:prop-kpp-off}, \ref{i3:prop-cov}, \ref{i3:prop-VGeo}, \ref{i3:prop-kpp-diag},
\ref{i3:prop-diag-rows}, \ref{i3:prop-rowtie}, \ref{i3:prop-unip} and Lemmas~\ref{i3:lem-K2group}, \ref{i3:lem-cnthyp} are proved
from these.  Theorem~\ref{i3:thm-linkIRR} (through Proposition~\ref{i3:prop-irr-row}) rests on Propositions~\ref{i3:prop-counting},
\ref{i3:prop-W-loc} (items (II)--(V); (I) is proved; Proposition~\ref{i3:prop-W} is proved from it) and
\ref{i3:prop-allrows-built}, and on the cited facts (HC), (GG) through Proposition~\ref{i3:prop-admLq-in}, which is proved
(Proposition~\ref{i3:prop-admLq} is proved from it, with (Sq) assuming (G)).  The structure of the formal proofs of all these
statements is described in \S\ref{i3:ssec-structure}.  Besides these, the
hypotheses are (SW), the slot hypotheses (Sel), (Rk), (Lp), (Ty), (Mn), (Inst)
for the first, and (SW), (G), $\mathrm{TuSize}$ and (AT) for the second.
\end{remark}


\subsection{Bessel uniqueness and the amplitude identity}\label{i3:ssec-grp-bessel}

We prove Propositions~\ref{i3:prop-bessel} and~\ref{i3:prop-amp}.

\subsubsection{Proof of Proposition~\ref{i3:prop-bessel}}\label{i3:ssec-pf-bessel}

Throughout, $\|\cdot\|=\|\cdot\|_{HS}$, $\mu$ is a Haar measure on $U$ and $A_B=2pq+1$, $A_G=n^2$, $A_S=n^3$,
$a_M=4n^3$.  Since $p,q\ge1$ (the signature of $H$ is $(k+r,r+1)$), $A_B\le A_G$, $A_G+1+A_B\le4n^3$ and
$a-A_G-1\ge A_S+1$ for $a\ge a_M$.

\medskip\noindent\emph{Step 0: numerics and elementary facts.}

\begin{lemma}\label{i3:lem-thr}
Put $\mathrm{thr}^*(\epsilon)=\max(npq+1,16nq_0+8\sqrt{2q_0})$.  Then $\mathrm{thr}^*(\epsilon)\le40n^3$.  Consequently, if
$h\ge120n^3$ then $\mathrm{thr}^*(\epsilon)+a_M<h$.
\end{lemma}

\begin{proof}
$p,q,q_0\le n$ give $npq+1\le2n^3$, $16nq_0\le16n^3$ and $8\sqrt{2q_0}\le16n\le16n^3$; and $40n^3+4n^3<120n^3$.
\end{proof}

\begin{lemma}\label{i3:lem-hs1}
For $g\in U$: $\|g\|\ge1$ and $\|g^{-1}\|\le n\|g\|$.  For $x\in M_n(\C)$, $\|x\|\le n\|g\|\,\|gx\|$.
\end{lemma}

\begin{proof}
$H^2=1$ and $H=H^*$, so $\|H\|_{op}=1$ and $1\le\|g^*\|_{op}\|g\|_{op}\le\|g\|^2$.  Next
$g^{-1}=Hg^*H$ and $\|H\|^2=n$ give the second bound, and $x=g^{-1}(gx)$ the third.
\end{proof}

\begin{lemma}\label{i3:lem-ktype}
If $\lambda$ is a Harish-Chandra parameter then $\Lambda_\lambda$ is integral and $\Lambda_b\le\Lambda_a$ whenever $a<b$
and $(a,b)$ is compact.  Hence there are an irreducible unitary representation $(\tau,V_\Lambda)$ of $K$ and a unit
vector $w_0\in V_\Lambda$ with $d\tau(e_{aa})w_0=\Lambda_aw_0$ and $d\tau(e_{ab})w_0=0$ for compact $a<b$.
\end{lemma}

\begin{proof}
For an index $a$ let $A_n,B_n$ (resp.\ $A_c,B_c$) count the $b>a$, $b<a$ with $(a,b)$ non-compact (resp.\ compact).
Then $2(\rho_n-\rho_c)(a)\equiv A_n+B_n+A_c+B_c=n-1\pmod2$, so $\Lambda_a\in\Z$.  For compact $a<b$ let $k\ge1$ count the
$c\in(a,b]$ with $\epsilon_c=\epsilon_a$ and $j$ the remaining $c\in(a,b)$; then $\rho_c(a)-\rho_c(b)=k$,
$\rho_n(a)-\rho_n(b)=j$, and $\lambda_a-\lambda_b\ge b-a=j+k$, so $\Lambda_a-\Lambda_b\ge2j\ge0$.  The restrictions of
$\Lambda$ to the two blocks are thus dominant integral weights of $U(p)$, $U(q)$; Theorem~\ref{i3:thm-knapp} gives
irreducible representations with these highest weights, their outer tensor product is an irreducible representation of
$K$, unitary for an invariant inner product, and the product of normalised highest weight vectors is $w_0$.
\end{proof}

\medskip\noindent\emph{Step 1: the model.}

Let $R$ be the right regular representation of $U_\epsilon$ on $L^2(U_\epsilon)$ and $E_{ab}$ its derivatives.  For a
space $\mathcal G$ of smooth vectors stable under the $E_{ab}$ put $N_j(x)=\sum_{l\le j}\sum_{|w|=l}\|E_wx\|$, the inner sum
over words $w$ in the index pairs, and $N_j^\oplus(f)=\sum_cN_j(f_c)$ on $\mathcal G^n$.  On $\mathcal G^n=\mathcal G\otimes\C^n$
let $(T_{ab}f)_c=E_{ab}f_c+\delta_{ca}f_b$ (the tensor product with the standard representation), let $K$ act by
$(k\cdot f)_c=\sum_dk_{cd}R(k)f_d$, and put
\[
 (\Omega f)_s=\sum_rE_{rs}f_r,\qquad \mathsf B(f,f')=\sum_c\epsilon'_c\langle f_c,f'_c\rangle,\qquad a_t=\lambda_t-\tfrac{n-1}2 .
\]
Write $\mu_\lambda=\lambda+\rho\in\Z^n$, $\mu_{n+1}=0$, $\Delta_i(z)$ for the leading principal $i\times i$ minor and
$\Phi_\lambda(z)=\prod_i\Delta_i(z)^{-(\mu_i-\mu_{i+1})}$, a holomorphic function where all leading minors are nonzero.

\begin{proposition}[formally verified, \cite{CharacterFamily}]\label{i3:prop-model}
Let $p,q\ge1$, $\lambda$ a Harish-Chandra parameter with gap $h>\mathrm{thr}^*(\epsilon)$, and $(\tau,V_\Lambda,w_0)$ as in
Lemma~\ref{i3:lem-ktype}.
\begin{enumerate}
\item[(a)] (The generating function.)  There is a smooth $\Psi:U_\epsilon\to\mathrm{End}(V_\Lambda)$, square integrable, with
 $\Psi(k_1gk_2)=\tau(k_1)\Psi(g)\tau(k_2)$, $\Psi(1)=1$ and $\sum_{a,b}E_{ab}E_{ba}\Psi=c(\lambda)\Psi$; it is obtained by
 averaging $\Phi_\lambda(k^{-1}gk)$ against the projection on $w_0$ over $K$, and is unique with these properties.  Let
 $\mathcal G_\lambda$ be the smallest $E_{ab}$-stable space containing the functions $g\mapsto\langle w_0,\Psi(g)v\rangle$
 ($v\in V_\Lambda$), $H_\lambda$ its closure in $L^2(U_\epsilon)$ and $\sigma_\lambda=R|_{H_\lambda}$.
\item[(b)] (The model is a discrete series.)  $H_\lambda$ is $R$-invariant, $\mathcal G_\lambda$ consists of smooth vectors
 of $\sigma_\lambda$, the $K$-finite vectors of $\sigma_\lambda$ are exactly $\mathcal G_\lambda$, every discrete series with
 parameter $(\epsilon,\lambda)$ is unitarily equivalent to $\sigma_\lambda$, and $\mathcal G_\lambda$ is dense in the smooth
 vectors of $\sigma_\lambda$ for every finite family of seminorms $u\mapsto\|d\sigma_\lambda(W)u\|$.  Moreover, for a
 discrete series $\sigma$ with parameter $(\epsilon,\lambda)$ and a smooth $w_0\ne0$, a linear map $E=\sum_i\beta_i\circ
 d\sigma(W_i)$ on smooth vectors ($\beta_i$ bounded) which vanishes on $\sigma(U)w_0$ vanishes on all smooth vectors.
\item[(c)] (Schur's lemma on the model.)  A sesquilinear form $b$ on $\mathcal G_\lambda$ with
 $b(E_{ac}x,y)=s_{ac}b(x,E_{ca}y)$ and $b(R(k)x,R(k)y)=b(x,y)$ is a multiple of $\langle\cdot,\cdot\rangle$.
\item[(d)] (Casimir on the tensor product.)  Every $f\in\mathcal G_\lambda^n$ is a finite sum of $\Omega$-eigenvectors in
 $\mathcal G_\lambda^n$, with eigenvalues among $a_1,\dots,a_n$; if $\lambda+e_s$ has two equal coordinates, $a_s$ is not
 an eigenvalue.
\item[(e)] (Shifted parameters.)  If $\lambda'=\lambda+e_s$ is a Harish-Chandra parameter of gap $h'>\mathrm{thr}^*(\epsilon)$,
 there is an injective linear $\iota:\mathcal G_{\lambda'}\to\mathcal G_\lambda^n$ with $\iota E_{ab}=T_{ab}\iota$,
 $\iota R(k)=k\cdot\iota$, whose image is the $a_s$-eigenspace of $\Omega$; and if $F\in\mathcal G_\lambda^n$, $z\in\mathcal
 G_{\lambda'}$ satisfy $\langle z,y\rangle=\sum_c\langle F_c,(\iota y)_c\rangle$ for all $y$, then
 $\langle z,R(g)y\rangle=\sum_c\langle F_c,((R\otimes\mathrm{std})(g)\iota y)_c\rangle$ for all $g$.
\end{enumerate}
\end{proposition}

\noindent\emph{What the proof involves.}  (a) is a computation with $\Phi_\lambda$: its decay along $\mathfrak p$ under
$h>8nq_0$, the Casimir and Gelfand equations, and Haar measure in polar coordinates; (b) uses a Casimir gap between the
layers $V_{k+1}=V_k+\sum_{(a,b)\ \mathrm{noncompact}}E_{ab}V_k$ (a Bochner-type inequality, the only place where
$h>npq+1$ enters), rigidity of data with the same lowest $K$-type, and an $L^2$-regularity lemma for Casimir
eigenfunctions; (c)--(e) are a trace computation for the eigenprojections of $\Omega$, positivity of an invariant form on
$\mathcal G_\lambda^n$ and a trace inequality comparing formal degrees, $d(\lambda+e_s)=r_s(\lambda)d(\lambda)$.  The
proof is long, and we do not reproduce it.  Its formally verified steps are outlined in \S\ref{i3:str-model}.

\medskip\noindent\emph{Step 2: polynomial decay by induction on the gap.}

\begin{lemma}\label{i3:lem-shift}
Let $\lambda$ be a Harish-Chandra parameter with gap $h$ and $s\le n$.  Then $\lambda+e_s$ has gap at least $h-1$; if it is
not a Harish-Chandra parameter, two of its coordinates coincide; and $a_1,\dots,a_n$ are distinct.
\end{lemma}

\begin{proof}
$|(\lambda+e_s)_a-(\lambda+e_s)_b|\ge|\lambda_a-\lambda_b|-1$.  The coordinates stay in $\frac{n-1}2+\Z$, and since
$\lambda_a-\lambda_b\ge1$ for $a<b$, $\lambda+e_s$ is non-increasing; if not strictly decreasing, two coordinates coincide.
The $a_t$ are distinct because $\lambda$ is strictly decreasing.
\end{proof}

\begin{lemma}\label{i3:lem-split}
Under the hypotheses of Proposition~\ref{i3:prop-model} there are linear maps $P_1,\dots,P_n$ of $\mathcal G_\lambda^n$,
commuting with the $T_{ab}$ and with $K$, with $\sum_sP_s=1$, $\Omega P_s=a_sP_s$ and
$N_j^\oplus(P_sf)\le C_jN_{j+n-1}^\oplus(f)$, such that (i) $P_s=0$ if $\lambda+e_s$ is not a Harish-Chandra parameter;
(ii) if $\lambda'=\lambda+e_s$ is one and $\mathrm{thr}^*(\epsilon)<h-1$, there are $\kappa\in\R\setminus\{0\}$, a linear
map $\mathsf q:\mathcal G_\lambda^n\to\mathcal G_{\lambda'}$ with $N_j(\mathsf qf)\le K_jN^\oplus_{j+n-1}(f)$, and
\[
 \sum_c\epsilon'_c\bigl\langle f'_c,((R\otimes\mathrm{std})(g)P_sf)_c\bigr\rangle=\kappa\,\langle\mathsf qf',R(g)\mathsf qf\rangle
 \qquad(f,f'\in\mathcal G_\lambda^n,\ g\in U_\epsilon).
\]
\end{lemma}

\begin{proof}
\emph{Properties of $\Omega$.}  With $T_{ab}=E_{ab}\otimes1+1\otimes e_{ab}$ one computes
$\sum_{a,b}T_{ab}T_{ba}=\mathrm{Cas}(E)\otimes1+1\otimes\mathrm{Cas}(e)+2\Omega$, so $\Omega$ commutes with the $T_{ab}$; it
commutes with $K$ because $\sum_{a,b}E_{ab}\otimes e_{ba}$ is invariant under simultaneous conjugation by a unitary $k$;
$N^\oplus_j(\Omega f)\le nN^\oplus_{j+1}(f)$; and $\mathsf B(\Omega f,f')=\mathsf B(f,\Omega f')$, by
$\langle E_{rs}x,y\rangle=s_{rs}\langle x,E_{sr}y\rangle$ and $\epsilon'_ss_{rs}=\epsilon'_r$.  Also $\mathsf B$ is
$K$-invariant and $\mathsf B(T_{ab}f,f')=s_{ab}\mathsf B(f,T_{ba}f')$.

\emph{Projections.}  The $a_t$ being distinct (Lemma~\ref{i3:lem-shift}), let $\ell_s$ be the Lagrange polynomials at the
$a_t$ and $P_s=\ell_s(\Omega)$.  The commutation properties, the seminorm bound and $\mathsf B$-symmetry follow from those
of $\Omega$; by Proposition~\ref{i3:prop-model}(d) every $f$ is a sum of eigenvectors with eigenvalues among the $a_t$, so
$\sum_sP_s=1$ and $\Omega P_s=a_sP_s$.  (i) If $\lambda+e_s$ is not a Harish-Chandra parameter, two coordinates coincide,
and $P_sf$ is an $a_s$-eigenvector, hence $0$ by Proposition~\ref{i3:prop-model}(d).

(ii) Let $\iota$ be as in Proposition~\ref{i3:prop-model}(e).  The form $b(x,y)=\mathsf B(\iota x,\iota y)$ on
$\mathcal G_{\lambda'}$ satisfies the hypotheses of Proposition~\ref{i3:prop-model}(c), so $b=\kappa\langle\cdot,\cdot\rangle$, and
$\kappa\in\R$ since $\mathsf B$ is Hermitian.  If $\kappa=0$, take $w\ne0$ and $G=(\epsilon'_c(\iota w)_c)_c$; then
$\mathsf B(\iota w,G)=\sum_c\|(\iota w)_c\|^2>0$, but $\mathsf B(\iota w,G)=\mathsf B(\iota w,P_sG)$ and $P_sG=\iota y$ for some
$y$, so $\mathsf B(\iota w,G)=\kappa\langle w,y\rangle=0$, a contradiction.  Define $\mathsf qf$ by $P_sf=\iota(\mathsf qf)$;
then $\mathsf qT_{ab}=E_{ab}\mathsf q$ and $\kappa\|\mathsf qf\|^2=\mathsf B(P_sf,P_sf)$, which with the bound for $P_s$ and
$N^\oplus_i(T_{ab}f)\le2N^\oplus_{i+1}(f)$ gives the bound for $\mathsf q$.  Finally, with $F_c=\epsilon'_cf'_c$ and
$z=\kappa\,\mathsf qf'$ one has, for $y\in\mathcal G_{\lambda'}$,
$\langle z,y\rangle=\mathsf B(P_sf',\iota y)=\mathsf B(f',\iota y)=\sum_c\langle F_c,(\iota y)_c\rangle$, and the last clause
of Proposition~\ref{i3:prop-model}(e) with $y=\mathsf qf$ is the stated identity.
\end{proof}

\begin{lemma}\label{i3:lem-induct}
Let $p,q\ge1$, $m\in\N$, and $\lambda$ a Harish-Chandra parameter with gap $h>\mathrm{thr}^*(\epsilon)+m$.  There is $C$
such that for all $x,y\in\mathcal G_\lambda$, $g\in U_\epsilon$ and index pairs $(a_t,b_t)_{t\le m}$,
\[
 |\langle y,R(g)x\rangle|\prod_{t=1}^m|g_{a_tb_t}|\le C\,N_{m(n-1)}(x)\,N_{m(n-1)}(y).
\]
\end{lemma}

\begin{proof}
Induction on $m$, simultaneously for all $(\lambda,h)$.  For $m=0$ this is Cauchy--Schwarz.  For the step, only the
component $c=a$ of $y\otimes e_a$ and $d=b$ of $x\otimes e_b$ contribute, so
\[
 \epsilon'_a\,g_{ab}\,\langle y,R(g)x\rangle=\sum_c\epsilon'_c\bigl\langle(y\otimes e_a)_c,((R\otimes\mathrm{std})(g)(x\otimes e_b))_c\bigr\rangle .
\]
Insert $\sum_sP_s=1$ in front of $x\otimes e_b$ (Lemma~\ref{i3:lem-split}).  The terms with $\lambda+e_s$ not a parameter
vanish.  For the others $\lambda'=\lambda+e_s$ has gap $\ge h-1>\mathrm{thr}^*(\epsilon)+m$ (Lemma~\ref{i3:lem-shift}), so
the shifted parameter is again a discrete-series parameter with smaller gap and the induction hypothesis applies to it;
by Lemma~\ref{i3:lem-split}(ii) the $s$-th term times $\prod_{t\ge1}|g_{a_tb_t}|$ is at most
$|\kappa|C'N_{m(n-1)}(\mathsf q(x\otimes e_b))N_{m(n-1)}(\mathsf q(y\otimes e_a))$, which is
$\le|\kappa|C'K_{m(n-1)}^2N_{(m+1)(n-1)}(x)N_{(m+1)(n-1)}(y)$.  Sum over $s$.
\end{proof}

\begin{lemma}\label{i3:lem-decay}
Let $\lambda$ be a Harish-Chandra parameter with gap $h>\mathrm{thr}^*(\epsilon)+m$, $m\in\N$, and $\sigma$ a discrete
series of $U$ with parameter $(\epsilon,\lambda)$.  There are $C$ and words $W_j$ in $\mathrm{Lie}(U)$ of length
$\le m(n-1)$ with
\[
 \|g\|^m|\langle v,\sigma(g)u\rangle|\le C\,\nu(u)\nu(v),\qquad \nu(u)=\sum_j\|d\sigma(W_j)u\|\quad(u,v\ \text{smooth},\ g\in U).
\]
In particular, for $b\le m$, $|\langle v,\sigma(g)u\rangle|\le C\nu(u)\nu(v)\|g\|^{-b}$.
\end{lemma}

\begin{proof}
Transport to $U_\epsilon$ (conjugation by $P$ preserves $\|\cdot\|$ and smooth vectors).  If $p=0$ or $q=0$ then
$\|g\|^2=n$ and Cauchy--Schwarz suffices.  Otherwise, since $\|g\|\le\sum_{a,b}|g_{ab}|$, expanding $\|g\|^m$ into
$n^{2m}$ monomials and applying Lemma~\ref{i3:lem-induct} to each gives the bound on $\mathcal G_\lambda$ with $N_{m(n-1)}$.
The words $E_w$ in the complexified $E_{ab}$ are finite linear combinations of words in real elements of
$\mathfrak g_\epsilon$ of the same length, so $N_{m(n-1)}\le\nu$ for suitable $W_j$.  By density of $\mathcal G_\lambda$ in the
smooth vectors of $\sigma_\lambda$ for $\nu$ (Proposition~\ref{i3:prop-model}(b)), an $e$-approximation and $e\to0$ extend
the bound to all smooth vectors of $\sigma_\lambda$; and a unitary equivalence $\sigma_\lambda\simeq\sigma$
(Proposition~\ref{i3:prop-model}(b)) maps smooth vectors to smooth vectors and intertwines the $d\sigma(W_j)$, so the
bound holds for $\sigma$.  The last clause follows from $\|g\|\ge1$ (Lemma~\ref{i3:lem-hs1}).
\end{proof}

\noindent By Lemma~\ref{i3:lem-thr}, under $h\ge120n^3$ Lemma~\ref{i3:lem-decay} applies with $m=a_M=4n^3$: every
coefficient of two smooth vectors of $\sigma$ is $O(\|g\|^{-4n^3})$.

\medskip\noindent\emph{Step 3: integrability against the Bessel character.}

\begin{theorem}[Integration in Cartan coordinates; Helgason, \emph{Groups and Geometric Analysis}, Academic Press 1984, Ch.~I, Thm.~5.8]\label{i3:thm-helgason}
Let $G$ be a connected semisimple Lie group with finite centre, $G=KA^+K$ a Cartan decomposition.  For suitable
normalisations of the Haar measures and $f\ge0$,
$\int_Gf\,dg=\int_K\int_{\mathfrak a^+}\int_Kf(k_1\exp(H)k_2)\prod_{\alpha\in\Sigma^+}(\sinh\alpha(H))^{m_\alpha}\,dk_1\,dH\,dk_2$.
\end{theorem}

\noindent It is applied to $SU(p,q)$, connected semisimple with finite centre.

\begin{lemma}\label{i3:lem-vol}
If $A\ge A_B=2pq+1$ then $y\mapsto\|x^{-1}y\|^{-A}$ is $\mu$-integrable on $U$ for every $x\in U$.
\end{lemma}

\begin{proof}
By left invariance take $x=1$.  $U(1)\times SU(p,q)\to U(p,q)$ has finite kernel and preserves $\|\cdot\|$, so it suffices to
bound volumes in $SU(p,q)$ ($q\le p$).  For $g=k_1\exp(H)k_2$, $\|g\|^2=(p-q)+2\sum_i\cosh(2t_i)$, so $\|g\|\le R$ forces
$0<t_i\le\log R$.  The positive restricted roots are $t_i\pm t_j$ (multiplicity $2$), $t_i$ (multiplicity $2(p-q)$) and $2t_i$
(multiplicity $1$), and $\sinh x\le e^x$ bounds the density of Theorem~\ref{i3:thm-helgason} by $e^{\sum_ic_it_i}$ with
$c_i=4(q-i)+2(p-q)+2$, $\sum_ic_i=2pq$.  Hence $\mu\{\|g\|\le R\}\le CR^{2pq}$, and a dyadic decomposition over
$2^j\le\|g\|<2^{j+1}$ (using $\|g\|\ge1$) gives $\int\|g\|^{-A}d\mu\le C'\sum_j2^{-j(A-2pq)}<\infty$.
\end{proof}

\begin{definition}\label{i3:def-beop}
For a unitary representation $\sigma$ of $U$, vectors $u,u'$, put
$b_{u,u'}(h,m)=\langle u',\sigma(mh)u\rangle\overline{\psi(m)}\pi(h)^{-1}\in\mathrm{End}(W)$ on $U_z\times N$ and, when it is
integrable for $\nu=\mu_W\times\mu_N$ ($\mu_W(U_z)=1$, $\mu_N$ a Haar measure on $N$),
$\beta(u,u')=\int b_{u,u'}\,d\nu$ (the \emph{Bessel operator}).  For $e,f\in W$, $e\otimes\bar f$ is $\eta\mapsto\langle
f,\eta\rangle e$.  Fix a unit vector $w$ with $|\langle w,\sigma(x)w\rangle|\le C_w\|x\|^{-a}$, $a\ge a_M$; the \emph{flat
vectors} are $V_{\rm fl}=\{u:|\langle w,\sigma(x)u\rangle|=O(\|x\|^{-a})\}$, and
$\mathcal D=\{\int_U\phi(g)\sigma(g)w\,d\mu:\phi\ \text{a test function}\}$, a test function being the restriction of a
compactly supported smooth function on $M_n(\C)$.  A \emph{multiplicity-one map} is a map $l$, linear on $V_{\rm fl}$, with
$l(\sigma(mh)u)=\psi(m)\pi(h)l(u)$ and $\beta(u,u')=l(u)\otimes\overline{l(u')}$ on $V_{\rm fl}$.
\end{definition}

\noindent Here $\psi=\psi_{z,c_{\rm std}}$, $c_{\rm std}=-2i$, and $z$ is supported on the middle block; the reduction to this
case is Lemma~\ref{i3:lem-torconj}.

\begin{proposition}[formally verified, \cite{CharacterFamily}]\label{i3:prop-GKin}
Let $\sigma$ be an irreducible unitary representation of $U$ with square-integrable coefficients and $w$, $a\ge a_M$ as in
Definition~\ref{i3:def-beop}.  Then:
\begin{enumerate}
\item[(c)] if $u,v$ satisfy $|\langle w,\sigma(x)u\rangle|=O(\|x\|^{-b_1})$ and $|\langle w,\sigma(x)v\rangle|=O(\|x\|^{-b_2})$ then
 $|\langle v,\sigma(x)u\rangle|=O(\|x\|^{-(\min(b_1,b_2)-A_G-1)})$;
\item[(d)] if $|\langle v,\sigma(x)u\rangle|=O(\|x\|^{-b})$ with $b\ge A_S+1$ then $b_{u,v}$ is $\nu$-integrable;
\item[(e)] $(h,m)\mapsto\|mh\|^{-b}$ is $\nu$-integrable for $b\ge A_B$;
\item[(f)--(h)] for $u,u'\in V_{\rm fl}$: $\langle\xi,\beta(u,u')\eta\rangle=\langle\beta(u',u)\xi,\eta\rangle$,
 $\Re\langle\xi,\beta(u,u)\xi\rangle\ge0$, and $\beta(\sigma(mh)u,u')=\psi(m)\pi(h)\beta(u,u')$;
\item[(GK)] (Gelfand--Kazhdan inputs) there are a map $\theta:U\to U$ with $\theta(U_z)\subset U_z$, a constant $c>0$, a class
 $\mathcal C$ of functions on $U$ (five times differentiable with derivatives $O(\|g\|^{-(a-A_G-1)})$) and a pairing
 $\mathcal B(f,F')$ on scalar times $\mathrm{End}(W)$-valued functions in $\mathcal C$, satisfying (GK1)--(GK5) of
 Lemma~\ref{i3:lem-GKsep}.
\end{enumerate}
\end{proposition}

\noindent\emph{What the proof involves.}  (c)--(h) are estimates and Fubini arguments using an explicit decomposition of
$U$ adapted to $U_zN$ (Iwasawa, Levi, a simultaneous Witt lemma) and Lemma~\ref{i3:lem-vol}; (GK) is the construction of the
involution $\theta$ of $U$ preserving $U_zN$ (an anti-automorphism composed with conjugation by a $z$-frame), the
decomposition of the big cell into $U_zN$-double cosets with an explicit iterated integration formula (Heisenberg levels,
the isotropic cone and a hyperboloid measure), and the regularised pairing $\mathcal B=c\,\mathcal B_\infty$ on the big
cell together with its convergence (this is where the numerical conditions of \S\ref{i3:ssec-gap} enter).  The
proofs are long, and we do not reproduce them.  Its formally verified steps are outlined in \S\ref{i3:str-GKin}.

\begin{proposition}\label{i3:prop-orbbd}
Let $k\ge1$, $\sigma$ a unitary representation of $U=U(H_{k,r})$, $z$ supported on the middle block with $T(z)>0$,
$c\in(\C^\times)^r$, $\pi$ a unitary representation of $U_z$ on a finite-dimensional space $W$, continuous on $U_z$
with $\pi(1)=1$, and $B$ a smooth Bessel map for $(z,\psi_{z,c},\sigma,\pi)$ (Definition~\ref{i3:def-bessel}).  Then
$g\mapsto\|B(\sigma(g)w)\|$ is bounded on $U$ for every smooth $w$.
\end{proposition}

\noindent The hypotheses $k\ge1$, $T(z)>0$, $z$ on the middle block, $c_a\ne0$ and the continuity of $\pi$ are those of
the formal statement.  The proof (\S\ref{i3:ssec-pf-orbbd}) writes $g=hm\,t\,k$ with $h\in U_z$, $m\in N$, $t$
in a torus and $k$ unitary; equivariance removes $hm$, and the remaining factor is controlled by the derivatives of
$\sigma$ along a complement of $\mathrm{Lie}(U_zN)$, which the torus contracts, except for one root direction which acts
through $d\psi$.  The proof uses the following notion of smooth vector: $v$ is smooth if near every point of $U$ the orbit
map $g\mapsto\sigma(g)v$ is the restriction of a $C^\infty$ map on an open subset of $M_n(\C)$.  This coincides with the
notion of Definition~\ref{i3:def-bessel} ($g\mapsto\sigma(g)v$ is $C^\infty$ on the Lie group $U$): $U$ is a closed embedded
submanifold of $M_n(\C)$, so restrictions of $C^\infty$ maps are $C^\infty$ on $U$, and every $C^\infty$ map on $U$ extends locally,
through a slice chart, to a $C^\infty$ map on an open subset of $M_n(\C)$.

\medskip\noindent\emph{Step 4: the reproducing integral and Schur orthogonality.}

\begin{lemma}\label{i3:lem-cw}
If $B:H^\infty\to W$ is linear with $\|B(v)\|\le C\sum_i\|d\sigma(W_i)v\|$, there are bounded linear $\beta_i:H\to W$ with
$B(v)=\sum_i\beta_i(d\sigma(W_i)v)$.  In particular $g\mapsto B(\sigma(g)w)$ is continuous for smooth $w$.
\end{lemma}

\begin{proof}
With $J(v)=(d\sigma(W_i)v)_i\in H^N$, $\ker J\subset\ker B$ and $Jv\mapsto Bv$ is bounded on $J(H^\infty)$; extend it to the
closure by continuity and to $H^N$ by orthogonal projection.
\end{proof}

\begin{lemma}\label{i3:lem-cutlim}
Let $B$ be as in Lemma~\ref{i3:lem-cw}, $w$ smooth with integrable coefficients $y\mapsto\langle v,\sigma(y)w\rangle$,
$\kappa\in\C$ with $\int_U\overline{\langle v,\sigma(y)w\rangle}\sigma(y)w\,d\mu=\kappa v$ for smooth $v$, and
$\sup_g\|B(\sigma(g)w)\|<\infty$.  Then
$\kappa B(v)=\int_U\overline{\langle v,\sigma(y)w\rangle}B(\sigma(y)w)\,d\mu(y)$ for smooth $v$.
\end{lemma}

\begin{proof}
Fix a real $\varphi\in C_c^\infty(M_n(\C))$ with $\varphi(0)=1$ and put
$v_N[\varphi;v']=\int\varphi(y/(N+1))\overline{\langle v',\sigma(y)w\rangle}\sigma(y)w\,d\mu$.  These are smooth, and
differentiating under the integral after the substitution $y\mapsto e^{tX}y$ gives
$d\sigma(X)v_N[\varphi;v']=v_N[L_X\varphi;v']+v_N[\varphi;d\sigma(X)v']$ with $(L_X\varphi)(y)=\frac d{dt}\varphi(e^{-tX}y)|_0$.  By
induction on the length of a word $W$ and dominated convergence,
$d\sigma(W)v_N[\varphi;v']\to\varphi(0)\kappa\,d\sigma(W)v'$, the terms with $L_X\varphi$ tending to $(L_X\varphi)(0)\kappa(\cdot)=0$.
Hence $B(v_N[\varphi;v])\to\kappa B(v)$ by Lemma~\ref{i3:lem-cw}.  On the other hand
$B(v_N[\varphi;v])=\int\varphi(y/(N+1))\overline{\langle v,\sigma(y)w\rangle}B(\sigma(y)w)\,d\mu$ (differentiation under a
compactly supported integral), which tends to the stated integral by dominated convergence, the coefficient being
integrable and $B(\sigma(y)w)$ bounded.
\end{proof}

\begin{lemma}\label{i3:lem-normal}
For $h\in U_z$ and $m\in N$: $h$ is the identity outside the middle block, $h^{-1}mh\in N$ and
$\psi_{z,c}(h^{-1}mh)=\psi_{z,c}(m)$.
\end{lemma}

\begin{proof}
$h$ fixes $e_a$ and $e^*_a$ ($a<r$), so its outer columns are identity columns, and $h^*Hh=H$ forces the outer rows to be
identity rows.  The defining conditions of $N$ (identity columns $j\le r+k$ below the diagonal and in the middle rows) are
then preserved by conjugation, by a direct entry-wise check; and $\ell_a(hmh^{-1})=\ell_a(m)$ because the outer rows of $h$
and the outer columns of $h^{-1}$ are those of the identity and $h^{-1}z=z$.
\end{proof}

\begin{lemma}\label{i3:lem-torconj}
There is a diagonal $\tau\in U$ commuting with $U_z$, normalising $N$, with
$\psi_{z,c}(\tau m\tau^{-1})=\psi_{z,c_{\rm std}}(m)$ for $m\in N$.
\end{lemma}

\begin{proof}
Put $T_j=\prod_{j\le l<r}c_{\rm std}/c_l$ (so $T_j=1$ for $j\ge r$, $c_aT_a=c_{\rm std}T_{a+1}$) and $\tau=\mathrm{diag}(d)$ with
$d_i=T_i$ ($i<r$), $d_i=1$ on the middle block, $d_i=\overline{T_{2r+k-i}}^{-1}$ ($i>r+k$).  Then $\bar d_id_j=1$ at the
nonzero entries of $H$, so $\tau\in U$; $\tau$ commutes with $U_z$ by Lemma~\ref{i3:lem-normal}; conjugation by a diagonal
matrix preserves the shape of $N$; and $\sum_ac_a\ell_a(\tau m\tau^{-1})=\sum_ac_aT_aT_{a+1}^{-1}\ell_a(m)=c_{\rm std}\sum_a\ell_a(m)$.
\end{proof}

\begin{lemma}\label{i3:lem-fold}
There is a measurable $\varrho:U\to[0,\infty)$ such that for every integrable $\varphi$ on $U$ (Banach-valued)
$\int_U\varphi\,d\mu=\int_U\varrho(y)\int_{U_z\times N}\varphi((mh)^{-1}y)\,d\nu(h,m)\,d\mu(y)$, the integrand being integrable.
\end{lemma}

\begin{proof}
Let $\beta_0(y)=\|y\|^{-A_B}$ and $O(y)=\int\beta_0(mhy)\,d\nu$.  Then $0<O(y)<\infty$ by Lemma~\ref{i3:lem-hs1} and
Proposition~\ref{i3:prop-GKin}(e), and $O$ is left $U_zN$-invariant: $mh\,m'h'=(m\cdot hm'h^{-1})(hh')$, and
$(h,m)\mapsto(hh',m\,hm'h^{-1})$ preserves $\nu$ (Lemma~\ref{i3:lem-normal}; $N$ is unimodular, being unipotent).  With
$\varrho=\beta_0/O$ one has $\int\varrho(sy)\,d\nu(s)=1$, so $\int\varphi=\iint\varrho(sy)\varphi(y)\,d\nu\,d\mu$, and the
measure-preserving change $(s,y)\mapsto(s,s^{-1}y)$ with Fubini gives the formula.
\end{proof}

\begin{lemma}\label{i3:lem-thmM}
In the setting of Definition~\ref{i3:def-beop} (with $a\ge a_M$) there is a multiplicity-one map $l$.
\end{lemma}

\noindent The proof, which contains the Gelfand--Kazhdan step, is given in Step 5.  We first show how it yields uniqueness
on orbits.

\begin{lemma}\label{i3:lem-fix}
Let $l$ be a multiplicity-one map, $d=d_\sigma$, and $F:U\to W$ measurable with $F(mhx)=\psi(m)\pi(h)F(x)$ and
$y\mapsto\|x^{-1}y\|^{-a}\|F(y)\|$ integrable for each $x$.  Then there is $c$ with
$\int_Ud\,\overline{\langle w,\sigma(y)w\rangle}F(xy)\,d\mu(y)=c\,l(\sigma(x)w)$ for all $x\in U$.
\end{lemma}

\begin{proof}
Let $\varphi_x(y)=d\,\overline{\langle w,\sigma(x^{-1}y)w\rangle}F(y)$, integrable by the decay of $w$; left invariance turns
the left side into $\int\varphi_x$.  For $s=mh$, the equivariance of $F$ gives
$\varphi_x(s^{-1}y)=d\,b_{\sigma(x)w,\sigma(y)w}(h,m)F(y)$.  The vectors $\sigma(x)w,\sigma(y)w$ are flat, so by
Proposition~\ref{i3:prop-GKin}(c),(d) these integrands are integrable and Lemma~\ref{i3:lem-fold} gives
$\int\varphi_x=\int_U\varrho(y)\,d\,\beta(\sigma(x)w,\sigma(y)w)F(y)\,d\mu(y)$.  Since
$\beta(\sigma(x)w,\sigma(y)w)F(y)=\langle l(\sigma(y)w),F(y)\rangle\,l(\sigma(x)w)$, the claim holds with
$c=\int\varrho(y)\,d\,\langle l(\sigma(y)w),F(y)\rangle\,d\mu(y)$, independent of $x$ (if $l(\sigma(x)w)=0$ for all $x$, both
sides vanish).
\end{proof}

\begin{lemma}\label{i3:lem-orbmult}
Let $l$ be a multiplicity-one map for $(\sigma,w)$ with $w$ smooth and all coefficients $\langle v,\sigma(x)w\rangle$ ($v$
smooth) $O(\|x\|^{-a})$, $a\ge a_M$.  If $B_m$ is a smooth Bessel map for $(z,\psi,\sigma,\pi)$, bounded on the orbit of $w$,
then there is $c$ with $B_m(\sigma(x)w)=c\,l(\sigma(x)w)$ for all $x\in U$.
\end{lemma}

\begin{proof}
$F(x)=B_m(\sigma(x)w)$ is continuous (Lemma~\ref{i3:lem-cw}), satisfies $F(mhx)=\psi(m)\pi(h)F(x)$ by
Lemma~\ref{i3:lem-normal} and the equivariance of $B_m$, and is bounded, so Lemma~\ref{i3:lem-vol} gives the integrability
needed in Lemma~\ref{i3:lem-fix}, which yields $c$.  By Theorem~\ref{i3:thm-godement} (pairing with an arbitrary vector),
$\int\overline{\langle v,\sigma(y)w\rangle}\sigma(y)w\,d\mu=d^{-1}v$, and Lemma~\ref{i3:lem-cutlim} applied to
$B_m\circ\sigma(x)$ with $\kappa=d^{-1}$, $v=w$ gives
$B_m(\sigma(x)w)=\int d\,\overline{\langle w,\sigma(y)w\rangle}F(xy)\,d\mu(y)=c\,l(\sigma(x)w)$.
\end{proof}

\medskip\noindent\emph{Step 5: the Gelfand--Kazhdan step.}

\begin{lemma}\label{i3:lem-sepminors}
If $\kappa,\alpha,\gamma:I\times I\to\C$ satisfy $\kappa(A,B)\overline{\kappa(E,F)}=\alpha(A,F)\gamma(B,E)$ for all
$A,B,E,F$, then $\kappa(A,B)\kappa(A',B')=\kappa(A,B')\kappa(A',B)$.
\end{lemma}

\begin{proof}
If $\kappa(E,F)\ne0$ for some $(E,F)$, then $\kappa(A,B)=\alpha(A,F)\gamma(B,E)/\overline{\kappa(E,F)}$, and both sides equal
$\alpha(A,F)\alpha(A',F)\gamma(B,E)\gamma(B',E)/\overline{\kappa(E,F)}^2$.
\end{proof}

\begin{lemma}\label{i3:lem-GKsep}
In the setting of Lemma~\ref{i3:lem-thmM}, let $\theta$, $c>0$, $\mathcal C$ and $\mathcal B$ be such that
\begin{enumerate}
\item[(GK1)] $\mathcal D$ is $\sigma(U)$-stable, $\mathcal D\subset V_{\rm fl}$, and for $u,v\in\mathcal D$ the functions
 $\langle v,\sigma(\cdot)u\rangle$, $\langle v,\sigma(\theta\,\cdot)u\rangle$ and $\langle v,\sigma(\cdot)u\rangle X_0$
 ($X_0\in\mathrm{End}(W)$) lie in $\mathcal C$;
\item[(GK2)] if $f=\langle y,\sigma(\cdot)x\rangle\in\mathcal C$, $\langle b_2,\sigma(\cdot)a_2\rangle\in\mathcal C$ and $b_{x,a_2}$,
 $b_{y,b_2}$ are integrable, then
 $\mathcal B(f,\langle b_2,\sigma(\cdot)a_2\rangle X_0)=\frac cd\,\mathrm{tr}\bigl(X_0^*\beta(y,b_2)^*\beta(x,a_2)\bigr)$;
\item[(GK3)] $\mathcal B(f\circ\theta,F')=\mathcal B(f,F')$;
\item[(GK4)] either $\int_U\langle y,\sigma(\theta g)x\rangle\overline{\langle y',\sigma(g)x'\rangle}\,d\mu=0$ for all vectors and then
 $\mathcal B(\langle b_1,\sigma(\theta\,\cdot)a_1\rangle,\langle b_2,\sigma(\cdot)a_2\rangle X_0)=0$ on $\mathcal D$; or there is a
 conjugate-linear isometric bijection $C$ of $H$ with $\langle y,\sigma(\theta g)x\rangle=\langle Cx,\sigma(g)Cy\rangle$,
 $C\sigma(h)=\sigma(\theta(h)^{-1})C$ ($h\in U_z$), and $b_{Cu,v}$ integrable for $u,v\in\mathcal D$;
\item[(GK5)] calling $(P_1,Q,P_2,Q_2)$ \emph{related} if $P_1^*\pi(h)Q=P_2^*\pi(\theta h)Q_2$ for $h\in U_z$: either
 $\langle\eta,Qx\rangle\langle P_1y,\xi\rangle=0$ for all related quadruples, or there is a conjugate-linear isometric bijection
 $C_0$ of $W$ with $\langle\eta,Qx\rangle\langle P_1y,\xi\rangle=\langle\eta,C_0P_2y\rangle\langle C_0Q_2x,\xi\rangle$ for all related
 quadruples.
\end{enumerate}
Then for $u_1,\dots,u_4\in\mathcal D$ and $e_1,\dots,e_4\in W$,
\begin{equation}\label{i3:eq-sep}
\langle e_1,\beta(u_1,u_2)e_2\rangle\langle e_3,\beta(u_3,u_4)e_4\rangle=\langle e_1,\beta(u_1,u_4)e_4\rangle\langle e_3,\beta(u_3,u_2)e_2\rangle .
\end{equation}
\end{lemma}

\begin{proof}
All $b_{u,v}$, $u,v\in\mathcal D$, are integrable (Proposition~\ref{i3:prop-GKin}(c),(d)).  In the first case of (GK4), for
$a_1,a_2\in\mathcal D$, (GK3), (GK4) and (GK2) give $0=\frac cd\mathrm{tr}(X_0^*\beta(a_1,a_2)^*\beta(a_1,a_2))$ for all $X_0$, so
$\beta=0$ on $\mathcal D$ and \eqref{i3:eq-sep} is trivial.  In the second case, applying (GK2) to
$f=\langle b_1,\sigma(\cdot)a_1\rangle$ and to $f\circ\theta=\langle Ca_1,\sigma(\cdot)Cb_1\rangle$ and comparing by (GK3),
\[
\beta(b_1,b_2)^*\beta(a_1,a_2)=\beta(Ca_1,b_2)^*\beta(Cb_1,a_2),
\]
and replacing $a_1$ by $\sigma(h)a_1$ (using $C\sigma(h)=\sigma(\theta(h)^{-1})C$ and the equivariance
$\beta(\sigma(h)u,u')=\pi(h)\beta(u,u')$; Proposition~\ref{i3:prop-GKin}(h) states it on flat vectors, but it holds,
by the substitution $h'\mapsto h'h^{-1}$ in the defining integral (invariance of the Haar measure of $U_z$), for every pair $(u,u')$ with $b_{u,u'}$ integrable, and
$Ca_1$ need not be flat; here $b_{Cu,v}$ is integrable by (GK4)),
\[
\beta(b_1,b_2)^*\pi(h)\beta(a_1,a_2)=\beta(Ca_1,b_2)^*\pi(\theta h)\beta(Cb_1,a_2)\qquad(h\in U_z).
\]
So $(\beta(b_1,b_2),\beta(a_1,a_2),\beta(Ca_1,b_2),\beta(Cb_1,a_2))$ is related.  In the first case of (GK5), taking
$(b_1,b_2)=(a_1,a_2)$ gives $|\langle\eta,\beta(a_1,a_2)x\rangle|^2=0$.  In the second, with $I=\mathcal D\times W$,
$\kappa((a_1,\eta),(a_2,x))=\langle\eta,\beta(a_1,a_2)x\rangle$ and the evident $\alpha,\gamma$ built from $C_0$, (GK5) is the
hypothesis of Lemma~\ref{i3:lem-sepminors}, whose conclusion is \eqref{i3:eq-sep}.
\end{proof}

\begin{lemma}\label{i3:lem-vanish}
Let $\Lambda:H\to\C$ with $y\mapsto\Lambda(\sigma(y)w)$ measurable, $|\Lambda(\sigma(y)w)|\le C\|y\|^b$ with $a-b\ge A_B$, and
$\Lambda(u)=d\int_U\langle\sigma(y)w,u\rangle\Lambda(\sigma(y)w)\,d\mu(y)$ for $u\in V_{\rm fl}$.  If $\Lambda=0$ on $\mathcal D$ then
$\Lambda=0$ on $V_{\rm fl}$.
\end{lemma}

\begin{proof}
$\mathcal D\subset V_{\rm fl}$ and $\sigma(x)w\in V_{\rm fl}$, by Lemma~\ref{i3:lem-hs1}.  For a test function $\phi$ and
$v_\phi=\int\phi(x)\sigma(x)w\,d\mu$, inserting the reproducing formula at $u=v_\phi$ and using Fubini (justified by
Lemma~\ref{i3:lem-vol} with exponent $a-b$) and the formula at $u=\sigma(x)w$ gives
$\int\phi(x)\Lambda(\sigma(x)w)\,d\mu(x)=\Lambda(v_\phi)=0$.  As $\phi$ is arbitrary, $\Lambda(\sigma(y)w)=0$ almost everywhere,
and then the reproducing formula gives $\Lambda=0$ on $V_{\rm fl}$.
\end{proof}

\begin{proof}[Proof of Lemma~\ref{i3:lem-thmM}]
By Proposition~\ref{i3:prop-GKin}(c),(d) all $b_{u,u'}$ with $u,u'\in V_{\rm fl}$ are integrable, and for $v\in V_{\rm fl}$,
$\|\beta(\sigma(y)w,v)\|\le C_v\|y\|^{b}$ with $b=a-A_G-1$ (bound $\|mhy\|^{-b}\le(n\|y\|)^b\|mh\|^{-b}$ and
Proposition~\ref{i3:prop-GKin}(e)).  \emph{Reproducing formula.}  For $u,v\in V_{\rm fl}$,
$\beta(u,v)=d\int_U\langle\sigma(y)w,u\rangle\beta(\sigma(y)w,v)\,d\mu(y)$: insert Theorem~\ref{i3:thm-godement} in the form
$\langle v,\sigma(s)u\rangle=d\int\langle\sigma(y)w,u\rangle\langle v,\sigma(s)\sigma(y)w\rangle\,d\mu(y)$ into $b_{u,v}$ and use
Fubini on $U_z\times N\times U$ (integrable by the decay bounds and Lemma~\ref{i3:lem-vol}).

\emph{Separation.}  By Proposition~\ref{i3:prop-GKin}(GK) and Lemma~\ref{i3:lem-GKsep}, \eqref{i3:eq-sep} holds on $\mathcal D$.

\emph{Extension.}  Let $\mathrm{Mn}(u_1,u_2,u_3,u_4)$ be the difference of the two sides of \eqref{i3:eq-sep}.  For fixed
$u_2,u_3,u_4$, $\Lambda=\mathrm{Mn}(\cdot,u_2,u_3,u_4)$ satisfies the hypotheses of Lemma~\ref{i3:lem-vanish} (the reproducing
formula is linear in the first slot), so vanishing on $\mathcal D$ in one variable extends to $V_{\rm fl}$.  Using the
antisymmetry $\mathrm{Mn}(u_3,u_2,u_1,u_4;e_3,e_2,e_1,e_4)=-\mathrm{Mn}(u_1,u_2,u_3,u_4;e_1,e_2,e_3,e_4)$ and the conjugation
symmetry $\mathrm{Mn}(u_2,u_1,u_4,u_3;e_2,e_1,e_4,e_3)=\overline{\mathrm{Mn}(u_1,\dots;e_1,\dots)}$ (from (f)), one moves the
variables one at a time from $\mathcal D$ to $V_{\rm fl}$; hence $\mathrm{Mn}=0$ on $V_{\rm fl}^4$.

\emph{Rank one.}  If $\kappa_0=\langle\xi_0,\beta(v_0,v_0)\xi_0\rangle>0$ for some flat $v_0$ and $\xi_0$, put
$l(u)=\kappa_0^{-1/2}\beta(u,v_0)\xi_0$; it is linear and equivariant by (h), and the vanishing minor at
$(u,u',v_0,v_0;\eta,\eta',\xi_0,\xi_0)$ together with (f) gives $\beta(u,u')=l(u)\otimes\overline{l(u')}$.  Otherwise
$\langle\xi,\beta(u,u)\xi\rangle=0$ for all flat $u$ and $\xi$ (it is real by (f) and $\ge0$ by (g)), the minor at
$(u,u',u',u;\eta,\eta',\eta',\eta)$ gives $|\langle\eta,\beta(u,u')\eta'\rangle|^2=0$, and $l=0$ works.
\end{proof}

\medskip\noindent\emph{Step 6: conclusion.}

\begin{proof}[Proof of Proposition~\ref{i3:prop-bessel}]
By hypothesis $z$ is supported on the middle block, and $\pi$ is a continuous irreducible unitary representation of the
compact connected group $U_z\cong U(k)$; hence (theorem of the highest weight) there are a unitary frame of the orthogonal
complement of $z$ in the middle block and a highest weight of $\pi$ with respect to the diagonal torus of that frame.  This
pair (frame, highest-weight representation), with Haar measures on $U$, $U_z$ and $N$, is the real-place datum used in the
setup of \cite{CharacterFamily}.  With $\tau$ as in
Lemma~\ref{i3:lem-torconj}, $v\mapsto B(\sigma(\tau)v)$ is a smooth Bessel map for $\psi_{z,c_{\rm std}}$ whenever $B$
is one for $\psi_{z,c}$, since $\tau hm=h(\tau m\tau^{-1})\tau$.  By Lemma~\ref{i3:lem-thr}, Lemma~\ref{i3:lem-decay} applies with
$m=a_M$, so all coefficients of smooth vectors are $O(\|g\|^{-a_M})$.  Let $w_0$ be a unit vector of the lowest $K$-type
(smooth, by Proposition~\ref{i3:prop-model}(b)); with $w=w_0$, $a=a_M$, Lemma~\ref{i3:lem-thmM} gives a multiplicity-one map $l$.

Let $B_1,B_2$ be smooth Bessel maps.  By Proposition~\ref{i3:prop-orbbd} each $B_i\circ\sigma(\tau)$ is bounded on the orbit
of $w_0$, and Lemma~\ref{i3:lem-orbmult} gives $c_i$ with $B_i(\sigma(\tau)\sigma(x)w_0)=c_i\,l(\sigma(x)w_0)$ for all $x$.
Choose $(a,b)\ne(0,0)$ with $ac_1+bc_2=0$.  Then $E=(aB_1+bB_2)\circ\sigma(\tau)$ is a continuous (Lemma~\ref{i3:lem-cw}) linear
map vanishing on $\sigma(U)w_0$.  By the last clause of Proposition~\ref{i3:prop-model}(b) (transported to $U$), $E$ vanishes on all smooth vectors.  Applied to
$\sigma(\tau^{-1})v$ this gives $aB_1(v)+bB_2(v)=0$.
\end{proof}

\subsubsection{Proof of Proposition~\ref{i3:prop-amp}}\label{i3:ssec-pf-amp}

\begin{proof}
Since $A$ is a good local Bessel datum for $\lambda'$, $\lambda'$ is a Harish-Chandra parameter; by Lemma~\ref{i3:lem-gap}
(applicable as $1000\le C_l$, $C_n\le C_l$) it has gap $h\ge120n^3$, and $h$ satisfies the two cell conditions of
\S\ref{i3:ssec-gap}.  Let $l,l'$ be as in Proposition~\ref{i3:prop-amp-in}(A1) and $x,y,x',y'$ as in (A2).
Since $\beta(x,y)=l(x)\otimes\overline{l(y)}\ne0$, $l(x)\ne0$; similarly $l'(x')\ne0$.  By Proposition~\ref{i3:prop-bessel}
(applied at $z$ and at $z'$, the hypothesis $h\ge120n^3$ holding by Lemma~\ref{i3:lem-gap}) the Bessel maps $B$ and $l$ are
linearly dependent; as $l\ne0$, $B=c_0l$ on smooth vectors for some $c_0\in\C$, and likewise $B'=c_0'l'$.  $K$-finite
vectors are smooth, so
\[
0\ne\langle Ba,\xi\rangle\langle\xi',B'b\rangle=\overline{c_0}\,c_0'\,\langle la,\xi\rangle\langle\xi',l'b\rangle ;
\]
hence $c_0,c_0'\ne0$ and $s=\langle la,\xi\rangle\langle\xi',l'b\rangle\ne0$.  Write $v_0=t\,w_0$ with $t=\|v_0\|>0$.  Then
$\Phi(g)=B(\sigma(g)v_0)=c_0t\,l(\sigma(g)w_0)$ and $\Phi'(g)=c_0't\,l'(\sigma(g)w_0)$ for $g\in U$, so by sesquilinearity of the
cell integral
\[
J_A(\nu,\gamma)=c_N^{-1}\mathcal J_{\nu,\gamma}[\Phi,\Phi']=K_1\,\mathcal J_{\nu,\gamma}\bigl[l(\sigma(\cdot)w_0),l'(\sigma(\cdot)w_0)\bigr],
\qquad K_1=c_N^{-1}(c_0t)\overline{(c_0't)}\ne0 .
\]
By (A3) with $(v,v')=(a,b)$, $s\,\mathcal J_{\nu,\gamma}[l(\sigma(\cdot)w_0),l'(\sigma(\cdot)w_0)]=d_\sigma^{-1}\mathrm{Orb}_{\gamma,\nu}(a,b,\xi,\xi')$,
with $d_\sigma>0$ by Theorem~\ref{i3:thm-godement}.  Therefore
$\mathrm{Orb}_{\gamma,\nu}(a,b,\xi,\xi')=d_\sigma\,s\,K_1^{-1}J_A(\nu,\gamma)$, and $\kappa=d_\sigma sK_1^{-1}\ne0$ does not depend
on $(\gamma,\nu)$.
\end{proof}


\subsection{Hilbert--Schmidt bounds off the diagonal}\label{i3:ssec-grp-S}

We prove Proposition~\ref{i3:prop-S-closed}, the one-level inequality, and the lemmas on which the proof rests.

\subsubsection{Proof of Proposition~\ref{i3:prop-S-closed}}\label{i3:S-ssec-proof}

The proof has four parts.
\begin{enumerate}
\item[(1)] An exact Hilbert-space computation (\S\ref{i3:S-sss-abstract}) bounds the norm of an operator
written as an absolutely convergent sum of \emph{terms} organised by \emph{rows}, through four row
estimates: the main row, the loops, the non-generic row, and a Gram bound.
\item[(2)] The main row (\S\ref{i3:S-sss-main}) is bounded by a recursion on lattice averages over the
levels, followed by an algebraic identity between volumes and masses (the \emph{mass formula}), whose
inputs come from Siegel--Weil.
\item[(3)] The term system of a block $(u,u')$ is constructed from the automorphic kernel
$\sum_\gamma\phi(x^{-1}\gamma y)$; nineteen structural components (\S\ref{i3:S-sss-components}) say that it has
the properties required in (1) and (2); the volume count (\S\ref{i3:S-sss-volume}) is the companion statement for
the diagonal of that kernel.
\item[(4)] The squared norm of the term operator $K_0(u,u')$ equals $\|K(u,u')\|_{HS}^2$ off the diagonal
(\S\ref{i3:S-sss-HS}).
\end{enumerate}
The assembly is \S\ref{i3:S-sss-assembly}.  The archimedean input, proportionality of orbital integrals to
Bessel integrals, is Proposition~\ref{i3:prop-amp}, and Bessel uniqueness under the gap is
Proposition~\ref{i3:prop-bessel}; both enter only through the components (C6), (C12), (C16), (C17).

\medskip\noindent\label{i3:S-sss-abstract}\emph{The abstract row bound.}

\begin{definition}[Term system]\label{i3:S-def-terms}
A \emph{term system} $X$ consists of: sets $\tau$ (terms) and $\mathrm{Rw}$ (rows); a finite set $I$; complex
inner product spaces $H$, $H_1$; maps $\mathrm{op}:\tau\to H$, $\mathrm{vec}:\tau\to H_1$, $\mathrm{row}:\tau\to\mathrm{Rw}$,
$\iota:\mathrm{Rw}\to I$; predicates $\mathrm{isId}$ and $\mathrm{gener}$ on $\tau$; a function $G:\tau\to\R$; a
matrix $\mathbf G\in\C^{I\times I}$ and a number $\gamma_1\in\R$; a function $s_0:\tau\times\tau\to\N$ (the
\emph{corank}); and a relation $\mathrm{scal}$ on $\tau\times\tau$ (\emph{scalar pairs}).  Put
\[
\mathrm{id}=\{\mathrm{isId}\},\quad\mathrm{deg}=\{G=0,\neg\mathrm{isId}\},\quad\mathrm{ng}=\{G\ne0,\neg\mathrm{gener}\},\quad
\mathrm{gen}=\{G\ne0,\mathrm{gener}\},\quad\mathrm{nz}=\{G\ne0\},
\]
$\mathrm{part}(S)=\sum_{t\in S}\mathrm{op}_t$, $K=\sum_t\mathrm{op}_t$,
$K_0=\mathrm{part}(\mathrm{id})+\mathrm{part}(\mathrm{gen})+\mathrm{part}(\mathrm{ng})$,
$v(S,i)=\sum_{t\in S,\,\iota(\mathrm{row}\,t)=i}\mathrm{vec}_t$ and $\Phi(S)=\gamma_1\sum_{i\in I}\|v(S,i)\|^2$ (sums
unconditional).  For $S\subset\tau$ let $\mathrm{all}(S)=\{(t,t')\in S^2:\mathrm{row}\,t=\mathrm{row}\,t'\}$ and inside it
$\mathrm{big}(S)=\{s_0\ge2\}$, $\mathrm{one}(S)=\{s_0=1\}$, $\mathrm{nsc}(S)=\{s_0=1,\neg\mathrm{scal}\}$,
$\mathrm{sc}(S)=\{s_0=1,\mathrm{scal}\}$.  For $Q\subset\tau^2$ put
$\Sigma(Q)=\gamma_1\sum_{p\in Q}\Re\langle\mathrm{vec}_{p_1},\mathrm{vec}_{p_2}\rangle$ and
$|\Sigma|(Q)=\gamma_1\sum_{p\in Q}|\langle\mathrm{vec}_{p_1},\mathrm{vec}_{p_2}\rangle|$.  The conditions used are:
\begin{itemize}
\item[(AC)] $\sum_t\|\mathrm{op}_t\|<\infty$, $\sum_t\|\mathrm{vec}_t\|<\infty$, $\gamma_1>0$;
\item[(ID)] $\mathrm{isId}\,t\Rightarrow G(t)=0$; \quad (IZ) $\mathrm{part}(\mathrm{id})=0$;
\item[(NE$_S$)] $\|\mathrm{part}(S)\|^2=\Re\sum_{i,j}\mathbf G_{ij}\langle v(S,i),v(S,j)\rangle$;
\item[(GG$_P$)] $P\gamma_1\mathbf 1-\mathbf G$ is positive semidefinite;
\item[(RO)] $\iota(\mathrm{row}\,t)=\iota(\mathrm{row}\,t')$ and $\mathrm{row}\,t\ne\mathrm{row}\,t'$ imply
 $\langle\mathrm{vec}_t,\mathrm{vec}_{t'}\rangle=0$;
\item[(CP)] $\mathrm{row}\,t=\mathrm{row}\,t'$ and $G(t)\ne0\ne G(t')$ imply $s_0(t,t')\ge1$;
\item[(RB$_{\varepsilon,d}$)] $|\Sigma|(\mathrm{big}(\mathrm{nz}))\le\varepsilon d$; \quad
(RS$_{\varepsilon,d}$) $|\Sigma|(\mathrm{nsc}(\mathrm{gen}))\le\varepsilon d$;
\item[(RM$_{A,d}$)] $\Sigma(\mathrm{sc}(\mathrm{gen}))\le Ad$; \quad
(RN$_{\varepsilon',d}$) $\Sigma(\mathrm{one}(\mathrm{ng}))\le\varepsilon'd$.
\end{itemize}
(RM) is the \emph{main row}, (RB) and (RS) the \emph{loops}, (RN) the \emph{non-generic row} and (GG) the
\emph{Gram bound}.  This is the notion of Definition~\ref{i3:def-terms}, with $\Gamma^\sharp,\Gamma_1$ written
$\mathbf G,\gamma_1$, together with the total $K$ and the further conditions (IZ), (RB)--(RN).
\end{definition}

\begin{lemma}[Four parts]\label{i3:S-lem-fourparts}
Under (ID) the sets $\mathrm{id},\mathrm{deg},\mathrm{ng},\mathrm{gen}$ partition $\tau$; under (ID), (AC) and completeness of $H$,
$K=K_0+\mathrm{part}(\mathrm{deg})$.
\end{lemma}

\begin{proof}
An identity term has $G=0$ and lies only in $\mathrm{id}$; any other term satisfies exactly one of $G=0$,
$G\ne0\wedge\mathrm{gener}$, $G\ne0\wedge\neg\mathrm{gener}$.  Each restricted family $\mathbf 1_S\mathrm{op}$ is
absolutely summable, hence summable in the complete space $H$, and sums are additive.
\end{proof}

\begin{lemma}\label{i3:S-lem-normsq}
If $H_1$ is a Hilbert space and $\sum_t\|f(t)\|<\infty$, then $(t,t')\mapsto\Re\langle f(t),f(t')\rangle$ is
summable on $\tau^2$ and $\|\sum_tf(t)\|^2=\sum_{(t,t')}\Re\langle f(t),f(t')\rangle$.
\end{lemma}

\begin{proof}
$|\Re\langle f(t),f(t')\rangle|\le\|f(t)\|\|f(t')\|$ is summable on $\tau^2$.  The continuous additive map
$x\mapsto\Re\langle w,x\rangle$ commutes with unconditional sums; applying this twice to
$\Re\langle\sum f,\sum f\rangle$ gives an iterated sum, equal to the double sum by Fubini for summable families.
\end{proof}

\begin{lemma}[Pair split]\label{i3:S-lem-pairsplit}
Assume (AC), (RO), (CP) and $H_1$ complete.  Then $\Phi(S)=\Sigma(\mathrm{all}(S))$ for every $S$, and
\[
\Phi(\mathrm{gen})=\Sigma(\mathrm{big}(\mathrm{gen}))+\Sigma(\mathrm{nsc}(\mathrm{gen}))+\Sigma(\mathrm{sc}(\mathrm{gen})),\qquad
\Phi(\mathrm{ng})=\Sigma(\mathrm{big}(\mathrm{ng}))+\Sigma(\mathrm{one}(\mathrm{ng})).
\]
\end{lemma}

\begin{proof}
By Lemma~\ref{i3:S-lem-normsq} applied to $t\mapsto\mathbf 1[t\in S,\iota\,\mathrm{row}\,t=i]\mathrm{vec}_t$ and summation over the
finite set $I$, $\Phi(S)/\gamma_1$ is the sum of $\Re\langle\mathrm{vec}_{p_1},\mathrm{vec}_{p_2}\rangle$ over the
pairs in $S^2$ with $\iota\,\mathrm{row}\,p_1=\iota\,\mathrm{row}\,p_2$.  By (RO) the pairs with different rows contribute
$0$, which gives $\Sigma(\mathrm{all}(S))$.  For $S=\mathrm{gen},\mathrm{ng}$, (CP) gives $s_0\ge1$ on $\mathrm{all}(S)$, so
$\mathrm{all}(S)$ is the disjoint union of $\{s_0\ge2\}$ and $\{s_0=1\}$, the latter split by $\mathrm{scal}$.
\end{proof}

\begin{lemma}\label{i3:S-lem-gram}
Let $I$ be finite, $a:I\to H_1$, $\mathbf G\in\C^{I\times I}$, $c\in\R$ with $c\mathbf 1-\mathbf G$ positive
semidefinite.  Then $\Re\sum_{i,j}\mathbf G_{ij}\langle a_i,a_j\rangle\le c\sum_i\|a_i\|^2$.  Consequently, under
(AC), (NE$_S$) and (GG$_P$), $\|\mathrm{part}(S)\|^2\le P\,\Phi(S)$.
\end{lemma}

\begin{proof}
Write $c\mathbf 1-\mathbf G=\sum_lv_lv_l^*$; then
$c\sum_i\|a_i\|^2-\sum_{i,j}\mathbf G_{ij}\langle a_i,a_j\rangle=\sum_l\|\sum_i\overline{v_l(i)}a_i\|^2\ge0$.  For the
consequence take $c=P\gamma_1$ and $a_i=v(S,i)$.
\end{proof}

\begin{theorem}[Abstract pair bound]\label{i3:S-thm-abstract}
Let $X$ be a term system with $H_1$ complete, and $P,d\ge0$.  Assume (AC), (NE$_{\mathrm{gen}}$),
(NE$_{\mathrm{ng}}$), (GG$_P$), (RO), (CP), (IZ), (RB$_{\varepsilon,d}$), (RS$_{\varepsilon,d}$), (RM$_{A,d}$),
(RN$_{\varepsilon',d}$).  Then
\[
\|K_0\|^2\le\Bigl(\sqrt{P(A+2\varepsilon)}+\sqrt{P(\varepsilon+\varepsilon')}\Bigr)^2d.
\]
\end{theorem}

\begin{proof}
For $Q\subset Q'$, $|\Sigma|(Q)\le|\Sigma|(Q')$ and $|\Sigma(Q)|\le|\Sigma|(Q)$.  Since $\mathrm{big}(\mathrm{gen})$ and
$\mathrm{big}(\mathrm{ng})$ lie in $\mathrm{big}(\mathrm{nz})$, Lemma~\ref{i3:S-lem-pairsplit} gives
$\Phi(\mathrm{ng})\le(\varepsilon+\varepsilon')d$ and $\Phi(\mathrm{gen})\le(A+2\varepsilon)d$, and
Lemma~\ref{i3:S-lem-gram} gives $\|\mathrm{part}(\mathrm{gen})\|^2\le P(A+2\varepsilon)d$,
$\|\mathrm{part}(\mathrm{ng})\|^2\le P(\varepsilon+\varepsilon')d$.  By (IZ), $K_0=\mathrm{part}(\mathrm{gen})+\mathrm{part}(\mathrm{ng})$, and the
triangle inequality concludes.
\end{proof}

\begin{definition}[Size numbers]\label{i3:S-def-SN}
For $k,r_S,f\in\N$ and $b,g_1,g_2,c,C_n\in\R$, $\mathrm{SN}(k,r_S,f,b,g_1,g_2,c,C_n)$ is the conjunction of:
$b\ge16$, $k\ge32$, $f\le k$, $4r_S+12\le k$, $(kb)^8\le g_j$ ($j=1,2$), $70\log b\le k$, $1\le5cf\log k$ and
$8C_n(2r_S+2)^5f\log b\le k$.
\end{definition}

\begin{lemma}[Numerics]\label{i3:S-lem-num}
\begin{enumerate}
\item If $b\ge1\le k$, $f\le k$, $(kb)^8\le g_j$, then $0\le P_g:=(1+3/g_1)^f(1+3/g_2)^f\le e^{6b^{-7}}$.
\item If $70\log b\le k$ then $e^{-k/5}\le b^{-14}$.
\item If $f\le k$, $k\ge32$, $4r_S+12\le k$, then $12f2^{-(k-2r_S-2)}\le f2^{-k/2}\le e^{-k/5}$.
\item If $1\le5cf\log k$, then $e^{-ckf\log k}\le e^{-k/5}$; if $8Cm^5f\log b\le k$, then
 $2^{-k/2}e^{Cm^5f\log b}\le e^{-k/5}$; if $4r_S+10\le k$, then $2^{-(k-2r_S-5)}\le2^{-k/2}$.
\item If $b\ge16$, $0\le P\le e^{6b^{-7}}$, $A\le(1+b^{-7})e^{b^{-14}}$ and $\varepsilon,\varepsilon'\le b^{-14}$, then
 $(\sqrt{P(A+2\varepsilon)}+\sqrt{P(\varepsilon+\varepsilon')})^2\le1+b^{-6}$.
\end{enumerate}
\end{lemma}

\begin{proof}
(1) $fb^7\le(kb)^8\le g_j$ and $(1+3/g)^f\le e^{3f/g}$.  (2) is immediate.  (3) With $j=k-2r_S-2$ one has
$2^{-k/2}\ge16\cdot2^{-j}$; and $\log k-\frac k2\log2\le-\frac k5$ for $k\ge32$ (from $\log x\le x-1$ at $x=k/32$).
(4) Direct.  (5) Put $u=b^{-7}\le10^{-3}$; then $e^{6u}\le1+7u$, $e^{u^2}\le1+u$, so $P(A+2\varepsilon)\le1+11u$ and
$P(\varepsilon+\varepsilon')\le3u^2$; the sum of square roots is $\le1+\frac{29}4u$ and $(1+\frac{29}4u)^2\le1+16u\le1+b^{-6}$.
\end{proof}

\begin{theorem}[Pair bound]\label{i3:S-thm-pairbound}
Let $X$ be a term system with $H_1$ complete.  Assume $\mathrm{SN}(k,r_S,f,b,g_1,g_2,c,C_n)$, $d\ge0$,
$0\le\delta\le b^{-7}$, $0\le e_1\le\exp(12f2^{-(k-2r_S-2)})$, (AC), (NE$_{\mathrm{gen}}$), (NE$_{\mathrm{ng}}$),
(GG$_{P_g}$), (RO), (CP), (IZ), (RB$_{\varepsilon,d}$) and (RS$_{\varepsilon,d}$) with $\varepsilon=e^{-ckf\log k}$,
(RM$_{(1+\delta)e_1,d}$), and (RN$_{\varepsilon',d}$) with $\varepsilon'=2^{-(k-2r_S-5)}e^{C_n(2r_S+2)^5f\log b}$.  Then
$\|K_0\|^2\le(1+b^{-6})d$.
\end{theorem}

\begin{proof}
Lemma~\ref{i3:S-lem-num} gives $P_g\le e^{6b^{-7}}$, $\varepsilon,\varepsilon'\le e^{-k/5}\le b^{-14}$ and
$(1+\delta)e_1\le(1+b^{-7})e^{b^{-14}}$; apply Theorem~\ref{i3:S-thm-abstract} and Lemma~\ref{i3:S-lem-num}(5).
\end{proof}

\medskip\noindent\label{i3:S-sss-main}\emph{The main row.}

The main row is controlled by a \emph{mass datum}, a finite list of positive numbers (local and global
volumes, formal degrees, densities) subject to identities.  The bound comes in two steps: an analytic
inequality for a quadratic form over lattice points (the recursion), and an exact algebraic evaluation of
the resulting quantity as $d_ud_{u'}/D_*$ (the mass formula).

\begin{definition}[Mass datum]\label{i3:S-def-mass}
A \emph{mass datum} $B$ over a finite set $P$ of auxiliary places consists of positive reals
$g_p,\mathrm{glk}_p,\mathrm{uk}_p,\mathrm{nn}_p,\mathrm{gg}_p$, local volumes
$\mu K'_p,\mu K^{\rm nat}_p,\mu K''_p,\mu S_p,\mu S'_p,\mu S''_p,\mu S'''_p,M_p,\lambda_p$ ($p\in P$) and their
counterparts at the remaining places (index $\infty$), global volumes
$\mathrm{vol}_G,\mathrm{vol}_{H_a},\mathrm{vol}_{H'_a},\mathrm{vol}_B,\mathrm{vol}_H,\mathrm{vol}_{H'},\mathrm{vol}_N,\mathrm{vol}_{N'},\mathrm{vol}_{ZZ},\mathrm{vol}_{PP}$,
positive constants $C_{oo},\deg,\dim\pi,\dim\pi',\rho_0,\rho'_0,d,d',n_0,\mu_H,\mu''_H,\mu_N$ (and primed
analogues), and real numbers $s_d,e_1,\delta,K_2,\mathrm{KV},\Sigma,K_{\rm top},J_N,J_C,J_{OO}$.  Put
\[
\gamma_1=\prod_p\frac{\mathrm{glk}_p}{\mathrm{uk}_p},\quad t=\prod_p\Bigl(1+\frac3{g_p}\Bigr),\quad
c_0=\frac{\mathrm{vol}_N\mathrm{vol}_H}{\dim\pi},\quad c'_0=\frac{\mathrm{vol}_{N'}\mathrm{vol}_{H'}}{\dim\pi'},
\]
$\mu K=\mu K_\infty\prod_p\mu K'_p$, $\mu K''=\mu K_\infty\prod_p\mu K''_p$, $MM=M_\infty\prod_pM_p$,
$\lambda=\lambda_\infty\prod_p\lambda_p$, $\mu S''=\mu S_\infty\prod_p\mu S''_p$,
$\mu S'''=\mu S'_\infty\prod_p\mu S'''_p$,
$D_*(B)=\deg\cdot\mathrm{vol}_G/\mu K''$, $AA=\mathrm{vol}_BC_{oo}(MM/\mathrm{vol}_{ZZ})/\mu K$ and $GS=t(d/\dim\pi)\mathrm{KV}$.
$B$ satisfies the \emph{mass hypotheses} if all entries listed as positive are positive, $s_d,\delta\ge0$, and:
\begin{itemize}
\item[(F)] $\Sigma=\rho_0\mathrm{KV}$; \ (Ma) $n_0\rho_0=AA$; \
(Ja) $J_N=\mathrm{vol}_{ZZ}J_C/\lambda$; \ (Nu) $J_{OO}=C_{oo}J_C$; \ (In) $J_{OO}\le1/\deg$;
\item[(Vo)] $\mathrm{vol}_B\mathrm{vol}_G=\mathrm{vol}_{H_a}\mathrm{vol}_{H'_a}$ and $\mathrm{vol}_N\mu_N=1$ (and primed);
\item[(Ms), (Co), (LS)] $\rho_0=\mathrm{vol}_{H_a}/(\mathrm{vol}_H\mu_H)$, $d=\dim\pi\,\mathrm{vol}_{H_a}/(\mathrm{vol}_H\mu''_H)$,
$\mu_H/\mu''_H=\gamma_1$, $\mu S''=\mu_N\mu''_H$ (and primed);
\item[(Or), (Ty), (Ix)] $M_\infty/\lambda_\infty=\mu K_\infty^2/(\mu S_\infty\mu S'_\infty)$,
$M_p/\lambda_p=\mu K'_p\mu K^{\rm nat}_p/(\mu S_p\mu S'_p)$, $\mu S_p=\frac{\mathrm{glk}_p}{\mathrm{uk}_p}\mu S''_p$,
$\mu S'_p=\frac{\mathrm{glk}_p}{\mathrm{uk}_p}\mu S'''_p$,
$\mu K'_p=\frac{\mathrm{gg}_p}{\mathrm{nn}_p\mathrm{uk}_p}\mu K''_p$, $\mu K^{\rm nat}_p/\mu K'_p=\mathrm{glk}_p\mathrm{nn}_p/\mathrm{gg}_p$
(for every $p\in P$).  The only archimedean identity is (Or).
\end{itemize}
The further relations used are (Ke) $\mathrm{KV}\le s_dn_0K_{\rm top}$, $K_{\rm top}\ge0$; (De) $s_d\le e_1$;
(TA) $c_0c'_0K_{\rm top}\le(1+\delta)J_N$; (2A) $K_2\le t\gamma_1\Sigma$.
\end{definition}

\begin{lemma}\label{i3:S-lem-sig}
If $B$ is positive and satisfies (F), (Ke), (Ma), (De), (TA), (Ja), (Nu), (In), then
$\Sigma\le(1+\delta)e_1(c_0c'_0)^{-1}\mathrm{vol}_B(MM/\lambda)/(\mu K\deg)$.
\end{lemma}

\begin{proof}
(Ja), (Nu), (In) give $J_N\le\mathrm{vol}_{ZZ}\lambda^{-1}C_{oo}^{-1}\deg^{-1}$, and (TA) gives
$K_{\rm top}\le(1+\delta)J_N/(c_0c'_0)$.  By (F), (Ke), (De), (Ma),
$\Sigma=\rho_0\mathrm{KV}\le e_1\rho_0n_0K_{\rm top}=e_1AA\,K_{\rm top}$.  Substitute.
\end{proof}

\begin{lemma}\label{i3:S-lem-columns}
Under positivity, (Vo), (LS), (Ms), (Co): $d=\gamma_1\rho_0\dim\pi$ and $c_0^{-1}=d\,\mu S''/\mathrm{vol}_{H_a}$ (and
primed).  Under positivity, (Or), (Ty), (Ix): $MM/(\lambda\,\mu K)=\mu K''/(\gamma_1\mu S''\mu S''')$.
\end{lemma}

\begin{proof}
Divide the column relation for $d$ by the mass relation for $\rho_0$ and use $\mu_H/\mu''_H=\gamma_1$; for
$c_0^{-1}$ combine $\mathrm{vol}_N\mu_N=1$ and $\mu S''=\mu_N\mu''_H$.  For the second identity multiply the local
ratio identities over all places and substitute (Ix): by (Ty), (Ix),
$\mu K^{\rm nat}_p/\mu S_p=\mu K''_p/\mu S''_p$ and $\prod_p\mu S'_p=\gamma_1\prod_p\mu S'''_p$, so
$MM/(\lambda\,\mu K)=\mu K_\infty\prod_p\mu K''_p\big/\bigl(\mu S_\infty\prod_p\mu S''_p\cdot\mu S'_\infty\prod_p\mu S'_p\bigr)
=\mu K''/(\gamma_1\mu S''\mu S''')$.
\end{proof}

\begin{theorem}[Mass formula]\label{i3:S-thm-massformula}
If $B$ satisfies the mass hypotheses, (Ke), (De), (TA) and (2A), then $GS=t\gamma_1\Sigma$ and
\[
K_2\le t\gamma_1\Sigma\le(1+\delta)e_1t\,\frac{dd'}{D_*(B)}.
\]
The same conclusion holds if (Ke), (De), (TA) are replaced by: $e_1>0$, $\mathrm{KV}\ge0$ and
$c_0c'_0\mathrm{KV}\le(1+\delta)e_1n_0J_N$.
\end{theorem}

\begin{proof}
By Lemma~\ref{i3:S-lem-columns} and (F), $GS=t\gamma_1\rho_0\mathrm{KV}=t\gamma_1\Sigma$.  Multiply
Lemma~\ref{i3:S-lem-sig} by $t\gamma_1$ and substitute Lemma~\ref{i3:S-lem-columns} and (Vo)
($\mathrm{vol}_B=\mathrm{vol}_{H_a}\mathrm{vol}_{H'_a}/\mathrm{vol}_G$): the right side is exactly $(1+\delta)e_1t\,dd'/D_*(B)$.  For the
variant, replace $s_d$ by $e_1$ and $K_{\rm top}$ by $\mathrm{KV}/(e_1n_0)$; the conclusions do not involve these entries.
\end{proof}

The identities (Vo), (Ms), (Co) are the place where Siegel--Weil enters the main row: the global volumes
$\mathrm{vol}_{H_a}$, $\mathrm{vol}_G$ of the stabilisers and of the unitary group, and the masses $\rho_0$, are
evaluated by (SW) (component (C17) below).

\emph{The recursion.}  A \emph{point of level $i$} is a finite list $p=((s_1,t_1),\dots,(s_m,t_m))$ with
$s_j\ge1$, $t_j\in E^\times$ and $\sum s_j=i$; $h\cdot(s,t)$ appends an entry.  A \emph{kernel datum} $D$ (of
rank $r_S$ and degree $f$) consists of: positive semidefinite kernels $K_r(i)$ on points of level $i$
($0\le i\le r_S+1$, $K_r(0)(\emptyset,\emptyset)=1$); exponents $e_i$; lattices
$\Lambda(b,s),\Lambda'(b,s),\Lambda^\flat(i)\subset E$ with volumes $\mathrm{vol}(b,s)$ and size vectors; a finite
set of words $w$ with positive weights and counts, containing the empty word (weight $1$, count $0$); Gram
forms $g_{\epsilon,i,w}(p,q)$ ($\epsilon>0$), positive semidefinite on finite sets; and, for a Hilbert space
$H$, vectors $V(p)\in H$ and coefficients $c_\rho(p)$ supported on level $r_S+1$; and an integer $K\ge4$ (the
\emph{word parameter}: words of count at most $4K(r_S+1-i)$ are used at level $i$, and $K$ is the exponent in the
lattice-sum bound below).  The quadratic
quantities are
\[
Q(T,\epsilon,i,w)=\mathrm{wt}(w)^{-1}\Re\sum_{p,q\in X_i}K_r(i)(p,q)\,g_{\epsilon,i,w}(p,q),
\]
where $X_i$ is the finite set of points of level $i$ of a \emph{tower} $T$ (finite sets of admissible
entries at each step).  The \emph{recursion hypotheses} $\mathrm{RH}(D,C_1,C_4,\Theta,j_n,V,c)$ consist of:
positivity of kernels, Gram forms and weights; $0\le e_i\le1$; $0\le\mathrm{vol}(b,s)$ for all $b,s$, and
$\mathrm{vol}(b,s)\le\Theta^{-2f}$ for $b,s\ge1$, $b+s\le r_S+2$; existence for every tower of a \emph{minimal theta domination} (kernels
$\Theta_{i,s,t}$, depending on $t$ only modulo $\Lambda(i-s+1,s)$, dominating the restrictions of $K_r(i)$, with
$(1+e_i)K_r(i-s)-\sum_t\Theta_{i,s,t}\succeq0$ over incongruent $t$); five analytic estimates (support,
wave, Poisson, lift, block with constant $C_4$) and the lattice-sum bound with constant
$e_p=(C_1\Theta^{-(K-5/2)})^f$; the energy estimate
$\lim_{\epsilon\to0^+}\Re g_{\epsilon,0,w}(\emptyset,\emptyset)/\mathrm{wt}(w)=j_n$; the cut-off condition
(for every finite set $Y$ of points of level $r_S+1$, for all sufficiently small $\epsilon$, depending on $Y$,
$g_{\epsilon,r_S+1,\emptyset}(p,q)=\langle V(q),V(p)\rangle$ for $p,q\in Y$ with $K_r(r_S+1)(p,q)\ne0$); and the
\emph{leading domination} in the sense of Proposition~\ref{i3:S-prop-limit}: $c_\rho$ is supported on points of level
$r_S+1$, $K_r(r_S+1)-\sum_{\rho\in R}c_\rho c_\rho^*\succeq0$ on finite sets for every finite $R$, and
$\sum_p|c_\rho(p)|\,\|V(p)\|<\infty$ for every $\rho$.  (The
precise form of the five estimates, of the lattice-sum bound and of the minimal domination is
Definition~\ref{i3:S-def-est} in \S\ref{i3:S-ssec-onelevel}; they enter only through
Proposition~\ref{i3:S-prop-onelevel}.  Positivity of the Gram forms, of the kernels and of the weights are
(GP), (KP), (WP) there.)

\begin{lemma}[Solving the recursion]\label{i3:S-lem-solve}
Let $0\le a\le\frac12$, $0\le\varepsilon\le a$, $r_Sa\le1$, $e_i\ge0$, $\sigma_i\ge0$, and suppose that for
$1\le i\le r_S+1$
\[
\sigma_i\le(1+\varepsilon)\sqrt{1+e_i}\,\sigma_{i-1}+\sum_{t=2}^ia^{t-1}\sigma_{i-t}.
\]
Then $\sigma_{r_S+1}^2\le\bigl(1+e^9(r_S+1)a\bigr)\prod_{i=1}^{r_S+1}(1+e_i)\,\sigma_0^2$.
\end{lemma}

\begin{proof}
Let $\lambda=1+\varepsilon+2a$ and $Q_k=\prod_{j\le k}\sqrt{1+e_j}$, nondecreasing and $\ge1$.  By strong induction
$\sigma_i\le\lambda^iQ_i\sigma_0$: with $M=\lambda^kQ_{k+1}\sigma_0$, the first term of the bound for $\sigma_{k+1}$ is at
most $(1+\varepsilon)M$, and the sum is at most $M\sum_{t\ge2}a^{t-1}\le2aM$ since $a\le\frac12$.  So
$\sigma_{r_S+1}^2\le\lambda^{2(r_S+1)}\prod(1+e_i)\sigma_0^2$.  With $x=(r_S+1)a\le\frac32$,
$\lambda^{2(r_S+1)}\le e^{6x}\le1+e^9x$ by convexity of $\exp$ on $[0,9]$.
\end{proof}

\begin{proposition}[One-level inequality]\label{i3:S-prop-onelevel}
Let $D$ be a kernel datum, $e_p\ge0$ and $C_4\ge1$.  Assume (GP), (KP), (WP), (Su), (Wa), (Po), (Li),
(Bl$_{C_4}$), (LS$_{e_p}$) of Definition~\ref{i3:S-def-est}, $0\le e_i\le1$ for $1\le i\le r_S+1$,
$\mathrm{vol}(b,s)\ge0$ for all $b,s$, and $\mathrm{vol}(b,s)\le1$ for $b,s\ge1$, $b+s\le r_S+2$.
Let $T$ be a tower having a minimal domination, $\epsilon>0$, $1\le i\le r_S+1$, and $\sigma:\N\to\R$ with, for
every $k<i$, $\sigma_k\ge0$ and $Q(T,\epsilon,k,w)\le\sigma_k^2$ for all words of count at most $4K(r_S+1-k)$.  Then
for every word $w$ of count at most $4K(r_S+1-i)$
\[
\sqrt{Q(T,\epsilon,i,w)}\le\sqrt{(1+3^fe_p)(1+2^fe_p)}\sqrt{1+e_i}\,\sigma_{i-1}
+\sum_{s=2}^i\Bigl((1+3^fe_p)(1+e_i)C_4^{fs}\prod_{s'=1}^{s-1}\mathrm{vol}(i-s+1,s')\Bigr)^{1/2}\sigma_{i-s}.
\]
\end{proposition}

\noindent The hypotheses are contained in $\mathrm{RH}$ (the volume bound
$\mathrm{vol}(b,s)\le\Theta^{-2f}\le1$ and $\mathrm{vol}\ge0$ are the clause ``sizes below $\Theta$'').  The proof is in
\S\ref{i3:S-ssec-onelevel}.

\begin{proposition}[Limit step]\label{i3:S-prop-limit}
Let $D$ be a kernel datum, $H$ a Hilbert space, $V$ and $c_\rho$ ($\rho\in\N$) as above, and $j_n,B\in\R$.  Assume
that the kernels $K_r(i)$ and the Gram forms are positive semidefinite on finite sets; the energy estimate
($\Re g_{\epsilon,0,w}(\emptyset,\emptyset)/\mathrm{wt}(w)\to j_n$ as $\epsilon\to0^+$, for every word $w$ of count at most
$N_{\rm top}=4K(r_S+1)$); the cut-off condition (for every finite set $Y$ of points of level $r_S+1$ there is
$\epsilon_Y>0$ such that for $0<\epsilon<\epsilon_Y$ and $p,q\in Y$ with $K_r(r_S+1)(p,q)\ne0$ one has
$g_{\epsilon,r_S+1,\emptyset}(p,q)=\langle V(q),V(p)\rangle$); and the leading domination ($c_\rho(p)\ne0$ only if $p$ is a
point of level $r_S+1$; $K_r(r_S+1)-\sum_{\rho\in R}c_\rho c_\rho^*$ is positive semidefinite on every finite set of
points of level $r_S+1$, for every finite $R\subset\N$; and $\sum_p|c_\rho(p)|\,\|V(p)\|<\infty$ for every $\rho$).  If for
every tower $T$, every $\epsilon>0$ and every $s_0\in\R$ with $\Re g_{\epsilon,0,w}(\emptyset,\emptyset)/\mathrm{wt}(w)\le s_0$ for
all words of count at most $N_{\rm top}$ one has $Q(T,\epsilon,r_S+1,\emptyset)\le Bs_0$, then
$\sum_{\rho\in R}\|\sum_pc_\rho(p)V(p)\|^2\le Bj_n$ for every finite $R\subset\N$.
\end{proposition}

\begin{proof}
\emph{A Gram inequality.}  Let $Y$ be finite and $k$ a kernel on $Y$ with $k-\sum_{\rho\in R}c_\rho c_\rho^*$ positive
semidefinite on $Y$.  For any positive semidefinite $m$ on $Y$, $\Re\sum_{p,q\in Y}m(p,q)\langle V(q),V(p)\rangle\ge0$: the
Gram matrix $(\langle V(p),V(q)\rangle)$ is positive semidefinite, so it equals $B_0^*B_0$ for a matrix $B_0$, and the double sum
becomes $\sum_jv_j^*mv_j\ge0$ with $v_j=(\overline{(B_0)_{jp}})_p$.  Applying this to $m=k-\sum_\rho c_\rho c_\rho^*$ and using
$\sum_{p,q}c_\rho(p)\overline{c_\rho(q)}\langle V(q),V(p)\rangle=\|\sum_pc_\rho(p)V(p)\|^2$ gives
\begin{equation}\label{i3:eq-rowsleform}
\sum_{\rho\in R}\Bigl\|\sum_{p\in Y}c_\rho(p)V(p)\Bigr\|^2\le\Re\sum_{p,q\in Y}k(p,q)\langle V(q),V(p)\rangle .
\end{equation}

\emph{Reduction to finite sets.}  For each $\rho$ the family $c_\rho(p)V(p)$ is absolutely summable, hence summable in
$H$, so along the directed set of finite sets $F$ of points, $\Phi(F):=\sum_{\rho\in R}\|\sum_{p\in F}c_\rho(p)V(p)\|^2$
tends to the left side of the claim.  It suffices to find, for every finite $F$, a finite $F'\supset F$ with
$\Phi(F')\le Bj_n$.

\emph{Choice of a tower.}  Let $F_0$ be the set of points of level $r_S+1$ in $F$ and $S\subset E^\times$ the finite set
of all labels $t_j$ occurring in the points of $F_0$.  Let $T$ be the tower whose admissible entries are $S$ at every
step; its level-$i$ set $X_i$ (points of level $i$ all of whose labels lie in $S$) is finite, since its points have
at most $i$ entries, taken from $\{1,\dots,i\}\times S$.  Put $Y=X_{r_S+1}$; it consists of points of level $r_S+1$
and contains $F_0$.  Let $F'=F\cup Y$.  Since $c_\rho$ vanishes off the points of level $r_S+1$, and those of $F'$ lie in
$Y$, $\Phi(F')=\sum_\rho\|\sum_{p\in Y}c_\rho(p)V(p)\|^2$, which by the leading domination and \eqref{i3:eq-rowsleform}
is at most $A:=\Re\sum_{p,q\in Y}K_r(r_S+1)(p,q)\langle V(q),V(p)\rangle$.

\emph{Passage to the limit.}  Let $\delta>0$.  There are finitely many words of count at most $N_{\rm top}$, so the
energy estimate gives, for all small $\epsilon>0$, $\Re g_{\epsilon,0,w}(\emptyset,\emptyset)/\mathrm{wt}(w)\le j_n+\delta$ for all of
them; the cut-off condition for $Y$ gives, for all small $\epsilon>0$,
$g_{\epsilon,r_S+1,\emptyset}(p,q)=\langle V(q),V(p)\rangle$ whenever $p,q\in Y$ and $K_r(r_S+1)(p,q)\ne0$.  Fix such an
$\epsilon$.  Terms with $K_r(r_S+1)(p,q)=0$ vanish on both sides and the empty word has weight $1$, so
$Q(T,\epsilon,r_S+1,\emptyset)=A$, and the hypothesis with $s_0=j_n+\delta$ gives $A\le B(j_n+\delta)$.  As $\delta\to0^+$,
$A\le Bj_n$; hence $\Phi(F')\le Bj_n$.
\end{proof}

\begin{theorem}[Lattice-average recursion]\label{i3:S-thm-recursion}
For $C_1,C_4\ge1$ there are $C,C_\Theta$ such that whenever $f\ge1$, the word parameter satisfies $K\ge4$,
$\Theta\ge C_\Theta$, $r_S(C/\Theta)^f\le1$ and $\mathrm{RH}(D,C_1,C_4,\Theta,j_n,V,c)$ hold, for every finite $R$
\[
\sum_{\rho\in R}\Bigl\|\sum_pc_\rho(p)V(p)\Bigr\|^2\le\Bigl(1+C(r_S+1)\bigl(\tfrac C\Theta\bigr)^f\Bigr)\prod_{i=1}^{r_S+1}(1+e_i)\cdot j_n.
\]
\end{theorem}

\begin{proof}
Put $C_0=\max(3C_1,2C_4)$, $C=\max(C_0,e^9)$, $C_\Theta=\max(2C_0,3C_1,1)$.  Let $\Theta\ge C_\Theta$ (so
$\Theta\ge1$) and $e_p=(C_1\Theta^{-(K-5/2)})^f$; as $K-\frac52\ge\frac32\ge1$ and $\Theta\ge1$,
$\Theta^{-(K-5/2)}\le\Theta^{-1}$; then $(1+3^fe_p)(1+2^fe_p)\le(1+(C_0\Theta^{-(K-5/2)})^f)^2$ and, for $s\ge2$, since
each $\mathrm{vol}(i-s+1,s')\le\Theta^{-2f}$ and $(1+3^fe_p)(1+e_i)\le4\le C_0^{2f(s-1)}C_4^{-fs}$, the $s$-th coefficient
of Proposition~\ref{i3:S-prop-onelevel} is at most $((C_0/\Theta)^f)^{s-1}$.  By Proposition~\ref{i3:S-prop-limit} it suffices
to bound $Q(T,\epsilon,r_S+1,\emptyset)$ by $(1+C(r_S+1)(C/\Theta)^f)\prod(1+e_i)s_0$ for each tower, $\epsilon$ and
admissible bound $s_0$ on the energies.  Note $s_0\ge0$: $X_0=\{\emptyset\}$ and $K_r(0)(\emptyset,\emptyset)=1$, so
$Q(T,\epsilon,0,w)=\Re g_{\epsilon,0,w}(\emptyset,\emptyset)/\mathrm{wt}(w)\ge0$ by (GP) and (WP), whence $0\le Q(T,\epsilon,0,w)\le s_0$.
Let $\sigma_k^2$ be the maximum of $\max(Q(T,\epsilon,k,w),0)$ over words of count
$\le4K(r_S+1-k)$; then $\sigma_0^2\le s_0$, and Proposition~\ref{i3:S-prop-onelevel} gives
the recursion of Lemma~\ref{i3:S-lem-solve} with $a=(C_0/\Theta)^f\le\frac12$ and
$\varepsilon=(C_0\Theta^{-(K-5/2)})^f\le(C_0/\Theta)^f=a$ (here $K\ge4$ is used).  Lemma~\ref{i3:S-lem-solve} gives
$Q\le\sigma_{r_S+1}^2\le(1+e^9(r_S+1)a)\prod(1+e_i)\sigma_0^2$, which is the required bound, since the empty word has
weight $1$ and $e^9\le C$, $a\le(C/\Theta)^f$.
\end{proof}

\begin{proposition}[formally verified, \cite{CharacterFamily}]\label{i3:S-prop-instance}
There is $C_1\ge1$ with the following property.  For every $C_A\ge1$, every $C_\Theta>0$ and all admissible
size constants there is an integer $k_0$ such that: if $E$ is an unramified CM field and $D$ the kernel datum
attached to a global size datum of the datum $\mathcal D$, with $k\ge k_0$ and the size hypotheses of the
standing hypotheses (with any closure constant $C'$ satisfying $8C_\Gamma\le C'$ and $C'm^5fL\le k$), then the
threshold $\Theta_h$ of $D$ (defined with $C_A$) satisfies $\Theta_h>4$, $\Theta_h\ge C_\Theta$, and, for every
$0\le C\le C_A$, $r_S(C/\Theta_h)^f\le1$ and $C(r_S+1)(C/\Theta_h)^f\le b^{-7}$; the sizes of $D$ are below $\Theta_h$; $D$ satisfies the lattice-sum bound
with $e_p=(C_1\Theta_h^{-(K-5/2)})^f$ and the Poisson step; and, if $D$ satisfies the local model (C1), its
kernels are positive semidefinite, it satisfies the support condition, $e_i=0$ for $1\le i\le r_S+1$, and every
tower has a minimal domination.
\end{proposition}

\noindent\emph{What the proof involves.}  Box counts in the ideals
$g^{-1}\mathfrak D^{-\nu}$ of $E$ and Poisson summation on $E\otimes\R/\Lambda$ give the lattice sums; the lower
bound $k^{c_*}\le\Theta_h$ comes from eight estimates on the local Bessel-type quantities (interlacing,
anti-concentration, threshold products); the local-model consequences come from five statements on the
explicit local kernels.  The threshold lower bound and the five local kernel statements are a long verification.  Its formally verified steps are outlined in \S\ref{i3:str-inst}.

\begin{proposition}[Densities]\label{i3:S-prop-dens}
Suppose $2r_S+f+7\le k$, $P$ is a set with $q:P\to\R$ and $\ell:P\to\N$ such that $\ell(p)$ is prime,
$\ell(p)\le q(p)$, at most $f$ elements of any finite subset of $P$ share a value of $\ell$, and every $x<s_d$ is
bounded by $\prod_{p\in S}\exp(4q(p)^{-(k-2r_S-2)})$ for some finite $S\subset P$.  Then
$s_d\le\exp(12f2^{-(k-2r_S-2)})$ and $s_d\le2$.
\end{proposition}

\begin{proof}
For $n\ge2$ and finite $S$, $\sum_Sq(p)^{-n}\le f\sum_{m\ge2}m^{-n}\le4f2^{-n}\sum_{m\ge2}m^{-2}\le3f2^{-n}$, using
$\sum_{m=2}^{M}m^{-2}\le\frac34$.  With $n=k-2r_S-2\ge f+5$ this bounds every product by
$\exp(12f2^{-n})$, hence $s_d\le\exp(12f2^{-n})$; and $18f\le2^{f+5}\le2^n$ gives $12f2^{-n}\le\frac23<\log2$.
\end{proof}

\begin{theorem}[Main row]\label{i3:S-thm-mainrow}
For $C_4\ge1$ and $C_{\rm thr}$ there are $C_A\ge\max(C_{\rm thr},1)$ and $0\le C\le C_A$ such that, for all
admissible size constants, there is $k_0$ with the following property.  Let $D$ be a kernel datum as in
Proposition~\ref{i3:S-prop-instance} with $k\ge k_0$, satisfying the local model, the real inputs with
constants $(C_4,j_n)$, the density input with density $s_d$ (Proposition~\ref{i3:S-prop-dens}), and the rows input
for coefficients $c$ and a mass datum $B$ with the mass hypotheses: $s_dK_r(r_S+1)-\sum_\rho c_\rho c_\rho^*\succeq0$,
$\mathrm{KV}\ge0$, and $c_0c'_0\mathrm{KV}\le n_0M$ whenever all $\sum_{\rho\in R}\|\sum_pc_\rho(p)V(p)\|^2\le M$.  Assume
$J_N=j_n$, $e_1=\max(s_d,1)$, $\delta=C(r_S+1)(C/\Theta_h)^f$.  Then
\[
\gamma_1\Sigma\le(1+\delta)e_1\frac{dd'}{D_*(B)},\qquad0\le\delta\le b^{-7},\qquad1\le e_1\le\exp(12f2^{-(k-2r_S-2)}).
\]
\end{theorem}

\begin{proof}
Proposition~\ref{i3:S-prop-instance} gives $C_1$, the sizes, lattice sums, Poisson step, positivity, support,
$e_i=0$ and minimal dominations, together with the numerical bounds on $\Theta_h$ and $\delta\le b^{-7}$, for $k_0$
large, with $C_\Theta$ and $C_A\ge C$ chosen from the constants $(C,C_\Theta)$ of Theorem~\ref{i3:S-thm-recursion}
for $(C_1,C_4)$; the kernel datum of the instance has word parameter $K\ge4$ and $f\ge1$
(\cite{CharacterFamily}), as Theorem~\ref{i3:S-thm-recursion} requires.  By
Proposition~\ref{i3:S-prop-dens}, $s_d\le2$ and $e_1\le\exp(12f2^{-(k-2r_S-2)})$.  Let $D^\sharp$ be $D$ with the top
kernel multiplied by $\lambda=\max(s_d,1)\in[1,2]$ and $e_{r_S+1}$ replaced by $\lambda-1$: all conditions of $\mathrm{RH}$
are preserved (nonnegative multiples of positive kernels are positive; the dominations at level $r_S+1$ are
multiplied by $\lambda$; the new leading domination is
$(s_dK_r-\sum c_\rho c_\rho^*)+(\lambda-s_d)K_r\succeq0$).  Theorem~\ref{i3:S-thm-recursion} for $D^\sharp$ gives
$\sum_\rho\|\sum_pc_\rho(p)V(p)\|^2\le(1+\delta)e_1j_n$, so the rows input gives
$c_0c'_0\mathrm{KV}\le(1+\delta)e_1n_0J_N$.  The variant of Theorem~\ref{i3:S-thm-massformula} with $K_2=t\gamma_1\Sigma$
gives $t\gamma_1\Sigma\le(1+\delta)e_1t\,dd'/D_*(B)$; divide by $t>0$.
\end{proof}

\begin{theorem}[Pair bound from the local data]\label{i3:S-thm-pbapplied}
Under the hypotheses of Theorem~\ref{i3:S-thm-mainrow}, let $X$ be a term system with $H_1$ complete such that
$\gamma_1(X)=\gamma_1(B)$, $\Sigma(\mathrm{sc}(\mathrm{gen}))=\gamma_1\Sigma_B$, the auxiliary places of $B$ are in bijection with
$\{1,2\}\times\{\text{infinite places of }E\}$ with $g_p=g_1$ or $g_2$ accordingly, $d=d_u$, $d'=d_{u'}$,
$D_*(B)=D_*$; assume (AC), (NE$_{\mathrm{gen}}$), (NE$_{\mathrm{ng}}$), (GG$_{t(B)}$), (RO), (CP), (IZ),
$\mathrm{SN}(k,r_S,f,b,g_1,g_2,c,C_n)$, and (RB), (RS), (RN) as in Theorem~\ref{i3:S-thm-pairbound} with
$d=d_ud_{u'}/D_*$.  Then $\|K_0\|^2\le(1+b^{-6})d_ud_{u'}/D_*$.
\end{theorem}

\begin{proof}
Theorem~\ref{i3:S-thm-mainrow} is (RM$_{(1+\delta)e_1,d}$) with the required bounds on $\delta,e_1$; through the
bijection of places $t(B)=(1+3/g_1)^f(1+3/g_2)^f=P_g$; and $d\ge0$.  Apply Theorem~\ref{i3:S-thm-pairbound}.
\end{proof}

\medskip\noindent\label{i3:S-sss-rows}\emph{The rows supplied by Siegel--Weil.}

\begin{definition}[Selections and rows]\label{i3:S-def-selection}\label{i3:S-def-rows}
A \emph{spectral selection} on a class $\mathcal C$ assigns to each $\mathcal F\in\mathcal C$ a threshold $x_s$ and,
for $x\ge x_s$ and $i\in I_x$, a \emph{spectral package} $(X,Y)$ of $\mathcal D(x,i)$: a Hida setting $X$ (an
automorphic instance of a unitary group $U$ of the hermitian space of $\mathcal D$ over $F$, a Haar measure and
fundamental domain, the right regular representation $\mathcal R$ on $L^2(U(F)\backslash U(\mathbb A))$, a level
subgroup $K_2$ with character $\chi_2$, given by an explicit formula in the setting (that $K_2$ is a group and $\chi_2$ a unitary
character of it is a property proved in Proposition~\ref{i3:prop-chi2}, not a datum), and idempotent $e=\overline{\chi_2}\mathbf 1_{K_2}/\mu_f(K_2)$,
discrete series $\sigma_v$ at the infinite places with formal degrees $d_v$ and unit lowest-$K$-type vectors
$w_v$), and a full package $Y$ over it (column spaces $C_u$, the Poincare-series maps $P_u$, and for each
pair of words the term system $X(u,u')=Y.t(u,u')$ of the kernel expansion of the $(u,u')$ block).  A
member is \emph{locally tied} if at every finite place its local data are isomorphic to the canonical local
model of $\mathcal D$.  The \emph{loop row with constant $C_0$} for the selection is: for every function $a_F$
of the family constants, eventually in the closure constants $(C',C_l)$ and in $x$, for all words $u,u'$
satisfying the counting hypotheses, whenever $(X,Y)$ is concretely tied and $X(u,u')$ satisfies (AC),
$X(u,u')$ satisfies (RB$_{\varepsilon_L,d}$) and (RS$_{\varepsilon_L,d}$) with $\varepsilon_L=e^{-\frac13kf\log k}$ and
$d=d_ud_{u'}/D_*$.  The \emph{rank row with constant $C_0$} is the family ranks hypothesis of
\cite{CharacterFamily}; we use it only through its consequences in Lemma~\ref{i3:S-lem-rankrow}.  Both rows
are closed statements in the sense of \S\ref{i3:ssec-setting} (``eventually in the closure constants'' means:
for $\mathcal K$-large families, for a suitable threshold system $\mathcal K$).
\end{definition}

\begin{proposition}[formally verified, \cite{CharacterFamily}]\label{i3:S-prop-NG}
Assume (SW).  For all admissible constants there are $C_n\ge0$ and thresholds $(C'_0,C_{l,0})$ such that for
$C'\ge C'_0$, $C_l\ge C_{l,0}$, every datum $\mathcal D$ satisfying the counting hypotheses at $(u,u')$, every
concretely tied package $(X,Y)$ with $X(u,u')$ satisfying (AC), the non-generic row
(RN$_{\varepsilon',d}$) holds with $\varepsilon'=2^{-(k-2r_S-5)}e^{C_n(2r_S+2)^5f\log b}$ and $d=d_ud_{u'}/D_*$.
\end{proposition}

\noindent\emph{What the proof involves.}  The pairs of corank one of non-generic terms are grouped into
classes, each class is bounded by a lattice sum over block minors and a mean amplitude, and the masses that
appear (stabiliser volumes and generic masses) are evaluated by (SW); the archimedean amplitudes use
Proposition~\ref{i3:prop-bessel}.  The reduction to leaf statements is long.  Its formally verified steps are outlined in \S\ref{i3:str-NG}.

\begin{lemma}[(FD)]\label{i3:S-lem-FD}
Let $U$ be the unitary group of an automorphic instance over $F$ ($F$ totally real, $E/F$ a CM extension, $U$ the unitary
group of a non-degenerate hermitian space $V$ over $E$), $G_\infty=\prod_{v\mid\infty}U(F_v)$, $\mu_f$ a Haar
measure on $U(\mathbb A_f)$, $\mathcal F$ a fundamental domain for $U(F)$ in $U(\mathbb A)$, and $\mathcal R$ the right
regular representation on $L^2=L^2(U(F)\backslash U(\mathbb A))$.  Let $K_f\subset U(\mathbb A_f)$ be the stabiliser of the
closure in $V\otimes_F\mathbb A_f$ of the self-dual lattice of the instance; it is an open compact subgroup.  Let $e$ be a
locally constant function on $U(\mathbb A_f)$ vanishing outside $K_f$, and $\mathcal R(e)$ the operator of right convolution
by $e$.
For each $v\mid\infty$ let $\sigma_v$ be an irreducible unitary representation of $U(F_v)$ with formal degree $d_v$
(for the fixed Haar measure of $U(F_v)$) and $w_v$ a unit vector whose coefficient
$\phi_v(g)=d_v\overline{\langle w_v,\sigma_v(g)w_v\rangle}$ is integrable, $\phi=\prod_v\phi_v$, and $T=\int_{G_\infty}\phi(g)\mathcal R(g)\,dg$.  Then the image of
$\mathcal R(e)\circ T$ is finite dimensional.
\end{lemma}

\noindent (The order $\mathcal R(e)\circ T$ is that of the formal statement; the two operators commute,
since $e$ lives on $U(\mathbb A_f)$ and $\phi$ on $G_\infty$.)  The proof, from Schur orthogonality, finiteness of the class number and finiteness
of multiplicities in the discrete spectrum, is in \S\ref{i3:S-ssec-leafproofs}.

\begin{lemma}\label{i3:S-lem-rankrow}
Let $C_0\ge1$ and let $\mathcal C$ be a class on which the rank row with constant $C_0$ is closed, i.e.\ the predicate
$\mathrm{Ranks}(\mathcal F;C_0,C_{YL}+1,\tfrac14,3,\tfrac12,C_{Tu})$ is closed on $\mathcal C$.  Put
$a(C_0)=e^{3/4}\cdot C_0e\ \ (\ge1)$.  Then the following three statements are closed on $\mathcal C$:
\begin{itemize}
\item[(i)] (frame row) $\mathrm{TFrame}(\mathcal F;C_{YL}+1,\tfrac14,3,\tfrac12,C_{Tu},C_0)$, with the additional
threshold $40\le C_l^{\mathcal F}$;
\item[(ii)] (size bound) $1\le a(C_0)$, and eventually in the closure constants $(C',C_l)$ and in $x$, every
member $\mathcal D(x,i)$ satisfying the standing hypotheses at $(C_0,C',C_l)$ satisfies $\mathrm{Sizes3b}(a(C_0))$, with
the additional threshold $40\le C_l^{\mathcal F}$;
\item[(iii)] (size row) eventually in $(C',C_l)$ and $x$, every member satisfying the standing hypotheses at
$(C_0,C',C_l)$ satisfies the size rows $\mathrm{SzRows}$ (the size hypotheses of
Proposition~\ref{i3:S-prop-instance}), with the additional threshold $1\le C_l^{\mathcal F}$.
\end{itemize}
Here $C_l^{\mathcal F}$ is the own closure constant of the family.
\end{lemma}

\noindent The proof is in \S\ref{i3:S-ssec-leafproofs};
it reduces the lemma to three bridge statements, which are formally verified.

\medskip\noindent\label{i3:S-sss-components}\emph{The nineteen structural components.}

Fix a datum $\mathcal D$ with host field $E$ ($f$ infinite places, maximal totally real subfield $F$), split
auxiliary primes $g_1\ne g_2$, ranks $(k,r_S)$, and a spectral package $(X,Y)$.  The \emph{concrete tie} of
$(X,Y)$ is the conjunction of: the measure tie and the standard additive character; the level subgroup
$K_2$ is normalised ($K(g)\subset K_2\subset K$, with $K$ the stabiliser in $G_f$ of the lattice, $\chi_2=1$ on $K(g)$,
$0<\mu_f(K_2)<\infty$); $d_u$ equals the
dimension of the column space $C_u$; the expected dimension of the package,
$D_*(X)=\bigl(\prod_vd_v\bigr)\mu(\mathcal F)\Re e(1)$, equals $D_*$; the archimedean coefficient
$g\mapsto\prod_v\langle w_v,\sigma_v(g_v)w_v\rangle$ decays and is integrable; and $\mathcal D$ is locally tied.  For
words $u,u'$ let $B(u,u')$ be the mass datum attached to the pair (built from base points, a pinned measure,
local spaces and uniformisers as in \cite{CharacterFamily}).

\begin{theorem}[The nineteen components]\label{i3:S-thm-components}
Assume (SW), (FD) and the conclusion of Proposition~\ref{i3:S-prop-NG}.  Then the following hold.
\begin{enumerate}
\item[(C1)] \emph{Local model.}  For every compatible pair of words, the kernel datum of the pair satisfies
the local model: (LM1) on every finite set the kernels $K_r(i)$ factor as $a(x)a(y)$ times a finite product of
explicit local kernels with class indicators; (LM2) the one-level data factor into local systems of top, pull
or outside type with $\Lambda(i-s+1,s)=\bigcap_pL_p$; (LM3) $K_r(r_S+1)$ vanishes on pairs whose last single
blocks differ; (LM4) $e_i=0$.
\item[(C2)] \emph{Density.}  For a locally tied datum with $2r_S+f+7\le k$ and words with
$T_u\equiv T_{u'}\pmod{g_S^2}$, the local densities are multipliable, the correction factors
$\alpha_p(u,u')$ are multipliable with $\prod_p\alpha_p>0$, and the density input of
Proposition~\ref{i3:S-prop-dens} holds with $s_d=\sup(\text{densities})/\prod_p\alpha_p$.
\item[(C3)] \emph{Rows input.}  The rows input of Theorem~\ref{i3:S-thm-mainrow} holds: summability of
$|c_\rho(p)|\|V(p)\|$, $s_dK_r(r_S+1)-\sum_\rho c_\rho c_\rho^*\succeq0$, $\mathrm{KV}\ge0$, and $c_0c'_0\mathrm{KV}\le n_0M$
whenever all row sums are $\le M$.
\item[(C4)] \emph{Non-generic row}: the conclusion of Proposition~\ref{i3:S-prop-NG}.
\item[(C5)] \emph{Degenerate terms.}  Eventually in the constants, under the $K_2$-hypotheses and the concrete
tie, $\mathrm{part}(\mathrm{deg})=0$ in $X(u,u')$ for $u\ne u'$.
\item[(C6)] \emph{Carrier.}  There are $C_4\ge1$ and $C_{A,0}$ such that for $C_A\ge\max(C_{A,0},1)$ and admissible
constants there is $k_0$ such that every tied datum with $k\ge k_0$ satisfying the size and counting
hypotheses has a full archimedean carrier: its kernel datum satisfies the real inputs with constants
$(C_4,J_N)$, and its global size datum has parameters $(k,r_S,f)$ of $\mathcal D$.
\item[(C7)] \emph{Kernel expansion.}  Under the concrete tie there is a linear $K^\flat:\mathrm{CM}_{u'}\to\mathrm{CM}_u$ with
$K^\flat=\sum_ta_t$ (unconditionally convergent sum of the term operators $a_t$ of $X(u,u')$) and
$\mathrm{Kform}(T,T')=\langle K^\flat\sharp_{u'}T',\sharp_uT\rangle$ for all $T\in C_u$, $T'\in C_{u'}$ (notation of the paragraph ``Column data'' before Proposition~\ref{i3:S-prop-transport}).
\item[(C8)] \emph{Fold.}  If $K_2$ is normalised and the archimedean coefficient decays, the unfolding of the
Poincare series of every word is legitimate.
\item[(C9)] \emph{Orbital integrals.}  If both folds are legitimate and fundamental domains exist, every
orbital integral of the pair converges absolutely.
\item[(C10)] \emph{Absolute convergence.}  Under the same hypotheses $X(u,u')$ satisfies (AC).
\item[(C11)] \emph{Finite column data.}  Under the gap hypothesis every full package has finite-dimensional
column spaces.
\item[(C12)] \emph{Decay.}  Under the gap hypothesis the archimedean coefficient of every Hida setting decays
and is integrable.
\item[(C13)] \emph{Choices.}  For standing $u\ne u'$ with $T_u\equiv T_{u'}\pmod{g_S^2}$ and $2r_S+f+7\le k$ there
are base points with nonsingular Gram matrix, a pinned measure, local spaces, irreducible uniformisers and a
row labelling $\kappa$ such that two eligible generic terms have the same label iff they form a scalar pair.
\item[(C14)] \emph{Term formula.}  Under local tying, the standing hypothesis and the congruence, the explicit
normalising constant $c_{\rm expl}$ is positive and the term formula of the pair holds.
\item[(C15)] \emph{Summability of rows.}  Under (C13), the concrete tie and $g_i\ge3$, the rows of $X(u,u')$
labelled by $\kappa$ are summable with constant $c_{\rm expl}\prod_p\alpha_p(u,u')$ against the point vectors.
\item[(C16)] \emph{Row form.}  Under the same hypotheses and the counting hypothesis,
$c_0c'_0\mathrm{KV}=n_0\sum_\rho\|\sum_pc_\rho(p)V(p)\|^2$ for $B(u,u')$.
\item[(C17)] \emph{Mass hypotheses.}  Under the gap hypothesis and (C13), $B(u,u')$ satisfies the mass hypotheses
of Definition~\ref{i3:S-def-mass}.
\item[(C18)] \emph{Link.}  $\gamma_1(X(u,u'))=\gamma_1(B)$, $\Sigma(\mathrm{sc}(\mathrm{gen}))=\gamma_1\Sigma_B$, a bijection of the
auxiliary places of $B$ with $\{1,2\}\times\{\text{infinite places of }E\}$ with $g_p=g_1,g_2$ accordingly,
$d=d_u$, $d'=d_{u'}$ and $D_*(B)=D_*$.
\item[(C19)] \emph{Split.}  $X(u,u')$ satisfies (AC), (NE$_{\mathrm{gen}}$), (NE$_{\mathrm{ng}}$), (GG$_{t(B)}$), (RO),
(CP), (IZ), and $\mathbf G_{ii}=\gamma_1$ for every $i$.
\end{enumerate}
\end{theorem}

The words ``eventually'', ``admissible constants'', ``standing'' have the meaning of the
standing hypotheses \eqref{i3:eq-standing}.  The \emph{gap hypothesis} on a datum is: $r\ge1801$ and
$b_{j+1}(\tau)-b_j(\tau)\ge16r^{19}$ for all $j$ and $\tau$.  It is not the conclusion of Lemma~\ref{i3:lem-gap}; it
follows from the standing hypotheses when $29\le C'$ and $16\le C_l$ (the node-gap clause gives
$b_{j+1}-b_j\ge C_lr^{19}f^5L\ge16r^{19}$, and $r\ge1801$ follows from $C'm^5fL\le k$ with $C'\ge29$ and the
other size clauses, as verified in \cite{CharacterFamily}).  The
proof of Theorem~\ref{i3:S-thm-components} is a list: components (C7), (C18), (C19) are deduced below from
smaller statements, (C4) is Proposition~\ref{i3:S-prop-NG}, and the others are formally verified propositions,
recorded with their sizes in Proposition~\ref{i3:S-prop-leaves}.

\begin{lemma}[Kernel expansion]\label{i3:S-lem-kexp}
Assume (C8)--(C10), that fundamental domains exist for the pair, that legitimate folding implies that the folded
kernels lie in $L^2$ when the archimedean coefficient is integrable, that the orbital integrals do not depend
on the fundamental domain, and the kernel expansion under all of these.  Then (C7) holds.
\end{lemma}

\begin{proof}
The concrete tie contains normalisation of $K_2$ and decay, so (C8) gives legitimate folds; with fundamental
domains (C9) gives convergent orbital integrals; integrability (part of the tie) puts the folded kernels in
$L^2$; independence of the fundamental domain follows; (C10) gives (AC).  These are the hypotheses of the
kernel expansion.
\end{proof}

\begin{lemma}[Typed bijection]\label{i3:S-lem-typbij}
Let $P$ be finite, $q:P\to\{g_1,g_2\}$ with $g_1\ne g_2$ and $\#q^{-1}(g_1)=\#q^{-1}(g_2)=f=\#\Omega$.  Then there
is a bijection $\beta:P\to\{1,2\}\times\Omega$ with $q(p)=g_1$ iff the first coordinate of $\beta(p)$ is $1$, and
$\prod_{p\in P}h(q(p))=h(g_1)^fh(g_2)^f$ for every $h$.
\end{lemma}

\begin{proof}
Choose bijections $q^{-1}(g_i)\to\Omega$ and put them together; split the product over $q^{-1}(g_1)\sqcup q^{-1}(g_2)$.
\end{proof}

\begin{proposition}[Typed count]\label{i3:S-prop-typed}
Let $E/F$ be a CM extension with $E$ totally complex, $f=[F:\Q]$, $R=\widehat{\mathcal O}_E=\mathcal O_E\otimes\widehat\Z$ with
the involution $x\mapsto\bar x$, $g_1\ne g_2$ primes each of which splits completely in $E$ (has $[E:\Q]$ primes above
it), $g=g_1g_2$, and $e\in R$ with $e^2=e$, $e\bar e=0$ and $1-e-\bar e\in gR$.  Put $K=\{x\in R:ex\in gR\}$.  A \emph{typed
reduction} of $K$ is a finite set $P$, finite fields $\mathbb F_v$ and ring homomorphisms $\mathrm{res}_v:R\to\mathbb F_v$
($v\in P$) such that $x\mapsto(\mathrm{res}_v(x))_v$ is surjective onto $\prod_v\mathbb F_v$ with kernel $K$, and
$|\mathbb F_v|\in\{g_1,g_2\}$ for every $v$.  Then for every typed reduction and each $i\in\{1,2\}$,
$\#\{v\in P:|\mathbb F_v|=g_i\}=f$; consequently $\prod_{v\in P}h(|\mathbb F_v|)=h(g_1)^fh(g_2)^f$ for every function $h$.
\end{proposition}

\noindent In \cite{CharacterFamily} the
statement is made for the typed reduction carried by a concretely tied package, whose data satisfy the
hypotheses above: the primes of the datum are distinct and split completely, and $1-e-\bar e\in gR$ follows from
the conditions on the datum by Lemma~\ref{i3:S-lem-onesub} below.  The proof is in \S\ref{i3:S-ssec-leafproofs}.

\emph{The data of $X(u,u')$ and $B(u,u')$ used below.}  We recall from the construction in
\cite{CharacterFamily} the parts of the term system and of the mass datum that enter (C18) and (C19); they are
definitions there.  (D1) A full package carries a typed reduction $(P,(\mathbb F_v),(\mathrm{res}_v))$ of
$K=\{x:ex\in gR\}$ (Proposition~\ref{i3:S-prop-typed}, with $e$ the idempotent of the Hida setting and $g=g_1g_2$)
and nontrivial additive characters $\psi_v$ of $\mathbb F_v$ (nontriviality is part of the definition of a full
package); the existence of a typed reduction with nontrivial characters, which is what makes full packages exist, is a
formally verified statement of \cite{CharacterFamily} (a written proof uses Dedekind's different theorem) and is listed
in Proposition~\ref{i3:S-prop-leaves}.  (D2) The index set of $X(u,u')$ is
$I=\prod_{v\in P}\mathrm{GL}_k(\mathbb F_v)$; the rows are triples $(\Lambda,\Lambda',i)$ with $\Lambda,\Lambda'$ in the finite row sets of
$u,u'$ and $i\in I$, and $\iota(\Lambda,\Lambda',i)=i$.  (D3) $H_1$ is an $\ell^2$-direct sum over the pairs $(\Lambda,\Lambda')$, and
$\mathrm{vec}_t$ has all coordinates zero except possibly the coordinate $(\Lambda_t,\Lambda'_t)$ of its row.  (D4)
$\mathbf G_{ij}=\overline{\gamma(j^{-1}i)}$ and $\gamma_1(X)=\Re\gamma(1)$, with $\gamma$ the function of
Proposition~\ref{i3:S-prop-GG} for $(\mathbb F_v,\psi_v)_{v\in P}$.  (D5) The auxiliary places of $B(u,u')$ are the primes
$\mathfrak p$ of $F$ dividing $g$, with $g_{\mathfrak p}$ the $g_i$ lying under $\mathfrak p$ ($=|\mathcal O_F/\mathfrak p|$ by
Lemma~\ref{i3:S-lem-splitprime}), $\mathrm{glk}_{\mathfrak p}=|\mathrm{GL}_k(\Z/g_{\mathfrak p})|$,
$\mathrm{uk}_{\mathfrak p}=|U_k(\Z/g_{\mathfrak p})|$; $\Sigma_B=\gamma_1(X)^{-1}\Sigma(\mathrm{sc}(\mathrm{gen}))$; $d=d_u$, $d'=d_{u'}$.
(D6) The corank $s_0$ and $G$ are defined from flags as in \S\ref{i3:S-ssec-leafproofs} (before
Corollary~\ref{i3:S-cor-CP}).

\begin{lemma}[Link]\label{i3:S-lem-link}
Let $g_1\ne g_2$ split completely in $E$, let $(X,Y)$ be concretely tied and $u\ne u'$, and assume
$D_*(B(u,u'))=D_*(X)$ (Proposition~\ref{i3:S-prop-dstar}).  Then (C18) holds.
\end{lemma}

\begin{proof}
For a prime $\ell$ put $a(\ell)=|\mathrm{GL}_k(\mathbb F_\ell)|/|U_k(\mathbb F_\ell)|$; it depends only on $\ell$ since finite fields
of equal cardinality are isomorphic.  \emph{$\gamma_1(X)=\gamma_1(B)$.}  By (D4) and Lemma~\ref{i3:S-lem-GGelem}(a),
$\gamma_1(X)=\gamma(1)=\prod_{v\in P}a(|\mathbb F_v|)$, which is $a(g_1)^fa(g_2)^f$ by Proposition~\ref{i3:S-prop-typed}.  By
Lemma~\ref{i3:S-lem-splitprime} each $g_i$ has exactly $f$ primes above it in $F$, each with residue field of $g_i$
elements, and no prime divides both; so by (D5) $\gamma_1(B)=\prod_{\mathfrak p}\mathrm{glk}_{\mathfrak p}/\mathrm{uk}_{\mathfrak p}=a(g_1)^fa(g_2)^f$.
\emph{$\Sigma$.}  By (D5) $\Sigma(\mathrm{sc}(\mathrm{gen}))=\gamma_1(X)\Sigma_B=\gamma_1(B)\Sigma_B$.  \emph{Bijection.}  By the same
count the map $\mathfrak p\mapsto g_{\mathfrak p}$ takes each value $g_1,g_2$ exactly $f$ times, and $E$ has $f$ infinite places;
Lemma~\ref{i3:S-lem-typbij} gives the bijection.  \emph{$d,d'$} by (D5).  \emph{$D_*$.}  $D_*(B)=D_*(X)$ by hypothesis,
and $D_*(X)=D_*$ is a clause of the concrete tie.
\end{proof}

\begin{proposition}[Bridge facts]\label{i3:S-prop-bridge}
Let $X$ be a Hida setting, $G(\mathbb A)=G_\infty\times G_f$ its adelic unitary group, $H(\mathbb A)$ the adelic unitary group of the
hermitian form of the datum, and $b:G(\mathbb A)\to H(\mathbb A)$ the adelic bridge (an isomorphism of topological groups compatible
with the matrix entries in the sense of (Br) below).  Let $K\subset G_f$ be the stabiliser of the lattice, $K_f\subset H(\mathbb A_f)$ the stabiliser of
$\widehat L$, and, for a prime $\mathfrak p$ of $F$, $K_{\mathfrak p}\subset H(F_{\mathfrak p})$ the stabiliser of the completed lattice
$L_{\mathfrak p}$.  Then:
\begin{enumerate}
\item[(B1)] $b(G(F))=H(F)$;
\item[(B2)] the finite part $b_f:G_f\to H(\mathbb A_f)$, $x\mapsto(b(1,x))_f$, is an isomorphism of topological groups;
\item[(B3)] $b_f(x)\in K_f$ if and only if $x\in K$;
\item[(B4)] $b_f$ is compatible with the projections to the local components $H(F_{\mathfrak p})$ and with matrix entries;
\item[(K)] $K$ and every $K_{\mathfrak p}$ are compact (each is the stabiliser of a lattice, hence closed and bounded).
\end{enumerate}
\end{proposition}

\noindent By (B3), $K$ and $K_f$ correspond under $b_f$.

\emph{The two models.}  $G_\infty=\prod_vU(J_{\varepsilon_v})$ (matrices over $\C$ in the basis adapted at $v$) and $G_f$ is the group of
$J$-unitary matrices over $\mathbb A_f\otimes_FE$; $G(F)$ is the set of $x$ for which there is $M\in U_J(E)$ with $x_v=P_v^{-1}v(M)P_v$ for
all $v$ and $x_f=M$.  $H(\mathbb A)$ is the group of $J$-unitary matrices over $\mathbb A_F\otimes_FE$, $H(F)$ those with entries in $E$, and
$K_f$ the set of $h$ whose component at every $\mathfrak p$ lies in $K_{\mathfrak p}$.  Let $\theta_v:\mathbb A_F\otimes_FE\to\C$ be the
projection to $F_v\otimes_FE\cong\C$ and $\tau:\mathbb A_F\otimes_FE\to\mathbb A_f\otimes_FE$ the finite projection; $a\mapsto((\theta_v(a))_v,\tau(a))$
is an isomorphism of topological rings ($\mathbb A_F=F_\infty\times\mathbb A_f$).  The bridge $b$ is an isomorphism of topological groups
with
\[
x_v=P_v^{-1}\,\theta_v(b(x))\,P_v\ \ (v\mid\infty),\qquad x_f=\tau(b(x))\qquad(x\in G(\mathbb A)),\tag{Br}
\]
$\theta_v$, $\tau$ applied entrywise.

\begin{proof}[Proof of Proposition~\ref{i3:S-prop-bridge}]
By (Br) and the ring isomorphism, $b(x)$ is the unique matrix whose images under $\theta_v$, $\tau$ are $P_vx_vP_v^{-1}$ and $x_f$.
(B1) If $x\in G(F)$ with matrix $M$, the matrix $M\in H(F)$ has the same images as $b(x)$, so $b(x)=M$; conversely if $b(x)=M$ has
entries in $E$, (Br) says $x\in G(F)$.  (B2) and (B4) By uniqueness, $b(1,x_f)$ is the matrix with infinite images $1$ and finite
image $x_f$, so $b_f(x_f)=x_f$ as a matrix over $\mathbb A_f\otimes_FE$: $b_f$ is the identity on matrices, a topological isomorphism of
$G_f$ onto $H(\mathbb A_f)$, compatible with entries and with the projections to the components at the primes $\mathfrak p$.  (B3) $\widehat L=
\prod_{\mathfrak p}L_{\mathfrak p}$ inside the restricted product $\mathbb A_f\otimes_FE=\prod'_{\mathfrak p}F_{\mathfrak p}\otimes_FE$, so a matrix
$x$ satisfies $x\widehat L=\widehat L$ if and only if $x_{\mathfrak p}L_{\mathfrak p}=L_{\mathfrak p}$ for every $\mathfrak p$; with (B2) this is (B3).
(K) $K_{\mathfrak p}$ is closed (the conditions $xL_{\mathfrak p}\subset L_{\mathfrak p}$, $x^{-1}L_{\mathfrak p}\subset L_{\mathfrak p}$ are closed)
and bounded (choose free lattices $L'\subset L_{\mathfrak p}\subset L''$; the entries of $x$ in a basis of $L'$ are bounded), hence
compact; $K$ is the closed subset of unitary matrices in $\prod_{\mathfrak p}\mathrm{Stab}(L_{\mathfrak p})$, which is compact by
Tychonoff and carries the restricted product topology.
\end{proof}

\begin{lemma}[Level volumes]\label{i3:S-lem-LV}
Let $(X,Y)$ be a package satisfying the measure tie.  Write $G(\mathbb A)=G_\infty\times G_f$ for the adelic unitary group of
$X$ with the product Haar measure $\mu=\mu_\infty\times\mu_f$, $\mathcal F\subset G(\mathbb A)$ the fundamental domain of $G(F)$
fixed by the package, $b:G(\mathbb A)\to H(\mathbb A)$ the adelic bridge, $b_f:G_f\to H(\mathbb A_f)$ its finite part
(Proposition~\ref{i3:S-prop-bridge}), and $K_X=b_f(K_2)$ the level unitary group.  Then
\[
\mathrm{vol}_G(B(u,u'))=\mu(\mathcal F)\qquad\text{and}\qquad \mu_{G_f}(K_X)=\mu_f(K_2).
\]
Here $\mathrm{vol}_G(B(u,u'))$ is the covolume of $H(F)$ in $H(\mathbb A)$ for $\mu_G$: the $\mu_G$-measure of any
fundamental domain of $H(F)$, and $0$ if there is none.
\end{lemma}

\begin{proof}
The measure tie gives $\mu_G=b_*\mu$ and $\mu_{G_f}=(b_f)_*\mu_f$.  \emph{Finite level.}  Since $b_f$ is injective (Proposition~\ref{i3:S-prop-bridge}(B2)),
$b_f^{-1}(b_f(K_2))=K_2$, so $\mu_{G_f}(K_X)=\mu_f(b_f^{-1}(b_fK_2))=\mu_f(K_2)$.  \emph{Covolume.}  Put $\mathcal F'=b(\mathcal F)$.
It is measurable since $b$ is a homeomorphism.  For $\gamma\in G(F)$ we have $b(\gamma x)=b(\gamma)b(x)$, so the translates
$\delta\mathcal F'$, $\delta\in H(F)=b(G(F))$ (Proposition~\ref{i3:S-prop-bridge}(B1)), are the images under $b$ of the translates $\gamma\mathcal F$.  The latter cover
$G(\mathbb A)$ up to a $\mu$-null set and are pairwise $\mu$-almost disjoint.  As $b$ carries $\mu$ to $\mu_G$, and $b^{-1}$ carries
$\mu_G$-null sets to $\mu$-null sets, $\mathcal F'$ is a fundamental domain of $H(F)$ for $\mu_G$.  Moreover
$\mu_G(\mathcal F')=\mu(b^{-1}(b\mathcal F))=\mu(\mathcal F)$.  The group $H(F)$ is countable, being in bijection with $G(F)$, and $\mu_G$
is a left-invariant Haar measure, being the image of one under a topological isomorphism.  Two fundamental domains of a
countable group acting by measure-preserving left translations have equal measure.  (Write each as the disjoint union of its
intersections with the translates of the other, and translate the pieces back.)  So the covolume is $\mu(\mathcal F)$.
\end{proof}

\begin{proposition}\label{i3:S-prop-KppVol}
Let $(X,Y)$ be a package whose level subgroup $K_2$ is normalised and is a subgroup, and whose character $\chi_2$ is a unitary
character of $K_2$.  Assume the measure tie, and that the sets $F_N(u)$, $F_N(u')$ of the full package used in $B(u,u')$ are strict
fundamental domains of the integral points of $N_u$, $N_{u'}$ at infinity.  Then $\mu K''(B(u,u'))=\mu_{G_f}(K_X)$.
\end{proposition}

\noindent Here (D7) the level volumes of $B(u,u')$ are, by construction, as follows.  $P$ is the set of primes $\mathfrak p$ of $F$
dividing $g=g_1g_2$; for $\mathfrak p\in P$, $\pi_{\mathfrak p}:H(\mathbb A_f)\to H(F_{\mathfrak p})$ is the projection, $\mu_{\mathfrak p}$ the
Haar measure on $H(F_{\mathfrak p})$ fixed by the package, $K_{\mathfrak p}\subset H(F_{\mathfrak p})$ the stabiliser of the completed lattice
$L_{\mathfrak p}$, and $K_f\subset H(\mathbb A_f)$ the stabiliser of $\widehat L$ (it corresponds to $K$ under $b_f$,
Proposition~\ref{i3:S-prop-bridge}(B3)); the measures $\mu_{\mathfrak p}$ and $\mu_{G_f}$ are Haar measures by the normalisation of
the package (not as a consequence of the measure tie).  Then
$\mu K''_{\mathfrak p}=\mu_{\mathfrak p}(K''_{\mathfrak p})$ with $K''_{\mathfrak p}=\pi_{\mathfrak p}(K_X)$, and the factor outside $P$ is
$\mu_{G_f}(K_f)/\prod_{\mathfrak p\in P}\mu_{\mathfrak p}(K_{\mathfrak p})$, so that
\[
\mu K''(B(u,u'))=\frac{\mu_{G_f}(K_f)}{\prod_{\mathfrak p\in P}\mu_{\mathfrak p}(K_{\mathfrak p})}\prod_{\mathfrak p\in P}\mu_{\mathfrak p}(K''_{\mathfrak p}).
\]
The proof (\S\ref{i3:S-ssec-kppvol}) uses that $K_2$ is a subgroup with $K(g)\subset K_2\subset K$, that $K_2$ is cut out inside $K$ by
conditions $(x-1)v\in M$ (Step~3(ii)), and that $\mu_{G_f}$, $\mu_{\mathfrak p}$ are Haar measures.  The measure tie and the strict
fundamental domains are carried by the formal statement but not used.

\begin{proposition}[$D_*$ identification]\label{i3:S-prop-dstar}
Let $(X,Y)$ be a package whose level subgroup $K_2$ is normalised and is a subgroup, and whose character $\chi_2$ is a unitary
character of $K_2$.  Assume the measure tie, and that the sets $F_N(u)$, $F_N(u')$ of the full package used in $B(u,u')$ are strict
fundamental domains of the integral points of $N_u$, $N_{u'}$ at infinity.  Then
\[
\text{(LV)}\qquad \mathrm{vol}_G(B(u,u'))=\mu(\mathcal F),\qquad \mu K''(B(u,u'))=\mu_f(K_2),
\]
and $D_*(B(u,u'))=D_*(X)$.
\end{proposition}

\begin{proof}
(LV) is Lemma~\ref{i3:S-lem-LV} combined with Proposition~\ref{i3:S-prop-KppVol}.
$e=\overline{\chi_2}\mathbf 1_{K_2}/\mu_f(K_2)$, and $\chi_2(1)=1$ since $\chi_2$ is a character, so $\Re e(1)=\mu_f(K_2)^{-1}$
($0<\mu_f(K_2)<\infty$ by normalisation).  By definition $D_*(B)=\deg\cdot\mathrm{vol}_G/\mu K''$ with $\deg=\prod_vd_v$,
and by (LV) this is $(\prod_vd_v)\mu(\mathcal F)/\mu_f(K_2)=(\prod_vd_v)\mu(\mathcal F)\Re e(1)=D_*(X)$.
\end{proof}

\begin{proposition}[Gelfand--Graev Gram bound]\label{i3:S-prop-GG}
Let $(\mathbb F_v,\psi_v)_{v\in P}$ be a finite family of finite fields with nontrivial additive characters,
$q_v=|\mathbb F_v|\ge3$, $G=\prod_v\mathrm{GL}_k(\mathbb F_v)$, $U_k$ the upper unitriangular group,
$\psi_{U,v}(u)=\psi_v(\sum_iu_{i,i+1})$, $\psi^0_{U,v}$ its extension by $0$ to $\mathrm{GL}_k(\mathbb F_v)$,
$\gamma(x)=\prod_v|U_k(\mathbb F_v)|^{-1}\sum_{y\in\mathrm{GL}_k(\mathbb F_v)}\psi^0_{U,v}(yx_vy^{-1})$ and $P_q=\prod_v(1+3/q_v)$.  Let $G$ act
freely on the right on a finite set $X$, $d:X\times X\to G$ with $j\cdot d(i,j)=i$ whenever $i,j$ lie in one orbit,
and $\mathbf G_{ij}=\overline{\gamma(d(i,j))}$ if $i,j$ lie in one orbit, $0$ otherwise.  Then $\gamma(1)>0$, $\mathbf G$ and its
transpose are positive semidefinite, and $P_q\gamma(1)\mathbf 1-\mathbf G\succeq0$.
\end{proposition}

\noindent The proof is in \S\ref{i3:S-ssec-leafproofs}; it uses one published theorem (multiplicity one of the
Gelfand--Graev representation).

\begin{lemma}[Split]\label{i3:S-lem-split}
Let $g_1\ne g_2$ split completely in $E$ with $g_1,g_2\ge3$.  Assume (C8)--(C10), the norm expansions
(NE$_{\mathrm{gen}}$), (NE$_{\mathrm{ng}}$) for $X(u,u')$ under the tie, and the hypotheses of Proposition~\ref{i3:S-prop-idcell} for $X(u,u')$, $u\ne u'$ (fundamental
domains, convergence of the orbital integrals, finite volume of $H_u$ at infinity, (Conj), (Ineq), (Stab)).  Then for a concretely tied package and standing $u\ne u'$, (C19)
holds and $t(B)=(1+3/g_1)^f(1+3/g_2)^f=P_g$.
\end{lemma}

\begin{proof}
(AC) is obtained from (C8)--(C10) as in Lemma~\ref{i3:S-lem-kexp}.  \emph{(RO) and the diagonal.}  If
$\iota(\mathrm{row}\,t)=\iota(\mathrm{row}\,t')$ and the rows differ, then by (D2) $(\Lambda_t,\Lambda'_t)\ne(\Lambda_{t'},\Lambda'_{t'})$, so by (D3)
$\mathrm{vec}_t,\mathrm{vec}_{t'}$ have disjoint coordinate supports and are orthogonal.  By (D4) $\mathbf G_{ii}=\overline{\gamma(1)}$, and
$\gamma(1)$ is a positive real number (Lemma~\ref{i3:S-lem-GGelem}(a)), so $\mathbf G_{ii}=\gamma(1)=\Re\gamma(1)=\gamma_1$.
\emph{(GG).}  Take $X=I=G$ acting on itself by right multiplication (free) and $d(i,j)=j^{-1}i$, so $j\,d(i,j)=i$; the
matrix of Proposition~\ref{i3:S-prop-GG} is $\overline{\gamma(j^{-1}i)}=\mathbf G_{ij}$, and $|\mathbb F_v|\in\{g_1,g_2\}$, so
$|\mathbb F_v|\ge3$.  Hence $t'\gamma_1\mathbf 1-\mathbf G\succeq0$ with $t'=\prod_v(1+3/|\mathbb F_v|)$, which is
$(1+3/g_1)^f(1+3/g_2)^f$ by Proposition~\ref{i3:S-prop-typed} with $h(m)=1+3/m$; the same count over the auxiliary
places of $B$ (Lemma~\ref{i3:S-lem-splitprime}, (D5)) gives $t(B)=\prod_{\mathfrak p}(1+3/g_{\mathfrak p})=t'$.  \emph{(CP)} is
Corollary~\ref{i3:S-cor-CP}.  \emph{(NE)} are hypotheses.  \emph{(IZ).}  The hypotheses of
Proposition~\ref{i3:S-prop-idcell} are assumed (they are among the leaves of Proposition~\ref{i3:S-prop-leaves}; the
fundamental domains are those of the hypotheses of Lemma~\ref{i3:S-lem-kexp}, and the convergence of the orbital
integrals is (C9)).  By Proposition~\ref{i3:S-prop-idcell} every identity term has vanishing operator, so $\mathrm{part}(\mathrm{id})=\sum_{\mathrm{id}}\mathrm{op}_t=0$ (the sum converges by (AC)).
\end{proof}

\begin{proposition}\label{i3:S-prop-idcell}
(Identity cell.)  Let $u\ne u'$.  Suppose that for every term $t$ of the block $X(u,u')$ the stabiliser $\Gamma_t$ has a
strict fundamental domain $\Omega_t$ and the orbital integral $\mathrm{OO}_t$ converges absolutely (its integrand is
$\mu_\infty$-integrable on $\Omega_t$), that the Bessel subgroup $H_u$ has finite volume at infinity, $\mu_{H,u}(H_u)<\infty$,
and that for every identity term $t$ with element $\gamma$: (Conj) $\gamma^{-1}S_u\gamma\subset S_{u'}$ with
$\psi_{u'}(\gamma^{-1}s\gamma)=\psi_u(s)$ for $s\in S_u$, and $\gamma^{-1}N_u\gamma\subset N_{u'}$; (Ineq) every linear
$T:V_\pi(u')\to V_\pi(u)$ with $T\pi_{u'}(\gamma^{-1}h\gamma)=\pi_u(h)T$ for $h\in H_u$ is zero; (Stab) the stabiliser $\Gamma_t$ is
a countable subgroup of $N_u$.  Then $\mathrm{OO}_t=0$ and $\mathrm{op}_t=0$ for every identity term $t$, and
$\mathrm{part}(\mathrm{id})=\sum_{t\ \mathrm{identity}}\mathrm{op}_t=0$.
\end{proposition}

\noindent The objects ($S_u$, $N_u$, $H_u$, $\psi_u$, $\pi_u$, $W_u$, $\Gamma_t$, $\Omega_t$, $\mathrm{OO}_t$) are those recalled at the
beginning of \S\ref{i3:ssec-pf-kppoff}; the proof, an unfolding along $N_uH_u$ followed by Schur orthogonality, is in
\S\ref{i3:S-ssec-pf-idcell}.  The hypotheses (Conj), (Ineq), (Stab) for $X(u,u')$ have written proofs (inequivalence via a lift and
tensor products), but the finite volume of $H_u$ at infinity does not; the formal development proves all of them (they are discharged
inside the formally verified off-diagonal bound, whose only hypotheses are those of Proposition~\ref{i3:prop-S-closed}), and they
are listed in Proposition~\ref{i3:S-prop-leaves}.

\begin{proposition}[formally verified, \cite{CharacterFamily}]\label{i3:S-prop-leaves}
Under (SW), (FD) and Proposition~\ref{i3:S-prop-NG}, the following hold as stated in Theorem~\ref{i3:S-thm-components}: the
components (C1), (C2), (C3), (C5), (C6), (C8)--(C17); the norm expansions used in Lemma~\ref{i3:S-lem-split}; for $u\ne u'$
standing, the hypotheses of Proposition~\ref{i3:S-prop-idcell} for $X(u,u')$, namely strict fundamental domains $\Omega_t$ and
integrability of the orbital integrand for every term, $\mu_{H,u}(H_u)<\infty$, and (Conj), (Ineq), (Stab) for every identity term;
the existence of a typed reduction with nontrivial characters $\psi_v$ (used in (D1)); and the hypotheses of
Lemma~\ref{i3:S-lem-kexp} besides (C8)--(C10).  (The bridge facts and the compactness of $K$ and $K_{\mathfrak p}$, used in
Lemma~\ref{i3:S-lem-LV} and \S\ref{i3:S-ssec-kppvol}, are Proposition~\ref{i3:S-prop-bridge}, proved above.)
\end{proposition}

\noindent\emph{What the proof involves.}  Each component is a separate development; the transitive dependencies of each are
large (shared with the rest of \cite{CharacterFamily}).  (C6) and (C16) use
Proposition~\ref{i3:prop-amp} (the archimedean orbital integrals are multiples of Bessel integrals) and
Proposition~\ref{i3:prop-bessel}; (C12) and (C17) use Proposition~\ref{i3:prop-bessel}; (C16) and (C17) use (SW)
(normalising constants, stabiliser volumes and generic masses); (C11) uses (FD) and (C12).  The hypotheses of
Proposition~\ref{i3:S-prop-idcell} are: fundamental domains (the same statement as in Lemma~\ref{i3:S-lem-kexp}), integrability
(component (C9), which needs both folds legitimate), positivity and finiteness of the volumes of the package, and three structural
statements (Conj), (Ineq), (Stab).  The bridge facts rest on the description of the rational points and of the finite-level compact
subgroup under the bridge.  The identity (LV), the $D_*$ identification, the identity cell and the transport are not in this list:
they are proved in \S\S\ref{i3:S-ssec-kppvol}, \ref{i3:S-ssec-pf-idcell}, \ref{i3:S-ssec-pf-transport}.

\begin{proof}[Proof of Theorem~\ref{i3:S-thm-components}]
(C4) is Proposition~\ref{i3:S-prop-NG}; (C7) is Lemma~\ref{i3:S-lem-kexp}; (C18) is Lemma~\ref{i3:S-lem-link} with
Propositions~\ref{i3:S-prop-typed}, \ref{i3:S-prop-dstar}; (C19) is Lemma~\ref{i3:S-lem-split} with
Propositions~\ref{i3:S-prop-GG}, \ref{i3:S-prop-idcell}; the rest is Proposition~\ref{i3:S-prop-leaves}.
\end{proof}

\medskip\noindent\label{i3:S-sss-volume}\emph{The volume count of the intertwiner spaces.}

The kernel blocks $K(u,u')$ are compressions of one operator: convolution with
\[
\phi_e(x)=\Bigl(\prod_vd_v\Bigr)\overline{\prod_v\langle w_v,\sigma_v(x_v)w_v\rangle}\,e(x_f),\qquad
\mathcal K(x,y)=\sum_{\gamma\in U(F)}\phi_e(x^{-1}\gamma y),
\]
on $L^2(U(F)\backslash U(\mathbb A))$, which (under the gap) is the orthogonal projection onto the
$e$-isotypic intertwiners from $\sigma_\infty=\bigotimes_v\sigma_v$.  Its trace is the dimension $D$ of the space of
intertwiners, and its main term is $M=\phi_e(1)\mu(\mathcal F)=(\prod_vd_v)\Re e(1)\mu(\mathcal F)=D_*(X)$: the
expected dimension of the concrete tie.  In the off-diagonal bound two facts about $\mathcal K$ are used: that
$D$ is finite (this is (FD), Lemma~\ref{i3:S-lem-FD}, entering (C11)), and that $D_*(X)$ is the expected value
(Proposition~\ref{i3:S-prop-dstar}).  The quantitative count $|D-M|\le e^{-\frac1{16}nf\log n}M$ is not used for
Proposition~\ref{i3:prop-S-closed}; it is the count bound of Proposition~\ref{i3:prop-conc1}.  We record its
structure since it is the same kernel.

\begin{definition}[Counting hypothesis]\label{i3:S-def-CH}
Let $n$ be the rank, $f$ the degree, $\mathrm{rd}$ the root discriminant, $(p,q)$ the signature at every infinite
place, $g$ the level, $g_0$ an auxiliary prime, $X(b)=nf\,\mathrm{rd}\,g\,b$, $N=n^2f$ and
$\tau_*=X(b)^{-11q}$.  $\mathrm{CH}_{\rm vol}(C_g,n_0,b)$ is: discrete series data with formal degrees $d_v$; $n_0\le n$;
$q^2\le n$; $\log\mathrm{rd}\le e^n$; $\log f\le\sqrt n$; $2\le b$, $\log b\le e^{nf\log n/16}$; $|\lambda_{v,a}|\le b^2$;
$1\le g\le b^{C_g}$; $g_0$ prime, $4n<g_0$, totally split in $E$; and every $\lambda_v$ has gap at least
$8\cdot10^{14}(C_g+4)n^{19}f^5\log(b\,\mathrm{rd})$.  For $x\in U(\mathbb A)$, $\tau(x)$ is the least normalised size of a
nonzero submodule of the lattice $\Lambda_x$.
\end{definition}

\noindent This is a different normalisation of the counting hypothesis
$\mathrm{CH}(n_0,A,b,h)$ of \S\ref{i3:ssec-counting} (with the gap constant made
explicit); to avoid a clash we write $\mathrm{CH}_{\rm vol}$.  The count used
in this section is Proposition~\ref{i3:prop-counting}; the two statements below
only record the structure of its proof.

\begin{proposition}[formally verified, \cite{CharacterFamily}]\label{i3:S-prop-regions}
Assume (SW), $\mathrm{CH}_{\rm vol}(C_g,n_0,b)$ with $n_0$ large in terms of $C_g$, and $q\ge1$.  Then: (a) $D<\infty$; (b) the
integral over $\{\tau\ge\tau_*\}\cap\mathcal F$ of $\sum_{\gamma\ne1\text{ near}}|\phi_e(x^{-1}\gamma x)|$ is at most
$e^{-\frac14nf\log n}M$; (c) the same integral over $\gamma$ not near is at most $e^{-n^2f}M$; (d)
$\int_{\{\tau<\tau_*\}}\mathcal K(x,x)\,d\mu\le De^{-n^2f}$; (e) $\mu(\{\tau<\frac12\}\cap\mathcal F)\le e^{-\frac14nf\log n}\mu(\mathcal F)$;
(f) $\mathcal K$ converges absolutely, $\mathcal K(x,x)\ge0$, $\int_{\mathcal F}\mathcal K(x,x)\,d\mu=D$ and
$|\mathcal K(x,x)-\phi_e(1)|$ is bounded by the near and far sums.  Here $\gamma$ is \emph{near} $x$ if
$\phi_e(x^{-1}\gamma x)\ne0$ and $d_{HS}((x^{-1}\gamma x)_v)\le10^{-6}n^{-8}f^{-2}$ at every $v$.
\end{proposition}

\noindent\emph{What the proof involves.}  Decay of lowest-$K$-type coefficients under the gap (a compression
formula), a lattice count for far terms, unipotence of near elements (from $4n<g_0$), Siegel-type bounds on
the measure of the region where $\Lambda_x$ has short submodules (flag measures), and the spectral expansion of
$\mathcal K$.

\begin{proposition}[Counting theorem; formally verified, \cite{CharacterFamily}]\label{i3:S-prop-count}
Assume (SW).  Under $\mathrm{CH}_{\rm vol}(C_g,n_0,b)$ with $n_0$ large in terms of $C_g$, $D<\infty$ and
$|D-M|\le e^{-\frac1{16}nf\log n}M$.
\end{proposition}

\noindent The statement is formally verified for all signatures.  For $q\ge1$ we deduce it from
Proposition~\ref{i3:S-prop-regions} below.  The case $q=0$ is not deduced here: the argument ``only $\gamma=1$
contributes'' would need the unipotence of the contributing $\gamma$ and the trace identity, which
Proposition~\ref{i3:S-prop-regions} provides only for $q\ge1$; for $q=0$ we rely on the formal statement.

\begin{proof}[Proof for $q\ge1$]
Let $\mathcal M=\{\tau\ge\tau_*\}\cap\mathcal F$.  Write
$\int_{\mathcal F}\mathcal K(x,x)-\phi_e(1)\mu(\mathcal F)=\int_{\mathcal M}(\mathcal K-\phi_e(1))+\int_{\mathcal F\setminus\mathcal M}\mathcal K-\phi_e(1)\mu(\mathcal F\setminus\mathcal M)$.
By Proposition~\ref{i3:S-prop-regions}(f),(b),(c) the first term is at most $(e^{-\Lambda/4}+e^{-n^2f})M$ in absolute value
($\Lambda=nf\log n$); the second lies in $[0,De^{-n^2f}]$ by (d); the third in $[-e^{-\Lambda/4}M,0]$ by (e), as
$\tau_*<\frac12$.  So $|D-M|\le(2e^{-\Lambda/4}+e^{-n^2f})M+e^{-n^2f}D$.  With $\beta=e^{-n^2f}\le\frac12$ this gives
$|D-M|\le\frac{2e^{-\Lambda/4}+2\beta}{1-\beta}M\le8e^{-\Lambda/4}M\le e^{-\Lambda/16}M$, since $n^2f\ge\Lambda\ge240$ for
$n\ge e^{15}$.
\end{proof}

\medskip\noindent\label{i3:S-sss-HS}\emph{From the term operator to Hilbert--Schmidt norms.}

For a linear map $A:W\to W'$ with $W$ finite dimensional, $\|A\|_{HS}^2=\sum_i\|Ae_i\|^2$ over an orthonormal
basis, and $\iota(A)=(Ae_i)_i\in W'^{\dim W}$ with the $\ell^2$ norm.  The kernel block is
$K(u,u')=(c_0(u)c_0(u'))^{-1/2}P_u^*P_{u'}:C_{u'}\to C_u$, and $K''(u,u')$ is the sum of its terms with $G=0$
that are not identity terms (the Kloosterman part).

\begin{lemma}\label{i3:S-lem-hs}
(1) $\|\iota(A)\|=\|A\|_{HS}$, and $\iota(A)=0$ implies $A=0$.  (2) If $V:W\to W_1$ is a linear isometric
isomorphism, or a surjective conjugate-linear map with $\langle Vx,Vy\rangle=\langle y,x\rangle$, and $g:W'\to W'_1$
preserves norms with $A_1(Vx)=g(Ax)$, then $\|A_1\|_{HS}=\|A\|_{HS}$.
\end{lemma}

\begin{proof}
(1) is the definition.  (2) The HS norm does not depend on the orthonormal basis; $(Ve_i)$ is an orthonormal
basis of $W_1$ and $\sum\|A_1Ve_i\|^2=\sum\|g(Ae_i)\|^2=\sum\|Ae_i\|^2$.
\end{proof}

\emph{Column data.}  We recall the objects of a full package $Y$ over a Hida setting $X$ that enter here.  For a
word $u$, $C_u$ is the column space: the space of linear maps $T$ from $V_\pi(u)$ to the functions
$H_{u,\infty}\times H_{u,f}\to\C$ such that every $T\xi$ is a form (left invariant under the rational points, right
$(K_u,\psi)$-equivariant) and $T(\pi(h)\xi)(y)=(T\xi)(y(h,1))$.  $\mathrm{Cl}_u$ is the finite set of column classes with representatives $h_\rho$, and
$\mathrm{CM}_u=V_\pi(u)^{\mathrm{Cl}_u}$ with the $\ell^2$ inner product.  The column map $\sharp_u:C_u\to\mathrm{CM}_u$ sends $T$ to
$(\rho\mapsto$ the Riesz representative in $V_\pi(u)$ of the functional $\xi\mapsto(T\xi)(1,h_\rho))$.  The package has
\emph{finite column data} if (among other clauses) every $\sharp_u$ is bijective and the Poincare-series maps
$P_u:C_u\to\mathcal H$ into the finite-dimensional space $\mathcal H$ of the package are linear; $C_u$ is then given the inner
product $\langle T,T'\rangle_{C_u}=\langle\sharp_uT',\sharp_uT\rangle$ (Lemma~\ref{i3:S-lem-sharp}).  The volumes of the package are
$\mathrm{vol}_N(u)=\mu_f(K_2)\,\mu_{N,u}(F_N(u))$ and $\mathrm{vol}_H(u)=\mu_{H,u}(H_u)$, with $F_N(u)$ the fundamental domain of $N_u$
fixed by $Y$, and the two normalising constants are
\[
c_{0,X}(u)=\frac{\mathrm{vol}_N(u)\,\mathrm{vol}_H(u)}{\dim V_\pi(u)},\qquad c_0(u)=\frac{\mathrm{vol}_N(u)\,\mathrm{vol}_H(u)}{d_u},\qquad
d_u=\prod_v\delta(\lambda_{u,v}),
\]
where $\delta(\lambda)$ is the dimension of an irreducible unitary representation of $U(k)$ of the highest weight attached to the
archimedean parameter $\lambda$ (the least such dimension, by definition; by Weyl's formula all such representations have the same
dimension).  The kernel form is $\mathrm{Kform}(T,T')=(c_{0,X}(u)c_{0,X}(u'))^{-1/2}\langle P_uT,P_{u'}T'\rangle$ and the kernel block
is $K(u,u')=(c_0(u)c_0(u'))^{-1/2}P_u^*P_{u'}:C_{u'}\to C_u$.  A \emph{draft datum} consists of a finite-dimensional Hilbert space
$\mathcal H$, finite-dimensional Hilbert spaces $C_u$, linear maps $P_u:C_u\to\mathcal H$, positive volumes, and further combinatorial
data; an isomorphism of draft data consists of linear isometric bijections $e_{\mathcal H}$, $e_u$ with
$e_{\mathcal H}P_u=P'_ue_u$, equal volumes and bijections of the remaining data.

\begin{proposition}[Transport]\label{i3:S-prop-transport}
Let $(X,Y)$ be a package with finite column data and $u,u'$ words.  Then:
\begin{enumerate}
\item[(a)] $\sharp_u$ is conjugate-linear, and $\langle\sharp_uT,\sharp_uT'\rangle=\langle T',T\rangle_{C_u}$;
\item[(b)] $\dim V_\pi(u)=d_u$;
\item[(c)] $c_{0,X}(u)=c_0(u)$;
\item[(d)] if $K^\flat:\mathrm{CM}_{u'}\to\mathrm{CM}_u$ is linear with $\mathrm{Kform}(T,T')=\langle K^\flat(\sharp_{u'}T'),\sharp_uT\rangle$ for all
$T\in C_u$, $T'\in C_{u'}$, then $K(u,u')=\sharp_u^{-1}K^\flat\sharp_{u'}$;
\item[(e)] for isomorphic draft data the Hilbert--Schmidt norms of the kernel blocks $(c_0(u)c_0(u'))^{-1/2}P_u^*P_{u'}$ coincide.
\end{enumerate}
\end{proposition}

\noindent The proof is in \S\ref{i3:S-ssec-pf-transport}; it uses one published theorem (Weyl's dimension formula).  Part (e) is the
``isomorphism of the package with the draft data of $\mathcal D$'' used below: the concrete tie provides such an isomorphism.

\begin{theorem}[Identification off the diagonal]\label{i3:S-thm-HS}
Let $(X,Y)$ be concretely tied, $u\ne u'$, and assume (C7), (AC) and $\mathrm{part}(\mathrm{deg})=0$ in $X(u,u')$.  Then
$K''(u,u')=0$ and $\|K(u,u')\|_{HS}^2=\|K_0(u,u')\|^2$.
\end{theorem}

\begin{proof}
By (C7) and Proposition~\ref{i3:S-prop-transport}(d), $K(u,u')=\sharp_u^{-1}K^\flat\sharp_{u'}$ with $K^\flat=\sum_ta_t$ and
$\mathrm{op}_t=\iota(a_t)$, and $\iota$ is continuous and linear, so $K=\iota(K^\flat)$.  The terms of the package are sums of terms of $X(u,u')$ over finite fibres, degenerate
terms corresponding to fibres of degenerate terms; regrouping an absolutely convergent sum along finite
fibres does not change it, so $K''$ is carried to the transport of $\sum_{\mathrm{deg}}a_s$, whose image under $\iota$
is $\mathrm{part}(\mathrm{deg})=0$; by Lemma~\ref{i3:S-lem-hs}(1) $K''(u,u')=0$.  By Lemma~\ref{i3:S-lem-fourparts},
$K=K_0+\mathrm{part}(\mathrm{deg})=K_0$.  Lemma~\ref{i3:S-lem-hs}(2) with $V=\sharp_{u'}$, $g=\sharp_u$ and
Proposition~\ref{i3:S-prop-transport} give $\|K(u,u')\|_{HS}^2=\|K^\flat\|_{HS}^2=\|\iota(K^\flat)\|^2=\|K_0(u,u')\|^2$.
\end{proof}

\medskip\noindent\label{i3:S-sss-assembly}\emph{Assembly.}

\begin{theorem}[Theorem S at a pair]\label{i3:S-thm-pair}
Assume (SW), (FD) and the conclusion of Proposition~\ref{i3:S-prop-NG}.  For every $C_{\rm thr}$ there are constants
such that, for admissible size constants, there are $k_0$, $C_n\ge0$, $C'_0$, $C_{l,0}$ with: for $C'\ge C'_0$,
$C_l\ge C_{l,0}$, every datum $\mathcal D$, package $(X,Y)$ and words $u\ne u'$ satisfying the size hypotheses with
$k\ge k_0$, the standing hypotheses at $(C',C_l)$, the gap hypothesis, local tying, the frame condition and the
loops (RB$_{\varepsilon_L,d}$), (RS$_{\varepsilon_L,d}$) with $d=d_ud_{u'}/D_*$,
\[
\|K_0(u,u')\|^2\le(1+b^{-6})\frac{d_ud_{u'}}{D_*}\quad\text{and}\quad\|K(u,u')\|_{HS}^2\le(1+b^{-6})\frac{d_ud_{u'}}{D_*}.
\]
\end{theorem}

\begin{proof}
By Theorem~\ref{i3:S-thm-components}, all nineteen components hold.  The gap hypothesis gives (C11) and (C12), hence
the concrete tie.  (C13) supplies base points, measures, local spaces, uniformisers and the row labelling;
(C14) the term formula; $g_i\ge(kb)^8\ge3$.  (C6) gives $C_4\ge1$ and the real inputs of the kernel datum with
constants $(C_4,J_N)$; (C1), (C2), (C3) give the local model, the density input and the rows input; (C16) gives
the row form, which with summability (C15) yields $\mathrm{KV}\ge0$ and the bound $c_0c'_0\mathrm{KV}\le n_0M$; (C17)
gives the mass hypotheses; $J_N=j_n$, $e_1=\max(s_d,1)$ and $\delta=C(r_S+1)(C/\Theta_h)^f$ hold by definition of
$B(u,u')$.  (C18) and (C19) give the identities and conditions of Theorem~\ref{i3:S-thm-pbapplied}; (C4) gives (RN)
with $\varepsilon'=2^{-(k-2r_S-5)}e^{C_n(2r_S+2)^5f\log b}$; the size numbers hold with $c=\frac13$ under the standing
hypotheses; the loops are hypotheses with $\varepsilon_L=e^{-\frac13kf\log k}$.  Theorem~\ref{i3:S-thm-pbapplied} gives the
first inequality.  (C5) gives $\mathrm{part}(\mathrm{deg})=0$, (C7) and (C10) the kernel expansion and (AC), and
Theorem~\ref{i3:S-thm-HS} gives the second.
\end{proof}

\begin{proof}[Proof of Proposition~\ref{i3:prop-S-closed}]
(FD) is Lemma~\ref{i3:S-lem-FD} and the conclusion of Proposition~\ref{i3:S-prop-NG} follows from (SW), so
Theorem~\ref{i3:S-thm-pair} applies.  The rank row and the loop row are closed statements; take $\varepsilon_1$ at most
theirs and, for $\varepsilon\le\varepsilon_1$, a threshold system containing theirs, those of
Lemma~\ref{i3:S-lem-rankrow} (the thresholds $40$, $40$, $1$), the thresholds $C'_0,C_{l,0}$ of Theorem~\ref{i3:S-thm-pair}, and the
Large conditions $29\le C_{k,\rm low}$, $29\le C_l^{\mathcal F}$, $1000+c_{\rm orb}\le C_l^{\mathcal F}$ (the last gives
$16\le C_l^{\mathcal F}$ and $1000\le C_l^{\mathcal F}$).  The constants of Theorem~\ref{i3:S-thm-pair} are chosen uniformly as
functions of the outer constants of $\mathcal F$ (with $C_0=3$, $(C_G,C_\Delta,c_T,C_b)=(C_\Gamma,5,\frac12,3)$), so these
thresholds form a threshold system.  Let $\mathcal F\in\mathcal C$ be large for it.  By Lemma~\ref{i3:S-lem-rankrow}, the
rank row gives the frame row, the size row and the size bound $\mathrm{Sizes3b}(a(3))$, $a(3)=3e^{7/4}$; the frame row, (FD) and the decay
(C12) show that the selection is admissible: eventually every selected package is concretely tied, and the member is
locally tied (hypothesis) and satisfies the frame condition.  At a member satisfying the standing hypotheses at
$(3,C_{k,\rm low},C_l^{\mathcal F})$, the gap hypothesis holds since $29\le C_{k,\rm low}$ and $16\le C_l^{\mathcal F}$
(\S\ref{i3:S-sss-components}).  For $x$ beyond all thresholds, $i\in I_x$ and $u\ne u'$: the size bound and the standing
hypotheses give the counting hypotheses; absolute convergence (C10) and the loop row give (RB), (RS) at
$\varepsilon_L$; $b\ge2$.  Theorem~\ref{i3:S-thm-pair}, applied with $C'=C_{k,\rm low}$ and $C_l=C_l^{\mathcal F}$, gives the
two inequalities.
\end{proof}

\begin{remark}\label{i3:S-rem-trust}
The proof rests on the following statements that are not proved in it: (SW), through the published Theorems~\ref{i3:thm-BHC} and \ref{i3:thm-SWT} (Proposition~\ref{i3:prop-SW}), and the class hypotheses (local tying, rank row, loop row); the
formally verified Propositions~\ref{i3:S-prop-instance}, \ref{i3:S-prop-NG}, \ref{i3:S-prop-leaves} (which includes the hypotheses of
Proposition~\ref{i3:S-prop-idcell}) and the bridge statements in the proof
of Lemma~\ref{i3:S-lem-rankrow}; the records of the setting recalled as (H1), (H2) in \S\ref{i3:S-ssec-pf-idcell}; the published
theorems (F1)--(F3) used for (FD), the Gelfand--Graev multiplicity one (Lemma~\ref{i3:S-lem-GGmult}) and Weyl's dimension formula
(used in Proposition~\ref{i3:S-prop-transport}(b)); and, through the components, Propositions~\ref{i3:prop-amp} and
\ref{i3:prop-bessel}.  The one-level inequality, the typed count, the Gram bound, (CP), Lemmas~\ref{i3:S-lem-link},
\ref{i3:S-lem-split} (from the definitions (D1)--(D7) of the construction), the level volumes (Lemma~\ref{i3:S-lem-LV},
Proposition~\ref{i3:S-prop-KppVol}), the $D_*$ identification (Proposition~\ref{i3:S-prop-dstar}), the identity cell
(Proposition~\ref{i3:S-prop-idcell}) and the transport (Proposition~\ref{i3:S-prop-transport}) are proved in
\S\S\ref{i3:S-ssec-onelevel}--\ref{i3:S-ssec-pf-transport}.  The largest single item is Proposition~\ref{i3:S-prop-leaves}.
Propositions~\ref{i3:S-prop-regions} and \ref{i3:S-prop-count} are not used in this proof.
\end{remark}

\subsubsection{Proof of the one-level inequality}\label{i3:S-ssec-onelevel}

\emph{Notation.}  A \emph{point} is a finite list $p=((s_1,t_1),\dots,(s_m,t_m))$ with $s_j\in\N$, $t_j\in E$; its
\emph{level} is $s_1+\dots+s_m$; it is a \emph{point of level $i$} if its level is $i$ and $s_j\ge1$, $t_j\ne0$ for all $j$.
$h\cdot(s,t)$ appends $(s,t)$ to $h$.  A \emph{tower} $T$ consists of finite sets $S_g(b,s)\subset E\setminus\{0\}$ ($b,s\ge1$);
$X_i$ is the (finite) set of points of level $i$ whose $j$-th entry satisfies $s_j\ge1$ and
$t_j\in S_g(s_1+\dots+s_{j-1}+1,s_j)$, and $X_{i,s}\subset X_i$ the set of those whose last entry has size $s$.  The kernel datum
has a finite word set $W_{\rm all}$ containing the empty word $\emptyset$ (count $0$, weight $1$); put $N_{\rm top}=4K(r_S+1)$,
$\mathrm{budget}(i)=4K(r_S+1-i)$, $W_b(n)=\{w\in W_{\rm all}:\mathrm{cnt}(w)\le n\}$.  For a finite set $Y$ of points put
$S(\epsilon,i,w,Y)=\Re\sum_{p,q\in Y}K_r(i)(p,q)g_{\epsilon,i,w}(p,q)$, so $Q(T,\epsilon,i,w)=S(\epsilon,i,w,X_i)/\mathrm{wt}(w)$.  For a kernel
$\kappa$ on points of level $i$ and $t,t'\in E$ put $\kappa_{t,t'}(h,h')=\kappa(h\cdot(s,t),h'\cdot(s,t'))$ and $\kappa_t=\kappa_{t,t}$; for a
kernel $\kappa$ on points of level $i-s$,
\[
N_{st}(\epsilon,i,s,w,Y,\kappa;t,t')=\mathrm{wt}(w)^{-1}\sum_{h,h'\in Y}\kappa(h,h')g_{\epsilon,i,w}(h\cdot(s,t),h'\cdot(s,t')),
\]
$N_0(\epsilon,k,w,Y,\kappa)=\mathrm{wt}(w)^{-1}\Re\sum_{h,h'\in Y}\kappa(h,h')g_{\epsilon,k,w}(h,h')$, and
$D_g(T,\epsilon,i,s,w)=\sum_{t\in S_g(i-s+1,s)}\Re N_{st}(\epsilon,i,s,w,X_{i-s},K_r(i)_t;t,t)$.  The datum also carries numbers
$\widehat N(\epsilon,i,w,Y,\kappa;l)\in\C$ ($l\in E$), size vectors $\rho^\flat_i,\rho_{b,s}\in\R^f$, functions $a_v:E\to\R$, and
$\mathrm{maj}(\rho,l)=\prod_{v=1}^f\min(1,(\rho_v/a_v(l))^K)$.  A \emph{probability vector} on a finite set $S$ is $\theta\ge0$ vanishing off
$S$ with $\sum\theta=1$.  A kernel $k$ is \emph{positive semidefinite on a finite set $Y$} ($k\succeq0$ on $Y$) if
$(k(y,y'))_{y,y'\in Y}$ is hermitian positive semidefinite.

\begin{definition}[Estimates, dominations]\label{i3:S-def-est}
Let $\epsilon>0$ and $Y$ a finite set of points of the stated level.
\begin{itemize}
\item[(GP), (KP), (WP)] For $\epsilon>0$ the Gram forms $g_{\epsilon,i,w}$ are $\succeq0$ on every finite set; the kernels $K_r(i)$
are $\succeq0$ on every finite set; $\mathrm{wt}(w)>0$.
\item[(Su)] For $1\le s\le i\le r_S+1$, points $h,h'$ of level $i-s$ and nonzero $t\ne t'$,
$K_r(i)(h\cdot(s,t),h'\cdot(s,t'))\ne0$ implies $(i,s)\ne(r_S+1,1)$ and $t^{-1}-t'^{-1}\in\Lambda'(i-s+1,s)$.
\item[(LS$_{e_p}$)] For $1\le i\le r_S+1$ and finite $L\subset\Lambda^\flat(i)\setminus\{0\}$, $\sum_{l\in L}\mathrm{maj}(\rho^\flat_i,l)\le e_p$; for
$1\le s\le i\le r_S+1$, $(i,s)\ne(r_S+1,1)$, and finite $L\subset\Lambda'(i-s+1,s)\setminus\{0\}$, $\sum_{l\in L}\mathrm{maj}(\rho_{i-s+1,s},l)\le e_p$.
\item[(Wa)] For $1\le i\le r_S+1$ and $w\in W_{\rm all}$ with $\mathrm{cnt}(w)+K\le N_{\rm top}$ there are probability vectors $\theta_A$ on
$W_b(\mathrm{cnt}(w)+K)$ ($A\subset\{1,\dots,f\}$) such that for all $Y$ of level $i-1$ and $\kappa\succeq0$ on $Y$:
$\widehat N(\epsilon,i,w,Y,\kappa;0)=N_0(\epsilon,i-1,w,Y,\kappa)$, and for every $l\in\Lambda^\flat(i)\setminus\{0\}$ there is $A$ with
$|\widehat N(\epsilon,i,w,Y,\kappa;l)|\le\mathrm{maj}(\rho^\flat_i,l)\sum_{w'}\theta_A(w')N_0(\epsilon,i-1,w',Y,\kappa)$.
\item[(Po)] For $1\le i\le r_S+1$, any $w$, $Y$ of level $i-1$, $\kappa\succeq0$ on $Y$, $t_0\in E$ and $B\in\R$ with
$\sum_{l\in L}|\widehat N(\epsilon,i,w,Y,\kappa;l)|\le B$ for every finite $L\subset\Lambda^\flat(i)\setminus\{0\}$: for every finite set $S$ of
nonzero $t\equiv t_0\bmod\Lambda(i,1)$, $\sum_{t\in S}\Re N_{st}(\epsilon,i,1,w,Y,\kappa;t,t)\le|\widehat N(\epsilon,i,w,Y,\kappa;0)|+B$.
\item[(Li)] For $1\le s\le i\le r_S+1$, $(i,s)\ne(r_S+1,1)$, and $w\in W_{\rm all}$ with $\mathrm{cnt}(w)+K\le N_{\rm top}$ there are
probability vectors $\theta_\pi$ on $W_b(\mathrm{cnt}(w)+K)^2$ ($\pi:\{1,\dots,f\}\to\{1,2,3\}$) such that for all $Y$ of level $i-s$,
finite $S\subset E$, $\kappa\succeq0$ on $\{h\cdot(s,t):h\in Y,t\in S\}$ and nonzero $t,t'\in S$ with $t^{-1}\ne t'^{-1}$ there is $\pi$
with $|N_{st}(\epsilon,i,s,w,Y,\kappa_{t,t'};t,t')|\le\mathrm{maj}(\rho_{i-s+1,s},t^{-1}-t'^{-1})\sum_{x}\theta_\pi(x)\frac12\bigl(\Re
N_{st}(\ldots,x_1,Y,\kappa_t;t,t)+\Re N_{st}(\ldots,x_2,Y,\kappa_{t'};t',t')\bigr)$.
\item[(Bl$_{C_4}$)] For $2\le s\le i\le r_S+1$ and $w\in W_{\rm all}$ with $\mathrm{cnt}(w)+(s+1)K\le N_{\rm top}$ and
$\mathrm{vol}(i-s+1,s)\le1$ there is a probability vector $\theta$ on $W_b(\mathrm{cnt}(w)+(s+1)K)$ such that for all $Y$ of level $i-s$,
$\kappa\succeq0$ on $Y$, $t_0\in E$ and finite sets $S$ of nonzero $t\equiv t_0\bmod\Lambda(i-s+1,s)$:
$\sum_{t\in S}\Re N_{st}(\epsilon,i,s,w,Y,\kappa;t,t)\le C_4^{fs}\prod_{s'=1}^{s-1}\mathrm{vol}(i-s+1,s')\sum_{w'}\theta(w')N_0(\epsilon,i-s,w',Y,\kappa)$.
\item[(TD)] A \emph{theta domination} is a family of kernels $\Theta_{i,s,t}$ such that for $1\le s\le i\le r_S+1$: $0\le e_i\le1$;
$\Theta_{i,s,t}-K_r(i)_t\succeq0$ on every finite set of points of level $i-s$, for $t\ne0$; $\Theta_{i,s,t}$ depends only on $t$ modulo
$\Lambda(i-s+1,s)$; and $(1+e_i)K_r(i-s)-\sum_{t\in S}\Theta_{i,s,t}\succeq0$ on every finite set of points of level $i-s$, for every
finite $S$ of pairwise incongruent elements modulo $\Lambda(i-s+1,s)$.  The theta kernels are also required to be
$\succeq0$ on finite sets (a clause of (TD), not of (KP)).
\item[(MD)] A \emph{minimal domination} for a tower $T$ is a family $\Theta_{i,s,t}$ such that for $1\le s\le i\le r_S+1$:
$\Theta_{i,s,t}-K_r(i)_t\succeq0$ on $X_{i-s}$ for $t\in S_g(i-s+1,s)$; $\Theta_{i,s,t}=\Theta_{i,s,t'}$ on $X_{i-s}^2$ for congruent
$t,t'\in S_g(i-s+1,s)$; and $(1+e_i)K_r(i-s)-\sum_{t\in S}\Theta_{i,s,t}\succeq0$ on $X_{i-s}$ for every $S\subset S_g(i-s+1,s)$ of pairwise
incongruent elements.
\end{itemize}
\end{definition}

\emph{Three elementary facts.}  (P1) If $k\succeq0$ on $Y'$ and $\phi:Y\to Y'$ is a map of finite sets, then
$k(\phi(\cdot),\phi(\cdot))\succeq0$ on $Y$ (it is $P^*AP$ with $P$ the $0/1$ matrix of $\phi$); in particular the zero extension of a
kernel $\succeq0$ on $X$ is $\succeq0$ on every finite set.  (P2) If $k_1,k_2\succeq0$ on $Y$, then so is $k_1k_2$ (Schur product
theorem), hence $\Re\sum_{y,y'}k_1(y,y')k_2(y,y')\ge0$.  (P3) If $M\succeq0$ on $Y$, $\|x\|_M=(\Re\sum M(y,y')x(y')\overline{x(y)})^{1/2}$ is a
seminorm on $\C^Y$.  Consequences: (N1) if $\kappa\succeq0$ on $Y$, then $\Re N_{st}(\ldots,\kappa;t,t)\ge0$ and $N_0(\ldots,\kappa)\ge0$ (by (GP),
(P1), (P2), (WP)); hence, $N_{st}$ being linear in $\kappa$, $\kappa_2-\kappa_1\succeq0$ on $Y$ implies
$\Re N_{st}(\ldots,\kappa_1;t,t)\le\Re N_{st}(\ldots,\kappa_2;t,t)$.  (N2) If $cK-\sum_{t\in R}\Theta_t\succeq0$ on $Y$, then
$\sum_{t\in R}N_0(\epsilon,k,w,Y,\Theta_t)\le c\,N_0(\epsilon,k,w,Y,K)$.

\begin{lemma}[Splitting by the last block]\label{i3:S-lem-st-split}
Under (GP), (KP), for every tower, $\epsilon>0$, $i\ge1$ and word $w$:
$\sqrt{S(\epsilon,i,w,X_i)}\le\sum_{s=1}^i\sqrt{S(\epsilon,i,w,X_{i,s})}$.
\end{lemma}

\begin{proof}
$M=K_r(i)g_{\epsilon,i,w}\succeq0$ on $X_i$ by (P2), and $\|\mathbf 1_A\|_M^2=S(\epsilon,i,w,A)$ for $A\subset X_i$.  A point of level $i\ge1$ is a
nonempty list whose last size $s$ satisfies $1\le s\le i$, so $\mathbf 1_{X_i}=\sum_{s=1}^i\mathbf 1_{X_{i,s}}$; apply (P3).
\end{proof}

\begin{lemma}[Schur step]\label{i3:S-lem-st-schur}
Assume (GP), (KP), (WP), (Su), (Li), (LS$_{e_p}$).  Let $1\le s\le i\le r_S+1$, $w\in W_{\rm all}$ with $\mathrm{cnt}(w)+K\le N_{\rm top}$, and
$M$ with $D_g(T,\epsilon,i,s,w')\le M$ for all $w'\in W_b(\mathrm{cnt}(w)+K)$.  Then
$S(\epsilon,i,w,X_{i,s})/\mathrm{wt}(w)\le D_g(T,\epsilon,i,s,w)+3^fe_pM$.
\end{lemma}

\begin{proof}
Put $Y=X_{i-s}$, $S=S_g(i-s+1,s)$, $\kappa=K_r(i)$.  \emph{Block decomposition.}  $h\cdot(s,t)\in X_i$ iff $h\in X_{\mathrm{level}(h)}$,
$\mathrm{level}(h)+s=i$, $s\ge1$, $t\in S_g(\mathrm{level}(h)+1,s)$; hence $X_{i,s}=\{h\cdot(s,t):h\in Y,t\in S\}$ with $(h,t)\mapsto h\cdot(s,t)$
injective, and
$S(\epsilon,i,w,X_{i,s})/\mathrm{wt}(w)=\sum_{t,t'\in S}\Re N_{st}(\epsilon,i,s,w,Y,\kappa_{t,t'};t,t')$.  Put
$a_x(t)=\Re N_{st}(\epsilon,i,s,x,Y,\kappa_t;t,t)\ge0$ (by (N1), as $\kappa_t\succeq0$ on $Y$ by (KP), (P1)); then $D_g(\cdot,x)=\sum_{t\in S}a_x(t)$,
$0\le D_g(\cdot,w)\le M$ (as $w\in W_b(\mathrm{cnt}(w)+K)$), and $e_p\ge0$ ((LS) with $L=\emptyset$).  \emph{Vanishing.}  If $t\ne t'$ and
$(i,s)=(r_S+1,1)$ or $t^{-1}-t'^{-1}\notin\Lambda'(i-s+1,s)$, then by (Su) $N_{st}(\ldots;t,t')=0$; in the case $(i,s)=(r_S+1,1)$ the
claim follows at once.  \emph{Off-diagonal kernel.}  Otherwise let $\rho=\rho_{i-s+1,s}$ and $m(t,t')=\mathrm{maj}(\rho,t^{-1}-t'^{-1})$ if
$t\ne t'$, $t^{-1}-t'^{-1}\in\Lambda'(i-s+1,s)$ and this value is $\ge0$, else $m(t,t')=0$.  For fixed $t$, $t'\mapsto t^{-1}-t'^{-1}$ is
injective, so (LS) gives $\sum_{t'}m(t,t')\le e_p$ and symmetrically $\sum_tm(t,t')\le e_p$.  With
$G(t,t')=\sum_\pi\sum_x\theta_\pi(x)\frac12(a_{x_1}(t)+a_{x_2}(t'))\ge0$, (Li) gives $\Re N_{st}(\ldots;t,t')\le m(t,t')G(t,t')$ for
$t\ne t'$ (if $\mathrm{maj}<0$ the left side is $\le0$).  \emph{Schur test.}  For each $x=(x_1,x_2)$,
$\sum_{t,t'}m(t,t')\frac12(a_{x_1}(t)+a_{x_2}(t'))\le\frac12e_pD_g(x_1)+\frac12e_pD_g(x_2)\le e_pM$; summing against $\theta_\pi$ and over
the $3^f$ maps $\pi$ bounds the off-diagonal part by $3^fe_pM$.
\end{proof}

\begin{lemma}[Diagonal step]\label{i3:S-lem-st-diag}
Assume (GP), (KP), (WP), (TD).  Let $1\le s\le i\le r_S+1$ and $R\subset E$ finite with every $t\in S_g(i-s+1,s)$ congruent modulo
$\Lambda(i-s+1,s)$ to some element of $R$.  Then
$D_g(T,\epsilon,i,s,w)\le\sum_{t_0\in R}\sum_{t\in S_g(i-s+1,s),\,t\equiv t_0}\Re N_{st}(\epsilon,i,s,w,X_{i-s},\Theta_{i,s,t_0};t,t)$.
\end{lemma}

\begin{proof}
All summands on the right are $\ge0$ by (N1).  For $t\in S_g(i-s+1,s)$ pick $t_0\in R$ with $t\equiv t_0$; by (TD) and (N1),
$\Re N_{st}(\ldots,K_r(i)_t;t,t)\le\Re N_{st}(\ldots,\Theta_{i,s,t};t,t)=\Re N_{st}(\ldots,\Theta_{i,s,t_0};t,t)$, a summand of the right side.
\end{proof}

\begin{lemma}[One-block step]\label{i3:S-lem-st-one}
Assume (GP), (KP), (WP), (TD), (Wa), (Po), (LS$_{e_p}$).  Let $1\le i\le r_S+1$, $w\in W_{\rm all}$ with $\mathrm{cnt}(w)+K\le N_{\rm top}$, $R$
finite and pairwise incongruent modulo $\Lambda(i,1)$, and $M$ with $N_0(\epsilon,i-1,w',X_{i-1},K_r(i-1))\le M$ for
$w'\in W_b(\mathrm{cnt}(w)+K)$.  Then the right side of Lemma~\ref{i3:S-lem-st-diag} for $s=1$ is at most
$(1+e_i)(N_0(\epsilon,i-1,w,X_{i-1},K_r(i-1))+2^fe_pM)$.
\end{lemma}

\begin{proof}
Put $Y=X_{i-1}$, $\Theta_{t_0}=\Theta_{i,1,t_0}$, $c=1+e_i$, and $G_A(t_0)=\sum_{w'}\theta_A(w')N_0(\epsilon,i-1,w',Y,\Theta_{t_0})\ge0$ with $\theta_A$
from (Wa).  For $t_0\in R$ and finite $L\subset\Lambda^\flat(i)\setminus\{0\}$, (Wa) and (LS) give
$\sum_{l\in L}|\widehat N(\ldots,\Theta_{t_0};l)|\le e_p\sum_AG_A(t_0)$ (only $l$ with $\mathrm{maj}\ge0$ contribute positively), so (Po) and
$\widehat N(\ldots;0)=N_0(\epsilon,i-1,w,Y,\Theta_{t_0})\ge0$ give
$\sum_{t\equiv t_0}\Re N_{st}(\epsilon,i,1,w,Y,\Theta_{t_0};t,t)\le N_0(\epsilon,i-1,w,Y,\Theta_{t_0})+e_p\sum_AG_A(t_0)$.  Sum over $t_0\in R$: by
(TD) and (N2), $\sum_{t_0}N_0(\ldots,w,Y,\Theta_{t_0})\le cN_0(\ldots,w,Y,K_r(i-1))$ and $\sum_{t_0}G_A(t_0)\le\sum_{w'}\theta_A(w')cM=cM$; there are $2^f$
sets $A$.
\end{proof}

\begin{lemma}[Multi-block step]\label{i3:S-lem-st-multi}
Assume (GP), (KP), (WP), (TD), (Bl$_{C_4}$) with $C_4\ge1$, and $\mathrm{vol}\ge0$.  Let $2\le s\le i\le r_S+1$, $w\in W_{\rm all}$ with
$\mathrm{cnt}(w)+(s+1)K\le N_{\rm top}$, $\mathrm{vol}(i-s+1,s)\le1$, $R$ pairwise incongruent modulo $\Lambda(i-s+1,s)$, and $M$ with
$N_0(\epsilon,i-s,w',X_{i-s},K_r(i-s))\le M$ for $w'\in W_b(\mathrm{cnt}(w)+(s+1)K)$.  Then the right side of
Lemma~\ref{i3:S-lem-st-diag} is at most $C_4^{fs}\prod_{s'=1}^{s-1}\mathrm{vol}(i-s+1,s')(1+e_i)M$.
\end{lemma}

\begin{proof}
Let $c_0=C_4^{fs}\prod_{s'<s}\mathrm{vol}(i-s+1,s')\ge0$ and $\theta$ from (Bl).  For each $t_0\in R$, (Bl) with $\kappa=\Theta_{i,s,t_0}$ bounds the
inner sum by $c_0\sum_{w'}\theta(w')N_0(\epsilon,i-s,w',X_{i-s},\Theta_{i,s,t_0})$.  Summing over $t_0$ and using (TD), (N2):
$\sum_{t_0}N_0(\ldots,\Theta_{i,s,t_0})\le(1+e_i)N_0(\ldots,K_r(i-s))\le(1+e_i)M$; $\theta$ is a probability vector.
\end{proof}

\begin{lemma}[Restricted datum]\label{i3:S-lem-Dp}
Let the kernels $K_r(i)$ be $\succeq0$ on finite sets, $0\le e_i\le1$, and $\Theta$ a minimal domination for $T$.  For each
$(b,s)$ choose $\sigma_{b,s}$ assigning to each class $c\in E/\Lambda(b,s)$ meeting $S_g(b,s)$ an element of $S_g(b,s)\cap c$.  Let $D_p$
be $D$ with $K_r(i)$ replaced by its restriction to $X_i\times X_i$ extended by zero, and
$\Theta^p_{i,s,t}(h,h')=\Theta_{i,s,\sigma(t)}(h,h')$ if $1\le s\le i\le r_S+1$, the class $[t]$ modulo $\Lambda(i-s+1,s)$ meets $S_g(i-s+1,s)$ and
$h,h'\in X_{i-s}$, and $0$ otherwise ($\sigma(t)=\sigma_{i-s+1,s}([t])$).  Then the kernels of $D_p$ are $\succeq0$ on finite sets and $D_p$
satisfies (TD).
\end{lemma}

\begin{proof}
Fix $1\le s\le i\le r_S+1$, $\Lambda=\Lambda(i-s+1,s)$, $S_g=S_g(i-s+1,s)$, $X'=X_{i-s}$.  (a) By the block decomposition, for $t\in S_g$,
$h\cdot(s,t)\in X_i$ iff $h\in X'$; for $t\notin S_g$, never.  (b) $K^p_r(i)\succeq0$ by (P1).  For $t_1\in S_g$,
$\Theta_{i,s,t_1}=(\Theta_{i,s,t_1}-K_r(i)_{t_1})+K_r(i)_{t_1}\succeq0$ on $X'$, so $\Theta^p$ is a zero extension of a kernel $\succeq0$ on
$X'$, hence $\succeq0$ on finite sets.  (c) \emph{Majorisation:} for $t\in S_g$, $\sigma(t)\equiv t$ gives $\Theta_{i,s,\sigma(t)}=\Theta_{i,s,t}$ on $X'^2$,
and by (a) $(K^p_r(i))_t$ is the zero extension of $K_r(i)_t|_{X'}$; so $\Theta^p_{i,s,t}-(K^p_r(i))_t$ is the zero extension of a kernel
$\succeq0$ on $X'$.  For $t\notin S_g$, $(K^p_r(i))_t=0$ and $\Theta^p\succeq0$.  \emph{Periodicity} holds by construction.
\emph{Domination:} for $S$ pairwise incongruent, let $S_0\subset S$ be the elements whose class meets $S_g$ and $S'=\sigma(S_0)$; $\sigma$ is
injective on $S_0$ and $S'$ is pairwise incongruent.  On $X'^2$, $\sum_{t\in S}\Theta^p_{i,s,t}=\sum_{t_1\in S'}\Theta_{i,s,t_1}$ and
$K^p_r(i-s)=K_r(i-s)$; off $X'^2$ both vanish.  So $(1+e_i)K^p_r(i-s)-\sum_S\Theta^p$ is the zero extension of a kernel $\succeq0$ on $X'$
by (MD).
\end{proof}

\begin{proof}[Proof of Proposition~\ref{i3:S-prop-onelevel}]
\emph{Reduction.}  Let $\Theta$ be a minimal domination for $T$ and $D_p$ as in Lemma~\ref{i3:S-lem-Dp}.  The quantities
$Q(T,\epsilon,k,w)$ and $D_g$ only involve entries indexed by $X_k\times X_k$, so they agree for $D$ and $D_p$; $D_p$ satisfies (KP)
and (TD); (Su) for $D_p$ follows from (Su) for $D$ (nonzero entries of $K^p_r$ are entries of $K_r$); the other hypotheses do
not involve the kernels.  So we may assume (TD).

Fix $w$ with $\mathrm{cnt}(w)\le\mathrm{budget}(i)$.  (a) For $1\le s\le i$ and $M$ with $D_g(T,\epsilon,i,s,w')\le M$ for
$w'\in W_b(\mathrm{cnt}(w)+K)$: $\mathrm{cnt}(w)+K\le N_{\rm top}$, so Lemma~\ref{i3:S-lem-st-schur} gives
$S(\epsilon,i,w,X_{i,s})/\mathrm{wt}(w)\le(1+3^fe_p)M$.  (b) $s=1$: let $R$ consist of one representative of each class of
$S_g(i,1)$ modulo $\Lambda(i,1)$.  For $\mathrm{cnt}(w')\le\mathrm{cnt}(w)+K$, every $w''\in W_b(\mathrm{cnt}(w')+K)$ has
$\mathrm{cnt}(w'')\le\mathrm{budget}(i)+2K\le\mathrm{budget}(i-1)$, so $N_0(\epsilon,i-1,w'',X_{i-1},K_r(i-1))=Q(T,\epsilon,i-1,w'')\le\sigma_{i-1}^2$.
Lemmas~\ref{i3:S-lem-st-diag}, \ref{i3:S-lem-st-one} with $M=\sigma_{i-1}^2$ give $D_g(T,\epsilon,i,1,w')\le(1+e_i)(1+2^fe_p)\sigma_{i-1}^2$; by (a),
$S(\epsilon,i,w,X_{i,1})/\mathrm{wt}(w)\le(1+3^fe_p)(1+2^fe_p)(1+e_i)\sigma_{i-1}^2$.  (c) $2\le s\le i$: with $R$ chosen likewise for
$\Lambda(i-s+1,s)$, for $\mathrm{cnt}(w')\le\mathrm{cnt}(w)+K$ and $w''\in W_b(\mathrm{cnt}(w')+(s+1)K)$,
$\mathrm{cnt}(w'')\le\mathrm{budget}(i)+(s+2)K\le\mathrm{budget}(i-s)$, so the $N_0$ at level $i-s$ is $\le\sigma_{i-s}^2$; also
$\mathrm{cnt}(w')+(s+1)K\le N_{\rm top}$ and $\mathrm{vol}(i-s+1,s)\le1$.  Lemmas~\ref{i3:S-lem-st-diag}, \ref{i3:S-lem-st-multi} and (a) give
$S(\epsilon,i,w,X_{i,s})/\mathrm{wt}(w)\le(1+3^fe_p)(1+e_i)C_4^{fs}\prod_{s'<s}\mathrm{vol}(i-s+1,s')\sigma_{i-s}^2$.  (d) By
Lemma~\ref{i3:S-lem-st-split} and $\mathrm{wt}(w)>0$, $\sqrt{Q}\le\sum_s\sqrt{S(\epsilon,i,w,X_{i,s})/\mathrm{wt}(w)}$; insert (b), (c) and use
$\sqrt{x\sigma^2}=\sqrt x\,\sigma$ for $\sigma\ge0$.
\end{proof}

\subsubsection{The lemmas (FD) and rank-row, the typed count, the Gram bound and (CP)}\label{i3:S-ssec-leafproofs}

\medskip\noindent\emph{Proof of Lemma~\ref{i3:S-lem-FD}.}
We use three published facts.
\begin{itemize}
\item[(F1)] Schur orthogonality (Theorem~\ref{i3:thm-godement}, in the convention fixed after it).
\item[(F2)] \emph{Finite class number} (Borel, \emph{Some finiteness properties of adele groups over number fields},
Publ.\ Math.\ IHES 16 (1963), 5--30, Thm.~5.1): for an open compact $K'\subset U(\mathbb A_f)$, $U(F)\backslash U(\mathbb A_f)/K'$ is finite (Thm.~5.1 is stated for $\prod_vU(\mathcal O_v)$; for any open compact $K'$ pass through $K'\cap\prod_vU(\mathcal O_v)$, of finite index in both).
\item[(F3)] \emph{Finite multiplicities in the discrete spectrum} (Borel--Garland, \emph{Laplacian and the discrete spectrum of
an arithmetic group}, Amer.\ J.\ Math.\ 105 (1983), 309--335, Thm.~2, stated for semisimple $G$; for semisimple groups Harish-Chandra, \emph{Automorphic forms on
semisimple Lie groups}, Lecture Notes in Math.\ 62 (1968), Thm.~1 and its corollary; reference to be confirmed; the reduction to reductive $G$ with anisotropic centre is standard but not stated there): let
$G$ be the group of real points of a reductive group over $\Q$ (here the restriction of scalars of $U$, so $G=G_\infty$) and
$\Gamma\subset G$ an arithmetic subgroup with $\mathrm{vol}(\Gamma\backslash G)<\infty$.  Then $L^2_{\rm disc}(\Gamma\backslash G)$,
the closed span of the irreducible closed $G$-invariant subspaces of $L^2(\Gamma\backslash G)$, is a Hilbert direct sum of
irreducible representations in which each irreducible unitary representation occurs with finite multiplicity.
\end{itemize}
Here $\mathrm{vol}(\Gamma\backslash G_\infty)<\infty$ for the groups $\Gamma_j$ below, because the centre of $U$ is the
anisotropic torus $U(1)_{E/F}$ (Borel--Harish-Chandra).

\emph{Step 1: $\sigma$ and the kernel $\phi$.}  Since $|\langle w_v,\sigma_v(g)w_v\rangle|\le1$, integrability of $\phi_v$
gives square integrability, so each $\sigma_v$ is square integrable; by (F1) its formal degree is unique, so the given $d_v$ is
it.  Put $\sigma=\bigotimes_v\sigma_v$ (Hilbert tensor product) and $w=\bigotimes_vw_v$, a unit vector.  The representation
$\sigma$ of $G_\infty=\prod_vU(F_v)$ is irreducible: its commutant is the tensor product of the commutants of the $\sigma_v$
(commutation theorem for tensor products of von Neumann algebras), i.e.\ $\C$.  Its coefficient
$\langle w,\sigma(g)w\rangle=\prod_v\langle w_v,\sigma_v(g_v)w_v\rangle$ is square integrable by Fubini, so (F1) applies to
$\sigma$; evaluating both sides of (F1) at $u_i=v_i=w$ and using Fubini gives $d_\sigma=d:=\prod_vd_v$.  Thus
$\phi(g)=d\,\overline{\langle w,\sigma(g)w\rangle}$, and $\phi$ is integrable.

\emph{Step 2: convolution identities.}  For $F\in L^1(G_\infty)$ write $\pi(F)=\int F(g)\pi(g)\,dg$ for a unitary
representation $\pi$, $F^*(g)=\overline{F(g^{-1})}$ (so $\pi(F^*)=\pi(F)^*$; $G_\infty$ is unimodular) and
$(\lambda_xF)(g)=F(x^{-1}g)$ (so $\pi(x)\pi(F)=\pi(\lambda_xF)$).  Note $\phi(g)=d\langle\sigma(g)w,w\rangle$.
\begin{itemize}
\item[(a)] $\sigma(\phi)v=\langle w,v\rangle w$ for all $v$ (the remark after Theorem~\ref{i3:thm-godement}); in particular
$\sigma(\phi)w=w$.
\item[(b)] $\phi^*=\phi$, since $\overline{\phi(g^{-1})}=d\langle w,\sigma(g^{-1})w\rangle=d\langle\sigma(g)w,w\rangle$.
\item[(c)] $\phi*F*\phi=\langle w,\sigma(F)w\rangle\,\phi$ for $F\in L^1(G_\infty)$; with $F$ an approximate identity this
gives $\phi*\phi=\phi$.  Indeed $(F*\phi)(g)=d\int F(h)\langle\sigma(h^{-1}g)w,w\rangle dh=d\langle\sigma(g)w,\sigma(F)w\rangle$, and
\[
(\phi*F*\phi)(g)=d^2\!\int\langle\sigma(h)w,w\rangle\langle\sigma(g)w,\sigma(h)\sigma(F)w\rangle\,dh
=d\,\langle w,\sigma(F)w\rangle\langle\sigma(g)w,w\rangle
\]
by (F1) with $(v_1,u_1,v_2,u_2)=(\sigma(g)w,\sigma(F)w,w,w)$.
\end{itemize}
By (b) and (c), for every unitary representation $\pi$ of $G_\infty$, $\pi(\phi)$ is a self-adjoint idempotent, i.e.\ an
orthogonal projection.

\emph{Step 3: $\mathrm{im}\,\pi(\phi)=\{Aw:A\in\mathrm{Hom}_{G_\infty}(\sigma,\pi)\}$, and $A\mapsto Aw$ is injective.}  If $A$
intertwines, $\pi(\phi)Aw=A\sigma(\phi)w=Aw$ by (a).  Conversely let $\pi(\phi)h=h$.  For $F,F'\in L^1$, by (c),
\[
\langle\pi(F')h,\pi(F)h\rangle=\langle h,\pi(\phi)\pi(F'^**F)\pi(\phi)h\rangle=\langle w,\sigma(F'^**F)w\rangle\|h\|^2
=\langle\sigma(F')w,\sigma(F)w\rangle\|h\|^2 .
\]
Hence $A_0:\sigma(F)w\mapsto\pi(F)h$ is well defined, linear and $\|h\|$ times an isometry on
$\mathcal D_w=\{\sigma(F)w:F\in L^1\}$.  $\mathcal D_w$ is $\sigma$-invariant ($\sigma(x)\sigma(F)w=\sigma(\lambda_xF)w$) and contains
$w=\sigma(\phi)w\neq0$, so its closure is $\mathcal H_\sigma$ by irreducibility.  The continuous extension $A$ of $A_0$ satisfies
$A\sigma(x)\sigma(F)w=\pi(\lambda_xF)h=\pi(x)A\sigma(F)w$, so $A\in\mathrm{Hom}_{G_\infty}(\sigma,\pi)$, and $Aw=A\sigma(\phi)w=\pi(\phi)h=h$.
If $Aw=0$ then $A\sigma(x)w=\pi(x)Aw=0$ for all $x$, and the span of $\sigma(G_\infty)w$ is dense, so $A=0$.  Therefore
$\dim\mathrm{im}\,\pi(\phi)=\dim\mathrm{Hom}_{G_\infty}(\sigma,\pi)$ (both possibly infinite).

\emph{Step 4: reduction to finitely many arithmetic quotients.}  $e$ is locally constant and vanishes outside the compact set
$K_f$, hence is compactly supported and bi-invariant under some open compact subgroup $K'$.  Then $\mathcal R(e)$ maps $L^2$
into the $K'$-invariant subspace $(L^2)^{K'}$, and commutes with $T=\mathcal R(\phi)$ (they act through the commuting factors
$U(\mathbb A_f)$, $G_\infty$), so $\mathcal R(e)T=T\mathcal R(e)$ and $\mathrm{im}(\mathcal R(e)T)\subset T((L^2)^{K'})$.  By (F2) write
$U(\mathbb A_f)=\bigsqcup_{j=1}^mU(F)x_jK'$; then $(L^2)^{K'}\cong\bigoplus_jL^2(\Gamma_j\backslash G_\infty)$ as
$G_\infty$-representations, with $\Gamma_j=U(F)\cap x_jK'x_j^{-1}$ arithmetic.  $T$ preserves $(L^2)^{K'}$ and restricts there to
$\pi(\phi)$ for $\pi=\bigoplus_jL^2(\Gamma_j\backslash G_\infty)$, so by Step~3,
$\dim T((L^2)^{K'})=\sum_j\dim\mathrm{Hom}_{G_\infty}(\sigma,L^2(\Gamma_j\backslash G_\infty))$.  Finally, for each $j$, a nonzero
$A\in\mathrm{Hom}(\sigma,L^2(\Gamma_j\backslash G_\infty))$ is, by Schur's lemma ($A^*A$ is a scalar), a multiple of an isometry, so
its image is a closed irreducible invariant subspace and lies in $L^2_{\rm disc}$.  Hence
$\mathrm{Hom}(\sigma,L^2(\Gamma_j\backslash G_\infty))=\mathrm{Hom}(\sigma,L^2_{\rm disc}(\Gamma_j\backslash G_\infty))$, whose dimension is the
multiplicity of $\sigma$ in $L^2_{\rm disc}$, finite by (F3).  \qed

\noindent The formal proof does not use (F3): it proves the needed finiteness from reduction theory (Siegel sets, cuspidality
at infinity, discrete-series bounds).  We use the published route; for (F2) and (F3) the exact statements are references to be
confirmed.

\medskip\noindent\emph{Proof of Lemma~\ref{i3:S-lem-rankrow}.}
\emph{Transfer of closedness.}  If $P$ is closed on $\mathcal C$ and there is a threshold function $\vartheta$ on the outer
constants such that every family satisfying $P$ and $C_l^{\mathcal F}\ge\vartheta$ satisfies $Q$, then $Q$ is closed on $\mathcal C$: keep
$\varepsilon_1$ and enlarge the threshold system by $\vartheta$.  Apply this to $P=\mathrm{Ranks}$ with
$\vartheta=40$, $40$, $1$ for (i), (ii), (iii), using the following implications, valid for every family $\mathcal F$

(i) If $40\le C_l^{\mathcal F}$ and $\mathrm{Ranks}$ holds, then $\mathrm{TFrame}$ holds; (ii) if moreover $1\le C_0$, then every member
satisfying the standing hypotheses satisfies $\mathrm{Sizes3b}(a(C_0))$, and $1\le a(C_0)$ since $e^{3/4}\ge1$ and $C_0e\ge1$;
(iii) if $1\le C_l^{\mathcal F}$ and $\mathrm{Ranks}$ holds, then (using $L\ge1$, so $1\le C_lL$) for $x$ beyond the threshold of
$\mathrm{Ranks}$ every member satisfying the standing hypotheses and the ties satisfies $\mathrm{SzRows}$: the mid-volume
component of $\mathrm{SzRows}$ is supplied by $\mathrm{Ranks}$ itself, the other three components by the bridge statements below.  These three implications are the formally verified part:
\begin{quote}
\textbf{Formally verified (bridge statements).}  (i) follows from the level-frame lemma with the critical level lower bound
(at $\epsilon_0=\frac13$, $C_V=0$) and the level formula for the family; (ii) from the critical level ratio and its bound;
(iii) from the Siegel mass ratio for unramified CM fields and the degree-ratio formula; the mid-volume hypothesis of the third is
taken from Ranks.  The statements are made under two instance hypotheses: that $d_\sigma(\tau)$ equals Harish-Chandra's formal degree
under the standing hypotheses, and the representation of totally positive elements by binary hermitian forms over unramified CM fields;
both are discharged in \cite{CharacterFamily}.
\end{quote}
\noindent The structure of the formal proofs of the three bridge implications is described in \S\ref{i3:str-bridge}.

\medskip\noindent\emph{Proof of Proposition~\ref{i3:S-prop-typed}.}
Write $R_v=R\otimes_{\widehat{\mathcal O}_F}\mathcal O_v$ for a prime $v$ of $F$, $\sigma$ the involution.

\begin{lemma}[Completely split primes]\label{i3:S-lem-splitprime}
Let $F\subset E$ be number fields and $p$ a prime with $[E:\Q]$ primes of $E$ above it.  Then $p$ has exactly $[F:\Q]$ primes
$v$ above it in $F$, and $|\mathcal O_F/v|=p$, $v(p)=1$ for each.
\end{lemma}

\begin{proof}
$\sum_{w\mid p}e(w|p)f(w|p)=[E:\Q]$ with $[E:\Q]$ summands forces $e=f=1$ above $p$ in $E$; by multiplicativity in $\Q\subset F\subset E$,
$e(v|p)=f(v|p)=1$.  Each $v$ has at most $[E:F]$ primes of $E$ above it, so $\#\{v\mid p\}\ge[F:\Q]$, and $\#\{v\mid p\}\le\sum e(v|p)f(v|p)=[F:\Q]$.
\end{proof}

\begin{lemma}[Local reduction]\label{i3:S-lem-localred}
Let $v$ be a prime of $F$, $A=E\otimes_FF_v$ (dimension $2$ over $F_v$) with the involution, $R_v\subset A$ a finitely generated
$\mathcal O_v$-submodule ring spanning $A$, stable under $\sigma$, and $\epsilon\in R_v$ with $\epsilon^2=\epsilon$, $\epsilon\sigma(\epsilon)=0$, $1-\epsilon-\sigma(\epsilon)=ct$ with
$|c|_v<1$, $t\in R_v$.  Then there is a ring homomorphism $\rho:A\to F_v$ with $\rho(\epsilon)=1$, $\rho(1\otimes a)=a$, $\rho(R_v)\subset\mathcal O_v$ and
$\epsilon x=\rho(x)\epsilon$ for all $x$.
\end{lemma}

\begin{proof}
$\epsilon'=\sigma(\epsilon)$ is idempotent with $\epsilon\epsilon'=0$, so $f=1-\epsilon-\epsilon'$ is idempotent; $f=f^n=c^nt^n\to0$ ($R_v$ is bounded), so $\epsilon+\epsilon'=1$.
Then $\epsilon,\epsilon'\ne0$ and they are independent, hence a basis of $A$; let $\rho(x)$ be the $\epsilon$-coordinate.  Multiplying
$x=\rho(x)\epsilon+\rho'(x)\epsilon'$ by $\epsilon$ gives $\epsilon x=\rho(x)\epsilon$, whence $\rho$ is a ring homomorphism with the stated values.  For $x\in R_v$
the characteristic polynomial of multiplication by $x$ has coefficients in $\mathcal O_v$ (it preserves the lattice $R_v$), and
$\rho(x)$ is a root of it ($x\epsilon=\rho(x)\epsilon$, $\epsilon\ne0$); so $\rho(x)$ is integral, hence in $\mathcal O_v$.
\end{proof}

\begin{lemma}\label{i3:S-lem-onesub}
Let $L$ be an $R$-lattice with hermitian form $h$ such that $h(z,w)\in gR$ for $z\in L$, $w\in gL$, let $z_0\in L$ with
$h(z_0,z_0)+1\in gR$, and let $(1-e-\bar e)y\in gL$ for all $y\in L$.  Then $1-e-\bar e\in gR$.
\end{lemma}

\begin{proof}
With $c=1-e-\bar e$: $c\,h(z_0,z_0)=h(z_0,cz_0)\in gR$ and $c(h(z_0,z_0)+1)\in gR$, so $c\in gR$.
\end{proof}

\begin{proof}[Proof of Proposition~\ref{i3:S-prop-typed}]
\emph{Step 1.}  For any typed reduction, $\prod_{v\in P}|\mathbb F_v|=|R/K|$ (surjective joint map with kernel $K$).
\emph{Step 2: a concrete typed reduction.}  Let $P_0$ be the set of primes $v$ of $F$ dividing $g$.  By
Lemma~\ref{i3:S-lem-splitprime}, $k_v=\mathcal O_F/v$ has $g_i$ elements for $v\mid g_i$, there are $f$ such $v$ for each $i$, and $v(g)=1$
(no $v$ divides both).  For $v\in P_0$ the image $\epsilon_v$ of $e$ satisfies $1-\epsilon_v-\sigma(\epsilon_v)=gt_v$, so
Lemma~\ref{i3:S-lem-localred} gives $\rho_v$; put $\mathrm{res}_v(x)=\rho_v(x_v)\bmod\mathfrak m_v$.  The joint map is surjective (lift $b_v$ to
$o_v\in\mathcal O_v$ and take $x$ with component $1\otimes o_v$ at $v$ and $0$ at the other places of $P_0$).  If $ex=gy$, $y\in R$, then
$\rho_v(x_v)=\rho_v(\epsilon_vx_v)=g\rho_v(y_v)\in\mathfrak m_v$.  Conversely, if all $\rho_v(x_v)\in\mathfrak m_v$, put $y=ex$: at $w\nmid g$, $g$ is a unit;
at $v\in P_0$, $y_v=\rho_v(x_v)\epsilon_v$ and $g^{-1}y_v\in R_v$.  So the kernel is $K$ and $|R/K|=g_1^fg_2^f$.  \emph{Step 3.}  For any
typed reduction with $n_i=\#\{v:|\mathbb F_v|=g_i\}$, $g_1^{n_1}g_2^{n_2}=g_1^fg_2^f$, and unique factorisation gives $n_1=n_2=f$.  The
product formula follows by grouping factors; $E$, totally complex of degree $2f$, has $f$ infinite places.
\end{proof}

\medskip\noindent\emph{Proof of Proposition~\ref{i3:S-prop-GG}.}
Let $\mathbb F$ be finite with $q$ elements, $G_k=\mathrm{GL}_k(\mathbb F)$, $U_k$ upper unitriangular ($|U_k|=q^{k(k-1)/2}$), $\psi$ an additive
character, $\gamma_\psi(x)=|U_k|^{-1}\sum_{y\in G_k}\psi^0_U(yxy^{-1})$, and for a representation $\pi$ of $G_k$ on $V$,
$R_\psi(\pi)=\sum_x\gamma_\psi(x)\pi(x)$.  A \emph{Whittaker functional} is $\ell\in V^*$ with $\ell(\pi(u)v)=\psi_U(u)\ell(v)$; $\pi$ is
\emph{generic} if one is nonzero.

\begin{lemma}\label{i3:S-lem-GGelem}
(a) $\psi_U$ is a character of $U_k$ and $\gamma_\psi(1)=|G_k|/|U_k|>0$.  (b) $\gamma_\psi(gxg^{-1})=\gamma_\psi(x)$, $\gamma_\psi(x^{-1})=\overline{\gamma_\psi(x)}$.
(c) $R_\psi(\pi)$ commutes with $\pi(G_k)$.
\end{lemma}

\begin{proof}
(a) $(uu')_{i,i+1}=u_{i,i+1}+u'_{i,i+1}$.  (b) Substitute $y\mapsto yg^{-1}$; and $\psi^0_U(z^{-1})=\overline{\psi^0_U(z)}$ (both $0$ off $U_k$).
(c) By (b), $\pi(g)R_\psi\pi(g)^{-1}=\sum_x\gamma_\psi(g^{-1}xg)\pi(x)=R_\psi$.
\end{proof}

\begin{lemma}\label{i3:S-lem-GGnum}
For real $q\ge3$ and $m\ge0$, $\prod_{i=1}^m(1-q^{-i})^{-1}\le1+3/q$.
\end{lemma}

\begin{proof}
With $x=1/q\le\frac13$, induction gives $p_m=\prod_{i\le m}(1-x^i)\ge1-\frac32x+\frac32x^{m+1}$: with $y=x^{m+1}$,
$(1-\frac32x+\frac32y)(1-y)-(1-\frac32x+\frac32xy)=\frac12y(1-3y)\ge0$.  So $p_m\ge1-\frac32x$ and
$(1-\frac32x)(1+3x)=1+\frac32x(1-3x)\ge1$.
\end{proof}

\begin{lemma}[Gelfand--Graev; cited]\label{i3:S-lem-GGmult}
If $\psi\ne1$ and $\pi$ is irreducible finite-dimensional, the Whittaker functionals on $\pi$ form a space of dimension $\le1$.
\end{lemma}

\noindent This is multiplicity one of the Gelfand--Graev representation $\mathrm{Ind}_{U_k}^{G_k}\psi_U$ (Frobenius reciprocity identifies
the space with $\mathrm{Hom}_{G_k}(\pi,\mathrm{Ind}\,\psi_U)$): I.~M.~Gelfand, M.~I.~Graev, Dokl.\ Akad.\ Nauk SSSR 147 (1962), 529--532; A.~V.~Zelevinsky,
LNM 869 (1981) (reference to be confirmed).  It is formally verified in \cite{CharacterFamily}.

\begin{lemma}\label{i3:S-lem-GGmir}
If $\psi\ne1$ and $\pi$ has a nonzero Whittaker functional $\ell$, then $\dim V\ge D_k=\prod_{i=1}^{k-1}(q^i-1)$.
\end{lemma}

\begin{proof}
We show by induction on $m\ge1$: for every representation of $G_m$ with nonzero Whittaker functional $\ell$ there are
$p_1,\dots,p_{D_m}\in G_m$ with last row $e_m$ such that the $\ell\circ\pi(p_a)$ are independent ($k=0$ is trivial).  $m=1$: $p_1=1$.  Step
$m\to m+1$: let $E(A)=\mathrm{diag}(A,1)$ and $c(x)=\bigl(\begin{smallmatrix}1&x\\0&1\end{smallmatrix}\bigr)$, $x\in\mathbb F^m$.  Then
$\psi_U(E(u))=\psi_U(u)$, $\psi_U(c(x))=\psi(x_m)$, and $E(A)c(x)=c(Ax)E(A)$.  By induction for $\pi\circ E$ get $p_a$; for each
$r\in\mathbb F^m\setminus\{0\}$ choose $g_r\in G_m$ with last row $r$, and put $f_{r,a}=\ell\circ\pi(E(p_ag_r))$ (last row of $p_ag_r$ is $r$).  Then
$f_{r,a}(\pi(c(x))v)=\psi(r\cdot x)f_{r,a}(v)$; for fixed $r$ the $f_{r,a}$ are independent (precomposition with an invertible operator);
and the characters $x\mapsto\psi(r\cdot x)$ are distinct for distinct $r$ (as $\psi\ne1$), so averaging against
$\psi(-s\cdot x)$ separates the $r$'s.  This gives $(q^m-1)D_m=D_{m+1}$ independent functionals, and $E(p_ag_r)$ has last row $e_{m+1}$.
\end{proof}

\begin{lemma}\label{i3:S-lem-GGconv}
If $\psi\ne1$ and $\pi$ is irreducible on $V$, then $R_\psi(\pi)=(|G_k|/\dim V)\mathrm{id}$ if $\pi$ is generic and $0$ otherwise.
\end{lemma}

\begin{proof}
$Q=|U_k|^{-1}\sum_u\psi_U(u)\pi(u)$ satisfies $Q\pi(u_0)=\psi_U(u_0)^{-1}Q$, so $Q^2=Q$ and $\mathrm{tr}\,Q=\mathrm{rank}\,Q$; and
$\sum_y\pi(y)^{-1}Q\pi(y)=R_\psi(\pi)$, so $\mathrm{tr}R_\psi=|G_k|\mathrm{rank}\,Q$.  With $s=\mathrm{diag}(1,-1,1,\dots)$, $\psi_U(sus)=\psi_U(u)^{-1}$; so
$f\mapsto f\circ Q\pi(s)$ maps $V^*$ into the Whittaker functionals, and $\ell\circ\pi(s)\circ Q=\ell\circ\pi(s)$ for Whittaker $\ell$.  By
Lemma~\ref{i3:S-lem-GGmult}, $\mathrm{rank}\,Q\le1$, with equality iff $\pi$ is generic.  By Lemma~\ref{i3:S-lem-GGelem}(c) and Schur's
lemma $R_\psi=c\,\mathrm{id}$, and taking traces $c\dim V=|G_k|\mathrm{rank}\,Q$.
\end{proof}

\begin{lemma}\label{i3:S-lem-GGplace}
If $\psi\ne1$, $q\ge3$, and $R_\psi(\pi)w=\mu w$ with $w\ne0$, then $\mu=\gamma_\psi(1)c$ with $0\le c\le1+3/q$.
\end{lemma}

\begin{proof}
$\ker(R_\psi-\mu)$ is invariant (Lemma~\ref{i3:S-lem-GGelem}(c)); take an irreducible $W$ in it.  By Lemma~\ref{i3:S-lem-GGconv}, $c=0$ or
$c=|U_k|/\dim W\le\prod_{i<k}q^i/(q^i-1)\le1+3/q$ by Lemmas~\ref{i3:S-lem-GGmir}, \ref{i3:S-lem-GGnum}.
\end{proof}

\begin{proof}[Proof of Proposition~\ref{i3:S-prop-GG}]
$\gamma(1)>0$ by Lemma~\ref{i3:S-lem-GGelem}(a).  Let $\pi$ act on $\C^X$ by $(\pi(g)a)(i)=a(i\cdot g)$ and $T=\sum_x\gamma(x)\pi(x)$.  Then
$T_{ij}=\sum_{i\cdot x=j}\gamma(x)$; by freeness this is $\gamma(d(i,j)^{-1})=\overline{\gamma(d(i,j))}=\mathbf G_{ij}$ if $i,j$ share an orbit, $0$
otherwise.  $\mathbf G$ is hermitian ($d(j,i)=d(i,j)^{-1}$ by freeness).  $T=\prod_vT_v$ with commuting
$T_v=R_{\psi_v}(\pi\circ\iota_v)$, all commuting with $\pi(G)$.  For an eigenvalue $\mu$ of $T$, take an irreducible $G$-invariant
$W\subset\ker(T-\mu)$; each $T_v$ acts on $W$ by a scalar $\mu_v$ (an eigenspace of $T_v$ in $W$ is invariant), and by
Lemma~\ref{i3:S-lem-GGplace} $\mu_v=\gamma_{\psi_v}(1)c_v$, $0\le c_v\le1+3/q_v$.  So every eigenvalue of the hermitian matrix $\mathbf G$ lies in
$[0,P_q\gamma(1)]$, which gives $\mathbf G\succeq0$ and $P_q\gamma(1)\mathbf 1-\mathbf G\succeq0$; the transpose of a positive semidefinite
hermitian matrix is positive semidefinite.
\end{proof}

\medskip\noindent\emph{The corank (CP).}
Let $(V,h)$ be a nondegenerate hermitian space over $E$.  A \emph{flag of norm $T$} is $F=(e_1,\dots,e_r;z;\tilde z)$ with
$e_1,\dots,e_r$ independent, $h(e_a,e_b)=0$, $h(e_a,z)=0$, $h(z,z)=-T$, and $\tilde z\in V$; its frame is $f=(e_1,\dots,e_r,z)$ and
$X_F=\sum Ee_a$.  For a lattice $\Lambda$, $F$ is a \emph{lift} if $z\in\Lambda$ and $\tilde z=z$.  For flags $F,F'$ of norms $T,T'$ let
$Z=(h(f_a,f'_b))$, $D=\det Z$, $D_r$ its leading $r\times r$ minor, $G(F,F')=D\bar D-TT'D_r\bar D_r$,
$U=(\Lambda\cap X_F)+\mathcal O_E\tilde z+(\Lambda\cap X_{F'})+\mathcal O_E\tilde z'$, $H'=(EU)^\perp$,
$\Theta(F,F',F'',\beta)=\{x\in H':\exists x''\in H'(F,F''),\ x-x''\in X_F,\ \beta x-x''\in X_{F''}\}$ and
$s_0=\max(k-\dim\Theta,0)$.  (D6): in $X(u,u')$, a term $t$ carries row lattices $(\Lambda_t,\Lambda'_t)$, which are full
$\mathcal O_E$-lattices in $V$ (so $\Lambda_t$ spans $V$ over $E$), a flag $F_t$ of norm $T_{u'}$
which is a lift for $\Lambda_t$, and an automorphism $\beta_t$; $u$ carries a base flag $F_0$ of norm $T_u$, a lift for every row lattice of
$u$; $G(t)=0$ iff $G(F_0,F_t)=0$; $s_0(t,t')=s_0(F_0,F_t,F_{t'},\beta_{t'}\beta_t^{-1};k)$ relative to $\Lambda_t$ when the row lattices of $t,t'$
agree, i.e.\ $(\Lambda_t,\Lambda'_t)=(\Lambda_{t'},\Lambda'_{t'})$, and $0$ otherwise; $\dim V=k+2r_S+1$.

\begin{lemma}\label{i3:S-lem-frameindep}
If $G(F,F')\ne0$, the $2r+2$ frame vectors of $F,F'$ are linearly independent.
\end{lemma}

\begin{proof}
Let $\sum a_if_i+\sum b_jf'_j=0$, $\ell=r+1$.  Pairing with $f_c$ and $f'_d$ gives $Zb=Ta_\ell\mathbf e_\ell$ and $Z^*a=T'b_\ell\mathbf e_\ell$; by the
adjugate, $Db_\ell=D_rTa_\ell$ and $\bar Da_\ell=\bar D_rT'b_\ell$.  If exactly one of $a_\ell,b_\ell$ is nonzero, or both, these give $G=0$
(in the last case $a_\ell b_\ell G=0$).  So $a_\ell=b_\ell=0$, $Zb=0$; if $D\ne0$ then $b=0$; if $D=0$ then $D_r\ne0$ and the first $r$ rows
give $b=0$.  Then $\sum_{a\le r}a_ae_a=0$ gives $a=0$.
\end{proof}

\begin{corollary}[(CP) for $X(u,u')$]\label{i3:S-cor-CP}
If $\Lambda$ spans $V$, $\dim V=k+2r+1$, $F,F'$ are lifts with $G(F,F')\ne0$, then $s_0(F,F',F'',\beta;k)\ge1$ for all $F'',\beta$.  Hence
$X(u,u')$ satisfies (CP).
\end{corollary}

\begin{proof}
The frame vectors lie in $W=EU$ ($ce_a\in\Lambda\cap X_F$ for some $c\ne0$; $z=\tilde z$).  For independent $u_1,\dots,u_m\in W$,
$v\mapsto(h(u_i,v))_i$ is onto $E^m$ (nondegeneracy) and kills $W^\perp$, so $\dim W^\perp\le\dim V-m$.  With $m=2r+2$
(Lemma~\ref{i3:S-lem-frameindep}), $\dim\Theta\le\dim H'\le k-1$.  For (CP): equal rows give equal row lattices, $G(t)\ne0$ gives
$G(F_0,F_t)\ne0$, and $F_0,F_t$ are lifts for $\Lambda_t$.
\end{proof}

\subsection{Kernel estimates off and on the diagonal}\label{i3:ssec-grp-kpp}

We prove Propositions~\ref{i3:prop-kpp-off}, \ref{i3:prop-kpp-diag}, \ref{i3:prop-W}, \ref{i3:prop-orbbd} and~\ref{i3:S-prop-KppVol}.

\subsubsection{Proof of Proposition~\ref{i3:prop-kpp-off}}\label{i3:ssec-pf-kppoff}

\emph{The objects.}  We recall from the construction of the term system $X(u,u')$ of a selected package $(X,Y)$
(Proposition~\ref{i3:prop-terms}) the facts used here.  At the infinite places the package fixes, for each word $u$,
a finite-dimensional Hilbert space $V_\pi(u)=\bigotimes_vV_{u,v}$ (tensor product over the infinite places $v$ of $F$,
with the product inner product on pure tensors, which span), a closed subgroup $S_u\subset U_\infty=\prod_vU(F_v)$ (the Bessel
subgroup: the elements acting unipotently on the flag $e_0,\dots,e_{r_S-1},z_u$ recalled below), a map $\pi_u$ from $S_u$ to the unitary
operators of $V_\pi(u)$, multiplicative, isometric and continuous on $S_u$ and trivial on its unipotent part $N_u$, a function $\psi_u$ of
modulus $1$ on $U_\infty$ (a character on $N_u$), and a continuous function $W_u:U_\infty\to V_\pi(u)$ (the archimedean Bessel vector) with
$W_u(sx)=\psi_u(s)\pi_u(s)W_u(x)$ for $s\in S_u$, $x\in U_\infty$.  Each factor $V_{u,v}$ carries an irreducible unitary
representation $\pi_{u,v}$ of $U(k)$ with highest weight determined by the archimedean weight $u(\tau_v)$ of the word, and
$\pi_u$ acts through these factors.  A term $t$ of $X(u,u')$ is a double coset with a representative
$\gamma_t\in U(F)$ and a pair of flags; it has a Gram determinant $G(t)$ (the function $G$ of
Definition~\ref{i3:def-terms}) and is \emph{relevant} or not according to the position of the two flags.  Its
\emph{orbital operator} is
\[
 \mathrm{OO}_t=\int_{\Omega_t}W_u(x)\otimes\overline{W_{u'}(\gamma_t^{-1}x)}\,d\mu_\infty(x)\ \in\mathrm{Hom}(V_\pi(u'),V_\pi(u)),
\]
with $\Omega_t$ a strict fundamental domain of the discrete group $\Gamma_t$ of rational points of the relative stabiliser $R_t$
(defined below; $\Omega_t$ exists and the integral converges absolutely by item (FD) of Proposition~\ref{i3:prop-kppoff-in}), and $\mathrm{op}_t$ is the image
of $\mathrm{OO}_t$ under a linear map (the assembly of the term over the row lattices); in particular
$\mathrm{OO}_t=0$ implies $\mathrm{op}_t=0$.  Finally $S_t\subset U_\infty$ denotes the pointwise stabiliser, at the infinite
places, of the isotropic part of the two flags of $t$.  A \emph{strict fundamental domain} of a subgroup $\Gamma$ is a measurable
$\Omega$ such that every $x$ has some $s\in\Gamma$ with $sx\in\Omega$, and $x\in\Omega$, $s\in\Gamma$, $sx\in\Omega$ imply $s=1$.

\emph{Linear algebra at infinity.}  We also recall the archimedean linear algebra of the construction.  Let $V=E^n$ with the
hermitian form $h$ of the datum ($n=2r_S+k+1$), conjugate-linear in the first variable.  For an infinite place $v$ (with a chosen
embedding $E\to\C$, also written $v$) write $x_v\in\C^n$ for the coordinate vector of $x\in V$ with $v$ applied entrywise (up to a fixed change of basis of $\C^n$ at
each $v$, which makes $h_v$ the standard form of its signature), $h_v$ for the
hermitian form on $\C^n$ with $h_v(x_v,y_v)=v(h(x,y))$, and $U(F_v)$ for its unitary group; for an $E$-subspace $A\subset V$ put
$A_v=\mathrm{span}_\C\{x_v:x\in A\}$.  The datum fixes $e_a,f_a\in V$ ($a<r_S$) with $h(e_a,e_b)=h(f_a,f_b)=0$, $h(e_a,f_b)=\delta_{ab}$,
and for each word $u$ a vector $z_u\perp e_a,f_a$ with $h(z_u,z_u)\ne0$; the flag of $u$ is $X_u=\mathrm{span}_E(e_a)$,
$X_u^+=X_u+Ez_u$.  At each $v$ the package fixes a frame $w_{u,v,b}\in\C^n$ ($b<k$), $h_v$-orthonormal and $h_v$-orthogonal to all
$e_{a,v},f_{a,v},z_{u,v}$.  For $s\in U_\infty$ let $m_v(s)\in M_k(\C)$, $m_v(s)_{ab}=h_v(w_{u,v,a},s_vw_{u,v,b})$.  The Bessel group $S_u$
contains every $s$ fixing all $e_{a,v}$ and $z_{u,v}$; $H_u\subset S_u$ is the pointwise stabiliser of all $e_{a,v},f_{a,v},z_{u,v}$,
$N_u\subset S_u$ is a subgroup, $\pi_u$ is multiplicative on $S_u$ and trivial on $N_u$, and (T) for $s\in H_u$,
$\pi_u(s)(\otimes_vx_v)=\otimes_v\pi_{u,v}(m_v(s))x_v$.  These facts (the frame relations, the description of $S_u$, $H_u$, $N_u$, the
properties of $\pi_u$, (T), and that each $\gamma_v$ is unitary) are part of the definition of the setting.  For a term $t$ with representative $\gamma=\gamma_t$, its flags are $F=(X_u,X_u^+)$ and
$F'=(\gamma X_{u'},\gamma X_{u'}^+)$; $\mathrm{Fix}_t=X_u^++\gamma X_{u'}^+$, and $S_t$ is the set of $h\in U_\infty$ with $h_vy=y$ for all
$y\in(\mathrm{Fix}_t)_v$ and all $v$.

\begin{proposition}[(Levi), (NR) and the first half of (FD) proved below; (NF) and the second half of (FD) formally verified, \cite{CharacterFamily}]\label{i3:prop-kppoff-in}
For every datum, tied package $(X,Y)$ and words $u,u'$:
\begin{enumerate}
\item[(NR)] (non-relevant character) if a term $t$ of $X(u,u')$ is not relevant, there is $r\in R_t$ with
$\psi_u(r)\ne\psi_{u'}(\gamma_t^{-1}r\gamma_t)$ ($R_t$ is defined below);
\item[(Levi)] for every $s\in S_u$ there is $s_0\in H_u$ with $s_0^{-1}s\in N_u$ and $m_v(s_0)=m_v(s)$ for every $v$ (applied at each
place; $S_u$, $N_u$, $H_u$ are defined place by place);
\item[(NF)] if $t$ is relevant and $G(t)=0$, then, with $\perp$ the orthogonal for $h$ in $V$,
$\mathrm{Fix}_t^\perp\cap X_u=X_u\cap\gamma_tX_{u'}$, $\mathrm{Fix}_t^\perp+X_u=(X_u^+)^\perp$ and
$\mathrm{Fix}_t^\perp+\gamma_tX_{u'}=(\gamma_tX_{u'}^+)^\perp$;
\item[(FD)] for a concretely tied package and every term $t$ of $X(u,u')$, $\Gamma_t$ has a strict fundamental domain $\Omega_t$ and
$x\mapsto W_u(x)\otimes\overline{W_{u'}(\gamma_t^{-1}x)}$ is integrable on $\Omega_t$.
\end{enumerate}
\end{proposition}

\noindent\emph{Structure of the remaining items.}  (NF) is linear algebra of the two flags over $E$, deduced in the formal proof from a normal form for pairs of
flags (new; its structure is described in \S\ref{i3:ssec-NF-structure}).  (FD) consists of the existence of $\Omega_t$ (its first half, proved in \S\ref{i3:ssec-pf-St-rest} from the
discreteness of $\Gamma_t$) and the absolute convergence of the orbital integrals (its second half, which uses both folds of the
pair and is formally verified; its structure is described in \S\ref{i3:str-FD}).  (Levi), the Levi decomposition of the Bessel group at one place, and (NR) are proved below.  From these, Lemma V and the geometry
of the stabiliser are proved below (Proposition~\ref{i3:prop-VGeo}).  The covariance of the orbital operators is
Proposition~\ref{i3:prop-cov}, proved from the smaller inputs of Proposition~\ref{i3:prop-cov-in}.

\begin{proof}[Proof of (Levi)]
Fix $v$ and drop it from the notation.  The vectors $e_a,f_a$ ($a<r_S$), $z_u$, $w_b$ ($b<k$) have a non-degenerate Gram
matrix (hyperbolic pairs, $h(z_u,z_u)\ne0$, an orthonormal frame, mutually orthogonal blocks), and there are
$2r_S+1+k=n$ of them, so they form a basis of $\C^n$.  Let $X=\mathrm{span}(e_a)$ and $W=\mathrm{span}(w_b)$; then
$X^\perp=X\oplus\C z_u\oplus W$ and $W^\perp=\mathrm{span}(e_a,f_a,z_u)$.  Let $s\in S_u$.  Since $(s-1)e_a\in X$, $sX=X$, and
$s$ being unitary, $sX^\perp=X^\perp$; also $sz_u\in z_u+X$.  For $w\in W$ write $sw=x+cz_u+Aw$ with $x\in X$, $Aw\in W$
(decomposition in $X^\perp$); then $c\,h(z_u,z_u)=h(z_u,sw)=h(sz_u,sw)-h(sz_u-z_u,sw)=h(z_u,w)-0=0$, since $sz_u-z_u\in X$ is
orthogonal to $sw\in X^\perp$.  Thus $sw\in Aw+X$, and $h(Aw,Aw')=h(sw,sw')=h(w,w')$ because $X$ is isotropic and orthogonal
to $X^\perp$; so $A$ is unitary on $W$, with matrix $(h(w_a,sw_b))_{ab}=m(s)$ in the orthonormal basis $(w_b)$.  Let $s_0$ be
the identity on $W^\perp$ and $A$ on $W$: it is unitary, fixes $e_a,f_a,z_u$ (so $s_0\in H_u$), and $m(s_0)=m(s)$.  Let
$n=s_0^{-1}s$.  Then $(n-1)e_a=(s-1)e_a\in E_{<a}$ and $(n-1)z_u=(s-1)z_u\in X$ (as $s_0^{-1}$ fixes $X$ pointwise); for
$y\perp X$, i.e.\ $y\in X^\perp$, $(n-1)y\in X$: this holds for $y\in X$ and $y=z_u$ by what precedes, and for $w\in W$,
$nw=s_0^{-1}(x+Aw)=x+w$.  Hence $n\in N_u$ at $v$; doing this at every $v$ gives $s_0$ and $s_0^{-1}s\in N_u$.
\end{proof}

\emph{Relevance.}  For a flag $F=(e_0,\dots,e_{r_S-1};z)$ in $V$ (the $e_a$ independent and isotropic, $z\perp e_a$, $h(z,z)\ne0$)
put $X_a=\mathrm{span}_E\{e_b:b<a\}$, $X=X_{r_S}$, $\nu_a=e_{a+1}$ ($a+1<r_S$), $\nu_{r_S-1}=z$, and let $\mathfrak n_F$ be the set of
$E$-linear $Y:V\to V$ with $h(Yx,y)+h(x,Yy)=0$, $YX_a\subset X_{a-1}$ ($1\le a\le r_S$) and $YX^\perp\subset X$.  For $Y\in\mathfrak n_F$,
$Y\nu_a\in X_{a+1}$; let $\ell_F(Y)=\sum_a\ell_{F,a}(Y)$ with $\ell_{F,a}(Y)$ the $e_a$-coordinate of $Y\nu_a$.  A pair $(F,F')$ is
\emph{Lie-relevant} if $\mathrm{tr}_{E/F}\ell_F(Y)=\mathrm{tr}_{E/F}\ell_{F'}(Y)$ for all $Y\in\mathfrak n_F\cap\mathfrak n_{F'}$.  In the
construction of the terms, $t$ is relevant if one of the flag pairs representing it is Lie-relevant; the pair
$(F_u,\gamma_tF_{u'})$, with $F_u=(e_a;z_u)$ and $\gamma_tF_{u'}=(\gamma_te_a;\gamma_tz_{u'})$, is one of them.

\begin{proof}[Proof of (NR)]
Let $t$ be not relevant, $\gamma=\gamma_t$, $F=F_u$, $F'=\gamma F_{u'}$.  Then $(F,F')$ is not Lie-relevant: there is
$Y\in\mathfrak n_F\cap\mathfrak n_{F'}$ with $\delta:=\lambda+\bar\lambda\ne0$, where $\lambda=\ell_F(Y)-\ell_{F'}(Y)$ (since
$\mathrm{tr}_{E/F}y=y+\bar y$).  Let $N$ be the number of infinite places of $F$, $c=g/(2N)\in\Q$ and $\sigma=c\,\delta^{-1}\in F$.
\emph{The element.}  $\sigma Y\in\mathfrak n_F\cap\mathfrak n_{F'}$ ($\bar\sigma=\sigma$).  An element $Y\in\mathfrak n_F$ is nilpotent: it maps
$X_a^\perp$ into $X_{a+1}^\perp$ (for $v\in X_a^\perp$, $x\in X_{a+1}$: $h(Yv,x)=-h(v,Yx)=0$), $X^\perp$ into $X$ and $X_a$ into
$X_{a-1}$.  So $n=\exp(\sigma Y)$ is a finite sum, an $E$-rational isometry (its adjoint is $\exp(-\sigma Y)$), and
$(n-1)X_a\subset X_{a-1}$, $(n-1)X^\perp\subset X$; likewise for $F'$, so $\gamma^{-1}n\gamma$ satisfies these conditions for
$F_{u'}$.  The conditions defining $N_u$ and $N_{u'}$ are $E$-linear, so the archimedean image $r=(n_v)_v$ lies in $N_u$ and
$\gamma^{-1}r\gamma$ in $N_{u'}$: $r\in R_t$.  \emph{The characters.}  As $Y^2\nu_a\in X_a$ and $h(f_a,\nu_a)=0$,
$h(f_a,n\nu_a)=\sigma h(f_a,Y\nu_a)=\sigma\ell_{F,a}(Y)$, and similarly $h(f_a,\gamma^{-1}n\gamma\nu'_a)=h(\gamma f_a,n\gamma\nu'_a)=
\sigma\ell_{F',a}(Y)$ ($\nu'_a$ for $u'$; $(\gamma f_a)$ is dual to $(\gamma e_a)$).  By the formula for $\psi_u$,
\[
\frac{\psi_u(r)}{\psi_{u'}(\gamma^{-1}r\gamma)}=\prod_v\exp\Bigl(2\pi i\,\frac{v(\sigma\lambda)+\overline{v(\sigma\lambda)}}{g}\Bigr)
=\prod_v\exp\Bigl(2\pi i\,\frac{v(\sigma\delta)}{g}\Bigr)=\exp\Bigl(2\pi i\,\frac{Nc}{g}\Bigr)=-1 ,
\]
since $\sigma\delta=c$ is rational.  Hence $\psi_u(r)\ne\psi_{u'}(\gamma^{-1}r\gamma)$.
\end{proof}

For a term $t$ with representative $\gamma=\gamma_t$ put $R_t=\{y\in N_u:\gamma^{-1}y\gamma\in N_{u'}\}$, where $N_u\subset U_\infty$ is the
unipotent part of the Bessel subgroup of $u$.  The term $t$ carries its row lattices $\Lambda=\Lambda_t$, $\Lambda'=\Lambda'_t$ (full
$\mathcal O_E$-lattices in $V$) and its representative is the archimedean image of an $E$-linear isometry $\gamma$ of $(V,h)$ with
$\gamma\Lambda'=\Lambda$ (we use the same letter for both).  For a full lattice $M\subset V$ let $N_M$ be the set of $E$-linear isometries
$n$ of $(V,h)$ with $nM=M$, $(n-1)X_{a}\subset X_{a-1}$ for $1\le a\le r_S$ (where $X_a=\mathrm{span}_E\{e_b:b<a\}$) and
$(n-1)X_u^\perp\subset X_u$.  Then
\[
\Gamma_t=\{n_\infty:\ n\in N_\Lambda,\ \gamma^{-1}n\gamma\in N_{\Lambda'}\},
\]
where $n_\infty=(n_v)_v$ is the archimedean image of $n$; it is a group, and $\Gamma_t\subset R_t$, since the conditions defining
$N_M$ are $E$-linear and pass to the spans at each $v$ (as in the proof of (NR) above).  $\Omega_t$ is a strict fundamental domain of
$\Gamma_t$.  For $v\in V_\pi(u)$, $w\in V_\pi(u')$ the operator $v\otimes\overline w$ is $\xi\mapsto\langle w,\xi\rangle v$; it is
linear in $v$ and conjugate-linear in $w$.

\begin{proposition}[(A), (S) proved below; (St) proved below modulo Proposition~\ref{i3:prop-St-rest}, formally verified, \cite{CharacterFamily}]\label{i3:prop-cov-in}
\begin{enumerate}
\item[(A)] (Averaging.)  Let $G$ be a second countable locally compact group with a left and right invariant Haar measure $\mu$,
$R\subset G$ a closed subgroup, $\Gamma\subset R$ a discrete subgroup with $R\subset\Gamma^{-1}K$ for a compact $K$ (i.e.\ for every $r\in R$
some $s\in\Gamma$ has $sr\in K$), $\chi:R\to\C$ continuous and multiplicative, trivial on $\Gamma$ and not identically $1$ on $R$, $f:G\to V$
($V$ a normed space) with $f(rx)=\chi(r)f(x)$ for $r\in R$, $x\in G$, and $\Omega$ a strict fundamental domain of $\Gamma$.  Then
$\int_\Omega f\,d\mu=0$ (Bochner integral, with the convention that it is $0$ if $f$ is not integrable on $\Omega$ or $V$
is not complete).
\item[(S)] (Shift.)  Let $G$ be a second countable topological group with a $\sigma$-finite left invariant measure $\mu$, $R$, $\Gamma$, $K$ as in
(A) with $\Gamma$ countable, $h\in G$ with $r\in R\iff h^{-1}rh\in R$, and $\nu$ a nonzero $\sigma$-finite measure on $G$, concentrated
on $R$, finite on compact sets, and invariant under $x\mapsto rx$ ($r\in R$), $x\mapsto x^{-1}$ and $x\mapsto h^{-1}xh$.  If $f:G\to V$ satisfies
$f(rx)=f(x)$ for $r\in R$, $\Omega$ is a strict fundamental domain of $\Gamma$ and $f$ is integrable on $\Omega$, then
$\int_\Omega f\,d\mu=\int_{h^{-1}\Omega}f\,d\mu$.
\item[(St)] (Structure.)  For every datum, tied package and words $u,u'$: for every term $t$, (i) and (ii) hold, and for every
relevant term $t$ with $G(t)=0$, (iii) and (iv) hold:
(i) $\mu_\infty$ is left and right invariant; $R_t$ is a closed subgroup of $U_\infty$ containing $\Gamma_t$, $\Gamma_t$ is countable and
discrete, and $R_t\subset\Gamma_t^{-1}K$ for a compact $K$;
(ii) $\psi_u$ and $\psi_{u'}$ are continuous unitary characters of $N_u$, $N_{u'}$, $\pi_u$ is trivial on $N_u$ (so
$W_u(sx)=\psi_u(s)W_u(x)$ for $s\in N_u$), and $\psi_u(s)=\psi_{u'}(\gamma^{-1}s\gamma)$ for $s\in\Gamma_t$;
(iii) every $h\in S_t$ lies in the Bessel subgroup of $u$ with $\psi_u(h)=1$, and $\gamma^{-1}h\gamma$ in that of $u'$ with
$\psi_{u'}(\gamma^{-1}h\gamma)=1$;
(iv) every $h\in S_t$ normalises $R_t$, and if $\psi_u(r)=\psi_{u'}(\gamma^{-1}r\gamma)$ for all $r\in R_t$ there is a measure $\nu$ as in (S)
for $(R_t,h)$ (the Haar-measure relation).
\end{enumerate}
\end{proposition}

\noindent (A) and (S) are the unfolding identities behind the Weil integration formula over $\Gamma\backslash R$ and
$R\backslash G$; we prove them directly.  (St) collects facts about the archimedean setting; most of it is proved below
from the explicit form of $N_u$ and $\psi_u$, and the remaining part is the smaller Proposition~\ref{i3:prop-St-rest}.

\emph{Explicit form of $S_u$, $N_u$, $\psi_u$.}  Write $E_{<a,v}=\mathrm{span}_\C\{e_{b,v}:b<a\}$ and $E_v=E_{<r_S,v}$.  The
Bessel subgroup is $S_u=\{s\in U_\infty:$ for all $v$ and $a<r_S$, $(s_v-1)e_{a,v}\in E_{<a,v}$ and $(s_v-1)z_{u,v}\in E_v\}$,
and $N_u$ is the set of $s\in S_u$ with, moreover, $(s_v-1)y\in E_v$ for every $y\in\C^n$ which is $h_v$-orthogonal to all
$e_{a,v}$.  With $g=g_0g_1$, $\nu_a=e_{a+1}$ for $a+1<r_S$ and $\nu_{r_S-1}=z_u$,
\[
\psi_u(s)=\prod_v\exp\Bigl(2\pi i\,\frac{c_v(s)+\overline{c_v(s)}}{g}\Bigr),\qquad c_v(s)=\sum_{a<r_S}h_v(f_{a,v},s_v\nu_{a,v}).
\]

\begin{proposition}[(a) proved; (c) proved modulo (NF) of Proposition~\ref{i3:prop-kppoff-in}; (b) proved modulo Proposition~\ref{i3:prop-levtr}, formally verified, \cite{CharacterFamily}]\label{i3:prop-St-rest}
For every datum, tied package and words $u,u'$, and every term $t$ of $X(u,u')$: (a) $\Gamma_t$ is discrete and there is a
compact $K$ with $R_t\subset\Gamma_t^{-1}K$; (b) $\psi_u(s)=\psi_{u'}(\gamma_t^{-1}s\gamma_t)$ for $s\in\Gamma_t$; (c) if $t$ is
relevant with $G(t)=0$, $h\in S_t$, and $\psi_u(r)=\psi_{u'}(\gamma_t^{-1}r\gamma_t)$ for all $r\in R_t$, there is a measure
$\nu$ as in (S) for $(R_t,h)$.
\end{proposition}

\noindent The proofs are in \S\ref{i3:ssec-pf-St-rest}.  (a) is the discreteness and cocompactness of the group of rational
points of the unipotent group $R_t$ preserving the two lattices of the term; it is an instance of the cocompactness of
arithmetic subgroups of unipotent groups (Borel--Harish-Chandra, Ann.\ of Math.\ 75 (1962); reference to be confirmed), which we
do not use: we prove it directly, by rational approximation along a filtration of the Lie algebra.  (c) takes for $\nu$ a Haar
measure of $R_t$ and decomposes $h=kn$ with $n\in R_t$ and $k$ normalising $R_t$ with relatively compact powers; it uses (NF).
(b) is reduced, by the product formula for the additive character and the multiplicativity of the level character, to the
transport of rational unipotent elements into the level subgroup (Proposition~\ref{i3:prop-levtr}), which is formally verified.

\begin{proof}[Proof of (St), apart from Proposition~\ref{i3:prop-St-rest}]
(i) \emph{Unimodularity.}  $\mu_\infty$ is a product of Haar measures of the groups $U(F_v)\cong U(p_v,q_v)$.  For such a
group $G\subset GL_n(\C)$ with Lie algebra $\mathfrak g$, the form $B(X,Y)=\mathrm{tr}(XY)$ is real on $\mathfrak g$ (from
$X^*=-JXJ^{-1}$) and non-degenerate (its complexification is the trace form of $\mathfrak{gl}_n(\C)=\mathfrak g\otimes\C$),
and $\mathrm{Ad}(x)$ preserves it; so $|\det\mathrm{Ad}(x)|=1$, the modular function is trivial, and $\mu_\infty$ is left
and right invariant.

\emph{$R_t$ is a closed subgroup.}  The conditions defining $N_u$ say that the linear maps $s_v-1$ send fixed vectors into
fixed subspaces, so $N_u$ is closed.  It is a subgroup: from $st-1=s(t-1)+(s-1)$ and the fact that each $s\in N_u$ maps
$E_{<a,v}$ and $E_v$ into themselves (since $(s_v-1)e_{b,v}\in E_{<b,v}$), and from $s^{-1}-1=-s^{-1}(s-1)$.  As
conjugation by $\gamma_t$ is a homeomorphism, $R_t=N_u\cap\gamma_tN_{u'}\gamma_t^{-1}$ is a closed subgroup.  The
elements of $\Gamma_t$ are the families $(n_v)_v$ of archimedean images of single matrices with entries in the countable
field $E$, so $\Gamma_t\subset R_t$ is countable.  The rest of (i) is Proposition~\ref{i3:prop-St-rest}(a).

(ii) \emph{$\psi_u$ is a continuous unitary character of $N_u$.}  $c+\bar c$ is real, so $|\psi_u|=1$; continuity is clear.
For $s,s'\in S_u$, $s'_v\nu_{a,v}-\nu_{a,v}\in E_{<a+1,v}$, and its $e_{a,v}$-coefficient is $h_v(f_{a,v},s'_v\nu_{a,v})$, as
$h_v(f_{a,v},\nu_{a,v})=0$; since $s_v$ is unitriangular on $E_{<a+1,v}$ with respect to $(e_{b,v})$, $s_v(s'_v-1)\nu_{a,v}$
has the same $e_{a,v}$-coefficient, and $h_v(f_{a,v},\cdot)$ is that coefficient on $E_{<a+1,v}$.  Hence
$c_v(ss')=h_v(f_a,s\nu_a)+h_v(f_a,s(s'-1)\nu_a)$ summed over $a$ equals $c_v(s)+c_v(s')$, and $\psi_u(ss')=\psi_u(s)\psi_u(s')$.
That $\pi_u$ is trivial on $N_u$ is part of the setting, so $W_u(sx)=\psi_u(s)\pi_u(s)W_u(x)=\psi_u(s)W_u(x)$ for $s\in N_u$.
The last clause of (ii) is Proposition~\ref{i3:prop-St-rest}(b).

(iii) Let $h\in S_t$.  Then $h_v$ fixes $e_{a,v}$ and $z_{u,v}$, so $h\in S_u$, and $c_v(h)=\sum_ah_v(f_{a,v},\nu_{a,v})=0$,
so $\psi_u(h)=1$.  Also $h_v$ fixes $\gamma_ve_{a,v}$ and $\gamma_vz_{u',v}$, so $\gamma^{-1}h\gamma$ fixes $e_{a,v}$ and
$z_{u',v}$; hence it lies in $S_{u'}$ with $\psi_{u'}(\gamma^{-1}h\gamma)=1$.

(iv) If a unitary $h$ fixes all $e_{a,v}$ and $z_{u,v}$, then $hN_uh^{-1}=N_u$: $(hsh^{-1}-1)e_{a,v}=h(s-1)e_{a,v}\in
hE_{<a,v}=E_{<a,v}$, similarly for $z_{u,v}$, and $y\perp e_{a,v}$ for all $a$ implies $h^{-1}y\perp h^{-1}e_{a,v}=e_{a,v}$.
Applied to $h\in S_t$ and to $\gamma^{-1}h\gamma$ (for $u'$), this shows that $h$ normalises $R_t$.  The Haar-measure
relation is Proposition~\ref{i3:prop-St-rest}(c).
\end{proof}

\begin{proof}[Proof of (A) and (S)]
\emph{Preliminaries.}  In both statements $\Gamma$ is countable (in (A) because a discrete subgroup of a second countable
group is countable).  For $x\in G$ the set $\{s\in\Gamma:sx\in K\}=\Gamma\cap Kx^{-1}$ is finite (discrete and contained in
a compact set; $\Gamma$ is closed, being discrete).  Fix a Borel linear order on $G$ (for instance, the lexicographic order of the sequences
$(\mathbf 1_{B_n}(x))_n$ for a countable base $(B_n)$ of the Hausdorff space $G$); then $D=\{r\in R\cap K:\ r\le sr$ for all $s\in\Gamma$ with $sr\in K\}$ is Borel and is a
strict fundamental domain for the left action of $\Gamma$ on $R$: every orbit $\Gamma r$ meets $K$ (hypothesis), in finitely
many points, and $D$ contains exactly the least one.  We use the elementary facts: (FD1) if $\Gamma$ (countable) acts by
measure preserving maps on a measure space and $D_1,D_2$ are strict fundamental domains, then
$\int_{D_1}\Phi=\int_{D_2}\Phi$ for every $\Gamma$-invariant $\Phi$ which is nonnegative or integrable on $D_1$
(decompose $D_1=\bigsqcup_s(D_1\cap sD_2)$ and translate); in particular all strict fundamental domains have the same
measure; (FD2) for $x\in G$, the set $\{q\in R:qx\in\Omega\}$ is a strict fundamental domain for $\Gamma$ acting on $R$ by
left multiplication.

\emph{(A).}  We may assume $V$ complete and $f$ integrable on $\Omega$.  Let $\nu$ be a left Haar measure of $R$ (a closed
subgroup, hence locally compact), viewed as a measure on $G$ carried by $R$.  Then $0<\nu(D)\le\nu(K)<\infty$ ($R$ is the
countable union of the $sD$).  \emph{$\nu$ is right invariant:} for $r\in R$, $Dr$ is a strict fundamental domain of
$\Gamma$ in $R$, and $\nu(Dr)=\Delta_R(r)\nu(D)$ with the modular function $\Delta_R$; by (FD1) $\nu(Dr)=\nu(D)$, so
$\Delta_R(r)=1$.  Consequently $\nu$ is also invariant under $x\mapsto x^{-1}$.  \emph{The character integrates to zero:}
extend $\chi$ by $0$ off $R$; it is left $\Gamma$-invariant.  Pick $r_0\in R$ with $\chi(r_0)\ne1$.  By right invariance
and (FD1) applied to $Dr_0$,
$\int_D\chi\,d\nu=\int_{Dr_0}\chi\,d\nu=\int_D\chi(rr_0)\,d\nu(r)=\chi(r_0)\int_D\chi\,d\nu$, so $\int_D\chi\,d\nu=0$.
\emph{Unfolding:} on $R\times G$ with $\nu\otimes\mu$ let $\Gamma$ act diagonally on the left, and put
$H(r,x)=\mathbf 1_\Omega(r^{-1}x)f(x)$.  $H$ is $\Gamma$-invariant ($(sr)^{-1}sx=r^{-1}x$, $f(sx)=\chi(s)f(x)=f(x)$), and
both $D\times G$ and $R\times\Omega$ are strict fundamental domains.  On $R\times\Omega$, by Fubini,
\[
\int_{R\times\Omega}H=\int_\Omega f(x)\,\nu\{r\in R:r^{-1}x\in\Omega\}\,d\mu(x)=\kappa\int_\Omega f\,d\mu,
\]
where $\kappa=\nu\{r:r^{-1}x\in\Omega\}=\nu\{q:qx\in\Omega\}=\nu(D)$ by inversion invariance, (FD2) and (FD1); the same
computation with $|f|$ shows that $H$ is integrable on $R\times\Omega$, so (FD1) applies.  On $D\times G$, substituting
$x=ry$ (left invariance of $\mu$) and using $f(ry)=\chi(r)f(y)$,
\[
\int_{D\times G}H=\int_D\int_\Omega f(ry)\,d\mu(y)\,d\nu(r)=\Bigl(\int_D\chi\,d\nu\Bigr)\int_\Omega f\,d\mu=0.
\]
Hence $\nu(D)\int_\Omega f\,d\mu=0$, and $\nu(D)>0$.

\emph{(S).}  We may assume $V$ complete.  Put $\Gamma'=h^{-1}\Gamma h$ and $\Omega'=h^{-1}\Omega$; since $h$ normalises $R$,
$\Gamma'\subset R$ is countable, $\Omega'$ is a strict fundamental domain of $\Gamma'$ in $G$, and $D'=h^{-1}Dh$ is a strict
fundamental domain of $\Gamma'$ in $R$.  By the hypotheses on $\nu$: $\nu$ is left $R$-invariant, so (FD1) applies to
$\Gamma$ and to $\Gamma'$; $\nu(D')=\nu(D)$ by conjugation invariance; $\nu(D)\le\nu(K)<\infty$; and $\nu(D)>0$, since
$\nu\ne0$ is carried by $R=\bigcup_s sD$.  By (FD2) and (FD1), $\nu\{r\in R:rx\in\Omega\}=\nu(D)$ for every $x$, and,
using inversion invariance, $\nu\{r\in R:r^{-1}y\in\Omega'\}=\nu\{q:qy\in\Omega'\}=\nu(D')=\nu(D)$ for every $y$.  The
function $f$ is integrable on $\Omega$; by Fubini (both measures $\sigma$-finite), the substitution $y=rx$ (left
invariance of $\mu$) and $f(r^{-1}y)=f(y)$,
\[
\nu(D)\int_{\Omega'}f\,d\mu=\int_R\int_{\Omega'}\mathbf 1_\Omega(rx)f(x)\,d\mu(x)\,d\nu(r)
=\int_R\int_\Omega\mathbf 1_{\Omega'}(r^{-1}y)f(y)\,d\mu(y)\,d\nu(r)=\nu(D)\int_\Omega f\,d\mu ,
\]
the same computation with $|f|$ justifying Fubini (it shows $\nu(D)\int_{\Omega'}|f|=\nu(D)\int_\Omega|f|<\infty$).  Divide
by $\nu(D)\in(0,\infty)$.
\end{proof}

\begin{proposition}[Covariance]\label{i3:prop-cov}
For every datum, tied package $(X,Y)$, words $u,u'$ and relevant term $t$ of $X(u,u')$ with $G(t)=0$:
$\pi_u(h)\circ\mathrm{OO}_t=\mathrm{OO}_t\circ\pi_{u'}(\gamma_t^{-1}h\gamma_t)$ for every $h\in S_t$.
\end{proposition}

\begin{proof}
Write $\gamma=\gamma_t$, $\Omega=\Omega_t$, $f(x)=W_u(x)\otimes\overline{W_{u'}(\gamma^{-1}x)}$, so $\mathrm{OO}_t=\int_\Omega f\,d\mu_\infty$, and
$\chi(r)=\psi_u(r)\overline{\psi_{u'}(\gamma^{-1}r\gamma)}$ for $r\in R_t$.  By (St)(ii) $\chi$ is continuous and multiplicative, and trivial on
$\Gamma_t$ (as $|\psi_{u'}|=1$).  For $r\in R_t$, $W_u(rx)=\psi_u(r)W_u(x)$ and
$W_{u'}(\gamma^{-1}rx)=W_{u'}((\gamma^{-1}r\gamma)\gamma^{-1}x)=\psi_{u'}(\gamma^{-1}r\gamma)W_{u'}(\gamma^{-1}x)$; the second slot being conjugate-linear,
$f(rx)=\chi(r)f(x)$.

\emph{Case 1: $\chi\ne1$ on $R_t$.}  By (A), with (St)(i), $\mathrm{OO}_t=0$, and both sides vanish.

\emph{Case 2: $\chi\equiv1$ on $R_t$.}  Fix $h\in S_t$ and put $h'=\gamma^{-1}h\gamma$.  By (St)(iii) and the transformation law of the Bessel
vectors, $W_u(hx)=\pi_u(h)W_u(x)$ and $W_{u'}(h'y)=\pi_{u'}(h')W_{u'}(y)$ for all $x,y$; the latter also holds for $h'^{-1}$, by (St)(iii)
applied to $h^{-1}\in S_t$ ($S_t$ is a group).  Let $F(x)=W_u(hx)\otimes\overline{W_{u'}(\gamma^{-1}x)}=\pi_u(h)\circ f(x)$; it is integrable
on $\Omega$ since $f$ is (Proposition~\ref{i3:prop-kppoff-in}(FD)).  For $r\in R_t$, $F(rx)=\pi_u(h)f(rx)=\chi(r)F(x)=F(x)$.  By (St)(iv) the hypotheses of (S) hold for $(R_t,\Gamma_t,h,F)$, so
\[
\int_\Omega F\,d\mu_\infty=\int_{h^{-1}\Omega}F\,d\mu_\infty=\int_\Omega F(h^{-1}x)\,d\mu_\infty(x),
\]
the second equality by left invariance of $\mu_\infty$.  Now $F(h^{-1}x)=W_u(x)\otimes\overline{W_{u'}(\gamma^{-1}h^{-1}x)}$ and
$\gamma^{-1}h^{-1}x=h'^{-1}\gamma^{-1}x$, so $W_{u'}(\gamma^{-1}h^{-1}x)=\pi_{u'}(h')^{-1}W_{u'}(\gamma^{-1}x)$.  Since $\pi_{u'}(h')$ is unitary,
$\langle\pi_{u'}(h')^{-1}w,\xi\rangle=\langle w,\pi_{u'}(h')\xi\rangle$, i.e.\ $v\otimes\overline{\pi_{u'}(h')^{-1}w}=(v\otimes\overline w)\circ\pi_{u'}(h')$.
Hence $F(h^{-1}x)=f(x)\circ\pi_{u'}(h')$.  Composition with a fixed operator commutes with the Bochner integral (finite-dimensional spaces), so
\[
\pi_u(h)\circ\mathrm{OO}_t=\int_\Omega F=\int_\Omega f(x)\circ\pi_{u'}(h')\,d\mu_\infty(x)=\mathrm{OO}_t\circ\pi_{u'}(h').\qedhere
\]
\end{proof}

\begin{proposition}[Lemma V and the stabiliser]\label{i3:prop-VGeo}
For every datum, tied package $(X,Y)$ and words $u,u'$:
\begin{enumerate}
\item[(V)] if a term $t$ of $X(u,u')$ is not relevant, then $\mathrm{OO}_t=0$ and $\mathrm{op}_t=0$;
\item[(Geo)] if $t$ is relevant and $G(t)=0$, there are $C_v\in U(k)$ such that for every
$M=(M_v)\in\prod_vU(k)$ some $h\in S_t$ acts by $\pi_u(h)=\bigotimes_v\pi_{u,v}(M_v)$ on $V_\pi(u)$ and by
$\pi_{u'}(\gamma_t^{-1}h\gamma_t)=\bigotimes_v\pi_{u',v}(C_vM_vC_v^*)$ on $V_\pi(u')$.
\end{enumerate}
\end{proposition}

\noindent The proof, from Propositions~\ref{i3:prop-kppoff-in} and \ref{i3:prop-cov-in}, is in \S\ref{i3:ssec-pf-VGeo}.

\begin{lemma}[Schur lemma for distinct parameters]\label{i3:lem-schurcore}
Let $\iota$ be a finite set, and for $i\in\iota$ let $\pi_i,\pi'_i$ be irreducible finite-dimensional unitary
representations of $U(k)$ on $V_i,V'_i$ with highest weights $\Lambda_i,\Lambda'_i$, and $C_i\in U(k)$.  Let
$V=\bigotimes_iV_i$, $V'=\bigotimes_iV'_i$, $\rho(M)=\bigotimes_i\pi_i(M_i)$ and $\rho'(M)=\bigotimes_i\pi'_i(C_iM_iC_i^*)$ for
$M\in G=\prod_iU(k)$.  If $\Lambda_{i_0}\ne\Lambda'_{i_0}$ for some $i_0$, then every linear $P:V'\to V$ with
$\rho(M)P=P\rho'(M)$ for all $M\in G$ is zero.
\end{lemma}

\begin{proof}
Restrict to the copy $U(k)\hookrightarrow G$ in the coordinate $i_0$ (with $M_i=1$ for $i\ne i_0$).  Then $V$ is, as a
representation of this $U(k)$, isomorphic to $V_{i_0}^{\oplus a}$ with $a=\prod_{i\ne i_0}\dim V_i$, and $V'$ to
$(\pi'_{i_0}\circ\mathrm{Ad}(C_{i_0}))^{\oplus a'}$; and $\pi'_{i_0}\circ\mathrm{Ad}(C_{i_0})$ is equivalent to $\pi'_{i_0}$, through
the operator $\pi'_{i_0}(C_{i_0})$.  $P$ intertwines these restrictions.  $\pi_{i_0}$ and $\pi'_{i_0}$ are inequivalent: an
equivalence carries weight vectors (for the diagonal torus) to weight vectors of the same weight, so the highest weight
is an invariant of the equivalence class, and $\Lambda_{i_0}\ne\Lambda'_{i_0}$.  Decompose $P$ into blocks
$P_{bc}:V'_{i_0}\to V_{i_0}$ according to the two direct sum decompositions: each block intertwines $\pi'_{i_0}$ (transported by
$\pi'_{i_0}(C_{i_0})$) with $\pi_{i_0}$, two inequivalent irreducible representations, so it vanishes by Schur's lemma (its
kernel and image are invariant subspaces).  Hence $P=0$.
\end{proof}

\noindent (The formal statement phrases the hypothesis through Harish-Chandra parameters $\lambda$, with highest weight
$\lambda+\rho_n-\rho_c$; since the shift $\rho_n-\rho_c$ is fixed, distinct parameters are the same as distinct highest
weights.)

\begin{proof}[Proof of Proposition~\ref{i3:prop-kpp-off}]
Take the $\varepsilon_1$ and threshold system of Proposition~\ref{i3:prop-terms}, add the threshold $3$, let $\mathcal F$ be
large and $x$ beyond the thresholds of the selection, and let a member satisfy the standing hypotheses at
$(3,C_{k,\rm low},C_l^{\mathcal F})$; as $C_{k,\rm low}\ge3\ge1$ these give the $K_2$-hypotheses of the datum at
$(C_{k,\rm low},C_l^{\mathcal F})$, and the selected package is concretely tied.  Let $u\ne u'$.  One of the
$K_2$-hypotheses (the separation of the words) gives $u(\tau)\ne u'(\tau)$ at every infinite place $\tau$, so the highest
weights of $\pi_{u,v}$ and $\pi_{u',v}$ differ at every $v$.

By Proposition~\ref{i3:prop-kppoff-in}(FD), for every term $t$ the strict fundamental domain $\Omega_t$ exists and
$\mathrm{OO}_t$ is defined.  Let $t$ be a degenerate term: $G(t)=0$ and $t$ is not the identity term.  If $t$ is not relevant, $\mathrm{op}_t=0$ by Proposition~\ref{i3:prop-VGeo}(V).
If $t$ is relevant, put $P=\mathrm{OO}_t$.  By Proposition~\ref{i3:prop-cov}, $\pi_u(h)P=P\pi_{u'}(\gamma_t^{-1}h\gamma_t)$ for $h\in S_t$; by
Proposition~\ref{i3:prop-VGeo}(Geo),
for every $M\in\prod_vU(k)$ there is such an $h$ with $\pi_u(h)=\rho(M)$ and $\pi_{u'}(\gamma_t^{-1}h\gamma_t)=\rho'(M)$, in the
notation of Lemma~\ref{i3:lem-schurcore} with $\pi_v=\pi_{u,v}$, $\pi'_v=\pi_{u',v}$.  Hence $\rho(M)P=P\rho'(M)$ for all $M$, and
Lemma~\ref{i3:lem-schurcore} gives $P=0$.  So $\mathrm{OO}_t=0$ and $\mathrm{op}_t=0$.
\end{proof}

\subsubsection{Proof of Proposition~\ref{i3:prop-kpp-diag}}\label{i3:ssec-pf-kppdiag}

We use the notation of (Ty) (\S\ref{i3:ssec-S}): $j(t)$, $\mathrm{lab}(t)$, $\Delta$, $S$, $H_b$, $H_1$, $d_\pi(u)$, $m(u)$.  Put
\[
A_1(j)=\Bigl(\frac Ck\Bigr)^{j(k-1-j)f}H_b(C)^{Cr_S^4}H_b(C_h)^{C_hr_S^4}H_b(C)^{\frac{j+1}2}H_b(C_c)^{C_cr_S},\qquad
A(c;j)=e^{c\,m^6fL}\Bigl(\frac ck\Bigr)^{j(k-1-j)f},
\]
and $B(c)=\bigl(\sum_{j=1}^{r_S}A(c;j)\bigr)^2$, where $L=\log b$ and $m=2r_S+2$.

\begin{proposition}[(Irr), (Num) proved below; (J), (Del), (Cos), (Dim) formally verified, \cite{CharacterFamily}]\label{i3:prop-kppdiag-in}
For all outer constants (with $c_b,c_T>0$) and $C_H\ge0$:
\begin{enumerate}
\item[(Irr)] at every tied package, a degenerate term $t$ of $X(u,u)$ which is not relevant has $\mathrm{vec}_t=0$;
\item[(J)] under the $K_2$-hypotheses, for every concretely tied package, every relevant degenerate term has $j(t)\ge1$;
\item[(Del)] there is $C^{(0)}_\Delta$ such that for every $C\ge C^{(0)}_\Delta$ there are thresholds $(C'_0,C_{l,0})$ such
 that, under the $K_2$-hypotheses at $(C',C_l)$ above them, for every concretely tied package, every relevant degenerate $t$ has
 $0<\Delta(j(t),\mathrm{lab}(t))\le H_b(C)\,S(H_1,j(t),\mathrm{lab}(t))^2$;
\item[(Cos)] there are $C_c\ge0$ and thresholds such that, under the $K_2$-hypotheses above them, for every concretely tied
 package, for $1\le j\le r_S$ and every
 finite set $T$ of labels each of which is $\mathrm{lab}(t)$ for some relevant degenerate $t$ with $j(t)=j$,
 $\sum_{\delta\in T}S(H_1,j,\delta)^{-3r_S}\le H_b(C_c)^{C_cr_S}$;
\item[(Dim)] at every tied package, $\Gamma_1\ge0$ and $\Gamma_1\,m(u)\,d_\pi(u)\le d_u$, $\Gamma_1$ being the constant of the
 term system $X(u,u)$;
\item[(Num)] for every $c\ge1$ there is $c_0(c)$ such that $2d\,B(c)\le d\,e^{-kf\log k}$ for every $d\ge0$ and every datum
 with $c_0(c)\,m^6L\le k\log k$.
\end{enumerate}
\end{proposition}

\noindent (J) is a statement about the geometry of the flag pairs (a relevant degenerate term with $j(t)=0$ would, by
the normal form of the flag pair, have equal flags and be the identity term; this uses that the levels $T_w$ are prime
to $g_0g_1$); (Del) and (Cos) are the local index bound and the count of coset labels; (Dim) is the dimension formula of
the column space.  These four are formally verified; their formally verified steps are outlined in \S\ref{i3:str-kppdiag}.  We now prove (Irr) and (Num).

\begin{proof}[Proof of (Irr)]
By construction the vector $\mathrm{vec}_t$ of a term in the row model is the image of its orbital operator
$\mathrm{OO}_t$ under a linear map (as for $\mathrm{op}_t$).  If $t$ is not relevant, $\mathrm{OO}_t=0$ by
Proposition~\ref{i3:prop-VGeo}(V), so $\mathrm{vec}_t=0$.
\end{proof}

\begin{proof}[Proof of (Num)]
We use that every datum has $b\ge2$ and $f\ge1$, and $m=2r_S+2$, $L=\log b$.  Let $c\ge1$ and $c_0(c)=\max(2,c^7,32c)$.  If
$r_S=0$ then $B(c)=0$; let $r_S\ge1$, so $m\ge4$, and assume $c_0m^6L\le k\log k$; as $L\ge\log2$, $k\log k\ge c_0m^6\log2$.
(a) $k\ge2m$: otherwise $k\log k<2m\log(2m)<4m^2\le m^6/64<c_0m^6\log2$.  Hence $k-1-j\ge k-1-r_S\ge\frac34k$ for $1\le j\le r_S$
(as $2m=4r_S+4$).  (b) $k\ge c^6$: otherwise $k\log k<6c^6\log c\le6c^7<c_0m^6\log2$.  Hence $c/k\le1$ and
$\log(k/c)\ge\frac56\log k$.  (c) $c\,m^6fL\le\frac c{c_0}fk\log k\le\frac1{32}fk\log k$.  Put $Q=fk\log k$; by (a), $Q\ge8\log8$.
By (a)--(c), for $1\le j\le r_S$,
\[
A(c;j)=e^{cm^6fL}\Bigl(\frac ck\Bigr)^{j(k-1-j)f}\le\exp\Bigl(\frac Q{32}-j\cdot\frac34k\cdot\frac56\log k\cdot f\Bigr)
=e^{-(\frac{5j}8-\frac1{32})Q}\le x^j,\qquad x=e^{-\frac{19}{32}Q}\le\frac12 .
\]
So $\sum_jA(c;j)\le2x$ and $2B(c)\le8x^2=8e^{-\frac3{16}Q}e^{-Q}\le e^{-Q}$, since $\frac3{16}Q\ge\frac32\log8>\log8$.  Multiply by
$d\ge0$.
\end{proof}

\begin{lemma}[Label sums]\label{i3:lem-labelsum}
Assume (Del) and (Cos) at $C\ge0$, $C_h\ge0$, $C_s\ge C_p+4$, and $k\ge r_S+3$.  For $1\le j\le r_S$ and every finite set $T$
of labels as in (Cos), $\sum_{\delta\in T}\kappa(j,\delta)\le A_1(j)$.
\end{lemma}

\begin{proof}
Let $\delta\in T$, $\delta=\mathrm{lab}(t)$ with $j(t)=j$; write $\Delta=\Delta(j,\delta)>0$ and $S=S(H_1,j,\delta)\ge1$.
Multiplying the two factors of $\kappa$,
$\kappa(j,\delta)=K_j\,\Delta^{\frac{j+1}2}S^{(C_p-C_s)k}$ with $K_j=(C/k)^{j(k-1-j)f}H_b(C)^{Cr_S^4}H_b(C_h)^{C_hr_S^4}\ge0$,
since $-\frac{k-1-j}2+\frac k2=\frac{j+1}2$.  By (Del), $\Delta^{\frac{j+1}2}\le H_b(C)^{\frac{j+1}2}S^{j+1}$.  As $S\ge1$
and $j+1+(C_p-C_s)k\le r_S+1-4k\le-3r_S$ (using $k\ge r_S+3$), $S^{j+1}S^{(C_p-C_s)k}\le S^{-3r_S}$.  Hence
$\kappa(j,\delta)\le K_jH_b(C)^{\frac{j+1}2}S^{-3r_S}$, and summing over $T$ with (Cos) gives
$\sum_T\kappa\le K_jH_b(C)^{\frac{j+1}2}H_b(C_c)^{C_cr_S}=A_1(j)$.
\end{proof}

\begin{lemma}[Finite sums]\label{i3:lem-finbound}
Under (Ty) at $(C_H,C_s,C,C_h,C_p)$, (Irr), (J), (Del), (Cos) and $k\ge r_S+3$, every finite set $Q$ of pairs in
$\mathrm{all}(\mathrm{deg})$ for $X(u,u)$ satisfies
$\sum_{(t,t')\in Q}\|\mathrm{vec}_t\|\|\mathrm{vec}_{t'}\|\le d_\pi(u)m(u)\bigl(\sum_{j=1}^{r_S}A_1(j)\bigr)^2$.
\end{lemma}

\begin{proof}
By (Irr), pairs with an irrelevant entry contribute $0$; drop them.  By (J) and $j(t)\le r_S$ (as $\dim_EI(t)\ge0$), the
remaining pairs are partitioned according to $(j(t),j(t'))\in[1,r_S]^2$.  Fix $(j',j'')$ and let $Q'$ be the corresponding
part.  Partition $Q'$ further by the labels $(\mathrm{lab}(t),\mathrm{lab}(t'))\in T'\times T''$, where $T'$, $T''$ are the
(finite) sets of first and second labels occurring.  (Ty) bounds the part with labels $(\delta',\delta'')$ by
$d_\pi m\,\kappa(j',\delta')\kappa(j'',\delta'')$; summing over $T'\times T''$ and using Lemma~\ref{i3:lem-labelsum} twice
(the summands are $\ge0$), the part $Q'$ contributes at most $d_\pi mA_1(j')A_1(j'')$.  Summing over $(j',j'')$ gives the
claim.
\end{proof}

\begin{proof}[Proof of Proposition~\ref{i3:prop-kpp-diag}]
\emph{Constants.}  At the outer constants $\mathbf c$ of the family, let $C_H,C_h,C_p$ be given by (Ty), $C_s=C_p+4$,
$C^{(0)}_T$ the constant of (Ty) at $C_s$, $C^{(0)}_\Delta$ that of (Del) at $C_H$, $C=\max(C^{(0)}_T,C^{(0)}_\Delta,0)$, $C_c$ that
of (Cos), and $C_f=1+C+C^2+C_h^2+C_c^2\ge1$.  The \emph{Kloosterman-diagonal threshold} is $\max(3,C'_0,C_{l,0})$, imposed on $C_{k,\rm low}$ and $C_l^{\mathcal F}$, where $(C'_0,C_{l,0})$ is the
coordinatewise maximum of the thresholds of (Ty) at $C$, (Del) at $C$, (Cos), of the bound $P_g\le2$ and, for $C_l$, of
$c_0(C_f)$ from (Num).  Let $\mathcal F$ be large, $x$ beyond the selection thresholds, and a member satisfy the standing
hypotheses at $(3,C',C_l)$, $C'=C_{k,\rm low}$, $C_l=C_l^{\mathcal F}$; as $C'\ge1$ these give the $K_2$-hypotheses at
$(C',C_l)$, the package is concretely tied, and the member satisfies the ties (a record of every family of the
class).  Also $k\ge m+3=2r_S+5\ge r_S+3$.

\emph{Summability.}  By Lemma~\ref{i3:lem-finbound}, the nonnegative family $\|\mathrm{vec}_t\|\|\mathrm{vec}_{t'}\|$ on
$\mathrm{all}(\mathrm{deg})$ has all finite partial sums at most $\mathsf B=d_\pi m(\sum_jA_1(j))^2$; it is therefore summable with
sum at most $\mathsf B$.

\emph{Chain of inequalities.}  By $|\langle a,b\rangle|\le\|a\|\|b\|$,
$|\Sigma|(\mathrm{all}(\mathrm{deg}))\le\Gamma_1\mathsf B\le d_u(\sum_jA_1(j))^2$ by (Dim).  Next $A_1(j)\le A(C_f;j)$: since
$0\le C\le C_f$ and $j(k-1-j)f\ge0$, $(C/k)^{j(k-1-j)f}\le(C_f/k)^{j(k-1-j)f}$; and the product of the four powers of $H_b$ is
$b^{m^2f(C^2r_S^4+C_h^2r_S^4+C(j+1)/2+C_c^2r_S)}\le b^{C_fm^6f}=e^{C_fm^6fL}$, because $r_S\le m$, $(j+1)/2\le m^4$ and
$C^2+C_h^2+C+C_c^2\le C_f$.  Hence $(\sum_jA_1(j))^2\le B(C_f)$.  The bound $P_g\le2$ holds above its threshold (it also follows
directly from $(kb)^8\le g_j$ and $f\le k$, as in the proof of Lemma~\ref{i3:lem-diagnum}: $P_g\le e^{6b^{-7}}\le2$).  So
$P_g|\Sigma|(\mathrm{all}(\mathrm{deg}))\le2d_uB(C_f)$.  Finally the $K_2$-hypotheses contain $C_lm^6L\le k\log k$, and
$C_l\ge c_0(C_f)$, so (Num) with $d=d_u$ gives $2d_uB(C_f)\le d_ue^{-kf\log k}$.
\end{proof}

\subsubsection{Proof of Proposition~\ref{i3:prop-W}}\label{i3:ssec-pf-W}

Let $K$ be a finite field, $g=|K|$, and $N\ge1$.  Let $\mathrm{Sq}_N\subset K^N$ be the set of $c$ with $P_c$ squarefree and
$P_c(0)\ne0$, and for $Q\subset K^N$ write $\mathrm{sq}(Q)=\sum_{c\in Q\cap\mathrm{Sq}_N}\mathrm{sq}(P_c)$.

\begin{proposition}[(I) proved in Lemma~\ref{i3:lem-sqtype}; (II)--(V) formally verified, \cite{CharacterFamily}]\label{i3:prop-W-loc}
Let $\bar K$ be an algebraic closure of $K$.
\begin{enumerate}
\item[(I)] If $N\ge1$, then $\mathrm{sq}(K^N)\ge1-2/g$.
\item[(II)] If $g\ge2^{12}$, then the $\mathrm{sq}$-mass of the $c\in\mathrm{Sq}_N$ for which some prime $l\notin B(g,N)$ has
 $x\mapsto$ (prime-to-$l$ part of $x$) not injective on the roots of $P_c$ in $\bar K$ is at most $23\log g/\sqrt g$.
\item[(III)] If $N\ge1$, $q\ne d$ are odd primes with $\max(q,d)\ge16$, $16qd\mid g-1$, and $c_0\in K^\times$, then the
 $\mathrm{sq}$-mass of the $c\in\mathrm{Sq}_N$ for which some prime $l\ne g$ and some root $x$ of $P_c$ have the same
 prime-to-$l$ part as $c_0$ is at most $\frac1{4\min(q,d)}+\frac{1+\log N}{g-1}$.
\item[(IV)] If $0<a<N$, the $\mathrm{sq}$-mass of the $c\in\mathrm{Sq}_N$ for which $P_c$ has a divisor of degree $a$ is at
 most $\frac78$.
\item[(V)] If $N\ge100$, the $\mathrm{sq}$-mass of the $c\in\mathrm{Sq}_N$ for which $P_c$ has no irreducible factor of prime
 degree $p$ with $N<2p$ and $p+3\le N$ is at most $1-\frac1{5\log N}$.
\end{enumerate}
\end{proposition}

\noindent (IV)
and (V) compare the degree types of squarefree polynomials, weighted by $\mathrm{sq}$, with the cycle types of random
permutations (the $\mathrm{sq}$-mass of a set of degree types is at most the proportion of
permutations of $\{1,\dots,N\}$ with those cycle types), and then count permutations with an invariant set of size $a$
($\le\frac78N!$) or without a cycle of prime length in $(N/2,N-3]$.  The whole is a long verification.
We prove the comparison with permutations and item (I) here; the formally verified steps of the rest are outlined in \S\ref{i3:str-Wloc}.

\begin{lemma}[Squarefree weights]\label{i3:lem-sqtype}
Let $K$ be a finite field with $g$ elements, $N\ge1$, and for $d\ge1$ let $I_d$ be the number of monic irreducible polynomials of degree
$d$ over $K$ with nonzero constant term.  For $c\in\mathrm{Sq}_N$ let the \emph{type} of $P_c$ be $(n_d)_{d\ge1}$, $n_d$ the number of its
irreducible factors of degree $d$.
\begin{enumerate}
\item[(a)] For every $(n_d)$ with $\sum_ddn_d=N$, the $\mathrm{sq}$-mass of the $c\in\mathrm{Sq}_N$ of type $(n_d)$ is
$\prod_d\binom{I_d}{n_d}(g^d-1)^{-n_d}$.
\item[(b)] For every set $A$ of types, the $\mathrm{sq}$-mass of the $c\in\mathrm{Sq}_N$ whose type lies in $A$ is at most the proportion of
permutations of $\{1,\dots,N\}$ whose cycle type lies in $A$.
\item[(c)] $\mathrm{sq}(K^N)\ge1-2/g$; this is item (I) of Proposition~\ref{i3:prop-W-loc}.
\end{enumerate}
\end{lemma}

\begin{proof}
(a) $c\mapsto P_c$ is a bijection of $K^N$ onto the monic polynomials of degree $N$.  By unique factorisation, a squarefree $P$ with
$P(0)\ne0$ of type $(n_d)$ is the same as a choice, for each $d$, of a set of $n_d$ distinct monic irreducibles of degree $d$ with
nonzero constant term, and then $\mathrm{sq}(P)=\prod_d(g^d-1)^{-n_d}$.  Summing over the choices gives the product.

(b) The roots of the polynomials counted by $I_d$ are elements of $\mathbb F_{g^d}^\times$, $d$ for each polynomial, and distinct
polynomials have disjoint sets of roots; so $dI_d\le g^d-1$.  Hence
$\binom{I_d}{n_d}(g^d-1)^{-n_d}\le I_d^{n_d}/(n_d!\,(g^d-1)^{n_d})\le1/(d^{n_d}n_d!)$.  By Cauchy's formula the number of permutations
of $\{1,\dots,N\}$ of cycle type $(n_d)$ is $N!/\prod_dd^{n_d}n_d!$.  Sum (a) over the types in $A$.

(c) For $c\in\mathrm{Sq}_N$, $\mathrm{sq}(P_c)=\prod_j(g^{\deg\Phi_j}-1)^{-1}\ge\prod_jg^{-\deg\Phi_j}=g^{-N}$.  Let $s_M$ be the number of
monic squarefree polynomials of degree $M$ ($s_0=1$).  Every monic $f$ of degree $M$ is uniquely $a^2b$ with $a,b$ monic and $b$
squarefree ($b$ is the product of the irreducible factors of odd multiplicity), so $g^M=\sum_{0\le k\le M/2}g^ks_{M-2k}$.  For $M\ge2$,
subtracting $g$ times the same identity for $M-2$ gives $s_M=g^M-g^{M-1}$; and $s_1=g$.  The squarefree monic $P$ of degree $N$ with
$P(0)=0$ are $XQ$ with $Q$ squarefree monic of degree $N-1$, at most $s_{N-1}\le g^{N-1}$ of them.  Hence $|\mathrm{Sq}_N|\ge
g^N-2g^{N-1}$ (for $N=1$, $|\mathrm{Sq}_1|=g-1$), and $\mathrm{sq}(K^N)\ge(g^N-2g^{N-1})g^{-N}=1-2/g$.
\end{proof}

\noindent By (b), items (IV) and (V) follow from the corresponding statements on permutations of $\{1,\dots,N\}$: the proportion of
permutations with an invariant set of size $a$ ($0<a<N$) is at most $\frac78$, and for $N\ge100$ the proportion of permutations
with a cycle of prime length $p$, $N<2p$, $p+3\le N$, is at least $\frac1{5\log N}$.  These two statements, and items (II), (III),
are formally verified.

\begin{proof}[Proof of Proposition~\ref{i3:prop-W}]
Write $\mu=P^{\otimes\iota}$, i.e.\ $\mu(x)=\prod_vP(x_v)$ for $x=(x_v)\in(K^N)^\iota$; $\mu\ge0$ and every set of the form
$\{x:x_{v}\in Q\}$ has mass $P(Q)=\sum_{c\in Q}P(c)$, since the other coordinates contribute $\prod_{w\ne v}\sum P=1$.
On $\mathrm{Sq}_N$, $P=\mathrm{sq}$, so $P(Q)=\mathrm{sq}(Q)$ for $Q\subset\mathrm{Sq}_N$, and by (I)
$P(K^N\setminus\mathrm{Sq}_N)=1-\mathrm{sq}(K^N)\le2/g$.

Let $\mathcal S=\{x:x_v\in\mathrm{Sq}_N\ \forall v\}$.  A tuple $x$ for which $(P_{x_v})_v$ is not good either lies outside
$\mathcal S$, or lies in $\mathcal S$ and violates one of (G3)--(G6) ((G1) and (G2) hold on $\mathcal S$; every $P_{x_v}$ is
monic of degree $N$).  We bound the five contributions.

\emph{Outside $\mathcal S$.}  $\mu(\mathcal S^c)\le\sum_v\mu(x_v\notin\mathrm{Sq}_N)\le2|\iota|/g$.

\emph{(G3) fails.}  Then for some $v$, $x_v$ lies in the exceptional set of (II).  By the union bound over $v$ and (II), the
mass is at most $|\iota|\cdot23\log g/\sqrt g$.

\emph{(G4) fails.}  Then for some $v$, $j$ and some sign, $x_v$ lies in the exceptional set of (III) for
$c_0=\chi(v,j)^{\pm1}$.  The union bound over $v$, $j$ and the two signs gives at most
$2|\iota||J|\bigl(\frac1{4\min(q,d)}+\frac{1+\log N}{g-1}\bigr)$.

\emph{(G5) fails.}  Then for some $a$ with $1\le a$, $2a\le N$ (so $0<a<N$), every $P_{x_v}$ has a divisor of degree $a$.  For
fixed $a$ this set of $x\in\mathcal S$ is a product of sets, of mass $\prod_v\mathrm{sq}(\cdot)\le(7/8)^{|\iota|}$ by (IV); there
are at most $N/2$ values of $a$.

\emph{(G6) fails.}  Then every $P_{x_v}$ lacks an irreducible factor of prime degree $p$ with $N<2p$, $p+3\le N$; this
product set has mass at most $(1-\frac1{5\log N})^{|\iota|}\le e^{-|\iota|/(5\log N)}$ by (V) and $1-y\le e^{-y}$.

Adding the five bounds gives the claim.
\end{proof}

\subsubsection{Proof of Proposition~\ref{i3:prop-orbbd}}\label{i3:ssec-pf-orbbd}

This subsection is self-contained.  The subgroups
$N$, $U_z$ and the Lie algebras $\mathfrak u$, $\mathfrak u_z$, $\mathfrak n$ are those fixed in the conventions paragraph
below (they agree with \S\ref{i3:ssec-bessel}).  Throughout, $k\ge1$.

\noindent\emph{Notation.}  $H_q$ is the Hermitian matrix with $(H_q)_{ij}=[i+j=2r+k]$ for ``outer'' $i$
($i\le r$ or $i\ge r+k$) and $(H_q)_{ij}=[i=j]$ otherwise; $H_q^2=1$, $H_q^*=H_q$.  $N_q\subset U(H_q)$ is
the analogue of $N$ for $H_q$, and $W_q\subset U(H_q)$ is the pointwise fixer of the $e_a,e^*_a$ ($a<r$), a copy of
$U(k,1)$ acting on the middle block (it does not fix $e_{w_0}-e_{w_L}$ in general); $w_0=r$, $w_L=r+k$;
$e_a=a$, $e^*_a=2r+k-a$ ($a<r$).  A \emph{torus vector} is $d\in\R_{>0}^n$ with $d_id_j=1$ for outer $i$ and
$i+j=2r+k$, $d_i=1$ for non-outer $i$, and $d_{w_0}\le1$; $\mathrm{dg}(d)$ is the diagonal matrix.  An
\emph{adapted basis} for $z$ is $\Gamma$ with $\Gamma^*H_{k,r}\Gamma=H_q$, $\Gamma e_a=e_a$, $\Gamma e^*_a=e^*_a$ and
$z=\rho\,\Gamma(e_{w_0}-e_{w_L})$, $\rho>0$.

\begin{lemma}[Iwasawa decomposition]\label{i3:ob-iwa}
Let $k\ge1$.  Every $x\in U(H_q)$ is $x=s\kappa$ with $s\in U(H_q)$ upper triangular and $\kappa\in U(H_q)$ unitary.
\end{lemma}
\begin{proof}
(This is the Iwasawa decomposition $G=PK$ for the minimal parabolic stabilising the standard isotropic
flag, cf.\ A.~W.~Knapp, \emph{Lie Groups Beyond an Introduction}, 2nd ed., Birkh\"auser 2002, Thm.~6.46; we
give the matrix proof.)  Put $A=H_q$; $A$ is the permutation matrix of the involution $p$ reversing the outer
indices and fixing the others, and $p$ reverses the grading $b(i)=2i-(2r+k)$ (outer), $b(i)=0$ (otherwise).
Let $g=Ax^*A$; then $g\in U(A)$ and $g$ is invertible with inverse $Ag^*A$ (as $g^*Ag=A$, $A^2=1$).  The
matrix $g^*g$ is positive definite, so $g^*g=B^*B$ with $B$ upper triangular with positive diagonal (Cholesky).
From $gAg^*=A$ one gets $(B^*B)A(B^*B)=A$, hence $W=B(AB^*A)$ satisfies $W^*W=1$.  Since $p$ reverses the
grading on the outer indices and fixes the middle ones, $AB^*A$ is block upper triangular for the partition of the
indices into the singletons $\{i\}$, $i$ outer, and the middle block, with positive outer diagonal entries and middle
block $B_{\rm mid}^*$ (the middle block of $B$ being $B_{\rm mid}$, upper triangular with positive diagonal).  Hence $W$ is block
upper triangular for the same partition, and unitary, so block diagonal.  Its outer diagonal entries are positive of
modulus $1$, so equal $1$; its middle block $B_{\rm mid}B_{\rm mid}^*$ is positive definite and unitary, hence $1$, so $B_{\rm mid}$ is a unitary upper triangular matrix
with positive diagonal, hence $B_{\rm mid}=1$.  Thus $W=1$, and $B^{-1}=AB^*A$ and $B^*AB=A$, and $k_0=gB^{-1}$ is unitary and lies
in $U(A)$; being unitary and $A$-unitary, $k_0$ commutes with $A$.  Hence
$x=Ag^*A=AB^*k_0^*A=(AB^*A)\,k_0^*$, and $s=AB^*A=B^{-1}$ (upper triangular, in $U(A)$) and $\kappa=k_0^*$ are as
required.
\end{proof}

\begin{lemma}[Levi decomposition]\label{i3:ob-levi}
Let $k\ge1$.  Every upper triangular $s\in U(H_q)$ is $s=m\,\mathrm{dg}(t)\,\mu\,\kappa$ with $m\in N_q$, $t$ a torus
vector with $t_{w_0}=1$, $\mu\in W_q$ and $\kappa\in U(H_q)$ unitary.
\end{lemma}
\begin{proof}
Let $d_i=1$ for $r\le i\le r+k$ and $d_i=s_{ii}$ otherwise.  The $(i,p(i))$ entry of $s^*H_qs=H_q$, together
with triangularity, gives $\bar d_id_{p(i)}=1$.  Let $\mu$ be the block of $s$ on indices $r\le i,j\le r+k$,
extended by the identity; this block preserves the restriction of $H_q$ (the off-block contributions vanish
by triangularity), so $\mu\in U(H_q)$, and $\mu$ fixes $e_a,e^*_a$, i.e.\ $\mu\in W_q$.  Put $t_i=|d_i|$ (a torus
vector with $t_{w_0}=1$) and $\kappa=\mathrm{diag}(d_i/|d_i|)$, unitary and in $U(H_q)$ because $\bar d_id_{p(i)}=1$.
Then $m=s(\mathrm{dg}(t)\mu\kappa)^{-1}$ lies in $U(H_q)$, is upper triangular with unit diagonal outside the middle block
and identity on the middle block, i.e.\ $m\in N_q$.
\end{proof}

\begin{lemma}[Cartan decomposition of $W_q$]\label{i3:ob-cartan}
Let $k\ge1$.  Every $\mu\in W_q$ is $\mu=h\,\mathrm{dg}(t)\,\kappa$ with $h\in W_q$ fixing $v_0=e_{w_0}-e_{w_L}$, $t$ a
torus vector with $t_{e_a}=1$ for all $a<r$, and $\kappa\in U(H_q)$ unitary.
\end{lemma}
\begin{proof}
(Cf.\ Knapp, \emph{op.~cit.}, Thm.~7.39, for $W_q\cong U(k,1)$.)  We have $\langle v_0,v_0\rangle_{H_q}=-2$.  Let
$y=\mu^{-1}v_0$; then $\langle y,y\rangle_{H_q}=-2$, and its Euclidean norm satisfies $N=\|y\|^2\ge2$, since
$2=-\langle y,y\rangle_{H_q}\le\|y\|\,\|H_qy\|=\|y\|^2$.  The function $u\mapsto u^{-2}+u^2$ is continuous on $(0,1]$,
equals $2$ at $u=1$ and tends to $\infty$ at $0$, so there is $u\in(0,1]$ with $u^{-2}+u^2=N$.  Let
$t=\mathrm{tor}(u)$ ($t_{w_0}=u$, $t_{w_L}=u^{-1}$, $1$ elsewhere) and $x=\mathrm{dg}(t)^{-1}v_0=u^{-1}e_{w_0}-ue_{w_L}$; $x$ and
$y$ have equal Euclidean norms and equal $H_q$-norms.  By Lemma~\ref{i3:ob-witt2} there is a unitary $K$
commuting with $H_q$ (so $K\in U(H_q)$) with $Kx=y$ and fixing every vector orthogonal to $x,y,H_qx,H_qy$,
in particular the $e_a,e^*_a$.  Put $\kappa=K^{-1}$ and $h=\mu\kappa^{-1}\mathrm{dg}(t)^{-1}$.  Then $h\in W_q$ and
$hv_0=\mu Kx=\mu y=v_0$.
\end{proof}

\begin{lemma}[Simultaneous Witt lemma]\label{i3:ob-witt2}
Let $H$ be Hermitian with $H^2=1$, and $x,y\in\C^n$ with $\|x\|=\|y\|$ and $\langle x,x\rangle_H=\langle y,y\rangle_H$.
There is a unitary $K$ with $HK=KH$, $Kx=y$ and $Kw=w$ for all $w$ orthogonal to $x,y,Hx,Hy$.
\end{lemma}
\begin{proof}
Write $x=x_++x_-$, $y=y_++y_-$ with $x_\pm=\frac12(x\pm Hx)$ in the $\pm1$ eigenspaces of $H$.  The two
hypotheses give $\|x_+\|^2+\|x_-\|^2=\|y_+\|^2+\|y_-\|^2$ and $\|x_+\|^2-\|x_-\|^2=\|y_+\|^2-\|y_-\|^2$, hence
$\|x_\pm\|=\|y_\pm\|$.  In each eigenspace choose a unitary $K_\pm$ with $K_\pm x_\pm=y_\pm$ which is the identity
on the orthogonal complement of $\{x_\pm,y_\pm\}$ (a Householder reflection in $x_\pm-y_\pm$ composed with a phase
if $\langle x_\pm,y_\pm\rangle\notin\R$; identity if $x_\pm=y_\pm$).  Then $K=K_+\oplus K_-$ commutes with $H$, $Kx=y$, and
fixes every vector orthogonal to $x_\pm,y_\pm$, i.e.\ to $x,y,Hx,Hy$.
\end{proof}

\begin{lemma}[Decomposition adapted to the Bessel subgroup]\label{i3:ob-Decomp}
Let $k\ge1$, $z$ supported on the middle block with $T(z)>0$ and $\Gamma$ an adapted basis.  Every $g\in U(H_{k,r})$ can be written
$g=h\,m\,(\Gamma\,\mathrm{dg}(d)\,R)$ with $h\in U_z$, $m\in N$, $d$ a torus vector, $R$ unitary and
$\Gamma\,\mathrm{dg}(d)R\in U(H_{k,r})$.
\end{lemma}
\begin{proof}
Let $\Gamma_0$ be unitary with $\Gamma_0^*H_{k,r}\Gamma_0=H_q$ (Lemma~\ref{i3:ob-bmu-struct}(a)) and $\Gamma'=\Gamma^{-1}$.  Then
$x=\Gamma'g\Gamma_0\in U(H_q)$.  Lemmas~\ref{i3:ob-iwa}, \ref{i3:ob-levi}, \ref{i3:ob-cartan} give $x=s\kappa_1$,
$s=m_0\,\mathrm{dg}(t)\,\mu\,\kappa_2$, $\mu=h_0\,\mathrm{dg}(t')\,\kappa_3$.  As $t$ is trivial on the middle block and $h_0\in W_q$
acts only there, $\mathrm{dg}(t)h_0=h_0\mathrm{dg}(t)$; products of torus vectors are torus vectors.  Put $h=\Gamma h_0\Gamma'$,
$m=\Gamma h_0^{-1}m_0h_0\Gamma'$, $d=tt'$, $R=\kappa_3\kappa_2\kappa_1\Gamma_0^*$.  Then
$g=\Gamma x\Gamma_0^*=\Gamma m_0h_0\,\mathrm{dg}(tt')R=h\,m\,(\Gamma\mathrm{dg}(d)R)$.  $h\in U_z$ since $h_0$ fixes $e_a,e^*_a$ and
$e_{w_0}-e_{w_L}$, which $\Gamma$ maps to $e_a,e^*_a,\rho^{-1}z$; $m\in N$ since conjugation by $\Gamma h_0^{\pm1}$ maps
$N_q$ to $N$: $\Gamma$ and $h_0$ fix the $e_a,e^*_a$, and moreover their outer rows (and those of $\Gamma^{-1}$, $h_0^{-1}$)
are identity rows, for $\Gamma$ by Lemma~\ref{i3:ob-B-frame} and for $h_0$ because it acts only on the middle block; so
the defining conditions of $N$ on outer rows and columns are preserved.  $R$ is a product of unitaries.
\end{proof}

\noindent\emph{Conventions for the rest of this subsection.}  Indices run over $\{0,\dots,n-1\}$,
$n=2r+k+1$; $e_i$ is the $i$-th standard basis vector; for $a<r$ we use the indices $a$ and $a^*=2r+k-a$ (so the
index written $e_a$ above is $a$), and we put $a+1:=r=w_0$ when $a=r-1$.  An index $i$ is \emph{strictly outer} if
$i<r$ or $i>r+k$.  $H=H_{k,r}$ has $H_{a a^*}=H_{a^* a}=1$ ($a<r$), $H_{ii}=1$ for $r\le i<r+k$, $H_{r+k,r+k}=-1$,
and all other entries $0$; $h(x,y)=y^*Hx$, so $T(z)=-\Re\,h(z,z)$, and $z$ lies in the prescribed subspace iff
$z_i=0$ for every strictly outer $i$.  $p$ is the involution with $p(i)=2r+k-i$ for outer $i$ and $p(i)=i$
otherwise, so that $(H_q)_{ij}=[j=p(i)]$.  For a matrix $X$ we write $\|X\|=\max_{i,j}|X_{ij}|$.  We put
\[
\mathfrak u=\{X:X^*H+HX=0\},\qquad \mathfrak u_z=\{X\in\mathfrak u:\ Xe_a=Xe_{a^*}=0\ (a<r),\ Xz=0\},
\]
$\mathfrak n=\{X\in\mathfrak u:\ X_{ia}=0$ whenever $a<r$ and $i\ge a$, and $X_{ij}=0$ whenever $r\le j\le r+k$ and
$i\ge r\}$; correspondingly
$N$ is the set of $g\in U$ with $g_{ia}=\delta_{ia}$ ($a<r$, $i\ge a$) and $g_{ij}=\delta_{ij}$ ($r\le j\le r+k$, $i\ge r$),
and $U_z=\{g\in U:\ ge_a=e_a,\ ge_{a^*}=e_{a^*}\ (a<r),\ gz=z\}$.  A position $(i,j)$ is \emph{upper} if $i<j$ and
$i$ or $j$ is outer, and \emph{central} if $r\le i,j\le r+k$.

Lemmas~\ref{i3:ob-B-calc}, \ref{i3:ob-B-chart}, \ref{i3:ob-B-reduce}, \ref{i3:ob-B-qbound},
\ref{i3:ob-B-abs} hold for $U=U(H_0)=\{g:g^*H_0g=H_0\}$ and
$\mathfrak u=\{X:X^*H_0+H_0X=0\}$ for any matrix $H_0$ with $H_0^2=1$ (for us $H_0=H_{k,r}$).  A unitary
representation $\sigma$ of $U$ on a Hilbert space $H$ is a map to bounded operators, multiplicative on $U$, with
$\sigma(1)=1$ and $\|\sigma(g)v\|=\|v\|$; $v$ is \emph{smooth} if every point of $U$ has an open neighbourhood $O$ in
$M_n(\C)$ and a $C^\infty$ map $F:O\to H$ with $F(g)=\sigma(g)v$ for $g\in O\cap U$.  For $X\in\mathfrak u$,
$e^{tX}\in U$, and $d\sigma(X)v=\frac{d}{dt}\sigma(e^{tX})v|_{t=0}$; for a word $W=(X_1,\dots,X_l)$ in $\mathfrak u$,
$d\sigma(W)=d\sigma(X_1)\cdots d\sigma(X_l)$, and $WY$ denotes the word with the letter $Y$ appended, so that
$d\sigma(WY)=d\sigma(W)d\sigma(Y)$.  Smooth vectors form a subspace stable under $\sigma(U)$ and the $d\sigma(X)$, and
$d\sigma(W)$ is linear on it.

\begin{lemma}[Structural inputs for boundedness]\label{i3:ob-bmu-struct}
Let $k\ge1$ and $z$ supported on the middle block with $T(z)>0$.
(a) An adapted basis $\Gamma$ exists, and there is a unitary $\Gamma_0$ with $\Gamma_0^*H_{k,r}\Gamma_0=H_q$.
(b) $\mathfrak u=\mathfrak u_z+\mathfrak n+\mathfrak q_\Gamma$, where $\mathfrak q_\Gamma=\{X\in\mathfrak u:(\Gamma^{-1}X\Gamma)_{ij}=0$ for
every upper position $(i,j)\}$; and $\exp(tX)\in U_z$ (resp.\ $N$) for $X\in\mathfrak u_z$ (resp.\ $\mathfrak n$).
(c) If $d$ is a torus vector, $R$ unitary, $M\ge1$ with $d_a\le Md_{a+1}$ for all $a<r$, and $A_{ij}=0$
for every upper position $(i,j)$, then $\|(\mathrm{dg}(d)R)^{-1}A(\mathrm{dg}(d)R)\|\le n^2\|A\|M^n$.
(d) For $c\in(\C^\times)^r$ there is $C_1$ such that for every $a<r$, torus vector $d$ and unitary $R$ with
$g'=\Gamma\mathrm{dg}(d)R\in U$ there is $Y\in\mathfrak u$ with $\|Y\|\le C_1$, $g'\exp(\tau Y)g'^{-1}\in N$ and
$\psi_{z,c}(g'\exp(\tau Y)g'^{-1})=\exp(i\tau d_a/d_{a+1})$ for all $\tau\in\R$.
(e) (Quotient bound) If $B:H\to W$ is Casselman--Wallach continuous and linear on smooth vectors, $X_j,Q_i\in\mathfrak u$
with $\mathfrak u=\mathrm{span}_\R\{X_j\}+\mathrm{span}_\R\{Q_i\}$, and $B(\sigma(\exp tX_j)v)=\chi_j(t)B(v)$ for smooth $v$, with
$\chi_j:\R\to\mathrm{End}(W)$ continuous and $\chi_j(0)=1$, then there are $C\ge0$ and $m$ with $\|B(v)\|\le C\rho$ whenever $v$ is
smooth and $\|d\sigma(Q_P)v\|\le\rho$ for all words $P$ in the indices $i$ of length $\le m$.
(f) (Absolute bound) In the situation of (e), for $C_0\ge0$, $c\in\N$ and smooth $w$ there is $K$ such that
$\|B(\sigma(g)w)\|\le K$ for every $g\in U$ and $M\ge1$ with $\|g^{-1}Q_ig\|\le C_0M^c$ for all $i$ and either $M=1$ or
there are $Y\in\mathfrak u$, $\|Y\|\le C_0$, and $\ell\in\R$, $|\ell|\ge M$, with $B(\sigma(g\exp\tau Y)u)=e^{i\ell\tau}B(\sigma(g)u)$ for
all $\tau$ and smooth $u$.
\end{lemma}
\begin{proof}
(a) is Lemma~\ref{i3:ob-B-adapt}; (b) is Lemmas~\ref{i3:ob-B-span} and \ref{i3:ob-B-exp}; (c) is Lemma~\ref{i3:ob-B-tor};
(d) is Lemma~\ref{i3:ob-B-root}; (e) is Lemma~\ref{i3:ob-B-qbound}; (f) is Lemma~\ref{i3:ob-B-abs} applied with the
constants $C,m$ of (e).
\end{proof}

\begin{proof}[Proof of Proposition~\ref{i3:prop-orbbd}]
Fix an adapted basis $\Gamma$ (Lemma~\ref{i3:ob-bmu-struct}(a)) and finite real bases $X_j$ of $\mathfrak u_z+\mathfrak n$
(the union of bases of $\mathfrak u_z$ and $\mathfrak n$) and $Q_i$ of $\mathfrak q_\Gamma$.  For $X_j\in\mathfrak u_z$ put
$\chi_j(t)=\pi(\exp tX_j)$; for $X_j\in\mathfrak n$ put $\chi_j(t)=\psi_{z,c}(\exp tX_j)\cdot1$.  By
Lemma~\ref{i3:ob-bmu-struct}(b) and equivariance (with $\psi_{z,c}(1)=1$, $\pi(1)=1$), $B(\sigma(\exp tX_j)v)=\chi_j(t)B(v)$;
$\chi_j$ is continuous because $\pi$ is continuous on $U_z$ and $\psi_{z,c}$ is continuous.
Lemma~\ref{i3:ob-bmu-struct}(e) gives $C,m$; let $C_1$ be as in (d), $C_2=n^2\sum_i\|\Gamma^{-1}Q_i\Gamma\|$, and let $K$
be given by (f) with $C_0=\max(C_1,C_2)$ and $c=n$.

Let $g\in U$ and write $g=hm\,g'$, $g'=\Gamma\mathrm{dg}(d)R$ (Lemma~\ref{i3:ob-Decomp}).  Equivariance,
$|\psi_{z,c}|=1$ and the isometry of $\pi$ give $\|B(\sigma(g)w)\|=\|B(\sigma(g')w)\|$.  If $M\ge1$ and
$d_a\le Md_{a+1}$ for all $a$, then by (c) applied to $A=\Gamma^{-1}Q_i\Gamma$,
$\|g'^{-1}Q_ig'\|\le C_0M^n$.  If all ratios $d_a/d_{a+1}$ are $\le1$, apply (f) with $M=1$.  Otherwise let
$M\ge1$ be the largest ratio, attained at $a_0$, and take $Y$ from (d) for $a_0$: then for smooth $u$,
$\sigma(g'\exp\tau Y)u=\sigma(g'\exp(\tau Y)g'^{-1})\sigma(g')u$ and equivariance with $h=1$ gives
$B(\sigma(g'\exp\tau Y)u)=e^{iM\tau}B(\sigma(g')u)$; apply (f).  In both cases $\|B(\sigma(g')w)\|\le K$.
\end{proof}

\begin{lemma}[Hermitian reflections]\label{i3:ob-B-move}
Let $H'$ be a Hermitian matrix, $h'(x,y)=y^*H'x$, and $u,z$ with $h'(u,u)=h'(z,z)\neq0$.  There is $\gamma$ with
$\gamma^*H'\gamma=H'$, $\gamma u=z$, and $\gamma x=x$ for every $x$ with $h'(x,u)=h'(x,z)=0$.
\end{lemma}
\begin{proof}
\emph{Step 1.} Suppose $c:=h'(u,u)-h'(u,z)\neq0$.  Put $v=u-z$ and $\gamma x=x-c^{-1}h'(x,v)\,v$.  Since $h'(u,u)=h'(z,z)$ is
real, $h'(v,v)=2h'(u,u)-h'(u,z)-\overline{h'(u,z)}=c+\bar c$, and $c=h'(u,v)$, so $\gamma u=u-v=z$.  For $x,y$,
expanding with $\overline{h'(y,v)}=h'(v,y)$,
\[
h'(\gamma x,\gamma y)=h'(x,y)-\big(\bar c^{-1}+c^{-1}-c^{-1}\bar c^{-1}(c+\bar c)\big)h'(x,v)h'(v,y)=h'(x,y),
\]
so $\gamma^*H'\gamma=H'$; and $\gamma x=x$ if $h'(x,v)=h'(x,u)-h'(x,z)=0$.
\emph{Step 2.} If $h'(u,u)-h'(u,z)=0$, apply Step 1 to $(u,-u)$ (here $c=2h'(u,u)\ne0$) and then to $(-u,z)$ (here
$c=h'(u,u)+h'(u,z)=2h'(u,u)\neq0$), and compose; vectors $h'$-orthogonal to $u$ and $z$ are fixed by both factors.
\end{proof}

\begin{lemma}[Adapted bases]\label{i3:ob-B-adapt}
Let $k\ge1$.  The matrix $\Gamma_0$ whose columns are $e_j$ ($j\ne r,r+k$), $(e_r+e_{r+k})/\sqrt2$ (column $r$) and
$(e_r-e_{r+k})/\sqrt2$ (column $r+k$) is unitary with $\Gamma_0^*H\Gamma_0=H_q$.  If $z_i=0$ for every strictly outer $i$ and
$T(z)>0$, an adapted basis $\Gamma$ for $z$ exists.
\end{lemma}
\begin{proof}
The columns of $\Gamma_0$ are orthonormal.  The entry $(\Gamma_0^*H\Gamma_0)_{ij}$ is $h(\mathrm{col}_j,\mathrm{col}_i)$.  If
$i,j\notin\{r,r+k\}$ it is $H_{ij}=(H_q)_{ij}$ (on these indices $H$ and $H_q$ agree).  Since $H$ restricted to
$\mathrm{span}(e_r,e_{r+k})$ is $\mathrm{diag}(1,-1)$ and $H$ does not mix this span with the other $e_j$,
$h(\mathrm{col}_r,\mathrm{col}_r)=\frac12(1-1)=0=h(\mathrm{col}_{r+k},\mathrm{col}_{r+k})$,
$h(\mathrm{col}_{r+k},\mathrm{col}_r)=h(\mathrm{col}_r,\mathrm{col}_{r+k})=\frac12(1+1)=1$, and the mixed entries with other
indices vanish; this is $H_q$.

Let $\rho=\sqrt{T(z)/2}>0$ and $u=\sqrt2\rho\,e_{r+k}$.  Then $h(u,u)=-2\rho^2=-T(z)$, and $h(z,z)$ is real (as $H$ is
Hermitian) with real part $-T(z)$, so $h(u,u)=h(z,z)\neq0$.  Lemma~\ref{i3:ob-B-move} gives $\gamma\in U$ with $\gamma u=z$
fixing every $x$ with $h(x,u)=h(x,z)=0$.  For strictly outer $i$, $h(e_i,u)=h(e_i,z)=0$, because the row of $H$ at a
non-strictly-outer index has support in the non-strictly-outer indices while $u,z$ vanish at strictly outer ones.  Put
$\Gamma=\gamma\Gamma_0$.  Then $\Gamma^*H\Gamma=\Gamma_0^*H\Gamma_0=H_q$; $\Gamma e_i=\gamma e_i=e_i$ for strictly outer $i$ (in
particular $i=a,a^*$); and since $\mathrm{col}_r-\mathrm{col}_{r+k}=\sqrt2e_{r+k}$,
$\rho\,\Gamma(e_{w_0}-e_{w_L})=\gamma(\sqrt2\rho e_{r+k})=\gamma u=z$.
\end{proof}

\begin{lemma}[Rows and columns of adapted bases]\label{i3:ob-B-frame}
Let $\Gamma$ be an adapted basis.  Then $\Gamma$ is invertible, and for every strictly outer $i$ the $i$-th column and the
$i$-th row of both $\Gamma$ and $\Gamma^{-1}$ are those of the identity matrix.
\end{lemma}
\begin{proof}
$H_q^2=1$ gives $(H_q\Gamma^*H)\Gamma=1$.  Columns: $\Gamma e_i=e_i$ by definition, hence $\Gamma^{-1}e_i=e_i$.  Rows: let $i'$ be the
strictly outer index with $i+i'=2r+k$.  The $(i',l)$ entry of $\Gamma^*H\Gamma=H_q$ is $[l=i]$; on the left it is
$\sum_b\overline{\Gamma_{bi'}}(H\Gamma)_{bl}=(H\Gamma)_{i'l}=\Gamma_{il}$, since column $i'$ of $\Gamma$ is $e_{i'}$ and row $i'$ of
$H$ is $e_i^{\,T}$.  So $e_i^{\,T}\Gamma=e_i^{\,T}$, hence $e_i^{\,T}\Gamma^{-1}=e_i^{\,T}$.
\end{proof}

\begin{lemma}[Splitting of the Lie algebra]\label{i3:ob-B-span}
Let $k\ge1$ and $\Gamma$ an adapted basis for $z$.  Every $Y\in\mathfrak u$ is $Y_1+Y_2+Y_3$ with $Y_1\in\mathfrak u_z$,
$Y_2\in\mathfrak n$, $Y_3\in\mathfrak q_\Gamma$.
\end{lemma}
\begin{proof}
\emph{Criterion.}  Since $(H_q)_{ij}=[j=p(i)]$, the $(i,j)$ entry of $X^*H_q+H_qX$ is $\overline{X_{p(j)i}}+X_{p(i)j}$; so
$X\in\mathfrak u(H_q)$ iff $\overline{X_{p(j)i}}+X_{p(i)j}=0$ for all $i,j$.  Consequently: if a set $P$ of positions satisfies
$(p(j),i)\in P\Rightarrow(p(i),j)\in P$, then $X\mapsto X\cdot1_P$ (entries outside $P$ replaced by $0$) preserves
$\mathfrak u(H_q)$; $X\mapsto(X_{p(i)p(j)})_{ij}$ preserves $\mathfrak u(H_q)$; and $a\cdot1\in\mathfrak u(H_q)$ if $\bar a+a=0$.
The sets of upper positions, of central positions and of upper non-central positions have this symmetry (a case check
using that $p$ reverses the order of the outer indices and fixes the others).

Let $Y'=\Gamma^{-1}Y\Gamma$; as $\Gamma^{-1}=H_q\Gamma^*H$, $Y'\in\mathfrak u(H_q)$.  Let $a=Y'_{w_0w_L}$; the criterion at
$(i,j)=(w_L,w_L)$ ($p(w_L)=w_0$) gives $\bar a+a=0$.  Let $U'=Y'\cdot1_{\rm upper}$, $U''=(U'_{p(i)p(j)})_{ij}$,
\[
B=(U'+U''+a\cdot1)\cdot1_{\rm central},\qquad A=Y'\cdot1_{\rm upper,\ non\text{-}central},\qquad C=Y'-B-A,
\]
all in $\mathfrak u(H_q)$.  (i) $B$ vanishes off central positions.  (ii) $B_{iw_0}=B_{iw_L}$ for all $i$: this is trivial for
non-central $i$; for $r\le i\le r+k$ one has $U'_{lw_0}=0$ for $l\ge r$ (an upper position $(l,r)$ needs $l<r$), so
$B_{iw_0}=U'_{p(i)w_L}+a[i=w_0]$ and $B_{iw_L}=U'_{iw_L}+a[i=w_L]$; for $i=w_0$ both equal $a$ ($U'_{w_Lw_L}=0$,
$U'_{w_0w_L}=a$), for $i=w_L$ both equal $a$, and otherwise $p(i)=i$ and both equal $U'_{iw_L}$.  (iii) $A$ vanishes off
upper non-central positions.  (iv) $C_{ij}=0$ at upper positions: if $(i,j)$ is not central, $C_{ij}=Y'_{ij}-0-Y'_{ij}$; if it is
central, then $i=r$ or $j=r+k$, so $(p(i),p(j))$ is not upper, $i\neq j$, and $B_{ij}=U'_{ij}=Y'_{ij}$, $A_{ij}=0$.

Put $Y_1=\Gamma B\Gamma^{-1}$, $Y_2=\Gamma A\Gamma^{-1}$, $Y_3=\Gamma C\Gamma^{-1}$; these lie in $\mathfrak u$ and sum to $Y$.
$Y_3\in\mathfrak q_\Gamma$ by (iv).  By Lemma~\ref{i3:ob-B-frame}, $Y_1e_a=\Gamma Be_a=0$ and $Y_1e_{a^*}=0$ by (i), and
$Y_1z=\rho\,\Gamma B(e_{w_0}-e_{w_L})=0$ by (ii); so $Y_1\in\mathfrak u_z$.  For $Y_2$: $(Y_2)_{ij}=\sum_{l,q}\Gamma_{il}A_{lq}(\Gamma^{-1})_{qj}$.
If $j=b<r$ and $i\ge b$: only $q=b$ contributes (column $b$ of $\Gamma^{-1}$ is $e_b$), $A_{lb}\neq0$ forces $l<b$, and then
$\Gamma_{il}=[i=l]=0$ (column $l$ of $\Gamma$ is $e_l$).  If $r\le j\le r+k$ and $i\ge r$: a strictly outer $q$ gives
$(\Gamma^{-1})_{qj}=[q=j]=0$; for the other $q$, $A_{lq}\neq0$ forces $l<r$ (an upper non-central position in a central column has
outer row index $<r$), and $\Gamma_{il}=0$ as before.  So $Y_2\in\mathfrak n$.
\end{proof}

\begin{lemma}[Exponentials]\label{i3:ob-B-exp}
For $X\in\mathfrak u_z$ and $t\in\R$, $e^{tX}\in U_z$; for $X\in\mathfrak n$, $e^{tX}\in N$.
\end{lemma}
\begin{proof}
For $X\in\mathfrak u$, $(tX)^*H=H(-tX)$ gives $(e^{tX})^*H=He^{-tX}$, so $e^{tX}\in U$.  If $Xv=0$ then every term of positive
degree in the exponential series kills $v$, so $e^{tX}v=v$; this gives $U_z$.  For $\mathfrak n$, fix a column $j$ and a
threshold $s$ (either $j=b<r$ and $s=b$, or $r\le j\le r+k$ and $s=r$).  In both cases $X_{il}=0$ for $i\ge s>l$ (as $l<r$ and
$i\ge l$) and $X_{ij}=0$ for $i\ge s$.  By induction on $m$, $(X^{m+1})_{ij}=\sum_lX_{il}(X^m)_{lj}=0$ for $i\ge s$ (terms with
$l<s$ vanish by the first property, those with $l\ge s$ by induction).  Hence $(e^{tX})_{ij}=\delta_{ij}$ for $i\ge s$.
\end{proof}

\begin{lemma}[Torus conjugation bound]\label{i3:ob-B-tor}
Statement (c) of Lemma~\ref{i3:ob-bmu-struct} holds for every $k\ge1$.
\end{lemma}
\begin{proof}
\emph{Step 1: $d_j\le Md_{j+1}$ for all $j<n-1$.}  For $j<r$ this is the hypothesis.  For $j\ge r+k$ put
$b=2r+k-j-1<r$; the torus relations give $d_jd_{b+1}=1$ and $d_{j+1}d_b=1$, hence
$d_j=d_jd_{j+1}d_b\le d_jd_{j+1}Md_{b+1}=Md_{j+1}$.  For $r\le j<r+k$: $d_j\le1$ (it is $d_{w_0}\le1$ or a central $1$) and
$d_{j+1}\ge1$ (if $j+1=r+k$, $d_{r+k}d_r=1$ with $d_r\le1$; otherwise $d_{j+1}=1$), so $d_j\le d_{j+1}\le Md_{j+1}$.
\emph{Step 2: if $(i,j)$ is not upper then $d_j\le M^nd_i$.}  If $i<j$, both are non-outer and $d_i=d_j=1$.  If $j\le i$,
iterating Step 1 gives $d_j\le M^{i-j}d_i\le M^nd_i$.
\emph{Step 3.}  $(\mathrm{dg}(d)R)^{-1}A\,\mathrm{dg}(d)R=R^*BR$ with $B_{ij}=d_i^{-1}A_{ij}d_j$.  At upper positions $B_{ij}=0$; elsewhere
$|B_{ij}|\le\|A\|M^n$ by Step 2.  Since $|R_{ij}|\le1$, each entry of $R^*BR$ is a sum of $n^2$ terms of size $\le\|B\|$.
\end{proof}

\begin{lemma}[Root vectors]\label{i3:ob-B-root}
Statement (d) of Lemma~\ref{i3:ob-bmu-struct} holds for every $k\ge1$, adapted basis $\Gamma$ (with $z=\rho\Gamma(e_{w_0}-e_{w_L})$)
and $c\in(\C^\times)^r$, with $C_1=\sum_{a<r}2|\kappa_a|$, $\kappa_a=c_a^{-1}$ ($a+1<r$), $\kappa_{r-1}=(\rho c_{r-1})^{-1}$.
\end{lemma}
\begin{proof}
Fix $a,d,R$ and write $E_{ij}$ for matrix units, $a'=p(a+1)=2r+k-a-1$, $\kappa=\kappa_a$, and
$E=\kappa E_{a,a+1}-\bar\kappa E_{a',a^*}$.  The four indices $a,a+1,a',a^*$ are outer and, as $k\ge1$, $a+1\le r<r+k\le a'$.
Using $H_qE_{ij}=E_{p(i)j}$ and $E_{ij}H_q=E_{ip(j)}$ for outer $i,j$, $E^*H_q+H_qE=\bar\kappa E_{a+1,a^*}-\kappa E_{a^*,a+1}
+\kappa E_{a^*,a+1}-\bar\kappa E_{a+1,a^*}=0$; and $E^2=0$ since $a+1\notin\{a,a'\}$ and $a^*\notin\{a,a'\}$.

From $g'^*Hg'=H$, $\Gamma^*H\Gamma=H_q$ and $\mathrm{dg}(d)^*H_q\mathrm{dg}(d)=H_q$ (torus relations) we get $R^*H_qR=H$.  Hence
$Y=R^*ER$ satisfies $Y^*H+HY=R^*(E^*H_q+H_qE)R=0$, $Y^2=R^*E^2R=0$, so $e^{\tau Y}=1+\tau Y$, and $\|Y\|\le2|\kappa|\le C_1$
(entries of $R$ have modulus $\le1$).  The torus relations $d_ad_{a^*}=1=d_{a'}d_{a+1}$ give $d_{a'}=\alpha d_{a^*}$ with
$\alpha=d_a/d_{a+1}$, hence $\mathrm{dg}(d)E=\alpha E\,\mathrm{dg}(d)$ and
\[
g'e^{\tau Y}g'^{-1}=1+\tau\,\Gamma\mathrm{dg}(d)E\,\mathrm{dg}(d)^{-1}\Gamma^{-1}=1+\alpha\tau M',\qquad M'=\Gamma E\Gamma^{-1}.
\]
By Lemma~\ref{i3:ob-B-frame}, $\Gamma e_a=e_a$ and $e_{a^*}^{\,T}\Gamma^{-1}=e_{a^*}^{\,T}$, so
$M'_{il}=[i=a]\,\kappa(\Gamma^{-1})_{a+1,l}-\bar\kappa\,\Gamma_{ia'}[l=a^*]$.
\emph{Membership in $N$.}  The matrix lies in $U$.  For $b<r$, $i\ge b$: $M'_{ib}=[i=a]\kappa[a+1=b]$ (column $b$ of $\Gamma^{-1}$ is $e_b$),
which vanishes as $i\ge b=a+1$ excludes $i=a$.  For $r\le j\le r+k$, $i\ge r$: $i\ne a$ and $j\ne a^*$, so $M'_{ij}=0$.
\emph{Character.}  First, $z_i=\rho(\Gamma_{iw_0}-\Gamma_{iw_L})=0$ for strictly outer $i$ (row $i$ of $\Gamma$ is $e_i^{\,T}$).  For
$b+1<r$: $\ell_b(1+\alpha\tau M')=\alpha\tau M'_{b,b+1}=\alpha\tau[b=a]\kappa$ (as $b+1\neq a^*$ and column $b+1$ of $\Gamma^{-1}$ is
$e_{b+1}$), so $c_b\ell_b=\alpha\tau[b=a]$.  For $b=r-1$: $\ell_b=z_b+\alpha\tau(M'z)_b=\alpha\tau[b=a]\kappa(\Gamma^{-1}z)_{r}$ (using
$z_b=z_{a^*}=0$), and $\Gamma^{-1}z=\rho(e_{w_0}-e_{w_L})$ gives $(\Gamma^{-1}z)_r=\rho$, so $c_b\ell_b=\alpha\tau[b=a]$.  Thus
$\sum_bc_b\ell_b=\alpha\tau\in\R$ and $\psi_{z,c}=\exp(i\alpha\tau)$.
\end{proof}

\begin{lemma}[Calculus of smooth vectors]\label{i3:ob-B-calc}
Let $\sigma$ be a unitary representation of $U$, $u,w,x$ smooth.
(i) There is an $\R$-linear $\Lambda:M_n(\C)\to H$ with $d\sigma(X)u=\Lambda(X)$ for $X\in\mathfrak u$.
(ii) For $y\in U$ and a word $W$ in $\mathfrak u$, the word $y^{-1}Wy$ (letters $y^{-1}X_ty$) lies in $\mathfrak u$, and
$d\sigma(W)\sigma(y)u=\sigma(y)\,d\sigma(y^{-1}Wy)u$.
(iii) For $p\in U$, $X\in\mathfrak u$ and a word $W$, $\tau\mapsto d\sigma(W)\sigma(pe^{\tau X})x$ has derivative
$d\sigma(W)\sigma(p)d\sigma(X)x$ at $\tau=0$.
(iv) For each $m$ there is $C\ge0$ with $\|d\sigma(W)w\|\le C\prod_t\|X_t\|$ for all words $W=(X_t)$ in $\mathfrak u$ of length $\le m$.
\end{lemma}
\begin{proof}
(i) Let $F$ be $C^\infty$ near $1$ with $F(g)=\sigma(g)u$ on $U$; since $e^{tX}\in U$ and $\frac{d}{dt}e^{tX}|_0=X$, $d\sigma(X)u=DF(1)[X]$;
take $\Lambda=DF(1)$.
(ii) $y^*H_0y=H_0$ and $H_0^2=1$ give $y^{-1}=H_0y^*H_0$ and $(y^{-1}Xy)^*H_0+H_0y^{-1}Xy=y^*(X^*H_0+H_0X)y$.  From
$e^{tX}y=ye^{ty^{-1}Xy}$, $\sigma(e^{tX})\sigma(y)v=\sigma(y)\sigma(e^{ty^{-1}Xy})v$; differentiating at $0$ ($\sigma(y)$ is bounded) gives
$d\sigma(X)\sigma(y)v=\sigma(y)d\sigma(y^{-1}Xy)v$ for smooth $v$.  Induction on the length of $W$ (apply this to
$v=d\sigma(y^{-1}W'y)u$) gives (ii).
(iii) Multiplying a local extension by a cut-off, choose $F\in C^\infty(M_n(\C),H)$ and an open $O\ni p$ with $F(g)=\sigma(g)x$ on
$O\cap U$, and put $\Phi(g,y)=F(gy)$.  For a letter $Y$ let $D_Y\Phi(g,y)=\partial_g\Phi(g,y)[gY]$ (the derivative along
$t\mapsto ge^{tY}$), and $D_W=D_{X_1}\cdots D_{X_l}$.  If $\Phi(g,y)=\sigma(g)v$ for $g\in U$ near a point $g_0$, then by induction on $W$,
$D_W\Phi(g,y)=\sigma(g)d\sigma(W)v$ for $g\in U$ near $g_0$.  For $y\in U$ near $p$, $\Phi(g,y)=\sigma(g)\sigma(y)x$ for $g\in U$ near $1$, so
$d\sigma(W)\sigma(y)x=D_W\Phi(1,y)$.  The operators $D_W$ (in $g$) commute with the derivative $D'$ along $y\mapsto ye^{\tau X}$ (equality of
mixed partial derivatives of the $C^\infty$ function $\Phi$), so the required derivative is $D_W\Psi(1,p)$ with $\Psi=D'\Phi$; and for $g\in U$
near $1$, $\Psi(g,p)=\frac{d}{d\tau}\sigma(gpe^{\tau X})x|_0=\sigma(g)\sigma(p)d\sigma(X)x$, so $D_W\Psi(1,p)=d\sigma(W)\sigma(p)d\sigma(X)x$.
(iv) Fix a real basis $b_1,\dots,b_N$ of $\mathfrak u$ and $\kappa\ge0$ such that every $X\in\mathfrak u$ is $\sum c_ib_i$ with $|c_i|\le\kappa\|X\|$.
We show by induction on $m$: for every word $P$ there is $C_P$ with $\|d\sigma(PW)w\|\le C_P\prod\|X_t\|$ for $|W|\le m$.  For $m=0$ take
$C_P=\|d\sigma(P)w\|$.  For $W=XW'$, by (i) applied to $d\sigma(W')w$ and linearity of $d\sigma(P)$,
$d\sigma(PXW')w=\sum_ic_i\,d\sigma(Pb_iW')w$, so $\|d\sigma(PXW')w\|\le\kappa\|X\|\sum_iC_{Pb_i}\prod\|W'\|$; take
$C_P=\max(\|d\sigma(P)w\|,\kappa\sum_iC_{Pb_i})$.
\end{proof}

\begin{lemma}[Product-of-exponentials charts]\label{i3:ob-B-chart}
Let $X_1,\dots,X_d,Q_1,\dots,Q_e\in\mathfrak u$ with $\mathfrak u=\mathrm{span}_\R\{X_j\}+\mathrm{span}_\R\{Q_i\}$, and
$s(t)=e^{t_1X_1}\cdots e^{t_dX_d}$ ($t\in\R^d$).  Then $s$ is $C^\infty$ with values in $U$, and there are $C^\infty$ maps
$Z_j:\R^d\to\mathfrak u$ with $Z_j(0)=X_j$, $s(t+\tau e_j)=s(t)e^{\tau Z_j(t)}$, and an open $O\ni0$ such that for every $Y\in\mathfrak u$
there are functions $a_j,b_i$, $C^\infty$ on $O$, with $s(t)^{-1}Ys(t)=\sum_ja_j(t)Z_j(t)+\sum_ib_i(t)Q_i$ on $O$.
\end{lemma}
\begin{proof}
The $Z_j$ by induction on $d$: write $s(t)=e^{t_1X_1}s'(t')$ with $t'=(t_2,\dots,t_d)$; put $Z_1(t)=s'(t')^{-1}X_1s'(t')$ and
$Z_{j}(t)=Z'_{j}(t')$ for $j\ge2$.  Then $s(t+\tau e_1)=e^{t_1X_1}s'(t')\,s'(t')^{-1}e^{\tau X_1}s'(t')=s(t)e^{\tau Z_1(t)}$, and the other
relations come from those of $s'$.  $Z_1\in\mathfrak u$ by Lemma~\ref{i3:ob-B-calc}(ii); it is $C^\infty$ since
$s'^{-1}=H_0s'^*H_0$; and $Z_j(0)=X_j$ as $s'(0)=1$.

Coefficients: on the real vector space $V=M_n(\C)$ let $M(t):\R^d\times\R^e\to V$, $(a,b)\mapsto\sum a_jZ_j(t)+\sum b_iQ_i$, which is
$C^\infty$ in $t$, takes values in $\mathfrak u$, and maps onto $\mathfrak u$ at $t=0$.  Choose a linear projection $\pi$ of $V$ onto
$\mathfrak u$ and a linear $r_0:V\to\R^d\times\R^e$ with $M(0)r_0=\pi$, and put $N(t)=M(t)r_0+(1-\pi)$; $N(0)=1$, so
$O=\{t:N(t)\text{ invertible}\}$ is open and contains $0$.  For $Y\in\mathfrak u$, $y(t)=s(t)^{-1}Ys(t)\in\mathfrak u$ is $C^\infty$; put
$(a(t),b(t))=r_0N(t)^{-1}y(t)$, $C^\infty$ on $O$.  With $v=N(t)^{-1}y(t)$: $M(t)r_0v+v-\pi v=y(t)$; applying $\pi$ (which fixes
$\mathfrak u$) gives $M(t)r_0v=y(t)$.
\end{proof}

\begin{lemma}[Reduction of derivatives along a chart]\label{i3:ob-B-reduce}
In the setting of Lemma~\ref{i3:ob-B-chart}, let $\sigma$ be a unitary representation of $U$, $W$ a word in $\mathfrak u$ and
$\varphi\in C_c^\infty(\R^d)$ real with support in $O$.  There is $C\ge0$ such that for all smooth $u$ and $\rho\in\R$ with
$\|d\sigma(Q_P)u\|\le\rho$ for every word $P$ in the indices $i$ of length $\le|W|$,
\[
\Big\|\int_{\R^d}\varphi(t)\,d\sigma(W)\sigma(s(t))u\,dt\Big\|\le C\rho .
\]
\end{lemma}
\begin{proof}
Induction on $W$ by appending the last letter, for all $\varphi$ at once.  For $W$ empty, $\|\sigma(s(t))u\|=\|u\|\le\rho$ and
$C=\int|\varphi|$.  Let the claim hold for $W$ and consider $WY$.  Take $a_j,b_i$ for $Y$ from Lemma~\ref{i3:ob-B-chart} and put
$f(t)=d\sigma(W)\sigma(s(t))u$, a $C^\infty$ function of $t$.  For $t\in O$, by Lemma~\ref{i3:ob-B-calc}(ii),(i) and linearity of $d\sigma(W)$,
\[
d\sigma(WY)\sigma(s(t))u=\sum_ja_j(t)\,d\sigma(W)\sigma(s(t))d\sigma(Z_j(t))u+\sum_ib_i(t)\,d\sigma(W)\sigma(s(t))d\sigma(Q_i)u,
\]
and by Lemma~\ref{i3:ob-B-calc}(iii) with $s(t+\tau e_j)=s(t)e^{\tau Z_j(t)}$, $d\sigma(W)\sigma(s(t))d\sigma(Z_j(t))u=\partial_jf(t)$.  Multiplying by
$\varphi$ (supported in $O$), integrating and integrating by parts,
\[
\int\varphi\,d\sigma(WY)\sigma(s)u=-\sum_j\int\partial_j(\varphi a_j)\,f+\sum_i\int\varphi b_i\,d\sigma(W)\sigma(s)\big(d\sigma(Q_i)u\big).
\]
The functions $\partial_j(\varphi a_j)$ and $\varphi b_i$ are $C^\infty$ with compact support in $O$.  The induction hypothesis bounds the
$j$-th term by $C_j\rho$, and the $i$-th term by $D_i\rho$ applied to the smooth vector $d\sigma(Q_i)u$, for which
$\|d\sigma(Q_P)d\sigma(Q_i)u\|=\|d\sigma(Q_{Pi})u\|\le\rho$ when $|P|\le|W|$.  Take $C=\sum C_j+\sum D_i$.
\end{proof}

\begin{lemma}[Quotient bound]\label{i3:ob-B-qbound}
Statement (e) of Lemma~\ref{i3:ob-bmu-struct} holds for every unitary representation $\sigma$ of $U$.
\end{lemma}
\begin{proof}
By Lemma~\ref{i3:lem-cw}, $B(v)=\sum_l\beta_l(d\sigma(W_l)v)$ on smooth vectors.  Take $s,O$ from Lemma~\ref{i3:ob-B-chart} and put
$\Phi(t)=\chi_1(t_1)\cdots\chi_d(t_d)$, continuous with $\Phi(0)=1$; by induction over the factors of $s(t)$,
$B(\sigma(s(t))v)=\Phi(t)B(v)$ for smooth $v$.  Choose $\delta>0$ with $\|\Phi(t)-1\|\le\frac12$ for $|t|\le\delta$ and the closed
$\delta$-ball inside $O$, and $\varphi\ge0$ in $C_c^\infty$ supported in that ball with $\int\varphi=1$.  Lemma~\ref{i3:ob-B-reduce} gives $C_l$
for $(W_l,\varphi)$; let $m=\max_l|W_l|$ and $C=2\sum_l\|\beta_l\|C_l$.  If $\|d\sigma(Q_P)v\|\le\rho$ for $|P|\le m$, then, as
$\sigma(s(t))v$ is smooth,
\[
\int\varphi(t)\Phi(t)B(v)\,dt=\sum_l\beta_l\Big(\int\varphi(t)\,d\sigma(W_l)\sigma(s(t))v\,dt\Big),
\]
of norm $\le\sum_l\|\beta_l\|C_l\rho$.  On the other hand this integral is $B(v)+\int\varphi(\Phi-1)B(v)$, and the second term has norm
$\le\frac12\|B(v)\|$; so $\|B(v)\|\le2\|\int\varphi\Phi B(v)\|\le C\rho$.
\end{proof}

\begin{lemma}[Absolute bound]\label{i3:ob-B-abs}
Let $\sigma$, $B$, $Q_i$ be as in Lemma~\ref{i3:ob-bmu-struct}(e) and let $C\ge0$, $m$ be such that $\|B(v)\|\le C\rho$ whenever
$\|d\sigma(Q_P)v\|\le\rho$ for all $|P|\le m$.  Then statement (f) of Lemma~\ref{i3:ob-bmu-struct} holds.
\end{lemma}
\begin{proof}
Let $C_w$ be the constant of Lemma~\ref{i3:ob-B-calc}(iv) for $w$ and length $m+cm$, $A=\max(1,C_0)$, and $K=CC_wA^{m+cm}$.  Fix $g,M$.
\emph{Claim:} for a word $R'$ in $\mathfrak u$ of length $\le cm$ with letters of norm $\le A$,
$\|B(\sigma(g)d\sigma(R')w)\|\le CC_wA^{m+cm}M^{cm}$.  Indeed, for $|P|\le m$, Lemma~\ref{i3:ob-B-calc}(ii) and unitarity give
$\|d\sigma(Q_P)\sigma(g)d\sigma(R')w\|=\|d\sigma(P'R')w\|$ with $P'=g^{-1}Q_Pg$, whose letters have norm $\le C_0M^c\le AM^c$; by
Lemma~\ref{i3:ob-B-calc}(iv) this is $\le C_w(AM^c)^mA^{cm}=C_wA^{m+cm}M^{cm}$, and the hypothesis on $B$ applies to the smooth vector
$\sigma(g)d\sigma(R')w$.  If $M=1$, the claim with $R'$ empty gives $\|B(\sigma(g)w)\|\le K$.  Otherwise, for smooth $u$, the function
$\tau\mapsto B(\sigma(ge^{\tau Y})u)=\sum_l\beta_l\,d\sigma(W_l)\sigma(ge^{\tau Y})u$ (Lemma~\ref{i3:lem-cw}) has derivative
$B(\sigma(g)d\sigma(Y)u)$ at $0$ by Lemma~\ref{i3:ob-B-calc}(iii), and it equals $e^{i\ell\tau}B(\sigma(g)u)$, whose derivative is
$i\ell B(\sigma(g)u)$.  Iterating, $B(\sigma(g)d\sigma(Y^{cm})w)=(i\ell)^{cm}B(\sigma(g)w)$.  The claim with $R'=Y^{cm}$ gives
$|\ell|^{cm}\|B(\sigma(g)w)\|\le CC_wA^{m+cm}M^{cm}$, and $M^{cm}\le|\ell|^{cm}$, $M^{cm}>0$ yield $\|B(\sigma(g)w)\|\le K$.
\end{proof}


\subsubsection{Proof of Proposition~\ref{i3:S-prop-KppVol}}\label{i3:S-ssec-kppvol}

Write $H_f=H(\mathbb A_f)$, $H_{\mathfrak p}=H(F_{\mathfrak p})$, and identify $G_f$ with $H_f$ through the bridge $b_f$, which is an
isomorphism of topological groups (Proposition~\ref{i3:S-prop-bridge}(B2)) compatible with the projections to the local components
and with matrix entries ((B4)), with $b_f(x)\in K_f\iff x\in K$ ((B3); $K$ the stabiliser of the lattice in $G_f$).  $K_2$ is a
subgroup by hypothesis (for the package this is Lemma~\ref{i3:lem-K2group}); normalisation gives $K(g)\subset K_2\subset K$, where $K(g)=\{x\in K:(x-1)\widehat L\subset g\widehat L\}$.  Hence
$K_X=b_f(K_2)$ is a subgroup of $K_f$ and $K''_{\mathfrak p}=\pi_{\mathfrak p}(K_X)$ a subgroup of $K_{\mathfrak p}$.  For an
idempotent $a\in\mathbb A_f$ with integral components and $x\in H_f$ put $x^{(a)}=ax+(1-a)$ (computed in the matrix ring over
$E\otimes_F\mathbb A_f$); since $a$ is fixed by the conjugation of $E/F$, $x^{(a)}$ is again unitary, and $x^{(a)}\in K_f$ if $x\in K_f$.
For orthogonal idempotents $a,a'$ one has $x^{(a)}x^{(a')}=x^{(a+a')}$, and $x^{(1)}=x$.  Let $e_{\mathfrak p}$ be the idempotent
equal to $1$ at $\mathfrak p$ and $0$ elsewhere, and $e=\sum_{\mathfrak p\in P}e_{\mathfrak p}$.

\emph{Step 1: openness.}  $K(g)$ is a neighbourhood of $1$: the lattice $\widehat L$ lies between $N^{-1}$ and $N$ times the
standard lattice for some integer $N\ge1$, so the set of $x$ such that all entries of $x-1$ and $x^{-1}-1$ lie in $gN^2$ times the
integral adeles, which is open, is contained in $K(g)$.  Hence the subgroup $K_X\supset b_f(K(g))$ is open.  Locally, let
$K_{\mathfrak p}(g)=\{y\in K_{\mathfrak p}:(y-1)L_{\mathfrak p}\subset gL_{\mathfrak p}\}$, open by the same argument.  For $y\in K_{\mathfrak p}(g)$
the element $\tilde y\in H_f$ equal to $y$ at $\mathfrak p$ and to $1$ elsewhere lies in $b_f(K(g))\subset K_X$ and $\pi_{\mathfrak p}(\tilde y)=y$, so
$K_{\mathfrak p}(g)\subset K''_{\mathfrak p}$ and $K''_{\mathfrak p}$ is open.

\emph{Step 2: local indices.}  $K_{\mathfrak p}$ is compact (Proposition~\ref{i3:S-prop-bridge}(K)) and $K''_{\mathfrak p}$ an open subgroup, so the index
$\ell_{\mathfrak p}=[K_{\mathfrak p}:K''_{\mathfrak p}]$ is finite and nonzero, and $K_{\mathfrak p}$ is the disjoint union of $\ell_{\mathfrak p}$ left
translates of $K''_{\mathfrak p}$, each of measure $\mu_{\mathfrak p}(K''_{\mathfrak p})$ by left invariance.  So
$\mu_{\mathfrak p}(K_{\mathfrak p})=\ell_{\mathfrak p}\,\mu_{\mathfrak p}(K''_{\mathfrak p})$, and $0<\mu_{\mathfrak p}(K''_{\mathfrak p})<\infty$
(an open nonempty set has positive Haar measure, a compact set finite measure).

\emph{Step 3: the index product.}  Let $\varphi:K_f\to\prod_{\mathfrak p\in P}K_{\mathfrak p}$, $x\mapsto(\pi_{\mathfrak p}x)_{\mathfrak p}$.
(a) $\varphi$ is surjective: for $k\in K_{\mathfrak p}$ the element equal to $k$ at $\mathfrak p$ and $1$ elsewhere lies in $K_f$, and the
product over $\mathfrak p\in P$ of such elements has prescribed components.  (b) $K_X=\varphi^{-1}\bigl(\prod_{\mathfrak p}K''_{\mathfrak p}\bigr)$.
The inclusion $\subset$ is the definition of $K''_{\mathfrak p}$.  Conversely let $x\in K_f$ with $\pi_{\mathfrak p}x\in K''_{\mathfrak p}$ for all
$\mathfrak p\in P$, and choose $w_{\mathfrak p}\in K_X$ with $\pi_{\mathfrak p}w_{\mathfrak p}=\pi_{\mathfrak p}x$.  (i) $x^{(e_{\mathfrak p})}$
depends only on $\pi_{\mathfrak p}x$, so $x^{(e_{\mathfrak p})}=w_{\mathfrak p}^{(e_{\mathfrak p})}$.  (ii) $K_X$ is stable under
$w\mapsto w^{(a)}$ for integral idempotents $a$: the defining conditions of $K_2$ are of the form $(x-1)v\in M$ for fixed vectors $v$ and
fixed submodules $M$ over the integral adeles of $E$ (the span of some frame vectors plus $g\widehat L$, see
\S\ref{i3:ssec-pf-unip}), together with $x\in K$; and $(w^{(a)}-1)v=a(w-1)v$.  So $x^{(e_{\mathfrak p})}\in K_X$ for each $\mathfrak p$,
and $x^{(e)}=\prod_{\mathfrak p}x^{(e_{\mathfrak p})}\in K_X$.  (iii) $x^{(1-e)}\in b_f(K(g))$: let $c\in\mathbb A_f$ be $0$ at the primes
dividing $g$ and $g^{-1}$ elsewhere; it is integral and $gc=1-e$, so $(x^{(1-e)}-1)v=(1-e)(x-1)v=g\,c(x-1)v\in g\widehat L$ for
$v\in\widehat L$.  Hence $x=x^{(e)}x^{(1-e)}\in K_X$.  From (a), (b),
\[
[K_f:K_X]=\Bigl[\prod_{\mathfrak p}K_{\mathfrak p}:\prod_{\mathfrak p}K''_{\mathfrak p}\Bigr]=\prod_{\mathfrak p\in P}\ell_{\mathfrak p},
\]
which is finite and nonzero by Step 2.

\emph{Step 4: global index.}  $K_X$ is an open subgroup of $K_f$ of finite index $[K_f:K_X]$; decomposing $K_f$ into left cosets as in
Step 2 gives $\mu_{G_f}(K_f)=[K_f:K_X]\,\mu_{G_f}(K_X)$, $\mu_{G_f}$ being a left Haar measure.

\emph{Conclusion.}  By (D7), Steps 2--4,
\[
\mu K''(B(u,u'))=\frac{[K_f:K_X]\,\mu_{G_f}(K_X)}{\prod_{\mathfrak p}\ell_{\mathfrak p}\mu_{\mathfrak p}(K''_{\mathfrak p})}\prod_{\mathfrak p}\mu_{\mathfrak p}(K''_{\mathfrak p})
=\mu_{G_f}(K_X),
\]
the cancellations being legitimate since $\prod_{\mathfrak p}\ell_{\mathfrak p}=[K_f:K_X]\ne0$ and $0<\mu_{\mathfrak p}(K''_{\mathfrak p})<\infty$.
All quantities are finite: $\mu_{G_f}(K_X)\le\mu_{G_f}(K_f)<\infty$, $K_f=b_f(K)$ being compact ((B2), (B3), (K)).\qed

\subsection{The level subgroup and the diagonal cells}\label{i3:ssec-grp-level}

We prove the facts on the level subgroup, state the counting hypotheses on the diagonal, and prove Propositions~\ref{i3:S-prop-idcell}, \ref{i3:S-prop-transport} and~\ref{i3:prop-VGeo}.

\subsubsection{The level subgroup: group property and unipotence}\label{i3:ssec-pf-unip}

This subsection proves the level-subgroup facts used in Proposition~\ref{i3:prop-rowtie}.  Let $V=E^n\otimes_F\mathbb A_f$ with
the hermitian form $\langle\,,\rangle$ (conjugate-linear in the first variable) extended adelically, $\widehat L\subset V$ the
completion of the self-dual lattice, $\widehat{\mathcal O}$ the integral adeles of $E$, $\sigma$ the conjugation, and $K$ the stabiliser
of $\widehat L$ in the unitary group $G_f$.  For a set $S\subset V$ and $g\in\N$ put $\mathrm{Mo}(S)=\widehat{\mathcal O}S+g\widehat L$, and
write $c\equiv0$ for $c\in g\widehat{\mathcal O}$.  Since $\widehat L$ is self-dual, $\langle\widehat L,\widehat L\rangle\subset\widehat{\mathcal O}$.

\begin{definition}[Frame]\label{i3:def-frame}
A \emph{frame modulo $g$} consists of vectors $e_a,f_a$ ($0\le a<r$), $z$, $\beta_i,\beta_i^\vee$ ($0\le i<k$) in $\widehat L$ and an
idempotent $\epsilon\in\widehat{\mathcal O}$ with $\epsilon\sigma(\epsilon)=0$, such that, modulo $g\widehat{\mathcal O}$:
$\langle e_a,e_b\rangle\equiv\langle f_a,f_b\rangle\equiv0$, $\langle e_a,f_b\rangle\equiv\delta_{ab}$; $\langle z,z\rangle\equiv-1$,
$\langle e_a,z\rangle\equiv\langle f_a,z\rangle\equiv0$; $z,e_a,f_a$ are orthogonal modulo $g$ to all $\beta_i,\beta_i^\vee$;
$\beta_i=\epsilon\beta_i$, $\beta_i^\vee=\sigma(\epsilon)\beta_i^\vee$, $\langle\beta_i^\vee,\beta_j\rangle\equiv\delta_{ij}\epsilon$; and
$\widehat L\subset\mathrm{Mo}(\text{all frame vectors})$.  Its \emph{level subgroup} $K_2$ is the set of $x\in K$ with
\[
(x-1)e_a\in\mathrm{Mo}(\{e_b:b<a\}),\qquad (x-1)z\in\mathrm{Mo}(\{e_b\}),\qquad (x-1)\beta_i\in\mathrm{Mo}(\{e_b\}\cup\{\beta_j:j<i\}).
\]
\end{definition}

The level subgroup of the package is of this form, with $g=g_0g_1$ (the two auxiliary primes, called $g_1,g_2$ in
Theorem~\ref{i3:S-thm-components}), the frame of the Hida setting, $r=r_S$ and $k$ the rank
parameter; the frame clauses are the integrality data of the package.  Note $\langle\beta_i,\beta_j\rangle=\sigma(\epsilon)\epsilon\langle\beta_i,\beta_j\rangle=0$
and likewise $\langle\beta_i^\vee,\beta_j^\vee\rangle=0$ exactly.

\begin{lemma}\label{i3:lem-K2group}
$K_2$ is a subgroup of $K$, and $K(g)\subset K_2$.
\end{lemma}

\begin{proof}
Every $x\in K$ maps $\widehat L$, $g\widehat L$ and every $\widehat{\mathcal O}$-module to modules of the same kind.  $K(g)\subset K_2$ and
$1\in K_2$ are clear.  \emph{Products.}  $(xy-1)v=x(y-1)v+(x-1)v$.  Each of the three sets $S$ occurring in the definition is a union of
earlier frame vectors whose own images under $x-1$ lie in $\mathrm{Mo}$ of a subset of $S$; hence $x$ maps $\mathrm{Mo}(S)$ into itself, and
$(xy-1)v\in\mathrm{Mo}(S)$.  \emph{Inverses.}  On $M=\mathrm{Mo}(S)$, $x$ induces on $M/g\widehat L$ a map $u$ with $(u-1)$ strictly lowering
the finite filtration by the sets of earlier frame vectors, so $u-1$ is nilpotent and $u^{-1}=\sum_{i\ge0}(1-u)^i$ preserves
$M/g\widehat L$; thus $x^{-1}M\subset M$.  Then $(x^{-1}-1)v=-x^{-1}(x-1)v\in x^{-1}\mathrm{Mo}(S)=\mathrm{Mo}(S)$.
\end{proof}

\emph{The level character.}  The Hida setting fixes an additive character $\psi_E$ of $E\otimes_F\mathbb A_f$ which is
trivial on $\widehat{\mathcal O}$.  For a frame modulo $g$ put $n_a=e_{a+1}$ for $a<r-1$ and $n_{r-1}=z$, and for $x\in G_f$
\[
A(x)=\sum_{a<r}\langle f_a,xn_a\rangle+\sum_{i<k-1}\langle\beta_i^\vee,x\beta_{i+1}\rangle,\qquad
\chi_2(x)=\psi_E\bigl(g^{-1}A(x)\bigr).
\]
This is the character of the level subgroup of the package.

\begin{proposition}[Level character]\label{i3:prop-chi2}
For every Hida setting, the character $\chi_2$ of the level subgroup is multiplicative on $K_2$ and $|\chi_2|=1$ on $K_2$; together with
Lemma~\ref{i3:lem-K2group}, $\chi_2$ is a unitary character of the group $K_2$.
\end{proposition}

\begin{proof}
Write $c\equiv0$ for $c\in g\widehat{\mathcal O}$.  If $c\equiv0$ then $\psi_E(g^{-1}c)=1$; so $\chi_2(xy)=\chi_2(x)\chi_2(y)$
follows from $A(xy)-A(x)-A(y)\equiv0$, by additivity of $\psi_E$.  For $f,n\in\widehat L$ and $x,y\in G_f$,
\[
\langle f,xyn\rangle-\langle f,xn\rangle-\langle f,yn\rangle=\langle f,(x-1)(y-1)n\rangle-\langle f,n\rangle .
\]
Apply this to the pairs $(f_a,n_a)$ and $(\beta_i^\vee,\beta_{i+1})$.  First, $\langle f_a,n_a\rangle\equiv0$ (as
$\langle f_a,e_{a+1}\rangle=\overline{\langle e_{a+1},f_a\rangle}\equiv0$ and $\langle f_a,z\rangle\equiv0$) and
$\langle\beta_i^\vee,\beta_{i+1}\rangle\equiv0$, by the frame congruences.  Next let $x,y\in K_2$.  By the definition of
$K_2$, $(y-1)n_a\in\mathrm{Mo}(\{e_b:b\le a\})$ (for $a=r-1$, $(y-1)z\in\mathrm{Mo}(\{e_b\})$), and since $x\in K$ preserves
$g\widehat L$ and $(x-1)e_b\in\mathrm{Mo}(\{e_c:c<b\})$, $(x-1)(y-1)n_a\in\mathrm{Mo}(\{e_c:c<a\})$.  Every element of this
module has $\langle f_a,\cdot\rangle\equiv0$: $\langle f_a,e_c\rangle\equiv0$ for $c\ne a$, the form is $\widehat{\mathcal
O}$-linear in the second variable, and $\langle\widehat L,g\widehat L\rangle\subset g\widehat{\mathcal O}$ by self-duality.
Likewise $(y-1)\beta_{i+1}\in\mathrm{Mo}(\{e_b\}\cup\{\beta_j:j\le i\})$ and $(x-1)(y-1)\beta_{i+1}\in\mathrm{Mo}(\{e_b\}\cup
\{\beta_j:j<i\})$, on which $\langle\beta_i^\vee,\cdot\rangle\equiv0$, since $\beta_i^\vee$ is orthogonal modulo $g$ to the
$e_b$ and $\langle\beta_i^\vee,\beta_j\rangle\equiv0$ for $j\ne i$.  Summing, $A(xy)-A(x)-A(y)\equiv0$.

For $|\chi_2|=1$: if $x\in K$ then $xn_a,x\beta_{i+1}\in\widehat L$, so $A(x)\in\widehat{\mathcal O}$; then
$\chi_2(x)^g=\psi_E(A(x))=1$, hence $|\chi_2(x)|=1$.
\end{proof}

\begin{proposition}[Unipotence modulo $g_0$]\label{i3:prop-unip}
Let a frame modulo $g$ be given with $2r+k+1\le n$, and let $g_0\mid g$.  Then for every $x\in K_2$ and $y\in\widehat L$,
$(x-1)^ny\in g_0\widehat L$.
\end{proposition}

\begin{proof}
\emph{(a) Reconstruction.}  For $w\in V$ put
\[
\mathrm{Rc}(w)=\sum_a\langle f_a,w\rangle e_a+\sum_a\langle e_a,w\rangle f_a-\langle z,w\rangle z+\sum_i\langle\beta_i^\vee,w\rangle\beta_i+\sum_i\langle\beta_i,w\rangle\beta_i^\vee .
\]
We claim $w-\mathrm{Rc}(w)\in g\widehat L$ for $w\in\widehat L$.  The set of such $w$ is an $\widehat{\mathcal O}$-module ($\mathrm{Rc}$ is
$\widehat{\mathcal O}$-linear in $w$), it contains $g\widehat L$ (as $\mathrm{Rc}(\widehat L)\subset\widehat L$), so by the spanning clause it
suffices to check the frame vectors.  For $w=e_b$: the coefficients are $\langle f_a,e_b\rangle\equiv\delta_{ab}$ (the conjugate of
$\langle e_b,f_a\rangle$), $\langle e_a,e_b\rangle\equiv0$, $\langle z,e_b\rangle\equiv0$, $\langle\beta^\vee_i,e_b\rangle\equiv\langle\beta_i,e_b\rangle\equiv0$,
so $\mathrm{Rc}(e_b)\equiv e_b$ modulo $g\widehat L$ (a coefficient in $g\widehat{\mathcal O}$ times a vector of $\widehat L$ lies in
$g\widehat L$).  The same computation applies to $f_b$, to $z$ (coefficient $-\langle z,z\rangle\equiv1$), to $\beta_j$ (coefficient
$\langle\beta_i^\vee,\beta_j\rangle\equiv\delta_{ij}\epsilon$ and $\epsilon\beta_j=\beta_j$; $\langle\beta_i,\beta_j\rangle=0$) and to
$\beta_j^\vee$ (coefficient $\langle\beta_i,\beta_j^\vee\rangle\equiv\delta_{ij}\sigma(\epsilon)$ and $\sigma(\epsilon)\beta^\vee_j=\beta^\vee_j$).

\emph{(b) Two tools.}  For $x\in G_f$ unitary, $\langle u,(x-1)w\rangle=\langle(x^{-1}-1)u,w\rangle$.  If $v\in\mathrm{Mo}(S)$, $u\in\widehat L$
and $\langle s,u\rangle\equiv0$ for $s\in S$, then $\langle v,u\rangle\equiv0$ (as $\langle g\widehat L,\widehat L\rangle\subset g\widehat{\mathcal O}$).
Also, $x\in K$ maps $g\widehat L$ into itself, so to show $(x-1)\mathrm{Mo}(S)\subset\mathrm{Mo}(T)$ it suffices to check $(x-1)s\in\mathrm{Mo}(T)$
for $s\in S$.  Fix $x\in K_2$; by Lemma~\ref{i3:lem-K2group}, $x^{-1}\in K_2$.

\emph{(c) The dual vectors.}  We claim
\[
(x-1)\beta_j^\vee\in\mathrm{Mo}(\{e_a\}\cup\{z\}\cup\{\beta^\vee_l:l>j\}),\qquad
(x-1)f_b\in\mathrm{Mo}(\{e_a\}\cup\{z\}\cup\{\beta_i\}\cup\{\beta^\vee_i\}\cup\{f_c:c>b\}).
\]
For the first, $w=(x-1)\beta_j^\vee\in\widehat L$ and $w\equiv\mathrm{Rc}(w)$; we inspect the five sums.  The $e_a$- and $z$-terms lie
in the span.  $\langle e_a,w\rangle=\langle(x^{-1}-1)e_a,\beta_j^\vee\rangle\equiv0$ by (b), since $(x^{-1}-1)e_a\in\mathrm{Mo}(\{e_c\})$
and $\langle e_c,\beta^\vee_j\rangle\equiv0$; so the $f_a$-terms lie in $g\widehat L$.  $\langle\beta_i^\vee,w\rangle=0$ exactly, because
$\beta_i^\vee=\sigma(\epsilon)\beta^\vee_i$ and $w=\sigma(\epsilon)w$ ($x$ is linear over $E\otimes\mathbb A_f$), and
$\langle\sigma(\epsilon)a,\sigma(\epsilon)b\rangle=\epsilon\sigma(\epsilon)\langle a,b\rangle=0$.  For the $\beta_i^\vee$-terms: if $i>j$ the
term is in the span; if $i\le j$, $\langle\beta_i,w\rangle=\langle(x^{-1}-1)\beta_i,\beta_j^\vee\rangle$ with
$(x^{-1}-1)\beta_i\in\mathrm{Mo}(\{e_c\}\cup\{\beta_l:l<i\})$, and $\langle e_c,\beta_j^\vee\rangle\equiv0$,
$\langle\beta_l,\beta_j^\vee\rangle\equiv0$ for $l<i\le j$; so the term lies in $g\widehat L$.  For the second, $w=(x-1)f_b$: the
$e_a$-, $z$-, $\beta_i$- and $\beta_i^\vee$-terms lie in the span; for the $f_a$-terms with $a\le b$,
$\langle e_a,w\rangle=\langle(x^{-1}-1)e_a,f_b\rangle\equiv0$ since $(x^{-1}-1)e_a\in\mathrm{Mo}(\{e_c:c<a\})$ and
$\langle e_c,f_b\rangle\equiv0$ for $c<a\le b$.

\emph{(d) A filtration.}  For $s\in\N$ let $G_s$ consist of: $e_b$ with $b+1\le s$; $z$ if $r+1\le s$; $\beta_j$ with $r+2+j\le s$;
$\beta_j^\vee$ with $r+1+k\le s+j$; $f_b$ with $2r+1+k\le s+b$.  Then $G_0=\emptyset$ and $G_N$ contains all frame vectors for
$N=2r+1+k$.  We claim $(x-1)\mathrm{Mo}(G_{s+1})\subset\mathrm{Mo}(G_s)$.  By (b) it suffices to check generators, using the definition
of $K_2$ for $e_b,z,\beta_j$ and (c) for $\beta^\vee_j,f_b$, and comparing indices:
for $e_b\in G_{s+1}$ the vectors $e_c$, $c<b$, satisfy $c+1\le b\le s$; for $z\in G_{s+1}$, every $e_c$ has $c+1\le r\le s$; for
$\beta_j\in G_{s+1}$, $e_c$ has $c+1\le r\le s$ and $\beta_l$, $l<j$, has $r+2+l\le r+1+j\le s$; for $\beta^\vee_j\in G_{s+1}$ we have
$s\ge r+k-j\ge r+1$, so all $e_c$ and $z$ are in $G_s$, and $\beta^\vee_l$, $l>j$, has $r+1+k\le s+1+j\le s+l$; for $f_b\in G_{s+1}$
we have $s\ge 2r+k-b\ge r+k+1$, so all $e_c$, $z$, $\beta_i$ ($r+2+i\le r+1+k\le s$) and $\beta^\vee_i$ ($r+1+k\le s$) are in $G_s$, and
$f_c$, $c>b$, has $2r+1+k\le s+1+b\le s+c$.

\emph{(e) Conclusion.}  Iterating (d), $(x-1)^m\mathrm{Mo}(G_{s+m})\subset\mathrm{Mo}(G_s)$; with $s=0$, $m=N$ and $\widehat L\subset\mathrm{Mo}(G_N)$
we get $(x-1)^N\widehat L\subset\mathrm{Mo}(\emptyset)=g\widehat L$.  As $n\ge N$ and $(x-1)g\widehat L\subset g\widehat L$,
$(x-1)^n\widehat L\subset g\widehat L\subset g_0\widehat L$.
\end{proof}

\noindent In Proposition~\ref{i3:prop-rowtie} this is applied with $g=g_0g_1$, the auxiliary prime $g_0$, and
$n=k+2r_S+1\ge 2r_S+k+1$.

\subsubsection{The counting hypotheses on the diagonal}\label{i3:ssec-pf-cnthyp}

\begin{definition}[Counting hypotheses]\label{i3:def-cnthyp}
Let $T_u$ be the auxiliary prime of the word $u$, $m=2r_S+2$, $L$ the size parameter and $\mathrm{vand}(u,\tau)$ the Vandermonde
product of the archimedean weights of $u$ at $\tau$.  A datum satisfies the \emph{counting hypotheses} at $(u,u')$ with constants
$(C',C_l,\mathbf c,C_0,A)$ if: it satisfies the size clauses (1)--(25) at $(C',C_l,\mathbf c)$ and the numerical clauses (N1)--(N14) at $C_0$
(\S\ref{i3:ssec-setting}), and $1\le C_0$; $\mathrm{Stand}(u,u')$ holds, namely $r_S\ge1$, $T_u,T_{u'}$ are primes $\ge3$ splitting completely in $E$, $g_1\ne g_2$
are primes splitting completely in $E$, $g_1g_2$ is prime to $T_uT_{u'}$ and $T_u\equiv T_{u'}\pmod{g_1g_2}$; $T_{u'}/T_u\le A$ and
$T_u/T_{u'}\le A$; $C_lm^6L\le k\log k$; $(\log f)^2\le Am^2L$ and $\log f\le Ar_S^2$; and for every $\tau$,
$\mathrm{vand}(u',\tau)\le A^k\mathrm{vand}(u,\tau)$ and $\mathrm{vand}(u,\tau)\le A^k\mathrm{vand}(u',\tau)$.
\end{definition}

\begin{lemma}\label{i3:lem-cnthyp}
If a datum satisfies the size clauses (1)--(25) at $(C',C_l,\mathbf c)$ and the numerical clauses (N1)--(N14) at $C_0\ge1$, with $2<C_l$,
$1\le C'$ and $A\ge1$, then it satisfies the counting hypotheses at $(u,u)$ for every word $u$.
\end{lemma}

\begin{proof}
Besides \eqref{i3:eq-standing} we use the standing clauses $\log f\le\sqrt{r_S}$ and $C_l\le m$, and the numerical data: $T_u$ splits
completely in $E$ and $T_u\equiv1\pmod{(g_1g_2)^2}$; $g_1,g_2$ split completely; $T_u\ne g_1,g_2$.  The bounds with $A$ at $u'=u$ are
$1\le A$ and $\mathrm{vand}\le A^k\mathrm{vand}$ (the Vandermonde products are positive, since the weights of $u$ at $\tau$ are
strictly decreasing).  \emph{$k\log k$.}  Using $C_l\le m$, $m\ge2$,
$m^4\le f$, $L>1$ (a size clause), $C'\ge1$ and $C'm^5fL\le k$: $C_lm^6L\le m^7L\le m^9L\le m^5fL\le C'm^5fL\le k\le k\log k$, as
$k\ge5$ (from $m\ge2$ and $k\ge m+3$) gives $\log k\ge1$.  \emph{$\log f$.}  As $f\ge1$, $(\log f)^2\le r_S\le m\le m^2L\le Am^2L$, and $\log f\le\sqrt{r_S}\le r_S^2\le Ar_S^2$.
\emph{$\mathrm{Stand}(u,u)$.}  $m\ge C_l>2$ and $m$ is even, so $m\ge4$ and $r_S=(m-2)/2\ge1$.  $T_u$ is prime (a field of the datum); if
$T_u=2$ then $2\equiv1$ modulo $(g_1g_2)^2\ge16$, impossible, so $T_u\ge3$.  $g_1\ne g_2$ are primes (host data).  $T_u$ is a prime different
from $g_1,g_2$, hence $g_1g_2$ is prime to $T_uT_u$; and trivially $T_u\equiv T_u$.
\end{proof}

\subsubsection{Proof of Proposition~\ref{i3:S-prop-idcell}}\label{i3:S-ssec-pf-idcell}

\emph{Structure used.}  Besides the objects recalled at the beginning of \S\ref{i3:ssec-pf-kppoff}, the construction of the Hida
setting provides, for each word $u$, the following (they are part of its definition or immediate from it).  (H1) $N_u$ and $H_u$ are contained in $S_u$, and
$nh\in S_u$ for $n\in N_u$, $h\in H_u$; $H_u$ is closed.  $\pi_u$ is defined on $S_u$, multiplicative there
($\pi_u(ab)=\pi_u(a)\pi_u(b)$ for $a,b\in S_u$), trivial on $N_u$, unitary on $S_u$, and $h\mapsto\pi_u(h)\xi$ is continuous; $|\psi_u|=1$ on
$S_u$; $W_u$ is continuous and $W_u(sx)=\psi_u(s)\pi_u(s)W_u(x)$ for $s\in S_u$, $x\in U_\infty$.  (H2) There are measures $\mu_{N,u}$ and $\mu_{H,u}$,
$\sigma$-finite, concentrated on $N_u$ resp.\ $H_u$, invariant under left translation by elements of $N_u$ resp.\ $H_u$, and $\nu_u$ on $U_\infty$, related by Weil's formula: for every measurable
$f:U_\infty\to[0,\infty]$,
\[
\int f\,d\mu_\infty=\int\Bigl(\iint f(nhx)\,d\mu_{N,u}(n)\,d\mu_{H,u}(h)\Bigr)d\nu_u(x).\tag{W}
\]
Strict fundamental domains are as defined at the beginning of \S\ref{i3:ssec-pf-kppoff}.

\begin{lemma}\label{i3:S-lem-ic-schur}
Let $G$ be a group with measurable left translations, $H\subset G$ a subset, $\mu$ a finite measure on $G$ with $\mu(G\setminus H)=0$,
invariant under left translation by elements of $H$, and $\pi,\sigma$ maps from $G$ to the operators of finite-dimensional Hilbert
spaces $V,V'$, inner-product preserving and multiplicative on $H$, with measurable matrix coefficients.  If every linear $T:V'\to V$ with $T\sigma(h)=\pi(h)T$ ($h\in H$) is zero, then for all $w,\xi\in V$, $w',\eta\in V'$,
\[
\int\overline{\langle\pi(h)w,\xi\rangle}\,\langle\sigma(h)w',\eta\rangle\,d\mu(h)=0 .
\]
\end{lemma}

\begin{proof}
Put $Tx=\int\langle\sigma(h)w',x\rangle\,\pi(h)w\,d\mu(h)$ for $x\in V'$; the integrand is bounded by $\|w'\|\|x\|\|w\|$ (unitarity), so $T$ is a
well-defined linear map.  For $k\in H$, using $\pi(k)\pi(h)=\pi(kh)$, $\langle\sigma(h)w',x\rangle=\langle\sigma(kh)w',\sigma(k)x\rangle$ and the
invariance of $\mu$ under $h\mapsto kh$,
\[
\pi(k)Tx=\int\langle\sigma(kh)w',\sigma(k)x\rangle\,\pi(kh)w\,d\mu(h)=\int\langle\sigma(h)w',\sigma(k)x\rangle\,\pi(h)w\,d\mu(h)=T\sigma(k)x .
\]
Hence $T=0$, and $0=\langle\xi,T\eta\rangle=\int\langle\sigma(h)w',\eta\rangle\langle\xi,\pi(h)w\rangle\,d\mu(h)$, which is the assertion since
$\langle\xi,\pi(h)w\rangle=\overline{\langle\pi(h)w,\xi\rangle}$.
\end{proof}

\begin{lemma}\label{i3:S-lem-ic-vol}
Let $\Gamma$ be a countable subgroup of a group $G$ with a measurable multiplication, $\mu$ a measure on $G$ invariant under left
translation by the elements of $\Gamma$, and $\Omega$ a strict fundamental domain of $\Gamma$.  Then $\mu\{n:ny\in\Omega\}=\mu(\Omega)$ for every
$y\in G$.
\end{lemma}

\begin{proof}
$A=\{n:ny\in\Omega\}=\Omega y^{-1}$ is again a strict fundamental domain (measurable as the preimage under a right translation, and the two
conditions are invariant under right translation).  Let $P,Q$ be strict fundamental domains.  For $n\in P$ there is exactly one $s\in\Gamma$
with $sn\in Q$: if $sn,s'n\in Q$ then $(s's^{-1})(sn)\in Q$ forces $s's^{-1}=1$.  So $P$ is the disjoint countable union of the
measurable sets $P\cap s^{-1}Q$, and by invariance
\[
\mu(P)=\sum_{s\in\Gamma}\mu(P\cap s^{-1}Q)=\sum_{s\in\Gamma}\mu(sP\cap Q)=\sum_{s\in\Gamma}\mu(Q\cap s^{-1}P)=\mu(Q),
\]
the third equality reindexing $s\mapsto s^{-1}$ and the last being the first with $P,Q$ exchanged.  Take $P=A$, $Q=\Omega$.
\end{proof}

\begin{lemma}\label{i3:S-lem-ic-weil}
If $g:U_\infty\to\C$ is measurable and $\mu_\infty$-integrable, then for $\nu_u$-almost every $x$ the function $(n,h)\mapsto g(nhx)$ is
$\mu_{N,u}\otimes\mu_{H,u}$-integrable, and $\int g\,d\mu_\infty=\int\bigl(\iint g(nhx)\,d\mu_{N,u}\,d\mu_{H,u}\bigr)d\nu_u(x)$.
\end{lemma}

\begin{proof}
For $g\ge0$ this is (W): (W) applied to $g$ shows that the inner integral is finite for almost every $x$ and integrable in $x$.  A real $g$ is
the difference of $\max(g,0)$ and $\max(-g,0)$, both integrable, and a complex $g$ is $\mathrm{Re}\,g+i\,\mathrm{Im}\,g$; the identity is
additive and the inner integrals split accordingly where both parts are integrable, which holds almost everywhere.
\end{proof}

\begin{proof}[Proof of Proposition~\ref{i3:S-prop-idcell}]
Fix an identity term $t$, and write $\gamma$, $\Gamma=\Gamma_t$, $\Omega=\Omega_t$, $N=N_u$, $H=H_u$, $S=S_u$, $\pi=\pi_u$, $\psi=\psi_u$,
$W=W_u$, $\mu_N=\mu_{N,u}$, $\mu_H=\mu_{H,u}$, $\nu=\nu_u$.  For $s\in S$ put $\sigma(s)=\pi_{u'}(\gamma^{-1}s\gamma)$, which is defined by (Conj),
and $W'(x)=W_{u'}(\gamma^{-1}x)$.  Since $s\mapsto\gamma^{-1}s\gamma$ is multiplicative, maps $S$ into $S_{u'}$ and $N$ into $N_{u'}$ (Conj),
$\sigma$ is multiplicative on $S$, trivial on $N$ and unitary on $S$.  For $s\in S$, by (H1) for $u'$ and (Conj),
\[
W'(sx)=W_{u'}\bigl((\gamma^{-1}s\gamma)(\gamma^{-1}x)\bigr)=\psi_{u'}(\gamma^{-1}s\gamma)\,\pi_{u'}(\gamma^{-1}s\gamma)W'(x)=\psi(s)\,\sigma(s)W'(x).
\]

Fix $\xi\in V_\pi(u)$, $\eta\in V_\pi(u')$, and let $g_0(x)=\overline{\langle W(x),\xi\rangle}\langle W'(x),\eta\rangle$.  Since
$(W(x)\otimes\overline{W'(x)})\eta=\langle W'(x),\eta\rangle W(x)$, we have $\langle\xi,\mathrm{OO}_t\eta\rangle=\int_\Omega g_0\,d\mu_\infty$, and
$g=\mathbf 1_\Omega g_0$ is integrable by hypothesis and measurable ($W,W'$ are continuous).  By Lemma~\ref{i3:S-lem-ic-weil},
\[
\langle\xi,\mathrm{OO}_t\eta\rangle=\int\Bigl(\iint\mathbf 1_\Omega(nhx)\,g_0(nhx)\,d\mu_N(n)\,d\mu_H(h)\Bigr)d\nu(x).
\]
For $n\in N$, $h\in H$ we have $nh\in S$, $\pi(nh)=\pi(n)\pi(h)=\pi(h)$, $\sigma(nh)=\sigma(h)$ and $|\psi(nh)|=1$, so
\[
g_0(nhx)=\overline{\psi(nh)}\psi(nh)\,\overline{\langle\pi(h)W(x),\xi\rangle}\langle\sigma(h)W'(x),\eta\rangle=c_x(h),\qquad
c_x(h)=\overline{\langle\pi(h)W(x),\xi\rangle}\langle\sigma(h)W'(x),\eta\rangle .
\]
As $\mu_N,\mu_H$ are concentrated on $N,H$, for almost every $x$ the inner double integral is (Fubini, the integrand being integrable)
$\int c_x(h)\,\mu_N\{n:n(hx)\in\Omega\}\,d\mu_H(h)$.  By (Stab), $\Gamma$ is a countable subgroup of $N$, so $\mu_N$ is invariant under left
translation by its elements, and Lemma~\ref{i3:S-lem-ic-vol} gives $\mu_N\{n:n(hx)\in\Omega\}=\mu_N(\Omega)$ for every $h$.  (If $\mu_N(\Omega)$
were infinite, integrability would force $c_x=0$ $\mu_H$-almost everywhere and the inner integral would vanish.)  So the inner integral is
$\mu_N(\Omega)\int c_x\,d\mu_H$.  Now $\mu_H$ is finite (hypothesis), concentrated on $H$ and left $H$-invariant, $\pi$ and $\sigma$ are
unitary and multiplicative on $H\subset S$ with continuous matrix coefficients, and (Ineq) says exactly that every $T$ with $T\sigma(h)=\pi(h)T$
on $H$ vanishes.  Lemma~\ref{i3:S-lem-ic-schur} with $w=W(x)$, $w'=W'(x)$ gives $\int c_x\,d\mu_H=0$.  Hence
$\langle\xi,\mathrm{OO}_t\eta\rangle=0$ for all $\xi,\eta$, i.e.\ $\mathrm{OO}_t=0$.

The term operator $\mathrm{op}_t$ is assembled from blocks indexed by pairs of column classes, each of which is a scalar multiple of
$\mathrm{OO}_t$ or zero, and the assembly is linear; so $\mathrm{op}_t=0$.  Summing over the identity terms, $\mathrm{part}(\mathrm{id})=0$.
\end{proof}

\subsubsection{Proof of Proposition~\ref{i3:S-prop-transport}}\label{i3:S-ssec-pf-transport}

\begin{lemma}\label{i3:S-lem-sharp}
$\sharp_u$ is additive and $\sharp_u(cT)=\bar c\,\sharp_uT$.  If $\sharp_u$ is injective, $\langle T,T'\rangle_{C_u}=\langle\sharp_uT',\sharp_uT\rangle$ is an
inner product on $C_u$, conjugate-linear in the first variable, and $\sharp_u$ preserves norms.
\end{lemma}

\begin{proof}
The functional $\xi\mapsto((T+T')\xi)(1,h_\rho)$ is the sum of those of $T$ and $T'$, and that of $cT$ is $c$ times that of $T$; the
Riesz representative of $c\ell$ is $\bar c$ times that of $\ell$.  Hence the first sentence.  The form is conjugate symmetric; it is linear in the
second variable since $\langle\sharp(cT'),\sharp T\rangle=\langle\bar c\,\sharp T',\sharp T\rangle=c\langle\sharp T',\sharp T\rangle$, and additive;
$\langle T,T\rangle=\|\sharp T\|^2\ge0$, with equality only if $\sharp T=0$, i.e.\ $T=0$ by injectivity.
\end{proof}

\noindent\emph{Weyl's dimension formula} (H.~Weyl, Math.\ Z.\ 23--24 (1925--26); see e.g.\ A.~W.~Knapp, \emph{Lie Groups Beyond an
Introduction}, Theorem~5.84, or W.~Fulton and J.~Harris, \emph{Representation Theory}, Corollary~24.6; it is
also formally verified, \cite{CharacterFamily}): an irreducible finite-dimensional representation of $U(k)$ with highest weight
$\Lambda=(\Lambda_1\ge\dots\ge\Lambda_k)$ has dimension $\prod_{a<b}(\Lambda_a-\Lambda_b+b-a)/(b-a)$.  In particular its dimension depends only on
$\Lambda$.

\begin{proof}[Proof of Proposition~\ref{i3:S-prop-transport}]
(a) is Lemma~\ref{i3:S-lem-sharp} (finite column data includes injectivity of $\sharp_u$).

(b) At each infinite place $v$, $V_{u,v}$ carries an irreducible unitary representation of $U(k)$ with the highest weight attached to
$\lambda_{u,v}$; by Weyl's formula every such representation has the same dimension, so $\dim V_{u,v}=\delta(\lambda_{u,v})$.  It remains to
see $\dim V_\pi(u)=\prod_v\dim V_{u,v}$.  Choose orthonormal bases $(b_{v,j})_j$ of $V_{u,v}$ and put $e_i=\otimes_vb_{v,i_v}$ for
$i=(i_v)\in\prod_v\{1,\dots,\dim V_{u,v}\}$.  Since the inner product of pure tensors is the product of the inner products, the $e_i$ are
orthonormal, hence linearly independent.  They span: pure tensors span $V_\pi(u)$, and for a pure tensor $y=\otimes_vx_v$ the coefficients
$c_i=\langle e_i,y\rangle=\prod_v\langle b_{v,i_v},x_v\rangle$ satisfy
$\sum_i|c_i|^2=\prod_v\sum_j|\langle b_{v,j},x_v\rangle|^2=\prod_v\|x_v\|^2=\|y\|^2$.  With $s=\sum_ic_ie_i$ we have $\langle s,y-s\rangle=0$, so
$\|y\|^2=\|s\|^2+\|y-s\|^2=\|y\|^2+\|y-s\|^2$ and $y=s$.  So the $e_i$ form a basis and $\dim V_\pi(u)=\prod_v\dim V_{u,v}=d_u$.

(c) follows from the definitions and (b).

(d) Both $K(u,u')$ and $\sharp_u^{-1}K^\flat\sharp_{u'}$ are linear maps $C_{u'}\to C_u$ (the second is a composite of two conjugate-linear maps
and a linear one).  For $T\in C_u$, $T'\in C_{u'}$, by (c) and the definition of the adjoint,
\[
\langle T,K(u,u')T'\rangle_{C_u}=(c_0(u)c_0(u'))^{-1/2}\langle P_uT,P_{u'}T'\rangle=\mathrm{Kform}(T,T')=\langle K^\flat\sharp_{u'}T',\sharp_uT\rangle,
\]
while by the definition of $\langle\,,\rangle_{C_u}$,
$\langle T,\sharp_u^{-1}K^\flat\sharp_{u'}T'\rangle_{C_u}=\langle\sharp_u\sharp_u^{-1}K^\flat\sharp_{u'}T',\sharp_uT\rangle=\langle K^\flat\sharp_{u'}T',\sharp_uT\rangle$.
As $T$ is arbitrary, $K(u,u')T'=\sharp_u^{-1}K^\flat\sharp_{u'}T'$.

(e) Let $(e_{\mathcal H},e_u)$ be an isomorphism from the draft datum $s$ to $s'$.  The normalising constants agree, since the volumes are equal and
$d_u$ depends only on the word.  For $y\in\mathcal H$ and $x\in C_u$,
$\langle e_ux,e_uP_u^*y\rangle=\langle P_ux,y\rangle=\langle e_{\mathcal H}P_ux,e_{\mathcal H}y\rangle=\langle P'_ue_ux,e_{\mathcal H}y\rangle$, and
$e_u$ is onto, so $P'^*_ue_{\mathcal H}=e_uP_u^*$.  Hence, with $c=(c_0(u)c_0(u'))^{-1/2}$,
$K_{s'}(e_{u'}x)=cP'^*_uP'_{u'}e_{u'}x=cP'^*_ue_{\mathcal H}P_{u'}x=e_uK_s(x)$.  Lemma~\ref{i3:S-lem-hs}(2) with $V=e_{u'}$ and $g=e_u$ gives
$\|K_{s'}(u,u')\|_{HS}=\|K_s(u,u')\|_{HS}$.
\end{proof}

\subsubsection{Proof of Proposition~\ref{i3:prop-VGeo}}\label{i3:ssec-pf-VGeo}

\emph{(V).}  Let $t$ be a term that is not relevant, $\gamma=\gamma_t$, $\Omega=\Omega_t$, $f(x)=W_u(x)\otimes\overline{W_{u'}(\gamma^{-1}x)}$ and
$\chi(r)=\psi_u(r)\overline{\psi_{u'}(\gamma^{-1}r\gamma)}$ for $r\in R_t$.  By Proposition~\ref{i3:prop-cov-in}(St)(ii), which holds for every term,
$\chi$ is continuous and multiplicative, trivial on $\Gamma_t$, and $f(rx)=\chi(r)f(x)$ for $r\in R_t$ (the computation at the beginning of the
proof of Proposition~\ref{i3:prop-cov}).  By Proposition~\ref{i3:prop-kppoff-in}(NR) some $r\in R_t$ has $\psi_u(r)\ne\psi_{u'}(\gamma^{-1}r\gamma)$;
as $|\psi_{u'}|=1$, $\chi(r)=1$ would give $\psi_u(r)=\psi_{u'}(\gamma^{-1}r\gamma)$, so $\chi\ne1$ on $R_t$.  By (St)(i) and the averaging statement
(A) of Proposition~\ref{i3:prop-cov-in}, $\mathrm{OO}_t=\int_\Omega f\,d\mu_\infty=0$, and hence $\mathrm{op}_t=0$.

\emph{Three lemmas for (Geo).}  Fix an infinite place $v$ and write $h_v(\cdot,\cdot)=\langle\cdot,\cdot\rangle_v$.

\begin{lemma}\label{i3:lem-geo-cplx}
For $E$-subspaces $A,B\subset V$: $(A+B)_v=A_v+B_v$, $(A\cap B)_v=A_v\cap B_v$, and $(A^\perp)_v=\{y\in\C^n:\langle x,y\rangle_v=0\ \forall x\in A_v\}$.
\end{lemma}

\begin{proof}
The map $V\otimes_{E,v}\C\to\C^n$, $x\otimes c\mapsto cx_v$, is an isomorphism (it maps a basis to a basis), and it carries
$A\otimes\C$ onto $A_v$.  The first identity is clear.  For the second, choose a basis of $V$ containing bases of $A\cap B$, of a complement of $A\cap B$ in
$A$ and of a complement of $A\cap B$ in $B$ (these together are linearly independent, since $\dim(A+B)=\dim A+\dim B-\dim(A\cap B)$); in coordinates
for such a basis the identity is evident.  For the third, $A^\perp$ is the kernel of the $E$-matrix $x\mapsto(h(a_i,x))_i$ for a basis $(a_i)$ of $A$,
and since $v(h(a,x))=\langle a_v,x_v\rangle_v$ the right side is the kernel of the image matrix under $v$; the kernel of a matrix has the same
dimension over $E$ and over $\C$ (rank is invariant under field extension), and the image of the $E$-kernel lies in the $\C$-kernel and is linearly
independent over $\C$, so it spans it.
\end{proof}

\begin{lemma}\label{i3:lem-geo-frame}
If $y\in\C^n$ is $\langle\,,\rangle_v$-orthogonal to all $e_{a,v},f_{a,v},z_{u,v}$, then $y=\sum_b\langle w_{u,v,b},y\rangle_v\,w_{u,v,b}$.
\end{lemma}

\begin{proof}
The $n=2r_S+1+k$ vectors $e_{a,v},f_{a,v},z_{u,v},w_{u,v,b}$ have the Gram matrix given by $\langle e,e\rangle=\langle f,f\rangle=0$,
$\langle e_a,f_b\rangle=\delta_{ab}$, $z\perp e,f$, $\langle z,z\rangle\ne0$, $w\perp e,f,z$ and $\langle w_a,w_b\rangle=\delta_{ab}$; it is invertible, so they
form a basis.  Write $y=\sum\alpha_ae_a+\sum\beta_af_a+cz+\sum c_bw_b$.  Then $0=\langle e_a,y\rangle=\beta_a$, $0=\langle f_a,y\rangle=\alpha_a$,
$0=\langle z,y\rangle=c\langle z,z\rangle$, so $y=\sum c_bw_b$ and $c_b=\langle w_b,y\rangle$.
\end{proof}

\begin{lemma}\label{i3:lem-geo-smat}
Let $t_1,\dots,t_k\in\C^n$ be orthonormal for $\langle\,,\rangle_v$ and $M\in U(k)$.  Put $\kappa(y)=(\langle t_d,y\rangle_v)_d$ and
$\sigma(y)=y+\sum_c((M-1)\kappa(y))_ct_c$.  Then $\sigma\in U(F_v)$, $\sigma(y)=y$ if $y\perp t_1,\dots,t_k$, and $\langle t_a,\sigma t_b\rangle_v=M_{ab}$.
\end{lemma}

\begin{proof}
$\kappa(t_b)$ is the $b$-th unit vector, so $\sigma(t_b)=\sum_cM_{cb}t_c$ and $\langle t_a,\sigma t_b\rangle_v=M_{ab}$.  If $y\perp t$ then $\kappa(y)=0$ and
$\sigma(y)=y$.  Writing $y=y_0+\sum_d\kappa(y)_dt_d$ with $y_0\perp t$, $\sigma(y)=y_0+\sum_c(M\kappa(y))_ct_c$, so
$\langle\sigma y,\sigma y'\rangle_v=\langle y_0,y'_0\rangle_v+\kappa(y)^*M^*M\kappa(y')=\langle y,y'\rangle_v$; the map with $M^*$ in place of $M$ is
inverse to $\sigma$.
\end{proof}

\begin{proof}[Proof of (Geo)]
Let $t$ be relevant with $G(t)=0$, $\gamma=\gamma_t$, $\mathrm{Fix}=\mathrm{Fix}_t$, $X=X_u$, $X^+=X_u^+$, $X'=\gamma X_{u'}$, $X'^+=\gamma X_{u'}^+$,
and fix $v$; write $\langle\,,\rangle$ for $\langle\,,\rangle_v$, $w_b=w_{u,v,b}$ and $w'_b=\gamma_vw_{u',v,b}$.  Since $\gamma_v$ is unitary and
maps $x_v$ to $(\gamma x)_v$, the $w'_b$ are orthonormal and orthogonal to $X'^+_v$, and the $w_b$ are orthogonal to $X^+_v$.

\emph{Step 1 (complex form of (NF)).}  (1) If $y\perp X^+_v$ there is $x\in X_v$ with $y-x\perp\mathrm{Fix}_v$; (2) if $y\perp X'^+_v$ there is $x\in X'_v$
with $y-x\perp\mathrm{Fix}_v$; (3) if $y\in X_v$ and $y\perp\mathrm{Fix}_v$ then $y\in(X\cap X')_v$.  Indeed, by Lemma~\ref{i3:lem-geo-cplx} and (NF),
$(X^+_v)^\perp=((X^+)^\perp)_v=(\mathrm{Fix}^\perp+X)_v=(\mathrm{Fix}^\perp)_v+X_v$ and $(\mathrm{Fix}^\perp)_v=(\mathrm{Fix}_v)^\perp$, which gives (1), and
(2) is the same with $X'$; and $X_v\cap(\mathrm{Fix}_v)^\perp=(X\cap\mathrm{Fix}^\perp)_v=(X\cap X')_v$ gives (3).

\emph{Step 2 (corrected frames).}  By Step 1 choose $x_b\in X_v$ and $x'_b\in X'_v$ with $\tilde w_b=w_b-x_b\perp\mathrm{Fix}_v$ and
$\tilde w'_b=w'_b-x'_b\perp\mathrm{Fix}_v$.  As $X_v,X'_v\subset\mathrm{Fix}_v$, $w_a\perp X_v$ and $w'_a\perp X'_v$:
$\langle\tilde w_a,\tilde w_b\rangle=\langle w_a,\tilde w_b\rangle=\langle w_a,w_b\rangle=\delta_{ab}$ and likewise $\langle\tilde w'_a,\tilde w'_b\rangle=\delta_{ab}$.

\emph{Step 3 ($K$ is unitary).}  Let $K_{ca}=\langle\tilde w_c,\tilde w'_a\rangle$.  Fix $a$, let $\xi=\sum_c\langle f_{c,v},\tilde w'_a\rangle e_{c,v}$ and
$y=\tilde w'_a-\xi$.  Then $y\perp e_{c,v}$ (as $e_c\in X\subset\mathrm{Fix}$ and $\langle e,e\rangle=0$), $y\perp f_{c,v}$ (as $\langle f_c,e_d\rangle=\delta_{cd}$),
and $y\perp z_{u,v}$ (as $z_u\in X^+\subset\mathrm{Fix}$ and $z\perp e$).  By Lemma~\ref{i3:lem-geo-frame}, $y=\sum_b\langle w_b,y\rangle w_b$, and
$\langle w_b,y\rangle=\langle w_b,\tilde w'_a\rangle=\langle\tilde w_b+x_b,\tilde w'_a\rangle=K_{ba}$ since $w_b\perp e$ and $x_b\in\mathrm{Fix}_v\perp\tilde w'_a$.
So $\tilde w'_a=\xi+\sum_cK_{ca}(\tilde w_c+x_c)$, and $\eta=\tilde w'_a-\sum_cK_{ca}\tilde w_c=\xi+\sum_cK_{ca}x_c$ lies in $X_v$ and is orthogonal to
$\mathrm{Fix}_v$.  By Step 1(3), $\eta\in(X\cap X')_v\subset\mathrm{Fix}_v$, so $\langle\eta,\tilde w'_b\rangle=0$ and
\[
\sum_c\overline{K_{ca}}K_{cb}=\Bigl\langle\sum_cK_{ca}\tilde w_c,\tilde w'_b\Bigr\rangle=\langle\tilde w'_a-\eta,\tilde w'_b\rangle=\delta_{ab}.
\]
So $K^*K=1$ and $C_v=K^*\in U(k)$.

\emph{Step 4 (the element $h$).}  Given $M=(M_v)$, let $h_v=\sigma$ be the element of Lemma~\ref{i3:lem-geo-smat} for $(\tilde w_c)$ and $M_v$, and
$h=(h_v)_v$.  Every $y\in\mathrm{Fix}_v$ is orthogonal to all $\tilde w_c$, so $h_vy=y$; thus $h\in S_t$.  In particular $h$ fixes all $e_{a,v}$ and
$z_{u,v}$, so $h\in S_u$, and $m_v(h)=M_v$: indeed $h_vx_b=x_b$, $\langle w_a,x_b\rangle=0$,
$\langle x_a,h_v\tilde w_b\rangle=\langle h_v^{-1}x_a,\tilde w_b\rangle=\langle x_a,\tilde w_b\rangle=0$, hence
$\langle w_a,h_vw_b\rangle=\langle\tilde w_a,h_v\tilde w_b\rangle=M_{v,ab}$.  By (Levi) there is $s_0\in H_u$ with $s_0^{-1}h\in N_u$ and $m_v(s_0)=M_v$;
since $\pi_u$ is multiplicative on $S_u$ and trivial on $N_u$, $\pi_u(h)=\pi_u(s_0)\pi_u(s_0^{-1}h)=\pi_u(s_0)$, and by (T)
$\pi_u(h)(\otimes x_v)=\otimes\pi_{u,v}(M_v)x_v$.  As pure tensors span, $\pi_u(h)=\bigotimes_v\pi_{u,v}(M_v)$.

\emph{Step 5 (the conjugate).}  Put $h'=\gamma^{-1}h\gamma$.  Since $\gamma e_a\in X'$ and $\gamma z_{u'}\in X'^+$ lie in $\mathrm{Fix}$, $h'$ fixes all
$e_{a,v}$ and $z_{u',v}$, so $h'\in S_{u'}$.  Its matrix is $m_v(h')_{ab}=\langle w'_a,h_vw'_b\rangle$ ($\gamma_v$ unitary).  Now $h_vx'_b=x'_b$ and
$h_v\tilde w'_b=\tilde w'_b+\sum_c((M_v-1)\kappa(\tilde w'_b))_c\tilde w_c$ with $\kappa(\tilde w'_b)=(K_{db})_d$, so
$h_vw'_b=w'_b+\sum_c((M_v-1)K)_{cb}\tilde w_c$.  Since $\langle w'_a,\tilde w_c\rangle=\langle\tilde w'_a+x'_a,\tilde w_c\rangle=\overline{K_{ca}}$,
\[
m_v(h')_{ab}=\delta_{ab}+\sum_c\overline{K_{ca}}\,((M_v-1)K)_{cb}=(1+K^*(M_v-1)K)_{ab}=(K^*M_vK)_{ab}=(C_vM_vC_v^*)_{ab}.
\]
By (Levi) and (T) for $u'$ as in Step 4, $\pi_{u'}(\gamma^{-1}h\gamma)=\bigotimes_v\pi_{u',v}(C_vM_vC_v^*)$.  The matrices $C_v$ do not depend on $M$.
\end{proof}


\subsection{Pattern groups and the Steinberg restriction}\label{i3:ssec-grp-St}

We prove Propositions~\ref{i3:prop-admLq-in} and~\ref{i3:prop-St-rest}, together with the first half of (FD).

\subsubsection{Proof of Proposition~\ref{i3:prop-admLq-in}}\label{i3:ssec-pf-admLq-in}

Throughout, $g$ is a prime, $G_n=GL_n(\mathbb F_g)$, rows and columns are indexed by $0,\dots,n-1$, $U_n\subset G_n$ is
the upper unitriangular group and $M_n=n(n-1)/2=\log_g|U_n|$.  For a subgroup $V\subset U_n$ and a homomorphism
$\phi:V\to\mathbb F_g$ put $e(V,\phi)=|V|^{-1}\sum_{u\in V}\psi(\phi(u))^{-1}u\in\C[G_n]$; thus $e_U(\psi)=e(U_r,\phi_U)$ with
$\phi_U(u)=\sum_{i<r-1}u_{i,i+1}$, and $e_J(\psi)=e(U_J,\phi_J)$ with $U_J=\{u\in U_r:u_{a,s+k}=0$ for $s\le a<s+k\}$
and
\[
\phi_J(u)=\sum_{i<s-1}u_{i,i+1}+[s\ge1]\bigl(u_{s-1,s+k}+u_{s+k,s+k+1}\bigr)+\sum_{i<s-1}u_{s+k+1+i,\,s+k+2+i}
+\sum_{a<k-1}u_{s+a,s+a+1}.
\]
For a class function $\varphi$ on $G_n$ write $\mathrm{wh}(\varphi)=\mathrm{tr}(e_U(\psi)|\varphi)$ (with $e_U$ of $G_n$;
$\mathrm{wh}=\varphi(1)$ on $G_0=1$).

\begin{lemma}\label{i3:lem-eVphi}
If $V\subset U_n$ is a subgroup and $\phi:V\to\mathbb F_g$ is additive, then $e(V,\phi)$ is an idempotent with identity
coefficient $|V|^{-1}$.
\end{lemma}

\begin{proof}
$\lambda=\psi\circ\phi$ is a character of $V$, and
$e^2=|V|^{-2}\sum_{u,u'}\lambda(uu')^{-1}uu'=|V|^{-1}\sum_y\lambda(y)^{-1}y=e$, since for fixed $u$ the map $u'\mapsto
uu'$ is a bijection of $V$.  The identity occurs only for $u=1$, with coefficient $|V|^{-1}\lambda(1)^{-1}=|V|^{-1}$.
\end{proof}

\begin{proof}[Proof of (J)]
Let $r=k+2s+1$.  For $u,u'\in U_r$ and $i<j$, $(uu')_{ij}=\sum_{i\le b\le j}u_{ib}u'_{bj}$.  (a) $U_J$ is a
subgroup: it contains $1$, and for $s\le a<s+k$, $(uu')_{a,s+k}=\sum_{a\le b\le s+k}u_{ab}u'_{b,s+k}$, where the terms
with $a\le b<s+k$ vanish because $u'_{b,s+k}=0$ ($s\le b<s+k$) and the term $b=s+k$ is $u_{a,s+k}=0$; it is finite, hence
closed under inverses.  (b) $\phi_J$ is additive on $U_J$: a superdiagonal entry is additive on $U_r$,
$(uu')_{i,i+1}=u_{i,i+1}+u'_{i,i+1}$; and for the jump entry,
$(uu')_{s-1,s+k}=u'_{s-1,s+k}+\sum_{s\le b<s+k}u_{s-1,b}u'_{b,s+k}+u_{s-1,s+k}=u_{s-1,s+k}+u'_{s-1,s+k}$ on $U_J$.
(c) $U_J$ is cut out of $U_r$ by the vanishing of the $k$ coordinates $u_{a,s+k}$, $s\le a<s+k$, all strictly above the
diagonal, and the coordinates above the diagonal are free on $U_r$; so $|U_J|=g^{M_r-k}$.  Lemma~\ref{i3:lem-eVphi} gives
(J): $e_J(\psi)$ is idempotent with identity coefficient $g^{k}g^{-M_r}$.
\end{proof}

\emph{Jacquet restriction and derivatives.}  For $n=a+b$ let $P_{a,b}\subset G_n$ be the block upper triangular parabolic
with Levi $G_a\times G_b$ and unipotent radical $N_{a,b}$ ($|N_{a,b}|=g^{ab}$).  For a class function $\chi$ on $G_n$ put
\[
\mathrm{res}_{a,b}\chi(h_1,h_2)=g^{-ab}\sum_{x\in P_{a,b},\ x\mapsto(h_1,h_2)}\chi(x)=\frac1{|N_{a,b}|}\sum_{y\in N_{a,b}}
\chi\bigl(\mathrm{diag}(h_1,h_2)y\bigr),
\]
and for $0\le i\le n$ define the class function $\chi^{[i]}$ on $G_i$ by $\chi^{[i]}(h)=\mathrm{wh}(\mathrm{res}_{i,n-i}
\chi(h,\cdot))$ (Whittaker trace in the second variable).  If $\chi$ is the character of a representation $V$, then
$\mathrm{res}_{a,b}\chi$ is the character of the representation of $G_a\times G_b$ on the $N_{a,b}$-invariants $V^{N_{a,b}}$
(the operator $|N_{a,b}|^{-1}\sum_y\rho(y)$ is the projection onto $V^{N_{a,b}}$, and it commutes with the Levi).
Hence $\mathrm{res}_{a,b}\chi=\sum_l\alpha_l\otimes\beta_l$ with irreducible characters $\alpha_l,\beta_l$, and
\[
\mathrm{wh}(\chi^{[i]})=\sum_l\mathrm{tr}(e_U|\alpha_l)\,\mathrm{tr}(e_U|\beta_l)\in\Z_{\ge0},\qquad
\mathrm{wh}(\chi^{[0]})=\mathrm{wh}(\chi),\tag{$\dagger$}
\]
since the trace of an idempotent in a representation is the rank of an idempotent operator (fact (F3) in the proof of
Proposition~\ref{i3:prop-admLq}), and $\mathrm{res}_{0,n}\chi(1,\cdot)=\chi$.

\emph{Harish-Chandra induction.}  For class functions $\alpha$ on $G_a$, $\beta$ on $G_b$, $n=a+b$, let
$\alpha\circ\beta(x)=|P_{a,b}|^{-1}\sum_{y\in G_n,\ y^{-1}xy\in P_{a,b}}(\alpha\otimes\beta)(\overline{y^{-1}xy})$, the
bar denoting the Levi component; this is the character of $\mathrm{Ind}_{P_{a,b}}^{G_n}$ of the inflation of
$\alpha\otimes\beta$, and $\sigma_1\circ\cdots\circ\sigma_t$ is defined by iteration from the right.

\begin{proposition}\label{i3:prop-mirder}
Let $\psi\ne1$, $r=k+2s+1$ with $k\ge1$, and let $e'_J(\psi)=e(U',\phi_U)$ with $U'=\{u\in U_r:u_{a,k}=0$ for $a<k\}$.
For every class function $\chi$ on $G_r$:
\begin{enumerate}
\item[(Ex)] $\mathrm{tr}(e_J(\psi)|\chi)=\mathrm{tr}(e'_J(\psi)|\chi)$;
\item[(Der)] $\mathrm{tr}(e'_J(\psi)|\chi)=\sum_{i<k}\mathrm{wh}(\chi^{[i]})\,(g^{k-i}-g^{k-i-1})+\mathrm{wh}(\chi^{[k]})$.
\end{enumerate}
Moreover:
\begin{enumerate}
\item[(BZ)] if $\sigma_1,\dots,\sigma_t$ are cuspidal irreducible characters of $G_{d_1},\dots,G_{d_t}$ with
$\mathrm{wh}(\sigma_j)=1$ and $\tau=\sigma_1\circ\cdots\circ\sigma_t$ on $G_n$, $n=\sum d_j$, then $\mathrm{wh}(\tau)=1$ and
$\mathrm{wh}(\tau^{[i]})=\#\{I\subset\{1,\dots,t\}:\sum_{j\in I}d_j=i\}$ for $0\le i\le n$.
\end{enumerate}
\end{proposition}

\noindent (Ex) is proved in \S\ref{i3:sss-ex} by a chain of $s$ root exchanges; (Der) is proved in \S\ref{i3:sss-der}
from the finite-field form of the Bernstein--Zelevinsky filtration of the restriction to the mirabolic subgroup, which we
prove there, together with the multiplicativity of the Whittaker trace under Harish-Chandra induction (from the Bruhat
decomposition) and the value of the Whittaker trace of the Gelfand--Graev characters; (BZ) is proved in \S\ref{i3:sss-bz}
from Zelevinsky's theorem that the Harish-Chandra restriction is multiplicative (Theorem~\ref{i3:thm-zel}).

\begin{proof}[Proof of (Mir)]
Let $\chi$ be the character of a representation of $G_r$; it is a class function.  By (Ex) and (Der),
\[
\mathrm{tr}(e_J|\chi)=\mathrm{wh}(\chi)(g^k-g^{k-1})+\sum_{0<i<k}\mathrm{wh}(\chi^{[i]})(g^{k-i}-g^{k-i-1})
+\mathrm{wh}(\chi^{[k]}),
\]
using $\mathrm{wh}(\chi^{[0]})=\mathrm{wh}(\chi)$ from $(\dagger)$.  By $(\dagger)$ every $\mathrm{wh}(\chi^{[i]})$ is a
nonnegative integer, and $g^{k-i}-g^{k-i-1}\ge0$.  Taking real parts gives
$(g^k-g^{k-1})\mathrm{Re}\,\mathrm{tr}(e_U|\chi)\le\mathrm{Re}\,\mathrm{tr}(e_J|\chi)$.
\end{proof}

\emph{Published facts for (Sq).}  We use, for $G_n$ with $g\ge5$:
\begin{enumerate}
\item[(HC)] (Harish-Chandra theory for $GL_n(\mathbb F_q)$; Green 1955, Howlett--Lehrer, Invent.\ Math.\ 58 (1980);
references to be confirmed against the source.)  Every irreducible character of $G_n$ is a constituent of
$\sigma_1\circ\cdots\circ\sigma_t$ for cuspidal $\sigma_j$ of $G_{d_j}$, the multiset $\{\sigma_j\}$ being unique; if the
$\sigma_j$ are pairwise distinct, $\sigma_1\circ\cdots\circ\sigma_t$ is irreducible.
\item[(GG)] (Gelfand--Graev, Lemma~\ref{i3:S-lem-GGmult}.)  For irreducible $\chi$, $\mathrm{wh}(\chi)\in\{0,1\}$; and a cuspidal
irreducible character $\sigma$ of $G_d$ with $\sigma(1)=\prod_{i=1}^{d-1}(g^i-1)$ has $\mathrm{wh}(\sigma)=1$.
\end{enumerate}
\emph{The labels of the datum.}  The normalised-types compatibility fixes a \emph{cuspidal labelling}: a map $\ell$ which assigns
to every irreducible cuspidal representation $V$ of $G_d$ ($d\ge1$) a polynomial $\ell_d(V)\in\mathbb F_g[X]$, such that
(Lab1) $\ell_d(V)$ depends only on the isomorphism class of $V$; (Lab2) $\ell_d(V)$ is monic, irreducible, of degree $d$, and
$\ne X$ (so its constant term is nonzero); (Lab3) $\ell_d(V)=\ell_d(V')$ implies $V\cong V'$.  The label of an irreducible
representation $\tau$ of $G_r$ is $\prod_c\ell(V_c)$ for a family $(V_c)$ of irreducible cuspidal representations of
$G_{d_c}$, $\sum d_c=r$, such that $\tau$ is a constituent of $V_1\circ\cdots\circ V_t$; by (HC) the multiset $\{V_c\}$ is
unique up to isomorphism, so the label is well defined.  Write $\mathrm{Irr}_d$ for the set of monic irreducible polynomials of
degree $d$ over $\mathbb F_g$ with nonzero constant term.

\begin{lemma}[Comparison of labels]\label{i3:lem-labcomp}
Assume (G).  For every $d\ge1$ there is a bijection $s_d$ of $\mathrm{Irr}_d$ such that every irreducible cuspidal representation
$V$ of $G_d$ has character $\chi_{s_d(\ell_d(V))}$.
\end{lemma}

\begin{proof}
By (G), $P\mapsto\chi_P$ is a bijection of $\mathrm{Irr}_d$ onto the irreducible cuspidal characters of $G_d$; for each $P$ choose a
representation $V_P$ with character $\chi_P$, and put $p(P)=\ell_d(V_P)$, an element of $\mathrm{Irr}_d$ by (Lab2).  If
$p(P)=p(Q)$ then $V_P\cong V_Q$ by (Lab3), so $\chi_P=\chi_Q$ and $P=Q$.  So $p$ is an injective map of the finite set $\mathrm{Irr}_d$
into itself, hence bijective; let $s_d=p^{-1}$.  If $V$ is irreducible cuspidal, its character is $\chi_P$ for some $P$, so
$V\cong V_P$ (a representation of a finite group is determined by its character), $\ell_d(V)=\ell_d(V_P)=p(P)$ by (Lab1), and
$s_d(\ell_d(V))=P$.
\end{proof}

\noindent Extend $s=(s_d)$ multiplicatively to the monic polynomials with nonzero constant term: $s(\prod_jP_j)=\prod_js(P_j)$.  It is a
bijection preserving degrees, the degrees of the irreducible factors, squarefreeness and $\mathrm{sq}$.  By the lemma, an irreducible
$\tau$ has label $f$ (for the datum) iff its cuspidal support consists of representations with characters $\chi_{s(P_j)}$, $P_j$ the
irreducible factors of $f$; that is, iff $s(f)$ is its label in Green's sense (the product of the $Q_j$ with $\sigma_j=\chi_{Q_j}$ in
(HC)).  So (ii), (iv) for the labels of the datum are equivalent to (ii), (iv) for Green's labels, and we use Green's labels.  (HC)
gives in addition: for squarefree $f=P_1\cdots P_t$ (distinct $P_j$), the irreducible characters with Green label $f$ are exactly
$\tau_f=\chi_{P_1}\circ\cdots\circ\chi_{P_t}$.

\begin{proof}[Proof of (Sq)]
Let $f=P_1\cdots P_t$ be a squarefree label of degree $r$, $d_j=\deg P_j$, and $\tau=\tau_f$.  By (HC) the series
idempotent is $\varepsilon_f=z_\tau$, so by (F2) of the proof of Proposition~\ref{i3:prop-admLq}, for $e=e_U,e_J$,
$\mathrm{wt}(e,f)=\frac{\tau(1)}{|G_r|}\mathrm{Re}\,\mathrm{tr}(e|\tau)/e(1)$.  By (GG) each $\chi_{P_j}$ has
$\mathrm{wh}=1$, so (BZ) applies: $\mathrm{tr}(e_U|\tau)=1$ and $\mathrm{wh}(\tau^{[i]})=c_i:=\#\{I:\sum_{j\in I}d_j=i\}$.

(ii) The degree of the induced character is $\tau(1)=[G_r:P]\prod_j\chi_{P_j}(1)$, $P$ the parabolic of type $(d_j)$.
With $|G_n|=g^{M_n}\prod_{i=1}^n(g^i-1)$ and $|P|=g^{M_r}\prod_j\prod_{i=1}^{d_j}(g^i-1)$,
$\tau(1)=\prod_{i=1}^r(g^i-1)\big/\prod_j(g^{d_j}-1)$.  Hence
$\mathrm{wt}(e_U,f)=\frac{\tau(1)}{|G_r|}\cdot1\cdot g^{M_r}=\prod_j(g^{d_j}-1)^{-1}=\mathrm{sq}(f)$.

(iv) By (Ex), (Der) and (BZ), $\mathrm{tr}(e_J|\tau)=\sum_{i<k}c_i(g^{k-i}-g^{k-i-1})+c_k\le\sum_{i\le k}c_ig^{k-i}
\le g^k\sum_{I}g^{-\sum_{j\in I}d_j}=g^k\prod_j(1+g^{-d_j})\le g^k(1+\frac1g)^t\le g^k(1+\frac1g)^r$.  With
$e_J(1)=g^kg^{-M_r}$ from (J) and $e_U(1)=g^{-M_r}$,
$\mathrm{wt}(e_J,f)=\frac{\tau(1)}{|G_r|}g^{M_r}g^{-k}\mathrm{tr}(e_J|\tau)\le(1+\frac1g)^r\,\mathrm{wt}(e_U,f)$.
\end{proof}

\noindent Proposition~\ref{i3:prop-admLq-in} is thus proved in this subsection from (G) and the cited (HC), (GG),
Theorem~\ref{i3:thm-zel} and the Bruhat decomposition.  The formal development proves its own versions of (HC) and (GG) for
$GL_n(\mathbb F_g)$, $g\ge5$; they are special cases of the published statements.

\medskip\noindent\label{i3:sss-ex}\emph{Pattern groups and the proof of (Ex).}

A \emph{pattern} on $\{0,\dots,n-1\}$ is a strict partial order $\Phi$ (an irreflexive transitive relation), viewed as a
set of positions $(a,b)$.  Put $U_\Phi=\{x\in G_n:x_{ab}=\delta_{ab}$ unless $(a,b)\in\Phi\}$, and for a set $E$ of
positions put $\sigma_E(x)=\sum_{(a,b)\in E}x_{ab}$.  For a class function $\chi$ on $G_n$ and a subgroup $V$ with a
character $\lambda$, the element $e=|V|^{-1}\sum_{u\in V}\lambda(u)^{-1}u$ satisfies
$\mathrm{tr}(e|\chi)=|V|^{-1}\sum_{u\in V}\lambda(u)^{-1}\chi(u)$.

\begin{lemma}\label{i3:lem-patgp}
Let $\Phi$ be a pattern.  (a) $U_\Phi$ is a subgroup of $G_n$ of order $g^{|\Phi|}$.  (b) If $E$ is a set of positions
off the diagonal containing no $(a,c)$ for which there is $b$ with $(a,b),(b,c)\in\Phi$, then $\sigma_E$ is additive on
$U_\Phi$; thus $\lambda_E=\psi\circ\sigma_E$ is a character of $U_\Phi$, and we write $e(\Phi,E)=|U_\Phi|^{-1}\sum_{u\in
U_\Phi}\lambda_E(u)^{-1}u$.
\end{lemma}

\begin{proof}
Let $x,x'$ have the shape of $U_\Phi$.  In $(xx')_{ac}=\sum_bx_{ab}x'_{bc}$ a term with $b\notin\{a,c\}$ is non-zero
only if $(a,b),(b,c)\in\Phi$, hence $(a,c)\in\Phi$; for $a=c$ such a $b$ would give $(a,a)\in\Phi$, so
$(xx')_{aa}=1$.  Hence $xx'$ has the shape of $U_\Phi$, and for $a\ne c$,
$(xx')_{ac}=x_{ac}+x'_{ac}+\sum_{b:(a,b),(b,c)\in\Phi}x_{ab}x'_{bc}$, the last sum being empty when $(a,c)\in E$; with
$\sigma_E(1)=0$ this gives (b).  After renumbering the indices along a linear order extending $\Phi$, matrices of this
shape are unitriangular, hence invertible; a finite set of invertible matrices containing $1$ and closed under
products is a group.  The entries at the positions of $\Phi$ are free, which gives the order.
\end{proof}

\noindent Thus $e_U(\psi)=e(\Phi_<,E_<)$ with $\Phi_<=\{a<b\}$ and $E_<=\{(a,a+1)\}$; $e'_J(\psi)=e(\Phi',E_<)$ with
$\Phi'=\Phi_<\smallsetminus\{(a,k):a<k\}$; and $e_J(\psi)=e(\Phi_J,E_J)$ with $\Phi_J=\Phi_<\smallsetminus\{(a,s+k):s\le
a<s+k\}$ and $E_J$ the set of positions whose entries are summed in $\phi_J$.

\begin{lemma}[Renumbering]\label{i3:lem-renum}
Let $\pi$ be a permutation of $\{0,\dots,n-1\}$ and $\pi\Phi=\{(\pi a,\pi b):(a,b)\in\Phi\}$.  Under the hypotheses of
Lemma~\ref{i3:lem-patgp}(b), $\mathrm{tr}(e(\Phi,E)|\chi)=\mathrm{tr}(e(\pi\Phi,\pi E)|\chi)$ for every class function
$\chi$ on $G_n$.
\end{lemma}

\begin{proof}
For the permutation matrix $w$ with $we_j=e_{\pi(j)}$ one has $(wxw^{-1})_{\pi a,\pi b}=x_{ab}$, so $x\mapsto wxw^{-1}$ maps
$U_\Phi$ onto $U_{\pi\Phi}$ with $\sigma_{\pi E}(wxw^{-1})=\sigma_E(x)$; and $\chi(wxw^{-1})=\chi(x)$.
\end{proof}

\begin{lemma}[Exchange]\label{i3:lem-exch}
Let $G$ be a finite group, $H_1,H_2$ subgroups with $|H_1|=|H_2|$, and $\lambda_j$ a character of $H_j$, with
$\lambda_1=\lambda_2$ on $H_1\cap H_2$.  Assume: for every $h\in H_2\smallsetminus H_1$ there is $a\in H_1$ with
$h^{-1}ah\in H_1$ and $\lambda_1(h^{-1}ah)\ne\lambda_1(a)$; and for every $h\in H_1\smallsetminus H_2$ there is $b\in H_2$
with $h^{-1}bh\in H_2$ and $\lambda_2(h^{-1}bh)\ne\lambda_2(b)$.  Then for every class function $\varphi$ on $G$,
$|H_1|^{-1}\sum_{a\in H_1}\lambda_1(a)^{-1}\varphi(a)=|H_2|^{-1}\sum_{b\in H_2}\lambda_2(b)^{-1}\varphi(b)$.
\end{lemma}

\begin{proof}
Let $T_j=\sum_{a\in H_j}\lambda_j(a)^{-1}\varphi(a)$ and, for $h\in G$, $F(h)=\sum_{a\in H_1}\lambda_1(a)^{-1}\varphi(ah)$.
If $h\in H_1$, the substitution $a\mapsto ah^{-1}$ gives $F(h)=\lambda_1(h)T_1$.  If $h\in H_2\smallsetminus H_1$, take
$a_0$ as in the hypothesis and $a_1=h^{-1}a_0h$; the substitution $a\mapsto a_1^{-1}aa_0$ and
$\varphi(a_1^{-1}aa_0h)=\varphi(a_1^{-1}(ah)a_1)=\varphi(ah)$ give $F(h)=\lambda_1(a_1)\lambda_1(a_0)^{-1}F(h)$, so
$F(h)=0$.  Hence $\sum_{b\in H_2}\lambda_2(b)^{-1}F(b)=|H_1\cap H_2|\,T_1$, using $\lambda_1=\lambda_2$ on $H_1\cap H_2$.  By
symmetry $\sum_{a\in H_1}\lambda_1(a)^{-1}\sum_{b\in H_2}\lambda_2(b)^{-1}\varphi(ba)=|H_1\cap H_2|\,T_2$.  The two double
sums agree because $\varphi(ab)=\varphi(ba)$.  Since $|H_1\cap H_2|\ge1$ and $|H_1|=|H_2|$, the claim follows.
\end{proof}

\begin{lemma}[Root exchange]\label{i3:lem-rootexch}
Let $\Phi_1,\Phi_2$ be patterns on $\{0,\dots,n-1\}$ whose union $\Phi$ is transitive, let $E$ satisfy the hypothesis of
Lemma~\ref{i3:lem-patgp}(b) for $\Phi_1$ and for $\Phi_2$, and let $W$ be a set of indices and $p,q\notin W$ with $(p,q)\in
E$, such that:
\begin{enumerate}
\item[(i)] $\Phi_1\smallsetminus\Phi_2=\{(p,w):w\in W\}$ and $\Phi_2\smallsetminus\Phi_1=\{(w,q):w\in W\}$;
\item[(ii)] for $w\in W$: if ($x=p$ or $(x,p)\in\Phi$) and ($y=w$ or $(w,y)\in\Phi$), then $(x,y)\in\Phi_1$, and
$(x,y)\in E$ only if $(x,y)=(p,q)$;
\item[(iii)] for $w\in W$: if ($x=w$ or $(x,w)\in\Phi$) and ($y=q$ or $(q,y)\in\Phi$), then $(x,y)\in\Phi_2$, and
$(x,y)\in E$ only if $(x,y)=(p,q)$.
\end{enumerate}
Then $\mathrm{tr}(e(\Phi_1,E)|\chi)=\mathrm{tr}(e(\Phi_2,E)|\chi)$ for every class function $\chi$ on $G_n$.
\end{lemma}

\begin{proof}
We apply Lemma~\ref{i3:lem-exch} to $H_j=U_{\Phi_j}$ and $\lambda_j=\lambda_E$ inside $U_\Phi$.  By (i),
$|\Phi_1|=|\Phi_1\cap\Phi_2|+|W|=|\Phi_2|$, so $|H_1|=|H_2|$ by Lemma~\ref{i3:lem-patgp}(a); the two characters are given by
the same formula.  Let $h\in H_2\smallsetminus H_1$.  Since $h\in U_\Phi$ and $h\notin U_{\Phi_1}$, $h$ has a non-zero
entry at a position of $\Phi\smallsetminus\Phi_1=\Phi_2\smallsetminus\Phi_1$: $h_{wq}\ne0$ for some $w\in W$.  For
$t\in\mathbb F_g$ let $a_0=1+tE_{pw}\in H_1$ (as $(p,w)\in\Phi_1$).  Then $h^{-1}a_0h=1+t\,(h^{-1}e_p)(e_w^{\mathsf T}h)$
has $(x,y)$-entry $\delta_{xy}+t\,(h^{-1})_{xp}h_{wy}$, and $(h^{-1})_{xp}h_{wy}\ne0$ forces ($x=p$ or $(x,p)\in\Phi$) and
($y=w$ or $(w,y)\in\Phi$), as $h,h^{-1}\in U_\Phi$.  By (ii), $h^{-1}a_0h\in H_1$ and
$\sigma_E(h^{-1}a_0h)=t\,(h^{-1})_{pp}h_{wq}=t\,h_{wq}$, while $\sigma_E(a_0)=0$ since $(p,w)\notin E$ (by (ii) with $x=p$,
$y=w$, as $w\ne q$).  Choosing $t$ with $\psi(th_{wq})\ne1$ gives $\lambda_E(h^{-1}a_0h)\ne\lambda_E(a_0)$.  If $h\in
H_1\smallsetminus H_2$, then $h^{-1}\in H_1\smallsetminus H_2$ has $(h^{-1})_{pw}\ne0$ for some $w\in W$; with
$b_0=1+tE_{wq}\in H_2$, (iii) gives in the same way $h^{-1}b_0h\in H_2$, $\sigma_E(h^{-1}b_0h)=t\,(h^{-1})_{pw}$ and
$\sigma_E(b_0)=0$.  Lemma~\ref{i3:lem-exch} and the trace formula above give the claim.
\end{proof}

\begin{proof}[Proof of (Ex)]
Let $W=\{a:s\le a<s+k\}$ and, for $0\le i\le s$, let $\Phi^{(i)}$ be the set of $(a,b)$ with: $a<b$ both in $W$; or $a<b$
both outside $W$; or $a<i$ and $b\in W$; or $a\in W$ and ($i<b<s$, or $b=s+k$ and $i<s$, or $b>s+k$).  Let $E$ be the set
of positions $(a,a+1)$ with $a+1<s$, with $a\ge s+k$, or with $a,a+1\in W$, together with $(s-1,s+k)$ if $s\ge1$.  The
following are direct verifications on the indices.
(1) Each $\Phi^{(i)}$ is a pattern and $E$ satisfies Lemma~\ref{i3:lem-patgp}(b) for it.
(2) $\Phi^{(s)}=\Phi_J$ and $\sigma_E=\phi_J$ on $U_J$ (the chain $0,\dots,s-1,s+k,\dots,r-1$ and the chain of $W$).
(3) For the permutation $\pi$ with $\pi(a)=s+a$ ($a<k$), $\pi(a)=a-k$ ($k\le a<k+s$), $\pi(a)=a$ ($a\ge k+s$), one has
$\pi\Phi'=\Phi^{(0)}$ and $\pi E'=E$ for $E'=\{(a,a+1):a+1\ne k\}$, and $\sigma_{E'}=\phi_U$ on $U'$ because $u_{k-1,k}=0$
there.
(4) For $1\le i\le s$, the patterns $\Phi_1=\Phi^{(i)}$, $\Phi_2=\Phi^{(i-1)}$ satisfy the hypotheses of
Lemma~\ref{i3:lem-rootexch} with $W$ as above, $p=i-1$, and $q=i$ if $i<s$, $q=s+k$ if $i=s$: their union is the set of
$(a,b)$ with $a<b$ in the same part ($W$ or its complement), or $a<i$ and $b\in W$, or $a\in W$, $b\notin W$ and $b\ge i$,
which is transitive; $\Phi^{(i)}\smallsetminus\Phi^{(i-1)}=\{(i-1,w):w\in W\}$ and
$\Phi^{(i-1)}\smallsetminus\Phi^{(i)}=\{(w,q):w\in W\}$; and in (ii), (iii) the hypotheses force $x\le i-1$, $x\notin W$,
and $y\in W$ or $y\ge i$, $y\notin W$, resp.\ $x\in W$ or $x<i$, and $y\ge q$, $y\notin W$, whence the membership, and the
only position of $E$ among them is $(i-1,i)$ for $i<s$, resp.\ $(s-1,s+k)$ for $i=s$, that is $(p,q)$.

By (2), Lemma~\ref{i3:lem-rootexch} applied $s$ times by (4), and Lemma~\ref{i3:lem-renum} with (3),
$\mathrm{tr}(e_J|\chi)=\mathrm{tr}(e(\Phi^{(s)},E)|\chi)=\dots=\mathrm{tr}(e(\Phi^{(0)},E)|\chi)=
\mathrm{tr}(e(\Phi',E')|\chi)=\mathrm{tr}(e'_J|\chi)$.
\end{proof}

\medskip\noindent\label{i3:sss-der}\emph{The mirabolic filtration and the proof of (Der).}

For a subgroup $H$ of a finite group $G$ and a class function $\varphi$ on $H$ put
$\mathrm{Ind}_H^G\varphi(x)=|H|^{-1}\sum_{y\in G,\ y^{-1}xy\in H}\varphi(y^{-1}xy)$; this is the induced class function,
which is linear in $\varphi$, transitive in stages, satisfies $\mathrm{Ind}_H^G(\varphi\cdot\xi|_H)=(\mathrm{Ind}_H^G\varphi)\xi$
for a linear character $\xi$ of $G$, and commutes with inflation from quotients by a normal subgroup contained in $H$
(Serre, \emph{Linear representations of finite groups}, \S\S3.3, 7.1--7.2; section numbers to be confirmed against the
source).  We also use the following special case of Mackey's formula, proved directly: \emph{if $G=HK$ for subgroups $H,K$ and
$\varphi$ is a class function on $H$, then $(\mathrm{Ind}_H^G\varphi)|_K=\mathrm{Ind}_{H\cap K}^K(\varphi|_{H\cap K})$.}
Indeed each $y\in G$ is $y=kh$ in exactly $|H\cap K|$ ways, and for $x\in K$,
$\sum_{y:y^{-1}xy\in H}\varphi(y^{-1}xy)=|H\cap K|^{-1}\sum_{h\in H}\sum_{k:k^{-1}xk\in H}\varphi(h^{-1}k^{-1}xkh)
=|H\cap K|^{-1}|H|\sum_{k:k^{-1}xk\in H\cap K}\varphi(k^{-1}xk)$.

\emph{Notation.}  For $m\le n$ let $\iota:G_m\to G_n$, $h\mapsto\mathrm{diag}(h,1_{n-m})$, and $\mathrm{tl}_m(x)$ the upper
left $m\times m$ block.  For $0\le i\le n$ let $R^n_i\subset P_{i,n-i}$ be the subgroup with lower right block in
$U_{n-i}$, $V^n_i\subset R^n_i$ the subgroup with upper left block $1$, so that $R^n_i=\iota(G_i)\ltimes V^n_i$ and
$|V^n_i|=g^{i(n-i)+M_{n-i}}$, and $\theta^n_i(x)=\psi(\sum_{i\le j<n-1}x_{j,j+1})$, a character of $R^n_i$ (the entries
$x_{j,j+1}$, $j\ge i$, lie in the unitriangular block and are additive) trivial on $\iota(G_i)$.  Unfolding $\mathrm{res}$
and $\mathrm{wh}$,
\[
\chi^{[i]}(h)=|V^n_i|^{-1}\sum_{v\in V^n_i}\chi(\iota(h)v)\,\theta^n_i(v)^{-1}.\tag{D0}
\]
For $1\le m\le n$ and $F:G_n\to\C$ define $F^{\downarrow m}:G_m\to\C$ by
$F^{\downarrow m}(p)=|V^n_m|^{-1}\sum_{v\in V^n_m}F(\iota(p)v)\,\vartheta(\iota(p)v)^{-1}$ with
$\vartheta(x)=\psi(\sum_{m-1\le j<n-1}x_{j,j+1})$.  Let $\mathrm{Mir}_n=\{x\in G_n:$ last row $(0,\dots,0,1)\}$, the
mirabolic subgroup; $R^n_i\subset\mathrm{Mir}_n$ for $i<n$.  Call $F$ \emph{mirabolic-central} if $F(pxp^{-1})=F(x)$ for
$p,x\in\mathrm{Mir}_n$; then $F|_{\mathrm{Mir}_n}$ is a linear combination of characters of representations of
$\mathrm{Mir}_n$.  For $i<n$ and a class function $\rho$ on $G_i$ put
$I^n_i(\rho)=\mathrm{Ind}_{R^n_i}^{\mathrm{Mir}_n}(\rho\otimes\theta^n_i)$, where
$(\rho\otimes\theta^n_i)(\iota(h)v)=\rho(h)\theta^n_i(v)$ is a class function on $R^n_i$.  Finally
$\Gamma_m=\mathrm{Ind}_{U_m}^{G_m}\psi_U$ with $\psi_U(u)=\psi(\sum_ju_{j,j+1})$ is the Gelfand--Graev character
($\Gamma_0=1$).

\begin{lemma}\label{i3:lem-tail}
Let $1\le m\le n$.  (a) If $F$ is mirabolic-central (in particular, if $F$ is a class function) on $G_n$, then
$F^{\downarrow m}$ is mirabolic-central on $G_m$.  (b) For $i<m$, $(F^{\downarrow m})^{[i]}=F^{[i]}$ (derivatives defined by
(D0) for any $F$).
\end{lemma}

\begin{proof}
Let $Q'=\iota(\mathrm{Mir}_m)\ltimes V^n_m$.  For $p\in\mathrm{Mir}_m$ the row $m-1$ of $\iota(p)$ is $e_{m-1}$, so on $Q'$
the entries $x_{j,j+1}$, $j\ge m-1$, are additive and $\vartheta$ is a character of $Q'$.  (a) For $q\in\mathrm{Mir}_m$,
$\iota(q)$ normalises $V^n_m$, and the substitution $v\mapsto\iota(q)v\iota(q)^{-1}$ gives
$F^{\downarrow m}(qpq^{-1})=|V^n_m|^{-1}\sum_vF(\iota(q)\iota(p)v\iota(q)^{-1})\vartheta(\iota(q)\iota(p)v\iota(q)^{-1})^{-1}=
F^{\downarrow m}(p)$, since $\iota(q)\in\mathrm{Mir}_n$ when $m<n$ (and $F^{\downarrow n}=F$).  (b) Every $y\in V^n_i$ is
uniquely $y=\iota(w)v$ with $w\in V^m_i$, $v\in V^n_m$, and, the row $m-1$ of $w$ being $e_{m-1}$ ($i\le m-1$),
$\theta^n_i(y)=\theta^m_i(w)\,\psi(\sum_{j\ge m-1}v_{j,j+1})$, while for $p=\iota(h)w$ also
$\vartheta(\iota(p)v)=\psi(\sum_{j\ge m-1}v_{j,j+1})$.  Inserting the definition of $F^{\downarrow m}$ into (D0) for
$(F^{\downarrow m})^{[i]}$ gives (D0) for $F^{[i]}$, with $|V^n_i|=|V^m_i||V^n_m|$.
\end{proof}

\begin{lemma}[One step]\label{i3:lem-mirstep}
Let $m\ge1$ and $F$ mirabolic-central on $G_{m+1}$.  Let $Q=\{x\in\mathrm{Mir}_{m+1}:\mathrm{tl}_m(x)\in\mathrm{Mir}_m\}$ and
$\xi(x)=\psi(x_{m-1,m})$, a character of $Q$.  Then for $x\in\mathrm{Mir}_{m+1}$,
\[
F(x)=F^{[m]}(\mathrm{tl}_m x)+\mathrm{Ind}_Q^{\mathrm{Mir}_{m+1}}\bigl((F^{\downarrow m}\circ\mathrm{tl}_m)\cdot\xi\bigr)(x).
\]
\end{lemma}

\begin{proof}
Both sides are linear in $F|_{\mathrm{Mir}_{m+1}}$, so we may take $F$ the character of a representation $\tau$ of
$\mathrm{Mir}_{m+1}=\iota(G_m)\ltimes V$, $V=V^{m+1}_m\cong\mathbb F_g^m$ (last column).  The characters of $V$ are
$\xi_c(v)=\psi(c\cdot v)$, $c$ a row vector; $\iota(h)$ maps $\xi_c$ to $\xi_{ch^{-1}}$, so $G_m$ permutes the non-trivial
ones transitively, and the stabiliser of $\xi=\xi_{e_{m-1}}$ is $\mathrm{Mir}_m$; its stabiliser in $\mathrm{Mir}_{m+1}$ is
$Q$.  Decompose $\tau=\tau^V\oplus\bigoplus_{c\ne0}\tau_{\xi_c}$ into $V$-isotypic parts; the group permutes the
$\tau_{\xi_c}$ transitively with stabiliser $Q$, so $\bigoplus_{c\ne0}\tau_{\xi_c}\cong\mathrm{Ind}_Q^{\mathrm{Mir}_{m+1}}
\tau_\xi$ (Serre, \emph{op.\ cit.}, \S7.3; to be confirmed).  The projection onto $\tau^V$ is $|V|^{-1}\sum_v\tau(v)$, so
$\mathrm{ch}\,\tau^V(x)=|V|^{-1}\sum_vF(xv)=F^{[m]}(\mathrm{tl}_mx)$ by (D0) (for $x=\iota(h)v_0$, $xv$ runs over
$\iota(h)V$).  The projection onto $\tau_\xi$ is $|V|^{-1}\sum_v\xi(v)^{-1}\tau(v)$, so for $q\in Q$,
$\mathrm{ch}\,\tau_\xi(q)=|V|^{-1}\sum_v\xi(v)^{-1}F(qv)$.  Write $v_c\in V$ for the element with last column
$(c,1)$, so $v_cv_{c'}=v_{c+c'}$, and $q=\iota(p)v_{c_0}$ with $p\in\mathrm{Mir}_m$.  Since the row $m-1$ of $p$ is
$e_{m-1}$, $(\iota(p)v_c)_{m-1,m}=c_{m-1}$, so $\xi(q)=\psi(c_{0,m-1})$ and $\vartheta(\iota(p)v_c)=\psi(c_{m-1})$.  The
substitution $c\mapsto c-c_0$ gives $\mathrm{ch}\,\tau_\xi(q)=\psi(c_{0,m-1})\,|V|^{-1}\sum_c\psi(c_{m-1})^{-1}F(\iota(p)v_c)
=\xi(q)F^{\downarrow m}(p)$.  Adding the two
characters gives the claim.
\end{proof}

\begin{lemma}[Mirabolic filtration]\label{i3:lem-mirfilt}
Let $n\ge1$ and $F$ mirabolic-central on $G_n$.  Then $F(x)=\sum_{i<n}I^n_i(F^{[i]})(x)$ for $x\in\mathrm{Mir}_n$.
\end{lemma}

\noindent Here, for $F$ the character of a representation of $\mathrm{Mir}_n$, each $F^{[i]}$ is a character of $G_i$
(that of the space of $(V^n_i,\theta^n_i)$-equivariant vectors), and the identity extends by linearity.

\begin{proof}
Induction on $n$.  For $n=1$, $\mathrm{Mir}_1=\{1\}$ and $I^1_0(F^{[0]})(1)=F^{[0]}=F(1)$.  Let $n=m+1$ with $m\ge1$.  For
$i=m$, $R^{m+1}_m=\mathrm{Mir}_{m+1}$, $\theta^{m+1}_m=1$, and $I^{m+1}_m(\rho)(x)=\rho(\mathrm{tl}_mx)$ as $\rho$ is a class
function; this is the first term of Lemma~\ref{i3:lem-mirstep}.  By Lemma~\ref{i3:lem-tail}(a) and the induction hypothesis,
$F^{\downarrow m}=\sum_{i<m}I^m_i((F^{\downarrow m})^{[i]})$ on $\mathrm{Mir}_m$, and $(F^{\downarrow m})^{[i]}=F^{[i]}$ by
Lemma~\ref{i3:lem-tail}(b).  For $i<m$, $R^{m+1}_i=\iota(R^m_i)\ltimes V$ with $V\trianglelefteq Q$, and
$\theta^{m+1}_i(x)=\theta^m_i(\mathrm{tl}_mx)\,\xi(x)$ on $R^{m+1}_i$; hence, by the properties of induction recalled
above, $(I^m_i(\rho)\circ\mathrm{tl}_m)\cdot\xi=\mathrm{Ind}_{R^{m+1}_i}^Q(\rho\otimes\theta^{m+1}_i)$ and
$\mathrm{Ind}_Q^{\mathrm{Mir}_{m+1}}((I^m_i(\rho)\circ\mathrm{tl}_m)\cdot\xi)=I^{m+1}_i(\rho)$.  Lemma~\ref{i3:lem-mirstep}
gives the claim.
\end{proof}

\begin{lemma}\label{i3:lem-blockres}
Let $i\le k$ and $\rho$ a class function on $G_i$.  Then $I^{k+1}_i(\rho)(\iota(h))=(\rho\circ\Gamma_{k-i})(h)$ for $h\in
G_k$.
\end{lemma}

\begin{proof}
$R^{k+1}_i$ contains the last-column group, so $\mathrm{Mir}_{k+1}=R^{k+1}_i\,\iota(G_k)$, and $R^{k+1}_i\cap\iota(G_k)=
\iota(R^k_i)$, on which $\theta^{k+1}_i\circ\iota=\theta^k_i$ (the entry $(k-1,k)$ of $\iota(r)$ is $0$).  By the special
case of Mackey's formula, $I^{k+1}_i(\rho)\circ\iota=\mathrm{Ind}_{R^k_i}^{G_k}(\rho\otimes\theta^k_i)$.  Inducing in stages
through $P_{i,k-i}=(G_i\times G_{k-i})\ltimes N_{i,k-i}$, with $R^k_i=(G_i\times U_{k-i})\ltimes N_{i,k-i}$ and $\rho\otimes
\theta^k_i$ trivial on $N_{i,k-i}$, gives $\mathrm{Ind}_{P_{i,k-i}}^{G_k}$ of the inflation of
$\rho\otimes\mathrm{Ind}_{U_{k-i}}^{G_{k-i}}\psi_U=\rho\otimes\Gamma_{k-i}$, which is $\rho\circ\Gamma_{k-i}$.
\end{proof}

\noindent For $a+b=n$, an $(a,b)$-shuffle is a permutation $\sigma$ of $\{0,\dots,n-1\}$ increasing on $\{0,\dots,a-1\}$
and on $\{a,\dots,n-1\}$; $w_\sigma$ is the permutation matrix with $w_\sigma e_j=e_{\sigma(j)}$.

\begin{theorem}[Bruhat decomposition; Borel, \emph{Linear Algebraic Groups}, Thm.~14.12; reference to be confirmed]\label{i3:thm-bruhat}
$G_n$ is the disjoint union of the double cosets $Bw_\sigma B$, $\sigma\in S_n$, $B$ the upper triangular group.
\end{theorem}

\noindent Consequently every $x\in G_n$ is $x=uw_\sigma p$ with $u\in U_n$, $p\in P_{a,b}$ and $\sigma$ an $(a,b)$-shuffle,
$\sigma$ unique: indeed $Bw_\sigma B=U_nw_\sigma B$ ($B=TU_n$ and $w_\sigma^{-1}Tw_\sigma=T$), $P_{a,b}$ is the union of the
$BwB$ with $w\in S_a\times S_b$, and each coset $\sigma(S_a\times S_b)$ contains exactly one shuffle.

\begin{lemma}[Whittaker traces of products]\label{i3:lem-whmult}
Let $\psi\ne1$ and $\alpha,\beta$ class functions on $G_a,G_b$.  Then $\mathrm{wh}(\alpha\circ\beta)=\mathrm{wh}(\alpha)\,
\mathrm{wh}(\beta)$.
\end{lemma}

\begin{proof}
Let $n=a+b$, $P=P_{a,b}$, $U=U_n$ and $\varphi(x)=\alpha(\mathrm{tl}_ax)\beta(x')$ on $P$, $x'$ the lower right block; it is
a class function on $P$.  By definition,
$\mathrm{wh}(\alpha\circ\beta)=|U|^{-1}|P|^{-1}\sum_{y\in G_n}S(y)$ with
$S(y)=\sum_{u\in U\cap yPy^{-1}}\psi_U(u)^{-1}\varphi(y^{-1}uy)$.  For $y'=u_0yp_0$ the substitution $u\mapsto u_0uu_0^{-1}$
shows $S(y')=S(y)$ ($\psi_U$ is a character of $U$, $\varphi$ a class function on $P$).  By the consequence of
Theorem~\ref{i3:thm-bruhat}, $|Uw_\sigma P|=|U||P|/|U\cap w_\sigma Pw_\sigma^{-1}|$ and
$\mathrm{wh}(\alpha\circ\beta)=\sum_\sigma S(w_\sigma)/|U\cap w_\sigma Pw_\sigma^{-1}|$ over the $(a,b)$-shuffles.  If
$\sigma$ has $i<a\le j$ with $\sigma(j)=\sigma(i)+1=:\ell+1$, the group $X=\{1+tE_{\ell,\ell+1}\}\subset U$ satisfies
$w_\sigma^{-1}Xw_\sigma=\{1+tE_{ij}\}\subset N_{a,b}$; so $X\subset U\cap w_\sigma Pw_\sigma^{-1}$, $\varphi(w_\sigma^{-1}uxw_\sigma)=
\varphi(w_\sigma^{-1}uw_\sigma)$ for $x\in X$, and $\psi_U(u(1+tE_{\ell,\ell+1}))=\psi_U(u)\psi(t)$; summing over $t$ gives
$S(w_\sigma)=0$.  A shuffle without such a pair has all of $\sigma(\{a,\dots,n-1\})$ before $\sigma(\{0,\dots,a-1\})$, i.e.
$\sigma=\sigma_0$ with $\sigma_0(j)=b+j$ ($j<a$), $\sigma_0(j)=j-a$ ($j\ge a$).  Then $U\cap w_{\sigma_0}Pw_{\sigma_0}^{-1}=
\{\mathrm{diag}(u_2,u_1):u_2\in U_b,u_1\in U_a\}$, $w_{\sigma_0}^{-1}\mathrm{diag}(u_2,u_1)w_{\sigma_0}=\mathrm{diag}(u_1,u_2)$, and
$\psi_U(\mathrm{diag}(u_2,u_1))=\psi_U(u_2)\psi_U(u_1)$; so $S(w_{\sigma_0})/|U_aU_b|=\mathrm{wh}(\alpha)\mathrm{wh}(\beta)$.
\end{proof}

\begin{lemma}\label{i3:lem-whGG}
$\mathrm{wh}(\Gamma_0)=1$ and $\mathrm{wh}(\Gamma_m)=g^m-g^{m-1}$ for $m\ge1$, if $\psi\ne1$.
\end{lemma}

\begin{proof}
Let $\delta_n$ be the indicator of $1\in G_n$; $\mathrm{wh}(\delta_n)=|U_n|^{-1}=g^{-M_n}$.  By (D0),
$\delta_{k+1}^{[i]}=g^{-(i(k+1-i)+M_{k+1-i})}\delta_i$.  Lemmas~\ref{i3:lem-mirfilt} and \ref{i3:lem-blockres} for the class
function $\delta_{k+1}$ at $\iota(h)$, and Lemma~\ref{i3:lem-whmult}, give
$g^{-M_k}=\mathrm{wh}(\delta_k)=\sum_{i\le k}g^{-(i(k+1-i)+M_{k+1-i}+M_i)}\mathrm{wh}(\Gamma_{k-i})=g^{-M_{k+1}}\sum_{j\le
k}\mathrm{wh}(\Gamma_j)$.  Hence $\sum_{j\le k}\mathrm{wh}(\Gamma_j)=g^k$ for all $k\ge0$; take differences.
\end{proof}

\begin{proof}[Proof of (Der)]
Every $u\in U'$ is uniquely $u=\iota(\iota(h))v$ with $h\in U_k$, $v\in V^r_{k+1}$ (the upper left $(k+1)$-block of $u$ is
$\mathrm{diag}(h,1)$), and $\phi_U(u)=\sum_{j<k-1}h_{j,j+1}+\sum_{j\ge k}u_{j,j+1}$.  Hence
$\mathrm{tr}(e'_J|\chi)=\mathrm{wh}(T)$ with $T(h)=\chi^{\downarrow(k+1)}(\iota(h))$.  By Lemma~\ref{i3:lem-tail}(a),
$F=\chi^{\downarrow(k+1)}$ is mirabolic-central, so by Lemmas~\ref{i3:lem-mirfilt}, \ref{i3:lem-blockres} and
\ref{i3:lem-tail}(b), $T=\sum_{i\le k}\chi^{[i]}\circ\Gamma_{k-i}$; each $\chi^{[i]}$ is a class function on $G_i$
(conjugation by $\iota(h')$ permutes the fibres of $P_{i,r-i}\to G_i\times G_{r-i}$).  Lemmas~\ref{i3:lem-whmult} and
\ref{i3:lem-whGG} give (Der).
\end{proof}

\medskip\noindent\label{i3:sss-bz}\emph{The derivatives of a product of cuspidals: proof of (BZ).}

Let $\mathcal R_n$ be the space of class functions on $G_n$ ($\mathcal R_0=\C$), with the product $\circ$, and for $c+d=n$
the restriction $\mathrm{res}_{c,d}:\mathcal R_n\to\mathcal R_c\otimes\mathcal R_d$ (class functions on $G_c\times G_d$).

\begin{theorem}[Zelevinsky, \emph{Representations of finite classical groups}, LNM 869 (1981), Ch.~III (the representation
ring of the groups $GL_n(\mathbb F_q)$ is a positive self-adjoint Hopf algebra); reference to be confirmed]\label{i3:thm-zel}
The product $\circ$ on $\bigoplus_n\mathcal R_n$ is associative and commutative, and the restriction is multiplicative: for
$\alpha\in\mathcal R_a$, $\beta\in\mathcal R_b$ and $c+d=a+b$,
$\mathrm{res}_{c,d}(\alpha\circ\beta)=\sum\mathrm{res}_{a_1,a_2}(\alpha)\circ\mathrm{res}_{b_1,b_2}(\beta)$, the sum over
$a_1+a_2=a$, $b_1+b_2=b$, $a_1+b_1=c$, where $(\alpha_1\otimes\alpha_2)\circ(\beta_1\otimes\beta_2)=(\alpha_1\circ\beta_1)
\otimes(\alpha_2\circ\beta_2)$.
\end{theorem}

\noindent Zelevinsky states this for the Grothendieck groups of representations, with induction and restriction normalised
as here (characters of $\mathrm{Ind}_P^G$ of inflations, and of $N$-invariants); it extends to class functions by
linearity.

\begin{proof}[Proof of (BZ)]
A cuspidal irreducible character $\sigma$ of $G_d$ ($d\ge1$) satisfies $\mathrm{res}_{c,d-c}\sigma=0$ for $0<c<d$: it is
the character of the $N_{c,d-c}$-invariants, a space of dimension $\mathrm{res}_{c,d-c}\sigma(1,1)=0$; and
$\mathrm{res}_{0,d}\sigma=1\otimes\sigma$, $\mathrm{res}_{d,0}\sigma=\sigma\otimes1$.  By Theorem~\ref{i3:thm-zel} and induction on
$t$, $\mathrm{res}_{i,n-i}\tau=\sum_{I}\sigma_I\otimes\sigma_{I^c}$ over $I\subset\{1,\dots,t\}$ with
$\sum_{j\in I}d_j=i$, where $\sigma_I$ is the product of the $\sigma_j$, $j\in I$ (empty product $1$).  By
Lemma~\ref{i3:lem-whmult}, $\mathrm{wh}(\sigma_I)=\prod_{j\in I}\mathrm{wh}(\sigma_j)=1$; in particular
$\mathrm{wh}(\tau)=1$, $\tau^{[i]}=\sum_I\sigma_I\,\mathrm{wh}(\sigma_{I^c})=\sum_I\sigma_I$, and
$\mathrm{wh}(\tau^{[i]})=\#\{I:\sum_{j\in I}d_j=i\}$.
\end{proof}



\subsubsection{Proof of Proposition~\ref{i3:prop-St-rest} and of the first half of (FD)}\label{i3:ssec-pf-St-rest}

We use the notation of \S\ref{i3:ssec-pf-kppoff}: $V=E^n$ with the hermitian form $h$, the vectors $e_a,f_a$ ($a<r_S$), $z_u$, the
term $t$ with row lattices $\Lambda,\Lambda'$ and representative $\gamma$ ($\gamma\Lambda'=\Lambda$), $N_M$, $\Gamma_t$, $R_t$,
$\mathrm{Fix}_t$, $S_t$.  Write $r=r_S$.

\emph{A graded basis.}  Let $W=\{y\in V:h(e_a,y)=h(f_a,y)=0\ \forall a\}$; as the $e_a,f_a$ span hyperbolic planes, $V=\mathrm{span}(e_a)\oplus
W\oplus\mathrm{span}(f_a)$.  Fix an $E$-basis $(w'_i)$ of $W$ and let $(b_i)_{i\in\mathcal K}$ be the $E$-basis of $V$ formed by the $e_a$, the
$w'_i$ and the $f_a$, with \emph{levels} $\ell(e_a)=a+1$, $\ell(w'_i)=r+1$, $\ell(f_a)=2r+1-a$; so $1\le\ell\le2r+1$.  For $m\ge0$ let
$\mathfrak p_m$ be the set of square matrices $Z$ indexed by $\mathcal K$ (over $E$ or over $\C$) with $Z_{ij}=0$ unless $\ell(j)\ge\ell(i)+m$;
thus $Z\in\mathfrak p_m$ lowers levels by at least $m$, $\mathfrak p_m\mathfrak p_{m'}\subset\mathfrak p_{m+m'}$, and $\mathfrak p_{2r+1}=0$.  At an
infinite place $v$ let $B_v$ be the matrix of the vectors $(b_i)_v$ and, for $s\in U_\infty$, $M_v(s)=B_v^{-1}s_vB_v$; let $\Xi$ be the
matrix of $\gamma$ in the basis $(b_i)$ (with entries in $E$) and $\Xi_v$ its image under $v$, so that $M_v(\gamma)=\Xi_v$.  Let $J$ be the
Gram matrix $J_{ij}=h(b_i,b_j)$ and $J_v$ its image.  From $h(e_a,e_b)=h(f_a,f_b)=0$, $h(e_a,f_b)=\delta_{ab}$ and $W\perp e_a,f_a$:
$J_{ij}\ne0$ only if $\ell(i)+\ell(j)=2r+2$, and $J^{-1}$ has the same shape.  For a matrix $A$ put $A^\#=J_v^{-1}A^*J_v$ (the adjoint
for $h_v$; over $E$, $A^\#=J^{-1}\bar A^{\rm T}J$); then $A$ is the matrix of an isometry iff $A^\#A=1$, $(\mathfrak p_m)^\#=\mathfrak p_m$
by the shape of $J$, and $(\Xi^{-1}A\Xi)^\#=\Xi^{-1}A^\#\Xi$ as $\Xi$ is an isometry.  Put $\mathrm{pr}(A)=\frac12(A-A^\#)$; it is
skew ($\mathrm{pr}(A)^\#=-\mathrm{pr}(A)$), it preserves each $\mathfrak p_m$, commutes with conjugation by $\Xi$, and fixes skew matrices.
For $Z\in\mathfrak p_1$ put $\mathrm{cay}(Z)=1+Z\sum_{j<2r+2}(Z/2)^j=(1+Z/2)(1-Z/2)^{-1}$ (as $Z^{2r+1}=0$).  Then: (C1) if $Z\in\mathfrak p_m$
($m\ge1$), $\mathrm{cay}(Z)-1\in\mathfrak p_m$ and $\mathrm{cay}(Z)-1-Z\in\mathfrak p_{2m}$; (C2) $\mathrm{cay}(-Z)=\mathrm{cay}(Z)^{-1}$;
(C3) if $Z^\#=-Z$ then $\mathrm{cay}(Z)^\#=(1+Z/2)^{-1}(1-Z/2)=\mathrm{cay}(Z)^{-1}$, so $\mathrm{cay}(Z)$ is the matrix of an isometry;
(C4) $\Xi^{-1}\mathrm{cay}(Z)\Xi=\mathrm{cay}(\Xi^{-1}Z\Xi)$.

\begin{lemma}\label{i3:lem-Npat}
Let $v$ be an infinite place and $s\in U(F_v)$.  Then $s$ satisfies the conditions defining $N_u$ at $v$ iff $M_v(s)-1\in\mathfrak p_1$.
The same holds over $E$: an $E$-linear isometry $n$ whose matrix $A$ has $A-1\in\mathfrak p_1$ satisfies $(n-1)X_a\subset X_{a-1}$
and $(n-1)X_u^\perp\subset X_u$.  Consequently $s\in R_t$ iff $M_v(s)-1\in\mathfrak p_1$ and $\Xi_v^{-1}M_v(s)\Xi_v-1\in\mathfrak p_1$ for every $v$.
\end{lemma}

\begin{proof}
Write $y=\sum\alpha_ae_a+w+\sum\beta_af_a$ ($w\in W_v$); then $h_v(e_a,y)=\beta_a$, so $y\perp e_{b,v}$ for all $b$ iff $y$ lies in the
span of the basis vectors of level $\le r+1$, and $E_v=\mathrm{span}(e_{a,v})$ is the span of those of level $\le r$; also
$E_{<a,v}$ is the span of those of level $\le a$, and $z_u\in W$ has level $r+1$.  If $M_v(s)-1\in\mathfrak p_1$, the conditions
$(s_v-1)e_{a,v}\in E_{<a,v}$, $(s_v-1)z_{u,v}\in E_v$ and $(s_v-1)y\in E_v$ for $y\perp e_{b,v}$ follow at once.  Conversely, the
conditions give the required vanishing in the columns of the $e_a$ and the $w'_i$; for the column of $f_a$ we must show that the
$f_b$-coefficient $h_v(e_b,(s_v-1)f_a)$ of $(s_v-1)f_a$ vanishes for $b\le a$.  As $N_u$ is a group ((St)(i)),
$(s_v^{-1}-1)e_b\in E_{<b,v}$, and $h_v(e_b,(s_v-1)f_a)=h_v(s_v^{-1}e_b,f_a)-\delta_{ab}=h_v((s_v^{-1}-1)e_b,f_a)=0$ since $h(e_c,f_a)=0$ for
$c<b\le a$.  The rational statement is the first half of the same argument over $E$ (with $X^\perp_u=\mathrm{span}(e_a)\oplus W$).  The
last assertion follows, since $M_v(\gamma^{-1}s\gamma)=\Xi_v^{-1}M_v(s)\Xi_v$ and the conditions for $N_{u'}$ involve the same $e_a$.
\end{proof}

\begin{proof}[Proof of (a): discreteness]
Let $\Lambda_0=\bigoplus_i\mathcal O_Eb_i$.  As $\Lambda$ and $\Lambda_0$ are full lattices there is an integer $m\ge1$ with
$m\Lambda_0\subset\Lambda\subset m^{-1}\Lambda_0$.  If $n\in N_\Lambda$ has matrix $A$, then $n(mb_i)\in\Lambda\subset m^{-1}\Lambda_0$, so
$m^2(A-1)$ has entries in $\mathcal O_E$.  Let $U$ be the set of $s\in U_\infty$ with $|m^2(M_v(s)-1)_{ij}|<1$ for all $v,i,j$; it is an open
neighbourhood of $1$.  If $n_\infty\in U\cap\Gamma_t$, every entry $\alpha$ of $m^2(A-1)$ is an algebraic integer with $|v(\alpha)|<1$ for
every infinite place $v$ of $F$; the embeddings of the CM field $E$ are the $v$ and their conjugates, so $|N_{E/\Q}(\alpha)|<1$, hence
$N_{E/\Q}(\alpha)=0$ and $\alpha=0$.  Thus $A=1$ and $n_\infty=1$: $U\cap\Gamma_t=\{1\}$, and the subgroup $\Gamma_t$ is discrete.
\end{proof}

\begin{lemma}[Rational approximation]\label{i3:lem-ratapprox}
For every $m\ge1$ and every integer $d\ge1$ there is $C\ge0$ such that: for every family $(P_v)_v$ of complex matrices with
$P_v\in\mathfrak p_m$, $\Xi_v^{-1}P_v\Xi_v\in\mathfrak p_1$ and $P_v^\#=-P_v$, there is a matrix $Y$ with entries in $d\mathcal O_E$ with
$Y\in\mathfrak p_m$, $\Xi^{-1}Y\Xi\in\mathfrak p_1$, $Y^\#=-Y$, and $|(v(Y)-P_v)_{ij}|\le C$ for all $v,i,j$.
\end{lemma}

\begin{proof}
The first two conditions are a system $L(Y)=0$ of linear equations with coefficients in $E$.  Let $\Pi$ be an $E$-linear projection
of the space of matrices onto $\ker L$; its image under $v$ is a projection onto $\ker L_v$, since $\ker(L)\otimes_{E,v}\C=\ker L_v$
(base change of a linear system along the field embedding $v$).  Choose an integer $c\ge1$ such that $c\Pi$ has entries in $\mathcal O_E$
and $c J$, $cJ^{-1}$ have entries in $\mathcal O_E$, and put $D=2c^3d$.  Since $\mathcal O_E$ is a full lattice in
$E\otimes_\Q\R\cong\prod_v\C$, there is $R_0$ such that every $(z_v)_v\in\prod_v\C$ has some $y\in\mathcal O_E$ with $|v(y)-z_v|\le R_0$
for all $v$.  Given $(P_v)$, choose $y$ with entries in $\mathcal O_E$ and $|v(y_{ij})-(P_v)_{ij}/D|\le R_0$, and put $Y_0=D\,\Pi(y)$ and
$Y=\mathrm{pr}(Y_0)$.  Then $Y_0$ has entries in $2c^2d\,\mathcal O_E$ and solves $L$; $Y$ has entries in $d\mathcal O_E$ (as $\mathcal
O_E$ is stable under conjugation), is skew, and solves $L$ (pr preserves $\mathfrak p_m$ and commutes with conjugation by $\Xi$).
Since $P_v$ solves $L_v$ and is skew, $v(Y)-P_v=\mathrm{pr}_v\bigl(v(\Pi)(D\,v(y)-P_v)\bigr)$, whose entries are bounded by a constant
depending only on $\Pi$, $J$, $D$ and $R_0$.
\end{proof}

\begin{lemma}[Membership]\label{i3:lem-ratmem}
There is an integer $d\ge1$ such that for every matrix $Y$ with entries in $d\mathcal O_E$, $Y\in\mathfrak p_1$, $\Xi^{-1}Y\Xi\in\mathfrak p_1$
and $Y^\#=-Y$, the $E$-linear map $n$ with matrix $\mathrm{cay}(Y)$ lies in $N_\Lambda$, $\gamma^{-1}n\gamma\in N_{\Lambda'}$, and
$M_v(n_\infty)=\mathrm{cay}(v(Y))$.  In particular $n_\infty\in\Gamma_t$.
\end{lemma}

\begin{proof}
By (C3) $n$ is an isometry, and by (C1) and Lemma~\ref{i3:lem-Npat} it satisfies the flag conditions.  Choose $m\ge1$ with
$m\Lambda_0\subset\Lambda,\Lambda'\subset m^{-1}\Lambda_0$; an $E$-linear $\phi$ whose matrix has $\phi-1$ with entries in $m^2\mathcal O_E$ satisfies
$(\phi-1)\Lambda\subset(\phi-1)m^{-1}\Lambda_0\subset m\Lambda_0\subset\Lambda$, so $\phi\Lambda\subset\Lambda$ (likewise for $\Lambda'$).  Choose
$c\ge1$ with $c\Xi$, $c\Xi^{-1}$ integral, and $d=2^{2r+2}m^2c^2$.  If $Y=dA$ with $A$ integral, then
$\mathrm{cay}(\pm Y)-1=\pm Y\sum_j(\pm Y/2)^j$ has entries in $m^2\mathcal O_E$; by (C2) $n\Lambda\subset\Lambda$ and $n^{-1}\Lambda\subset\Lambda$,
so $n\Lambda=\Lambda$ and $n\in N_\Lambda$.  The matrix of $\gamma^{-1}n\gamma$ is $\Xi^{-1}\mathrm{cay}(Y)\Xi=\mathrm{cay}(\Xi^{-1}Y\Xi)$ by (C4),
and $\Xi^{-1}Y\Xi$ has entries in $2^{2r+2}m^2\mathcal O_E$, lies in $\mathfrak p_1$ and is skew; the same argument gives
$\gamma^{-1}n\gamma\in N_{\Lambda'}$.  The last formula is the image of $\mathrm{cay}(Y)$ under $v$.
\end{proof}

\begin{lemma}[One step]\label{i3:lem-onestep}
Let $m\ge1$, $A\in\mathfrak p_m$ with $(1+A)^\#(1+A)=1$, $Y\in\mathfrak p_m$ skew, and $P=\mathrm{pr}(A)$.  Then
$\mathrm{cay}(Y)^{-1}(1+A)\,\mathrm{cay}(P-Y)^{-1}-1\in\mathfrak p_{2m}$.
\end{lemma}

\begin{proof}
Modulo $\mathfrak p_{2m}$ (a two-sided ideal in the algebra of matrices of $\mathfrak p_0$), (C1) and (C2) give $\mathrm{cay}(Z)^{-1}=\mathrm{cay}(-Z)\equiv1-Z$
for $Z\in\mathfrak p_m$, so the product is $\equiv(1-Y)(1+A)(1-P+Y)\equiv1+A-P=1+\frac12(A+A^\#)$.  Unitarity gives
$A+A^\#=-A^\#A\in\mathfrak p_{2m}$.
\end{proof}

\begin{proof}[Proof of (a): cocompactness]
For $0\le j\le 2r+1$ let $R^j$ be the set of $x\in R_t$ with $M_v(x)-1\in\mathfrak p_{j+1}$ for every $v$; so $R^0=R_t$
(Lemma~\ref{i3:lem-Npat}) and $R^{2r}=\{1\}$ (as $\mathfrak p_{2r+1}=0$).  Fix $d$ as in Lemma~\ref{i3:lem-ratmem} and, for each $j$, $C_j$ as
in Lemma~\ref{i3:lem-ratapprox} for $(m,d)=(j+1,d)$.  Let $S_j$ be the set of $s\in U_\infty$ such that, for each $v$,
$M_v(s)=\mathrm{cay}(Z_v)$ for some $Z_v\in\mathfrak p_1$ with all entries of modulus $\le C_j$; it is compact (the image of a compact set
of matrices under a continuous map, intersected with the closed subgroup $U_\infty$ of $\prod_vGL_n(\C)$).

\emph{Step: for $j<2r$, $R^j\subset\Gamma_tR^{j+1}S_j$.}  Let $x\in R^j$, $A_v=M_v(x)-1\in\mathfrak p_{j+1}$ and $P_v=\mathrm{pr}(A_v)$.  Then
$P_v\in\mathfrak p_{j+1}$, $P_v$ is skew, and $\Xi_v^{-1}P_v\Xi_v=\mathrm{pr}(\Xi_v^{-1}M_v(x)\Xi_v-1)\in\mathfrak p_1$ by Lemma~\ref{i3:lem-Npat}.  Lemma~\ref{i3:lem-ratapprox}
gives $Y$; Lemma~\ref{i3:lem-ratmem} gives $g\in\Gamma_t$ with $M_v(g)=\mathrm{cay}(v(Y))$.  Put $Z_v=P_v-v(Y)$: it lies in $\mathfrak p_{j+1}$, is
skew, has $\Xi_v^{-1}Z_v\Xi_v\in\mathfrak p_1$ and entries of modulus $\le C_j$.  By (C3) there is $c\in U_\infty$ with $M_v(c)=\mathrm{cay}(Z_v)$,
and $c\in S_j$; by (C1), (C4) and Lemma~\ref{i3:lem-Npat}, $c\in R_t$.  Let $y=g^{-1}xc^{-1}\in R_t$.  By Lemma~\ref{i3:lem-onestep} with
$m=j+1$, $M_v(y)-1\in\mathfrak p_{2j+2}\subset\mathfrak p_{j+2}$, so $y\in R^{j+1}$ and $x=gyc$.

\emph{Conclusion.}  By descending induction on $j$, $R^j\subset\Gamma_tK_j$ with $K_{2r}=\{1\}$ and $K_j=K_{j+1}S_j$ compact (for
$x=gyc$ as above, $y=g'k$ with $g'\in\Gamma_t$, $k\in K_{j+1}$, so $x=(gg')(kc)$).  Hence $R_t=R^0\subset\Gamma_tK$ with $K=K_0$
compact, i.e.\ every $r\in R_t$ has $s=g^{-1}\in\Gamma_t$ with $sr\in K$.
\end{proof}

\begin{proof}[Proof of the first half of (FD)]
$U_\infty$ is a second countable locally compact group and $\Gamma_t$ is a discrete subgroup by (a).  Choose an open $U\ni1$ with
$\Gamma_t\cap UU^{-1}=\{1\}$, and a sequence $(x_i)_{i\ge0}$ with $U_\infty=\bigcup_iUx_i$ (second countability).  Put
$\Omega_t=\bigcup_i\bigl(Ux_i\setminus\Gamma_t\bigcup_{j<i}Ux_j\bigr)$, a Borel set as $\Gamma_t$ is countable.  For $x\in U_\infty$ let $i$ be
least with $x\in\Gamma_tUx_i$, and $x=sy$ with $s\in\Gamma_t$, $y\in Ux_i$; then $y\notin\Gamma_tUx_j$ for $j<i$, so $s^{-1}x=y\in\Omega_t$.  If
$y,sy\in\Omega_t$ with $s\in\Gamma_t$, the indices $i$ of $y$ and of $sy$ agree (each lies in $\Gamma_t$ times the other's piece), so
$y,sy\in Ux_i$ and $s\in UU^{-1}$, whence $s=1$.  Thus $\Omega_t$ is a strict fundamental domain of $\Gamma_t$.
\end{proof}

\begin{proof}[Proof of (c), from (NF) of Proposition~\ref{i3:prop-kppoff-in}]
Let $t$ be relevant with $G(t)=0$ and $\eta\in S_t$ (the element called $h$ in (c)); the hypothesis on the characters is not used.
Put $I=X_u\cap\gamma X_{u'}$ and $\mathrm{Fix}=\mathrm{Fix}_t$.  For $E$-subspaces $A,B$ of $V$ and an infinite place $v$, $(A+B)_v=A_v+B_v$,
$(A\cap B)_v=A_v\cap B_v$ and $(A^\perp)_v=(A_v)^{\perp}$ for $h_v$ (base change along the field embedding $v$).  So (NF) gives at every $v$:
(N1) every $y\in\C^n$ orthogonal to $(X_u^+)_v$ is congruent modulo $(X_u)_v$ to a vector orthogonal to $\mathrm{Fix}_v$; (N3) a vector
of $(X_u)_v$ orthogonal to $\mathrm{Fix}_v$ lies in $I_v$.  Note $I_v\subset(X_u)_v$ is isotropic, $I_v\subset\mathrm{Fix}_v$, and
$e_a,z_u,\gamma e_a,\gamma z_{u'}\in\mathrm{Fix}$.  Fix $v$ and drop it from the notation; $w_b$ ($b<k$) is the frame of the package at $v$.

\emph{Adapted frame.}  $w_b\perp e_a,f_a,z_u$, so by (N1) there is $x_b\in(X_u)_v$ with $\tilde w_b:=w_b-x_b\perp\mathrm{Fix}_v$.  Then
$h_v(\tilde w_a,\tilde w_b)=h_v(w_a,\tilde w_b)=h_v(w_a,w_b)=\delta_{ab}$ (as $x_a\in\mathrm{Fix}_v\perp\tilde w_b$ and $w_a\perp X_u$).

\emph{Decomposition.}  If $y\perp\mathrm{Fix}_v$, then $y-\sum_bh_v(\tilde w_b,y)\tilde w_b\in I_v$.  Indeed $y\perp e_a$ for all $a$, so in the
basis $e_a,f_a,z_u,w_b$ (see the proof of (Levi)) $y$ has no $f_a$-component: $y=s+cz_u+\sum_bd_b\tilde w_b$ with $s\in(X_u)_v$ (absorbing the
$x_b$).  Pairing with $z_u\in\mathrm{Fix}_v$ gives $c\,h(z_u,z_u)=0$, so $c=0$; pairing with $\tilde w_b$ gives $d_b=h_v(\tilde w_b,y)$; and
$s=y-\sum_bd_b\tilde w_b\in(X_u)_v$ is orthogonal to $\mathrm{Fix}_v$, so $s\in I_v$ by (N3).

\emph{Splitting $\eta$.}  $\eta_v$ is unitary and fixes $\mathrm{Fix}_v$ pointwise, hence preserves $\mathrm{Fix}_v^\perp\ni\tilde w_b$.  Let
$Q_{ab}=h_v(\tilde w_a,\eta_v\tilde w_b)$; by the decomposition, $\eta_v\tilde w_b=\sum_aQ_{ab}\tilde w_a+\varrho_b$ with $\varrho_b\in I_v$.  Expanding
$h_v(\eta_v\tilde w_a,\eta_v\tilde w_b)=\delta_{ab}$ and using $I_v\perp\tilde w_c$ and $h_v(I_v,I_v)=0$ gives $Q^*Q=1$.  Let $\theta_v$ act by $Q$ on
$\mathrm{span}(\tilde w_b)$ (in the orthonormal basis $\tilde w_b$) and as the identity on its orthogonal complement (a complement, the
span being positive definite); $\theta_v$ is unitary.  Doing this at every $v$ gives $\theta\in U_\infty$.  It fixes $\mathrm{Fix}_v\perp\tilde w_b$
pointwise, so $\theta\in S_t$, and by (iv) of the proof of (St) $\theta R_t\theta^{-1}=R_t$.  For $m\in\Z$, $\theta^m$ acts by $Q^m$ on the
same span, so all powers of $\theta$ lie in the image $C$ of the compact group $\prod_vU(k)$ under the continuous map
$(Q'_v)_v\mapsto$ (the element acting by $Q'_v$ on $\mathrm{span}(\tilde w_{b})$ and trivially on its complement); $C$ is compact.

\emph{The unipotent part.}  Let $n=\theta^{-1}\eta$.  Then $n_v$ fixes $\mathrm{Fix}_v$ pointwise, and $n_v\tilde w_b=\theta_v^{-1}(\sum_aQ_{ab}\tilde w_a+\varrho_b)
=\tilde w_b+\varrho_b$ (as $\varrho_b\in\mathrm{Fix}_v\perp\tilde w_c$).  Let $y\perp I_v$.  Then $(n_v-1)y\perp\mathrm{Fix}_v$ (for $p\in\mathrm{Fix}_v$,
$h_v(p,n_vy)=h_v(n_vp,n_vy)=h_v(p,y)$), and $h_v(\tilde w_b,(n_v-1)y)=0$, because $\tilde w_b=n_v(\tilde w_b-\varrho_b)$ gives
$h_v(\tilde w_b,n_vy)=h_v(\tilde w_b-\varrho_b,y)=h_v(\tilde w_b,y)$.  By the decomposition, $(n_v-1)y\in I_v$.  Hence $n\in N_u$: $n_v$ fixes
$e_{a,v}$ and $z_{u,v}$, and $y\perp e_{a,v}$ for all $a$ implies $y\perp I_v$, so $(n_v-1)y\in I_v\subset E_v$.  And $\gamma^{-1}n\gamma\in N_{u'}$:
it fixes $e_{a,v}$ and $z_{u',v}$ (as $\gamma e_a,\gamma z_{u'}\in\mathrm{Fix}$), and if $y\perp e_{a,v}$ for all $a$ then $\gamma_vy\perp(\gamma X_{u'})_v\supset I_v$,
so $(n_v-1)\gamma_vy\in I_v\subset(\gamma X_{u'})_v$ and $(\gamma_v^{-1}n_v\gamma_v-1)y\in E_v$.  Thus $n\in R_t$ and $\eta=\theta n$.

\emph{The measure.}  $R_t$ is a closed subgroup of $U_\infty$ ((St)(i)), hence locally compact and second countable; let $\mu_R$ be a left
Haar measure of $R_t$ and $\nu$ its image in $U_\infty$.  Then $\nu$ is $\sigma$-finite, nonzero, carried by $R_t$, finite on compact sets
and invariant under $x\mapsto rx$ ($r\in R_t$).  By (a) and (St)(i), $R_t$ contains a countable discrete subgroup $\Gamma_t$ with
$R_t\subset\Gamma_t^{-1}K$, $K$ compact; the argument of the proof of (A) (``$\nu$ is right invariant'') uses only this and shows that
$R_t$ is unimodular; so $\nu$ is right invariant and invariant under $x\mapsto x^{-1}$.  Conjugation by $n\in R_t$ preserves $\nu$ (left
and right invariance).  For $\theta$: $\alpha(x)=\theta^{-1}x\theta$ is a topological automorphism of $R_t$, so $\alpha_*\mu_R=c\,\mu_R$ for some
$c>0$ (uniqueness of Haar measure); let $V_0\subset R_t$ be a compact neighbourhood of $1$, so $0<\mu_R(V_0)<\infty$.  For $m\in\Z$,
$\alpha^m(V_0)=\theta^{-m}V_0\theta^m\subset R_t\cap CV_0C$, a compact set $L$, and $\mu_R(\alpha^m(V_0))=c^{-m}\mu_R(V_0)\le\mu_R(L)<\infty$ for all
$m\in\Z$; hence $c=1$.  Since $x\mapsto\eta^{-1}x\eta$ is the composite of conjugation by $\theta$ and by $n$, $\nu$ is invariant under it.
So $\nu$ has all the properties required in (S) for $(R_t,\eta)$.
\end{proof}

\emph{The finite-adelic characters.}  We use the notation of \S\ref{i3:ssec-pf-unip}: $V_f=V\otimes_F\mathbb A_f$ with the adelic form
$\langle\,,\rangle$, $\widehat L$, $K$, $\mathrm{Mo}$, and the frame modulo $g=g_0g_1$ of the Hida setting, which we now write
$e^\circ_a,f^\circ_a,z^\circ,\beta_i,\beta_i^\vee$ to distinguish it from the rational vectors $e_a,f_a,z_u$ (viewed in $V_f$); $K_2$ is its level
subgroup.  For $x\in GL(V)$ rational, $x_f$ denotes its image in $G_f$.  The character $\psi_E$ of the setting satisfies, besides
$\psi_E(\widehat{\mathcal O})=1$, the product formula $\psi_E(y)\psi_\infty(y)=1$ for $y\in E$, where $\psi_\infty(y)=\prod_v\exp(2\pi
i(v(y)+\overline{v(y)}))$.  For $x\in G_f$ put, with $\nu_a$ as in the formula for $\psi_u$ and $n^\circ_a=e^\circ_{a+1}$ ($a<r-1$),
$n^\circ_{r-1}=z^\circ$,
\[
\psi_{f,u}(x)=\psi_E\Bigl(g^{-1}\sum_{a<r}\langle f_a,x\nu_a\rangle\Bigr),\qquad \chi_N(x)=\psi_E\Bigl(g^{-1}\sum_{a<r}\langle f^\circ_a,xn^\circ_a\rangle\Bigr),
\]
and let $K^N$ be the set of $x\in K$ with $(x-1)e^\circ_a\in\mathrm{Mo}(\{e^\circ_b:b<a\})$ and $(x-1)z^\circ\in\mathrm{Mo}(\{e^\circ_b\})$.  So
$K_2\subset K^N$, and $\chi_N$ is the first factor of $\chi_2$.  The proofs of Lemma~\ref{i3:lem-K2group} and of the multiplicativity in
Proposition~\ref{i3:prop-chi2}, restricted to the vectors $e^\circ_a,z^\circ$ and to the first sum, show: $K^N$ is a subgroup of $K$ and $\chi_N$
is multiplicative on $K^N$, with $\chi_N(1)=1$.  For each row lattice $\Lambda$ of a word $u$ fix a column class $\rho(\Lambda)$ of $u$ whose
lattice is $\Lambda$ (the package provides one), and put $a_\Lambda=h_{\rho(\Lambda)}x_u\in G_f$, where $x_u\in K$ is the element of the
Hida setting which fixes all $e_a,f_a$ and satisfies $x_u^{-1}z_u-z^\circ\in g\widehat L$.

\begin{proposition}[Level transport; formally verified, \cite{CharacterFamily}]\label{i3:prop-levtr}
For every datum, full package and words $u,u'$:
\begin{enumerate}
\item[(LT)] for every row lattice $\Lambda$ of $u$ and every $n\in N_\Lambda$, the element $y=a_\Lambda^{-1}n_fa_\Lambda$ lies in $K_2$ and
$\chi_N(y)=\psi_{f,u}(n_f)$;
\item[(Car)] for every term $t$ of $X(u,u')$ with row lattices $(\Lambda,\Lambda')$ and representative $\gamma$,
$\kappa_t=a_\Lambda^{-1}\gamma_fa_{\Lambda'}$ lies in $K^N$.
\end{enumerate}
\end{proposition}

\noindent (LT) says that a rational element of the unipotent radical preserving a row lattice is, after transport by the class
representative, an element of the level subgroup, on which the level character recovers the finite part of the Bessel character;
its formal proof is a long computation with the integral frame of the setting.  (Car) is the compatibility of the term's
representative with the flags modulo $g$ (the admissibility of the term).  Its formally verified steps are outlined in \S\ref{i3:str-levtr}.

\begin{proof}[Proof of (b), from Proposition~\ref{i3:prop-levtr}]
Let $s\in\Gamma_t$, $s=n_\infty$ with $n\in N_\Lambda$ and $n'=\gamma^{-1}n\gamma\in N_{\Lambda'}$; then $\gamma^{-1}s\gamma=n'_\infty$.
\emph{Product formula.}  For a rational isometry $m$ and a word $w$ put $c_w(m)=\sum_ah(f_a,m\nu_a)\in E$ ($\nu_a$ for $w$).  At each $v$,
$c_v(m_\infty)=v(c_w(m))$ in the formula for $\psi_w$, so $\psi_w(m_\infty)=\psi_\infty(c_w(m)/g)$; and $\psi_{f,w}(m_f)=\psi_E(c_w(m)/g)$, the
adelic form extending $h$.  Hence $\psi_w(m_\infty)\psi_{f,w}(m_f)=1$.  \emph{Level.}  By (LT), $y=a_\Lambda^{-1}n_fa_\Lambda$ and
$y'=a_{\Lambda'}^{-1}n'_fa_{\Lambda'}$ lie in $K_2\subset K^N$, with $\chi_N(y)=\psi_{f,u}(n_f)$ and $\chi_N(y')=\psi_{f,u'}(n'_f)$.  Since
$n'_f=\gamma_f^{-1}n_f\gamma_f$, $y=\kappa_ty'\kappa_t^{-1}$, and $\kappa_t\in K^N$ by (Car).  As $\chi_N$ is multiplicative on the group $K^N$,
$\chi_N(y)=\chi_N(\kappa_t)\chi_N(y')\chi_N(\kappa_t)^{-1}=\chi_N(y')$.  Therefore $\psi_{f,u}(n_f)=\psi_{f,u'}(n'_f)$, and by the product formula
$\psi_u(s)=\psi_{f,u}(n_f)^{-1}=\psi_{f,u'}(n'_f)^{-1}=\psi_{u'}(\gamma^{-1}s\gamma)$.
\end{proof}


\subsection{Outline of the formally verified steps}\label{i3:ssec-structure}

We outline, for each statement of this section that is formally verified, the chain of steps of its formal proof, read
from the top levels of the formal development.  The outline is not a proof.  Each step is classified as (L), (N) or (R) as
explained in the introduction to this annex; for an (L) step we state what the formal statement needs beyond the published
one, and for an (N) step we indicate the idea of its proof.  When we say that we do not describe the proof of a step, the
step enters only through its statement.  We begin with item (NF) of Proposition~\ref{i3:prop-kppoff-in}.

\medskip\noindent\label{i3:ssec-NF-structure}\emph{The normal form theorem (item (NF) of Proposition~\ref{i3:prop-kppoff-in}).}

Item (NF) of Proposition~\ref{i3:prop-kppoff-in} is formally verified, \cite{CharacterFamily}.  It is the specialisation to the pair $(F_u,\gamma_tF_{u'})$ of a \emph{normal form theorem} for pairs of
flags, stated over an arbitrary field $K$ with an involution and an embedding $\sigma:K\to\C$ commuting with conjugation.  For
flags $F=(e_0,\dots,e_{r-1};z)$, $F'=(e'_0,\dots,e'_{r-1};z')$ of norms $T,T'$ in a finite-dimensional hermitian space, with $X$, $X'$
their isotropic parts, $I=X\cap X'$ and $V_0=X+X'+Kz+Kz'$ (so $V_0=\mathrm{Fix}_t$ for the pair of the term), the hypotheses are:
$\sigma(T),\sigma(T')>0$; $F$ and $F'$ are \emph{positive at $\sigma$} (every $w$ orthogonal to all vectors of the flag and not in its
isotropic part has $\sigma(h(w,w))>0$); $(F,F')$ is Lie-relevant; and $G(F,F')=0$.  The conclusion is the existence of $j\le r$
and $c\in K^\times$ with: $X=X_j\oplus I$, $X'=X'_j\oplus I$, $z'-cz\in X+X'$, $c\bar cT=T'$, the minor $D_j\ne0$, the radical of $h$ on
$V_0$ is $I$, $\dim V_0=r+j+1$, an element of $X$ (or $X'$) is orthogonal to $V_0$ iff it lies in $I$, and every $w$ orthogonal to
the vectors of $F$ (resp.\ $F'$) is congruent modulo $X$ (resp.\ $X'$) to a vector orthogonal to $V_0$.  The last three clauses are
the three identities of (NF).  That the pair of a term satisfies the positivity hypotheses at a suitable embedding is checked in
the formal proof from the signature of the datum; we do not describe this check here.  The steps are the following.
\begin{enumerate}
\item[(R)] \emph{Anisotropy.}  Positivity at $\sigma$ implies that an isotropic vector orthogonal to the flag lies in its
isotropic part.
\item[(N)] \emph{Case $I\ne0$: alignment of levels.}  For $x\in I$ with leading level $a$ in $F$ (coefficient $\alpha$) and $a'$ in $F'$
(coefficient $\alpha'$), and $v$ orthogonal to $e_b$ ($b\le a$), the element $Y=x\wedge v$, $Yy=h(v,y)x-h(x,y)v$, is nilpotent, skew,
lies in $\mathfrak n_F$, and its root coordinate is $\ell_F(Y)=\alpha h(e_{a+1},v)$; the same for $F'$.  Lie-relevance, applied to such
$Y$, forces $\alpha e_{a+1}-\alpha'e'_{a'+1}\in X_{a+1}+X'_{a'+1}$.  Iterating along a basis of $I$ adapted to the levels gives $j$ and
$c$ with the shape $X=X_j\oplus I$, $X'=X'_j\oplus I$, matching leading coefficients in ratio $c$, and $z'-cz\in X+X'$; we describe this iteration only in outline.
\item[(N)] \emph{Case $I=0$, $D_r\ne0$.}  Projecting $z$, $z'$ orthogonally to $X+X'$ (possible since the pairing between $X$ and $X'$
given by $D_r$ is nondegenerate) gives $z_1=z-x$, $z'_1=z'-x'$ with $h(z_1,z_1)=-T$, $h(z'_1,z'_1)=-T'$ and
$G=|D_r|^2(|h(z_1,z'_1)|^2-TT')$.  By positivity, $(X+X')^\perp$ has exactly one negative direction at $\sigma$.  So $G=0$ is the case
of equality in the reverse Cauchy--Schwarz inequality, and $z'_1=cz_1$; the normal form holds with $j=r$.
\item[(L)] \emph{Reverse Cauchy--Schwarz.}  For a hermitian form of signature $(m,1)$ on a complex vector space and vectors
$a,b$ with $h(a,a)<0$, $h(b,b)<0$: $|h(a,b)|^2\ge h(a,a)h(b,b)$, with equality iff $a,b$ are proportional (classical; reference to be confirmed).
The formal version is over $K$ with $\sigma$; the proportionality constant is in $K$ because the equality case is a linear
dependence over $K$.
\item[(N)] \emph{Case $I=0$, $D_r=0$ is impossible} under anisotropy and $G=0$.
\item[(R)] \emph{Consequences of the normal form.}  $D_j\ne0$; $c\bar cT=T'$ (compare norms of $z'$ and $cz$ modulo the isotropic
parts); the radical of $V_0$ is $I$ and $\dim V_0=r+j+1$ (Gram matrix computation with $D_j\ne0$); an element of $X$ orthogonal to
$V_0$ lies in $I$; the complement statement; and the symmetric statements, by exchanging $(F,F',c)$ with $(F',F,c^{-1})$.
\end{enumerate}
Only the reverse Cauchy--Schwarz inequality is in the literature; the alignment of levels by relevance (the second item) is the
new point, and the remaining steps are linear algebra.



\medskip\noindent\label{i3:str-model}\emph{The discrete-series model (Proposition~\ref{i3:prop-model}).}

The formal proof lives in a library on the discrete series of $U(p,q)$ organised in seven parts: highest weights of $K$,
the generating function $\Psi$ (``Casimir coefficients''), existence of the model, a Bochner-type inequality, rigidity of
lowest $K$-types, tensor products with the standard representation, and formal degrees.  The decay lemma
(Lemma~\ref{i3:lem-decay}) is obtained from (a)--(e) exactly as in \S\ref{i3:ssec-pf-bessel}.  The steps:
\begin{enumerate}
\item[(M1)] (L) \emph{The lowest $K$-type.}  Theorem of the highest weight (Theorem~\ref{i3:thm-knapp}) for
$K\cong U(p)\times U(q)$, after Lemma~\ref{i3:lem-ktype}.  Nothing extra is needed.
\item[(M2)] (N) \emph{The generating function, item (a).}  $\Psi$ is obtained by averaging $\Phi_\lambda(k^{-1}gk)$,
$\Phi_\lambda=\prod_i\Delta_i^{-(\mu_i-\mu_{i+1})}$, against the projection on $w_0$ over $K$, as in the statement.  The formal proof checks: the
$K$-biequivariance and $\Psi(1)=1$ (R); smoothness (R); the Casimir equation $\sum E_{ab}E_{ba}\Psi=c(\lambda)\Psi$, a
differentiation of the product of powers of leading minors (a Capelli-type computation: the minors $\Delta_i$ are highest
weight vectors for the two-sided action of $GL_n(\C)$, and the Casimir acts on such a product by the value of the
Casimir at the weight $\mu_\lambda-\rho$); the Gelfand equation (the convolution identity characterising $\tau$-spherical
functions); a decay estimate along $\exp\mathfrak p$, obtained from lower bounds for the leading minors of $\exp X$,
$X\in\mathfrak p$, by moment estimates in polar coordinates, under $h>8nq_0$; square integrability from the decay and
the integration formula in Cartan coordinates (Theorem~\ref{i3:thm-helgason}, (L)); and uniqueness from the restriction
to $\exp\mathfrak p$ and the Cartan decomposition $U_\epsilon=K\exp\mathfrak p$ (L: e.g.\ Knapp, \emph{op.\ cit.}, Thm.~6.31,
reference to be confirmed).  Background (L): a discrete series with lowest $K$-type $\tau$ has the $\tau$-spherical function
$g\mapsto P\sigma(g)P$ ($P$ the projection on the lowest $K$-type), and $\tau$-spherical functions satisfying the Casimir
and Gelfand equations determine the representation they generate (R.~Godement, \emph{A theory of spherical functions I},
Trans.\ AMS 73 (1952) 496--556; reference to be confirmed).  What is new is the explicit formula through $\Phi_\lambda$ and its
decay with the explicit threshold; outside the holomorphic case (where $\Phi_\lambda$ reduces to a power of one
determinant) we do not know a published source for it (reference to be confirmed).
\item[(M3)] (L) with an (N) proof, \emph{item (b): the model is the discrete series.}  Formally: the closure $H_\lambda$ is
invariant and nonzero, $\sigma_\lambda$ is irreducible, square integrable, has a $K$-type of highest weight $\Lambda_\lambda$
and satisfies the Vogan minimality; any discrete series with parameter $(\epsilon,\lambda)$ is equivalent to it.  The
formal route: on any discrete series with parameter $(\epsilon,\lambda)$ the Casimir acts on the lowest $K$-type by
$c(\lambda)$ (Vogan minimality excludes highest weight vectors of weight $\Lambda_\lambda-\alpha$, $\alpha$ noncompact, in the
first layer); the layers $V_{k+1}=V_k+\sum E_{ab}V_k$ have a Casimir gap (the Bochner-type inequality; the only place where
$h>npq+1$ enters); and two data with the same lowest $K$-type and Casimir value are rigidly equivalent.  The published
statement this reproves is: \emph{for a regular Harish-Chandra parameter $\lambda$ of $U(p,q)$ there is, up to equivalence,
exactly one discrete series with infinitesimal character $\lambda$ and lowest $K$-type $\Lambda_\lambda=\lambda+\rho_n-\rho_c$,
which occurs with multiplicity one} (Harish-Chandra, Theorem~\ref{i3:thm-hc} and its extension to reductive groups;
Schmid, Ann.\ of Math.\ 102 (1975) 535--564; Hecht--Schmid, Invent.\ Math.\ 31 (1975) 129--154; Vogan, Ann.\ of Math.\ 109
(1979) 1--60; reference to be confirmed).  Beyond it the formal statement needs: (i) that Definition~\ref{i3:def-ds} (Vogan
minimality of $\Lambda_\lambda$ among the $K$-types present) singles out $\pi_\lambda$; this is the cited characterisation, whose
exact form we have not matched with the definition; (ii) the identification of the $K$-finite vectors with
$\mathcal G_\lambda$ (R, from admissibility: each $K$-type occurs with finite multiplicity in an irreducible unitary
representation, Harish-Chandra, Trans.\ AMS 75 (1953) 185--243, reference to be confirmed); (iii) density of the $K$-finite vectors
in the smooth vectors for every finite family of seminorms (L: Harish-Chandra; e.g.\ G.~Warner, \emph{Harmonic analysis on
semi-simple Lie groups I}, Springer 1972, \S4.4; reference to be confirmed); (iv) the last clause of (b), which follows from (iii)
and irreducibility: the closure in $H^\infty$ of the span of $\sigma(U)w_0$ is a closed invariant subspace of the smooth
representation, it is stable under the $K$-isotypic projections, hence contains all $K$-finite vectors (finite multiplicities),
hence is $H^\infty$; a linear map of the form $\sum\beta_i\circ d\sigma(W_i)$ is continuous for the smooth topology (R).
\item[(M4)] (L) \emph{Item (c), Schur's lemma on the model.}  An invariant sesquilinear form on an irreducible admissible
$(\mathfrak g,K)$-module which is hermitian-symmetric in the sense $b(E_{ac}x,y)=s_{ac}b(x,E_{ca}y)$ is a multiple of the
invariant inner product (Schur's lemma for irreducible admissible $(\mathfrak g,K)$-modules, e.g.\ Wallach, \emph{Real reductive
groups I}, \S3.3; reference to be confirmed).  Extra: none, given (M3).
\item[(M5)] (L) \emph{Item (d), Casimir on the tensor product.}  On $\mathcal G_\lambda\otimes\C^n$ one has
$\sum_{a,b}T_{ab}T_{ba}=\mathrm{Cas}(E)\otimes1+1\otimes\mathrm{Cas}(e)+2\Omega$ (proof of Lemma~\ref{i3:lem-split}), so $\Omega$ is,
up to constants, the Casimir of the tensor product; by Kostant's theorem
(B.~Kostant, \emph{On the tensor product of a finite and an infinite dimensional representation}, J.\ Funct.\ Anal.\ 20 (1975)
257--285; reference to be confirmed) a module with infinitesimal character $\lambda$ tensored with $\C^n$ is a direct sum of its
generalised eigenspaces for the infinitesimal characters $\lambda+e_s$, on which $\Omega$ acts by $a_s$ (here diagonalisably,
since the $a_s$ are distinct).  The formal proof establishes the eigen-decomposition directly by resolvent identities and a
matching of poles.  The clause ``$a_s$ is not an eigenvalue if $\lambda+e_s$ is singular'': under the gap a singular
$\lambda+e_s$ lies on the wall of a \emph{compact} simple root ($\lambda_{s-1}=\lambda_s+1$ with $\epsilon_{s-1}=\epsilon_s$), and the
translation of a discrete series to such a wall vanishes (Zuckerman, Ann.\ of Math.\ 106 (1977) 295--308; Knapp--Zuckerman,
Ann.\ of Math.\ 116 (1982) 389--455; reference to be confirmed).
\item[(M6)] (L) with (N) parts, \emph{item (e), shifted parameters.}  That the $a_s$-eigenspace is a model of the discrete
series with parameter $\lambda+e_s$ is Zuckerman's translation principle within a Weyl chamber (same reference).  The formal
statement needs in addition the explicit intertwiner $\iota$ with the model $\mathcal G_{\lambda+e_s}$ and the transport of
the pairing identity to all $g$ (R, by continuity and the equivariance of $\iota$), and the positivity
$\kappa\ne0$ of the induced form, which is proved in Lemma~\ref{i3:lem-split}.  The formal proof also derives
$d(\lambda+e_s)=r_s(\lambda)d(\lambda)$ from a trace inequality for the eigenprojections; this ratio is the quotient of two
instances of Harish-Chandra's formal degree formula (Theorem~\ref{i3:thm-hc}), (L).
\end{enumerate}
(M1), (M3), (M4), (M5) and most of (M6) are published representation theory of $U(p,q)$; the formal
development reproves them under the gap hypothesis.  The genuinely new parts are (M2) (the explicit generating function and
its decay) and the effective polynomial decay of Lemma~\ref{i3:lem-decay} for \emph{smooth} vectors with explicit Sobolev
seminorms and the explicit range $m<h-\mathrm{thr}^*(\epsilon)$.  The exponents of the matrix coefficients of the discrete
series are known (D.~Mili\v ci\'c, Duke Math.\ J.\ 44 (1977) 59--88; reference to be confirmed), and for $K$-finite vectors they give
decay of the order of the gap; the uniformity over smooth vectors with explicit seminorms is what is extra.  Roughly two thirds
of the core is published mathematics.

\medskip\noindent\label{i3:str-GKin}\emph{The analytic and Gelfand--Kazhdan inputs (Proposition~\ref{i3:prop-GKin}).}

The formal proof is a multiplicity-one library for Bessel models of $U$ together with a library of integrals over $N$ and
$U_zN$ (iterated Heisenberg integrals and tail integrals).
\begin{enumerate}
\item[(G1)] (R) given (L), \emph{item (c).}  Reproducing formula $\langle v,\sigma(x)u\rangle=d_\sigma\int_U\langle v,\sigma(y)w\rangle
\langle w,\sigma(y^{-1}x)u\rangle\,d\mu(y)$ for a unit vector $w$, a consequence of Schur orthogonality (Theorem~\ref{i3:thm-godement});
then the convolution of $\|\cdot\|^{-b_1}$ and $\|\cdot\|^{-b_2}$ is bounded by $\|\cdot\|^{-(\min(b_1,b_2)-A_G-1)}$ using
$\|y^{-1}x\|\ge\|x\|/(n\|y\|)$ (Lemma~\ref{i3:lem-hs1}) and Lemma~\ref{i3:lem-vol}.
\item[(G2)] (N), \emph{items (d), (e).}  The Haar measure of $N$ is written as an iterated integral over Heisenberg levels, and
$\{(h,m):\|mh\|\le R\}$ has $\nu$-measure $O(R^{A_B})$ ($U_z$ compact).  This is an explicit volume computation; the
integrability thresholds $A_S$, $A_B$ are what the later numerical conditions need.
\item[(G3)] (R), \emph{items (f)--(h)}: hermitian symmetry and positivity of $\beta$ by Fubini and the reproducing formula;
equivariance by invariance of the measures.
\item[(G4)] (N) with an (L) model, \emph{the involution $\theta$.}  $\theta$ is an anti-automorphism of $U$ of transpose type
composed with conjugation by an element of the $z$-frame, chosen so that $\theta(U_z)=U_z$, $\theta(N)=N$ and
$\psi\circ\theta=\psi$ on $N$.  This is the involution of the Gelfand--Kazhdan method for Bessel models, as in the published
multiplicity-one theorems (Aizenbud--Gourevitch--Rallis--Schiffmann, Ann.\ of Math.\ 172 (2010) 1407--1434; Jiang--Sun--Zhu,
GAFA 20 (2010) 690--709; Sun--Zhu, Ann.\ of Math.\ 175 (2012) 23--44); there it is used to prove invariance of
$U_zN$-biequivariant distributions, which the formal route avoids.
\item[(G5)] (L)+(R), \emph{the dichotomy of (GK4), (GK5).}  Either $\sigma\circ\theta^{-1}$ (twisted to a representation) is not
equivalent to the conjugate of $\sigma$, and then the cross coefficients are orthogonal by Schur orthogonality
(Theorem~\ref{i3:thm-godement}); or it is, and Schur's lemma gives the conjugate-linear isometry $C$.  The same for $\pi$ on $U_z$.
\item[(G6)] (N), \emph{the regularised pairing (GK2), (GK3).}  The big cell is decomposed into $U_zN$-double cosets with an
explicit iterated integration formula (Heisenberg levels, the isotropic cone and a hyperboloid measure); $\mathcal B$ is
defined as $c\,\mathcal B_\infty$ by integration over the double-coset space; (GK2) is an unfolding, (GK3) the
$\theta$-invariance of the double-coset measure, and the absolute convergence uses the decay class $\mathcal C$ and the
numerical conditions of \S\ref{i3:ssec-gap}.  This is the longest part.
\end{enumerate}
The published multiplicity-one theorems (Jiang--Sun--Zhu, Sun--Zhu) imply Proposition~\ref{i3:prop-bessel}
for all parameters, so the whole of Proposition~\ref{i3:prop-GKin} could be bypassed by citing them, provided their notion of
Bessel functional (continuous on the Casselman--Wallach globalisation) is matched with Definition~\ref{i3:def-bessel}; the
statement there is for the pair $(U(n+1),U(n))$ and its Bessel generalisations, and the match of the Bessel data of
Definition~\ref{i3:def-bessel} with theirs remains to be checked.  Within the formal route, (G1), (G3), (G5) are published
or routine, (G2), (G4), (G6) are new; roughly one third is published.

\medskip\noindent\label{i3:str-orbbd}\emph{The inputs of Proposition~\ref{i3:prop-orbbd}.}

Proposition~\ref{i3:prop-orbbd} is proved in \S\ref{i3:ssec-pf-orbbd}.  Its structural inputs are classified as follows.
(L) Iwasawa decomposition $G=PK$ for the minimal parabolic (Lemma~\ref{i3:ob-iwa}; Knapp, \emph{op.\ cit.}, Thm.~6.46 and its
reductive form; the written proof is the matrix version in the basis where $P$ is upper triangular).  (L) Cartan decomposition
of $U(k,1)$ (Lemma~\ref{i3:ob-cartan}; Knapp, Thm.~7.39).  (L) Witt's theorem in the form of Lemmas~\ref{i3:ob-witt2} and
\ref{i3:ob-B-move} (reference to be confirmed; the written proofs are explicit reflections).  (L) Calculus of smooth vectors,
Lemma~\ref{i3:ob-B-calc} (G\aa rding; e.g.\ Warner, \emph{op.\ cit.}, \S4.4; reference to be confirmed).  (L) Canonical coordinates of
the second kind, Lemma~\ref{i3:ob-B-chart} (standard Lie theory; reference to be confirmed).  (R) Levi decomposition
(Lemma~\ref{i3:ob-levi}), adapted bases, splitting of the Lie algebra, exponentials.  (N) the torus conjugation bound
(Lemma~\ref{i3:ob-B-tor}), the root vectors with explicit $\psi$-frequency (Lemma~\ref{i3:ob-B-root}) and the quotient and absolute
bounds (Lemmas~\ref{i3:ob-B-qbound}, \ref{i3:ob-B-abs}), which together are the argument.  We give written proofs of the (L) items
because their statements are the matrix forms in the fixed basis, which the textbook statements give only after a
translation of the same length.

\medskip\noindent\label{i3:str-HCfdeg}\emph{The formal degree (Proposition~\ref{i3:thm-HC-fdeg}).}

\begin{enumerate}
\item[(F1)] (R) \emph{Normalisation.}  The compact-dual measure of the datum is a fixed multiple of the measure used in the model
library; the scale tie identifies them.
\item[(F2)] (L) via the model, \emph{existence}: Proposition~\ref{i3:prop-model}(b) with the gap $h>\max(npq+1,16nq_0+8\sqrt{2q_0})$,
implied by $h>32Npq+2$.  Published: Harish-Chandra's existence theorem.
\item[(F3)] (L) \emph{The ratio step.}  $d(\lambda+e_s)/d(\lambda)=r_s(\lambda)=\dim V(\lambda+e_s)/\dim V(\lambda)$, from
(M6) of \S\ref{i3:str-model}.  Hence $d(\lambda_t)/\dim V(\lambda_t)$ is constant along the family $\lambda_t=\lambda+t\,\delta$, $t\in\N$,
$\delta=(n-1,n-2,\dots,0)$ (a sum of unit shifts; each $\lambda_t$ is a Harish-Chandra parameter of gap $\ge h+t$).  Published: the quotient of two instances of
Harish-Chandra's formula.
\item[(F4)] (N) \emph{The constant, by a limit.}  As $t\to\infty$ the coefficient $\langle w,\Psi_{\lambda_t}(\exp X)w\rangle$ concentrates
at $X=0$; after the substitution $X\mapsto t^{-1/2}X$ the integral $\int|\langle w,\sigma(g)w\rangle|^2d\mu$ converges, by a
Gaussian comparison on $\mathfrak p$ and the density at the identity of the Haar measure in exponential coordinates, to
an explicit Gaussian integral, which is matched with the compact-dual normalisation; lower and upper bounds for the
$\liminf$ and $\limsup$ are proved separately.  The constant comes out as $1$.
\item[(F5)] (L) Weyl's dimension formula for $U(N)$ (see \S\ref{i3:S-ssec-pf-transport}) and the independence of $w$ (Theorem~\ref{i3:thm-godement}).
\end{enumerate}
The statement is a special case of Harish-Chandra's formula in the compact-dual normalisation, as explained in
Remark~\ref{i3:rem-HC-fdeg}; only the matching of definitions and normalisations is missing for a citation.  The formal proof is a
new proof of a published formula ((F4) is the only non-routine new step).  Essentially all of the core is published.

\medskip\noindent\label{i3:str-mass}\emph{The mass ratio (Proposition~\ref{i3:prop-massratio}).}

\begin{enumerate}
\item[(Ma1)] (N) as organised, (L) in substance, \emph{the inductive mass formula.}  For every $m\ge1$, $\mathrm{mass}_I(m)=c_1\prod_{j=2}^mR(E,j)$
with a constant $c_1>0$ independent of $m$; the first statement is the quotient of two instances.  The induction splits off a
rank-one sublattice spanned by a vector of norm $1$ (possible by the representation of totally positive elements by binary
hermitian forms, from Landherr's theorem, Theorem~\ref{i3:thm-landherr}) and compares the masses of $L$ and $L'=v^\perp$ by the
ratio of local densities.
\item[(Ma2)] (L) \emph{Local densities} of the unimodular lattice $O_{E,\mathfrak p}^M$ representing $1$: at a prime unramified in
$E/F$ the ratio is $(1-\eta(\mathfrak p)^MN\mathfrak p^{-M})$; at infinity the factor is $(M-1)!/\pi^M$ per place (volume of the
unit sphere).  These are standard computations (e.g.\ Y.~Hironaka, \emph{Local zeta functions on hermitian forms and its
application to local densities}, J.\ Number Theory 71 (1998) 40--64; G.~Shimura, \emph{Euler products and Eisenstein series},
CBMS 93, 1997; reference to be confirmed).
\item[(Ma3)] (L) \emph{The Siegel mass formula for hermitian lattices} (Siegel; Braun, Ann.\ of Math.\ 41 (1940); in exact form
Gan--Hanke--Yu, Duke Math.\ J.\ 107 (2001) 103--133; reference to be confirmed): the mass of a genus is an explicit product of
local densities and a volume.  The formal statement uses only the ratio of two consecutive ranks, for which the global
constant cancels; what is needed beyond the published formula is the matching of the normalisations of the local
densities with the regrouped Euler product $L(E,M)$ (done in the remark after the statement).
\item[(Ma4)] (R) given (L), \emph{the second statement under (VT)}: the two-stage induction described after the statement, from
the volume ratio (a consequence of the rank-one Siegel--Weil formula, Theorem~\ref{i3:thm-SWT}) and the rank-two genus mass
(an instance of (Ma1)).
\end{enumerate}
The proposition is published in substance; only the matching of normalisations is taken from the formal verification.

\medskip\noindent\label{i3:str-amp}\emph{The amplitude identity (Proposition~\ref{i3:prop-amp-in}).}

\begin{enumerate}
\item[(A1)] (R) given Lemma~\ref{i3:lem-thmM}, \emph{the Bessel pair}: the multiplicity-one map of Lemma~\ref{i3:lem-thmM}, applied
at $z$ and at $z'$, gives $l$, $l'$ with $\beta=l\otimes\bar l$ on flat vectors; $\mathcal D_0$ consists of smooth vectors because
$K$-finite vectors of $\sigma$ are smooth (L: Harish-Chandra, as in (M3)(ii) of \S\ref{i3:str-model}).
\item[(A2)] (N) \emph{Non-vanishing}: a computation on the span of the derivatives of the lowest vector, showing that $\beta$ is
nonzero on $\mathcal D_0\times\mathcal D_0$; we do not describe it in detail.
\item[(A3)] (L)+(N), \emph{the amplitude identity}: the orbital integral over a cell is unfolded with Schur orthogonality
(Theorem~\ref{i3:thm-godement}, (L)), the inversion invariance of the Haar measure of the relative stabiliser (R: it is a closed
subgroup of a unipotent group, hence unimodular), and Fubini; the absolute convergence of the cell integrals (N) is an explicit
exponent count which uses the second numerical condition on $h$ of \S\ref{i3:ssec-gap}.
\end{enumerate}
About half of the proof is published or routine; the cell convergence estimate and (A2) are new.

\medskip\noindent\label{i3:str-count}\emph{The counting theorem (Proposition~\ref{i3:prop-counting}).}

The formal proof derives the count from a \emph{volume step} (a recursive bound for the covolumes of hermitian lattices) and
then runs a trace-formula argument with explicit error terms.
\begin{enumerate}
\item[(C0)] (L)+(R) \emph{Volume step from Siegel--Weil.}  The mean of the one-vector theta series over $U(F)\backslash U(\mathbb A)$
equals the product of local densities (rank-one case of (SW), Theorem~\ref{i3:thm-SWT}); a monotone limit over test functions
gives the mean of the counting function of vectors of given norm, and from it the ratio of covolumes of $W$ and of the complement
of a vector (the volume step).
\item[(C1)] (N) \emph{Decay with gap of lowest-$K$-type coefficients.}  The coefficient of a lowest-$K$-type vector is compressed to
the generating function $\Psi$ of \S\ref{i3:str-model} (a strict lowest $K$-type gives $P\sigma(g)P=\Psi(g)$), whose decay under the
gap gives a pointwise bound $|\phi_e(x^{-1}\gamma x)|\le\ldots$ in terms of $\|x^{-1}\gamma x\|$.  Background (L): the lowest-$K$-type
coefficient is the $\tau$-spherical function (Godement), as in (M2).
\item[(C2)] (N) \emph{Absolute volume bound} for hermitian lattices of small covolume, by induction on the rank from the volume
step; and (C3) (N) \emph{flag probability bounds}: the measure of the set of $x\in U(F)\backslash U(\mathbb A)$ for which $x$ has an
isotropic flag of small height, again from the volume step.
\item[(C4)] (L)+(R) \emph{Finite multiplicities}: the space of intertwiners is finite dimensional, as in (FD) (Lemma~\ref{i3:S-lem-FD}).
\item[(C5)] (N) \emph{The terms of the geometric side.}  $D_{\rm count}(I)$ is the trace of the kernel
$\sum_\gamma\phi_e(x^{-1}\gamma x)$ over $\mathcal F_I$; the identity contributes $D^*(I)$; the other $\gamma$ are split into
near-unipotent elements, far elements in a moderate region, elements with a rare flag, and a deep region, and each class is
bounded by (C1)--(C3) and elementary lattice counts.
\item[(C6)] (R) Assembly with the explicit constants.
\end{enumerate}
The method is the limit-multiplicity method of the trace formula (DeGeorge--Wallach, Ann.\ of Math.\ 107
(1978) 133--150; Savin, Invent.\ Math.\ 95 (1989) 149--159; with explicit error terms Shin--Templier, Invent.\ Math.\ 203 (2016)
1--177; reference to be confirmed).  The published results count in the level or weight aspect with fixed group, or with error
terms polynomial in the level; the statement here (rank $n\to\infty$, uniform in the field, relative error
$e^{-\frac1{16}nf\log n}$) is not among them to our knowledge.  (C0) and (C4) are published; (C1)--(C3), (C5) are new as
quantitative statements, so that roughly one quarter of the proof is published.

\medskip\noindent\label{i3:str-allrows}\emph{The all-rows datum (Proposition~\ref{i3:prop-allrows-built}).}

\begin{enumerate}
\item[(R1)] (N) \emph{Row labels}: from the finite-dimensionality of the rows, the realisation $X$ and the three ties, the labelled
row basis spans exactly the good part of the row space.
\item[(R2)] (R) \emph{Reduction}: the reduction $K\to\prod GL_r(\mathbb F_{g_k})$ at the $2f$ places above $g_0g_1$ is surjective
with kernel $K(g)$ (at a place of $F$ split in $E$ the local stabiliser is $GL_r$ of the local ring, whose reduction is surjective;
we do not describe the formal proof of this step in detail).
\item[(R3)] (R) given (J) of Proposition~\ref{i3:prop-admLq-in}: the cylinder idempotents $e*\varepsilon_{E_S}$ are idempotent, of positive
type, supported on elements unipotent modulo $g_{1-k}$, and bi-invariant under $K(g)$, i.e.\ admissible.
\item[(R4)] (R) \emph{Additivity} of the dimensions over label tuples (orthogonal central idempotents of the finite Hecke algebra
$\C[K/K(g)]$).
\item[(R5)] Use of Proposition~\ref{i3:prop-counting} (formally verified) for the counts of the cylinders; (R) the identification of
the counting problems $I_{k,S}$.
\end{enumerate}
The proof is mostly routine given Propositions~\ref{i3:prop-admLq-in} and \ref{i3:prop-counting}; (R1) is the new step.

\medskip\noindent\label{i3:str-terms}\emph{The term systems (Proposition~\ref{i3:prop-terms}).}

The formal statement assembles the following components (numbered as in Theorem~\ref{i3:S-thm-components}).
(R) The identification $j$ and the identities $j(\sum\mathrm{op}_t)=K(u,u')$, $j(\mathrm{part}(\mathrm{deg}))=K''(u,u')$, and
$j(\mathrm{part}(\mathrm{id}))=\mathrm{id}$ for $u=u'$: the kernel expansion (C7) and the transport of
Proposition~\ref{i3:S-prop-transport} (proved).  The kernel expansion itself rests on: legitimacy of the folding of the
Poincar\'e series (C8), the existence of fundamental domains for the stabilisers (proved, \S\ref{i3:ssec-pf-St-rest}), absolute
convergence of the orbital integrals (C9) (the integrability half of (FD), \S\ref{i3:str-FD}), their independence of the
fundamental domain, the well-definedness of the amplitudes, and (AC) (C10).  (AC), (ID), (RO), (CP), (NE), (GG) are the
conditions of (C19); for $u\ne u'$ their written proofs are Lemma~\ref{i3:S-lem-split} and Proposition~\ref{i3:S-prop-GG}, from the
Gelfand--Graev multiplicity one (L, Lemma~\ref{i3:S-lem-GGmult}); the formal statement covers also $u=u'$ by the same argument.  (N) The Parseval identity (NE): the finite-field Gram matrix at the auxiliary primes is the Gelfand--Graev
matrix of (D4).  (R) The orthogonality (RO) of distinct rows with the same index holds because $H_1$ is an orthogonal direct sum over the
pairs of row lattices (item (D3) before Lemma~\ref{i3:S-lem-link}).  This is a structural statement; its
published inputs are the Gelfand--Graev multiplicity one and, through (C9), the decay of discrete-series coefficients; its new content is the construction of the Poincar\'e-series model and the convergence statements (C8)--(C10),
which are shared with Proposition~\ref{i3:prop-S-closed}.

\medskip\noindent\label{i3:str-Wloc}\emph{The local estimates (Proposition~\ref{i3:prop-W-loc}).}

For a squarefree $P=\prod_j\Phi_j$ with $P(0)\ne0$, $\mathrm{sq}(P)=\prod_j(g^{\deg\Phi_j}-1)^{-1}$; $|GL_N(K)|\,\mathrm{sq}(P)$ is the
number of regular semisimple elements of $GL_N(K)$ with characteristic polynomial $P$ (the centraliser of such an element is
$\prod_j\mathbb F_{g^{\deg\Phi_j}}^\times$), so $\mathrm{sq}(K^N)$ is the proportion of regular semisimple elements.
\begin{enumerate}
\item[(W1)] (R; proved in Lemma~\ref{i3:lem-sqtype}) \emph{Exact type weights.}  For a degree type $(n_d)$ with $\sum dn_d=N$, the $\mathrm{sq}$-mass of the squarefree
polynomials of that type is $\prod_d\binom{I_d}{n_d}(g^d-1)^{-n_d}$, $I_d$ the number of monic irreducible polynomials of degree $d$
with nonzero constant term.
\item[(W2)] (N; proved in Lemma~\ref{i3:lem-sqtype}) \emph{Comparison with permutations}: the $\mathrm{sq}$-mass of every set of degree types is at most the proportion of
permutations of $\{1,\dots,N\}$ with those cycle types.  From (W1), since $I_d\le(g^d-1)/d$ gives
$\binom{I_d}{n_d}(g^d-1)^{-n_d}\le\frac1{d^{n_d}n_d!}$, and $\prod_d\frac1{d^{n_d}n_d!}$ is the proportion of permutations of cycle type
$(n_d)$ (Cauchy's formula, L).  This is the finite-field analogue of the cycle index comparison (J.~Fulman, Bull.\ AMS 39 (2002)
51--85; reference to be confirmed).
\item[(W3)] (L)+(R), \emph{item (I)}: $\mathrm{sq}(K^N)\ge1-2/g$; proved in Lemma~\ref{i3:lem-sqtype}, from the count of squarefree
polynomials (classical) and $\mathrm{sq}(P)\ge g^{-N}$, which is also the formal route.  The qualitative statement (the proportion of
regular semisimple elements of $GL_N(\mathbb F_g)$ is $1-1/g+O(g^{-2})$) is in G.~E.~Wall, Bull.\ Austral.\ Math.\ Soc.\ 60 (1999)
253--284, and J.~Fulman, P.~Neumann, C.~Praeger, Mem.\ AMS 830 (2005) (reference to be confirmed).
\item[(W4)] (N) \emph{Items (II), (III)}: counts of irreducible polynomials whose roots have non-injective prime-to-$l$ parts, or share
the prime-to-$l$ part of a given element: at most $g^{3\delta/4}$, resp.\ $g^{\delta/4}$, irreducibles of degree $\delta$ (when
$l^{4v_l(g^\delta-1)}\le g^\delta$); the number of prime factors of $g^\delta-1$ is at most $\delta\log_2g$ (R); summation over $\delta$
with geometric series bounds (R).  We do not describe the proofs of the two irreducible counts here.
\item[(W5)] (L)+(N) \emph{Items (IV), (V)} by (W2) (the formal proof of (IV) uses in addition two subset-sum statements,
whose proofs we do not describe): (IV) the proportion of permutations with an invariant set of size $a$, $0<a<N$, is
at most $\frac78$ (N, elementary; the asymptotics are in Eberhard--Ford--Green, Duke Math.\ J.\ 165 (2016) 2961--3001, reference to
be confirmed); (V) a permutation has a cycle of a given length $p>N/2$ with probability exactly $1/p$, and these events are disjoint, so the
proportion with a cycle of prime length in $(N/2,N-3]$ is $\sum_p1/p\ge\frac1{5\log N}$ for $N\ge100$ (L: explicit prime counts,
Rosser--Schoenfeld, Illinois J.\ Math.\ 6 (1962) 64--94; reference to be confirmed; the formal proof uses its own Chebyshev-type bound).
\end{enumerate}
This is elementary finite-field combinatorics: the qualitative statements are classical (cycle index, Wall), and the explicit
constants are new and short.  About half of it is published in substance.

\medskip\noindent\label{i3:str-NG}\emph{The non-generic row (Proposition~\ref{i3:S-prop-NG}).}

\begin{enumerate}
\item[(NG1)] (L)+(R) \emph{Masses from (SW).}  The stabiliser volumes and generic masses entering the row are evaluated by (SW)
(Theorems~\ref{i3:thm-BHC}, \ref{i3:thm-SWT}), as in the mass hypotheses (C17).
\item[(NG2)] (N) \emph{Choice of constants and pinned data}: base points, a pinned measure, local spaces; existence of a
pinning and the form of the non-generic terms.
\item[(NG3)] (N) \emph{Local input}: positivity of the local amplitudes and the theta window, from the kernel of the local model.
\item[(NG4)] (N) \emph{Point-sum input}: for each class of corank-one pairs, a lattice sum over block minors is bounded by
$e^{C_B(2r_S+2)^5f\log b}$ times the mean amplitude; it rests on a uniform joint norm bound, uniform amplitude bounds (through
Proposition~\ref{i3:prop-bessel}) and uniform pointwise bounds.
\item[(NG5)] (R) Assembly: grouping the corank-one pairs into classes and summing the class bounds.
\end{enumerate}
The published input is (SW); the rest is new.  We describe (NG2)--(NG4) only at the level of their statements.

\medskip\noindent\label{i3:str-inst}\emph{The instance (Proposition~\ref{i3:S-prop-instance}).}

\begin{enumerate}
\item[(I1)] (L) \emph{Poisson step.}  For a lattice $\Lambda$ in a euclidean space $V$, $\varphi\ge0$ with absolutely summable Fourier
transform on the dual lattice and $\sum_{l\in\Lambda}\varphi(l+x)=\mathrm{covol}(\Lambda)^{-1}\sum_{\mu\in\Lambda^*}\hat\varphi(\mu)e^{2\pi i\langle\mu,x\rangle}$:
every finite partial sum of $\varphi(l+x)$ is at most $\mathrm{covol}(\Lambda)^{-1}(\hat\varphi(0)+\sum_{\mu\ne0}|\hat\varphi(\mu)|)$.  This is
the Poisson summation formula (e.g.\ E.~M.~Stein, G.~Weiss, \emph{Introduction to Fourier analysis on euclidean spaces},
Princeton 1971, Ch.~VII, Cor.~2.6; reference to be confirmed) and a positivity remark; nothing extra.
\item[(I2)] (N)+(R) \emph{Lattice sums}: box counts in the ideals $g^{-1}\mathfrak D^{-\nu}$ of $E$, with explicit constants.
\item[(I3)] (N) \emph{Sizes and the threshold}: $\Theta_h\ge k^{c_*}$, from eight estimates on the local Bessel-type quantities
(interlacing, anti-concentration, threshold products) and monotonicity of the thresholds.
\item[(I4)] (N) \emph{Local-model consequences}: positivity of the kernels, the support condition, $e_i=0$ and minimal dominations,
from five statements on the explicit local kernels, whose proofs we do not describe.
\end{enumerate}
Only (I1) is published.

\medskip\noindent\label{i3:str-bridge}\emph{The bridge implications of Lemma~\ref{i3:S-lem-rankrow}.}

\begin{enumerate}
\item[(B-i)] (N)+(R) \emph{Frame row from Ranks}: a lower bound for the critical level of the family (at $\epsilon_0=\frac13$, $C_V=0$),
from explicit inequalities between the size parameters ($k\ge5$, $\log k\ge1$, $L\ge1$, $b\ge1$, $f\ge2$, $k^2\le g$), and the level
formula of the family.
\item[(B-ii)] (R) \emph{Size bound}: the ratio of critical levels is at most $e^{3/4}C_0e$, from the same inequalities.
\item[(B-iii)] (L)+(R) \emph{Size row}: the Siegel mass ratio for unramified CM fields (the first statement of
Proposition~\ref{i3:prop-massratio}; published in substance, \S\ref{i3:str-mass}) and the degree-ratio formula
$d_\sigma(\tau)=\dim V(\lambda(\tau))$ (Proposition~\ref{i3:thm-HC-fdeg}; published in substance, \S\ref{i3:str-HCfdeg}), together with
positivity of Vandermonde products (R).
\end{enumerate}
This is bookkeeping around two statements which are published in substance; the explicit inequalities are new but
elementary.

\medskip\noindent\label{i3:str-kppdiag}\emph{Items (J), (Del), (Cos), (Dim) of Proposition~\ref{i3:prop-kppdiag-in}.}

\begin{enumerate}
\item[(J)] (N) \emph{via the normal form.}  A relevant degenerate term has a normal form with index $j$ and scalar $c$ (the normal
form theorem of \S\ref{i3:ssec-NF-structure}); if $j=0$ the normal form forces the isotropic parts and the last vectors of the two
flags to agree, so the flags are equal; since the levels $T_w$ are prime to $g_0g_1$, a term with equal flags is the identity
term, which is excluded.  So $j(t)\ge1$.  The only non-routine input is the normal form theorem.
\item[(Del)] (L)+(N) \emph{Local index bound.}  $\Delta(j,\mathrm{lab})$ is a product of local norms of leading minors; each local norm is
bounded by $H_b(C)\,S(H_1,j,\mathrm{lab})^2$ via a bound for the norm of an algebraic integer by its archimedean sizes (product
formula and integrality, L) and the height bound of the datum (R).
\item[(Cos)] (L)+(N) \emph{Coset labels.}  The labels $\delta$ are tuples of algebraic integers (leading minors of integral Gram
matrices of unimodular lattices, R) with $\delta_0=1$ and $\delta_l=0$ for $l>j$; the sum $\sum_\delta S(H_1,j,\delta)^{-3r_S}$ is a
lattice sum over $O_E$, bounded by a norm lower bound for nonzero algebraic integers ($|N_{E/\Q}x|\ge1$, L) and a convergence bound
for lattice sums of negative powers of the archimedean sizes (Epstein-zeta type, L, e.g.\ via counting lattice points in boxes;
reference to be confirmed).  The explicit constant $21+4|C_H|$ is new.
\item[(Dim)] (N) \emph{Dimension formula}: from a transport of column spaces, a stabiliser bound and an index inequality (three
statements, whose proofs we do not describe).
\end{enumerate}
Thus (J) reduces to the normal form theorem, and (Del), (Cos) are explicit number-theoretic estimates with classical
inputs.

\medskip\noindent\label{i3:str-Sq}\emph{Item (Sq) of Proposition~\ref{i3:prop-admLq-in}.}

The proof of (Sq) in \S\ref{i3:ssec-pf-admLq-in} uses (HC), (GG) and the comparison of labels.  (HC) and (GG) are published (L), as
stated there; the formal development proves its own versions of them (for $GL_n(\mathbb F_g)$, $g\ge5$), which are special cases of the
published statements.  The comparison of labels is Lemma~\ref{i3:lem-labcomp}, proved in \S\ref{i3:ssec-pf-admLq-in}; it is (R) given (G) and the
three properties of the cuspidal labelling of the datum.  So (Sq) is proved in this paper from (G), the cited (HC), (GG) and
Theorem~\ref{i3:thm-zel}; Proposition~\ref{i3:prop-admLq-in} is therefore proved here.

\medskip\noindent\label{i3:str-levtr}\emph{Level transport (Proposition~\ref{i3:prop-levtr}).}

\begin{enumerate}
\item[(T1)] (R) \emph{Projection}: the orthogonal projection away from $\mathrm{span}(e_a,f_a)$ preserves the self-dual lattice
(unimodularity of the hyperbolic planes, R), in the global and in the adelic models.
\item[(T2)] (R) \emph{Splitting of $K$-membership}: an element of $N_{f,u}\cap K$ is determined by its action on the hyperbolic part and
on the complement.
\item[(T3)] (L)+(R) \emph{B\'ezout step}: there is $C$ prime to $g$ with $C\,e_a$, $C\,f_a$ in the adelic lattice (multiples of the
frame vectors), and if $\psi_E(Cy)=1$ and $\psi_E(gy')=1$-type congruences hold then $\psi_E(y)=1$ (B\'ezout for coprime $C,g$
and additivity of $\psi_E$; L: the conductor of $\psi_E$ is the inverse different, R).
\item[(T4)] (N) \emph{Congruences}: for $n\in N_{f,u}\cap K$, $\kappa=x_u^{-1}nx_u$ lies in $K$, and the hermitian pairings
$h(f_a,\kappa\,\mathrm{nxt}(e,z)_a)$ and their transports differ by elements of $g\widehat O$ (after multiplication by $C^2$), which with
(T3) gives $\chi_N(\kappa)=\psi_{f,u}(n_f)$; the membership $\kappa\in K_2$ follows from the same congruences.  This is a long
explicit computation with the integral frame of the Hida setting.
\item[(T5)] (R) \emph{(Car)}: the representative of an admissible term is compatible with the flags modulo $g$, from the definition of
admissibility and the congruence of the frames.
\end{enumerate}
These are elementary lattice congruences; no published theorem beyond B\'ezout and the conductor of the standard character
enters, and the length of the proof lies in the bookkeeping, not in the ideas.

\medskip\noindent\label{i3:str-FD}\emph{The integrability half of (FD).}

The second half of item (FD) of Proposition~\ref{i3:prop-kppoff-in} (integrability of the orbital integrands for every term) is
derived in the formal development from a single \emph{lattice sum} statement: for words $u,u'$ with legitimate folds and fundamental
domains, the integral over $\Omega_t\times(\text{finite adeles})$ of $|\phi|$ times the sum over the second Poincar\'e series is finite.
Its proof:
\begin{enumerate}
\item[(D1)] (N) \emph{Slab cover}: the domain is covered by slabs adapted to the double-coset unfolding (regrouping of the terms).
\item[(D2)] (N) \emph{Pointwise majorant by columns}: on each slab the integrand is bounded by a product of column decay factors, from
the decay of the archimedean coefficients (Lemma~\ref{i3:lem-decay}, via the gap).
\item[(D3)] (R) given Proposition~\ref{i3:prop-levtr} \emph{invariance}: the integrand is invariant under the rational elements of the
stabiliser (this is where level transport enters).
\item[(D4)] (N)+(R) \emph{Counting}: the number of rational points per slab is bounded by a lattice-point count.
\item[(D5)] (R) \emph{Neatness}: the unipotent elements of $K_u$ give a neat subgroup, so the unfolding has no multiplicities.
\end{enumerate}
The argument is standard in shape (unfolding plus majorants, as in the convergence of Poincar\'e series; Borel, \emph{Introduction
aux groupes arithm\'etiques}, Hermann 1969, for the reduction-theoretic background; reference to be confirmed), and new in its explicit
majorants.  We describe (D1), (D2), (D4) only at the level of their statements.



\medskip\noindent\textbf{Status.}
\emph{Proved here:} Theorem~\ref{i3:thm-main} from the statements listed below; in particular the reduction of the
off-diagonal bound to its leaves (Proposition~\ref{i3:prop-S-closed} and the one-level inequality), Bessel uniqueness
(Proposition~\ref{i3:prop-bessel}), the boundedness of Bessel maps on orbits (Proposition~\ref{i3:prop-orbbd}), the
passage to almost orthogonality (Theorem~\ref{i3:thm-linkS-slots}), the irreducibility bound
(Proposition~\ref{i3:prop-irr-row}), the level counts (Propositions~\ref{i3:prop-V}, \ref{i3:prop-part5} and
\ref{i3:prop-countF}), the unitary Siegel--Weil statement (SW) (Proposition~\ref{i3:prop-SW}), Proposition~\ref{i3:prop-admLq-in},
and the proved items of Propositions~\ref{i3:prop-kppoff-in}, \ref{i3:prop-cov-in}, \ref{i3:prop-St-rest},
\ref{i3:prop-kppdiag-in} and \ref{i3:prop-W-loc}, as indicated in their headers.
\emph{Published theorems, quoted:} Theorems~\ref{i3:thm-hc}, \ref{i3:thm-godement}, \ref{i3:thm-knapp}, \ref{i3:thm-BHC},
\ref{i3:thm-SWT}, \ref{i3:thm-landherr}, \ref{i3:thm-helgason}, \ref{i3:thm-bruhat} and \ref{i3:thm-zel}, the facts (F2), (F3)
of \S\ref{i3:S-ssec-leafproofs}, and Green's parametrisation (G), which enters as a hypothesis; Theorem~\ref{i3:thm-shortprimes}
is proved in Section~\ref{i4:sec}.
\emph{Formally verified, \cite{CharacterFamily}:} Propositions~\ref{i3:prop-amp-in}, \ref{i3:prop-massratio},
\ref{i3:thm-HC-fdeg}, \ref{i3:prop-counting}, \ref{i3:prop-allrows-built}, \ref{i3:prop-terms}, \ref{i3:prop-model},
\ref{i3:prop-GKin}, \ref{i3:S-prop-instance}, \ref{i3:S-prop-NG}, \ref{i3:S-prop-leaves}, \ref{i3:S-prop-regions},
\ref{i3:S-prop-count} and \ref{i3:prop-levtr}, Lemma~\ref{i3:lem-dataB}, and the items so marked in
Propositions~\ref{i3:prop-kppoff-in}, \ref{i3:prop-cov-in}, \ref{i3:prop-St-rest}, \ref{i3:prop-kppdiag-in} and
\ref{i3:prop-W-loc}.  The hypotheses (G), (AT), (VT) and the slot hypotheses of Theorem~\ref{i3:thm-linkS-slots} are stated
where they are used; Proposition~\ref{i3:prop-slots-SW} discharges part of the slot hypotheses.



\section{The windowed Selmer bound}\label{i5:sec}
\providecommand{\Sha}{\mathrm{III}}
\providecommand{\lgth}{\operatorname{lg}}

This section treats the fifth arithmetic input of
Theorem~\ref{M:thm:alpoge-families}: a bound, uniform over a window of scales,
for the total $l$-adic size of the Selmer groups of the Tate twists
$\Q_l/\Z_l(-e)$ over the CM abelian fields attached to a Dirichlet character
and a pair of large primes.  We first state the bound
(\S\ref{i5:ssec-statement}) and the Galois input it uses
(\S\ref{i5:ssec-galois}), then reduce it, by elementary and
Galois-cohomological arguments written in full, to a small number of
propositions: a Soul\'e-type effective bound for $H^2$, four automorphic
counts, an Euler-system bound for the minus part, and an Iwasawa-theoretic
bound at a finite set of exceptional primes.  Statements marked
\emph{formally verified} are proved in \cite{CharacterFamily}; we state them
exactly and say what their proofs involve, but do not reproduce them.
Published theorems are stated in ``Theorem (Author)'' form, with the
hypotheses checked.

The main result is Theorem~\ref{i5:thm-vwin}.  Subsections~\ref{i5:ssec-statement}--\ref{i5:ssec-bab}
contain the definitions and statements: the bound itself, the Galois and cohomological
inputs, the reduction to three bounds and to characters over $\Q$, the prime-by-prime
cases, the Poitou--Tate step, the automorphic counts, the minus part and the exceptional
primes.  Subsection~\ref{i5:ssec-proof} proves Theorem~\ref{i5:thm-vwin} from these
statements, and \S\S\ref{i5:ssec-pf-minus}--\ref{i5:x5-ssec-arch} prove them, apart from the
formally verified cores, whose definitions are collected in \S\ref{i5:x5-ssec-defs}.
Subsection~\ref{i5:x6-ssec} outlines the formally verified steps, and the section ends
with a summary of its status.

\medskip
\noindent\emph{Conventions.} For a module $X$ over a local ring $R$,
$\lgth_R X\in\N\cup\{\infty\}$ is its length; for $a\in\N\cup\{\infty\}$,
$a^\natural=a$ if $a<\infty$ and $a^\natural=0$ otherwise.  Galois cohomology
is continuous cochain cohomology.  For a prime $l$, a \emph{coefficient ring}
is a finite extension $\mathcal O$ of $\Z_l$ inside $\overline\Q_l$, with
uniformizer $\varpi$, valuation $v_{\mathcal O}$ and
$e_{\mathcal O}=v_{\mathcal O}(l)=\lgth_{\mathcal O}(\mathcal O/l)$.  We write
$G=G_\Q$ and $\chi$ for the $l$-adic cyclotomic character.

\subsection{Statement of the windowed bound}\label{i5:ssec-statement}

\begin{definition}[$l$-data and Selmer groups]\label{i5:def-ldata}
Let $F$ be a number field, $\Omega$ an algebraic closure and $l$ a prime.  An
\emph{$l$-datum} $D$ for $(F,\Omega)$ consists of the set $S$ of places of $F$
above $l$, for each $v\in S$ a valuation ring of $\Omega$ above $v$ with
decomposition group $D_v\subset G_F$, and the factorisation
$\chi_D\colon G_{F,S}\to\Z_l^\times$ of the cyclotomic character through
$G_{F,S}=G_F/N_S$, where $N_S$ is the closed normal subgroup generated by the
inertia groups at the finite places outside $S$.  Put
\[
 \Sha^1_D(F,\Q_l/\Z_l(-e))=\ker\Bigl(H^1(G_{F,S},\Q_l/\Z_l(-e))\to\prod_{v\in S}H^1(D_v,\Q_l/\Z_l(-e))\Bigr),
\]
the Selmer group with the strict condition above $l$ and the unramified
condition elsewhere, and $\operatorname{sha}_D(e)=\lgth_{\Z_l}\Sha^1_D$.  Put
also $V_D(e)=\lgth_{\Z_l}H^2(G_{F,S},\Z_l(e+1))$ and
$\Sha^2_D(F,\Z_l(e+1))=\ker\bigl(H^2(G_{F,S},\Z_l(e+1))\to\prod_{v\in S}H^2(D_v,\Z_l(e+1))\bigr)$.
\end{definition}

An $l$-datum always exists: choose valuation rings above each $v\mid l$; the
cyclotomic character is unramified outside $l$, so it factors through
$G_{F,S}$.  Archimedean places are not included in $\Sha^1_D$; this is
immaterial below, since all fields that occur are CM (Lemma~\ref{i5:lem-cm}).

\begin{definition}[the windowed bound]\label{i5:def-vwin}
(1) For a number field $E$, $t\in\N$ and $B\in\R$, $\mathrm{VV}(E,t,B)$ means:
for every algebraic closure of $E$ and every family $(D_p)_p$ of $p$-data,
$\operatorname{sha}_{D_p}(e)<\infty$ for all primes $p$ and $0<e<t$, and for
every finite set $\mathcal S$ of primes
$\sum_{l\in\mathcal S}\log l\sum_{e=1}^{t-1}\operatorname{sha}_{D_l}(e)^\natural\le B$.

(2) Let $\psi$ be a primitive Dirichlet character of conductor $N_\psi$ and
order $o_\psi$.  For $e,q,d\in\N$, a number field $E$ is a
\emph{$(\psi,e,q,d)$-field} if there are $N$ with $2^{e+2},N_\psi,q,d\mid N$,
an embedding $j\colon E\hookrightarrow\Q(\zeta_N)$ and a primitive $N$-th root
of unity $\zeta$ such that $\sigma_a\colon\zeta\mapsto\zeta^a$ fixes $j(E)$ if
and only if $a\equiv\pm1\bmod 2^{e+2}$, $\psi(a)=1$ and
$(\frac aq)=(\frac ad)=1$.

(3) For $\varepsilon,Y>0$, $(q,d)$ is a \emph{window pair}
$(q,d)\in\mathrm{WP}(\varepsilon,Y)$ if $q\equiv3\bmod8$ and $d\equiv7\bmod16$
are primes, $\tfrac12Y^{\varepsilon/5}<\log q\le Y^{\varepsilon/5}$,
$\tfrac12Y^{\varepsilon/5}(\log Y)^2<\log d\le Y^{\varepsilon/5}(\log Y)^2$,
$q\mid d+1$, and every odd prime $p\le Y^{\varepsilon/5}$ divides $d+1$.  An
integer $e$ is \emph{on page $x$}, $\mathrm{Page}_\varepsilon(x,e)$, if
$x^{2+8\varepsilon}\le2^e\le2x^{2+8\varepsilon}$.

(4) $\mathrm{VWinAt}(\psi,C_t)$: for all $0<\varepsilon<1/52$ and every
$e\colon\R\to\N$ with $\mathrm{Page}_\varepsilon(x,e(x))$ for $x\ge1$, there
are $J\ge1$ and $X_0$ such that for all $X\ge X_0$ and every
$(q,d)\in\mathrm{WP}(\varepsilon,X^{1/J})$ there is $Y$ with $X^{1/J}\le Y$,
$Y^2\le X$, such that for every $x\in[Y,Y^2]$ and every
$(\psi,e(x),q,d)$-field $E$,
$\mathrm{VV}\bigl(E,\lfloor x^{C_t}\rfloor,\sqrt{\exp(x^\varepsilon)}\bigr)$.
$\mathrm{VWin}$ is the statement that $\mathrm{VWinAt}(\psi,401)$ holds for every
primitive $\psi$.
\end{definition}

A function $e$ as in (4) exists: for $x\ge1$ choose $n$ with
$2^n\le x^{2+8\varepsilon}<2^{n+1}$ and put $e(x)=n+1$.  If $\mathrm{VV}(E,t,B)$,
$t'\le t$ and $B\le B'$, then $\mathrm{VV}(E,t',B')$, since all terms are
non-negative.

\begin{theorem}[windowed Selmer bound]\label{i5:thm-vwin}
$\mathrm{VWin}$ holds, given:
\textup{(i)} the published theorems of
\S\ref{i5:ssec-galois}, Kato's Theorem~\ref{i5:x-thm:kato},
the Ferrero--Washington theorem~\ref{i5:thm-fw}, Stickelberger's theorem,
Chebotarev's density theorem, the classical results cited in
\S\ref{i5:ssec-pf-shak} and \S\ref{i5:x5d-ssec}, the theorems of Quillen,
Milnor--Moore and Borel (and the case $i=1$ via Dirichlet's unit theorem)
used for Proposition~\ref{i5:x5-prop-bodd}, and van der Kallen's homological
stability theorem with Bass's stable-rank bound, used in
Lemma~\ref{i5:x5-lem-k3}(a);
\textup{(ii)} the formally verified Propositions~\ref{i5:x5-prop-U1},
\ref{i5:x5-prop-U234} and \ref{i5:x5-prop-U4} (for $\mathrm{OneCharU}$),
the two counting statements derived from $(\mathrm{C}_K)$, $(\mathrm{C}_\Q)$,
$(\mathrm N)$ stated after Proposition~\ref{i5:x-prop:u3count} (for
$\mathrm{NumberU3U}$, themselves conditional on the families of published
statements listed there), \ref{i5:x5-prop-sp4core}
(for $\mathrm{Count}_{\rm II}$, $\mathrm{Count}_{\rm III}$), and
\ref{i5:x5-prop-Vcore}, Lemma~\ref{i5:x5-lem-Vopen} and
\ref{i5:x5-prop-Vhost} (for $\mathrm{Count}_{\rm V}$; the cores are stated
exactly with the definitions of \S\ref{i5:x5-ssec-defs}),
together with the purely algebraic lemma quoted in the proof of
Lemma~\ref{i5:x5-lem-irrcyc}.  The cores of (ii) are verified relative to
the published statements of Theorem~\ref{i5:thm-pt} (in (i)), which the formal
development itself proves, and to
Corollary~\ref{i5:cor-lemmaK}: $\mathrm{OneCharU}$ (through (U3$'$), (U4)),
$\mathrm{Count}_{\rm II}$, $\mathrm{Count}_{\rm III}$ (through
Proposition~\ref{i5:x5-prop-sp4core}(f)) and $\mathrm{Count}_{\rm V}$ (through
Proposition~\ref{i5:x5-prop-Vhost}(d)) use Corollary~\ref{i5:cor-lemmaK}.
There is no circularity: Corollary~\ref{i5:cor-lemmaK} follows from
$\mathrm{ShaK}$ (Proposition~\ref{i5:prop-shak}), whose proof uses none of
these cores.
Everything else, including Propositions~\ref{i5:prop-psibound},
\ref{i5:prop-stick} and \ref{i5:prop-thmbab}, is proved in this section.  In
words: for the CM field $E$ attached to
$(\psi,q,d)$ and a page $e(x)$,
\[
 \sum_l\log l\sum_{1\le e<x^{401}}\lgth\Sha^1(E,\Q_l/\Z_l(-e))\le\sqrt{\exp(x^\varepsilon)} .
\]
\end{theorem}

\begin{remark}[the trivial character]\label{i5:rem-trivial}
For $\psi=\mathbf 1$ ($N_\psi=o_\psi=1$) the field is explicit.  The fixed
field of $\{a\equiv\pm1\bmod2^{e+2}\}$ in $\Q(\zeta_{2^{e+2}})$ is the real
dyadic field $\Q(\zeta_{2^{e+2}})^+$, of degree $2^e$; since
$q\equiv d\equiv3\bmod4$, the fixed field of the squares in $(\Z/q)^\times$
(resp.\ $(\Z/d)^\times$) is $\Q(\sqrt{-q})$ (resp.\ $\Q(\sqrt{-d})$).  Hence
\[
 E=\Q(\zeta_{2^{e+2}})^+\bigl(\sqrt{-q},\sqrt{-d}\bigr),\qquad [E:\Q]=2^{e+2},
\]
the three fields being linearly disjoint (they ramify at disjoint sets of
primes); $E$ is CM with $E^+=\Q(\zeta_{2^{e+2}})^+(\sqrt{qd})$.  The
exceptional set of Definition~\ref{i5:def-Bchi} is then $B_{\mathbf1}=\{3,7\}$
($p_K=7$, $h^*=h(\Q(\sqrt{-7}))=1$, $p_1=3$), and the conductor of $E$ is
$N_h=2^{e+2}qd$ (Lemma~\ref{i5:lem-hostcond}).  The reader may keep this case
in mind throughout; no step simplifies in an essential way, except that all
characters $\psi'$ below then have order dividing $2^{e+1}$ and Case~IV can
only occur for $l\in\{3,7\}$.
\end{remark}

\begin{remark}[role]\label{i5:rem-role}
$\mathrm{VWin}$ supplies the Selmer hypothesis of the existence input
(Section~\ref{i4:sec}) and, through it, bounds the Selmer sum entering condition (P1) of
Definition~\ref{M:def:package}.  It is a hypothesis of
Theorem~\ref{M:thm:alpoge-families}, not a condition on a single package.
\end{remark}

\subsection{The Galois and cohomological inputs}\label{i5:ssec-galois}

The proof uses Galois representations attached to regular algebraic polarized
cuspidal representations of $\mathrm{GL}_a$ over CM fields, through a weak form
of local--global compatibility.  We state this weak form and derive it from the
published theorems.

\medskip\noindent\emph{Local parameters.}
Let $E$ be a number field, $w$ a finite place of norm $q_w$, $W_{E_w}$ the
Weil group (inside $G_E$ via a place of $\overline\Q$ above $w$), and
$\mathrm{Irr}(E_w)$ the irreducible smooth representations of
$\mathrm{GL}_n(E_w)$, $n\ge1$.  For $\sigma\in W_{E_w}$, $\deg\sigma\in\Z$ is
defined by $\sigma(x)\equiv x^{q_w^{\deg\sigma}}$ on the residue field (so
arithmetic Frobenius has degree $1$).

\begin{definition}[local Artin maps]\label{i5:def-artin}
A \emph{global Artin map} for $E$ is a continuous open surjection
$\mathrm{art}\colon\mathbb A_E^\times\to G_E^{\rm ab}$ with kernel the closure
of $E^\times(E\otimes\R)^{\times,0}$, mapping $\mathcal O_v^\times$ onto the image
of inertia and uniformizers to geometric Frobenius elements, and such that
for every finite place $v$ the closure of $\mathrm{art}(E_v^\times)$ is the
image of the decomposition group at $v$.  A \emph{local
Artin map} at $w$ is a surjective homomorphism $a\colon W_{E_w}\to E_w^\times$,
continuous on inertia, with $\mathrm{ord}_w a(\sigma)=-\deg\sigma$, such that
every smooth character of $W_{E_w}$ has the form $\chi\circ a$ and
$\mathrm{art}(\iota_wa(\sigma))=\sigma|_{E^{\rm ab}}$ for every global Artin
map.
\end{definition}

\begin{lemma}\label{i5:lem-artin}
Every local Artin map at $w$ equals $\mathrm{Art}_{E_w}^{-1}$, the inverse of
local reciprocity normalised so that uniformizers go to geometric Frobenius.
\end{lemma}

\begin{proof}
By local--global compatibility of class field theory,
$u(\sigma)=a(\sigma)\mathrm{Art}_{E_w}^{-1}(\sigma)^{-1}$ lies in the kernel of
$\mathrm{art}_E\circ\iota_w$.  Let $x$ be in this kernel.  Acting on
$E(\mu_M)$ with $M=q_w^k-1$, $x$ acts as $\zeta\mapsto\zeta^{q_w^{-\mathrm{ord}\,x}}$,
so $\mathrm{ord}_w x=0$.  Let $U\subset\mathcal O_w^\times$ be open of index $m$.
By Chevalley's theorem (Chevalley, J.\ Math.\ Soc.\ Japan 3 (1951), 36--44, Th\'eor\`eme~1) there is
$\mathfrak a$ prime to $w$ such that every unit $\equiv1\bmod\mathfrak a$ is an
$m$-th power of a unit.  Since $\iota_w(x)$ lies in the closure of
$E^\times(E\otimes\R)^{\times,0}$, there is $\epsilon\in E^\times$ with
$\epsilon^{-1}\iota_w(x)$ in the open subgroup
$U\times\prod_{v\mid\mathfrak a}(1+\mathfrak a\mathcal O_v)\times\prod_{\text{other }v}\mathcal O_v^\times\times(E\otimes\R)^\times$.
Comparing components, $\epsilon$ is a unit $\equiv1\bmod\mathfrak a$, hence an
$m$-th power, hence in $U$; so $x\in U$.  As $U$ is arbitrary, $x=1$.
\end{proof}

\begin{definition}[good classes; $\mathrm{rec}^T$-maps]\label{i5:def-recT}
A supercuspidal $y\in\mathrm{Irr}(E_w)$ is \emph{good} if it is a character of
$\mathrm{GL}_1$, or $E_w\cong\Q_g$ and $y$ is a character twist of a
representation with nonzero vectors fixed by $1+M_n(\mathfrak p)$ (depth zero);
$x\in\mathrm{Irr}(E_w)$ is good if its cuspidal support consists of good
supercuspidals.  A map $\mathrm{rc}$ from $\mathrm{Irr}(E_w)$ to
Frobenius-semisimple Weil--Deligne representations over $\C$ is a
\emph{$\mathrm{rec}^T$-map} if, with $a$ a local Artin map:
$\dim\mathrm{rc}(x)=n$ on $\mathrm{GL}_n$; $\mathrm{rc}(\chi)=\chi\circ a$ for
characters; $\mathrm{rc}$ of a Langlands quotient of segments
$[\rho,\dots,\rho|\det|^{k-1}]$ is $\bigoplus\mathrm{rc}(\rho)\otimes\mathrm{Sp}_k$;
supercuspidals go to irreducible representations; $\mathrm{rc}$ commutes with
$|\det|^s$-twists; $|\det\mathrm{rc}(x)(\sigma)|=|\omega_x(a(\sigma))|$; and,
for good $x$ only: $\mathrm{rc}$ commutes with character twists and
contragredients, $\det\mathrm{rc}(x)=\omega_x\circ a$, depth zero implies
tame; for arbitrary $x$ and good $y$, $\mathrm{rc}(x)=\mathrm{rc}(y)$ implies
$x=y$; and for the depth-zero
supercuspidal $\pi_\theta$ over $\Q_g$ attached to a regular character $\theta$
of $\mathbb F_{g^n}^\times$, $\mathrm{rc}(\pi_\theta)$ has trivial monodromy and
the characteristic polynomial of inertia is
$\prod_{j<n}(X-\theta(z^{g^j}))$, $z$ the reduction of $a'(t)$ for the local
Artin map $a'$ of the unramified extension of degree $n$.
\end{definition}

\begin{theorem}[Harris--Taylor 2001, Thm.~A; Henniart 2000]\label{i5:thm-llc}
For a $p$-adic field $K$ there is a unique family of bijections $\mathrm{rec}_K$
from $\mathrm{Irr}(\mathrm{GL}_n(K))$ to $n$-dimensional Frobenius-semisimple
Weil--Deligne representations, sending $\chi$ to
$\chi\circ\mathrm{Art}_K^{-1}$, compatible with $L$- and $\varepsilon$-factors of
pairs, central characters, contragredients and character twists, sending
supercuspidals to irreducible representations, and compatible with the
Langlands classification (segment $\mapsto\mathrm{rec}_K(\rho)\otimes\mathrm{Sp}_k$).
[Harris--Taylor, Ann.\ of Math.\ Stud.\ 151 (2001); Henniart, Invent.\ Math.\ 139 (2000), 439--455; reference to be confirmed, including the normalisations.]
\end{theorem}

\begin{theorem}[Aubert--Baum--Plymen--Solleveld, \emph{Depth and the local Langlands correspondence}, Progr.\ Math.\ 319 (2016), 17--41, Thm.~2.9 (with $m=1$, $D=F$)]\label{i5:thm-abps}
$\pi$ has depth zero if and only if $\mathrm{rec}_K(\pi)$ is trivial on wild
inertia.
\end{theorem}

\begin{theorem}[explicit level zero: Henniart, Math.\ Nachr.\ 158 (1992), 7--26; Bushnell--Henniart, J.\ Amer.\ Math.\ Soc.\ 18 (2005), 685--710; reference to be confirmed]\label{i5:thm-lz}
For $K\cong\Q_g$, $\mathrm{rec}_K(\pi_\theta)\cong
\mathrm{Ind}_{W_{K_n}}^{W_K}(\tilde\theta\circ\mathrm{Art}_{K_n}^{-1})$ with
$\tilde\theta$ tamely ramified and equal to the inflation of $\theta$ on
$\mathcal O_{K_n}^\times$ up to an unramified twist, all normalised
geometrically.
\end{theorem}

\begin{proposition}[local parameters exist]\label{i5:prop-recT}
For every number field $E$, finite place $w$ and place of $\overline\Q$ above
it, $\mathrm{rec}_{E_w}$ is a $\mathrm{rec}^T$-map.
\end{proposition}

\begin{proof}
Take $a=\mathrm{Art}_{E_w}^{-1}$ (Lemma~\ref{i5:lem-artin} shows there is no
other choice).  All axioms except the last two are clauses of
Theorem~\ref{i5:thm-llc}; injectivity holds since $\mathrm{rec}$ is bijective.
``Depth zero implies tame'' is Theorem~\ref{i5:thm-abps}.  For the type law,
restrict the induced representation of Theorem~\ref{i5:thm-lz} to inertia:
since $K_n/K$ is unramified, $I_K=I_{K_n}$, and by Mackey the restriction is
$\bigoplus_{j<n}(\tilde\theta\circ\mathrm{Art}_{K_n}^{-1})^{\mathrm{Frob}^j}$, whose
eigenvalues at $t$ are the $\theta(z^{g^j})$; unramified twists do not change
it.
\end{proof}

\begin{remark}
\cite{CharacterFamily} proves Proposition~\ref{i5:prop-recT} independently, with
a witness that agrees with a Langlands-type map only on good classes.  The
axioms of Definition~\ref{i5:def-recT} constrain $\mathrm{rc}$ on good classes
only; this is all that is used below.
\end{remark}

\medskip\noindent\emph{Galois representations.}

\begin{definition}[weak compatibility]\label{i5:def-away}
Here and below, Weil--Deligne representations are taken Frobenius-semisimple and up to
isomorphism; for a continuous $r$, $\mathrm{WD}(r|_{G_{E_w}})$ (resp.\ $\mathrm{WD}_{\rm pst}(r|_{G_{E_w}})$ at $w\mid l$)
denotes the Frobenius-semisimplification of the Weil--Deligne representation attached to
$r|_{G_{E_w}}$ by Grothendieck's monodromy theorem (resp.\ by Fontaine's functor $D_{\rm pst}$).
For Weil--Deligne representations $x,y$ of $W_{E_w}$, $\mathrm{Away}(x,y)$
means: (i) $\dim x=\dim y$; (ii) if $y$ has abelian image and trivial
monodromy, then $x\cong y$; (iii) if $E_w\cong\Q_g$ with $g>16(\dim y)^2$ or
$\dim y=1$, and $y$ is tame, has trivial monodromy and is \emph{tamely
square-free} (every repeated root of the characteristic polynomial of a
generator of tame inertia is a $(q_w-1)$-st root of unity), then $x$ has
trivial monodromy and the same characteristic polynomials as $y$ on inertia.
\end{definition}

\begin{theorem}[Barnet-Lamb--Gee--Geraghty--Taylor, Ann.\ of Math.\ (2) 179 (2014), Thm.~2.1.1, collecting Shin 2011, Chenevier--Harris 2013, Caraiani 2012; at $w\mid l$ beyond the Iwahori case: Barnet-Lamb--Gee--Geraghty--Taylor, \emph{Local-global compatibility for $l=p$} I, II, and Caraiani 2014]\label{i5:thm-blggt}
Let $E$ be CM and $(\pi,\chi)$ a regular algebraic polarized cuspidal
automorphic representation of $\mathrm{GL}_a(\mathbb A_E)$, $\iota\colon
\overline\Q_l\cong\C$.  There is a continuous semisimple
$r=r_{l,\iota}(\pi)\colon G_E\to\mathrm{GL}_a(\overline\Q_l)$ with
$r^c\cong r^\vee\otimes\epsilon_l^{1-a}r_{l,\iota}(\chi)|_{G_E}$, such that
for all finite $w$,
$\iota\mathrm{WD}(r|_{G_{E_w}})^{\rm F\text{-}ss}\cong\mathrm{rec}_{E_w}(\pi_w\otimes|\det|^{(1-a)/2})$
(with $\mathrm{WD}$ formed by Fontaine's functor if $w\mid l$); for $w\mid l$,
$r|_{G_{E_w}}$ is potentially semistable, crystalline if $\pi_w$ is unramified,
with $\tau$-Hodge--Tate weights $\{a_{\iota\tau,i}+a-i\}$ determined by the
infinitesimal character of $\pi$ at $\iota\tau$.
\end{theorem}

\begin{remark}[what each source supplies]\label{i5:rem-blggt-src}
BLGGT, Thm.~2.1.1 (arXiv:1010.2561, checked; its conventions, $\epsilon_l$ of
Hodge--Tate weight $-1$, geometric $\mathrm{Art}$, twist $|\det|^{(1-a)/2}$,
are those above) gives: the polarisation; the displayed isomorphism at every
$w\nmid l$ (its part (2)); de Rham with the stated Hodge--Tate weights at every
$w\mid l$ (part (3)); and the displayed isomorphism at $w\mid l$, with
semistability, and crystallinity when $\pi_w$ is unramified, only when $\pi_w$
has an Iwahori-fixed vector (part (4)). Potential semistability at every
$w\mid l$ follows from de Rham (Berger, Invent.\ Math.\ 148 (2002), 219--284).
For general $\pi_w$ at $w\mid l$: Barnet-Lamb--Gee--Geraghty--Taylor,
\emph{Local-global compatibility for $l=p$}, I (Ann.\ Fac.\ Sci.\ Toulouse
Math.\ 21 (2012), 57--92) and II (Ann.\ Sci.\ ENS 47 (2014), 165--179), give
the displayed isomorphism up to semisimplification, that is, with monodromy
forgotten; Caraiani, \emph{Monodromy and local-global compatibility for $l=p$},
Algebra Number Theory 8 (2014), 1597--1646, adds the monodromy operator for
$a$ even and extends the compatibility to Frobenius-semisimplification. With
these, the theorem as stated is covered at $w\mid l$ for $a$ even. For $a$ odd
and $\pi_w$ without Iwahori-fixed vectors, agreement of the monodromy
operators at $w\mid l$ is not covered by the statements checked (bibliographic
data of the two $l=p$ papers taken from a secondary reference list).
\end{remark}

\begin{proposition}[Galois input]\label{i5:prop-row0b}
Let $E$, $\pi$, $\iota$ be as in Theorem~\ref{i5:thm-blggt} and
$\sigma_w=\pi_w\otimes|\det|^{(1-a)/2}$.  There are a continuous semisimple
$r\colon G_E\to\mathrm{GL}_a(\overline\Q_l)$, an integer $wt$ and multisets
$HT(\tau)$ such that: (1) for every finite $w\nmid l$,
$\mathrm{Away}(\mathrm{WD}(r|_{G_{E_w}}),\mathrm{rc}(\sigma_w))$ for
$\mathrm{rc}=\mathrm{rec}_{E_w}$; (2) $HT(\tau)=\{\frac{a-1}2-z\}$, $z$ running
over the infinitesimal character of $\pi$ at $\iota\tau$, and
$HT(\tau\circ c)=wt-HT(\tau)$; (3) at $w\mid l$, $r|_{G_{E_w}}$ is crystalline
with jumps $HT(\tau)$ if $\pi_w$ is spherical, potentially semistable with jumps
$HT(\tau)$ in general, and
$\mathrm{Away}(\mathrm{WD}_{\rm pst}(r|_{G_{E_w}}),\mathrm{rec}_{E_w}(\sigma_w))$.
\end{proposition}

\begin{proof}
Take $r=r_{l,\iota}(\pi)$.  By Theorem~\ref{i5:thm-blggt},
$x=\mathrm{WD}(r|_{G_{E_w}})^{\rm F\text{-}ss}\cong\mathrm{rec}_{E_w}(\sigma_w)$
(at $w\mid l$ with Fontaine's functor; for the sources at $w\mid l$ see
Remark~\ref{i5:rem-blggt-src}).  With the convention of Definition~\ref{i5:def-away}, $\mathrm{WD}(r|_{G_{E_w}})$ is the
Frobenius-semisimplification, so Theorem~\ref{i5:thm-blggt} reads
$\mathrm{WD}(r|_{G_{E_w}})\cong\mathrm{rec}_{E_w}(\sigma_w)$, and clauses (i)--(iii) of $\mathrm{Away}$
follow at once: (i) and (ii) from the isomorphism, (iii) because isomorphic Weil--Deligne
representations have the same monodromy operator and the same characteristic polynomials on
inertia. No statement about semisimplicity of Frobenius on the unsemisimplified representation
is used. At $w\mid l$ the isomorphism, including the monodromy operator, is
part (4) of Theorem~\ref{i5:thm-blggt} when $\pi_w$ has an Iwahori-fixed
vector; for the remaining $\pi_w$ see Remark~\ref{i5:rem-blggt-src}.  The conventions agree by Lemma~\ref{i5:lem-artin} (geometric
normalisation of $\mathrm{Art}$) and because the twist $|\det|^{(1-a)/2}$ is the
one of BLGGT.  For (2), with Hodge--Tate weights counted on covariant
$D_{\rm dR}$ (so $\Q_l(1)$ has weight $-1$, as in BLGGT) and the infinitesimal
character written as $\lambda+\rho$, the weights $a_{\iota\tau,i}+a-i$ equal
$\frac{a-1}2-z_i$ with $z_i=-a_{\iota\tau,i}+i-\frac{a+1}2$; the relation for
$\tau\circ c$ follows from Clozel's purity lemma (Clozel 1990, Lemme~4.9;
reference to be confirmed): $\mathrm{inf}_{\bar\tau}(\pi)=w_0-\mathrm{inf}_\tau(\pi)$,
so $wt=a-1-w_0$.
\end{proof}

\begin{remark}[the form used by the counts]\label{i5:rem-row0b-all}
In \cite{CharacterFamily} the Galois input is stated in a stronger-looking
form: clause (1), and the conditional clause of (3), are asserted for
\emph{every} $\mathrm{rec}^T$-map $\mathrm{rc}$ at $w$, not only for
$\mathrm{rec}_{E_w}$.  The formal proof obtains this form from its own
theta-lift construction of $r$.  We have checked how the counts of
\S\ref{i5:ssec-counts} use it, and no step fixes a particular $\mathrm{rc}$.
Every map $\mathrm{rc}$ at which the strong form is applied comes from an
existence statement: the existence of a $\mathrm{rec}^T$-map at $w$
(Proposition~\ref{i5:prop-recT}, which is proved formally with an independent
witness), the existence of a local Langlands correspondence
(Theorem~\ref{i5:thm-llc}), or a choice made from one of these and transported
along isomorphisms of local fields.  Some intermediate statements are
quantified over all $\mathrm{rec}^T$-maps.  Each is derived from the strong form
at one chosen map, sometimes with transfer lemmas that use only the axioms of
Definition~\ref{i5:def-recT} on good classes.  The final statement of this
section mentions no local Langlands map, so every such quantifier is
eventually instantiated at a chosen map.  The map $\mathrm{rec}_{E_w}$
satisfies each of these existence statements (Theorem~\ref{i5:thm-llc},
Proposition~\ref{i5:prop-recT}), and by its uniqueness it is compatible with
isomorphisms of local fields.  If it is chosen every time, the argument applies
the Galois input only with $\mathrm{rc}=\mathrm{rec}_{E_w}$, that is, in the
form of Proposition~\ref{i5:prop-row0b}, which follows from
Theorem~\ref{i5:thm-blggt} (Barnet-Lamb--Gee--Geraghty--Taylor).  The strong
form for arbitrary $\mathrm{rec}^T$-maps is therefore not needed, and we do not
claim it from the literature.  This observation rests on how $\mathrm{rc}$ is
introduced where the formal proof invokes the strong form; we have not
followed every such use.

Clause (3) of Proposition~\ref{i5:prop-row0b} is asserted unconditionally,
which is what Theorem~\ref{i5:thm-blggt} gives.  The formal statement is
weaker: it asserts potential semistability, the jumps and $\mathrm{Away}$ at
$w\mid l$ only when $\mathrm{rc}(\sigma_w)$ is $\Q_p$-big, tame and
square-free, or abelian, together with trivial monodromy in those cases.
\end{remark}

\begin{remark}
Proposition~\ref{i5:prop-row0b} is a strict weakening of full local--global
compatibility: it asserts equality of Weil--Deligne representations only when
$\mathrm{rc}(\pi_w)$ has abelian image and trivial monodromy, and equality of
inertial characteristic polynomials in tame square-free situations over $\Q_g$,
$g>16a^2$.  \cite{CharacterFamily} proves it independently, through a
theta-lift construction of $r$.  In the reduction below it enters only through
the automorphic counts of \S\ref{i5:ssec-counts}.
\end{remark}

\medskip\noindent\label{i5:ssec-coh}We next recall the Galois cohomology that we use.

Let $F$ be a number field, $l$ a prime, $M$ a finite discrete
$\Z/l^r$-module with continuous action, $S\supseteq S_l$ finite,
$M^\vee=\mathrm{Hom}(M,\mu_{l^r})$.

\begin{theorem}[Poitou--Tate duality and local duality; Neukirch--Schmidt--Wingberg, \emph{Cohomology of Number Fields}, Ch.~VII--VIII; Milne, \emph{Arithmetic Duality Theorems}, Ch.~I; NSW Thms.~7.1.8, 7.2.6, 7.2.15, 7.3.1, 8.3.18, 8.3.20, 8.6.7, 8.6.10; Milne, \emph{ADT} I.2.6]\label{i5:thm-pt}
(1) If $l\ne2$ or $F$ is totally imaginary, the groups $H^i(G_{F,S},M)$ are
finite, vanish for $i\ge3$, the nine-term Poitou--Tate sequence is exact, and
$\Sha^1(G_{F,S},M)\cong\Sha^2(G_{F,S},M^\vee)^\vee$.
(2) For a finite place $v$, $H^i(D_v,M)$ is finite, vanishes for $i\ge3$,
$H^2(D_v,\mu_{l^r})\cong\Z/l^r$, and cup product
$H^i(D_v,M)\times H^{2-i}(D_v,M^\vee)\to\Z/l^r$ is perfect.
(3) If $v\nmid l$ and inertia acts trivially on $M$, the unramified subgroups
of $H^1(D_v,M)$ and $H^1(D_v,M^\vee)$ are exact annihilators.
\end{theorem}

The hypotheses are met: all global fields to which (1) is applied are CM or
contain $\mu_l$ with $l$ odd, and (2), (3) are applied to finite modules.  We
also use the Hochschild--Serre five-term sequence and inflation--restriction
for profinite groups (standard; Neukirch--Schmidt--Wingberg, Ch.~I).

\subsection{Reduction to the three bounds}\label{i5:ssec-assembly}

\medskip\noindent\emph{Plus and minus parts.}

\begin{definition}\label{i5:def-pm}
Let $E$ be CM with complex conjugation $c$, $l$ odd, $D$ an $l$-datum.  Put
$\Sha^\pm_D(e)=\Sha^1_D\cap\ker(c\mp(-1)^e)$ and
$\operatorname{sha}^\pm_D(e)$ their lengths.  $\mathrm{VV}^\pm(E,t,B)$ is
Definition~\ref{i5:def-vwin}(1) with $\operatorname{sha}^\pm$ and $l$ odd.
$\mathrm{VV}_l(L,t,B)$ means: for every $l$-datum $D$ on $L$,
$\operatorname{sha}_D(e)<\infty$ for $0<e<t$ and
$\sum_{e<t}\operatorname{sha}_D(e)^\natural\le B$.
\end{definition}

\begin{lemma}\label{i5:lem-split}
For $E$ CM and $l$ odd, $\operatorname{sha}_D(e)=\operatorname{sha}^+_D(e)+\operatorname{sha}^-_D(e)$.
Consequently, $\mathrm{VV}_2(E,t,B_2)$, $\mathrm{VV}^-(E,t,B_-)$ and
$\mathrm{VV}^+(E,t,B_+)$ imply $\mathrm{VV}(E,t,B_2\log2+B_-+B_+)$.
\end{lemma}

\begin{proof}
A lift of $c$ permutes the decomposition groups above $l$ up to conjugacy, so
$c$ preserves $\Sha^1_D$, and $c^2=1$.  As $2$ is a unit, $x\mapsto\frac12(x\pm
s\,cx)$ ($s=(-1)^e$) split $\Sha^1_D$ into the two eigenspaces, and length is
additive.  The second claim follows by splitting $\mathcal S$ into $\{2\}$ and
the odd primes.
\end{proof}

\medskip\noindent\emph{The window statements.}

\begin{definition}\label{i5:def-windows}
\begin{enumerate}
\item $(e,t)$, with $e,t\colon\R\to\N$, is \emph{admissible}
$\mathrm{adm}(A_1,C_t)$ if for $x\ge2$: $e(x)\ge2$, $2^{e(x)}\le x^{A_1}$,
$2\le t(x)\le x^{C_t}$.
\item An \emph{abelian datum} is $M=(m_M,H)$ with $H\le(\Z/m_M)^\times$;
$\deg M=[(\Z/m_M)^\times:H]$; $M_\psi=(N_\psi,\ker\psi)$; $\mathrm{cyc}(m)=(m,1)$.
$(M,e,q,d)$-fields are defined as in Definition~\ref{i5:def-vwin}(2) with
``$a\bmod m_M\in H$'' in place of $\psi(a)=1$; $M$ \emph{has $\mu_l$} if
$l\mid m_M$ and $H\subset\{a\equiv1\bmod l\}$.  $\mathfrak f(E)$ is the
conductor of an abelian field.  An \emph{$l$-layer of height $j$} is
$L'=L(\mu_{l^k})$ for some $k$ with $[L':L]=l^j$.  (When $\mu_l\subset L$,
resp.\ $\mu_4\subset L$ if $l=2$, which holds in every use below, these are
exactly the fields $L\subset L'\subset L(\mu_{l^\infty})$ with $[L':L]=l^j$.)
\item $\mathrm{PP}(m;q,d)$: $q\equiv3\bmod8$, $d\equiv7\bmod16$ primes, $\gcd(qd,m)=1$.
\item A \emph{constant bundle} $K=(R_0,(R_l),C_K)$ consists of functions
$R_0(M),R_l(M)\in\N$ and $C_K>0$; for a finite set $B'$ of primes and abelian
data $L(l)$ put $R=R_0(M)+\sum_{l\in B'}R_l(L(l))$ and
$J_K=2^{R+2}(2C_KC_t/\varepsilon)^R$.
\item $B_a(x)$: every $(M,e(x),q,d)$-field $E$ satisfies
$\mathrm{VV}_2(E,t(x),\exp(x^{4\varepsilon/5}))$.  $B_b(x)$: for every
$l\in B'$, every $L(l)$-field $L_0$ and every $l$-layer $L_1/L_0$ of height
$\lfloor\log_lR_l(L(l))\rfloor+3$, $\mathrm{VV}_l(L_1,t(x),\exp(x^{4\varepsilon/5}))$.
$\mathrm{BWin}(X,Y)$: $X^{1/J_K}\le Y$, $Y^2\le X$, and $B_a(x)$, $B_b(x)$ for
all $x\in[Y,Y^2]$.
\item $\mathrm{ThmBab}(K,C_t)$: for all $M$ and all finite $B',L$ with each
$l\in B'$ odd and $L(l)$ having $\mu_l$, and all $0<\varepsilon<1/52$,
$A_1$, there is $X_0$ such that for all admissible $(e,t)$, $X\ge X_0$ and
$\mathrm{PP}(m_M;q,d)$ with $\log(qd)\le X^{\varepsilon/(4J_K)}$, some $Y$
satisfies $\mathrm{BWin}(X,Y)$.
\item $\beta_C(o,n,N,t)=3nt^2\log(tN)+C\sum_{l\mid o,\ l\text{ odd}}l\,n^2t^2\log(tlN)^2$.
$\mathrm{ThmBc}(C_t)$: there is $C$ such that for all $M$, admissible $(e,t)$,
$x\ge2$, $\mathrm{PP}(m_M;q,d)$ and CM $(M,e(x),q,d)$-fields $E$,
$\mathrm{VV}^-(E,t(x),\beta_C(\deg M,[E:\Q],\mathfrak f(E),t(x)))$.
$\mathrm{ThmBcFree}$: the same for all integers $e,t\ge2$ in place of
$e(x),t(x)$.
\item $\mathrm{HVAt}(\psi,C;x)$: every CM $(\psi,e(x),q,d)$-field $E$ satisfies
\[\mathrm{VV}^+\bigl(E,t(x),(\mathfrak f(E)t(x))^C+C\exp(x^{4\varepsilon/5})\bigr).\]
\item Let $B_\psi$ be as in Definition~\ref{i5:def-Bchi} and
$L_\psi(l)=\mathrm{cyc}(lN_\psi)$, $J_\psi=J_K$ for $(M_\psi,B_\psi,L_\psi)$.
$\mathrm{ThmHV}(K,C_t)$: for every $\psi$ there are $C$ and primes $r_1,r_2$
such that for all $0<\varepsilon<1/52$ and $A_1$ there is $X_0$: for admissible $(e,t)$,
$X\ge X_0$, $\mathrm{PP}(N_\psi;q,d)$ with $2X^{C_t}+1<q,d$, the numbers
$N_\psi,o_\psi,p_K,h^*,p_1,r_1,r_2$ all $<q,d$, $\log(qd)\le X^{\varepsilon/(4J_\psi)}$,
and every $Y$ with $\mathrm{BWin}(X,Y)$: $\mathrm{HVAt}(\psi,C;x)$ for $x\in[Y,Y^2]$.
\item $\mathrm{JointWin}(C_t)$: for all $\psi$, $0<\varepsilon<1/52$, $A_1$ and
admissible $(e,t)$ there are $J>0$, $C$, $P$, $X_0$ such that for $X\ge X_0$ and
$\mathrm{PP}(N_\psi;q,d)$ with $2X^{C_t}+1<q,d$, $P<q,d$,
$\log(qd)\le X^{\varepsilon/(4J)}$, there is $Y$ with $X^{1/J}\le Y$,
$Y^2\le X$ and $B_a(x)\wedge\mathrm{HVAt}(\psi,C;x)$ on $[Y,Y^2]$.
\end{enumerate}
\end{definition}

\begin{definition}[the exceptional set]\label{i5:def-Bchi}
For primitive $\psi$: $p_K$ is the least prime $\equiv3\bmod4$, $\ge7$, not
dividing $N_\psi$; $h^*$ the class number of $\Q(\sqrt{-p_K})$; $p_1$ the least
prime $\equiv1$ modulo the odd part of $o_\psi$ with $p_1\nmid2N_\psi p_K$; and
$B_\psi=\{3\}\cup\{\text{odd primes dividing }o_\psi N_\psi p_Kh^*p_1\}$.
\end{definition}

\medskip\noindent\emph{Elementary lemmas.}

\begin{lemma}\label{i5:lem-cm}
If $q\equiv3\bmod4$ is prime, every $(\psi,e,q,d)$-field is CM.
\end{lemma}
\begin{proof}
$E$ is abelian, so all complex conjugations restrict to $\sigma_{-1}$, which
does not fix $E$ since $(\frac{-1}q)=-1$; the fixed field of $\sigma_{-1}$ is
totally real of index $2$.
\end{proof}

\begin{lemma}\label{i5:lem-cond}
A $(\psi,e,q,d)$-field $E$ satisfies $1\le\mathfrak f(E)\le2^{e+2}N_\psi qd$ and
$[E:\Q]\le\mathfrak f(E)$.
\end{lemma}
\begin{proof}
The defining conditions depend on $a$ modulo
$\mathrm{lcm}(2^{e+2},N_\psi,q,d)$, so $E$ embeds in that cyclotomic field;
and $[E:\Q]\le\varphi(\mathfrak f(E))$.
\end{proof}

\begin{lemma}[asymptotics]\label{i5:lem-asym}
(1) For all $C,o$ there is $K$ (e.g.\ $K=3+|C|\sum_{l\mid o\text{ odd}}l^3$) with
$\beta_C(o,n,N,t)\le Ke^{8W}$ whenever $N,t\ge1$, $n\le N$, $\log(Nt)\le2W$.
(2) For $\varepsilon>0$ and $K,C$ there is $x_0$ such that for $x\ge x_0$ and
$1\le u\le e^{2x^{\varepsilon/4}}$,
$e^{x^{4\varepsilon/5}}\log2+Ke^{8x^{\varepsilon/4}}+u^C+Ce^{x^{4\varepsilon/5}}\le\sqrt{\exp(x^\varepsilon)}$.
(3) For $0<\varepsilon<1/52$, $C_t$, $m$ there is $x_0$ such that for
$x\ge x_0$, $1\le N\le2^{e+2}mqd$, $\mathrm{Page}_\varepsilon(x,e)$,
$1\le t\le x^{C_t}$ and $\log(qd)\le x^{\varepsilon/4}$: $\log(Nt)\le2x^{\varepsilon/4}$.
(4) There is $Y_0$ with $\log(qd)\le Y^{\varepsilon/4}$ for $Y\ge Y_0$ and
$(q,d)\in\mathrm{WP}(\varepsilon,Y)$.
(5) For $\varepsilon,J>0$, $C_t$, $P$ there is $X_0$ such that $X\ge X_0$ and
$(q,d)\in\mathrm{WP}(\varepsilon,X^{1/J})$ imply $2X^{C_t}+1<q,d$ and $P<q,d$.
(6) If $\mathrm{Page}_\varepsilon(x,e(x))$ for $x\ge1$, $C_t\ge1$ and
$t(x)=\lfloor x^{C_t}\rfloor$, then $(e,t)$ is $\mathrm{adm}(4,C_t)$.
\end{lemma}
\begin{proof}
(1) $nt^2\le(Nt)^2\le e^{4W}$, $\log(tN)\le e^{2W}$, and
$\log(tlN)\le l e^{2W}$.  (2) Every term is $\exp(O(x^{\max(4\varepsilon/5,\varepsilon/4)}))$
and $\max(4\varepsilon/5,\varepsilon/4)<\varepsilon$.
(3) $\log(Nt)\le\log(8m)+(2+8\varepsilon+C_t)\log x+x^{\varepsilon/4}$ and
$\log x=o(x^{\varepsilon/4})$.  (4) $\log(qd)\le Y^{\varepsilon/5}(1+(\log Y)^2)$.
(5) $\log q,\log d>\frac12X^{\varepsilon/(5J)}$.  (6) For $x\ge2$,
$2^{e(x)}\ge x^2\ge4$, $2^{e(x)}\le2x^3\le x^4$, and $2\le x\le x^{C_t}$.
\end{proof}

\begin{proposition}\label{i5:prop-vwinat}
For $C_t\ge1$, $\mathrm{JointWin}(C_t)$ and $\mathrm{ThmBc}(C_t)$ imply
$\mathrm{VWinAt}(\psi,C_t)$ for every $\psi$.  Moreover
$\mathrm{ThmBc}(C_t)\iff\mathrm{ThmBcFree}$.
\end{proposition}

\begin{proof}
Fix $\varepsilon$, $e$; put $t(x)=\lfloor x^{C_t}\rfloor$, admissible by
Lemma~\ref{i5:lem-asym}(6).  JointWin gives $J,C,P,X_J$; let $C_b$ be the
constant of ThmBc and $K$ that of Lemma~\ref{i5:lem-asym}(1) for
$(C_b,o_\psi)$, and take thresholds $x_T,x_N,Y_0$ from (2), (3), (4).  Put
$J'=\max(J,1)$, $P'=\max(P,N_\psi)$, $X_B$ from (5),
$L=\max(x_T,x_N,Y_0,2)$, $X_0=\max(X_J,X_B,L^{J'})$.  For $X\ge X_0$ and
$(q,d)\in\mathrm{WP}(\varepsilon,X^{1/J'})$: by (5), $\mathrm{PP}(N_\psi;q,d)$
and the size conditions hold, and by (4) $\log(qd)\le X^{\varepsilon/(4J)}$.
JointWin yields $Y$.  For $x\in[Y,Y^2]$ and a $(\psi,e(x),q,d)$-field $E$
(CM by Lemma~\ref{i5:lem-cm}), $B_a$, HVAt and ThmBc give the three bounds of
Lemma~\ref{i5:lem-split}; by Lemma~\ref{i5:lem-cond} and (3),
$\log(\mathfrak f(E)t(x))\le2x^{\varepsilon/4}$, so by (1) the minus bound is
$\le Ke^{8x^{\varepsilon/4}}$, and (2) with $u=\mathfrak f(E)t(x)$ bounds the
total by $\sqrt{\exp(x^\varepsilon)}$.

For the equivalence: ($\Leftarrow$) admissibility gives $e(x),t(x)\ge2$.
($\Rightarrow$) given $e,t\ge2$, put $x=\max(2^e,t)$ and use the functions equal
to $e,t$ at $x$ and to $2$ elsewhere, which are $\mathrm{adm}(2,C_t)$.
\end{proof}

\begin{proposition}\label{i5:prop-jointwin}
If $\mathrm{ThmBab}(K,C_t)$ and $\mathrm{ThmHV}(K,C_t)$ then
$\mathrm{JointWin}(C_t)$.
\end{proposition}
\begin{proof}
Each $l\in B_\psi$ is odd and $\mathrm{cyc}(lN_\psi)$ has $\mu_l$, so ThmBab
applies to $(M_\psi,B_\psi,L_\psi)$.  Take $J=J_\psi$, $C$ from ThmHV,
$P=N_\psi+o_\psi+p_K+h^*+p_1+r_1+r_2$; ThmBab gives $Y$ with BWin, and ThmHV
applies to the same $Y$.
\end{proof}

\medskip\noindent\emph{Hosts.}

\begin{definition}[hosts]\label{i5:def-host}
A \emph{host} $h=(\psi,e,q_1,d)$ is a primitive $\psi$, $e\ge2$ and
$\mathrm{PP}(N_\psi;q_1,d)$.  Put $q=N_\psi$, $o=o_\psi$,
$N_h=\mathrm{lcm}(2^{e+2},q)q_1d$ and
$H_h=\{a\in(\Z/N_h)^\times: a\equiv\pm1\ (2^{e+2}),\ \psi(a)=1,\ a\text{ a square mod }q_1\text{ and mod }d\}$.
$\mathrm{Big}_h(t)$ means $2t+1<q_1,d$ and each of $q,o,p_K,h^*,p_1$ is
$<q_1$ and $<d$.  $X_h(\mathcal O,e)$ is the set of characters
$\psi'\colon(\Z/N_h)^\times\to\mathcal O^\times$ trivial on $H_h$ with
$\psi'(-1)=(-1)^e$.  A number field $E$ \emph{has level $(N,H)$} if it embeds
in $\Q(\zeta_N)$ as the fixed field of $\{\sigma_u:u\in H\}$.
\end{definition}

\begin{lemma}\label{i5:lem-level}
(i) A subfield of $\Q(\zeta_N)$ has level $(N,H)$ for a unique $H$.  (ii) If
$M\mid N$, $E$ has level $(M,H)$ iff it has level $(N,\pi^{-1}H)$, $\pi$ the
reduction.  (iii) A $(\psi,e,q_1,d)$-field has level $(N_h,H_h)$.
\end{lemma}
\begin{proof}
(i) and (ii) are Galois theory of $\Q(\zeta_N)$ (all embeddings of the abelian
field $E$ have the same image).  (iii) The four conditions depend only on
$a\bmod N_h$ ($q_1,d$ prime, so $(\frac ap)=1$ iff $a$ is a square), and
$N_h$ divides the $N$ of the definition; apply (ii).
\end{proof}

\begin{lemma}[conductor]\label{i5:lem-hostcond}
For a host $h$ and a $(\psi,e,q_1,d)$-field $E$, $\mathfrak f(E)=N_h$.
\end{lemma}
\begin{proof}
$\mathfrak f(E)\le N_h$ by Lemma~\ref{i5:lem-level}.  Conversely, let
$E\hookrightarrow\Q(\zeta_M)$ and suppose $v_p(M)<v_p(N_h)$ for some prime $p$.
We first find $c$ prime to $N_h$, $c\equiv1\bmod N_h/p$, with $c\notin H_h$.
By the Chinese remainder theorem it suffices to find $R\mid N_h$ with
$p^{v_p(N_h)}\mid R$ and $w\equiv1\bmod R/p$, prime to $p$, such that every
$c\equiv w\bmod R$ violates a condition of $H_h$: for $p\in\{q_1,d\}$ take
$R=p$, $w=-1$ (not a square, as $p\equiv3\bmod4$); for other odd $p$, or $p=2$
with $v_2(N_h)\le v_2(q)$, take $R=q$ and $w\equiv1\bmod q/p$ with $\psi(w)\ne1$,
which exists because $\psi$ is primitive; otherwise $p=2$, $v_2(N_h)=e+2$, and
$R=2^{e+2}$, $w=1+2^{e+1}\not\equiv\pm1$ (as $e\ge2$).  Now put
$L=\mathrm{lcm}(M,N_h)$.  By Lemma~\ref{i5:lem-level}, $E$ has level
$(L,\pi_M^{-1}H')=(L,\pi_{N_h}^{-1}H_h)$.  The Chinese remainder theorem gives
$v$ prime to $L$ with $v\equiv c\bmod N_h$ and $v\equiv1\bmod L/p$; as
$M\mid L/p$, $v\in\pi_M^{-1}H'$, hence $c\in H_h$, a contradiction.  So
$N_h\mid M$.
\end{proof}

\begin{definition}[character pieces]\label{i5:def-V}
For $\psi'$ an $\mathcal O$-valued Dirichlet character of level $N$, viewed as
a character of $G$, and $e\ge0$, let
$\rho_{e,\psi'}=K_{\mathcal O}/\mathcal O$ with $G$ acting by
$\psi'^{-1}\chi^{-e}$.  For $n\ge1$ let $\mathcal D(n)$ be the family of the
decomposition groups $D_w$, $w\mid n$, and of all inertia groups $I_w$; for a
closed $U\le G$ and a $G$-module $A$,
$\mathrm{Sel}_U(A;\mathcal D)=\{x\in H^1(U,A):x|_{U\cap D'}=0\ \forall D'\in\mathcal D\}$.
Put $V_l(e,\psi')=\lgth_{\mathcal O}\mathrm{Sel}_G(\rho_{e,\psi'};\mathcal D(l))$,
$V'_{l,N}(e,\psi')=\lgth_{\mathcal O}\mathrm{Sel}_G(\rho_{e,\psi'};\mathcal D(lN))\le V_l(e,\psi')$,
$\mathrm{Vn}(e,\psi')=V'^\natural_{l,N}/e_{\mathcal O}$, and
$\mathrm{VLe}(e,\psi',B)$: $V'<\infty$ and $\mathrm{Vn}\log l\le B$.
\end{definition}

\begin{definition}[plus-part statements]\label{i5:def-pointhv}
$\mathrm{LayerLe}(\psi,l,j,t,B)$: for every $\mathrm{cyc}(lN_\psi)$-field
$L_0$ and $l$-layer $L_1/L_0$ of height $j$, $\mathrm{VV}_l(L_1,t,B)$.
$\mathrm{PointHV}$: for every $\psi$ and $j\colon\N\to\N$ there are $C$ and
primes $r_1,r_2$ such that for every host $h$ with character $\psi$ and
$r_1,r_2<q_1,d$, every $t\ge2$ with $\mathrm{Big}_h(t)$, every $B\ge0$ with
$\mathrm{LayerLe}(\psi,l,j(l),t,B)$ for all $l\in B_\psi$, and every CM
$(\psi,e,q_1,d)$-field $E$: $\mathrm{VV}^+(E,t,(N_ht)^C+CB)$.
$\mathrm{SumIn}$ (resp.\ $\mathrm{SumOut}$) is the same statement restricted
to one prime $l\in B_\psi$ (resp.\ to sets of odd primes outside $B_\psi$, with
bound $(N_ht)^C$).
\end{definition}

\begin{proposition}\label{i5:prop-point}
$\mathrm{SumIn}\wedge\mathrm{SumOut}\Rightarrow\mathrm{PointHV}\Rightarrow\mathrm{ThmHV}(K,C_t)$
for every $K$ and $C_t>0$.
\end{proposition}
\begin{proof}
First implication: split a set of odd primes into its parts in and outside
$B_\psi$; with $M=N_ht\ge2$, the constants add up, using $M^k\ge k+1$.
Second: apply PointHV with $j(l)=\lfloor\log_lR_l(L_\psi(l))\rfloor+3$ and
$X_0=2^{J_\psi}$.  For $x\in[Y,Y^2]$, $x\ge2$, admissibility gives
$e(x),t(x)\ge2$ and $t(x)\le X^{C_t}$, so $\mathrm{Big}_h(t(x))$ holds for the
host $h=(\psi,e(x),q,d)$; $B_b$ is exactly LayerLe with
$B=\exp(x^{4\varepsilon/5})$; and $N_h=\mathfrak f(E)$ by
Lemma~\ref{i5:lem-hostcond}.
\end{proof}

\subsection{Reduction to characters over $\Q$}\label{i5:ssec-red}

For odd $l$ and $N,e\ge1$ put $\nu(l,N,e)=\max(v_l(N),1)+v_l(e)$ if $l\mid N$
or $(l-1)\mid e$, and $0$ otherwise; $\omega(N)$ the number of prime factors of
$N$; $\mathrm{loss}(l,N,n,e)=(n+n^2)\nu+\omega(N)nv_l(n)$ and
$\mathrm{lossAll}=\mathrm{loss}+\omega(N)v_l(n)$.  All modules are torsion
$\mathcal O$-modules whose cyclic submodules have finite length.

\begin{lemma}\label{i5:lem-lens}
(1) $\lgth(Z/aZ)\le\lgth Z[a]$ for $a\in\mathcal O$.
(2) If a finite abelian group $Q$ acts on $P$ and $\mathcal O$ contains the
$\exp(Q)$-th roots of unity, $\lgth P\le\sum_\phi\lgth P^{(\phi)}$ over the
characters $\phi\colon Q\to\mathcal O^\times$.
(3) Let $U'\trianglelefteq U$ be closed, $U'$ open in $U$, $[U,U]\subset U'$, and
let $U$ act on $A$ by scalars $c(u)$.  If $U$ is topologically generated by
$U'$ and a finite set $T$, then
$\lgth\ker(H^1(U,A)\to H^1(U',A))\le|T|\lgth H^0(U,A)$ and
$\lgth\operatorname{coker}(H^1(U,A)\to H^1(U',A)^U)\le[U:U']^2\lgth H^0(U,A)$.
If moreover $U'$ acts trivially and $[U:U']\mid n$, the kernel has length
$\le n\,v_l(n)e_{\mathcal O}$.
\end{lemma}
\begin{proof}
(1) Reduce to finitely generated submodules, of finite length, where kernel
and cokernel of multiplication by $a$ have equal length.  (2) For commuting
$A,B$, $\lgth\ker AB\le\lgth\ker A+\lgth\ker B$; apply this to
$\prod_j(g_i-\zeta_i^j)=0$ for generators $g_i$ of $Q$.  (3) Kernel: classes in
the kernel are represented by cocycles $z$ vanishing on $U'$, determined by
their values on $T$; since $[u,g]\in U'$ acts trivially,
$(c(u)-1)z(g)=(c(g)-1)z(u)$, so $(c(u)-1)z$ is a coboundary.  With $a$ a
generator of the ideal $(c(u)-1)$, (1) bounds the kernel by
$\lgth Z_0[a]\le|T|\lgth A[a]$ and $A[a]=H^0(U,A)$.  Cokernel: it injects into
$H^2(U/U',A^{U'})$ (five-term sequence), which is killed by $a$ (the abelian
quotient acts trivially on its own cohomology and by scalars), and $2$-cocycles
are determined by $[U:U']^2$ values.  The last claim: the kernel is
$H^1(U/U',A)$, killed by $l^{v_l(n)}$, and its cocycles inject into
$A[l^{v_l(n)}]^{|U/U'|}$.
\end{proof}

\begin{lemma}\label{i5:lem-local}
Let $l$ be odd, $e\ge1$, $\psi'$ of level $N$.  (1) For $w\mid l$ and every
subgroup $U\supseteq I_w$, and for $U=G$,
$\lgth H^0(U,\rho_{e,\psi'})\le\nu(l,N,e)e_{\mathcal O}$; and $\nu\le2Ne$.
(2) If $\psi'^n=1$, $p\ne l$, $w\mid p$, then
$\lgth\ker(H^1(D_w,\rho)\to H^1(I_w,\rho))\le v_l(n)e_{\mathcal O}$, and the map
is injective if $\psi'(I_w)=1$.
\end{lemma}
\begin{proof}
(1) Pass to $U'=I_w\cap\ker(G\to(\Z/N)^\times)$, on which $\psi'=1$ and
$\chi(U')=1+l^k\Z_l$ with $k=\max(v_l(N),1)$ if $l\mid N$, $\chi(U')=\Z_l^\times$
otherwise ($\Q_l(\zeta_{l^n})/\Q_l$ is totally ramified, $\Q_l(\zeta_{N'})$
unramified for $l\nmid N'$).  The invariants have length
$\min_uv_{\mathcal O}(\chi(u)^e-1)$: zero if $l\nmid N$, $(l-1)\nmid e$ (take $u$ a
root of unity of order $l-1$), and $k+v_l(e)$ otherwise, by
$v_l((1+x)^m-1)=v_l(x)+v_l(m)$ for $l$ odd, $v_l(x)\ge1$.
(2) The kernel is $\rho^{I_w}/(\varphi-1)$.  If $\psi'(I_w)=1$, $\varphi$ acts
by $\psi'(\varphi)^{-1}p^{-e}\ne1$ on the divisible module $\rho$, so the
quotient is $0$.  Otherwise $\rho^{I_w}\subset\rho[\zeta-1]$ for a root of unity
$\zeta\ne1$ of order dividing $n$, of length $\le e_{\mathcal O}\le v_l(n)e_{\mathcal O}$
if the order is an $l$-power and $0$ otherwise.
\end{proof}

\begin{proposition}[character decomposition]\label{i5:prop-red}
Let $E$ be CM of level $(N,H)$, $n=[(\Z/N)^\times:H]$, $l$ odd, $\mathcal O$
containing the $\varphi(N)$-th roots of unity, $D$ an $l$-datum on $E$,
$e\ge1$, and $T$ the set of characters of $(\Z/N)^\times$ trivial on $H$ with
$\psi'(-1)=(-1)^e$.  Then
\[
 e_{\mathcal O}\operatorname{sha}^+_D(e)\le\sum_{\psi'\in T}\bigl(V'_{l,N}(e,\psi')+e_{\mathcal O}\mathrm{lossAll}(l,N,n,e)\bigr),
\]
without the $\mathrm{lossAll}$ terms if $l\nmid n$; and $\mathrm{lossAll}\le(Ne)^6$
if $2\le N$, $n\le N$.
\end{proposition}

\begin{proof}
(a) \emph{Model over $\Q$.}  Conjugation by an extension of an isomorphism of
fields transports $l$-data, inertia, the cyclotomic character and complex
conjugation, so $\operatorname{sha}^\pm_D$ is independent of the choices.  By
inflation--restriction ($N_S$ acts trivially and its classes vanish exactly when
they vanish on the generating inertia groups, a closed normal condition),
$\Sha^1_D=\mathrm{Sel}_{G_E}(\Q_l/\Z_l(-e);\mathcal D(l))$, compatibly with a lift
$c\in G$ of complex conjugation; tensoring with $\mathcal O$ multiplies lengths
by $e_{\mathcal O}$.  Hence
$e_{\mathcal O}\operatorname{sha}^+_D(e)=\lgth X^{c=(-1)^e}$ with
$X=\mathrm{Sel}_{G_E}(M_e;\mathcal D(l))$, $M_e=\rho_{e,\mathbf1}$.

(b) \emph{Sign to characters.}  $G$ acts on $X$ through
$Q=(\Z/N)^\times/H$, $c\mapsto-1$.  By Lemma~\ref{i5:lem-lens}(2),
$\lgth X^{c=\epsilon}\le\sum_\phi\lgth(X^{c=\epsilon})^{(\phi)}$, and the term
vanishes unless $\phi(-1)=\epsilon$, since $\pm2$ is a unit.  So
$e_{\mathcal O}\operatorname{sha}^+\le\sum_{\psi'\in T}\lgth X^{(\psi')}$.

(c) \emph{Case $l\nmid n$.}  On $U=G_E$, $M_e=\rho_{e,\psi'}$ and
$X^{(\psi')}=\mathrm{Sel}_U(\rho;\mathcal D(l))^G$.  Restriction from $G$ is
bijective onto it: $\mathrm{cor}\circ\mathrm{res}=n$ and
$\mathrm{res}\circ\mathrm{cor}=n$ on invariants, $n$ is a unit, and
$n^{-1}\mathrm{cor}$ preserves the local conditions because
$\ker(H^1(D')\to H^1(D'\cap U))$ is killed by a divisor of $n$.  So
$\lgth X^{(\psi')}=V_l(e,\psi')$.

(d) \emph{Places dividing $N$.}  For $p\mid N$, $p\ne l$, fix $w_p\mid p$.  The
map $x\mapsto(x|_{D_{w_p}})_p$ lands in
$\prod_p\ker(H^1(D_{w_p})\to H^1(I_{w_p}))$ and its kernel lies in
$\mathrm{Sel}_G(\rho;\mathcal D(lN))$ (conjugation acts trivially on
$H^1(G,-)$).  By Lemma~\ref{i5:lem-local}(2),
$V_l\le V'_{l,N}+\omega(N)v_l(n)e_{\mathcal O}$, with equality of $V_l$ and
$V'_{l,N}$ if $l\nmid n$.

(e) \emph{General case.}  Replace (c) by
$\lgth X^{(\psi')}\le V_l+\mathrm{loss}\cdot e_{\mathcal O}$: with $Y$ the preimage
of $\mathrm{Sel}_{G_E}(\rho)^G$ in $H^1(G,\rho)$, the map
$y\mapsto(y|_{D_{w_0}},(y|_{I_{w_p}})_p)$ ($w_0\mid l$) has kernel in
$\mathrm{Sel}_G(\rho;\mathcal D(l))$ (inertia at $w\nmid lN$ lies in $G_E$), so
$\lgth X^{(\psi')}$ is at most $V_l$ plus the lengths of the local restriction
kernels and of $\operatorname{coker}(H^1(G,\rho)\to H^1(G_E,\rho)^G)$.
Lemmas~\ref{i5:lem-lens}(3) and \ref{i5:lem-local}(1) bound these by
$n\nu e_{\mathcal O}$, $nv_l(n)e_{\mathcal O}$ each, and $n^2\nu e_{\mathcal O}$.
Finally, with $M=Ne\ge2$: $\nu\le2M$, $\omega(N),n,v_l(n)\le M$, so
$\mathrm{lossAll}\le6M^3\le M^6$.
\end{proof}

\subsection{The prime-by-prime cases}\label{i5:ssec-cases}

\begin{definition}[components and cases]\label{i5:def-cases}
For a character $\psi'$ of level $N$ and a prime $r$, $\mathrm{comp}_r\psi'$ is
$\psi'$ composed with the projection of $(\Z/N)^\times$ onto its $r$-primary
factor.  Write $\psi'=\theta_0\kappa$ with $\theta_0$ of order prime to $l$ and
$\kappa$ of $l$-power order, and $\chi'_l=\theta_0(\mathrm{comp}_l\theta_0)^{-1}$.
Let $c_2(\psi')$ be the exponent of the conductor of $\mathrm{comp}_2\psi'$,
$\theta_1=\theta_0(\mathrm{comp}_2\theta_0)^{-1}$, $f_0$ the order of $l$
modulo $\mathrm{lcm}(4,\operatorname{ord}\theta_1)$, $n_0=v_2(l^{f_0}-1)$ and
$R_1=16o(lq)^24^{n_0}$.  For a host $h$, odd $l$, $e\ge1$ and
$\psi'\in X_h(\mathcal O,e)$:
Case~I: $l\notin B_\psi$, $(l-1)\nmid2e$.
Case~Z: $l\notin B_\psi$, $(l-1)\mid2e$, $\psi'=\mathbf1$.
Case~II: $e\ge2$, $\operatorname{ord}\chi'_l$ even, and $l\in B_\psi$ or $(l-1)\mid2e$.
Case~III: $l\notin B_\psi$, $(l-1)\mid2e$, $\operatorname{ord}\psi'$ odd, $\psi'\ne\mathbf1$.
Case~IV: $l\in B_\psi$, $e\ge2$, $\operatorname{ord}\chi'_l$ odd.
Case~V: $l\in B_\psi$, $e=1$; subcase V if $R_1\le2^{c_2}$, Vs if $R_1>2^{c_2}$.
\end{definition}

\begin{lemma}[host arithmetic]\label{i5:lem-hostarith}
Let $h$ be a host and $\psi'$ trivial on $H_h$.
(1) The primes dividing $N_h$ are $2$, those of $q$, $q_1$, $d$;
$v_{q_1}(N_h)=v_d(N_h)=1$; $16\mid N_h$; $[(\Z/N_h)^\times:H_h]\mid2^{e+2}o$;
$\operatorname{ord}\psi'\mid2^{e+1}o$.
(2) For odd $r$: $\operatorname{ord}\mathrm{comp}_r\psi'\mid2o$; it divides $o$ if
$r\mid q$; $(\mathrm{comp}_r\psi')^2=1$ if $r\nmid q$.
(3) If $\mathrm{Big}_h(t)$: every $l\in B_\psi$ is $<q_1,d$; $p_K,p_1\nmid N_h$;
every odd $l\notin B_\psi$ has $l\ge5$, $l\nmid oqh^*[(\Z/N_h)^\times:H_h]$,
$l\ne p_K,p_1$, and if $0<e<t$, $(l-1)\mid2e$ then $e\ge2$, $l\le2t$, $l\nmid N_h$.
(4) (Exhaustion) If $\mathrm{Big}_h(t)$ and $0<e<t$: for $l\notin B_\psi$,
$\kappa=1$, $\chi'_l=\psi'$ when $(l-1)\mid2e$, and one of I, Z, II, III holds;
for $l\in B_\psi$ one of II, IV, V holds.
(5) If $l\ne q_1,d$ is odd and $\operatorname{ord}\chi'_l$ is odd, the
conductor of $\psi'$ divides $q$.
\end{lemma}

\begin{proof}
(1) $H_h$ is the kernel of a map to a group of order $2^e\cdot o\cdot2\cdot2$;
$u^{2^{e+1}o}$ lies in $H_h$ for every $u$.  (2) With $v$ the $r$-projection of
$u$: if $r\nmid q$, $v^2\in H_h$; if $r\mid q$, $v^o\in H_h$.  (3) Elements of
$B_\psi$ divide one of $o,q,p_K,h^*,p_1$ or equal $3<7\le p_K$, all $<q_1,d$;
$p_K,p_1$ are odd, prime to $q$ and $<q_1,d$, so do not divide $N_h$.  An odd
$l\notin B_\psi$ is $\ne3$ and divides none of the five numbers, nor the index by
(1); if $(l-1)\mid2e$ then $4\le l-1\le2e$ and $l\le2t-1<q_1,d$.
(4) For $l\notin B_\psi$, $l\nmid\operatorname{ord}\psi'$ by (1),(3); if
$(l-1)\mid2e$, $l\nmid N_h$ so $\mathrm{comp}_l\psi'=1$.  The cases are then a
case distinction on $(l-1)\mid2e$, $\psi'=1$ and the parity of the order; for
$l\in B_\psi$ on $e=1$ and the parity of $\operatorname{ord}\chi'_l$.
(5) $\mathrm{comp}_2\psi'=\mathrm{comp}_2\chi'_l$ has odd and $2$-power order;
for $p\in\{q_1,d\}$, $\mathrm{comp}_p\psi'$ has odd order and square $1$.  So the
components at $2,q_1,d$ are trivial, and $v_p(N_h)=v_p(q)$ elsewhere.
\end{proof}

\begin{definition}[counts]\label{i5:def-counts}
For a case $P$, $\mathrm{Count}_P$ means: for every $\psi$ there is $C$ such
that for every host $h$ with character $\psi$, $t$ with $\mathrm{Big}_h(t)$, odd
$l$, $\mathcal O$, $0<e<t$ and $\psi'\in X_h(\mathcal O,e)$ in Case $P$:
$\mathrm{VLe}(e,\psi',(lN_he)^C)$.  This defines $\mathrm{Count}_{\rm II}$,
$\mathrm{Count}_{\rm III}$, $\mathrm{Count}_{\rm V}$, $\mathrm{Count}_{\rm Vs}$,
$\mathrm{Count}_{\rm Z}$.  $\mathrm{Count}_{\rm I}$: for every $\psi$ there are
$C$ and primes $r_1,r_2$ such that for hosts with $r_1,r_2<q_1,d$,
$\mathrm{Big}_h(t)$, $0<e<t$, coefficient rings $(\mathcal O_l)$ and finite sets
$\mathcal L$ of odd primes in Case~I, all $V'<\infty$ and
$\sum_{l\in\mathcal L}\log l\sum_{\psi'\in X_h(\mathcal O_l,e)}\mathrm{Vn}(e,\psi')\le(N_ht)^C$.
$\mathrm{Count}_{\rm IV}$: for $\psi$, $l\in B_\psi$, $j$ there is $C$ such that
for hosts, $t\ge2$ with $\mathrm{Big}_h(t)$, $B$ with
$\mathrm{LayerLe}(\psi,l,j,t,B)$ and $\psi'$ trivial on $H_h$ with
$\operatorname{ord}\chi'_l$ odd: $V'(e,\psi')<\infty$ for $0<e<t$ and
$\sum_{e<t}\mathrm{Vn}(e,\psi')\le B+(N_ht)^C$.
\end{definition}

\begin{proposition}\label{i5:prop-sum}
(1) $\mathrm{Count}_{\rm I}\wedge\mathrm{Count}_{\rm Z}\wedge\mathrm{Count}_{\rm II}\wedge\mathrm{Count}_{\rm III}\Rightarrow\mathrm{SumOut}$.
(2) $\mathrm{Count}_{\rm II}\wedge\mathrm{Count}_{\rm V}\wedge\mathrm{Count}_{\rm Vs}\wedge\mathrm{Count}_{\rm IV}\Rightarrow\mathrm{SumIn}$.
\end{proposition}

\begin{proof}
Put $M=N_ht$, and choose $\mathcal O_l$ with the $\varphi(N_h)$-th roots of unity.
(1) For $l\notin B_\psi$, $l\nmid[(\Z/N_h)^\times:H_h]$
(Lemma~\ref{i5:lem-hostarith}(3)), so Proposition~\ref{i5:prop-red} (first form,
with $T=X_h$) gives
$\sum_e\operatorname{sha}^+(e)^\natural\le\sum_e\sum_{\psi'}\mathrm{Vn}(e,\psi')$.
For fixed $e$, the Case-I primes contribute $\le M^{C_1}$; the others satisfy
$(l-1)\mid2e$, so $l\le2t$, $lN_he\le M^3$, and each of the at most $2t$ such
primes contributes $\le N_hM^{3K}$ ($K$ the largest of the other constants).
Summing over $e<t\le M$ gives $\le M^{C_1+3K+4}$.
(2) Fix $l\in B_\psi$; $M\ge3$, $\log l\le M$.  In cases II, V, Vs,
$\mathrm{Vn}\log l\le M^{2C_m}$.  In Case IV, $\psi'$ has conductor dividing
$q$ (Lemma~\ref{i5:lem-hostarith}(5)), so there are at most $q$ such
characters, each with $\sum_e\mathrm{Vn}\le B+M^{C_4}$.  The second form of
Proposition~\ref{i5:prop-red} adds $N_htM^6$.  Altogether
$\log l\cdot q(B+M^{C_4})+tN_hM^{2C_m}+\log l\cdot N_htM^6\le M^C+CB$ with
$C=C_4+2C_m+9+lN_\psi$.
\end{proof}

\medskip\noindent\emph{Case Z.}

\begin{lemma}\label{i5:lem-kw}
(1) Let $l$ be odd and $\phi\colon G\to\Z/l$ continuous with $\phi(I_w)=0$ for
$w\nmid l$.  Then $\phi$ vanishes on $G_{\Q(\mu_{l^2})}$.
(2) (Control) Let $U\le G$ be open normal, $\theta=\psi'\chi^n$ with $\psi'$
trivial on $U$ and $n\ne0$, $U_\infty=U\cap\ker\chi$.  Then
$\mathrm{Sel}_U(\rho(\theta);\mathcal D(l))\to H^1(U_\infty,\rho(\theta))$ is injective.
\end{lemma}
\begin{proof}
(1) By Kronecker--Weber (Neukirch V (1.10)) and the ramification
condition, $\phi$ factors through the cyclic group $(\Z/l^k)^\times$; the kernel
of reduction to $(\Z/l^2)^\times$ consists of $l$-th powers.  (2) It suffices
that $H^1(U/U_\infty,\rho)=0$.  $U/U_\infty\cong\Delta\times\Z_l$ with $\Delta$ of
order prime to $l$; if $\theta(\Delta)\ne1$ the invariants vanish; otherwise
$H^1=\rho/(\theta(\gamma)-1)\rho=0$, since $\theta(\gamma)=\chi(\gamma)^n\ne1$ and
$\rho$ is divisible.
\end{proof}

\begin{lemma}[normalised cocycles]\label{i5:lem-sht}
Let $\alpha\colon G\to\mathcal O^\times$, $\delta\in G$ with
$\alpha(\delta)-1$ a unit, and $F\ge1$ with $\alpha(I_w)=1$ for $w\nmid F$ and,
for $w\mid F$, $w\nmid l$, some $g\in I_w$ with $\alpha(g)-1$ a unit.  For
$a\ne0$ let $\mathrm{ShT}(a)$ be the cocycles $c\colon G\to\rho(\alpha)$ with
$c(\delta)=0$, $ac=0$, $c|_{I_w}=0$ for $w\nmid F$, and $c|_{D_w}=dm$ with
$am=0$ for $w\mid F$.  If $\lgth\mathrm{ShT}(a)\le B$ for all $a$, then
$\lgth\mathrm{Sel}_G(\rho(\alpha);\mathcal D(l))\le B$.
\end{lemma}
\begin{proof}
Each class has a unique representative with $c(\delta)=0$.  A Selmer class has
finite range, so for suitable $a$ its normalised representative lies in
$\mathrm{ShT}(a)$ (at $w\mid F$, $w\nmid l$, one solves for $m$ using the unit
$\alpha(g)-1$, and at $w\mid l$ uses the cocycle relation to move between
conjugate places).  So the Selmer group is an increasing union of images of
the $\mathrm{ShT}(\varpi^k)$.
\end{proof}

\begin{proposition}[Case Z]\label{i5:prop-countZ}
If $l\ge5$, $e\ge2$ is even and $(l-1)\mid2e$, then $V_l(e,\mathbf1)=0$.
Hence $\mathrm{Count}_{\rm Z}$ holds with $C=0$.
\end{proposition}
\begin{proof}
In Case Z, $l\ge5$ (since $3\in B_\psi$) and parity forces $e$ even.  The
Selmer group $\mathrm{Sel}$ of $A=\rho_{e,\mathbf1}$ is $l$-primary, so it
suffices that $\mathrm{Sel}[l]=0$.

If $(l-1)\mid e$, $G$ acts trivially on $A[l]$, so an $l$-torsion class is
represented by a homomorphism $G\to A[l]$ vanishing on inertia away from $l$;
by Lemma~\ref{i5:lem-kw}(1) it vanishes on $\ker\chi$, hence is $0$ by
Lemma~\ref{i5:lem-kw}(2).

Otherwise $l=2k+1$ with $k$ even and $e/k$ odd, so $l\equiv1\bmod4$.  Let
$K=\Q(\sqrt l)$; $G_K$ acts trivially on $A[l]$ (Euler's criterion).  Pick
$\delta$ in inertia at $l$ with $\chi(\delta)$ a generator mod $l$, so
$\alpha(\delta)-1$ is a unit for $\alpha=\chi^{-e}$ and $\delta\notin G_K$.  We
apply Lemma~\ref{i5:lem-sht} with $F=l$, $B=0$.  For $c\in\mathrm{ShT}(a)$ with
$lc=0$, $c|_{G_K}$ is a homomorphism to $A[l]$, unramified at all finite places
(at $w\mid l$ one shows $lm=0$ using an inertia element with $\chi=$ generator).
By class field theory it factors through the narrow class group of $K$, of order
dividing $2h_K$ (Hilbert class field, Neukirch VI (6.9)), and
$h_K<l$ by the Minkowski bound and a divisor-sum estimate
($h_K\le\sqrt l(1+\log\sqrt l)<l$; Minkowski: Neukirch I (6.2) and the proof of I (6.3)).  So
$c|_{G_K}=0$, and $c(g)=c(\delta)+\alpha(\delta)c(\delta^{-1}g)=0$ for
$g\notin G_K$.
\end{proof}

\subsection{The Poitou--Tate step and the effective bound for $H^2$}\label{i5:ssec-pt}

\begin{proposition}[Poitou--Tate step]\label{i5:prop-shaleV}
If $l\ne2$ or $F$ has no real place, then for every $l$-datum $D$ and $e\in\Z$,
$\lgth_{\Z_l}\Sha^2_D(F,\Z_l(e+1))=\operatorname{sha}_D(e)$; hence
$\operatorname{sha}_D(e)\le V_D(e)$, and
$V_D(e)\le\operatorname{sha}_D(e)+\sum_{v\in S}\lgth H^0(D_v,\Q_l/\Z_l(-e))$.
\end{proposition}

\begin{proof}
Put $\Gamma=G_{F,S}$, $T=\Z_l(e+1)$, $T_r=T/l^r$, $A=\Q_l/\Z_l(-e)$, so
$\mathrm{Hom}(T_r,\mu_{l^r})=A[l^r]$.

\emph{Step 0: passage to $\Z_l$-coefficients.}  Let $\Gamma'$ be $\Gamma$ or a
$D_v$.  By Theorem~\ref{i5:thm-pt} the $H^i(\Gamma',T_r)$ are finite and
$H^3(\Gamma',T_1)=0$.  Then $H^2(\Gamma',T)$ is finitely generated and
$H^2(\Gamma',T)\to H^2(\Gamma',T_r)$ is surjective with kernel $l^rH^2$.
Indeed: a cocycle reducing to a coboundary is $l^r$ times a cocycle ($T$
torsion-free), which gives the kernel; a lift $y$ of a cocycle mod $l^r$ has
$dy=l^rc$, and $H^3(T_1)=0$ lets one write $c=d(\sum l^ib_i)$ by successive
approximation, which gives surjectivity; for finite generation, $X/lX$ is
finite and $\bigcap l^rX=0$, the latter because a cocycle in
$l^rZ^2+B^2$ for all $r$ is a coboundary (choose compatible primitives at
each level, which is possible since the sets of primitives modulo coboundaries
are finite non-empty torsors under $H^1(\Gamma',T_r)$, and an inverse limit of
non-empty finite sets is non-empty); then a separated $\Z_l$-module with finite
$X/lX$ is finitely generated (embed it in its completion and use Nakayama).

\emph{Step 1.}  Let $f\colon X=H^2(\Gamma,T)\to P=\prod_vH^2(D_v,T)$ and $f_r$
its analogue mod $l^r$.  For $r$ large we claim: (i) $f(x)\in l^rP$ implies
$f(x)\in l^rf(X)$, and (ii) if $a\in A$ has $da|_{D_v}=db_v|_{D_v}$ with
$b_v\in A[l^r]$ for all $v$, then $da=db$ with $b\in A[l^r]$.  If $e\ne0$,
$\chi(D_v)$ is open, so $H^0(D_v,A)$ is killed by some $l^{c_0}$, and by local
duality $H^2(D_v,T_r)$ has order $\le l^{c_0}$; hence $P$ is finite, killed by
some $l^m$, and (i), (ii) hold for $r\ge c_0+m$.  If $e=0$, (ii) is trivial
($b=0$), and the exact sequence $H^2(T_r)\to\prod H^2(D_v,T_r)\to H^0(A[l^r])^\vee\to0$
shows $\#\operatorname{coker}f/l^r=l^r$ for all $r$, so
$\operatorname{coker}f\cong\Z_l$ is torsion-free, which gives (i).

\emph{Step 2.}  By (i) and Step 0, $\pi_X$ maps $\ker f$ onto $\ker f_r$ with
kernel $N_r=\ker f\cap l^rX$; so $\#(\ker f/N_r)=\#\Sha^2(T_r)=\#\Sha^1(A[l^r])$
by Theorem~\ref{i5:thm-pt}(1).  Let $P_r\subset\Sha^1_D(F,A)$ be the classes
represented by cocycles killed by $l^r$ that are local coboundaries of
elements of $A[l^r]$; the $P_r$ increase and exhaust $\Sha^1_D$ (cocycles have
finite range), and by (ii) $\Sha^1(A[l^r])\to P_r$ is bijective.

\emph{Step 3.}  If $(N_r)$ decreases to $0$, $(P_r)$ exhausts $Y$, and
$\#(X'/N_r)=\#P_r<\infty$ for large $r$, then $\lgth X'=\lgth Y$: if $\lgth Y$
is infinite both are unbounded; if it is $n$, then $P_r=Y$ for large $r$, the
surjections $X'/N_{r+1}\to X'/N_r$ between groups of order $l^n$ are
bijective, so $N_r$ is eventually $0$.  Apply this to $X'=\ker f=\Sha^2_D$ and
$Y=\Sha^1_D$.

For the last inequality, the exact sequence
$0\to\Sha^2_D\to H^2(\Gamma,T)\to\prod_vH^2(D_v,T)$ and local duality in the
limit, $H^2(D_v,T)\cong H^0(D_v,A)^\vee$, give
$V_D(e)\le\operatorname{sha}_D(e)+\sum_v\lgth H^0(D_v,A)$.
\end{proof}

\begin{lemma}\label{i5:lem-loch0}
For $F$ of degree $n$, $l$ prime, $e\ge1$:
$\sum_{v\mid l}\lgth H^0(D_v,\Q_l/\Z_l(-e))\le n(2+\lfloor\log_ln\rfloor+e)$.
\end{lemma}
\begin{proof}
$\chi(D_v)$ has index $i_v\le n$ in $\Z_l^\times$, so contains
$(1+l^2)^{i_v}$, and $v_l((1+l^2)^{i_ve}-1)=2+v_l(i_v)+v_l(e)$.  There are at most
$n$ places.
\end{proof}

\begin{definition}[effective $H^2$ bound]\label{i5:def-lemmaK}
Put $\mathrm{kB}_C(e,n,D)=\exp((2en)^C\log(2D))$.  $\mathrm{LemmaK}(C_K)$: for
every prime $l$, totally imaginary $L\supset\mu_l$ ($\mu_4$ if $l=2$), $l$-datum
$D$ and $e\ge1$, $H^2(G_{L,S},\Z_l(e+1))$ is finite and
$\log\#H^2\le\mathrm{kB}_{C_K}(e,[L:\Q],|d_L|)$.  $\mathrm{ShaK}(C)$: under the
same hypotheses $\operatorname{sha}_D(e)<\infty$ and
$\operatorname{sha}_D(e)\log l\le\mathrm{kB}_C(e,[L:\Q],|d_L|)$.
\end{definition}

\begin{proposition}\label{i5:prop-shak}
$\mathrm{ShaK}(C)$ holds for some $C$.
\end{proposition}

The proof is in \S\ref{i5:ssec-pf-shak}.  It passes to a layer $L'\supset L$
with $\mu_{l^c}\subset L'$ maximal and $2e+1<l^c$, $2e+1<r_2(L')$.  Along the
cyclotomic layers $F/L'$ it identifies $l^{\operatorname{sha}}$ with the number
of eigen-homomorphisms $\mathrm{Cl}(F)\to\Z/l^\nu$ for the twisted Galois action
(Kummer theory and descent along $L'(\mu_{l^\infty})$); these reductions are
proved there.  The counting step SCount and the torsion bound for
$H_m(\mathrm{GL}_N(\mathcal O_{L'}),\Z)$ (STors) are also proved there
(Propositions~\ref{i5:x-prop:scount}, \ref{i5:x-prop:stors}).  The counting
step SFew (Proposition~\ref{i5:x-prop:sfew}) is reduced in \S\ref{i5:x5-ssec}
to two formally verified propositions; so Proposition~\ref{i5:prop-shak} is
proved up to SFew.

\begin{corollary}\label{i5:cor-lemmaK}
$\mathrm{LemmaK}(C_K)$ holds with $C_K=\max(C,0)+9$.
\end{corollary}
\begin{proof}
By Proposition~\ref{i5:prop-shaleV} and Lemma~\ref{i5:lem-loch0},
$V_D(e)\le s+n(2+\lfloor\log_ln\rfloor+e)<\infty$, $s=\operatorname{sha}_D(e)$,
so $H^2$ is finite of order $l^{V_D(e)}$.  As $l\le n+1$,
$\log\#H^2\le\mathrm{kB}_C+4n\log(2D)(2e)(2e+3)\log(n+1)$, and with $x=2en\ge2$,
$y=x^{\max(C,0)}\log(2D)$ the right side is $\le(x^9-1)y+e^y\le e^{x^9y}$.
\end{proof}

\begin{theorem}[Soul\'e 1979, Invent.\ Math.\ 55, Th\'eor\`eme~5 (IV.3.2), with Borel 1974]\label{i5:thm-soule}
For a number field $L$, an odd prime $l$ and $e\ge1$,
$H^2(\mathcal O_L[1/l],\Z_l(e+1))$ is finite.
\end{theorem}

Theorem~\ref{i5:thm-soule} gives the finiteness in LemmaK but not the
effective size, which is the content of Proposition~\ref{i5:prop-shak}; we are
not aware of a published statement in this form.

\medskip\noindent\label{i5:ssec-IVVs}We now treat Cases IV and Vs.

\begin{lemma}\label{i5:lem-upV}
Let $L/\Q$ be finite abelian, $D$ an $l$-datum on $L$, $e\ge1$, $\psi'$ with
$G_L\le\ker\psi'$.  Then
$V_l(e,\psi')\le e_{\mathcal O}\operatorname{sha}_D(e)+[L:\Q]\lgth H^0(G,\rho_{e,\psi'})$.
If $L$ contains $\mu_l$, $l$ odd, then $V_D(e)\le2\,\mathrm{kB}_{C_K}(e,[L:\Q],|d_L|)$,
and if $L\subset\Q(\mu_m)$ then $\log|d_L|\le[L:\Q]\log m$.
\end{lemma}
\begin{proof}
Restrict to $G_L$; as in step (a) of Proposition~\ref{i5:prop-red}, the target
has length $e_{\mathcal O}\operatorname{sha}_D(e)$, and the kernel is bounded by
Lemma~\ref{i5:lem-lens}(3) with $T$ lifts of $\mathrm{Gal}(L/\Q)$.  The second claim is
Corollary~\ref{i5:cor-lemmaK} ($\lgth\le2\log\#$); the third is the
conductor--discriminant formula (reference to be confirmed).
\end{proof}

\begin{proposition}[Case IV]\label{i5:prop-countIV}
$\mathrm{Count}_{\rm IV}$ holds with $C=lN_\psi l^j+3$.
\end{proposition}
\begin{proof}
By Lemma~\ref{i5:lem-hostarith}(5), $\psi'$ has conductor dividing $N_\psi$, so
it is trivial on $G_{L_1}$ for $L_0=\Q(\zeta_{lN_\psi})$ and the layer
$L_1=\Q(\zeta_{lN_\psi l^j})$, with $n=[L_1:\Q]\le lN_\psi l^j$.  LayerLe gives
$\sum_{e<t}\operatorname{sha}_D(e)^\natural\le B$ on $L_1$.  By
Lemmas~\ref{i5:lem-upV} and \ref{i5:lem-local}(1),
$\mathrm{Vn}(e,\psi')\le\operatorname{sha}_D(e)^\natural+n\cdot2N_he$; summing,
$\le B+2nN_ht^2\le B+(N_ht)^{n+3}$.
\end{proof}

\begin{proposition}[Case Vs]\label{i5:prop-countVs}
$\mathrm{Count}_{\rm Vs}$ holds.
\end{proposition}
\begin{proof}
Let $M=\max B_\psi$.  In Case Vs, $\operatorname{ord}\psi'\le2o\,2^{c_2}<2oR_1$,
and $R_1\le16o(lq)^2l^{2\varphi(4o)}$ since $f_0\mid\varphi(4o)$ (the order of
$\theta_1$ divides $2o$); so
\[\operatorname{ord}(\psi')\,l\le n_V:=2o\cdot16o(MN_\psi)^2M^{2\varphi(4o)}M\]
(we write $n_V$ to avoid a clash with $n_0=v_2(l^{f_0}-1)$ of Definition~\ref{i5:def-cases}).
Let $L\subset\Q(\mu_{lN_h})$ be the abelian field with $\mu_l\subset L$,
$G_L\le\ker\psi'$, $[L:\Q]\le n_V$ (fixed field of the intersection of the
kernels of $\psi'$ and of reduction mod $l$).  By Lemma~\ref{i5:lem-upV} and
Proposition~\ref{i5:prop-shaleV},
$V'(1,\psi')\le e_{\mathcal O}V_D(1)+n_V\cdot2N_he_{\mathcal O}$, and with
$x=lN_h$, $V_D(1)\le2x^A$, $A=(2n_V)^{\lceil C_K\rceil}(n_V+1)$.  Hence
$\mathrm{Vn}\log l\le(2x^A+2n_VN_h)l\le x^{A+n_V+3}$.
\end{proof}

\subsection{The automorphic counts}\label{i5:ssec-counts}

The remaining counts are where the Galois input enters.  In Case~I the
argument is an Eisenstein congruence for $\mathrm U(3)$; in Cases~II and~III a
$\mathrm{GSp}_4$ argument; in Case~V a $\mathrm U(3)$ lifting argument over
$\Q(i)$.  We prove the reduction of Case~I to a single count and state the
counts.

\begin{definition}\label{i5:def-gd}
For an imaginary quadratic $K$, an odd prime $l$, $e\ge1$ and a character
$\psi'$ of even level $N$, $\mathrm{GD}(K,l,e,\psi')$ means:
$l\nmid\operatorname{ord}\psi'$, $(\mathrm{comp}_l\psi')^2=1$,
$\psi'(-1)=(-1)^e$, $l\nmid|d_K|h_K$, $\gcd(d_K,N)=1$,
$\ker(\chi\bmod l)\cap\ker\chi'_l\not\subset G_K$, and $(l-1)\nmid2e$.
\end{definition}

\begin{theorem}[$\mathrm{CorGK}$; the Case~I count]\label{i5:prop-corg}
For every imaginary quadratic field $K$ there are $c_5,A\in\N$ such that for
all $e,N\ge1$, $H\le(\Z/N)^\times$, distinct primes $p,p'\nmid30N|d_K|$,
coefficient rings $(\mathcal O_l)_l$ and finite sets $\mathcal L$ of primes each
of which admits a character $\psi\colon(\Z/N)^\times\to\mathcal O_l^\times$ with
$\ker\psi=H$ and $\mathrm{GD}(K,l,e,\psi)$: all the $V'(e,\psi)$ below are finite
and
\[
\sum_{l\in\mathcal L}\log l\sum_{\ker\psi=H}\mathrm{Vn}(e,\psi)\le2c_5\,\varphi([(\Z/N)^\times:H])\,(N(e+1))^A(\log p+\log p'),
\]
the inner sum being over the $\mathcal O_l$-valued characters of $(\Z/N)^\times$.
\end{theorem}

The proof is in \S\ref{i5:ssec-pf-counts}.  It bounds the Selmer length of
$\psi\chi^e$ by the $l$-adic size of the congruences between the Eisenstein
trace $\Psi+1+\Psi^{-1}$ (restricted to $G_K$) and the Frobenius traces of
irreducible conjugate self-dual automorphic three-dimensional representations
of $G_K$ of weights $(e,0,-e)$ and bounded level.  A single non-zero integer
controls these congruences for all $l$ at once.  The section proves the
reduction to two statements.  $\mathrm{OneCharU}$ is the Selmer bound for one
character, an automorphy lifting and rigidity theorem in a residually reducible
$\mathrm U(3)$ situation.  $\mathrm{NumberU3U}$ is the counting estimate.  Both
are recorded as formally verified (Propositions~\ref{i5:x-prop:onechar},
\ref{i5:x-prop:numberu3u}).

\begin{proposition}\label{i5:prop-countI}
Theorem~\ref{i5:prop-corg} implies $\mathrm{Count}_{\rm I}$.
\end{proposition}
\begin{proof}
Fix $\psi$. By Definition~\ref{i5:def-Bchi}, $p_K=p_K(\psi)$ is a prime $\equiv3\pmod4$, $p_K\ge7$, $p_K\nmid N_\psi$; let $K=\Q(\sqrt{-p_K})$, so $d_K=-p_K$.
Take $c_5,A$ from Theorem~\ref{i5:prop-corg} for $K$, primes $r_1>30N_\psi p_K$ and $r_2>r_1$, and $k\in\N$ with $2c_5(\log r_1+\log r_2)\le16^k$. We show the bound of
$\mathrm{Count}_{\mathrm I}$ with $C=k+1+A$ and these $r_1,r_2$.

Let $h$ be a host with character $\psi$, $r_1,r_2<q_1,d$, $\mathrm{Big}_h(t)$, $0<e<t$, coefficient rings $\mathcal O_l$, and $\mathcal L$ a finite set of odd primes
in Case~I. The prime divisors of $N_h=\mathrm{lcm}(2^{e+2},N_\psi)q_1d$ are $2$, those of $N_\psi$, $q_1$ and $d$; as $r_i>30N_\psi p_K$ and $r_i<q_1,d$,
neither $r_i$ divides $30N_h|d_K|$, and $r_1\ne r_2$. Let $T$ be the set of kernels of the characters in $\bigcup_{l\in\mathcal L}X_h(\mathcal O_l,e)$. For
$H'\in T$ apply Theorem~\ref{i5:prop-corg} to $H'$ and the set of $l\in\mathcal L$ for which some $\psi'\in X_h(\mathcal O_l,e)$ has kernel $H'$; such $\psi'$ satisfy GD by
Lemma~\ref{i5:x-lem:gd-host}. This gives finiteness of all $V'$, and, grouping the characters of $X_h(\mathcal O_l,e)$ by kernel (each group lies in
$\mathrm{orb}_{\mathcal O_l}(H')$ and $\mathrm{Vn}\ge0$),
\[
\sum_{l\in\mathcal L}\log l\sum_{\psi'\in X_h(\mathcal O_l,e)}\mathrm{Vn}(e,\psi')\le2c_5(\log r_1+\log r_2)(N_h(e+1))^A\sum_{H'\in T}\varphi([(\Z/N_h)^\times:H']).
\]
The quotients by the $H'\in T$ are cyclic, so by Lemma~\ref{i5:x-lem:orb}(c) the last sum is $\le\#(\Z/N_h)^\times\le N_h$. Finally $16\mid N_h$, $e+1\le t$ and
$2c_5(\log r_1+\log r_2)\le16^k\le N_h^k$ give $2c_5(\log r_1+\log r_2)N_h(N_h(e+1))^A\le N_h^{k+1}(N_ht)^A\le(N_ht)^{k+1+A}$.
\end{proof}


$\mathrm{Count}_{\rm II}$, $\mathrm{Count}_{\rm III}$ and $\mathrm{Count}_{\rm V}$ are Propositions~\ref{i5:prop-c23} and \ref{i5:prop-c5}, which are stated and proved in \S\ref{i5:x5-ssec}; their reductions to smaller formally verified automorphy-lifting statements are given in \S\ref{i5:x5-ssec}.

\begin{remark}
Propositions~\ref{i5:x-prop:onechar}, \ref{i5:prop-c23} and \ref{i5:prop-c5}
rest on new mathematics (automorphy lifting in residually reducible $\mathrm U(3)$
and $\mathrm{GSp}_4$ situations).  Their reductions are proved in
\S\ref{i5:x5-ssec}; the remaining lifting cores are stated there as formally
verified propositions, stated exactly with the definitions of
\S\ref{i5:x5-ssec-defs}; their architecture is in \S\ref{i5:x5-ssec-arch}.  Proposition~\ref{i5:x-prop:numberu3u} is
deduced from a smaller formally verified count, Proposition~\ref{i5:x-prop:u3count}.
\end{remark}

\subsection{The minus part and the exceptional primes}\label{i5:ssec-minus}

For a Dirichlet character $\chi_0$ of level $N$ and $k\ge1$ put
$\mathrm{den}(N,k)=N\prod_{p\le k+1}p$, $a_k(\chi_0)=\mathrm{den}(N,k)B_{k,\chi_0}$
(generalised Bernoulli number of the associated primitive character) and
$\mathrm{bnd}(N,k)=16N^{k+1}k!$.

\begin{proposition}[$\mathrm{PsiBound}$]\label{i5:prop-psibound}
Let $E$ be a CM abelian field with a compatible surjection
$\rho\colon(\Z/N)^\times\to\mathrm{Gal}(E/\Q)$, $l$ an odd prime with
$l\nmid[E:\Q]$, $\mathcal O$ a coefficient ring, $e\ge1$, and
$\psi\colon\mathrm{Gal}(E/\Q)\to\mathcal O^\times$ with
$\psi\rho(-1)=(-1)^{e+1}$.  Then the $\psi$-isotypic part of
$\mathcal O\otimes\Sha^1(E,\Q_l/\Z_l(-e))$ has $\mathcal O$-length at most
$v_{\mathcal O}(a_{e+1}(\psi\rho))$.
\end{proposition}

This is the unconditional (divisibility) half of a main conjecture.  It is
proved in \S\ref{i5:ssec-pf-minus} (Theorem~\ref{i5:x-thm:thmbc}).  The proof
reduces to the Selmer group of $\alpha=(\psi\rho)^{-1}\chi^{-e}$ over $\Q$ and
disposes of the case $\alpha\equiv\chi$ modulo the maximal ideal directly.
Otherwise it runs Kolyvagin's argument with the Euler system of Stickelberger
(Jacobi-sum) elements: derivative classes at auxiliary primes chosen by
Chebotarev, congruences that evaluate the bottom class against
$B_{e+1,\psi}$, local duality at the auxiliary primes, and the reciprocity law
from Theorem~\ref{i5:thm-pt}.

\begin{proposition}[primes dividing the degree]\label{i5:prop-stick}
There is $C\ge0$ such that for every CM field $E\subset\Q(\zeta_N)$, odd $l$,
$l$-datum $D$ and $e\ge1$: $\operatorname{sha}^-_D(e)<\infty$ and
$\operatorname{sha}^-_D(e)^\natural\log l\le Cl[E:\Q]^2(e+1)\log(elN)^2$.
\end{proposition}

This is Proposition~\ref{i5:x-prop:stickbound}, proved in
\S\ref{i5:ssec-pf-minus}.

It combines an annihilator $l^r$ of $\Sha^-$ with
$r\log l\le30[E:\Q](e+1)\log(elN)$, coming from a Stickelberger-type element
at level $Nl$ whose sign-determinant divides a non-zero integer of controlled
size, with the $l$-rank bound
$\lgth H^1[l]\le8l[E:\Q]\log(lN)$ obtained by Kummer theory over $E(\mu_l)$ and
the Minkowski bound for the class number.

\begin{lemma}[Bernoulli numbers]\label{i5:lem-bern}
(1) $a_k(\chi_0)$ is an algebraic integer, $|a_k(\chi_0)|\le\mathrm{bnd}(N,k)$
for complex $\chi_0$, and $a_k(\chi_0)\ne0$ if $\chi_0(-1)=(-1)^k$.
(2) If $X$ is a Galois-stable set of characters of level $N$ with
$\chi(-1)=(-1)^k$, then $M=\prod_{\chi\in X}a_k(\chi)$ is a non-zero integer,
and for every $l$ and $T\subseteq X$,
$\sum_{\chi\in T}v_l(a_k(\chi))\le v_l(M)$, with
$\sum_lv_l(M)\log l\le|X|\log\mathrm{bnd}(N,k)$.
\end{lemma}
\begin{proof}
(1) Expand $B_{k,\chi_0}=\sum_j\binom kjB_jf^{j-1}\sum_a\chi_0(a)a^{k-j}$; the
denominators of $B_j$ divide $\prod_{p\le k+1}p$ (von Staudt--Clausen;
reference to be confirmed) and that of $f^{j-1}$ divides $N$.  The bound follows
from $|B_k(x)|\le4k!/(2\pi)^k$ on $[0,1]$ and $\prod_{p\le k+1}p\le4^{k+1}$.
Non-vanishing: $L(1-k,\chi_0)=-B_{k,\chi_0}/k$ is non-zero in this parity by
the functional equation and $L(k,\bar\chi_0)\ne0$.  (2) $M$ is Galois-invariant
and integral, hence in $\Z$; all factors are $l$-integral, and
$\sum_lv_l(M)\log l=\log|M|$.
\end{proof}

\begin{proposition}\label{i5:prop-thmbc}
Propositions~\ref{i5:prop-psibound} and \ref{i5:prop-stick} imply
$\mathrm{ThmBcFree}$.
\end{proposition}
\begin{proof}
Let $E$ be a CM $(M,e,q,d)$-field, $F=\mathfrak f(E)$, $n=[E:\Q]$.  The exponent
of $\mathrm{Gal}(E/\Q)$ divides $2^{e+1}\deg M$, so an odd prime dividing $n$
divides $\deg M$.  For odd $l\nmid n$, decompose $\mathcal O\otimes\Sha^-$ into
isotypic parts (possible as $l\nmid n$); the minus condition selects the $\psi$
with $\psi\rho(-1)=(-1)^{e'+1}$, and Proposition~\ref{i5:prop-psibound} with
Lemma~\ref{i5:lem-bern}(2) gives
$\sum_{l\nmid n}\log l\operatorname{sha}^-(e')^\natural\le n\log\mathrm{bnd}(F,e'+1)\le n(2t+5)\log(tF)$.
Summing over $e'<t$ gives $\le3nt^2\log(tF)$.  For $l\mid n$,
Proposition~\ref{i5:prop-stick} gives $\le Cln^2t^2\log(tlF)^2$ per prime.  The
total is $\beta_C(\deg M,n,F,t)$.
\end{proof}

\medskip\noindent\label{i5:ssec-bab}We turn to the exceptional primes.

\begin{theorem}[Ferrero--Washington, Ann.\ of Math.\ (2) 109 (1979), 377--395, with Washington, GTM 83, 2nd ed., Prop.~13.23; reference to be confirmed]\label{i5:thm-fw}
For $K$ abelian over $\Q$ containing $\mu_l$ ($\mu_4$ if $l=2$) there is $r$
with $\#\mathrm{Cl}(K_n)/l\le l^r$ for all layers $K_n$ of the cyclotomic
$\Z_l$-extension.
\end{theorem}
The published theorem is $\mu=0$ for every prime $l$, including $l=2$; the
form stated is equivalent to it by Washington, GTM 83, Prop.~13.23.

\begin{proposition}\label{i5:prop-thmbab}
Assume Theorems~\ref{i5:thm-pt} and \ref{i5:thm-fw} and
Corollary~\ref{i5:cor-lemmaK}.  Then there is a constant bundle $K$ with
$\mathrm{ThmBab}(K,C_t)$ for all $C_t\ge1$.
\end{proposition}

The proof is in \S\ref{i5:ssec-pf-minus} (Theorem~\ref{i5:x-thm:thmbab}).  Over a base field $L_0\supset\mu_l$
with $\mu_{l^{c_0}}\subset L_0$ maximal, descent along $L_0(\mu_{l^\infty})$
identifies $\operatorname{sha}(L_1,l,e)$, for every layer $L_1/L_0$ of height
$m$, with a twisted coinvariant length $\operatorname{tw}(e,m)$ on the Iwasawa
module $X_\infty$; Theorem~\ref{i5:thm-fw} makes $X_\infty$ generated by $R$
elements, on which a generator $\gamma$ is unipotent mod $l$, so
$\operatorname{sha}$ grows by at most $R$ per layer and is large only for $e$
in a set $\mathrm{Exc}$ of at most $R$ ``centres''.  For $l=2$ the base field is
the compositum of the field of $M$ with $\Q(i)$, $\Q(\sqrt{-q})$,
$\Q(\sqrt{-d})$, and every $(M,e,q,d)$-field $E$ satisfies
$\operatorname{sha}(E,2,k)\le1+\operatorname{sha}(L_1,2,k)$ for a $2$-layer of
height $\le e+2$.  Finally a pigeonhole choice of the window gives $Y$ with
$X^{1/J}\le Y$, $Y^2\le X$ such that no centre $e\le x^{C_t}$ exceeds
$x^{\varepsilon/(2C_K)}$ for $x\in[Y,Y^2]$; LemmaK then bounds the remaining
terms by $\exp(x^{4\varepsilon/5})$.  This is where the window (rather than a
single scale) is needed.

\subsection{Proof of Theorem~\ref{i5:thm-vwin}}\label{i5:ssec-proof}

\begin{proof}
$\mathrm{Count}_{\rm Z}$, $\mathrm{Count}_{\rm IV}$ and $\mathrm{Count}_{\rm Vs}$
are Propositions~\ref{i5:prop-countZ}, \ref{i5:prop-countIV} and
\ref{i5:prop-countVs} (the last using Corollary~\ref{i5:cor-lemmaK});
$\mathrm{Count}_{\rm I}$ is Propositions~\ref{i5:prop-corg} and
\ref{i5:prop-countI}; $\mathrm{Count}_{\rm II}$, $\mathrm{Count}_{\rm III}$,
$\mathrm{Count}_{\rm V}$ are Propositions~\ref{i5:prop-c23} and
\ref{i5:prop-c5} (stated and proved in \S\ref{i5:x5-ssec}).  By Proposition~\ref{i5:prop-sum}, SumIn and SumOut hold;
by Proposition~\ref{i5:prop-point}, $\mathrm{ThmHV}(K,401)$ for every $K$.
Proposition~\ref{i5:prop-thmbab} gives $K$ with $\mathrm{ThmBab}(K,401)$, so
$\mathrm{JointWin}(401)$ by Proposition~\ref{i5:prop-jointwin}.
Proposition~\ref{i5:prop-thmbc} gives ThmBcFree, hence $\mathrm{ThmBc}(401)$
(Proposition~\ref{i5:prop-vwinat}), and Proposition~\ref{i5:prop-vwinat}
gives $\mathrm{VWinAt}(\psi,401)$ for every $\psi$.
\end{proof}


\subsection{Stickelberger elements, the minus part and the exceptional primes}\label{i5:ssec-pf-minus}

We prove Propositions~\ref{i5:prop-psibound}, \ref{i5:prop-stick} and \ref{i5:prop-thmbab}. Lemmas~\ref{i5:x-lem:bern-int} and
\ref{i5:x-lem:bern-prod} restate Lemma~\ref{i5:lem-bern} in the form used here.

\begin{theorem}[minus part]\label{i5:x-thm:thmbc}
Poitou--Tate duality and unramified orthogonality imply $\mathrm{PsiBound}$, and $\mathrm{PsiBound}$ implies $\mathrm{ThmBcFree}$. Here $\mathrm{PsiBound}$
is the statement: let $E$ be a CM abelian field with a surjection $\rho\colon(\Z/N)^\times\to\mathrm{Gal}(E/\Q)$ compatible with the cyclotomic action,
$l$ an odd prime with $l\nmid[E:\Q]$, $\mathcal O$ a coefficient ring, $e\ge1$, and $\psi\colon\mathrm{Gal}(E/\Q)\to\mathcal O^\times$ with
$\psi\rho(-1)=(-1)^{e+1}$. Then the $\mathcal O$-length of the $\psi$-part of the strict Selmer group $\Sha^1$ of $\mathcal O\otimes\Q_l/\Z_l(-e)$ over $E$
is at most $v_{\mathcal O}\big(\mathrm{den}(N,e+1)\,B_{e+1,\psi\rho}\big)$, where $B_{k,\chi}$ is the generalised Bernoulli number of the primitive
character attached to $\chi$ and $\mathrm{den}(N,k)=N\prod_{p\le k+1}p$.
\end{theorem}

\begin{lemma}[Bernoulli numbers]\label{i5:x-lem:bern-int}
Let $\chi$ be a Dirichlet character of level $N$ with values in a field of characteristic $0$, $k\ge1$, and $a_k(\chi)=\mathrm{den}(N,k)B_{k,\chi}$. Then
$a_k(\chi)$ is an algebraic integer; for complex $\chi$, $|a_k(\chi)|\le\mathrm{bnd}(N,k):=16N^{k+1}k!$; $a_k(\chi)\ne0$ if $\chi(-1)=(-1)^k$; and
$a_k$ at level $N$ divides $a_k$ at level $N'$ when $N\mid N'$.
\end{lemma}
\begin{proof}
Let $\chi_0$ be primitive of conductor $f\mid N$. Then $B_{k,\chi_0}=\sum_{j=0}^{k}\binom kjB_jf^{j-1}\sum_{a=1}^f\chi_0(a)a^{k-j}$; by von Staudt--Clausen
the denominator of $B_j$ divides $\prod_{p\le k+1}p$, and $f^{j-1}$ contributes at most $f\mid N$; so $a_k$ is integral. From
$|B_k(x)|\le4k!/(2\pi)^k$ on $[0,1]$ (Fourier expansion for $k\ge2$, direct for $k=1$) and $\prod_{p\le k+1}p\le4^{k+1}$ (Erd\H{o}s) we get
$|a_k|\le N\cdot4^{k+1}\cdot4N^kk!/(2\pi)^k\le16N^{k+1}k!$. Non-vanishing: $L(1-k,\chi_0)=-B_{k,\chi_0}/k$ is non-zero for $\chi_0(-1)=(-1)^k$ by the
functional equation (no $\Gamma$-pole at $1-k$ in this parity, $L(k,\overline{\chi_0})\ne0$, using $L(1,\chi)\ne0$ for $k=1$; $\chi_0=1$, $k=1$ is
excluded by parity). Finally $\mathrm{den}(N',k)=(N'/N)\mathrm{den}(N,k)$.
\end{proof}

\begin{lemma}[product of Bernoulli numbers]\label{i5:x-lem:bern-prod}
Let $k\ge1$ and $X$ a $\mathrm{Gal}(\overline\Q/\Q)$-stable set of Dirichlet characters of level $N$ with $\chi(-1)=(-1)^k$. Then
$M=\prod_{\chi\in X}a_k(\chi)$ is a non-zero integer with $\log|M|\le|X|\log\mathrm{bnd}(N,k)$; for every prime $l$ and $T\subseteq X$ (viewed
$l$-adically), $\sum_{\chi\in T}v_l(a_k(\chi))\le v_l(M)$; and $\sum_lv_l(M)\log l\le|X|\log\mathrm{bnd}(N,k)$.
\end{lemma}
\begin{proof}
$\sigma(a_k(\chi))=a_k(\sigma\chi)$, so $M$ is a Galois-invariant algebraic integer, hence in $\Z$; the rest is Lemma~\ref{i5:x-lem:bern-int} and
integrality of all factors in any $l$-adic embedding.
\end{proof}

\begin{proposition}[primes not dividing the degree]\label{i5:x-prop:primebound}
Assume $\mathrm{PsiBound}$. Let $E\subset\Q(\zeta_N)$ be CM, $(D_p)_p$ a family of $p$-data and $e\ge1$. Then $\operatorname{sha}^-_{D_p}(e)<\infty$ for odd
$p\nmid[E:\Q]$, and for every finite set $\mathcal S$ of such primes
$\sum_{l\in\mathcal S}\log l\cdot\operatorname{sha}^-_{D_l}(e)^\natural\le[E:\Q]\log\mathrm{bnd}(N,e+1)$.
\end{proposition}
\begin{proof}
Fix $l\in\mathcal S$ and $\mathcal O$ containing all character values of $\mathrm{Gal}(E/\Q)$. As $l\nmid[E:\Q]$, $\mathcal O\otimes\Sha^-$ is the sum of its
$\psi$-parts, and the minus condition selects the $\psi$ with $\psi\rho(-1)=(-1)^{e+1}$ ($\rho(-1)$ is complex conjugation). With
$\lgth_{\mathcal O}(\mathcal O\otimes M)=e_{\mathcal O}\lgth_{\Z_l}M$ and $\mathrm{PsiBound}$, $\operatorname{sha}^-\log l\le\sum_{\chi\in T}v_l(a_{e+1}(\chi))\log l$ for the
set $T$ of such $\chi=\psi\rho$; in particular it is finite. The set $X$ of complex characters of $(\Z/N)^\times$ trivial on $\ker\rho$ with
$\chi(-1)=(-1)^{e+1}$ is Galois stable with $|X|\le[E:\Q]$; Lemma~\ref{i5:x-lem:bern-prod} bounds the sum over $l\in\mathcal S$.
\end{proof}

\medskip\noindent\emph{Higher Stickelberger elements.}
Let $M\ge2$, $\Gamma$ a finite abelian group, $\iota\colon(\Z/M)^\times\to\Gamma$ a homomorphism, $[g]$ the image of $g\in\Gamma$ in a group ring,
and $\langle b\rangle\in\{0,\dots,M-1\}$ the representative of $b\in\Z/M$. For $k\ge1$ and $a$ prime to $M$ put $E_k(b)=M^{k-1}B_k(\langle b\rangle/M)/k$ and
\[
\theta_k=-\sum_{b\in(\Z/M)^\times}E_k(b)[\iota(b)]^{-1},\qquad
y_k(a)=\big(1-a^k[\iota(a)]^{-1}\big)\theta_k=-\sum_{b}\big(E_k(b)-a^kE_k(a^{-1}b)\big)[\iota(b)]^{-1}.
\]
For $a$ odd and prime to $M$ put $j_a(b)=(a\langle a^{-1}b\rangle-\langle b\rangle)/M\in\Z$, $E^{(1)}_a(b)=\tfrac{a-1}2-j_a(b)$, and for $e\ge0$
$\tilde y_e(a)=-\sum_b\langle b\rangle^eE^{(1)}_a(b)[\iota(b)]^{-1}\in\Z[\Gamma]$; superscripts $(M)$ indicate the level. For $\chi\colon\Gamma\to\C^\times$ let
$\chi_\iota(b)=\chi(\iota(b))^{-1}$ and $B_{k,\chi_\iota}=M^{k-1}\sum_b\chi_\iota(b)B_k(\langle b\rangle/M)$ (computed at level $M$). Put
$\mathrm{denY}(M,k)=Mk(k+1)!$, $\mathrm{sizeY}(M,a,k)=2a^kM^kk!$, $\mathrm{bdX}=\mathrm{denY}\cdot\mathrm{sizeY}$, $\mathrm{dk}(k)=2k\prod_{p\le k+1}p$.

\begin{lemma}[signed determinants]\label{i5:x-lem:BC-signdet}
Let $A$ be a commutative ring, $\Gamma$ finite abelian, $c\in\Gamma$ with $c^2=1\ne c$, $\varepsilon=\pm1$, $g_1,\dots,g_s$ representatives of
$\Gamma/\langle c\rangle$, and for $z\in A[\Gamma]$ put $\det_{c,\varepsilon}(z)=\det\big(z(g_ig_j^{-1})+\varepsilon z(cg_ig_j^{-1})\big)_{i,j}$. Then
(1) $\det_{c,\varepsilon}(z)\in([c]-\varepsilon,z)$; so if $cx=\varepsilon x$ and $zx=0$ in an $A[\Gamma]$-module then $\det_{c,\varepsilon}(z)x=0$;
(2) under $A\to\Omega$, a field containing the $\exp(\Gamma)$-th roots of unity, $\det_{c,\varepsilon}(z)\mapsto\prod_{\chi(c)=\varepsilon}\chi(z)$;
(3) $\det_{c,\varepsilon}$ commutes with ring maps and $\det_{c,\varepsilon}(tz)=t^s\det_{c,\varepsilon}(z)$.
\end{lemma}
\begin{proof}
$R=A[\Gamma]/([c]-\varepsilon)$ is free on $[g_1],\dots,[g_s]$, and the matrix is that of multiplication by $z$ on $R$; the adjugate identity gives
(1). Over $\Omega$, $R\cong\prod_{\chi(c)=\varepsilon}\Omega$ and $m_z$ is diagonal with entries $\chi(z)$, giving (2). (3) is clear.
\end{proof}

\begin{lemma}[higher Stickelberger determinants]\label{i5:x-lem:BC-ydet}
Let $a\ge2$ be prime to $M$, $k\ge2$, $c=\iota(-1)\ne1$, $\varepsilon=(-1)^k$. Then (1) $\mathrm{denY}(M,k)y_k(a)\in\Z[\Gamma]$; (2)
$\chi(y_k(a))=(1-a^k\chi_\iota(a))(-B_{k,\chi_\iota}/k)$, non-zero if $\chi(c)=\varepsilon$, with $|\chi(y_k(a))|\le\mathrm{sizeY}(M,a,k)$; (3)
$X:=\mathrm{denY}^{|\Gamma|/2}\det_{c,\varepsilon}(y_k(a))$ is a non-zero integer with $|X|\le\mathrm{bdX}(M,a,k)^{|\Gamma|/2}$.
\end{lemma}
\begin{proof}
(1) $kE_k(b)=\sum_i\binom kiB_i\langle b\rangle^{k-i}M^{i-1}$ and von Staudt--Clausen. (2) $\chi(\theta_k)=-B_{k,\chi_\iota}/k$. If $\chi(c)=\varepsilon$, then
$B_{k,\chi_\iota}=B_{k,\chi_0}\prod_{p\mid M}(1-\chi_0(p)p^{k-1})\ne0$ by Lemma~\ref{i5:x-lem:bern-int} ($k\ge2$), and $|a^k\chi_\iota(a)|>1$. The size
bound follows from $|B_k(x)|\le4k!/(2\pi)^k$. (3) $X=\det_{c,\varepsilon}(\mathrm{denY}\,y_k(a))$ is an integer determinant; apply
Lemma~\ref{i5:x-lem:BC-signdet}(2) and (2).
\end{proof}

\begin{lemma}[change of level; first-moment congruence]\label{i5:x-lem:BC-ylevel}
Let $l\mid N$, $a$ odd prime to $N$, $e\ge1$, $k=e+1$, $\iota\colon(\Z/N)^\times\to\Gamma$, $M_m=Nl^m$, $\iota_m=\iota\circ\mathrm{red}$. Then
(1) $y_k(a)^{(M_m)}=y_k(a)^{(N)}$; (2) at any level $M$ with $a$ odd prime to $M$, $\mathrm{dk}(k)(y_k(a)-\tilde y_e(a))\in M\Z[\Gamma]$; (3)
$z:=y_k(a)^{(N)}\in\Z_l[\Gamma]$, and for each $r$, $\tilde y_e(a)^{(M_m)}\equiv z\bmod l^r$ for $m$ large.
\end{lemma}
\begin{proof}
(1) The fibre over $b$ is $\{\langle b\rangle+jN\}_{j<l^m}$; Raabe's formula $\sum_{j<L}B_k((x+j)/L)=L^{1-k}B_k(x)$ shows $\theta_k$ is level-independent.
(2) With $x=\langle a^{-1}b\rangle$, $ax=\langle b\rangle+Mj_a(b)$; expanding $kE_k(t)=t^k/M-\frac k2t^{k-1}+\sum_{i\ge2}\binom kiB_it^{k-i}M^{i-1}$ gives
$k(E_k(b)-a^kE_k(x))=k\langle b\rangle^{k-1}(\tfrac{a-1}2-j_a(b))+MR$, $R$ with denominator dividing $2\prod_{p\le k+1}p$. (3) For
$m\ge r+v_l(\mathrm{dk}(k))$, (1) and (2) give $y_k(a)^{(N)}-\tilde y_e(a)^{(M_m)}\in l^r\Z_{(l)}[\Gamma]$; $r=0$ gives integrality.
\end{proof}

\begin{lemma}[Stickelberger]\label{i5:x-lem:BC-clkill}
For $M\ge3$, $K=\Q(\mu_M)$ and $a$ odd prime to $M$, $\sum_{b}E^{(1)}_a(b)\sigma_b^{-1}$ annihilates $\mathrm{Cl}(K)$.
\end{lemma}
\begin{proof}
The norm element kills $\mathrm{Cl}(K)$. Since $Mj_a(b)=a\langle a^{-1}b\rangle-\langle b\rangle$, $\sum_bj_a(b)\sigma_b^{-1}=(a\sigma_{a^{-1}}-1)\Theta$ with
$\Theta=\sum_t\frac{\langle t\rangle}M\sigma_t^{-1}\in\Q[\mathrm{Gal}(K/\Q)]$.  $\Theta$ itself is not integral, but
$(a\sigma_{a^{-1}}-1)\Theta=\sum_bj_a(b)\sigma_b^{-1}$ is, and Stickelberger's theorem (Washington, GTM 83, Thm.~6.10: if
$\beta\in\Z[\mathrm{Gal}(K/\Q)]$ and $\beta\Theta$ is integral then $\beta\Theta$ annihilates $\mathrm{Cl}(K)$; here $\beta=a-\sigma_a$ up to the unit
$\sigma_{a^{-1}}$, cf.\ Washington Lemma~6.9; proved in the formalization via Gauss sums) shows that it kills $\mathrm{Cl}(K)$.
\end{proof}

\begin{lemma}[annihilation of the Selmer group]\label{i5:x-lem:BC-annih}
Let $E\subset\Q(\mu_N)$ be Galois over $\Q$, $\rho\colon(\Z/N)^\times\to\mathrm{Gal}(E/\Q)$ the cyclotomic action, $l\mid N$, $D$ an $l$-datum on $E$
with $l$ odd or the mod-$4$ cyclotomic character trivial on $G_E$, $e\ge1$, $a$ odd prime to $N$, and $z\in\Z_l[\mathrm{Gal}(E/\Q)]$ the element
$y_{e+1}(a)^{(N)}$ for $\iota=\rho^{-1}$. Then $z\cdot\varphi(N)x=0$ for $x\in\Sha^1_D(E,\Q_l/\Z_l(-e))$.
\end{lemma}
\begin{proof}
Let $A=\Q_l/\Z_l(-e)$, $U_\infty=G_{E(\mu_{l^\infty})}$ and $\mathrm{res}\colon H^1(G_{E,S},A)\to\mathrm{Hom}(U_\infty,\Q_l/\Z_l)$.
(i) $\mathrm{res}$ is injective: its kernel is $H^1$ of a procyclic group with generator $\gamma$ of infinite order, i.e.\
$A/(\chi_{\mathrm{cyc}}(\gamma)^{-e}-1)A=0$.
(ii) If $\mathrm{res}(y)$ vanishes on $U_\infty\cap G_{\Q(\mu_N)}$ then $\varphi(N)y=0$, since $u^{\varphi(N)}$ lies there for $u\in U_\infty$.
(iii) Let $l^rx=0$. $\mathrm{res}(x)$ is unramified everywhere and trivial on decomposition groups above $l$; by compactness it extends to such a
$\psi$ on $G_{E(\mu_{l^n})}$. Let $m\ge\max(n,r)+2$, $M=Nl^m$, $K=\Q(\mu_M)$, $u\in U_\infty\cap G_{\Q(\mu_N)}\subset G_K$, $\sigma_b\in G_\Q$ acting on
$\mu_M$ by $b$, $w_b=\sigma_b^{-1}u\sigma_b$. Since $\mathrm{res}(\rho(b)x)(u)=\chi_{\mathrm{cyc}}(\sigma_b)^{-e}\mathrm{res}(x)(w_b)$ and
$\chi_{\mathrm{cyc}}(\sigma_b)\equiv\langle b\rangle\bmod l^m$,
$\mathrm{res}(\tilde y_e(a)^{(M)}x)(u)=-\sum_bE^{(1)}_a(b)\psi(w_b)$. By class field theory $\psi|_{G_K}=\psi''\circ\mathrm{Art}_K$ with
$\psi''$ on $\mathrm{Cl}(K)$, and $\mathrm{Art}_K(w_b)=\sigma_b^{-1}\mathrm{Art}_K(u)$; the sum vanishes by Lemma~\ref{i5:x-lem:BC-clkill}.
(iv) By Lemma~\ref{i5:x-lem:BC-ylevel}(3), $zx=\tilde y_e(a)^{(M)}x$ for $m$ large; conclude by (iii) and (ii).
\end{proof}

\begin{lemma}[size of $\mathrm{bdX}$]\label{i5:x-lem:BC-bdx}
For $N\ge1$, $l\ge2$, $e\ge1$: $\log\mathrm{bdX}(Nl,2Nl+1,e+1)+\log(Nl)\le30(e+1)\log(elN)$.
\end{lemma}
\begin{proof}
With $T=3elN$, $M=Nl$, $k=e+1$: $M,k\le T$, $a\le T^3$, $(k+1)!\le T^{2k+2}$, $k!\le T^k$, so $M\,\mathrm{bdX}\le T^{7k+6}\le T^{10k}$, and
$\log T\le3\log(elN)$.
\end{proof}

\begin{proposition}[primes dividing the degree]\label{i5:x-prop:stickbound}
There is $C\ge0$ such that for every CM field $E\subset\Q(\zeta_N)$, prime $l$ (odd, or $l=2$ with the mod-$4$ cyclotomic character trivial on $G_E$),
$l$-datum $D$ and $e\ge1$: $\operatorname{sha}^-_D(e)<\infty$ and $\operatorname{sha}^-_D(e)^\natural\log l\le C\,l\,[E:\Q]^2(e+1)\log(elN)^2$.
\end{proposition}
\begin{proof}
(a) \emph{Annihilation.} Put $N'=Nl$, $a=2N'+1$, $\Gamma=\mathrm{Gal}(E/\Q)$, $c=\rho'(-1)\ne1$, $\varepsilon=(-1)^{e+1}$, $z$ as in
Lemma~\ref{i5:x-lem:BC-annih} at level $N'$. For $x\in\Sha^-_D(e)$, $x'=\varphi(N')x$ has $cx'=\varepsilon x'$, $zx'=0$, so
$\det_{c,\varepsilon}(z)x'=0$ (Lemma~\ref{i5:x-lem:BC-signdet}). This determinant divides in $\Z_l$ the integer $X$ of Lemma~\ref{i5:x-lem:BC-ydet}(3),
hence divides $l^{v_l(X)}$, $l^{v_l(X)}\le\mathrm{bdX}^{[E:\Q]/2}$. With $l^{v_l\varphi(N')}\le N'$ and Lemma~\ref{i5:x-lem:BC-bdx}, $l^r$ kills
$\Sha^-_D(e)$ with $r\log l\le30[E:\Q](e+1)\log(elN)$.
(b) \emph{$l$-rank.} $\lgth H^1(G_{E,S},\Q_l/\Z_l(-e))[l]\le8l[E:\Q]\log(lN)$: over $F=E(\mu_l)$ (degree prime to $l$ over $E$) the twist is trivial mod $l$,
and by Kummer theory $H^1[l]$ injects into an $l$-Selmer group of order $\le l^{2[F:\Q]}h_F$, with $\log h_F\le3[F:\Q]\log(lN)$ (Minkowski,
$|d_F|\le(lN)^{[F:\Q]}$) and $[F:\Q]\le l[E:\Q]$.
(c) If $a^r$ kills $P\subset X$ then $\lgth P\le r\lgth X[a]$. Combine.
\end{proof}

\medskip\noindent\emph{Notation for the Euler-system argument.}
$l$ is odd, $\mathcal O$ a coefficient ring with maximal ideal $\mathfrak m$, uniformiser $\varpi$, valuation $v$; $W=\mathcal O\otimes\Q_l/\Z_l$;
$G=G_\Q$; $\chi$ the cyclotomic character; $G_n=G_{\Q(\mu_n)}$ (so commutators lie in $G_n$). A \emph{place} $A$ is a valuation ring of $\overline\Q$
above $p$, with $D_A\supset I_A$ and Frobenius elements $\varphi_A$; $A$ \emph{divides} $n$ if $p\mid n$. If $A\nmid n$ then $I_A\subset G_n$ and $\varphi_A$
acts on $\mu_n$ by $p$; $D_A$ is topologically generated by $I_A$ and $\varphi_A$. For $\gamma\colon G\to\mathcal O^\times$ a \emph{$\gamma$-cocycle} is a
continuous $c\colon G\to W$ with $c(gh)=c(g)+\gamma(g)c(h)$; $\partial m(g)=(\gamma(g)-1)m$. For $\alpha\colon G\to\mathcal O^\times$ put $\beta=\chi\alpha^{-1}$.

\begin{lemma}[the module $W$]\label{i5:x-lem:BC-W}
$W=K/\mathcal O$; its submodules are the $\varpi^{-n}\mathcal O/\mathcal O$. Hence: (i) each $w$ has annihilator $(a)$, $a\ne0$; (ii) $W[a]=\mathcal Ow_0\cong\mathcal O/a$
for any $w_0$ with annihilator $(a)$; (iii) of two elements one is a multiple of the other; (iv) if $f\colon X\to W$ is linear and $ef(X)=0$ then
$\lgth X\le\lgth\ker f+v(e)$.
\end{lemma}

\begin{definition}\label{i5:x-def:BC-sht}
For continuous $\alpha\colon G\to\mathcal O^\times$, $\mathrm{Sel}(\alpha)\subset H^1(G,W(\alpha))$ consists of the classes vanishing on every $I_A$ and on
$D_A$ for $A\mid l$. For $\delta\in G$ with $\alpha(\delta)-1\in\mathcal O^\times$, $F\ge1$, $a\ne0$, $\mathrm{ShT}(\alpha,\delta,F,a)$ is the module of
$\alpha$-cocycles $c$ with (T1) $c(\delta)=0$; (T2) $ac=0$; (T3) $c|_{I_A}=0$ for $A\nmid F$; (T4) $c|_{D_A}=\partial m$ with $m\in W[a]$ for $A\mid F$.
$S_n\subset\mathrm{ShT}$ is the submodule with $c|_{D_A}=0$ for all $A\mid n$.
\end{definition}

\begin{lemma}[from cocycles to the Selmer group]\label{i5:x-lem:BC-shasht}
Let $\alpha(\delta)-1\in\mathcal O^\times$, $\alpha(I_A)=1$ for $A\nmid F$, and for $A\mid F$, $A\nmid l$ let some $g\in I_A$ have $\alpha(g)-1$ a unit. If
$\lgth\mathrm{ShT}(\alpha,\delta,F,a)\le n$ for all $a\ne0$, then $\lgth\mathrm{Sel}(\alpha)\le n$.
\end{lemma}
\begin{proof}
For an $\alpha$-cocycle: (a) if $c=\partial m$ at $\sigma^{-1}x\sigma$ then $c(x)=\partial(\alpha(\sigma)m-c(\sigma))(x)$; (b) if $c=\partial m$ at $g$
and at $d^{-1}gd$ with $\alpha(g)-1$ a unit, then $c(d)=\partial m(d)$. Given $[c]\in\mathrm{Sel}(\alpha)$, subtract $\partial((\alpha(\delta)-1)^{-1}c(\delta))$ for
(T1); take $a$ killing $m_0$ (with $c|_{D_{A_0}}=\partial m_0$, $A_0\mid l$) and the finitely many values of $c$ for (T2); (T3) from $\alpha|_{I_A}=1$;
(T4) at $A\mid l$ by (a) and conjugacy, at $A\mid F$, $A\nmid l$ by (b) since $D_A$ normalises $I_A$. So $\mathrm{Sel}(\alpha)$ is the directed union of
the images of $\mathrm{ShT}(\alpha,\delta,F,\varpi^k)$.
\end{proof}

\begin{lemma}[the exceptional case]\label{i5:x-lem:BC-exc}
If $\alpha\equiv\chi\bmod\mathfrak m$ on $G$ and $\alpha(I_A)=1$ for $A\nmid l$, then $\mathrm{Sel}(\alpha)=0$.
\end{lemma}
\begin{proof}
Let $n(g)$ give the action on $\mu_l$, so $\alpha\equiv n\bmod\varpi$; take $t\in I_{A_0}$, $A_0\mid l$, acting by $-1$ on $\mu_l$, so $\alpha(t)-1$ is a
unit. By Lemma~\ref{i5:x-lem:BC-shasht} ($F=l$, $\delta=t$) and induction on $v(a)$ it suffices that $c\in\mathrm{ShT}(\alpha,t,l,a)$, $\varpi c=0$, forces
$c=0$. (T4) and $c(t)=0$ give $c|_{D_{A_0}}=0$. For additive $\lambda\colon W[\varpi]\to\Z/l$, $y=\lambda\circ c$ is a cocycle into $\mu_l$; by Kummer theory
$\zeta^{y(g)}=g\beta/\beta$ with $\beta^l=b\in\Q^\times$. As $y$ vanishes on an inertia group at every $p$, $l\mid\mathrm{ord}_p(b)$ for all $p$, so
$b=y_0^l$ ($l$ odd), $\beta=\zeta^jy_0$ and $y=(n-1)j$; $y(t)=0$ gives $j=0$.
\end{proof}

\begin{lemma}[reciprocity]\label{i5:x-lem:BC-rec}
Assume, for $F=\Q$, $l$ odd, $r\ge1$, $\mu=\Z/l^r(1)$:
(PT) for $S\ni l$ finite and $w\in H^2(G_S,\mu)$, $\sum_{p\in S}\mathrm{inv}_p(w|_{D_p})=0$, each $\mathrm{inv}_p$ injective;
(UO) for $A\mid p\ne l$, $f\in H^1(D_A,\mu)$ is unramified iff $f\cup x=0$ for all unramified $x\in H^1(D_A,\Z/l^r)$.
Let $\alpha$ be continuous, $w_0\in W$ with annihilator $(a)$, $\{\alpha\equiv1\bmod a\}$ open, $l\mid N$, $q\ne l$ prime, $c$ an $\alpha$-cocycle and $n$
a $\beta$-cocycle with values in $W[a]$, with: (i) for $A\nmid Nq$, $\alpha\equiv1\bmod a$ and $c=n=0$ on $I_A$; (ii) for $A\mid N$, $A\nmid q$,
$c|_{D_A}=\partial m$, $m\in W[a]$; (iii) for some $A_0\mid q$, $\alpha\equiv\beta\equiv1\bmod a$ on $D_{A_0}$ and $c|_{I_{A_0}}=0$. Then for
$g\in D_{A_0}$, $\tau\in I_{A_0}$ and $n(\tau)=yw_0$: $yc(g)=0$.
\end{lemma}
\begin{proof}
Identify $W[a]=B=\mathcal O/a$, a $\Z/l^r$-algebra; then $\alpha'\beta'=\chi\bmod l^r$. If $n'(\tau)c'(g)\ne0$ pick $\lambda\colon B\to\Z/l^r$ not killing it,
and put $w(g,h)=\lambda(n'(g)\alpha'(g)c'(h))$, a $2$-cocycle with values in $\mu$. By (i) it is inflated from $G_S$, $S$ the primes dividing $Nq$; by
(ii) it is a coboundary on $D_p$ for $p\in S\smallsetminus\{q\}$; by (PT), $w|_{D_{A_0}}=\partial b$. On $D=D_{A_0}$, $c',n'$ are homomorphisms,
$c'|_I=0$ and $w=\lambda(n'(g)c'(h))$. Let $\sigma$ be a Frobenius and $f(h)=\lambda(n'(h)c'(\sigma))$. For unramified $x$, $x(h)c'(\sigma)=x(\sigma)c'(h)$
(both agree on $I$ and at $\sigma$), so $f\cup x=x(\sigma)\partial b$ is a coboundary; by (UO) $f$ is unramified, so $f(\tau)=0$. Thus the homomorphism
$h\mapsto\lambda(n'(\tau)c'(h))$ vanishes on $I$ and at $\sigma$, hence at $g$: contradiction.
\end{proof}

\begin{lemma}[extending equivariant homomorphisms]\label{i5:x-lem:BC-ext}
If $\nu\colon G_F\to W$ is a continuous homomorphism with $\nu(gug^{-1})=\gamma(g)\nu(u)$ and $\gamma(\delta')-1$ is a unit, then
$n(g)=-(\gamma(\delta')-1)^{-1}\nu([g,\delta'])$ is a $\gamma$-cocycle extending $\nu$, with values in $\mathcal O\nu(G_F)$.
\end{lemma}
\begin{proof}
$[gh,\delta']=g[h,\delta']g^{-1}[g,\delta']$, and $\nu([u,\delta'])=-(\gamma(\delta')-1)\nu(u)$ for $u\in G_F$.
\end{proof}

\medskip\noindent\emph{Data.} In Lemmas~\ref{i5:x-lem:BC-step}--\ref{i5:x-lem:BC-cheb}, Proposition~\ref{i5:x-prop:BC-euler} and Theorem~\ref{i5:x-thm:BC-count}:
(D0) $F=N_0l^m$, $m\ge1$, $l\nmid N_0$; (D1) $\alpha\equiv1\bmod l^m$ on $G_F$; (D2) $\alpha(g)\equiv-1\bmod l^m$ if $g$ acts by $-1$ on $\mu_F$; (D3)
$\alpha(g_0)-\chi(g_0)\in\mathcal O^\times$ for some $g_0$. Then $\beta\equiv1\bmod l^m$ on $G_F$ and $\beta(g_0)-1$ is a unit. A homomorphism
$\nu\colon G_F\to W$ is \emph{$\gamma$-admissible of level $n$} if continuous, $\nu(gug^{-1})=\gamma(g)\nu(u)$, $l^m\nu=0$, $\nu|_{I_A}=0$ for $A\nmid n$.
Fix a prime $r\equiv1\bmod F$; a prime $q$ is \emph{admissible} if $q\equiv1\bmod F$, $q\ne r$, $l\nmid\mathrm{ord}(r\bmod q)$; $d$ is admissible if it is
a square-free product of admissible primes.

\begin{lemma}[Kolyvagin's local step]\label{i5:x-lem:BC-step}
Assume (PT), (UO), (D0)--(D3). Let $F\mid N$, $q\equiv1\bmod F$ prime, $\nu$ $\beta$-admissible of level $Nq$, $a_0\mid a\mid l^m$, $w_0$ with
annihilator $(a)$, $a\nu=0$, and $c\in\mathrm{ShT}(\alpha,\delta,F,a_0)$ with $c|_{D_A}=0$ for $A\mid N$, $A\nmid Fq$. Then for $A\mid q$, $g\in D_A$,
$\tau\in I_A$, $\nu(\tau)=yw_0$: $yc(g)=0$.
\end{lemma}
\begin{proof}
Extend $\nu$ by Lemma~\ref{i5:x-lem:BC-ext} ($\delta'=g_0$) and apply Lemma~\ref{i5:x-lem:BC-rec}: $q\ne l$; (i) since $I_A\subset G_F$ for $A\nmid Nq$ and
(T3); (ii) by (T4) or the hypothesis; (iii) since $D_A\subset G_F$ for $A\mid q$ ($q\equiv1\bmod F$).
\end{proof}

\begin{lemma}[elements of $\mathrm{ShT}$]\label{i5:x-lem:BC-shtnu}
Assume (D0)--(D3), $a\mid l^m$, $c\in\mathrm{ShT}(\alpha,\delta,F,a)$. (1) $c|_{G_F}$ is $\alpha$-admissible of level $F$; (2) $c|_{G_F}=0$ implies $c=0$;
(3) if $q\equiv1\bmod F$, $A\mid q$, $\varphi$ a Frobenius at $A$ and $c(\varphi)=0$, then $c|_{D_{A'}}=0$ for all $A'\mid q$.
\end{lemma}
\begin{proof}
(1) by (D1). (2) Expand $c(g\delta)=c([g,\delta]\delta g)$ to get $(\alpha(\delta)-1)c(g)=0$. (3) $D_A\subset G_F$ is generated by $I_A$ and $\varphi$;
conjugate.
\end{proof}

\begin{lemma}[Chebotarev choice]\label{i5:x-lem:BC-cheb}
Assume (D0)--(D3). Let $\nu$ be $\beta$-admissible, $c$ $\alpha$-admissible with $c|_{G_F}\ne0$, $T\ne0$. Then there are an admissible $q\nmid T$, $A\mid q$
and a Frobenius $\varphi$ at $A$ with $c(\varphi)\ne0$ and $\nu(G_F)=\mathcal O\nu(\varphi)$.
\end{lemma}
\begin{proof}
\emph{Kummer lemma.} Let $\varrho^{l^m}=r$, $H=G_F\cap\mathrm{Stab}(\varrho)$. If $f\colon G_F\to X$ is additive, $f(gug^{-1})=\gamma(g)f(u)$, $l^mf=0$ and
$f|_H=0$, then $(\gamma(g)-\chi(g))f=0$: $gug^{-1}\equiv u^{\chi(g)}\bmod H$ by the action on $\varrho$.
Let $B=\ker c|_{G_F}$, $u_0$ with $\nu(G_F)=\mathcal O\nu(u_0)$ (Lemma~\ref{i5:x-lem:BC-W}(iii)), $K=\{u:\nu(u)\in\mathfrak m\nu(u_0)\}$ ($K=B$ if
$\nu=0$), $C=G_{Fl}$. (a) $H\not\subset B$ by the Kummer lemma with $g_0$. (b) $H\not\subset K$ by the Kummer lemma for $\nu\bmod\mathfrak m\nu(u_0)$ and
$g$ acting by $-1$, using (D2). (c) $H\not\subset C$: for $u$ acting by $1+F$ on $\mu_{Fl}$ and $g$ by $-1$, $gug^{-1}u\in H$ acts by $(1+F)^2\not\equiv1$.
(d) $x^{l^m}\in B$. \emph{Three-subgroup lemma:} if $H$ lies in none of $A,B,C$ and $x^k\in B$ for $x\in H$, $k$ odd, then $H\not\subset A\cup B\cup C$
(take $x\in H\cap A\smallsetminus(B\cup C)$, $y\in H\cap B\smallsetminus(A\cup C)$; then $xy\in C$ and $x\cdot xy$ avoids all three, since $x^2\in B$
would give $x\in B$). So $V=H\smallsetminus(K\cup B\cup C)$ is open and non-empty, and Chebotarev (degree-one form, proved in the formalization) gives
$q\nmid TFlr$, $A$, $\varphi\in V$. Then $c(\varphi)\ne0$, $\nu(\varphi)$ generates, $q\equiv1\bmod F$, $l\nmid(q-1)/F$, and $\varphi\varrho=\varrho$ gives
$\mathrm{ord}(r\bmod q)\mid N_0(q-1)/F$: $q$ is admissible.
\end{proof}

\begin{lemma}[Stickelberger elements]\label{i5:x-lem:BC-X}
For $M\ge1$ let $K_M=\Q(\mu_M)$, $\sigma_t\colon\zeta\mapsto\zeta^t$, and for $b<M$ put
$\theta_{M,b}(t)=\lfloor\langle t\rangle b/M\rfloor-\lfloor\langle-t\rangle b/M\rfloor$, $\Theta_{M,b}=\sum_t\theta_{M,b}(t)\sigma_t^{-1}$. Let $\mathfrak a\ne0$ be an
ideal of $K_M$.
\begin{enumerate}
\item There is $X\in K_M^\times$ with $X\overline X=1$ and $(X)=\mathfrak a^{\Theta_{M,b}}$, unique up to $\mu_{2M}$.
\item (Norm relation) For a prime $q\nmid M$, $b<q$, $\mathfrak a'$ an ideal of $K_{qM}$, $X'$ as in (1) for $(qM,b,\mathfrak a')$ and $X$ for
$(M,b,N\mathfrak a')$: $N_{K_{qM}/K_M}X'=\zeta X/\sigma_q^{-1}X$, $\zeta\in\mu_{2M}$.
\item (Congruence) If moreover $\mathfrak a'$ is prime to $q$ and $b<M$, then $X'\equiv\zeta'\sigma_q^{-1}(X)$ modulo every prime above $q$, $\zeta'\in\mu_{2qM}$.
\end{enumerate}
\end{lemma}
\begin{proof}
Note $\theta_{M,b}(-t)=-\theta_{M,b}(t)$. (1) By Chebotarev in the Hilbert class field choose a degree-one prime $\mathfrak p\nmid M$ in the class of
$\mathfrak a^{-1}$, $p=N\mathfrak p$. With $\omega$ the $M$-th power residue character mod $\mathfrak p$ and $J=\prod_{j=1}^{b-1}J(\omega,\omega^j)$, Stickelberger's
factorisation gives $(\overline J/J)=\mathfrak p^{\Theta_{M,b}}$ and $|J|^2=p^{b-1}$. Writing $\mathfrak a\mathfrak p=(y)$ and
$Y=\prod_t\sigma_t^{-1}(y)^{\theta_{M,b}(t)}$, $X=YJ/\overline J$ works. Uniqueness: a quotient is a unit with all conjugates of absolute value $1$, a root
of unity (Kronecker). (2) The distribution relation $\sum_{t'\mapsto t}\theta_{qM,b}(t')=\theta_{M,b}(t)-\theta_{M,b}(q^{-1}t)$ ($b<q$) shows both sides generate
$(N\mathfrak a')^{(1-\sigma_q^{-1})\Theta_{M,b}}$; conclude by Kronecker.
(3) Fix $\mathfrak Q\mid q$ in $K'=K_{qM}$, reduction $\varrho$ to $k=\mathcal O_{K'}/\mathfrak Q$, $K=K_M$, $N=N_{K'/K}$, $g=\sigma_q^{-1}$. Use (F1) $a^q=c^q$
implies $a\equiv c$ ($q$-th powers are injective on $k$) and (F2) $g(\xi)^q=\xi$ for $\xi\in\mu_M$.
\emph{(a) Jacobi sums.} For a degree-one prime $\mathfrak P\nmid qM$ of $K'$ with $\mathfrak P_0=\mathfrak P\cap K$, the residue characters satisfy
$\omega'^q=\omega$ (both are $\equiv x^{(p-1)/M}$ mod $\mathfrak P$ in $\mu_{qM}$). By (F1), (F2), $\omega'(x)\equiv g(\omega(x))$ and the same for inverses,
hence $J'\equiv g(J)$, $\overline{J'}\equiv g(\overline J)$; $J'$ divides $p^{b-1}\notin\mathfrak Q$, so $\overline{J'}/J'\equiv g(\overline J/J)$, and by uniqueness
in (1), $X_{\mathfrak P}\equiv\zeta_3g(X_{\mathfrak P_0})$, $\zeta_3\in\mu_{2qM}$.
\emph{(b) Principal ideals.} For $y$ prime to $q$, $\varrho\circ\tau$ depends only on $\tau|_{\mu_M}$ (as $\mathcal O_{K'}=\Z[\zeta_{qM}]$ and (F1)), so
$f(t)=\varrho(\sigma_{t'}^{-1}y)\ne0$ is well defined; for $s$ above $u=q\bmod M$, $\varrho\circ\sigma_s=\mathrm{Frob}_q\circ\varrho$, giving $f(t)=f(ut)^q$.
The distribution relation then yields $\varrho(\prod_{t'}\sigma_{t'}^{-1}(y)^{\theta_{qM,b}(t')})=\prod_t(f(ut)^{q-1})^{\theta_{M,b}(t)}$, and since
$g\sigma_t^{-1}(Ny)$ is a product of $q-1$ conjugates each reducing to $f(ut)$,
$\prod_{t'}\sigma_{t'}^{-1}(y)^{\theta_{qM,b}(t')}\equiv g(\prod_t\sigma_t^{-1}(Ny)^{\theta_{M,b}(t)})$, both $\mathfrak Q$-units.
\emph{(c) Combination.} Choose $\mathfrak P\nmid qM$ of degree one in the class of $\mathfrak a'^{-1}$, $\mathfrak P\mathfrak a'=(y)$, and let $A$, $B$ be the two
products of (b). Then $A/X'$, $B/X$ are as in (1) for $\mathfrak P$, $\mathfrak P_0$, so (a) gives $A/X'\equiv\zeta_3g(B)/g(X)$ and (b) gives $A\equiv g(B)$.
As $X'$ and $A$ are $\mathfrak Q$-units, so are $g(B)$, $g(X)$, and the identity
$x-\zeta_3^{-1}z=x(A-B_1)/A-(A/x-\zeta_3B_1/z)xz\zeta_3^{-1}/A$ ($x=X'$, $z=g(X)$, $B_1=g(B)$) gives (3) with $\zeta'=\zeta_3^{-1}$.
\end{proof}

\begin{proposition}[the Euler system of Stickelberger elements]\label{i5:x-prop:BC-euler}
Assume (D0)--(D3), $r\equiv1\bmod F$ prime, $A_0\mid r$, $b<F$, lifts $\sigma_t\in G$ of $t\in(\Z/F)^\times$, $w\in W$ with $l^mw=0$, and
$z_b=\sum_t\theta_{F,b}(t)\alpha(\sigma_t)^{-1}$. Then there are $\nu_d\colon G\to W$, $d$ admissible, with
(E1) $\nu_d|_{G_F}$ $\beta$-admissible of level $Frd$;
(E2) for $qd$ admissible, $q\nmid d$, $A\mid q$, $\varphi$ a Frobenius at $A$: $\nu_{qd}(\tau)=u\nu_d(\varphi)$ for some $\tau\in I_A$, $u\in\mathcal O^\times$;
(E3) $\nu_1(\tau_0)=u_0z_bw$ for some $\tau_0\in I_{A_0}$, $u_0\in\mathcal O^\times$.
\end{proposition}
\begin{proof}
Put $M=l^m$, $\zeta\in\mu_M$ primitive.
\emph{Step 1.} For admissible $d$ let $\mathfrak r_d$ be the prime of $K_{Fd}$ below $A_0$ and $X(d)$ as in Lemma~\ref{i5:x-lem:BC-X}(1) for
$(Fd,b,\mathfrak r_d)$: a unit away from $r$, with $\mathrm{ord}_{A_0}(\sigma_tX(1))=\theta_{F,b}(t)\mathrm{ord}_{A_0}(r)$. For $qd$ admissible,
$N\mathfrak r_{qd}=\mathfrak r_d^s$ with $s\mid\mathrm{ord}(r\bmod qFd)$ prime to $l$, and Lemma~\ref{i5:x-lem:BC-X}(2),(3) give, for $\tau$ generating
$\mathrm{Gal}(K_{Fqd}/K_{Fd})$ and $\varphi$ acting on $\mu_{Fd}$ by $q$, $\prod_{i<q-1}\tau^iX(qd)=\zeta_1(X(d)/\varphi^{-1}X(d))^s$ and
$X(qd)\equiv\zeta_2\varphi^{-1}(X(d))^s$ modulo $A\mid q$, $\zeta_i\in\mu_{2Fqd}$.
\emph{Step 2.} With $g_0\zeta=\zeta^{a_*}$, $P(u)=g_0(u)/u^{a_*}$ is multiplicative, Galois-equivariant on abelian extensions and kills $\mu_M$. Put
$x(d)=P(X(d))^{2N_0d}$; the relations become exact: $N_qx(qd)=(x(d)/\varphi^{-1}x(d))^{qs}$, $x(qd)\equiv\varphi^{-1}(x(d))^{qs}$ mod $A\mid q$.
\emph{Step 3.} Fix $\tau_q$ trivial on $\mu_k$ for $q\nmid k$ and generating on $\mu_q$, $D_qu=\prod_{i=0}^{q-2}\tau_q^i(u)^i$, $D_d=\prod_{q\mid d}D_q$. Then
$\tau_q(D_qu)/D_qu=u^{q-1}/N_qu$, and by induction on the number of prime factors, $g(D_dx(d))/D_dx(d)\in K_{Fd}^{\times M}$ for $g\in G_F$, since
$(\tau_q-1)D_dx(d)=D_{d/q}(x(d))^{q-1}((\varphi_q^{-1}-1)D_{d/q}x(d/q))^{qs}$ and $M\mid q-1$.
\emph{Step 4 (Kummer classes).} With $Y_d^M=D_dx(d)$ and $z_g^M=g(D_dx(d))/D_dx(d)$, $\zeta^{\kappa_d(g)}=g(PY_d)/(PY_d\,P(z_g))$ defines a continuous
homomorphism $\kappa_d\colon G_F\to\Z/M$, unramified outside $Frd$. \emph{Kolyvagin's formula:} for $qd$ admissible, $A\mid q$, $t\in I_A$ acting on $\mu_q$
as $\tau_q$, $\varphi$ a Frobenius at $A$: $\kappa_{qd}(t)=-qs\kappa_d(\varphi)$. (Here $tD_{qd}x(qd)/D_{qd}x(qd)=(w^kz^{qs})^M$ with $w=D_dx(qd)$,
$q-1=Mk$, $z^M=\varphi^{-1}(D_dx(d))/D_dx(d)$; reduce mod $A$ using Step 2, that $t$ is trivial on $A$-units and $\varphi^{-1}(V)^q\equiv V$, and that
$\mu_M$ injects mod $A$.) \emph{Initial value:} with $\varrho^M=r$, $t\in I_{A_0}$, $t\varrho=\zeta\varrho$, and $z_g=1$ for $d=1$, the tame symbol gives
$\kappa_1(\sigma_c^{-1}t\sigma_c)c\equiv K(c)\bmod M$, $K(c)=2N_0(\theta'(c\bar g_0)-a_*\theta'(c))$, $\theta'(c)=\theta_{F,b}(c\bar g_0)-a_*\theta_{F,b}(c)$.
\emph{Step 5.} Put $\nu_d(u)=\sum_t\kappa_d(\sigma_t^{-1}u\sigma_t)\beta(\sigma_t)w$; it is $\beta$-equivariant, giving (E1). (E2) follows from Kolyvagin's
formula applied at the conjugate places $\sigma_c^{-1}A$. (E3): as $\beta(\sigma_c)\equiv c\alpha(\sigma_c)^{-1}$, $\nu_1(t)=Z(K)w$ with
$Z(k)=\sum_ck(c)\alpha(\sigma_c)^{-1}$, and $Z(c\mapsto k(c\bar g_0)-a_*k(c))\equiv(\alpha(g_0)-a_*)Z(k)$ by (D1); twice this gives
$u_0=2N_0(\alpha(g_0)-a_*)^2$, a unit by (D3) and $l\nmid2N_0$.
\end{proof}

\begin{theorem}[the Euler-system bound]\label{i5:x-thm:BC-count}
Assume (PT), (UO), (D0)--(D3), $b<F$, lifts $\sigma_t$, $z=z_b$, $\alpha(\delta)-1\in\mathcal O^\times$ and $a_0\ne0$ with $a_0z\mid l^m$. Then
$\lgth\mathrm{ShT}(\alpha,\delta,F,a_0)\le v(z)$.
\end{theorem}
\begin{proof}
Write $zk=l^m$, so $a_0\mid k$; choose $w$ with annihilator $(l^m)$, $r$, $A_0$ and $\nu_d$ as in Proposition~\ref{i5:x-prop:BC-euler}; let $(a_d)$ be the
annihilator of $\nu_d(G_F)$. $S_n$ refers to $\mathrm{ShT}(\alpha,\delta,F,a_0)$.
\emph{Key step.} Let $n=r\cdot$(admissible $d$), $qd$ admissible, $A,\varphi,\tau,u$ as in (E2), $a\in(a_0)$, i.e.\ $a_0\mid a$, where $(a_0)$ is the annihilator of $\nu_{qd}(\tau)$. Then
$a_{qd}=ae$, (1) $ec(\varphi)=0$ for $c\in S_n$, and (2) $\lgth S_n\le\lgth S_{nq}+v(e)$. For (1), Lemma~\ref{i5:x-lem:BC-step} gives $yc(\varphi)=0$ where
$\nu_{qd}(\tau)=yw_0$, $W[a_{qd}]=\mathcal Ow_0$; if $v(y)\le v(a_{qd})$ then $e\in y\mathcal O$, otherwise $\nu_{qd}(\tau)=0$, $a_0$ is a unit and $c=0$. For (2),
$c\mapsto c(\varphi)$ is killed by $e$ with kernel in $S_{nq}$ (Lemma~\ref{i5:x-lem:BC-shtnu}(3)); apply Lemma~\ref{i5:x-lem:BC-W}(iv).
\emph{Base.} By (E3) the same argument with $n=1$, $q=r$, $\nu_1$, $A_0$ gives $a_1=ke_0$ and $\lgth S_1\le\lgth S_r+v(e_0)$.
\emph{Induction.} Let $\mathcal L$ be the admissible $d$ with $a_0\mid a_d$, $j(d)=v(a_d)\le J=v(l^m)$. Claim: $\lgth S_{rd}+j(d)\le J$ on $\mathcal L$. If
$0\ne c\in S_{rd}$, $c|_{G_F}$ is non-zero and $\alpha$-admissible (Lemma~\ref{i5:x-lem:BC-shtnu}); Lemma~\ref{i5:x-lem:BC-cheb} ($T=d$) gives $q$, $A$,
$\varphi$ with $c(\varphi)\ne0$ and $\nu_d(G_F)=\mathcal O\nu_d(\varphi)$, so by (E2) the key step applies with $a=a_d$: $v(e)>0$, $qd\in\mathcal L$, and
$\lgth S_{rd}+j(d)\le\lgth S_{rqd}+j(qd)$ with $j(qd)>j(d)$. Descending induction on $J-j(d)$ proves the claim.
\emph{Conclusion.} At $d=1$: $\lgth S_r+v(k)+v(e_0)\le v(z)+v(k)$, so $\lgth\mathrm{ShT}\le v(z)$.
\end{proof}

\begin{lemma}[twisted Stickelberger values]\label{i5:x-lem:BC-zcong}
Let $l^m\mid F$, $m,e\ge1$, $k=e+1$, $b$ prime to $F$, $\psi$ a character modulo $F$ with values in $\overline\Q_l$, and $B_{k,\psi}$ computed at level $F$.
Then
$\big|\mathrm{dk}(k)\big(b^e\psi(b)\sum_t\theta_{F,b}(t)\langle t\rangle^e\psi(t)-2(\psi(b)b^k-1)B_{k,\psi}/k\big)\big|_l\le|l^m|_l$.
\end{lemma}
\begin{proof}
With $E'(x)=F^{k-1}B_k(x/F)/k$ and $\langle t\rangle b=\langle bt\rangle+Fj$, $\theta_{F,b}(t)=2j-(b-1)$, and the expansion of
Lemma~\ref{i5:x-lem:BC-ylevel}(2) gives
$\mathrm{dk}(k)[b^e\langle t\rangle^e\theta_{F,b}(t)+2(E'(\langle bt\rangle)-b^kE'(\langle t\rangle))]\in F\Z$. Multiply by $\psi(bt)$ and sum, using
$\sum_t\psi(bt)E'(\langle bt\rangle)=B_{k,\psi}/k$ and $\sum_t\psi(bt)E'(\langle t\rangle)=\psi(b)B_{k,\psi}/k$.
\end{proof}

\begin{lemma}[choice of auxiliary data]\label{i5:x-lem:BC-zeval}
Let $\psi$ be primitive of conductor $f=f_0l^s$ ($l\nmid f_0$), values in $\mathcal O$, order prime to $l$, $\psi(-1)=(-1)^{e+1}$, $e\ge1$; put
$\alpha=(\chi^e\psi)^{-1}$, assume (D3), and let $a=\mathrm{den}(f,e+1)B_{e+1,\psi}$. For every $a_0\ne0$ there are $m\ge\max(1,s)$, $F=f_0l^m$, $b<F$ and lifts
$\sigma_t$ with (D0)--(D3), $a_0z_b\mid l^m$ and $z_b\mid a$.
\end{lemma}
\begin{proof}
$T=|B_{e+1,\psi}/(e+1)|_l>0$ by parity; choose $m$ with $|l^m|<|a_0||\mathrm{dk}(e+1)|T$. (D1), (D2) follow from $\chi\equiv1\bmod l^m$ on $G_F$ and
$\chi^e\psi(g)\equiv(-1)^e(-1)^{e+1}$ for $g$ acting by $-1$. Take $b=\langle g_0\bmod F\rangle$. Then $z_b\equiv\sum_t\theta_{F,b}(t)\langle t\rangle^e\psi(t)$,
$\psi(b)b^{e+1}-1$ is a unit multiple of $\alpha(g_0)-\chi(g_0)$, and the level-$F$ Bernoulli number is $B_{e+1,\psi}(1-\psi(l)l^e)^{[l\nmid f]}$, so
Lemma~\ref{i5:x-lem:BC-zcong} and the ultrametric inequality ($l$ odd) give $|z_b|=T$. Hence $a_0z_b\mid l^m$ and $|a|\le T$.
\end{proof}

\begin{proof}[Proof of Theorem~\ref{i5:x-thm:thmbc}]
\emph{$\mathrm{PsiBound}\Rightarrow\mathrm{ThmBcFree}$.} Let $M$, $e,t\ge2$, $\mathrm{PP}(m_M;q,d)$, $E$ a CM $(M,e,q,d)$-field, $F=\mathfrak f(E)$, $n=[E:\Q]$. An
odd $l\mid n$ divides $\deg M$, since $\sigma_a^{2^{e+1}\deg M}$ fixes $E$ for all $a$ prime to $m_M$. Finiteness for $0<e'<t$ is
Proposition~\ref{i5:x-prop:stickbound} or \ref{i5:x-prop:primebound}. Over $l\nmid n$, Proposition~\ref{i5:x-prop:primebound} and
$\log\mathrm{bnd}(F,e'+1)\le(2t+5)\log(tF)$ give at most $3nt^2\log(tF)$; each $l\mid\deg M$ contributes at most $Cln^2t^2\log(tlF)^2$ by
Proposition~\ref{i5:x-prop:stickbound}. The total is $\le\beta_C(\deg M,n,F,t)$.

\emph{(PT)$+$(UO)$\Rightarrow\mathrm{PsiBound}$.} As $l\nmid[E:\Q]$, the $\psi$-part is $\mathrm{Sel}(\alpha)$ with $\alpha=(\psi\rho)^{-1}\chi^{-e}$ (restriction to
$G_E$ and isotypic parts, as in step~(c) of Proposition~\ref{i5:prop-red}). Passing to the primitive character changes $a_{e+1}$ by $N/f$
(Lemma~\ref{i5:x-lem:bern-int}), so let $\psi$ denote $\psi\rho$, primitive of conductor $f$, order prime to $l$, $\psi(-1)=(-1)^{e+1}$, and
$a=a_{e+1}(\psi)$. Values $\psi(g)\ne1$ have $\psi(g)-1$ a unit, and each $A\mid f$ has $g\in I_A$ with $\psi(g)\ne1$; so $\alpha(I_A)=1$ for $A\nmid fl$
and the last hypothesis of Lemma~\ref{i5:x-lem:BC-shasht} holds.
\emph{Exceptional case} ($\alpha\equiv\chi\bmod\mathfrak m$): for $A\nmid l$, $g\in I_A$, $\psi(g)\equiv1$, hence $\psi(g)=1$; Lemma~\ref{i5:x-lem:BC-exc} gives
$\mathrm{Sel}(\alpha)=0$.
\emph{Main case} ((D3) holds): Lemma~\ref{i5:x-lem:BC-zeval} with $a_0=1$ gives $F_1$; take $\delta$ acting by $-1$ on $\mu_{F_1}$, so $\alpha(\delta)-1$ is a
unit by (D2). For $a_0\ne0$, Lemma~\ref{i5:x-lem:BC-zeval} and Theorem~\ref{i5:x-thm:BC-count} give
$\lgth\mathrm{ShT}(\alpha,\delta,F,a_0)\le v(z_b)\le v(a)$, and $\mathrm{ShT}(\alpha,\delta,F,a_0)=\mathrm{ShT}(\alpha,\delta,fl,a_0)$ ($F$ and $fl$ have the same prime
divisors). Lemma~\ref{i5:x-lem:BC-shasht} gives $\lgth\mathrm{Sel}(\alpha)\le v(a)$.
\end{proof}

\begin{remark}[why no main conjecture is needed]\label{i5:x-rem:minus}
Under Proposition~\ref{i5:prop-shaleV}, the $\psi$-part of the minus part is dual to the $\psi^{-1}$-part of $\Sha^2(\mathcal O_E[1/l],\Z_l(e+1))$, in the
parity where $B_{e+1,\psi}\ne0$. The \emph{equality} of its order with $\#\mathcal O/(B_{e+1,\psi}/(e+1))$ (up to Euler factors) follows from the main
conjecture (Mazur--Wiles 1984; Wiles 1990, \emph{Ann. of Math.} 131). Theorem~\ref{i5:x-thm:thmbc} needs only the \emph{upper bound}, which Kolyvagin's
method gives directly: the argument above is the Euler system of Jacobi-sum Stickelberger elements, projected to the $\chi\alpha^{-1}$-eigenspace mod
$l^m$, as in Kolyvagin (1990, Grothendieck Festschrift II) and Rubin (1990, appendix to Lang, GTM 121; 2000, Ann.\ of Math.\ Studies 147). Its only
duality inputs are (PT) and (UO); Stickelberger and Chebotarev are proved in the formalization; the twisted form (coefficients
$\mathcal O(\psi^{-1}\chi^{-e})$, strict at $l$, normalisation $\mathrm{den}(N,e+1)$) is proved above rather than cited. Coates--Sinnott-type
annihilation (1974) enters only in Proposition~\ref{i5:x-prop:stickbound}, for $l\mid[E:\Q]$, where a crude bound suffices.
\end{remark}

\medskip\noindent\emph{Notation for the proof of Theorem~\ref{i5:x-thm:thmbab}.}
For a number field $F$ put $\operatorname{sha}(F,l,e)=\lgth_{\Z_l}\Sha^1_D(F,\Q_l/\Z_l(-e))$ for any $l$-datum $D$ (independent of $D$: changing $A_v$
conjugates $D_v$). For a $\Z_l$-module $X$ and $f\in\mathrm{End}(X)$ let $\operatorname{cl}(f)=\lgth(X/fX)$; write $f\in l^k$ if $f(X)\subset l^kX$;
$\gamma$ is \emph{unipotent mod $l$} if $(\gamma-1)^N\in l$ for some $N$; $X$ is \emph{$R$-generated} if $\Z_l^R\twoheadrightarrow X$; $u\in\Z_l$ has
\emph{level $c$} if $v_l(u-1)=c$ ($c\ge2$ if $l=2$). Put $\operatorname{tw}(e,m)=\operatorname{cl}((u^e\gamma)^{l^m}-1)$ and $j_0=j_0(l,R)=\lfloor\log_lR\rfloor+3$.

\begin{lemma}[elementary facts]\label{i5:x-lem:bab-elem}
Let $X$ be $R$-generated, $g,w\in\mathrm{End}(X)$. (1) $\operatorname{cl}(l)\le R$. (2) If $a\mid b$, $\operatorname{cl}(g^a-1)\le\operatorname{cl}(g^b-1)$. (3) If $g$ is
unipotent mod $l$, $(g-1)^R\in l$. (4) $(g-1)^{l^k}\in l\iff g^{l^k}-1\in l$. (5) $w-1\in l^k\Rightarrow w^n-1\in l^k$; $w-1\in l\Rightarrow w^l-1\in l^2$.
(6) If $g$ is unipotent mod $l$ and $l\mid a-1$, $ag$ is unipotent mod $l$.
\end{lemma}
\begin{proof}
(1) $X/lX$ is a quotient of $\mathbb F_l^R$. (2) $g^a-1$ divides $g^b-1$. (3) $\bar g-1$ is nilpotent on a space of dimension $\le R$.
(4) $(T-1)^{l^k}\equiv T^{l^k}-1\bmod l$. (5) Factor $w^n-1$; expand $w^l-1=\sum_i\binom li(w-1)^i$. (6) Expand $(a(g-1)+(a-1))^N$.
\end{proof}

\begin{lemma}\label{i5:x-lem:bab-step}
If $X$ is $R$-generated and $w-1\in l^2$, then $\operatorname{cl}(w^l-1)\le\operatorname{cl}(w-1)+R$.
\end{lemma}
\begin{proof}
$w^l-1=(w-1)\Phi$, $\Phi=\sum_{i<l}w^i=l+D$ with $D\in l^2$, so $\operatorname{cl}(w^l-1)\le\operatorname{cl}(w-1)+\lgth(X/\Phi X)$. From $lX\subset\Phi X+l\cdot lX$
and closedness of $\Phi X$, $lX\subset\Phi X$, so $\lgth(X/\Phi X)\le R$.
\end{proof}

\begin{lemma}[growth along the tower]\label{i5:x-lem:bab-growth}
If $X$ is $R$-generated and $g$ unipotent mod $l$, then $\operatorname{cl}(g^{l^m}-1)\le\operatorname{cl}(g^{l^{j_0}}-1)+R(m-j_0)^+$. Hence, if $u$ has level $c$,
$\operatorname{tw}(e,m)\le\operatorname{tw}(e,j_0)+R(m-j_0)^+$.
\end{lemma}
\begin{proof}
With $k=j_0-2$, $R\le l^k$, Lemma~\ref{i5:x-lem:bab-elem}(3)--(5) give $g^{l^n}-1\in l^2$ for $n\ge k+1$; iterate Lemma~\ref{i5:x-lem:bab-step} from
$n=j_0$, and use (2) for $m\le j_0$. For the second claim, $u^e\gamma$ is unipotent mod $l$ by (6).
\end{proof}

\begin{lemma}[exceptional centres]\label{i5:x-lem:bab-centres}
Let $X$ be $R$-generated, $\gamma$ unipotent mod $l$, $u$ of level $c$. There are $s_1,\dots,s_R\in\Z_l$ such that $|s_i-e|_l\ge l^{-B}$ for all $i$
implies $\operatorname{tw}(e,m)\le R(\max(m,j_0)+c+B)$.
\end{lemma}
\begin{proof}
\emph{Matrix lift.} For $h\in\mathrm{End}(X)$ choose $S\in M_R(\Z_l)$ with $\pi S=h\pi$; then $\operatorname{cl}(ah-1)\le v_l\det(aS-1)=\sum_iv_l(az_i-1)$ (Smith
normal form; $z_i$ the eigenvalues of $S$ in $\mathbb C_l$).
\emph{Lifting the exponent.} $v_l(u^n-1)=c+v_l(n)$, so with $u_j=u^{l^j}$, $|u_j^a-u_j^b|=l^{-(j+c)}|a-b|$ for $a,b\in\Z_l$ (continuity of
$x\mapsto u_j^x$).
\emph{One centre per eigenvalue.} For $z\in\mathbb C_l$ there is $s\in\Z_l$ with $l^{-(j+c)}|s-e|\le|u_j^ez-1|$ for all $e$: trivial if $|z|\ne1$;
otherwise take $s$ minimising $|u_j^xz-1|$ on $\Z_l$ and use the ultrametric inequality.
Apply this to $h=\gamma^{l^{j_0}}$, $a=u_{j_0}^e$: $\operatorname{tw}(e,j_0)\le\sum_i(j_0+c+v_l(s_i-e))$; conclude with Lemma~\ref{i5:x-lem:bab-growth}.
\end{proof}

\begin{definition}[Iwasawa module of a tower]\label{i5:x-def:bab-tower}
Let $L_0\supset\mu_l$ ($\mu_4$ if $l=2$), $L_\infty=L_0(\mu_{l^\infty})$. For $N\in\{1,l\}$, $D_\infty^{(N)}$ is the module of continuous
$\varphi\colon G_{L_\infty}\to\Q_l/\Z_l$ vanishing on all inertia groups and, if $N=l$, on all decomposition groups above $l$;
$X_\infty^{(N)}=\mathrm{Hom}(D_\infty^{(N)},\Q_l/\Z_l)$, with $G_{L_0}$ acting by conjugation, twisted by $\chi_{\mathrm{cyc}}^e$ on the $e$-th twist.
$\mathrm{IsTower}(l,L_0,R)$: $L_0$ is totally complex, contains $\mu_l$ ($\mu_4$), $R\ge1$, and $X_\infty^{(l)}$ is $R$-generated.
\end{definition}

\begin{lemma}[procyclicity]\label{i5:x-lem:bab-procyc}
A closed subgroup $V\le\Z_l^\times$ (with $V\subset1+4\Z_2$ if $l=2$) is topologically generated by one element.
\end{lemma}
\begin{proof}
$\Z_l^\times=\mu_{l-1}\times(1+l\Z_l)$ with $1+l\Z_l\cong\Z_l$ ($1+4\Z_2\cong\Z_2$); $V$ is the product of a cyclic prime-to-$l$ part and a closed
subgroup of $\Z_l$.
\end{proof}

\begin{lemma}[cocycles over a procyclic quotient]\label{i5:x-lem:bab-cob}
Let $\chi\colon\Gamma\to\Z_l^\times$ be continuous on a profinite $\Gamma$, $M$ a discrete $\Gamma$-module with $\ker\chi$ acting trivially, and $\gamma$
with $\chi(\Gamma)\subset\overline{\chi(\gamma)^{\Z}}$ and $\gamma-1$ surjective on $M$. Then every continuous cocycle $z$ with $z(\ker\chi)=0$ is a
coboundary.
\end{lemma}
\begin{proof}
Subtract $(g-1)m$ with $(\gamma-1)m=z(\gamma)$; the new $z'$ vanishes on $\gamma^{\Z}$ and $\ker\chi$, and is constant on cosets of $\ker\chi$, so its
(closed) zero set $Z$ has $\chi(Z)\supset\overline{\chi(\gamma)^\Z}\supset\chi(\Gamma)$, i.e.\ $Z=\Gamma$.
\end{proof}

\medskip\noindent\emph{Notation.} For $F$ a number field, $\chi=\chi_{\mathrm{cyc}}$, $M_e=\Q_l/\Z_l(-e)$, $U=G_{F(\mu_{l^\infty})}=\ker\chi$, let $\mathcal H_F(e)$ be the
module of continuous homomorphisms $\varphi\colon U\to\Q_l/\Z_l$ with (i) $\varphi(U\cap I_w)=0$ for all finite $w$; (ii) $\varphi(U\cap D_w)=0$ for $w\mid l$;
(iii) $\varphi(huh^{-1})=\chi(h)^{-e}\varphi(u)$ for $h\in G_F$.

\begin{lemma}[inflation--restriction with local conditions]\label{i5:x-lem:bab-infres}
Assume $e\ne0$ and $\mu_4\subset F$ if $l=2$. For an $l$-datum $D$ on $F$, inflation followed by restriction to $U$ maps $\Sha^1_D(F,M_e)$
isomorphically onto $\mathcal H_F(e)$; so $\operatorname{sha}(F,l,e)=\lgth\mathcal H_F(e)$.
\end{lemma}
\begin{proof}
(a) \emph{Inflation} from $G_{F,S}=G_F/N_S$ is injective with image the classes vanishing on all $I_w$, $w\nmid l$: such a cocycle vanishes on all
conjugates of these $I_w$ (as $z(hnh^{-1})=hz(n)$), hence on $N_S$, hence factors. Local conditions at conjugate places are equivalent. So
$\Sha^1_D$ is identified with the group $\Sha_F(e)$ of classes vanishing on $D_w$ ($w\mid l$) and $I_w$ ($w\nmid l$ finite). This step uses no
hypothesis on $e$ or $\mu_4$.
(b) For $z$ in $\Sha_F(e)$, $z|_U$ is a homomorphism satisfying (iii) by the cocycle rule, and (i), (ii) since coboundaries vanish on elements of $U$.
(c) \emph{Injectivity.} By Lemma~\ref{i5:x-lem:bab-procyc}, $\chi(G_F)=\overline{\chi(\gamma)^\Z}$; $\chi(\gamma)^{-e}\ne1$ as $\chi(G_F)$ is infinite; so
Lemma~\ref{i5:x-lem:bab-cob} applies.
(d) \emph{Surjectivity.} Let $C=\overline{\gamma^\Z}$, so $G_F=CU$. Write $C=C_l\times C'$; $\chi|_{C_l}$ is injective (infinite image in
$\Z_l$), so $C\cap U\subset C'$, which $\varphi$ kills. Then $z(cu)=c\varphi(u)$ is well defined, continuous ($C\times U\to G_F$ is a quotient map) and,
by (iii), a cocycle extending $\varphi$. It vanishes on $I_w\subset U$ for $w\nmid l$ by (i), and is a coboundary on $D_w$, $w\mid l$, by
Lemma~\ref{i5:x-lem:bab-cob} applied to $D_w$ (with $\chi(D_w)$ infinite and (ii)). By (a) its class lies in $\Sha^1_D$.
\end{proof}

\begin{proposition}[descent]\label{i5:x-prop:bab-descent}
Let $\mu_l$ ($\mu_4$) $\subset L_0$ and $c_0$ maximal with $\mu_{l^{c_0}}\subset L_0$.
(1) $\chi(G_{L_0})=1+l^{c_0}\Z_l$; so some $\gamma\in G_{L_0}$ has $u=\chi(\gamma)$ a topological generator, of level $c_0$.
(2) If $X_\infty^{(l)}$ is finitely generated, $\gamma$ is unipotent mod $l$ on it.
(3) For $e\ge1$ and every $l$-layer $L_1/L_0$ of height $m$, $\operatorname{sha}(L_1,l,e)=\operatorname{tw}(e,m)$ on $X=X_\infty^{(l)}$.
\end{proposition}
\begin{proof}
(1) $\chi(G_{L_0})$ is a closed subgroup of $1+l^{c_0}\Z_l\cong\Z_l$ not inside $1+l^{c_0+1}\Z_l$. (2) The action factors through the pro-$l$ group
$\mathrm{Gal}(L_\infty/L_0)$, whose image in $\mathrm{Aut}(X/lX)$ is an $l$-group. (3) For $F=L_1$, $U=G_{L_\infty}$ and (i), (ii) define
$D_\infty^{(l)}$; by Lemma~\ref{i5:x-lem:bab-infres}, $\operatorname{sha}(L_1,l,e)=\lgth(D_\infty^{(l)}(-e))^{G_{L_1}}$. Stabilisers are open, so the
invariants are the kernel of $\gamma^{l^m}\star-1$ ($\gamma^{l^m}$ generates $\mathrm{Gal}(L_\infty/L_1)$). On $X$, $\gamma$ acts as $u^e\gamma$ in the
$e$-th twist, and $(u^e\gamma)^{l^m}-1$ is the transpose of $\gamma^{-l^m}\star-1$; by Pontryagin duality kernel and cokernel have the same length.
\end{proof}

\begin{definition}[$\mathrm{ABound}$, $\mathrm{SecondBound}$]\label{i5:x-def:abound}
$\mathrm{SecondBound}(l,R,L_0)$: for $e\ge1$, $l$-layers $L_1/L_0$ of height $m$ and $L_2/L_0$ of height $j_0(l,R)$, and an $l$-datum $D$ on $L_2$,
$\operatorname{sha}(L_1,l,e)\le V_D(e)+R(m-j_0)^+$.
$\mathrm{ABound}(l,R,L_0)$: $\mathrm{SecondBound}(l,R,L_0)$, and there are a finite $K/\Q_l$, $c$ with $\mu_{l^c}\subset L_0$, $R'\le R$ and $s_i\in K$
($i\le R'$) with $\operatorname{sha}(L_1,l,e)\le R'(\max(m,j_0)+c+B)$ whenever $|s_i-e|\ge l^{-B}$ for all $i$.
\end{definition}

\begin{proposition}[tower bound]\label{i5:x-prop:abound}
Assume PTD, local finiteness and local Tate duality. Then $\mathrm{IsTower}(l,L_0,R)\Rightarrow\mathrm{ABound}(l,R,L_0)$.
\end{proposition}
\begin{proof}
Proposition~\ref{i5:x-prop:bab-descent} gives $\gamma,u,c=c_0$ and $\operatorname{sha}(L_1,l,e)=\operatorname{tw}(e,m)$. The centres come from
Lemma~\ref{i5:x-lem:bab-centres} ($K=\Q_l$, $R'=R$). SecondBound: Lemma~\ref{i5:x-lem:bab-growth} gives
$\operatorname{tw}(e,m)\le\operatorname{sha}(L_2,l,e)+R(m-j_0)^+$, and $\operatorname{sha}(L_2,l,e)\le V_D(e)$ by Proposition~\ref{i5:prop-shaleV} ($L_2$ totally complex).
\end{proof}

\medskip\noindent\emph{Notation.} For closed $V\le G_F$ and a set $\Sigma$ of primes, $\mathrm{Ch}_l(V,\Sigma)$ is the group of continuous $\varphi\colon V\to\Z/l$ with
$\varphi(V\cap I_w)=0$ for all finite $w$ of residue characteristic $\notin\Sigma$. For $V=G_{F(\mu_{l^\infty})}$, $\mathrm{Ch}_l(V,\emptyset)\cong
D_\infty^{(1)}[l]\supseteq D_\infty^{(N)}[l]$.

\begin{lemma}[adding a tame prime]\label{i5:x-lem:bab-tame}
Let $p\ne l$ with $I_w\subset V$ for $w\mid p$, and the places above $p$ in at most $k$ $V$-orbits. If $\#\mathrm{Ch}_l(V,\Sigma)\le n$ then
$\#\mathrm{Ch}_l(V,\Sigma\cup\{p\})\le nl^k$.
\end{lemma}
\begin{proof}
A homomorphism $I_w\to\Z/l$ factors through the $\Z_l(1)$-quotient of tame inertia, so is determined by its value at one $\gamma_w$. The map
$\varphi\mapsto(\varphi(\gamma_{w_i}))_{i\le k}$ (orbit representatives) has fibres that are cosets of $\mathrm{Ch}_l(V,\Sigma)$, since a homomorphism is
conjugation invariant.
\end{proof}

\begin{lemma}[index two]\label{i5:x-lem:bab-idx2}
Let $U\le V$ be open of index $\le2$ in a profinite $V$, $S$ a $V$-conjugation-stable set with $S\cap V\subset U$, and $\mathrm{Ch}(H)$ the continuous
$H\to\mathbb F_2$ vanishing on $H\cap S$. If $\#\mathrm{Ch}(V)\le n$ then $\#\mathrm{Ch}(U)\le n^2$.
\end{lemma}
\begin{proof}
Assume $U\ne V$, fix $\sigma\in V\smallsetminus U$, $c(u)=\sigma^{-1}u\sigma$ and $f(\varphi)=\varphi+\varphi\circ c$. Then $f\circ f=0$, so
$\#\mathrm{Ch}(U)\le(\#\ker f)^2$, and $\ker f=P=\{\varphi\circ c=\varphi\}$. To see $\#P\le n$: if $\varphi\in P$ and $\varphi(\sigma^2)=0$, then
$\psi(u)=\psi(\sigma u)=\varphi(u)$ extends $\varphi$ to $\mathrm{Ch}(V)$; the sign character $\varepsilon$ ($\varepsilon(U)=0$) lies in $\mathrm{Ch}(V)$. Fixing
$\varphi_0\in P$ with $\varphi_0(\sigma^2)=1$ (if any), send $\varphi$ to the extension of $\varphi$ or $\varphi+\varphi_0$, plus $\varepsilon$ in the second case;
evaluation at $\sigma$ and restriction to $U$ recover $\varphi$, so this is injective.
\end{proof}

\begin{lemma}[orbits and inertia over $K_M(\mu_{2^\infty})$]\label{i5:x-lem:bab-orb}
Let $M$ be an abelian datum, $K_M$ an $M$-field, $V_M=G_{K_M(\mu_{2^\infty})}$, $p\nmid2m_M$. (1) $I_w\subset V_M$ for $w\mid p$. (2) If $p\equiv3\bmod8$
(resp.\ $7\bmod16$), the places above $p$ form at most $2\deg M$ (resp.\ $4\deg M$) $V_M$-orbits.
\end{lemma}
\begin{proof}
(1) $K_M(\mu_{2^\infty})$ is unramified at $p$. (2) The number of orbits is $[G_\Q:V_MD_{w_0}]\le[G_\Q:G_{K_M}][\Z_2^\times:\overline{\langle p\rangle}]$, and
$\overline{\langle p\rangle}$ has index $2$ (resp.\ $4$): $-p$ generates $1+4\Z_2$ (resp.\ $1+8\Z_2$).
\end{proof}

\begin{lemma}[restriction along a quadratic step]\label{i5:x-lem:bab-shatwo}
If $[E':E]\le2$ and $k\in\Z$, then $\operatorname{sha}(E,2,k)\le1+\operatorname{sha}(E',2,k)$.
\end{lemma}
\begin{proof}
By Lemma~\ref{i5:x-lem:bab-infres}(a), restriction maps $\Sha_E$ to $\Sha_{E'}$ with kernel inside $\mathcal K=\ker(H^1(G_E,M_k)\to H^1(G_{E'},M_k))$. For
$E'\ne E$ and $\gamma\in G_E\smallsetminus G_{E'}$, $[z]\mapsto z(\gamma)\bmod(\gamma-1)A$, $A=M_k^{G_{E'}}$, embeds $\mathcal K$ in a subquotient of
$\Q_2/\Z_2$, and $2\mathcal K=0$ (from $z(\gamma^2)=0$). Every subquotient of $\Q_2/\Z_2$ has $2$-torsion of length $\le1$.
\end{proof}

\begin{lemma}[the base field and the layer]\label{i5:x-lem:bab-comp}
Let $\mathrm{PP}(m_M;q,d)$ and $L_0$ a $\mathrm{base}_2(M,q,d)$-field. (1) $\mu_4\subset L_0$, $[L_0:\Q]\le8\deg M$, $L_0\hookrightarrow\Q(\zeta_{4m_Mqd})$.
(2) For $e\ge0$, $L_1=L_0(\mu_{2^{e+2}})$ is a $2$-layer of height $m\le e+2$, and every $(M,e,q,d)$-field $E$ has $E\subset L_1$, $[L_1:E]\le2$.
\end{lemma}
\begin{proof}
(1) An $A$-field has degree $\deg A$ and lies in $\Q(\zeta_{m_A})$; $\deg(A\vee B)\le\deg A\deg B$, and $\mathrm{cyc}(4)$, $\Q(\sqrt{-q})$, $\Q(\sqrt{-d})$ have
degree $2$; $a\equiv1\bmod4$ on $H_{\mathrm{base}_2}$. (2) In $(\Z/N)^\times$, $L_1$ is fixed by $\{a\in H'_M:(a/q)=(a/d)=1,\ a\equiv1\ (2^{e+2})\}$ and $E$
by the same group with $a\equiv\pm1$; the index is $\le2$.
\end{proof}

\begin{proposition}[towers from Ferrero--Washington]\label{i5:x-prop:fw-towers}
Assume Ferrero--Washington as in Remark~\ref{i5:x-rem:fw}.
(1) For odd $l$ and an abelian datum $A$ having $\mu_l$ there is $R$ with $\mathrm{IsTower}(l,L_0,R)$ for every $A$-field $L_0$.
(2) For every abelian datum $M$ there is $R$ such that for every $\mathrm{PP}(m_M;q,d)$ the $\mathrm{base}_2(M,q,d)$-field $L_0$ (compositum of $M$, $\Q(i)$,
$\Q(\sqrt{-q})$, $\Q(\sqrt{-d})$) satisfies $\mathrm{IsTower}(2,L_0,R)$ and $\mathrm{TwoBase}(M,q,d,L_0)$: the three properties of
Lemma~\ref{i5:x-lem:bab-comp}(1), and for $e\ge2$ and every $(M,e,q,d)$-field $E$ a $2$-layer $L_1/L_0$ of height $\le e+2$ with
$\operatorname{sha}(E,2,k)\le1+\operatorname{sha}(L_1,2,k)$ for $k\ge1$.
\end{proposition}
\begin{proof}
(1) Ferrero--Washington gives $r$ with $\#\mathrm{Cl}(K_n)/l\le l^r$ for all layers $K_n=K(\mu_{l^n})$. By class field theory $D_\infty^{(1)}[l]$ is the union
of the duals of $\mathrm{Cl}(K_n)/l$, so has dimension $\le r$; $X_\infty^{(l)}$ is a quotient of $X_\infty^{(1)}$, hence $r$-generated by topological
Nakayama. Take $R=\max(1,r)$; all $A$-fields are isomorphic and totally complex.
(2) The same argument at $l=2$ for an $(M\vee\mathrm{cyc}(4))$-field gives $g=g(M)$ with $\#\mathrm{Ch}_2(V_M,\emptyset)\le2^g$. Here
$V_0=G_{L_0(\mu_{2^\infty})}=V_M\cap V_q\cap V_d$, $V_p=G_{\Q(\sqrt{-p})}$.
\emph{Step 1.} Lemmas~\ref{i5:x-lem:bab-orb} and \ref{i5:x-lem:bab-tame} (for $q$, then $d$) give $\#\mathrm{Ch}_2(V_M,\{q,d\})\le2^{g+6\deg M}$.
\emph{Step 2.} With $S$ the inertia groups at residue characteristics $\notin\{q,d\}$ ($S\cap V_M\subset V_q$, as $\Q(\sqrt{-q})$ is unramified outside $q$),
Lemma~\ref{i5:x-lem:bab-idx2} twice gives $\#\mathrm{Ch}_2(V_0,\{q,d\})\le2^{4(g+6\deg M)}$.
\emph{Step 3.} $D_\infty^{(2)}(L_0)[2]\subset D_\infty^{(1)}(L_0)[2]=\mathrm{Ch}_2(V_0,\emptyset)$, so as in (1) $\mathrm{IsTower}(2,L_0,R)$ with
$R=4(g+6\deg M)$, independent of $q,d$.
\emph{Step 4.} TwoBase: Lemmas~\ref{i5:x-lem:bab-comp} and \ref{i5:x-lem:bab-shatwo}.
\end{proof}

\begin{lemma}[window on a logarithmic scale]\label{i5:x-lem:bab-window}
Let $R\in\N$, $0<\eta\le C_t$, $X>1$, $\mathrm{Exc}\subset\N$ with $|\mathrm{Exc}|\le R$. There is $Y$ with $X^{1/J}\le Y$, $Y^2\le X$,
$J=2^{R+2}(C_t/\eta)^R$, such that for $x\in[Y,Y^2]$ and $e\in\mathrm{Exc}$: $e\le x^{C_t}\Rightarrow e\le x^\eta$.
\end{lemma}
\begin{proof}
With $U=\log X$, $B=C_t/\eta$, $y_k=U/(2(2B)^k)$ ($k\le R$): $y$ is bad for $a$ only if $a\in(y/B,2y]$, and these intervals are disjoint for the $y_k$.
So some $y_k$ is good for all $a\in\{\log e/C_t\}$; take $Y=e^{y_k}$.
\end{proof}

\begin{lemma}[cyclotomic layers]\label{i5:x-lem:bab-layerdisc}
(1) If $\mu_{l^k}\subset L$ then $l^k\le2[L:\Q]$. (2) If $L\hookrightarrow\Q(\zeta_N)$ then $\log|d_L|\le[L:\Q]\log N$. (3) Let $n_0=(n_1+1)l^j$,
$c_0=\log2+n_0\log(2n_0N_0)+n_0$. If $\log P\le w$, $w\ge1$, $[L_0:\Q]\le n_1$, $L_0\hookrightarrow\Q(\zeta_{N_0P})$ and $L_1/L_0$ is an $l$-layer of height
$j$, then $[L_1:\Q]\le n_0$ and $\log(2|d_{L_1}|)\le c_0w$.
\end{lemma}
\begin{proof}
(1) $\varphi(l^k)\mid[L:\Q]$. (2) $d_L^{[K:L]}\mid d_K$ and $|d_K|\le N^{\varphi(N)}$ for $K=\Q(\zeta_N)$. (3) $L_1\hookrightarrow\Q(\zeta_{N'})$ with
$N'\le2n_0N_0P$ by (1); apply (2).
\end{proof}

\begin{lemma}[logarithmic bookkeeping]\label{i5:x-lem:bab-logs}
Let $l\ge2$, $J>0$, $C_t\ge0$, $X\ge1$, $x\ge X^{1/J}$. (i) $\lfloor\log_l\lfloor X^{C_t}\rfloor\rfloor\le2C_tJ\log x$; (ii) $X^{\varepsilon/(4J)}\le x^{\varepsilon/4}$;
(iii) if $x\ge1$ and $2^e\le x^{A_1}$ then $e+2\le(2|A_1|+2)(1+\log x)$.
\end{lemma}
\begin{proof}
$\log X\le J\log x$ and $\log2>\frac12$.
\end{proof}

\begin{lemma}[final comparison]\label{i5:x-lem:bab-final}
For $\mathrm{pt}_C(x)=C(1+\log x)+2\exp(Cx^{3\varepsilon/4})$ there is $x_0\ge2$ with $x^{C_t}\mathrm{pt}_C(x)\le\exp(x^{4\varepsilon/5})$ for $x\ge x_0$.
\end{lemma}
\begin{proof}
$x^{C_t}\mathrm{pt}_C(x)\le\exp(\log3+(|C_t|+|C|)\log x+|C|+|C|x^{3\varepsilon/4})$, and each term is eventually $\le\frac14x^{4\varepsilon/5}$.
\end{proof}

\begin{lemma}[summation]\label{i5:x-lem:bab-vvat}
If $b\ge0$, $t\le T$, $Tb\le B$ and $\operatorname{sha}(F,l,e)\le b$ is finite for $1\le e<t$, then $\mathrm{VV}_l(F,t,B)$.
\end{lemma}

\begin{lemma}[LemmaK at fixed level]\label{i5:x-lem:bab-klen}
Assume $\mathrm{LemmaK}(C_K)$, $C_K>0$. If $L\supset\mu_l$ ($\mu_4$), $D$ an $l$-datum on $L$, $e\ge1$, then $V_D(e)\le2\,\mathrm{kB}_{C_K}(e,[L:\Q],|d_L|)<\infty$,
$\mathrm{kB}_C(e,n,D)=\exp((2en)^C\log(2D))$. If moreover $[L:\Q]\le n_0$, $\log(2|d_L|)\le c_0x^{\varepsilon/4}$, $x\ge1$, $e\le x^{\varepsilon/(2C_K)}$, then
$\mathrm{kB}\le\exp((2n_0)^{C_K}c_0x^{3\varepsilon/4})$.
\end{lemma}
\begin{proof}
$L$ is totally complex, so LemmaK gives $V_D(e)\log l\le\mathrm{kB}$; and $(2e[L:\Q])^{C_K}\le(2n_0)^{C_K}x^{\varepsilon/2}$.
\end{proof}

\begin{lemma}[pointwise bound]\label{i5:x-lem:bab-point}
Assume $\mathrm{LemmaK}(C_K)$, $\mathrm{ABound}(l,R,L_0)$ with $c$, $s_i$ as there. For $Z\in\N$ let $k=\lfloor\log_lZ\rfloor$ and
$\mathrm{Exc}(Z)=\{1\le e\le Z:\exists i,\ |s_i-e|<l^{-k}\}$; then $|\mathrm{Exc}(Z)|\le R$. Let $L_2/L_0$ be a layer of height $j_0$ with $[L_2:\Q]\le n_0$,
$\log(2|d_{L_2}|)\le c_0x^{\varepsilon/4}$, and suppose $e'\le x^{\varepsilon/(2C_K)}$ for $e'\in\mathrm{Exc}(Z)$ with $e'\le x^{C_t}$. Then for $1\le e\le Z$,
$e\le x^{C_t}$, and every layer $L_1/L_0$ of height $m$,
$\operatorname{sha}(L_1,l,e)\le R(m+j_0+c+k)+2\exp((2n_0)^{C_K}c_0x^{3\varepsilon/4})$.
\end{lemma}
\begin{proof}
Two elements of $[1,Z]$ near the same $s_i$ differ by a multiple of $l^{k+1}>Z$. For $e\notin\mathrm{Exc}(Z)$ use ABound with $B=k$; for
$e\in\mathrm{Exc}(Z)$ use SecondBound and Lemma~\ref{i5:x-lem:bab-klen}.
\end{proof}

\begin{lemma}[uniform pointwise bound]\label{i5:x-lem:bab-point2}
With $C=R(a+j_0+c_1+2C_tJ)+1+(2n_0)^{C_K}c_0$, in the situation of Lemma~\ref{i5:x-lem:bab-point} with $c\le c_1$, $Z=\lfloor X^{C_t}\rfloor$,
$X^{1/J}\le x\le X$ and $m\le a(1+\log x)$: $\operatorname{sha}(L_1,l,e)+1\le\mathrm{pt}_C(x)$.
\end{lemma}
\begin{proof}
$e\le Z$, and $k\le2C_tJ\log x$ by Lemma~\ref{i5:x-lem:bab-logs}(i); insert in Lemma~\ref{i5:x-lem:bab-point}.
\end{proof}

\begin{theorem}[ThmBab]\label{i5:x-thm:thmbab}
Assume PTD, local finiteness, local Tate duality, Ferrero--Washington (for $K$ abelian over $\Q$ containing $\mu_l$ ($\mu_4$ if $l=2$) there is $r$
with $\#\mathrm{Cl}(K_n)/l\le l^r$ for all cyclotomic $\Z_l$-layers $K_n$), and $\mathrm{LemmaK}(C_K)$. Then there is a constant bundle $K$ with
$\mathrm{ThmBab}(K,C_t)$ for all $C_t\ge1$.
\end{theorem}
\begin{proof}
Propositions~\ref{i5:x-prop:abound} and \ref{i5:x-prop:fw-towers} give $R_l(A)$ with $\mathrm{ABound}(l,R_l(A),L_0)$ for all $A$-fields ($l$ odd, $A$ having
$\mu_l$), and $R_0(M)$ with $\mathrm{TwoBase}$ and $\mathrm{ABound}(2,R_0(M),L_0)$ for each base field. As $\mathrm{kB}_C$ increases in $C$, assume $C_K\ge1$. Put
$K=(R_0,R_\cdot,C_K)$, and for data $(M,B',L)$: $R=R_0(M)+\sum_{l\in B'}R_l(L(l))$, $J=2^{R+2}(2C_KC_t/\varepsilon)^R$
(Definition~\ref{i5:def-windows}), $\eta=\varepsilon/(2C_K)\le C_t$.

\emph{Part A ($l=2$).} By Lemma~\ref{i5:x-lem:bab-layerdisc}(3) ($j=j_0$, $n_1=8\deg M$, $N_0=4m_M$, $P=qd$, $w=x^{\varepsilon/4}$, admissible by
Lemma~\ref{i5:x-lem:bab-logs}(ii)), layers $L_2/L_0$ of height $j_0$ satisfy the hypotheses of Lemma~\ref{i5:x-lem:bab-point} with $n_0,c_0$ depending on
$(M,R)$; and $c\le c_1=16\deg M$ by Lemma~\ref{i5:x-lem:bab-layerdisc}(1). For an $(M,e(x),q,d)$-field $E$, TwoBase gives $L_1$ of height
$m\le(2|A_1|+2)(1+\log x)$ (Lemma~\ref{i5:x-lem:bab-logs}(iii)). If $x\ge2$, $X^{1/J}\le x\le X$ and $\mathrm{Exc}_A=\mathrm{Exc}(\lfloor X^{C_t}\rfloor)$ is
$\eta$-good at $x$, Lemma~\ref{i5:x-lem:bab-point2} gives $\operatorname{sha}(E,2,k)\le1+\operatorname{sha}(L_1,2,k)\le\mathrm{pt}_C(x)$ for $1\le k<t(x)\le x^{C_t}$, and
Lemmas~\ref{i5:x-lem:bab-vvat}, \ref{i5:x-lem:bab-final} give $\mathrm{VV}_2(E,t(x),\exp(x^{4\varepsilon/5}))$ for $x\ge x_0$: this is $B_a$ at $x$.

\emph{Part B ($l\in B'$).} Fix an $L(l)$-field $L_0$ and a layer $L_2$ of height $j_0$; $n_0=[L_2:\Q]$, $c_0=\log(2|d_{L_2}|)$. All $L(l)$-fields and their
layers of height $j_0$ are isomorphic to these. With $\mathrm{Exc}_l$ the exceptional set of $L_0$, Lemma~\ref{i5:x-lem:bab-point2} ($m=a=j_0$) and
Lemmas~\ref{i5:x-lem:bab-vvat}, \ref{i5:x-lem:bab-final} give $\mathrm{VV}_l(L_1,t(x),\exp(x^{4\varepsilon/5}))$ for $x\ge x_0$ whenever $\mathrm{Exc}_l$ is
$\eta$-good at $x$.

\emph{Window.} Let $x_1$ dominate all thresholds and $X_0=\max(2,x_1^J)$. For $X\ge X_0$, $\mathrm{Exc}=\mathrm{Exc}_A\cup\bigcup_l\mathrm{Exc}_l$ has $\le R$
elements, and Lemma~\ref{i5:x-lem:bab-window} gives $Y$ with $X^{1/J}\le Y$, $Y^2\le X$, such that $\mathrm{Exc}$ is $\eta$-good on $[Y,Y^2]$. Parts A and B
apply at every $x\in[Y,Y^2]$: this is $\mathrm{BWin}(K,C_t,\varepsilon,M,B',L,e,t,q,d;X,Y)$.
\end{proof}

\begin{remark}\label{i5:x-rem:fw}
The Ferrero--Washington hypothesis is $\mu=0$ for the cyclotomic $\Z_l$-extension of an abelian field: Ferrero--Washington (1979, \emph{Ann. of
Math.} 109, Main Theorem), which covers every prime $l$, including $l=2$.  The form used in Theorem~\ref{i5:thm-fw} (boundedness of
$\#\mathrm{Cl}(K_n)/l$) is equivalent to $\mu=0$ by Washington, \emph{Introduction to Cyclotomic Fields}, Prop.~13.23 (2nd ed.).
\end{remark}

\subsection{Proof of Proposition~\ref{i5:prop-shak}}\label{i5:ssec-pf-shak}

\medskip\noindent\emph{Notation.}
$n_L=[L:\Q]$, $D_L=|d_L|$, $r_2(L)$ the number of complex places, $h_L$ the class number,
$\mathrm{gens}(L)=2n_L+\lfloor\log_2h_L\rfloor+2$. For $N,m\in\N$, $\mathrm{tc}(L,N,m)=\#\big(H_m(\mathrm{GL}_N(\mathcal O_L),\Z)_{\mathrm{tors}}\big)$
(group homology with trivial coefficients). $\operatorname{sha}(L,l,E)=\lgth_{\Z_l}\Sha^1_D(L,\Q_l/\Z_l(-E))$ for any $l$-datum $D$ (it does not depend
on $D$: changing $A_v$ conjugates $D_v$). $\mathrm{kB}_C(e,n,D)=\exp\big((2en)^C\log(2D)\big)$, so LemmaK$(C)$ reads
$\log\#H^2(G_{L,S},\Z_l(e+1))\le\mathrm{kB}_{C}(e,n_L,D_L)$.
A \emph{layer} $\mathrm{Lay}(l,c,\nu,\zeta,\gamma)$ for $F/L$ means: $F/L$ is Galois with group generated by $\gamma$ of order $l^{\nu-c}$,
$\zeta\in\mathcal O_F^\times$ is a primitive $l^\nu$-th root of unity with $\gamma\zeta=\zeta^{1+l^c}$, and $1\le c\le\nu$ ($c\ge2$ if $l=2$).
For $a\in\N$: $\alpha\in F^\times$ is \emph{$a$-twisted} as in Definition~\ref{i5:x5d-def-val}(5) (i.e.\ $\operatorname{ord}_P(\alpha)=\mathfrak d(\gamma P)-a\,\mathfrak d(P)$ for all $P$, for some divisor $\mathfrak d$);
$\mathrm{nrm}_a(\alpha)=\prod_{j<\operatorname{ord}\gamma}\gamma^j(\alpha)^{a^j}$; and $\mathrm{EqHom}(\gamma,a,l^\nu)$ is the group of maps $\varphi$ from the
primes of $F$ to $\Z/l^\nu$ with $\sum_P\operatorname{ord}_P(\alpha)\varphi(P)=0$ for all $\alpha\in F^\times$ and $\varphi(\gamma^{-1}P)=a\varphi(P)$, i.e.\ the
$a$-eigen-homomorphisms $\mathrm{Cl}(F)\to\Z/l^\nu$.

\begin{definition}[the intermediate statements]\label{i5:x-def:lemmaK-steps}
\begin{enumerate}
\item $\mathrm{ShaK}(C)$: for every prime $l$, every totally complex $L\supset\mu_l$ ($\mu_4$ if $l=2$) and $E\ge1$:
$\operatorname{sha}(L,l,E)<\infty$ and $\operatorname{sha}(L,l,E)\log l\le\mathrm{kB}_C(E,n_L,D_L)$.
\item $\mathrm{SLayer}$: for $L\supset\mu_l$ ($\mu_4$) and $E\ge1$ there are a totally complex $L'\supset L$ and $c,d$ with $\mu_{l^c}\subset L'$,
$\mu_{l^{c+1}}\not\subset L'$ ($c\ge1$, $c\ge2$ if $l=2$), $2E+1<l^c$, $2E+1<r_2(L')$, $n_{L'}=d\,n_L$, $d\le(n_L+1)(4E+4)$,
$l^c\le2n_{L'}$, $D_{L'}\le D_L^d(l^c)^{n_{L'}}$, and $\operatorname{sha}(L,l,E)\le\operatorname{sha}(L',l,E)$.
\item $\mathrm{SSha}$: if $L$ is totally complex, $\mu_{l^c}\subset L$ maximal as in (2), $E\ge1$, and for every layer $\mathrm{Lay}(l,c,\nu,\zeta,\gamma)$ for
$F/L$ the group $\mathrm{EqHom}(\gamma,(1+l^c)^E,l^\nu)$ is finite of order $\le B$, then $l^{\operatorname{sha}(L,l,E)}\le B$.
\item $\mathrm{SFew}$: for $L$ totally complex, a layer for $F/L$, $2e+3<l^c$ and $2e+3<r_2(L)$, with $a=(1+l^c)^{e+1}$: every set $S$ of
$a$-twisted elements of $F^\times$ such that $\mathrm{nrm}_a(\alpha)^s\equiv\mathrm{nrm}_a(\alpha')^s$ modulo $(F^\times)^{l^\nu}$ forces
$\alpha=\alpha'$, where $s=l^{2e+2+(2e+3)(c+2e+3)}$, has
$|S|\le(l^\nu)^{r_2(L)}\,\mathrm{tc}(L,4e+10,2e+3)\,\mathrm{tc}(L,4e+10,2e+2)$.
\item $\mathrm{SCount}$: for $L$ totally complex, a layer, $E\ge1$, if the statement in (4) holds with exponent $s=l^b$ and bound $B$ for
$a=(1+l^c)^E$, then $\mathrm{EqHom}(\gamma,a,l^\nu)$ is finite and
$\#\mathrm{EqHom}\cdot(l^\nu)^{r_2(L)}\le B\,l^{(b+v_l(E))\mathrm{gens}(L)+c+v_l(E)}$.
\item $\mathrm{STors}$: there is $C$ with $\log\mathrm{tc}(L,4E+6,m)\le\mathrm{kB}_C(E,n_L,D_L)$ for $L$ totally complex, $E\ge1$, $m\le2E+1$.
\end{enumerate}
\end{definition}


A number field $K$ is \emph{$l$-admissible} if $\mu_l\subset K$ ($\mu_4\subset K$
if $l=2$); such a field is totally complex.  Statements (2) and (3) of
Definition~\ref{i5:x-def:lemmaK-steps} are proved below
(Propositions~\ref{i5:x-prop:K4-slayer} and \ref{i5:x-prop:K4-ssha}), together
with the assembly (Proposition~\ref{i5:x-prop:snum}).

Statement (4) of Definition~\ref{i5:x-def:lemmaK-steps} (SFew) is Proposition~\ref{i5:x-prop:sfew}, reduced in \S\ref{i5:x5-ssec}; statements (5) and (6) are proved here.


\medskip\noindent\emph{SCount.}
Divisors of $F$ are written additively, $\mathrm{Div}(F)=\bigoplus_P\Z$ over the
maximal ideals $P$ of $\mathcal O_F$, $\mathrm{div}\colon F^\times\to\mathrm{Div}(F)$,
$\mathrm{Cl}(F)=\mathrm{Div}(F)/\mathrm{div}(F^\times)$; for an automorphism
$\sigma$, $\operatorname{ord}_P(\sigma^{-1}\alpha)=\operatorname{ord}_{\sigma P}(\alpha)$.
The argument is first carried out for an abstract module with a twisted
endomorphism.

\begin{definition}[twisting data]\label{i5:x-def:K4-datum}
A \emph{twisting datum} consists of abelian groups $Y$, $I$ with $I$ torsion
free, endomorphisms $g$ of $Y$ and $g_I$ of $I$, a homomorphism $d\colon Y\to I$
with $dg=g_Id$, an integer $n\ge1$ with $g^n=1$, $g_I^n=1$, and an integer
$a\ge2$.  Put $\delta=g-a$, $\delta_I=g_I-a$,
$\mathrm{nm}=\sum_{j=0}^{n-1}a^jg^{n-j}$, $Y_{\rm t}$ the torsion of $Y$,
$U=d^{-1}(\delta_II)$, and
\[
W=U/\big((\delta Y+Y_{\rm t})\cap U\big),\quad
E_Q=\ker d/\big(\delta(\ker d)+Y_{\rm t}\big),\quad
C_Q=I/(dY+\delta_II).
\]
For $m\ge1$ let $Y\langle m\rangle=Y/(mY+Y_{\rm t})$ and
$\mathrm{Sel}(m)=\{[y]\in Y\langle m\rangle:\ gy-y\in mY+Y_{\rm t},\ dy\in mI\}$.
\end{definition}

\begin{lemma}\label{i5:x-lem:K4-abstract}
Let $(Y,I,g,g_I,d,n,a)$ be a twisting datum.
\begin{enumerate}
\item $\delta_I$ is injective, and $d(Y_{\rm t})=0$.
\item If $I/dY$ and $E_Q$ are finite, then $C_Q$ and $W$ are finite and $\#W=\#E_Q\cdot\#C_Q$.
\item If $a^n-1=qu$ with $q,u\in\Z$, then for every $s\in\N$ there is a
homomorphism $f\colon W\to Y\langle q\rangle$ with $f([y])=[s\,\mathrm{nm}(y)]$,
and $us$ kills $\ker f$.
\item If $l\mid a-1$ and $\mathrm{Sel}(l)$ is finite, then $\#W[l]\le\#\mathrm{Sel}(l)$.
\end{enumerate}
\end{lemma}

\begin{proof}
(1) If $g_Ix=ax$ then $x=g_I^nx=a^nx$ with $a^n-1\ge1$, so $x=0$; and $I$ is torsion free.

(2) By (1) there is a unique $\psi\colon U\to I$ with $\delta_I\psi(u)=du$.  Let
$\bar\delta$ be the endomorphism of the finite group $\bar I=I/dY$ induced by
$\delta_I$.  Since $\psi(\delta y+t)=dy$ for $t\in Y_{\rm t}$, $\psi$ induces
$\Psi\colon W\to\ker\bar\delta$, which is onto (if $\delta_Ix=dy$ then $y\in U$,
$\psi(y)=x$).  Its kernel is the image of $\ker d$ (if $\psi(u)=dy$ then
$u-\delta y\in\ker d$), and $\ker d\to W$ has kernel
$(\delta(\ker d)+Y_{\rm t})$, because $d(\delta y+t)=0$ forces $\delta_I(dy)=0$,
hence $dy=0$.  So $0\to E_Q\to W\to\ker\bar\delta\to0$ is exact, and
$\#\ker\bar\delta=\#\operatorname{coker}\bar\delta=\#C_Q$.

(3) Telescoping with $g^n=1$ gives $\mathrm{nm}\,\delta=\delta\,\mathrm{nm}=(1-a^n)g=-qu\,g$,
so $s\,\mathrm{nm}(\delta z)\in qY$; with $\mathrm{nm}(Y_{\rm t})\subset Y_{\rm t}$ the map
is well defined on $W$.  If $f([y])=0$, put $v=sy$, so $\mathrm{nm}(v)=qz+t$,
$t\in Y_{\rm t}$.  Applying $\delta$: $q(-u\,gv-\delta z)\in Y_{\rm t}$, and $q\ne0$,
so $u\,gv\in R:=\delta Y+Y_{\rm t}$.  As $g^{n-1}$ commutes with $\delta$,
$uv=g^{n-1}(u\,gv)\in R$, i.e.\ $(us)[y]=0$.

(4) For $w=[u]\in W[l]$ write $lu=\delta y+t$, $t\in Y_{\rm t}$, and
$a-1=lc'$.  Then $gy-y=l(u+c'y)-t$, and if $du=\delta_Ii$ then
$\delta_I(dy)=l\,du=\delta_I(li)$, so $dy=li$ by (1); thus $[y]\in\mathrm{Sel}(l)$.
If $w,w'$ give the same class, $y-y'=lz+t_0$, then
$l(u-u'-\delta z)\in Y_{\rm t}$, so $w=w'$.  Hence $w\mapsto[y]$ is injective.
\end{proof}

\begin{lemma}\label{i5:x-lem:K4-count}
(1) For a finite abelian group $H$ and $\nu\ge0$, $\#\mathrm{Hom}(H,\Z/l^\nu)\le l^{v_l(\#H)}$.
(2) Let $W$ be finite abelian, $f\colon W\to Z$ a homomorphism, $s,g\in\N$ and
$u\in\Z$ with $l\nmid u$, such that $ul^s$ kills $\ker f$ and $\#W[l]\le l^g$.
Then $l^{v_l(\#W)}\le\#f(W)\cdot l^{sg}$.
\end{lemma}

\begin{proof}
(1) Write $\#H=l^em$, $l\nmid m$.  Restriction to $mH$ (of order dividing $l^e$)
is injective on $\mathrm{Hom}(H,\Z/l^\nu)$ since $m$ is invertible mod $l^\nu$,
and $\#\mathrm{Hom}(mH,\Z/l^\nu)\le\#mH$.
(2) With $K=\ker f$, $l^sK$ is killed by $u$, so
$l^{v_l\#K}\le\#K[l^s]\le\#W[l^s]\le(\#W[l])^s\le l^{sg}$; and
$l^{v_l\#f(W)}\le\#f(W)$.
\end{proof}

\begin{lemma}[arithmetic of a layer]\label{i5:x-lem:K4-layer}
Let $L$ be totally complex, $\mathrm{Lay}(l,c,\nu,\zeta,\gamma)$ a layer for $F/L$, $n=l^{\nu-c}=[F:L]$.
\begin{enumerate}
\item $F/L$ is unramified outside $l$.
\item $F$ is totally complex, and $\gamma$ permutes the $n\,r_2(L)$ archimedean
places of $F$ in $r_2(L)$ orbits of length $n$.
\item For $a\ge2$, $(Y,I,g,g_I,d,n,a)=(F^\times,\mathrm{Div}(F),\gamma^{-1},\
\mathfrak d\mapsto\mathfrak d(\gamma\,\cdot),\ \mathrm{div},n,a)$ is a twisting datum.
$U$ is the group of $a$-twisted elements, $\mathrm{nm}=\mathrm{nrm}_a$,
$I/dY=\mathrm{Cl}(F)$ is finite, $E_Q$ is finite and
$\#E_Q\cdot(a-1)=(a^n-1)^{r_2(L)}$.
\end{enumerate}
\end{lemma}

\begin{proof}
(1) If $P\nmid l$ and $\gamma^j$ lies in the inertia group of $P$, then
$\zeta^{(1+l^c)^j}\equiv\zeta\bmod P$; the $l^\nu$-th roots of unity are distinct
mod $P$, so $l^\nu\mid(1+l^c)^j-1$, hence $l^{\nu-c}\mid j$ by
Lemma~\ref{i5:x-lem:K4-lte} and $\gamma^j=1$.
(2) Each archimedean place of $L$ is complex, so its decomposition group in
$\mathrm{Gal}(F/L)$ is trivial and it has $n$ extensions, permuted simply
transitively.
(3) The first claims are the definitions ($g^{n-j}=\gamma^j$), and $\ker d=\mathcal O_F^\times$.
Finiteness of $\mathrm{Cl}(F)$ is the finiteness of the class number.  Let $X$
be the set of archimedean places of $F$ and
$\lambda(u,m)=(\log\|u\|_w+m)_{w\in X}$ on $\mathcal O_F^\times\times\Z$.  By
Dirichlet's unit theorem the log-embedding of $\mathcal O_F^\times$ is a full
lattice in the trace-zero hyperplane with kernel $\mu_F$; since
$(1,\dots,1)$ is not in that hyperplane, $\Lambda=\operatorname{im}\lambda$ is a
full lattice in $\R^X$ and $\ker\lambda=\mu_F\times0$.  With $\pi^*$ the
permutation of coordinates induced by $\gamma$, $\lambda$ intertwines
$\delta\times(1-a)$ with $\pi^*-a$.  By (2), $\pi^*$ is a sum of $r_2(L)$ cyclic
permutation matrices of size $n$, so $|\det(\pi^*-a)|=(a^n-1)^{r_2(L)}$, which is
the index $[\Lambda:(\pi^*-a)\Lambda]$.  Transporting through $\lambda$, this
index is the order of $E_Q\times\Z/(a-1)$.
\end{proof}

\begin{lemma}\label{i5:x-lem:K4-eqhom}
In Lemma~\ref{i5:x-lem:K4-layer}(3), $\varphi\mapsto(\mathfrak d\mapsto\sum_P\mathfrak d(P)\varphi(P))$
embeds $\mathrm{EqHom}(\gamma,a,l^{\nu'})$ into $\mathrm{Hom}(C_Q,\Z/l^{\nu'})$;
hence $\#\mathrm{EqHom}(\gamma,a,l^{\nu'})\le l^{v_l(\#C_Q)}$ when $C_Q$ is finite.
\end{lemma}

\begin{proof}
The two defining conditions of $\mathrm{EqHom}$ say that the functional kills
$\mathrm{div}(F^\times)$ and $\delta_I\mathrm{Div}(F)$; it determines $\varphi$ on
the basis vectors $\mathbf 1_P$.  Conclude by Lemma~\ref{i5:x-lem:K4-count}(1).
\end{proof}

\begin{lemma}[Selmer bounds]\label{i5:x-lem:K4-sel}
(1) For a number field $L$,
$\mathrm{Sel}_L(l)=\{x\in L^\times/L^{\times l}:\ l\mid\operatorname{ord}_p(x)\ \forall p\nmid l\}$
has order $\le l^{2n_L+\lfloor\log_2h_L\rfloor}$.
(2) In Lemma~\ref{i5:x-lem:K4-layer}(3), $\#\mathrm{Sel}(l)\le l^2\,\#\mathrm{Sel}_L(l)\le l^{\mathrm{gens}(L)}$.
\end{lemma}

\begin{proof}
(1) Let $H$ be the preimage of $\mathrm{Sel}_L(l)$ in $L^\times$,
$H_0=\{y:(y)=\mathfrak a^l\}$ and $U_0=\mathcal O_L^\times L^{\times l}$.  Then
$[U_0:L^{\times l}]\le l^{r_1+r_2}\le l^{n_L}$ (Dirichlet), $y\mapsto[\mathfrak a]$
embeds $H_0/U_0$ into $\mathrm{Cl}(L)[l]$, of order $\le l^{\lfloor\log_2h_L\rfloor}$,
and $y\mapsto(\operatorname{ord}_py\bmod l)_{p\mid l}$ embeds $H/H_0$ into
$(\Z/l)^{\le n_L}$.

(2) Here $Y\langle l\rangle=F^\times/F^{\times l}\mu_F$; let $\rho$ generate
$\mu_F$.  A class in $\mathrm{Sel}(l)$ is represented by $\alpha$ with all
$\operatorname{ord}_P\alpha$ divisible by $l$ and $\gamma\alpha=\alpha v^l\rho^r$,
$0\le r<l$.  Taking $N_{F/L}$, $N(v)^l=N(\rho)^{-r}$, so the \emph{type}
$(r,N(v))$ takes at most $l^2$ values.  If $\alpha,\alpha'$ have the same type,
$\beta=\alpha/\alpha'$ satisfies $\gamma\beta=\beta y^l$ with $N(y)=1$; by
Hilbert's Theorem~90 for the cyclic extension $F/L$, $y=w/\gamma w$, so
$x_0=\beta w^l\in L^\times$, and by (1) of Lemma~\ref{i5:x-lem:K4-layer},
$\operatorname{ord}_px_0=\operatorname{ord}_Px_0\equiv0\bmod l$ for $p\nmid l$.
So $[\alpha]=[\alpha']+[x_0]$ with $x_0\in\mathrm{Sel}_L(l)$.  The last
inequality is (1), as $\mathrm{gens}(L)=2n_L+\lfloor\log_2h_L\rfloor+2$.
\end{proof}

\begin{proposition}[SCount]\label{i5:x-prop:scount}
Statement (5) of Definition~\ref{i5:x-def:lemmaK-steps} holds.
\end{proposition}

\begin{proof}
Let $a=(1+l^c)^E$, $n=l^{\nu-c}$, $f_E=v_l(E)$, and take the twisting datum of
Lemma~\ref{i5:x-lem:K4-layer}(3).  By Lemma~\ref{i5:x-lem:K4-lte}, $a\ge2$,
$v_l(a-1)=c+f_E$ and $a^n-1=l^\nu\cdot l^{f_E}m$ with $l\nmid m$.  By
Lemmas~\ref{i5:x-lem:K4-layer} and \ref{i5:x-lem:K4-abstract}(2), $W$ and $C_Q$
are finite with $\#W=\#E_Q\#C_Q$, and by Lemma~\ref{i5:x-lem:K4-eqhom}
$\#\mathrm{EqHom}\le l^{v_l\#C_Q}$.  Taking $v_l$ of
$\#E_Q(a-1)=(a^n-1)^{r_2(L)}$ gives $v_l\#E_Q+c+f_E=r_2(L)(\nu+f_E)\ge\nu r_2(L)$.

Lemma~\ref{i5:x-lem:K4-abstract}(3) with $q=l^\nu$, $u=l^{f_E}m$, $s=l^b$ gives
$f\colon W\to Y\langle l^\nu\rangle$, $f([\alpha])=[\mathrm{nrm}_a(\alpha)^{l^b}]$,
with kernel killed by $m\,l^{b+f_E}$.  Choosing one $a$-twisted $\alpha_z$ in each
fibre over $z\in f(W)$ gives a set $S$ as in statement (4) (two of its elements
with $l^\nu$-congruent $l^b$-th norm powers have the same image, hence are equal),
so $\#f(W)\le B$.  Since $l\mid a-1$,
Lemmas~\ref{i5:x-lem:K4-abstract}(4) and \ref{i5:x-lem:K4-sel}(2) give
$\#W[l]\le l^{\mathrm{gens}(L)}$, and Lemma~\ref{i5:x-lem:K4-count}(2) gives
$l^{v_l\#W}\le B\,l^{(b+f_E)\mathrm{gens}(L)}$.  Hence
\[
\#\mathrm{EqHom}\cdot l^{\nu r_2(L)}\le l^{v_l\#C_Q+v_l\#E_Q+c+f_E}
=l^{v_l\#W}l^{c+f_E}\le B\,l^{(b+f_E)\mathrm{gens}(L)+c+f_E}.\qedhere
\]
\end{proof}

\medskip\noindent\emph{STors.}
For a set $X$, $\Z[X]$ is the free abelian group on $X$, with
$\|f\|_1=\sum_x|f(x)|$; a $G$-set $X$ makes $\Z[X]$ a permutation module.
$B_{\rm tors}$ is the torsion subgroup of an abelian group $B$.

\begin{proposition}[STors]\label{i5:x-prop:stors}
Let $L$ be totally complex, $n=n_L$, $D=D_L$, $e\ge1$, $1\le N\le4e+6$ and
$0\le m\le2e+1$.  Then $H_m(\mathrm{GL}_N(\mathcal O_L),\Z)_{\rm tors}$ is finite
and $\log\#H_m(\mathrm{GL}_N(\mathcal O_L),\Z)_{\rm tors}\le\mathrm{kB}_{24}(e,n,D)$.
In particular statement (6) of Definition~\ref{i5:x-def:lemmaK-steps} holds with $C=24$.
\end{proposition}

The proof has three parts: (A) a torsion bound for a group acting on an
acyclic complex of permutation modules with few orbits and small stabilisers;
(B) such a complex for $\mathrm{GL}_N(\mathcal O_L)$, the ordered nerve of the
normal cones of the vertices of the Ryshkov polyhedron of Hermitian forms,
with bounds on orbits and stabilisers; (C) the numerics.

\subparagraph{(A) Dimension shifting.}

\begin{lemma}\label{i5:x-lem:K5-elem}
(1) If $A\to B\to C$ is exact, $\#A\le a$ and $\#C_{\rm tors}\le c$, then
$\#B_{\rm tors}\le ac$; if $A,C$ are finite, $\#B\le\#A\#C$.
(2) If $\#\iota\le A$ and $\mathcal L\subset\Z[\iota]$ consists of vectors of
$\ell^1$-norm $\le b$ ($b\ge1$), then $\#(\Z[\iota]/\langle\mathcal L\rangle)_{\rm tors}\le b^A$.
(3) For a finite group $S$ and $j\ge0$, $\#H_{j+1}(S,\Z)\le\#S^{\#S^{j+1}}$.
(4) If $G$ acts on $X$ with $\le A$ orbits and stabilisers of order $\le s$, then
$\#H_{j+1}(G,\Z[X])\le s^{As^{j+1}}$.
\end{lemma}

\begin{proof}
(1) $B_{\rm tors}\to C_{\rm tors}$ has kernel inside the image of $A$.
(2) Let $\mathcal L_0\subset\mathcal L$ be a maximal independent subset, and
complete it by standard basis vectors to a basis $v_1,\dots,v_r$ of
$\Q[\iota]$.  By expansion of the determinant,
$[\Z[\iota]:\langle v_i\rangle]=|\det(v_i)|\le\prod\|v_i\|_1\le b^A$.  The torsion
of the cokernel is $\mathrm{sat}\langle\mathcal L_0\rangle/\langle\mathcal L\rangle$,
a quotient of $\mathrm{sat}\langle\mathcal L_0\rangle/\langle\mathcal L_0\rangle$,
which injects into $\Z[\iota]/\langle v_i\rangle$.
(3) Norm followed by augmentation is multiplication by $\#S$ and factors through
$H_{j+1}(S,\Z[S])=0$ (Shapiro), so $\#S$ kills $H_{j+1}(S,\Z)$, a quotient of the
free group of standard $(j+1)$-cycles, of rank $\le\#S^{j+1}$.
(4) Split $X$ into orbits; by Shapiro's lemma the orbit $Gx$ contributes
$H_{j+1}(G_x,\Z)$, bounded by (3), and (1) gives the product bound.
\end{proof}

\begin{proposition}[dimension shifting]\label{i5:x-prop:K5-shift}
Let $G$ act on sets $X_p$ ($p\ge0$) with equivariant face maps, $d_p$ the
alternating sums and $\varepsilon\colon\Z[X_0]\to\Z$ the augmentation.  If
$\Z[X_{m+1}]\to\cdots\to\Z[X_0]\to\Z\to0$ is exact and for $p\le m$ the set
$X_p$ has $\le A$ orbits and stabilisers of order $\le s$, then
$\#H_m(G,\Z)_{\rm tors}\le(m+2)^As^{mAs^m}$.  If $A\le e^a$ and $s\le e^\sigma$,
then $\log\#H_m(G,\Z)_{\rm tors}\le\exp(a+(m+1)\sigma+2m)$.
\end{proposition}

\begin{proof}
We prove more generally, for a $G$-module $M$ with a resolution of this kind
whose differential $d_{m+1}\colon\Z[X_{m+1}]\to\Z[X_m]$ has columns of $\ell^1$-norm $\le b$, that
$\#H_m(G,M)_{\rm tors}\le b^As^{mAs^m}$, by induction on $m$.  For $m=0$,
$M_G$ is the cokernel of $\Z[X_1/G]\to\Z[X_0/G]$ (right exactness of
coinvariants), and Lemma~\ref{i5:x-lem:K5-elem}(2) applies.  For $m+1$, let
$K=\ker\varepsilon$; the shifted complex resolves $K$ to level $m$, and the
sequence $H_{m+1}(G,\Z[X_0])\to H_{m+1}(G,M)\to H_m(G,K)$ with
Lemma~\ref{i5:x-lem:K5-elem}(4),(1) gives
$s^{As^{m+1}}b^As^{mAs^m}\le b^As^{(m+1)As^{m+1}}$.  Here $b=m+2$.  For the
logarithmic form, $A\log(m+2)+mAs^m\log s\le As^{m+1}(2m+1)\le e^{a+(m+1)\sigma+2m}$.
\end{proof}

\subparagraph{(B) The complex.}
Fix $L$, $n$, $D$, $N$ as in Proposition~\ref{i5:x-prop:stors} and put $d=nN$,
$I_0=D^Nd^d$, $R_0=N^NI_0^2$.  Let $\sigma_w$, $w\in W$, be representatives of
the $n/2$ complex embeddings, $\Lambda=\mathcal O_L^N$, $G=\mathrm{GL}_N(\mathcal O_L)$,
and $E=\prod_w\mathrm{Herm}_N(\C)$.  For $h\in E$ and $y\in\Lambda$ put
$q_h(y)=\sum_w\sigma_w(y)^*h_w\sigma_w(y)$; $G$ acts on $E$ by
$(gh)_w=\sigma_w(g^{-1})^*h_w\sigma_w(g^{-1})$, so $q_{gh}(y)=q_h(g^{-1}y)$.
Let $\mathrm{Psd}$, $\mathrm{Pos}$ be the tuples of positive semidefinite,
resp.\ definite, forms, and
\[
P=\{h\in E:\ q_h(y)\ge1\ \forall y\in\Lambda\smallsetminus0\}
\]
the Ryshkov polyhedron; $\mathrm{Min}(h)=\{y:q_h(y)=1\}$.  A \emph{vertex} is
$v\in P$ such that $k=0$ is the only $k\in E$ with $q_k=0$ on $\mathrm{Min}(v)$.
Let $\Omega=\{x\in E^*:\ x>0\text{ on }\mathrm{Psd}\smallsetminus0\}$, for a
vertex $v$ let $\mathcal N(v)=\{x\in E^*:\ x(v)\le x(h)\ \forall h\in P\}$, and
\[
X_p=\{(v_0,\dots,v_p)\ \text{vertices}:\ \Omega\cap\textstyle\bigcap_j\mathcal N(v_j)\ne\emptyset\},
\]
with faces deleting entries.  $G$ acts on $X_\bullet$ compatibly.

\begin{lemma}\label{i5:x-lem:K5-P}
(1) $\#\{a\in\mathcal O_L:\ |\sigma_w(a)|\le R\ \forall w\}\le(2R+1)^n$ for $R\in\N$.
(2) $P\ne\emptyset$, $P\subset\mathrm{Pos}$, $\Omega\ne\emptyset$, and $1\le I_0\prod_w\det h_w$ for $h\in P$.
(3) For $h\in\mathrm{Pos}$ and $B\in\R$ there are a finite $S\subset\Lambda$ and a
neighbourhood $U$ of $h$ with $q_k(y)>B$ for $k\in U$, $y\notin S$.
\end{lemma}

\begin{proof}
(1) Two distinct elements congruent mod $2R+1$ would differ by $(2R+1)x$ with
$0<|N(x)|<1$.
(2) The tuple of identity forms lies in $P$ (some conjugate of a nonzero
coordinate has modulus $\ge1$), and $x=\sum_w\operatorname{tr}$ lies in $\Omega$.
For $h\in P$ let $G_h$ be the real Gram matrix of $q_h$ on $\Lambda\cong\Z^d$
(via a $\Z$-basis of $\mathcal O_L$).  It is $\ge1$ on nonzero integral vectors,
hence positive semidefinite, and Minkowski's convex body theorem applied to
$G_h+\epsilon$ and $\epsilon\to0$ gives $1\le(4d)^d\det G_h$; in particular
$G_h\succ0$, and since $\R^d\cong\prod_w\C^N$ ($L$ totally complex), every
$h_w\succ0$.  Comparing with the $n$ embeddings gives
$2^d\det G_h=D^N(\prod_w\det h_w)^2$, whence $1\le2^dd^dD^N\Pi^2\le(I_0\Pi)^2$.
(3) There are $c>0$ and $U\ni h$ with $q_k(z)\ge c\|z\|^2$ on $U$ (sup norm);
take $R>\max(B/c,1)$ and $S$ the vectors with all conjugates of coordinates
$\le R$, finite by (1).
\end{proof}

\begin{lemma}[convexity]\label{i5:x-lem:K5-convex}
(a) If $h\in P$ and $q_k=0$ on $\mathrm{Min}(h)$, then $h+tk\in P$ for $|t|$ small.
(b) For $x\in\Omega$ and $c\in\R$, $\{h\in P:x(h)\le c\}$ is compact, and a compact
subset of $P$ contains finitely many vertices.
(c) Every $x\in\Omega$ lies in some $\mathcal N(v)$; for finite $s\subset\Omega$
only finitely many $\mathcal N(v)$ meet $\mathrm{conv}(s)$; the $\mathcal N(v)$ and
$\Omega$ are closed convex, resp.\ convex.
\end{lemma}

\begin{proof}
(a) Apply Lemma~\ref{i5:x-lem:K5-P}(3) at $h$ with $B=1$: off the finite $S$ the
inequality persists near $h$, on $S\smallsetminus\mathrm{Min}(h)$ it is strict, and
on $\mathrm{Min}(h)$, $q_{h+tk}=1$.
(b) $x\ge\mu\|\cdot\|$ on $\mathrm{Psd}$ for some $\mu>0$, and $P\subset\mathrm{Psd}$
is closed.  Near each point, vertices have $\mathrm{Min}(v)\subset S$, and a vertex
is determined by $\mathrm{Min}(v)$ (if $v'$ is a vertex equal to $1$ there, $v-v'$
vanishes on $\mathrm{Min}(v)$).
(c) By (b) $x$ attains its minimum on $P$ on a compact convex face, which has
an extreme point $e$ (Krein--Milman); by (a), if $k$ vanishes on $\mathrm{Min}(e)$
then $e\pm tk$ lie in the face, so $k=0$ and $e$ is a vertex.  If
$z\in\mathrm{conv}(s)\cap\mathcal N(v)$ and $h_0\in P$, then $x(v)\le x(h_0)$ for
some $x\in s$, so $v$ lies in a compact set as in (b).
\end{proof}

\begin{lemma}[nerve lemma]\label{i5:x-lem:K5-nerve}
Let $\mathcal X$ be a nonempty convex subset of a real normed space and
$(\mathcal A_i)_{i\in I}$ closed convex sets covering $\mathcal X$ such that for
every finite $s\subset\mathcal X$ only finitely many $\mathcal A_i$ meet
$\mathrm{conv}(s)$.  Let $\mathcal S_p$ be the set of $(i_0,\dots,i_p)$ with
$\mathcal X\cap\bigcap_j\mathcal A_{i_j}\ne\emptyset$.  Then the augmented complex
$\cdots\to\Z[\mathcal S_1]\to\Z[\mathcal S_0]\to\Z\to0$ is exact.
\end{lemma}

\begin{proof}
A cycle involves finitely many tuples; replacing $\mathcal X$ by the convex hull
$\mathcal K$ of a point of $\mathcal X$ and one common point for each tuple, we
may assume $\mathcal X=\mathcal K$ compact convex, meeting finitely many
$\mathcal A_i$.  Induct on the number $\Phi(\mathcal K)$ of finite $\sigma\subset I$
with $\mathcal K\cap\bigcap_\sigma\mathcal A_j\ne\emptyset$.  If some point of
$\mathcal K$ lies in all $\mathcal A_i$ meeting $\mathcal K$, prefixing an index
$a$ with that point in $\mathcal A_a$ is a contracting homotopy.  Otherwise take
$\sigma$ maximal with a common point $x_0\in\mathcal K$ and $i$ with
$\mathcal K\cap\mathcal A_i\ne\emptyset$, $x_0\notin\mathcal A_i$; then
$\mathcal K\cap\mathcal A_i$ and $\mathcal K\cap\bigcap_\sigma\mathcal A_j$ are disjoint
compact convex sets, separated (Hahn--Banach) by a hyperplane $f=u$.  The pieces
$\mathcal K_\le,\mathcal K_\ge,\mathcal K_=$ are compact convex and nonempty,
$\mathcal S(\mathcal K)=\mathcal S(\mathcal K_\le)\cup\mathcal S(\mathcal K_\ge)$ and
$\mathcal S(\mathcal K_=)=\mathcal S(\mathcal K_\le)\cap\mathcal S(\mathcal K_\ge)$ (a
segment between common points crosses $f=u$), and each piece has smaller $\Phi$
($\sigma$ is missing for $\mathcal K_\le,\mathcal K_=$, $\{i\}$ for $\mathcal K_\ge$).
Mayer--Vietoris concludes.
\end{proof}

\begin{corollary}\label{i5:x-cor:K5-acyc}
The augmented complex $\Z[X_\bullet]\to\Z$ is exact.
\end{corollary}

\begin{proof}
Lemma~\ref{i5:x-lem:K5-nerve} with $\mathcal X=\Omega$, $\mathcal A_v=\mathcal N(v)$,
by Lemmas~\ref{i5:x-lem:K5-P}(2) and \ref{i5:x-lem:K5-convex}(c).
\end{proof}

\begin{lemma}[spanning]\label{i5:x-lem:K5-span}
(a) If $M\subset\Lambda$ and $k=0$ is the only $k\in\mathrm{Psd}$ with $q_k=0$ on
$M$, then $M$ contains the columns of some $A\in M_N(\mathcal O_L)$ with $\det A\ne0$.
(b) For $(v_0,\dots,v_p)\in X_p$, $M=\bigcap_j\mathrm{Min}(v_j)$ satisfies (a).
\end{lemma}

\begin{proof}
(a) A nonzero linear form $f=\sum c_az_a$ vanishing on $M$ would give
$k_w=\overline{u_w}u_w^{\mathsf T}$, $u_w=\sigma_w(c)$, with
$q_k(y)=\sum_w|\sigma_w f(y)|^2=0$ on $M$.
(b) Let $x\in\Omega\cap\bigcap\mathcal N(v_j)$ and $c$ the barycentre of the
$v_j$; then $\mathrm{Min}(c)=M$.  If $0\ne k\in\mathrm{Psd}$ vanishes on $M$, then
$c-tk\in P$ for small $t>0$ by Lemma~\ref{i5:x-lem:K5-convex}(a), and
$x(c-tk)<x(c)=\min_Px$, a contradiction.
\end{proof}

\begin{lemma}[stabilisers]\label{i5:x-lem:K5-stab}
Every stabiliser $G_t$, $t\in X_p$, is finite of order $\le3^{nN^2}$.
\end{lemma}

\begin{proof}
$G_t\subset G_{v_0}$, which permutes the finite set $\mathrm{Min}(v_0)$; with $A$
from Lemma~\ref{i5:x-lem:K5-span}(a), $g\mapsto gA$ is injective, so $G_{v_0}$
is finite.  A finite subgroup of $M_N(\mathcal O_L)^\times$ injects into
$M_N(\mathcal O_L/3)$: if $u=1+3x$ has finite order $k$, then $x=0$, by strong
induction on $k$ using that $z=3zw$ forces $z=0$: for $3\nmid k$ expand
$(1+3x)^k=1$ and use a Bezout relation $ka+3b=1$; for $k=3k'$ apply induction to
$u^3$ and expand $u^3=1$.
\end{proof}

\begin{lemma}[norms and coordinates]\label{i5:x-lem:K5-nrm}
Let $h\in P$ and $A\in M_N(\mathcal O_L)$, $\det A\ne0$, with columns in
$\mathrm{Min}(h)$; $I=|N(\det A)|$.  Then $I\le I_0$, there is $B\in M_N(\mathcal O_L)$
with $AB=BA=I$, and for $y\in\mathrm{Min}(h)$ all conjugates of the coordinates of
$By$ have modulus $\le R_0$.
\end{lemma}

\begin{proof}
Put $h'_w=\sigma_w(A)^*h_w\sigma_w(A)\succ0$.  Its diagonal entries sum over $w$
to $q_h(Ae_a)=1$, so by Hadamard $\det h'_w\le1$; and
$\prod_w\det h'_w=I\prod_w\det h_w$.  With Lemma~\ref{i5:x-lem:K5-P}(2) this
gives $I\le I_0$ and $\det h'_w\ge1/I_0$.  $B=\pm\gamma\operatorname{adj}A$ with
$\gamma$ the product of the other conjugates of $\det A$.  For $y\in\mathrm{Min}(h)$
and $c=\sigma_w(A)^{-1}\sigma_w(y)$, $c^*h'_wc\le1$; by Cauchy--Schwarz all
entries of $h'_w$ and $h'_wc$ have modulus $\le1$, and Cramer's rule gives
$|c_a|\le N!\,I_0$, so $|\sigma_w((By)_a)|=I|c_a|\le R_0$.
\end{proof}

\begin{lemma}[orbits]\label{i5:x-lem:K5-orb}
$\#(X_p/G)\le I_0^{1+nN^2}\big((2R_0+1)^d+1\big)^{nN^2(p+1)}\le(2Dd)^{12(p+1)d^3N}$.
\end{lemma}

\begin{proof}
\emph{Hermite reduction.} If $A,A'$ have $|N\det|=I$ and the matrices $B,B'$ of
Lemma~\ref{i5:x-lem:K5-nrm} agree mod $I$, then $g=A'B/I\in G$ and $gA=A'$.
So up to $G$ there are $\le\sum_{I\le I_0}I^{nN^2}\le I_0^{1+nN^2}$ such $A$.
\emph{Vertices through $A$.} By Lemma~\ref{i5:x-lem:K5-nrm}, $y\mapsto By$
embeds the set of $y$ with $q_h(y)=1$ for some $h\in P$ having the columns of $A$
in $\mathrm{Min}(h)$ into a set of size $\le(2R_0+1)^d$
(Lemma~\ref{i5:x-lem:K5-P}(1)).  A vertex $v$ is determined by a subset of
$\mathrm{Min}(v)$ of size $\le\dim E\le nN^2$ whose functionals $q_\cdot(y)$
span those of $\mathrm{Min}(v)$; hence at most $((2R_0+1)^d+1)^{nN^2}$ vertices
have the columns of $A$ in their $\mathrm{Min}$.
\emph{Tuples.} For $t\in X_p$, Lemma~\ref{i5:x-lem:K5-span} gives $A$ with
columns in all $\mathrm{Min}(v_j)$ and $|N\det A|\le I_0$; moving $A$ to a
representative by $G$ moves $t$ to a tuple of vertices through it.  This gives
the first bound; the second is arithmetic with $\beta=2Dd$:
$I_0\le\beta^{2d}$, $2R_0+1\le\beta^{5d+2}$, and
$2d(1+dN)+((5d+2)d+1)dN(p+1)\le12(p+1)d^3N$.
\end{proof}

\begin{proof}[Proof of Proposition~\ref{i5:x-prop:stors}]
Apply Proposition~\ref{i5:x-prop:K5-shift} to $G$ and $X_\bullet$: exactness is
Corollary~\ref{i5:x-cor:K5-acyc}; for $p\le m$ there are
$\le(2Dd)^{12(m+1)d^3N}$ orbits (Lemma~\ref{i5:x-lem:K5-orb}) and stabilisers of
order $\le3^{nN^2}$ (Lemma~\ref{i5:x-lem:K5-stab}).  With
$a=12(m+1)d^3N\log(2Dd)$ and $\sigma=2nN^2$,
$\log\#H_m(G,\Z)_{\rm tors}\le\exp(a+(m+1)\sigma+2m)$.  Put $x=en\ge1$,
$\ell=\log(2D)>0.69$; then $N,d\le10x$, $m+1\le4x$, $\log(2Dd)\le16x\ell$, so
$a+(m+1)\sigma\le7681200\,x^6\ell$ and $2m\le12x^6\ell$, and the exponent is
$\le2^{24}x^{24}\ell=(2en)^{24}\log(2D)$.
\end{proof}

\begin{lemma}[class number bound]\label{i5:x-lem:classno}
For every number field $L$, $\log h_L\le4\,n_L\log(2D_L)\log2$; hence $\mathrm{gens}(L)\le2n_L+4n_L\log(2D_L)+2$.
\end{lemma}
\begin{proof}
If $n_L=1$, $h_L=1$. If $n_L\ge2$, every ideal class contains an integral ideal of norm at most the Minkowski bound
$\frac{n!}{n^n}(\frac4\pi)^{r_2}\sqrt{D_L}\le\sqrt{D_L}$ (Minkowski). The number of integral ideals of norm $\le Y$ is at most
$\sum_{m\le Y}\tau_n(m)\le Y(1+\log Y)^{n-1}$, where $\tau_n$ is the $n$-fold divisor function (each prime power $p^k$ has at most $\tau_n(p^k)$ ideals of
norm $p^k$; the bound on $\sum\tau_n$ follows by induction on $n$ from $\sum_{m\le Y}1/m\le1+\log Y$). Hence
$\log h_L\le\frac12\log D_L+(n-1)\log(1+\frac12\log D_L)\le n\log(2D_L)\le4\log2\cdot n\log(2D_L)$. Finally
$\lfloor\log_2h_L\rfloor\le\log h_L/\log2\le4n_L\log(2D_L)$.
\end{proof}

\begin{lemma}[three elementary inequalities]\label{i5:x-lem:num}
Let $u\ge2$, $Y\ge\frac12$, $C,C_0\in\R$.
\begin{enumerate}
\item If $P\le u^9Y$ and $t\le2\exp(u^CY)+P$, then $t\le\exp(u^{\max(C,9)+2}Y)$.
\item If $0\le t\le u$, $0\le b\le u+2u^2$, $0\le g\le u+2uY+2$, $0\le c\le u$, $0\le v\le u$, then $((b+v)g+c+v)\,t\le u^9Y$.
\item Let $C_0\ge0$, $E,n,D,n',D',d,w\ge1$ integers with $n'=dn$, $d\le(n+1)(4E+4)$, $w\le2n'$, $D'\le D^dw^{n'}$. Then
$\mathrm{kB}_{C_0}(E,n',D')\le\mathrm{kB}_{5C_0+11}(E,n,D)$.
\end{enumerate}
\end{lemma}
\begin{proof}
(1) Let $A=u^{\max(C,9)}Y\ge2^9/2$. Then $u^CY\le A$ and $P\le A$, so $t\le2e^A+A\le3e^A\le e^{A+2}\le e^{4A}\le e^{u^2A}$.
(2) As $uY\ge1$: $b+v\le4u^2$, $g\le u+2uY+2\le6uY$, $c+v\le2u\le2u^2Y$. So the left side is $\le(24u^3Y+2u^2Y)u\le28u^4Y\le u^9Y$, as $u^5\ge32$.
(3) Put $x=2En\ge2$. $n'=dn\le(n+1)(4E+4)n\le16En^2\le x^4$, so $2En'\le x^5$ and $(2En')^{C_0}\le x^{5C_0}$. Also
$\log(2D')\le\log2+d\log D+n'\log w\le d\log(2D)+n'\log(2n')\le n'\log(2D)+2n'^2\le n'\log(2D)+4n'^2\log(2D)\le5n'^2\log(2D)$,
using $d\le n'$, $\log(2n')\le2n'$ and $\log(2D)\ge\log2>\frac12$. Hence $(2En')^{C_0}\log(2D')\le5x^{5C_0+8}\log(2D)\le x^{5C_0+11}\log(2D)$.
\end{proof}

\begin{proposition}[assembly of $\mathrm{ShaK}$]\label{i5:x-prop:snum}
$\mathrm{SFew}\wedge\mathrm{SCount}\wedge\mathrm{SSha}\wedge\mathrm{SLayer}\wedge\mathrm{STors}\Rightarrow\mathrm{ShaK}(C)$ with
$C=5(\max(C_0,9)+2)+11$, $C_0$ the constant of STors.
\end{proposition}
\begin{proof}
Let $L\supset\mu_l$ ($\mu_4$), $E=k+1\ge1$. SLayer gives $L'$, $c$, $d$ with the listed properties and $\operatorname{sha}(L,l,E)\le\operatorname{sha}(L',l,E)$.
For every layer $\mathrm{Lay}(l,c,\nu,\zeta,\gamma)$ for $F/L'$, SFew (with $e=k$; its hypotheses $2k+3=2E+1<l^c$, $<r_2(L')$ hold) and SCount (with
$b=2k+2+(2k+3)(c+2k+3)$) give $\#\mathrm{EqHom}\le T_1T_2\,l^{X}$, where $T_1=\mathrm{tc}(L',4k+10,2k+3)$, $T_2=\mathrm{tc}(L',4k+10,2k+2)$,
$X=(b+f)g+c+f$, $f=v_l(E)$, $g=\mathrm{gens}(L')$. SSha gives $l^{s}\le T_1T_2l^X$ for $s=\operatorname{sha}(L',l,E)$, so
$s\log l\le\log T_1+\log T_2+X\log l$, and $\log T_i\le\exp(u^{C_0}Y)$ by STors, where $n'=n_{L'}$, $Y=\log(2D_{L'})\ge\log2>\frac12$, $u=2En'$.
Now $n'=2r_2(L')\ge2k+4$, so $u\ge2n'\ge4$ and $u\ge2E$; $c<l^c\le2n'$ and $l\le2n'$, so $c\le u$ and $\log l\le l\le u$; $f<E\le u$;
$b\le u+2u^2$ (as $(2k+3)(c+2k+3)\le n'\cdot3n'$); and $g\le2n'+4n'Y+2\le u+2uY+2$ by Lemma~\ref{i5:x-lem:classno}. Lemma~\ref{i5:x-lem:num}(2) gives
$X\log l\le u^9Y$, and Lemma~\ref{i5:x-lem:num}(1) gives $s\log l\le\exp(u^{C_1}Y)=\mathrm{kB}_{C_1}(E,n',D_{L'})$ with $C_1=\max(C_0,9)+2$. Finally Lemma~\ref{i5:x-lem:num}(3) with
$(n',D',d,w)=(n_{L'},D_{L'},d,l^c)$ gives $\le\mathrm{kB}_{5C_1+11}(E,n_L,D_L)$; and $\operatorname{sha}(L,l,E)\le s$ is finite.
\end{proof}

\begin{lemma}[lifting the exponent]\label{i5:x-lem:K4-lte}
Let $c\ge1$ ($c\ge2$ if $l=2$), $E\ge1$, $\nu\ge c$, $a=(1+l^c)^E$. Then $v_l(a-1)=c+v_l(E)$ and $v_l(a^{l^{\nu-c}}-1)=\nu+v_l(E)$. Moreover
$(1+l^c)^j\equiv1\bmod l^\nu$ if and only if $l^{\nu-c}\mid j$.
\end{lemma}
\begin{proof}
For $l$ odd this is the classical lifting-the-exponent lemma ($v_l((1+x)^m-1)=v_l(x)+v_l(m)$ for $v_l(x)\ge1$) applied to $x=l^c$ and $m=E$, $m=El^{\nu-c}$, $m=j$. For $l=2$ and $v_2(x)\ge2$ the same two steps work:
$(1+x)^2-1=x(2+x)$ with $v_2(2+x)=1$, and for odd $m$, $(1+x)^m-1=x\big(m+\sum_{i\ge2}\binom mix^{i-1}\big)$ with odd second factor; and
$v_2(x)$ stays $\ge2$ along the induction.
\end{proof}


\medskip\noindent\emph{Descent along $K(\mu_{l^\infty})$ and SSha.}

\medskip\noindent\emph{Notation.} Let $K$ be an $l$-admissible number field, $\Omega\supset K$ an algebraic closure, $\chi\colon G_K\to\Z_l^\times$ the cyclotomic
character, $U_\infty=\ker\chi=G_{K(\mu_{l^\infty})}$ and, for $\nu\ge1$, $U_\nu=\chi^{-1}(1+l^\nu\Z_l)=G_{K(\mu_{l^\nu})}$. For a place $w$ of $\Omega$
let $I_w\subset D_w\subset G_K$ be its inertia and decomposition groups. For $E\in\Z$, $M_E=\Q_l/\Z_l$ with $g$ acting by $\chi(g)^{-E}$. For a closed
normal subgroup $V\subset G_K$ with $U_\infty\subset V$ let $\mathcal E_V(K,E)$ be the $\Z_l$-module of continuous homomorphisms $\varphi\colon V\to\Q_l/\Z_l$ with
(i) $\varphi(I_w\cap V)=0$ for every finite place $w$; (ii) $\varphi(D_w\cap V)=0$ for every $w\mid l$; (iii) $\varphi(gug^{-1})=\chi(g)^{-E}\varphi(u)$
for $g\in G_K$, $u\in V$. Write $\mathcal E_\infty=\mathcal E_{U_\infty}$, $\mathcal E_\nu=\mathcal E_{U_\nu}$. Thus $\mathcal E_\infty(K,E)$ is the module of
$G_K$-invariants of $D^{(l)}_\infty(-E)$ in Definition~\ref{i5:x-def:bab-tower}. For a class $z\in H^1(G_{K,S},M_E)$ let $r(z)$ be the restriction of (the
inflation of) $z$ to $U_\infty$; as $U_\infty$ acts trivially on $M_E$, $r(z)$ is a well-defined homomorphism.

\begin{lemma}[the cyclotomic quotient]\label{i5:x-lem:K4-procyc}
Let $K$ be $l$-admissible and $c_K$ maximal with $\mu_{l^{c_K}}\subset K$.
\begin{enumerate}
\item $\chi(G_K)=1+l^{c_K}\Z_l\cong\Z_l$. If $\tilde\gamma\in G_K$ and $u=\chi(\tilde\gamma)$ has level exactly $c_K$, then $u$ topologically generates $\chi(G_K)$,
$u^{-E}\ne1$ for $E\ne0$, and for $\nu\ge c_K$, $\mathrm{Gal}(K(\mu_{l^\nu})/K)$ is cyclic of order $l^{\nu-c_K}$ generated by $\tilde\gamma$.
\item Let $C$ be the closure of $\langle\tilde\gamma\rangle$. Then $G_K=C\cdot U_\infty$, and every continuous homomorphism from $C\cap U_\infty$ to an
$l$-primary torsion group vanishes.
\item For every $w\mid l$, $\chi(D_w)$ is an open subgroup of $\chi(G_K)$; in particular it contains an element $u_w$ with $u_w^{-E}\ne1$ topologically
generating it, for $E\ne0$.
\end{enumerate}
\end{lemma}
\begin{proof}
(1) The equality is Proposition~\ref{i5:x-prop:bab-descent}(1). The group $1+l^{c_K}\Z_l$ is torsion free ($c_K\ge1$, $\ge2$ if $l=2$), so $u^{-E}\ne1$. The
order statement follows from $\mathrm{Gal}(K(\mu_{l^\nu})/K)\cong\chi(G_K)/(1+l^\nu\Z_l)$ and Lemma~\ref{i5:x-lem:K4-lte}.
(2) $\chi(C)$ is closed and contains $u$, so $\chi(C)=\chi(G_K)$ and $G_K=CU_\infty$. The procyclic group $C$ is the product $C_l\times C'$ of its pro-$l$
part and its prime-to-$l$ part. $\chi|_C$ kills $C'$ (no nontrivial continuous map from a prime-to-$l$ profinite group to the pro-$l$ group $\Z_l$) and
maps the procyclic pro-$l$ group $C_l$ onto $\Z_l$, hence isomorphically. So $C\cap U_\infty=\ker(\chi|_C)=C'$, which has no nontrivial continuous
homomorphism to an $l$-primary group.
(3) $\chi(D_w)$ is the image of the local cyclotomic character of $K_w$, which is open in $\Z_l^\times$ (the extension $K_w(\mu_{l^\infty})/K_w$ is infinite
and the image of local inertia has finite index, cf.\ the proof of Lemma~\ref{i5:lem-loch0}). It is a closed subgroup of $1+l^{c_K}\Z_l\cong\Z_l$, hence
$1+l^k\Z_l$ for some $k$, generated by any element of level $k$; such an element has infinite order.
\end{proof}

\begin{lemma}[descent]\label{i5:x-lem:K4-descent}
Let $K$ be $l$-admissible, $D$ an $l$-datum for $(K,\Omega)$ and $E\ne0$. Then $r$ is injective on $H^1(G_{K,S},M_E)$ and maps
$\Sha^1_D(K,M_E)$ isomorphically onto $\mathcal E_\infty(K,E)$. In particular $\operatorname{sha}(K,l,E)=\lgth_{\Z_l}\mathcal E_\infty(K,E)$.
\end{lemma}
\begin{proof}
Write $z$ also for a continuous cocycle on $G_K$ (inflated). For $k\in U_\infty$ (acting trivially) the cocycle identity applied to $gk=(gkg^{-1})g$ gives
\[
z(gkg^{-1})=\chi(g)^{-E}z(k)\qquad(g\in G_K,\ k\in U_\infty).\tag{$\ast$}
\]
\emph{Injectivity.} Let $r(z)=0$. Take $\tilde\gamma$, $u$ as in Lemma~\ref{i5:x-lem:K4-procyc}(1). As $u^{-E}-1\ne0$ and $\Q_l/\Z_l$ is divisible, there is $x$ with
$(u^{-E}-1)x=z(\tilde\gamma)$; replacing $z$ by $z-\partial x$ we may assume $z(\tilde\gamma)=0$ (still $z|_{U_\infty}=0$). Then $z(gk)=z(g)$ and
$z(kg)=z(g)$ for $k\in U_\infty$, so $z$ factors through $G_K/U_\infty\cong\chi(G_K)$; it vanishes on $\tilde\gamma^j$ for all $j\in\Z$ by the cocycle identity, hence on the
dense subgroup they generate, hence everywhere by continuity. So $z$ is a coboundary.

\emph{Image in $\mathcal E_\infty$.} Let $z\in\Sha^1_D$ and $\varphi=r(z)$; $\varphi$ is a continuous homomorphism and satisfies (iii) by ($\ast$). If $w$ is a finite place
not above $l$ then $I_w\cap U_\infty$ maps to $1$ in $G_{K,S}$, so $\varphi=0$ there. If $w\mid l$, then $D_w=hD_vh^{-1}$ for the chosen $D_v$ of $D$ and some $h$; on
$D_v$, $z(k')=k'y-y$ for some $y$, so $z(k')=0$ for $k'\in D_v\cap U_\infty$, and by ($\ast$) $z(hk'h^{-1})=0$. So $\varphi\in\mathcal E_\infty(K,E)$.

\emph{Surjectivity.} Let $\varphi\in\mathcal E_\infty(K,E)$, and $C$ as in Lemma~\ref{i5:x-lem:K4-procyc}(2). Define $z(\kappa k)=\kappa\cdot\varphi(k)$ for $\kappa\in C$,
$k\in U_\infty$. This is well defined: if $\kappa k=\kappa'k'$ then $\kappa'^{-1}\kappa=k'k^{-1}\in C\cap U_\infty$, on which $\varphi$ vanishes by
Lemma~\ref{i5:x-lem:K4-procyc}(2), and $\kappa'^{-1}\kappa$ acts trivially. It is continuous ($C\times U_\infty\to G_K$ is a closed continuous surjection of compact
spaces, hence a quotient map), and it is a cocycle: $\kappa k\kappa_1k_1=\kappa\kappa_1(\kappa_1^{-1}k\kappa_1)k_1$ and
$\varphi(\kappa_1^{-1}k\kappa_1)=\kappa_1^{-1}\cdot\varphi(k)$ by (iii). Clearly $z|_{U_\infty}=\varphi$. Local conditions: for a finite place $w\nmid l$,
$I_w\subset U_\infty$ ($K(\mu_{l^\infty})/K$ is unramified outside $l$), so $z|_{I_w}=0$; hence $z$ vanishes on the closed normal subgroup $N_S$ generated by these
inertia groups (on $N_S$, $z$ is a continuous homomorphism, conjugation-equivariant by ($\ast$)), so $z$ is inflated from $G_{K,S}$. For $v\in S$, $z|_{D_v}$ vanishes on
$D_v\cap U_\infty=\ker(\chi|_{D_v})$ by (ii), so the injectivity argument applied to the procyclic quotient $\chi(D_v)$ with the generator $u_v$ of
Lemma~\ref{i5:x-lem:K4-procyc}(3) shows that $z|_{D_v}$ is a coboundary. So $[z]\in\Sha^1_D$ and $r[z]=\varphi$.

Finally, the length of $\Sha^1_D(K,M_E)$ is $\operatorname{sha}(K,l,E)$ (Definition~\ref{i5:def-ldata} and the notation before Definition~\ref{i5:x-def:lemmaK-steps}).
\end{proof}

\begin{lemma}[descent to a finite layer]\label{i5:x-lem:K4-finite}
Let $K$ be $l$-admissible, $E\ne0$, $\varphi\in\mathcal E_\infty(K,E)$. There is $n_0$ such that for every $\nu\ge n_0$ there is $\psi\in\mathcal E_\nu(K,E)$ with
$\psi|_{U_\infty}=\varphi$ and $l^\nu\psi=0$.
\end{lemma}
\begin{proof}
We use repeatedly: $U_\infty=\bigcap_\nu U_\nu$ with $U_\nu$ closed and decreasing; hence if $Z\subset G_K$ is compact and $Z\cap U_\infty=\emptyset$, then
$Z\cap U_\nu=\emptyset$ for $\nu$ large.

(a) \emph{Extension as a homomorphism.} $\ker\varphi$ is open in $U_\infty$, so $\ker\varphi\supset U_\infty\cap V$ for an open normal subgroup $V$ of $G_K$.
$U_\infty V$ is open and contains $\bigcap_\nu U_\nu$, so (compactness) $U_n\subset U_\infty V$ for some $n$; then $U_n=U_\infty\cdot(V\cap U_n)$, and
$\psi_0(kv)=\varphi(k)$ ($k\in U_\infty$, $v\in V\cap U_n$) is a well-defined continuous homomorphism on $U_n$ extending $\varphi$ (using
$kvk_1v_1=kk_1(k_1^{-1}vk_1)v_1$ and normality of $V$).

(b) \emph{Local conditions.} Let $Z=\bigcup_{w\mid l}D_w$; it is compact (finitely many conjugacy classes of the compact groups $D_v$, $v\mid l$, under the
compact group $G_K$). $\psi_0$ is locally constant, so $\{x\in Z\cap U_n:\psi_0(x)\ne0\}$ is compact and, by (ii) for $\varphi$, disjoint from $U_\infty$; so it is
empty after restricting to some $U_{n_1}$. Inertia groups $I_w$ with $w\nmid l$ lie in $U_\infty$, where $\psi_0=\varphi$ vanishes on them; those with
$w\mid l$ lie in $Z$.

(c) \emph{Equivariance.} The map $(g,x)\mapsto\psi_0(gxg^{-1})-\chi(g)^{-E}\psi_0(x)$ on $G_K\times U_{n_1}$ is locally constant and vanishes for $x\in U_\infty$ by
(iii) for $\varphi$. The set where it is nonzero is compact and disjoint from $G_K\times U_\infty$, so it is empty on $G_K\times U_{n_2}$ for some $n_2$.

(d) \emph{Torsion.} $\psi_1=\psi_0|_{U_{n_2}}$ is locally constant on a compact group, so takes finitely many values and is killed by some $l^{t}$. Put
$n_0=\max(n_2,t)$; for $\nu\ge n_0$ the restriction $\psi=\psi_1|_{U_\nu}$ lies in $\mathcal E_\nu(K,E)$, extends $\varphi$, and $l^\nu\psi=0$.
\end{proof}

\begin{lemma}[class field theory at a finite layer]\label{i5:x-lem:K4-cft}
Let $K$ be $l$-admissible, $E\in\Z$, $\nu\ge1$, $F=K(\mu_{l^\nu})$ (so $U_\nu=G_F$, and $F$ is totally complex). Let $\mathrm{art}\colon G_F\to\mathrm{Cl}(F)$ be
the Artin map of the Hilbert class field. Then $\theta\mapsto\theta\circ\mathrm{art}$ is a bijection from the group $\mathcal C(F,E)$ of homomorphisms
$\theta\colon\mathrm{Cl}(F)\to\Q_l/\Z_l$ with $\theta([P])=0$ for $P\mid l$ and $\theta(\tau c)=\chi(g)^{-E}\theta(c)$ whenever $g\in G_K$ and $\tau=g|_F$, onto
$\mathcal E_\nu(K,E)$.
\end{lemma}
\begin{proof}
Class field theory for the Hilbert class field $H/F$ (Neukirch 1999, Ch.~VI, \S6):
$H/F$ is the maximal abelian extension unramified at all places, and the Artin map gives $\mathrm{Cl}(F)\cong\mathrm{Gal}(H/F)$ with $[P]\mapsto\mathrm{Frob}_P$.
Consequently $\mathrm{art}$ is surjective and continuous, kills every inertia group, maps $D_w$ onto $\langle[P]\rangle$ for $w\mid P$, and
$\mathrm{art}(gxg^{-1})=\tau(\mathrm{art}(x))$ for $g\in G_K$, $\tau=g|_F$ (transport of structure). Hence $\theta\circ\mathrm{art}\in\mathcal E_\nu(K,E)$ for
$\theta\in\mathcal C(F,E)$, and the map is injective by surjectivity of $\mathrm{art}$. Conversely let $\psi\in\mathcal E_\nu(K,E)$; its kernel is open and its fixed field
is abelian over $F$ and unramified at all finite places by (i), and at the archimedean places because $F$ is totally complex; so it lies in $H$ and
$\psi=\theta\circ\mathrm{art}$ for a homomorphism $\theta$. Then $\theta([P])=0$ for $P\mid l$ by (ii) and $\mathrm{art}(D_w)=\langle[P]\rangle$, and the equivariance of
$\theta$ follows from (iii) and surjectivity of $\mathrm{art}$.
\end{proof}

\begin{lemma}[the layer and the comparison map]\label{i5:x-lem:K4-layex}
Let $L$ be a number field, $c\ge1$ ($c\ge2$ if $l=2$) with $\mu_{l^c}\subset L$, $\mu_{l^{c+1}}\not\subset L$, $\nu\ge c$, $F=L(\mu_{l^\nu})$, $E\ge1$.
\begin{enumerate}
\item There are $g\in G_L$ with $\chi(g)=1+l^c$, a primitive $l^\nu$-th root of unity $\zeta\in\mathcal O_F^\times$, and $\gamma=g|_F$ such that
$\mathrm{Lay}(l,c,\nu,\zeta,\gamma)$ holds for $F/L$.
\item For such $g,\gamma$, $\theta\mapsto\big(P\mapsto\iota^{-1}\theta([P])\big)$, with $\iota\colon\Z/l^\nu\cong l^{-\nu}\Z_l/\Z_l$, is an injection from
$\{\theta\in\mathcal C(F,E):l^\nu\theta=0\}$ into $\mathrm{EqHom}(\gamma,(1+l^c)^E,l^\nu)$.
\end{enumerate}
\end{lemma}
\begin{proof}
(1) $L$ is $l$-admissible with $c_L=c$, so by Lemma~\ref{i5:x-lem:K4-procyc}(1) there is $g$ with $\chi(g)=1+l^c$, and $\gamma=g|_F$ satisfies $\gamma\zeta=\zeta^{1+l^c}$ for
any primitive $l^\nu$-th root $\zeta$. The order of $\gamma$ is the order of $1+l^c$ mod $l^\nu$, i.e.\ $l^{\nu-c}$ (Lemma~\ref{i5:x-lem:K4-lte}). Since
$\zeta^{l^{\nu-c}}\in\mu_{l^c}\subset L$, $[F:L]\le l^{\nu-c}$; so $F/L$ is Galois with group generated by $\gamma$, of order $l^{\nu-c}$.
(2) Let $\varphi_\theta(P)=\iota^{-1}\theta([P])$. For $\alpha\in F^\times$, $\sum_P\operatorname{ord}_P(\alpha)\varphi_\theta(P)=\iota^{-1}\theta([(\alpha)])=0$. As
$g^{-1}|_F=\gamma^{-1}$, $\theta([\gamma^{-1}P])=\chi(g^{-1})^{-E}\theta([P])=(1+l^c)^E\theta([P])$, the last step because $l^\nu\theta=0$ and
$\chi(g)\equiv1+l^c$ mod $l^\nu$. So $\varphi_\theta\in\mathrm{EqHom}$. Injectivity: the $[P]$ generate $\mathrm{Cl}(F)$.
\end{proof}

\begin{proposition}[SSha]\label{i5:x-prop:K4-ssha}
Statement (3) of Definition~\ref{i5:x-def:lemmaK-steps} holds.
\end{proposition}
\begin{proof}
Let $L$, $l$, $c$, $E\ge1$, $B$ be as in (3); $L$ is $l$-admissible with $c_L=c$. Fix an $l$-datum $D$ and put $\mathcal W=r(\Sha^1_D(L,M_E))\subset\mathcal E_\infty(L,E)$
(Lemma~\ref{i5:x-lem:K4-descent}). Let $T\subset\mathcal W$ be finite. By Lemma~\ref{i5:x-lem:K4-finite} choose $\nu\ge\max(c,2)$ such that each $\varphi\in T$ extends to
$\psi_\varphi\in\mathcal E_\nu(L,E)$ with $l^\nu\psi_\varphi=0$. Let $F=L(\mu_{l^\nu})$ and $g,\zeta,\gamma$ as in Lemma~\ref{i5:x-lem:K4-layex}(1). By
Lemma~\ref{i5:x-lem:K4-cft}, $\psi_\varphi=\theta_\varphi\circ\mathrm{art}$ with $\theta_\varphi\in\mathcal C(F,E)$, and $l^\nu\theta_\varphi=0$ because
$\theta\mapsto\theta\circ\mathrm{art}$ is injective. The composite $\varphi\mapsto\theta_\varphi\mapsto\varphi_{\theta_\varphi}$ (Lemma~\ref{i5:x-lem:K4-layex}(2)) is injective
($\theta_\varphi$ determines $\psi_\varphi$, which determines $\varphi=\psi_\varphi|_{U_\infty}$) from $T$ into $\mathrm{EqHom}(\gamma,(1+l^c)^E,l^\nu)$, whose
order is $\le B$ by hypothesis (applied to the layer of Lemma~\ref{i5:x-lem:K4-layex}(1)). So $\#T\le B$ for every finite $T\subset\mathcal W$; hence $\mathcal W$ is finite
with $\#\mathcal W\le B$. A finite $\Z_l$-module has length $k$ with $\#=l^k$; since $r$ is injective (Lemma~\ref{i5:x-lem:K4-descent}),
$\operatorname{sha}(L,l,E)=\lgth\mathcal W=k<\infty$ and $l^k=\#\mathcal W\le B$.
\end{proof}

\medskip\noindent\emph{SLayer.}

\begin{lemma}[discriminant of a cyclotomic step]\label{i5:x-lem:K4-disc}
Let $L'=L(\zeta)$ with $\zeta$ a primitive $N$-th root of unity. Then $|d_{L'}|\le|d_L|^{[L':L]}\,N^{[L':\Q]}$.
\end{lemma}
\begin{proof}
By the transitivity of discriminants and $\mathfrak d_{L'/L}=N_{L'/L}(\mathfrak D_{L'/L})$ (Neukirch 1999, Ch.~III, \S2),
$|d_{L'}|=\mathcal N(\mathfrak D_{L'/L})\,|d_L|^{[L':L]}$, $\mathcal N$ the absolute norm of the different $\mathfrak D_{L'/L}$. Let $f\in\mathcal O_L[X]$ be the minimal polynomial
of $\zeta$; since $L'=L(\zeta)$, $f'(\zeta)\in\mathfrak D_{L'/L}$ (Serre, \emph{Corps locaux}, Ch.~III, \S6). As $\mathcal O_L$ is integrally closed, $X^N-1=f h$ in
$\mathcal O_L[X]$; differentiating at $\zeta$, $N\zeta^{N-1}=f'(\zeta)h(\zeta)$, so $N=f'(\zeta)h(\zeta)\zeta\in\mathfrak D_{L'/L}$. Hence
$\mathcal N(\mathfrak D_{L'/L})$ divides $\mathcal N(N\mathcal O_{L'})=N^{[L':\Q]}$.
\end{proof}

\begin{lemma}[monotonicity of $\operatorname{sha}$ in a cyclotomic layer]\label{i5:x-lem:K4-mono}
Let $L$ be $l$-admissible, $E\ge1$ and $L'=L(\mu_{l^k})$ for some $k$. Then $\operatorname{sha}(L,l,E)\le\operatorname{sha}(L',l,E)$ in $\N\cup\{\infty\}$.
\end{lemma}
\begin{proof}
Embed $L'\subset\Omega$. Since $L'(\mu_{l^\infty})=L(\mu_{l^\infty})$, both fields have the same $U_\infty$, and $L'$ is $l$-admissible. The inertia and
decomposition groups of $L'$ are $I_w\cap G_{L'}$, $D_w\cap G_{L'}$, and $U_\infty\subset G_{L'}$, so conditions (i), (ii) are the same for $L$ and $L'$, while
(iii) for $L'$ only involves $g\in G_{L'}\subset G_L$. Hence $\mathcal E_\infty(L,E)\subset\mathcal E_\infty(L',E)$, and by Lemma~\ref{i5:x-lem:K4-descent} (for $L$ and for $L'$)
$\operatorname{sha}(L,l,E)=\lgth\mathcal E_\infty(L,E)\le\lgth\mathcal E_\infty(L',E)=\operatorname{sha}(L',l,E)$.
\end{proof}

\begin{proposition}[SLayer]\label{i5:x-prop:K4-slayer}
Statement (2) of Definition~\ref{i5:x-def:lemmaK-steps} holds.
\end{proposition}
\begin{proof}
Let $L$ be $l$-admissible, $E\ge1$, $c_0=c_L$, $n=n_L$. Let $m\ge1$ be minimal with $l^m\ge4E+4$, so $l^m\le l(4E+4)$. Put $L'=L(\mu_{l^{c_0+m}})$,
$d=l^m$, and $c=c_{L'}$ (maximal with $\mu_{l^c}\subset L'$), so $c\ge c_0+m$.
\begin{itemize}
\item $[L':L]=l^m=d$ by Lemma~\ref{i5:x-lem:K4-procyc}(1); so $n_{L'}=dn$.
\item $L'\supset L$ is $l$-admissible, hence totally complex; $c\ge c_0\ge1$ ($\ge2$ if $l=2$), $\mu_{l^c}\subset L'$, $\mu_{l^{c+1}}\not\subset L'$.
\item $2E+1<4E+4\le l^m\le l^c$; and $2r_2(L')=n_{L'}=l^mn\ge4E+4$, so $2E+1<r_2(L')$.
\item $l\le n+1$: for $l$ odd, $\Q(\mu_l)\subset L$ gives $l-1\mid n$; for $l=2$, $n\ge2$. So $d\le l(4E+4)\le(n+1)(4E+4)$.
\item $l^c\le2n_{L'}$ by Lemma~\ref{i5:x-lem:bab-layerdisc}(1).
\item $|d_{L'}|\le|d_L|^{d}(l^{c_0+m})^{n_{L'}}\le|d_L|^d(l^c)^{n_{L'}}$ by Lemma~\ref{i5:x-lem:K4-disc} with $N=l^{c_0+m}$.
\item $\operatorname{sha}(L,l,E)\le\operatorname{sha}(L',l,E)$ by Lemma~\ref{i5:x-lem:K4-mono}.\qedhere
\end{itemize}
\end{proof}




\begin{proof}[Proof of Proposition~\ref{i5:prop-shak}]
SLayer and SSha are Propositions~\ref{i5:x-prop:K4-slayer} and
\ref{i5:x-prop:K4-ssha}; SFew, SCount and STors are
Propositions~\ref{i5:x-prop:sfew}, \ref{i5:x-prop:scount} and
\ref{i5:x-prop:stors}.  Proposition~\ref{i5:x-prop:snum} gives $\mathrm{ShaK}(C)$.
\end{proof}

\subsection{The Case~I count}\label{i5:ssec-pf-counts}



\medskip\noindent\emph{Notation.} $K$ is an imaginary quadratic field with discriminant $d_K$ and class number $h_K$; $w$ is a finite place of $K$ of residue
characteristic $p$ and norm $q_w$. For a Dirichlet character $\psi$ of level $N$ with values in a coefficient ring $\mathcal O$ and $e\ge1$, put
$\Psi_{e,\psi}=\psi\,\chi_{\mathrm{cyc}}^{e}\colon G_\Q\to\mathcal O^\times$ ($\psi$ viewed as a Galois character through $\Q(\mu_N)$) and
\[
t_{e,\psi}(\sigma)=\Psi_{e,\psi}(\sigma)+1+\Psi_{e,\psi}(\sigma)^{-1}\qquad(\sigma\in G_K),
\]
the trace of $\Psi_{e,\psi}\oplus1\oplus\Psi_{e,\psi}^{-1}$. For $H\le(\Z/N)^\times$ let $\mathrm{orb}_{\mathcal O}(H)$ be the set of $\mathcal O$-valued characters
$\psi$ of $(\Z/N)^\times$ with $\ker\psi=H$. For $q\ge1$, $z\in\C^\times$ put $\mathrm{eis}_3(q,e,z)=q^{-e}z+1+q^ez^{-1}$.
Fix $\iota\colon\overline{\Q}_l\cong\C$. $\mathrm{Adm}(e,N)$ is the set of continuous representations $\rho\colon G_K\to\mathrm{GL}_3(\overline{\Q}_l)$ which are
irreducible, conjugate self-dual, of Hodge--Tate weights $(e,0,-e)$, automorphic via $\iota$, of level dividing $N$ (resp.\ $N|d_K|$ for those
arising from an endoscopic $2+1$ datum) and regular-ordinary at $l$.
Precisely:

\begin{definition}[level, admissibility]\label{i5:x-def:bab-adm}
Let $F$ be a number field, $\rho\colon G_F\to\mathrm{GL}_n(\overline{\Q}_l)$, $M\ge1$, $m=(m_1,\dots,m_n)\in\Z^n$. Say a character is \emph{of level $M$}
if it is trivial on $G_{F(\mu_M)}$.
\begin{enumerate}
\item[(a)] $\rho$ has \emph{level $(M;m)$} if $\rho$ is continuous; for every finite $w\nmid l$, $\rho|_{I_w}$ is diagonalisable with diagonal entries
the restrictions of characters of level $M$; and for every $w\mid l$ there are a place $A\mid w$ of $\overline{\Q}$, characters $\chi_i$ of level $M$
and characters $\kappa_i$ of $D_A$ such that $\rho|_{D_A}$ is conjugate to an upper triangular representation with diagonal $(\kappa_1,\dots,\kappa_n)$,
$\kappa_i=\chi_{\mathrm{cyc}}^{m_i}\chi_i$ on $I_A$, and either $\kappa_i\ne\kappa_j\chi_{\mathrm{cyc}}$ on $D_A$ for all $i,j$, or $\rho|_{D_A}$ is
potentially crystalline.
\item[(b)] $\rho$ is \emph{regular-ordinary at $l$ of level $(M;m)$} if for every $w\mid l$ there is $A\mid w$ with a flag as in (a) satisfying
$\kappa_i\ne\kappa_j\chi_{\mathrm{cyc}}$ on $D_A$ for all $i,j$.
\item[(c)] (stable) $\rho\colon G_K\to\mathrm{GL}_3(\overline{\Q}_l)$ is \emph{stably admissible} for $(e,N)$ if it is irreducible, has level
$(N;e,0,-e)$, is regular-ordinary at $l$ of level $(N;e,0,-e)$, and $\rho\otimes\chi_{\mathrm{cyc}}^{-1}$ is attached via $\iota$ to a regular algebraic
cuspidal representation $\pi$ of $\mathrm{GL}_3(\mathbb A_K)$ with $\pi^\vee\cong\pi^c$ (matching of Frobenius characteristic polynomials at almost all
places, after the normalisation $q_w^{(n-1)/2}$), together with the automorphy witness of (d).
\item[(d)] (endoscopic) $\rho$ is \emph{endoscopically admissible} for $(e,N|d_K|)$ if there is an irreducible $W\colon G_\Q\to\mathrm{GL}_2(\overline{\Q}_l)$
with $W\otimes\chi_{\mathrm{cyc}}^{-e}$ attached via $\iota$ to a newform of weight $2e+1$, satisfying the automorphy witness over every CM quadratic
extension, of level $(N|d_K|;e,-e)$, with $\det W=\chi_{\mathrm{cyc}}^{0}\cdot(\text{a character of finite level})$, and $\rho\cong W|_{G_K}\oplus1$
(the conjugating matrix places $W$ in the corner entries and $1$ in the middle).
\end{enumerate}
$\mathrm{Adm}(e,N)$ is the set of stably admissible or endoscopically admissible $\rho$.
\end{definition}

\begin{definition}[automorphy witness]\label{i5:x-def:witA}
Let $F$ be a CM field and $r\colon G_F\to\mathrm{GL}_n(\overline{\Q}_l)$ continuous. For a finite place $v$ of $F$ and a place $A\mid v$ of $\overline{\Q}$ with
decomposition group $D_A$, a \emph{flag} at $A$ is a tuple of characters $\kappa_1,\dots,\kappa_n\colon D_A\to\overline{\Q}_l^\times$ such that $r|_{D_A}$ is
conjugate to an upper triangular representation with diagonal $(\kappa_1,\dots,\kappa_n)$.
\begin{enumerate}
\item $r$ is \emph{unlinked at $v$} if for every $A\mid v$ there is a flag at $A$ with $\kappa_i\ne\kappa_j\chi_{\mathrm{cyc}}$ on $D_A$ for all $i,j$.
\item $r$ is \emph{potentially unlinked at $v$} if for every $A\mid v$ there is a flag at $A$ with $\kappa_i^m\ne(\kappa_j\chi_{\mathrm{cyc}})^m$ on $D_A$
for all $i,j$ and all $m\ge1$.
\item For a cuspidal automorphic representation $\pi$ of $\mathrm{GL}_n(\mathbb A_F)$, $t\in\R$ and a finite place $v$ with uniformiser $\varpi_v$ and
$q_v=Nv$, the \emph{window condition} $\mathrm{W}(\pi,t,v)$ says: the cuspidal support of $\pi_v$ consists of characters $\xi$ of $F_v^\times$, and
each such $\xi$ satisfies $q_v^{\,t-1/2}<|\xi(\varpi_v)|<q_v^{\,t+1/2}$.
\item The \emph{automorphy witness} $\mathrm{Wit}(r)$ is the conjunction of: (i) for every $v\mid l$, if $r$ is unlinked at $v$ then $r$ is potentially
unlinked at $v$; (ii) for every regular algebraic cuspidal $\pi$ of $\mathrm{GL}_n(\mathbb A_F)$ to which $r$ is attached via $\iota$ (equality of
Frobenius characteristic polynomials with Satake parameters at almost all places), there is $t\in\R$ such that $\mathrm{W}(\pi,t,v)$ holds for every
finite $v$ with $v\nmid l$ or $r$ potentially unlinked at $v$.
\end{enumerate}
In Definition~\ref{i5:x-def:bab-adm}(c) the automorphy predicate is: there is a regular algebraic cuspidal $\pi$ with $\pi^\vee\cong\pi^c$ to which
$r=\rho\otimes\chi_{\mathrm{cyc}}^{-1}$ is attached, and, $r$ being continuous, $\mathrm{Wit}(r)$; in (d) the requirement ``over every CM quadratic
extension'' means $\mathrm{Wit}(W|_{G_{K'}})$ for every CM quadratic $K'/\Q$ with $W|_{G_{K'}}$ irreducible (no existence of an attached
automorphic representation over $K'$ is required).
\end{definition}

\begin{remark}
Clause (ii) is uniform in $v$: one real number $t$ bounds the absolute values of all unramified-character Satake parameters of $\pi$ at once
(up to the factor $q_v^{\pm1/2}$), at all places away from $l$ and at the potentially unlinked places above $l$; the condition that the cuspidal support
consists of characters forces $\pi_v$ to be a subquotient of a principal series at those places.
\end{remark}

The condition $\mathrm{GD}$ is that of Definition~\ref{i5:def-gd}.

\begin{definition}\label{i5:x-def:gd-onechar}
\begin{enumerate}
\item $\mathrm{OneCharU}$: for every imaginary quadratic $K$, odd $l$, $\mathcal O$, $N$, $e$, $\psi$ with $\mathrm{GD}(K,l,e,\psi)$, and every $\sigma\in G_K$:
if $\operatorname{tr}\rho(\sigma)\ne t_{e,\psi}(\sigma)$ for every $\rho\in\mathrm{Adm}(e,N)$, then $V'(e,\psi)<\infty$ and there is a finite set
$X\subset\mathrm{Adm}(e,N)$ of pairwise non-isomorphic representations with
$\mathrm{Vn}(e,\psi)\le2\sum_{\rho\in X}v_l\big(\operatorname{tr}\rho(\sigma)-t_{e,\psi}(\sigma)\big)$.
\item $\mathrm{NumberU3U}$: for every imaginary quadratic $K$ there are $c,A\in\N$ such that for all $e,N\ge1$, every prime $p\ge7$ with $p\nmid Nd_K$,
every place $w\mid p$, every finite set $\mathrm{Orb}$ and map $z\colon\mathrm{Orb}\to\mu_\infty(\C)$ such that for every $\tau\in\mathrm{Aut}(\C)$ there is a bijection $u\colon\mathrm{Orb}\to\mathrm{Orb}$
with $z(u(o))=\tau(z(o))$ for all $o\in\mathrm{Orb}$, there is
$NN\in\Z\smallsetminus\{0\}$ with $\log|NN|\le c\,\#\mathrm{Orb}\,(N(e+1))^A\log p$ such that for every odd $l\ne p$, every Frobenius $\sigma\in G_K$ at $w$ and
every family $(X_o)_{o\in\mathrm{Orb}}$ of finite sets of pairwise non-isomorphic elements of $\mathrm{Adm}(e,N)$: $\operatorname{tr}\rho(\sigma^{-1})\ne\iota^{-1}\mathrm{eis}_3(q_w,e,z_o)$
for all $o$, $\rho\in X_o$, and $|NN|_l\le\prod_o\prod_{\rho\in X_o}|\operatorname{tr}\rho(\sigma^{-1})-\iota^{-1}\mathrm{eis}_3(q_w,e,z_o)|_l$.
\item $\mathrm{PropGK}$: for every number field $K$ there are $c_5,A\in\N$ such that for $e\ge1$, $N\ge1$, $H\le(\Z/N)^\times$ and a prime $p\nmid30N|d_K|$ there is
$NN\ne0$ with $\log|NN|\le c_5\varphi([(\Z/N)^\times:H])(N(e+1))^A\log p$ such that for every odd $l\ne p$ and $\mathcal O$: if some $\psi\in\mathrm{orb}_{\mathcal O}(H)$
satisfies GD, then $V'(e,\psi)<\infty$ for all $\psi\in\mathrm{orb}_{\mathcal O}(H)$ and $\sum_{\psi\in\mathrm{orb}_{\mathcal O}(H)}\mathrm{Vn}(e,\psi)\le2v_l(NN)$.
\item $\mathrm{CorGK}$: for every number field $K$ there are $c_5,A$ such that for $e,N,H$, distinct primes $p,p'\nmid30N|d_K|$, coefficient rings
$(\mathcal O_l)_l$ and a finite set $\mathcal L$ of primes each of which admits a GD character in $\mathrm{orb}_{\mathcal O_l}(H)$: all $V'(e,\psi)$ are finite and
$\sum_{l\in\mathcal L}\log l\sum_{\psi\in\mathrm{orb}_{\mathcal O_l}(H)}\mathrm{Vn}(e,\psi)\le2c_5\varphi([(\Z/N)^\times:H])(N(e+1))^A(\log p+\log p')$.
\end{enumerate}
\end{definition}

\begin{lemma}\label{i5:x-lem:orb}
Let $G$ be a finite abelian group, $H\le G$, and $R$ a domain containing the $\#G$-th roots of unity.
(a) If $\psi,\psi'\colon G\to R^\times$ have kernel $H$, then $\psi'=\psi^j$ with $\gcd(j,[G:H])=1$, and $\operatorname{ord}\psi=[G:H]$.
(b) The number of characters with kernel exactly $H$ is $\varphi([G:H])$ if $G/H$ is cyclic and $0$ otherwise.
(c) If $T$ is a set of subgroups $H$ with $G/H$ cyclic, then $\sum_{H\in T}\varphi([G:H])\le\#G$.
\end{lemma}
\begin{proof}
A character with kernel $H$ is a faithful character of $G/H$, whose image is a finite subgroup of $R^\times$, hence cyclic; so $G/H$ is cyclic of order
$\operatorname{ord}\psi$, and faithful characters of a cyclic group of order $n$ are the $\psi^j$, $\gcd(j,n)=1$, $\varphi(n)$ in number (they exist since $R$
contains the $n$-th roots of unity). This gives (a) and (b). For (c), the sets of characters with kernel exactly $H$, $H\in T$, are disjoint subsets of
the dual group of $G$, which has $\#G$ elements.
\end{proof}

\begin{lemma}\label{i5:x-lem:gd-kernel}
If $\psi$ satisfies $\mathrm{GD}(K,l,e,\psi)$ and $\ker\psi'=\ker\psi$, then $\psi'$ satisfies $\mathrm{GD}(K,l,e,\psi')$.
\end{lemma}
\begin{proof}
By Lemma~\ref{i5:x-lem:orb}(a), $\psi'=\psi^j$ with $j$ prime to $\operatorname{ord}\psi=\operatorname{ord}\psi'$. The conditions not involving $\psi$ are unchanged, and
$l\nmid\operatorname{ord}\psi'$. The component of $\psi'$ at $l$ is the $j$-th power of that of $\psi$, so its square is trivial. If $\psi(-1)=1$ then
$\psi'(-1)=1$; if $\psi(-1)=-1$ then $\operatorname{ord}\psi$ is even, $j$ is odd and $\psi'(-1)=-1$. Finally $\chi'_l$ is multiplicative in the character, so
$\chi'_l(\psi')=\chi'_l(\psi)^j$ with $j$ prime to its order, and the two have the same kernel.
\end{proof}

\begin{lemma}\label{i5:x-lem:sumlog}
For $n\in\Z\smallsetminus\{0\}$ and a finite set $\mathcal L$ of primes, $\sum_{l\in\mathcal L}v_l(n)\log l\le\log|n|$.
\end{lemma}
\begin{proof}
The prime powers $l^{v_l(n)}$, $l\in\mathcal L$, are pairwise coprime divisors of $n$, so their product divides $n$ and is $\le|n|$; take logarithms.
\end{proof}

\begin{lemma}\label{i5:x-lem:propGK}
$\mathrm{NumberU3U}\wedge\mathrm{OneCharU}\Rightarrow\mathrm{PropGK}$.
\end{lemma}
\begin{proof}
If $K$ is not imaginary quadratic, no character satisfies GD, and $NN=1$, $c_5=A=0$ work. Let $K$ be imaginary quadratic, $c,A$ as in NumberU3U, $c_5=c$,
and let $e,N,H,p$ be given. Then $p\nmid N$, $p\nmid d_K$ and, since $p\nmid30$, $p\ge7$. Choose $w\mid p$; $q_w$ is a power of $p$, hence prime to $N$.
Let $\mathrm{Orb}$ be the set of complex characters of $(\Z/N)^\times$ with kernel $H$, and $z_\chi=\chi(q_w)^{-1}$, a root of unity; the image of $z$ is permuted
by $\mathrm{Aut}(\C)$ since $\tau\circ\chi$ has the same kernel as $\chi$. NumberU3U gives $NN$, and $\#\mathrm{Orb}\le\varphi([(\Z/N)^\times:H])$ by
Lemma~\ref{i5:x-lem:orb}(b), whence the bound on $\log|NN|$.

Let $l\ne p$ be odd and $\psi_0\in\mathrm{orb}_{\mathcal O}(H)$ satisfy GD. By Lemma~\ref{i5:x-lem:gd-kernel} every $\psi\in\mathrm{orb}_{\mathcal O}(H)$ satisfies GD, and
$E\colon\psi\mapsto\iota\circ\psi$ is a bijection $\mathrm{orb}_{\mathcal O}(H)\to\mathrm{Orb}$ (Lemma~\ref{i5:x-lem:orb}(a)). Let $\sigma\in G_K$ be a Frobenius at a place
of $\overline{\Q}$ above $w$. Since $w\nmid lN$, $\sigma$ acts on $\mu_N$ and on $\mu_{l^\infty}$ by $\zeta\mapsto\zeta^{q_w}$; hence
$\Psi_{e,\psi}(\sigma)=q_w^e\psi(q_w)$ and $\iota\,t_{e,\psi}(\sigma^{-1})=\mathrm{eis}_3(q_w,e,z_{E\psi})$. The first assertion of NumberU3U, applied to
singletons, shows $\operatorname{tr}\rho(\sigma^{-1})\ne t_{e,\psi}(\sigma^{-1})$ for every $\rho\in\mathrm{Adm}(e,N)$. Hence OneCharU applies to $(\psi,\sigma^{-1})$:
$V'(e,\psi)<\infty$ and there are finite sets $X_\psi$ with $\mathrm{Vn}(e,\psi)\le2\sum_{\rho\in X_\psi}v_l(\operatorname{tr}\rho(\sigma^{-1})-t_{e,\psi}(\sigma^{-1}))$.
The second assertion of NumberU3U for $X_o=X_{E^{-1}o}$ gives that all these differences are nonzero and $|NN|_l\le\prod_o\prod_\rho|\cdot|_l$; taking
$-\log_l$, $\sum_o\sum_\rho v_l(\cdot)\le v_l(NN)$. Summing the bounds for $\mathrm{Vn}(e,\psi)$ over $\psi$ and reindexing by $E$ gives
$\sum_{\psi}\mathrm{Vn}(e,\psi)\le2v_l(NN)$.
\end{proof}

\begin{lemma}\label{i5:x-lem:corGK}
$\mathrm{PropGK}\Rightarrow\mathrm{CorGK}$.
\end{lemma}
\begin{proof}
Take $c_5,A$ from PropGK. If $\mathcal L=\emptyset$ there is nothing to prove; otherwise $e\ge1$ (part of GD). PropGK at $p$ and at $p'$ gives
$NN,NN'\ne0$. Let $l\in\mathcal L$. If $l\ne p,p'$ both apply; if $l=p$ then $l\ne p'$ and the one at $p'$ applies; symmetrically if $l=p'$. In each case
the $V'(e,\psi)$ are finite and, valuations being $\ge0$, $\sum_\psi\mathrm{Vn}(e,\psi)\le2v_l(NN)+2v_l(NN')$. Multiplying by $\log l$, summing over $\mathcal L$
and using Lemma~\ref{i5:x-lem:sumlog} gives the bound $2\log|NN|+2\log|NN'|$, and PropGK bounds both logarithms.
\end{proof}

\begin{lemma}\label{i5:x-lem:gd-host}
Let $h=(\psi,e,q_1,d)$ be a host with $\mathrm{Big}_h(t)$, $0<e<t$, $l$ an odd prime in Case~I, $\psi'\in X_h(\mathcal O,e)$, and $K=\Q(\sqrt{-p_K(\psi)})$.
Then $\mathrm{GD}(K,l,e,\psi')$ holds.
\end{lemma}
\begin{proof}
We check the conditions of GD in turn. $N_h$ is even since $2^{e+2}\mid N_h$. Since $l\notin B_\psi$ (Case~I), Lemma~\ref{i5:lem-hostarith} gives
$l\nmid\operatorname{ord}\psi'$ and that the component of $\psi'$ at $l$ has trivial square (in its proof for $l\notin B_\psi$). The parity $\psi'(-1)=(-1)^e$ is the
definition of $X_h(\mathcal O,e)$. $K$ is imaginary quadratic and, as $p_K\equiv3\pmod4$, $d_K=-p_K$; by Lemma~\ref{i5:lem-hostarith}, $l\ne p_K$ and $l\nmid h^*=h_K$, so
$l\nmid|d_K|h_K$. Next $p_K\nmid N_h$: $p_K$ is odd, $p_K\nmid N_\psi$ by definition, and $p_K\in B_\psi$ so $p_K<q_1,d$ by Lemma~\ref{i5:lem-hostarith}; hence
$\gcd(|d_K|,N_h)=1$. For the non-containment condition, let $I$ be an inertia group at a place of $\overline{\Q}$ above $p_K$. Since $p_K\ne l$, $I$ acts
trivially on $\mu_l$; since $p_K\nmid N_h$ and $\chi'_l$ factors through $\mathrm{Gal}(\Q(\mu_{N_h})/\Q)$, $I\subseteq\ker\chi'_l$. So
$I\subseteq\ker(\chi_{\mathrm{cyc}}\bmod l)\cap\ker\chi'_l$. As $p_K$ ramifies in $K$, $I\not\subseteq G_K$, which is the required non-containment. Finally
$(l-1)\nmid2e$ is part of Case~I.
\end{proof}

\begin{definition}[$\mathrm{KatoTwoW}$]\label{i5:x-def:katotwow}
$\mathrm{KatoTwoW}$ is the statement: let $l$ be a prime, $\iota\colon\overline{\Q}_l\cong\C$, $K$ a number field, $e,N\in\N$,
$\theta,\psi\colon G_\Q\to\overline{\Q}_l^\times$ homomorphisms and $W\colon G_\Q\to\mathrm{GL}_2(\overline{\Q}_l)$, such that, with
$\chi=\chi_{\mathrm{cyc}}$ and $\Psi=\chi^e\psi$:
\begin{enumerate}
\item $l$ is odd, $e\ge1$, $(l-1)\nmid2e$;
\item $K$ is imaginary quadratic, $l\nmid|d_K|h_K$; $N\ge1$ is even and $\gcd(|d_K|,N)=1$;
\item $\ker\theta=G_K$ and $\theta^2=1$;
\item $\psi$ is trivial on $G_{\Q(\mu_N)}$, has finite order prime to $l$, $\psi^2$ is trivial on the inertia groups above $l$, and
$\chi(c)^e\psi(c)=1$ for some complex conjugation $c$;
\item $W$ has level $(N|d_K|;\,e,-e)$ in the sense of Definition~\ref{i5:x-def:bab-adm}(a) below, and at some place above $l$ the inertia
acts, on an open subgroup of inertia, with characteristic polynomial $(X-\chi^{e})(X-\chi^{-e})$;
\item $\det W=\theta$; $\operatorname{tr}W(s)=\Psi(s)+\Psi(s)^{-1}$ for $s$ in inertia at $l$; and
$|\operatorname{tr}W(g)-(\theta(g)\Psi(g)+\Psi(g)^{-1})|<1$ for all $g\in G_\Q$ (residual congruence);
\item (ordinary flag) at some place $w\mid l$, $W|_{D_w}$ is upper triangular with diagonal characters $\kappa_0,\kappa_1$, there is a character
$\chi'$ trivial on $G_{\Q(\mu_l)}$ with $\kappa_0=\chi^{e}\chi'$, $\kappa_1=\chi^{-e}\chi'$ on $I_w$, and $I_w\subset\ker\theta$;
\item $W$ restricted to $G_{K(\mu_{l^\infty})}$ is irreducible;
\item there is a newform $f$ of weight $2e+1$ such that $W\otimes\chi^{-e}$ is attached to $f$ via $\iota$, and $W$ satisfies the automorphy
witness condition over every CM quadratic extension (Definition~\ref{i5:x-def:bab-adm}(d)).
\end{enumerate}
Then $H^1_f(\Q,W)=0$ and $H^1_f(\Q,W\otimes\theta)=0$ (Bloch--Kato Selmer groups with $\overline{\Q}_l$-coefficients, i.e.\ every continuous
cocycle satisfying the Bloch--Kato local conditions is a coboundary).
\end{definition}

$\mathrm{OneCharU}$ is Proposition~\ref{i5:x-prop:onechar}, whose reduction is given in \S\ref{i5:x5-ssec}.


\medskip\noindent\emph{The Kato-type vanishing.}

We show that $\mathrm{KatoTwoW}$ (Definition~\ref{i5:x-def:katotwow}) is a
consequence of Kato's finiteness theorem.  Recall Kato's normalisation
(Kato 2004, \S\S8, 14).  Let $f$ be a normalised newform of weight $k\ge2$
on $\Gamma_1(M)$ with nebentypus $\varepsilon$ and coefficient field
$F=\Q(a_n(f))$, and let $\lambda$ be a finite place of $F$ above $l$.  Kato's
$V_{F_\lambda}(f)$ is the $2$-dimensional $F_\lambda$-linear representation of
$G_\Q$ cut out by $f$ in $H^1_{\mathrm{et}}(Y_1(M)_{\overline\Q},
\mathrm{Sym}^{k-2})\otimes F_\lambda$; it is unramified at every prime
$p\nmid Ml$, and there
\begin{equation}\label{i5:x-eq:kato-frob}
\det\bigl(1-\mathrm{Fr}_p^{-1}u\mid V_{F_\lambda}(f)\bigr)=1-a_p(f)u+\varepsilon(p)p^{k-1}u^2
\end{equation}
with $\mathrm{Fr}_p$ the arithmetic Frobenius (Kato 2004, 14.10).  For a
finite extension $K$ of $\Q$ and a $G_\Q$-stable $\mathcal O_\lambda$-lattice
$T$ of a representation $V$, $\mathrm{Sel}(K,T)\subseteq H^1(K,V/T)$ is the
Bloch--Kato Selmer group: the classes whose restriction to every place $v$ of
$K$ lies in the image of $H^1_f(K_v,V)$, where $H^1_f(K_v,V)=0$ for $v$
archimedean, is the unramified subgroup for $v\nmid l$, and is
$\ker(H^1(K_v,V)\to H^1(K_v,B_{\mathrm{crys}}\otimes V))$ for $v\mid l$
(Kato 2004, 14.1).

\begin{theorem}[Kato 2004, \emph{Ast\'erisque} 295, Thm.~14.2(1)]\label{i5:x-thm:kato}
Let $f$ be a normalised newform of weight $k\ge2$ and $K$ a finite abelian
extension of $\Q$.  Let $r$ be an integer with $1\le r\le k-1$ and $r\ne k/2$.
Then for every finite place $\lambda$ of $F$ and every $G_\Q$-stable
$\mathcal O_\lambda$-lattice $T$ of $V_{F_\lambda}(f)(r)$, $\mathrm{Sel}(K,T)$
is finite.
\end{theorem}

(See Ast\'erisque 295 (2004), \S14.  Part (2) of the theorem, which needs
$k$ even and $L(f,\chi,k/2)\ne0$, is not used.)  No $L$-value hypothesis enters part (1); this is what makes it
applicable here, since our weight $k=2e+1$ is odd and $r=e$ is never the
centre.

\begin{lemma}\label{i5:x-lem:sel-fin-vanish}
Let $V$ be a finite-dimensional $F_\lambda$-representation of $G_K$,
unramified outside a finite set, with lattice $T$.  If $\mathrm{Sel}(K,T)$ is
finite then $H^1_f(K,V)=0$ (the classes in $H^1(K,V)$ that lie in
$H^1_f(K_v,V)$ at every place).  For a finite extension $E/F_\lambda$,
$H^1_f(K,V\otimes_{F_\lambda}E)=H^1_f(K,V)\otimes_{F_\lambda}E$.
\end{lemma}

\begin{proof}
Classes of $H^1_f(K,V)$ are unramified outside the finite set $S$ of places
above $l$ and of ramification of $V$, so $H^1_f(K,V)\subseteq
H^1(G_{K,S},V)=H^1(G_{K,S},T)\otimes\Q_l$, and $H^1(G_{K,S},T)$ is a finitely
generated $\mathcal O_\lambda$-module.  The map $H^1(K,V)\to H^1(K,V/T)$ sends
$H^1_f(K,V)$ into $\mathrm{Sel}(K,T)$ (by the definition of the local
conditions of $\mathrm{Sel}$), and its image is divisible; a divisible
subgroup of a finite group is $0$.  So $H^1_f(K,V)$ lies in the image of
$H^1(G_{K,S},T)$, a finitely generated $\mathcal O_\lambda$-module, and a
$\Q_l$-vector space inside such a module is $0$.  For the second assertion,
$V\otimes E\cong V^{[E:F_\lambda]}$ as $F_\lambda[G_K]$-modules and all local
conditions are additive.
\end{proof}

\begin{proposition}[Kato-type vanishing]\label{i5:x-prop:katotwow}
$\mathrm{KatoTwoW}$ holds.  In fact its conclusion follows from conditions
(1) (only $e\ge1$), (3), the continuity of $W$ contained in (5), and the first
half of (9) of Definition~\ref{i5:x-def:katotwow}, by
Theorem~\ref{i5:x-thm:kato}.
\end{proposition}

\begin{proof}
\emph{Coefficients.}  $G_\Q$ is compact and $\overline\Q_l^2$ is the countable
union of the closed subspaces $E^2$, $E$ running over finite extensions of
$\Q_l$ inside $\overline\Q_l$; by Baire's theorem the image of $W$, and the
image of any continuous cocycle $c\colon G_\Q\to\overline\Q_l^2$, lie in
$\mathrm{GL}_2(E)$, resp.\ $E^2$, for one finite $E$ (a preimage $c^{-1}(E^2)$
has interior, hence contains a coset of an open subgroup, and finitely many
cosets cover $G_\Q$).  So it suffices to prove the vanishing with coefficients
in every sufficiently large finite $E$, and a coboundary over $E$ is a
coboundary over $\overline\Q_l$.  The local conditions descend to $E$: they
are stated with $\overline\Q_l$-coefficients (the witness of a coboundary on
inertia, resp.\ in $W\otimes B_{\mathrm{cris}}$, has entries in
$\overline\Q_l$, resp.\ in $\overline\Q_l\otimes B_{\mathrm{cris}}$), and applying an
$E$-linear retraction $\overline\Q_l\to E$ to the witness gives a witness over
$E$, because the retraction commutes with the Galois action on $W$ (whose
matrices have entries in $E$) and on $B_{\mathrm{cris}}$.

\emph{Identification with Kato's representation.}  Let $k=2e+1\ge3$ and let
$f$ be the newform of (9).  Via $\iota^{-1}$, $F=\Q(a_n(f))$ embeds in
$\overline\Q_l$; let $\lambda$ be the induced place and $V'=V_{F_\lambda}(f)
\otimes_{F_\lambda}\overline\Q_l$.  Put $W'=W\otimes\chi^{-e}$.  By (9)
($W'$ attached to $f$), for almost all primes $p$, $W'$ is unramified at $p$
and the characteristic polynomial of $W'(\sigma^{-1})$, for every Frobenius
$\sigma$ at $p$, is $\iota^{-1}$ of $X^2-a_p(f)X+\varepsilon(p)p^{k-1}$; by
\eqref{i5:x-eq:kato-frob} the same holds for $V'$.  Hence
$\operatorname{tr}W'=\operatorname{tr}V'$ on all Frobenius elements at almost
all primes; these are dense in $G_\Q$ (Chebotarev, Neukirch VII (13.4)), and
traces are continuous (by the continuity of $W$ in (5)), so
$\operatorname{tr}W'=\operatorname{tr}V'$.  ``Attached'' in (9) includes, in
the formal definition, that $W'$ is semisimple.  By the Brauer--Nesbitt
theorem (reference to be confirmed) $V'^{\mathrm{ss}}\cong W'$; and $V'$ is
irreducible, $V_{F_\lambda}(f)$ being absolutely irreducible for a newform of
weight $\ge2$ (Ribet; reference to be confirmed).  Hence
\[
 W\cong V_{F_\lambda}(f)(e)\otimes_{F_\lambda}\overline\Q_l ,
\]
where $(e)$ is the twist by $\chi^e$.  Note $1\le e\le k-1=2e$ and $e\ne k/2$
as $k$ is odd.

\emph{$H^1_f(\Q,W)=0$.}  Theorem~\ref{i5:x-thm:kato} with $K=\Q$, $r=e$ gives
$\mathrm{Sel}(\Q,T)$ finite for a lattice $T$ of $V_{F_\lambda}(f)(e)$;
Lemma~\ref{i5:x-lem:sel-fin-vanish} gives $H^1_f(\Q,V_{F_\lambda}(f)(e)
\otimes E)=0$ for every finite $E$, and $H^1_f$ is invariant under
isomorphism of representations.

\emph{$H^1_f(\Q,W\otimes\theta)=0$.}  By (3), $\theta$ is the quadratic
character of the imaginary quadratic field $K$ of (2), which is abelian over
$\Q$.  Theorem~\ref{i5:x-thm:kato} over $K$ and
Lemma~\ref{i5:x-lem:sel-fin-vanish} give $H^1_f(K,W)=0$.  Since $\theta$ is
trivial on $G_K$, restriction is a map $H^1(\Q,W\otimes\theta)\to H^1(K,W)$.
It sends $H^1_f$ into $H^1_f$: at a place $v$ of $K$ above $p\ne l$, the
inertia group of $v$ is contained in that of $p$, so unramified classes
restrict to unramified classes; above $l$, the crystalline condition
$\ker(H^1(\cdot,V)\to H^1(\cdot,B_{\mathrm{crys}}\otimes V))$ is functorial
for restriction to the open subgroup $G_{K_v}\subseteq G_{\Q_l}$.  The
condition over $\Q$ is imposed at the place $v_0$ of $K$ above $l$ fixed by
the embedding defining $G_{\Q_l}\subseteq G_\Q$; if $l$ splits in $K$, the
other place above $l$ is $gv_0$ for some $g\in G_\Q\smallsetminus G_K$, its
decomposition group is $gG_{K_{v_0}}g^{-1}$, and conjugation by $g$ carries
the restricted class at $v_0$ to the restricted class at $gv_0$ and preserves
$H^1_f$ (the crystalline condition is invariant under conjugation), so the
condition holds there too.  It is
injective, since corestriction composed with restriction is multiplication by
$[K:\Q]=2$, invertible on a $\overline\Q_l$-vector space.  Hence
$H^1_f(\Q,W\otimes\theta)=0$.

The local conditions in Definition~\ref{i5:x-def:katotwow} are exactly the
Bloch--Kato ones used above (cocycles which are coboundaries on every inertia
group away from $l$, and crystalline classes at the place above $l$; there are
no conditions at infinity, and $H^1$ of $G_\R$ with $l$ odd vanishes anyway).
\end{proof}

\begin{remark}\label{i5:x-rem:kato-diff}
Comparison with the formal proof.  \cite{CharacterFamily} does not use
Theorem~\ref{i5:x-thm:kato}: it reconstructs Kato's Euler system (zeta
elements from Siegel units and their explicit reciprocity law) and runs Kolyvagin's argument.  The proof above
replaces all of that by Kato's published theorem.  The comparison of
hypotheses is: conditions (2) (except that $K$ is imaginary quadratic, used
only to know $K/\Q$ is abelian), (4), (5) except the continuity of $W$,
(6), (7), (8), the condition $(l-1)\nmid2e$ and the parity of $l$ in
(1), and the automorphy witness in (9) are \emph{not needed} for the
conclusion; they are needed by the formal route (ordinary normal form, an
$l$-ordinary source of level prime to $l$) and by the other uses of these
conditions in \S\ref{i5:ssec-counts}.  Nothing in Kato's statement is missing
from (1)--(9): Kato needs only a newform of weight $\ge2$, an abelian base
field and $1\le r\le k-1$, $r\ne k/2$.  The two points to keep in mind are the
normalisations: ``attached'' in (9) matches characteristic polynomials of
\emph{geometric} Frobenius with the Hecke polynomial, which is Kato's
\eqref{i5:x-eq:kato-frob}, so $W\otimes\chi^{-e}$ is Kato's $V(f)$ itself and
not its dual; and the Hodge--Tate weights $(e,-e)$ of $W$ then correspond to
$r=e$.

Two points should be noted.  First, the comparison of the statement of
Theorem~\ref{i5:x-thm:kato} with Kato 2004, 14.1 (definition of the Selmer group),
14.2(1) and 14.10 (Frobenius characteristic polynomial), is a reference to be
confirmed.  Second, the identification of the period ring used in the formal
statement of the crystalline condition with Fontaine's $B_{\mathrm{cris}}$ rests on
its interface in the formal development (a ring isomorphism compatible with the
Galois action, Frobenius and filtration); we do not describe its construction.
\end{remark}

\medskip\noindent\emph{The counting estimate $\mathrm{NumberU3U}$.}

The estimate splits into an elementary norm argument, proved here, and a
count of automorphic representations, which is formally verified.

\begin{definition}[trace data]\label{i5:x-def:trdata}
Let $q,e\in\N$.  A \emph{trace datum} for $(q,e)$ is a finite set $S$ with a
map $t\colon S\to\C$ such that: for every $\tau\in\mathrm{Aut}(\C)$ there is
a permutation $u_\tau$ of $S$ with $t(u_\tau s)=\tau(t(s))$; $|t(s)|\le\sqrt
q+2$; and $q^et(s)$ is an algebraic integer.  Given a group $\Gamma$,
$g\in\Gamma$, a prime $l$, $\iota\colon\overline\Q_l\cong\C$ and a set
$\mathcal A$ of representations $\Gamma\to\mathrm{GL}_3(\overline\Q_l)$, the
datum \emph{classifies $\mathcal A$ at $g$} if there is
$\mathrm{cl}\colon\mathcal A\to S$ with $\mathrm{cl}(\rho)=\mathrm{cl}(\rho')
\Rightarrow\rho\cong\rho'$ and $\operatorname{tr}\rho(g)=\iota^{-1}
t(\mathrm{cl}(\rho))$.
\end{definition}

\begin{lemma}\label{i5:x-lem:eisdiff}
Let $q\ge7$, $e\ge1$, $z$ a root of unity and $t\in\C$ with $|t|\le\sqrt q+2$
and $q^et$ integral.  Then $t\ne\mathrm{eis}_3(q,e,z)$, and
$F=q^e(t-\mathrm{eis}_3(q,e,z))$ is a non-zero algebraic integer with
$\log|F|\le2e\log q+1$.
\end{lemma}

\begin{proof}
Put $P=q^e\ge7$.  As $|z|=1$, $|\mathrm{eis}_3(q,e,z)|\ge P-1-P^{-1}>\sqrt
P+2\ge\sqrt q+2\ge|t|$, the middle inequality because $P-\sqrt P-3-P^{-1}>0$
for $P\ge7$.  Next $F=q^et-z-P-P^2z^{-1}$ is a sum of algebraic integers
($z^{-1}$ is a power of $z$).  Finally $|F|\le P(\sqrt q+2)+1+P+P^2\le2P^2$,
since $P(\sqrt P+3)+1\le P^2$ for $P\ge7$; so $\log|F|\le2e\log q+\log2$.
\end{proof}

\begin{lemma}\label{i5:x-lem:intprod}
Let $Y$ be finite and $F\colon Y\to\C$ with algebraic integer values such that
every $\tau\in\mathrm{Aut}(\C)$ permutes the family $(F(y))_y$.  Then
$\prod_yF(y)\in\Z$.
\end{lemma}

\begin{proof}
The product is an algebraic integer fixed by $\mathrm{Aut}(\C)$, hence in
$\Q$, hence in $\Z$.
\end{proof}

\begin{proposition}[norm bound]\label{i5:x-prop:normbound}
Let $q\ge7$, $e\ge1$, $(S,t)$ a trace datum for $(q,e)$, $\mathrm{Orb}$ a
finite set and $z\colon\mathrm{Orb}\to\mu_\infty(\C)$ whose family of values is
permuted by $\mathrm{Aut}(\C)$.  There is $NN\in\Z\smallsetminus\{0\}$ with
$\log|NN|\le\#\mathrm{Orb}\cdot\#S\cdot(2e\log q+1)$ such that for every prime
$l\nmid q$, every $\iota$, $\Gamma$, $g$, every set $\mathcal A$ classified by
$(S,t)$ at $g$, and every family $(X_o)_{o\in\mathrm{Orb}}$ of finite sets of
pairwise non-isomorphic elements of $\mathcal A$: for $\rho\in X_o$,
$\operatorname{tr}\rho(g)\ne\iota^{-1}\mathrm{eis}_3(q,e,z_o)$, and
\[
 |NN|_l\le\prod_{o}\prod_{\rho\in X_o}\bigl|\operatorname{tr}\rho(g)-\iota^{-1}\mathrm{eis}_3(q,e,z_o)\bigr|_l .
\]
\end{proposition}

\begin{proof}
Put $F(o,s)=q^e(t(s)-\mathrm{eis}_3(q,e,z_o))$ and $NN=\prod_{o,s}F(o,s)$.
Each factor is a non-zero algebraic integer with $\log|F|\le2e\log q+1$
(Lemma~\ref{i5:x-lem:eisdiff}); since
$\tau\,\mathrm{eis}_3(q,e,w)=\mathrm{eis}_3(q,e,\tau w)$, $\tau$ permutes the
factors, so $NN\in\Z\smallsetminus\{0\}$ (Lemma~\ref{i5:x-lem:intprod}) with
the stated bound.  For $\rho\in X_o$, $\operatorname{tr}\rho(g)=\iota^{-1}
t(\mathrm{cl}\rho)\ne\iota^{-1}\mathrm{eis}_3$.  All factors have
$|\iota^{-1}F|_l\le1$; $\mathrm{cl}$ is injective on each $X_o$ (its elements
are pairwise non-isomorphic), so dropping the factors with
$s\notin\mathrm{cl}(X_o)$ gives $|NN|_l\le\prod_o\prod_{\rho\in X_o}
|\iota^{-1}F(o,\mathrm{cl}\rho)|_l$, and $|\iota^{-1}F(o,\mathrm{cl}\rho)|_l=
|\operatorname{tr}\rho(g)-\iota^{-1}\mathrm{eis}_3(q,e,z_o)|_l$ as $|q|_l=1$.
\end{proof}

\begin{proposition}[counting admissible representations; formally verified, \cite{CharacterFamily}]\label{i5:x-prop:u3count}
Let $K$ be imaginary quadratic.  There are $c_1,A_1,c_2,A_2\in\N$ such that
for all $e,N\ge1$, every prime $p\nmid Nd_K$ and every place $w\mid p$ of $K$,
with $q_w$ the norm of $w$:
(i) there is a trace datum $D_1$ for $(q_w,e)$ with $\#D_1\le
c_1(N(e+1))^{A_1}$ which, for every odd prime $l\ne p$, every $\iota$ and every
Frobenius $\sigma\in G_K$ at $w$, classifies at $\sigma^{-1}$ the stably
admissible elements of $\mathrm{Adm}(e,N)$;
(ii) there is a trace datum $D_2$ for $(q_w,e)$ with $\#D_2\le
c_2(N|d_K|(e+1))^{A_2}$ classifying in the same sense the endoscopically
admissible elements (level $N|d_K|$).
\end{proposition}
\noindent\emph{The counting statements.}  Parts (i) and (ii) are derived below from
two formally verified counting statements.  An \emph{eigenvalue datum}
$(S,\beta)$ for $(n,q,m,\mathrm{w})$ ($m\in\Z^n$, $\mathrm{w}\in\Z$) is a finite
set $S$ with multisets $\beta(s)\subset\C$ of size $n$ such that every
$\tau\in\mathrm{Aut}(\C)$ permutes the family $(\beta(s))_s$ (applied
elementwise), $|b|=q^{\mathrm{w}/2}$ for $b\in\beta(s)$, and $q^ab$ (resp.\
$q^ab^{-1}$) is an algebraic integer whenever $-a\le m_i$ for all $i$ (resp.\
$m_i\le a$ for all $i$).  It \emph{classifies at $w$} a set $\mathcal A$ of
representations $G_F\to\mathrm{GL}_n(\overline\Q_l)$ if there is $\mathrm{cl}\colon
\mathcal A\to S$ with $\mathrm{cl}(\rho)=\mathrm{cl}(\rho')\Rightarrow\rho\cong
\rho'$ and such that for every $\rho\in\mathcal A$ and every valuation ring $A\mid
w$, $\rho$ is trivial on $I_A$ and $\iota\operatorname{charpoly}\rho(\sigma^{-1})
=\prod_{b\in\beta(\mathrm{cl}\rho)}(X-b^{-1})$ for every Frobenius $\sigma$ at
$A$.  For a CM field $F$, $M\ge1$ and $m\in\Z^n$, $\mathrm{Adm}^K_{M,m}$ is the
set of $\rho$ that are irreducible, satisfy $\mathrm{Aut}_\iota(\rho)$, have level
$(M;m)$ (Definition~\ref{i5:x-def:bab-adm}(a)), and are \emph{odd}: at every
$w\mid l$, $\rho$ is unlinked or potentially unlinked
(Definition~\ref{i5:x-def:witA}(1),(2)).  For $F=\Q$, $\mathrm{Adm}^\Q_{M,m,
\mathrm{w}}$ is the set of $\rho$ that are irreducible, satisfy
$\mathrm{Aut}^{(4)}_\iota$ (Definition~\ref{i5:x5-def-gsp4}, for general $n$),
have level $(M;m)$, satisfy ${}^t\rho(s)P\rho(s)=\theta(s)P$ for some
$P\in\mathrm{GL}_n$ and some $\theta=\bar\chi^{\mathrm{w}}\theta_0$ with
$\theta_0$ of finite level, and either ($l\mid M$ and $\rho$ unlinked at every
$w\mid l$) or $m_i\ne m_j+1$ for all $i,j$.

\begin{itemize}
\item[$(\mathrm{C}_K)$] For every CM field $K$ and $n\ge2$ there are $c,A$ such
that for every strictly decreasing $m\in\Z^n$ with $m_i+m_{n-1-i}=1-n$, every
$H\ge\max|m_i|$, $M\ge1$, prime $p\nmid M$, $p\nmid d_K$ and place $w\mid p$,
there is an eigenvalue datum for $(n,q_w,m,1-n)$ of size $\le c(M(H+1))^A$
classifying $\mathrm{Adm}^K_{M,m}$ at $w$ for every odd prime $l\ne p$ and every
$\iota$.
\item[$(\mathrm{C}_\Q)$] For every $n\ge2$ there are $c,A$ such that for every
strictly decreasing $m\in\Z^n$ and $\mathrm{w}$ with $m_i+m_{n-1-i}=\mathrm{w}$,
$H\ge\max|m_i|$, $M\ge1$ and prime $p\nmid M$, there is an eigenvalue datum for
$(n,p,m,\mathrm{w})$ of size $\le c(M(H+1))^A$ classifying
$\mathrm{Adm}^\Q_{M,m,\mathrm{w}}$ at $p$ for every odd $l\ne p$ and $\iota$.
\item[$(\mathrm{N})$] If $e\ge1$ and $W\colon G_\Q\to\mathrm{GL}_2(\overline\Q_l)$ is
irreducible, $W\otimes\bar\chi^{-e}$ is attached to a newform of weight $2e+1$,
$W$ satisfies the witness over every CM quadratic field, has level
$(M;e,-e)$ and $\det W$ is $\bar\chi^0$ times a character of finite level, then
$W\in\mathrm{Adm}^\Q_{M,(e,-e),0}$.
\end{itemize}
$(\mathrm{C}_K)$ and $(\mathrm{C}_\Q)$ are formally verified in
\cite{CharacterFamily} through the bundle of published statements described below
together with two level-compatibility statements (Satake parameters of the
attached $\pi$ at places $w\nmid l$ and membership of $\pi$ in the finite set of
cuspidal representations of level $M$ and weight $m$) and an admissibility
restriction statement from $\Q$ to $K$, all three themselves derived formally
from the same bundle; $(\mathrm{N})$ is formally verified from the bundle.
From $(\mathrm{C}_K)$ with $n=3$ and $m=(e-1,-1,-e-1)$, \cite{CharacterFamily}
derives (formally) the following statement: \emph{for every CM field $K$ there
are $c_1,A_1$ such that for $e,N\ge1$, $p\nmid Nd_K$ and $w\mid p$ there is a
trace datum $D_1$ for $(q_w,e)$, $\#D_1\le c_1(N(e+1))^{A_1}$, classifying at
$\sigma^{-1}$, for every odd $l\ne p$, every $\iota$ and every Frobenius $\sigma$
at $w$, the set $\mathcal A_3(e,N)$ of $\rho\colon G_K\to\mathrm{GL}_3(\overline\Q_l)$ that
are irreducible, of level $(N;e,0,-e)$, odd, and with
$\mathrm{Aut}_\iota(\rho\otimes\bar\chi^{-1})$} (the twist moves the weights by
$1$ and $\mathrm{tr}$ of the datum is the sum of the inverses of the
eigenvalues).  From $(\mathrm{C}_\Q)$ with $n=2$, $m=(e,-e)$, $\mathrm{w}=0$ and
$(\mathrm N)$ it derives: \emph{there are absolute $c_2,A_2$ such that for every
number field $K$, $e,M\ge1$, prime $p\nmid M$, $h\in\N$ and place $w\mid p$ of
$K$, there is a trace datum $D_2$ for $(p^h,e)$, $\#D_2\le c_2(M(e+1))^{A_2}$,
classifying at $\sigma^{-1}$, for every odd $l\ne p$, $\iota$ and every
$\sigma\in G_K$ which is a Frobenius at a place above $w$ for the norm $p^h$,
the endoscopically admissible elements for $(e,M)$} (Definition~\ref{i5:x-def:bab-adm}(d), with $M$ in place of $N|d_K|$).

\begin{proof}[Proof of Proposition~\ref{i5:x-prop:u3count} from the two derived statements]
(i) Take $c_1,A_1$ and $D_1$ from the first statement.  It suffices to show that
every stably admissible $\rho$ lies in $\mathcal A_3(e,N)$, since a datum
classifying a set classifies each subset (restrict $\mathrm{cl}$).
Irreducibility, the level and $\mathrm{Aut}_\iota(\rho\otimes\bar\chi^{-1})$ are
part of Definition~\ref{i5:x-def:bab-adm}(c).  Oddness: let $w\mid l$.  By (b)
of that definition there are a place $A\mid w$ and a flag
$(\kappa_1,\kappa_2,\kappa_3)$ of $\rho|_{D_A}$ with $\kappa_i\ne\kappa_j
\chi_{\rm cyc}$ on $D_A$ for all $i,j$.  We show that $\rho$ is unlinked at $w$
in the sense of Definition~\ref{i5:x-def:witA}(1), i.e.\ that such a flag exists
at \emph{every} $A'\mid w$.  Since $\Omega/K$ is normal, $A'=g(A)$ for some
$g\in G_K$, so $D_{A'}=gD_Ag^{-1}$.  If $h\in\mathrm{GL}_3$ makes $h\rho(s)h^{-1}$
upper triangular with diagonal $(\kappa_i(s))$ for $s\in D_A$, then
$h\rho(g)^{-1}$ does the same for $\rho|_{D_{A'}}$ with the characters
$\kappa'_i(s')=\kappa_i(g^{-1}s'g)$, because $h\rho(g)^{-1}\rho(s')\rho(g)h^{-1}
=h\rho(g^{-1}s'g)h^{-1}$.  As $\chi_{\rm cyc}$ takes values in an abelian group,
$\chi_{\rm cyc}(g^{-1}s'g)=\chi_{\rm cyc}(s')$, so $\kappa'_i\ne\kappa'_j
\chi_{\rm cyc}$ on $D_{A'}$.  Hence $\rho$ is unlinked at every $w\mid l$, in
particular odd, and $\rho\in\mathcal A_3(e,N)$.  The bound $\#D_1\le
c_1(N(e+1))^{A_1}$ is the stated one.  (The inclusion uses none of the
published hypotheses; the formal inclusion contains an argument of the bundle
which is not used.)

(ii) Let $K$ be imaginary quadratic and $M=N|d_K|$.  Take the absolute
constants $c_2,A_2$ of the second statement.  Let $p\nmid Nd_K$ and $w\mid p$;
then $p\nmid M$, and $q_w=p^h$ with $h=f(w/p)\in\{1,2\}$.  The second statement
with this $K$, $M$, $h$ gives $D_2$, a trace datum for $(p^h,e)=(q_w,e)$, with
$\#D_2\le c_2(N|d_K|(e+1))^{A_2}$, classifying at $\sigma^{-1}$ the
endoscopically admissible elements for $(e,N|d_K|)$, for every odd $l\ne p$,
every $\iota$ and every Frobenius $\sigma\in G_K$ at $w$ (a Frobenius at $w$ for
$K$ has norm $q_w=p^h$).  This is (ii), with constants independent of $K$.
\end{proof}

\noindent\emph{The published inputs.}  The two counting statements
$(\mathrm{C}_K)$, $(\mathrm{C}_\Q)$, the statement $(\mathrm N)$ and the
level and restriction statements are derived formally from the following
published statements, which the formal development isolates as separate
propositions and proves inside the development.  We list them by
family, with their number.
The bundle consists of $65$ statements: (H1) nine on automorphic
representations, (H2) eighteen of the local representation theory of
$\mathrm{GL}_n$ over $p$-adic fields and of finite groups of Lie type, (H3)
fifteen global statements on automorphic forms, Galois representations and
class field theory, (H4) twelve of Galois cohomology, (H5) eleven on CM abelian
varieties and special values.  Each is listed, with its usual attribution, in
Table~\ref{i5:x6-tab-bundle} of \S\ref{i5:x6-sssec-count}; the formal proofs of
the cores of \S\ref{i5:x5-ssec} use, in addition, the five statements (H6) of
the same table.  We do not claim that every one of the $65$ is used on the route to
(i) and (ii).  The formally verified steps of
$(\mathrm C_K)$, $(\mathrm C_\Q)$ and $(\mathrm N)$ are outlined in
\S\ref{i5:x6-sssec-count}.

\begin{proposition}[$\mathrm{NumberU3U}$]\label{i5:x-prop:numberu3u}
$\mathrm{NumberU3U}$ holds, given Proposition~\ref{i5:x-prop:u3count}.
\end{proposition}

\begin{proof}
Let $c_1,A_1,c_2,A_2$ be as in Proposition~\ref{i5:x-prop:u3count},
$d=|d_K|$, $c=5(c_1+c_2d^{A_2})$ and $A=A_1+A_2+1$.  Let $e,N\ge1$, $p\ge7$
prime to $Nd$, $w\mid p$, and $\mathrm{Orb},z$ as in the statement.  Then
$q_w=p^h$ with $h\le2$.  The disjoint union $D=D_1\sqcup D_2$ is a trace datum
classifying all of $\mathrm{Adm}(e,N)$; with $Y=N(e+1)$,
$\#D\le c_1Y^{A_1}+c_2d^{A_2}Y^{A_2}\le(c_1+c_2d^{A_2})Y^{A_1+A_2}$.  As
$q_w\ge7$, Proposition~\ref{i5:x-prop:normbound} gives $NN\ne0$ with
$\log|NN|\le\#\mathrm{Orb}\,\#D\,(2e\log q_w+1)$, and since $\log p\ge1$,
$\log q_w\le2\log p$ and $e+1\le Y$, $2e\log q_w+1\le5Y\log p$; so
$\log|NN|\le c\,\#\mathrm{Orb}\,Y^A\log p$.  For odd $l\ne p$ we have $l\nmid
q_w$, and the second part of Proposition~\ref{i5:x-prop:normbound} with
$g=\sigma^{-1}$ gives the two remaining assertions.
\end{proof}

\begin{proof}[Proof of Theorem~\ref{i5:prop-corg}]
By Propositions~\ref{i5:x-prop:onechar} and \ref{i5:x-prop:numberu3u} and
Lemma~\ref{i5:x-lem:propGK}, $\mathrm{PropGK}$ holds; by
Lemma~\ref{i5:x-lem:corGK}, $\mathrm{CorGK}$ holds.
\end{proof}


\subsection{Deformation data for the automorphic cores}\label{i5:x5-ssec-defs}

This subsection contains the definitions needed to state the five formally
verified cores of \S\ref{i5:x5-ssec} exactly (Propositions~\ref{i5:x5-prop-U1},
\ref{i5:x5-prop-U234}, \ref{i5:x5-prop-sp4core}, \ref{i5:x5-prop-Vcore},
\ref{i5:x5-prop-Vhost}).  They follow the formal definitions of
\cite{CharacterFamily}; nothing in them is new mathematics, and some of them
are phrased in a way suited to formalisation rather than to readability (for
instance, local conditions are imposed through the existence of a conjugating
matrix rather than through a sub-functor).  We state them in this form so that
the statements below are exactly the formal ones.

\medskip\noindent\emph{Conventions.} $l$ is a prime and $\overline\Q_l$ an
algebraic closure of $\Q_l$ (not completed); the formal statements use the
same field $\overline\Q_l$ for the target of the embeddings $j$, for
coefficients and for $\iota$.
$|\cdot|_l$ is the absolute value of $\overline\Q_l$ with $|l|_l=l^{-1}$ and
$v_l(x)=-\log|x|_l/\log l$.  A \emph{coefficient ring} is a subring
$\mathcal O\subset\overline\Q_l$ containing $\Z_l$, finite over $\Z_l$ and a discrete
valuation ring; $k_{\mathcal O}$ is its residue field and $e_{\mathcal O}$ its
ramification index.  $\Omega$ is an algebraic closure of $\Q$, $G_F=\mathrm{Gal}
(\Omega/F)$ with the Krull topology, and for a finite set $S$ of primes
$G_{\Q,S}=G_\Q/N_S$ as in Definition~\ref{i5:def-ldata}, with projection
$\pi_S$.  For a valuation ring $A$ of $\Omega$, $D_A$ and $I_A$ are its
decomposition and inertia groups.  $\chi\colon G_\Q\to\Z_l^\times$ is the
cyclotomic character, $\chi_{\mathcal O}$ the same character with values in
$\mathcal O^\times$ and $\bar\chi$ the same character with values in
$\overline\Q_l^\times$.  For a Dirichlet character $\psi$ with values in $\mathcal O$,
$\Psi_{e,\psi}=\chi_{\mathcal O}^{\,e}\psi\colon G_\Q\to\mathcal O^\times$.  An
$\mathcal O$-algebra $R$ is \emph{complete local over $\mathcal O$} if $R$ is
local and noetherian, $\mathfrak m_R$-adically complete, $\mathcal O\to R$ is a
local homomorphism and $R=\mathcal O+\mathfrak m_R$.

\begin{definition}[the group $\mathcal G_m$]\label{i5:x5-def-Gm}
Let $R$ be a commutative ring and $m\ge1$; index rows and columns by
$0,\dots,m-1$.  Let $J=J_m$ have entries $J_{ij}=1$ if $j=m-1-i$ and $0$
otherwise, and $\theta(g)=J\,{}^t\!g^{-1}J$ for $g\in\mathrm{GL}_m(R)$.
$\mathcal G_m(R)=\mathrm{GL}_m(R)\rtimes\{\pm1\}$, with $-1$ acting by $\theta$.
For $x=(A,\varepsilon)$ write $\operatorname{mat}x=A$ and $\nu(x)=\varepsilon$
(the multiplier).  $\mathcal G_m(R)$ acts on $M_m(R)$ by
$\mathrm{Ad}(A,1)X=AXA^{-1}$ and $\mathrm{Ad}(A,-1)X=A\,(-J\,{}^tXJ)\,A^{-1}$.
For a group $\Gamma$ and a homomorphism $r\colon\Gamma\to\mathcal G_m(R)$: a
\emph{cocycle} is $\gamma\colon\Gamma\to M_m(R)$ with
$\gamma(gh)=\gamma(g)+\mathrm{Ad}\,r(g)\,\gamma(h)$; the \emph{coboundary} of
$X\in M_m(R)$ is $g\mapsto X-\mathrm{Ad}\,r(g)X$.  For $h\in\mathrm{GL}_m(R)$,
$\operatorname{conj}_h r$ is $r$ followed by conjugation by $(h,1)$ in
$\mathcal G_m(R)$.  For a ring map $f$, $f_*r$ is $r$ followed by $f$ on
entries.  $r$ \emph{has multiplier} $\eta\colon\Gamma\to\{\pm1\}$ if
$\nu\circ r=\eta$.

For $m=3$ put $w(a)=1-a$ ($a=0,1,2$), $\operatorname{wt}(a,b)=w(a)-w(b)$,
$\operatorname{tor}(t)=\mathrm{diag}(t,1,t^{-1})$ (so
$\theta(\operatorname{tor}t)=\operatorname{tor}t$), and for $\eta\colon\Gamma\to
\{\pm1\}$, $\Psi\colon\Gamma\to R^\times$ let
$r_0^{\eta,\Psi}(g)=(\operatorname{tor}\Psi(g),\eta(g))$.  With $A(g)=
\operatorname{mat}r(g)$: $r$ is \emph{upper} (resp.\ \emph{lower}) if
$A(g)_{ab}=0$ for all $g\in\ker\eta$ and $\operatorname{wt}(a,b)<0$ (resp.\
$>0$); \emph{block} if $A(g)_{ab}=0$ for $g\in\ker\eta$ and
$\operatorname{wt}(a,b)$ odd.  $\operatorname{corner}(A)$ is the $2\times2$
submatrix on the indices $\{0,2\}$.  For $H\le\Gamma$, a map $a\colon\Gamma\to
M_m(R)$ is \emph{irreducible on $H$} if the only $R$-submodules
$W\subseteq R^m$ with $a(g)W\subseteq W$ for all $g\in H$ are $0$ and $R^m$;
$a,a'$ are \emph{$H$-isomorphic} if $Ma(g)=a'(g)M$ for some
$M\in\mathrm{GL}_m(R)$ and all $g\in H$.  $r$ is of \emph{type~I} if
$\operatorname{mat}\circ r$ is irreducible on $\ker\eta$, of \emph{type~II} if
$r$ is block and $\operatorname{corner}\circ\operatorname{mat}\circ r$ is
irreducible on $\ker\eta$.  For $c,g_1\in\Gamma$ put $g_1^*=cg_1^{-1}c$ and
$Z=A(g_1)-A(g_1^*)$; $r$ is \emph{normalised at $(c,g_1)$} if $A(c)=1$ and
$Z_{01}=Z_{10}=0$.  $(\eta,\Psi,c,g_1)$ is a \emph{marking} if $c^2=1$,
$\eta(c)=-1$, $\Psi(c)=1$, $\eta(g_1)=1$ and $\Psi(g_1)^2-1\in R^\times$.
\end{definition}

\begin{definition}[frames and their points]\label{i5:x5-def-frame}
Let $\mathcal O$ be a local ring.  A \emph{frame} $F$ over $\mathcal O$
consists of topological groups $G,\Gamma$, a homomorphism $\pi\colon G\to
\Gamma$, characters $\eta\colon\Gamma\to\{\pm1\}$ and
$\Psi\colon\Gamma\to\mathcal O^\times$, an element $c\in\Gamma$, an index set
$\iota$ with subgroups $D_v,I_v\le G$ ($v\in\iota$), an index $v_l\in\iota$ and
a subset $S_b\subseteq\iota$.  Put $T=\{v_l\}\cup S_b$, $\eta_v=\eta\circ
\pi|_{D_v}$, $\Psi_v=\Psi\circ\pi|_{D_v}$, $I'_v=I_v\cap D_v$, and
$N_v=\{s\in I'_v:\eta_v(s)=1,\ \Psi_v(s)\equiv1\bmod\mathfrak m_{\mathcal O}\}$.
Let $R$ be a local $\mathcal O$-algebra and $r\colon D_v\to\mathcal G_3(R)$.
\begin{enumerate}
\item[(T)] $r$ is \emph{triangular at $v$} if there is $h\in\mathrm{GL}_3(R)$,
$h\equiv1\bmod\mathfrak m_R$, such that $r'=\operatorname{conj}_{h^{-1}}r$ is
upper (for $\eta_v$) and $\operatorname{mat}r'(s)_{aa}=\Psi_v(s)^{w(a)}$ for all
$s\in I'_v$ with $\eta_v(s)=1$ and all $a$.
\item[(D)] $r$ is \emph{diagonal at $v$} if moreover $r'$ can be taken lower as
well.
\item[(U)] $r$ is \emph{unramified at $v$} if $r(s)=1$ for $s\in N_v$.
\end{enumerate}
A homomorphism $r\colon\Gamma\to\mathcal G_3(R)$ is a \emph{point of $F$ at
$g_1\in\Gamma$} if: $r$ has multiplier $\eta$; $\operatorname{mat}r(g)\equiv
\operatorname{tor}\Psi(g)\bmod\mathfrak m_R$ entrywise; $r\circ\pi|_{D_v}$ is
triangular at $v_l$, diagonal at every $v\in S_b$, unramified at every
$v\notin T$; $r$ is normalised at $(c,g_1)$; and every entry of
$\operatorname{mat}r$ is locally constant modulo $\mathfrak m_R^n$ for every
$n$.  A point $u$ with values in $R_u$ is \emph{universal at $g_1$} if for
every $R'$ complete local over $\mathcal O$ and every point $r$ of $F$ at $g_1$
with values in $R'$ there is a unique $\mathcal O$-algebra map $f\colon R_u\to
R'$ with $f_*u=r$.  $g_1$ is \emph{marked} if $\eta(g_1)=1$ and
$\Psi(g_1)^2-1\in\mathcal O^\times$.

\emph{Weights.}  For a point $u$ with values in $R$ and $j\in\Z$, $R_j(u)$ is
the set of $x\in R$ such that $f'(x)=t^jf(x)$ whenever $R'$ is complete local
over $\mathcal O$, $f,f'\colon R\to R'$ are $\mathcal O$-algebra maps, $t\in
R'^\times$ and $f'_*u=\operatorname{conj}_{\operatorname{tor}t}(f_*u)$.  Each
$R_j(u)$ is an $\mathcal O$-module, $R_iR_j\subseteq R_{i+j}$, the entry
$\operatorname{mat}u(g)_{ab}$ lies in $R_{\operatorname{wt}(a,b)}(u)$, and
$S(u)=\sum_jR_j(u)$ is an $\mathcal O$-subalgebra.  $R_0(u)$ is the
\emph{weight-zero part}.  A \emph{point of $u$} with values in an
$\mathcal O$-algebra $K$ is a pair $(\xi,r)$ of an $\mathcal O$-algebra map
$\xi\colon S(u)\to K$ and $r\colon\Gamma\to\mathcal G_3(K)$ with
$\nu(r(g))=\nu(u(g))$ and $\operatorname{mat}r(g)_{ab}=\xi(\operatorname{mat}
u(g)_{ab})$.

\emph{Characteristic-zero points.}  For a field $K\supseteq\mathcal O$,
$r\colon\Gamma\to\mathcal G_3(K)$ is a \emph{characteristic-zero point of $F$
at $g_1$} if there are a subring $A\subseteq K$ containing $\mathcal O$ which is
a discrete valuation ring finite over $\mathcal O$, $t\in K^\times$ and a point
$r'$ of $F$ at $g_1$ with values in $A$ such that $r'=\operatorname{conj}_
{\operatorname{tor}t}r$ in $\mathcal G_3(K)$.

\emph{Rigidity.}  For a topological field $K\supseteq\mathcal O$ and
$r\colon\Gamma\to\mathcal G_3(K)$, $r$ is \emph{rigid for $F$} if every
continuous cocycle $y\colon\Gamma\to M_3(K)$ for $r$ is a coboundary as soon as:
$y(\pi(s))=0$ for $s\in N_v$, $v\notin T$; and the deformation
$r_\varepsilon=((1+\varepsilon y)\operatorname{mat}r,\nu\circ r)$ of $r$ to
$K[\varepsilon]$ satisfies condition (T) at $v_l$ and condition (D) at every
$v\in S_b$, with $h\in\mathrm{GL}_3(K[\varepsilon])$ arbitrary (no congruence
condition on $h$).
\end{definition}

\begin{definition}[host data and host frames]\label{i5:x5-def-hostframe}
Let $\mathcal O$ be a coefficient ring, $N\ge1$, $\psi\colon(\Z/N)^\times\to
\mathcal O^\times$, $e,o_1\in\N$; $\theta_0,\kappa,c_2$ are as in
Definition~\ref{i5:def-cases}, $S_b(l,\psi)$ is the set of primes $p\mid N$,
$p\ne l$, with $\mathrm{comp}_p\kappa\ne1$, and $T(l,\psi)=\{l\}\cup S_b(l,\psi)$.
$\mathrm{HostData}(l,e,\psi,o_1)$ means: $l\ne2$; $e\ge1$; $N$ even;
$\psi(-1)=(-1)^e$; $o_1\ge1$; $(\mathrm{comp}_r\theta_0)^{2o_1}=1$ for every
odd prime $r$; $8o_1v^2<2^{c_2(\psi)}$ for every $v\in T(l,\psi)$; and
$c_2(\psi)\ge4$.  Let $\eta_4\colon G_\Q\to\{\pm1\}$ be the character of
$\Q(i)/\Q$.  A \emph{host frame} for $(e,\psi)$ is a finite set $S$ of primes,
valuation rings $A_p$ of $\Omega$ ($p\in S$) and a frame $F$ over $\mathcal O$
with $G=G_\Q$, $\Gamma=G_{\Q,S}$, $\iota=S$, such that: $S$ is the set of
primes dividing $lN$; $A_p$ lies over $p$; $\pi=\pi_S$; $\eta\circ\pi=\eta_4$;
$\Psi\circ\pi=\Psi_{e,\psi}$; $c=\pi(\sigma)$ for some complex conjugation
$\sigma$; $D_p=D_{A_p}$ and $I_p=I_{A_p}$; $v_l=l$; and $S_b=S_b(l,\psi)$.
\end{definition}

\begin{definition}[the conditions $\mathrm{B}'$]\label{i5:x5-def-Bprime}
Let $K$ be a CM field and $\iota\colon\overline\Q_l\cong\C$.  For $n\ge1$ and a
homomorphism $\rho\colon G_K\to\mathrm{GL}_n(\overline\Q_l)$, $\mathrm{Aut}_\iota(\rho)$ is
the automorphy predicate of Definition~\ref{i5:x-def:witA}: there is a regular
algebraic cuspidal $\pi$ of $\mathrm{GL}_n(\mathbb A_K)$ with $\pi^\vee\cong\pi^c$
to which $\rho$ is attached via $\iota$, and, if $\rho$ is continuous,
$\mathrm{Wit}(\rho)$ holds.  For
$a\colon G_\Q\to M_3(\overline\Q_l)$ and $j\in\Z$, $\mathrm{AutOn}_{K,\iota,j}(a)$ means:
for some embedding of $K$ in $\Omega$ there is a homomorphism
$\rho\colon G_K\to\mathrm{GL}_3(\overline\Q_l)$ with $\rho(s)=a(s)$ for $s\in G_K$ and
$\mathrm{Aut}_\iota(\rho\otimes\bar\chi^{\,j})$.  For $r\colon G_{\Q,S}\to
\mathcal G_3(\overline\Q_l)$ put $a_r=\operatorname{mat}\circ r\circ\pi_S$;
$\mathrm{B}'_{K,\iota}(r)$ means: (i) $a_r$ is irreducible on every open
subgroup of $G_\Q$ contained in $\ker\eta_4$; (ii) $a_r$ is irreducible on
$\ker\eta_4\cap\ker\chi$; (iii) $\mathrm{AutOn}_{K,\iota,-1}(a_r)$.
\end{definition}

\begin{definition}[OneCharU data]\label{i5:x5-def-dat}
Let $K$ be an imaginary quadratic field contained in $\Omega$, $\mathcal O$ a
coefficient ring, $N\ge1$, $e\ge1$ and $\psi\colon(\Z/N)^\times\to\mathcal
O^\times$.  A \emph{datum} $D$ for $(K,\mathcal O,e,\psi)$ consists of:
$\mathrm{GD}(K,l,e,\psi)$ (Definition~\ref{i5:def-gd}); $\eta\colon G_\Q\to\{\pm1\}$
with $\ker\eta=G_K$; a complex conjugation $c_D$; an embedding $j\colon\Omega
\to\overline\Q_l$ and the corresponding continuous embedding
$\iota_D\colon G_{\Q_l}\to G_\Q$ (with $\sigma\circ j=j\circ\iota_D(\sigma)$); a
finite extension $K_v/\Q_l$ with $\eta(\iota_D g)=1\iff g\in G_{K_v}$; such that
$\eta\circ\iota_D$ is trivial on the inertia group $I_l\le G_{\Q_l}$,
$\tau=(\mathrm{comp}_l\psi)\circ\iota_D$ has open kernel, $\Psi=\Psi_{e,\psi}$ is
continuous, and $\Psi(\iota_Dg)=\chi(g)^e\tau(g)$ for $g\in I_l$.  Put
$\bar r_0=r_0^{\eta,\bar\Psi}$, $\bar\Psi=\Psi\bmod\mathfrak m_{\mathcal O}$.  A
\emph{lift} (of $\bar r_0$, or of $\bar r_0\circ\iota_D$) with values in $R$
complete local over $\mathcal O$ is a homomorphism $r$ to $\mathcal G_3(R)$
reducing to $\bar r_0$ (multiplier included) whose entries are locally constant
modulo each $\mathfrak m_R^n$.

\emph{The ordinary local condition.}  Let $(R^{\rm loc},u^{\rm loc})$ be a
universal lift of $\bar r_0\circ\iota_D$ (no condition).  Put $H_L=\ker(\eta
\circ\iota_D)\cap\ker\tau$, $\psi_a(s)=\Psi(\iota_Ds)^{w(a)}$ (indices $a$ mod
$3$), and choose (once and for all, by choice) $s_0\in H_L\cap I_l$ with
$\chi(s_0)$ of infinite order and $\phi_0\in H_L$ such that $I_l\cap H_L$ and
$\phi_0$ topologically generate $H_L$ and $\phi_0\in\mathrm{Frob}^mI_l$ for some
$m\ge1$.  For a lift $r$ put $\mathrm{sh}_a(s)=\operatorname{mat}r(s)-
\psi_a(s)$, $\mathrm{cc}(a)=(\psi_a(s_0)-\psi_{a+1}(s_0))(\psi_a(s_0)-
\psi_{a+2}(s_0))$, $\mathrm{tr}(a)=\operatorname{tr}(\mathrm{sh}_{a+1}(s_0)
\mathrm{sh}_{a+2}(s_0)\operatorname{mat}r(\phi_0))$,
$\mathrm{dd}(a)=\mathrm{cc}(a+1)\mathrm{tr}(a)-\chi(\phi_0)\mathrm{cc}(a)
\mathrm{tr}(a+1)$ and $\delta=l\cdot\mathrm{dd}(0)\mathrm{dd}(1)$.  Let
$I^{\rm ord}\subseteq R^{\rm loc}$ be the ideal generated (for $r=u^{\rm loc}$)
by the coefficients of $\operatorname{charpoly}(\operatorname{mat}r(s))-
\prod_a(X-\psi_a(s))$, $s\in H_L\cap I_l$, and by the entries of
$\mathrm{sh}_0(s_0')\mathrm{sh}_1(s_1')\mathrm{sh}_2(s_2')$,
$s_i'\in H_L\cap I_l$; let $I_Y=\{x:\delta^nx\in I^{\rm ord}\text{ for some }n\}$,
and let $I_t$ be the intersection of those minimal primes of $\sqrt{I_Y}$ that
are contained in the kernel of every $\mathcal O$-algebra map
$x\colon R^{\rm loc}\to\mathcal O$ with $x_*u^{\rm loc}=r_0^{\eta,\Psi}\circ
\iota_D$.

\emph{The global ring.}  Let $\mathrm{Kill}_D$ be the set of
$g\in I_A$, $A$ a valuation ring of $\Omega$ not above $l$, with $\bar r_0(g)=1$.
For $g_1\in G_\Q$, a lift $r$ of $\bar r_0$ is \emph{$D$-admissible at $g_1$} if it is normalised
at $(c_D,g_1)$, $r(g)=1$ for $g\in\mathrm{Kill}_D$, and $r\circ\iota_D=f_*
u^{\rm loc}$ for some $\mathcal O$-algebra map $f$ with $I_t\subseteq\ker f$.  A
\emph{global ring} for $(D,g_1)$ is a universal pair $(R_D,u_D)$ for
$D$-admissible lifts (the representing property as in
Definition~\ref{i5:x5-def-frame}).  $g_1$ is \emph{marked} if $(\eta,\Psi,c_D,
g_1)$ is a marking.

\emph{Points.}  $E(u)$ is the $\mathcal O$-subalgebra of $R_D$ generated by all
entries of $\operatorname{mat}u_D$, and $W_0(u)$ the $\mathcal O$-subalgebra
generated by all products $\prod_i\operatorname{mat}u_D(g_i)_{a_ib_i}$ with
$\sum_i\operatorname{wt}(a_i,b_i)=0$.  A \emph{point} is a pair $(\xi,r)$ with
$\xi\colon E(u)\to\overline\Q_l$ an $\mathcal O$-algebra map and $r\colon G_\Q\to\mathcal
G_3(\overline\Q_l)$ with $\nu(r(g))=\nu(u_D(g))$ and $\operatorname{mat}r(g)_{ab}=
\xi(\operatorname{mat}u_D(g)_{ab})$.  An $\mathcal O$-algebra map $f\colon R_D
\to\overline\Q_l$ is \emph{local} if $|f(x)|_l<1$ for $x\in\mathfrak m_{R_D}$; it is
\emph{two-signed} if moreover $f_*u_D$ is neither upper nor lower.  For $r$ with
multiplier $\eta$, $\rho_K(r)=\operatorname{mat}\circ r|_{G_K}$.  For
$W\colon G_\Q\to\mathrm{GL}_2(\overline\Q_l)$, $W$ \emph{extends} $r$ if
$W(g)=\operatorname{corner}\operatorname{mat}r(g)$ for $g\in\ker\eta$ and
$\operatorname{tr}W(c_D)=0$; $\mathrm{Mod}_\iota(W)$ means: $W\otimes\bar
\chi^{-e}$ is attached via $\iota$ to a newform of weight $2e+1$, and $W$
satisfies the automorphy witness over every CM quadratic extension, as in
Definition~\ref{i5:x-def:bab-adm}(d).
\end{definition}

\begin{definition}[Bloch--Kato vanishing]\label{i5:x5-def-BK}
Let $F$ be a number field and $V$ a topological $\overline\Q_l$-vector space with a
$\overline\Q_l$-linear action $\rho$ of $G_F$.  A continuous cocycle $c\colon G_F\to V$ is
\emph{Bloch--Kato} if for every finite $w\nmid l$ there is a valuation ring
$A\mid w$ with $c|_{I_A}$ a coboundary, and for every embedding $\tau\colon F\to\overline\Q_l$ there is
an extension $j\colon\Omega\to\overline\Q_l$ of $\tau$ such that the
restriction of $c$ to the decomposition group of $j$ becomes a coboundary in
$V\otimes_{\Q_l}B_{\mathrm{cris}}$, i.e.\ there is $x\in V\otimes_{\Q_l}
B_{\mathrm{cris}}$ with $c(g)\otimes1=gx-x$ for all $g$ in that decomposition
group.  $\mathrm{BK}(\rho)$ means that every continuous Bloch--Kato cocycle is a
coboundary.  For $W\colon G_\Q\to\mathrm{GL}_2(\overline\Q_l)$, a CM quadratic $K$ and
$\theta\colon G_\Q\to\overline\Q_l^\times$, a representation $R$ of $G_\Q$ on
$M_2(\overline\Q_l)$ is an \emph{adjoint of $W|_{G_K}$ twisted by $\theta$} if there are
a complex conjugation $c$ and $P\in\mathrm{GL}_2(\overline\Q_l)$ with
${}^tW(s)PW(s')=\theta(s)P$ whenever $s,s'\in G_K$ and $s'=csc^{-1}$,
$R(s)X=W(s)XW(s)^{-1}$ for $s\in G_K$, and $R(c)X=-P^{-1}\,{}^tX\,P$.
\end{definition}

\begin{definition}[$\mathrm{GSp}_4$ data]\label{i5:x5-def-gsp4}
Index rows and columns of $4\times4$ matrices by $0,\dots,3$.  Put
$s=(1,1,-1,-1)$, $\lambda=(1,0,1,0)$, $\mu=(1,1,0,0)$,
$\operatorname{wt}(i,j)=(\lambda_i-\lambda_j,\mu_i-\mu_j)\in\Z^2$,
$\operatorname{adj}(X)_{ij}=s_is_jX_{3-j,3-i}$; $g$ is $\mathrm{GSp}$ with
multiplier $\nu$ if $\operatorname{adj}(g)g=\nu$, and $\mathfrak{sp}_4=\{X:
\operatorname{adj}X=-X\}$; $X$ is \emph{block} if $X_{ij}=0$ whenever
$\lambda_i\ne\lambda_j$.  A $\mathrm{GSp}_4$-\emph{datum} $D$ over a coefficient
ring $\mathcal O$ consists of a finite set $S$ of primes, $e_D\in\N$,
characters $\psi_D,\xi,\omega\colon G_\Q\to\mathcal O^\times$, $c_D\in G_\Q$,
a valuation ring $A_p$ of $\Omega$ for each prime $p$, elements
$\sigma_p,\phi_D,g_0\in G_\Q$ and $s_0\in D_{A_l}$.  Put
$\Psi_D=\chi_{\mathcal O}^{e_D}\psi_D$, $\alpha=\chi_{\mathcal O}\omega^{-1}\xi$,
$\beta=\Psi_D\alpha^{-1}$, $\vartheta=\alpha\Psi_D$,
$(\mathrm{ps}_0,\dots,\mathrm{ps}_3)=(\vartheta,\Psi_D,\alpha,1)$,
$r_0^D=\mathrm{diag}(\mathrm{ps}_i)$, $\bar r_0^D$ its reduction,
$S_b(D)=\{p\in S:p\ne l,\ l\mid\#\psi_D(I_{A_p})\}$ and $T(D)=\{l\}\cup S_b(D)$.
A valuation ring $A$ is \emph{away from} a set $U$ of primes if it lies over no
$p\in U$.  A function $f$ on $G_\Q$ is \emph{$T$-unramified with value $a$} if
$f(g)=a$ for every $g\in I_A$ with $A$ away from $T(D)$ and $\bar r_0^D(g)=1$.
For $\gamma\colon G_\Q\to\mathcal O^\times$ with reduction $\bar\gamma$,
$\mathrm{XVan}(\gamma)$ means: every locally constant $z\colon G_\Q\to
k_{\mathcal O}$ with $z(gh)=z(g)+\bar\gamma(g)z(h)$, $T$-unramified with value
$0$, and of the form $g\mapsto\bar\gamma(g)m_v-m_v$ on $D_{A_v}$ for each $v\in
T(D)$, is of the form $g\mapsto\bar\gamma(g)m-m$.

$D$ is \emph{good} if: $l\ne2$; $S$ consists of primes and contains $2$ and
$l$; $e_D\ge2$; $A_p$ lies over $p$; $c_D$ is a complex conjugation;
$\ker\psi_D$ and $\ker\xi$ are open; $I_A\le\ker\psi_D\cap\ker\xi$ for $A$ away
from $S$; $\Psi_D(c_D)=1$; $\omega^{l-1}=1$ and $\bar\omega=\bar\chi_{\mathcal O}$;
$\xi^m=1$ for some $m$ prime to $l$; $I_{A_v}\le\ker\xi$ for $v\in T(D)$;
$\xi(c_D)=-1$; $\sigma_v$ is a Frobenius at $A_v$ for $v\in T(D)$;
$\chi(\sigma_l)=1$; $\overline{\xi\Psi_D^i}(\sigma_v)\ne1$ for $v\in T(D)$ and
$i\in\Z$; $\mathrm{XVan}(\gamma)$ and $\mathrm{XVan}(\gamma^{-1})$ for
$\gamma\in\{\alpha,\beta,\vartheta\}$; $\phi_D=\lim_n\sigma_l^{\,l^{n!}}$;
$\bar\alpha(g_0)\ne1$ and $\bar\Psi_D(g_0)\ne1$; $s_0\in I_{A_l}\cap
\ker\psi_D\cap\ker\omega$ and $\chi(s_0)$ has infinite order.

\emph{Points.}  For a lift $r\colon G_\Q\to M_4(R)$ of $\bar r_0^D$ ($R$ complete
local over $\mathcal O$, continuity as above), the conditions are:
$\mathrm{Mult}$: $r(g)$ is $\mathrm{GSp}$ with multiplier $\vartheta(g)$ for all
$g$; $\mathrm{Unr}$: $r$ is $T$-unramified with value $1$; $\mathrm{Nrm}$:
$r(\phi_D)$ is block and $h=r(g_0)+\vartheta(g_0)r(g_0^{-1})$ has
$h_{02}=h_{20}=0$; at $l$, the condition $P_l$ on lifts of $\bar r_0^D|_{D_{A_l}}$:
$\mathrm{GSp}$ with multiplier $\vartheta$, and, with
$\mathrm{pr}_i=\prod_{j\ne i}(r(s_0)-\mathrm{ps}_j(s_0))$:
$\operatorname{charpoly}r(s_0)=\prod_j(X-\mathrm{ps}_j(s_0))$,
$\mathrm{pr}_j\,r(g)\,\mathrm{pr}_i=0$ for $i<j$, and
$\mathrm{pr}_i(r(s)-\mathrm{ps}_i(s))\mathrm{pr}_i=0$ for $s\in I_{A_l}$; at
$p\in S_b(D)$, the condition $P_p$: $\mathrm{GSp}$ with multiplier $\vartheta$ and
$r(g)=1$ for $g\in I_{A_p}\cap\ker\psi_D$.  A local lift $r_v$ \emph{lies on a
component through $r_0$} for $P_v$ if there are a complete local domain $B$
over $\mathcal O$, a lift $\rho$ with values in $B$ satisfying $P_v$, an
$\mathcal O$-algebra map $x\colon B\to\mathcal O$ with $x_*\rho=r_0^D|_{D_{A_v}}$
and $y\colon B\to R$ with $y_*\rho=r_v$.  $r$ is a \emph{point of $D$} if
$\mathrm{Mult}$, $\mathrm{Unr}$, $\mathrm{Nrm}$ hold and $r|_{D_{A_v}}$ lies on a
component through $r_0$ for $P_v$, $v\in T(D)$.  For an $\mathcal O$-algebra $K$,
$r\colon G_\Q\to M_4(K)$ is a \emph{$K$-point} if $r=y_*\rho$ for a point
$\rho$ with values in some $R$ and an $\mathcal O$-algebra map $y\colon R\to K$.
$r$ is \emph{irreducible} if its image spans $M_4$.

\emph{Properties.}  $\mathrm{Fin}_0(D)$: for every point $\rho$ with values in
$R$, the $\mathcal O$-subalgebra generated by the products
$\prod_i\rho(g_i)_{a_ib_i}$ with $\sum_i\operatorname{wt}(a_i,b_i)=(0,0)$ is
finite over $\mathcal O$.  $\mathrm{NoEnd}(D)$: for no $\overline\Q_l$-point $r$ with all
$r(g)$ block do the matrices $(r(g)_{ab})_{a,b\in\{0,2\}}$ span $M_2(\overline\Q_l)$.
$\mathrm{AllIrr}_D(P)$: every irreducible $\overline\Q_l$-point satisfies $P$.
$\mathrm{Rig}_D(r)$: every $\overline\Q_l[\varepsilon]$-point $r'$ with $r'\equiv r\bmod
\varepsilon$ has $r'(g)-r(g)=\varepsilon(Xr(g)-r(g)X)$ for some
$X\in\mathfrak{sp}_4(\overline\Q_l)$ and all $g$.  $a_g(r)=\operatorname{tr}r(g)-
\operatorname{tr}r_0^D(g)$.  $\operatorname{sha}(D)=\lgth_{\mathcal O}
\mathrm{Sel}_G(\rho_{e_D,\psi_D};\mathcal D(n_D))$, $n_D=\prod_{p\in S}p$
(notation of Definition~\ref{i5:def-V}), and $\mathrm{Vn}(D)=
\operatorname{sha}(D)/e_{\mathcal O}$.  $\mathrm{Ineq}_D(g)$:
$\operatorname{sha}(D)<\infty$ and there is a finite set $\Xi$ of irreducible
$\overline\Q_l$-points, pairwise not conjugate under $\mathrm{GL}_4(\overline\Q_l)$, with
$\mathrm{Vn}(D)\log l\le2\sum_{r\in\Xi}(-\log|a_g(r)|_l)$.  For $\iota\colon
\overline\Q_l\cong\C$, $\mathrm{Aut}^{(4)}_\iota(r)$: if $r$ is continuous then for every
CM quadratic $K$ with $r|_{G_K}$ irreducible, $r|_{G_K}$ is attached via $\iota$
to a regular algebraic cuspidal representation of $\mathrm{GL}_4(\mathbb A_K)$
and $\mathrm{Wit}(r|_{G_K})$ holds.

\emph{Data attached to a host.}  For a host $h$, $\psi'\in X_h(\mathcal O,e)$ and
$n',p_1,m_0\in\N$, $\mathrm{OfChar}(D,e,\psi',n',p_1,m_0)$ means: $e_D=e$,
$\psi_D=\psi'$, $S$ is the set of primes dividing $2lN_hp_1$, and
$\xi=\xi_0\xi'$ with $\xi'$ a primitive odd character of conductor $2^{n'}$ and
$\xi_0$ a character mod $p_1$ of order $m_0$.  $\mathrm{OfCase}(D,h,e,\psi')$
means: Case~II and $\mathrm{OfChar}(D,e,\psi',n_{\rm II},1,1)$, or Case~III and
$\mathrm{OfChar}(D,e,\psi',n_{\rm III},p_1(\chi),\operatorname{ord}\psi')$.
\end{definition}

The local condition at $l$ in Definitions~\ref{i5:x5-def-frame}
and~\ref{i5:x5-def-dat} is an ordinarity condition (a filtration with diagonal
characters $\Psi^{1},1,\Psi^{-1}$ on inertia, in the frame case after
conjugation by a matrix congruent to $1$; in the case of a datum, the Zariski
closure of the ordinary locus through the base points), and the condition at
the places $p\in S_b$ is diagonalisability with the same diagonal characters on
inertia.  \emph{No} condition of Steinberg type occurs anywhere: at the places
$v\notin T$ the lifts are unramified on $N_v$ (resp.\ trivial on
$\mathrm{Kill}_D$, resp.\ $T$-unramified).

\subsection{Reductions for SFew and the automorphic counts}\label{i5:x5-ssec}

This subsection proves the reduction steps behind Propositions
\ref{i5:x-prop:sfew}, \ref{i5:x-prop:onechar}, \ref{i5:prop-c23} and
\ref{i5:prop-c5}.  In each case what remains is stated as a smaller
proposition, formally verified in \cite{CharacterFamily}: for SFew the
construction of detecting cycles and a rank bound for the homology of
$\mathrm{GL}_n(\mathcal O_L)$; for the three counts an automorphy lifting
theorem in a residually reducible situation, a rigidity statement and a
length estimate for the deformation ring (stated exactly with the
definitions of \S\ref{i5:x5-ssec-defs}), plus, for Counts II and III, an exactly stated
arithmetic estimate.

\subsubsection{SFew}

Notation is that of \S\ref{i5:ssec-pf-shak}.  For a group $G$ and a ring $A$,
$C_m(G;A)$ is the free $A$-module on $G^m$ with the inhomogeneous bar
differential (trivial coefficients), $Z_m,B_m$ its cycles and boundaries, and
$\rho_q$ reduction modulo $q$.  We write
$\mathrm{st}\colon\mathrm{GL}_n\hookrightarrow\mathrm{GL}_{n'}$ ($n\le n'$) for
stabilisation and $\psi_d(g)=1_d\otimes g\in\mathrm{GL}_{dn}$ for the
block-diagonal $d$-fold sum.

\begin{definition}\label{i5:x5-def-prim}
Let $R$ and $A$ be commutative rings, $m\ge0$ and $B$ an $A$-module.  A
family of $A$-linear maps $\Phi_n\colon C_{m+1}(\mathrm{GL}_n(R);A)\to B$
is \emph{primitive} (over $A$) if (a) $\Phi_n(B_{m+1})=0$; (b)
$\Phi_{n'}(\mathrm{st}_*z)=\Phi_n(z)$ for $n\le n'$ and cycles $z$; (c)
$\Phi_{dn}(\psi_{d*}z)=d\,\Phi_n(z)$ for cycles $z$.  Unless stated
otherwise $A=\Z/q$ \textup{(}so $B$ is a $\Z/q$-module\textup{)}.  A family of
homomorphisms $\varphi_n\colon H_m(\mathrm{GL}_n(R);\Z)\to B$ is primitive if
$\varphi_{n'}\circ\mathrm{st}_*=\varphi_n$ and
$\varphi_{dn}\circ\psi_{d*}=d\varphi_n$.
For $l\ge1$ and $x\in H_m(\mathrm{GL}_n(R);\Z)$ put $z_0(x)=x$ and
$z_{t+1}(x)=\psi_{l*}z_t(x)-l^{t+2}\mathrm{st}_*z_t(x)$, so that
$z_t(x)\in H_m(\mathrm{GL}_{l^tn}(R);\Z)$.
\end{definition}

The remaining cores are the following three statements.  Fix a finite field
$k$ and $q\ge1$.  An \emph{additive family} of degree $m+1$ is a family of
$\Z/q$-linear maps $\theta_V\colon C_{m+1}(\mathrm{GL}(V);\Z/q)\to\Z/q$, $V$
running over finite-dimensional $k$-spaces, vanishing on boundaries, invariant
under transport by isomorphisms $V\cong W$ on cycles, and with
$\theta_{V\oplus W}((\rho\oplus\sigma)_*z)=\theta_V(\rho_*z)+\theta_W(\sigma_*z)$
for all homomorphisms $\rho\colon G\to\mathrm{GL}(V)$,
$\sigma\colon G\to\mathrm{GL}(W)$ and cycles $z\in Z_{m+1}(G;\Z/q)$.  For
$k\in G$ the prism operator is
$\mathrm{prism}_k[g_1|\dots|g_m]=\sum_{i=0}^m(-1)^i[g_1|\dots|g_i|k|k^{-1}g_{i+1}k|\dots|k^{-1}g_mk]$,
and $z_e=\sum_{j\in(\Z/q)^{e+1}}[j_0|1|j_1|1|\dots|j_e|1]\in C_{2e+2}(\Z/q;\Z/q)$
is the standard cycle.

\begin{definition}\label{i5:x5d-def-val}
Let $k$ be a finite field.
\textup{(1)} For a commutative ring $R$, a ring map $r\colon R\to k$ and $n\ge0$,
$\mathrm{red}_{r,n}\colon\mathrm{GL}_n(R)\to\mathrm{GL}(k^n)$ applies $r$ to the
entries; $\mathrm{scal}_V\colon k^\times\to\mathrm{GL}(V)$ is the scalar action.
\textup{(2)} Let $q\ge1$, $\bar\zeta\in k^\times$ with $\bar\zeta^q=1$ and $x\in\N$.
Then $\chi_x\colon\Z/q\to k^\times$ is the homomorphism $j\mapsto\bar\zeta^{x\tilde\jmath}$,
$\tilde\jmath\in\{0,\dots,q-1\}$ representing $j$. In
$z_e=\sum_{j\in(\Z/q)^{e+1}}[j_0|1|j_1|1|\dots|j_e|1]$ the entries in even
positions are the class of $1$ (a generator of $\Z/q$).
\textup{(3)} Let $v\colon k^\times\to\Z/q$ be a homomorphism and $e\ge0$. An additive
family $\theta$ of degree $2e+3$ with values in $\Z/q$ \emph{has value
$(\bar\zeta,v)$} if
$\theta_V\big(\mathrm{prism}_g((\mathrm{scal}_V\circ\chi_x)_*z_e)\big)=x^{e+1}v(\det g)$
for all finite $k$-spaces $V$, all $x\in\N$ and all $g\in\mathrm{GL}(V)$.
\textup{(4)} For number fields $L\subset F$, a ring map $r\colon\mathcal O_F\to k$,
$\alpha\in F^\times$ and $u\in k^\times$, $\alpha$ \emph{reduces to $u$ under $r$}
if $\alpha=a/c$ with $a,c\in\mathcal O_F$, $r(c)\ne0$ and $r(a)=u\,r(c)$.
\textup{(5)} For $\sigma\in\mathrm{Gal}(F/L)$ and a maximal ideal $P$ of $\mathcal O_F$,
$\sigma P=\sigma(P)$. For $a\in\N$, $\alpha\in F^\times$ is \emph{$a$-twisted} if
there is a divisor $\mathfrak d$ (a finitely supported function from the primes
of $F$ to $\Z$) with $\operatorname{ord}_P(\alpha)=\mathfrak d(\gamma P)-a\,\mathfrak d(P)$
for all $P$.
\end{definition}

\begin{proposition}[detection and construction; proved in \S\ref{i5:x5d-ssec}]\label{i5:x5-prop-detect-construct}
In \textup{(i-a)} and \textup{(i-b)} let $k$ be a finite field.
\textup{(i-a)} If $q\ge1$, $\zeta\in k^\times$
is a primitive $q$-th root of unity, $v\colon k^\times\to\Z/q$ a homomorphism and
$e\ge0$, there is an additive family of degree $2e+3$ with values in $\Z/q$ that
has value $(\zeta,v)$.
\textup{(i-b)} Let $A$ be a commutative ring,
$\theta$ an additive family of degree $m+1$ with values in $A$ \textup{(}the
definition of additive family with $\Z/q$ replaced by $A$\textup{)}, $R$ a
commutative ring and $r\colon R\to k$ a ring map. Then
$\Phi_n=\theta_{k^n}\circ(\mathrm{red}_{r,n})_*$ is a primitive family
\textup{(}Definition~\ref{i5:x5-def-prim}\textup{)} over $A$ with values in the
$A$-module $B=A$.
In particular this holds for $A=B=\Z/q$.
\textup{(ii)} Let
$\mathrm{Lay}(l,c,\nu,\zeta,\gamma)$ be a layer for $F/L$, $e\in\N$ with $2e+3<l^c$,
$a=(1+l^c)^{e+1}$, and $\alpha\in F^\times$ $a$-twisted. There are $N\in\N$, a cycle
$y\in Z_{2e+3}(\mathrm{GL}_N(\mathcal O_L);\Z/l^\nu)$ and $0\ne b\in\mathcal O_F$ such
that for every finite field $k$, every ring map $r\colon\mathcal O_F\to k$ with $r|_{\mathcal O_L}$
surjective and $r(b)\ne0$, every homomorphism $v\colon k^\times\to\Z/l^\nu$ and
every additive family $\theta$ of degree $2e+3$ with values in $\Z/l^\nu$ having
value $(r(\zeta),v)$, there is $u\in k^\times$ such that $\mathrm{nrm}_a(\alpha)$
reduces to $u$ under $r$ and
\[
\theta_{k^N}\big((\mathrm{red}_{r|_{\mathcal O_L},N})_*y\big)=\mathrm{pe}\cdot v(u),\qquad
\mathrm{pe}=\prod_{0\le j\le2e+2,\ j\ne e+1}\big((1+l^c)^{e+1}-(1+l^c)^j\big).
\]
\textup{(}No primitivity of $r(\zeta)$ is assumed; $r(\zeta)^{l^\nu}=1$ holds automatically.
If $r(\zeta)$ is primitive, the families of \textup{(i-a)} for the root $r(\zeta)$ qualify.\textup{)}
\end{proposition}

Statement (i) is a Soul\'e-type evaluation of cycles on the homology of
$\mathrm{GL}(V)$ over a finite field; (ii) builds, from the divisor relation
$\operatorname{ord}_P(\alpha)=\mathfrak d(\gamma P)-a\,\mathfrak d(P)$, a module over
$\mathcal O_L$ with an automorphism whose prism cycle reduces at each residue
field to the required eigen-decomposition.  Complete proofs are given in
\S\ref{i5:x5d-ssec}; all three statements are also formally verified in
\cite{CharacterFamily}.

\begin{proposition}[rank bound]\label{i5:x5-prop-bodd}
Let $L$ be totally complex, $k$ even with $k+1<r_2(L)$, $l\ge1$ and
$n\ge2(k+1)+4$.  The $\Z$-span of
$\{z_k(x):x\in H_{k+1}(\mathrm{GL}_n(\mathcal O_L);\Z)\}$ has rank $\le r_2(L)$.
\end{proposition}

Here $z_k$ is as in Definition~\ref{i5:x5-def-prim}.  The statement is
formally verified in \cite{CharacterFamily} by a different route; the proof
below instead derives it from published theorems on the rational homology of
$\mathrm{GL}(\mathcal O_L)$.
The hypotheses ``$k$ even'' and ``$k+1<r_2(L)$'' are not needed for it.

\begin{proof}
We use three published facts.
\emph{(Q) Quillen} \textup{(}Quillen, \emph{Cohomology of groups},
Actes Congr\`es Internat.\ Math.\ Nice 1970, vol.~2, 47--51, reference to be confirmed, for the plus construction
and $K_i=\pi_iB\mathrm{GL}^+$; Loday, \emph{K-th\'eorie alg\'ebrique et
repr\'esentations de groupes}, Ann.\ Sci.\ \'Ec.\ Norm.\ Sup.\ 9 (1976), \S1.2,
Th.~1.2.6, for the H-space structure induced by block sum\textup{)}. $H_*(\mathrm{GL}(R);\Q)=\varinjlim_nH_*(\mathrm{GL}_n(R);\Q)$
equals $H_*(B\mathrm{GL}(R)^+;\Q)$, and $B\mathrm{GL}(R)^+$ is a homotopy
commutative H-space whose product is induced by block sum $\oplus$; so
$\mathcal H=H_*(\mathrm{GL}(R);\Q)$ is a connected graded commutative and
cocommutative Hopf algebra, the coproduct induced by the diagonal and the
product by $\oplus$, and $K_i(R)=\pi_i(B\mathrm{GL}(R)^+)$ for $i\ge1$.
Moreover: the stabilisation map
$s_m\colon H_*(\mathrm{GL}_m(R);\Q)\to\mathcal H$ is a map of coalgebras (it is
induced by a map of spaces, and the coproducts come from the diagonals via the
K\"unneth isomorphism over $\Q$); and for $a,b\ge0$ the product of $\mathcal H$
restricted to $s_a(H_*(\mathrm{GL}_a))\otimes s_b(H_*(\mathrm{GL}_b))$ is
$s_{a+b}\circ(\oplus)_*$, independently of the positions of the blocks, since
permutation matrices act on $\mathcal H$ by inner automorphisms, hence
trivially.
\emph{(MM) Milnor--Moore, Cartan--Serre} \textup{(}Milnor--Moore,
\emph{On the structure of Hopf algebras}, Ann.\ of Math.\ 81 (1965), 211--264, Appendix (Theorem),
with Cor.~4.18 and Thm.~5.18\textup{)}.  For such an H-space the rational
Hurewicz map is an isomorphism of $\pi_*\otimes\Q$ onto the primitives
$P\subset\mathcal H$, and $\mathcal H$ is the symmetric algebra
$\mathrm{Sym}(P)$ (graded sense).
\emph{(B) Borel} \textup{(}Borel, \emph{Stable real cohomology of
arithmetic groups}, Ann.\ Sci.\ \'Ec.\ Norm.\ Sup.\ 7 (1974), 235--272, Prop.~12.2\textup{)}.  For a number
field $L$ with $r_1$ real and $r_2$ complex places and $i>1$,
$\dim_\Q K_i(\mathcal O_L)\otimes\Q$ is
$r_1+r_2$ for $i\equiv1\bmod4$, $r_2$ for $i\equiv3\bmod4$, and $0$ for
even $i$.  The case $i=1$ is not Borel's theorem: $K_1(\mathcal O_L)\otimes\Q
\cong\mathcal O_L^\times\otimes\Q$ \textup{(}Bass--Milnor--Serre, Publ.\ Math.\ IH\'ES 33 (1967), Cor.~4.3(a): $SK_1(\mathcal O_L)=0$\textup{)}, of dimension $r_1+r_2-1$ by
Dirichlet's unit theorem; this case occurs \textup{(}$k=0$\textup{)}.  For $L$
totally complex ($r_1=0$) the dimension is $\le r_2(L)$ for every $i\ge1$.

Let $M=l^kn$ and $s\colon H_{k+1}(\mathrm{GL}_M(\mathcal O_L);\Q)\to
\mathcal H_{k+1}$ the stabilisation.  By Lemma~\ref{i5:x5-lem-k3}(a) (with
$i=k+1$, $M\ge n\ge2(k+1)+4$) each stabilisation step is an isomorphism in
degree $k+1$, so $s$ is injective.  The rank of a $\Z$-module is the
$\Q$-dimension of its tensor product with $\Q$, and
$H_{k+1}(\,\cdot\,;\Z)\otimes\Q=H_{k+1}(\,\cdot\,;\Q)$; hence the rank in
question is the dimension of the $\Q$-span of the $s(z_k(x))$, and it suffices
to show $s(z_k(x))\in P_{k+1}$, since $\dim P_{k+1}=\dim K_{k+1}(\mathcal
O_L)\otimes\Q\le r_2(L)$ by (MM) and (B).

\emph{The operations.} On $\mathrm{GL}_m$, $\psi_d(g)=1_d\otimes g$ is, up to
conjugation by a permutation matrix, the composite of the diagonal
$\mathrm{GL}_m\to\mathrm{GL}_m^d$ and block sum; inner automorphisms act
trivially on group homology, and $\psi_d$ commutes with stabilisation up to
such a conjugation.  By the two facts stated after (Q) ($s$ is a coalgebra
map, and the product on stabilised classes is $s\circ(\oplus)_*$), it follows
that $s\circ\psi_{d*}=[d]\circ s$, where
$[d]=\mu^{(d)}\circ\Delta^{(d)}$ is the $d$-th convolution power of the
identity of $\mathcal H$, and $s\circ\mathrm{st}_*=s$.  As $\mathcal H$ is
commutative, $[d]$ is an algebra map; on a primitive $p$,
$\Delta^{(d)}p=\sum_j1\otimes\cdots\otimes p\otimes\cdots\otimes1$ gives
$[d]p=dp$.  Hence $[d]$ acts on $\mathrm{Sym}^m(P)$ by $d^m$.

\emph{Conclusion.} Write $s(x)=\sum_{m\ge1}x_m$ with $x_m\in\mathrm{Sym}^m(P)$
of degree $k+1$ ($m\ge1$ as the degree is positive).  Since $P$ is
concentrated in degrees $\ge1$, $x_m=0$ for $m>k+1$.  The recursion
$z_{t+1}=\psi_{l*}z_t-l^{t+2}\mathrm{st}_*z_t$ gives
$s(z_{t+1})=([l]-l^{t+2})s(z_t)$, so
\[
s(z_k(x))=\sum_{m=1}^{k+1}\prod_{t=0}^{k-1}\bigl(l^m-l^{t+2}\bigr)\,x_m .
\]
For $2\le m\le k+1$ the factor with $t=m-2\in[0,k-1]$ vanishes.  Hence
$s(z_k(x))=\prod_{t<k}(l-l^{t+2})\,x_1\in P_{k+1}$.
\end{proof}

\begin{lemma}[stability, Bockstein, counting]\label{i5:x5-lem-k3}
\textup{(a)} If $R$ is a Dedekind domain, $i\ge0$, $n\ge2i+4$ and $n'\ge n$,
then $\mathrm{st}_*\colon H_i(\mathrm{GL}_n(R);\Z)\to H_i(\mathrm{GL}_{n'}(R);\Z)$
is an isomorphism; so for $L$ totally complex
$\#H_i(\mathrm{GL}_{n'}(\mathcal O_L),\Z)_{\rm tors}=\mathrm{tc}(L,n,i)$.
\textup{(b)} A primitive family $\Phi$ on $(m+1)$-chains induces a primitive
family $\varphi_n([z])=\Phi_n(\rho_qz)$ on $H_{m+1}(\mathrm{GL}_n(R);\Z)$.
\textup{(c)} If $l^\nu B=0$ and $\varphi$ is primitive with values in $B$, then
for all $n,t$ there is $u\in\Z$ with $l^t\varphi_n(x)=\varphi_{l^tn}(u\,z_t(x))$
for all $x$.
\textup{(d)} For a group $G$ there is a homomorphism
$\beta\colon Z_{k+1}(G;\Z/q)\to H_k(G;\Z)_{\rm tors}$ such that every
$y\in\ker\beta$ is $\rho_qz$ for an integral cycle $z$.
\textup{(e)} If $\beta\colon Y\to T$, $g\colon Y\to B$ are homomorphisms of
abelian groups, $T$ finite, and finite subsets of $g(\ker\beta)$ have at most
$c$ elements, then finite subsets of $g(Y)$ have at most $c\,\#T$ elements.
\textup{(f)} If $H$ is abelian with finite torsion, $J\le H$ has rank $\le r$
and $f\colon H\to B$ with $qB=0$, then finite subsets of $f(J)$ have at most
$q^r\#H_{\rm tors}$ elements.
\end{lemma}

\begin{proof}
(a) We use only the following fact: \emph{for a Dedekind domain $R$, $i\ge0$
and $n\ge2i+4$, the stabilisation map
$H_i(\mathrm{GL}_n(R);\Z)\to H_i(\mathrm{GL}_{n+1}(R);\Z)$ is bijective.}  It
follows from van der Kallen's homological stability theorem (van der Kallen,
\emph{Invent.\ Math.}\ 60 (1980)) together with Bass's bound on the stable rank
of a noetherian ring of dimension $\le1$: van der Kallen, Thm.~4.11 with $e=\max(1,\mathrm{sdim})=1$
(sdim $=1$ by Bass's stable range theorem, quoted in van der Kallen \S2.2), giving bijectivity for $n\ge2i+1$;
the formalization proves the fact itself.
Compose, and restrict to torsion; finiteness of the torsion is part of STors.
(b) Reduction mod $q$ maps cycles to cycles and boundaries to boundaries and
commutes with $\mathrm{st}_*$, $\psi_{d*}$.
(c) By induction, $\varphi(z_t(x))=l^t(1-lc_t)\varphi_n(x)$ with $c_0=0$,
$c_{t+1}=c_t+l^t-l^{t+1}c_t$: primitivity turns the recursion into
multiplication by $l-l^{t+2}$.  Take $u=\sum_{i<\nu}(lc_t)^i$, so that
$u(1-lc_t)\equiv1$ modulo $l^\nu$.
(d) Lift $y$ to an integral chain $x$; then $\partial x=qw$ with $w$ a cycle
and $q[w]=0$; put $\beta(y)=[w]$.  Changing the lift changes $w$ by a boundary;
if $[w]=0$, $w=\partial v$ and $z=x-qv$ works.
(e) Partition a finite $S\subset g(Y)$ by the value of $\beta$ on chosen
preimages; translating a fibre by one of its elements embeds it in
$g(\ker\beta)$.
(f) The preimages generate a subgroup of rank $s'\le r$; writing its elements as
$\sum c_iv_i+\tau$ with $\tau$ torsion and reducing $c_i$ mod $q$ bounds the image
by $q^{s'}\#H_{\rm tors}$.
\end{proof}

\begin{proposition}[counting]\label{i5:x5-prop-upper}
Let $L$ be totally complex, $k$ even with $k+1<r_2(L)$, $q=l^\nu$, $B$ a
$\Z/q$-module and $\Phi$ a primitive family on $(k+1)$-chains of
$\mathrm{GL}_\bullet(\mathcal O_L)$.  Every finite set $S$ of elements
$l^k\Phi_n(y)$, $y\in Z_{k+1}(\mathrm{GL}_n(\mathcal O_L);\Z/q)$, has
$\#S\le q^{r_2(L)}\mathrm{tc}(L,2k+6,k+1)\,\mathrm{tc}(L,2k+6,k)$.
\end{proposition}

\begin{proof}
Stabilise all $y$ to one $n\ge2k+6$ (this does not change $\Phi$).  With $\beta$
from Lemma~\ref{i5:x5-lem-k3}(d), $T=H_k(\mathrm{GL}_n;\Z)_{\rm tors}$ has order
$\mathrm{tc}(L,2k+6,k)$ by (a), so by (e) it suffices to bound finite subsets of
$l^k\Phi_n(\ker\beta)$ by $q^{r_2(L)}\mathrm{tc}(L,2k+6,k+1)$.  For
$y=\rho_qz\in\ker\beta$, (b) and (c) give
$l^k\Phi_n(y)=\varphi_M(u\,z_k([z]))$, $M=l^kn$.  So these values lie in
$\varphi_M(J)$, $J$ the span of the $z_k(x)$ in
$H_{k+1}(\mathrm{GL}_M(\mathcal O_L);\Z)$, of rank $\le r_2(L)$ by
Proposition~\ref{i5:x5-prop-bodd}; (a) and (f) conclude.
\end{proof}

\begin{proposition}[power detection]\label{i5:x5-prop-powdet}
Let $L\subset F$ be number fields, $q=l^\nu$, $\zeta\in\mathcal O_F$ a
primitive $q$-th root of unity, $\beta\in F^\times$ and $0\ne b\in\mathcal O_F$.
If at every maximal ideal $P\not\ni b$ of $\mathcal O_F$ with
$\mathcal O_L\to\mathcal O_F/P$ surjective, $\beta$ is a $P$-unit whose residue
is a $q$-th power, then $\beta\in(F^\times)^q$.
\end{proposition}

\begin{proof}
Let $x^q=\beta$ and $E=F(x)$, Galois over $F$.  Given $g\in\mathrm{Gal}(E/F)$,
Chebotarev's theorem (Neukirch VII (13.4), with the density-one property of
degree-one primes) gives a prime $v\nmid bq$ of $F$ of prime norm $p$,
unramified in $E$, and $w\mid v$ with Frobenius $g$.  Since $\mathcal O_F/v=\mathbb F_p$,
the hypothesis applies at $v$: $x\equiv c\bmod w$ with $c=\zeta^is\in\mathcal O_F$
($X^q-s^q$ splits over $\mathcal O_F$).  Then
$g(x)-x\equiv(x-c)^p-(x-c)\equiv0\bmod w$, while $g(x)=\eta x$ with $\eta^q=1$;
so $\eta\equiv1\bmod v$, and $\eta\ne1$ would give $q\equiv\sum_i\eta^i=0\bmod v$.
Hence $g(x)=x$ for all $g$, and $x\in F$.
\end{proof}

\begin{proposition}[SFew]\label{i5:x-prop:sfew}
Statement (4) of Definition~\ref{i5:x-def:lemmaK-steps} holds, given
Propositions~\ref{i5:x5-prop-detect-construct} and \ref{i5:x5-prop-bodd}.
\end{proposition}

\begin{proof}
Let $S$ be as in (4), $q=l^\nu$, $s=l^{2e+2+(2e+3)(c+2e+3)}$.  For $\alpha\in S$
take $(N_\alpha,y_\alpha,b_\alpha)$ from (ii), put $b^*=l\prod_\alpha b_\alpha$,
and let $I$ be the set of maximal ideals $P\not\ni b^*$ of $\mathcal O_F$ with
$\mathcal O_L\to\mathcal O_F/P=k_P$ surjective.  For $P\in I$, $\zeta$ stays a
primitive $q$-th root of unity in $k_P$ ($l\notin P$), and $k_P^\times$ is
cyclic of order divisible by $q$, so there is $v_P\colon k_P^\times\to\Z/q$ with
kernel $(k_P^\times)^q$.  Let $\theta^P$ come from (i-a) and $\Phi^P$ from
(i-b).

\emph{The constant.} If $0<d<2e+3$ then $l^{c+2e+3}\nmid(1+l^c)^d-1=l^c(d+l^ct)$,
since otherwise $l^{\min(c,2e+3)}\le d<2e+3<l^c$; each factor of $\mathrm{pe}$
is a unit times such a difference.  Hence $\mathrm{pe}\cdot x=l^{(2e+3)(c+2e+3)}$
for some $x\in\Z/q$.

\emph{Injectivity.} Put $f(\alpha)=(l^{2e+2}\Phi^P_{N_\alpha}(y_\alpha))_{P\in I}$.
If $f(\alpha)=f(\alpha')$, then for each $P$, with $u,u'$ the residues of
$\mathrm{nrm}_a(\alpha),\mathrm{nrm}_a(\alpha')$, (ii) gives
$l^{2e+2}\mathrm{pe}\,v_P(u)=l^{2e+2}\mathrm{pe}\,v_P(u')$; multiplying by $x$,
$s\,v_P(u)=s\,v_P(u')$, so $(u/u')^s$ is a $q$-th power in $k_P$.
Proposition~\ref{i5:x5-prop-powdet} with $b=b^*$ makes
$\mathrm{nrm}_a(\alpha)^s\mathrm{nrm}_a(\alpha')^{-s}$ a $q$-th power in $F$, and
the separation hypothesis gives $\alpha=\alpha'$.

\emph{Count.} The product family $(\Phi^P)_{P\in I}$ is primitive with values in
$(\Z/q)^I$, so Proposition~\ref{i5:x5-prop-upper} with $k=2e+2$ (even,
$k+1<r_2(L)$, $2k+6=4e+10$) gives
$\#S=\#f(S)\le q^{r_2(L)}\mathrm{tc}(L,4e+10,2e+3)\mathrm{tc}(L,4e+10,2e+2)$.
\end{proof}

\subsubsection{OneCharU}

\medskip\noindent\emph{Set-up.} Let $K$ be imaginary quadratic with character $\eta$, $l$
odd, $e\ge1$, $N$ even, $\psi$ of level $N$ with $\mathrm{GD}(K,l,e,\psi)$,
$\Psi=\chi^e\psi$ and $t_{e,\psi}(\sigma)=\Psi(\sigma)+1+\Psi(\sigma)^{-1}$.
Let $\mathcal G_3=\mathrm{GL}_3\rtimes\{\pm1\}$, $-1$ acting by
$g\mapsto J\,{}^tg^{-1}J^{-1}$ ($J$ antidiagonal).  The base representation is
$r_0(g)=(\mathrm{diag}(\Psi(g),1,\Psi(g)^{-1}),\eta(g))$; fix a complex
conjugation $c$, a place above $l$ and an element $g_1\in G_\Q$.  Let $R$ be the
complete local noetherian $\mathcal O$-algebra pro-representing lifts $r$ of
$\bar r_0$ with multiplier $\eta$ that are trivial on the inertia elements away
from $l$ acting trivially through $\bar r_0$, satisfy the ordinary local
condition at $l$, and are \emph{normalised at $(c,g_1)$}: $A(c)=1$ and
$Z=A(g_1)-A(cg_1^{-1}c)$ has $Z_{01}=Z_{10}=0$ ($A$ the matrix component).  A
point $f$ of $R$ gives $r_f$ and $\rho_f=r_f|_{G_K}$; it is of \emph{type I} if
$\rho_f$ is irreducible, of \emph{type II} if $\rho_f$ is block diagonal with an
irreducible corner block on $\{e_0,e_2\}$.  It is \emph{two-signed} if it is
neither upper nor lower triangular on $G_K$.  We call $g_1$ a \emph{marking} if
$c^2=1$, $\eta(c)=-1$, $\Psi(c)=1$, $\eta(g_1)=1$ and
$\Psi(g_1)^2-1\in\mathcal O^\times$.

\begin{lemma}\label{i5:x5-lem-unram}
\textup{(a)} If $l\nmid\operatorname{ord}\psi$, then $\psi$ takes values in
$\Z_l[\mu_m]$, $m=\operatorname{ord}\psi$, an unramified coefficient ring.
\textup{(b)} If $\psi,\psi'$ have values in $\mathcal O,\mathcal O'$ with equal
images in $\overline\Q_l$, then $V'(e,\psi)=\infty\iff V'(e,\psi')=\infty$ and
$\mathrm{Vn}(e,\psi)=\mathrm{Vn}(e,\psi')$.
\end{lemma}

\begin{proof}
(a) $X^m-1$ is separable mod $l$.  (b) Pass to a common $\mathcal O''$; then the
Galois module becomes a direct sum of $[\mathcal O'':\mathcal O]$ copies, the
Selmer group (a kernel of restriction maps) becomes
$\mathcal O''\otimes_{\mathcal O}\mathcal S$ by flatness, and lengths multiply by
$e_{\mathcal O''}/e_{\mathcal O}$; dividing by the ramification index gives the
claim.
\end{proof}

\begin{lemma}[shape of points]\label{i5:x5-lem-shape}
Let $g_1$ be a marking and $f$ a two-signed point.  Then $\rho_f$ is of type I or
II, and every point $f'$ with $\rho_{f'}\cong\rho_f$ is conjugate to $f$ by an
element of the diagonal torus $\mathrm{tor}(t)=\mathrm{diag}(t,1,t^{-1})$.
\end{lemma}

\begin{proof}
Since $r(cgc)$ acts through the outer automorphism,
$A(cg^{-1}c)=J\,{}^tA(g)J$ for $g\in G_K$; together with the normalisation
this forces $Z=\mathrm{diag}(z,0,-z)$, and at $r_0$,
$z=\Psi(g_1)-\Psi(g_1)^{-1}$ is a unit.  As $r_f\equiv r_0$ modulo the maximal
ideal, $z(r_f)$ is a unit and $z(r_{f'})\ne-z(r_f)$ ($l\ne2$).
\emph{Shape.} A $G_K$-stable subspace is stable under $\mathrm{diag}(z,0,-z)$,
which has distinct eigenvalues, so it is spanned by basis vectors; the symmetry
$A_{ij}=0\Rightarrow A_{2-j,2-i}=0$ on $G_K$ reduces the six cases to: upper
triangular, lower triangular, or block diagonal with corner $\{0,2\}$, in which
case the corner is irreducible (type II) or, by the same argument with
$\mathrm{diag}(z,-z)$, triangular.  A two-signed point is not triangular.
\emph{Conjugacy.} If $MA=A'M$ on $G_K$, then $MZ=Z'M$, and since no diagonal entry
of $Z$ equals a different-index entry of $Z'$, $M$ is diagonal; applying the
relation to $cg^{-1}c$ shows $A'=\mathrm{tor}(t)A\,\mathrm{tor}(t)^{-1}$ on $G_K$
for suitable $t$ (with $t^2=m_0/m_2$ in the case $m_0m_2\ne m_1^2$, where the
entries $A_{01},A_{10},A_{12},A_{21}$ vanish).  This extends to $c$ since
$\mathrm{tor}(t)$ is fixed by the outer automorphism.
\end{proof}

The cores are the following, in the language of
Definition~\ref{i5:x5-def-dat} (for a datum $D$ the set-up above is the case
$\eta=\eta_D$, $c=c_D$; the ring $R$ of the set-up is a global ring
$(R_D,u_D)$).

\begin{proposition}[(U1); formally verified, \cite{CharacterFamily}]\label{i5:x5-prop-U1}
Let $K$ be imaginary quadratic, $\mathcal O$ a coefficient ring with
$e_{\mathcal O}=1$, $N\ge1$, $e\ge1$, $\psi\colon(\Z/N)^\times\to\mathcal
O^\times$, $D$ a datum for $(K,\mathcal O,e,\psi)$ and $\iota\colon\overline\Q_l\cong\C$.
There is $g_1\in G_\Q$ such that $(\eta,\Psi,c_D,g_1)$ is a marking and, for
every universal lift $(R^{\rm loc},u^{\rm loc})$ and every global ring
$(R_D,u_D)$ for $(D,g_1)$:
\textup{(F)} $W_0(u_D)$ is a finitely generated $\mathcal O$-module;
\textup{(A)} for every point $(\xi,r)$ of type~I, with multiplier $\eta$,
$\mathrm{Aut}_\iota(\rho_K(r)\otimes\bar\chi^{-1})$;
\textup{(W)} for every point $(\xi,r)$ of type~II with
$\operatorname{mat}r(g)_{11}=1$ for $g\in\ker\eta$, and every
$W\colon G_\Q\to\mathrm{GL}_2(\overline\Q_l)$ extending $r$, $\mathrm{Mod}_\iota(W)$.
\end{proposition}

\begin{proposition}[(U2), (U3); formally verified, \cite{CharacterFamily}]\label{i5:x5-prop-U234}
Let $K,\mathcal O,N,e,\psi,D$ be as in Proposition~\ref{i5:x5-prop-U1} but with
$e_{\mathcal O}$ arbitrary, $g_1$ marked, $(R_D,u_D)$ a global ring for
$(D,g_1)$, $f\colon R_D\to\overline\Q_l$ a local $\mathcal O$-algebra map, $r=f_*u_D$,
and $\iota\colon\overline\Q_l\cong\C$.
\textup{(U2a)} If $r$ has multiplier $\eta$, is of type~I and
$\mathrm{Aut}_\iota(\rho_K(r)\otimes\bar\chi^{-1})$, then
$\mathrm{BK}(\mathrm{Ad}\,r)$ (Definition~\ref{i5:x5-def-BK}, $F=\Q$,
$V=M_3(\overline\Q_l)$).
\textup{(U2b)} If $r$ is of type~II and $W\colon G_\Q\to\mathrm{GL}_2(\overline\Q_l)$ is
continuous, extends $r$, has $\det W=\eta$ and $\mathrm{Mod}_\iota(W)$, then
$\mathrm{BK}(R)$ for every $\theta$ with $\theta=\eta$ and every adjoint $R$ of
$W|_{G_K}$ twisted by $\theta$.
\textup{(U3)} If $r$ has multiplier $\eta$, is of type~I and
$\mathrm{Aut}_\iota(\rho_K(r)\otimes\bar\chi^{-1})$, then $\rho_K(r)$ is stably
admissible for $(e,N)$ (Definition~\ref{i5:x-def:bab-adm}(c)).
\end{proposition}

\begin{proposition}[(U3$'$), (U4): type~II admissibility and the length estimate; formally verified, \cite{CharacterFamily}]\label{i5:x5-prop-U4}
Let $K,\mathcal O,N,e,\psi,D,\iota$ be as in Proposition~\ref{i5:x5-prop-U1}
(so $e_{\mathcal O}=1$), let $g_1$ be as there, and let $(R_D,u_D)$ be a global
ring for $(D,g_1)$.
\textup{(U3$'$)} If $f\colon R_D\to\overline\Q_l$ is local, $r=f_*u_D$ is of type~II, and
the conclusion (W) of Proposition~\ref{i5:x5-prop-U1} holds for the extensions
of $r$, then $\rho_K(r)$ is endoscopically admissible for $(e,N|d_K|)$
(Definition~\ref{i5:x-def:bab-adm}(d)).
\textup{(U4)} Assume the conclusions of Propositions~\ref{i5:x5-prop-U1}
and~\ref{i5:x5-prop-U234} and of Proposition~\ref{i5:x-prop:katotwow}, the
published Galois-cohomological statements of Theorem~\ref{i5:thm-pt}, and
Corollary~\ref{i5:cor-lemmaK}.  Let
$\sigma\in G_K$ be such that $\operatorname{tr}\operatorname{mat}f_*u_D(\sigma)
\ne t_{e,\psi}(\sigma)$ for every two-signed $f$.  Then $V'(e,\psi)<\infty$ and
there is a finite set $T$ of two-signed $f$, such that $f_*u_D$ and
$f'_*u_D$ are not conjugate by $\operatorname{tor}(t)$, $t\in\overline\Q_l^\times$, for
$f\ne f'$ in $T$, with
\[
\mathrm{Vn}(e,\psi)\le2\sum_{f\in T}v_l\bigl(\operatorname{tr}\operatorname{mat}
f_*u_D(\sigma)-t_{e,\psi}(\sigma)\bigr).
\]
\end{proposition}

\noindent\emph{Structure of the proofs.}  For Propositions~\ref{i5:x5-prop-U1},
\ref{i5:x5-prop-U234} and \ref{i5:x5-prop-U4}, which steps are published
theorems, which are re-derived in \cite{CharacterFamily}, and which are new, see
\S\ref{i5:x6-sssec-U1}--\S\ref{i5:x6-sssec-U4}.

These statements are exact.  Proposition~\ref{i5:x5-prop-U1} is a single formal
statement; (U2a), (U2b), (U3) are three formal statements; (U3$'$) and (U4) are
not single formal declarations but are assembled, inside the formal proof of
$\mathrm{OneCharU}$, from formally verified lemmas (the existence of a datum,
of the universal rings and of a graded model, the shape of two-signed points,
and the two commutative-algebra estimates of Proposition~\ref{i5:x5-prop-UM}
below); \S\ref{i5:x5-ssec-arch} gives their architecture.  In (U2a) and (U3) the
hypothesis ``$\mathrm{Aut}_\iota$'' is supplied by (A) of
Proposition~\ref{i5:x5-prop-U1}; in (U4) the rigidity at every two-signed point
needed by the length estimate comes from (U2a), (U2b) and
Proposition~\ref{i5:x-prop:katotwow} (for the type-II points).


\begin{proposition}[$\mathrm{OneCharU}$]\label{i5:x-prop:onechar}
$\mathrm{OneCharU}$ holds, given Propositions~\ref{i5:x5-prop-U1} and
\ref{i5:x5-prop-U234}, \ref{i5:x5-prop-U4} and \ref{i5:x-prop:katotwow}.
\end{proposition}

\begin{proof}
\emph{Unramified coefficients.} $\mathrm{GD}$ includes $l\nmid\operatorname{ord}\psi$;
by Lemma~\ref{i5:x5-lem-unram} we may replace $\mathcal O$ by $\Z_l[\mu_m]$
without changing $\mathrm{GD}$, $t_{e,\psi}$, finiteness of $V'$ or
$\mathrm{Vn}$.  So $e_{\mathcal O}=1$.
\emph{Points.} Choose a datum $D$ for $(K,\mathcal O,e,\psi)$ (its existence
from $\mathrm{GD}$ is a formally verified lemma), take $g_1$ from
Proposition~\ref{i5:x5-prop-U1} and a global ring $(R_D,u_D)$ (formally verified
to exist).  A two-signed point is of type I or II (Lemma~\ref{i5:x5-lem-shape});
at type-I points (A) of Proposition~\ref{i5:x5-prop-U1} and (U3) give stable
admissibility, at type-II points (W) and (U3$'$) give endoscopic admissibility; so
$\rho_f\in\mathrm{Adm}(e,N)$.  For $\sigma\in G_K=\ker\eta$ the trace of $r_0$
is $t_{e,\psi}(\sigma)$, so the trace difference at $f$ is
$\operatorname{tr}\rho_f(\sigma)-t_{e,\psi}(\sigma)$, non-zero by the hypothesis
of $\mathrm{OneCharU}$.
\emph{Bound.} (U4) (whose hypotheses (A), (W), (U2a), (U2b) and
Proposition~\ref{i5:x-prop:katotwow} hold) gives $V'<\infty$ and $T$.  If $f,f'\in T$ and
$\rho_f\cong\rho_{f'}$, Lemma~\ref{i5:x5-lem-shape} makes them torus
conjugate, so $f=f'$.  Thus $X=\{\rho_f:f\in T\}$ consists of pairwise
non-isomorphic elements of $\mathrm{Adm}(e,N)$ in bijection with $T$, and the
bound of (U4) is the bound of $\mathrm{OneCharU}$.
\end{proof}

\subsubsection{$\mathrm{Count}_{\rm II}$ and $\mathrm{Count}_{\rm III}$}

Notation is that of Definition~\ref{i5:def-cases}: $h$ a host with character
$\psi$, $\psi'\in X_h(\mathcal O,e)$.  To $(h,l,\psi')$ are attached exponents
$n_{\rm II},n_{\rm III}$ (in Case II, $n_{\rm II}$ is $\lceil\log_2\rceil$ of a
maximum one of whose terms is $2^{c_2(\psi')+6}o_hl^3q_h^2$).

\begin{lemma}\label{i5:x5-lem-transfer}
Let $f\colon\mathcal O\hookrightarrow\mathcal O'$ and $\psi''=f\circ\psi'$.  Then
$\psi''$ has the same order, kernel, conductor, components, $c_2$, $n_0$,
parity and exponents $n_{\rm II},n_{\rm III}$; $\psi''\in X_h(\mathcal O',e)$ and
lies in the same case.  Moreover $\mathrm{VLe}(e,\psi'',B)\Rightarrow\mathrm{VLe}(e,\psi',B)$.
\end{lemma}

\begin{proof}
Composition with an injective map commutes with all operations defining these
invariants and preserves kernels.  The last assertion is
Lemma~\ref{i5:x5-lem-unram}(b).
\end{proof}

\begin{proposition}[$\mathrm{GSp}_4$ cores; formally verified, \cite{CharacterFamily}]\label{i5:x5-prop-sp4core}
Notation of Definition~\ref{i5:x5-def-gsp4}.  Let $h$ be a host with
$\mathrm{Big}_h(t)$, $l\ne2$ prime, $\mathcal O$ a coefficient ring, $0<e<t$ and
$\psi'\in X_h(\mathcal O,e)$.
\begin{enumerate}
\item[(a)] (automorphy, finiteness) For every good $D$ with
$\mathrm{OfCase}(D,h,e,\psi')$: $\mathrm{Fin}_0(D)$, $\mathrm{NoEnd}(D)$, and
$\mathrm{AllIrr}_D(\mathrm{Aut}^{(4)}_\iota)$ for some $\iota\colon\overline\Q_l\cong\C$.
\item[(b)] (existence) If $\psi'$ is in Case~II or~III and $\mathcal O$ contains
the $2^{n_{\rm II}}$-th and $2^{n_{\rm III}}$-th roots of unity, there is a good
$D$ with $\mathrm{OfCase}(D,h,e,\psi')$.
\item[(c)] (comparison) For every good $D$ with $\mathrm{OfCase}(D,h,e,\psi')$,
$V'(e,\psi')\le\operatorname{sha}(D)$.
\item[(d)] (level) If $D$ is good and $\mathrm{OfChar}(D,e,\psi',n',p_1,m_0)$
with $c_2(\psi')\le n'$ and $m_0\mid o_h$, then $\psi_D,\xi,\omega$ are trivial on
$G_{\Q(\mu_M)}$, $M=l2^{n'}N_hp_1$, and their $k$-th powers are trivial,
$k=(l-1)2^{n'}\cdot2o_h$.
\end{enumerate}
Moreover, for every good $D$:
\begin{enumerate}
\item[(e)] (rigidity) if $\mathrm{OfCase}(D,h,e,\psi')$ for some host $h$ with
$\mathrm{Big}_h(t)$, $0<e<t$ and $\psi'\in X_h(\mathcal O,e)$, and $\psi_D,\xi,\omega$
are trivial on $G_{\Q(\mu_M)}$ for some $M\ge1$, then for every $\iota$ every irreducible
$\overline\Q_l$-point $r$ with $\mathrm{Aut}^{(4)}_\iota(r)$ satisfies $\mathrm{Rig}_D(r)$;
\item[(f)] (length estimate) if $\mathrm{Fin}_0(D)$, $\mathrm{NoEnd}(D)$ and
$\mathrm{AllIrr}_D(\mathrm{Rig}_D)$, then $\mathrm{Ineq}_D(g)$ for every $g\in
G_\Q$ with $\mathrm{AllIrr}_D(a_g\ne0)$;
\item[(g)] ($\mathrm{NumberSp4}$) there are $C,A\in\N$, independent of $D$, such
that: if $\psi_D,\xi,\omega$ are trivial on $G_{\Q(\mu_M)}$ ($M\ge1$) and have
trivial $k$-th powers ($k\ge1$), $p\ge5$ is a prime with $p\nmid M$, $p\notin S$,
and $\mathrm{AllIrr}_D(\mathrm{Aut}^{(4)}_\iota)$ for some $\iota$, then there is $\sigma\in G_\Q$
with $\mathrm{AllIrr}_D(a_\sigma\ne0)$ and
$\mathrm{Ineq}_D(\sigma)\Rightarrow\mathrm{Vn}(D)\log l\le
2Ck\,(M(e_D+1))^A\log p$.
\end{enumerate}
\end{proposition}

\noindent\emph{Structure of the proof:} \S\ref{i5:x6-sssec-sp4}.

The statement is exact; (a), (e) and the counting statement behind (g) are the
three substantial formal statements, (b), (c), (d), (f) are formally verified
lemmas of the formal proof of $\mathrm{Count}_{\rm II}$ and $\mathrm{Count}_{\rm
III}$.  Their inputs: (b) uses the Bernoulli-unit statement (the residual
non-vanishing of generalised Bernoulli numbers together with the statement that
an unramified residual cocycle for a character whose generalised Bernoulli
number is an $l$-adic unit is a coboundary); (f) uses the commutative-algebra
estimates (M), (M') of Proposition~\ref{i5:x5-prop-UM}, the existence of a
graded model, the published local theorems, and Soul\'e-type vanishing and
finiteness for $\psi^{-1}\chi^{-e}$, derived formally from
Corollary~\ref{i5:cor-lemmaK}, Poitou--Tate and the local theorems; the models
and universal rings use the published Galois-cohomological statements of
Theorem~\ref{i5:thm-pt} (including Greenberg--Wiles and unramified
orthogonality for $\mathrm{GSp}_4$).  Note that (c) is an
inequality, and that the automorphy predicate in (a) is
$\mathrm{Aut}^{(4)}_\iota$, which asks for cuspidal representations of
$\mathrm{GL}_4$ over CM quadratic fields (with the automorphy witness), not for
automorphy for $\mathrm{GSp}_4$ over $\Q$.  (a) is an automorphy lifting
statement for a residual representation which is a sum of four characters; the
transfer to $\mathrm{GL}_4$ and the attached Galois representations (Gee--Ta\"ibi;
Taylor, Laumon, Weissauer) are named for comparison only (references to
be confirmed).


\begin{lemma}[arithmetic estimate]\label{i5:x5-lem-c23arith}
For a positive integer $m$ let $p_0(m)$ be the least prime $\ge5$ not dividing
$m$.  For every character $\chi$ (of the kind attached to hosts) and all
$C,A\in\N$ there is $C'\in\N$ such that: for every host $h$ with character
$\chi$, every odd prime $l$, every $\mathcal O$, every $e\ge2$, every
$\psi'\in X_h(\mathcal O,e)$ and all $n',p_1$ with $(n',p_1)=(n_{\rm II},1)$ or
$(n',p_1)=(n_{\rm III},p_1(\chi))$ with $p_1(\chi)$ the prime of
Definition~\ref{i5:def-Bchi} for the character $\chi$ of $h$,
\[
2C\,(l-1)2^{n'}(2o_h)\,\bigl(l2^{n'}N_hp_1(e+1)\bigr)^A\log p_0(lN_hp_1)\le(lN_he)^{C'} .
\]
\end{lemma}

Here $e$ is independent of the exponent of the host, $\log$ is the
natural logarithm, and, exactly as in the formally verified version in
\cite{CharacterFamily},
$n_{\rm II}=\lceil\log_2\max(lN_h4^{n_0},\,2^{c_2+6}o_hl^3q_h^2)\rceil$ and
$n_{\rm III}=\lceil\log_2\max(lq_hp_14^{n_0},\,8l^2)\rceil$, with
$n_0=n_0(\psi')$, $c_2=c_2(\psi')$ of Definition~\ref{i5:def-cases}, $q_h=N_\chi$,
$o_h=o_\chi$ and $p_1=p_1(\chi)$.

\begin{proof}
Put $X=lN_he$, $f=\varphi(4o_h)$ and $P=p_1(\chi)$; all constants below depend
only on $\chi$.  \emph{Inputs.} $l,N_h,e\le X$; $16\mid N_h$
(Lemma~\ref{i5:lem-hostarith}(1)), so $X\ge16$; $q_h\mid N_h$ and $o_h\le
q_h$, so $o_h,q_h\le X$; $2^{c_2}$ is the conductor of $\mathrm{comp}_2\psi'$,
which divides $N_h$, so $2^{c_2}\le X$; every $k\in\N$ satisfies $k\le2^k\le
X^k$.  Finally $2^{n_0}<l^f$: the order $o$ of $\theta_1$ divides $2o_h$
(its odd components have order dividing $2o_h$ by
Lemma~\ref{i5:lem-hostarith}(2)) and is prime to $l$, so
$M=\mathrm{lcm}(4,o)$ divides $4o_h$ and is prime to $l$; the order $f_0$ of
$l$ mod $M$ divides $\varphi(M)$, which divides $f$; hence
$2^{n_0}\le l^{f_0}-1<l^f$.  In particular $4^{n_0}\le X^{2f}$.

\emph{Bound for $2^{n'}$.} For $m\ge1$, $2^{\lceil\log_2m\rceil}\le2m$.  In
Case II, $lN_h4^{n_0}\le X^{2+2f}$ and $2^{c_2+6}o_hl^3q_h^2\le64X^7$, so with
$X^2\ge128$, $2^{n_{\rm II}}\le\max(2X^{2+2f},128X^7)\le X^{2f+10}$.  In
Case III, $lq_hP4^{n_0}\le X^{1+q_h+P+2f}$ and $8l^2\le X^5$, so
$2^{n_{\rm III}}\le X^{11+q_h+P+2f}$.  Thus $2^{n'}\le X^a$ with
$a=(2f+10)+(11+q_h+P+2f)$.

\emph{Assembly.} With $B=1+P$ we have $p_1\le B\le X^B$, and: $2C\le X^{2C}$;
$(l-1)2^{n'}(2o_h)\le X^{1+a+2o_h}$; $l2^{n'}N_hp_1(e+1)\le X^{1+a+1+B+2}$
(as $e+1\le X^2$); and, by Bertrand's postulate applied to $lN_hp_1+5$,
$p_0(lN_hp_1)\le2lN_hp_1+10\le X^{4+2+B}$, so $\log p_0\le
p_0\le X^{B+6}$.  Multiplying,
the left side is at most $X^{C'}$ with
$C'=2C+(1+a+2o_h)+(a+B+4)A+(B+6)$, which depends only on $\chi$, $C$, $A$.
\end{proof}

\begin{proposition}\label{i5:prop-c23}
$\mathrm{Count}_{\rm II}$ and $\mathrm{Count}_{\rm III}$ hold, given
Proposition~\ref{i5:x5-prop-sp4core}.
\end{proposition}

\begin{proof}
\emph{Main estimate.} Assume first that $\mathcal O$ contains the roots of
unity of Proposition~\ref{i5:x5-prop-sp4core}(b), and let $\psi'$ be in Case~II
or~III.  By (b) there is a good $D$ with $\mathrm{OfCase}(D,h,e,\psi')$; by (a),
$\mathrm{Fin}_0(D)$, $\mathrm{NoEnd}(D)$ and $\mathrm{AllIrr}_D(\mathrm{Aut}^{(4)}
_\iota)$ for some $\iota$.  Unwinding $\mathrm{OfCase}$ gives
$\mathrm{OfChar}(D,e,\psi',n',p_1,m_0)$ with $(n',p_1,m_0)=(n_{\rm II},1,1)$ or
$(n_{\rm III},p_1(\chi),\operatorname{ord}\psi')$; in both cases $c_2(\psi')\le
n'$ (in Case~II because $2^{c_2}$ is at most the second term of the maximum
defining $n_{\rm II}$; in Case~III because $\operatorname{ord}\psi'$ is odd, so
$\mathrm{comp}_2\psi'=1$ and $c_2=0$) and $m_0\mid o_h$ (in Case~III by
Lemma~\ref{i5:lem-hostarith}).  So (d) applies with $M=l2^{n'}N_hp_1$ and
$k=(l-1)2^{n'}\cdot2o_h$.  By (e) (with this $M$) every irreducible
$\overline\Q_l$-point is rigid, i.e.\ $\mathrm{AllIrr}_D(\mathrm{Rig}_D)$.  Let $p_0$ be
the least prime $\ge5$ not dividing $lN_hp_1$.  Then $p_0\nmid2$ and $p_0\nmid
lN_hp_1$, so $p_0\nmid M$ and $p_0\notin S=\{p\mid2lN_hp_1\}$.  By (g) there is
$\sigma$ with $\mathrm{AllIrr}_D(a_\sigma\ne0)$ and
$\mathrm{Ineq}_D(\sigma)\Rightarrow\mathrm{Vn}(D)\log l\le2Ck(M(e+1))^A\log p_0$;
by (f), $\mathrm{Ineq}_D(\sigma)$ holds; in particular
$\operatorname{sha}(D)<\infty$.  By (c), $V'(e,\psi')\le\operatorname{sha}(D)<
\infty$, so
\[
\mathrm{Vn}(e,\psi')\log l\le\mathrm{Vn}(D)\log l\le2C(l-1)2^{n'}(2o_h)
\bigl(l2^{n'}N_hp_1(e+1)\bigr)^A\log p_0 .
\]
Here $e=e_D\ge2$ ($D$ is good).  Lemma~\ref{i5:x5-lem-c23arith} above, applied
with the constants $C,A$ of (g), bounds the right side by $(lN_he)^{C'}$ with
$C'$ depending only on the character of $h$.  Hence
$\mathrm{VLe}(e,\psi',(lN_he)^{C'})$.

\emph{General $\mathcal O$.} Let $\mathcal O'\supseteq\mathcal O$ contain the
$2^{n_{\rm II}+n_{\rm III}}$-th roots of unity.  By
Lemma~\ref{i5:x5-lem-transfer}, $\psi''=f\circ\psi'$ is in the same case with
the same exponents, the main estimate applies to it, and $\mathrm{VLe}$
transfers back to $\psi'$.  The constant $C'$ depends only on the character of
$h$, which is the shape required by $\mathrm{Count}_{\rm II}$ and
$\mathrm{Count}_{\rm III}$.
\end{proof}

\subsubsection{$\mathrm{Count}_{\rm V}$}

In subcase V of Case~V (that is, $l\in B_\psi$, $e=1$ and $R_1\le2^{c_2}$,
Definition~\ref{i5:def-cases}) one works with host frames
(Definition~\ref{i5:x5-def-hostframe}) over $K=\Q(i)$, whose absolute Galois
group is $\ker\eta_4$.  In this subsection $K$ is a CM field with $[K:\Q]=2$
containing a square root of $-1$, i.e.\ $K=\Q(i)$.

\begin{proposition}[unitary lifting over $\Q(i)$; formally verified, \cite{CharacterFamily}]\label{i5:x5-prop-Vcore}
Let $l$ be a prime, $\iota\colon\overline\Q_l\cong\C$, $\mathcal O$ a coefficient ring,
$N\ge1$, $e,o_1\in\N$ and $\psi\colon(\Z/N)^\times\to\mathcal O^\times$ with
$\mathrm{HostData}(l,e,\psi,o_1)$.  For every host frame $(S,(A_p),F)$ for
$(e,\psi)$ there is a marked $g_1\in G_{\Q,S}$ such that for every point $u$ of
$F$ universal at $g_1$, with values in $R$:
\textup{(F)} the weight-zero part $R_0(u)$ is a finitely generated
$\mathcal O$-module;
\textup{(A)} for every point $(\xi,r)$ of $u$ with values in $\overline\Q_l$ such that
$r$ is of type~I (for $\eta$), $\mathrm{AutOn}_{K,\iota,-1}(a_r)$.
\end{proposition}

\begin{lemma}[irreducibility on open subgroups; formally verified, \cite{CharacterFamily}]\label{i5:x5-lem-Vopen}
In the situation of Proposition~\ref{i5:x5-prop-Vcore}, let $u$ be universal at a
marked $g_1$ with $R_0(u)$ finite over $\mathcal O$.  Then for every point
$(\xi,r)$ of $u$ with values in $\overline\Q_l$, $r$ is a characteristic-zero point of
$F$ at $g_1$; and every characteristic-zero point $r$ at $g_1$ of type~I has
$a_r$ irreducible on every open subgroup of $G_\Q$ contained in $\ker\eta_4$.
\end{lemma}

\begin{proposition}[host inputs; formally verified, \cite{CharacterFamily}]\label{i5:x5-prop-Vhost}
Let $l,\mathcal O,N,e,\psi,o_1$ be as in Proposition~\ref{i5:x5-prop-Vcore} and
let $(S,(A_p),F)$ be a host frame for $(e,\psi)$.
\begin{enumerate}
\item[(a)] (rigidity) For every $\iota$, every marked $g_1$ and every
characteristic-zero point $r$ of $F$ at $g_1$ with values in $\overline\Q_l$ and
$\mathrm{B}'_{K,\iota}(r)$, $r$ is rigid for $F$.
\item[(b)] (number) There is $C\in\N$, depending on nothing (for fixed
$K=\Q(i)$), such that for every $\iota$ and every marked $g_1$ there is
$\varphi\in G_{\Q,S}$ with: $\operatorname{tr}\operatorname{mat}r(\varphi)\ne
\operatorname{tr}\operatorname{tor}\Psi(\varphi)$ for every characteristic-zero
point $r$ at $g_1$ of type~I with $\mathrm{B}'_{K,\iota}(r)$; and for every
finite family $r_1,\dots,r_n$ of such points, pairwise not
$\ker\eta$-isomorphic,
\[
\Bigl(\sum_iv_l\bigl(\operatorname{tr}\operatorname{mat}r_i(\varphi)-
\operatorname{tr}\operatorname{tor}\Psi(\varphi)\bigr)\Bigr)\log l\le(lNe)^C.
\]
\item[(c)] (existence) For a host $h$, an odd prime $l$, $e\ge1$ and
$\psi'\in X_h(\mathcal O,e)$ with $R_1\le2^{c_2(\psi')}$,
$\mathrm{HostData}(l,e,\psi',o_h)$ holds (with $N=N_h$), and
$\mathrm{XData}(l,\psi',q_h)$ holds: $q_h\ge1$, the conductor of
$\theta_1^2$ divides $q_h$, and $2lq_h4^{n_0}\le2^{c_2(\psi')}$
(notation of Definition~\ref{i5:def-cases}, for $\psi'$); for any
$\mathrm{HostData}$, a host frame exists; and for every host frame and marked
$g_1$, a point of $F$ universal at $g_1$ exists.
\item[(d)] (abstract bound) Assume that there is $q\ge1$ with
$\mathrm{XData}(l,\psi,q)$ (as in (c)).  Let $g_1$ be marked, $u$ universal at $g_1$ with
$R_0(u)$ finite, $\varphi\in G_{\Q,S}$, and $P$ a property of homomorphisms
$G_{\Q,S}\to\mathcal G_3(\overline\Q_l)$ such that every point $(\xi,r)$ of $u$ of type~I
has $P(r)$, and every characteristic-zero point $r$ at $g_1$ of type~I with
$P(r)$ is rigid and has $\operatorname{tr}\operatorname{mat}r(\varphi)\ne
\operatorname{tr}\operatorname{tor}\Psi(\varphi)$.  Then there are
characteristic-zero points $r_1,\dots,r_n$ at $g_1$ of type~I with $P(r_i)$,
pairwise not $\ker\eta$-isomorphic, and $V\in\N$ with $V'(e,\psi)\le V$ and
$V\le2e_{\mathcal O}\sum_iv_l\bigl(\operatorname{tr}\operatorname{mat}
r_i(\varphi)-\operatorname{tr}\operatorname{tor}\Psi(\varphi)\bigr)$.
\end{enumerate}
\end{proposition}

\noindent\emph{Structure of the proofs:} \S\ref{i5:x6-sssec-V} (Propositions~\ref{i5:x5-prop-hostend}, \ref{i5:x5-prop-Vcore}, Lemma~\ref{i5:x5-lem-Vopen}) and \S\ref{i5:x6-sssec-Vhost} (Proposition~\ref{i5:x5-prop-Vhost}).

These statements are exact.  Proposition~\ref{i5:x5-prop-Vcore}, Lemma~\ref{i5:x5-lem-Vopen}
(two formal statements), (a), (b) and the three parts of (c) are formal
statements; (d) is assembled in the formal proof of $\mathrm{Count}_{\rm V}$ from
the formally verified abstract bound for frames, the comparison
$V'(e,\psi)\le V$ for host frames, two formally verified vanishing statements
for host frames, the commutative-algebra estimates of
Proposition~\ref{i5:x5-prop-UM}, and the published Galois-cohomological inputs
of Theorem~\ref{i5:thm-pt} (local and global finiteness, local duality, local
and global Euler characteristic formulas, Poitou--Tate); see
\S\ref{i5:x5-ssec-arch}.  The inputs of the two vanishing statements are
stated explicitly: the first (a vanishing statement for the cocycles of the base point $r_0$
with values in the weight pieces $j\le0$) is a Soul\'e-type vanishing and finiteness statement,
derived formally from Corollary~\ref{i5:cor-lemmaK}, Poitou--Tate and the
local theorems of Theorem~\ref{i5:thm-pt}; the second (the one using
$\mathrm{XData}$) is derived formally from the Bernoulli-unit statement:
the residual non-vanishing of generalised Bernoulli numbers together with the
statement that an unramified residual cocycle for a character whose
generalised Bernoulli number $B_{1,\chi^{-1}}$ is an $l$-adic unit is a
coboundary.  The existence of a point universal at $g_1$ in (c) uses global
finiteness (Theorem~\ref{i5:thm-pt}).  The hypothesis $R_1\le2^{c_2}$ enters
only through (c), which supplies both $\mathrm{HostData}$ and
$\mathrm{XData}$.
Note that in (b) the constant is uniform and the bound is in terms of the
level $N$ of $\psi$ (for a host, $N=N_h$).

\noindent\emph{Comparison with published automorphy lifting theorems.}  The
residual representations are $\bar r_0=r_0^{\eta,\bar\Psi}$, i.e.\ on $G_K$ the
sum $\bar\Psi\oplus1\oplus\bar\Psi^{-1}$ of three characters
(Propositions~\ref{i5:x5-prop-U1}, \ref{i5:x5-prop-Vcore}), resp.\ the sum
$\bar\vartheta\oplus\bar\Psi_D\oplus\bar\alpha\oplus1$ of four characters
(Proposition~\ref{i5:x5-prop-sp4core}).  (1) Clozel--Harris--Taylor, Thorne
(2012), Barnet-Lamb--Gee--Geraghty--Taylor (2014) and
Boxer--Calegari--Gee--Pilloni all require a residually irreducible (big,
adequate or vast) image; none applies.  (2) The residually reducible theorems
are Thorne (2015) and Allen--Newton--Thorne (\emph{Automorphy lifting for
residually reducible $l$-adic Galois representations, II}, Compositio Math.\ 156 (2020), 2399--2422, arXiv:1912.11269).  The table lists their hypotheses: for
Allen--Newton--Thorne as in their Thm.~1.1 (the table omits hypotheses that do not change the verdicts);
for Thorne (2015) as recorded in the introduction of Allen--Newton--Thorne (reference to be confirmed for Thorne's statements),
against Propositions~\ref{i5:x5-prop-U1} and~\ref{i5:x5-prop-Vcore}; ``holds''
and ``fails'' refer to the exact statements above.
\begin{center}\small
\begin{tabular}{p{0.34\textwidth}p{0.17\textwidth}p{0.17\textwidth}p{0.2\textwidth}}
\toprule
hypothesis & (U1) & (Vcore) & $\mathrm{GSp}_4$ (a)\\
\midrule
$l>3$, $l\nmid n$ (ANT; $n=3$) & holds: $\mathrm{GD}$ forces $(l-1)\nmid2e$, so $l\ge5$ & fails: only $l\ne2$, $l=3$ allowed & not applicable ($\mathrm{GSp}_4$)\\
number of constituents of $\bar\rho^{ss}$: exactly $2$ (Thorne), any (ANT) & $3$: Thorne fails, ANT count fine & $3$: same & $4$ characters\\
constituents pairwise non-isomorphic & holds, by the marking $g_1$ produced by the proposition (part of its conclusion): $\bar\Psi(g_1)^2\ne1$ gives $\bar\Psi\ne1,\bar\Psi^{\pm2}\ne1$ on $G_K$ & same & not determined\\
same after restriction to $G_{K(\zeta_l)}$; $\bar\rho^{ss}$ primitive, no quotient of order $l$ (ANT) & not determined & not determined & --\\
each constituent conjugate self-dual (with the right sign) & not determined & not determined & --\\
$\rho$ ordinary & holds (local condition at $l$) & holds (condition (T)) & ordinary-type condition $P_l$\\
Steinberg-type place $\tilde v_0\nmid l$ for $\rho$ & fails: no such place; lifts trivial on $\mathrm{Kill}_D$ & fails: unramified on $N_v$, diagonal at $S_b$ & --\\
ordinary automorphic $\pi$ with $\overline{r(\pi)}^{ss}\cong\bar\rho^{ss}$, Steinberg at $\tilde v_0$ & fails: not assumed & fails: not assumed & --\\
conclusion: automorphy of one given $\rho$ & stronger: all type-I points of a ring, finiteness (F), type-II modularity (W) & stronger: all type-I points, finiteness (F) & finiteness, no endomorphisms, automorphy\\
\bottomrule
\end{tabular}
\end{center}
So neither published theorem applies as stated: Thorne (2015) because there are
three (or four) constituents, both because no Steinberg-type place and no
automorphic lift with the same residual semisimplification are available in
these deformation problems, and Allen--Newton--Thorne in addition because
$l=3$ is allowed in Proposition~\ref{i5:x5-prop-Vcore}.  Whether the
automorphy conclusions (A) could be deduced from Allen--Newton--Thorne after
adding auxiliary Steinberg places (a level-raising argument) we do not
discuss.  (3) For $\mathrm{GSp}_4$ with residual representation a sum of
characters we know no published lifting theorem.  The cores are therefore taken
from the formal verification.

\begin{lemma}\label{i5:x5-lem-irrcyc}
Let $(S,(A_p),F)$ be a host frame for $(e,\psi)$ with
$\mathrm{HostData}(l,e,\psi,o_1)$, and $g_1$ marked.  Every
characteristic-zero point $r$ of $F$ at $g_1$ of type~I has $a_r$ irreducible on
$\ker\eta_4\cap\ker\chi$ (condition (ii) of Definition~\ref{i5:x5-def-Bprime}).
\end{lemma}

\begin{proof}
We use the following purely algebraic statement (formally verified,
\cite{CharacterFamily}; no topology enters).  \emph{Let $\Gamma$ be a group,
$\kappa$ an algebraically closed field of characteristic $0$,
$\rho\colon\Gamma\to\mathrm{GL}_m(\kappa)$ a homomorphism and $H'$ a normal
subgroup of $\Gamma$ containing every commutator $aba^{-1}b^{-1}$.  If $\rho$ is
irreducible and $\rho|_{H'}$ is reducible, there is a homomorphism
$\phi\colon\Gamma\to\kappa^\times$ with $\phi\ne1$, $H'\subseteq\ker\phi$,
$\phi^m=1$, and $\phi\otimes\rho\cong\rho$.}

Put $H=\ker\eta_4$ and $H'=H\cap\ker\chi$; since $\chi$ is abelian, $H'$ is
normal in $H$ and contains all commutators of $H$.  By
Lemma~\ref{i5:x5-lem-Vopen}, $a_r$ is irreducible on every open subgroup of
$G_\Q$ contained in $H$ (condition (i)); in particular on $H$.  Suppose $a_r$ is
reducible on $H'$.  The statement above with $\Gamma=H$,
$\kappa=\overline\Q_l$, $m=3$ gives $\phi\ne1$ trivial on $H'$ with $\phi^3=1$
and $M\in\mathrm{GL}_3(\overline\Q_l)$ with $Mr(g)=\phi(g)r(g)M$ on $H$.
\emph{$\ker\phi$ is open.} $\phi$ factors through $\chi(H)\cong H/H'$ and kills
cubes, so $\ker\phi\supseteq\{g\in H:\chi(g)\in\chi(H)^3\}$.  As $H$ is open of
finite index in $G_\Q$ and $\chi$ is continuous and surjective onto
$\Z_l^\times$, $\chi(H)$ is a closed subgroup of finite index, hence open, in
$\Z_l^\times$; and $\chi(H)^3$ is open in $\chi(H)$ since
$\Z_l^\times$ is a finite group times a copy of $\Z_l$ and cubing on $\Z_l$ has
open image.  So $\ker\phi$ contains an open subgroup and is open.
By (i) $r$ is irreducible on $\ker\phi$, and $M$ commutes with $r(\ker\phi)$,
so $M$ is scalar by Schur's lemma; then $\phi(g)=1$ for all $g\in H$, a
contradiction.
\end{proof}

\begin{proposition}\label{i5:prop-c5}
$\mathrm{Count}_{\rm V}$ holds, given Propositions~\ref{i5:x5-prop-Vcore} and
\ref{i5:x5-prop-Vhost} and Lemma~\ref{i5:x5-lem-Vopen}.
\end{proposition}

\begin{proof}
Let $C$ be the constant of Proposition~\ref{i5:x5-prop-Vhost}(b) (for
$K=\Q(i)$); we prove $\mathrm{Count}_{\rm V}$ with $C+1$.  Let
$h,t,l,\mathcal O,e,\psi'$ be in subcase V, i.e.\ $l\in B_\psi$, $e=1$ and
$R_1\le2^{c_2}$.  By Proposition~\ref{i5:x5-prop-Vhost}(c),
$\mathrm{HostData}(l,1,\psi',o_h)$ holds with $N=N_h$ and there is a host frame
$(S,(A_p),F)$.  Fix any $\iota$.  Proposition~\ref{i5:x5-prop-Vcore} gives a
marked $g_1$, and (c) a point $u$ universal at $g_1$; by (F), $R_0(u)$ is finite.
By (c), $\mathrm{XData}(l,\psi',q_h)$ holds, which is the additional
hypothesis of (d).
Let $P=\mathrm{B}'_{K,\iota}$.  Every point $(\xi,r)$ of $u$ of type~I has
$P(r)$: $r$ is a characteristic-zero point with (i) by
Lemma~\ref{i5:x5-lem-Vopen}, (ii) by Lemma~\ref{i5:x5-lem-irrcyc}, and (iii) by
(A).  Let $\varphi$ be given by (b) for $(\iota,g_1)$.  Every
characteristic-zero type-I point with $P$ is rigid by (a) and has non-zero
trace difference at $\varphi$ by (b).  So (d) applies and gives points
$r_1,\dots,r_n$ (characteristic-zero, type~I, with $P$, pairwise not
$\ker\eta$-isomorphic) and $V$ with $V'(e,\psi')\le V\le2e_{\mathcal O}s$,
$s=\sum_iv_l(\operatorname{tr}\operatorname{mat}r_i(\varphi)-
\operatorname{tr}\operatorname{tor}\Psi(\varphi))$; by (b), $s\log l\le
(lN_he)^C$.  In particular $V'(e,\psi')<\infty$, and since $lN_he\ge2$,
\[
\mathrm{Vn}(e,\psi')\log l=\frac{V'}{e_{\mathcal O}}\log l\le2s\log l\le2(lN_he)^C\le(lN_he)^{C+1},
\]
which is $\mathrm{VLe}(e,\psi',(lN_he)^{C+1})$.
\end{proof}


\subsubsection{Proof of Proposition~\ref{i5:x5-prop-detect-construct}}\label{i5:x5d-ssec}

We prove the three statements in the form of
Proposition~\ref{i5:x5-prop-detect-construct}. Throughout, $A$ is a commutative ring
and $G,H$ are groups. The bar differential is
$\partial[g_1|\dots|g_{m+1}]=[g_2|\dots|g_{m+1}]+\sum_{i=1}^{m}(-1)^i[\dots|g_ig_{i+1}|\dots]+(-1)^{m+1}[g_1|\dots|g_m]$.
A homomorphism $f$ induces $f_*[g_1|\dots|g_m]=[f(g_1)|\dots|f(g_m)]$, and
$\kappa_k(g)=k^{-1}gk$. ``Cycle'' and ``boundary'' refer to $\partial$. For an
isomorphism $\varphi\colon V\to W$, $\iota_\varphi(g)=\varphi g\varphi^{-1}$.

\medskip\noindent\emph{The homogeneous model.}
For $n\ge-1$ let $P_n$ be the free $A$-module on $G^{n+1}$ ($P_{-1}=A$), with
$\partial^h(x_0,\dots,x_n)=\sum_j(-1)^j(x_0,\dots,\widehat{x_j},\dots,x_n)$
($\partial^h=0$ on $P_{-1}$). For $a\in G$ the cone is
$c_a(x_0,\dots,x_n)=(a,x_0,\dots,x_n)$. A map of sets $f\colon G\to G$ acts by
$f_*(x_0,\dots,x_n)=(f(x_0),\dots,f(x_n))$ ($f_*=\mathrm{id}$ on $P_{-1}$), and
$H_f(x_0,\dots,x_n)=\sum_{i=0}^n(-1)^i(x_0,\dots,x_i,f(x_i),\dots,f(x_n))$ ($H_f=0$ on
$P_{-1}$). Let $r_k(x)=xk$, $\lambda_h(x)=hx$,
$\pi(x_0,\dots,x_n)=[x_0^{-1}x_1|\dots|x_{n-1}^{-1}x_n]$ and
$s[g_1|\dots|g_n]=(1,g_1,g_1g_2,\dots,g_1\cdots g_n)$.

\begin{lemma}\label{i5:x5d-lem-homog}
For $n\ge-1$, $a,k\in G$ and maps $f\colon G\to G$:
\textup{(1)} $\partial^hc_ay=y-c_a\partial^hy$; $\partial^hf_*=f_*\partial^h$;
$H_f(c_ay)=c_a\big(f_*(c_ay)-H_fy\big)$.
\textup{(2)} $\partial^hH_f+H_f\partial^h=f_*-\mathrm{id}$ on $P_n$.
\textup{(3)} For $n\ge0$: $\pi s=\mathrm{id}$, $\pi\lambda_h=\pi$,
$s(\pi x)=\lambda_{x_0^{-1}}x$ for basis elements $x$, $\pi\partial^h=\partial\pi$,
$\pi H_{r_k}=\mathrm{prism}_k\pi$ and $\pi(r_k)_*=(\kappa_k)_*\pi$.
\end{lemma}

\begin{proof}
(1) The face of $(a,x_0,\dots,x_n)$ omitting $a$ is $y$ and the others are $-c_a$
of the faces of $y$. The second identity is clear. In $H_f(a,x_0,\dots)$ the
$i=0$ term is $(a,f(a),f(x_0),\dots)=c_af_*(c_ay)$, and the term $i+1$ is $-c_a$
of the term $i$ of $H_fy$.
(2) Induction on $n$; for $n=-1$ both sides vanish. A basis element of $P_n$ is
$c_ay$. By (1), $\partial^hH_f(c_ay)=f_*(c_ay)-H_fy-c_a\big(\partial^hf_*(c_ay)-\partial^hH_fy\big)$
and $H_f\partial^h(c_ay)=H_fy-c_a\big(f_*(c_a\partial^hy)-H_f\partial^hy\big)$.
Since $\partial^hf_*(c_ay)=f_*y-f_*(c_a\partial^hy)$, the sum is
$f_*(c_ay)-c_a\big(f_*y-(\partial^hH_f+H_f\partial^h)y\big)$, which equals
$f_*(c_ay)-c_ay$ by the induction hypothesis.
(3) The first three identities are telescoping. With $g_i=x_{i-1}^{-1}x_i$:
omitting $x_0$ gives $[g_2|\dots]$, omitting $x_j$ ($1\le j\le n$) merges
$g_j,g_{j+1}$ with sign $(-1)^j$, omitting $x_{n+1}$ drops $g_{n+1}$ with sign
$(-1)^{n+1}$; this is $\partial\pi$. The image under $\pi$ of
$(x_0,\dots,x_i,x_ik,\dots,x_nk)$ is
$[g_1|\dots|g_i|k|k^{-1}g_{i+1}k|\dots|k^{-1}g_nk]$, and
$(x_{i-1}k)^{-1}x_ik=\kappa_k(g_i)$.
\end{proof}

\begin{lemma}\label{i5:x5d-lem-chain}
\textup{(1)} $\partial\partial=0$; $\partial f_*=f_*\partial$, $(gf)_*=g_*f_*$ and
$f_*\mathrm{prism}_k=\mathrm{prism}_{f(k)}f_*$ for homomorphisms $f,g$.
\textup{(2)} If $G$ is trivial, every cycle in $C_{m+1}(G;A)$ is a boundary.
\textup{(3)} \textup{(}prism identity\textup{)} For $z\in C_{m+1}(G;A)$,
$\partial\,\mathrm{prism}_kz+\mathrm{prism}_k\partial z=(\kappa_k)_*z-z$. In
particular, if $f\colon C\to G$ is a homomorphism, $k$ commutes with $f(C)$ and
$z$ is a cycle, then $\mathrm{prism}_k(f_*z)$ is a cycle.
\textup{(4)} \textup{(}multiplicativity\textup{)} For $k,k'\in G$ and a cycle
$z\in C_{m+1}(G;A)$,
$\mathrm{prism}_{kk'}z-\mathrm{prism}_kz-\mathrm{prism}_{k'}((\kappa_k)_*z)$ is a
boundary.
\textup{(5)} $z_e$ is a cycle in $C_{2e+2}(\Z/q;\Z/q)$.
\end{lemma}

\begin{proof}
(1) $\partial^h\partial^h=0$ (each double omission occurs twice with opposite
signs), so $\partial\partial\pi=\pi\partial^h\partial^h=0$ by
Lemma~\ref{i5:x5d-lem-homog}(3), and $\pi$ is onto. The other identities hold on
basis elements because $f(k^{-1}gk)=f(k)^{-1}f(g)f(k)$.
(2) $C_n=A\,e_n$, $e_n=[1|\dots|1]$, and $\partial e_{n+1}=\varepsilon_ne_n$ with
$\varepsilon_n=1+\sum_{j=0}^{n}(-1)^{j+1}$, which is $0$ for $n$ even and $1$ for
$n$ odd. If $z=ae_{m+1}$ is a cycle, $\varepsilon_ma=0$, and
$\partial(ae_{m+2})=\varepsilon_{m+1}ae_{m+1}=(1-\varepsilon_m)ae_{m+1}=z$.
(3) Write $z=\pi y$ and apply $\pi$ to Lemma~\ref{i5:x5d-lem-homog}(2) for
$f=r_k$, using Lemma~\ref{i5:x5d-lem-homog}(3). If $k$ commutes with $f(C)$,
$\kappa_k\circ f=f$, and the identity applied to the cycle $f_*z$ gives
$\partial\,\mathrm{prism}_kf_*z=0$.
(4) Put $H=H_{r_k}$, $H'=H_{r_{k'}}$, $H''=H_{r_{kk'}}$, $\varrho=(r_k)_*$, so
$(r_{k'})_*\varrho=(r_{kk'})_*$. Define $D\colon P_n\to P_{n+2}$ by $D=0$ on
$P_{-1}$ and $D(c_ay)=c_a\big(H'\varrho(c_ay)+Dy\big)$. We claim
\[
(\ast)\qquad\partial^hDx-D\partial^hx=-(H''x-Hx-H'\varrho x)\qquad(x\in P_n).
\]
For $n=-1$ both sides vanish. For $x=c_ay$, Lemma~\ref{i5:x5d-lem-homog}(1),(2)
give $\partial^hDx=H'\varrho x+Dy-c_a\big((r_{kk'})_*x-\varrho x-H'\varrho\partial^hx+\partial^hDy\big)$
and $D\partial^hx=Dy-c_a\big(H'\varrho c_a\partial^hy+D\partial^hy\big)$; since
$\partial^hx=y-c_a\partial^hy$ the difference is
$H'\varrho x-c_a\big((r_{kk'})_*x-\varrho x-H'\varrho y\big)-c_a(\partial^hDy-D\partial^hy)$,
and inserting $(\ast)$ for $y$ it equals
$H'\varrho x-c_a((r_{kk'})_*x-H''y)+c_a(\varrho x-Hy)=H'\varrho x-H''x+Hx$ by
Lemma~\ref{i5:x5d-lem-homog}(1). Next, $\partial^h,H,H',\varrho$ commute with
$\lambda_h$ and $\lambda_hc_a=c_{ha}\lambda_h$, so $D\lambda_h=\lambda_hD$ by
induction; with $\pi\lambda_h=\pi$ and $s\pi x=\lambda_{x_0^{-1}}x$ this gives
$\pi D(s\pi w)=\pi Dw$ for all $w$. For a cycle $z$ let $y=s(z)$. Then
$\pi D\partial^hy=\pi Ds(\partial z)=0$, and applying $\pi$ to $(\ast)$ gives
$\partial(\pi Dy)=-\big(\mathrm{prism}_{kk'}z-\mathrm{prism}_kz-\mathrm{prism}_{k'}(\kappa_k)_*z\big)$.
(5) Write $z_e=\sum_j[\Gamma(j)]$. The first faces do not depend on $j_0$, so they
sum to $q$ times a chain, i.e.\ $0$. For $t=0,\dots,e$ let $A_t(j)$ be the face
merging positions $2t+1,2t+2$ (sign $-1$; it replaces $j_t,1$ by $j_t+1$) and
$B_t(j)$ the face with sign $+1$ at position $2t+2$ (for $t<e$ it replaces
$1,j_{t+1}$ by $1+j_{t+1}$; for $t=e$ it drops the last entry). With $\delta_t$
the $t$-th unit vector and $\delta_{e+1}=0$, $B_t(j+\delta_t-\delta_{t+1})=A_t(j)$,
and $j\mapsto j+\delta_t-\delta_{t+1}$ is a bijection of $(\Z/q)^{e+1}$; so the
remaining faces cancel.
\end{proof}

\medskip\noindent\emph{Transfers.}
Let $X$ be a finite set with a right $G$-action, $M$ an abelian group and
$c\colon G\times X\to M$ a cocycle: $c(gh,x)=c(g,x)c(h,x\cdot g)$. Put
$T_c[g_1|\dots|g_m]=\sum_{x\in X}[c(g_1,x)|c(g_2,x\cdot g_1)|\dots|c(g_m,x\cdot g_1\cdots g_{m-1})]$.

\begin{lemma}\label{i5:x5d-lem-tr}
\textup{(1)} $\partial T_c=T_c\partial$.
\textup{(2)} If $f\colon G\to H$ is a homomorphism, $Y$ a finite right $H$-set, $c'$
a cocycle on $H\times Y$ and $p\colon X\to Y$ satisfies $p(x\cdot g)=p(x)\cdot f(g)$,
$c(g,x)=c'(f(g),p(x))$ and $\#p^{-1}(y)=1$ in $A$ for all $y$, then $T_c=T_{c'}f_*$.
\textup{(3)} For cocycles $c,c'$ on $X,X'$, $T_{c\sqcup c'}=T_c+T_{c'}$.
\textup{(4)} If $b\colon X\to M$ and $c'(g,x)=b(x)^{-1}c(g,x)b(x\cdot g)$, then
$T_{c'}z-T_cz$ is a boundary for every cycle $z$.
\end{lemma}

\begin{proof}
(1) With $Q_i=g_1\cdots g_i$ and $t_i(x)=c(g_i,x\cdot Q_{i-1})$, the cocycle identity
gives $c(g_ig_{i+1},x\cdot Q_{i-1})=t_i(x)t_{i+1}(x)$, so merges and the last
face commute with $T_c$; the first face of the $x$-summand is the $x\cdot g_1$
summand of $T_c[g_2|\dots]$, and $x\mapsto x\cdot g_1$ is a bijection.
(2) The $x$-summand of $T_c$ equals the $p(x)$-summand of $T_{c'}f_*$; group by
fibres. (3) is clear.
(4) Let $W=M^X\rtimes G$, $(n,g)(n',g')=(n\cdot{}^gn',gg')$,
${}^gn'(x)=n'(x\cdot g)$; it acts on $X$ through $G$ and $c_W((n,g),x)=n(x)$ is a
cocycle. For a cocycle $d$, $\rho_d(g)=(d(g,\cdot),g)$ is a homomorphism and
$T_d=T_{c_W}(\rho_d)_*$ by (2). For $\beta=(b,1)$,
$\beta^{-1}\rho_c(g)\beta=\rho_{c'}(g)$, so $\rho_{c'}=\kappa_\beta\rho_c$. With
$w=(\rho_c)_*z$, Lemma~\ref{i5:x5d-lem-chain}(3) gives
$T_{c'}z-T_cz=T_{c_W}(\partial\,\mathrm{prism}_\beta w)=\partial T_{c_W}\mathrm{prism}_\beta w$.
\end{proof}

For a finite $k$-space $V$, $\lambda\colon V\smallsetminus0\to k^\times$ is
\emph{homogeneous} if $\lambda(uv)=u\lambda(v)$ (choose $v_\ell$ on each line $\ell$
and put $\lambda(uv_\ell)=u$). For $\rho\colon G\to\mathrm{GL}(V)$, $G$ acts on
the right on the set $\mathbb P(V)$ of lines by $[v]\cdot g=[\rho(g)^{-1}v]$, and
$c_{\rho,\lambda}(g,[v])=\lambda(v)/\lambda(\rho(g)^{-1}v)$ is a well-defined
cocycle (the quotients telescope). For homogeneous $\lambda,\lambda'$,
$c_{\rho,\lambda'}$ arises from $c_{\rho,\lambda}$ as in
Lemma~\ref{i5:x5d-lem-tr}(4) with $b([v])=\lambda(v)/\lambda'(v)$.

\begin{lemma}[transfer to $k^\times$]\label{i5:x5d-lem-trans}
Let $k$ be a finite field, $\#k=1$ in $A$, and $t\colon C_{m+1}(k^\times;A)\to A$
linear and vanishing on boundaries. There is an additive family $\theta$ of
degree $m+1$ with values in $A$ such that
$\theta_k((\mathrm{scal}_k)_*z)=t(z)$ for every cycle $z$.
\end{lemma}

\begin{proof}
Fix homogeneous $\lambda_V$ for each $V$ and put
$\theta_V=t\circ T_{c_{\mathrm{id},\lambda_V}}$. \emph{Key formula:}
$\theta_V(\rho_*z)=t(T_{c_{\rho,\lambda}}z)$ for every $\rho\colon G\to\mathrm{GL}(V)$,
homogeneous $\lambda$ and cycle $z$; indeed
$T_{c_{\rho,\lambda_V}}=T_{c_{\mathrm{id},\lambda_V}}\rho_*$ by
Lemma~\ref{i5:x5d-lem-tr}(2) ($p=\mathrm{id}$) and $t$ kills
$T_{c_{\rho,\lambda}}z-T_{c_{\rho,\lambda_V}}z$ by Lemma~\ref{i5:x5d-lem-tr}(4).
(a) $\theta_V(\partial y)=t(\partial Ty)=0$.
(b) For $\varphi\colon V\to W$ use the key formula with $\rho=\iota_\varphi$,
$\lambda=\lambda_V\circ\varphi^{-1}$, then Lemma~\ref{i5:x5d-lem-tr}(2) with
$f=\mathrm{id}$ and the bijection $[w]\mapsto[\varphi^{-1}w]$:
$\theta_W((\iota_\varphi)_*z)=t(T_{c_{\mathrm{id},\lambda_V}}z)=\theta_V(z)$.
(c) On $V\oplus W$ take $\lambda(v,w)=\lambda_V(v)$ for $v\ne0$ and $\lambda_W(w)$
for $v=0$. The map $\mathbb P(V\oplus W)\to\mathbb P(V)\sqcup\mathbb P(W)$,
$[v,w]\mapsto[v]$ ($v\ne0$), $[0,w]\mapsto[w]$, is equivariant, carries
$c_{\rho\oplus\sigma,\lambda}$ to $c_{\rho,\lambda_V}\sqcup c_{\sigma,\lambda_W}$, and its
fibres have $\#W=1$ (in $A$) or one element; the key formula and
Lemma~\ref{i5:x5d-lem-tr}(2),(3) give (c).
Finally, for $V=k$, $\rho=\mathrm{scal}_k$, $\lambda(v)=v$: $\mathbb P(k)$ is a point
and $c(u,\cdot)=u$, so $T_c=\mathrm{id}$ and
$\theta_k((\mathrm{scal}_k)_*z)=t(z)$.
\end{proof}

\begin{lemma}[the cocycle on $k^\times$]\label{i5:x5d-lem-coc}
Let $k$ be finite, $q\ge1$, $\zeta\in k^\times$ of order $q$, $v\colon k^\times\to\Z/q$
a homomorphism, $e\ge0$, and $\chi_x$ as in Definition~\ref{i5:x5d-def-val}. There
is a linear $t\colon C_{2e+3}(k^\times;\Z/q)\to\Z/q$ vanishing on boundaries with
$t(\mathrm{prism}_g(\chi_x)_*z_e)=x^{e+1}v(g)$ for all $x\in\N$, $g\in k^\times$.
\end{lemma}

\begin{proof}
For $q=1$ take $t=0$. Let $q\ge2$, $N=\#k^\times$, $\omega$ a generator of $k^\times$
and $y(a)\in\Z$ with $\omega^{y(a)}=a$. Since $\zeta$ has order $q$,
$r:=\gcd(N,y(\zeta))$ satisfies $qr=N$. Write $r=\alpha N+\beta y(\zeta)$ and
$\ell(a)=\beta y(a)$. Then $N\mid\ell(a)+\ell(b)-\ell(ab)$ and
$\ell(\zeta^x)\equiv rx\pmod N$. Put $U(a,b)=(\ell(a)+\ell(b)-\ell(ab))/N\in\Z$; it is
symmetric and $U(b,c)-U(ab,c)+U(a,bc)-U(a,b)=0$. Let $u=U\bmod q$, define
$F_0=1$, $F_1(g_1)=v(g_1)$,
$F_{n+2}(g_1,\dots,g_{n+2})=u(g_1,g_2)F_n(g_3,\dots,g_{n+2})$, and let $t$ be the
linear extension of $F_{2e+3}$.
\emph{$F_n\circ\partial=0$.} For $n=0$: $1-1$; for $n=1$:
$v(g_2)-v(g_1g_2)+v(g_1)=0$. For $n+2$: the first face and the merges at positions
$1,2$ contribute $F_n(g_4,\dots)\big(u(g_2,g_3)-u(g_1g_2,g_3)+u(g_1,g_2g_3)\big)$;
the remaining faces contribute $u(g_1,g_2)$ times the non-first faces of
$F_n\circ\partial$ at $(g_3,\dots)$, i.e.\ $-u(g_1,g_2)F_n(g_4,\dots)$ by induction;
the sum vanishes by the cocycle identity.
\emph{Insertion.} For $h\in(k^\times)^{2E}$,
$\sum_{i=0}^{2E}(-1)^iF_{2E+1}(h_1,\dots,h_i,g,h_{i+1},\dots)=F_{2E}(h)v(g)$: for $E=0$
this is $v(g)$; for $h=(a,b,h')$ the terms $i=0,1$ cancel by symmetry of $u$ and
the others give $u(a,b)F_{2E-2}(h')v(g)$ by induction. As $k^\times$ is abelian,
$t(\mathrm{prism}_g[h])=F_{2e+2}(h)v(g)$.
\emph{Evaluation.} With $\chi=\chi_x$,
$t(\mathrm{prism}_g(\chi_x)_*z_e)=\big(\sum_{j\in\Z/q}u(\chi(j),\chi(1))\big)^{e+1}v(g)$,
and by telescoping over $j\mapsto j+1$,
$N\sum_jU(\chi(j),\chi(1))=q\,\ell(\zeta^x)=q\,r(x+qm)$ for some $m\in\Z$; dividing by
$N=qr$ the sum is $\equiv x\pmod q$.
\end{proof}

\begin{lemma}\label{i5:x5d-lem-det}
Let $k$ be a finite field with an element $a\ne0,1$, $M$ an abelian group,
$w\colon k^\times\to M$ a homomorphism and $F_V\colon\mathrm{GL}(V)\to M$
homomorphisms with $F_W(\varphi g\varphi^{-1})=F_V(g)$,
$F_{V\oplus W}(g\oplus h)=F_V(g)+F_W(h)$ and $F_k(\mathrm{scal}_k(u))=w(u)$. Then
$F_V(g)=w(\det g)$.
\end{lemma}

\begin{proof}
By induction via $k^{n+1}\cong k\oplus k^n$,
$F_{k^n}(\mathrm{diag}(d_i))=\sum_iw(d_i)$. For $i\ne j$ let $T(x)=1+xE_{ij}$ and $D$
diagonal with $a$ at $i$, $1$ elsewhere: $DT(x)D^{-1}=T(ax)$ and
$T(x)T(y)=T(x+y)$, so $F(T(ax))=F(T(x))$; with $x=c/(a-1)$, $T(ax)=T(x)T(c)$ and
$F(T(c))=0$. Every invertible matrix over a field is a product of an invertible
diagonal matrix and transvections $T(c)$ (Gaussian elimination), and
$\det T(c)=1$. A general $V$ is isomorphic to some $k^n$.
\end{proof}

\medskip\noindent\emph{Proof of \textup{(i-a)}.}
Since $\zeta$ has order $q$, $q\mid\#k-1$, so $\#k=1$ in $\Z/q$. Let $t$ be as
in Lemma~\ref{i5:x5d-lem-coc} and $\theta$ as in Lemma~\ref{i5:x5d-lem-trans}
($A=\Z/q$, $m+1=2e+3$). If $q=1$ there is nothing to prove; let $q\ge2$, so
$\zeta\ne0,1$. Fix $x$, put $w_V=(\mathrm{scal}_V\circ\chi_x)_*z_e$ (a cycle,
Lemma~\ref{i5:x5d-lem-chain}(1),(5)) and $F_V(g)=\theta_V(\mathrm{prism}_gw_V)$.
Scalars commute with $g$, so $\mathrm{prism}_gw_V$ is a cycle and
$(\kappa_g)_*w_V=w_V$ (Lemma~\ref{i5:x5d-lem-chain}(3)). $F_V$ is a homomorphism by
Lemma~\ref{i5:x5d-lem-chain}(4) and (a). Since $\iota_\varphi\circ\mathrm{scal}_V=\mathrm{scal}_W$,
$(\iota_\varphi)_*\mathrm{prism}_gw_V=\mathrm{prism}_{\varphi g\varphi^{-1}}w_W$, and (b)
gives conjugation invariance. For direct sums let
$f=(\mathrm{scal}_V\chi_x,\mathrm{scal}_W\chi_x)\colon\Z/q\to\mathrm{GL}(V)\times\mathrm{GL}(W)$
and apply (c) with the two projections to the cycle
$\mathrm{prism}_{(g,h)}f_*z_e$: $F_{V\oplus W}(g\oplus h)=F_V(g)+F_W(h)$. For
$u\in k^\times$,
$\mathrm{prism}_{\mathrm{scal}_k(u)}w_k=(\mathrm{scal}_k)_*\mathrm{prism}_u(\chi_x)_*z_e$
with $\mathrm{prism}_u(\chi_x)_*z_e$ a cycle, so
$F_k(\mathrm{scal}_k(u))=t(\mathrm{prism}_u(\chi_x)_*z_e)=x^{e+1}v(u)$.
Lemma~\ref{i5:x5d-lem-det} with $w(u)=x^{e+1}v(u)$ gives
$F_V(g)=x^{e+1}v(\det g)$. \qed

\medskip\noindent\emph{Proof of \textup{(i-b)}.}
\begin{lemma}[trivial summands]\label{i5:x5d-lem-triv}
Let $\theta$ be additive of degree $m+1$, $z\in C_{m+1}(H;A)$ a cycle, $\mu$ a
representation of $H$ on $V$, $P$ a finite $k$-space and
$\varphi\colon V\oplus P\to U$ an isomorphism. Then
$\theta_U((\iota_\varphi\circ(\mu\oplus1_P))_*z)=\theta_V(\mu_*z)$.
\end{lemma}

\begin{proof}
The trivial representation $1_P$ factors through the trivial group, where the
image of $z$ is a boundary (Lemma~\ref{i5:x5d-lem-chain}(2)); so $\theta_P(1_*z)=0$.
Conclude with (b) and (c).
\end{proof}

Boundaries go to boundaries, so $\Phi_n$ kills them. If $\varphi\colon W\to V$
intertwines $\sigma$ with $\rho$, then $\rho=\iota_\varphi\circ\sigma$ and (b)
gives $\theta_W(\sigma_*z)=\theta_V(\rho_*z)$ on cycles. Under
$k^{n'}\cong k^n\oplus k^{n'-n}$, $\mathrm{red}_{r,n'}\circ\mathrm{st}$ becomes
$\mathrm{red}_{r,n}\oplus1$, so Lemma~\ref{i5:x5d-lem-triv} gives
$\Phi_{n'}(\mathrm{st}_*z)=\Phi_n(z)$. Under $k^{dn}\cong(k^n)^d$,
$\mathrm{red}_{r,dn}\circ\psi_d$ becomes the $d$-fold sum $\rho^{\oplus d}$ of
$\rho=\mathrm{red}_{r,n}$; by induction on $d$ (for $d=0$,
Lemma~\ref{i5:x5d-lem-triv}; then (c) and $(k^n)^{d+1}\cong k^n\oplus(k^n)^d$),
$\theta_{(k^n)^d}(\rho^{\oplus d}_*z)=d\,\theta_{k^n}(\rho_*z)$. \qed

\medskip\noindent\emph{Proof of \textup{(ii)}: algebraic lemmas.}
Write $\mathrm{ng}(M)$ for the least number of generators of an abelian group $M$
($\infty$ if none is finite). Group homology is computed by the complex
$C_*(G;A)$, and the long exact homology sequences and Shapiro's lemma below
are applied to the homology of this bar complex (e.g.\ Brown,
\emph{Cohomology of Groups}, Ch.~III; to be confirmed against the source).

\begin{lemma}\label{i5:x5d-lem-gen}
\textup{(1)} If $M_1\xrightarrow fM_2\xrightarrow gM_3$ is exact,
$\mathrm{ng}(M_2)\le\mathrm{ng}(M_1)+\mathrm{ng}(M_3)$.
\textup{(2)} If $\mathrm{ng}(M)\le r$, $l\ge2$ and $x_0,\dots,x_r\in M$, there are
$n_t\in\Z$, not all divisible by $l$, with $\sum n_tx_t=0$.
\textup{(3)} Let $G$ be a group, $p\colon G\to\Z/q$ surjective, $K=\ker p$, coefficients
a trivial module $\Z/r$. Then $\mathrm{ng}H_1(G)\le\mathrm{ng}H_1(K)+\mathrm{ng}H_0(K)$
and $\mathrm{ng}H_{n+2}(G)\le\mathrm{ng}H_{n+2}(K)+\mathrm{ng}H_{n+1}(K)+\mathrm{ng}H_n(G)$.
\textup{(4)} $\mathrm{ng}H_{2e}(\Z/q\times\Z/q;\Z/r)\le2e+1$ for $e\ge0$.
\end{lemma}

\begin{proof}
(1) $g(M_2)$ is a subgroup of a quotient of $\Z^{\mathrm{ng}(M_3)}$, hence generated
by $\le\mathrm{ng}(M_3)$ elements (subgroups of $\Z^s$ are free of rank $\le s$);
lifts of these together with $f$ of generators of $M_1$ generate $M_2$.
(2) Lift the $x_t$ to $\Z^r$; $r+1$ vectors there satisfy a nontrivial integer
relation; divide by $l$ as long as all coefficients are divisible by $l$ (this
terminates); map the relation to $M$.
(3) Let $W$ be the functions $\Z/q\to\Z/r$ with $(gF)(j)=F(j+p(g))$. The sequence
$0\to\Z/r\to W\xrightarrow TW\xrightarrow s\Z/r\to0$ (constants;
$(TF)(j)=F(j+1)-F(j)$; $s(F)=\sum_jF(j)$) is exact ($\ker T$ = constants; if
$\sum F=0$, $E(j)=\sum_{i<j}F(i)$ has $TE=F$). As $p$ is onto,
$W\cong\mathrm{Ind}_K^G\Z/r$, so $H_n(G;W)\cong H_n(K;\Z/r)$ (Shapiro's lemma, e.g.\ Brown,
\emph{Cohomology of Groups}, III.6; to be confirmed against the source). Split the sequence at $Z=\ker s$; the long exact sequences
give exact $H_{n+2}(G;W)\to H_{n+2}(G)\to H_{n+1}(G;Z)$,
$H_{n+1}(G;W)\to H_{n+1}(G;Z)\to H_n(G)$, $H_1(G;W)\to H_1(G)\to H_0(G;Z)$ with
$H_0(G;W)\to H_0(G;Z)$ onto; apply (1).
(4) $H_0=\Z/r$. With $p$ the first projection, $K\cong\Z/q$; applying (3) to $K$
itself with $p$ bijective ($H_{\ge1}$ of the trivial group vanish) gives
$\mathrm{ng}H_n(\Z/q)\le1$; then (3) gives
$\mathrm{ng}H_{2e+2}\le1+1+\mathrm{ng}H_{2e}$.
\end{proof}

For $t\in\N$ let $d_t\colon\Z/q\to\Z/q\times\Z/q$, $x\mapsto(x,tx)$.

\begin{lemma}[relation lemma]\label{i5:x5d-lem-rel}
Let $l$ be prime, $q=l^\nu$, $e\ge0$. There are $n_0,\dots,n_{2e+3}\in\Z/q$ and
$t_0$ with $n_{t_0}$ a unit such that $\sum_tn_t(d_t)_*z_e$ is a boundary in
$C_{2e+2}(\Z/q\times\Z/q;\Z/q)$.
\end{lemma}

\begin{proof}
The $(d_t)_*z_e$ are cycles (Lemma~\ref{i5:x5d-lem-chain}(1),(5)).
$H_{2e+2}(\Z/q\times\Z/q;\Z/q)$ needs $\le2e+3$ generators
(Lemma~\ref{i5:x5d-lem-gen}(4)), so Lemma~\ref{i5:x5d-lem-gen}(2) gives integers
$n_t$, some $n_{t_0}$ prime to $l$, with $\sum n_t[(d_t)_*z_e]=0$ in
$H_{2e+2}$ computed by $C_*$, i.e.\ $\sum n_t(d_t)_*z_e$ is a boundary. Reduce mod $q$.
\end{proof}

For a finite abelian group $X$ let $\Z[X]$ be its group ring (basis $[x]$) and
$\pi_t\colon\Z[(\Z/q)^2]\to\Z[\Z/q]$ the ring map induced by $(x,y)\mapsto x+ty$.

\begin{lemma}[eta lemma]\label{i5:x5d-lem-eta}
Let $l$ be prime, $c,\nu,M\in\N$, $M<l^c$, $q=l^\nu$, $\kappa_0=1+l^c$,
$\prod_{j<M}(X-\kappa_0^j)=\sum_{i=0}^Mp_iX^i$. For each $t_0\in\{0,\dots,M\}$ there
is $\eta\in\Z[(\Z/q)^2]$ with $\pi_t(\eta)=0$ for $t\in\{0,\dots,M\}\smallsetminus\{t_0\}$
and $\pi_{t_0}(\eta)=\sum_ip_i[\kappa_0^i]$.
\end{lemma}

\begin{proof}
\emph{Step 1.} For $w\in\N$ and $\kappa=1+w$, $(X^w-1)^M$ divides
$\sum_ip_iX^{\kappa^i}$ in $\Z[X]$ (with $p_i$ formed with $\kappa$). Let
$\mathcal E(f)(X)=f(X^{\kappa})$, a ring endomorphism; then
$\sum_ip_iX^{\kappa^i}=\big(\prod_{j<M}(\mathcal E-\kappa^j)\big)(X)$, and it suffices
that $(X^w-1)^j\mid x$ implies $(X^w-1)^{j+1}\mid\mathcal E(x)-\kappa^jx$. Write
$x=(X^w-1)^jy$ and $\mathcal E(X^w-1)=(X^w-1)\Sigma$, $\Sigma=\sum_{i\le w}X^{wi}$;
then $\mathcal E(x)-\kappa^jx=(X^w-1)^j[(\Sigma^j-\kappa^j)\mathcal E(y)+\kappa^j(\mathcal E(y)-y)]$,
and $X^w-1$ divides $\Sigma-\kappa=\sum_{i\le w}(X^{wi}-1)$ and
$\mathcal E(y)-y$ (as $X^{n\kappa}-X^n=X^n(X^{nw}-1)$).
\emph{Step 2.} Under $X\mapsto[1]$, $w=l^c$: $([l^c]-1)^M$ divides
$R:=\sum_ip_i[\kappa_0^i]$ in $\Z[\Z/q]$.
\emph{Step 3.} For $t\ne t_0$ in $\{0,\dots,M\}$, $0<|t_0-t|<l^c$, so $l^c$ is a
multiple of $t_0-t$ in $\Z/q$; and $[a]-1\mid[ac']-1$. There are $M$ such $t$, so
$R=\prod_{t\ne t_0}([t_0-t]-1)\cdot Q$ for some $Q$.
\emph{Step 4.} With $\iota(x)=(x,0)$ and $\Sigma_0=\sum_{t\ne t_0}t$, put
$\eta=\prod_{t\ne t_0}([(0,1)]-[(t,0)])\cdot\iota([-\Sigma_0]Q)$. Since
$\pi_{t'}([(0,1)]-[(t,0)])=[t']-[t]$, $\pi_{t'}(\eta)=0$ for $t'\ne t_0$, and
$\pi_{t_0}(\eta)=\prod_{t\ne t_0}[t]([t_0-t]-1)[-\Sigma_0]Q=R$.
\end{proof}

\begin{lemma}\label{i5:x5d-lem-ord}
If $l$ is prime, $1\le c\le\nu$ \textup{(}$c\ge2$ if $l=2$\textup{)}, then $1+l^c$ has
order $l^{\nu-c}$ in $(\Z/l^\nu)^\times$.
\end{lemma}

\begin{proof}
If $y=1+l^mu$ with $m\ge c$, $l\nmid u$, the binomial expansion gives
$y^l=1+l^{m+1}u'$ with $u'\equiv u\pmod l$: the middle terms are divisible by
$l^{2m+1}$ and the last by $l^{lm}$, and $lm\ge m+2$ under the hypotheses. Hence
$(1+l^c)^{l^s}=1+l^{c+s}u_s$ with $l\nmid u_s$.
\end{proof}

\medskip\noindent\emph{Lattices.}
Let $F$ be a number field. For a divisor $\mathfrak a$, $I(\mathfrak a)=\{x:\operatorname{ord}_Px\ge\mathfrak a(P)\ \forall P\}$;
for a finite set $J$ and divisors $(\mathfrak a_j)$,
$\Lambda(\mathfrak a)=\bigoplus_jI(\mathfrak a_j)\subset F^J$.

\begin{lemma}\label{i5:x5d-lem-approx}
Given finitely many primes $P\in S$ and $n_P\in\Z$, there is $x\in F^\times$ with
$\operatorname{ord}_Px=n_P$ for $P\in S$ and $\operatorname{ord}_Px\ge0$ otherwise.
\end{lemma}

\begin{proof}
Let $s=\max(\{0\}\cup\{-n_P:P\in S\})$ and $0\ne d\in\Z\cap\prod_{P\in S}P$. Let $S'=S\cup\{P:P\mid d\}$,
$n'_P=n_P+s\operatorname{ord}_Pd\ge0$ for $P\in S$ and $n'_P=s\operatorname{ord}_Pd$ for
$P\in S'\smallsetminus S$. By the Chinese remainder theorem (Neukirch I (3.6)) choose
$y\in\mathcal O_F$ with $y\equiv y_P\bmod P^{n'_P+1}$, $y_P\in P^{n'_P}\smallsetminus P^{n'_P+1}$;
then $x=y/d^s$ works.
\end{proof}

\begin{lemma}[Steinitz-type lattice change]\label{i5:x5d-lem-stein}
Let $\mathfrak a,\mathfrak b$ be $J$-families of divisors and $\beta\in F^\times$ with
$\operatorname{div}\beta=\sum_j(\mathfrak b_j-\mathfrak a_j)$, and $J\ne\emptyset$ or
$\beta=1$. There is $g\in\mathrm{GL}_J(F)$ with $\det g=\beta$ and
$g(\Lambda(\mathfrak a))=\Lambda(\mathfrak b)$.
\end{lemma}

\begin{proof}
(a) \emph{For divisors $d,e$ there are $p,q',r,s\in F$, $p\ne0$, $ps-q'r=1$,
$p\in I(d)$, $q'\in I(e)$, $r\in I(-e)$, $s\in I(-d)$.} Let $C=I(d)$, $D=I(e)$,
$0\ne p\in C\cap\mathcal O_F$, $\mathfrak i=pC^{-1}$ (integral). By
Lemma~\ref{i5:x5d-lem-approx} choose $q'$ with $\operatorname{ord}_Pq'=\operatorname{ord}_PD$ for
$P\mid\mathfrak i$ or $P$ in the support of $D$, and $\ge0$ elsewhere. Then
$q'D^{-1}$ is integral and prime to $\mathfrak i$, so
$pC^{-1}+q'D^{-1}=\mathcal O_F$: $1=ps+q'r'$ with $s\in C^{-1}$, $r'\in D^{-1}$; put
$r=-r'$.
(b) For $i\ne j$ and a divisor $d$, apply (a) to $(d,\mathfrak a_i+d-\mathfrak a_j)$
and let $g$ fix the other coordinates and send
$(x_i,x_j)\mapsto(px_i+q'x_j,rx_i+sx_j)$; then $\det g=1$ and
$g\Lambda(\mathfrak a)=\Lambda(\mathfrak a+de_i-de_j)$, by the valuation bounds for
$g$ and for $g^{-1}\colon(y_i,y_j)\mapsto(sy_i-q'y_j,-ry_i+py_j)$.
(c) $\mathrm{diag}(\beta\text{ at }j_0,1,\dots)$ maps $\Lambda(\mathfrak a)$ onto
$\Lambda(\mathfrak a+\operatorname{div}(\beta)e_{j_0})$.
(d) If $J=\emptyset$ take $g=1$. Otherwise apply (c), then (b) with
$(i,j_0,\mathfrak b_i-\mathfrak a'_i)$ for each $i\ne j_0$; the total is preserved, so
coordinate $j_0$ ends at $\mathfrak b_{j_0}$; the determinant is $\beta$.
\end{proof}

\medskip\noindent\emph{Setting.}
From now on $L\subset F$, $F/L$ Galois with group $\Gamma$, $\iota\colon\mathcal O_L\to\mathcal O_F$,
$q=l^\nu$, $\zeta\in\mathcal O_F^\times$, $\zeta^q=1$. Then
$\operatorname{ord}_{\sigma P}(\sigma x)=\operatorname{ord}_P(x)$, and $b\in\sigma P\Leftrightarrow b\in P$
for $b\in\mathcal O_L$. For $0\ne b\in\mathcal O_L$ let
$\mathcal O_{F,b}=\{x:\operatorname{ord}_Px\ge0\ \forall P\not\ni b\}$ ($\Gamma$-stable),
$\alpha$ a \emph{$b$-unit} if $\operatorname{ord}_P\alpha=0$ for all $P\not\ni b$,
$\mathcal G_J=\mathrm{GL}_L(F^J)$, $\mathcal G_{J,b}=\{g:g\mathcal O_{F,b}^J=\mathcal O_{F,b}^J\}$,
$\mathcal H_{J,b}=\{g\in\mathcal G_{J,b}:g\mathcal O_F^J=\mathcal O_F^J\}$.
Fix $r\colon\mathcal O_F\to k$ to a field; ``reduces to'' is as in
Definition~\ref{i5:x5d-def-val}(4), extended to $x\in F$, $u\in k$.
Reductions are unique ($a/c=a'/c'$ gives $r(a)r(c')=r(a')r(c)$), $a\in\mathcal O_F$
reduces to $r(a)$, and reductions are additive and multiplicative.
$r$ \emph{splits} if (S1) for every $w\colon\Gamma\to k$ there is
$a\in\mathcal O_F$ with $r(\sigma a)=w(\sigma)$ for all $\sigma$, and (S2) every
$a\in\mathcal O_F$ with $r(\sigma a)=0$ for all $\sigma$ lies in
$\mathfrak p\mathcal O_F$, $\mathfrak p=\ker(r\iota)$.

\begin{lemma}\label{i5:x5d-lem-res}
Let $0\ne b\in\mathcal O_L$, $r(\iota b)\ne0$.
\textup{(1)} Each $x\in\mathcal O_{F,b}$ has $b^nx\in\mathcal O_F$ for some $n$, so it
reduces; each $b$-unit reduces to an element of $k^\times$.
\textup{(2)} If $\mathfrak a,\mathfrak b$ are supported on primes containing $b$ and
$g\in\mathcal G_J$ maps $\Lambda(\mathfrak a)$ onto $\Lambda(\mathfrak b)$, then
$g\in\mathcal G_{J,b}$.
\textup{(3)} If $g\in\mathrm{GL}_F(F^I)$ maps $\mathcal O_{F,b}^I$ onto itself, then $\det g$ is a
$b$-unit, and for every $\sigma\in\Gamma$ there is $\bar g\in\mathrm{GL}_I(k)$ such
that $\sigma(\det g)$ reduces to $\det\bar g$ and, whenever $\sigma(w_i)$ reduces to
$u_i$ for all $i$, $\sigma((gw)_i)$ reduces to $(\bar gu)_i$.
\textup{(4)} Given $v\colon k^\times\to\Z/q$, let $U$ be the group of $b$-units and
$f(x)=v(u)$ if $x$ reduces to $u\in k^\times$ \textup{(}$f=0$ otherwise\textup{)}.
Then $U$ is $\Gamma$-stable, $f(xy)=f(x)+f(y)$ on $U$, and $f(1)=0$.
\end{lemma}

\begin{proof}
(1) Take $n=\max(0,-\operatorname{ord}_Px)$ over the finitely many $P\ni b$. For a $b$-unit,
$\alpha$ and $\alpha^{-1}$ reduce and the product of the reductions is $1$.
(2) For $x\in\mathcal O_{F,b}^J$ some $b^nx\in\Lambda(\mathfrak a)$; then
$b^ngx=g(b^nx)\in\Lambda(\mathfrak b)$, and $\mathfrak b_j(P)=0=\operatorname{ord}_Pb$ for
$P\not\ni b$; so $gx\in\mathcal O_{F,b}^J$; likewise for $g^{-1}$.
(3) The matrices $A,B$ of $g,g^{-1}$ have entries in $\mathcal O_{F,b}$ and $AB=1$,
so $\det A$ is a unit of $\mathcal O_{F,b}$. The map $\rho_\sigma\colon\mathcal O_{F,b}\to k$,
$x\mapsto$ reduction of $\sigma x$, is a ring map by (1); take $\bar g=\rho_\sigma(A)$.
(4) $\operatorname{ord}_P$ is additive and $\operatorname{ord}_P(\sigma\alpha)=\operatorname{ord}_{\sigma^{-1}P}\alpha$;
additivity of $f$ follows from (1) and multiplicativity of reductions.
\end{proof}

\begin{lemma}\label{i5:x5d-lem-split}
\textup{(1)} There is $0\ne b_0\in\mathcal O_L$ such that every $r$ with $r\iota$
surjective and $r(\iota b_0)\ne0$ splits.
\textup{(2)} For every divisor $\mathfrak d$ and $0\ne b_0\in\mathcal O_L$ there is
$0\ne b\in\mathcal O_L$, $b_0\mid b$, with $\iota b\in P$ whenever $\mathfrak d(P)\ne0$.
\textup{(3)} If $r$ splits, $r\iota$ is onto and $r(\iota b)\ne0$, then
$\varrho\colon\mathcal O_{F,b}^J\to V:=k^{\Gamma\times J}$,
$\varrho(x)(\sigma,j)=$ reduction of $\sigma(x_j)$, is surjective, and there is a
unique homomorphism $\mu\colon\mathcal G_{J,b}\to\mathrm{GL}(V)$ with
$\varrho(gx)=\mu(g)\varrho(x)$.
\textup{(4)} In \textup{(3)}, $\mathcal O_F^J$ is a finitely generated projective
$\mathcal O_L$-module with $\mathcal O_L$-linear $\mathcal H_{J,b}$-action, and
$\varrho|_{\mathcal O_F^J}$ is $r\iota$-semilinear, surjective, with kernel in
$\mathfrak p\mathcal O_F^J$ and $\varrho(hx)=\mu(h)\varrho(x)$.
\end{lemma}

\begin{proof}
(1) Take an $L$-basis $(B_i)$ of $F$, its trace-dual basis $(B_i^\vee)$ and
$0\ne m\in\Z$ with $e_i=mB_i^\vee$, $f_i=mB_i\in\mathcal O_F$; $b_0=m^2$. Then
$b_0x=\sum_i\mathrm{Tr}_{F/L}(xf_i)e_i$; applying $\sigma$ and Dedekind's
independence of characters (reference to be confirmed) gives
$\sum_i\tau(f_i)\sigma(e_i)=\delta_{\tau\sigma}b_0$. (S1): choose $s_i\in\mathcal O_L$ with
$r\iota(s_i)=\sum_\tau r(b_0)^{-1}w(\tau)r(\tau f_i)$; then $a=\sum s_ie_i$ works.
(S2): $b_0a=\sum c_ie_i$ with $c_i=\mathrm{Tr}(af_i)\in\mathfrak p$; with
$b'\in\mathcal O_L$, $r(b')=r(b_0)^{-1}$, $a=(1-b'b_0)a+b'\sum c_ie_i\in\mathfrak p\mathcal O_F$.
(2) $b=b_0\prod_{P\in\mathrm{supp}\,\mathfrak d}\mathrm N(P)$.
(3) Surjectivity is (S1) coordinatewise. If $\varrho(y)=0$, write $y_j=b^{-n}a_j$;
then $r(\sigma a_j)=0$ for all $\sigma$, so $a_j\in\mathfrak p\mathcal O_F$ by (S2), and
$y=\sum_sm_sw_s$ with $m_s\in\mathfrak p$, $w_s\in\mathcal O_{F,b}^J$. Then
$\varrho(gy)=\sum r(m_s)\varrho(gw_s)=0$. Hence $\mu(g)\varrho(x):=\varrho(gx)$ is well
defined, $k$-linear ($r\iota$ onto), multiplicative, invertible, unique.
(4) $\mathcal O_F^J$ is finitely generated and torsion free over the Dedekind ring
$\mathcal O_L$, hence projective (reference to be confirmed). The rest follows from (3),
(S1), (S2).
\end{proof}

\begin{lemma}[free model]\label{i5:x5d-lem-free}
Let $R$ be a commutative ring, $M$ a finitely generated projective $R$-module with
an action $\alpha$ of a group $\mathcal H$. There are $N$ and
$\lambda\colon\mathcal H\to\mathrm{GL}_N(R)$ such that for every surjection
$r\colon R\to k$ onto a field, every representation $\mu$ of $\mathcal H$ on a
$k$-space $V$ and every surjective $r$-semilinear $\pi\colon M\to V$ with
$\ker\pi\subseteq(\ker r)M$ and $\pi\alpha(h)=\mu(h)\pi$, there are $P\subseteq k^N$
and an isomorphism $\varphi\colon V\oplus P\to k^N$ with
$\mathrm{red}_{r,N}\circ\lambda=\iota_\varphi\circ(\mu\oplus1_P)$.
\end{lemma}

\begin{proof}
Choose $s\colon R^N\to M$ onto and $i$ with $si=1$, and put
$\lambda(h)=i\alpha(h)s+(1-is)$; using $si=1$, $\lambda$ is a homomorphism. Define
$\bar\imath(\pi x)=r(ix)$ and $\bar s(r(y))=\pi(sy)$ (well defined by the kernel
condition and semilinearity); $\bar s\bar\imath=1$. Let $P=\ker\bar s$,
$\varphi(v,p)=\bar\imath(v)+p$. Writing $r(y)=\bar\imath(w_1)+w_2$, $w_2\in P$, one
has $w_1=\pi(sy)$ and
$r(\lambda(h)y)=\bar\imath(\mu(h)w_1)+w_2=\varphi(\mu(h)w_1,w_2)$.
\end{proof}

\begin{lemma}[chain lemma]\label{i5:x5d-lem-cchain}
Let $C,D,G$ be groups, $\mathcal H\le G$, $\rho\colon D\to G$, $d_t\colon C\to D$,
$n_t\in A$, $g_t\in G$ \textup{(}$t$ in a finite set\textup{)}, $t_0$ an index and
$z\in C_{m+1}(C;A)$ a cycle. Assume $\sum_tn_t(d_t)_*z$ is a boundary, $g_t$
commutes with $\rho(d_t(C))$, $\rho(D)\subseteq\mathcal H$,
$g_{t_0}^{-1}\rho(D)g_{t_0}\subseteq\mathcal H$ and $g_{t_0}^{-1}g_t\in\mathcal H$. Then
there is a cycle $y\in C_{m+2}(\mathcal H;A)$ with
$\mathrm{incl}_*y-\sum_tn_t\mathrm{prism}_{g_t}((\rho d_t)_*z)$ a boundary in
$C_{m+2}(G;A)$.
\end{lemma}

\begin{proof}
Let $\sum_tn_t(d_t)_*z=\partial\mathbf b$, $\rho_H=\rho$, $\sigma_H=\kappa_{g_{t_0}}\circ\rho$
(both into $\mathcal H$), $k_t=g_{t_0}^{-1}g_t$; then
$\kappa_{k_t}\sigma_Hd_t=\rho_Hd_t$. Put
$y=(\sigma_H)_*\mathbf b-(\rho_H)_*\mathbf b+\sum_tn_t\mathrm{prism}_{k_t}((\sigma_Hd_t)_*z)$. By
Lemma~\ref{i5:x5d-lem-chain}(3),
$\partial\,\mathrm{prism}_{k_t}((\sigma_Hd_t)_*z)=(\rho_Hd_t)_*z-(\sigma_Hd_t)_*z$,
so $\partial y=0$. In $G$, Lemma~\ref{i5:x5d-lem-chain}(3) for $\rho_*\mathbf b$ gives
$(\kappa_{g_{t_0}}\rho)_*\mathbf b-\rho_*\mathbf b\equiv\sum_tn_t\mathrm{prism}_{g_{t_0}}w_t$,
$w_t=(\rho d_t)_*z$, and Lemma~\ref{i5:x5d-lem-chain}(4) gives
$\mathrm{prism}_{g_t}w_t\equiv\mathrm{prism}_{g_{t_0}}w_t+\mathrm{prism}_{k_t}((\kappa_{g_{t_0}})_*w_t)$,
modulo boundaries; subtracting gives the claim.
\end{proof}

\medskip\noindent\emph{Blocks.}
A \emph{block datum} on $J$ is $(s,\omega,g')$: $s\colon J\to\Gamma$, $\omega\colon J\to X$,
$g'_x\in\mathrm{GL}_F(F^{\omega^{-1}(x)})$. $g\in\mathcal G_J$ \emph{has block form}
$(s,\omega,g')$ if $s_j((gw)_j)=\big(g'_{\omega_j}((s_iw_i)_{\omega_i=\omega_j})\big)_j$. For
$\chi\colon J\to\Z/q$, divisors $\mathfrak b$ and $\delta\colon\Z/q\to F^\times$, $g$
\emph{carries} $(\chi,\mathfrak b,\delta)$ if $g(\mathcal O_F^J)=\Lambda(\mathfrak b)$ and
$g$ has a block form with $X=\Z/q$, $s_j(\zeta^{\chi_j})=\zeta^{\omega_j}$ and
$\det g'_x=\delta(x)$. Let $\mathrm{rep}_\chi(c)=\mathrm{diag}(\zeta^{\chi_jc})_j$, and on
$J=(\Z/q)^2$, $\mathrm{rep}(c_1,c_2)=\mathrm{diag}(\zeta^{j_1c_1+j_2c_2})_j$.

\begin{lemma}\label{i5:x5d-lem-block}
\textup{(1)} Every block datum is the block form of some $g\in\mathcal G_J$.
\textup{(2)} A $g$ carrying $(\chi,\mathfrak b,\delta)$ commutes with $\mathrm{rep}_\chi(\Z/q)$.
\textup{(3)} If moreover $g\in\mathcal G_{J,b}$, each $g'_x$ maps
$\mathcal O_{F,b}^{\omega^{-1}(x)}$ onto itself.
\textup{(4)} If $\mathfrak c_j(s_jP)=\mathfrak b_j(P)$ and each $g'_x$ maps
$\mathcal O_F^{\omega^{-1}(x)}$ onto $\Lambda((\mathfrak c_j)_{\omega_j=x})$, then
$g(\mathcal O_F^J)=\Lambda(\mathfrak b)$.
\textup{(5)} $\mathrm{rep}\circ d_t=\mathrm{rep}_{\chi_t}$, $\chi_t(j)=j_1+tj_2$;
$\mathrm{rep}(x)$ preserves every $\Lambda(\mathfrak a)$ and $\mathcal O_{F,b}^J$; if
$\zeta$ is primitive there is $\kappa\colon\Gamma\to\N$ with
$\sigma\zeta=\zeta^{\kappa(\sigma)}$.
\end{lemma}

\begin{proof}
(1) $g=S^{-1}(\bigoplus_xg'_x)S$ with $S(w)_j=s_j(w_j)$.
(2) At coordinate $j$, $s_j(\zeta^{\chi_jc}w_j)=\zeta^{\omega_jc}s_j(w_j)$, and $g'_x$ is
$F$-linear; $s_j$ is injective.
(3) For $w$ on block $x$ put $W_j=s_j^{-1}(w_j)$ if $\omega_j=x$, else $0$; then
$(s_j(gW)_j)_{\omega_j=x}=g'_x(w)$, and $\mathcal O_{F,b}$ is $\Gamma$-stable;
so $g'_x(\mathcal O_{F,b}^{\omega^{-1}(x)})\subseteq\mathcal O_{F,b}^{\omega^{-1}(x)}$.
The same applies to $g^{-1}\in\mathcal G_{J,b}$, which has block form
$(s,\omega,g'^{-1})$; hence $g'_x$ maps $\mathcal O_{F,b}^{\omega^{-1}(x)}$ onto itself.
(4) $\operatorname{ord}_P((gw)_j)\ge\mathfrak b_j(P)$ for all $P$ iff
$\operatorname{ord}_Q(s_j(gw)_j)\ge\mathfrak c_j(Q)$ for all $Q$.
(5) Direct computation; $\zeta$ is a unit; $\sigma\zeta$ is a $q$-th root of unity.
\end{proof}

\begin{lemma}[carrying matrices]\label{i5:x5d-lem-carry}
Let $\chi\colon J\to\Z/q$, $\eta\in\Z^J$, $\mathfrak d$ a divisor.
\textup{(1)} If $\sum_{\chi_j=x}\eta_j=0$ for all $x$, some $g\in\mathcal G_J$ carries
$(\chi,\eta\mathfrak d,1)$.
\textup{(2)} Assume $\mathrm{Lay}(l,c,\nu,\zeta,\gamma)$, $\kappa_0=1+l^c$, some $\chi_j=1$,
$\sum_j\eta_j[\chi_j]=\sum_{i=0}^Mp_i[\kappa_0^i]$ in $\Z[\Z/q]$, and $\beta\in F^\times$
with $\operatorname{ord}_P\beta=\sum_ip_i\mathfrak d(\gamma^iP)$. Then some $g$ carries
$(\chi,\eta\mathfrak d,\delta)$ with $\delta(1)=\beta$, $\delta(x)=1$ for $x\ne1$.
\end{lemma}

\begin{proof}
(1) $s\equiv1$, $\omega=\chi$; Lemma~\ref{i5:x5d-lem-stein} on each block (from $0$ to
$(\eta_j\mathfrak d)_{\chi_j=x}$, determinant $1$), then
Lemma~\ref{i5:x5d-lem-block}(1),(4).
(2) By Lemma~\ref{i5:x5d-lem-ord}, $\kappa_0$ has order $l^{\nu-c}=\operatorname{ord}\gamma$, so
$T(\kappa_0^i)=\gamma^i$ is well defined; put $T(x)=1$, $B(x)=x$ if $x$ is not a
power of $\kappa_0$, and $B(\kappa_0^i)=1$. Since $\gamma^i\zeta=\zeta^{\kappa_0^i}$,
$T(x)^{-1}(\zeta^x)=\zeta^{B(x)}$. Take $s_j=T(\chi_j)^{-1}$, $\omega_j=B(\chi_j)$,
$\mathfrak c_j(P)=\eta_j\mathfrak d(T(\chi_j)P)$. Applying the additive map
$[y]\mapsto[B(y)=x]\,\mathfrak d(T(y)P)$ to the given identity yields
$\sum_{\omega_j=x}\mathfrak c_j(P)=[x=1]\sum_ip_i\mathfrak d(\gamma^iP)=\operatorname{ord}_P\delta(x)$.
The block $x=1$ is nonempty; Lemma~\ref{i5:x5d-lem-stein} gives $g'_x$ with
$\det g'_x=\delta(x)$, and Lemma~\ref{i5:x5d-lem-block}(1),(4) give $g$.
\end{proof}

\begin{lemma}[evaluation]\label{i5:x5d-lem-val}
\textup{(1)} Let $\theta$ have value $(\bar\zeta,v)$, $V\cong\bigoplus_{b\in B}U_b$,
$\sigma$ a representation of $\Z/q$ acting on $U_b$ by the scalars $\chi_{x_b}$, and
$g\in\mathrm{GL}(V)$ acting on $U_b$ by $g'_b$. Then
$\theta_V(\mathrm{prism}_g(\sigma_*z_e))=\sum_bx_b^{e+1}v(\det g'_b)$.
\textup{(2)} Let $k$ be finite, $r$ split with $r\iota$ onto and $r(\iota b)\ne0$,
$\mu$ as in Lemma~\ref{i5:x5d-lem-split}(3), $\zeta$ primitive,
$\bar\zeta=r(\zeta)$, $\theta$ with value $(\bar\zeta,v)$, and $f$ as in
Lemma~\ref{i5:x5d-lem-res}(4). If $g\in\mathcal G_{J,b}$ carries $(\chi,\mathfrak b,\delta)$,
then
$\theta_V(\mathrm{prism}_{\mu(g)}((\mu\,\mathrm{rep}_\chi)_*z_e))=\sum_{\sigma\in\Gamma}\sum_{x\in\Z/q}(\kappa(\sigma)x)^{e+1}f(\sigma\delta(x))$.
\end{lemma}

\begin{proof}
(1) Let $\mathcal R\colon\Z/q\times\prod_b\mathrm{GL}(U_b)\to\mathrm{GL}(V)$,
$(c,(h_b))\mapsto\bigoplus_b\chi_{x_b}(c)h_b$. Then
$\mathrm{prism}_g(\sigma_*z_e)=\mathcal R_*Z$ with $Z=\mathrm{prism}_{(1,g')}(\mathrm{inl}_*z_e)$, a
cycle (Lemma~\ref{i5:x5d-lem-chain}(3)). By (b), (c) and induction on $|B|$ (for
$B=\emptyset$ use Lemma~\ref{i5:x5d-lem-chain}(2)),
$\theta_V(\mathcal R_*Z)=\sum_b\theta_{U_b}(\mathrm{prism}_{g'_b}((\mathrm{scal}\circ\chi_{x_b})_*z_e))$.
(2) Let $(s,\omega,g')$ be the block datum. By Lemmas~\ref{i5:x5d-lem-block}(3) and
\ref{i5:x5d-lem-res}(3) there are $\bar g_{x,\sigma}$ with $\sigma\delta(x)$ reducing to
$\det\bar g_{x,\sigma}$. Decompose $V=\bigoplus_{(\sigma,x)}k^{\omega^{-1}(x)}$ via
$w\mapsto(w(\sigma s_j,j))_{\omega_j=x}$. For $w$ choose $a\in\mathcal O_F^J$ with
$r(\tau a_j)=w(\tau,j)$. At $(\sigma s_j,j)$, $\mu(\mathrm{rep}_\chi(c))w$ is the reduction of
$\sigma s_j(\zeta^{\chi_jc}a_j)=\zeta^{\kappa(\sigma)\omega_jc}\sigma s_j(a_j)$, so
$\mu\,\mathrm{rep}_\chi$ acts on block $(\sigma,x)$ by $\chi_{\kappa(\sigma)x}$; and
$\mu(g)w$ is the reduction of $\sigma(g'_x((s_ia_i))_j)$, so $\mu(g)$ acts there by
$\bar g_{x,\sigma}$. Apply (1); $v(\det\bar g_{x,\sigma})=f(\sigma\delta(x))$.
\end{proof}

\begin{lemma}\label{i5:x5d-lem-beta}
Let $\gamma\in\Gamma$, $\kappa_0\in\N$, $E<M$, and $\operatorname{ord}_P\alpha=\mathfrak d(\gamma P)-\kappa_0^E\mathfrak d(P)$.
Let $Q=\prod_{j<M,j\ne E}(X-\kappa_0^j)$, $\prod_{j<M}(X-\kappa_0^j)=\sum p_iX^i$, and
$\alpha^{p(\gamma^{-1})}=\prod_i\gamma^{-i}(\alpha)^{c'_i}$ for $p=\sum c'_iX^i$.
\textup{(1)} $\operatorname{ord}_P\alpha^{Q(\gamma^{-1})}=\sum_ip_i\mathfrak d(\gamma^iP)$.
\textup{(2)} Under $\mathrm{Lay}(l,c,\nu,\zeta,\gamma)$, with $U\subseteq F^\times$ a $\Gamma$-stable
subgroup, $f$ additive on $U$, $\alpha\in U$, $E'\in\N$ and $a=(1+l^c)^{E'}$:
$\sum_\sigma\kappa(\sigma)^{E'}f(\sigma\alpha^{p(\gamma^{-1})})=p(a)f(\mathrm{nrm}_a\alpha)$.
\end{lemma}

\begin{proof}
(1) $\operatorname{ord}_P\gamma^{-i}\alpha=\mathfrak d(\gamma^{i+1}P)-\kappa_0^E\mathfrak d(\gamma^iP)$,
so the order is $\sum_i[(X-\kappa_0^E)Q]_i\,\mathfrak d(\gamma^iP)$.
(2) $\kappa$ is multiplicative mod $q$ and $\kappa(\gamma^j)\equiv(1+l^c)^j$. The left
side is $\sum_ic'_i\sum_\sigma\kappa(\sigma)^{E'}f(\sigma\gamma^{-i}\alpha)$; with
$\tau=\sigma\gamma^{-i}$, $\kappa(\sigma)^{E'}=a^i\kappa(\tau)^{E'}$, giving $p(a)S$,
$S=\sum_{j<\operatorname{ord}\gamma}a^jf(\gamma^j\alpha)=f(\mathrm{nrm}_a\alpha)$.
\end{proof}

\medskip\noindent\emph{Proof of \textup{(ii)}.}
Let $J=(\Z/q)^2$, $q=l^\nu$, $\kappa_0=1+l^c$, $M=2e+3$, $\kappa$ as in
Lemma~\ref{i5:x5d-lem-block}(5), and $\mathfrak d$ with
$\operatorname{ord}_P\alpha=\mathfrak d(\gamma P)-a\mathfrak d(P)$.
\emph{Data.} Lemma~\ref{i5:x5d-lem-split}(1),(2) give $b_0$ and $b$. Lemma~\ref{i5:x5d-lem-rel}
gives $n_t$ ($t\le M$) and $t_0$; Lemma~\ref{i5:x5d-lem-eta} ($M<l^c$) gives $\eta$.
Let $\beta=\alpha^{Q(\gamma^{-1})}$ with $(E,M)=(e+1,2e+3)$; by
Lemma~\ref{i5:x5d-lem-beta}(1), $\operatorname{ord}_P\beta=\sum_ip_i\mathfrak d(\gamma^iP)$.
\emph{Matrices.} Since $\sum_j\eta_j[\chi_t(j)]=\pi_t(\eta)$,
Lemma~\ref{i5:x5d-lem-carry}(1) gives $g_t$ carrying $(\chi_t,\eta\mathfrak d,1)$ for
$t\ne t_0$, and Lemma~\ref{i5:x5d-lem-carry}(2) ($\chi_{t_0}(1,0)=1$) gives $g_{t_0}$
carrying $(\chi_{t_0},\eta\mathfrak d,\delta_{t_0})$, $\delta_{t_0}(1)=\beta$. By
Lemma~\ref{i5:x5d-lem-res}(2) all $g_t\in\mathcal G:=\mathcal G_{J,b}$. Put
$\mathcal H=\mathcal H_{J,b}$, $\rho=\mathrm{rep}$. By Lemma~\ref{i5:x5d-lem-block}(5),(2),
$g_t$ commutes with $\rho(d_t(\Z/q))$; $\rho(x)\in\mathcal H$;
$g_{t_0}^{-1}\rho(x)g_{t_0}\in\mathcal H$ since $\rho(x)$ preserves
$\Lambda(\eta\mathfrak d)=g_{t_0}\mathcal O_F^J$; and $g_{t_0}^{-1}g_t\in\mathcal H$
since both map $\mathcal O_F^J$ onto $\Lambda(\eta\mathfrak d)$.
\emph{Chains.} Lemma~\ref{i5:x5d-lem-cchain} ($C=\Z/q$, $D=(\Z/q)^2$, $z=z_e$) gives a
cycle $y'\in C_{2e+3}(\mathcal H;\Z/q)$ with
$\mathrm{incl}_*y'\equiv\sum_tn_t\mathrm{prism}_{g_t}((\rho d_t)_*z_e)$ modulo boundaries.
Lemmas~\ref{i5:x5d-lem-split}(4) and \ref{i5:x5d-lem-free} ($R=\mathcal O_L$,
$M=\mathcal O_F^J$) give $N$ and $\lambda\colon\mathcal H\to\mathrm{GL}_N(\mathcal O_L)$,
independent of $k$ and $r$. Put $y=\lambda_*(n_{t_0}^{-1}y')$ and $b_\alpha=\iota(b)\ne0$.
\emph{Evaluation.} Let $r,v,\theta$ be as in the statement. As $b_0\mid b$, $r$
splits; let $\mu,U,f$ be as above. For $P\not\ni b$, also $\gamma P\not\ni b$, so
$\alpha\in U$ and $\mathrm{nrm}_a\alpha\in U$ reduces to some $u\in k^\times$, with
$f(\mathrm{nrm}_a\alpha)=v(u)$. By Lemmas~\ref{i5:x5d-lem-free} and
\ref{i5:x5d-lem-triv}, $\theta_{k^N}((\mathrm{red}_{r\iota,N})_*y)=n_{t_0}^{-1}\theta_V(\mu_*\mathrm{incl}_*y')
=n_{t_0}^{-1}\sum_tn_t\theta_V(\mathrm{prism}_{\mu(g_t)}((\mu\rho d_t)_*z_e))$, and by
Lemma~\ref{i5:x5d-lem-val}(2) the $t$-th term is
$\sum_\sigma\sum_x(\kappa(\sigma)x)^{e+1}f(\sigma\delta_t(x))$. For $t\ne t_0$ it
vanishes ($f(1)=0$); for $t_0$ only $x=1$ survives, giving
$\sum_\sigma\kappa(\sigma)^{e+1}f(\sigma\beta)=Q(a)f(\mathrm{nrm}_a\alpha)$ by
Lemma~\ref{i5:x5d-lem-beta}(2), and $Q(a)=\mathrm{pe}$. Hence the value is
$\mathrm{pe}\cdot v(u)$. \qed

\medskip\noindent\emph{Formal verification.}
Statements (i-a), (i-b) and (ii) of Proposition~\ref{i5:x5-prop-detect-construct}
are formally verified in \cite{CharacterFamily}, through formal counterparts of
Lemmas~\ref{i5:x5d-lem-chain}, \ref{i5:x5d-lem-trans}, \ref{i5:x5d-lem-coc},
\ref{i5:x5d-lem-rel}, \ref{i5:x5d-lem-eta} and \ref{i5:x5d-lem-stein}.  The
formal version of Lemma~\ref{i5:x5d-lem-eta} has the extra hypothesis
$c\le\nu$; the lemma as written above is slightly more general, and its proof
does not use that hypothesis.


\subsection{The automorphic cores}\label{i5:x5-ssec-arch}

For each core we list the formally verified statements its formal proof is
built from, and prove the assembly steps.  Statements marked ``formally
verified'' are stated exactly (with the definitions of
\S\ref{i5:x5-ssec-defs} and below) unless the contrary is said.

\subsubsection{Two estimates from commutative algebra}

Let $\mathcal O$ be a complete discrete valuation ring of characteristic $0$,
$A=\bigoplus_{j\in\Z}A_j$ a $\Z$-graded $\mathcal O$-algebra and
$x\colon A\to\mathcal O$ an $\mathcal O$-algebra map.  Put
$\mathrm{sat}(x)=\{f\in A:cf\in(\ker x)^2\text{ for some }0\ne c\in\mathcal O\}$,
$\mathrm{Cot}_j=$ the image of $A_j\cap\ker x$ in $A/\mathrm{sat}(x)$, and for a
family $F=(F_i)$ in $A$, $\mathrm{cotlen}_j(F)=\lgth_{\mathcal O}\bigl(\mathrm{Cot}_j
/(\mathrm{Cot}_j\cap\langle F_i\rangle)\bigr)\in\N\cup\{\infty\}$, where
$\langle F_i\rangle$ is the $\mathcal O$-span of the images of the $F_i$; $F$
\emph{spans $\mathrm{Cot}_j$ up to torsion} if every element of $\mathrm{Cot}_j$
has a non-zero multiple in $\langle F_i\rangle$.  A \emph{point derivation} at
$x$ (with values in an $\mathcal O$-module $Q$ that is a module over the target
of $x$) is an $\mathcal O$-linear $D$ with $D(fg)=x(f)D(g)+x(g)D(f)$; it has
\emph{weight $j$} if it kills $A_i$ for $i\ne j$.

\begin{proposition}[(M) formally verified, \cite{CharacterFamily}; (M$'$) proved in \S\ref{i5:x6-sssec-M}; cotangent estimates]\label{i5:x5-prop-UM}
\textup{(M$'$)} Assume moreover $k_{\mathcal O}$ finite.  Let $j\in\Z$ with
$x(A_j)=0$, $F=(F_i)$ a family in $A_j$, $0\ne c\in\mathcal O$, $M$ an
$\mathcal O$-module and $D\colon M\to\mathrm{Hom}_{\mathcal O}(A,\mathcal O/c)$
an injective additive map such that $D(\mu)(F_i)=0$ for all $\mu,i$, and every
$D(\mu)$ is the reduction mod $c$ of a point derivation of weight $j$ at $x$
with values in $\mathcal O$.  Then $\lgth_{\mathcal O}M\le\mathrm{cotlen}_j(F)$.

\textup{(M)} Let $\bar K\supseteq\mathcal O$ be an algebraically closed field
with a function $v\colon\bar K\to\R$ which is additive on products of non-zero
elements, satisfies $v(x+y)\ge\min(v(x),v(y))$ for non-zero $x,y,x+y$, and
restricts to the normalised valuation on $\mathcal O\smallsetminus\{0\}$.
Assume: $A_0$ is complete local over $\mathcal O$, $A$ is generated over $A_0$
by finitely many elements of degrees $\pm1,\pm2$, and $x(A_j)=0$ for $j\ne0$.
Let $I$ be the ideal generated by families $F_0,\dots,F_r\in A_{-1}$,
$G_1,\dots,G_s\in A_{-2}$, $E_1,\dots,E_{r'}\in A_1$, $E'_1,\dots,E'_{s'}\in
A_2$ and $Z_1,\dots,Z_m\in A_0\cap\ker x$, and let $a\in A$.  Assume (M1)
$r\ge1$; $x$ is a smooth point of dimensions $(s,r,m,r'+1,s'+1)$ in degrees
$(-2,-1,0,1,2)$, i.e.\ there is a graded $\mathcal O$-algebra map
$\phi\colon A\to\mathrm{Frac}(\mathcal O)[[t_\tau]]$ to a power series ring in
that many variables of those degrees, with constant terms given by $x$, every
variable reached by some $\phi(f)$ alone among the linear terms, and
$x(f)=0$ with vanishing linear terms only for $f\in\mathrm{sat}(x)$; and
$(F_i)$, $(G_i)$, $(Z_i)$ span $\mathrm{Cot}_{-1},\mathrm{Cot}_{-2},\mathrm{Cot}_0$
up to torsion; (M2) the image of $A_0$ in $A/I$ is a finitely generated
$\mathcal O$-module; (M3) for every $\bar K$-point $x_1\ne x$ of $A/I$ killing all
$A_j$, $j\ne0$, there is no point of $A/I$ agreeing with $x_1$ on $A_0$ that
kills $A_{<0}$ but not $A_{>0}$, or there is none that kills $A_{>0}$ but not
$A_{<0}$; (M4) at every \emph{two-signed} $\bar K$-point $\xi$ of $A/I$ (one
that kills neither $A_{<0}$ nor $A_{>0}$) any two point derivations vanishing
on $I$ are linearly dependent; (M5) $a\in A_0$, $x(a)=0$, and $\xi(a)\ne0$ at
every two-signed $\bar K$-point $\xi$ of $A/I$.  Then there is a finite set $T$
of two-signed $\bar K$-points of $A/I$ meeting every orbit of the
$\bar K^\times$-action $(t\cdot\xi)(f)=t^j\xi(f)$ ($f\in A_j$) on two-signed
points; $\mathrm{cotlen}_{-1}(F)<\infty$; and for every such $T$,
$\mathrm{cotlen}_{-1}(F)\le2\sum_{\xi\in T}v(\xi(a))$.
\end{proposition}

(M) is where the factor $2$ in all length estimates comes from; both
statements are pure commutative algebra.  The proof of (M$'$) and the structure
of the formal proof of (M) are in \S\ref{i5:x6-sssec-M}.

\subsubsection{(U1)}\label{i5:x5-sssec-archU1}

Fix a datum $D$ (Definition~\ref{i5:x5-def-dat}) and a valuation ring $B$ of
$\Omega$ over $l$.  For a local ring $A$ over $\mathcal O$, a homomorphism
$r\colon G_\Q\to\mathcal G_3(A)$ is an \emph{ordinary point at $(B,g_1)$} if: it
has multiplier $\eta$ and reduces to $\bar r_0$; its entries are locally
constant mod every $\mathfrak m_A^n$; it is normalised at $(c_D,g_1)$; $r(s)=1$
for every $s\in I_{B'}$, $B'$ not over $l$, with $\eta(s)=1$, $\Psi(s)=1$; and
$r|_{D_B}$ is triangular in the sense of Definition~\ref{i5:x5-def-frame}(T)
(with $D_B,I_B$).  $\mathrm{FinEnd}(B,g_1)$: there is a subgroup $U\le G_\Q$ of
finite index such that for every complete discrete valuation ring $A$ over
$\mathcal O$ with finite residue field in which $\mathfrak m_{\mathcal O}$ maps to
$0$ (so $A$ has characteristic $l$), every ordinary point $r$ at $(B,g_1)$ with
values in $A$ and every $g\in U\cap\ker\eta$, the characteristic polynomial of
$\operatorname{mat}r(g)$ is that of $\operatorname{tor}\Psi(g)$.
$\mathrm{AutEnd}_3(B,g_1)$: for every subring $A\subseteq\overline\Q_l$, a discrete
valuation ring finite over $\mathcal O$, and every ordinary point $r$ at
$(B,g_1)$ with values in $A$ whose image in $\mathcal G_3(\overline\Q_l)$ is irreducible
on every open subgroup of $\ker\eta$, $\mathrm{Aut}_\iota(\rho_K(r)\otimes\bar
\chi^{-1})$.  $\mathrm{AutEnd}_2(B,g_1)$: for the same $A$ and every ordinary
point $r$ of type~II with $\operatorname{mat}r(g)_{11}=1$ on $\ker\eta$, every
$W$ extending $r$ (with $c_D$) satisfies $\mathrm{Mod}_\iota(W)$.

\begin{proposition}[residually reducible lifting at ordinary points; formally verified, \cite{CharacterFamily}]\label{i5:x5-prop-gdend}
For every datum $D$, $\iota$ and $B$ over $l$ there is $g_1$ with
$(\eta,\Psi,c_D,g_1)$ a marking, $\mathrm{FinEnd}(B,g_1)$,
$\mathrm{AutEnd}_3(B,g_1)$ and $\mathrm{AutEnd}_2(B,g_1)$.
\end{proposition}

\begin{lemma}[formally verified, \cite{CharacterFamily}]\label{i5:x5-lem-U1steps}
Let $D$ be a datum with $e_{\mathcal O}=1$, $g_1$ marked, $B$ the valuation ring
of $\Omega$ defined by $j$ (i.e.\ $x\in B\iff|j(x)|\le1$), and $(R_D,u_D)$ a
global ring for $(D,g_1)$.  Granting the published local theorems of
\cite{CharacterFamily} (local Tate duality, finiteness, Euler characteristic and
the Bloch--Kato local statements over $\Q_l$ and its finite extensions):
\textup{(1)} $\mathrm{FinEnd}(B,g_1)$ implies that $W_0(u_D)$ is finite over
$\mathcal O$;
\textup{(2)} if $W_0(u_D)$ is finite, then for every point $(\xi,r)$ there are a
subring $A\subseteq\overline\Q_l$, a discrete valuation ring finite over $\mathcal O$,
$t\in\overline\Q_l^\times$ and an ordinary point $r'$ at $(B,g_1)$ with values in $A$ with
$r'=\operatorname{conj}_{\operatorname{tor}t}r$ in $\mathcal G_3(\overline\Q_l)$;
\textup{(3)} an ordinary point $r$ at $(B,g_1)$ with values in such an $A$ which
is of type~I in $\mathcal G_3(\overline\Q_l)$ is irreducible on every open subgroup of
$\ker\eta$.
\end{lemma}

The hypothesis $e_{\mathcal O}=1$ is used by none of (1)--(3); (U1) holds for
every coefficient ring, and $e_{\mathcal O}=1$ is only part of the formal
statement.

\begin{proof}[Proof of Proposition~\ref{i5:x5-prop-U1} from Proposition~\ref{i5:x5-prop-gdend}, Lemma~\ref{i5:x5-lem-U1steps} and the published local theorems]
Let $B$ be the valuation ring defined by $j$; it lies over $l$.  Take $g_1$ from
Proposition~\ref{i5:x5-prop-gdend}.  (F) is Lemma~\ref{i5:x5-lem-U1steps}(1).
(A): let $(\xi,r)$ be a point of type~I with multiplier $\eta$, and take
$A,t,r'$ from Lemma~\ref{i5:x5-lem-U1steps}(2).  On $\ker\eta$ the matrix parts
satisfy $\operatorname{mat}r'(g)=\operatorname{tor}(t)\operatorname{mat}r(g)
\operatorname{tor}(t)^{-1}$ (for $\nu=1$ conjugation by $(\operatorname{tor}t,
1)$ is ordinary conjugation), so $r'$ is of type~I, hence irreducible on open
subgroups of $\ker\eta$ by Lemma~\ref{i5:x5-lem-U1steps}(3), and
$\mathrm{AutEnd}_3$ gives $\mathrm{Aut}_\iota(\rho_K(r')\otimes\bar\chi^{-1})$.
Since $\rho_K(r')=\operatorname{tor}(t)\rho_K(r)\operatorname{tor}(t)^{-1}$ and
$\mathrm{Aut}_\iota$ (attachment and the witness) depends only on the
isomorphism class, $\mathrm{Aut}_\iota(\rho_K(r)\otimes\bar\chi^{-1})$.
(W): let $(\xi,r)$ be of type~II with $\operatorname{mat}r_{11}=1$ on $\ker\eta$
and $W$ extending $r$.  Take $A,t,r'$ as before and put $W'=\mathrm{diag}(t,t^{-1})
W\mathrm{diag}(t,t^{-1})^{-1}$.  Conjugation by $\operatorname{tor}t$ preserves
the block shape and multiplies the corner by $\mathrm{diag}(t,t^{-1})$ on both
sides, so $r'$ is of type~II, $\operatorname{mat}r'_{11}=1$ on $\ker\eta$ and
$W'$ extends $r'$ (the trace at $c_D$ is unchanged).  $\mathrm{AutEnd}_2$ gives
$\mathrm{Mod}_\iota(W')$, and $\mathrm{Mod}_\iota$ is invariant under
isomorphism, so $\mathrm{Mod}_\iota(W)$.
\end{proof}

What is formally verified is Proposition~\ref{i5:x5-prop-gdend} and
Lemma~\ref{i5:x5-lem-U1steps}.  The formal proof of
Proposition~\ref{i5:x5-prop-gdend} is a residually reducible lifting argument of
Skinner--Wiles type (base change to an auxiliary abelian CM field, Hida theory,
Eisenstein-congruence seeds, propagation of automorphy along the components of
the ordinary deformation space, and solvable descent); it contains no
Taylor--Wiles patching.  Its structure is described in \S\ref{i5:x6-sssec-U1}.

\subsubsection{(U2)--(U4) and $\mathrm{OneCharU}$}

(U2a), (U2b) and (U3) are single formal statements, derived in
\cite{CharacterFamily} from its bundle of published hypotheses (including an
adjoint lifting statement derived there from cohomological inputs) and from
the level and witness computations; the structure of their proofs is
described in \S\ref{i5:x6-sssec-U234}.  (U3$'$) is a formally verified lemma
reading the endoscopic admissibility off (W) and the local conditions.  The
length estimate (U4) is assembled from four formally verified steps:
\begin{enumerate}
\item[(S1)] (graded model) There are a graded $\mathcal O$-algebra $A$ with
$r\ge1$ and a smooth point $x$ of the dimensions of (M1), families
$F,G,E,E',Z$, and a surjection $q\colon A\to R_D$ with kernel the ideal $I$ they
generate, compatible with the base point ($x_0\circ q=x$, $x_{0*}u_D=r_0$), with
each $\operatorname{mat}u_D(g)_{ab}$ the image of an element of
$A_{\operatorname{wt}(a,b)}$, the weight pieces of $R_D$ topologically spanned by
monomials in the entries, every two-signed $\overline\Q_l$-point of $A/I$ conjugate
under the torus action to $f\circ q$ for a local $f$; if $f,f'$ are local
and $f'_*u_D=\operatorname{conj}_{\operatorname{tor}t}f_*u_D$ then $f'\circ
q=t\cdot(f\circ q)$ (by the density clause, a local $f$ being continuous for the
$\mathfrak m$-adic topology); and every point
derivation at such a point vanishing on $I$ coming from a deformation of
$f_*u_D$ to $\overline\Q_l[\varepsilon]$.  (Statement summarised; the formal
definition of a model has fifteen clauses.)
\item[(S2)] (hypotheses (M1)--(M5)) The spanning clauses of (M1) follow from
the rigidity of $r_0$ in weights $\le0$, which is Bloch--Kato vanishing for the
weight-$j$ pieces of $\mathrm{ad}\,r_0$, $j\le0$ (formally verified from
Theorem~\ref{i5:thm-pt} and Corollary~\ref{i5:cor-lemmaK}); (M2) follows from
(F); (M3) from the formally verified statement that a local point which is both
upper and lower is $r_0$; (M4) from the rigidity of every
two-signed point, which is (U2a) at type-I points and (U2b) together with
Proposition~\ref{i5:x-prop:katotwow} at type-II points; (M5) for
$a=\operatorname{tr}\operatorname{mat}u_D(\sigma)-t_{e,\psi}(\sigma)$ from the
hypothesis of (U4).
\item[(S3)] (Selmer to cotangent) $V'(e,\psi)\le\mathrm{cotlen}_{-1}(F)$.  It
is deduced from (M$'$) with $j=-1$, applied to finite-length submodules of the
Selmer group of Definition~\ref{i5:def-V}, each class giving a point derivation
of weight $-1$ at $x$ (the class is the $\Psi^{-1}$-entry of a first-order
deformation of $r_0$); the conclusion is exact, the description of the proof is
a summary.
\item[(S4)] (cotangent to points) (M) gives a finite set $T_A$ meeting every
torus orbit of two-signed points, $\mathrm{cotlen}_{-1}(F)<\infty$, and
$\mathrm{cotlen}_{-1}(F)\le2\sum_{\xi\in T'}v(\xi(a))$ for every finite set
$T'$ of two-signed points meeting every orbit.
\end{enumerate}
Steps (S1)--(S3) use the published Galois-cohomological statements of
Theorem~\ref{i5:thm-pt} (local theorems, Poitou--Tate, global finiteness and
Euler characteristic) and Corollary~\ref{i5:cor-lemmaK}; so does (U3$'$).

\begin{proof}[Proof of (U4) from (S1)--(S4)]
By (S2) the hypotheses of (M) hold for the model of (S1) and
$a=\operatorname{tr}\operatorname{mat}u_D(\sigma)-t_{e,\psi}(\sigma)$, which
lies in $A_0$ (diagonal entries have weight $0$) and is killed by $x$ (at $r_0$
the trace is $t_{e,\psi}(\sigma)$, since $\sigma\in G_K=\ker\eta$).  (S4) gives
$T_A$ and $\mathrm{cotlen}_{-1}(F)<\infty$, and (S3) gives $V'(e,\psi)\le
\mathrm{cotlen}_{-1}(F)<\infty$.  Each $\xi\in T_A$ is torus-conjugate to some
$f_\xi\circ q$ with $f_\xi$ local and two-signed (S1).  Let $T$ consist of one
$f_\xi$ for each torus-conjugacy class of the $(f_\xi)_*u_D$, $\xi\in T_A$; it
is finite, and its elements are pairwise not torus-conjugate.  Put
$T'=\{f\circ q:f\in T\}$.  Every two-signed point of $A/I$ is torus-conjugate
to some $f'\circ q$ with $f'$ local (S1), hence to some $\xi\in T_A$, hence to
$f_\xi\circ q$; if $f\in T$ is the chosen representative of the class of
$(f_\xi)_*u_D$, then $f_*u_D=\operatorname{conj}_{\operatorname{tor}t}
(f_\xi)_*u_D$ for some $t$, so $f\circ q=t\cdot(f_\xi\circ q)$ by (S1).  Thus
$T'$ meets every orbit, and $f\mapsto f\circ q$ is injective since $q$ is
surjective.  Applying (M) to $T'$, with $v=v_l$ and $e_{\mathcal O}=1$,
\[
\mathrm{Vn}(e,\psi)=V'(e,\psi)\le\mathrm{cotlen}_{-1}(F)\le2\sum_{f\in T}
v_l\bigl(f(q(a))\bigr)=2\sum_{f\in T}v_l\bigl(\operatorname{tr}\operatorname{mat}
f_*u_D(\sigma)-t_{e,\psi}(\sigma)\bigr).
\]
\end{proof}

The proof of $\mathrm{OneCharU}$ from (U1)--(U4) is
Proposition~\ref{i5:x-prop:onechar}.

\subsubsection{$\mathrm{Count}_{\rm V}$}

For a host frame the ordinary points and the conditions $\mathrm{FinEnd}$,
$\mathrm{AutEnd}_3$ are defined as in \S\ref{i5:x5-sssec-archU1} with
$\eta=\eta_4$, $\Psi=\Psi_{e,\psi}$, $c_D=\sigma$ a complex conjugation and $B$
a valuation ring over $l$.

\begin{proposition}[host endomorphisms; formally verified, \cite{CharacterFamily}]\label{i5:x5-prop-hostend}
Let $K=\Q(i)$, $l$ prime, $\iota$, $\mathcal O$, $N$, $e$, $\psi$, $o_1$ with
$\mathrm{HostData}(l,e,\psi,o_1)$, $\sigma$ a complex conjugation and $B$ over
$l$.  There is $g_1\in G_\Q$ with $(\eta_4,\Psi_{e,\psi},\sigma,g_1)$ a marking,
$\mathrm{FinEnd}(B,g_1)$ and $\mathrm{AutEnd}_3(B,g_1)$.
\end{proposition}

\begin{lemma}[formally verified, \cite{CharacterFamily}]\label{i5:x5-lem-Vsteps}
Let $(S,(A_p),F)$ be a host frame for $(e,\psi)$ with
$\mathrm{HostData}(l,e,\psi,o_1)$, $c=\pi(\sigma)$.  (Part (3) uses, for the
frame, the class-number-one statement and the frame conditions, both
consequences of $\mathrm{HostData}$.)
\textup{(1)} If $r$ is a point of $F$ at $\pi(g_1)$ with values in a local ring,
then $r\circ\pi$ is an ordinary point at $(A_l,g_1)$.
\textup{(2)} Let $\mathcal O$ have finite residue field with $2\in\mathcal
O^\times$, $A$ a domain, local over $\mathcal O$ with $\mathfrak
m_{\mathcal O}A=0$, $(\eta,\Psi,c,g_1)$ a marking, $r$ with multiplier $\eta$,
$r\equiv r_0\bmod\mathfrak m_A$, normalised at $(c,g_1)$, such that on a
finite-index subgroup $U$ the characteristic polynomials on $U\cap\ker\eta$ are
those of $r_0$.  Then $r$ is upper or lower.
\textup{(3)} If $u$ is universal at a marked $g_1$ and every point of $F$ at
$g_1$ with values in a complete discrete valuation ring of characteristic $l$
over $\mathcal O$ with finite residue field is upper or lower, then $R_0(u)$ is
finite over $\mathcal O$.
\end{lemma}

\begin{proof}[Proof of Proposition~\ref{i5:x5-prop-Vcore} from Proposition~\ref{i5:x5-prop-hostend}, Lemmas~\ref{i5:x5-lem-Vsteps} and~\ref{i5:x5-lem-Vopen}]
Write $c=\pi(\sigma)$ and $B=A_l$, which lies over $l$.  Take $g_1$ from
Proposition~\ref{i5:x5-prop-hostend}.  Since $\eta\circ\pi=\eta_4$ and
$\Psi\circ\pi=\Psi_{e,\psi}$, $\pi(g_1)$ is marked.  Let $u$ be universal at
$\pi(g_1)$.  (F): by Lemma~\ref{i5:x5-lem-Vsteps}(3) it suffices that every
point $r$ with values in a complete discrete valuation ring $A$ of characteristic
$l$ is upper or lower.  By (1), $r\circ\pi$ is an ordinary point, so
$\mathrm{FinEnd}$ gives $U$ with the characteristic polynomial condition on
$U\cap\ker\eta_4$; (2), applied to $r\circ\pi$ and the marking
$(\eta_4,\Psi_{e,\psi},\sigma,g_1)$, shows that $r\circ\pi$ is upper or lower, and
since $\pi$ is surjective so is $r$.  (A): let $(\xi,r)$ be a point of type~I.
By Lemma~\ref{i5:x5-lem-Vopen} $r$ is a characteristic-zero point, i.e.\
$r'=\operatorname{conj}_{\operatorname{tor}t}r$ is a point of $F$ with values in
a finite discrete valuation ring $A\subseteq\overline\Q_l$, and $a_r$ is irreducible on
open subgroups of $\ker\eta_4$; so is $a_{r'}$ (conjugate on $\ker\eta_4$).  By
(1), $r'\circ\pi$ is ordinary at $(A_l,g_1)$, and $\mathrm{AutEnd}_3$ gives
$\mathrm{Aut}_\iota(\rho\otimes\bar\chi^{-1})$ for $\rho=a_{r'}|_{G_K}$; as in
the proof of (U1), this transfers to $a_r|_{G_K}$, which is
$\mathrm{AutOn}_{K,\iota,-1}(a_r)$.
\end{proof}

For Proposition~\ref{i5:x5-prop-Vhost}: (a) is derived formally, as (U2a), from
the bundle of published hypotheses of \cite{CharacterFamily} and the adjoint
lifting statement derived there; (b) is
derived from $\mathrm{NumberU3U}$ (Proposition~\ref{i5:x-prop:numberu3u}) by
choosing a suitable prime $p\ge7$, $p\nmid lN$, and $\varphi$ a Frobenius
at a place above $p$, and showing that $a_r|_{G_K}$ is stably admissible of
level $N$ for every characteristic-zero type-I point with $\mathrm{B}'$; (d) is
the abstract bound for frames (the frame version of (S1)--(S4), with (M),
(M$'$), the vanishing statement for the weight $\le0$ cocycles of the frame, the
class-number-one statement for the frame, and the published Galois cohomology
of Theorem~\ref{i5:thm-pt}) followed by the formally verified comparison of the
frame Selmer length with $V'(e,\psi)$.  The proof of $\mathrm{Count}_{\rm V}$
from Propositions~\ref{i5:x5-prop-Vcore} and~\ref{i5:x5-prop-Vhost} is
Proposition~\ref{i5:prop-c5}.

\subsubsection{$\mathrm{Count}_{\rm II}$ and $\mathrm{Count}_{\rm III}$}

Proposition~\ref{i5:x5-prop-sp4core}(a) is the $\mathrm{GSp}_4$ analogue of
Propositions~\ref{i5:x5-prop-gdend} and~\ref{i5:x5-prop-hostend} combined with the
finiteness criterion of Lemma~\ref{i5:x5-lem-Vsteps}(3); (e) is the
$\mathrm{GSp}_4$ analogue of (U2a); (f) is the $\mathrm{GSp}_4$ analogue of
(S1)--(S4), with weights in $\Z^2$ graded through the first coordinate and
(M), (M$'$) used as above; (g) is a count of admissible
$4$-dimensional representations of $G_\Q$ of level $M$ and weights
$(e+1,e,1,0)$ (formally verified from the published hypotheses), followed by a
norm argument as in Proposition~\ref{i5:x-prop:normbound} (in its rank-four
form, with $\mathrm{eis}_4$).  The assembly is
the proof of Proposition~\ref{i5:prop-c23}.


\subsection{Outline of the formally verified steps}\label{i5:x6-ssec}

We outline, for each formally verified statement of \S\ref{i5:x5-ssec} and
\S\ref{i5:x5-ssec-arch}, how its formal proof in \cite{CharacterFamily} is
organised.  The outline is not a proof, and nothing below is used as a proof
step elsewhere, except Proposition~\ref{i5:x5-prop-UM}(M$'$), which is proved in
full in \S\ref{i5:x6-sssec-M}.  Each step is classified as (L), (N) or (R) as
explained in the introduction to this annex.  For an (L) step we say whether the formal development isolates it as one of the published
statements of Table~\ref{i5:x6-tab-bundle} (or derives it from them in a few
lines), each of which it proves as a separate proposition, or
\emph{re-derives} it inside a longer argument; and we say
where the formal statement differs from the published one; for an (N) step we
say what is constructed and why it works.  The outline follows the formal
statements and the top levels of their proofs; when we say that we do not
describe a part, it enters only through its statement.  The identification of
a formal statement with a published theorem is ours (the formal text names no
authors).  Several steps are shared by several cores; they are
described once, in \S\ref{i5:x6-sssec-common}.

Two facts shape everything below.  First, in every core the residual
representation is a sum of characters ($\bar\Psi\oplus1\oplus\bar\Psi^{-1}$ on
$G_K$, resp.\ four characters for $\mathrm{GSp}_4$), so the automorphy lifting
theorems for residually irreducible (adequate, big, enormous) images do not
apply, and the relevant literature is the residually reducible one: Skinner and
Wiles (Publ.\ Math.\ IH\'ES 89 (1999), 5--126), Thorne, \emph{J. Amer. Math. Soc.} 28 (2015), 785--870,
and Allen--Newton--Thorne, Compositio Math.\ 156 (2020), 2399--2422, arXiv:1912.11269.  Second, the
deformation rings carry an action of the diagonal torus of the residual
representation, hence a $\Z$-grading, and the quantitative content of the
section (the length estimates) is extracted from this grading; this part has
no counterpart in the automorphy lifting literature.

\medskip\noindent\label{i5:x6-sssec-common}\emph{Steps shared by several cores.}

\medskip\noindent\emph{(C1) Universal rings (L, re-derived).}
\emph{Statement.} Let $\Gamma$ be a profinite group satisfying Mazur's
finiteness condition $\Phi_l$ (every open subgroup has finitely many continuous
homomorphisms to $\Z/l$), $\mathcal O$ a complete discrete valuation ring with
residue field $k$ of characteristic $l$, and $\bar r$ a continuous
representation of $\Gamma$ over $k$.  Then the functor of lifts of $\bar r$ to
complete local noetherian $\mathcal O$-algebras with residue field $k$ is
representable; so is any subfunctor cut out by the local conditions of
Definitions~\ref{i5:x5-def-frame}, \ref{i5:x5-def-hostframe},
\ref{i5:x5-def-dat} and~\ref{i5:x5-def-gsp4}, which are given by closed
conditions on matrix entries.  This is Mazur's representability theorem for
framed deformations (Mazur 1989; reference to be confirmed).  The formal proof is a
direct construction (the lifts to the dual numbers form a finite set, they are
separated by finitely many group elements, and the ring is built level by
level), with $\bar r$ the reducible residual representation itself.
\emph{Input.} Condition $\Phi_l$ for $G_{\Q,S}$ is deduced from global
finiteness of $H^1(G_{\Q,S},M)$ for finite $l$-torsion $M$ (re-proved in the formal development, (H4.5)
of Table~\ref{i5:x6-tab-bundle}).

\medskip\noindent\emph{(C2) Graded presentations (L, re-derived, with (N) bookkeeping).}
Conjugation by $\operatorname{tor}(t)=\mathrm{diag}(t,1,t^{-1})$, $t\in\mathcal
O^\times$, acts on the universal ring; in the limit it makes the ring the
completion of a $\Z$-graded $\mathcal O$-algebra in which the matrix entry
$(a,b)$ has weight $\operatorname{wt}(a)-\operatorname{wt}(b)\in\{0,\pm1,\pm2\}$.
The formal development constructs (\emph{graded model}, (S1) of
\S\ref{i5:x5-ssec-arch}) a presentation
\[
R\cong \mathcal O[[z_\tau]]/(F,G,E,E',Z)
\]
with homogeneous variables and homogeneous relations, in which the numbers of
variables and of relations of weight $j$ are controlled by
$\dim H^1$ and $\dim H^2$ of the weight-$j$ piece of $\mathrm{ad}\,r_0$ with
the local conditions.  This is the classical presentation of a deformation ring
with $h^1$ generators and $h^2$ relations (Mazur; B\"ockle for the relative
form; reference to be confirmed), made torus-equivariant.  The count of relations
uses the Greenberg--Wiles formula, Poitou--Tate and the global Euler
characteristic formula ((H4.1), (H4.6), (H4.8), re-proved in the formal development).  The new piece of
bookkeeping is the exact defect in each weight: in weights $\pm1$ and $0$ the
defect is the dimension of the $(-1)$-eigenspace of complex conjugation on the
weight piece, and as a result the weight $-1$ part is over-determined by
exactly one equation (``$r+1$ equations in $r$ variables'' in (M1)).  That
single excess equation is what makes the weight $-1$ cotangent module have
finite length, and it is the source of the Selmer bound.

\medskip\noindent\emph{(C3) Vanishing of adjoint Selmer groups at automorphic points (L, re-derived).}
\emph{Formal statement.} Let $L$ be a CM field containing no non-trivial
$l$-th root of unity and satisfying a largeness condition on its degree, let
$2\le n\le4$, and let $x$ be a point of the universal ordinary
pseudo-character (determinant) ring of an $n$-dimensional conjugate self-dual
residual representation of $G_L$, which is automorphic (attached to a regular
algebraic conjugate self-dual cuspidal $\pi$), locally generic and ordinary.
Then every continuous adjoint cocycle at $x$ satisfying the ordinary local
conditions is a coboundary; i.e.\ the ordinary (Bloch--Kato) adjoint Selmer
group at $x$ vanishes.
\emph{Method.} An $R^{\rm ord}=\mathbb T^{\rm ord}$ argument at the point $x$:
the ordinary Hecke algebra (Hida theory for definite unitary groups over $L$,
derived formally from the bundle) and Taylor--Wiles
auxiliary data give a patched ring over a power series ring, and a counting
argument shows that every derivation at $x$ vanishes on traces.  This is the
form of the statement proved by Geraghty (ordinary automorphy lifting), Allen
(\emph{Deformations of polarized automorphic Galois representations and adjoint
Selmer groups}, Duke Math.\ J.\ 165 (2016), 2407--2460) and Newton--Thorne (\emph{Adjoint Selmer groups of automorphic
Galois representations of unitary type}, J.\ Eur.\ Math.\ Soc.\ 25 (2023), 1919--1967); Geraghty, \emph{Modularity lifting
theorems for ordinary Galois representations}, Math.\ Ann.\ 373 (2019), 1341--1427 (reference to be confirmed).
\emph{Mismatch.} The published theorems we know are stated under hypotheses on
the residual image (adequacy or irreducibility) which fail here; the formal
statement assumes instead that $x$ itself is automorphic and locally generic,
and restricts to $n\le4$ and to fields without $l$-th roots of unity.  We do not
know whether a published statement covers this case, and we describe the
patching argument behind it only in outline.  It is used for (U2a), (U2b),
Proposition~\ref{i5:x5-prop-sp4core}(e) and Proposition~\ref{i5:x5-prop-Vhost}(a).

\medskip\noindent\emph{(C4) Soul\'e-type vanishing in non-positive weight (L, re-derived, explicit).}
\emph{Statement.} For $e\ge1$, the Bloch--Kato Selmer groups of the
weight-$j$ pieces of $\mathrm{ad}\,r_0$, $j<0$, vanish, i.e.\
$H^1_f(\Q,\overline\Q_l(\psi^{-1}\chi^{-e}))=0$ and its analogues, and the
corresponding integral groups are finite with an explicit bound.  This is the
vanishing of Soul\'e (Theorem~\ref{i5:thm-soule}), with an explicit bound;
here it is derived from Corollary~\ref{i5:cor-lemmaK} (proved in this
section), Poitou--Tate and the local theorems of Theorem~\ref{i5:thm-pt}.  In
weight $0$ the corresponding statement (every weight-$0$ cocycle of $r_0$
satisfying the local conditions is a coboundary) is a class-number statement:
for host frames it is the triviality of the class group of $\Q(i)$, (R).

\medskip\noindent\emph{(C5) Herbrand-type vanishing (L, re-derived) and a Bernoulli-unit criterion (N).}
\emph{Statement.} If $\chi$ is a Dirichlet character with
$B_{1,\chi^{-1}}$ an $l$-adic unit, every everywhere unramified residual
$\bar\chi$-cocycle is a coboundary.  This is the Herbrand direction of the
Herbrand--Ribet theorem (Stickelberger's theorem annihilates the $\chi$-part of
the class group by $B_{1,\chi^{-1}}$; Washington, \emph{Introduction to
Cyclotomic Fields}, Ch.~6; reference to be confirmed), re-derived through a
Bernoulli sum.  It is combined with a new, explicit non-vanishing criterion:
\emph{if $\varphi$ has odd conductor $f$ and order dividing $l^d-1$, and $\xi$ is
primitive of conductor $2^n$ with $f\cdot4^{v_2(l^d-1)}\le2^n$ and $\varphi\xi$
odd, then $|B_{1,\varphi\xi}|_l=1$}, proved by a direct computation with
character sums.  The criterion is in the spirit of Washington's theorem on the
$l$-part of class numbers in the cyclotomic $\Z_2$-extension (reference to
be confirmed), but with an explicit conductor threshold; this threshold is the origin
of the conditions $\mathrm{XData}$ and of the exponents $n_{\rm II},n_{\rm III}$.

\medskip\noindent\emph{(C6) Ribet-type passage from Selmer classes to tangent vectors (L-type idea, (R) execution).}
An extension class of $\Psi$ by $1$ gives a first-order deformation of $r_0$
supported in weight $-1$; restricting to $b$-torsion classes gives an injection
of the Selmer group of $\Psi$ (Definition~\ref{i5:def-V}) into weight-$(-1)$
point derivations at the base point that kill the weight-$(-1)$ relations.
This is (S3) of \S\ref{i5:x5-ssec-arch}; with (M$'$) it gives
$V'(e,\psi)\le\mathrm{cotlen}_{-1}(F)$.  The idea is Ribet's (reference to
be confirmed); the formal execution is bookkeeping with the presentation (C2).

\medskip\noindent\label{i5:x6-sssec-M}\emph{The commutative algebra: (M) and (M$'$).}

\begin{proof}[Proof of Proposition~\ref{i5:x5-prop-UM}\textup{(M$'$)}]
Let $q=\#k_{\mathcal O}$.  Since $x(1)=1$ and $x(A_j)=0$, we have $j\ne0$.  Fix
$\mu\in M$ and a point derivation $\tilde D$ of weight $j$ at $x$ with values in
$\mathcal O$ whose reduction is $D(\mu)$.  (i) $\tilde D$ kills $\mathrm{sat}(x)$:
for $f,g\in\ker x$, $\tilde D(fg)=x(f)\tilde D(g)+x(g)\tilde D(f)=0$, so
$\tilde D((\ker x)^2)=0$; if $cf\in(\ker x)^2$ with $c\ne0$, then
$c\tilde D(f)=0$ in the torsion-free $\mathcal O$, so $\tilde D(f)=0$.  (ii)
$A_j\subseteq\ker x$, so $\tilde D|_{A_j}$ factors through
$\mathrm{Cot}_j=A_j/(A_j\cap\mathrm{sat}(x))$; reducing modulo $c$ and using
$D(\mu)(F_i)=0$, $D(\mu)$ induces an $\mathcal O$-linear map
$\varphi_\mu\colon N\to\mathcal O/c$, $N=\mathrm{Cot}_j/(\mathrm{Cot}_j\cap
\langle F_i\rangle)$.  (iii) $D(\mu)$ is determined by $\varphi_\mu$: it kills
$A_i$ for $i\ne j$ (weight $j$), and $A=\bigoplus_iA_i$.  As $D$ is additive and
injective, $\mu\mapsto\varphi_\mu$ is an injective additive map
$M\to\mathrm{Hom}_{\mathcal O}(N,\mathcal O/c)$.  If $N$ has infinite length
there is nothing to prove.  Otherwise $\#\mathrm{Hom}_{\mathcal O}(N,\mathcal
O/c)\le\#N=q^{\lgth N}$, by induction on $\lgth N$: for $0\to N'\to N\to
k_{\mathcal O}\to0$ the sequence $0\to\mathrm{Hom}(k_{\mathcal O},\mathcal
O/c)\to\mathrm{Hom}(N,\mathcal O/c)\to\mathrm{Hom}(N',\mathcal O/c)$ is exact
and $\mathrm{Hom}(k_{\mathcal O},\mathcal O/c)$, the socle of $\mathcal O/c$,
has at most $q$ elements.  Hence $M$ is finite, so it is an $\mathcal O$-module
of finite length with $\#M=q^{\lgth M}$, and $\lgth_{\mathcal O}M\le\lgth N=
\mathrm{cotlen}_j(F)$.
\end{proof}

\noindent\emph{Structure of the formal proof of (M) (N).}  This is the new
commutative algebra of the section; it is derived together with a calculus of
multigraded Hilbert functions.  The steps are the following.
\begin{enumerate}
\item (R) By (M1) one chooses coordinates so that $F_i\equiv\varpi^{c_i}y_i$
modulo saturated squares for a basis $(y_i)$ of $\mathrm{Cot}_{-1}$ modulo
torsion and the excess equation lies in the squares; then
$\mathrm{cotlen}_{-1}(F)=\sum_ic_i$.
\item (N) One forms a $\Z^3$-graded Rees-type algebra over $A_0$ built from the
equations, and compares Hilbert--Samuel-type multiplicities.  Lower bound: in
large degree the multiplicity of the quotient by the chart kernel and by $a$
is at least $(\sum_ic_i)\,2^{r+r'}$.
\item (R) Additivity of multiplicities in exact sequences with the
non-zero-divisor $a$ moves this to the multiplicity of $A/(I,a)$.
\item (N) Upper bound: the multiplicity of $A/(I,a)$ is at most
$2^{r+r'+1}\sum_{\xi\in T}v(\xi(a))$, the sum over torus-orbit representatives
of two-signed points.  Here (M3) excludes contributions from one-signed points,
(M4) makes each orbit reduced (its tangent space is the orbit direction), so the
local multiplicity at an orbit is $v(\xi(a))$; (M2) and (M5) make the
quotient finite over $\mathcal O$ and $a$ a non-zero-divisor.
\item (R) Dividing by $2^{r+r'}$ gives $\mathrm{cotlen}_{-1}(F)\le
2\sum_{\xi\in T}v(\xi(a))$; the factor $2$ comes from the two signs of the
torus action.
\end{enumerate}
In spirit this is a numerical criterion bounding a tangent length by a
congruence quantity, but it is not Wiles' numerical criterion (no complete
intersection, no $\eta$-invariant) and not a Fitting-ideal argument.  We know
no published statement of this form.  We describe the internal lemmas of steps 2 and 4
only at the level of their statements.

\medskip\noindent\label{i5:x6-sssec-U1}\emph{Proposition~\ref{i5:x5-prop-gdend} and (U1).}

The proof of (U1) from Proposition~\ref{i5:x5-prop-gdend} and
Lemma~\ref{i5:x5-lem-U1steps} is in \S\ref{i5:x5-sssec-archU1}.  The formal
proof of Proposition~\ref{i5:x5-prop-gdend} has the following steps.
\begin{enumerate}
\item[(G1)] (R) Parity and ramification: $\Psi(c_D)=1$, $c_D^2=1$, and every
place at which $\eta$ ramifies divides $lNd_K$; $G_K=\ker\eta$.
\item[(G2)] (N, with (L) inputs re-proved in the formal development) \emph{An auxiliary field and a
frame.}  One constructs a CM field $L\supseteq K$, abelian over $\Q$, and an
abstract deformation problem (``frame'') over $L$ with the following
properties: $l\ne2$; genericity of the residual characters ($\omega\ne
\bar\Psi^{\pm1}$, $\omega\ne\bar\Psi^{\pm2}\delta$, $\bar\Psi(g_0)^2\ne1$,
$\omega$ not a power of $\bar\Psi$); bounds $\le C_0$ for the relevant
$H^1$-ranks and class groups; an Euler-characteristic inequality from the
Greenberg--Wiles formula; and \emph{large $l$-adic degree}: the local degrees
$d_v$ above $l$ satisfy $\sum_vd_v>3C_0+7$ and $d_v>5$.  The cohomological
inputs (Poitou--Tate, Euler characteristic, local duality) are
re-proved in the formal development.  The
role of the large degree is the same as in Skinner--Wiles and Thorne: it gives
Krull-dimension lower bounds for the ordinary deformation space that beat the
dimension of the reducible locus.  The genericity conditions replace the
image hypotheses of the published theorems.
\item[(G3)] (L, re-derived from the bundle) \emph{Hida theory.} For $n=2,3$
there is an ideal $J$ of the universal ordinary pseudo-character ring $B_0$
over the Iwasawa algebra $\Lambda\cong\mathcal O[[X_1,\dots,X_{n[L:\Q]}]]$ such
that $B_0/J$ is finite over $\Lambda$, every minimal prime of $B_0/J$ has
dimension $1+n[L:\Q]$, and the classical points are dense and automorphic.  This
is Hida theory for definite unitary groups (Hida; Geraghty; reference to
be confirmed).  It is derived formally from the bundle (base change to unitary
groups, Galois representations, local Langlands); we do not describe the derivation.
\item[(G4)] (N, with (L) ingredients) \emph{Seeds.}  For each component one
needs one automorphic point congruent to the residual representation.
For $n=2$: a cusp form congruent to an Eisenstein series, constructed from
explicit Eisenstein series, their constant terms and the Deligne--Serre lifting
lemma; this is Ribet's construction (Ribet, Invent.\ Math.\ 34 (1976), 151--162),
re-derived.  For $n=3$: the conjugate self-dual twist of the adjoint (a
Gelbart--Jacquet type lift) of an Eisenstein-congruent two-dimensional
representation, restricted to $L$.  The adjoint lift is derived formally from
the bundle by a pigeonhole argument over the finitely many automorphic
representations of fixed level, together with the descent to unitary groups
(H6.1) and base change, which are
re-proved in the formal development.  We do not describe the internal
congruence-transfer step.
\item[(G5)] (N) \emph{Propagation of automorphy.}  Let $\mathcal S$ be the
ordinary lifting space of the frame.  If the Hida quotient is finite over
$\Lambda$, if modularity propagates downward from ``nice'' type-I primes and
along chains of ``thick'' links of primes, and if one type-I prime of dimension
$1+n[L:\Q]$ lies over the Hida quotient (a seed, (G4)), then \emph{every}
type-I prime of $\mathcal S$ lies over the Hida quotient.  The proof uses the
connectedness of $\mathrm{Spec}$, a Krull-dimension upper bound for the
reducible locus, and local charts and flatness; rank $3$ uses the rank-$2$
statement for the strata of shape $1+2$.  This is the Skinner--Wiles
mechanism (every component meets a modular one through points of
characteristic $l$), here for $\mathrm{GL}_3$ over a CM field with three
residual constituents; it is the main new part of the core, together with the frames, charts and
dimension bounds on which it rests.
\item[(G6)] (R) \emph{Two runs.}  The argument is run twice, with auxiliary
split places $v_0\ne v_0'$ of prime norm at which $\bar\Psi$ is trivial and
$\omega$ is not; the two outputs are glued to get local-global compatibility
at every place.  These places are a neatness device; they are not Steinberg
places, and no Steinberg condition is imposed on the deformations.
\item[(G7)] (L; known theorem, re-proved in the formal development) \emph{Descent.}  Automorphy over $L$ descends to $K$
by cyclic base change of prime degree and its descent (Arthur--Clozel;
(H6.2) of Table~\ref{i5:x6-tab-bundle}), using that $L/K$ is abelian and the
irreducibility over $L$; the local conditions needed for descent (finite image
of inertia away from $l$, ordinary flag at $l$, unramifiedness preserved) are
(R) except one step derived from the linked/unlinked analysis, which we do not describe.
For type-II points the two-dimensional descent from $K$ to $\Q$ is (H6.3),
re-proved in the formal development, and the passage to a newform is derived from the bundle.
\item[(G8)] (N) \emph{Finiteness} $\mathrm{FinEnd}$: with $U=G_L$, every
characteristic-$l$ ordinary point has the residual characteristic polynomial.
This follows from (G3) and (G5): the Hida quotient is finite over $\Lambda$ and
all type-I components lie over it, so the off-triangular primes are contained
in an ideal with Artinian quotient modulo $\lambda$.  It is not a Chevalley
argument and not a lattice (Ribet) argument.
\end{enumerate}
\emph{Lemma~\ref{i5:x5-lem-U1steps}.} (1) (N) $\mathrm{FinEnd}\Rightarrow$ the
weight-zero algebra is finite: a graded finiteness criterion, described in
(V5) of \S\ref{i5:x6-sssec-V}; it uses the published local theorems (re-proved in the formal development).
(2) (N, elementary) a torus-rescaling argument, described in (V6).
(3) (N, elementary) irreducibility on open subgroups from a regular inertia
element, as in (V7).

\emph{What is new, and the comparison.}  The input that is genuinely new is the
combination (G2)+(G5)+(G8): a residually reducible lifting theorem for three
residual constituents with no Steinberg place and no automorphic lift assumed,
whose conclusion is automorphy of \emph{all} type-I points of the ring together
with finiteness of the weight-zero algebra.  For (U1) the condition $\mathrm{GD}$
forces $(l-1)\nmid2e$, hence $l\ge5$.  The comparison with Thorne (2015) and
Allen--Newton--Thorne is the table in \S\ref{i5:x5-ssec}.

\medskip\noindent\label{i5:x6-sssec-U234}\emph{(U2a), (U2b), (U3) and (U3$'$).}

\emph{(U2a).}
(1) (R) The point factors through a finite discrete valuation ring and is
torus-conjugate to an ordinary point (Lemma~\ref{i5:x5-lem-U1steps}(2)); the
frame over $L$ of (G2) is attached to it.
(2) (R) Level and unlinkedness at $l$ from (U3), hence local genericity.
(3) (R) Restriction to $G_L$; automorphy transfers to $L$ by base change
(re-proved in the formal development).
(4) (L, re-derived) (C3) with $n=3$ over $L$: the adjoint Selmer group over $L$
vanishes.
(5) (R) Descent of cocycles from $L$ to $\Q$ (restriction to a finite-index
subgroup is injective on $H^1$ with characteristic-zero coefficients) and the
place-by-place comparison of local conditions with
Definition~\ref{i5:x5-def-BK}, using the published local theorems.

\emph{(U2b).}  (L; known theorem, re-proved in the formal development) A modular $W$ is attached to a newform, and its
restriction to $G_K$ is automorphic by quadratic base change ((H1.8), (H1.9));
then (C3) with $n=2$.  The off-diagonal pieces of $\mathrm{Ad}\,r$ at a type-II
point ($W$ and $W\otimes\theta$) are handled by
Proposition~\ref{i5:x-prop:katotwow} (Kato's Euler system), proved in this
section.

\emph{(U3).}  (R) The level, the shape at $l$ and the automorphy witness are
read off the local conditions of the point; the unlinkedness at $l$ comes from
the linked/unlinked analysis of the same frame machinery as (G5), which we
describe only at its top level.  \emph{(U3$'$)} (R): block form $\mathrm{diag}(W,1)$ and the
local conditions.

\medskip\noindent\label{i5:x6-sssec-U4}\emph{The lemmas behind (U4).}

The assembly of (U4) from (S1)--(S4) is proved in \S\ref{i5:x5-ssec-arch}.  The
four steps:
(S1) graded model: (C2), with (C1).
(S2) hypotheses of (M): spanning in weights $\le0$ from (C4); (M2) from (F) of
(U1); (M3) (R): a local point which is both upper and lower triangular is the
base point; (M4) from (U2a) at type-I points and from (U2b) together with
Proposition~\ref{i5:x-prop:katotwow} at type-II points; (M5) is the hypothesis.
(S3) (C6) and (M$'$), proved in \S\ref{i5:x6-sssec-M}.
(S4) (M), (N).
So the new content of (U4) is (M) and the equivariant bookkeeping of (C2); the
remaining inputs are (C3), (C4), Kato's theorem and the published local and
global Galois cohomology.

\medskip\noindent\label{i5:x6-sssec-sp4}\emph{Proposition~\ref{i5:x5-prop-sp4core}: the $\mathrm{GSp}_4$ cores.}

The residual representation is $\bar\vartheta\oplus\bar\Psi_D\oplus\bar\alpha
\oplus1$ with Hodge--Tate weights $(e+1,e,1,0)$; the assembly of
$\mathrm{Count}_{\rm II}$, $\mathrm{Count}_{\rm III}$ from (a)--(g) is the proof
of Proposition~\ref{i5:prop-c23}.

\emph{(a) Automorphy and finiteness.}  The method is a residually reducible
potential automorphy argument for $\mathrm{GL}_4$ over CM fields, of the same
type as (G2)--(G8); it is not a reduction to the three-dimensional core, not a
lattice argument, and not Hida theory for $\mathrm{GSp}_4$.
(a1) (L, re-derived) Hida theory for unitary groups over a tower of CM fields
(as (G3)), and a propagation statement for the universal ordinary
pseudo-character ring (a prime containing a nice automorphic point contains the
Hecke ideal; the ideal propagates along chains of thick primes), built on the
patching modules of \cite{CharacterFamily}, whose internals we do not describe.
(a2) (N) Seeds.  Case II: base change and automorphic induction of a Hilbert
Eisenstein series of weight one, a quadratic real-to-complex lift and cyclic
base change.  Case III: $\rho\otimes\tau$ with $\rho$ an automorphic
two-dimensional representation over $\Q$ and $\tau$ induced from a character of
an auxiliary imaginary quadratic field, so that automorphy reduces to base
change and automorphic induction (re-proved in the formal development) and does not use a tensor product
functoriality; a Ribet lattice and auxiliary weight-four characters are used.
We describe this step only in outline.
(a3) (N) Connectedness with explicit constants: once the degree $d$ of the
tower satisfies $30(C+1)^2\le d$, the reducible loci have Krull dimension at
most $d+20C+1$ and tangent rank bounded likewise, so components join through
irreducible points (Skinner--Wiles and Thorne dimension counting, re-derived
with explicit numbers).
(a4) (N) $\mathrm{Fin}_0$ and $\mathrm{NoEnd}$: points of the bad kind cannot be
pro-modular, because pro-modular points are pure.  We describe this step at the
level of the assembly only.
(a5) (L; known theorem, re-proved in the formal development) Descent to $\Q$ and to CM quadratic fields by
Arthur--Clozel base change ((H6.2), (H6.3), (H1.8)).

\emph{(b) Existence of a good datum} (N): an odd primitive character $\xi$ of
$2$-power conductor $2^{n'}$, $n'\ge c_2(\psi')+6$, with the vanishing
conditions supplied by (C5).
\emph{(c)} (R): unfolding of $\operatorname{sha}(D)$ and the comparison of
Selmer lengths (C6).
\emph{(d)} (R): divisibility bookkeeping.
\emph{(e) Rigidity} (L, re-derived): (C3) with $n=4$ over an auxiliary CM
quadratic field and a twist making the restriction polarized; the tangent vector
becomes a Bloch--Kato class, which is a coboundary, projected back to
$\mathfrak{sp}_4$.
\emph{(f) Length estimate} (N): the $\mathrm{GSp}_4$ form of (C2), with weights
in $\Z^2$ graded through the first coordinate, (M), (M$'$), and (C4) for
$\psi^{-1}\chi^{-e}$.
\emph{(g) Counting}: the count of cuspidal cohomological eigensystems of level
$M$ and weight $(e+1,e,1,0)$ is $(\mathrm C_\Q)$ of
\S\ref{i5:x6-sssec-count}; the Ramanujan-type bound for regular algebraic
cuspidal representations (from purity, re-proved in the formal development) gives
$|t|\le4p^{(e+1)/2}\sqrt p$, hence $t\ne\mathrm{eis}_4$ and
$|t-\mathrm{eis}_4|\le p^{2(e+1)}$; the norm argument is that of
Proposition~\ref{i5:x-prop:normbound} in rank four (R).

New, in this core: (a2)--(a4) (a residually reducible automorphy theorem for
four residual characters of $\mathrm{GSp}_4$, with no Steinberg place), the
explicit Bernoulli criterion of (C5), and (f).

\medskip\noindent\label{i5:x6-sssec-V}\emph{Propositions~\ref{i5:x5-prop-hostend}, \ref{i5:x5-prop-Vcore} and Lemma~\ref{i5:x5-lem-Vopen}.}

\emph{Proposition~\ref{i5:x5-prop-hostend}.}
(V1) (R) $\sigma^2=1$, $\Psi(\sigma)=1$; the places where $\eta_4$ ramifies
divide $lN$; $G_{\Q(i)}=\ker\eta_4$.
(V2) (N) A frame over an auxiliary abelian CM field $L$, as in (G2).  The
host-specific point is how the genericity conditions are met for $l=3$: the
condition $c_2(\psi)\ge4$ of $\mathrm{HostData}$ gives an element of $2$-adic
inertia on which $\omega$ and $\psi$ are trivial but $\eta_4=-1$, and an element
$s$ of the decomposition group at $l$ with $\eta(s)=1$ and $\bar\Psi(s)^2\ne1$.
Thus the non-quadraticity of $\bar\Psi$ comes from the dyadic part of $\psi$,
not from $\omega$, which is what removes the usual restriction $l\ge5$.
(V3) (N) (G3)--(G6) applied to this frame.  (V4) (R) finiteness from (G8) with
$U=G_L$.  (V4$'$) (L; known theorem, re-proved in the formal development) descent from $L$ to $\Q(i)$, as (G7).
We have not verified that every genericity condition inside (G2)--(G6) holds for
$l=3$; the parts described here need only $l$ odd.

\emph{Proposition~\ref{i5:x5-prop-Vcore}} (the proof from
Proposition~\ref{i5:x5-prop-hostend} and Lemma~\ref{i5:x5-lem-Vsteps} is in
\S\ref{i5:x5-ssec-arch}).  The three parts of Lemma~\ref{i5:x5-lem-Vsteps}:
(1) (R) restriction of the frame conditions to decomposition groups.
(2) (N, elementary) \emph{Characteristic-$l$ points are triangular.}  If $A$ is a
local domain of characteristic $l$, $r$ normalised and marked, and on a
finite-index subgroup the characteristic polynomials are those of $r_0$, then
they are those of $r_0$ everywhere on $\ker\eta$ (pass from $g^n$ to $g$), and
over $\mathrm{Frac}(A)$ the shape analysis of Lemma~\ref{i5:x5-lem-shape}
excludes type~I (the characteristic polynomials split into characters), so $r$
is upper or lower triangular.  Only $l$ odd is used.
(V5) (N) \emph{Graded finiteness criterion} (Lemma~\ref{i5:x5-lem-Vsteps}(3), also
Lemma~\ref{i5:x5-lem-U1steps}(1)): \emph{let $R=\mathcal O[[z_i]]/I$ with
homogeneous $I$ and weighted variables; if at every point of $R$ with values in
a complete discrete valuation ring of characteristic $l$ the weight-zero
variables vanish and either all negative-weight or all positive-weight
variables vanish, then $R_0$ is finite over $\mathcal O$.}  On a triangular
point the diagonal characters equal the residual ones (by the class-number-one
statement), so the weight-zero variables vanish.  Proof of the criterion:
every weight-zero monomial lies in $\sqrt{\varpi R}$ (if a prime containing
$\varpi$ missed it, a dimension-one prime above it, normalised by
Krull--Akizuki, would give a characteristic-$l$ point at which a mixed-sign
monomial does not vanish, which is impossible); by homogeneity this contracts to
$R_0$, so $R_0/\varpi$ is finite, and $R_0$ is finite by the complete Nakayama
lemma (Matsumura, \emph{Commutative Ring Theory}, Thm.~8.4).  The tools are
standard; the criterion, in this graded form, is new as far as we know (its
spirit is that of the finiteness of residually reducible pseudo-deformation
rings).

\emph{Lemma~\ref{i5:x5-lem-Vopen}.}
(V6) (N, elementary) \emph{Torus rescaling.}  If $R_0$ is finite, every
$\overline\Q_l$-point $\xi$ of the graded ring has $|\xi(f)|<1$ on $R_0\cap
\mathfrak m$; the mixed products $z_i^{|w_h|}z_h^{|w_i|}$ lie there, which
gives $t\in\overline\Q_l^\times$ with $|t^{w_i}\xi(z_i)|<1$ for all $i$; the
substitution $z_i\mapsto t^{w_i}\xi(z_i)$ kills the homogeneous $I$ and defines
a point of $R$ with values in a finite discrete valuation ring, equal to
$\xi$ up to $\operatorname{tor}(t)$.
(V7) (N, elementary) \emph{Irreducibility on open subgroups.}  An element $u$
of inertia at $l$, a power of an element with cyclotomic value $1+l$ lying in
the normal core of the open subgroup, has three distinct eigenvalues
$\Psi(u),1,\Psi(u)^{-1}$ by the ordinary condition; if $r$ were irreducible on
$\ker\eta$ but reducible on the open subgroup, the invariant subspaces would be
coordinate subspaces and one finds $g$ with $\operatorname{tr}r(gx^n)=0$ for all
$n$, $x$ upper triangular at $l$; reducing gives
$\bar\Psi(gx^n)+1+\bar\Psi(gx^n)^{-1}=0$ for all $n$, which forces
$\bar\Psi(x)^2=1$, a contradiction.  Valid for $l=3$.

New in this core: (V2) (the dyadic mechanism replacing $l\ge5$) and (V5)--(V7),
on top of the shared (G3)--(G7).

\medskip\noindent\label{i5:x6-sssec-Vhost}\emph{Proposition~\ref{i5:x5-prop-Vhost}.}

\emph{(a) Rigidity.} (R) Transport of a characteristic-zero point to the
ordinary point over $\Q(i)$ (V6); local genericity from unlinkedness; (L,
re-derived) (C3) with $n=3$, the same statement as in (U2a); (R) the
Bloch--Kato conditions place by place (at $l$ by the published local theorems,
re-proved in the formal development); (R) averaging over the index-two subgroup $G_{\Q(i)}$ to pass from
$G_{\Q(i)}$ to $G_\Q$ ($2$ is invertible).

\emph{(b) Number.} (R) Choose a prime $p\equiv3\bmod4$ with $7\le p\le12lN$ and
$p\nmid lN$ (Dirichlet; proved), so that $p$ is inert in $\Q(i)$ with norm
$p^2$; take $\varphi=\sigma^{-1}$ for a Frobenius $\sigma$ at a place above
$p$; identify $\operatorname{tr}\operatorname{tor}\Psi(\varphi)$ with
$\mathrm{eis}_3(p^2,e,z_0)$.  (L/N) Apply $\mathrm{NumberU3U}$
(Proposition~\ref{i5:x-prop:numberu3u}, resting on
Proposition~\ref{i5:x-prop:u3count}) after showing that $a_r|_{G_{\Q(i)}}$ is
stably admissible for every characteristic-zero type-I point with
$\mathrm{B}'$ (irreducibility from (V7), automorphy from $\mathrm{B}'$, level
from the frame conditions, unlinkedness from the frame analysis; we do not
describe the last two).  (R) Sum over the two embeddings of $\Q(i)$ and collect
constants.

\emph{(c) Existence.} (R) $\mathrm{HostData}$ and $\mathrm{XData}$ from the
arithmetic of hosts (the inequality $2lq_h4^{n_0}\le R_1\le2^{c_2}$); (R) the
frame ($S$ the primes dividing $lN$, places, a complex conjugation); (L,
re-derived) the universal ring, (C1), with $\Phi_l$ from global finiteness
(re-proved in the formal development).

\emph{(d) Abstract bound.} (N) The frame version of (U4): a weighted
presentation of the universal frame ring (C2); vanishing of tangent
coordinates of non-positive weight (C4); two-signed points are of type~I, using
the class-number-one statement and the Herbrand-type vanishing (C5) (this is
where $\mathrm{XData}$ enters); at rigid points the tangent space is the orbit
direction; then the cotangent bound (M) and the Selmer comparison (C6), (M$'$).
The final comparison of the frame Selmer length with $V'(e,\psi)$ is (R).

\medskip\noindent\label{i5:x6-sssec-count}\emph{Proposition~\ref{i5:x-prop:u3count}: the counting statements and the published inputs.}

\emph{$(\mathrm C_K)$.}
(K1) (L, re-derived) For $m$ with $|m_i|\le H$ there is an irreducible
algebraic representation $\Xi$ of $\mathrm{Res}_{K/\Q}\mathrm{GL}_n$ with the
matching infinitesimal character and $\dim\Xi\le((n+1)^2(H+1))^{n^2[K:\Q]}$
(Weyl's dimension formula; Knapp, \emph{Lie Groups Beyond an Introduction},
Thm.~5.84; reference to be confirmed).
(K2) (L; known theorem, re-proved in the formal development) Each regular algebraic cuspidal $\pi$ of that weight is
$\Xi$-cohomological (H1.1).
(K3) (L; known theorem, re-proved in the formal development) The number of pairwise non-isomorphic such $\pi$ with
$K(M)$-fixed vectors is at most $\dim H^*(Y_{K(3M)},\Xi)$ (H1.3).
(K4) (L/N, re-derived) $\dim H^*(\Gamma,\Xi)\le c\,M^A\dim\Xi$ for congruence
subgroups $\Gamma$ of level $M$ of $\mathrm{GL}_n(\mathcal O_F)$, uniformly in
$M$: finiteness of adelic double cosets, the index bound
$[\Gamma(1):\Gamma(3M)]\le(3M)^{[F:\Q]n^2}$, Shapiro's lemma, and, for each group
commensurable with $\mathrm{GL}_n(\mathcal O_F)$, a bounded exact complex of
permutation modules with finite stabilisers (from an acyclic simplicial complex
of building type; for $F$ totally real by passing to a CM extension and taking
fixed points of complex conjugation).  Qualitatively this is the fact that
these groups are virtually of type FL (Borel--Serre; reference to be confirmed); the
explicit polynomial dependence on the level is the new part.
(K5) (L; known theorem, re-proved in the formal development) Satake parameters at $w\nmid M$ (local components
and the Satake isomorphism); stability of the set under $\mathrm{Aut}(\C)$ by
Clozel's rationality (H1.2), with conjugate self-duality from Gelfand--Kazhdan
and strong multiplicity one (H3.7), (H3.4); purity (H3.9); and integrality, whose proof we do not
describe.
(K6) (L) Two representations with the same Frobenius characteristic
polynomials almost everywhere are isomorphic: Chebotarev (Neukirch VII (13.4))
and Brauer--Nesbitt (reference to be confirmed); re-derived.
(K7) (L; known theorem, re-proved in the formal development) For admissible $\rho$ attached to $\pi$: $\pi$ has the right
weight and level and the Frobenius eigenvalues are the Satake parameters, by
local-global compatibility (H3.8) and Hodge--Tate weights; (R) assembly.

\emph{$(\mathrm C_\Q)$} (N, an elementary device): take $n^2+1$ imaginary quadratic
fields unramified at $p$ (Dirichlet primes $\equiv3\bmod4$); an irreducible
$n$-dimensional $r$ is isomorphic to its twist by at most $n^2$ quadratic
characters, so $r|_{G_{K_i}}$ is irreducible for some $i$, hence admissible
over $K_i$; apply $(\mathrm C_{K_i})$ at a place above $p$ of residue degree
$f_i$ and take $f_i$-th roots of the eigenvalues.

\emph{$(\mathrm N)$} (R, with (L) inputs re-proved in the formal development): newform to cuspidal
representation of $\mathrm{GL}_2(\mathbb A_\Q)$ (H1.9), quadratic base change
(H1.8), duality of two-dimensional representations up to the determinant, and
a weight-gap check.

\emph{The published inputs.}  Table~\ref{i5:x6-tab-bundle} lists the $65$
statements of the bundle (H1)--(H5) and the five further statements (H6) used
by the cores of \S\ref{i5:x5-ssec}.  The statements are stated informally (the
exact statements are in \cite{CharacterFamily}); the attributions are ours, and
``reference to be confirmed'' marks a reference we have not confirmed.  Unless
stated otherwise local statements are for all non-archimedean local fields of
characteristic $0$, and global Galois-cohomological statements for finite
$\Z/l^r[G]$-modules over number fields.

{\small
\begin{longtable}{p{0.07\textwidth}p{0.47\textwidth}p{0.38\textwidth}}
\caption{The published statements isolated and proved as separate propositions in \cite{CharacterFamily}.}\label{i5:x6-tab-bundle}\\
\toprule
 & statement (stated informally) & usual attribution; reference\\
\midrule
\endfirsthead
\toprule
 & statement & attribution; reference\\
\midrule
\endhead
\multicolumn{3}{l}{(H1) automorphic representations (9)}\\
H1.1 & regular algebraic cuspidal $\pi$ over a totally complex field is cohomological for the $\Xi$ with opposite infinitesimal character & Clozel 1990, Lemme 3.14 (with Borel--Wallach); reference to be confirmed\\
H1.2 & $\mathrm{Aut}(\C)$-conjugate of a regular algebraic cuspidal $\pi$ is cuspidal with conjugated infinity type & Clozel 1990, Thm.~3.13 (rationality); reference to be confirmed\\
H1.3 & pairwise non-isomorphic $\Xi$-cohomological cuspidal $\pi$ with $U$-fixed vectors number at most $\dim H^*(Y_U,\Xi)$ & Borel 1983 (cuspidal cohomology injects); reference to be confirmed\\
H1.4 & Galois representations for regular algebraic polarized cuspidal $\mathrm{GL}_2$ over totally real fields, with local-global compatibility, Hodge--Tate weights, crystalline at spherical places above $l$ & Eichler--Shimura, Deligne, Carayol, Taylor, Blasius--Rogawski; crystallinity T.~Saito; reference to be confirmed\\
H1.5 & strong multiplicity one for cuspidal $\mathrm{GL}_a$ over totally real $F$ & Jacquet--Shalika II, Thm.~4.4; Piatetski-Shapiro\\
H1.6 & invariant pairing $\pi\times\pi$ twisted by transpose-inverse, totally real $F$ & Gelfand--Kazhdan (reference to be confirmed)\\
H1.7 & admissibility of cuspidal representations (finite-dimensional isotypic fixed vectors), totally real $F$ & Harish-Chandra, Borel--Jacquet; reference to be confirmed\\
H1.8 & quadratic base change for regular algebraic cuspidal $\mathrm{GL}_2$ from totally real $F$ to a CM quadratic $E$ (or $\pi\cong\pi\otimes\eta_{E/F}$) & Langlands 1980, Arthur--Clozel 1989 (Ch.~3; reference to be confirmed)\\
H1.9 & a newform of weight $k\ge2$ gives a cuspidal $\pi$ of $\mathrm{GL}_2(\mathbb A_\Q)$ with matching Hecke polynomials & adelisation of newforms (Jacquet--Langlands, Gelbart); reference to be confirmed\\
\midrule
\multicolumn{3}{l}{(H2) local representation theory (18)}\\
H2.1 & irreducible smooth is admissible; Jacquet functors and induction preserve finite length; induction commutes with contragredient & Bernstein--Zelevinsky 1976/77 (reference to be confirmed)\\
H2.2 & product of irreducible unitarisable representations is irreducible & Bernstein 1984; reference to be confirmed\\
H2.3 & geometric lemma (filtration of Jacquet modules of induced representations) & Bernstein--Zelevinsky 1977, Lemma 2.12\\
H2.4 & $\bigoplus_nR(\mathrm{GL}_n)$ is the polynomial ring on segments & Zelevinsky 1980, Cor.~7.5\\
H2.5 & products with totally disjoint (unlinked up to $|\det|^{\pm1}$) cuspidal supports are irreducible & Bernstein--Zelevinsky / Zelevinsky 1980, Thm.~4.2 \\
H2.6 & Jacquet's lemma: $K(1)(N\cap K_0)$-fixed vectors map bijectively to Levi-fixed vectors of the Jacquet module & Jacquet, Casselman; Moy--Prasad; reference to be confirmed\\
H2.7 & tempered iff product of square-integrable & Jacquet 1977, Zelevinsky; reference to be confirmed\\
H2.8 & essentially square-integrable iff $L(\Delta)$ for a segment & Bernstein--Zelevinsky, Zelevinsky 1980 Thm.~9.3 (Bernstein)\\
H2.9 & square-integrable representations are attached to segments & same; reference to be confirmed\\
H2.10 & generic irreducible is a product of segment representations with unlinked segments & Zelevinsky 1980, Thm.~9.7 (non-degenerate $=$ generic)\\
H2.11 & uniqueness of the multisegment; $L(m+n)$ is a subquotient of $L(m)\times L(n)$ & Langlands classification (Zelevinsky, Silberger); reference to be confirmed\\
H2.12 & every irreducible is $L(m)$ & Langlands classification; reference to be confirmed\\
H2.13 & standard modules have irreducible cosocle & Langlands quotient theorem (Silberger); reference to be confirmed\\
H2.14 & local converse theorem for $\mathrm{GL}_{m+1}$ with twists by $\mathrm{GL}_m$ & Henniart, Invent.\ Math.\ 113 (1993), 339--350; reference to be confirmed\\
H2.15 & depth zero is preserved by parabolic induction & Moy--Prasad; reference to be confirmed\\
H2.16 & Brauer's induction theorem in degree-zero form & Brauer; Serre, \emph{Linear Representations}; reference to be confirmed\\
H2.17 & units $\equiv1\bmod a$ are $m$-th powers for suitable $a$ prime to $b$ & Chevalley, J.\ Math.\ Soc.\ Japan 3 (1951), 36--44, Th\'eor\`eme~1\\
H2.18 & cuspidal characters of $\mathrm{GL}_d(\mathbb F_g)$: degree, values on elliptic elements, parametrisation by regular characters & Green 1955 (reference to be confirmed)\\
\midrule
\multicolumn{3}{l}{(H3) global automorphic forms, Galois representations, class field theory (15)}\\
H3.1 & rapid decay of cusp forms of uniform moderate growth & Gelfand--Piatetski-Shapiro, Harish-Chandra; Moeglin--Waldspurger I.2.18 (reference to be confirmed)\\
H3.2 & finiteness of automorphic forms of fixed level, $K_\infty$-type and $Z(\mathfrak g)$-character & Harish-Chandra 1968; reference to be confirmed\\
H3.3 & the same for $\mathrm{GL}_a$ over totally complex fields & Harish-Chandra; reference to be confirmed\\
H3.4 & rigidity of isobaric sums; bound $q^{-1/2}<|\alpha|<q^{1/2}$ for unitary cuspidal Satake parameters & JS I Cor.~2.5 and JS II Thm.~4.4\\
H3.5 & discrete decomposition of the cuspidal spectrum; cusp forms bounded & Gelfand--Graev--Piatetski-Shapiro, Godement; reference to be confirmed\\
H3.6 & base change from unitary groups $U(n)$ to $\mathrm{GL}_n$ over CM $E$ with $E/E^+$ unramified & Labesse 2011, Mok, Kaletha--Minguez--Shin--White; reference to be confirmed\\
H3.7 & Gelfand--Kazhdan duality over totally complex fields & Gelfand--Kazhdan (reference to be confirmed)\\
H3.8 & Galois representations for regular algebraic polarized cuspidal $\mathrm{GL}_a$ over CM fields with local-global compatibility at all finite places & Harris--Taylor, Shin, Chenevier--Harris, Caraiani, BLGGT; Theorem~\ref{i5:thm-blggt} and Remark~\ref{i5:rem-blggt-src}\\
H3.9 & purity of algebraic cuspidal representations & Clozel 1990, Lemme 4.9; reference to be confirmed\\
H3.10 & regular algebraic implies cohomological & Clozel 1990, Lemme 3.14; reference to be confirmed\\
H3.11 & global Artin reciprocity & Artin, Tate; Neukirch VI (5.5)\\
H3.12 & ray class fields & Takagi, Artin; Neukirch VI; reference to be confirmed\\
H3.13 & $l$-adic characters of algebraic Hecke characters are de Rham and potentially crystalline above $l$ & Weil, Serre (\emph{Abelian $l$-adic representations}); reference to be confirmed\\
H3.14 & Weil--Deligne representation of that character at $w$ is the local component & Serre, Fontaine; reference to be confirmed\\
H3.15 & de Rham implies potentially semistable & Berger 2002 (Andr\'e, Kedlaya, Mebkhout); reference to be confirmed\\
\midrule
\multicolumn{3}{l}{(H4) Galois cohomology (12)}\\
H4.1 & normalised local invariants; Poitou--Tate nine-term sequence; $\Sha$-duality; $H^n=0$, $n\ge3$ ($l$ odd or $F$ totally complex) & Poitou, Tate; Theorem~\ref{i5:thm-pt}; reference to be confirmed\\
H4.2 & finiteness of local cohomology & Tate; reference to be confirmed\\
H4.3 & local Tate duality & Tate; NSW Thm.~7.2.6\\
H4.4 & local Euler characteristic formula & Tate; NSW Thm.~7.3.1\\
H4.5 & finiteness of $H^i(G_S,M)$ & Tate; reference to be confirmed\\
H4.6 & global Euler characteristic formula & Tate; reference to be confirmed\\
H4.7 & unramified classes are exact annihilators of each other & Tate; Milne, \emph{ADT} I.2.6; NSW 7.2.15\\
H4.8 & Greenberg--Wiles formula & Wiles 1995, Greenberg; reference to be confirmed\\
H4.9 & finiteness over finite extensions of $\Q_l$ & Tate; reference to be confirmed\\
H4.10 & local Tate duality over finite extensions of $\Q_l$ & Tate; NSW (note)\\
H4.11 & local Euler characteristic over finite extensions of $\Q_l$ & Tate; NSW (note)\\
H4.12 & local Bloch--Kato dimension formula for $H^1_f$ and duality & Bloch--Kato 1990, Cor.~3.8.4 (reference to be confirmed)\\
\midrule
\multicolumn{3}{l}{(H5) CM abelian varieties and special values (11)}\\
H5.1 & deformation theory of abelian schemes over nilpotent PD-thickenings & Grothendieck--Messing 1972; reference to be confirmed\\
H5.2 & de Rham cohomology of abelian schemes: freeness, Hodge filtration, base change, duality & Oda 1969, Berthelot--Breen--Messing; reference to be confirmed\\
H5.3 & dual abelian scheme, Poincar\'e bundle; crystalline equals de Rham over PD-thickenings & Grothendieck, Raynaud, Berthelot; reference to be confirmed\\
H5.4 & period lattices span $H_{\rm dR}$; polarisation pairing is $\pm2\pi i$ times the Riemann form & Riemann relations, Deligne; reference to be confirmed\\
H5.5 & periods of smooth families vary holomorphically (Gauss--Manin) & Katz--Oda 1968; reference to be confirmed\\
H5.6 & N\'eron--Ogg--Shafarevich; CM abelian varieties have potential good reduction & Serre--Tate 1968; reference to be confirmed\\
H5.7 & main theorem of complex multiplication for elliptic curves & Shimura 1971, Thm.~5.4; reference to be confirmed\\
H5.8 & $K(j(\mathcal O_K))$ is the Hilbert class field and related facts & Weber, Hasse, Deuring; Neukirch VI \S6\\
H5.9 & rationality and $l$-integrality of partial zeta values at $1-k$ & Siegel--Klingen; Deligne--Ribet 1980; reference to be confirmed\\
H5.10 & functional equation of partial zeta functions & Hecke\\
H5.11 & functional equation of Hecke $L$-functions with CM infinity type, explicit $\varepsilon$ & Hecke, Tate 1950 \\
\midrule
\multicolumn{3}{l}{(H6) further statements used by the cores (5)}\\
H6.1 & weak descent of conjugate self-dual cuspidal $\mathrm{GL}_n$, $2\le n\le4$, to $U(n)$ (agreement at almost all split places) & Labesse 2011, Mok, Kaletha--Minguez--Shin--White; reference to be confirmed\\
H6.2 & cyclic base change of prime degree for regular cuspidal $\mathrm{GL}_n$ between totally complex fields: existence, descent of invariant $\pi$, fibres & Arthur--Clozel 1989 (reference to be confirmed)\\
H6.3 & descent of $\mathrm{GL}_2$ from imaginary quadratic fields to $\Q$, or automorphy over $\Q$ of odd $W$ automorphic over an abelian CM field & Langlands, Arthur--Clozel solvable descent; reference to be confirmed\\
H6.4 & existence of a totally definite quaternion algebra over totally real $F$ of even degree, unramified at finite places & Albert--Brauer--Hasse--Noether; reference to be confirmed\\
H6.5 & cusp forms which are eigenvectors of a fixed Lie algebra element lie in the span of cuspidal subrepresentations & Gelfand--Graev--Piatetski-Shapiro; reference to be confirmed\\
\bottomrule
\end{longtable}
}

The informal statements of H4.2, H6.4 and of part of H3.13 and H3.8 are less precise than the others.



\medskip\noindent\textbf{Status.}
\emph{Proved here:} Theorem~\ref{i5:thm-vwin} from the inputs listed in its statement; the reductions of
\S\S\ref{i5:ssec-assembly}--\ref{i5:ssec-bab}; Propositions~\ref{i5:prop-shak}, \ref{i5:prop-psibound}, \ref{i5:prop-stick},
\ref{i5:prop-thmbab} and the Case~I count (Propositions~\ref{i5:prop-corg}, \ref{i5:prop-countI}); the reductions of
$\mathrm{SFew}$, $\mathrm{OneCharU}$, $\mathrm{Count}_{\rm II}$, $\mathrm{Count}_{\rm III}$ and $\mathrm{Count}_{\rm V}$ to their cores
(\S\ref{i5:x5-ssec}); the assembly steps of \S\ref{i5:x5-ssec-arch}; and part (M$'$) of Proposition~\ref{i5:x5-prop-UM}.
\emph{Published theorems, quoted:} Theorems~\ref{i5:thm-llc}, \ref{i5:thm-abps}, \ref{i5:thm-lz}, \ref{i5:thm-blggt},
\ref{i5:thm-pt}, \ref{i5:thm-soule}, \ref{i5:thm-fw} and \ref{i5:x-thm:kato}, Stickelberger's theorem, Chebotarev's density
theorem, and the classical results cited in \S\ref{i5:ssec-pf-shak} and \S\ref{i5:x5d-ssec}, as listed in
Theorem~\ref{i5:thm-vwin}(i).
\emph{Formally verified, \cite{CharacterFamily}:} the cores listed in Theorem~\ref{i5:thm-vwin}(ii), namely
Propositions~\ref{i5:x5-prop-U1}, \ref{i5:x5-prop-U234}, \ref{i5:x5-prop-U4}, \ref{i5:x5-prop-sp4core},
\ref{i5:x5-prop-Vcore}, \ref{i5:x5-prop-Vhost}, Lemma~\ref{i5:x5-lem-Vopen} and the counting statements of
Proposition~\ref{i5:x-prop:u3count}, together with the statements they are built from in \S\ref{i5:x5-ssec-arch}
(Propositions~\ref{i5:x5-prop-gdend}, \ref{i5:x5-prop-hostend}, part (M) of Proposition~\ref{i5:x5-prop-UM}, and
Lemmas~\ref{i5:x5-lem-U1steps}, \ref{i5:x5-lem-Vsteps}).  The formal development also proves the published statements of
Table~\ref{i5:x6-tab-bundle} on which these cores rest.




\section{Existence of families with uniform constants}\label{i4:sec}

This section proves Theorem~\ref{M:thm:alpoge-families} from five inputs.
Four of them are written elsewhere and are used here as hypotheses, with
their exact statements: the loss bound, the divisor bound, the level counts
and Theorem~S, and the windowed Selmer bound. The fifth is the existence
argument itself. Its logical part fixes the constants in the right order
and is proved here in full. Its arithmetic part builds hosts, columns and
bundles of automorphic origin. Where that construction is too long for this
section, or where we do not give a written proof, we state the result
and label it ``formally verified''. The arithmetic part also rests
on a large number of published theorems, most of which are not stated in
this paper. Everything that is not proved here, formally verified statements
and published inputs alike, is listed in Remark~\ref{i4:rem-trust} and summarised
at the end of the section. For each formally verified statement the
structure of its formal proof is described in
\S\ref{i4:ssec-outline}, with every step classified as a
published theorem, a routine step or a new statement.

The main results are Theorem~\ref{i4:thm-top} and Theorem~\ref{i4:thm-final}, which
gives Theorem~\ref{M:thm:alpoge-families}.  The section is organised as follows.
Subsections~\ref{i4:ssec-defs}--\ref{i4:ssec-arith} contain the definitions and statements: hosts,
families and connected classes; the four inputs, which are used as hypotheses
(Propositions~\ref{i4:hyp-loss}, \ref{i4:hyp-div}, \ref{i4:hyp-S} and \ref{i4:hyp-selmer}, proved in
Sections~\ref{i1:sec}, \ref{i2:sec}, \ref{i3:sec} and \ref{i5:sec}); thresholds and closed properties;
connectedness from design existence; design existence; the closed estimates on the working
class; the hypothesis structure; and the arithmetic construction.
Subsection~\ref{i4:ssec-top} proves the two main theorems from these statements.
Subsections~\ref{i4:ssec-grp-hosts}--\ref{i4:ssec-grp-norm} contain the proofs of the statements:
hosts and kits, the supply constructions, the division side, realisation, the automorphic
construction, and norm primes in short intervals.  Subsection~\ref{i4:ssec-outline} outlines the
formally verified steps, and the section ends with a summary of its status.

\medskip
\noindent\emph{Conventions.} Throughout, $\psi=1$ is the trivial character.
We record the simplifications it brings.
\begin{enumerate}
\item The level attached to $\psi$ keeps only the trivial character. Its
weighted prime sum is $G$ itself, and there is no exceptional character.
\item The host group is $U=E_0\times A$ with $E_0=(\Z/2)^2$ and $A$ finite
abelian of order $n$. The twist half of the requirement $\mathcal G_1$
(Definition~\ref{i4:def-G1}) contains the clause $o=1$ for the integer
$o=|A|/\mathrm{ord}(a_0)$ of the frame. So the translation
subgroup $A_0$ is all of $A$, and $A$ is cyclic, generated by $a_0$, of
$2$-power order.
\item The characters of $U$ that are trivial on $A$ are the four characters
of $E_0$. The trivial one is kept. The other three are $\chi_q$, $\chi_d$
and $\chi_q\chi_d$ (through the Artin map, by the normalisation
$s_1\circ\mathrm{Art}=(\frac\cdot q)$, $s_2\circ\mathrm{Art}=(\frac\cdot d)$
of Definition~\ref{i4:def-host}), and these are the characters whose prime
sums have to be bounded.
\item The equidistribution condition on the top set is vacuous. Its target
share is $4/(3o+1)=1$.
\end{enumerate}
We write $u=x^\varepsilon$ and $b=\exp(x^\varepsilon)$. A statement holds
\emph{eventually} in a family if it holds for all $x\ge x_0$ and all
indices $i\in I(x)$, for some $x_0$.

\subsection{Hosts, families and connected classes}\label{i4:ssec-defs}

\begin{definition}[host]\label{i4:def-host}
Let $\varepsilon\in(0,1/52)$ and $x\ge1$. A \emph{host} at scale $x$
consists of the following data.
\begin{itemize}
\item A finite abelian group $A$ of order $n$.
\item Primes $q\equiv3\bmod8$ with $q\ne3$, and $d\equiv7\bmod16$, and an integer $Q_S$
with $2^{\nu+2}\mid Q_S$ and $\gcd(Q_S,qd)=1$, with $(\frac{-d}q)=1$ and
$(\frac{-d}l)=1$ for every odd prime $l\mid Q_S$.
\item A CM field $E$, Galois over $\Q$, contained in $\Q(\zeta_{N_E})$ with
$N_E=Q_Sqd$, together with an isomorphism
$\mathrm{Gal}(E/\Q)\cong U=E_0\times A$, and containing $\sqrt{-d}$ and
$\sqrt{-q}$, with $\sigma(\sqrt{-d})=s_2(\sigma)\sqrt{-d}$ and
$\sigma(\sqrt{-q})=s_1(\sigma)\sqrt{-q}$, where $\sigma\in\mathrm{Gal}(E/\Q)$
is written $((s_1(\sigma),s_2(\sigma)),a)$ in the multiplicative notation
below.
\item The Artin map $\mathrm{Art}:(\Z/N_E)^\times\to U$, which has the
normalisations (H1)--(H5) of Proposition~\ref{i4:prop-rd}, so that
$s_1\circ\mathrm{Art}=(\frac\cdot q)$ and $s_2\circ\mathrm{Art}=(\frac\cdot d)$
(by (H1): $\sqrt{-q}\in\Q(\zeta_q)$ and $\sqrt{-d}\in\Q(\zeta_d)$ are
quadratic Gauss sums, on which $\zeta\mapsto\zeta^u$ acts by $(\frac uq)$
and $(\frac ud)$).
\item The \emph{arithmetic data}: for every prime $p$, the inertia subgroup
$I_p\subset U$, which is the image under $\mathrm{Art}$ of the units
congruent to $1$ modulo the prime-to-$p$ part of $N_E$, and the Frobenius
element $\mathrm{Fr}_p\in U$, which is the image under $\mathrm{Art}$ of a
unit congruent to $p$ modulo the prime-to-$p$ part of $N_E$ and to $1$
modulo its $p$-part. By the decomposition law in cyclotomic fields these
are the inertia group and a Frobenius element of $p$ in
$\mathrm{Gal}(E/\Q)\cong U$.
\item Level primes $g_0,g_1,p_{de},p'_{de},q_0$ and ranks $(k,m,f,t,r)$
with $f=2n$.
\end{itemize}
The host is \emph{good} if a list of explicit conditions holds. The list
is: validity (with constant $C_8$); arithmeticity (the four clauses listed
at the start of the proof of Proposition~\ref{i4:prop-kits}); the lower
bound on the second torus set; $4L_w|T_2|\le n_0$, where $L_w$ is the
window length of item 9 of (E1), $T_2$ the thick set and $n_0$ the rank
parameter with $n=o\,n_0$; the $(q,d)$-bound, the box bound, the thinness
bound and the bad-prime bound, with constants $C_5,c_6,C_b$;
$|\mathrm{dlog}|\le x^\varepsilon$; the two finite sets
$\mathrm{BB}(r,g_0)$ and $\mathrm{BB}(r,g_1)$ (the
cyclotomic norm sets at the two generator primes) are disjoint; for
$i=0,1$ there is $\pi\in\mathcal O_E$ with $\pi\in K=\Q(\sqrt{-d})$,
$\pi\bar\pi=g_i$ and $\pi\equiv1\bmod N_E$;
$p_{de}\equiv p'_{de}\equiv1\bmod\mathrm{lcm}(15,N_E)$; $q_0$ is a prime
with $q_0\equiv1\bmod N_E$, $q_0\le N_E^{C_8}$ and
$q_0\notin\{p_{de},p'_{de}\}$; and the level primes are exactly
$\{g_0,g_1,p_{de},p'_{de},q_0\}$. A host also carries a real constant
$C_{T_2}$, the constant of its thick-set bound.

\emph{Coordinates on $U$.} Two notations for $E_0$ are used. In the
multiplicative notation of Proposition~\ref{i4:prop-rd},
$E_0=\{\pm1\}^2$ and an element of $U$ is $((s_1,s_2),a)$; the
\emph{$d$-sign} $s_2$ is the action on $\sqrt{-d}$. In the additive
notation $E_0=(\Z/2)^2$ of the formal development, the $d$-sign is the
\emph{first} coordinate. The CM type $\Phi\subset U$ is the set of elements
with trivial $d$-sign ($s_2=1$, equivalently first additive coordinate
$0$), that is $\Phi=\mathrm{Gal}(E/\Q(\sqrt{-d}))$; it has $|\Phi|=2n$
elements.
\end{definition}

For a primitive Dirichlet character $\chi\ne1$ we use the prime sum
$G_\chi(y)=\sum_{n\le y}\Lambda(n)\chi(n)(y/n-1)+y\,L'(1,\chi)/L(1,\chi)$,
which appears in clause (2) of Theorem~\ref{M:thm:alpoge-families}.

\begin{definition}[the requirement for the trivial character]\label{i4:def-G1}
Let $H$ be a host at scale $x$ with exponent $\varepsilon$, and let $t$ be
its rank parameter. The \emph{requirement} $\mathcal G_1(H)$ is the
conjunction of the following two conditions.
\begin{itemize}
\item (Twists.) The integer $o$ of the frame equals $1$. For each of the
three nontrivial characters $\chi$ of $U$ that are trivial on $A$, we have
$|G_\chi(x)|\le x^{1/2+\varepsilon}$ and $|L'(1,\chi)/L(1,\chi)|\le4\log x$.
Here $\chi$ is identified with the primitive Dirichlet character cut out
by the Artin map. (The remaining clauses of the twist condition of the
formal development, namely that the kept characters lie in $X^+$, the clause
on the exceptional character and the shares of the top set, are part of
the conclusions of Propositions~\ref{i4:Sprop-field} and~\ref{i4:Sprop-host}
below, and are supplied there in the proof of
Proposition~\ref{i4:prop-supply}.)
\item (Selmer.) $\mathrm{VV}(E,t,\sqrt b)$. This means: for every prime
$l$ and every $0<e<t$ the $\Z_l$-length
$\operatorname{sha}_l(e)=\operatorname{length}\mathrm{Sha}^1(E,\Q_l/\Z_l(-e))$
is finite, and for every finite set $\mathcal S$ of primes
$\sum_{l\in\mathcal S}\log l\sum_{e=1}^{t-1}\operatorname{sha}_l(e)\le\sqrt b$.
\end{itemize}
For $\lambda>1$ the statement $\mathrm{HostExW}(\lambda,\mathcal G_1)$ is
the following. For every $\varepsilon\in(0,1/52)$ there are constants
$C_{T_2}\in\R$, $C_8\ge1$ and $C_5,c_6,C_b>0$, a threshold $x_0$, and a set
of scales whose $x$-coordinates contain intervals $[W,W^\lambda]$ with $W$
arbitrarily large, such that every such scale with $x\ge x_0$ carries valid
ranks, a prime $q_0$ and a host $H$ which is good with the constants
$C_8,C_5,c_6,C_b$ and the auxiliary prime $q_0$, whose thick-set constant is
$C_{T_2}$, and which satisfies $\mathcal G_1$.
\end{definition}

\begin{definition}[families, classes, windows]\label{i4:def-family}
Fix an exponent $\varepsilon$. A \emph{family} $\mathcal F$ of exponent
$\varepsilon$ consists of the following data.
\begin{enumerate}
\item A threshold $x_1>1$.
\item For each real $x$, a finite index set $I(x)$.
\item A set $\mathrm{Win}\subset\R$ that is unbounded above, with
$I(x)\ne\emptyset$ for $x\in\mathrm{Win}$.
\item For each $i\in I(x)$, a column datum $D(x,i)$ (a host with its frame
and its columns, indexed by words) and size data $L(x,i)$ at scale $x$ and
exponent $\varepsilon$.
\end{enumerate}
The data are subject to the following uniformity conditions. At fixed $x$,
all $D(x,i)$ share their ranks, the primes $q,d$, $Q_S$, the Artin table,
the level primes, the word set, the thick set $T_2$, the root discriminant
$\mathrm{rd}_E$, the constant $C_{T_2}$ and the top bound $b_{\rm top}$.
The integer coordinates of the top set do not depend on $i$. The constants
$J,C_\Gamma,C_c,C_t$ and a constant $C_o$ do not depend on $(x,i)$, and
$o\le C_o$, $Q_S\le C_on$ at every point. Finally the \emph{cover} condition holds: for every $i_0$ and every
translation $u_0$ there is $i$ with translation $u_0$ that carries the same
numerators and weights as $i_0$ at equal coordinates.

The family also carries three lists of constants, $c_{\rm pre}$,
$c_{\rm out}$ and $c_{\rm mid}(C_H)$, and a list $c_{\rm last}$. The middle
list depends on a real parameter $C_H$. The family \emph{has windows of
ratio $\lambda$} if for every $y$ there is $W\ge y$ with
$[W,W^\lambda]\subset\mathrm{Win}$. A \emph{family class} is a property of
families. The class $\mathcal G_1.\mathrm{cls}(\lambda)$ consists of the
families with windows of ratio $\lambda$ whose every point $(x,i)$ with
$x>x_1$ satisfies $\mathcal G_1$.
\end{definition}

\begin{definition}[chain-tied; connected]\label{i4:def-chain}
Let $\mathcal F$ have exponent $\varepsilon$, and let
$e_D,e_L,K,c_{HH},e_K,C_L,x\in\R$. The family is \emph{chain-tied at $x$}
with these constants if there are a base index $i_0\in I(x)$ and a map
$\mathrm{pt}:A\to I(x)$ such that the following holds. Here $A$ is the
group of the host $H$ of $D(x,i_0)$. The columns at the translation
$v\in A$ are those of $D(x,\mathrm{pt}(v))$, which shares its frame and
level data with $D(x,i_0)$. From these there are data $q,d$, $T,S\subset A$
of sizes $r,s$, $m\ge0$, $N$, $(a_t)_{t\in T}$, $w_{v,j}$ for $v\in A$ and
$1\le j\le m$, and $a,b,b',k,\eta,X',\theta$, which satisfy properties
(1)--(5) of Theorem~\ref{M:thm:alpoge-families}(b) for $H$. The column
package in property (5) is read from $H$ along the identity of $A$, and its
constants satisfy $e_D^{\rm pkg}\le e_D$, $e_L^{\rm pkg}\le e_L$ and
$c_{HH}^{\rm pkg}\le c_{HH}$. A family class $\mathcal E$ is
\emph{connected} if the following holds. There are reals $e_D,e_L$ and
$\varepsilon_1>0$ such that for every $\varepsilon\in(0,\varepsilon_1)$
there are a family $\mathcal F\in\mathcal E$ of exponent $\varepsilon$ and
constants $K,c_{HH},e_K,C_L,x_0$ for which $\mathcal F$ is chain-tied at
every $x\in\mathrm{Win}(\mathcal F)$ with $x\ge x_0$.
\end{definition}

\begin{definition}[Selmer hypothesis on a family]\label{i4:def-hypV}
$\mathrm{HypV}(\mathcal F)$ means: eventually, $\mathrm{VV}(E,t,\sqrt b)$
holds at every point $(x,i)$, where $E$ and $t$ are the field and rank
parameter of $D(x,i)$.
\end{definition}

\begin{lemma}\label{i4:lem-hypV-cls}
For every $\lambda>1$, every member of $\mathcal G_1.\mathrm{cls}(\lambda)$
satisfies $\mathrm{HypV}$.
\end{lemma}
\begin{proof}
Every point beyond $x_1$ satisfies $\mathcal G_1$. The second conjunct of
$\mathcal G_1$ is the Selmer bound of Definition~\ref{i4:def-hypV}, so
$\mathrm{HypV}$ holds with threshold $x_1+1$.
\end{proof}

\subsection{The four inputs used as hypotheses}\label{i4:ssec-hyps}

The family estimates that the construction must satisfy are called
\emph{links}. Each is required eventually in the family, with constants
that may depend on the family. We name the links and state the four inputs
in the form in which they are used.
\begin{itemize}
\item $(\mathrm{LN})$, the local index link $\mathrm{LinkN}$ of
Definition~\ref{i2:def-LinkN}. There are $C_N$ and $x_0$ such that for all
$x\ge x_0$, all $i\in I(x)$ with $D(x,i)$ Sha-finite and all lifts $j$:
if $(\mathrm{vv}_D(l))_l$ is summable, then
$\eta\Sigma_{L(x,i)}(j)\le d_{Z_0}(j)\bigl(C_N(fmbL^2+(\log x)^2N_c)+\sum'_l\mathrm{vv}_D(l)\bigr)$.
The notation is that of the divisor bound; $\sum'_l$ is the sum if the
sequence is summable and $0$ otherwise.
\item $(\mathrm{SLR})$, small-level rank equality. Eventually
$d_{Z_0}(j)=D_0$.
\item $(\mathrm{DivN})$, the conclusion of the divisor bound.
\item $(\mathrm{LEta})$, twist coverage. This is the monodromy, cover and
Hecke clauses of the Galois component of the tied datum at each point.
\item $(\mathrm{LIrr})$, the irreducibility bound. Eventually
$0\le\delta_p\le\exp(-x^{\varepsilon_z/7})$.
\item $(\mathrm{LDiv})$, the division bound with constants
$(C_{\rm abs},C_H)$. There are $C,x_0$ such that at every point with
$x\ge x_0$ the following nine conditions hold:
(1) the single-map condition (the division lattices map into one
coefficient lattice); (2) the frame-number bound (hh) with constant $C_H$:
$s+\mu_{\max}(RR)\le x^{C_H}b$ (notation of Lemma~\ref{i4:Llem-hh});
(3) the vanishing of the gain and of the excess above the levels
(Lemma~\ref{i4:lem-zeroabove}); (4) the Bessel identity, that is
$\mathrm{BuRational}$ together with the clause c5 of
Definition~\ref{i2:def-DKCB}; (5) $\mathrm{RRFull}$; (6) the height-map
bound: the height of every nonzero division map $B_u$ is at most $h(u)$;
(7) condition (P1) with constant $C$ and exponent $C_{\rm abs}$
(Proposition~\ref{i2:prop-P1}); (8) the c2 bound
$\bigl|\sum_w\mathrm{sz}_w(E(u,j,w))-(\sum_ta_t)\log(g_0g_1)\bigr|\le CB_0$
for every word $u$ and lift $j$; (9) the variation bound
$\bigl|(h(u)+\sum_ta_t\log T_u)-(h(u')+\sum_ta_t\log T_{u'})\bigr|\le
C(|T|aL_a+mB_0)$ for all words $u,u'$.
\item $(\mathrm{LG})$, the prime-by-prime comparison of the loss bound.
\item $(\mathrm{LS}(C_0))$, almost-orthogonality of the kernel blocks:
this is $\mathrm{LinkS}(C_0)$ of Definition~\ref{i3:def-linkS}. Eventually
there are $\eta_0,\eta_1\in\R$ and $e_0\ge0$ with
$\eta_0,\eta_1,\sqrt{e_0}\le\exp(-x^{\varepsilon/8})$ such that the
following implication holds: if $C_0^{-1}\le d_u\Gamma^2/D_*\le C_0$ for
every word $u$, then (i) $\|K(u,u')\|_{HS}^2\le(1+\eta_0)\,d_ud_{u'}/D_*$
for all $u\ne u'$; (ii) $\|K(u,u)-\mathrm{id}\|_{HS}^2\le e_0\,d_u$ for all
$u$; (iii) $D_{\rm all}\le(1+\eta_1)D_*$.
\item $(\mathrm{LMu})$, the slope count. For some $C$, eventually
$N'_\mu(u)\le C\,X''_n$ at every word $u$. The quantity $N'_\mu$ is
recalled in the proof of Proposition~\ref{i4:prop-mu}.
\item $(\mathrm{LGl})(C_{\rm abs},C_H,C_0,e_D,e_L)$, the global link. It
says: if the chosen data with $(C_0,C_H)$ (Definition~\ref{i4:def-chosen})
and $(\mathrm{LG})$, $(\mathrm{LS}(C_0))$, $(\mathrm{LIrr})$,
$(\mathrm{LDiv})$ and $(\mathrm{LMu})$ all hold, then there are
$K,c_{HH},e_K,C_L,x_0$ such that the family is chain-tied with
$(e_D,e_L,K,c_{HH},e_K,C_L)$ at every $x\ge x_0$ in $\mathrm{Win}$.
\end{itemize}

\begin{proposition}[the divisor bound; hypothesis]\label{i4:hyp-div}
Let $\mathcal F$ have exponent $\varepsilon>0$ and satisfy $\mathrm{HypV}$,
$(\mathrm{LN})$ with constant $C_N$, and $(\mathrm{SLR})$. Then there is
$x_0$ such that for all $x\ge x_0$, $i\in I(x)$ and all lifts $j$,
\[
 \eta\Sigma_{L(x,i)}(j)\le D_0\,(|C_N|+1)\,x^{\varepsilon}(B_0+N_c).
\]
This conclusion is $(\mathrm{DivN})$.
\end{proposition}

\begin{proposition}[the loss bound; hypothesis]\label{i4:hyp-loss}
\begin{enumerate}
\item[(a)] Let $\mathcal C$ be a class of families. If $\mathcal C$ is
closed (Definition~\ref{i4:def-closed}) for group compatibility and for the
second admissibility condition, then $\mathcal C$ is closed for
$(\mathrm{LG})$: there is $\varepsilon_1>0$ such that for every
$\varepsilon\in(0,\varepsilon_1]$ there is a threshold datum $\kappa$ such
that $(\mathrm{LG})$ holds for every family of $\mathcal C$ of exponent
$\varepsilon$ which is $\kappa$-large at some $C_H$.
\item[(b)] (Theorem~\ref{i1:thm-loss}.) Let $\mathcal F$ satisfy
$\mathrm{ChosenTT}$ (the top-set condition), $\mathrm{ChosenBadp}$ (the
bad-prime condition of the small-shares record, as produced in
Lemma~\ref{i4:Llem-loss}) and $(\mathrm{LG})$ with constant $C$. Put
$K=1+|C|+C_{\rm ord}$ with $C_{\rm ord}=12$. Then for all large $x$ and all
$i$ such that $\mathrm{ChW}(X_{x,i})$ holds (the weight vector has small
shares),
\[
 I_O(L_{x,i})\le(3K+13)\,|T|\,m\,b,\qquad
 \sum_t\sum_j P(a_t,t,w_j)-g(L_{x,i})\le(3K+13)\,|T|\,m\,b
\]
for every reading of $L_{x,i}$ and every admissible cutoff. In other
words, when small-share weight vectors exist at every large scale, (P7)
holds with $c_{\rm loss}=3K+13$ and $e_L=0$.
\end{enumerate}
\end{proposition}

\begin{proposition}[the level counts and Theorem S; hypothesis]\label{i4:hyp-S}
On the working class $D_3(\mathcal A)$ of \S\ref{i4:ssec-links}:
\begin{enumerate}
\item[(a)] the implication $\mathrm{TuSize}\Rightarrow(\mathrm{LS}(3))$ is
closed (Theorem~S, part (b), with the slot record of the hypothesis
structure);
\item[(b)] $(\mathrm{LIrr})$ is closed (Theorem~S, part (c), with the
irreducibility row and Green's hypothesis);
\item[(c)] for every member, condition (P4) holds at every point beyond a
threshold, with a constant $c_{\rm count}$ (the level counts, Theorem~S
part (d)).
\end{enumerate}
\end{proposition}
\noindent Section~\ref{i3:sec} states (a) and (b) for every family of the working class
beyond the closure threshold ($8C_\Gamma\le C_{k,\rm low}$); in the language
of Definition~\ref{i4:def-closed} below this says exactly that the two
properties are closed on $D_3(\mathcal A)$.

\begin{proposition}[the windowed Selmer bound; hypothesis]\label{i4:hyp-selmer}
$\mathrm{VWinAt}(1,401)$ holds. Explicitly: for all
$0<\varepsilon<1/52$ and every $e:\R\to\N$ with
$x^{2+8\varepsilon}\le2^{e(x)}\le2x^{2+8\varepsilon}$ for $x\ge1$, there are $J\ge1$
and $X_0$ with the following property. For all $X\ge X_0$ and every window
pair $(q,d)$ for $(\varepsilon,X^{1/J})$ there is $Y\ge X^{1/J}$ with
$Y^2\le X$ such that for every $x\in[Y,Y^2]$ and every $(1,e(x),q,d)$-field
$E$, $\mathrm{VV}(E,\lfloor x^{401}\rfloor,\sqrt{\exp(x^\varepsilon)})$
holds.
\end{proposition}

The passage from the Selmer bound to hosts is not part of Section~\ref{i5:sec}. It is
the following statement, proved in the development.

\begin{proposition}[host supply]\label{i4:prop-supply}
$\mathrm{VWinAt}(1,401)$ implies $\mathrm{HostExW}(3/2,\mathcal G_1)$.
\end{proposition}
\noindent The proof (\S\ref{i4:ssec-proof-supply}) is an assembly of the
five constructions stated exactly in
Propositions~\ref{i4:Sprop-field}--\ref{i4:Sprop-tw} below (the field of
the character, the window data, valid ranks, the host, the twist bounds).
The twist bounds are proved in \S\ref{i4:Sst-tw} and the window data in
\S\ref{i4:Sst-win} modulo the existence of good window pairs
(Proposition~\ref{i4:Sst-prop-pair}); the field of the character, the
valid ranks and the host are formally verified, and their formally verified
steps are outlined in \S\S\ref{i4:Sst-ssec} and~\ref{i4:Sst-outline}.

\begin{definition}[window data]\label{i4:Sdef-win}
Let $\varepsilon\in(0,1/52)$, $J\ge1$, $c\ge1$, $1<x$, $Y>0$, $M\ge1$, and $q,d,e\in\N$.
A primitive Dirichlet character $\chi$ \emph{induces} a function $f$ modulo
$N$ if $\chi(a)=f(a)$ for all integers $a$ prime to $N$.
\begin{enumerate}
\item $\mathrm{PairGood}_M(\varepsilon,Y,x,q,d)$ means:
(i) $(q,d)\in\mathrm{WP}(\varepsilon,Y)$;
(ii) $\gcd(M,qd)=1$;
(iii) for every Dirichlet character $\xi$ modulo $M$, every
$(i,j)\in\{0,1\}^2\smallsetminus\{(0,0)\}$, every primitive Dirichlet
character $\chi$ inducing $\xi\chi_q^i\chi_d^j$ modulo $2Mqd$, and all
$y\in[Y,x]$: $|G_\chi(y)|\le y^{1/2+\varepsilon/2}$ and
$|L'/L(1,\chi)|\le4\log y$;
(iv) for every real primitive character $\varrho$ of modulus dividing $8M$,
every $(i,j)\ne(0,0)$ and every primitive $\chi$ inducing
$\varrho\chi_q^i\chi_d^j$ modulo $2Mqd$: $L(s,\chi)\ne0$ for
$\Rel s>1-1/400$ and
$|\Iml s|\le8M\exp(Y^{\varepsilon/5})\bigl(\exp(Y^{\varepsilon/5}(\log Y)^2)\bigr)^j$.
We write $\mathrm{PairGood}=\mathrm{PairGood}_1$.
\item $\mathrm{WinDatum}(\varepsilon,J,c,x,Y,q,d,e)$ means: $x^{1/J}\le Y\le x$;
$\mathrm{PairGood}(\varepsilon,Y,x,q,d)$; $x^{2+8\varepsilon}\le c2^e\le2x^{2+8\varepsilon}$;
and every $(1,e,q,d)$-field $E$ satisfies
$\mathrm{VV}(E,\lfloor x^{400}\rfloor,\sqrt{\exp(x^\varepsilon)})$.
\item A host datum $H$ over a frame $\varphi$ satisfies $\mathrm{KerChi}$ if
$\sigma_a$ ($a$ prime to $N=Q_Sqd$) lies in the kernel of the Artin map
exactly when $a\equiv\pm1\bmod2^{\nu+2}$ and $(\frac aq)=(\frac ad)=1$, where
$\nu=\nu(\varphi)$; and $\mathrm{TwistForm}$ if every character $\theta$ of
$U$ trivial on the first factor is, through the Artin map, of the form
$\chi_q^i\chi_d^j$ with $\theta(\iota)=(-1)^j$.
\end{enumerate}
\end{definition}

\begin{proposition}[the field of the character; formally verified, \cite{CharacterFamily}]\label{i4:Sprop-field}
There are $c,e_0\in\N$, $c\ge1$, such that for all $e\ge e_0$ and all
$q,d,b\in\N$ with $q\equiv3\bmod8$, $d\equiv7\bmod16$ primes,
$(\frac{-d}q)=1$, $q\ne3$, $b\ge4$, and $(\frac{-d}p)=1$ for every odd prime
$p\le b$, there are a frame $\varphi$ and host axioms $H$ over it with
$\nu(\varphi)=e$, $\mathrm{ord}(\varphi)=c\,2^e$, $H.q=q$, $H.d=d$, $o=1$,
the field rows $\mathrm{FieldRows}(H,b)$, $Q_S\mid2^{e+2}$,
$Q_S\le4\,\mathrm{ord}(\varphi)$, $\mathrm{KerChi}$, $\mathrm{TwistForm}$,
and the two character clauses of the requirement: the kept characters
(those $\theta$ with $\chi_\theta=1$, see \S\ref{i4:Sst-tw}) lie in $X^+$,
and the $q$-sign character $((s_1,s_2),a)\mapsto s_1$ of $U$ is not kept.
\end{proposition}
\noindent The structure of its formal proof is described in \S\ref{i4:Sst-field}.

\begin{proposition}[window data]\label{i4:Sprop-win}
Let $\varepsilon\in(0,1/52)$ and $c\ge1$, and assume $\mathrm{VWinAt}(1,401)$.
There is $J\ge1$ such that for every $y$ there is $Z\ge y$ such that for
every $x\in[Z,Z^{3/2}]$ there are $Y,q,d,e$ with
$\mathrm{WinDatum}(\varepsilon,J,c,x,Y,q,d,e)$.
\end{proposition}
\noindent The proof, complete modulo the existence of good window pairs (Proposition~\ref{i4:Sst-prop-pair}, formally verified), is in \S\ref{i4:Sst-win}.
\noindent (The pair existence statement used in the formal proof is
proved there unconditionally.)

\begin{proposition}[valid ranks; formally verified, \cite{CharacterFamily}]\label{i4:Sprop-ranks}
For $\varepsilon\in(0,1/52)$ there is $x_0$ such that for every scale of
exponent $\varepsilon$ with $x\ge x_0$, every frame $\varphi$ with
$x^{2+8\varepsilon}\le\mathrm{ord}(\varphi)\le2x^{2+8\varepsilon}$ and all host
axioms over $\varphi$, there are ranks $R$ with frame $\varphi$ that are
valid at the scale and admitted for the host axioms. Moreover valid ranks
satisfy $t\le x^{400}$ once $x\ge2^{100}$.
\end{proposition}
\noindent The proof, complete given two clauses of the host axioms, is in \S\ref{i4:Sst-ranks}.

\begin{proposition}[the host; formally verified, \cite{CharacterFamily}]\label{i4:Sprop-host}
For $\varepsilon\in(0,1/52)$ and $C_0\ge1$ there are $C_{T_2},C_8,c_6,C_b,x_0$
with $C_0\le C_8$, $c_6,C_b>0$, such that for every scale $s$ of exponent
$\varepsilon$ with $x\ge x_0$, valid admitted ranks $R$, host axioms $H_0$
over $R$, $M\in\N$ and $Y$ with: $J(s)\le C_0$, $o\le C_0$,
$Q_S\le C_0n$, $0<M\le C_0$, $Q_S\mid\mathrm{lcm}(2^{\nu+2},M)$,
$x^{1/J}\le Y\le x$, $\mathrm{PairGood}_M(\varepsilon,Y,x,H_0.q,H_0.d)$
and $\mathrm{FieldRows}(H_0,\lfloor Y^{\varepsilon/5}\rfloor)$, there are a
host $H$ extending $H_0$ and a prime $q_0$ with constant $C_{T_2}$, $H$ good
with constants $(C_8,4,c_6,C_b)$ and auxiliary prime $q_0$, and the shares
of the top set bounded by $x^{-1+\varepsilon}$.
\end{proposition}
\noindent The structure of its formal proof is described in \S\ref{i4:Sst-host}.
\noindent (It rests on Proposition~\ref{i4:thm-shortnorm} and the
smallest-prime bound Theorem~\ref{i4:thm-linnik}.)

\begin{proposition}[the twist bounds]\label{i4:Sprop-tw}
Let $H$ be host data over a frame with $\mathrm{TwistForm}$, every prime
divisor of $Q_S$ equal to $2$, $1\le x$, $Y\le x$ and
$\mathrm{PairGood}(\varepsilon,Y,x,q,d)$. Then for the three nontrivial
characters $\chi$ of $U$ trivial on $A$: $|G_\chi(y)|\le y^{1/2+\varepsilon}$
at $y=x$ (the pointwise bound of the bounded set) and
$|L'/L(1,\chi)|\le4\log x$.
\end{proposition}
\noindent The proof is in \S\ref{i4:Sst-tw}.
\noindent The formal statement is the more general
Lemma~\ref{i4:Sst-lem-tw}, about the bounded set of the host data; the
three characters above lie in that set (this uses the surjectivity of the
Artin map, not Conventions (2) and (3)). Both are proved in
\S\ref{i4:Sst-tw}.

\subsection{Thresholds and closed properties}\label{i4:ssec-closed}

The constants of the construction are fixed in stages. First come the
outer constants of a family. Then come the middle constants, whose lower
bounds may depend on the outer ones. Last come the remaining constants,
whose lower bounds may depend on the outer and middle ones. The following
definition makes this precise.

\begin{definition}[threshold data, largeness, closedness]\label{i4:def-closed}
A \emph{threshold datum} $\kappa=(\kappa_{\rm pre},\kappa_t,\kappa_{\rm mid},\kappa_{\rm last})$
consists of a real number and three functions from finite lists of reals to
$\R$. We write $\kappa\le\kappa'$ for the componentwise order and
$\kappa\vee\kappa'$ for the componentwise maximum. A family $\mathcal F$ is
\emph{$\kappa$-large at $C_H$} if the following three conditions hold:
\begin{itemize}
\item every entry of $c_{\rm pre}$ is at least $\kappa_{\rm pre}$;
\item every entry of $c_{\rm mid}(C_H)$ is at least
$\kappa_{\rm mid}(c_{\rm out})$;
\item every entry of $c_{\rm last}$ is at least
$\kappa_{\rm last}(c_{\rm out}\frown c_{\rm mid}(C_H))$.
\end{itemize}
A class $\mathcal C$ is \emph{closed for} a property $P(\mathcal F,C_H)$ if
there is $\varepsilon_1>0$ with the following property. For every
$\varepsilon\in(0,\varepsilon_1]$ there is a threshold datum $\kappa$ such
that $P(\mathcal F,C_H)$ holds for every $\mathcal F\in\mathcal C$ of
exponent $\varepsilon$ and every $C_H$ at which $\mathcal F$ is
$\kappa$-large.
\end{definition}

\begin{lemma}\label{i4:lem-closed}
\begin{enumerate}
\item[(a)] $\mathcal F$ is $(\kappa\vee\kappa')$-large at $C_H$ if and only
if it is both $\kappa$-large and $\kappa'$-large at $C_H$. If
$\kappa\le\kappa'$, then $\kappa'$-largeness implies $\kappa$-largeness.
\item[(b)] If $\mathcal C$ is closed for $P$ and for $Q$, then it is closed
for $P\wedge Q$.
\item[(c)] If $\mathcal C$ is closed for $P$ and $P\Rightarrow Q$ holds for
every family of positive exponent, then $\mathcal C$ is closed for $Q$.
\item[(d)] If $\mathcal C$ is closed for $P$ and
$\mathcal C'\subset\mathcal C$, then $\mathcal C'$ is closed for $P$.
\item[(e)] A property that holds for every member of $\mathcal C$ beyond a
member-dependent threshold in $x$ is closed on $\mathcal C$. One may take
$\varepsilon_1=1$ and the zero threshold datum.
\end{enumerate}
\end{lemma}
\begin{proof}
Each largeness condition has the form $k\le c$. Part (a) follows because
$\max(k,k')\le c$ holds exactly when $k\le c$ and $k'\le c$. For (b), take
$\min(\varepsilon_P,\varepsilon_Q)$ and the threshold datum
$\kappa_P\vee\kappa_Q$, and apply (a). Parts (c)--(e) are immediate from
the definition.
\end{proof}

\subsection{Connectedness from design existence}\label{i4:ssec-desWin}

\begin{definition}[chosen data; design existence]\label{i4:def-chosen}
For a family $\mathcal F$ and reals $C_0,C_H$, the statement
$\mathrm{ChosenE}(\mathcal F,C_0,C_H)$ is the following. There are
$C_5,c_6,C_b>0$ such that, eventually, the following conditions hold:
\begin{itemize}
\item the lower bound on the thick set;
\item the condition on the top set;
\item the bad-prime bound with constant $C_b$;
\item the $(q,d)$-bound with constant $C_5$;
\item the thinness bound with constant $c_6$;
\item the \emph{proviso} with constant $C_0$: $C_0^{-1}\le d_u\Gamma^2/D_*\le C_0$
for every word $u$;
\item the bound $C_0\Gamma(L)^2x^{C_H}+1$ on the number of words.
\end{itemize}
Further:
\begin{itemize}
\item $\mathrm{TuSize}(\mathcal F)$ means that, eventually, the level
primes $T_u$ satisfy their lower bound $b^{m/2}\le T_u$ and
$\log T_u\le C_{T_u}mL$, for every word $u$;
\item $\mathrm{DGamSize}(\mathcal F)$ means that for some $C$, eventually,
the stabiliser denominators satisfy
$\log D_\Gamma(u)\le C\,fN_m\log(fN_m)$, where $N_m=k+m(t+1)$.
\end{itemize}
\emph{Designs exist} for $(\mathcal C,C_0)$ if the following holds. There is
$\varepsilon_0>0$ such that for every $\varepsilon\in(0,\varepsilon_0]$ and
\emph{every} threshold datum $\kappa$ there are $\mathcal F\in\mathcal C$
of exponent $\varepsilon$ and $C_H\ge0$ with the following properties:
$\mathcal F$ is $\kappa$-large at $C_H$, and $\mathrm{ChosenE}(\mathcal
F,C_0,C_H)$, $\mathrm{TuSize}(\mathcal F)$ and $\mathrm{DGamSize}(\mathcal F)$
hold.
\end{definition}

We also use four auxiliary classes and one implication.
\begin{itemize}
\item $\mathrm{WithSm}(\mathcal C)$ adds to $\mathcal C$ the smoothness of
the stabiliser denominators.
\item $\mathrm{WithTw}(\mathcal C)$ adds the twist condition.
\item $\mathrm{Des}(\mathcal C,C_0)$ is $\mathrm{WithTw}(\mathrm{WithSm}(\mathcal C))$
together with ``$\mathrm{ChosenE}(C_0,C_H)$ for some $C_H$'',
$\mathrm{TuSize}$ and $\mathrm{DGamSize}$.
\item $(\mathrm{LDivt})$ is the implication
$(\mathrm{LEta})\wedge(\mathrm{DivN})\Rightarrow(\mathrm{LDiv})$.
\end{itemize}

\begin{lemma}[no gain above the levels]\label{i4:lem-zeroabove}
Let $\mathcal F$ be any family. At every point $(x,i)$ with $x>x_1$, for
every word $u$, every lift $j$ and every prime $w$ of the coefficient field
whose residue characteristic $p_w$ satisfies $p_w>X''_n$ (where
$X''_n=X'xn$ and $X'=\exp(x^{\varepsilon/3})$),
$p_w\neq T_u$, $p_w\notin\{g_0,g_1\}$, $p_w\neq d$ and $p_w\nmid D_\Gamma(u)$, the gain
$\mathrm{gain}(j,w)$ and the excess term $E(u,j,w)$ vanish.
\end{lemma}
\begin{proof}
For $x>x_1$ the size data of $\mathcal F$ at $(x,i)$ come with a support
$S$ (a finite set of primes $w$) and the following four clauses of the
family record.
\begin{itemize}
\item[(Z1)] $\mathrm{gain}(j,w)=0$ for $w\notin S$.
\item[(Z2)] If $w\in S$ and $p_w$ is neither of residue class zero nor
of split type, then $\mathrm{gain}(j,w)=0$ for every $j$.
\item[(Z3)] (fill clause) If $w\in S$ and $p_w$ is of residue class zero or of
split type, then $\mathrm{gain}(j,w)=\mathrm{gain}^P(p_w,j)$ for every
lift $j$, where $\mathrm{gain}^P(p,j)$ is the explicit fill value, which
is defined to be $0$ whenever $p>X''_n$.
\item[(Z4)] (excess definition) For every word $u$, lift $j$ and prime
$w$, with $p=p_w$: $E(u,j,w)=E_T(u,w)$ if $p=T_u$; otherwise
$E(u,j,w)=v_p(D_\Gamma(u))+e(u,j,w)$, where $e(u,j,w)$ is the term at
$g_0$ if $p=g_0$, the term at $g_1$ if $p=g_1$, the term at $d$ if
$p$ is ramified of class $d$ (the only such prime is $d$), the term at
residue-class-zero or split-type primes if $p$ is of one of these two
kinds \emph{and} $p\le X''_n$, and $e(u,j,w)=0$ in every other case.
\end{itemize}
Let $p_w$ satisfy the hypotheses. For the gain: if $w\notin S$ use (Z1);
if $w\in S$ and $p_w$ is neither of residue class zero nor of split type,
use (Z2); otherwise (Z3) gives $\mathrm{gain}(j,w)=\mathrm{gain}^P(p_w,j)=0$
because $p_w>X''_n$. For the excess: $p_w\ne T_u$, so the second case of
(Z4) applies; $v_{p_w}(D_\Gamma(u))=0$ since $p_w\nmid D_\Gamma(u)$; and
none of the four nonzero cases of $e$ occurs, because $p_w\ne g_0,g_1$,
$p_w\ne d$ and $p_w>X''_n$. Hence $E(u,j,w)=0$.
\end{proof}
\noindent No hypothesis on the family beyond $x>x_1$ is used. The lemma is used in
Proposition~\ref{i4:prop-mu}(c), as conjunct (3) of $(\mathrm{LDiv})$ in
Lemma~\ref{i4:lem-bridge}, and in Proposition~\ref{i4:Lprop-D}(f).

\begin{proposition}[the size inputs]\label{i4:prop-mu}
Let $\mathcal F$ have exponent $\varepsilon\in(0,1/52)$, and write
$X''_n=X'xn$ with $X'=\exp(x^{\varepsilon/3})$.
\begin{enumerate}
\item[(a)] $\mathrm{TuSize}$ implies that, eventually,
$\log T_u\le2X''_n$.
\item[(b)] $\mathrm{DGamSize}$ with constant $C$ implies that, eventually,
$\log D_\Gamma(u)\le10|C|X''_n$.
\item[(c)] Suppose every lattice $\nu_u$ is eventually full and of
positive rank, and (a) and (b) hold. Then $(\mathrm{LMu})$ holds. In
particular $(\mathrm{LDiv})$ is not needed.
\end{enumerate}
\end{proposition}
\begin{proof}
(a) For large $x$ we have $C_{T_u}\le x^{\varepsilon/100}$. Then
$\log T_u\le m\,x^{\varepsilon+\varepsilon/100}\le mx\le2nx\le2X''_n$,
using $m\le2n$.

(b) The size record contains the following facts: $N_m\le\eta/m$,
$\log\eta\le x^{\varepsilon/4}$, and $n\le x^3$. They give
$N_m\le X'$ and $\log(fN_m)\le\log2+3\log x+x^{\varepsilon/4}\le4x$.
Hence $0\le fN_m\log(fN_m)\le8X'xn\le10X''_n$.

(c) The quantity bounded by $(\mathrm{LMu})$ is
\[
N'_\mu(u)=\Bigl|\mu(\nu_u)-\mu(\nu^H_u)-\sum_w2\ell_we_w\,\mathrm{sg}_u(w)\Bigr|,
\]
where $\nu_u=I_u\nu^H_u$ and $I_u=\prod_w\mathfrak P_w^{-\lfloor e_w\mathrm{sg}_u(w)\rfloor}$.
Here $w$ runs over the primes of the coefficient field, with weights
$\ell_w=\log\mathrm N\mathfrak P_w/[LL:\Q]$, and $\mathrm{sg}_u(w)$ is the
gain at $(u,w)$. We bound it in three steps.
\begin{itemize}
\item The slope of a full lattice changes under a fractional ideal by
$\mu(I\Lambda)=\mu(\Lambda)-2\log\mathrm N(I)/[LL:\Q]$. Hence
$N'_\mu(u)\le2\sum_{\mathrm{sg}_u(w)\ne0}\ell_w$, because each floor
differs from its argument by less than $1$.
\item By Lemma~\ref{i4:lem-zeroabove} the gain vanishes at primes
$p_w>X''_n$ outside $\{T_u,g_0,g_1,d\}$ and the prime divisors of
$D_\Gamma(u)$. Since $\sum_{w\mid p}\ell_w\le\log p$, Chebyshev's bound
gives
\[
\sum_{\mathrm{sg}_u(w)\ne0}\ell_w\le X''_n\log4+\log T_u+\log g_0+\log g_1+\log d+\log D_\Gamma(u).
\]
\item The size record gives $\log g_j,\log d\le x\le X''_n$. Together with
(a) and (b) this yields
$N'_\mu(u)\le2(\log4+3+2+10|C|)X''_n$.\qedhere
\end{itemize}
\end{proof}

\begin{theorem}[connectedness from design existence]\label{i4:thm-desWin}
Let $\mathcal C,\mathcal C',\mathcal E$ be family classes and let
$C_{\rm abs},C_0,e_D,e_L\in\R$. Assume the divisor bound
(Proposition~\ref{i4:hyp-div}) and the following:
\begin{enumerate}
\item designs exist for $(\mathrm{WithTw}(\mathrm{WithSm}(\mathcal C')),C_0)$;
\item $\mathcal C'\subset\mathcal C\cap\mathcal E$;
\item every member of $\mathcal E$ satisfies $\mathrm{HypV}$;
\item $\mathrm{Des}(\mathcal C,C_0)$ is closed for each of the following
nine properties:
\begin{itemize}
\item $(\mathrm{LN})$; $(\mathrm{LEta})$; $(\mathrm{LIrr})$;
$(\mathrm{SLR})$; $(\mathrm{LG})$;
\item $\mathrm{TuSize}\wedge\mathrm{DGamSize}\Rightarrow(\mathrm{LDivt})$;
\item $\mathrm{ChosenE}(C_0,C_H)\Rightarrow$ (every $\nu_u$ is eventually
full and of positive rank);
\item $\mathrm{TuSize}\Rightarrow(\mathrm{LS}(C_0))$;
\item $\mathrm{TuSize}\wedge\mathrm{DGamSize}\Rightarrow(\mathrm{LGl})(C_{\rm abs},C_H,C_0,e_D,e_L)$.
\end{itemize}
\end{enumerate}
Then $\mathcal E$ is connected, with the exponents $e_D,e_L$.
\end{theorem}
\begin{proof}
By Lemma~\ref{i4:lem-closed}(b), $\mathrm{Des}(\mathcal C,C_0)$ is closed
for the conjunction of the nine properties. Let $\varepsilon_1$ be its
exponent, let $\varepsilon_3$ be the exponent of (1), and put
$\varepsilon_*=\min(\varepsilon_1,\varepsilon_3)$. Fix
$\varepsilon\in(0,\varepsilon_*)$, and let $\kappa$ be the threshold datum
given by closedness at $\varepsilon$.

Now apply (1) to \emph{this} $\kappa$. It gives $\mathcal F$ and
$C_H\ge0$ such that $\mathcal F$ is $\kappa$-large at $C_H$ and
$\mathrm{ChosenE}$, $\mathrm{TuSize}$ and $\mathrm{DGamSize}$ hold. By (2),
$\mathcal F\in\mathrm{Des}(\mathcal C,C_0)$, so the nine properties hold
at $(\mathcal F,C_H)$. By (2) and (3), $\mathcal F\in\mathcal E$ and
$\mathrm{HypV}(\mathcal F)$ holds.

With $(\mathrm{LN})$ and $(\mathrm{SLR})$, the divisor bound gives
$(\mathrm{DivN})$. With $(\mathrm{LEta})$, the implication
$(\mathrm{LDivt})$ then gives $(\mathrm{LDiv})$. The seventh closed
property, with $\mathrm{ChosenE}(C_0,C_H)$, gives that every $\nu_u$ is
eventually full and of positive rank; $\mathrm{TuSize}$ and
$\mathrm{DGamSize}$ give (a) and (b) of Proposition~\ref{i4:prop-mu}; so
Proposition~\ref{i4:prop-mu}(c) gives $(\mathrm{LMu})$. All hypotheses of
$(\mathrm{LGl})$ now hold. It yields $K,c_{HH},e_K,C_L,x_0$ such that
$\mathcal F$ is chain-tied at every $x\ge x_0$ in $\mathrm{Win}(\mathcal F)$.
\end{proof}

\begin{remark}
The threshold datum is produced by the finitely many estimates
\emph{before} the family is chosen, and design existence is then applied
to it. This is the order of quantifiers that
Theorem~\ref{M:thm:alpoge-families} requires.
\end{remark}

\subsection{Design existence}\label{i4:ssec-design}

\begin{theorem}[Minkowski's bound; Minkowski, J.\ reine angew.\ Math.\ 101 (1887), 196--202; Serre, Thm.~1]\label{i4:thm-mink}
For $N\in\N$ put
\[
M(N)=\prod_{p\le N+1}p^{r(N,p)},\qquad r(N,p)=\sum_{j=0}^N\Bigl\lfloor\frac N{p^j(p-1)}\Bigr\rfloor .
\]
For every $N$, the order of every finite subgroup of $\mathrm{GL}_N(\Q)$
divides $M(N)$.
\end{theorem}
\noindent This is Minkowski's theorem in its classical form: the
$p$-adic exponent of Minkowski's number is
$\sum_{j\ge0}\lfloor N/(p^j(p-1))\rfloor$ for every prime $p$ (including
$p=2$), and the truncations $j\le N$ and $p\le N+1$ lose nothing, since
$p^j(p-1)\ge2^{N+1}>N$ for $j>N$ and $\lfloor N/(p-1)\rfloor=0$ for
$p>N+1$. (Check: $M(1)=2$, $M(2)=24$.) Serre's survey on bounds for the orders of finite subgroups of $G(k)$,
Theorem~1 (in \emph{Group Representation Theory}, EPFL Press, Lausanne 2007, 405--450; arXiv:1011.0346), bounds $v_\ell(|A|)$ by Minkowski's
exponent for every finite subgroup $A$ of $\mathrm{GL}_N(\Q)$ and every
prime $\ell$; this is the divisibility form above. (The statement is also proved in the formal
development; we treat it as cited literature by choice, since it is
classical, and not because it is assumed there.) It is applied to $\Gamma$-stabilisers,
which are finite subgroups of $\mathrm{GL}_N(\Q)$; no other hypothesis
is needed.

\begin{lemma}\label{i4:lem-logmink}
$M(N)$ divides $(2N)!$. Hence, if $y\ge2$, then $\log M(2y)\le12\,y\log y$.
Consequently, if eventually $D_\Gamma(u)\mid M(2fN_m)$ for every word $u$,
then $\mathrm{DGamSize}$ holds with constant $12$.
\end{lemma}
\begin{proof}
For $p\le N+1$ we have $p^j(p-1)\ge p^{j+1}/2$. Hence
$\lfloor N/(p^j(p-1))\rfloor\le\lfloor2N/p^{j+1}\rfloor$, and summing over
$j$ gives $r(N,p)\le v_p((2N)!)$ by Legendre's formula. Therefore
$M(2y)\le(4y)!\le(4y)^{4y}$, and
$\log M(2y)\le4y(2\log2+\log y)\le12y\log y$. For the last claim, note
that $fN_m\ge2$.
\end{proof}

\begin{definition}[design choices]\label{i4:def-choices}
Let $c_T:\R^2\to\R$. The statement $\mathrm{DesignChoices}(\mathcal C,C_0,c_T)$
is the following. There is $\varepsilon_0>0$ such that for every
$\varepsilon\in(0,\varepsilon_0]$ and every threshold datum $\kappa$ there
are $\mathcal F\in\mathcal C$ of exponent $\varepsilon$ and
$C_H\ge0$, $C_5,C_b>0$ with the following properties. First, $\mathcal F$
is $\kappa$-large at $C_H$. Second, the following conditions hold
eventually:
\begin{itemize}
\item the thick-set, top-set, bad-prime ($C_b$), $(q,d)$ ($C_5$) and code
bounds;
\item the box and $(q,d)$-row conditions;
\item $D_\Gamma(u)\mid M(2fN_m)$ for every word $u$;
\item $c_T(\varepsilon,C_{YL})\le C_{T_u}$;
\item the tame data are distinct;
\item $C_F=C_H$;
\item the word-count bound;
\item $C_{YL}\le9$.
\end{itemize}
The \emph{working class} of a class $\mathcal A$ is
$\mathrm{WithTw}(\mathrm{WithSm}(\mathrm{WithChN}(\mathcal A)))$.
Here $\mathrm{WithChN}$ adds to $\mathcal A$ the following eventual
conditions: $D_\Gamma\mid M(2fN_m)$, distinct tame data, the box and
$(q,d)$-row conditions, and $C_{YL}\le9$.
\end{definition}

\begin{theorem}[design existence from design choices]\label{i4:thm-choices}
Assume $\mathrm{DesignChoices}(\mathcal C,C_0,c_T)$. Assume also that
$\mathcal C$ is closed for each of the following three properties:
\begin{itemize}
\item $c_T(\varepsilon,C_{YL})\le C_{T_u}\Rightarrow\mathrm{TuSize}$;
\item the proviso with constant $C_0$;
\item ``boxes, $(q,d)$-rows and a $(q,d)$-bound imply a thinness bound with
some $c_6>0$''.
\end{itemize}
Then designs exist for $(\mathcal C,C_0)$.
\end{theorem}
\begin{proof}
Let $\varepsilon$ be at most both exponents. Let $\kappa$ be arbitrary,
and let $\kappa'$ be the datum of the closedness. Apply the design choices
to $\kappa\vee\kappa'$. The family $\mathcal F$ obtained is both
$\kappa$-large and $\kappa'$-large (Lemma~\ref{i4:lem-closed}(a)), so the
three closed properties hold for it:
\begin{itemize}
\item the first gives $\mathrm{TuSize}$;
\item the second gives the proviso;
\item the third gives $c_6$ and the thinness bound.
\end{itemize}
Together with the eventual bounds of the choices, this is
$\mathrm{ChosenE}(\mathcal F,C_0,C_H)$. Finally, $\mathrm{DGamSize}$
follows from Lemma~\ref{i4:lem-logmink}.
\end{proof}

The proviso is where the class number formula enters. We prove the
inequality, since it is the only place where the constant $C_0=3$ is
fixed.

\begin{lemma}[the proviso with constant $e$]\label{i4:lem-proviso}
At a point of a family, let $u$ be a word with trivial class-$0$
stabilisers. Make the following assumptions.
\begin{itemize}
\item The class number formula holds:
$d_u=\dim M_0(u)=d^*_u(T_u^k-1)^f$, with $d^*_u>0$.
\item With $Y=(D_*\Gamma^{-2}/d^*_u)^{1/(fk)}$, the level prime satisfies
$Y\le T_u\le(1+\frac1{fk})Y$.
\item $T_u\ge2$ and $f\le k$.
\end{itemize}
Then $\tfrac12\le d_u\Gamma^2/D_*\le e$. In particular the proviso holds
with $C_0=3$.
\end{lemma}
\noindent In this lemma $d_u=\dim M_0(u)$. It is the quantity called
$d^S_u$ in Lemma~\ref{i4:Llem-lvl} and in Proposition~\ref{i4:Lprop-forms},
where $d_u$ denotes the rank of the division lattice $\nu_u$.
\begin{proof}
For the upper bound,
$(T_u^k-1)^f\le T_u^{fk}\le(1+\frac1{fk})^{fk}Y^{fk}\le eY^{fk}$, so
$d_u\le eD_*\Gamma^{-2}$.

For the lower bound, $2f\le2^k\le T_u^k$. By Bernoulli's inequality,
$(1-T_u^{-k})^f\ge1-fT_u^{-k}\ge\frac12$. Hence
$(T_u^k-1)^f\ge\frac12T_u^{fk}\ge\frac12Y^{fk}$.
\end{proof}

\begin{lemma}[the lowest-space identity]\label{i4:lem-lowest}
Let $G$ be the unitary group of a rank-$k$ hermitian form over $E/F$, and
let $U\subset G(\mathbb A_{F,f})$ be open compact with a character $\chi_U$.
Let $y_c$ ($c\in\mathrm{cls}_0$) represent $G(F)\backslash G(\mathbb A_{F,f})/U$,
and suppose $G(F)\cap y_cUy_c^{-1}=1$ for all $c$. Let $V_\lambda$ be the
span of the highest-weight polynomial functions of the frames for weights
$\lambda=(\lambda_\tau)_\tau$. Let $M_0$ be the space of functions
$\varphi(x,z)$ satisfying both of the following:
\begin{itemize}
\item $\varphi(\gamma x\kappa,\gamma z)=\chi_U(\kappa)\varphi(x,z)$;
\item $\varphi(x,\cdot)\in V_\lambda$.
\end{itemize}
Then
$\dim M_0=|\mathrm{cls}_0|\prod_\tau\dim V(\lambda_\tau)$.
\end{lemma}
\begin{proof}
The restriction map $\varphi\mapsto(\varphi(y_c,\cdot))_c$ goes into
$\bigoplus_cV_\lambda$. Every $x$ can be written as $\gamma y_c\kappa$ with
$c$ unique. Because the stabilisers are trivial, $\gamma$ and $\kappa$ are
unique as well. So the restriction map is injective. Conversely, given
$(\varphi_c)_c$, the formula
$\varphi(\gamma y_c\kappa,z)=\chi_U(\kappa)\varphi_c(\gamma^{-1}z)$ is well
defined. It lies in $M_0$ because $V_\lambda$ is stable under the action on
frames. Finally $V_\lambda$ is the tensor product over the places $\tau$
of the irreducible representations $V(\lambda_\tau)$, so
$\dim V_\lambda=\prod_\tau\dim V(\lambda_\tau)$; Weyl's dimension formula
gives each factor.
\end{proof}

\begin{proposition}[zero estimates and the explicit formula]\label{i4:thm-litthin}
There are constants $C_1,c,C>0$ such that the following three statements
hold. Zeros are counted with multiplicity. Characters are Dirichlet
characters; a primitive character of conductor $1$ is the trivial one,
with $L(s,1)=\zeta(s)$. For $\chi$ primitive of conductor $N$ and $u\ge1$ put
$\psi_0(u,\chi)=\sum_{n\le u}\Lambda(n)\chi(n)-\delta_{N=1}\,u$.
\begin{itemize}
\item[$(\mathrm A_6)$]
\begin{enumerate}
\item[(i)] Let $Q\ge3$, and count the zeros $\rho$ of $L(s,\chi)$ for all
primitive $\chi$ of conductor $\le Q$, with $\Rel\rho\ge1-1/400$ and
$|\Iml\rho|\le Q$. Their number is at most $C_1Q^{1/4}$.
\item[(ii)] Let $\chi$ be primitive of conductor $N\ge3$, and let $\rho$ be
a zero of $L(s,\chi)$ with
$\Rel\rho\ge1-c/\log\max(N,N|\Iml\rho|)$. Then $\chi$ is real, $\rho$ is
real and simple, and $\rho$ is the only zero in that region.
\item[(iii)] Every zero $\rho$ of $\zeta$ with $0<\Rel\rho<1$ satisfies
$\Rel\rho\le1-c/\log(|\Iml\rho|+3)$.
\end{enumerate}
\item[$(\mathrm A_7)$] Let $\chi$ be primitive of conductor $N$ and
$2\le T\le u$. The set of zeros with $0<\Rel\rho<1$, $|\Iml\rho|<T$ is finite, and
\[
\Bigl|\psi_0(u,\chi)+\sum_{|\Iml\rho|<T}\frac{u^\rho}\rho\Bigr|\le
C\Bigl(\frac{u\log^2(NuT)}T+\log(Nu)+\sum_{|\Iml\rho|<1}\frac1{|\rho|}\Bigr).
\]
\item[$(\mathrm A_{11})$] Let $\chi$ be primitive of conductor $N$ and
$t\in\R$. The number of zeros with $0<\Rel\rho<1$, $|\Iml\rho-t|\le1$ is finite and at most
$C\log(N(|t|+2))$.
\end{itemize}
Here the sums in $(\mathrm A_7)$ run over zeros with $0<\Rel\rho<1$,
and in $(\mathrm A_6)$(i), (ii) a zero is any zero of $L(s,\chi)$ in the
stated region.
\end{proposition}
Each item follows from a classical theorem (all in Davenport,
\emph{Multiplicative Number Theory}, and Montgomery, \emph{Topics in
Multiplicative Number Theory}, LNM 227 (1971); numbers as given below). We
state them and derive the items. Throughout, $\chi$ is primitive of
conductor $N$, $\sigma=\Rel s$, $t=\Iml s$, and $c_\ast,C_\ast$ are
absolute constants.

\begin{theorem}[classical zero estimates; Davenport, \emph{Multiplicative Number Theory}, 3rd ed., GTM 74, Chs.~13--17, 19; Montgomery, LNM 227, Thm.~12.2; reference to be confirmed]\label{i4:thm-classzeros}
\begin{enumerate}
\item[(Z0)] $L(s,\chi)\ne0$ for $\sigma\ge1$ (for $\chi=1$: for $\sigma\ge1$, $s\ne1$).
\item[(Z1)] (de la Vall\'ee Poussin) $\zeta(s)\ne0$ for
$\sigma\ge1-c_\ast/\log(|t|+2)$.
\item[(Z2)] (Gronwall, Landau, Titchmarsh; Davenport, Ch.~14) If $\chi$ is a nonprincipal character modulo $q$, then
$L(s,\chi)$ has no zero with $\sigma\ge1-c_\ast/\log(q(|t|+2))$, except,
when $\chi$ is real, at most one zero, which is real and simple.
\item[(Z3)] The number of zeros $\rho$ of $L(s,\chi)$ with $0<\Rel\rho<1$
and $|\Iml\rho-t|\le1$ is at most $C_\ast\log(N(|t|+2))$.
\item[(Z4)] (truncated explicit formula) For $2\le T\le u$,
\[
\psi(u,\chi)=\delta_{N=1}u-\sum_{|\Iml\rho|<T}\frac{u^\rho}\rho-(1-a)\log u-b(\chi)
+\sum_{m\ge1}\frac{u^{a-2m}}{2m-a}+R(u,T),
\]
with $a\in\{0,1\}$ the parity of $\chi$, the sum over zeros with
$0<\Rel\rho<1$, $|R(u,T)|\le C_\ast(u\log^2(NuT)/T+\log u)$, and
$|b(\chi)|\le C_\ast\log(N+2)+\sum_{|\Iml\rho|<1}1/|\rho|$ (for $N=1$,
$b=\log2\pi$ and the term $(1-a)\log u$ is absent; the sum over $m$ then
equals $-\frac12\log(1-u^{-2})$).
\item[(Z5)] (zero density, Montgomery, Thm.~12.2; reference to be confirmed,
statement to be confirmed against the source) For $Q,T\ge2$, writing
$Z(\sigma)=\sum_{q\le Q}\sideset{}{^*}\sum_{\chi\bmod q}N(\sigma,T,\chi)$,
\[
Z(\sigma)\le
C_\ast(Q^2T)^{3(1-\sigma)/(2-\sigma)}\log^9(QT)\quad(\tfrac12\le\sigma\le\tfrac45),\qquad
Z(\sigma)\le
C_\ast(Q^2T)^{2(1-\sigma)/\sigma}\log^{14}(QT)\quad(\tfrac45\le\sigma\le1),
\]
where $\sum^*$ runs over primitive characters and $N(\sigma,T,\chi)$
counts the zeros with $\Rel\rho\ge\sigma$, $|\Iml\rho|\le T$.
\end{enumerate}
\end{theorem}

\begin{proof}[Proof of Proposition~\ref{i4:thm-litthin}]
Take $c=\min(c_\ast/2,\,\tfrac12\log3)$.

$(\mathrm A_6)$(i). By (Z0) every zero counted has $\Rel\rho<1$. Apply
the second range of (Z5) with $T=Q$ and $\sigma=399/400\ge4/5$: the
exponent is $2(1/400)/(399/400)=2/399$, so the count is at most
$C_\ast(Q^3)^{2/399}\log^{14}(Q^2)=C_\ast Q^{2/133}\log^{14}(Q^2)\le C_1Q^{1/4}$
for all $Q\ge3$, with $C_1$ absolute.

$(\mathrm A_6)$(ii). Let $N\ge3$ and put $\ell(t)=\log\max(N,N|t|)\ge\log3$.
Since $N(|t|+2)\le3N\max(1,|t|)$, we have
$\log(N(|t|+2))\le\log3+\ell(t)\le2\ell(t)$, hence
$c/\ell(t)\le c_\ast/\log(N(|t|+2))$. So every zero with
$\Rel\rho\ge1-c/\ell(\Iml\rho)$ lies in the region of (Z2) (with $q=N$),
and (since $c\le\frac12\log 3$) has $\Rel\rho\ge\frac12$, so it is not a
trivial zero. By (Z2) $\chi$ is real, $\rho$ is real and simple, and it
is the only zero of $L(s,\chi)$ in the region of (Z2), a fortiori the
only one in the smaller region.

$(\mathrm A_6)$(iii). If $0<\Rel\rho<1$ and
$\Rel\rho>1-c/\log(|\Iml\rho|+3)$, then since $\log(|t|+3)\ge\log(|t|+2)$
and $c\le c_\ast$, $\rho$ lies in the region of (Z1), a contradiction.

$(\mathrm A_{11})$ is (Z3).

$(\mathrm A_7)$. Finiteness follows from (Z3), covering $|t|<T$ by
finitely many unit intervals. In (Z4), $\psi_0(u,\chi)=\psi(u,\chi)-\delta_{N=1}u$
up to the term $\Lambda(u)\chi(u)$ when $u$ is an integer, which is at
most $\log u$. Moving the sum over zeros to the left, the remaining terms
are bounded by $\log u$ (or $|\tfrac12\log(1-u^{-2})|\le1$ when $N=1$), $|b(\chi)|$, $\sum_{m\ge1}u^{a-2m}/(2m-a)\le\sum_{m\ge1}u^{1-2m}\le1$,
and $|R|$. Since $u\ge2$ and $N\ge1$, $\log u$, $\log(N+2)$ and $2$ are
all at most $4\log(Nu)$, and $\log(Nu T)\ge\log u$; so the total is at
most $C(u\log^2(NuT)/T+\log(Nu)+\sum_{|\Iml\rho|<1}1/|\rho|)$ with $C$
absolute.
\end{proof}

\begin{proposition}[the working class]\label{i4:prop-working}
Let $\mathcal A$ be a family class with a twist selection whose members
satisfy the seven duties. Then:
\begin{enumerate}
\item the working class of $\mathcal A$ is closed for the rank ties, that
is, for the conclusion of Lemma~\ref{i4:lem-lowest} at every word;
\item under the class number formula at all points, the working class is
closed for the proviso with constant $3$;
\item every class is closed for the thinness implication of
Theorem~\ref{i4:thm-choices};
\item $\mathrm{DesignChoices}(\mathrm{WithTw}(\mathcal A),3,c_T)$ implies
the same statement for the working class.
\end{enumerate}
\end{proposition}
\begin{proof}
(1) Each member carries, beyond a threshold, a tied datum. This datum
supplies the data of Lemma~\ref{i4:lem-lowest} at every word with trivial
stabilisers. The level primes split completely in $E$, because of the
level conditions and the congruence on $T_u$. Apply
Lemma~\ref{i4:lem-closed}(e).

(2) Lemma~\ref{i4:lem-proviso}, combined with (1), gives the proviso with
constant $e<3$. Its hypotheses hold at every point with $x>x_1$, for the
following reasons. The class number formula is the hypothesis of (2). The
window $Y\le T_u\le(1+\frac1{fk})Y$ is a field of the column data of every
member: the level prime is chosen in that window by the level choice of
the realisation (\S\ref{i4:ssec-proof-realise}, step 3), and the
identifications of $d^*_u$ and $D_*$ with the quantities defining $Y$ are
fields of the family at every point with $x>x_1$. $T_u\ge2$ because $T_u$
is prime, and $f\le k$ for large $x$ by validity of the ranks
($f\le4x^{2+8\varepsilon}<x^{9/2}\le k$). Trivial class-$0$ stabilisers are
part of the data supplied in (1). The proviso is monotone in $C_0$.

(3) By Proposition~\ref{i4:thm-litthin}, at a host with boxes and
$(q,d)$-rows we have $|G_\chi(y)|\le y\exp(-cx^{\varepsilon/8})$ for every
class character $\chi$ and every $y\ge\exp(x^{\varepsilon/3})/3$. The
argument is the truncated explicit formula $(\mathrm A_7)$ combined with
the zero-free region and the zero-density estimate $(\mathrm A_6)$ and the
local zero count $(\mathrm A_{11})$; the box condition is what excludes an
exceptional zero, via $(\mathrm A_6)$(ii). This uses
$\log N_E\le2x^{\varepsilon/5}\log^2x$, which comes from the $(q,d)$-rows.
The bound holds for every family beyond a threshold.

(4) Each of the conditions added by $\mathrm{WithChN}$ and by
$\mathrm{WithSm}$ is among the eventual conditions of the design choices.
So the same witnesses serve.
\end{proof}

\begin{theorem}[design existence for the working class]\label{i4:thm-design}
Let $\mathcal A$ be a family class with a twist selection whose members
satisfy the seven duties, and let $c_T:\R^2\to\R$. Assume:
\begin{itemize}
\item $\mathrm{DesignChoices}(\mathrm{WithTw}(\mathcal A),3,c_T)$;
\item the working class is closed for
$c_T(\varepsilon,C_{YL})\le C_{T_u}\Rightarrow\mathrm{TuSize}$;
\item the working class is closed for the class number formula at all
points.
\end{itemize}
Then designs exist for (the working class of $\mathcal A$, $3$).
\end{theorem}
\begin{proof}
Combine Proposition~\ref{i4:prop-working}(2)--(4) with
Theorem~\ref{i4:thm-choices}.
\end{proof}

\subsection{The closed estimates on the working class}\label{i4:ssec-links}

Write $D_3(\mathcal A)=\mathrm{Des}(\mathrm{WithChN}(\mathcal A),3)$ for a
class $\mathcal A$ carrying a hypothesis structure
(Definition~\ref{i4:def-hyps} below). Lemma~\ref{i4:lem-closed} lets us
treat each of the nine properties of Theorem~\ref{i4:thm-desWin}(4)
separately.

\begin{lemma}[rank equality at level three]\label{i4:lem-rankeq}
Let $\mathcal Z$ and $M_j$ be finite-dimensional inner product spaces. Let
$\theta_P:\mathcal Z\to M_j$ be linear, with adjoint $\theta_T$, and put
$Z_j=\theta_T(M_j)$. Make the following assumptions:
\begin{itemize}
\item $M_j$ has a basis $(\phi_k)_{k\le d_j}$;
\item $d_j=\dim S_0$, where $S_0=\bigoplus_cE_c$;
\item $\dim E_c=\sum_{c(\sigma)=c}n_0(\sigma)$, where $n_0(\sigma)$ is the
size of the fibre of a map $\mathrm{ft}_0:\{1,\dots,D_0\}\to\Sigma$;
\item there are an invertible matrix $\Theta$ and a surjective matrix
$\ell$ such that every $\sum_k(\Theta\ell y)_k\phi_k$ lies in the image of
$\theta_P$.
\end{itemize}
Then $\theta_T$ is injective and $\dim Z_j=D_0$.
\end{lemma}
\begin{proof}
Since $\Theta\ell$ is surjective, every $\phi_k$ lies in the image of
$\theta_P$. So $\theta_P$ is onto, and its adjoint is injective: if
$\theta_T w=0$ with $w=\theta_Pv$, then
$\|w\|^2=\langle v,\theta_Tw\rangle=0$. Therefore
$\dim Z_j=d_j=\sum_c\dim E_c=\sum_\sigma n_0(\sigma)=D_0$.
\end{proof}

\begin{proposition}[closed estimates from the inputs]\label{i4:prop-links-proved}
Let $\mathcal A$ carry a hypothesis structure, and assume the loss bound
and the level counts and Theorem~S. Then $D_3(\mathcal A)$ is closed for
each of the following:
\begin{itemize}
\item $(\mathrm{LEta})$;
\item $(\mathrm{SLR})$;
\item $(\mathrm{LIrr})$;
\item $\mathrm{TuSize}\Rightarrow(\mathrm{LS}(3))$;
\item $(\mathrm{LG})$.
\end{itemize}
\end{proposition}
\begin{proof}
We take the five properties in turn.

$(\mathrm{LEta})$: its three clauses are clauses of the Galois component
of the tied datum, which the twist selection provides beyond a threshold.
Apply Lemma~\ref{i4:lem-closed}(e).

$(\mathrm{SLR})$: beyond the selection threshold, the selection supplies
at every point a level-three realisation. Its rows give the hypotheses of
Lemma~\ref{i4:lem-rankeq}:
\begin{itemize}
\item the basis comes from the theta spaces;
\item the level-zero rows follow from primitivity of the realisations;
\item the $\Lambda'$-rows are a field of the structure;
\item the divisor rows are a field of the structure;
\item the zero-locus field gives the surjectivity of $\ell$.
\end{itemize}
The lemma gives $\dim Z_{k+1}(\mu_0)=D_0$.

$(\mathrm{LIrr})$ and $\mathrm{TuSize}\Rightarrow(\mathrm{LS}(3))$: these
are Proposition~\ref{i4:hyp-S}(b) and (a).

$(\mathrm{LG})$: a member lies in the second-layer class
(Proposition~\ref{i4:prop-working}). Its tied data contain a group tie,
which is group compatibility. $\mathrm{ChosenE}(3,C_H)$ contains the
bounds on the tops and the bad-prime bound. With the Mertens-type estimate
$\sum_{p\le e^T}\log p/p\le2T(T+2)$, these give the second admissibility
condition for $x\ge(16C_b+2)^{3/\varepsilon}$. Proposition~\ref{i4:hyp-loss}(a)
gives $(\mathrm{LG})$.
\end{proof}

The three remaining properties are where the column package is built. We
state them here. Their proof, in \S\ref{i4:ssec-proof-links}, reduces them
to the division bridge, to the package lemmas, to the kernel join and to
explicit constant estimates, all proved there. What is left is a set of
pointwise leaves, stated as Proposition~\ref{i4:prop-links-leaves}.

\begin{proposition}[closed estimates on the division and global side]\label{i4:prop-links-formal}
Let $\mathcal A$ carry a hypothesis structure, and assume Green's
hypothesis, the loss bound and the level counts and Theorem~S. Assume also
Proposition~\ref{i4:prop-links-leaves}. Then $D_3(\mathcal A)$ is closed
for each of the following:
\begin{enumerate}
\item[(a)] $\mathrm{TuSize}\wedge\mathrm{DGamSize}\Rightarrow(\mathrm{LDivt})$
with $C_{\rm abs}=2$;
\item[(b)] $\mathrm{ChosenE}(3,C_H)\Rightarrow$ (every $\nu_u$ is
eventually full and of positive rank);
\item[(c)] $\mathrm{TuSize}\wedge\mathrm{DGamSize}\Rightarrow(\mathrm{LGl})(2,C_H,3,2,0)$.
\end{enumerate}
In particular the exponents are $e_D=2$ and $e_L=0$.
\end{proposition}

\begin{proposition}[the leaves of the division and global links]\label{i4:prop-links-leaves}
Let $\mathcal A$ carry a hypothesis structure. The following hold on
$D_3(\mathcal A)$, each beyond a member-dependent threshold.
\begin{enumerate}
\item[(L0)] (\emph{Clause inputs.}) $D_3(\mathcal A)$ is closed, under
$\mathrm{TuSize}\wedge\mathrm{DGamSize}$, for each of the clauses adm2,
cut, c2, c3, the frame-number bound (hh), $\mathrm{BuRational}$ and
$\mathrm{RRFull}$ of $\mathrm{DKCB}(\cdot,2)$ (Definition~\ref{i2:def-DKCB}),
and for the inputs of c5: along the selection, eventually the gap
hypothesis and the finiteness property $\mathrm{FinX}$ of the selected
spectral datum hold, the multiplicity condition of the selected indefinite
realisation holds, and the base-vector image is full.
\item[(L1)] (\emph{Arithmetic clause.}) At every point satisfying the
hypotheses of the wiring step (\S\ref{i4:ssec-proof-links}, step (C)),
with the constants $c_0,c_e$ of \eqref{i4:eq-c0ce} and the integer $d_g$ of
Lemma~\ref{i4:lem-dgam}, the arithmetic clause of the division package
holds at every word.
\item[(L2)] (\emph{Package lemma of the division bridge.}) At the same
points, the base package with full Hecke data of the wiring step, the tied
rows, the columns of the realisation, the divisor rows, the
decomposition-point rows, the $\kappa$-rows, the height formula and the
frame tie yield a division kernel package with constants
$(\max(0,50c_0+10,60c_e),2)$.
\item[(L3)] (\emph{Level-zero leaves.}) At every word in the Hida setting,
the concrete forms at level zero imply the slope line, the length bound of
the second-branch inclusion, the integral Hecke condition and the ramified
Hecke condition at level zero.
\item[(L4)] (\emph{Family-level steps.}) Under $\mathrm{TuSize}$, the
eleven-item minimal global record implies the fifth family part, the two
slope conclusions, the rank condition of the ramified-rank family, the
count, the lattice condition and the frame bound $\mathrm{FrameF}(C_H)$,
with the inputs listed in step (5) of \S\ref{i4:ssec-proof-links}; and
the core link: these items, with the losses of
Proposition~\ref{i4:hyp-loss}(b), $(\mathrm{LIrr})$ and $(\mathrm{LDiv})$,
give constants $K,c_{HH},e_K,C_L,x_0$ such that the family is chain-tied
with $(C_{\rm abs},0,K,c_{HH},e_K,C_L)$ at every $x\ge x_0$ in
$\mathrm{Win}$.
\end{enumerate}
\end{proposition}

\noindent The proof is in \S\ref{i4:ssec-proof-Lleaves}. It is complete
except for the five formally verified Propositions~\ref{i4:Lprop-sel},
\ref{i4:Lprop-D}, \ref{i4:Lprop-pkg}, \ref{i4:Lprop-forms} and
\ref{i4:Lprop-fam} below, which are strictly smaller than (L0)--(L4): the
assembly of the clauses, the arithmetic clause, the reduction of the
package lemma to the division inputs, the height bound, the gauge form
bound, the commensurability and fullness of the division lattices, the
rank positivity, the congruence length bound and the loss requirements are
all proved in \S\ref{i4:ssec-proof-Lleaves}. We use the notation of the
size record: at a point $(x,i)$ with $x>x_1$ we write $L=x^\varepsilon$,
$b=e^L$, $a=xb$, $X'=\exp(x^{\varepsilon/3})$, $b'=a/X'$; $f=2n,k,m,t,r$ for
the ranks, $N_m=k+m(t+1)$, $B_0=fmbL$; $\mathrm{Ga}$, $\mathrm{Fo}$, $E_e(u)$
for the gauge, frame and excess terms at a word $u$; $\Lambda^0_u$ for the
level-zero lattice and $\nu_u$ for the division lattice at $u$, of rank
$d_u$; $RR$ for the row lattice of the division datum and $\kappa_R$ for its
archimedean row constants. Put
\[
M_t=\max(C_t,5)+3,\qquad P=1+f+k+m+t+N_m+L\quad\text{(the size polynomial)},
\]
$E_b=4M_t+1$ and $E_c=9M_t+2$, $E_d=9M_t+1$.

\begin{proposition}[the selection clauses; (a) formally verified, \cite{CharacterFamily}]\label{i4:Lprop-sel}
\begin{enumerate}
\item[(a)] $D_3(\mathcal A)$ is closed for each of the properties
``$\mathrm{TuSize}\wedge\mathrm{DGamSize}\Rightarrow$ clause'', where the
clause is adm2, cut, c2 or c3 of $\mathrm{DKCB}(\cdot,2)$;
\item[(b)] every class $\mathcal C$ of families all of whose members
satisfy $\mathrm{TuSize}$ is closed for the property ``eventually, the gap
hypothesis holds at the datum $D(x,i)$ of every point'' ($\mathrm{GapEv}$);
\item[(c)] for every datum $D$ satisfying the gap hypothesis and every
spectral tie $T$ of $D$, $\mathrm{FinX}(T)$ holds.
\end{enumerate}
\end{proposition}
\noindent The structure of the formal proof of (a) is described in \S\ref{i4:Lst-sel}.
\noindent\emph{Remark.} In (b) the hypothesis that every member
of the class satisfies $\mathrm{TuSize}$ is needed. It gives
what (L0) uses: if $\mathcal C$ is the class of members of
$D_3(\mathcal A)$ satisfying $\mathrm{TuSize}$, a large member of
$D_3(\mathcal A)$ that satisfies $\mathrm{TuSize}$ is a large member of
$\mathcal C$, so $D_3(\mathcal A)$ is closed for
``$\mathrm{TuSize}\Rightarrow\mathrm{GapEv}$''. (c) has no
hypothesis besides the gap: Wu decay and
finite dimensionality of rows are consequences. Its proof:
the gap hypothesis gives Wu decay and integrability of the archimedean
coefficient (the decay estimates of
Section~\ref{i3:sec}, \S\ref{i3:ssec-bessel} and Lemma~\ref{i3:lem-decay}); finite
dimensionality of the rows is
Proposition~\ref{i4:A-prop-inputs-leaves}(M5); and
these two facts give $\mathrm{FinX}$.
\noindent What remains: (a) is the admissible twist link (the period
valuation bound along the class, from the six period fields, the field
$\mathrm{aP10}$, the generator-prime field, the service field and the
coprimality of the level primes to the stabiliser discriminant), the
$\mathcal O_{LL}$-stability of the lattices $\tilde M_\mu$, and the second
complement; (b) and the implication are part of Theorem~S. The
closedness wrappers are short; the content is
in Sections~\ref{i1:sec} and~\ref{i3:sec}.

\begin{proposition}[the analytic size estimates of the division side; formally verified, \cite{CharacterFamily}, except the first claim of (e)]\label{i4:Lprop-D}
In (a), (b), (e), (f) let $\mathcal F\in D_3(\mathcal A)$ satisfy
$\mathrm{TuSize}$ and $\mathrm{DGamSize}$; (c) and (d) are statements
about the class $D_3(\mathcal A)$.
\begin{enumerate}
\item[(a)] (\emph{maximal slope}) eventually, for every $i$ and word $u$,
$\mu_{\max}(\Lambda^0_u)\le x^{E_b}b$;
\item[(b)] (\emph{slope lemma}) if moreover the level conditions,
$\mathrm{MuRank}$, $(\mathrm{LG})$ and $(\mathrm{LMu})$ hold, there are $C,x_4$
such that for $x\ge x_4$, every $i$ and $u$, at every point where the
definitional identities of the point hold,
$\mu_{\max}(\nu_u)\le\mu_{\max}(\Lambda^0_u)+\mathrm{Ga}+\mathrm{Fo}-E_e(u)+C\,m\,B_0$
(the statement also gives a lower bound for $\mu(\nu_u)$, which is not
used here);
\item[(c)] (\emph{rows}) if the selection carries rows 14 (primitivity),
16 (divisor condition), 17 (decomposition-point condition) and 25
(definite rows, closed on $D_3(3,\mathcal A)$) of the row field, then
$D_3(\mathcal A)$ is closed for the property: $\mathrm{TuSize}$ and
$\mathrm{DGamSize}$ imply that eventually $\mathrm{BuRational}$ holds and
eventually $\mathrm{RRFull}$ holds;
\item[(d)] (\emph{sharp functional bound}) if the selection carries rows
12 (bounded rows), 14, 16 and 25 of the row field, then there is a
function $C_{F0}$ of the outer constants such that $D_3(\mathcal A)$ is
closed, with the threshold datum all of whose middle entries equal
$\max(E_b,E_c,C_{F0},E_d)+1$ (so a member is large at $C_H$ when every
entry of $c_{\rm mid}(C_H)$ is at least this number), for the property:
$\mathrm{DIVinN}$ implies that eventually
$\mu_{\max}(RR)+\mathrm{mean}_e\log\kappa_R(e)\le x^{C_{F0}}b$;
\item[(e)] (\emph{pointwise size bounds}) at every point with $x>x_1$ and
$x\ge9$ (where the size conditions $\mathrm{SizesAt}$ and the parameter
fills are fields of the family, and Mertens' bound
$\sum_{p<y}\log p/p\le2\log y$ is used for $\mathrm{Ga}$): $P\le x^{M_t}$; $\mathrm{Ga}\le64P^4b$ and
$\mathrm{Fo}\le256P^9b$; and, if $\log T_u\le C_{T_u}mL$ for all $u$,
\[
-\bigl(h_6(u)+\mathrm{mean}_e\log\kappa_R(e)\bigr)\le(500+2|C_{T_u}|+30|C_{sz}|)P^9b,
\]
where $h_6(u)$ is the explicit height term of the point;
\item[(f)] (\emph{slope count}) if moreover every $\nu_u$ is eventually full
and $\mathrm{MuRank}$ holds, then $(\mathrm{LMu})$ holds.
\end{enumerate}
\end{proposition}
\noindent The structure of its formal proof is described in \S\ref{i4:Lst-D}.
\noindent\emph{Remark.} (a) needs a twist
selection of the family (members of $D_3(\mathcal A)$ carry one) and
$\mathrm{TuSize}$ only; $E_b=\mathrm{EbMu}=4(\max(C_t,5)+3)+1=4M_t+1$.
(b) has hypotheses
$\mathrm{LvlF}$, $\mathrm{TuSize}$, $\mathrm{MuRank}$, $(\mathrm{LG})$,
$(\mathrm{LMu})$, as stated. In (c), (d) the row hypotheses (rows 14, 16, 17, 25 for (c); rows 12, 14, 16, 25
for (d)) are named explicitly; they are hypotheses of the statements, outside the closed
property. In (e) the term $\mathrm{FDef}$ (frame definition of the point) is also needed; it is a field of the family at every point with $x>x_1$; at
family points the size conditions, the fills and Mertens' bound supply the remaining estimates.
\noindent What remains: (a) is a column-value argument for stabilisers
and an upper slope bound for the level-zero lattice; (b) is the slope
comparison between $\nu_u$ and $\Lambda^0_u$ through the places bound and
the Mertens-type sums over tops; (c) is period descent along the rows; (d)
is a pointwise bound for the row lattice plus threshold bookkeeping; (e)
are explicit estimates term by term; (f) is the slope count from the
size bounds of Proposition~\ref{i4:prop-mu}(a), (b) without using
$(\mathrm{LDiv})$, the vanishing above the levels being
Lemma~\ref{i4:lem-zeroabove}.

\noindent Parts (b) and (c) of Proposition~\ref{i4:Lprop-sel} and the
first claim $P\le x^{M_t}$ of Proposition~\ref{i4:Lprop-D}(e) are proved in
\S\ref{i4:ssec-proof-selD}, modulo the following smaller statement.

\begin{proposition}[gap leaves; formally verified, \cite{CharacterFamily}]\label{i4:Lprop-gapleaves}
\begin{enumerate}
\item[(G1)] (\emph{Theorem~S, gap row}) For every class $\mathcal C$, if
$\mathcal C$ is closed for the residual eventuality $\mathrm{ResidEv}$, then
there is $\varepsilon_1>0$ such that for every $\varepsilon\in(0,\varepsilon_1]$ there is a
threshold datum $\kappa$ with: every family $\mathcal F$ of exponent
$\varepsilon$ that belongs to $\mathcal C$, satisfies $\mathrm{TuSize}$
and is $\kappa$-large at $C_H$ satisfies $\mathrm{GapEv}$.
\item[(G2)] Every class is closed for $\mathrm{ResidEv}$.
\item[(G3)] (\emph{Wu decay}) at every datum satisfying the gap hypothesis,
every Hida setting has Wu decay and integrable archimedean coefficient.
\item[(G4)] (\emph{finiteness from decay}) if the rows are finite
dimensional ($\mathrm{RowsFinDim}$), then for every datum, every Hida
setting $X$ with Wu decay and integrable archimedean coefficient, and every
full datum over $X$ (in particular every spectral tie), $\mathrm{FinX}$
holds.
\end{enumerate}
\end{proposition}
\noindent The structure of its formal proof is described in \S\ref{i4:Lst-gap}.

\begin{proposition}[division inputs at every word; formally verified, \cite{CharacterFamily}]\label{i4:Lprop-pkg}
At a point as in (L2), let $P_k$ be the core package of the wiring step.
For every tied division datum $I$ compatible with $P_k$ (and equal, after
transport, to the datum of the selection), the column bridge built from
$I$ is tied, and at every word there are division inputs with chain ties
and first columns, built from: the bound of Lemma~\ref{i4:Llem-etatw}, the
twist record, the base rank $D_0$, the lift exponent at the twisted places,
the arithmetic clause (L1) with $d_g$, and the row data given by the row
field.
\end{proposition}
\noindent The structure of its formal proof is described in \S\ref{i4:Lst-pkg}.
\noindent What remains is the construction of the division-input records
from these data and the identification of the column bridge.
(The statement is stated informally; the exact statement is in \cite{CharacterFamily}.
The objects ``core package'', ``tied division datum'' and ``division
inputs with chain ties and first columns'' have some twenty fields that we do
not define here.)

\begin{proposition}[the level-zero leaves; formally verified, \cite{CharacterFamily}]\label{i4:Lprop-forms}
This is (L3): at every word $u$ in the Hida setting, the concrete forms at
level zero imply (1) the slope line
$\mu(\Lambda^{\rm ram}_u)=\mu(\Lambda^0_u)-f\log(T-1)/(T-2)$, $T=T_u$;
(2) whenever $n_b=b_{\rm top}(\tau)\ge0$, the second-branch inclusion $T^{n_b}S_{0,w}\subseteq B_{0,w,v}$ with
length $\ell(B_{0,w,v}/T^{n_b}S_{0,w})=e_w\,d^S_u(\mathrm{wsum}_\tau+k\,b_{\rm top}(\tau))/k$;
(3) the integral Hecke condition $U_v(B_{0,w,v})\subseteq S_{0,w}$,
$D_v(S_{0,w})\subseteq B_{0,w,v}$; (4) the ramified Hecke condition: $U_v$
and $D_v$ preserve the ramified subspace and $U_vD_v=T^{N_0-1}=T^{k-1}$ on it. Here
$w$ is a place over $T$ seen by the form, $v\in VV_w$, $\tau=\tau_w(v)$ and
$e_w$ is the ramification index of $w$.
\end{proposition}
\noindent The structure of its formal proof is described in \S\ref{i4:Lst-forms}.
\noindent\emph{Remark.} The
hypothesis $b_{\rm top}(\tau)\ge0$ of (2) is needed; it is supplied by
Proposition~\ref{i4:Lprop-fam}(a) at every top.
\noindent What remains: (1) is the ramified slope formula from character
finiteness and admissible stabilisers; (2) comes from the rationality of the
level-zero forms; (3) is integrality at level zero (a local-global
comparison of lattices away from a place); (4) is a diamond-matching
argument.

\begin{proposition}[the family-level items; formally verified, \cite{CharacterFamily}]\label{i4:Lprop-fam}
Let $\mathcal F$ satisfy $\mathrm{TuSize}$ and the minimal global record.
Then:
\begin{enumerate}
\item[(a)] (\emph{setting}) eventually, at every word, the Hida setting
holds: $2fN_m+1<T_u$, $4a_{\max}(u)<T_u$, the alcove condition, and
$b_{\rm top}(\tau)\ge0$ at every top;
\item[(b)] (\emph{items}) let $C_{\rm abs}\ge0$, $C_H>0$ and $C_0\ge3$;
under the hypotheses of $(\mathrm{LGl})$ with these constants and the
small-shares record, eventually: the fifth family part, the two slope
conclusions, the rank condition of the ramified-rank family, the count, the
logarithmic-derivative bound, the maximal and normalised slopes, the
lattice condition, the real frame and $\mathrm{FrameF}(C_H)$, the nodes, the
character bound and the windows;
\item[(c)] (\emph{core link}) let $C_{\rm abs}\ge0$, $C_H>0$, $C_5$ and
$c_6>0$. If the family satisfies the count, the two slope conclusions
($\mathrm{SlopesM}$, $\mathrm{SlopesN}$), the lattice condition, the
nodes, $\mathrm{FrameF}(C_H)$, the losses $\mathrm{LossFS}(\cdot,0)$, the
character bound and the windows (the items of (b)), and moreover
$\mathrm{ChosenTT}$, $\mathrm{ChosenQD}(C_5)$, $\mathrm{ChosenThin}(c_6)$,
$(\mathrm{LIrr})$ and $(\mathrm{LDiv})$ with constants $(C_{\rm abs},C_H)$,
then there are $K,c_{HH},e_K,C_L,x_0$ such that the family is chain-tied
with $(C_{\rm abs},0,K,c_{HH},e_K,C_L)$ at every $x\ge x_0$ in
$\mathrm{Win}$.
\end{enumerate}
\end{proposition}
\noindent The structure of its formal proof is described in \S\ref{i4:Lst-fam}.
\noindent\emph{Remark.} In (c), $e_L=0$. The hypotheses $C_{\rm abs}\ge0$, $C_H>0$,
$c_6>0$, $\mathrm{ChosenTT}$, $\mathrm{ChosenQD}$, $\mathrm{ChosenThin}$,
$(\mathrm{LIrr})$ and $(\mathrm{LDiv})$ are all needed. In the use (proof
of (L4)) they are available: the chosen clauses are clauses of
$\mathrm{ChosenE}$, and $(\mathrm{LIrr})$, $(\mathrm{LDiv})$ are among the
inputs of $(\mathrm{LGl})$, which also require $C_H>0$ and $C_0\ge3$.
(b) is the list of items produced from
$\mathrm{TuSize}$, the setting, $\mathrm{LvlF}$ and the hypotheses of the
minimal link statement (the mass ratio, base volume $=$ mass, the
two-sided $\Sigma$-degree bound, the open slope parts, column values,
the free transports with slope line, and the three level-zero conditions
of Proposition~\ref{i4:Lprop-forms}).
\noindent What is not written out here are the family-level steps and
the five properties of chain-tied data derived from these items.


\subsection{The hypothesis structure}\label{i4:ssec-struct}

The class $\mathcal A$ is produced from bundles of automorphic origin. We
describe the logical frame and prove the reduction. The construction
itself is stated in \S\ref{i4:ssec-arith}.

\begin{definition}[eventual statements; slots; rows]\label{i4:def-ev}
A \emph{parameter} $p=(\varepsilon,C,c_D,c_h,C_5,c_6,C_b)$ consists of an
exponent, a record of constants and five reals. It is \emph{positive} if
$C_5,c_6,C_b>0$.

A statement $Q(p)$ holds \emph{for eventually all parameters} if the
following holds. There is $\varepsilon_0>0$ such that, for all
$\varepsilon\in(0,\min(\varepsilon_0,1/52))$, all $C_8\ge1$ and all
$C_5,c_6,C_b$, there is a threshold datum $\kappa_0$ with this property:
for every $\kappa_0$-large record $C$ there is $c_1$ such that for all
$c_D\ge c_1$ there is $c_2$ such that for all $c_h\ge c_2$, $Q(p)$ holds.

A \emph{slot} is a predicate on points of scale contexts. A \emph{linkage}
$\mathrm{Lk}$ relates points to bundles over them.
\begin{itemize}
\item $\mathrm{ExB}(Q,\mathrm{Lk})$ says: eventually, every $Q$-point
carries, for every $c$, a linked bundle.
\item A \emph{row} for a property $H$ of (class, selection), written
$\mathrm{Row}_Q(H)$, says: there is a property $E$ of bundles such that
eventually every linked bundle at a $Q$-point has $E$, and $H(\mathcal
A,\sigma)$ holds for every selection $\sigma$ whose bundles have $E$.
\item $\mathrm{LargeBase}(KP,Q,R)$ says: eventually, every core point
satisfying $Q$ over a base satisfying $KP$ also satisfies $R$.
\item $\mathrm{ExBL}(KP,Q,Q',\mathrm{Lk})$ says: there is $R$ with
$\mathrm{LargeBase}(KP,Q,R)$ and $\mathrm{ExB}(Q'\wedge R,\mathrm{Lk})$.
\item $\mathrm{RowL}(H)$ says: there is $R$ with $\mathrm{LargeBase}(KP,Q,R)$
and $\mathrm{Row}_{Q'\wedge R}(H)$.
\end{itemize}
\end{definition}

\begin{lemma}[elimination]\label{i4:lem-elim}
\begin{enumerate}
\item[(a)] Eventual statements are closed under conjunction and are
monotone under implication.
\item[(b)] $\mathrm{RowL}(H)\wedge\mathrm{RowL}(H')\Rightarrow\mathrm{RowL}(H\wedge H')$.
\item[(c)] $\mathrm{ExBL}$ and $\mathrm{RowL}(H)$ together give one $R$
with $\mathrm{LargeBase}(KP,Q,R)$, $\mathrm{ExB}(Q'\wedge R)$ and
$\mathrm{Row}_{Q'\wedge R}(H)$.
\item[(d)] Suppose $\mathrm{ExB}(Q)$ and $\mathrm{Row}_Q(H)$ hold. Then
there are a class $\mathrm{cls}(E)$ and a selection $\sigma_E$ on it,
consisting of bundles with property $E$, such that
$H(\mathrm{cls}(E),\sigma_E)$ holds.
\end{enumerate}
\end{lemma}
\begin{proof}
(a) Take the minimum of the exponents, the join of the threshold data, and
the maxima of $c_1$, $c_2$ and $x_0$.

(b) and (c): the large-base conditions combine by (a). A row for a slot
restricts to a row for a stronger slot. Two rows combine into a row for
$E\wedge E'$.

(d) Let $\mathrm{cls}(E)$ be the class of families linked to a parameter
$p$ at whose points every $c$ admits a $p$-built bundle with $E$. Choosing
one such bundle at every point defines $\sigma_E$, and the row gives
$H(\mathrm{cls}(E),\sigma_E)$.
\end{proof}

\begin{definition}[hypothesis structure]\label{i4:def-hyps}
A \emph{hypothesis structure} consists of a family class $\mathcal A$, a
selection $\sigma$ of bundles on $\mathcal A$, a function $c_T$, and the
following fields.
\begin{enumerate}
\item[(i)] $\mathrm{DesignChoices}(\mathrm{WithTw}(\mathcal A\wedge
\mathcal G.\mathrm{cls}(\lambda)),3,c_T)$ holds for every $\lambda>1$ and
every requirement $\mathcal G$ with $\mathrm{HostExW}(\lambda,\mathcal G)$.
\item[(ii)] The seven duties and lattice stability hold for every member.
\item[(iii)] The working class is closed for
$c_T\le C_{T_u}\Rightarrow\mathrm{TuSize}$ and for the class number
formula.
\item[(iv)] The thin-scale input (Proposition~\ref{i4:thm-litthin}), the
volume ratio and the rank-two genus mass hold.
\item[(v)] The selected bundles have the following pointwise properties:
bounded rows, twist records, primitivity, the $\Lambda'$-rows, divisor
rows, decomposition points, instance ranks and the realisation join.
\item[(vi)] Wu decay, finite dimensionality of rows, and the three
positivity statements hold.
\item[(vii)] $D_3(\mathcal A)$ is closed for the following: the torsion
frame; the loop row and Theorem~S; the point count; the base-vector
images; the definite rows; the six period conditions; the level-prime
valuations; the tuple bound; the second complement; zero-locus
surjectivity; the real-place ties; and the irreducibility count.
\item[(viii)] $D_3(\mathcal A)$ is closed for $(\mathrm{LN})$.
\end{enumerate}
\end{definition}

\begin{theorem}[class selection and the structure]\label{i4:thm-classSel}
Assume the Siegel mass ratio for CM fields and the volume ratio. Let
$KP,Q,Q'$, a linkage and $c_T$ be given. Assume the following:
\begin{itemize}
\item the four retwisting conditions (base, first, export, second);
\item $\mathrm{ExBL}(KP,Q,Q',\mathrm{Lk})$;
\item $\mathrm{RowL}(H)$ for each of the 26 row conditions that produce
fields (v)--(viii), for the counting row and for the level row.
\end{itemize}
Assume also the inputs in (iv) and (vi). Then a hypothesis structure
exists.
\end{theorem}
\begin{proof}
By Lemma~\ref{i4:lem-elim}(b), the rows combine into one row
$\mathrm{RowL}(H)$, where $H$ is their conjunction. By (c) and (d) we
obtain $R$, $E$, the class $\mathcal A=\mathrm{cls}(E)$ and the selection
$\sigma=\sigma_E$, with $H(\mathcal A,\sigma)$.

Field (i) is the realisation step: Theorem~\ref{i4:thm-realise} applied
with the slot $Q'\wedge R$ and the property ``every $c$ admits a built
bundle with $E$''. Fields (v)--(viii) are conjuncts of $H$. They pass to
the subclasses by Lemma~\ref{i4:lem-closed}(d). Field (ii), the seven
duties and lattice stability, comes from the property $QO$ of the
realisation step: every member of $\mathcal A=\mathrm{cls}(E)$ is linked to
a parameter at whose points $QO$ holds, and $QO$ contains the duties and
the $\mathcal O_{LL}$-stability of the lattices
(\S\ref{i4:ssec-proof-realise}, step 1). Fields (iii), (iv)
and (vi) are the hypotheses, together with the level row, which takes
$c_T(\varepsilon,y)=\max(11,16+3\max(5,y+1))$.
\end{proof}

\subsection{The arithmetic construction}\label{i4:ssec-arith}

The following results supply the hypotheses of
Theorem~\ref{i4:thm-classSel}. They form the automorphic core of the
existence argument: hosts and level primes in short windows, Hermitian
lattices of automorphic origin with prescribed data, and the 26 rows. The
construction itself is written out in \S\ref{i4:ssec-proof-Aleaves};
what remains of it is stated there as smaller formally verified
propositions.

\begin{theorem}[Hecke; Neukirch VII (6.1), (8.5), (8.6)]\label{i4:thm-hecke}
The Hecke $L$-function of a Gr\"ossencharakter $\chi$ of a number field has
a meromorphic continuation and a functional equation. It is entire unless
$\chi=|\cdot|^{iy}$ for some real $y$.
\end{theorem}
\noindent This theorem is background. The arithmetic of the hosts uses the
following two statements about primes.

\begin{theorem}[least prime in a progression; Linnik, Mat.\ Sb.\ 15 (1944); effective form Heath-Brown, Proc.\ London Math.\ Soc.\ (3) 64 (1992), 265--338]\label{i4:thm-linnik}
There are absolute constants $c_0,L_0\ge1$ such that for all $M\ge1$ and
all $a$ prime to $M$ the least prime $p\equiv a\pmod M$ satisfies
$p\le c_0M^{L_0}$.
\end{theorem}
\noindent We use it with $a=1$. Then $A=L_0+\log_2c_0+1$ gives a prime
$p\equiv1\pmod M$ with $p\le M^A$ for every $M\ge2$.

\begin{proposition}[norm primes in short intervals]\label{i4:thm-shortnorm}
There are constants $A,c,c'>0$ with the following property. Let $d\equiv
3\pmod 4$ be a prime with $d\ge100$, let $M\ge3$ be an integer, and let
$X\ge1$ and $h$ be reals with
\[
\log(4dM^2)\le c\sqrt{\log X},\qquad e^{-c\sqrt{\log X}}\le h\le1 .
\]
Then at least $c'hX/((4dM^2)^A\log X)$ primes $p$ with
$\lfloor X\rfloor<p\le\lfloor X(1+h)\rfloor$ can be written as
$4p=A_0^2+dB_0^2$ with integers $A_0,B_0$ such that $M\mid B_0$ and
$2M\mid A_0-2-B_0$.
\end{proposition}
\noindent The congruences say exactly that $\pi=(A_0+B_0\sqrt{-d})/2$
satisfies $\pi\equiv1\pmod{M}$. The proof is in
\S\ref{i4:ssec-proof-shortnorm}: an elementary reduction to a lower bound for
the trivial ray class in short intervals (Lemma~\ref{i4:lem-trivshort}),
which is proved there from four zero statements for $L$-families
(Proposition~\ref{i4:prop-lfam-zeros}) and the construction of the ray class
family (Proposition~\ref{i4:prop-lfam-ray}). The zero statements are
proved in \S\ref{i4:ssec-proof-lfam-zeros}; the construction of the ray
class family is proved in \S\ref{i4:ssec-proof-lfamray}, with its leaves
(Proposition~\ref{i4:prop-ray-leaves}) proved in \S\ref{i4:Rst-ray}.

\begin{proposition}[the root discriminant]\label{i4:prop-rd}
Let $\nu\ge2$, let $A$ be a finite abelian group of order $n$ with a
homomorphism $h:A\to H_\nu=(\Z/2^{\nu+2})^\times/\{\pm1\}$, let $Q_S\ge1$
and let $q,d$ be primes with $q\equiv3\pmod8$, $d\equiv7\pmod{16}$,
$2^{\nu+2}\mid Q_S$ and $(Q_S,qd)=1$. Put $N_E=Q_Sqd$ and
$G=\{\pm1\}^2\times A$. Let $E$ be a Galois number field with an
isomorphism $\ell:\mathrm{Gal}(E/\Q)\to G$, and let $\sqrt{-d}\in E$ satisfy
$\sigma(\sqrt{-d})=s_2(\ell(\sigma))\sqrt{-d}$, where $s_2$ is the second
sign. Let $L=\Q(\zeta)$ with $\zeta$ a primitive $N_E$-th root of unity,
$\tau_u(\zeta)=\zeta^u$, let $\iota:E\to L$ be a field embedding and
$\mathrm{Art}:(\Z/N_E)^\times\to G$ a homomorphism such that
\begin{itemize}
\item[(H1)] $\iota(\ell^{-1}(\mathrm{Art}(u))x)=\tau_u(\iota(x))$ for all
$u$ and $x\in E$;
\item[(H2)] $\mathrm{Art}(-1)=((-1,-1),1)$;
\item[(H3)] $h$ maps the $A$-component of $\mathrm{Art}(u)$ to the class of
$u$ in $H_\nu$;
\item[(H4)] $u\equiv1\pmod{Q_Sq}$ implies $\mathrm{Art}(u)\in\{1,((1,-1),1)\}$;
\item[(H5)] $u\equiv1\pmod{Q_Sd}$ implies $\mathrm{Art}(u)^2=1$.
\end{itemize}
Then, with $f=2n$,
\[
2^{\nu+1-2^{-\nu}}\sqrt d\le\mathrm{rd}_E\le\frac{Q_S}2\sqrt{qd},\qquad
2f\le\mathrm{rd}_E^2 .
\]
\end{proposition}
\noindent Every host has these data, so the bounds hold for every host. The
proof is in \S\ref{i4:ssec-proof-rd}. It is a conductor--discriminant
computation. The logarithm of the root discriminant is written as a sum
over the divisors of $N_E$ that counts how far the kernel of the Artin map
is from containing the congruence subgroups; the hypotheses (H2)--(H5)
then control the terms at $2$, $q$ and $d$.

\begin{proposition}[the kit parameters]\label{i4:prop-kits}
Let $\varepsilon\in(0,1/52)$, and let $C$ be a standard record of
constants. Eventually, at valid size data whose host is arithmetic, there
is a parameter choice whose constants are those of $C$ and whose window
constant is $10$. Moreover, the ten groups of size statements (E3) hold
eventually. These bound $k,m,f,\eta_{sz},|T|,\rho_P,b_T,\Delta,L_w,j_0$,
the node sizes and the word count by explicit powers of $x$ and $b$.
\end{proposition}
\noindent The proof is in \S\ref{i4:ssec-proof-kits}. It consists of
twenty-one elementary eventual inequalities for the first record and ten
groups for the third. All of them follow from four shared comparison
lemmas, the explicit rank bounds and Proposition~\ref{i4:prop-rd}.

\begin{theorem}[realisation]\label{i4:thm-realise}
Assume the Siegel mass ratio and the volume ratio. Suppose the base,
first, export and second retwisting conditions hold for $(KP,Q,Q')$, and
that a property $QO(p)$ holds eventually at every $Q'$-point. Assume
Minkowski's bound (Theorem~\ref{i4:thm-mink}),
Theorem~\ref{i4:thm-linnik}, Proposition~\ref{i4:thm-shortnorm}, and
Proposition~\ref{i4:prop-realise-leaves}. Then
$\mathrm{DesignChoices}(\mathrm{WithTw}(\mathcal A\wedge\mathcal G.\mathrm{cls}(\lambda)),3,c_T)$
holds for $\mathcal A=\mathrm{cls}(QO)$, for every $\lambda>1$ and every
$\mathcal G$ with $\mathrm{HostExW}(\lambda,\mathcal G)$.
\end{theorem}
\noindent The proof is in \S\ref{i4:ssec-proof-realise}. It places hosts
and level primes in the windows, glues the kits of
Proposition~\ref{i4:prop-kits} into families, builds the layered data in
two passes, and identifies volumes with masses. The elementary estimates
it needs that are not written out here are the following.

\begin{definition}[place data]\label{i4:def-placedata}
Let $x\ge1$, $\Delta\ge1$, $a,B>0$, $k\ge2$, $m\ge1$. Sequences
$u=(u_i)_{i<k}$, $v=(v_i)_{i<k-1}$, $b=(b_j)_{j<m}$, $w=(w_j)_{j<m}$ form
\emph{place data} if: (P1) $0\le v_i\le4k\Delta$; (P2) $v_i-v_{i+1}=4\Delta$;
(P3) $v_i<u_i$ and $u_{i+1}<v_i$; (P4) $-b_0<u_{k-1}$; (P5) either
$u_0\le2a$ and $-a\le u_{k-1}$, or $4k\Delta\le u_0-v_0\le8k\Delta$ and
$u_{k-1}=v_{k-2}-\frac12$; (P6) $b_j<w_j$ and $w_j-b_j\le b_j/x$; (P7)
$w_j<b_{j+1}$ and $b_{j+1}-b_j\ge1$; (P8) $4k\Delta\le b_0/x$ and
$w_{m-1}\le2B$. For a sequence $y$ of length $N$ put
$\dim_N(y)=\prod_{i<j<N}(y_i-y_j)/\prod_{i<N}i!$, and
$\mathrm{row}(v,b)=(v_0,\dots,v_{k-2},-b_0,\dots,-b_{m-1})$.
\end{definition}

\begin{proposition}[size leaves of the realisation]\label{i4:prop-realise-leaves}
Let $\varepsilon\in(0,1/52)$.
\begin{enumerate}
\item[(R1)] (\emph{Logarithmic ratio.}) Let $C_0\ge0$ and
$C=30+(2+2C_0)/\varepsilon$. Let $x\ge1$, $L=x^\varepsilon$, $B=e^L$,
$a=xB$, integers $f,m\ge1$, $k\ge2$ with $k\le x^{C_0}$, reals $v_E\ge\pi$
with $\log v_E\le L$, $\Delta\ge1$, $g_1,g_2\ge1$, place data
$(u_\rho,v_\rho,b_\rho,w_\rho)$ for $\rho<f$, and positive reals
$\mathrm{mass}_k,\mathrm{mass}_r,I_\sigma,I_\pi,\Lambda(M)$ $(k<M<k+m)$ with
$|\log\Lambda(M)|\le f\log\zeta(M)$. Put $r=k-1+m$ and assume
\[
\mathrm{mass}_r=\mathrm{mass}_k\prod_{M=k+1}^{r}\frac{((M-1)!(v_E/\pi)^M)^f}{\sqrt{(2v_E)^f}}\Lambda(M),\qquad
2^{-2f}(g_1g_2)^{fm(2k+m-1)/2}\le\frac{I_\sigma}{I_\pi},
\]
$D_*\ge(2/3)^f\mathrm{mass}_rI_\sigma\prod_\rho\dim_r(\mathrm{row}(v_\rho,b_\rho))$
and $d_*=\mathrm{mass}_kI_\pi\prod_\rho\dim_k(u_\rho)$. Then
\[
\sum_{j<m}\frac1f\sum_{\rho<f}\log\frac{w_{\rho,j}v_E}{\pi}+m(\log g_1+\log g_2)
-C\Big(L+\frac{m^2L}k\Big)\le\frac1{fk}\log\frac{D_*}{d_*}.
\]
\item[(R2)] (\emph{Place data.}) For a standard record and $c_D,c_h\ge0$
there is $x_0$ such that at every scale of exponent $\varepsilon$ with
$x\ge x_0$ and all choice data satisfying the choice conditions, the
column type of every word at every archimedean place is place data for
$(x,\Delta,a,B,k,m)$, and its window weights satisfy $w_j\ge1$.
\item[(R3)] (\emph{Block lengths.}) For a standard record the nine
block-length statements of Lemma~\ref{i4:lem-E2parts} hold eventually.
\item[(R4)] (\emph{Window weights.}) The class of all families is closed
for: ``eventually, a weight vector for $(D(x,i),L(x,i))$ exists in the
windows of the family''.
\item[(R5)] (\emph{Class decomposition.}) For a totally definite hermitian
space $W$ of a column datum, with standard lattice $L$, adelic unitary group
$G$, $\Gamma=U(W)(F)$ and $K$ the product of the local stabilisers of $L$
and the compact real groups, there are $g_1,\dots,g_n$ with
$G=\bigsqcup_i\Gamma g_iK$, each $\Gamma_i=\Gamma\cap g_iKg_i^{-1}$ finite,
and $\sum_i|\Gamma_i|^{-1}=\mathrm{mass}(L)$.
\end{enumerate}
\end{proposition}
\noindent The proofs are in \S\ref{i4:ssec-proof-Rleaves}. (R1) is a
Vandermonde computation; (R2) and (R3) are comparisons of explicit powers
of $x$; (R5) is the double-coset description of the genus, with Borel's
finiteness theorem. (R4) is proved there modulo one smaller statement,
Proposition~\ref{i4:Rprop-window}.


\begin{theorem}[Landherr, Abh.\ Math.\ Sem.\ Hamburg 11 (1936) 245--248]\label{i4:thm-landherr}
Two nondegenerate hermitian spaces over a quadratic extension $E/F$ of
number fields are isometric if and only if they have the same dimension,
the same discriminant in $F^\times/N_{E/F}E^\times$, and the same signature
at every real place of $F$ that does not split in $E$.
\end{theorem}

\begin{theorem}[mass of a standard lattice]\label{i4:thm-genusmass}
Let $E/F$ be a CM extension that is unramified at all finite places, and
let $L$ be a unimodular lattice in a totally definite hermitian space $W$
of rank $n$. Then $\mathrm{mass}(L)=\mathrm{mass}(\mathrm{std}_n)$.
\end{theorem}
\begin{proof}
Let $d$ be a Gram determinant of $W$.
\begin{itemize}
\item At the real places $d>0$.
\item At a finite place $v$, $d$ differs by a norm from the unit Gram
determinant of $L_v$. A unit is a norm in the unramified extension
$E_v/F_v$.
\end{itemize}
By the Hasse norm theorem (Neukirch VI (4.5)) $d$ is a
global norm. By Theorem~\ref{i4:thm-landherr}, $W\cong E^n$.

Locally, every unimodular lattice over an unramified $E_v$ has an
orthonormal basis. To see this, note first that some vector $x$ has
$h(x,x)$ a unit: otherwise the trace would map into $\mathfrak p$, contrary
to unimodularity. Rescale so that $h(x,x)=1$, split off $x$, and induct.
Hence the image of $L$ and $\mathrm{std}_n$ lie in one genus. The mass
depends only on the genus up to isometry.
\end{proof}

\begin{theorem}[Green, Trans.\ AMS 80 (1955) 402--447; reference to be confirmed]\label{i4:thm-green}
The cuspidal irreducible characters of $\mathrm{GL}_d(\mathbb F_g)$ are the
characters attached by Green's construction to the regular characters of
$\mathbb F_{g^d}^\times$.
\end{theorem}
\noindent It is applied to the residue fields $\mathbb F_g$ at the level primes $g$.

The automorphic construction rests on the following published theorems,
in addition to Theorems~\ref{i4:thm-hecke}, \ref{i4:thm-landherr} and
\ref{i4:thm-green}. We state the form in which they are used, and mark the references to be confirmed.

\begin{theorem}[Bernstein--Zelevinsky, Ann.\ Sci.\ ENS (4) 10 (1977), Thm.~2.9]\label{i4:thm-bz}
An irreducible admissible representation of $\mathrm{GL}_n$ over a
$p$-adic field has a cuspidal support that is unique up to permutation.
\end{theorem}

\begin{theorem}[Bernstein--Zelevinsky finiteness; reference to be confirmed]\label{i4:thm-bzfin}
Let $k$ be a nonarchimedean local field. Every irreducible smooth
representation of $\mathrm{GL}_n(k)$ is admissible; Jacquet modules of
representations of finite length have finite length; parabolic induction
takes representations of finite length to representations of finite
length; and parabolic induction commutes with taking contragredients.
\end{theorem}
\noindent Theorem~\ref{i4:thm-bzfin} is the finiteness that is used for
the cuspidal supports of a given level and for the finiteness of types;
Theorem~\ref{i4:thm-bz} is the uniqueness statement.

\begin{theorem}[Serre, \emph{Linear Representations of Finite Groups}, 2.3 Thm.~3 and 2.6 Thm.~8]\label{i4:thm-serre}
The irreducible characters of a finite group are orthonormal, and every
representation is the direct sum of its isotypic components, cut out by the
central idempotents.
\end{theorem}

In the next two theorems $E/F$ is a CM extension, $W$ and $V$ are
hermitian spaces over $E$ of dimensions $n\ge1$ and $m$, $r$ is the Witt
index of $V$, $[U(V)]=U(V)(F)\backslash U(V)(\mathbb A_F)$, and $\omega$ is
a Weil representation of $U(W\oplus(-W))\times U(V)$ attached to an
additive character $\psi$ and a splitting character. For a theta function
$\theta(g,h;f)$ of a vector $f$ we write
\[
I(g,f)=\frac1{\mathrm{vol}([U(V)])}\int_{[U(V)]}\theta(g,h;f)\,dh
\]
for its \emph{mean} over $[U(V)]$, for any Haar measure on
$U(V)(\mathbb A_F)$. The identities are stated with the mean, so they do not
depend on the normalisation of the measure.

\begin{theorem}[Siegel--Weil formula in Weil's range; A.~Weil, Acta Math.\ 113 (1965) 1--87; for unitary groups A.~Ichino, Math.\ Z.\ 247 (2004) 241--277]\label{i4:thm-sw}
Assume $r=0$ or $n+r<m$. Then the theta integral converges absolutely:
$\int_{[U(V)]}\sum_\xi|f(\xi h)|\,dh$ is bounded by a continuous seminorm of
the Schwartz function $f$, and $h\mapsto\theta(g,h;f)$ is integrable on
$[U(V)]$ for every $g$. If moreover $m>2n$, then the Siegel--Weil section
$\Phi_f$ attached to $f$ is left invariant under the rational points of the
Siegel parabolic, its Siegel Eisenstein series converges absolutely at
$s_0=(m-n)/2$, and
\[
E(g,s_0,\Phi_f)=I(g,f)\qquad\text{for all }g .
\]
\end{theorem}

\begin{theorem}[Siegel--Weil formula in the second range; A.~Ichino, Math.\ Z.\ 255 (2007) 721--729; holomorphy: T.~Ikeda, Compos.\ Math.\ 103 (1996) 183--218]\label{i4:thm-sw2}
Assume $n<m\le2n$ and $n+r<m$. Let $f$ be a Fock vector (archimedean
component the Gaussian attached to a majorant and to Cartan involutions
compatible with the maximal compact subgroup of an Iwasawa decomposition).
Then $h\mapsto\theta(g,h;f)$ is integrable on $[U(V)]$ for every $g$, the
Siegel Eisenstein series $E(g,s,\Phi_f)$ attached to $f$ is holomorphic at
$s_0=(m-n)/2$, and its value there is $E(g,s_0,\Phi_f)=I(g,f)$.
\end{theorem}
\noindent Only these two ranges are used (Lemma~\ref{i4:A-lem-swrange});
no regularised formula is needed, because in both the theta integral
converges absolutely.

\begin{theorem}[local theta correspondence and Howe duality; Waldspurger, Israel Math.\ Conf.\ Proc.\ 2 (1990) 267--324 ($p\ne2$); Gan--Takeda, J.\ Amer.\ Math.\ Soc.\ 29 (2016) 473--493; archimedean: Howe, J.\ Amer.\ Math.\ Soc.\ 2 (1989) 535--552; reference to be confirmed]\label{i4:thm-howe}
For a unitary dual pair over a local field, every irreducible smooth
representation $\pi$ of one member that occurs as a quotient of the Weil
representation has a unique irreducible quotient $\theta(\pi)$ of its big
theta lift, and $\pi\mapsto\theta(\pi)$ is injective. At an unramified
place the lift of an unramified representation is unramified, with Satake
parameter given by the standard recipe.
\end{theorem}
\noindent The unramified case is the one we use; it is also
proved in the formal development.

The Deligne--Lusztig parametrisation is not used here. The cuspidal
representations of $\mathrm{GL}_d(\mathbb F_g)$ enter only through Green's
parametrisation (Theorem~\ref{i4:thm-green}).

\begin{theorem}[bundles and rows]\label{i4:thm-assembly}
Assume the published inputs listed in Remark~\ref{i4:rem-trust}(ii), and
Proposition~\ref{i4:prop-assembly-leaves}. Let $c_{\rm Dep}$ and the
evaluation property $A$ be those of
Proposition~\ref{i4:prop-assembly-leaves}(A4), and let $p$ be the parameter
system with origin $c_{\rm Dep}$ extended by $A$. Then the hypotheses of
Theorem~\ref{i4:thm-classSel} hold for suitable $KP,Q,Q',\mathrm{Lk}$, with
$(\mathrm{LN})$ as the norm condition (field (viii)). More precisely:
\begin{enumerate}
\item[(a)] (Large unramified exponents.) The following hold for $p$: the
Siegel mass ratio for CM fields, the four retwisting conditions, and
$\mathrm{ExBL}$.
\item[(b)] (All rows.) The 26 row conditions hold as $\mathrm{RowL}$
statements.
\item[(c)] (Open origin.) The open parameter statements hold at the
origin. These are primitivity, the $\Lambda'$-rows, the divisor rows,
decomposition points, the selection jump, the tie, the labels, goodness,
geometry, multiplicity, the constant $\kappa$ and definite forms. They are
monotone along stronger extensions, so they hold at $p$.
\item[(d)] Wu decay, finite dimensionality of rows, the positivity
statements and the volume ratio hold. They follow from the Siegel--Weil
formulas (Theorems~\ref{i4:thm-sw} and~\ref{i4:thm-sw2}).
\end{enumerate}
\end{theorem}
\noindent The proof is in \S\ref{i4:ssec-proof-assembly}. The parts are
proved in the following order, which is the logical one: (A2) gives (c);
(c) and (A3) give the norm row; (c), the norm row and (A4) give (b) through
Theorem~\ref{i4:A-thm-allrows}; (a) and (d) are (A1), which does not use
(b) or (c).

\begin{proposition}[leaves of the automorphic construction]\label{i4:prop-assembly-leaves}
Let $p_0$ be the parameter system with origin $c_{\rm Dep}$, and $p$ its
extension by the evaluation property $A$ of (A4).
\begin{enumerate}
\item[(A1)] (\emph{Members.}) Statements (a) and (d) of
Theorem~\ref{i4:thm-assembly} hold.
\item[(A2)] (\emph{Bundle-level inputs of the open rows.}) For every
bundle over a witness point of $p_0$ carrying a bundle witness: it has a
comparison datum; its realisation satisfies the second multiplicity
condition $\mathrm{MI}^*_2$; given $\mathrm{MI}^*_2$, the pairings on the
definite base and the division-row pairings vanish; at points with a
parameter choice the $\Lambda$-rows and the $\kappa$-rows hold; the local
Langlands correspondence gives the goodness and label rows; the ray class
field statement gives weak geometric compatibility; the defect row on the
point-evaluation side holds; the row conditions at the companion hold; and
the join of the linked realisation gives the selection jump
(Proposition~\ref{i4:A-prop-open}, (O1)--(O9)).
\item[(A3)] (\emph{Norm condition.}) Under the hypotheses of
Propositions~\ref{i4:A-prop-at23} and~\ref{i4:A-prop-galois}, the norm
condition for $(\mathrm{LN})$ is linked at $p_0$.
\item[(A4)] (\emph{Functional condition and covolume bound.}) There are
$c_{\rm Dep}$ and an evaluation property $A$ with $\mathrm{EvLgG}(A)$ such
that the functional condition and the covolume bound are linked at $p$
(Proposition~\ref{i4:A-prop-funcov}).
\end{enumerate}
\end{proposition}
\noindent The proof is in \S\ref{i4:ssec-proof-Aleaves}. It is complete
modulo the published theorems listed in Remark~\ref{i4:rem-trust}(ii) and
the following exactly stated, formally verified propositions:
Proposition~\ref{i4:A-prop-inputs-leaves} (the leaves of the member
inputs) and Proposition~\ref{i4:A-prop-inputs}(b)--(d),
Proposition~\ref{i4:A-prop-stages} (the stages of the member statement),
Proposition~\ref{i4:A-prop-rowleaves} (the pointwise leaves of the rows),
Proposition~\ref{i4:A-prop-open} (the bundle-level leaves of (A2)),
Proposition~\ref{i4:A-prop-split} (the split-case conclusion behind the row bounds at $2$ and $3$),
Proposition~\ref{i4:A-prop-galois} (the Galois upgrade of the norm row),
Proposition~\ref{i4:A-prop-sern} (the series input) and
Proposition~\ref{i4:A-prop-funcov} (the functional and covolume rows).
What is proved in full is the construction of the lifts as far as it is
used, the row calculus, the reduction of each of the 26 rows to a
pointwise or bundle-level statement, the rows that are formal (twist,
the five open rows, the norm row), and the assembly of (A1)--(A4).

\begin{corollary}\label{i4:cor-structure}
Assume the published inputs listed in Remark~\ref{i4:rem-trust}(ii), and
Theorems~\ref{i4:thm-classzeros} and~\ref{i4:thm-mink}, and
the formally verified Propositions~\ref{i4:Rprop-kappaI}, \ref{i4:Rprop-num},
\ref{i4:Sprop-field}, \ref{i4:Sprop-ranks}, \ref{i4:Sprop-host},
\ref{i4:Sst-prop-pair}, \ref{i4:A-prop-inputs-leaves},
\ref{i4:A-prop-inputs}(b)--(d), \ref{i4:A-prop-stages}, \ref{i4:A-prop-rowleaves},
\ref{i4:A-prop-open}, \ref{i4:A-prop-split}, \ref{i4:A-prop-galois},
\ref{i4:A-prop-sern} and~\ref{i4:A-prop-funcov}, together with the
non-endoscopic third-layer hypotheses and the compatibility of geometric
Galois representations. Then a hypothesis structure exists.
\end{corollary}
\begin{proof}
Proposition~\ref{i4:prop-assembly-leaves} holds by
\S\ref{i4:ssec-proof-Aleaves}, so Theorem~\ref{i4:thm-assembly} supplies
the hypotheses of Theorem~\ref{i4:thm-classSel}: the Siegel mass ratio is
(a), the volume ratio, Wu decay, finite dimensionality and the positivity
statements are (d), and the rows and $\mathrm{ExBL}$ are (a), (b). In
field (iv), the thin-scale input is Proposition~\ref{i4:thm-litthin} and
the genus mass is Theorem~\ref{i4:thm-genusmass} with $n=2$. The
realisation step used in field (i) (Theorem~\ref{i4:thm-realise}) needs
Theorem~\ref{i4:thm-mink}, Theorem~\ref{i4:thm-linnik},
Proposition~\ref{i4:thm-shortnorm} and Proposition~\ref{i4:Rprop-window}.
\end{proof}

\subsection{The top-level theorem and the proof of Theorem~\ref{M:thm:alpoge-families}}\label{i4:ssec-top}

\begin{theorem}[uniform existence of connected families]\label{i4:thm-top}
Assume the loss bound, the divisor bound, and the level counts and
Theorem~S, and the statements listed in Remark~\ref{i4:rem-trust}. Let
$\lambda>1$, and let $\mathcal G$ be a requirement with
$\mathrm{HostExW}(\lambda,\mathcal G)$ all of whose class members satisfy
$\mathrm{HypV}$. Then $\mathcal G.\mathrm{cls}(\lambda)$ is connected, with
$e_D=2$ and $e_L=0$.
\end{theorem}
\begin{proof}
Let $(\mathcal A,\sigma,c_T)$ be the hypothesis structure of
Corollary~\ref{i4:cor-structure}. Apply Theorem~\ref{i4:thm-desWin} with
\[
\mathcal C=\mathrm{WithChN}(\mathcal A),\quad
\mathcal C'=\mathrm{WithChN}(\mathcal A\wedge\mathcal G.\mathrm{cls}(\lambda)),\quad
\mathcal E=\mathcal G.\mathrm{cls}(\lambda),
\]
and $C_{\rm abs}=2$, $C_0=3$, $e_D=2$, $e_L=0$. We check its four
hypotheses.
\begin{itemize}
\item Hypothesis (2) is immediate.
\item Hypothesis (3) is the assumption on $\mathcal G$.
\item Hypothesis (1) is Theorem~\ref{i4:thm-design} applied to
$\mathcal A\wedge\mathcal G.\mathrm{cls}(\lambda)$. Its first hypothesis is
field (i) of the structure. The duties and the two closedness hypotheses
restrict to the subclass by Lemma~\ref{i4:lem-closed}(d).
\item In hypothesis (4), closedness for $(\mathrm{LN})$ is field (viii).
Five of the remaining properties are
Proposition~\ref{i4:prop-links-proved}, and the last three are
Proposition~\ref{i4:prop-links-formal}. Green's hypothesis is
Theorem~\ref{i4:thm-green}.\qedhere
\end{itemize}
\end{proof}

\medskip\noindent\label{i4:ssec-final}We now deduce Theorem~\ref{M:thm:alpoge-families}.

\begin{theorem}\label{i4:thm-final}
Assume the loss bound (Proposition~\ref{i4:hyp-loss}), the divisor bound
(Proposition~\ref{i4:hyp-div}), the level counts and Theorem~S
(Proposition~\ref{i4:hyp-S}), and the windowed Selmer bound
(Proposition~\ref{i4:hyp-selmer}), and the statements listed in
Remark~\ref{i4:rem-trust}. Then Theorem~\ref{M:thm:alpoge-families}
holds.
\end{theorem}
\begin{proof}
The proof has five steps.

\emph{Step 1: hosts.} By Proposition~\ref{i4:hyp-selmer},
$\mathrm{VWinAt}(1,401)$ holds. By the host supply
(Proposition~\ref{i4:prop-supply}), $\mathrm{HostExW}(3/2,\mathcal G_1)$
holds.

\emph{Step 2: connectedness.} By Lemma~\ref{i4:lem-hypV-cls}, every member
of $\mathcal G_1.\mathrm{cls}(3/2)$ satisfies $\mathrm{HypV}$.
Theorem~\ref{i4:thm-top} with $\lambda=3/2$ and $\mathcal G=\mathcal G_1$
shows that this class is connected. So there are $e_D,e_L$ and
$\varepsilon_1'>0$ (one may take $e_D=2$, $e_L=0$) such that for every
$\varepsilon\in(0,\varepsilon_1')$ there are a family $\mathcal F$ of
exponent $\varepsilon$ in the class and constants $K,c_{HH},e_K,C_L,x_0$
with $\mathcal F$ chain-tied at every $x\ge x_0$ of
$\mathcal X=\mathrm{Win}(\mathcal F)$.

\emph{Step 3: the range of $\varepsilon$.} Put
$\varepsilon_1=\min(\varepsilon_1',1/52)$. This loses nothing: a family of
exponent $\varepsilon$ exists only for $\varepsilon<1/52$, because the
scales require it. The quantities $e_D$, $e_L$ and $\varepsilon_1$ are
chosen before $\varepsilon$. The constants $K,c_{HH},e_K,C_L,x_0$ are
chosen with the family, before $x$.

\emph{Step 4: clause (a) and the signs.} The family has windows of ratio
$3/2$, which is clause (a). The size windows of the package contain
$e_D^{\rm pkg},e_L^{\rm pkg}\ge0$. Since $\mathcal X$ is unbounded, the
family is chain-tied at some $x$, so one may replace $e_D,e_L$ by
$\max(e_D,0)$ and $\max(e_L,0)$. With the values $2$ and $0$ this is
automatic.

\emph{Step 5: clause (b).} Let $x\in\mathcal X$ with $x\ge x_0$.
Definition~\ref{i4:def-chain} gives a base member $i_0$, its host $H$ with
group $(\Z/2)^2\times A$, the map $\mathrm{pt}$ on all of $A$, and data
satisfying (1)--(5) for $H$. Here $A$ is the full translation group by
convention (2), so the package has one lift for each element of $A$. In
property (2) the characters are the three nontrivial characters trivial on
$A$ (convention (3)). The bounds of property (2),
$|G_\chi(x)|\le x^{1/2+\varepsilon}$ and $|L'(1,\chi)/L(1,\chi)|\le C_L\log x$,
are part of the chain-tied data at $i_0$. The second does not come from
the bound with constant $4$ for $L'/L$ in $\mathcal G_1$.

The sizes (1), weights (3), thin scales and cost (4), and the package (5)
are clauses of the chain-tied statement. In each case the package
constants are at most $e_D$, $e_L$ and $c_{HH}$. Clause (b) is thus
obtained by definition: ``chain-tied'' (Definition~\ref{i4:def-chain}) is
defined as satisfying (1)--(5) of Theorem~\ref{M:thm:alpoge-families}(b), so
the whole content of (b) lies in the core link,
Proposition~\ref{i4:Lprop-fam}(c), which is formally verified. None of
$K,c_{HH},e_K,C_L,x_0$ depends on $x$ or on the window. This is
Theorem~\ref{M:thm:alpoge-families}.
\end{proof}

\begin{remark}[statements not proved in this section]\label{i4:rem-trust}
Beyond the four inputs used as hypotheses, the proof above rests on the
following.
\begin{enumerate}
\item[(i)] \emph{Formally verified statements} \cite{CharacterFamily}:
\begin{itemize}
\item hosts: the field of the character, the valid ranks and the host
(Propositions~\ref{i4:Sprop-field}, \ref{i4:Sprop-ranks}
and~\ref{i4:Sprop-host}) and the existence of good window pairs
(Proposition~\ref{i4:Sst-prop-pair}), from which, with the twist bounds
and the window data (Propositions~\ref{i4:Sprop-tw} and~\ref{i4:Sprop-win},
proved in \S\S\ref{i4:Sst-tw}, \ref{i4:Sst-win}), the host supply
(Proposition~\ref{i4:prop-supply}) is proved;
\item division and global leaves: Propositions~\ref{i4:Lprop-gapleaves},
\ref{i4:Lprop-sel}(a), \ref{i4:Lprop-D} (except the first claim of (e)),
\ref{i4:Lprop-pkg}, \ref{i4:Lprop-forms} and \ref{i4:Lprop-fam};
\item realisation: Propositions~\ref{i4:Rprop-kappaI} and~\ref{i4:Rprop-num},
from which the leaves of the window count
(Proposition~\ref{i4:Rprop-window-leaves}) and
Proposition~\ref{i4:Rprop-window} are proved;
\item the automorphic construction: Propositions~\ref{i4:A-prop-inputs-leaves},
\ref{i4:A-prop-inputs}(b)--(d), \ref{i4:A-prop-stages},
\ref{i4:A-prop-rowleaves}, \ref{i4:A-prop-open},
\ref{i4:A-prop-split} (from which the row bounds at $2$ and $3$,
Proposition~\ref{i4:A-prop-at23rows}, are proved), \ref{i4:A-prop-galois},
\ref{i4:A-prop-sern} and~\ref{i4:A-prop-funcov}.
\end{itemize}
The analytic number theory is proved in full modulo published theorems:
norm primes in short intervals (Proposition~\ref{i4:thm-shortnorm}) from
the zero statements of Proposition~\ref{i4:prop-lfam-zeros} and the
ray class family of Proposition~\ref{i4:prop-lfam-ray} with its leaves
(Proposition~\ref{i4:prop-ray-leaves}). The pointwise row leaves and the
bundle-level leaves of the automorphic construction are by far the
largest. For every formally verified proposition the structure of its
formal proof (which steps are published theorems, which are routine, which
are new) is described in \S\ref{i4:ssec-outline}; those
accounts are not proofs.
\item[(ii)] \emph{Published theorems.} Stated in this section: Hecke
(Theorem~\ref{i4:thm-hecke}); the least prime in a progression
(Theorem~\ref{i4:thm-linnik}); Landherr
(Theorem~\ref{i4:thm-landherr}) and the Hasse norm theorem (Neukirch VI
(4.5)); Green (Theorems~\ref{i4:thm-green} and~\ref{i4:A-thm-green}); Bernstein--Zelevinsky
(Theorems~\ref{i4:thm-bz} and~\ref{i4:thm-bzfin}); Serre
(Theorem~\ref{i4:thm-serre}); the two Siegel--Weil formulas
(Theorems~\ref{i4:thm-sw} and~\ref{i4:thm-sw2});
Howe duality (Theorem~\ref{i4:thm-howe});
Minkowski's bound (Theorem~\ref{i4:thm-mink}); the
classical zero estimates and truncated explicit formula
(Theorem~\ref{i4:thm-classzeros}; reference to be confirmed for the second range of (Z5)), from which
Proposition~\ref{i4:thm-litthin} is derived; Landau's lemma on the
logarithmic derivative and the truncated Perron formula
(Lemmas~\ref{i4:lem-landau} and~\ref{i4:lem-perron}), from which Proposition~\ref{i4:prop-lfam-zeros} is proved; the
arithmetic of imaginary quadratic fields, Poisson summation for the
Gaussian and Minkowski's bound for ideal classes
(Theorems~\ref{i4:Rst-thm-quad}, \ref{i4:Rst-thm-poisson}
and~\ref{i4:Rst-thm-minkideal}; Neukirch), from which
Proposition~\ref{i4:prop-ray-leaves} is proved. The further published
theorems stated in \S\ref{i4:ssec-outline} serve only to
classify the steps of formally verified proofs. Used with a
statement in the text but without a theorem environment: the classical
facts on cyclotomic fields and differents in Lemma~\ref{i4:lem-rd-sigma},
and the non-vanishing of Hecke $L$-functions on $\Rel s\ge1$
(Lemma~\ref{i4:A-lem-chainL}); reference to be confirmed.

Not stated in this paper (the formal development isolates them as separate
propositions and proves each inside the development; reference to be confirmed): the published inputs of the third place
listed in Remark~\ref{i4:A-rem-pub}; and moreover Clozel's cohomological
criterion; $\ell$-adic characters and parameters attached to Hecke
characters; Chebotarev in degree one; Galois representations attached to
polarised cuspidal representations, with local-global compatibility above
$\ell$, purity, and the de Rham to potentially semistable theorem; base
change for definite unitary groups; the crystalline and integral lattice
theorems of Bhatt--Scholze; the doubling method and its local zeta
integrals; the formal degree of compact duals; Bernstein's irreducibility
of unitary products and the irreducibility of unlinked products; the
geometric-lemma filtration; injectivity of the depth-zero Jacquet functor;
Jacquet--Shalika; finiteness of automorphic forms on $\mathrm{GL}_n$;
Zelevinsky's polynomial ring; Artin characters of ray classes; the generic
unitary dual and standard cosocles of $\mathrm{GL}_n$; Green's regular
characters; existence and covers of discrete series and Harish-Chandra
modules; the Weil representation of unitary pairs; strong approximation for
$\mathrm{SU}$; Harish-Chandra analyticity; the trace class property of
cusp forms; Schur orthogonality for square-integrable representations;
the non-endoscopic third-layer hypotheses; compatibility of geometric
Galois representations; the hypotheses $\mathrm{Hypotheses}_2$ of
Proposition~\ref{i4:A-prop-at23} (Fontaine's $B_{\rm cris}$, global Artin
reciprocity, local Bloch--Kato, local Tate duality, the local Euler
characteristic, local finiteness, the Bhatt--Scholze integral comparison);
depth-zero support; the local theta inputs, the agreement of theta
integrals in the Weil tensor models and the local matching of
Proposition~\ref{i4:A-prop-funcov}; the covolume bound for $X$; and the
existence of a local Langlands datum.
\end{enumerate}
The remaining arguments of this section are proved here: the logical
layer, which fixes the constants in the stated order; the size, proviso,
lowest-space, rank-equality and genus-mass arguments; the
root-discriminant estimate; the kit parameters; the reductions of the
division and global links to their leaves; realisation from its size
leaves; and the reduction of the 26 rows and of the open-origin rows to
their leaves.
\end{remark}

\subsection{Root discriminants, kits and the host supply}\label{i4:ssec-grp-hosts}

\subsubsection{The root-discriminant estimate}\label{i4:ssec-proof-rd}

Let $N\ge1$. Write $\phi$ for Euler's function and $\Lambda$ for the von
Mangoldt function. For $c\mid N$ let
$W_c=\{u\in(\Z/N)^\times:u\equiv1\pmod c\}$, the kernel of reduction to
$(\Z/c)^\times$. For a subgroup $U\subset(\Z/N)^\times$ put
\[
\varrho_U(c)=\frac{|W_c\setminus U|}{|W_c|},\qquad
\Sigma(U)=\sum_{m\mid N}\Lambda(m)\,\varrho_U(N/m),
\]
and for $u\in\Z/N$ let $o(u)$ be the additive order of $u-1$.

\begin{lemma}\label{i4:lem-rd-W}
Let $c\mid N$.
\begin{enumerate}
\item[(a)] $|W_c|\phi(c)=\phi(N)$.
\item[(b)] For $m\mid N$, $o(u)\mid m$ if and only if $u\equiv1\pmod{N/m}$.
\item[(c)] If $a,b\mid N$ and $W_a\cup W_b\subset U$, then
$W_{\gcd(a,b)}\subset U$.
\item[(d)] $0\le\varrho_U(c)\le1-\phi(c)/\phi(N)$. If $W_c\not\subset U$
then $\varrho_U(c)\ge1/2$. If $uv\in U$ for all $u,v\in W_c\setminus U$,
then $\varrho_U(c)\le1/2$.
\end{enumerate}
\end{lemma}
\begin{proof}
(a) Reduction is onto with kernel $W_c$. (b) With $N=mk$, $m(u-1)=0$ in
$\Z/N$ means $k\mid u-1$. (c) Given $u\in W_{\gcd(a,b)}$, the Chinese
remainder theorem gives $v\in(\Z/N)^\times$ with $v\equiv1\pmod a$ and
$v\equiv u\pmod b$. Then $v\in W_a$ and $v^{-1}u\in W_b$, so $u\in U$.
(d) The first bound holds because $1\in W_c\cap U$. If $u_0\in
W_c\setminus U$, multiplication by $u_0$ maps $W_c\cap U$ injectively into
$W_c\setminus U$, and under the last hypothesis it maps $W_c\setminus U$
injectively into $W_c\cap U$.
\end{proof}

\begin{lemma}\label{i4:lem-rd-sigma}
\begin{enumerate}
\item[(a)] $\sum_{u\in(\Z/N)^\times\setminus U}\Lambda(o(u))/\phi(o(u))=\Sigma(U)$.
\item[(b)] Let $\zeta$ be a primitive $N$-th root of unity, $L=\Q(\zeta)$,
$E\subset L$ a subfield and $U=\{u:\tau_u|_E=\mathrm{id}\}$. Then
$\log|d_E|/[E:\Q]=\Sigma(U)$.
\end{enumerate}
\end{lemma}
\begin{proof}
(a) For $1\ne m\mid N$ one has
$\sum_{m\mid m'\mid N}\Lambda(m')\phi(N/m')=\phi(N)\Lambda(m)/\phi(m)$. Both
sides vanish unless $m=p^j$ is a prime power. If it is, write
$N=p^vM$ with $p\nmid M$. The left side is
$\log p\,\phi(M)\sum_{i\le v-j}\phi(p^{v-j-i})=\log p\,\phi(M)p^{v-j}$, which
is the right side. Apply this with $m=o(u)$ (note $o(u)\ne1$ for
$u\notin U$), exchange the sums, and use Lemma~\ref{i4:lem-rd-W}(b): the
$u\notin U$ with $o(u)\mid m$ are the elements of $W_{N/m}\setminus U$.

(b) We use the following classical facts.
\begin{itemize}
\item $\mathcal O_L=\Z[\zeta]$, and $u\mapsto\tau_u$ is an isomorphism
$(\Z/N)^\times\cong\mathrm{Gal}(L/\Q)$ (Washington, \emph{Introduction to
Cyclotomic Fields}; reference to be confirmed).
\item If $\mathcal O_{F'}=\mathcal O_F[\alpha]$, the different of $F'/F$ is
generated by $g'(\alpha)$, $g$ the minimal polynomial of $\alpha$. The
different is multiplicative in towers, and its absolute norm is the
discriminant (Neukirch, \emph{Algebraic Number Theory}, Ch.~III;
reference to be confirmed).
\item $\log|\Phi_m(1)|=\Lambda(m)$ for $m>1$. This follows from
$\prod_{1<j\mid m}\Phi_j(1)=m$ by induction.
\end{itemize}
For any subfield $F\subset L$ we have $\mathcal O_L=\mathcal O_F[\zeta]$,
and the minimal polynomial of $\zeta$ over $F$ is
$\prod_{\sigma\in\mathrm{Gal}(L/F)}(X-\sigma\zeta)$. The tower formula
gives
\[
\log|d_L|=\sum_{1\ne\sigma\in\mathrm{Gal}(L/F)}\log|\mathrm N_{L/\Q}(\zeta-\sigma\zeta)|+[L:F]\log|d_F| .
\]
Subtract the case $F=E$ from the case $F=\Q$. This gives
$[L:E]\log|d_E|=\sum_{u\notin U}\log|\mathrm N_{L/\Q}(\zeta-\zeta^u)|$.
Now $\zeta-\zeta^u=\zeta^u(\zeta^{1-u}-1)$, the root $\zeta^{1-u}$ has
order $o(u)$, and $\mathrm N(\zeta^u)$ has absolute value $1$. Hence
$\log|\mathrm N_{L/\Q}(\zeta-\zeta^u)|=\phi(N)\Lambda(o(u))/\phi(o(u))$.
Divide by $\phi(N)=[L:E][E:\Q]$ and apply (a).
\end{proof}

\begin{lemma}\label{i4:lem-rd-bounds}
\begin{enumerate}
\item[(a)] (Lower bound.) Let $e\ge0$, $2^{e+2}\mid N$, let $d$ be an odd
prime dividing $N$, suppose every $u\in U$ is $\equiv\pm1\pmod{2^{e+2}}$,
and suppose $\varrho_U(N/d)\ge1/2$. Then
$\Sigma(U)\ge(e+1-2^{-e})\log2+\frac12\log d$.
\item[(b)] (Upper bound.) Let $2\mid N$, and let $q\ne d$ be odd primes
dividing $N$ with $\varrho_U(N/q),\varrho_U(N/d)\le1/2$. Then
$\Sigma(U)\le\log N-\log2-\frac12\log q-\frac12\log d$.
\item[(c)] (Index bound.) $\log[(\Z/N)^\times:U]\le2\Sigma(U)$.
\end{enumerate}
\end{lemma}
\begin{proof}
(a) Write $N=2^vM$ with $M$ odd, and $W'=W_{2^{e+2}M}$. For $0\le i\le e+2$,
every $u\in W_{2^iM}\cap U$ is $\equiv1\pmod M$ and $\equiv\pm1\pmod{2^{e+2}}$.
For $i\ge2$ the sign must be $+$, so $W_{2^iM}\cap U\subset W'$. For
$i<2$ it lies in at most two cosets of $W'$. With
Lemma~\ref{i4:lem-rd-W}(a) this gives
$\varrho_U(2^iM)\ge1-(\phi(2^i)+[i<2])/2^{e+1}$. Keep only the terms
$m=2^{v-i}$, $0\le i\le e+1$, and $m=d$ of $\Sigma(U)$; all terms are
nonnegative. Since $\sum_{i\le e+1}\phi(2^i)=2^{e+1}$, the $2$-terms add up
to at least $(e+1-2^{-e})\log2$, and the $d$-term is at least $\frac12\log d$.

(b) The terms $m=q,d$ are at most $\frac12\log q$ and $\frac12\log d$. Every
other term is at most $\Lambda(m)(1-\phi(N/m)/\phi(N))$ by
Lemma~\ref{i4:lem-rd-W}(d). Now $\sum_{m\mid N}\Lambda(m)=\log N$. The
identity in the proof of Lemma~\ref{i4:lem-rd-sigma}(a), with $m=2$, gives
$\sum_{2\mid m\mid N}\Lambda(m)\phi(N/m)/\phi(N)=\log2$, and the even $m$
are among the other terms.

(c) Let $c$ be the least divisor of $N$ with $W_c\subset U$. By
Lemma~\ref{i4:lem-rd-W}(c), every $f\mid N$ with $W_f\subset U$ is a
multiple of $c$. The index is at most $\phi(N)/|W_c|=\phi(c)\le c$, and
$\log c=\sum\Lambda(m)$ over the $m\mid N$ with $c\nmid N/m$. For each such
$m$, $W_{N/m}\not\subset U$, so $\varrho_U(N/m)\ge1/2$ and
$\Lambda(m)\le2\Lambda(m)\varrho_U(N/m)$.
\end{proof}

\begin{proof}[Proof of Proposition~\ref{i4:prop-rd}]
Put $N=N_E$ and $U=\ker\mathrm{Art}$.

\emph{Surjectivity and the second sign.} Every $\sigma\in\mathrm{Gal}(E/\Q)$
extends through $\iota$ to some $\tau_u$, and (H1) with injectivity of
$\iota$ gives $\mathrm{Art}(u)=\ell(\sigma)$; so $\mathrm{Art}$ is onto.
The quadratic Gauss sum $g_d=\sum_a(\frac ad)\zeta^{aN/d}$ satisfies
$g_d^2=-d$ (as $d\equiv3\pmod4$) and $\tau_u g_d=(\frac ud)g_d$. Since
$\iota(\sqrt{-d})=\pm g_d$, (H1) gives
$s_2(\mathrm{Art}(u))=(\frac ud)$.

\emph{The four properties of $U$.}
\begin{itemize}
\item $\tau_u$ fixes $\iota(E)$ if and only if $u\in U$, by (H1). So
$\log\mathrm{rd}_E=\Sigma(U)$ by Lemma~\ref{i4:lem-rd-sigma}(b).
\item $W_{Q_Sq}\not\subset U$. Otherwise the kernel $U'$ of
$s_2\circ\mathrm{Art}$ would contain $W_{Q_Sq}$ and $W_d$ (the Jacobi
symbol is $1$ on $W_d$), hence all of $(\Z/N)^\times$ by
Lemma~\ref{i4:lem-rd-W}(c). This contradicts (H2).
\item $\varrho_U(N/d)\le1/2$ and $\varrho_U(N/q)\le1/2$. On
$W_{Q_Sq}\setminus U$, (H4) forces $\mathrm{Art}=((1,-1),1)$, an element of
order $2$, so products of two such elements lie in $U$. The group
$W_{Q_Sd}$ embeds into $(\Z/q)^\times$, so it is cyclic with a generator
$g$. By (H5) $\mathrm{Art}(g)^2=1$, so every element of
$W_{Q_Sd}\setminus U$ is an odd power of $g$ with image $\mathrm{Art}(g)$.
In both cases Lemma~\ref{i4:lem-rd-W}(d) applies.
\item Every $u\in U$ is $\equiv\pm1\pmod{2^{\nu+2}}$, by (H3).
\end{itemize}

\emph{Conclusion.}
\begin{itemize}
\item Lemma~\ref{i4:lem-rd-bounds}(a) with $e=\nu$ and the second
property (which gives $\varrho_U(N/d)\ge1/2$, as $N/d=Q_Sq$) yields the
lower bound.
\item Lemma~\ref{i4:lem-rd-bounds}(b) yields the upper bound, since
$\log N=\log Q_S+\log q+\log d$, and $q\ne d$ by their residues modulo $8$.
\item The index of $U$ is $|G|=4n=2f$, so
Lemma~\ref{i4:lem-rd-bounds}(c) gives $2f\le\mathrm{rd}_E^2$.\qedhere
\end{itemize}
\end{proof}

\subsubsection{The kit parameters}\label{i4:ssec-proof-kits}

\emph{Notation.} At a scale of exponent $\varepsilon$ write
\[
L=x^\varepsilon,\quad b=e^L,\quad a=xb,\quad X'=e^{x^{\varepsilon/3}},\quad
b'=a/X',\quad w''_n=b'/(xn),\quad X''_n=xnX',\quad\delta_1=x^{-\varepsilon/8},
\]
and $\theta=\varepsilon/100$. The ranks are $n$, $f=2n$, even $k,m$ with
$m\ge2$ and $m+3\le k$, odd $t$, $r=k+m-1$ and $N_m=k+m(t+1)$. Validity of
the ranks means
\[
x^{2+8\varepsilon}\le n\le2x^{2+8\varepsilon},\quad
\sqrt{x\delta_1}\le m\le2\sqrt{x\delta_1},\quad
x^{9/2+10\varepsilon}\le k\le2x^{9/2+10\varepsilon},
\]
and $r+t=n_P+T_On_W$ with $T_O=\lceil x^{377}\rceil$. The other
quantities ($\Delta=\lceil x^{c_D}\rceil$, $h_0=2\lceil x^{c_h}\rceil$,
$\eta_{sz}=mh_0$, $S_c=4k\Delta$, the host primes $q,d,p_{de},p'_{de},g_0,g_1$
and the conductor $Q_E=Q_Sqd$) are those of the size data; their validity
conditions are recalled where used. ``Eventually'' means: for every
$\varepsilon\in(0,1/52)$ and standard $C$ there is $c_1$ such that for
$c_D\ge c_1$ there is $c_2$ such that for $c_h\ge c_2$ there is $x_0$ with
the property at all valid size data with $x\ge x_0$. Eventual statements
are closed under conjunction (take maxima of the thresholds). In every
statement below the thresholds are explicit; we
name them when they are not $0$.

\begin{lemma}[shared estimates]\label{i4:lem-kits-shared}
\begin{enumerate}
\item[(a)] (Comparison.) For $0\le\alpha<\beta$, $\kappa>0$, reals $c,C$ and
$p\in\N$: $c\log x+Cx^\alpha(\log x)^p+C\le\kappa x^\beta$ for large $x$.
\item[(b)] (Scale.) For real $c$ and $\kappa>0$, for large $x$:
$x^c\le X'^\kappa$; $x^cX'\le b^\kappa$; $b'x^c\le b$; $b^{1-\kappa}x^c\le b'$;
for $1\le n\le x^3$, $b^{1-\kappa}x^c\le w''_n\le b'$; and
$\exp(cx^{\varepsilon/5}(\log x)^2)x^c\le X'^\kappa$.
\item[(c)] (Ranks.) For $x\ge1000^{1000}$ and valid ranks: $f\le x^3$,
$m\le x$, $k\le r\le\min(x^5,2k)$, $x^{377}\le t\le x^{401}$,
$N_m\le x^{402}$, $1\le L\le x$.
\item[(d)] (Host.) For a valid host with constant $C_8\ge1$ and large $x$:
$\log q\le x^{\varepsilon/5}$, $\log d\le x^{\varepsilon/5}(\log x)^2$,
$\log Q_E\le2x^{\varepsilon/5}(\log x)^2$,
$\log p_{de},\log p'_{de}\le2C_8x^{\varepsilon/5}(\log x)^2$,
$1\le\mathrm{rd}_E\le Q_E$, and $\log\mathrm{rd}_E\le L$.
\item[(e)] (Top count.) There is $C_N$, depending only on $\varepsilon$,
such that for large $x$ the top count satisfies $1\le N_1\le C_N\log x$,
and for $\Lambda\ge1$ the quantity $\Lambda_G=\Lambda N_1^2\log a$
satisfies $x^\varepsilon\le\Lambda_G\le2\Lambda C_N^2(\log x)^2x^\varepsilon$.
\end{enumerate}
\end{lemma}
\begin{proof}
(a) Each summand is $o(x^\beta)$. (b) Take logarithms: $\log X'=x^{\varepsilon/3}$,
$\log b=x^\varepsilon$, $\log b'=\log x+x^\varepsilon-x^{\varepsilon/3}$, and
$xn\le x^4$. Each clause becomes an instance of (a). (c) Since
$x^{1/1000}\ge1000$, a constant at most $1000$ costs at most $x^{1/1000}$.
Then $f\le4x^{2+8\varepsilon}\le x^{2.156}$, $r\le2k\le x^{4.694}$ and
$m\le x^{0.501}$, using $\varepsilon<1/52$. The block-length definitions give
$n_W\le x^{22.9}$, and $r+t=8fr+4+T_On_W$ gives
$T_O\le t\le13x^{400.9}$. Hence $N_m\le x^{402}$. (d) Host validity
places $q,d$ in $(\sqrt{X_i},X_i]$ with $\log X_1=Y^{\varepsilon/5}$,
$\log X_2=Y^{\varepsilon/5}(\log Y)^2$ for some $Y\le x$, and gives
$\log Q_E\le2Y^{\varepsilon/5}(\log Y)^2$ and $p_{de},p'_{de}\le Q_E^{C_8}$.
Proposition~\ref{i4:prop-rd} gives $1\le\mathrm{rd}_E\le\frac{Q_S}2\sqrt{qd}\le Q_E$,
and $2x^{\varepsilon/5}(\log x)^2\le L$ by (a). (e) The top count is
$N_1=\lceil\max(6eG'\delta_1,(h+3)\log_2x)\rceil+1$ with
$G'=(M+5)(6r_G+1)\lceil1/\delta_1\rceil$ and $M=\lceil(h+3)10^{6r_G+1}\log x\rceil$;
both entries of the maximum are $O(\log x)$. Finally
$\log a=\log x+x^\varepsilon\in[x^\varepsilon,2x^\varepsilon]$.
\end{proof}

\noindent\emph{Further notation.} The proof below also uses the following
quantities of the size data (cap-exist-kits, \S\S1--2).
\begin{itemize}
\item $h=2+8\varepsilon$ (the exponent of $n$; it is \emph{not} the node
gap $h_0$) and $r_G=\lceil8(h+2)/\varepsilon\rceil+1$; these enter only the
top count $N_1$ of Lemma~\ref{i4:lem-kits-shared}(e), whose constant is
$C_N=12e\big((h+3)10^{6r_G+1}+6\big)(6r_G+1)+2(h+3)+2$.
\item $r_S=m/2-1$ (an integer, $m$ being even), so $m=2r_S+2$.
\item $\delta_2=|T_2|/n$, the density of the type set $T_2$ of the host, and
$N_*=2n\delta_1$.
\item For real $C$: $A(C)=C\,r\big(f+f^2(k-1)^3L^2/m\big)$;
$\rho_P(C)=4N_*x^\varepsilon\Lambda_G/k+C\log x\,\Lambda_G/k$;
$\mathrm{gap}(C)=Ckm^4+N_mm^{10}$.
\item $l_A=kx^\varepsilon/(m\delta_2)$; $D_l=nt^2(N_m+t)L^2+kn^2(S_c+t+L)$
and $l_*=mb/D_l$ (the symbol $D_l$ is local to the kits and is unrelated to the proviso quantity $D_*$); $h_0=\eta_{sz}/m$.
\item $c_{LL}$, $C_\rho$, $C_{gap}$ (and $C_{kLow},C_{k6},C_{fL},C_{\Delta G},
C_w,C_\Delta,C_{Count},C_{h0r},C_{h0k},C_{h0N},C_Y$) are entries of the
standard record $C$; for a standard record each of them and $c_{LL}^{-1}$
is $\ge1$. The remaining constants are $C_{nW}=10$, $C_{Lw}=11\cdot405$,
$C_{j0}=405$, $c_w=9/10$, $C_{YL}=9$, $C_{\eta F}=6$, $C_{pde}=C_8\ge1$,
$C_F\ge1$; $\eta_F=6x^{-C_F}$.
\end{itemize}
In particular $A(10)\ge10rf$ (item 4), and in item 4 $c_{LL}\,r\ge c_{LL}x\ge1$
once $x\ge c_{LL}^{-1}$, so $\log\log b=\varepsilon\log x\le\log r\le
c_{LL}rf\log r$; and $(\log f)^2\le36(\log x)^2\le r_S$ because
$r_S\ge m/2-1\ge x^{1/4}/2-1$.

\begin{proof}[Proof of Proposition~\ref{i4:prop-kits}]
\emph{The first record (E1).} The parameter choice takes as its constants
those of $C$; they are positive for a standard record. Its window constant
is $C_{nW}=10$. Its four arithmetic fields are the four clauses of
arithmeticity of the host:
\begin{itemize}
\item each $g_i$ is a norm $\pi\bar\pi$;
\item $g_i\equiv1\pmod{Q_E}$;
\item $p_{de}$ and $p'_{de}$ split completely;
\item the level primes split in $K$.
\end{itemize}
Its estimate fields are twenty-one eventual inequalities. We list them
with their proofs; $x_r=1000^{1000}$.
\begin{enumerate}
\item $x^{2C_t+5}<q,d$; $x^{\varepsilon_z/6}\le\log q$;
$Q_E^{36}t^{150}\le b^{1/2}$; $\log p_{de},\log p'_{de}\le L$. Host
validity places $q,d$ in $(\sqrt{X_i},X_i]$ with $\log X_1=Y^{\varepsilon/5}$,
$\log X_2=Y^{\varepsilon/5}(\log Y)^2$, where $Y\ge x^{1/J}$ for the pairs used
(they are window pairs for $(\varepsilon,X^{1/J})$ with $x\le X$,
Proposition~\ref{i4:hyp-selmer}); hence
$\log q,\log d\ge x^{\varepsilon_z/5}/2$, and
$\varepsilon_z=\varepsilon/J\ge\varepsilon/C_{pde}$. (The constant
$C_{pde}$ is the constant $C_8$ of the host.) Then
$Q_E^{36}t^{150}\le\exp(222x^{\varepsilon/5}(\log x)^2)\le b^{1/2}$, using
$t<q$. Finally $2C_{pde}x^{\varepsilon/5}(\log x)^2\le L$.
\item $C_{kLow}m^5fLx^\varepsilon\le k\le x^{1/2}nR_0$ and $\sqrt{k/m}\le R_0$,
where $R_0=nx^{-\varepsilon}$. Put $w=x^{1/2}$ and $y=x^{\varepsilon/16}$.
Then $n\in[w^4y^{128},2w^4y^{128}]$, $m\le2w/y$ and
$k\in[w^9y^{160},2w^9y^{160}]$. Each inequality compares monomials in $w,y$,
and holds once $y\ge128|C_{kLow}|+2$.
\item $m^2L\le k$, $mx^{1/2}\log x\le k$, $C_{k6}m^6L\le k\log k$,
$C_{fL}L\le f\log k$. With $y=x^{1/2}\ge\max(8,64C_{k6},C_{fL})$:
$m\le2y$, $k\ge y^9$, $f\ge y^4$.
\item $\log f\le\sqrt{r_S}$, $(\log f)^2\le m^2L$,
$\log\log b\le c_{LL}rf\log r$, $(4f+10)r\le A(10)$. Here $\log f\le6\log x$
and $m\ge x^{1/4}$, $m\le2r_S+3$; also $\log\log b=\varepsilon\log x\le\log r$,
and $A(10)\ge10rf$.
\item[5--6.] The $\eta_{sz}$-bounds; these are items (1) and (2) of the group
$\eta_{sz}$ below.
\setcounter{enumi}{6}
\item $|TT|\le N_*$; $\rho_P(C_\rho)\le x^{-2}$;
$b_T(t)\le b'+N_*x^\varepsilon\Lambda_Gb'/k$. The first and third are
validity, with $N_*=2n\delta_1$. For the second, $n\delta_1\le2x^{2+8\varepsilon}$
and $\Lambda_G\le A_0(\log x)^2x^\varepsilon$ by
Lemma~\ref{i4:lem-kits-shared}(e). So both summands of $\rho_P$ are at most
$\frac12x^{5/2+10\varepsilon}/k$, and $x^{5/2+10\varepsilon}\le x^{-2}k$.
\item $4S_c\le\Lambda_Gb'/k$, $4C_{\Delta G}k^2\Delta\le\Lambda_Gb'$,
$\log\Delta\le\frac12x^{\varepsilon/4}$. Here $b'\ge x^{c_D+11}$ for large
$x$, while $16C_{\Delta G}k^2\Delta\le128C_{\Delta G}x^{c_D+10}$.
\item For every integer $l\ge11$ there is $j$ with $l_A\le(l-1)^j$ and
$2j+2\le L_w$; $L_w\le C_{Lw}\log x$; $j_0\le C_{j0}\log x$. Here $l_A\le x^7$,
so $j=\lceil7\log x\rceil$ works. The value $L_w=10\lceil405\log x\rceil+2\lceil30/\varepsilon\rceil+22$
and $nN_m/m\le4x^{404}$ give the rest, for $\log x\ge60/\varepsilon+34$.
\item The exactness clauses (shape of the Harish-Chandra parameters,
$\eta_F=6x^{-C_F}$, $p_{de}\equiv p'_{de}\equiv1\pmod{15}$, trivial
conductor ideal). These are read off from validity.
\item $2b/l_*+t/2+3\le|b_i(\tau)-e_j(\tau)|$ ($c_1=402$, $c_2=7c_D$). By (c),
$D_l\le12y^6$ with $y=x^{c_D}$, while $|b_i-e_j|\ge h_0/2\ge x^{7c_D}$.
\item $w_j(\tau)\le X''_nw_{j'}(\tau')$ and $b_{m-1}(\tau)\le2b_0(\tau)$.
These are item (9) of the third record.
\item $k^{C_w}\le w''_n$ and $b^{9/10}\le w''_n$. These follow from (b)
with $(c,\kappa)=(5|C_w|,1/10)$.
\item $k^{C_\Delta}(C_\Delta m)^3/\eta_F^{C_\Delta}$, $C_\Delta mN_m$,
$k^{C_\Delta}(C_\Delta m^4L)^4$, $N_mm^{10}$ and $(N_mfL)^{C_{Count}}$ are all
$\le\Delta$, for $c_1=5C_\Delta+C_FC_\Delta+406C_{Count}+412$. Each left
side is at most $x^{c_1}$, by (c).
\item[15--16.] The node and gap bounds by $h_0$. These are item (6) of the
third record, with $c_2=\max(1300+c_D,406|C_{Count}|,1210,3+402C_{h0N})$.
\setcounter{enumi}{16}
\item $\mathrm{gap}(C_{gap})$ is less than $v_i-v_{i+1}$, $v_{k-2}+b_0$ and
$b_{j+1}-b_j$ ($c_1=c_2=430$). We have $\mathrm{gap}\le x^{430}$, while each
right side is at least $\min(4\Delta,b_1,h_0)\ge x^{430}$.
\item[18--20.] $r(b'+S_c)\log x\cdot l_A\le\delta_2mb$, the two $l_*$-bounds,
and the word count $\le2^{k-3}$. These are items (5), (7) and (10) of the
third record.
\setcounter{enumi}{20}
\item There is $Y=(kb)^8$ with $N_m^4Q_E^{C_Y}\le Y^{1/4}$,
$\log Y\le9L$ and $Y\le g_i\le2Y$. Indeed $N_m^4\le x^{1608}\le b$,
$Q_E^{C_Y}\le b$, and $\log Y=8(\log k+L)\le9L$.
\end{enumerate}

\emph{The third record (E3): ten groups.} Throughout, (c) and (d) of
Lemma~\ref{i4:lem-kits-shared} are used without comment.
\begin{enumerate}
\item \emph{Sizes.} $n,m,k$ lie within a factor $x^{\pm\theta}$ of
$x^{2+8\varepsilon}$, $\sqrt{x\delta_1}$ and $x^{9/2+10\varepsilon}$. Also
$1\le t\le x^{C_t+\theta}$, $x^{-\theta}n\delta_1\le|TT|\le N_*$,
$\delta_2\ge x^{-\varepsilon}$ and $L_w\le(\log x)^2$. This holds for
$\log x\ge200C_{pde}/\varepsilon+4455$, since then $x^\theta\ge2C_{pde}$.
\item \emph{$k$ against the other ranks.} Fourteen inequalities hold,
among them
\[
x^\theta m^5fLx^\varepsilon\le k,\quad
x^\theta m^3L\log\mathrm{rd}_E\le k(\log k)^2,\quad
(\log f)^2\le m^2L,\quad f\le r^2,
\]
once $x^{\varepsilon/1000}\ge100+8|C_\Gamma|$. Each compares exponents of
$x$; for example the first left side is at most
$x^{\theta+5(1/2-\varepsilon/16+\nu)+2+8\varepsilon+\varepsilon/100+\nu+2\varepsilon}$
with $\nu=\varepsilon/1000$, which is at most $x^{9/2+10\varepsilon}$.
\item \emph{Host.} $Q_E^{36}t^{150}\le b^{1/2}$;
$x^{\varepsilon_z/6}\le\log q\le x^{\varepsilon/5+\theta}$;
$\log d\le x^{\varepsilon/5+\theta}$; $|\lambda|\le x^\varepsilon$;
$1\le\mathrm{rd}_E\le\frac{Q_S}2\sqrt{qd}$; $\log p_{de},\log p'_{de}\le L$.
These follow as in item (1) of (E1), with Proposition~\ref{i4:prop-rd}.
\item \emph{$\Lambda_G$ and the band.} $x^\varepsilon\le\Lambda_G\le x^{\varepsilon+\theta}$,
and $\frac{4N_*x^\varepsilon\Lambda_G}k+\frac{x^\theta\log x\,\Lambda_G}k\le x^{-2}$.
This is (e) together with the computation of item (7) of (E1).
\item \emph{$\eta_{sz}$} ($c_2=401$).
\begin{itemize}
\item $t\le\eta_{sz}\le b/(t^3fk^3)$, $\log\eta_{sz}\le x^{\varepsilon/4}$,
$t\eta_{sz}\le w''_n/x$ and $m(t+1)\eta_{sz}\le w''_n$. These follow from
$\eta_{sz}\le x^{c_h+3}$ and $x^{c_h+1224}\le w''_n\le b$, by (b).
\item $m(S_c+\eta_{sz})\le w''_nL$, $N_m\eta_{sz}\le\delta_2mw''_n$,
$\eta_{sz}^2rL\le\delta_2m(w''_n)^2$ and $kS_c\le w''_n$. These follow
from $M^{10}\le w''_n$ with $M=x^{c_D+c_h+402}$.
\item $r(b'+S_c)\log x\cdot l_A\le\delta_2mb$. This follows from
$b'\le b/(8x^{14})$, $S_c\le b/(8x^{14})$ and $l_A\le2x^7$.
\end{itemize}
\item \emph{Nodes} ($c_2=1300+c_D$). With
$D_l=nt^2(N_m+t)L^2+kn^2(S_c+t+L)$ and $l_*=mb/D_l$, each of
$12b/l_*+2t+8$, $ktL$, $x^\theta r^{19}f^5L$, $4N_m^3+2t+6$,
$x^\theta km^{9/2}$, $2(r+mt)^3+t+3$ and $N_mm^{10}$ is at most
$200x^{1299}x^{c_D}\le h_0$.
\item \emph{$l_*$.} $b'+S_c+\eta_{sz}+b/l_*<l_*$ and
$2\eta_{sz}+t+8b/l_*+7<l_*$. With $P=Q^7$ and $Q=x^{c_D+c_h+402}$, all of
$D_l,S_c,\eta_{sz},t$ are at most $P$, and $b'P^2\le b$, $P^2\le b'$ by (b).
Then $b/l_*\le P$ and $l_*\ge100b'$.
\item \emph{$\Delta$} ($c_1=412$). $N_mm^{10}\le\Delta$, $x^\theta mN_m\le\Delta$,
$\log\Delta\le\frac12x^{\varepsilon/4}$, $kS_c\le w''_n$,
$x^\theta mk\Delta\le w''_n$, $4S_c\le\Lambda_Gb'/k$,
$4x^\theta k^2\Delta\le\Lambda_Gb'$ and $x^\theta m\log(\mathrm{rd}_E/2)\le k$.
These follow from $16xk^2\Delta\le x^{c_D+12}\le w''_n\le b'$ and
$\log\mathrm{rd}_E\le L$.
\item \emph{Windows.} $b^{1/2}\le w''_n$, $b_{m-1}(\tau)\le2b_0(\tau)$,
and $w_j(\tau)\le X''_nw_{j'}(\tau')$. At every place $w''_n\le b_0(\tau)$
and $w_j(\tau)\le2b$, while $\eta_{sz}\le w''_n$. Hence, with $b_0(\tau)$ the first base value
(written $b_1(\tau)$ in step 2 of \S\ref{i4:ssec-proof-realise}),
$b_{m-1}=b_0+(m-1)h_0\le2b_0$ and $w_j\le2b\le xb=X''_nw''_n$.
\item \emph{Level primes and words.} $(kb)^8\le g_i\le2g_j$ and
$\log g_i\le x^\theta L$. The number of words
$\lceil3\exp(C_\Gamma m^5fL)^2x^{C_F}\rceil$ is at most $2^{k-3}$, because
$m^5fx^\varepsilon\le128x^{9/2+9\varepsilon}$ and
$k-3\ge2(256C_\Gamma C_{pde}+C_F+2)x^{9/2+9\varepsilon}$ for
$x^\varepsilon\ge2(256C_\Gamma C_{pde}+C_F+2)+3$.
\end{enumerate}
This completes the proof.
\end{proof}

\subsubsection{The host supply}\label{i4:ssec-proof-supply}

\begin{lemma}\label{i4:Slem-small}
(a) If $(q,d)\in\mathrm{WP}(\varepsilon,Y)$ then $(\frac{-d}q)=1$ and
$(\frac{-d}p)=1$ for every odd prime $p\le Y^{\varepsilon/5}$.
(b) If host data satisfy $\mathrm{KerChi}$, the field $E$ is a
$(1,\nu,q,d)$-field.
\end{lemma}
\begin{proof}
(a) If $p\mid d+1$ then $-d\equiv1\bmod p$. Apply this to $p=q$ and to the
odd $p\le Y^{\varepsilon/5}$. (b) $E\subset\Q(\zeta_N)$ with $N=Q_Sqd$,
divisible by $2^{\nu+2},q,d$; by Galois theory $\sigma_a$ fixes $E$ iff
$\sigma_a$ is in the kernel of the Artin map, which is the stated
condition by $\mathrm{KerChi}$.
\end{proof}

\begin{proof}[Proof of Proposition~\ref{i4:prop-supply}]
Let $\varepsilon\in(0,1/52)$. Take $c,e_0$ from
Proposition~\ref{i4:Sprop-field} and $J$ from
Proposition~\ref{i4:Sprop-win} (with this $c$). Put
$C_0=J+1+8\ge1$ (so $J,o=1,8\le C_0$), and take
$C_{T_2},C_8,c_6,C_b,x_A$ from Proposition~\ref{i4:Sprop-host}; $x_B$
from Proposition~\ref{i4:Sprop-ranks} (both parts). Put $K=5$ and $A=K^{5/\varepsilon}$,
so $A^{\varepsilon/5}=K$, and choose $x_0\ge x_A,x_B,2^{100},A^J$ with
$c2^{e_0}<x_0$. The constants of $\mathrm{HostExW}$ are
$(C_{T_2},C_8,C_5=4,c_6,C_b,x_0)$, with $C_8\ge C_0\ge1$; the set of
scales is: $J(s)=J$ and $x\in\mathrm{Win}=\{x:\exists Y,q,d,e,\ \mathrm{WinDatum}\}$,
which contains $[Z,Z^{3/2}]$ for arbitrarily large $Z$ by
Proposition~\ref{i4:Sprop-win}.

Let $s$ be such a scale with $x\ge x_0$, and $\mathrm{WinDatum}(\varepsilon,J,c,x,Y,q,d,e)$.
Then $Y\ge x^{1/J}\ge A$, so $Y^{\varepsilon/5}\ge K=5$. By
Lemma~\ref{i4:Slem-small}(a), $q,d$ satisfy the host hypotheses and
$(\frac{-d}p)=1$ for odd $p\le b:=\lfloor Y^{\varepsilon/5}\rfloor\ge4$; and
$q\ne3$ since $\log q>Y^{\varepsilon/5}/2\ge5/2>\log3$. Moreover $e\ge e_0$:
otherwise $c2^e\le c2^{e_0}<x_0\le x\le x^{2+8\varepsilon}\le c2^e$. So
Proposition~\ref{i4:Sprop-field} gives $\varphi,H_0$ with $\nu=e$,
$\mathrm{ord}=c2^e\in[x^{2+8\varepsilon},2x^{2+8\varepsilon}]$, $o=1$,
$\mathrm{FieldRows}(H_0,b)$, $Q_S\mid2^{e+2}$, $Q_S\le4\,\mathrm{ord}$,
$\mathrm{KerChi}$, $\mathrm{TwistForm}$. Proposition~\ref{i4:Sprop-ranks}
gives valid admitted ranks $R$ with frame $\varphi$, so
$n\ge x^{2+8\varepsilon}$, and then $Q_S\le4\cdot2x^{2+8\varepsilon}\le8n\le C_0n$.
Proposition~\ref{i4:Sprop-host} with $M=1$ gives $H$ and $q_0$, good with
the required constants and constant $C_{T_2}$.

It remains to check $\mathcal G_1(H)$. \emph{Twists:} $o=1$ by
Proposition~\ref{i4:Sprop-field}; the kept characters and the
exceptional-character clause by Proposition~\ref{i4:Sprop-field}; the
shares of the top set by Proposition~\ref{i4:Sprop-host}; and, since every
prime divisor of $Q_S$ is $2$, Proposition~\ref{i4:Sprop-tw} gives the two
bounds of Definition~\ref{i4:def-G1} at $y=x$. \emph{Selmer:} by
Lemma~\ref{i4:Slem-small}(b) $E$ is a $(1,e,q,d)$-field, so
$\mathrm{VV}(E,\lfloor x^{400}\rfloor,\sqrt{\exp(x^\varepsilon)})$; since
$t\le x^{400}$ (Proposition~\ref{i4:Sprop-ranks}, $x\ge2^{100}$) and
$\mathrm{VV}(E,t',B)$ is monotone (decreasing in $t'$, increasing in $B$),
$\mathrm{VV}(E,t,\sqrt b)$ holds with $b=\exp(x^\varepsilon)$.
\end{proof}

\subsection{The supply constructions}\label{i4:Sst-ssec}

This subsection concerns the five constructions
Propositions~\ref{i4:Sprop-field}--\ref{i4:Sprop-tw} used in the proof of
Proposition~\ref{i4:prop-supply}. Two of them are short and are proved
here in full (the twist bounds, \S\ref{i4:Sst-tw}, and the valid ranks
modulo two of the host axioms, \S\ref{i4:Sst-ranks}); one is proved here
modulo a single analytic statement, the existence of good window pairs
(\S\ref{i4:Sst-win}); the formally verified steps for the field of the
character and for the host are outlined in \S\ref{i4:Sst-outline}. Steps are
classified as (L), (N) or (R) as explained in the introduction to this annex.
Throughout, $\psi=1$ (Conventions (1)--(4)), and the
notation $\chi_q(a)=(\frac aq)$, $\chi_d(a)=(\frac ad)$ (Jacobi symbols) is
used for functions on the integers. A primitive Dirichlet character $\chi$
\emph{induces} a function $f$ \emph{modulo} $M$ if $\chi(a)=f(a)$ for every
integer $a$ prime to $M$.

The genuinely new mathematics in the supply is concentrated in three
places: the existence of good window pairs
(Proposition~\ref{i4:Sst-prop-pair}, a mean-square argument over quadratic
twists by primes, step (P3) below), the choice of the top set of a host
(step (H6)), and a variant of Behrend's construction (step (H5)). The rest
is cyclotomic class field theory, classical results on zeros of Dirichlet
$L$-functions (most of them re-derived in \cite{CharacterFamily}), the
short-interval bound Proposition~\ref{i4:thm-shortnorm}, the
smallest-prime bound Theorem~\ref{i4:thm-linnik}, and bookkeeping.

\subsubsection{The twist bounds}\label{i4:Sst-tw}

Let $H$ be host data over a frame $\varphi$, with group
$U=E_0\times A$, $E_0=\{\pm1\}^2$, and Artin map $\mathrm{Art}$ on
$(\Z/N_E)^\times$, $N_E=Q_Sqd$. Write $A_0\subset A$ for the subgroup
generated by $a_0$ and $\iota_d=((1,-1),0)\in U$ for the element with
$q$-sign $1$, $d$-sign $-1$ and trivial $A$-component. A character
$\theta$ of $U$ is \emph{trivial on the first factor} if it is trivial on
$\{(1,1)\}\times A_0$; then $\theta(\iota_d)=\pm1$, and we write $X^+$,
$X^-$ for the two sets of such characters with $\theta(\iota_d)=1$,
resp.\ $-1$. For a character $\theta$ of $U$ let $\chi_\theta$ be the
primitive Dirichlet character inducing $\theta\circ\mathrm{Art}$ modulo
$N_E$. The character $\theta$ is \emph{kept} if $\chi_\theta=1$. The
\emph{bounded set} is
\[
\mathcal B(H)=\{\theta\in X^+:\ \theta\ \text{not kept}\}\cup X^- .
\]

\begin{lemma}[twist bounds, general form]\label{i4:Sst-lem-tw}
Let $H$ be host data with $\mathrm{TwistForm}$ such that every prime
divisor of $Q_S$ is $2$. Let $\varepsilon>0$, $1\le x$, $Y\le x$, and
assume $\mathrm{PairGood}(\varepsilon,Y,x,q,d)$ for $(q,d)=(H.q,H.d)$.
Then for every $\theta\in\mathcal B(H)$,
\[
|G_{\chi_\theta}(x)|\le x^{1/2+\varepsilon},\qquad
|L'/L(1,\chi_\theta)|\le4\log x .
\]
\end{lemma}

\begin{proof}
Let $\theta\in\mathcal B(H)$. In both cases $\theta$ is trivial on the
first factor, so $\mathrm{TwistForm}$ gives $i,j\in\{0,1\}$ with
$\theta(\iota_d)=(-1)^j$ and $\theta(\mathrm{Art}(a))=\chi_q(a)^i\chi_d(a)^j$
for all $a$ prime to $N_E$. We claim $(i,j)\ne(0,0)$. If $\theta\in X^-$
then $(-1)^j=-1$, so $j=1$. If $\theta\in X^+$ and $i=j=0$, then
$\theta\circ\mathrm{Art}$ is trivial on the residues prime to $N_E$, so
$\chi_\theta=1$ and $\theta$ is kept, which is excluded. Since every prime
divisor of $Q_S$ is $2$, an integer prime to $2qd$ is prime to $N_E$;
hence $\chi_\theta$ induces $\chi_q^i\chi_d^j$ modulo $2qd$. Clause (iii)
of $\mathrm{PairGood}$ (Definition~\ref{i4:Sdef-win}), applied at $y=x\in[Y,x]$,
gives $|G_{\chi_\theta}(x)|\le x^{1/2+\varepsilon/2}\le x^{1/2+\varepsilon}$
(as $x\ge1$) and $|L'/L(1,\chi_\theta)|\le4\log x$.
\end{proof}

\begin{proof}[Proof of Proposition~\ref{i4:Sprop-tw}]
Let $\theta$ be one of the three nontrivial characters of $U$ trivial on
$A$. It is trivial on $\{(1,1)\}\times A_0$, so $\theta\in X^+\cup X^-$.
If $\theta\in X^+$, then $\theta$ is not kept: the Artin map is surjective
(it is the restriction $\mathrm{Gal}(\Q(\zeta_{N_E})/\Q)\to\mathrm{Gal}(E/\Q)$
followed by the identification $\mathrm{Gal}(E/\Q)\cong U$, and restriction is
surjective by Galois theory), so $\theta\circ\mathrm{Art}\ne1$ and
$\chi_\theta\ne1$. Hence $\theta\in\mathcal B(H)$, and
Lemma~\ref{i4:Sst-lem-tw} gives the two bounds.
\end{proof}

\noindent\emph{Classification.} Every step is (R). Lemma~\ref{i4:Sst-lem-tw}
is the exact form of the formal statement (it is stated for the bounded
set, without the conventions $o=1$, $\psi=1$ that identify that set with the
three characters); both it and Proposition~\ref{i4:Sprop-tw} are proved
here in full.

\subsubsection{Window data}\label{i4:Sst-win}

The formal proof of Proposition~\ref{i4:Sprop-win} combines the windowed
Selmer bound (the hypothesis $\mathrm{VWinAt}(1,401)$,
Proposition~\ref{i4:hyp-selmer}) with the following statement, which
carries all the analysis.

\begin{proposition}[good window pairs; formally verified, \cite{CharacterFamily}]\label{i4:Sst-prop-pair}
For $0<\varepsilon<1/52$ and $A\ge1$ there is $Y_1$ such that for every
$Y\ge Y_1$ there are $q,d$ with $\mathrm{PairGood}(\varepsilon,Y,Y^A,q,d)$.
\end{proposition}
\noindent (The formal statement is proved for $\mathrm{PairGood}$ with an
arbitrary auxiliary modulus $M\ge1$; only $M=1$ is used here.)

\begin{proof}[Proof of Proposition~\ref{i4:Sprop-win} from Proposition~\ref{i4:Sst-prop-pair}]
Let $\varepsilon\in(0,1/52)$ and $c\ge1$. Choose $e:\R\to\N$ with
$x^{2+8\varepsilon}\le2^{e(x)}\le2x^{2+8\varepsilon}$ for $x\ge1$ (take
$2^{e(x)}$ the least power of $2$ that is $\ge x^{2+8\varepsilon}$).
$\mathrm{VWinAt}(1,401)$ gives $J\ge1$ and $X_0$. Proposition~\ref{i4:Sst-prop-pair}
with $A=J$ gives $Y_1$. Let $k=c^{1/(2+8\varepsilon)}\ge1$, so
$k^{2+8\varepsilon}=c$. Given $y$, put
$M=\max(1,y,Y_1,k^3,k^{400})$ and $X=\max(X_0,M^J)$. Then
$X^{1/J}\ge M\ge Y_1$, so there are $q,d$ with
$\mathrm{PairGood}(\varepsilon,X^{1/J},X,q,d)$; in particular $(q,d)$ is a
window pair for $(\varepsilon,X^{1/J})$, and $\mathrm{VWinAt}$ gives $W$ with
$X^{1/J}\le W$, $W^2\le X$ and the Selmer bound
$\mathrm{VV}(E,\lfloor x'^{401}\rfloor,\sqrt{\exp(x'^\varepsilon)})$ for all
$x'\in[W,W^2]$ and all $(1,e(x'),q,d)$-fields $E$.

Put $Z=kW\ge M\ge y$. Since $k^3\le M\le W$ we have
$Z^{3/2}=(kW)^{3/2}\le(W^4)^{1/2}=W^2$. Let $x\in[Z,Z^{3/2}]$ and write
$x=kx'$; then $x'\in[W,W^2]$ and $x\le W^2\le X$. Take
$Y=X^{1/J}$ and $e=e(x')$. Then: $x^{1/J}\le X^{1/J}=Y$ and
$Y\le W\le x$; $\mathrm{PairGood}(\varepsilon,Y,x,q,d)$ follows from
$\mathrm{PairGood}(\varepsilon,Y,X,q,d)$ since $x\le X$ (the clauses only
shrink); $x^{2+8\varepsilon}=c\,x'^{2+8\varepsilon}\le c2^e\le
2c\,x'^{2+8\varepsilon}=2x^{2+8\varepsilon}$. Finally, a $(1,e,q,d)$-field $E$
satisfies $\mathrm{VV}(E,\lfloor x'^{401}\rfloor,\sqrt{\exp(x'^\varepsilon)})$,
and since $x^{400}=k^{400}x'^{400}\le x'^{401}$ (as $k^{400}\le M\le x'$)
and $x'\le x$, monotonicity of $\mathrm{VV}$ (decreasing in the first
argument, increasing in the second) gives
$\mathrm{VV}(E,\lfloor x^{400}\rfloor,\sqrt{\exp(x^\varepsilon)})$. This is
$\mathrm{WinDatum}(\varepsilon,J,c,x,Y,q,d,e)$.
\end{proof}

\noindent\emph{Structure of the proof of Proposition~\ref{i4:Sst-prop-pair}.}
Fix $\varepsilon$, $A$, and put
\[
z=Y^{\varepsilon/5},\qquad z'=Y^{\varepsilon/5}(\log Y)^2,\qquad
\tau=\varepsilon/2,\qquad K=A .
\]
For a primitive character $\chi$ and $Q\ge1$ say that $\chi$ is
\emph{$Q$-boxed} if $\chi\ne1$, its conductor is at most $Q$, and $L(s,\chi)\ne0$
for $\Rel s>1-1/400$, $|\Iml s|\le Q$. Say that $\chi$ is
\emph{$(Q,C)$-good} if it is $Q$-boxed, $|L'/L(1,\chi)|\le C\log Y$, and
$|G_\chi(y)|\le y^{1/2+\tau}$ for all $y\in[Y,Y^K]$. Let
$P_1$ be the set of primes $q\equiv3\bmod8$ with $e^{z/2}<q\le e^z$ and,
for $q\in P_1$, $P_2(q)$ the set of primes $d\equiv7\bmod16$ with
$e^{z'/2}<d\le e^{z'}$ and $d\equiv-1$ modulo $q$ and modulo every odd
prime $p\le z$. The steps are:

\begin{enumerate}
\item[(P1)] (R) Every $q\in P_1$ and $d\in P_2(q)$ form a window pair
$(q,d)\in\mathrm{WP}(\varepsilon,Y)$, and $q<d$.

\item[(P2)] (N) \emph{The identity for $G_\chi$ and the truncation.} For a
primitive $\chi\ne1$ and $y\ge1$,
\[
G_\chi(y)=-\sum_{n\ge1}\Lambda(n)\chi(n)\min(1,y/n),
\]
because $L'/L(1,\chi)=-\sum_n\Lambda(n)\chi(n)/n$ (the prime number
theorem for $\chi$, (L)). Truncating at $Z=Y^{400K}$ gives
$G_\chi(y)=-S_Z(y,\chi)+y\sum_{n>Z}\Lambda(n)\chi(n)/n$ with
$S_Z(y,\chi)=\sum_{n\le Z}\Lambda(n)\min(1,y/n)\chi(n)$. The formal
statement is: there is an absolute $C$ such that a $Q$-boxed $\chi$ with
$Y^{\varepsilon/6}\le\log Q\le Y^{\varepsilon/4}$ and
$|S_Z(y,\chi)|\le y^{1/2+\tau}/2$ for all $y\in[Y,Y^K]$ is $(Q,C)$-good, for
$Y$ large in terms of $\varepsilon,\tau,K$. Its two analytic ingredients are
a bound for the tail $y\sum_{n>Z}\Lambda(n)\chi(n)/n$ and the bound
$|L'/L(1,\chi)|\le C\log Y$, both deduced from the zero-free box by contour
arguments of the kind of Landau's lemma (Lemma~\ref{i4:lem-landau}); we do not
describe their formal proofs here.

\item[(P3)] (N) \emph{Mean square over a family of quadratic twists.} Let
$8\mid m_0$, $a_0$ prime to $m_0$, and $U$ with
$\tfrac14(\log Y)^2\log m_0\le\log U$ and
$Y^{\varepsilon/6}\le\log U\le Y^{\varepsilon/5}(\log Y)^2$, and assume every
prime factor of $m_0$ is at most $Y^{\varepsilon/5}$ or larger than
$Y^{400K}$. Let $P=P(m_0,a_0,U)$ be the set of primes $p\equiv a_0\bmod m_0$
with $U^{1/2}<p\le U$. Then for $Y\ge Y_0(\varepsilon,K)$, every
$y\in[Y,Y^K]$ and every $\kappa:\N\to\C$ with $|\kappa|\le1$,
\[
\sum_{q\in P}\Big|\sum_{n\le Z}\Lambda(n)\min(1,y/n)\kappa(n)\chi_q(n)\Big|^2
\le C\,y\log y\,|P| .
\]
Idea: open the square. The inner sum $\sum_{q\in P}\chi_q(n_1n_2)$ is, by
quadratic reciprocity (the residue of $q$ modulo $8$ is fixed by $a_0$), a
sign times $\sum_{q\in P}(\frac qs)$ with $s$ essentially the odd squarefree kernel of
$n_1n_2$ prime to $m_0$. The terms with $s=1$ form the diagonal, which
contributes $\ll y\log y\,|P|$ because $\sum_n\Lambda(n)^2\min(1,y/n)^2\ll
y\log y$ and $n_1n_2$ is then a square up to primes dividing $m_0$. For
$1<s\le Z^2$ the sum over $q$ is at most $Y^{-B}|P|$ by (P4), with one
possible exceptional $s$, whose bookkeeping we do not describe. The large size of the family ($\log|P|$ is of order at least $Y^{\varepsilon/6}$)
against the length $Z=Y^{400K}$ of the sums is what makes the diagonal
dominate.

\item[(P4)] (L) \emph{Primes in progressions with a quadratic twist.} For
odd squarefree $1<s\le Y^{A'}$ prime to $m_0$, outside at most one
exceptional value, $|\sum_{q\in P}(\frac qs)|\le Y^{-B}|P|$; and
$|P|\ge U^{1-o(1)}$. This is the prime number theorem for arithmetic
progressions to the modulus $m_0s$, in the form with the classical
zero-free region and at most one exceptional real character (Landau--Page;
Davenport, Ch.~14 and 20; (Z2), (Z4) of Theorem~\ref{i4:thm-classzeros}).
Extra in the formal version: explicit uniformity in a modulus $m_0$ as large
as $\exp(4\log U/(\log Y)^2)$, where the saving is
$\exp(-c(\log Y)^2/4)$. In \cite{CharacterFamily} it is re-derived,
together with Siegel's theorem and the Siegel--Walfisz theorem (Davenport,
Chs.~21--22); we do not describe the derivation here.

\item[(P5)] (L) \emph{Zero density.} There is an absolute $C$ such that
for $Q\ge3$ at most $CQ^{1/4}$ primitive characters of conductor at most
$Q$ have a zero with $\Rel s>1-1/400$, $|\Iml s|\le Q$. This is the
consequence $(\mathrm A_6)$(i) of the zero-density estimate (Z5) of
Theorem~\ref{i4:thm-classzeros} derived in the proof of
Proposition~\ref{i4:thm-litthin}. In \cite{CharacterFamily} it is
re-derived (mollifier, zero detection, duality and a large-sieve bound);
we do not describe this derivation.

\item[(P6)] (N) \emph{Most twists are good (``stage'').} With
$m_0,a_0,U,P$ as in (P3), let $\Xi$ be a set of at most $R$ primitive
characters of conductor at most $F$, $1\le F\le U^{1/2}$. Then for every
$\eta>0$ and $Y\ge Y_0(\varepsilon,\tau,K,\eta,R)$, the set of $q\in P$
for which some $\xi\in\Xi$ and some primitive $\chi$ inducing $\xi\chi_q$
is not $(FU,C)$-good has at most $\eta|P|$ elements. Proof: such a $q$ is
bad either because $\chi$ is not $FU$-boxed or because
$|S_Z(y,\chi)|>y^{1/2+\tau}/2$ for some $y\in[Y,Y^K]$ (by (P2)). Distinct
$q$ give distinct characters (the conductor is $fq$ with $f\le F<q$), so by
(P5) at most $C(FU)^{1/4}\le CU^{3/8}$ primes are bad for the first
reason, a negligible proportion by (P4). For the second reason, $S_Z(y)$ is
piecewise differentiable in $y$; on blocks $[Y',W]$ with $W\le2Y'$ the
supremum of $|S_Z|^2$ is bounded by its mean plus an integral involving
$\partial_yS_Z$, the mean square (P3) bounds the average over $q$,
and Markov's inequality, with the saving $y^{2\tau}/\log y$ against the
$O(K\log Y)$ blocks, bounds the proportion of bad $q$ by $\eta/2$.

\item[(P7)] (R) \emph{The two stages.} Apply (P6) with $(m_0,a_0,U)=(8,3,e^z)$
and $\Xi$ the primitive characters inducing the characters modulo $8$: all
but a proportion $\eta$ of $q\in P_1$ are good. For such $q$, apply (P6)
with $m_0=16\cdot q\prod_{3\le p\le z}p$, $a_0$ given by the Chinese
remainder theorem ($a_0\equiv7\bmod16$, $a_0\equiv-1$ modulo the odd part),
$U=e^{z'}$, and $\Xi$ the primitive characters inducing $\xi\chi_q^a$
($\xi$ modulo $8$, $a\in\{0,1\}$): all but a proportion $\eta$ of
$d\in P_2(q)$ are good. The hypotheses of (P3) hold since
$\log m_0\le3z$ (Chebyshev) and $\log U=z'=z(\log Y)^2$. Take $\eta=1/2$ and
choose $q$, then $d$.

\item[(P8)] (N, short) \emph{The constant $4$.} For a $Q$-boxed $\chi$ with
$\log Q$ large, $|L'/L(1,\chi)|\le1000\log\log Q+C$. Here
$\log Q\le3Y^{\varepsilon/5}(\log Y)^2$, so
$1000\log\log Q\le200\varepsilon\log Y+O(\log\log Y)<4\log Y$ for large $Y$,
because $200\varepsilon<200/52<4$. This is where the bound $\varepsilon<1/52$
and the constant $4$ of $\mathrm{PairGood}$ come from. The bound
$1000\log\log Q+C$ is proved by a contour argument in the box, which we do
not describe.

\item[(P9)] (R) \emph{Assembly.} For the chosen $(q,d)$ and $M=1$:
clause (i) of $\mathrm{PairGood}$ is (P1); (ii) is trivial; (iii) follows from
$(Q,C)$-goodness with $\tau=\varepsilon/2$ on $[Y,Y^A]$ and from (P8) since
$\log y\ge\log Y$; (iv) is the box part of goodness, with
$Q=8e^z$ for $j=0$ (a character inducing $\rho\chi_q$ modulo $2qd$ already
induces it modulo $8q$) and $Q=8e^{z+z'}$ for $j=1$, which are the heights
$8\exp(Y^{\varepsilon/5})\exp(Y^{\varepsilon/5}(\log Y)^2)^j$ of
Definition~\ref{i4:Sdef-win}.
\end{enumerate}

\noindent What rests on \cite{CharacterFamily}: the mean-square bound (P3),
the stage (P6) and the analytic bounds of (P2) and (P8) are new in this
form (the methods are classical: explicit formula, large-sieve-type
orthogonality over a family of quadratic characters, zero density, Markov);
(P4), (P5) are re-derivations of published theorems with explicit
uniformity.

\subsubsection{Valid ranks}\label{i4:Sst-ranks}

Recall (\S\ref{i4:ssec-proof-kits}) that ranks over a frame $\varphi$ with
$|A|=n$ are integers $k,m,t$ with $k,m$ even, $t$ odd, $m\ge2$,
$m+3\le k$, $k+m+4\le t$, $3m<t$, $m\le f=2n\le k$ and
$m^2\le k+m(t+1)$; $r=k+m-1$. They are \emph{valid} at a scale $x$ of
exponent $\varepsilon$ if
$x^{2+8\varepsilon}\le n\le2x^{2+8\varepsilon}$,
$\sqrt{x\delta_1}\le m\le2\sqrt{x\delta_1}$,
$x^{9/2+10\varepsilon}\le k\le2x^{9/2+10\varepsilon}$ and $r+t=n_P+T_On_W$,
with $n_P=8fr+4$, $T_O=\lceil x^{377}\rceil$,
$n_W=8\lceil\frac{10}8r(f+f^2(k-1)^3L^2/m)\rceil$ (so $n_W$ does not depend
on $t$). Ranks $R$ are \emph{admitted} for host axioms $H_0$ over their
frame if, for every $b$, every choice of level data (level primes, the
primes $p_{de},p'_{de}$, the thick set $T_2$, the window length $L_w$, the
constant $C_{T_2}$) and whenever $\mathrm{FieldRows}(H_0,b)$ holds and
$q,d>4tr+1$, there is a host over $R$ extending $H_0$ with these choices
and split bound $b$.

\begin{proof}[Proof of Proposition~\ref{i4:Sprop-ranks}]
\emph{Existence.} Let $x\ge16$ and $\varphi$ with
$x^{2+8\varepsilon}\le n\le2x^{2+8\varepsilon}$. Since
$\sqrt{x\delta_1}=x^{1/2-\varepsilon/16}\ge2$ and every interval $[a,2a]$
with $a\ge2$ contains an even integer, choose even $m$ with
$\sqrt{x\delta_1}\le m\le2\sqrt{x\delta_1}$ and even $k$ with
$x^{9/2+10\varepsilon}\le k\le2x^{9/2+10\varepsilon}$. Put $r=k+m-1$,
$f=2n$ and $t=n_P+T_On_W-r=8fr+4+T_On_W-r$. Then $m\le2x\le f$ and
$f\le4x^{3}\le x^{9/2}\le k$ (as $x\ge16$), $m+3\le k$, and $t\ge8r+4-r\ge
r+5=k+m+4$ (as $f\ge1$), so $t>3m$ and $m^2\le m(t+1)$. Since $8\mid n_W$
and $r$ is odd, $t$ is odd. So $(k,m,t)$ are ranks over $\varphi$, and they
are valid by construction.

\emph{Admission.} Given the choices, a host extending $H_0$ consists of the
data of $H_0$, the choices, and three further items: a square root
$\sqrt{-q}\in E$ on which $\mathrm{Gal}(E/\Q)$ acts through the $q$-sign
(part of $\mathrm{FieldRows}$); an embedding $E\to\C$ in the CM type (the
CM type is nonempty, as it has $2n$ elements); and an ideal $\mathfrak b$
of $\mathcal O_E$ with $\mathfrak D_{E/K}=\mathfrak b\bar{\mathfrak b}$,
$\mathfrak b+\bar{\mathfrak b}=\mathcal O_E$, stable under
$\mathrm{Gal}(E/K)$, where $K=\Q(\sqrt{-d})$ and $\mathfrak D_{E/K}$ is the
relative different. The last two are among the host axioms (the CM-type and
different clauses), so all ranks are admitted.

\emph{The bound on $t$.} Let $x\ge2^{100}$ and put $y=x^{1/100}\ge2$. With
$\varepsilon<1/52$: $n\le2x^{2+8\varepsilon+\varepsilon/100}\le2y^{216}$, so
$f\le4y^{216}$; $k\le2y^{470}$, $r\le2k\le4y^{470}$; $m\ge x^{1/2-\varepsilon/16}\ge y^{49}$;
$L^2=x^{2\varepsilon}\le y^4$. Hence
$f^2(k-1)^3L^2/m\le16\cdot8\,y^{432+1410+4-49}=128y^{1797}$,
$f+f^2(k-1)^3L^2/m\le132y^{1797}$, and $n_W\le\frac{10}{8}\cdot8\cdot
4y^{470}\cdot132y^{1797}+8\le5288y^{2267}$. Also $T_O\le2y^{37700}$. So
\[
t\le8fr+4+T_On_W\le128y^{686}+4+10576y^{39967}\le10708y^{39967}\le
y^{40000}=x^{400},
\]
using $y^{33}\ge2^{33}\ge10708$.
\end{proof}

\noindent\emph{Classification.} Existence and the bound on $t$: (R).
Admission: (R), given two clauses of the host axioms (the CM type and the
factorisation of the relative different), which are conclusions of
Proposition~\ref{i4:Sprop-field}; the different clause is discussed in
(F4) below. The proposition is proved here in full relative to the
definition of host axioms.

\subsection{The division side}\label{i4:ssec-grp-div}

\subsubsection{The closed estimates on the division side}\label{i4:ssec-proof-links}

By Lemma~\ref{i4:lem-closed} the three properties may be treated
separately, and a conjunction of closed properties is closed. We use the
division kernel constant bounds $\mathrm{DKCB}(\mathcal F,C_{\rm abs})$ of
Definition~\ref{i2:def-DKCB}. They consist of a package clause (pkg), a
frame-number bound (hh), and the clauses adm2 (admissibility), cut, c2,
c3, c5 (complements), $\mathrm{BuRational}$, $\mathrm{RRFull}$ and
$L_a\ge0$.

\begin{lemma}[the division bridge]\label{i4:lem-bridge}
Let $C_{\rm abs}\ge1$. $\mathrm{DKCB}(\mathcal F,C_{\rm abs})$ implies
$(\mathrm{LDivt})$ for $\mathcal F$. Consequently, if the bounds are closed on a class under
$\mathrm{TuSize}\wedge\mathrm{DGamSize}$, so is $(\mathrm{LDivt})$.
\end{lemma}
\begin{proof}
Assume $(\mathrm{LEta})$ and $(\mathrm{DivN})$. Let $M$ be the maximum of
the thresholds of all clauses and of $x_1+1$, and let
$C=\max(C_1,C_2,C_3,0)$ be the maximum of the package and complement
constants. At $x\ge M$, the package at $(x,i)$ gives the word facts. Its
ties give $\mu_{CC}=\mu^\vee$, $h_u-\overline{\log\kappa}=h(u)$ and
$2F^R_o=F$. Its divisor inequality is the display of
Definition~\ref{i2:def-DKCB}. Since $L_a\ge0$, $|T|a\ge0$,
$m(B_0+N_c)\ge0$ and $x^\varepsilon\le x^{C_{\rm abs}\varepsilon}$ (as
$x\ge1$ and $C_{\rm abs}\ge1$), the right side is at most
$\mathrm{MAIN}(u)+C|T|aL_a+Cm(B_0+N_c)x^{C_{\rm abs}\varepsilon}$, which is
(P1) with exponent $C_{\rm abs}$ (Proposition~\ref{i2:prop-P1}). We now
obtain the nine conjuncts of $(\mathrm{LDiv})$ at every point with
$x\ge M$:
\begin{itemize}
\item the single-map condition (1), the height-map bound (6) and (P1) (7)
are the word facts of the package, as just explained;
\item the frame-number bound (2) is the clause (hh), which applies since
$(\mathrm{DivN})$ holds;
\item the vanishing above the levels (3) is Lemma~\ref{i4:lem-zeroabove},
since $M\ge x_1+1$;
\item the Bessel identity (4) is $\mathrm{BuRational}$ together with c5;
\item $\mathrm{RRFull}$ (5) is the clause of that name;
\item the c2 bound (8) is the clause c2;
\item the variation bound (9) is the clause c3, which is exactly this
bound for the explicit height $h(u)$ of the point.
\end{itemize}
Each bound is monotone in its constant, so all hold with the common
constant $C$.
\end{proof}

\begin{lemma}[group order]\label{i4:lem-gzero}
Let $G$ be a finite group of order $\prod_{k\in\{0,1\}}|\mathrm{GL}_r(\Z/g_k)|^f$
with $g_0,g_1\le2Y$. Then $\log|G|\le2fr^2\log(2Y)$.
\end{lemma}
\begin{proof}
$|\mathrm{GL}_r(\Z/g)|\le g^{r^2}$.
\end{proof}

\begin{lemma}[the denominator $d_g$]\label{i4:lem-dgam}
Let a point have valid ranks (so that $2r^2\le t^{150}$ and
$N_m\le3t^2$) and size data with $(kb)^8\le g_i\le2g_j$,
$\log g_i\le x^\theta L$ and $(Q_Sqd)^{36}t^{150}\le\sqrt b$. Let
$\log D_\Gamma(u)\le C\,fN_m\log(fN_m)$ for every word $u$. Let the selected
datum have group $G$ as in Lemma~\ref{i4:lem-gzero}, and suppose every
$|\Gamma_c|$ divides $D_\Gamma(u_0)$. Let $d_g$ be the least common multiple
of the integers $D_G=|P|\,\mathrm{lcm}_c|\Gamma_c|$ of the level lattices.
Then for every $N_c\ge0$, with $U=L(fmbL+N_c)$,
\[
\log d_g\le(2+9C)\,U .
\]
\end{lemma}
\begin{proof}
By Lagrange's theorem, $d_g$ divides $|G|D_\Gamma(u_0)$. Next,
$\log(2g_0)\le2L^2$ and $2r^2\le t^{150}\le\sqrt b\le mb$. With
Lemma~\ref{i4:lem-gzero} this gives $\log|G|\le2fr^2\log(2g_0)\le2U$.
Put $a_0=fN_m$. Since $f\le Q_Sqd$ (the Galois group of the host has order
$4n$ and is a quotient of $(\Z/Q_Sqd)^\times$) and $N_m\le3t^2$, we get
$a_0\le3\sqrt b$ and $\log a_0\le3L$. Hence
$\log D_\Gamma(u_0)\le C\cdot9bL\le9CU$.
\end{proof}

\begin{lemma}[kernel join]\label{i4:lem-kerjoin}
Let $\mathcal E,\mathcal E_P$ be finite-dimensional Hilbert spaces, and let
$e_R:\mathcal E\to\mathcal E_P^+\subset\mathcal E_P$ be $c_s$-semilinear and
onto, with $\|e_RF\|=c\|F\|$ for some $c>0$. Here $c_s$ is the identity or
complex conjugation. Let $V$ be a linear map of $\mathcal E_P$ preserving
$\mathcal E_P^+$ isometrically. Let $c_u$ and $U_\mu$ be isometries, and
let $\alpha\ne0$. Let $B_1,B_2$ be bilinear forms on
$\C^{a}\times\C^{n}$ that agree on the pairs of standard basis vectors.
Suppose
\begin{itemize}
\item the period is $\mathrm{per}(\varphi,c_Fy)=B_1(c_u\varphi,y)$;
\item $B_2(x,y)=c_s(\alpha\langle s_E(Ve_R(\iota c_Fy)),P(U_\mu c_s(x))\rangle)$;
\item $\langle s_E(e_RF),P(c_s(c_u\varphi))\rangle=c_s(\kappa\langle\overline{j(F)},P_u(T)\rangle)$
with $\kappa\ne0$.
\end{itemize}
Then there are a linear isometry $V'$ of $\mathcal E$, an isometry $U'$ and
$a\ne0$ with
$\mathrm{per}(\varphi,c_Fy)=a\langle\overline{j(V'\iota c_Fy)},P_u(T)\rangle$,
for every column $T$ whose dictionary image is the ramified vector of
$U'\varphi$.
\end{lemma}
\begin{proof}
Define $V'$ by $e_R(V'F)=V(e_RF)$. It is well defined and linear, because
$e_R$ is injective, $c_s$-semilinear and onto $\mathcal E_P^+$, and $V$
preserves $\mathcal E_P^+$. It is isometric, because
$c\|V'F\|=\|Ve_RF\|=c\|F\|$. Put
$U'=c_u^{-1}c_sU_\mu c_sc_u$ and $a=c_s(\alpha)\kappa$. Then
$\mathrm{per}(\varphi,c_Fy)=B_2(c_u\varphi,y)$, and substituting the
definition of $V'$ and the third identity gives the claim.
\end{proof}

\begin{proof}[Proof of Proposition~\ref{i4:prop-links-formal}]
\emph{(a)} By Lemma~\ref{i4:lem-bridge} it suffices that the clauses of
$\mathrm{DKCB}$ with $C_{\rm abs}=2$ are closed on $D_3(\mathcal A)$ under
$\mathrm{TuSize}\wedge\mathrm{DGamSize}$. The clauses adm2, cut, c2, c3,
hh, $\mathrm{BuRational}$, $\mathrm{RRFull}$ and the inputs of c5 are
(L0) of Proposition~\ref{i4:prop-links-leaves}, and $L_a\ge0$ holds at every
point. It remains to treat c5 and the package clause.

\emph{The clause c5.} Fix a member and a point beyond the thresholds of the
gap hypothesis, of $\mathrm{FinX}$, of the multiplicity condition and of
the full base-vector image (all from (L0) and the hypothesis structure).
The selection joins the realisation at the point to its kernel. The
multiplicity condition $\mathrm{MI}^*_2$ provides $V$, scalars $a_\mu\ne0$,
matrices $U_\mu$ and normalising constants $\zeta_\mu$ with
$\zeta_\mu(\mu_0)=1$. Taking $t=\mu_0$ in $\mathrm{MI}^*_2$, together with
the character tie and the row tie of the join, shows that the
forms $B_1,B_2$ of Lemma~\ref{i4:lem-kerjoin} agree on basis vectors. So
the lemma gives the kernel join.

The real spectral ties give two more facts. First, the column dictionary
is bijective and pinned, so the column join holds with $W=\mathrm{id}$.
Second, the spectral row of $F$ is $\overline{j(F)}$. By the full
base-vector image, every $x\in\mathcal E$ is $\iota q$ for a unique $q$ in
the orthogonal complement of the kernel part. Since the Bessel functional
is $B_u(\varphi,x)=\mathrm{per}_u(\varphi,q)$, the kernel join gives
$B_u(\varphi,x)=a\,\beta_S(u)(U'\varphi,V'x)$. This is c5.

\emph{The package clause.} Fix a member with $C_{sz}\ge0$,
$\mathrm{DGamSize}$ with constant $C_d$, the row field (constant $C$), the
single-map condition and $\mathrm{DIVinN}$ (constant $C_e$). Put
$c_D=2+9\max(C_d,0)$ and
\begin{equation}\label{i4:eq-c0ce}
c_0=\max\bigl(28,\ 2(1+|C_c|)+1,\ C_{sz}(3+2C_{YL}+\max(C_t,5)+3)\bigr),\qquad
c_e=\max\bigl(c_0,\ 2C+C_e+c_D+3\bigr).
\end{equation}
Take $x$ beyond the thresholds of these fields and beyond an explicit scale
threshold depending on $(\varepsilon,C_t,C_c,C_H,c_e)$. At such a point:
\begin{itemize}
\item[(i)] The eight conditions of the selected indefinite realisation
form a primitive indefinite package with full Hecke data.
\item[(ii)] Lemma~\ref{i4:lem-dgam}, applied with $C=\max(C_d,0)$, gives
$\mathrm{DGamSize}$ for the selected level-three datum with constants
$(1,c_D)$.
\item[(iii)] The arithmetic clause holds at every word, by (L1).
\item[(iv)] The package lemma (L2) then gives a division kernel package
with constants $(\max(0,50c_0+10,60c_e),2)$.
\end{itemize}
This is the package clause with $C_{\rm abs}=2$.

\emph{(b)} The real-place second-layer ties are closed on $D_3(\mathcal A)$
(field (vii)). Beyond their threshold each point carries tied volumes,
column-value data and level-zero lattice data. From these:
\begin{itemize}
\item the free $T$-part and the slope line, which (L3) supplies at every
word, give the division side conditions;
\item the second-branch length and the two Hecke conditions at level zero,
also from (L3), give the congruence length bound.
\end{itemize}
Together with the eventual setting provided by $\mathrm{TuSize}$, these are
the level conditions: $d_u=d^S_u(1-1/(T-1))^f$ for every word, and the
length bound.

The level conditions make every $\nu_u$ commensurable with the level-zero
lattice, which spans the space of $\nu_u$. Hence every vector has a nonzero
integer multiple in $\nu_u$, which is fullness. The proviso of
$\mathrm{ChosenE}(3,C_H)$ gives $d^S_u>0$, so $d_u>0$, which is positive
rank. Since $\mathrm{TuSize}$ is a conjunct of $\mathrm{Des}$, it is closed
on $D_3(\mathcal A)$, and Lemma~\ref{i4:lem-closed} gives (b).

\emph{(c)} The ties, the volume ratio and the genus mass in rank two
(Theorem~\ref{i4:thm-genusmass}) give the eleven-item minimal global record
on $D_3(\mathcal A)$. The hypothesis $C_H>0$ of
Proposition~\ref{i4:Lprop-fam}(b), (c) is obtained by closedness and
conjoined with the record: every class is closed for $C_H>0$, with the
threshold datum whose middle constants equal $1$, since a member large
for it has every entry of $c_{\rm mid}(C_H)$, in particular $C_H$, at
least $1$. Four of its items hold for every family:
\begin{itemize}
\item the rank-two mass ratio;
\item the two-sided bound on the $\sigma$-degree;
\item the lower open part of the slope;
\item the maximal open part of the slope.
\end{itemize}
Base volume equals mass by the unfolding argument
(\S\ref{i4:ssec-proof-realise}, step 6). The remaining items are the
column values, the free $T$-part, the slope line, the second-branch
length and the two Hecke conditions at level zero; these come from the
ties and (L3).

Now fix a member with $\mathrm{TuSize}$, $\mathrm{DGamSize}$ and the
hypotheses of $(\mathrm{LGl})$. The steps, in order:
\begin{enumerate}
\item The level conditions, commensurability and fullness hold by (b).
\item $\mathrm{MuRank}$ holds by the proviso.
\item (L4) gives the family-level items and the frame bound.
\item The losses $\mathrm{LossFS}(\cdot,0)$ come from
Proposition~\ref{i4:hyp-loss}(b) together with the wall clause, which is a
duty of the class.
\item The core link of (L4) gives chain-tied constants with
$e_D=C_{\rm abs}=2$ and $e_L=0$.
\end{enumerate}
This is $(\mathrm{LGl})(2,C_H,3,2,0)$.
\end{proof}

\subsubsection{The division and global leaves}\label{i4:ssec-proof-Lleaves}

We prove Proposition~\ref{i4:prop-links-leaves} item by item. Throughout,
$\mathcal F\in D_3(\mathcal A)$ and $(x,i)$ is a point with $x>x_1$; the
size record then holds at $(x,i)$, and we use from it: $0<\varepsilon<1/52$,
$x\ge1$, $b\ge1$, $m\ge2$, $k\ge m+3$, $t\ge r+5\ge5$, $f\ge1$, and the
bounds $n\le x^{2+8\varepsilon+\varepsilon/100}$,
$k\le x^{9/2+10\varepsilon+\varepsilon/100}$, $m\le x^{1/2+\varepsilon/100}$,
$t\le x^{C_t}$, and $Q_S\le4n\,x^{\varepsilon/100}$.

\medskip\noindent\emph{Two elementary bounds.}

\begin{lemma}[absorbing the size polynomial]\label{i4:Llem-bpol}
If $x\ge\max(9,C)$ and $P\le x^{M_t}$, then $CP^eb\le x^{eM_t+1}b$ for every
integer $e\ge0$.
\end{lemma}
\begin{proof}
$1\le P\le x^{M_t}$ gives $P^e\le x^{eM_t}$, and $C\le x$; multiply by
$b\ge0$.
\end{proof}

\begin{lemma}[factorial threshold]\label{i4:Llem-fact}
Let $e\ge0$ and $N=\lceil(e+1)/\varepsilon\rceil$. If $x\ge N!$ and
$x\ge1$, then $x^e\le b$.
\end{lemma}
\begin{proof}
$x\ge N!$ gives $x^e\le x^{e+1}/N!\le x^{\varepsilon N}/N!\le
e^{x^\varepsilon}=b$, the last step by one term of the exponential series.
\end{proof}

\medskip\noindent\emph{(L0): the clauses of the division kernel bounds.}

\begin{lemma}[gauge form bound]\label{i4:Llem-gafo}
If $\mathcal F$ satisfies c2, there is $x_0$ such that for $x\ge x_0$,
$x>x_1$, every $i$ and every word $u$,
$\mathrm{Ga}+\mathrm{Fo}-E_e(u)\le x^{E_c}b$.
\end{lemma}
\begin{proof}
The clause c2 gives $C,x_a$ with
$\bigl|\sum_{w}\mathrm{sz}_w(E(u,j,w))-(\sum_ta_t)\log\prod_{g\in\mathcal G}g\bigr|\le
CB_0$ for every lift $j$, where $E_e(u)=\sum_j\sum_w\mathrm{sz}_w(E(u,j,w))$,
$a_t>0$ are the window weights and $\mathcal G$ is a finite set of primes.
The subtracted term is $\ge0$, so each inner sum is $\ge-CB_0$, and summing
over the $m$ lifts, $-E_e(u)\le CmB_0$. Since $m,f,L\le P$,
$CmB_0\le|C|P^4b$. Let $x\ge x_0=\max(x_a,256,2+|C|)$ and $Q=x^{9M_t+1}$.
By Proposition~\ref{i4:Lprop-D}(e) and Lemma~\ref{i4:Llem-bpol} (with
$(C,e)=(64,4)$ and $(256,9)$), $\mathrm{Ga}\le Qb$ and $\mathrm{Fo}\le Qb$;
and $|C|P^4b\le|C|x^{4M_t}b\le|C|Qb$. Adding,
$\mathrm{Ga}+\mathrm{Fo}-E_e(u)\le(2+|C|)Qb\le x^{9M_t+2}b$.
\end{proof}

\begin{lemma}[height bound from component bounds]\label{i4:Llem-hh}
Let $\mathcal F$ satisfy: $\max(E_b,E_c,C_{F0},E_d)+1\le C_H$; the slope
lemma of Proposition~\ref{i4:Lprop-D}(b); $\mathrm{MuRank}$; eventual
fullness of every $\nu_u$; the bounds of Proposition~\ref{i4:Lprop-D}(a),
(d), (e) and Lemma~\ref{i4:Llem-gafo}; and $\mathrm{TuSize}$. Then the
frame-number bound (hh) with constant $C_H$ holds: eventually
\[
s+\mu_{\max}(RR)\le x^{C_H}b,\qquad
s=R^{-1}\sum_ud_u\bigl(\mu(\nu_u)-h_6(u)\bigr),\quad R=\sum_ud_u .
\]
\end{lemma}
\begin{proof}
\emph{Step 1.} For every real $C$ there is $x_C$ such that for $x\ge x_C$
and all $i,u$,
\begin{equation}\label{i4:Leq-hhrest}
\mu_{\max}(\Lambda^0_u)+\mathrm{Ga}+\mathrm{Fo}-E_e(u)+CmB_0-h_6(u)+\mu_{\max}(RR)\le x^{C_H}b.
\end{equation}
Put $E=\max(E_b,E_c,C_{F0},E_d)$ and take $x$ beyond the thresholds of the
quoted bounds and $\ge\max(9,|C|,500+2|C_{T_u}|+30|C_{sz}|)$. Then each of
the following five quantities is at most $x^Eb$: $\mu_{\max}(\Lambda^0_u)$;
$\mathrm{Ga}+\mathrm{Fo}-E_e(u)$; $CmB_0\le|C|P^4b\le x^{4M_t+1}b$
(Lemma~\ref{i4:Llem-bpol}); $-h_6(u)-\mathrm{mean}\log\kappa_R\le x^{E_d}b$
(Proposition~\ref{i4:Lprop-D}(e), the hypothesis $\log T_u\le C_{T_u}mL$
being a clause of $\mathrm{TuSize}$, and Lemma~\ref{i4:Llem-bpol} with
$e=9$); and $\mu_{\max}(RR)+\mathrm{mean}\log\kappa_R\le x^{C_{F0}}b$. The
mean cancels in the sum, so the left side of \eqref{i4:Leq-hhrest} is at
most $5x^Eb\le x^{E+1}b\le x^{C_H}b$.

\emph{Step 2.} Let $C$ be the constant of the slope lemma and take $x$
beyond its threshold, those of $\mathrm{MuRank}$, of fullness and of
Step 1. By $\mathrm{MuRank}$, $d_u>0$, so $\nu_u\ne0$; it is finitely
generated, so $\mu(\nu_u)\le\mu_{\max}(\nu_u)$. With the slope lemma and
\eqref{i4:Leq-hhrest}, $\mu(\nu_u)-h_6(u)+\mu_{\max}(RR)\le x^{C_H}b$ for
every $u$. Multiply by $d_u\ge0$, sum, and divide by $R\ge d_{u_0}>0$.
\end{proof}

\begin{proof}[Proof of (L0)]
By Lemma~\ref{i4:lem-closed} it suffices to treat each clause and to
combine. The clauses adm2, cut, c2, c3 are
Proposition~\ref{i4:Lprop-sel}(a). The inputs of c5: the gap hypothesis is
Proposition~\ref{i4:Lprop-sel}(b) applied to the members of
$D_3(\mathcal A)$ satisfying $\mathrm{TuSize}$ (as in the remark
after that proposition), and $\mathrm{FinX}$ follows from it pointwise,
with the same threshold, by Proposition~\ref{i4:Lprop-sel}(c); the multiplicity condition of the selected indefinite
realisation is a conjunct of the row field; and the eventual fullness of
the base-vector image is the field ``base-vector images'' of (vii).
$\mathrm{BuRational}$ and $\mathrm{RRFull}$ are
Proposition~\ref{i4:Lprop-D}(c), which uses only the row field; the row
field is closed on $D_3(\mathcal A)$, so the implication is closed.

There remains (hh). The real-place ties, the row field, the closedness of
Proposition~\ref{i4:Lprop-D}(d), $(\mathrm{LG})$ (closed by
Proposition~\ref{i4:prop-links-proved}, so a fortiori under
$\mathrm{TuSize}\wedge\mathrm{DGamSize}$) and $\mathrm{TuSize}$ are jointly
closed (Lemma~\ref{i4:lem-closed}). Fix a member with these properties,
$\mathrm{DGamSize}$, c2 and $\mathrm{DIVinN}$, and let $C_{F0}$ come from
Proposition~\ref{i4:Lprop-D}(d). The ties with $\mathrm{TuSize}$ give the
level conditions, $\mathrm{NuFullF}$ and $\mathrm{MuRank}$
(Lemma~\ref{i4:Llem-lvl} below, together with the proviso of
$\mathrm{ChosenE}(3,C_H)$); with $\mathrm{TuSize}$ and
$\mathrm{DGamSize}$ they give $(\mathrm{LMu})$
(Proposition~\ref{i4:Lprop-D}(f)). Now Proposition~\ref{i4:Lprop-D}(b) gives
the slope lemma, Proposition~\ref{i4:Lprop-D}(a) the $E_b$-bound,
Lemma~\ref{i4:Llem-gafo} the $E_c$-bound and
Proposition~\ref{i4:Lprop-D}(d) the largeness and the sharp functional
bound; Lemma~\ref{i4:Llem-hh} gives (hh).
\end{proof}

\medskip\noindent\emph{(L1): the arithmetic clause.}

At a point of the wiring step we use the following data: the level primes
$g_0,g_1$ with $Y\le g_i\le2Y$, $(kb)^8\le Y$, $\log Y\le C_{YL}L$; the
host root discriminant $\mathrm{rd}_E\ge1$; $n_V=r+t$;
$L_a=\log(a\,\mathrm{rd}_E)$ and $L'=\log(b\,\mathrm{rd}_E\,g_0g_1n_V)$; the
group $G$ of the tied division datum, of order
$\prod_{k'\in\{0,1\}}|\mathrm{GL}_r(\Z/g_{k'})|^f$; the cloud budget
$N_c=|T_2|\,k\,b\,L^2/n^2\ge0$; and $U=L(B_0+N_c)$.

The \emph{arithmetic clause} at a word has six components: (1) the sum over
the lifts of the record's level-six terms is at most that of the scale;
(2) $\log|G|\le\lambda_0$; (3) the level-three bound, i.e.\ a sum identity
for the level-three entries and a bound for the error entries; (4) the
$Q$-clause $3\delta Q\le1$, where $\delta=8fr/\min(g_0,g_1)$ and
$Q=\max(1,B+x^{C_H}b+28\log|G|)$ with $B$ the boundary term of the core;
(5) the scale's group term is at most $\lambda_0$; (6) $c_0\ge28$. The
\emph{fixed scale} is the scale built from the record of the point: its
level-six entries are those of the record shifted by $C_{sz}fn_V^2L'$, its
level-three entries are
$\ell_3(j)=\sum_\rho q_j\,w_{j\rho}\log\Omega+C\ell_{3,u}$ with
$q_j=(k+j-1)-(m-j)$, $\ell_{3,u}=B_0\log\log x$, its error unit is $U$, its
group term is $\log|G|$, and its frame entries are the record's frame
values. We take $\lambda_0=\log|G|$.

\begin{lemma}[numerical facts]\label{i4:Llem-num}
At a point of the wiring step, with $x$ beyond the scale threshold:
\begin{enumerate}
\item[(a)] $P\le x^{M_t}$, $P^e\le X'\le b$ for $0\le e\le8$, $\log x\le L$;
\item[(b)] $0\le L_a\le4P$ and $C_{sz}L'\le c_0L_a$;
\item[(c)] $\log|G|\le2mU$, and $m\log|G|\le b$;
\item[(d)] the weights satisfy $0\le w_j(\tau)\le2b$ at every real place;
\item[(e)] $U\le2P^4b$ and $B_0\le P^3b$.
\end{enumerate}
\end{lemma}
\begin{proof}
(a) For $P\le x^{M_t}$, put $X=x^{M_t-1}$; each of $f,k,m,t,L$ and
$m(t+1)$ is at most $X$ by the size record, so $P\le8X\le x^{M_t}$. For
$y\ge1$, $\log y\le(2/e)y^{e/2}$; applied with $y=x$, $e=\varepsilon/3$ and
the scale condition $48M_t/\varepsilon\le x^{\varepsilon/6}$, it gives
$x^{8M_t}\le X'$, and $X'\le b$; similarly $\log x\le L$.
(b) $L_a=\log x+L+\log\mathrm{rd}_E\ge0$, and $\log\mathrm{rd}_E\le L$
because $\mathrm{rd}_E\le\frac{Q_S}2\sqrt{qd}\le2n\,x^{\varepsilon/100}\sqrt{qd}\le2x^3\sqrt{qd}$
(Proposition~\ref{i4:prop-rd}, the size-record bounds on $Q_S$ and $n$,
and $2+8\varepsilon+\varepsilon/50<3$) with $\log q,\log d\le
x^{\varepsilon/2}$; so $L_a\le3L\le4P$ after (a). Expanding,
$L'\le(L+\log\mathrm{rd}_E)+2\log2+2C_{YL}L+M_t\log x\le(3+2C_{YL}+M_t)L_a$,
and $C_{sz}(3+2C_{YL}+M_t)\le c_0$ by \eqref{i4:eq-c0ce}.
(c) Lemma~\ref{i4:lem-gzero} with $Y=g_0$ gives
$\log|G|\le2fr^2\log(2g_0)$; $\log(2g_0)\le2L^2$ and $2r^2\le b\le mb$, so
$\log|G|\le2L\cdot fmbL\le2mU$. Also $m\log|G|\le4mfr^2L^2\le4P^6\le P^7\le b$.
(d) At a top place $w_j\le3b'+\eta$, at a thick place $w_j\le b+\eta+1$,
at a thin place $w_j\le b'+\eta$, with $\eta\le b/2$ and $b'\le b/4$ (as
$4x\le P^2\le X'$); each is $\le2b$.
(e) $B_0=fmbL\le P^3b$; $N_c\le P^3b$ because $|T_2|\le n$, $k,L\le P$ and
$n\ge1$; hence $U\le2P^4b$.
\end{proof}

\begin{lemma}[the fixed scale]\label{i4:Llem-scale}
At a point of the wiring step, for every core over the fixed scale with
constant $C\ge0$:
\begin{enumerate}
\item[(a)] $\sum_j\ell_3(j)=M+mC\ell_{3,u}$, where
$M=\sum_{j,\rho}q_jw_{j\rho}\log\Omega\le16P^4b$;
\item[(b)] if $\mathrm{bd}(j)\le c_{13}U$ for every $j$, $\lambda_0\le
c_G\,mU$, $\log d_g\le c_DU$ and $2C+c_G+c_{13}+c_D+1\le c_e$, then
$mC\ell_{3,u}+\lambda_0+\sum_j\bigl(CU+\mathrm{bd}(j)+\log d_g\bigr)+1\le c_e\,mU$;
\item[(c)] $2\,\mathrm{Fo}(P)=\mathrm{Fo}$, where $\mathrm{Fo}(P)$ is the
frame number of the core;
\item[(d)] the boundary term satisfies $B\le M+2\mathrm{Fo}(P)+c_emU\le b^2$.
\end{enumerate}
\end{lemma}
\begin{proof}
(a) Reindex the lifts. For the bound, $|q_j|\le k+m\le2P$,
$0\le w_{j\rho}\le2b$ (Lemma~\ref{i4:Llem-num}(d)),
$|\log\Omega|\le L_a\le4P$ (the formula
$\log\Omega=\gamma+\log(2\pi/\mathrm{rd}_E)-\Re(L'/L)(1,\chi_d)$ with
$|\Re(L'/L)(1,\chi_d)|\le L$ from the size record), and there are at most
$f\le P$ places for each of the $m\le P$ lifts.
(b) $\log\log x\le\log x\le L$ gives $C\ell_{3,u}\le CB_0L\le CU$; and
$1\le mU$ since $m\ge2$, $L\ge1$, $B_0\ge1$. Sum the five bounds.
(c) Write $\mathrm{Fo}=\sum_j\sum_{w\in\mathrm{supp}}2\ell_we_wF(j,w)$. The
base side of the core is untwisted, so the frame entries of the core vanish
off the twisted places; the record's frame values vanish off the support;
the frame entries of the core are the record's frame values; and log weights
and ramification indices agree under the identification of places. So both
sums run over the same set and differ by the factor $2$.
(d) The first inequality is (a), (b) and the definition of $B$. By (a),
$M\le16P^4b$; by (c) and Proposition~\ref{i4:Lprop-D}(e),
$2\mathrm{Fo}(P)=\mathrm{Fo}\le256P^9b$; by Lemma~\ref{i4:Llem-num}(e),
$c_emU\le2c_eP^5b$. So $B\le KP^9b$ with $K=272+2c_e$. Finally $KP^9b\le
b^2$: for $x\ge\max(9,K)$ and $x\ge\lceil(9M_t+2)/\varepsilon\rceil!$,
Lemma~\ref{i4:Llem-bpol} gives $KP^9b\le x^{9M_t+1}b$, and
Lemma~\ref{i4:Llem-fact} gives $x^{9M_t+1}\le b$.
\end{proof}

\begin{proof}[Proof of (L1)]
Take $x$ beyond the scale threshold, which contains the lower bounds on
$x$ used in Lemmas~\ref{i4:Llem-num} and~\ref{i4:Llem-scale} and $x\ge
N!$ for $N=\lceil(\max(C_H,0)+1)/\varepsilon\rceil$. Components (1), (2),
(5) and (6): (1) is the identity of the level-six entries summed over the
lifts (the shift is the same on both sides); (2) and (5) hold with
$\lambda_0=\log|G|$; (6) is $c_0\ge28$, part of \eqref{i4:eq-c0ce}.

(3): the sum identity is Lemma~\ref{i4:Llem-scale}(a). For the error bound
apply Lemma~\ref{i4:Llem-scale}(b) with $c_{13}=C_e$ (the block bounds
from $\mathrm{DIVinN}$, which are $\le C_eU$), $c_G=2$ (by
Lemma~\ref{i4:Llem-num}(c), $\lambda_0\le2mU$) and $c_D=2+9\max(C_d,0)$ (by
Lemma~\ref{i4:lem-dgam}, $\log d_g\le c_DU$); the constant condition
$2C+2+C_e+c_D+1\le c_e$ is in \eqref{i4:eq-c0ce}.

(4): it suffices that $24frQ\le\min(g_0,g_1)$, since then $3\delta
Q=24frQ/\min(g_0,g_1)\le1$. Now $B\le b^2$ by
Lemma~\ref{i4:Llem-scale}(d); $x^{C_H}\le b\le b^2$ by
Lemma~\ref{i4:Llem-fact}; and $28\log|G|\le56mU\le112P^5b\le KP^9b\le b^2$
by Lemma~\ref{i4:Llem-num}(c), (e). So $Q\le3b^2$. Also $fr\le P^2$, so
$72fr\le KP^9b\le b^2$. Hence $24frQ\le72frb^2\le b^4\le(kb)^8\le Y\le
\min(g_0,g_1)$.
\end{proof}

\medskip\noindent\emph{(L2): the package lemma of the division bridge.}

\begin{lemma}[the $\eta$-sum at the twisted places]\label{i4:Llem-etatw}
Let $S_0$ be a base side with twist function $\mu$ and twisted places
$T_w$, and let $\eta_i(t,y)\ge0$ ($0\le i<m$) be the lift exponents of the
family. Assume the numerical tie $\eta_{i+1}(y)=\eta_i(\mu(y),y)$ for
$y\in T_w$, the twist record ($T_w\subseteq\mathrm{supp}$ and
$\mu_y\ne\mu_0$ on $T_w$), $(\mathrm{DivN})$ in the form
$\eta\Sigma(j)=\sum_{w\in\mathrm{supp},\,\mu_w\ne\mu_0}l_we_w\eta_j(w)\le
D_0C_eU$ for every lift $j$, and $D_0=d_Z$. Then
$\sum_{y\in T_w}l_ye_y\eta_i(\mu(y),y)\le C_eU\,d_Z$ for $0\le i<m$.
\end{lemma}
\begin{proof}
By the numerical tie each term is $l_ye_y\eta_{i+1}(y)$, the term of
$\eta\Sigma(i+1)$ at $y$; by the twist record the sum runs over a subset of
the index set of $\eta\Sigma(i+1)$, whose terms are $\ge0$.
\end{proof}

\begin{lemma}[the package face]\label{i4:Llem-face}
In the situation of Proposition~\ref{i4:Lprop-pkg}, take for the package
face the core package $P_k$, the word map, the identifications $e_{N,u}$
of coordinate spaces with division spaces, the lattices $N_C$ and the tied
column bridge of the tied datum, the records $R_u$ of the numerical datum,
$N_u=e_{N,u}^{-1}(\Lambda_u)$ ($\Lambda_u$ the division lattice of $u$),
and the division inputs of Proposition~\ref{i4:Lprop-pkg}. Then:
\begin{enumerate}
\item[(i)] $N_C(w_u)=I_u\cdot N_u$, $I_u=\prod_{i\in\mathrm{supp}}\mathfrak
p_i^{-\lfloor u_i(\sigma_u(i))\rfloor}$;
\item[(ii)] $h_u(R_u)-\mathrm{mean}\log\kappa=h(u)$;
\item[(iii)] $2F_o(P_k)=F_o$;
\item[(iv)] $\mathrm{rhs}(R_u)\le\mathrm{MAIN}(u)+C|T|aL_a+Cm(B_0+N_c)x^{2\varepsilon}$
for every $C\ge\max(0,50c_0+10,60c_e)$.
\end{enumerate}
\end{lemma}
\begin{proof}
(i) Apply the injective map $e_{N,u}$. The lattice tie gives
$e_{N,u}(N_C(w_u))=L_u$ with $L_u=I'_u\Lambda_u$, $I'_u$ the product over
the support; the shift ideal of the record shift is the product over all
places, but the shift vanishes off the support, so $I'_u=I_u$; and
$e_{N,u}(I_uN_u)=I_u\Lambda_u$.

(ii) The record ties identify the numerical datum with the one built from
the size data. The record height is
$h_u(R_u)=h_A+\sum_jh_{B,j}-\sum_jh_{X,j}$, and term by term: $h_A$ is the
size term $\mathrm{SZ}(u)$ (same $a_t$, $S_c$, $\eta$, $L_a$); $h_{B,j}$ is
the kernel place mean bound $K_j(u)$ (the kernel place terms agree place by
place, as the record's places correspond bijectively to the embeddings);
$h_{X,j}$ is the lift lower bound $\Lambda_j$, because for a partition with
parts $A_0\ge\dots\ge A_{p-1}\ge1$ the hook length formula in the form
\[
\prod_{\rm hooks}h\cdot\prod_{i<l<p}(A_i-A_l+l-i)=\prod_{i<p}(A_i+p-1-i)!
\]
makes the logarithm of the hook product equal to the size-data log-hook
term. Hence $h_u(R_u)=\mathrm{SZ}(u)+\sum_jK_j(u)-\sum_j\Lambda_j$, which by
the height formula is $h(u)+\mathrm{mean}_e\log\kappa_R(e)$, and
$\kappa_R(\sigma)=\kappa(v_\sigma)$.

(iii) is Lemma~\ref{i4:Llem-scale}(c).

(iv) By definition
$\mathrm{rhs}(R_u)=\mathrm{MAIN}(u)+(50c_0+10)|T|aL_a+60c_e\,m\,x^\varepsilon(B_0+N_c)$.
Use $|T|aL_a\ge0$, $m(B_0+N_c)\ge0$ and $x^\varepsilon\le x^{2\varepsilon}$.
\end{proof}

\begin{proof}[Proof of (L2)]
Let the division kernel datum be the one of the selection, with selected
rows, column datum, column face, scale and column tie, and let $P_k$ be
the base package with full Hecke data of the wiring step. The package has
the theta space, base data and column data of the selected datum (these
are fields of the realisation); so its base rows are the selected rows at
the same level subspace $Z_s$, its column face is that of the selected
datum (by the column tie), and the decomposition-point rows and the
$\kappa$-rows move to it.

(1) Transport the column tie along the division-rows hypothesis: this
gives a column $C'$, lattices $N'$ and a tied datum $I'$ with $C'=C$,
$I'=I$. (2) The base side of $P_k$ equals that of the datum (same rows at
the same $Z_s$), and the column faces and periods at equal $Z_s$ agree.
(3) Transport $(N',I')$ to a tied datum $(N'',I'')$ for $P_k$ along the
identifications of (2). (4) Under these identifications the tuples, the
evaluation maps and the word maps of $I''$ and $I$ coincide. Steps (1),
(3), (4) are identifications of equal data and involve no further argument.
(5) Proposition~\ref{i4:Lprop-pkg} at $(N'',I'')$ gives the tie of the
column bridge, the division inputs, the chain ties and the first columns;
the arithmetic clause it requires is (L1), and the $\eta$-sum input is
Lemma~\ref{i4:Llem-etatw}. Lemma~\ref{i4:Llem-face} gives the four
remaining conditions of the package face, and the three ties of the
package (tuple, evaluation, word) are (4), composed for the words with the
word tie of the datum. This is a division kernel package with constants
$(\max(0,50c_0+10,60c_e),2)$.
\end{proof}

\medskip\noindent\emph{(L3): the level-zero leaves, and what follows from them.}

Item (L3) is Proposition~\ref{i4:Lprop-forms}. We prove here the two
consequences that the proof of Proposition~\ref{i4:prop-links-formal}(b)
draws from it: the division side conditions and the congruence length
bound. Fix a word $u$ in the Hida setting with level prime $T=T_u$ and a
place $w$ over $T$; $\mathcal O_w$ is a discrete valuation ring with
uniformiser $\varpi$, $T=\epsilon\varpi^{e_w}$ with $\epsilon$ a unit, and
$\ell$ denotes $\mathcal O_w$-length.

\begin{lemma}[lengths under products]\label{i4:Llem-H1}
Let $A$ be an $\mathcal O_w$-module of finite rank and $g,h$ endomorphisms
with $g(A),h(A)\subseteq A$ of finite colength. Then
$\ell(A/gh(A))\le\ell(A/g(A))+\ell(A/h(A))$. Consequently, for
$G=g_1\cdots g_r$, $\ell(A/G(A))\le\sum_j\ell(A/g_j(A))$ and
$\ell(A/G^m(A))\le m\,\ell(A/G(A))$.
\end{lemma}
\begin{proof}
$0\to g(A)/gh(A)\to A/gh(A)\to A/g(A)\to0$ is exact, and $g$ induces a
surjection $A/h(A)\to g(A)/gh(A)$. Induct.
\end{proof}

\begin{lemma}[twisted inclusion bound]\label{i4:Llem-H2}
Let $S$ be free of finite rank, $B$ a submodule, and $U,D$ endomorphisms
with $\varpi^bS\subseteq B$, $U(B)\subseteq S$, $D(S)\subseteq B$ and
$UD=\varpi^e$ on $S$. Then
$\ell(S/\varpi^bU(S))\le\ell(B/\varpi^bS)+e\,\mathrm{rk}\,S$. The same holds
with $S,B$ replaced by $S\cap P,B\cap P$ for a saturated submodule $P$
stable under $U$ and $D$, if $UD=\varpi^e$ only on $S\cap P$ (with the
right side unchanged).
\end{lemma}
\begin{proof}
Put $T_1=U(\varpi^bS)\subseteq T_2=U(B)\subseteq S$. The map $y\mapsto Uy$
is a surjection $B\to T_2/T_1$ vanishing on $\varpi^bS$, so
$\ell(T_2/T_1)\le\ell(B/\varpi^bS)$; and $\varpi^eS=UD(S)\subseteq T_2$, so
$\ell(S/T_2)\le\ell(S/\varpi^eS)=e\,\mathrm{rk}\,S$. For the second form:
$S\cap P$ is free of rank $\le\mathrm{rk}\,S$ (a submodule of a free module
over a principal ideal domain), and inclusion induces an injection
$(B\cap P)/\varpi^b(S\cap P)\to B/\varpi^bS$, since $x=\varpi^bs\in P$ with
$s\in S$ forces $s\in P$ by saturation.
\end{proof}

\begin{lemma}[congruence length bound]\label{i4:Llem-cong}
Assume the four conclusions of Proposition~\ref{i4:Lprop-forms} at $u$.
Let $S_{\rm ram}=S_{0,w}\cap P$ ($P$ the ramified subspace), and let
$N_\lambda(w)=S_{\rm ram}$ if $w$ is not seen by the form, and
$N_\lambda(w)=G^m(S_{\rm ram})$ with $G=\prod_{v\in VV_w}T^{n_b(v)}U_v$
otherwise. Then $N_\lambda(w)\subseteq S_{\rm ram}$ and
\[
\ell\bigl(S_{\rm ram}/N_\lambda(w)\bigr)\le m\sum_{v\in VV_w}e_w\,d^S_u\Bigl(\frac{a_{\max}}k+2\,\mathrm{Wt}_u(\tau)+k+b_{\rm top}(\tau)\Bigr),\qquad\tau=\tau_w(v).
\]
\end{lemma}
\begin{proof}
If $w$ is not seen, $VV_w=\emptyset$ and the length is $0$. Otherwise fix
$v$, write $S=S_{0,w}$ (free of rank $\le d^S_u$), $B=B_{0,w,v}$,
$n_b=b_{\rm top}(\tau)\ge0$ (Hida setting) and $g_v=T^{n_b}U_v$.
\emph{Stability:} for $x\in S_{\rm ram}$, $T^{n_b}x\in B$ by (2), so
$g_vx\in U_v(B)\subseteq S$ by (3), and $g_vx\in P$ by (4).
\emph{Colength:} Lemma~\ref{i4:Llem-H2} (second form) with $b=e_wn_b$,
$e=e_w(k-1)$, $U=U_v$, $D=\epsilon^{-(k-1)}D_v$ (multiplying by a unit does
not change images) and (2) gives
$\ell(S_{\rm ram}/g_v(S_{\rm ram}))\le n_v+e_w(k-1)d^S_u$ with
$n_v=e_wd^S_u(\mathrm{wsum}_\tau+kb_{\rm top}(\tau))/k$. By
Lemma~\ref{i4:Llem-H1}, $\ell(S_{\rm ram}/N_\lambda(w))\le
m\sum_v(n_v+e_w(k-1)d^S_u)$. At a top, $\lambda_{0,0}=a_{\rm top}\le
a_{\max}$, $|\lambda_{0,a}|\le\mathrm{Wt}_u(\tau)$ for $a\ne0$ and
$b_{\rm top}=-\lambda_{0,k-1}\le\mathrm{Wt}_u(\tau)$; so
$(\mathrm{wsum}_\tau+kb_{\rm top})/k\le a_{\max}/k+2\mathrm{Wt}_u(\tau)$,
and $(k-1)\le k$ gives the bound.
\end{proof}

\begin{lemma}[level conditions, fullness, rank]\label{i4:Llem-lvl}
Let $\mathcal F$ satisfy the real-place ties, $\mathrm{TuSize}$, the proviso
of $\mathrm{ChosenE}(3,C_H)$ and the conclusions of
Proposition~\ref{i4:Lprop-forms} beyond a threshold, at every word in the
Hida setting, and let the Hida setting hold eventually
(Proposition~\ref{i4:Lprop-fam}(a)). Then eventually, at every word:
\begin{enumerate}
\item[(a)] $d_u=d^S_u(1-1/(T-1))^f$ and the congruence length bound
(the \emph{level conditions});
\item[(b)] $\nu_u$ is commensurable with $M_{\rm ram}=\iota_u^{-1}(L^0_u)$,
and the division lattice spans (fullness, $\mathrm{NuFullF}$);
\item[(c)] $d_u>0$ ($\mathrm{MuRank}$).
\end{enumerate}
\end{lemma}
\begin{proof}
(a) The free $T$-part (a field of the concrete forms) gives the integer
identity $d_u(T-1)^f=d^S_u(T-2)^f$, whence the first condition as $T\ge2$;
the second is Lemma~\ref{i4:Llem-cong}. The threshold is the maximum of
those of the setting, of the ties and of the leaves, and $x_1+1$.

(b) Write $X\le_wY$ if every $y\in X$ has $sy\in Y$ for some $s\notin w$.
At $w$ over $T$, the $T$-definition of $\nu_u$ gives $\iota_u(x)\in
N_\lambda(w)$ for $x\in\nu_u$, and $N_\lambda(w)\subseteq S_{\rm ram}(w)$ by
Lemma~\ref{i4:Llem-cong}; by the local spanning of $L^0_u$ some $s\notin w$
has $sx\in M_{\rm ram}$. At the other places $\nu_u\le_wM_{\rm ram}$ is part
of $N'$-definedness. A vector lying locally in a submodule $M$ at every
place lies in $M$ (the ideal $\{a:ax\in M\}$ lies in no maximal ideal), so
$\nu_u\subseteq M_{\rm ram}$, and $\nu_u$ is finitely generated as
$M_{\rm ram}$ is noetherian. Choose a prime $p\ne T$ that is not inert
($p=q$ if $T=2$, $p=2$ otherwise) and a place over it; there
$N'$-definedness gives $M_{\rm ram}\le_w\nu_u$, and taking the product of
the multipliers over a finite generating set gives $c\ne0$ with
$cM_{\rm ram}\subseteq\nu_u$. The division lattice is
$\mathrm{span}(I_u\nu_u)$ with $I_u$ a nonzero fractional ideal; it is
finitely generated, contains $a_0cM_{\rm ram}$ for any integral $0\ne
a_0\in I_u$, and $M_{\rm ram}$ spans because $L^0_u$ does. So it spans.

(c) Fullness gives $\mathrm{rank}\,\nu_u=d_u$; the proviso with $C_0=3>0$
gives $d^S_u>0$; the setting gives $T>2$; so
$d_u=d^S_u(1-1/(T-1))^f>0$.
\end{proof}

\medskip\noindent\emph{(L4): the family-level steps.}

\begin{lemma}[loss requirements]\label{i4:Llem-loss}
Let $\mathrm{ChosenE}(C_0,C_H)$ hold with constants $C_5,c_6,C_b$,
$C_b>0$, and assume the chosen wall. Then eventually the small-shares
record holds: thick density $x^{-\varepsilon}\le|T_2|/n$, the bad-prime
bound $\mathrm{ChBadp}$, and the shares; the remaining clauses are copied
from $\mathrm{ChosenE}$.
\end{lemma}
\begin{proof}
\emph{Thick density.} The thick-set clause gives, beyond a threshold,
$n_0e^{-C_{T2}S}/L_w\le|T_2|$ with $S=\sqrt{\log n_0\log L_w}$, $n=o\,n_0$,
$L_w\ge2$. Put $C_T=\max(C_{T2},0)+1$ and
$\kappa=\min(\varepsilon/2,\varepsilon^2/(12C_T^2))$. For $x\ge
\exp((4/\kappa)^2)$ one has $2\log\log x\le\kappa\log x$ (put $y=\log x$
and use $\log y\le2\sqrt y$). The size bounds give
$0\le\log n_0\le\frac94\log x$ and $0\le\log L_w\le2\log\log x\le\kappa\log
x$; hence $C_T^2S^2\le\frac3{16}\varepsilon^2(\log x)^2$, so
$C_{T2}S\le\frac9{20}\varepsilon\log x$. With $\log L_w\le\frac\varepsilon2\log
x$ this gives $x^{-\varepsilon}x^{\varepsilon/100}\le e^{-C_{T2}S}/L_w$, and
with $o\le x^{\varepsilon/100}$, $x^{-\varepsilon}n\le|T_2|$.

\emph{Bad primes.} The tops clause gives $\delta_1n\le2|\mathrm{TT}|$; the
old bad-prime clause gives
$\sum_{p\in\mathcal P}\frac{\log p}p|\mathrm{bad}(p)|\le4C_bx^{-2}\delta_1n\sum_{p\in\mathcal
P}\frac{\log p}p$, where every $p\in\mathcal P$ is below $\lceil X'\rceil$.
For $y\ge1$, $\sum_{p<\lceil y\rceil}\log p/p\le2\log y$: for $y\ge4$ by
Lemma~\ref{i1:lem-cheb}, and for $y<4$ by inspection. With $y=X'$ the right
side is at most $16C_bx^{\varepsilon/3}|\mathrm{TT}|x^{-2}\le
|\mathrm{TT}|x^{-2+\varepsilon}$ once $16C_b\le x^{2\varepsilon/3}$.

\emph{Shares.} Beyond the wall threshold, take an index $i_0$ (the window
set lies in the covered range); the point $(x,i_0)$ carries its top weights
in the admissible box, and the wall holds at every index. This is the
shares clause with witness $i_0$.
\end{proof}

\begin{lemma}[from the losses to the loss bounds]\label{i4:Llem-L8}
If $\mathrm{LossFS}(\mathcal F,0)$ holds with constant $C=3K+13$, then for
the model data built from the family at a large scale, with the chosen
weight vector, $I_O(u_0)\le C\,n_T\,m\,b$ and
$\mathrm{ASK}(u_0)-\mathrm{Ga}(u_0)\le C\,n_T\,m\,b$ for every $u_0$.
\end{lemma}
\begin{proof}
For each $u_0$ apply $\mathrm{LossFS}$ at the underlying point
$\mathrm{pt}(u_0)$, which carries the chosen weights; the model's top count,
number of lifts and scale parameters are those of the point, and
$x^{0\cdot\varepsilon}=1$.
\end{proof}

\begin{proof}[Proof of (L4)]
Let $\mathcal F$ satisfy $\mathrm{TuSize}$, the minimal global record and
the hypotheses of $(\mathrm{LGl})$. In order:
(1) the Hida setting holds eventually (Proposition~\ref{i4:Lprop-fam}(a));
(2) the level conditions, commensurability, $\mathrm{NuFullF}$ and
$\mathrm{MuRank}$ hold eventually (Lemma~\ref{i4:Llem-lvl}, the leaves
being Proposition~\ref{i4:Lprop-forms}); (3) the small-shares record holds
eventually (Lemma~\ref{i4:Llem-loss}); (4) the losses
$\mathrm{LossFS}(\cdot,0)$ follow from Proposition~\ref{i4:hyp-loss}(b),
applied with the tops and bad-prime clauses of (3) and $(\mathrm{LG})$,
together with the wall clause, and Lemma~\ref{i4:Llem-L8} turns them into
the loss bounds of the model data; (5) the family-level items are
Proposition~\ref{i4:Lprop-fam}(b); (6) the core link is
Proposition~\ref{i4:Lprop-fam}(c), with $e_D=C_{\rm abs}$ and $e_L=0$.
\end{proof}

\begin{proof}[Proof of Proposition~\ref{i4:prop-links-leaves}]
(L0) is proved above, modulo Propositions~\ref{i4:Lprop-sel},
\ref{i4:Lprop-D}, \ref{i4:Lprop-forms} and~\ref{i4:Lprop-fam}(a) (the last
two through Lemma~\ref{i4:Llem-lvl}); (L1) is proved above, modulo
Proposition~\ref{i4:Lprop-D}(e); (L2) modulo
Proposition~\ref{i4:Lprop-pkg}; (L3) is Proposition~\ref{i4:Lprop-forms};
(L4) is proved above, modulo Propositions~\ref{i4:Lprop-forms}
and~\ref{i4:Lprop-fam}.
\end{proof}

\begin{remark}\label{i4:Lrem-trust}
On the division and global side, the statements not proved here are
exactly the formally verified Propositions~\ref{i4:Lprop-sel}--\ref{i4:Lprop-fam}.
Of these, Proposition~\ref{i4:Lprop-forms} coincides with item
(L3); the others are proper parts of (L0), (L2) and (L4), and
(L1) depends on no formally verified statement beyond the
pointwise frame estimate $\mathrm{Fo}\le256P^9b$ of
Proposition~\ref{i4:Lprop-D}(e).
\end{remark}

\subsubsection{The selection clauses (b), (c) and the size polynomial bound}\label{i4:ssec-proof-selD}

\begin{proof}[Proof of Proposition~\ref{i4:Lprop-sel}(b)]
By (G2) $\mathcal C$ is closed for $\mathrm{ResidEv}$; (G1) gives $\varepsilon_1$
and, for $\varepsilon\le\varepsilon_1$, a threshold datum $\kappa$. Let
$\mathcal F\in\mathcal C$ of exponent $\varepsilon$ be $\kappa$-large at
$C_H$. By hypothesis $\mathcal F$ satisfies $\mathrm{TuSize}$, so (G1)
gives $\mathrm{GapEv}(\mathcal F)$. This is closedness for
``$\mathcal F\in\mathcal C\Rightarrow\mathrm{GapEv}$'', hence for $\mathrm{GapEv}$ on $\mathcal C$.
\end{proof}

\begin{proof}[Proof of Proposition~\ref{i4:Lprop-sel}(c)]
$\mathrm{RowsFinDim}$ holds by Proposition~\ref{i4:A-prop-inputs-leaves}(M5).
At a datum with the gap hypothesis, (G3) gives Wu decay and integrability
for the Hida setting of the tie, and (G4) gives $\mathrm{FinX}(T)$.
\end{proof}

\begin{proof}[Proof of the first claim of Proposition~\ref{i4:Lprop-D}(e)]
Let $c=\max(C_t,5)$ and $X=x^{c+2}$; $x\ge9$. The size conditions give
$n\le x^{2+8\varepsilon+\varepsilon/100}$, $k\le x^{9/2+10\varepsilon+\varepsilon/100}$,
$m\le x^{1/2+\varepsilon/100}$ (as $\sqrt{x\delta_1}\le x^{1/2}$), and the
$t$-window of the size datum gives $t\le x^{C_t}\le x^c$.
Hence $f=2n\le x\cdot x^{2+8\varepsilon+\varepsilon/100}\le X$, $k,m,t\le X$, $L=x^\varepsilon\le X$,
$1\le X$, and $m(t+1)\le x^{1/2+\varepsilon/100}\cdot x^{1+c}\le X$, so
$N_m=k+m(t+1)\le2X$. Thus $P\le8X\le x\cdot X=x^{c+3}=x^{M_t}$.
\end{proof}

\subsection{Realisation}\label{i4:ssec-grp-real}

\subsubsection{Proof of realisation}\label{i4:ssec-proof-realise}

We use the vocabulary of the retwisting conditions. A \emph{base} $B$
consists of the following data:
\begin{itemize}
\item a pre-datum $\bar B$ (host, ranks and prescribed data);
\item a coefficient field $L(B)$ with an embedding $L(B)\to\C$, whose image
is $L_\C(B)$;
\item fields $K_s(B)\subset F_0(B)\subset L_\C(B)$.
\end{itemize}
$B'$ \emph{enlarges} $B$ if $B'$ has the same pre-datum, the same $K_s$ and
the same $F_0$, and $L_\C(B)\subset L_\C(B')$.

A \emph{twist datum} $T$ has a finite set $\mathrm{pl}(T)$ of primes and
twisting places. It is \emph{untwisted} if it has no twisting places. It
is \emph{forced} if $\mathrm{pl}(T)$ contains $g_0$, $g_1$, $d$ and every
prime $\le X''_n$.

A \emph{core} over $(B,f,T)$ is the collection of column, row, division,
spectral and norm stages over $B$, with assembled column datum $D(c)$.

The four retwisting conditions read as follows.
\begin{itemize}
\item $\mathrm{BaseCond}(KP)$: bases with $KP$ exist with prescribed finite
algebraic sets in $F_0$ and in $L_\C$, and $KP$ is stable under
enlargement.
\item $\mathrm{FirstG}(KP,Q)$: generically, over a forced untwisted $T$,
there is a core satisfying $Q$.
\item $\mathrm{ExportF}(KP,Q)$: generically, such cores satisfy the export
condition.
\item $\mathrm{SecondG}(KP,Q,Q')$: generically, such a core with the export
condition can be retwisted. That is, there are an enlargement $B'$, a twist
datum $T'$ with the same primes, avoiding any prescribed set, and a core
over $(B',f,T')$ satisfying $Q$, the twist condition at all places and
$Q'$.
\end{itemize}
Here ``generically'' means: for all data outside finitely many exceptional
algebraic numbers and primes, chosen as functions of $(\bar B,f,K_s,F_0)$.

\emph{Step 1: two passes.} Assume that volumes equal masses (step 6). Then,
eventually in contexts, there are a column datum $D$ and a layer $Y$ over
$D$ with the record, ledger and twist conditions, coming from a core
satisfying $Q'$, with all tuples used and with $\mathcal O_{LL}$-stable
lattices $\tilde M_\mu$.

Indeed, forcedness holds generically with the exceptional set of primes
$\{g_0,g_1,d\}\cup\{p\le X''_n\}$. Intersect the four generic statements:
take the unions of their exceptional sets. Then proceed as follows.
\begin{itemize}
\item $\mathrm{BaseCond}$ gives a base $B$ with $KP$ that contains the
exceptional algebraic numbers.
\item The untwisted datum on the exceptional primes is generic and forced.
\item $\mathrm{FirstG}$ gives a core $c$ with $Q$.
\item $\mathrm{ExportF}$ gives the export condition for $c$.
\item $\mathrm{SecondG}$ gives a retwisted core $c'$ with $Q'$ and the
twist condition at all places.
\end{itemize}
The layer over $D(c')$ is read off from $c'$. Its two size fields are the
size formula
\[
d^*_u=\mathrm{mass}(\mathrm{std}_k)(\iota^{\rm col}_{g_0}\iota^{\rm col}_{g_1})^f\prod_\tau\dim_k(u_\tau)
\]
and the target $D^*$, in which the volume of the indefinite space of
signature $(P,Q)$ is replaced by $\mathrm{mass}(\mathrm{std}_r)$. This is
where volumes equal masses is used. The record conditions are the following.
\begin{itemize}
\item The untwisted part vanishes where $\mu=\mu_0$.
\item The set $\{\mu\ne\mu_0\}$ is finite.
\item Fullness holds because the level-zero lattice spans.
\item The tame condition holds by the following computation. Let $w$
generate $\mathbb F_g^\times$ and $\chi_j=w^{2(j+1)}$. If
$4mg^{1/4}<g-1$, then the $2m+1$ characters $1,\chi_j^{\pm1}$ are distinct,
and remain distinct in prime-to-$\ell$ parts for $\ell$ outside an explicit
finite set. The inequality follows from $g\ge(kb)^8$ and $k\ge m+3$.
\end{itemize}

\emph{Step 2: the Harish-Chandra parameters.} Let $\Delta=\lceil x^{c_D}\rceil$
and $h=\lceil x^{c_h}\rceil$ with $c_h\ge402$, and let $w$ be a binary word
of length $k-2$. Put
\[
v_i=(k-2-i)4\Delta,\qquad
u_{l+1}=v_l-\tfrac12-[w_l=1]\Delta,\qquad
b_j=b_1+2jh,\qquad e_j=b_j+h .
\]
Take $u_0=a_t+(k-1)/2$ and $u_{k-1}=-b_t-(k-1)/2$ at the top place of $t$,
and $u_0=(k-2)4\Delta+4k\Delta+\frac12$, $u_{k-1}=-\frac12$ elsewhere. The
base value $b_1(\tau)$ is
\begin{itemize}
\item $b_t+\lfloor b'/x\rfloor$ at the top place of $t$,
\item $\lceil b\rceil$ at a thick place,
\item $\lceil w''_n\rceil$ at a thin place,
\end{itemize}
corrected at one thick place $\tau_0$ by $0\le s<w_E$. The correction makes
$\prod_\tau\varphi_\tau(z)^{b_1(\tau)}=1$ on the roots of unity $z\in E$.
This is possible because $\varphi_{\tau_0}$ generates the character group
of the cyclic group $\mu(E)$.

These are Harish-Chandra parameters: $u_i\in\frac12+\Z$; the interlacing
$v_i<u_i$, $u_{i+1}<v_i$ holds, at the top place because $4k\Delta\le a$
and $k\le b'/x$ for large $x$; and $e_j-b_j=b_{j+1}-e_j=h\ge(t+3)/2$,
because $c_h\ge402$ and $t\le x^{401}$. The condition on
the $e_j$ holds because $(e_j(\tau))_\tau$ differs from $(b_1(\tau))_\tau$
by a constant and by the following lemma.

\begin{lemma}[roots of unity have norm one]\label{i4:lem-normmu}
Let $E$ be a host field and $K=\Q(\sqrt{-d})$. Then $\mathrm N_{E/K}(z)=1$
for every root of unity $z\in E$.
\end{lemma}
\begin{proof}
Let $z$ have order $m$.
\begin{itemize}
\item $d\nmid m$. Otherwise $E$ contains a primitive $d$-th root of unity,
and $d$ would have ramification index at least $d-1$ in $E$. But the
inertia groups above $d$ have order at most $2$ (their labels are $1$ and
$((1,-1),1)$), while $d\ge7$.
\item The element $\iota$ with label $((1,-1),1)$ is inertial above $d$.
So $\iota(z)\equiv z$ modulo a prime above $d$. If $\iota(z)\ne z$, then
$m\iota(z)^{m-1}\equiv0$ modulo that prime, so $d\mid m$, which was just
excluded. Hence $\iota(z)=z$.
\item Put $s=c\iota$, where $c$ is complex conjugation. Its label is
$((-1,1),1)$, so $s$ fixes $K$ and $s(z)=c(z)=z^{-1}$. Pairing $\sigma$
with $\sigma s$ in $\mathrm{Gal}(E/K)$ gives $\mathrm N_{E/K}(z)=1$.\qedhere
\end{itemize}
\end{proof}

\emph{Step 3: level primes.} Let $c_s>0$ and let $T_s>0$ satisfy
$c_smL\le\log T_s$. We claim there is a prime $T\in[T_s,(1+1/(fk))T_s]$
with $T=\pi\bar\pi$, where $\pi\equiv1\pmod{N(g_0g_1)^2}$ is integral in
$K$.

Apply Proposition~\ref{i4:thm-shortnorm} with $X=T_s$, $h=1/(fk)$ and
$M=N(g_0g_1)^2$. Then $\log(4dM^2)\le80x^\varepsilon$, since
$\log g_j\le9x^\varepsilon$. On the other hand
$c\sqrt{\log X}\ge c\sqrt{c_s}\,x^{1/6+\varepsilon/2}$, since
$m\ge x^{1/3}$, and $h\ge x^{-12}$. For large $x$ the hypotheses hold.

Then $(T-1)/N=\alpha+\bar\alpha+N\alpha\bar\alpha$ with $\pi=1+N\alpha$, so
$T\equiv1\pmod N$. Moreover $T\ge X'$, so $T\ne p_{de},p'_{de}$. With
$T^\flat_u=(D_*/d_u)^{1/(fk)}\Gamma_S^{-2/(fk)}$, this gives a \emph{level
choice} for every $(d_u,D_*)$ satisfying the window inequality.

\emph{Step 4: the auxiliary prime and hosts in windows.} Let $r(H)$ be the
least prime not dividing $N_H(p_{de}-1)(p'_{de}-1)$, and let $q_0(H)$ be
the least prime $\equiv1\pmod{r(H)N_H}$.
\begin{itemize}
\item Bertrand's postulate gives $r(H)N_H\le N_H^{3+2C_8}$.
\item Theorem~\ref{i4:thm-linnik} then gives $q_0(H)\le N_H^{(3+2C_8)A}$.
So $\mathrm{HostExW}$ with constant $C_8$ yields $\mathrm{HostExW}$ with
the constant $(3+2C_8)A$ and auxiliary prime $q_0(H)$. The new level prime
$q_0(H)$ splits in $K$ and differs from $p_{de},p'_{de}$.
\item Since $\log N_H\le2x^{\varepsilon/5}(\log x)^2$, we get
$q_0(H)<X'\le b\le g_i$, by Lemma~\ref{i4:lem-kits-shared}(b).
\end{itemize}

Object existence follows. At large scales, for every good host with
window weights, every injective code of length
$\lceil3\exp(C_\Gamma m^5fL)^2x^{C_F}\rceil$ and every base word, there is
a valid point with $QO(\kappa)$. It is assembled as follows.
\begin{itemize}
\item The choice data come from step 2.
\item The degrees, the target and the window inequality come from step 5.
\item The level choice comes from step 3.
\item The layered data come from step 1, at the context defined by the
level choice.
\item The size data are explicit. Their window, depth and cloud clauses
follow from the shape of step 2 and from Lemma~\ref{i4:lem-kits-shared}
by a direct comparison of powers of $x$.
\item The ledger point is produced as follows. Enlarge the support to all
places of residue characteristic at most
$B=\lceil X''_n\rceil+g_0+g_1+d+\sum_u(T_u+D_\Gamma(u))$, and to all places
where $\mu\ne\mu_0$. Then glue the local lattices prescribed at finitely
many places into one $\mathcal O$-lattice.
\end{itemize}
For the gluing: if $M$ is a lattice and $A_w$, $w\in S$, are $w$-saturated
modules with $n_wM\subset A_w$ and $n_wA_w\subset M_{(w)}$, then
\[
N=\bigcap_{w\in S}A_w\cap\bigcap_{w\notin S}M_{(w)}
\]
has $N_{(w)}=A_w$ for $w\in S$ and $N_{(w)}=M_{(w)}$ for $w\notin S$. The
proof clears one denominator away from one place in a Dedekind domain.

Combining object existence, $\mathrm{HostExW}$ and (R4) gives kits at all
scales of the system. Each kit is a host together with one point for every
base word, all having $QO$, the wall property and $\mathcal G$.

\emph{Step 5: degrees and target.} Fix a level prime $g$ with
$\mathbb F=\Z/g$. Then
\[
\iota^{\rm col}_g=\frac{|\mathrm{GL}_k(\mathbb F)|}{|N_k(\mathbb F)|},\qquad
\iota^{\rm row}_g=\frac{|\mathrm{GL}_r(\mathbb F)|}{|P(\mathbb F)|},
\]
with $|N_k(\mathbb F)|=g^{k(k-1)/2}$, $|P(\mathbb F)|=g^{r(r-1)/2-k}$ and
$|\mathrm{GL}_n(\mathbb F)|=g^{n^2}\theta_n$, where
$\theta_n=\prod_{i\le n}(1-g^{-i})\in[\frac12,1]$. Hence
\[
\frac{I_\sigma}{I_\pi}\ge2^{-2f}(g_0g_1)^{fm(2k+m-1)/2}.
\]
By the Siegel mass ratio the masses telescope, and $|\log L_E(M)|\le f\log\zeta(M)$.

(R2) makes the column types place data, and (R1) gives
\[
\frac1{fk}\log\frac{D_*}{d_u}\ge m(\log g_0+\log g_1)-C_1\Big(L+\frac{m^2L}k\Big)\ge15mL,
\]
using $\log g_i\ge8L$ and, for large $x$, $C_1(L+m^2L/k)\le mL$. Since $\frac2{fk}\log\Gamma_S=2C_\Gamma m^5L/k\le mL$,
this is the window inequality with $c_s=1$ and $D_*=D_S(P)$.

The degrees are positive integers, for three reasons.
\begin{itemize}
\item Reduction modulo a prime above $g$ is injective on
$\mathrm{Aut}(M)$. This is Minkowski's lemma: $(1+g^jB)^\ell=1$ forces
$B\equiv0$, using $g\ge3$.
\item Every prime factor of $|\mathrm{Aut}(M)|$ is at most $2fk+1<g$.
Hence $|\mathrm{Aut}(M)|$ divides $\iota^{\rm col}_g$, and
$\mathrm{mass}(\mathrm{std}_k)\iota^{\rm col}_g$ is an integer.
\item $\dim_k(u)=\pm\det\binom{u_i-1/2}{j}$ is a positive integer.
\end{itemize}

\emph{Step 6: volumes equal masses.} Let $W$ be definite, and let $\mu$,
$\mathcal F$ be a Haar measure and a fundamental domain from the tied
volumes, with $\mu(K)=1$. Write $G=\bigsqcup\Gamma g_iK$, as in (R5). Let
$N_i(x)$ be the number of $\gamma$ with $\gamma x\in g_iK$. Then
$N_i=|\Gamma_i|$ on $\Gamma g_iK$, and $N_i=0$ elsewhere. Unfolding gives
$\int_{\mathcal F}N_i\,d\mu=\mu(g_iK)=1$. Summing over $i$,
$\mathrm{vol}(W)=\sum_i|\Gamma_i|^{-1}=\mathrm{mass}(L_0(W))$, and this
mass is $\mathrm{massI}(D,\mathrm{rank}\,W)$ by
Theorem~\ref{i4:thm-genusmass}.

For signature $(P,Q)$ with $P\ge2$, induct on $Q$. The volume ratio with
sign $-1$ multiplies the volume by $\rho(E,M)$. The volume ratio with sign
$+1$, applied to definite spaces, shows that $\mathrm{massI}$ satisfies the
same recursion.

\begin{proof}[Proof of Theorem~\ref{i4:thm-realise}]
We first glue kits into families. Choose a standard record $C$ that is
large for the given threshold data and has the prescribed $C_{T_2}$ and
$C_{pde}$. Its last entries are chosen after its middle entries, which
are chosen after its outer entries. Then take $c_D,c_h$ beyond the
thresholds of object existence and of Proposition~\ref{i4:prop-kits}.

For every scale $s$ of the system with $s.x\ge x_0$, choose a kit $k(s)$.
Define a family as follows. Its index set at $x$ is that of $k(s(x))$ if
$x\in\mathrm{Win}$ and $x>x_1$; its points are the kit points; its
constants are those of $C$, with the window index replaced by $C_{pde}$.
The uniformity requirements of a family at fixed $x$ are the kit axioms:
the same frame, ranks, primes and level data, compatible word sets, and
every base word covered. Every pointwise property of the kit points passes
to the family. Proposition~\ref{i4:prop-kits}, with (R3), supplies the
numerology at every point, and hence $C_{YL}=9$ and the window property.

The family is linked to $\kappa=(\varepsilon,C,c_D,c_h,C_5,c_6,C_b)$, has
$QO(\kappa)$ at all points and lies in $\mathcal G.\mathrm{cls}(\lambda)$.
Pointwise, it has:
\begin{itemize}
\item the twist condition;
\item the $(q,d)$-bound;
\item the box and $(q,d)$-row conditions, from the goodness of the host;
\item $D_\Gamma\mid M(2fN_m)$, by Theorem~\ref{i4:thm-mink};
\item distinct tame data;
\item the thick-set, distinguished-set, bad-prime and code bounds;
\item $C_F=C_H$ and $C_H\ge0$, by the choice of the record;
\item the word-count bound, by group 10 of (E3).
\end{itemize}
Setting $C_{T_u}=c_T(\varepsilon,9)$ gives the design choices.
\end{proof}

\subsubsection{The size leaves of realisation}\label{i4:ssec-proof-Rleaves}

For a sequence $y=(y_i)_{i<N}$ write
$V_N(y)=\prod_{i<j<N}(y_i-y_j)$, so that $\dim_N(y)=V_N(y)/\prod_{i<N}i!$.

\medskip\noindent\emph{(R1): the Vandermonde ratio.}

\begin{lemma}[Vandermonde ratio at one place]\label{i4:lem-vandplace}
Let $x\ge1$, $\varepsilon>0$, $B=e^{x^\varepsilon}$, $a=xB$, and let
$(u,v,b,w)$ be place data (Definition~\ref{i4:def-placedata}). Then
$V_k(u)>0$, $V_{k-1+m}(\mathrm{row}(v,b))>0$ and
\[
(k-1)\sum_{j<m}\log b_j-2k\log(8a)-2k(1+\log k)\ \le\
\log V_{k-1+m}(\mathrm{row}(v,b))-\log V_k(u).
\]
\end{lemma}
\begin{proof}
By (P2) and (P3), $u_0>v_0>u_1>v_1>\dots>v_{k-2}>u_{k-1}$, so $V_k(u)>0$.
By (P6)--(P8), $b_0\ge4k\Delta x\ge8$, $b$ is strictly increasing with
$b_j-b_i\ge j-i$, and $b_j<w_j\le w_{m-1}\le2B$. The row sequence is
$v_0>\dots>v_{k-2}\ge0>-b_0>\dots>-b_{m-1}$, so its Vandermonde product is
positive and equals
\[
V_{k-1}(v)\prod_{i<k-1,\,j<m}(v_i+b_j)\prod_{i<j<m}(b_j-b_i)\ \ge\
V_{k-1}(v)\prod_{j<m}b_j^{\,k-1},
\]
since $v_i\ge0$ and $b_j-b_i\ge1$. It remains to show
$\log V_{k-1}(v)-\log V_k(u)\ge-2k\log(8a)-2k(1+\log k)$.

\emph{Outer pairs.} Every difference $u_i-u_j$ $(i<j)$ with $i=0$ or
$j=k-1$ is at most $8a$. In the first case of (P5), $u_0-u_{k-1}\le3a$. In
the second case, (P8), (P6) and $x\ge1$ give $4k\Delta\le b_0/x<w_0\le2B\le
2a$; then $u_0-u_j<(u_0-v_0)+(v_0-v_j)\le12k\Delta$ for $j\le k-2$,
$u_0-u_{k-1}=(u_0-v_0)+4\Delta(k-2)+\frac12\le12k\Delta$, and
$u_i-u_{k-1}<v_{i-1}-v_{k-2}+\frac12\le4k\Delta$ for $1\le i\le k-2$; all
are at most $6a$. There are $2k-3$ outer pairs.

\emph{Inner pairs.} For $1\le i<j\le k-2$, interlacing gives
$u_i-u_j<v_{i-1}-v_j=4\Delta(j-i+1)$, while the pair $(i-1,j-1)$ of $v$
contributes $4\Delta(j-i)$. The map $(i,j)\mapsto(i-1,j-1)$ is injective
into the pairs of $v$, and every pair of $v$ contributes a factor at least
$4\Delta\ge1$. Hence
\[
\log V_{k-1}(v)-\log V_k(u)\ge-(2k-3)\log(8a)-\sum_{d=1}^{k-3}(k-2-d)\log\Big(1+\frac1d\Big)
\ge-2k\log(8a)-k(1+\log k),
\]
by $\log(1+1/d)\le1/d$ and $\sum_{d\le k}1/d\le1+\log k$.
\end{proof}

\begin{proof}[Proof of (R1)]
All quantities are positive by Lemma~\ref{i4:lem-vandplace}. Put
$\ell_v=\log(v_E/\pi)\in[0,L]$ and $K=(2+2C_0)/\varepsilon$.

\emph{Factorials cancel.} Since $\sum_{M=k+1}^r\log(M-1)!=\sum_{i<r}\log i!
-\sum_{i<k}\log i!$, the mass relation and the definition of $\dim$ give
\[
\log\frac{D_*}{d_*}\ge f\log\tfrac23+\sum_{M=k+1}^{r}\Big(fM\ell_v-\tfrac f2\log(2v_E)+\log\Lambda(M)\Big)
+\log\frac{I_\sigma}{I_\pi}+\sum_\rho\Big(\log V_r(\mathrm{row}_\rho)-\log V_k(u_\rho)\Big).
\]

\emph{Elementary terms.} $f\log\frac23\ge-f$. As $(m-1)(k+1)\ge km-k$,
$\sum_MfM\ell_v\ge fkm\ell_v-fkL$. As $\log(2v_E)\le1+L\le2L$, the second
sum is at least $-fmL$. For $M\ge3$, $\zeta(M)<2$, so $\log\Lambda(M)\ge-f$
and $\sum_M\log\Lambda(M)\ge-fmL$. As $(2k+m-1)/2\ge k$ and $g_1g_2\ge1$,
$\log(I_\sigma/I_\pi)\ge-2f+fkm\log(g_1g_2)$.

\emph{One place.} Fix $\rho$ and drop it. By (P6)--(P8),
$\log b_j\le\log w_j\le1+L$, and $w_j<b_{j+1}$, $b_0\ge1$ give the
telescoping bound $\sum_j(\log w_j-\log b_j)\le\log w_{m-1}-\log b_0\le1+L$.
Also $\log x\le L/\varepsilon$, $\log(8a)\le7+L/\varepsilon+L$ and
$\log k\le C_0L/\varepsilon$. By Lemma~\ref{i4:lem-vandplace},
\[
k\sum_j\log w_j=(k-1)\sum_j\log b_j+\sum_j\log b_j+k\sum_j(\log w_j-\log b_j)
\le\big(\log V_r(\mathrm{row})-\log V_k(u)\big)+kL(20+K)+2mL,
\]
using $L\ge1$ to absorb constants.

\emph{Conclusion.} Summing over $\rho$,
$\log(D_*/d_*)\ge k\sum_{\rho,j}\log w_{\rho,j}+fkm\ell_v+fkm\log(g_1g_2)
-f(3+kL+4mL+kL(20+K))$. Divide by $fk$; with $3/k\le2L$ and
$4mL/k\le4m^2L/k$ the error is at most $(23+K)L+4m^2L/k\le C(L+m^2L/k)$.
\end{proof}

\medskip\noindent\emph{(R2): place data.}

\begin{lemma}[place data from the shape of a parameter]\label{i4:lem-placeshape}
Let $k\ge m+3$, $m\ge2$, $t\ge0$, and let $(u,v,b,w)$ be a Harish-Chandra
parameter of ranks $(k,m,t)$: $v_i<u_i$ and $u_{i+1}<v_i$ for $i<k-1$,
$-b_0<u_{k-1}$, $w_j-b_j\ge(t+3)/2$ and $b_{j+1}-w_j\ge(t+3)/2$. Let
$x,\Delta\ge1$, $h_0\ge0$, $a,B,b_1\in\R$, and assume:
$v_i=(k-2-i)4\Delta$; either $u_0\le2a$ and $-a\le u_{k-1}$, or
$u_0=(k-2)4\Delta+4k\Delta+\frac12$ and $u_{k-1}=-\frac12$;
$b_j=b_1+jh_0$, $w_j=b_j+h_0/2$; $x(4k\Delta+h_0)\le b_1$; and
$b_1+mh_0\le2B$. Then $(u,v,b,w)$ is place data for $(x,\Delta,a,B,k,m)$
and $w_j\ge1$.
\end{lemma}
\begin{proof}
Note $4k\Delta\ge1$, $b_j\ge b_1$, $x\cdot4k\Delta\le b_1$, $xh_0\le b_1$.
(P1), (P2) are read off from $v_i=(k-2-i)4\Delta$; (P3), (P4) are
conditions of the parameter. (P5): in the second alternative
$u_0-v_0=4k\Delta+\frac12\in[4k\Delta,8k\Delta]$ and
$u_{k-1}=v_{k-2}-\frac12$. (P6): $w_j-b_j=h_0/2\le b_1/x\le b_j/x$ and
$w_j-b_j>0$. (P7): $b_{j+1}-b_j\ge t+3\ge1$. (P8): $4k\Delta\le b_0/x$ and
$w_{m-1}\le b_1+mh_0\le2B$. Finally $w_j\ge b_1\ge x\cdot4k\Delta\ge1$.
\end{proof}

We recall the choice conditions used (\S\ref{i4:ssec-proof-realise},
step 2), with $\Delta=\lceil x^{c_D}\rceil$, $h_0=2\lceil x^{c_h}\rceil$.
\emph{Window:} for $t$ in the top set, $a\le a_t\le a+a/(kx^{13})$ and
$b'\le b_t\le b'+2n\delta_1x^\varepsilon\Lambda_Gb'/k$. \emph{Shape:} the
column type at every archimedean place $\tau$ is a Harish-Chandra parameter
with $v_i=(k-2-i)4\Delta$; off the top places $u_0=(k-2)4\Delta+4k\Delta+
\frac12$, $u_{k-1}=-\frac12$; $b_j=b_1(\tau)+jh_0$, $w_j=b_j+h_0/2$, where
$b_1(\tau)=b_t+\lfloor b'/x\rfloor$ at the top place of $t$,
$\lceil B\rceil\le b_1(\tau)<\lceil B\rceil+w_E$ at a thick place ($w_E$ the
number of roots of unity in $E$), and $b_1(\tau)=\lceil w''_n\rceil$ at a
thin place. \emph{Top:} at the top place of $t$, $u_0=a_t+(k-1)/2$ and
$u_{k-1}=-b_t-(k-1)/2$.

\begin{proof}[Proof of (R2)]
\emph{Scale estimates.} By Lemma~\ref{i4:lem-kits-shared}(b),(c),(e), for
large $x$: $f\le x^3$, $m\le x$, $x^{9/2+10\varepsilon}\le k\le x^5$,
$n\le2x^{2+8\varepsilon}$, $\Delta\le2x^{c_D}$, $h_0\le4x^{c_h}$; $w''_n$
exceeds any fixed power of $x$; and $X'\ge3x$. Hence, beyond a threshold,
\[
x(4k\Delta+h_0)\le w''_n\le b',\quad 3b'\le B,\quad mh_0+8f^2+2\le B,\quad k\le B,\quad 2B\le a.
\]
By Lemma~\ref{i4:lem-kits-shared}(e) and $(\log x)^2\le4x$,
$\Lambda_G\le x^{2+\varepsilon}/4$ for large $x$, so
$2n\delta_1x^\varepsilon\Lambda_G\le x^{4+10\varepsilon}\le k$ and the
window condition gives $b'\le b_t\le2b'$. Finally $w_E\le2[E:\Q]^2=8f^2$.

\emph{The base point.} We have $w''_n\le b_1(\tau)\le B+8f^2+2$: at a top
place $w''_n\le b'\le b_1\le3b'\le B$; at a thick place
$B\le b_1<B+1+8f^2$; at a thin place $b_1\in[w''_n,w''_n+1)$ with
$w''_n\le B$.

\emph{Conclusion.} The hypotheses of Lemma~\ref{i4:lem-placeshape} hold:
$x(4k\Delta+h_0)\le w''_n\le b_1$ and $b_1+mh_0\le2B$. For the alternative
on $u$: off the top places it is a choice condition. At the top place of
$t$, $a_t\le\frac32a$ and $k\le B\le a/2$ give $u_0\le2a$, and
$b_t\le2b'\le B\le a/2$ gives $-u_{k-1}\le a$.
\end{proof}

\medskip\noindent\emph{(R3): block lengths.}

We use the notation of \S\ref{i4:ssec-proof-kits}. In addition, put
\[
A=10\,r\Big(f+\frac{f^2(k-1)^3L^2}{m}\Big),\quad n_W=8\Big\lceil\frac A8\Big\rceil,\quad
z=frL,\quad P_P=C_Pz^{12},\quad D_p=f\binom r2,\quad n_P=8fr+4,\quad s=r+6,
\]
and let $\delta_2\ge x^{-\varepsilon}$ be the type density of the host. A
standard record has $C_t=401$, $C_p=1$, $c_{TO}=377$, $C_P\ge1000$ and
$C_b,C_{hT},C_{bN}\ge1$. Rank validity includes $r+t=n_P+T_On_W$.

\begin{lemma}[explicit rank bounds]\label{i4:lem-rankbounds}
For $x\ge1000^{1000}$ and valid ranks: $f\le x^{2.156}$,
$k\le r\le\min(x^{4.694},2k)$, $x^{0.498}\le m\le x^{0.501}$,
$x^{13/2}\le frL\le x^{6.87}$, $f\binom r2\le x^{12}$, $kf/m\le x^7$,
$n_W\le x^{22.9}$, $x^{377}\le t\le x^{401}$, $N_m\le x^{402}$.
\end{lemma}
\begin{proof}
Since $x^{1/1000}\ge1000$, a factor at most $1000$ costs at most
$x^{1/1000}$. With $\varepsilon<1/52$: $f\le4x^{2+8\varepsilon}$,
$r\le2k\le4x^{9/2+10\varepsilon}$, $x^{1/2-\varepsilon/16}\le m\le
2x^{1/2-\varepsilon/16}$, $L\le x^{0.02}$; the first five bounds are
products of these. Next $n_W\le A+8$, and
$f^2r(k-1)^3L^2/m\le x^{22.7}$, so $n_W\le x^{22.9}$. From
$r+t=8fr+4+T_On_W$ and $n_W\ge1$, $t\ge T_O\ge x^{377}$ and
$t\le13x^{400.9}\le x^{401}$. Finally $N_m=k+m(t+1)\le x^{402}$.
\end{proof}

\begin{lemma}[the nine block-length statements]\label{i4:lem-E2parts}
For $\varepsilon\in(0,1/52)$ and standard $C$, each of the following holds
eventually; $x_r=1000^{1000}$.
\begin{enumerate}
\item[(1)] $P_P,D_p,T_O\le x^{401}$.
\item[(2)] $A\le n_W<A+8$, $8\mid n_W$; $T_O=\lceil x^{377}\rceil$; $s$ is
odd and $r+5<s\le r+7$; $fr(r+1)\le n_W$, $(k-1)f+2r\le n_W$, $n_P\le n_W$.
\item[(3)] $(2fm)^4n_W\le P_P$ and $n_W\le P_P$.
\item[(4)] $P_P^4D_p^3kfx/(m\delta_2)\le T_O$ and $D_pk/m\le T_O$.
\item[(5)] $(D_p+2)n_W+2r\le T_O$, $40(D_p+2)\le T_O$, $11D_p^2\le T_O$,
$r+2(2\cdot5^4r)(D_p+1)\le T_O$.
\item[(6)] $T_OP_Px^{C_b}\le b$, $T_OP_P\le b^{1/2}$,
$P_P^3T_O^2x^{C_b}\le b$, $P_P^3T_O^2\le b^{1/2}$.
\item[(7)] For $c_D\ge805$: $T_O^2\le\Delta$ and
$(T_O+D_p+2)(T_O+D_p+2+2N_m)\le\Delta$.
\item[(8)] For $c_h\ge1000$: $4T_O^2+2t\le h_0$,
$P_PD_pfm^2N_mT_O\le h_0$, $P_P(fmk)^3k/m\le h_0$,
$C_{hT}T_ON_mf^2m^2kD_pP_P\le h_0$.
\item[(9)] $T_OD_p^2(\log b)x^{C_b}\le b$, $T_OD_p^2\log b\le b^{1/2}$,
$C_{bN}rmkP_P^2D_p^3\le b$.
\end{enumerate}
\end{lemma}
\begin{proof}
All are comparisons of explicit powers of $x$ through
Lemma~\ref{i4:lem-rankbounds}, except (2), which is arithmetic, and (6),
(9), which compare powers of $x$ with $b=\exp(x^\varepsilon)$ through
Lemma~\ref{i4:lem-kits-shared}(a),(b). We give the exponents.
(1) $P_P\le x\cdot x^{6.87\cdot12}\le x^{401}$ for $x\ge C_P$;
$D_p\le x^{12}$; $T_O\le x^{378}$.
(2) The first two are properties of $8\lceil A/8\rceil$; $k,m$ even make
$r$ odd and $s=r+6$ odd. With $j=k-1$, $r+1\le2j$, $f\ge m$, $L\ge1$:
$n_W\ge A\ge10fr+10frj^3$, which dominates $fr(r+1)\le2frj$,
$jf+2r\le2fr$ and $8fr+4\le10fr$.
(3) $fm\le z$ and $n_W\le10fr+10f^2r^4L^2+8\le28z^4$, so
$(2fm)^4n_W\le448z^8\le1000z^{12}\le P_P$.
(4) With $\delta_2^{-1}\le x$ and $C_P^4\le x$: the left side is at most
$x^{1+330+36+7+1+1}=x^{376}\le T_O$; and $D_pk\le x^{17}$.
(5) With $y=x^{23}$, $D_p,n_W,r\le y$ and $T_O\ge y^3$; each left side is
at most $10^4y^2\le y^3$.
(6) $z\le x^{14}$, so $P_P^3T_O^2x^{C_b}\le4C_P^3x^{1258+C_b}\le
x^{1259+C_b}\le b^{1/2}$ for large $x$; the other three follow from
$T_OP_P\le P_P^3T_O^2$ and $b^{1/2}\le b$.
(7) With $y=x^{402}$: $\Delta\ge x^{805}\ge35y^2$, while $T_O^2\le4y^2$
and the product is at most $5y\cdot7y$.
(8) $T_O\le x^{378}$, $t\le x^{401}$, $P_P\le x^{85}$, $N_m\le x^{402}$:
the four left sides are at most $x^{757},x^{882},x^{117},x^{891}$, and
$h_0\ge x^{1000}$.
(9) Every left side is at most $Kx^N$ for fixed $K,N$ (using $\log b\le x$,
$D_p\le32x^{15}$, $z\le8x^{11}$), and $Kx^N\le\exp(x^\varepsilon/2)$ for
large $x$ by Lemma~\ref{i4:lem-kits-shared}(a).
\end{proof}

\begin{proof}[Proof of (R3)]
Eventual statements are closed under conjunction; take the maxima of the
thresholds in Lemma~\ref{i4:lem-E2parts}.
\end{proof}

\medskip\noindent\emph{(R4): window weights.}

The property in (R4) is derived from three closed properties of a family:
(i) the closed form of the first window theorem, which for every
admissible middle constant $\Lambda\ge T_0$ supplies a constant
$C(\Lambda)>0$ and a bound $T_2(\Lambda)>0$; (ii) the sign conditions
(middle window constant $\Lambda>0$, $C_e\ge0$); (iii) the logarithmic
window bound, which holds once the window constant $C_N$ exceeds a
constant $C_\star(\varepsilon)$.

\begin{definition}[weight requirement]\label{i4:Rdef-weight}
Fix a point $(D,L)$ of a family, with scale $x$, exponent $\varepsilon$
and the family constants $K$. Let $\mathcal A$ and $\mathcal B$ be the
finite boxes of integer vectors on the tops $T$ given by the
two windows of the point, and $\mathcal P$ the finite set of
\emph{candidate primes} ($p<X'=\exp(x^{\varepsilon/3})$ satisfying the
candidate condition of the window count). For $(a,b)\in\mathcal A\times\mathcal
B$ and $p\in\mathcal P$ let
\[
\Sigma(a,b,p)=m\sum_{t_0\in\mathrm{GR}(p)}\kappa_{\rm loc}\bigl(p,T_{t_0},a|_{t_0},b|_{t_0}\bigr),
\]
where: $h_p\in U$ is the Frobenius step at $p$, of additive order $o_p$;
$\mathrm{GR}(p)$ is the set of good representatives $t_0$ of the
$h_p$-orbits (a fixed set of representatives, one in each orbit,
restricted to the good ones, see below); $T_{t_0}=\{j<o_p: t_0+jh_p\in T\}$ (the tops on the orbit, in
orbit coordinates) and $a|_{t_0}(j)=a(t_0+jh_p)$, likewise $b|_{t_0}$; and
\[
\kappa_{\rm loc}(p,T',a',b')=\begin{cases}\min\bigl(1,\;K_{II}\,|T'|\,r_G^2\,k\,x^{-A}\bigr)&p<p_0(k),\\
\kappa_I(p,o_p,T',a',b')&p\ge p_0(k),\end{cases}
\]
with $K_{II}$ an entry of the family constants, $r_G$ the band length of the
scale, $A$ the scale exponent, $p_0(k)$ the band threshold, and
$\kappa_I$ the type-I excess of the point: beyond the first threshold of
the family, $\kappa_I(p,o,T',a',b')$ is the least admissible excess, that is
the infimum of the $\kappa\in[0,1]$ for which the local lattice property of
excess $\kappa$ holds (\S\ref{i4:Rst-kappaI}); in particular
$0\le\kappa_I\le1$. A prime
$p<X'$ is a candidate if it is prime, its prime class (a datum of the
instance) is the inert class, $x^{\varepsilon/4}\le o_p$, and, if
$p\ge p_0(k)$, $p$ is not in the band $\mathrm{Band}(\rho_B)$, where for
$\rho'>0$ the band $\mathrm{Band}(\rho')$ is the set of $p$ with
$p_0(k)\le p<X'$ and $X'(1-\rho')\le p^i\le X'(1+\rho')$ for some $i\ge1$; the
\emph{narrow band} is $\mathrm{Band}(\rho)$, with $\rho\le\rho_B$ the two band
ratios of the point.

\emph{Good representatives.} Let $n'=o_p$ and $T'=T_{t_0}\subset\{0,\dots,n'-1\}$.
With $h=2+8\varepsilon$, $K_G=r_G-1=\lceil8(h+2)/\varepsilon\rceil$,
$u=\mathrm{good}_U(\delta_1)=\lceil1/\delta_1\rceil$, and $M$, $G_p=G'$, $N_1$ as in
Lemma~\ref{i4:lem-kits-shared}(e), the representative $t_0$ is \emph{good} if:
(a) every cyclic window $\{i,i+1,\dots,i+L_2-1\}\bmod n'$ contains at most
$K_G$ elements of $T'$, where $L_2=2\max_p\lceil\log_p((kx^A)^4)\rceil$, the
maximum over the primes of inert class;
(b) every cyclic window of length $G'$ contains at most $N_1$ elements of $T'$;
(c) with $\ell=(6r_G+1)u$ and $S=\lfloor n'/\ell\rfloor$, for every $s_0<S$
there is $j\le M$ such that the block $s=(s_0+j)\bmod S$ is \emph{marked}:
for $1\le c\le6r_G+1$, the cell $[s\ell+(c-1)u,\,s\ell+cu)$ contains no
element of $T'$ if $c$ is odd and exactly one element of $T'$, which is
even, if $c$ is even.
The windows are $\mathcal A=\mathrm{box}(W_A)$, $\mathcal B=\mathrm{box}(W_B)$,
the integer vectors supported on $T$ with entries in the integer intervals
$W_A=[\lceil a\rceil,\lfloor a+a/(kx^{A/2})\rfloor]$ and
$W_B=[\lceil b'\rceil,\lfloor b'+N_*x^\varepsilon\Lambda b'/k\rfloor]$; here
$a$ and $b'=a/X'$ are the scale quantities of \S\ref{i4:ssec-proof-kits},
not weight vectors. A
\emph{weight vector for $(D,L)$ in the windows} is a pair
$(a,b)\in\mathcal A\times\mathcal B$ with
\[
\mathrm{WReq}(a,b):\qquad\sum_{p\in\mathcal P}\frac{\log p}{p}\Bigl(2\cdot
\mathbf 1\bigl[\Sigma(a,b,p)>x^{-\varepsilon}\bigr]+4\,\Sigma(a,b,p)\Bigr)\le
x^{-9}.
\]
\end{definition}

\begin{proposition}[weights from the three window properties]\label{i4:Rprop-window}
Let $\mathcal F$ be a family of exponent $\varepsilon$ such that
\begin{enumerate}
\item[(i)] (\emph{first window theorem at the family}) there are $C_H,C\ge0$
with $C_H\le k_I$, $C_H\le D_I$, $2C_H\le\rho$, $2C\le\kappa_1$, $C\le\kappa_2$
(entries of the family constants $K$) such that the second conclusion of the
first window theorem holds for $\mathcal F$ with constants
$(\Lambda,C_H,C)$, $\Lambda$ the middle window constant of $K$;
\item[(ii)] $\Lambda>0$;
\item[(iii)] the logarithmic window bound holds eventually at every point.
\end{enumerate}
Then for some $x_0$, for all $x\ge x_0$ and $i\in I(x)$, a weight vector
for $(D(x,i),L(x,i))$ exists in the windows of $\mathcal F$.
\end{proposition}
\noindent (The condition $C_e\ge0$ is not needed here; it is used only in the
closedness argument for (R4) below.) The proof is in
\S\ref{i4:ssec-proof-window}; it is complete modulo
Proposition~\ref{i4:Rprop-window-leaves}.

\begin{proposition}[leaves of the window count]\label{i4:Rprop-window-leaves}
Let $\mathcal F$ satisfy (i)--(iii) of Proposition~\ref{i4:Rprop-window}.
There is $x_0$ such that at every point $(D,L)=(D(x,i),L(x,i))$
with $x\ge x_0$, writing $m$ for the rank and $|T|$ for the number of tops:
\begin{enumerate}
\item[(W0)] $\mathcal A\ne\emptyset$ and $\mathcal B\ne\emptyset$;
\item[(W1)] $\Sigma(a,b,p)\ge0$ for all $a,b$ and all $p\in\mathcal P$;
\item[(W2)] for every $p\in\mathcal P$,
$\sum_{a\in\mathcal A}\sum_{b\in\mathcal B}\Sigma(a,b,p)\le|\mathcal A|\,|\mathcal B|\,m\,|T|\,x^{-12}$;
\item[(W3)] $\sum_{p\in\mathcal P}\log p/p\le x^{\varepsilon/3}(2+x^{\varepsilon/3})$;
\item[(W4)] $(x^{\varepsilon}+1)\,m\,|T|\,x^{-12}\cdot x^{\varepsilon/3}(2+x^{\varepsilon/3})\le x^{-9}/4$.
\end{enumerate}
\end{proposition}
\noindent The proof (\S\ref{i4:ssec-proof-wleaves}) is complete modulo
Propositions~\ref{i4:Rprop-kappaI} and~\ref{i4:Rprop-num}.

\begin{proposition}[the type-I excess; formally verified, \cite{CharacterFamily}]\label{i4:Rprop-kappaI}
Let $\mathcal F$ satisfy (i) of Proposition~\ref{i4:Rprop-window}. There is
$C_m\ge0$ such that, eventually, at every point, for every prime $p$ with
$p_0(k)\le p<X'$ that is not in the narrow band $\mathrm{Band}(\rho)$, every
good representative $t_0$ with $3G_p\le o_p$, at which
the level bound $L_2<\mathrm{good}_U(\delta_1)$ and the hypothesis $\mathrm{HypI}(p,o_p,|T_{t_0}|)$ hold: $\kappa_I(p,o_p,T_{t_0},a',b')\ge0$ for all
$a',b'$, and its mean over $\mathrm{box}_{T_{t_0}}(W_A)\times\mathrm{box}_{T_{t_0}}(W_B)$ is at most
$C_m|T_{t_0}|N_1/x^{A/2}$.
\end{proposition}
\noindent The structure of its formal proof is described in \S\ref{i4:Rst-kappaI}.

\begin{proposition}[eventual numerics; formally verified, \cite{CharacterFamily}]\label{i4:Rprop-num}
Let $\mathcal F$ satisfy (i)--(iii) of Proposition~\ref{i4:Rprop-window},
and let $C_m$ be as in Proposition~\ref{i4:Rprop-kappaI}. Then eventually at every
point: $K_{II}r_G^2kx^{-A}\le x^{-12}$; $C_mN_1/x^{A/2}\le x^{-12}$;
$L_2<\mathrm{good}_U(\delta_1)$; $3G_p\le x^{\varepsilon/4}$;
$1\le a/(kx^{A/2})$; $1\le N_*x^\varepsilon\Lambda b'/k$; the band
ratio $\rho\le\rho_B$; and $\mathrm{HypI}(p,o_p,|T_{t_0}|)$ at every
candidate $p\ge p_0(k)$ and good $t_0$. Moreover the size bounds
$m\le x^{1/2+\varepsilon/100}$ and $|T|\le2x^{2+8\varepsilon+\varepsilon/100}$ hold.
\end{proposition}
\noindent The structure of its formal proof is described in \S\ref{i4:Rst-num}.

\begin{proof}[Proof of (R4) from Proposition~\ref{i4:Rprop-window}]
By Lemma~\ref{i4:lem-closed} and Proposition~\ref{i4:Rprop-window} it
suffices that the class of all families is closed for each of (i)--(iii);
each is witnessed with exponent $1$. (ii): take the threshold datum with
middle threshold $1$ and the other thresholds $0$; a large member has
$\Lambda\ge1>0$, and $C_e=1$ is a fixed entry. (iii): take the last
threshold $C_\star(\varepsilon)$. (i): take middle threshold $T_0$ and as
last threshold $l\mapsto\max_{c\in l}\max(2T_2(c),2C(c))$. For a large
member $\Lambda\ge T_0$ and every last-list constant is at least
$2T_2(\Lambda)$ and $2C(\Lambda)$, and these are the clauses of the first
theorem with $C_H=T_2(\Lambda)$ and constant $C(\Lambda)$.
\end{proof}

\medskip\noindent\emph{(R5): class decomposition.}

\begin{proof}[Proof of (R5)]
\emph{Reduction to finite adeles.} $K$ is open, hence measurable. The
projection $\pi:G\to G_f=U(W)(\mathbb A_{F,f})$ is onto, injective on
$\Gamma$, and $K=\pi^{-1}(K_f)$ with $K_f$ the stabiliser of
$\hat L=\prod_vL_v$. So $\Gamma gK\mapsto\Gamma\pi(g)K_f$ is a bijection of
double coset spaces preserving the groups $\Gamma\cap gKg^{-1}$.

\emph{Lattices.} For $g\in G_f$ let $\Lambda_g=\{x\in W:x\in g_vL_v\text{
for all finite }v\}$, a lattice with completions $g_vL_v$; $\Lambda_1=L$
and $\Lambda_{\gamma gk}=\gamma\Lambda_g$. Conversely, if
$\Lambda_g=\gamma\Lambda_{g'}$, then $g\hat L$ and $\gamma g'\hat L$ are both
the closure of the $\hat{\mathcal O}_E$-span of $\Lambda_g$ (strong
approximation for $\mathcal O_F$: every integral finite adele is $r+Dy$
with $r\in\mathcal O_F$, $y$ integral, for any $D\ne0$), so
$\Gamma gK_f=\Gamma g'K_f$. Every $\Lambda_g$ lies in the genus of $L$, and
every lattice $M$ of the genus is some $\Lambda_g$ (choose $g_v$ with
$g_vL_v=M_v$, and $g_v=1$ where $M_v=L_v$). Hence $\Gamma gK_f\mapsto
[\Lambda_g]$ is a bijection onto the isometry classes in the genus of $L$.

\emph{Finiteness.} $W$ is totally definite, so
$\mathrm{Tr}_{E/\Q}(c\,h(x,x))$ is positive definite for suitable
$c\in\mathcal O_F\smallsetminus\{0\}$. There are $d,e\ne0$ with
$d\,h(\Lambda_g,\Lambda_g)\subset\mathcal O_E$ and
$e\Lambda_g^\#\subset\Lambda_g$ for all $g$ (true for $L$, and local
isometries preserve both). Such lattices have bounded covolume in the
Euclidean space $(W,\mathrm{Tr}(c\,h))$ and bounded denominators, so by
reduction theory for positive definite lattices (Hermite's finiteness,
reference to be confirmed) they form finitely many isometry classes. So the
genus, hence $\Gamma\backslash G/K$, is finite.

\emph{Stabilisers.} $\Gamma\cap gK_fg^{-1}=\mathrm{Aut}(\Lambda_g)$. Its
matrices in a basis have entries in a fixed fractional ideal and preserve a
positive definite form, so they are bounded at every archimedean place; a
discrete bounded set is finite.

\emph{Mass.} By the bijection, $\sum_i|\Gamma_i|^{-1}=\sum_{[M]}
|\mathrm{Aut}(M)|^{-1}=\mathrm{mass}(L)$. One $g_i$ in each double coset
gives the disjoint decomposition.
\end{proof}

\subsubsection{The window weights and the member hypotheses}\label{i4:ssec-proof-r4c}

\medskip\noindent\label{i4:ssec-proof-window}\emph{Proof of Proposition~\ref{i4:Rprop-window}.}

\begin{proof}
Let $x_0$ be as in Proposition~\ref{i4:Rprop-window-leaves} and fix a
point $(D,L)=(D(x,i),L(x,i))$ with $x\ge x_0$; we use the notation of
Definition~\ref{i4:Rdef-weight}. For $(a,b)\in\mathcal A\times\mathcal B$
put
\[
Z(a,b)=\sum_{p\in\mathcal P}\frac{\log p}{p}\Bigl(\mathbf 1\bigl[\Sigma(a,b,p)>x^{-\varepsilon}\bigr]+\Sigma(a,b,p)\Bigr).
\]

\emph{Step 1 (Markov).} Let $S\ge0$. If $S>x^{-\varepsilon}$ then
$\mathbf 1[S>x^{-\varepsilon}]=1<x^{\varepsilon}S$; otherwise the indicator
is $0\le x^\varepsilon S$. By (W1) this applies to $S=\Sigma(a,b,p)$, and
since $\log p/p\ge0$,
\[
Z(a,b)\le\sum_{p\in\mathcal P}\frac{\log p}{p}\,(x^{\varepsilon}+1)\,\Sigma(a,b,p).
\]

\emph{Step 2 (averaging).} Sum over $(a,b)\in\mathcal A\times\mathcal B$,
exchange the finite sums and apply (W2) at each $p$, then (W3) and (W4):
\[
\sum_{a,b}Z(a,b)\le|\mathcal A||\mathcal B|\,(x^\varepsilon+1)\,m|T|x^{-12}
\sum_{p\in\mathcal P}\frac{\log p}{p}
\le|\mathcal A||\mathcal B|\,(x^\varepsilon+1)\,m|T|x^{-12}\,x^{\varepsilon/3}(2+x^{\varepsilon/3})
\le|\mathcal A||\mathcal B|\,\frac{x^{-9}}4 .
\]
The right side is the sum over $\mathcal A\times\mathcal B$ of the
constant $x^{-9}/4$. Since $\mathcal A$ and $\mathcal B$ are nonempty (W0),
some $(a,b)$ has $Z(a,b)\le x^{-9}/4$ (a sum of finitely many reals bounded
by the sum of constants has a term bounded by the constant; apply this
first over $a$, then over $b$).

\emph{Step 3 (weight requirement).} For this $(a,b)$ and each $p$, with
$\iota\in\{0,1\}$ the indicator and $S=\Sigma(a,b,p)$, we have
$2\iota+4S\le4(\iota+S)$. Multiplying by $\log p/p\ge0$ and summing,
the left side of $\mathrm{WReq}(a,b)$ is at most $4Z(a,b)\le x^{-9}$.
So $(a,b)$ is a weight vector in the windows.
\end{proof}

\medskip\noindent\emph{Proof of Proposition~\ref{i4:A-prop-inputs}(a).}

\begin{proof}
We go through the 21 conjuncts of the member hypotheses in the order of
the statement. The ray class field statement and the principal ideal
theorem follow from (P1) by (M1). The statements (P2), (P4), (P5), (P6),
(P7), (P11), (P1), (P3) are hypotheses. The Weil representation for
unitary dual pairs is (M2). Existence of discrete series, the
infinitesimal-character statement, existence in the volume
normalisation and the uniqueness for $U(p,q)$ are (M3), from (P8).
$\mathrm{RowsFinDim}$ is (M5). The formal degree formula in the volume
normalisation is (P10). The Siegel mass ratio is (M6). The three
statements for every Harish-Chandra parameter are (M4), from (P9), (P8)
and the uniqueness just obtained. This is the full list.
\end{proof}


\subsubsection{The leaves of the window count}\label{i4:ssec-proof-wleaves}

\begin{proof}[Proof of Proposition~\ref{i4:Rprop-window-leaves}]
Take $C_m$ from Proposition~\ref{i4:Rprop-kappaI} and $x_0$ beyond the
thresholds of Propositions~\ref{i4:Rprop-kappaI} and~\ref{i4:Rprop-num}
and beyond $48^{1/(1/2-9.69\varepsilon)}$ (used in (W4)).

(W0) $W_A=[\lceil y\rceil,\lfloor y+w\rfloor]$ with $w=a/(kx^{A/2})\ge1$ is
nonempty because $\lceil y\rceil<y+1\le y+w$; likewise $W_B$. A box over a
nonempty interval is nonempty (coordinates off $T$ are $0$).

(W1) For $p<p_0(k)$, $\kappa_{\rm loc}=\min(1,\cdot)\ge0$, as the
family constants are nonnegative, so $K_{II}\ge0$.
For $p\ge p_0(k)$ and $t_0$ good, the hypotheses of
Proposition~\ref{i4:Rprop-kappaI} hold: $3G_p\le x^{\varepsilon/4}\le o_p$;
$L_2<\mathrm{good}_U$; $\mathrm{HypI}$; $p_0(k)\le p<X'$; and $p$ is not
in the narrow band: a candidate is not in $\mathrm{Band}(\rho_B)$, and
$\mathrm{Band}(\rho)\subset\mathrm{Band}(\rho_B)$ as $\rho\le\rho_B$
(Proposition~\ref{i4:Rprop-num}). So $\kappa_I\ge0$, and $\Sigma=m\sum\kappa_{\rm loc}\ge0$.

(W2) \emph{Marginals.} The map $j\mapsto t_0+jh_p$ ($j<o_p$) is injective,
so for every function $F$ of the restriction,
$\sum_{a\in\mathcal A}F(a|_{t_0})=c_A\sum_{a'\in\mathrm{box}_{T_{t_0}}(W_A)}F(a')$ and
$|\mathcal A|=c_A|\mathrm{box}_{T_{t_0}}(W_A)|$, with $c_A$ the number of
completions off the orbit; the same for $\mathcal B$. Hence the average of
$\kappa_{\rm loc}(p,T_{t_0},a|_{t_0},b|_{t_0})$ over $\mathcal A\times\mathcal B$ equals its
mean over the local boxes. \emph{Mean bound.} For $p<p_0(k)$ the value is
constant and $\le K_{II}|T_{t_0}|r_G^2kx^{-A}\le|T_{t_0}|x^{-12}$; for
$p\ge p_0(k)$ the mean is $\le C_m|T_{t_0}|N_1/x^{A/2}\le|T_{t_0}|x^{-12}$
(Propositions~\ref{i4:Rprop-kappaI}, \ref{i4:Rprop-num}). Here
$|T_{t_0}|$ is the number of tops on the orbit of $t_0$ (the orbit of a
representative is the image of $T_{t_0}$). \emph{Sum.} Hence
$\sum_{a,b}\Sigma(a,b,p)=m\sum_{t_0}\sum_{a,b}\kappa_{\rm loc}\le|\mathcal A||\mathcal B|\,m\,x^{-12}\sum_{t_0\in\mathrm{GR}(p)}|T_{t_0}|$,
and the orbits of distinct representatives are disjoint, so
$\sum_{t_0}|T_{t_0}|\le|T|$.

(W3) Every candidate is $<X'=e^u$, $u=x^{\varepsilon/3}$. For $1\le n<X'$,
$\log n<u$, so $\sum_{p\in\mathcal P}\log p/p\le u\sum_{1\le n\le M-1}1/n$
with $M=\lceil X'\rceil$; the harmonic sum is $\le1+\log(M-1)\le1+u$. So the
sum is $\le u(1+u)\le u(2+u)$.

(W4) With $m\le x^{1/2+\varepsilon/100}$, $|T|\le2x^{2+8\varepsilon+\varepsilon/100}$,
$x^\varepsilon+1\le2x^\varepsilon$ and $u(2+u)\le3u^2$ (as $u\ge1$), the
left side is at most
$12x^{\varepsilon+\varepsilon/100+1/2+2+8\varepsilon+\varepsilon/100-12+2\varepsilon/3}$; the exponent
is at most $-19/2+9.69\varepsilon<-9$ for $\varepsilon<1/52$, so
$48\,x^{-19/2+9.69\varepsilon}\le x^{-9}$ as soon as
$x\ge48^{1/(1/2-9.69\varepsilon)}$, which holds for $x\ge x_0$; this gives
the bound $x^{-9}/4$.
\end{proof}

\subsection{The automorphic construction}\label{i4:ssec-grp-aut}

\subsubsection{Proof of the automorphic construction}\label{i4:ssec-proof-assembly}

Parts (a) and (d) are (A1) of Proposition~\ref{i4:prop-assembly-leaves}.
We prove (c) from (A2), and then (b) from (c), (A3) and (A4).

\emph{Linked rows.} A bundle $b$ is \emph{linked} at a point if there are
the following data:
\begin{itemize}
\item a witness point $q$, consisting of a base, a frame, a twist datum, a
column choice, a level, an object and a witness that contains one
norm-choice datum;
\item a bundle $b'$ over $q$ with a bundle witness;
\item an identification of $b$ with $b'$.
\end{itemize}
A row condition $H$ is \emph{linked} at a parameter system if, eventually,
some property $E$ holds for all linked bundles at built good points, and
$H$ holds for every selection of $E$-bundles.

Two principles follow from the definitions.
\begin{itemize}
\item[(P)] A property that holds for every bundle over a witness point with
a bundle witness gives a linked row.
\item[(I)] If $A\Rightarrow R$ holds for every such bundle and the row of
$A$ is linked, then the row of $R$ is linked.
\end{itemize}

\emph{Monotonicity.} Let $N'$ be a norm-choice extension stronger than $N$:
every $N'$-datum yields an $N$-datum at the same arguments. Replace the
norm-choice datum in a build record, or in a witness, by the corresponding
$N$-datum, and copy everything else. This shows that cores built (resp.\
good) at $p(N')$ are built (resp.\ good) at $p(N)$, and that bundles linked
at $p(N')$ are linked at $p(N)$. So the property $E$ of a row linked at
$p(N)$ holds for all linked bundles at $p(N')$, and every $E$-selection for
$p(N')$ is one for $p(N)$. Hence linked rows, and in particular the twelve
open statements, pass from $p(N)$ to $p(N')$.

\emph{The twelve rows at the origin.} Let $p_0$ be the system with origin
$c_{\rm Dep}$; in items 1--12, ``linked'' means linked at $p_0$.
\begin{enumerate}
\item \emph{Primitivity.} The real input of a bundle witness carries
primitivity of the realised rows and type separation. Apply (P).
\item \emph{Decomposition points.} The row conditions at the companion
(Proposition~\ref{i4:A-prop-open}(O8)) give the decomposition-point
condition at every bundle over a witness point with a bundle witness.
Apply (P).
\item \emph{The tie.} Every linked bundle has a comparison datum, by (A2).
\item \emph{The selection jump.} By
Proposition~\ref{i4:A-prop-open}(O9), applied with rows 1, 2 and 3 of this
list (primitivity, decomposition points and the tie, all linked at $p_0$),
the selection-jump row is linked at $p_0$. In words: the comparison datum
pins the realisation, and the join of the linked realisation gives the
join with the kernel.
\item \emph{Multiplicity.} By (A2), the condition $\mathrm{MI}^*_2$ holds at
every bundle. Its derivation uses Theorems~\ref{i4:thm-sw}
and~\ref{i4:thm-howe}, together with the doubling method and its
unramified computation.
\item \emph{Divisor rows.} Apply (I) with $A=\mathrm{MI}^*_2$: given
$\mathrm{MI}^*_2$, the pairings vanish (O3), which gives the level-three
divisor rows; the fourth clause of (O8) turns these into the divisor
condition.
\item \emph{$\Lambda'$-rows.} The $\Lambda$-rows hold at points with a
parameter choice (O4), and the first clause of (O8) turns them into the
$\Lambda'$-rows of the companion. The finiteness of types uses
Theorems~\ref{i4:thm-bzfin} and~\ref{i4:thm-green}.
\item \emph{The constant $\kappa$.} By (A2), keeping the slot of row (7).
The closedness requirement is met with exponent $1$ and the zero threshold
datum.
\item[9--10.] \emph{Goodness and the labels.} By (A2). Green's
parametrisation (Theorem~\ref{i4:thm-green}) supplies Green's hypothesis.
\setcounter{enumi}{10}
\item \emph{Geometry.} By (A2), from the ray class field statement, the
divisor rows and the tie.
\item \emph{Definite forms.} This open statement is the defect row on the
point-evaluation side, (O7). The written proof uses the orthogonality and
isotypic decomposition of Theorem~\ref{i4:thm-serre}.
\end{enumerate}
Each row is linked at $p_0$, so the open statements hold at the origin.
By monotonicity they hold at every stronger extension, in particular at
$p$. This is (c).

\emph{Part (b).} By (A3) the norm row for $(\mathrm{LN})$ is linked at
$p_0$, hence at $p$ by monotonicity. By (A4) the functional condition and
the covolume bound are linked at $p$. With (c), all hypotheses of
Theorem~\ref{i4:A-thm-allrows} at $p$ hold (its published hypotheses are
in Remark~\ref{i4:rem-trust}(ii), and surjectivity onto $X$ and the lower
bound $\lambda_{\rm low}$ are Proposition~\ref{i4:A-prop-funcov}(d)). It
gives the 26 rows, with $NN$ the norm property of (A3). This is (b).
\qed

\subsubsection{The row bounds at $2$ and $3$}\label{i4:ssec-proof-at23rows}

We use Definition~\ref{i4:A-def-cost}. For a prime $l$ and $h\in\N$ put
$\lambda_N(l,h)=s_l+\lfloor\log_l2h\rfloor+\lfloor\log_lN\rfloor+1$ and
\[
e_L(l,h)=32(2hs_l+1)+\lambda_N,\qquad c_{0}(l,h,n')=20(n'+1)\lambda_N,\qquad
c_B(l,h,n',t)=40(t+1)(n'+1)\lambda_N ,
\]
with $\lambda_N=\lambda_N(l,h)$. At a twisted place $k$ with prime $l=l_k$,
ring $\mathcal O_k$ and local level $s_0(k)$, a \emph{declaration} is a pair
$(\mathrm{col},\mathrm{row})$ of disjoint finite sets of places of $E$
over $l$ covering them; its rows are \emph{conjugate} to its columns if
$v\in\mathrm{row}$ exactly when $v$ is conjugate to some $u\in\mathrm{col}$.
For $\pi\in Z^{\rm good}_k$ and parameters $(e',c_0,c_B)$ the
\emph{split bound} is
\begin{multline*}
\mathrm{bs}(e',c_0,c_B)=\sum_{u\in\mathrm{col}}\bigl(e'+n_u m(e'+\delta)+2\delta+h^W_u+n_ume'+rc_0\bigr)
+s_\pi+\sum_{0<e<t}\mathrm{sha}_l(e)+v_l(h_E)\\
+t\sum_{w\in\mathrm{row}}(n_w+2)c_B+4n_V(|\mathrm{col}|+1)\delta+t\,|\mathrm{col}|\,h_0
+t\sum_{v\in Q_k}\ell\bigl(\mathcal O_k/(\zeta_v-1)\bigr),
\end{multline*}
where $n_u=e_uf_u$ is the local degree, $\delta=\sum_{0<j<t}(s_0+v_l(j))$,
$h^W_u$ the length of the invariants of $W'$ at $u$, $h_0$ the maximum over
$u\in\mathrm{col}$ and $0<j<t$ of the length (in units of $v(l)$) of the
invariants of the decomposition group at $u$ on $K_k/\mathcal O_k(j)$
($K_k$ the fraction field of $\mathcal O_k$, the twist by the $j$-th power
of the cyclotomic character),
$\ell$ the length in units of $v(l)$, and $\zeta_v$ the value of the tame
character at the place $v\in Q_k$.

\begin{proposition}[the split-case conclusion; formally verified, \cite{CharacterFamily}]\label{i4:A-prop-split}
Assume $\mathrm{Hypotheses}_2$ and a setup as in
Proposition~\ref{i4:A-prop-at23} at a datum with size datum $L$ and
parameter choice. Let $k$ be a twisted place, $l=l_k$, with: (H1) $l\ne
g_0,g_1$; (H2) $l$ splits in $K=\Q(\sqrt{-d})$ or $l\in$ small; (H3) if
$l\notin$ small, every place of $E$ over $l$ is unramified over $\Q$; (H4)
if $l=2$ then $4\mid N$; (H5) $l=2$, or $l$ is odd with $(\frac{-d}l)=1$.
Then there is a declaration whose rows are conjugate to its columns such
that for every $\pi\in Z^{\rm good}_k$:
\begin{enumerate}
\item[(S)] if $\pi$ has finite $s$-length, $\mathrm{sha}_l(e)<\infty$ for
$0<e<t$, and $\zeta_v\ne1$ for $v\in Q_k$, then
$v_{cc}(k,\pi)\le\mathrm{bs}\bigl(e_L(l,h_N),\,c_0(l,h_{\rm hub},4nn_V),\,c_B(l,\lfloor2b\rfloor,4nn_V,t)\bigr)$,
where $h_N\le h_B$ is the norm height of $L$ and $h_{\rm hub}=\lfloor2b\rfloor+1$;
\item[(W)] $h^W_u\le r\,c_0(l,h_{\rm hub},4nn_V)$ for $u\in\mathrm{col}$;
\item[(Z)] $\zeta_v\ne1$ for $v\in Q_k$.
\end{enumerate}
\end{proposition}
\noindent The structure of its formal proof is described in \S\ref{i4:Gst-split}.
\noindent In the formal proof the declaration comes from the place types
at $l$ (using (H5)), the local conditions at $l$ (lift, cyclotomic
invariants, $T$-block) from $\mathrm{Hypotheses}_2$, and (S), (W), (Z) from
the admitted package at $\pi$ (Fontaine, jump windows, Bloch--Kato
annihilators, Tate finiteness, the dual and closing statements). The
comparison of the last term of $\mathrm{bs}$ with $\epsilon(l)$ is not part
of this proposition: the killing condition of the setup contains, at every
twisted place, the bound
$t\sum_{v\in Q_k}\ell(\mathcal O_k/(\zeta_v-1))\le\epsilon(l)$, $\epsilon(l)=l^{-2}$
(clause (K2) of Definition~\ref{i4:A-def-at23}).

\begin{lemma}[invariant lengths]\label{i4:A-lem-h0}
Let $k$ be a twisted place, $l=l_k$, with $s_0=s_0(k)\ge1$ and $s_0\ge2$ if
$l=2$. Then for every declaration, $h_0\le s_0+\lfloor\log_lt\rfloor$.
\end{lemma}
\begin{proof}
Let $u\in\mathrm{col}$ and $0<j<t$. The inertia group at $u$ contains $g$
whose cyclotomic value is $\beta=1+l^{s_0}$, and $g$ acts on
$K_k/\mathcal O_k(j)$ by $\beta^{-j}$. An invariant $y$ satisfies
$\beta^{-j}y=y$, hence $(\beta^j-1)y=0$. Since $s_0\ge1$ ($s_0\ge2$ for
$l=2$), $v_l(\beta^j-1)=s_0+v_l(j)$ (lifting the exponent), so $a=\beta^j-1\ne0$
and the invariants lie in the $a$-torsion of $K_k/\mathcal O_k$, which is
$a^{-1}\mathcal O_k/\mathcal O_k\cong\mathcal O_k/(a)$; so their length is at
most that of $\mathcal O_k/(a)$, which is $s_0+v_l(j)$ in units of
$v(l)$. Finally $v_l(j)\le\lfloor\log_lj\rfloor\le\lfloor\log_lt\rfloor$.
\end{proof}

\begin{lemma}[numerics of the split bound]\label{i4:A-lem-numsplit}
Let $L$ carry a parameter choice, $l$ prime, and write $\Lambda=\Lambda_l$
(Definition~\ref{i4:A-def-cost}, so the parameter $y$ there is $y=h_B$). Then:
(a) $h_B\le6b$, $b\le h_B$, $N\le h_B$, $t\le h_B$, $1\le h_B$;
(b) $s_l+1\le\Lambda$, $1\le\Lambda$, and $\lfloor\log_ln\rfloor\le\Lambda$ if $n\le8h_B$;
(c) if $h\le4h_B$ then $\lambda_N(l,h)\le3\Lambda$,
$e_L(l,h)\le1000h_Bs_l$, $c_0(l,h,n')\le1000a\Lambda$ when $1\le a$, $n'\le8a$,
and $c_B(l,h,n',t)\le1000ta\Lambda$ when $n'\le2a$, $t\ge1$.
\end{lemma}
\begin{proof}
(a) The parameter choice gives $\eta_L\le b/(t^3fk^3)\le b$, $t\le\eta_L$, and
$S_L<\ell^*\le b$ where $\ell^*=mb/(\cdots)$ has denominator $\ge m$; so
$S_L,\eta_L,t\le b$ and $h_B\le6b$ (as $b\ge1$). Also $N^{36}t^{150}\le
b^{1/2}$, so $N\le b$. (b) $\log_l(2h_B)\ge0$ since $l\ge2$, $h_B\ge1$; and
$\lfloor\log_ln\rfloor\le\log_l(8h_B)=\log_l4+\log_l(2h_B)\le2+\log_l(2h_B)\le\Lambda+0$ using $s_l\ge1$.
(c) By (b) applied to $2h\le8h_B$ and $N\le8h_B$: $s_l+1\le\Lambda$,
$\lfloor\log_l2h\rfloor\le\Lambda$ and $\lfloor\log_lN\rfloor\le\Lambda$, so
$\lambda_N\le3\Lambda$.
With $\Lambda\le6h_Bs_l$ (from $\log_l(2h_B)\le3h_B$) and $hs_l\le4h_Bs_l$:
$e_L\le64\cdot4h_Bs_l+32+18h_Bs_l\le1000h_Bs_l$. Then
$c_0\le20\cdot9a\cdot3\Lambda\le1000a\Lambda$ and
$c_B\le40\cdot2t\cdot3a\cdot3\Lambda=720ta\Lambda$.
\end{proof}

\begin{lemma}[the split bound]\label{i4:A-lem-bsplit}
Let $c\ge1$, $\Lambda\ge s_0+1$, $y\ge0$, and suppose
$e'\le cy$, $c_0\le cfn_V\Lambda$, $h^W_u\le crfn_V\Lambda$ ($u\in\mathrm{col}$),
$c_B\le ctfn_V\Lambda$, $h_0\le c\Lambda$. Then
$\mathrm{bs}(e',c_0,c_B)\le s_\pi+6cfmy+30cf^2t^2n_V\Lambda+\sum_{0<e<t}\mathrm{sha}_l(e)+v_l(h_E)+t\sum_{v\in Q_k}\ell(\mathcal O_k/(\zeta_v-1))$.
\end{lemma}
\begin{proof}
Put $A=\sum_{\mathrm{col}}n_u$, $A'=\sum_{\mathrm{row}}n_w$; both are at most
$[E:\Q]=2f$, and $|\mathrm{col}|\le A$, $|\mathrm{row}|\le A'$. Also
$0\le\delta\le t\Lambda$ (as $v_l(j!)$-type sums give
$\delta\le t(s_0+1)$). The column sum is at most
$A\,m(2e'+\delta)+|\mathrm{col}|(e'+2\delta+2crfn_V\Lambda)$. With $1\le m\le
f$, $r\le t$, $1\le t\le n_V$, each of the terms $Am\cdot2e'\le4cfmy$,
$|\mathrm{col}|e'\le2cfmy$, $Am\delta\le2cf^2t^2n_V\Lambda$,
$2|\mathrm{col}|\delta\le4(\cdot)$, $2|\mathrm{col}|crfn_V\Lambda\le4(\cdot)$,
$t(A'+2|\mathrm{row}|)c_B\le6(\cdot)$, $4n_V(|\mathrm{col}|+1)\delta\le12(\cdot)$,
$t|\mathrm{col}|h_0\le2(\cdot)$, with $(\cdot)=cf^2t^2n_V\Lambda$, is
elementary; they add up to $6cfmy+30cf^2t^2n_V\Lambda$.
\end{proof}

\begin{proof}[Proof of Proposition~\ref{i4:A-prop-at23rows}(a),(b)]
Fix a setup and a twisted place $k$, $l=l_k$, and $\pi\in Z^{\rm good}_k$.

\emph{(a), $l=2$.} $2^{\nu+2}\mid Q_S$ gives $4\mid Q_S\mid N$ (H4) and $2\in$
small. Each $g_i\equiv1\bmod N$, so $g_i\notin$ small and $l\ne g_i$ (H1);
(H2), (H3), (H5) are immediate. Also $s_0(k)=v_2(Q_S)\ge2$, and
$s_0=s_l$ since $v_2(N)=v_2(Q_S)\ge1$.

\emph{(b), $l=3\notin$ small.} $3\ne q,d$ and $3\ne g_i$ by the
definition of a host; $3\le60\le$ split bound; $3\nmid N$ (the prime
divisors of $N$ are the small primes, $q$ and $d$). Hence $l$ splits in $K$
or is small (H2), places over $3$ are unramified (H3), and
$(\frac{-d}3)=1$ (H5), all by the host axioms for primes $\le$ the split
bound other than $q,d$. As $3\notin$ small, $s_0=1=s_l$.

In both cases Proposition~\ref{i4:A-prop-split} gives a declaration with
(S), (W), (Z). The setup supplies the hypotheses of (S): finite
$s$-length from clause (K1) of the killing condition, finite
$\mathrm{sha}_l(e)$ from the finiteness of the Shafarevich groups (which
covers $l=2$), and (Z). Lemma~\ref{i4:A-lem-h0}
gives $h_0\le s_l+\lfloor\log_lt\rfloor$. Apply
Lemma~\ref{i4:A-lem-numsplit} with $h_1=h_N\le h_B$, $h_2=h_{\rm hub}\le2b+1\le4h_B$,
$h_3=\lfloor2b\rfloor\le4h_B$, $n'=4nn_V=2fn_V$, $a=fn_V\ge1$:
$e'\le1000h_Bs_l$, $c_0\le1000fn_V\Lambda$, $h^W_u\le1000rfn_V\Lambda$,
$c_B\le1000tfn_V\Lambda$, $h_0\le2\Lambda\le1000\Lambda$. By
Lemma~\ref{i4:A-lem-bsplit} with $c=1000$, $y=h_Bs_l$, and (S),
\[
v_{cc}(k,\pi)\le s_\pi+6000fmh_Bs_l+30000f^2t^2n_V\Lambda_l+VV(l)+v_{hE}(l)+t\sum_{v\in Q_k}\ell(\mathcal O_k/(\zeta_v-1)),
\]
and clause (K2) of the killing condition bounds the last term by
$\epsilon(l)$. This is
$\mathrm{Bd}(X.\rho,k,\pi,\mathrm{cost}(L,6000,30000,l))$.
\end{proof}

\begin{proof}[Proof of Proposition~\ref{i4:A-prop-at23rows}(c)]
Write $s=s_l$, $h=h_B$. Since $v_l(N)=v_l(Q_S)$ and $2^{v_l(Q_S)}\le
l^{v_l(Q_S)}\le Q_S\le n^2\le x^{2C_0}$, we get $v_l(Q_S)\le2C_0\log
x/\log2\le3C_0\log x$, so $s\le3C_0\log x$ (using $C_0\log x\ge1$). With
$h\le6b$: $Afmhs\le18AC_0\,fmb\log x\le C\,fmb\log x$. Next
$\log_l(2h)\le2\ell+6$, because $\log(2h)\le\log(12b)\le\ell+3$ and
$\log l\ge1/2$. With $\log x\le x$, $\ell\ge1$, $x>1$:
$\Lambda_l\le3C_0\log x+2\ell+7\le(3C_0+9)x\ell$. Using $f\le k\le r\le t\le x^{C_0}$
(ranks) and $n_V\le2t$:
$Bf^2t^2n_V\Lambda_l\le2B(3C_0+9)\,x^{5C_0+1}\ell^2\le C\,x^{C}\ell^2$, as
$x^{5C_0+1}\le x^C$ and $C\ge2B(3C_0+9)\ge0$. Adding gives
$\mathrm{cost}\le\mathrm{mid}_F(L,C)$.
\end{proof}

\subsection{The automorphic construction in detail}\label{i4:ssec-proof-Aleaves}

This subsection proves Proposition~\ref{i4:prop-assembly-leaves}. We first
describe the objects (\S\ref{i4:A-sss-chain}--\S\ref{i4:A-sss-local}),
then the bookkeeping of rows (\S\ref{i4:A-sss-calc}), then the 26 rows
(\S\ref{i4:A-sss-rows}), and finally the proofs of (A1)--(A4)
(\S\ref{i4:A-sss-A1}--\S\ref{i4:A-sss-A4}).

Throughout, ``eventually'' means: for all values of the staged parameters
$(\varepsilon,C_8,K_0,C,c_D,c_h)$ in the staged range and all scale
contexts with $x\ge x_0$, for some threshold $x_0$. Eventual statements are
closed under finite conjunction: take the minimum of the exponents, the
join of the thresholds and the maximum of the lower bounds. We use this
silently. We write $p_0$ for the parameter system with origin $c_{\rm Dep}$
(no Selmer extension, geometric closure of all geometries) and $p_A$ for
its extension by an evaluation property $A$.

\subsubsection{The even-rank theta-lift chain}\label{i4:A-sss-chain}

Let $D$ be a column datum, $E/F$ its CM extension, $\psi$ a standard
additive character of $F\backslash\mathbb A_F$ and $S$ a finite set of
places of $F$ containing the archimedean ones, the level primes and the
places ramified in $E$. For a hermitian space $V$ over $E$ write
$[U(V)]=U(V)(F)\backslash U(V)(\mathbb A_F)$, with Tamagawa measure.

\begin{definition}[chain data]\label{i4:A-def-chain}
A \emph{chain datum} consists of integers $n\ge1$, $t\ge n+2$, $m\ge1$ and
Hecke characters $\chi_0,\dots,\chi_{m-1}$ of $\mathbb A_E^\times$ that are
conjugate self-dual, unramified outside $S$, and whose archimedean
exponents are pairwise distinct. A \emph{stage} $j\le m$ is a hermitian
space $V_j$ over $E$ with
\[
\dim V_j=n+jt,\qquad \text{Witt index of }V_j<t,
\]
together with a self-dual lattice, an open compact subgroup
$J_j\subset U(V_j)(\mathbb A_{F,f})$ that is hyperspecial outside $S$, and
archimedean discrete-series $K$-types. Let $\mathcal C_j$ be the space of
maps from the $K$-type to cusp forms on $U(V_j)$ of level $J_j$
(Definition~\ref{i4:A-def-cusp}). A \emph{step} $j\to j+1$ is a Weil
representation $\omega_j$ of $U(V_j)\times U(V_{j+1})$ attached to $\psi$
and the splitting pair $(\chi_j,\chi_j)$, with a Fock vacuum vector
$\phi_j$ spherical outside $S$. The \emph{upward lift} and the
\emph{downward map} are
\[
\theta_j^\uparrow F(h)=\int_{[U(V_j)]}\overline{\theta(g,h;\phi_j)}\,F(g)\,dg,
\qquad \theta_j^\downarrow=(\theta_j^\uparrow)^*:\mathcal C_{j+1}\to\mathcal C_j ,
\]
and the \emph{chain map} is $\Theta_k=\theta_0^\downarrow\circ\dots\circ
\theta_{k-1}^\downarrow:\mathcal C_k\to\mathcal C_0$. For a multiset $B$ of
unramified Satake parameters put
\[
\mathrm{seg}(\chi,t)=\{\chi|\cdot|^{(t-1)/2-l}:0\le l<t\},\qquad
B_k=B+\sum_{i<k}\mathrm{seg}(\chi_i,t).
\]
\end{definition}

\begin{definition}[cusp forms]\label{i4:A-def-cusp}
Let $V$ be a nondegenerate hermitian space over $E$. For a nonzero totally
isotropic subspace $X\subset V$ let $N_X\subset U(V)$ be the unipotent
radical of the stabiliser of $X$: the $g$ with $(g-1)X=0$,
$(g-1)V\subset X^\perp$ and $(g-1)X^\perp\subset X$. A function $\phi$ on
$U(V)(\mathbb A_F)$ is \emph{cuspidal} if
$\int_{N_X(F)\backslash N_X(\mathbb A_F)}\phi(ng)\,dn=0$ for every such
$X$ and every $g$. The space $\mathcal C_j$ is spanned by the linear maps
$T$ from the smooth vectors of the archimedean discrete-series model to
functions on $U(V_j)(\mathbb A_F)$ whose values are left
$U(V_j)(F)$-invariant, smooth, $\mathfrak z$-finite, of moderate growth,
cuspidal and right $J_j$-invariant, with $T$ equivariant for the
archimedean action and continuous for the Sobolev topology.
\end{definition}

The point of the chain is the following. A cusp form on the definite
member $U(V_0)$ with prescribed Satake parameters $B$ outside $S$ is
reached as the image of $\Theta_k$ of a cusp form on $U(V_k)$ whose Satake
parameters are $B_k$; the segments $\mathrm{seg}(\chi_i,t)$ are the
standard unramified contribution of a theta lift between spaces whose
dimensions differ by $t$ (the unramified case of
Theorem~\ref{i4:thm-howe}). The chain map is surjective, which is what
makes every prescribed definite form appear in a lattice of automorphic
origin. Surjectivity rests on three facts: an abstract adjoint criterion
(Lemma~\ref{i4:A-lem-adjsurj}), the non-vanishing of the $L$-values that
control each step (Lemma~\ref{i4:A-lem-chainL}), and the Siegel--Weil
formula at each step, whose hypotheses are checked in
Lemma~\ref{i4:A-lem-swrange}.

\begin{lemma}[adjoint surjectivity]\label{i4:A-lem-adjsurj}
Let $A$ be a finite-dimensional complex vector space with a positive
definite hermitian form, $C$ a complex vector space with a positive
definite hermitian form, and $\Theta:C\to A$ linear. Let $\mathcal
P_A\subset A\times I$ and $\mathcal P_C\subset C\times J$ be labelling
relations, $\sim_A$ an equivalence relation on $I$, $\sim_C$ a relation on
$J$, and $T:I\to J$ a map, such that:
\begin{itemize}
\item[(i)] the labelled vectors span $A$;
\item[(ii)] if $(u,B),(u',B')\in\mathcal P_A$, $B\sim_AB'$ and $u+au'\ne0$,
then $(u+au',B)\in\mathcal P_A$; and $(au,B)\in\mathcal P_A$ for $a\ne0$;
\item[(iii)] $T(B)\sim_CT(B')$ implies $B\sim_AB'$;
\item[(iv)] vectors of $C$ with pairwise non-$\sim_C$ labels are linearly
independent;
\item[(v)] for each $(u,B)\in\mathcal P_A$ there is $x\in C$ with
$(x,T(B))\in\mathcal P_C$ and $\langle y,x\rangle_C=\langle\Theta
y,u\rangle_A$ for all $y\in C$.
\end{itemize}
Then $\Theta$ is surjective.
\end{lemma}
\begin{proof}
Suppose $W=\Theta(C)\ne A$. As $A$ is finite dimensional and the form is
definite, there is $u\ne0$ orthogonal to $W$. By (i) and (ii),
$u=\sum_r v_r$ with $(v_r,B_r)\in\mathcal P_A$. Group the terms by the
$\sim_A$-class of the label and add them one at a time inside each class
using (ii), discarding a partial sum when it vanishes. Each class
contributes $0$ or a nonzero labelled vector, so $u=\sum_{i=1}^su_i$ with
$(u_i,B_i)\in\mathcal P_A$, $u_i\ne0$, $B_i\not\sim_AB_{i'}$ for $i\ne i'$,
and $s\ge1$. By (v) choose $x_i$ labelled $T(B_i)$ with $\langle
y,x_i\rangle_C=\langle\Theta y,u_i\rangle_A$. By (iii) the labels $T(B_i)$
are pairwise non-$\sim_C$, so by (iv) the $x_i$ are independent and
$x=\sum x_i\ne0$. But $\langle y,x\rangle_C=\langle\Theta y,u\rangle_A=0$
for every $y$; with $y=x$ this gives $x=0$, a contradiction.
\end{proof}

In the application, $A=\mathcal C_0$ (finite dimensional by
Harish-Chandra's finiteness theorem), $C=\mathcal C_k$, the labels are
multisets of Satake parameters outside $S$, $\sim_A$ and $\sim_C$ are
equality of labels, $T(B)=B_k$, (iv) is linear independence of Hecke
eigenforms with distinct eigensystems, and (v) is the statement that the
upward chain lift of an eigenform with parameters $B$ is nonzero with
parameters $B_k$ and is adjoint to $\Theta_k$. Condition (iii) is the
injectivity in Lemma~\ref{i4:A-lem-chainL}; the nonvanishing in (v) is
the Rallis inner product formula, whose $L$-value is the one of
Lemma~\ref{i4:A-lem-chainL} and whose Eisenstein side is identified by
Siegel--Weil.

\begin{lemma}[$L$-values along the chain]\label{i4:A-lem-chainL}
For every chain datum, every $k<m$ and every $B$ with
$|\alpha_v|\le q_v^{(n-1)/2}$ for all $\alpha\in B$ and all $v\notin S$, the
partial $L$-function $L^S(s+\frac12,B_k\times\chi_k^{-1})$ is holomorphic
and nonzero at $s=t/2$. Moreover $B\mapsto B_k$ is injective.
\end{lemma}
\begin{proof}
By definition of $B_k$,
\[
L^S(s+\tfrac12,B_k\times\chi_k^{-1})=L^S(s+\tfrac12,B\times\chi_k^{-1})
\prod_{i<k}\prod_{l<t}L^S\bigl(s+\tfrac12+\tfrac{t-1}2-l,\chi_i\chi_k^{-1}\bigr).
\]
At $s=t/2$ each local factor of the first term has
$|\alpha\chi_k^{-1}(\varpi_v)q_v^{-(t+1)/2}|\le q_v^{(n-1)/2-(t+1)/2}\le
q_v^{-2}$, because $t\ge n+2$; so the Euler product converges absolutely
and is nonzero. In the second term the arguments are the integers $t-l\ge1$.
For $i\ne k$ the archimedean exponents of $\chi_i$ and $\chi_k$ differ, so
$\chi_i\chi_k^{-1}$ is a Hecke character whose archimedean components are
not all of the form $|\cdot|^{iy}$, so it is not of the form $|\cdot|^{iy}$;
its completed $L$-function is entire (Theorem~\ref{i4:thm-hecke}) and nonzero on
$\Rel s\ge1$ (non-vanishing of Hecke $L$-functions on $\Rel s\ge1$;
reference to be confirmed). Removing the finitely many Euler factors at $S$
creates no zero on $\Rel s\ge1$, since each removed factor is
$(1-aq_v^{-s})^{-1}$ with $|a|=1$. Injectivity: $B_k$ determines $B$ by
cancelling the fixed multiset $\sum_{i<k}\mathrm{seg}(\chi_i,t)$.
\end{proof}

\begin{lemma}[Siegel--Weil ranges along the chain]\label{i4:A-lem-swrange}
For a chain datum and a step $j\to j+1$ put $n_j=\dim V_j$,
$m_j=\dim V_{j+1}=n_j+t$ and $r_j$ the Witt index of $V_{j+1}$. Then
$n_j<m_j$ and $m_j-r_j>n_j$ (Weil's convergence condition). Moreover
$m_0>2n_0$, and $n_j<m_j\le 2n_j$ for $1\le j<m$.
\end{lemma}
\begin{proof}
$m_j-n_j=t\ge3$, and $r_j<t$ by Definition~\ref{i4:A-def-chain}, so
$m_j-r_j>m_j-t=n_j$. For $j=0$: $m_0=n+t\ge2n+2>2n_0$. For $j\ge1$:
$n_j=n+jt\ge t$, hence $m_j=n_j+t\le2n_j$.
\end{proof}

\begin{remark}[on the Siegel--Weil input]\label{i4:A-rem-sw}
The Siegel--Weil formula is used for the doubled pair
$(U(V_j\oplus(-V_j)),U(V_{j+1}))$, that is, with $W=V_j$ and $V=V_{j+1}$,
$n=n_j$, $m=m_j$, $r=r_j$. Lemma~\ref{i4:A-lem-swrange} gives $n_j+r_j<m_j$
at every step, which is the convergence hypothesis of both
Theorems~\ref{i4:thm-sw} and~\ref{i4:thm-sw2}. At the first step
$m_0>2n_0$, so Theorem~\ref{i4:thm-sw} applies. At the later steps
$n_j<m_j\le2n_j$, so Theorem~\ref{i4:thm-sw2} applies; there the formula
is used only for Fock vectors, which is the class of vectors it covers.
Both identities are with the mean $I(g,f)$, so no choice of measure on
$[U(V_{j+1})]$ enters. Both theorems remain ``reference to be confirmed''.
\end{remark}

\begin{remark}[on Howe duality]\label{i4:A-rem-howe}
Only the unramified clause of Theorem~\ref{i4:thm-howe} is used: at a
place $v\notin S$ the lift of an unramified representation with Satake
multiset $B$ along a step of dimension gap $t$ is unramified with Satake
multiset $B+\mathrm{seg}(\chi_{j,v},t)$. The development proves this
clause internally (it is leaf (L6) of Proposition~\ref{i4:A-prop-inputs});
the full Howe duality statement is not used.
\end{remark}

\subsubsection{Local components}\label{i4:A-sss-local}

The cusp forms produced by the chain have the following local components.

\emph{Unramified places} $v\notin S$. The group $J_{j,v}$ is hyperspecial
and the local component is unramified; its Satake multiset is $B_j$ as in
Definition~\ref{i4:A-def-chain}. Finiteness of the possible cuspidal
supports of a given level uses Theorem~\ref{i4:thm-bzfin}; their
uniqueness is Theorem~\ref{i4:thm-bz}.

\emph{Level primes} $g$. At a level prime the residue field is $\mathbb
F_g$, and the local data are depth-zero types: inflations to a maximal
compact subgroup of cuspidal representations of $\mathrm{GL}_d(\mathbb
F_g)$. These are labelled by Green's parametrisation, in the following
precise form.

\begin{theorem}[Green, Trans.\ AMS 80 (1955) 402--447; reference to be confirmed]\label{i4:A-thm-green}
Let $g$ be a prime, $d>0$, and $K$ the field with $g^d$ elements. There is a
map $\theta\mapsto\tau(\theta)$ from characters of $K^\times$ to class
functions on $\mathrm{GL}_d(\mathbb F_g)$ such that: for $\theta$ regular
($\theta^{g^i}\ne\theta^{g^j}$ for $0\le i<j<d$), $\tau(\theta)$ is a
cuspidal irreducible character of degree $\prod_{i=1}^{d-1}(g^i-1)$, and
$\tau(\theta)(y)=(-1)^{d-1}\sum_{i<d}\theta(z^{g^i})$ whenever the
characteristic polynomial of $y$ is irreducible with root $z\in K$; for
regular $\theta,\theta'$, $\tau(\theta)=\tau(\theta')$ if and only if
$\theta'=\theta^{g^i}$ for some $i<d$; and every cuspidal irreducible
character of $\mathrm{GL}_d(\mathbb F_g)$ is $\tau(\theta)$ for some regular
$\theta$.
\end{theorem}
\noindent This is the precise form of Theorem~\ref{i4:thm-green}.
Hypotheses: none beyond $d>0$. Cuspidality is meant in the usual sense of
vanishing Harish-Chandra restriction to every proper standard Levi
subgroup, which is the sense used in the development.

\begin{lemma}[Green's parametrisation by prime labels]\label{i4:A-lem-greenlab}
Fix an injective multiplicative map $t:\overline{\mathbb F}_g^\times\to\C^\times$
and, for each $d>0$, an element $u_d\in\overline{\mathbb F}_g$ of order
$g^d-1$. For $x\in\overline{\mathbb F}_g^\times$ let $\theta_x(u_d^n)=t(x)^n$.
For $d>0$ there is a map $P\mapsto\tau_P$ from monic polynomials over
$\mathbb F_g$ to class functions on $\mathrm{GL}_d(\mathbb F_g)$ such that,
for every prime label $P$ of degree $d$ (monic, irreducible, $P(0)\ne0$),
$\tau_P$ is cuspidal of degree $\prod_{i=1}^{d-1}(g^i-1)$ with
$\tau_P(y)=(-1)^{d-1}\sum_{i<d}\theta_x(z^{g^i})$ for every root $x$ of $P$
and every $y$ with irreducible characteristic polynomial and root $z$;
$P\mapsto\tau_P$ is injective on prime labels of degree $d$; and every
cuspidal irreducible character is some $\tau_P$.
\end{lemma}
\begin{proof}
Let $K\subset\overline{\mathbb F}_g$ have $g^d$ elements; $K^\times$ is
cyclic, generated by $u_d$. For $x\in K^\times$ the character
$\mathrm{th}(x)=\theta_x|_{K^\times}$ is well defined, $x\mapsto\mathrm{th}(x)$
is injective (as $t$ is) and onto the characters of $K^\times$ (a
character is fixed by its value at $u_d$, a $(g^d-1)$-th root of unity, and
$t$ maps the $g^d-1$ roots of $X^{g^d-1}-1$ injectively into the $g^d-1$
such roots of unity), and $\mathrm{th}(x^m)=\mathrm{th}(x)^m$. Hence
$\mathrm{th}(x)$ is regular if and only if the $x^{g^i}$, $i<d$, are pairwise
distinct. For a prime label $P$ of degree $d$ fix a root $x_P\in K^\times$;
its conjugates are pairwise distinct, so $\mathrm{th}(x_P)$ is regular; put
$\tau_P=\tau(\mathrm{th}(x_P))$. Cuspidality and degree are given by
Theorem~\ref{i4:A-thm-green}. The value formula for another root
$x=x_P^{g^j}$ follows because the sum over the Frobenius orbit of $z$ does
not depend on the chosen root and $\theta_x(w^m)=\theta_{x^m}(w)$.
Injectivity: $\tau_P=\tau_{P'}$ gives $\mathrm{th}(x_{P'})=\mathrm{th}(x_P^{g^i})$,
so $P,P'$ share a root and are equal. Exhaustion: a cuspidal $\chi$ is
$\tau(\mathrm{th}(x))$ with $\mathrm{th}(x)$ regular; the minimal polynomial
$P$ of $x$ has degree $k\mid d$, and $k<d$ would contradict regularity; so
$P$ is a prime label of degree $d$, $x_P=x^{g^j}$, and $\tau_P=\chi$ by the
injectivity clause of Theorem~\ref{i4:A-thm-green}.
\end{proof}

The series idempotents of $\mathrm{GL}_r(\mathbb F_g)$ that cut out the
local types are sums of the central idempotents
$z_\chi=\frac{\chi(1)}{|G|}\sum_{x\in G}\chi(x^{-1})x$ over the irreducible
characters with a given Harish-Chandra label; their orthogonality and
completeness are Theorem~\ref{i4:thm-serre}. The Deligne--Lusztig
parametrisation is not used.

\emph{Archimedean places.} At each real place of $F$, $U(V_j)$ has a
signature $(p,q)$, and the archimedean component is a discrete series of
$U(p,q)$ with a given Harish-Chandra parameter $\lambda\in\R^{p+q}$ and
signs $\varepsilon$. The inputs are Harish-Chandra's classification and
existence of discrete series and the uniqueness of the discrete series with
given parameter (via the lowest $K$-type). For $V_0$ totally definite the
group is compact at infinity and the $K$-type is a finite-dimensional
representation. These inputs enter through
Proposition~\ref{i4:A-prop-inputs}(a); see Remark~\ref{i4:A-rem-pub}.

\begin{remark}[published inputs]\label{i4:A-rem-pub}
The development isolates the following published statements as separate
propositions, which intermediate formal lemmas take as arguments (``the
third-place hypotheses''), and proves each of them inside the development: global Artin reciprocity and Weber's ray
class field theorem; existence of CM elliptic curves with good reduction
outside a prime; the theory of Harish-Chandra modules; existence of
discrete-series covers; Harish-Chandra's finiteness of automorphic forms; rapid
decay of cusp forms; the Casselman--Wallach globalisation; the
classification and existence of discrete series; the formal degree formula;
the Siegel--Weil formulas of Theorems~\ref{i4:thm-sw}
and~\ref{i4:thm-sw2}; the Siegel--Weil formula in
the Schr\"odinger model; meromorphic continuation of Siegel Eisenstein
series; unramified doubling; Bernstein--Zelevinsky finiteness
(Theorem~\ref{i4:thm-bzfin}); Schur
orthogonality; Dixmier--Malliavin; Flath's tensor product theorem;
Zelevinsky's classification; the geometric lemma; Green's cuspidality
criterion; and the local Langlands correspondence for $\mathrm{GL}_n$, used
only in depth zero and for characters. Of these, only
Theorems~\ref{i4:thm-hecke}, \ref{i4:thm-bz}, \ref{i4:thm-bzfin},
\ref{i4:thm-serre}, \ref{i4:thm-sw}, \ref{i4:thm-sw2}, \ref{i4:thm-howe}
and \ref{i4:A-thm-green} are stated in
this paper; for the others the references are to be confirmed. Green's cuspidality criterion as used (a
character of $\mathrm{GL}_d(\mathbb F_g)$ is cuspidal if and only if its
averages over the unipotent radicals of the standard maximal parabolics
vanish) is a definition-level reformulation and is not
Theorem~\ref{i4:A-thm-green}; we state it separately (it too is
proved in the formal development).
\end{remark}

\subsubsection{The row calculus}\label{i4:A-sss-calc}

A \emph{slot} is a property of core points; a \emph{linkage} is a property
$\mathrm{Lk}(\beta,c,b)$ of a built point $\beta$, a real $c$ and a bundle
$b$ over $(\beta.D,\beta.O,c)$. A \emph{selection} on a family class
$\mathcal A$ is a choice of bundles over the points of the members of
$\mathcal A$.

\begin{definition}[rows]\label{i4:A-def-row}
Let $Q$ be a slot, $\mathrm{Lk}$ a linkage and $H$ a property of pairs
$(\mathcal A,\sigma)$ (a family class and a selection).
\begin{enumerate}
\item $\mathrm{Row}(Q,\mathrm{Lk},H)$ means: there is a bundle property $E$
such that eventually, for every built point $\beta$ in slot $Q$, every $c$
and every $b$ with $\mathrm{Lk}(\beta,c,b)$, $E(b)$ holds; and
$H(\mathcal A,\sigma)$ holds for every pair in which every member of
$\mathcal A$ has a linked parameter all of whose points are built in slot
$Q$, linked by $\mathrm{Lk}$, with selected bundles satisfying $E$.
\item For a base condition $KP$ and slots $Q,Q'$, and a property $r$ of
slots, $\mathrm{Lg}(KP,Q,Q';r)$ means: there is a slot $R$ with
$\mathrm{LargeBase}(KP,Q,R)$ and $r(Q'\wedge R)$. Here
$\mathrm{LargeBase}(KP,Q,R)$ says that eventually, for every base
satisfying $KP$, every core point in slot $Q$ is in slot $R$.
\end{enumerate}
A \emph{row statement} at a parameter system $p$ is
$\mathrm{Lg}(KP_x,\mathrm{Built}(p),\mathrm{Good}(p);\mathrm{Row}(\cdot,\mathrm{Lk}_p,H))$,
with $\mathrm{Lk}_p$ the bundle linkage of $p$.
\end{definition}

\begin{lemma}[row calculus]\label{i4:A-lem-calc}
\begin{enumerate}
\item[(a)] $\mathrm{LargeBase}(KP,Q,\mathrm{True})$ holds; so $r(Q'\wedge
\mathrm{True})$ implies $\mathrm{Lg}(KP,Q,Q';r)$.
\item[(b)] If $\mathrm{Lg}(KP,Q,Q';r)$ and $r(S)\Rightarrow r'(S)$ for every
slot $S$, then $\mathrm{Lg}(KP,Q,Q';r')$.
\item[(c)] $\mathrm{Lg}(KP,Q,Q';\mathrm{Row}(\cdot,\mathrm{Lk},H))$ and
$\mathrm{Lg}(KP,Q,Q';\mathrm{Row}(\cdot,\mathrm{Lk},H'))$ imply
$\mathrm{Lg}(KP,Q,Q';\mathrm{Row}(\cdot,\mathrm{Lk},H\wedge H'))$.
\item[(d)] If $H(\mathcal A,\sigma)$ holds for all pairs, then
$\mathrm{Row}(Q,\mathrm{Lk},H)$ holds for all $Q,\mathrm{Lk}$.
\end{enumerate}
\end{lemma}
\begin{proof}
(a) Every core point is in slot $\mathrm{True}$. (b) Keep $R$. (d) Take
$E=\mathrm{True}$. (c) Let $(R,E)$ and $(R',E')$ be witnesses and
$R''=R\wedge R'$; $\mathrm{LargeBase}(KP,Q,R'')$ is the conjunction of two
eventual statements. A built point in $Q'\wedge R''$ lies in $Q'\wedge R$
and in $Q'\wedge R'$, so eventually its linked bundles satisfy $E\wedge
E'$; a pair selected for $(Q'\wedge R'',E\wedge E')$ is selected for both
$(Q'\wedge R,E)$ and $(Q'\wedge R',E')$, since rows are monotone in the
slot. Hence $H\wedge H'$.
\end{proof}

\begin{lemma}[pointwise rows]\label{i4:A-lem-row0}
Let $Q$ be a slot and $\mathcal R(D,O)$ a property of a datum and an
object, and suppose that eventually every built point $\beta$ in slot $Q$
satisfies $\mathcal R(\beta.D,\beta.O)$. Let $H_{\mathcal R}(\mathcal
A,\sigma)$ say: for every family $F$ in the closure of $\mathcal A$ there is
$x_0$ such that $\mathcal R$ holds at the datum and object of every point
of $F$ at every scale $x\ge x_0$. Then $\mathrm{Row}(Q,\mathrm{Lk},
H_{\mathcal R})$ holds for every linkage $\mathrm{Lk}$, with closedness
threshold $0$ and $\varepsilon_1=1$.
\end{lemma}
\begin{proof}
Take $E(b)$ to be $\mathcal R$ at the point underlying $b$. The eventual
clause is the hypothesis. In a selected pair every member family has,
eventually, only built points of slot $Q$ satisfying $E$; so along every
family in the closure $\mathcal R$ holds eventually, which is
$H_{\mathcal R}$. Since $E$ ignores the bundle, the linkage plays no role.
\end{proof}

\begin{lemma}[bundle linkage gives reciprocity linkage]\label{i4:A-lem-linkR}
If $b$ is linked to a built point $\beta$ by the bundle linkage of $p$,
then its record $b.R$ is linked to $\beta$ by the reciprocity linkage.
\end{lemma}
\begin{proof}
The bundle linkage gives a witness point $q$, a bundle $b'$ over $q$ with a
bundle witness $w$ and an equality $(\beta.D,\beta.O,b)=(q.D,q.O,b')$. The
reciprocity linkage asks for $q$, a record compatible with $q$ and the
equality $(\beta.D,\beta.O,b.R)=(q.D,q.O,b'.R)$. Take $b'.R$ with the
compatibility recorded in $w$, and apply $(D,O,b)\mapsto(D,O,b.R)$.
\end{proof}

A core point is \emph{over true columns} if its construction record carries
a column view: a base, a character frame, a twisting datum, a weighted
column family $w$ with level $Z$, the object, a norm choice, and an
identification of the twist of the object with the twist read off the
true columns of $w$ at $Z$, together with the base sign condition and the
real tie of $w.D(Z)$.

\begin{lemma}\label{i4:A-lem-overcols}
Good cores lie over true columns.
\end{lemma}
\begin{proof}
A good core carries a built record as its first component, and a built
record yields a column view (Proposition~\ref{i4:A-prop-rowleaves}(P0)).
\end{proof}

\subsubsection{The 26 rows}\label{i4:A-sss-rows}

The conclusion (b) of Theorem~\ref{i4:thm-assembly} is that each of the
following 26 statements is a row statement at $p$ (Definition~\ref{i4:A-def-row})
for the bundle linkage $\mathrm{Lk}_p$, with base condition $KP_x$, slots
$\mathrm{Built}(p)$, $\mathrm{Good}(p)$, a norm property $NN$ and the
definite-side condition $\mathrm{DR}$. In each row, $\mathcal A$ is a
family class, $\sigma$ a selection, $b$ a selected bundle over a point
$(D,O)$ of a member family at scale $x$, and ``closed'' means closed under
the closure operation of family classes with the indicated $\varepsilon_1$
and threshold datum. Table~\ref{i4:A-tab-rows} lists the rows with their
exact content and sources; ``U'' marks a universal row (it holds for every
slot and linkage). Status: P = proved here, R = reduced here to a smaller
proposition, O = open row (from (A2)), H = hypothesis of the row theorem.

{\small
\begin{longtable}{>{\raggedright}p{0.4cm}>{\raggedright}p{2.0cm}>{\raggedright}p{5.8cm}>{\raggedright\arraybackslash}p{2.5cm}}
\caption{The 26 rows.}\label{i4:A-tab-rows}\\
\toprule
\# & Row & Exact statement $H(\mathcal A,\sigma)$ & Source, status\\
\midrule
\endfirsthead
\toprule
\# & Row & Exact statement & Source, status\\
\midrule
\endhead
1 & class count & for every family $F$, $x>x_1(F)$, every $i$ and every word $u$ satisfying the bigness and splitting conditions: $\#\mathrm{cls}_0(u)=\mathrm{mass}_I(D,k)\,\mathrm{idx}(g_0,g_1)\,(T_u^k-1)^f$ at $D=D_F(x,i)$; and $L_a\ge0$ for $x\ge x_1(F)+1$; closed with $\varepsilon_1=1$, threshold $0$ & Lemma~\ref{i4:A-lem-classnum}; U; R\\
2 & $j$-invariant & at every selected bundle, $\iota(V_0(\mu_0))$ is the whole target of the realisation map, for $x>x_1(F)+1$ & from row 14; R\\
3 & instance rank & at every point: adelic quotient volumes of hermitian spaces of equal rank do not depend on the signature, and genus mass, class finiteness and class mass hold at rank $\mathrm{rk}\,Y_0$ & Siegel--Weil; U; R\\
4 & second complement & on classes whose members satisfy the $T_u$-size and $D_\Gamma$-size conditions, the family second-complement inputs hold; closed with $\varepsilon_1=1/100$ & U; R\\
5 & level-prime valuations & at bundle points: for every word $u$, the sharp level-prime column valuations and the twisted column bound at $u$ & over true columns; R\\
6 & arithmetic & at bundle points: a constant $C_1$ with the used-tuple bound and the family size bound at type primes, eventually & delivery over true columns; R\\
7 & reference chain & at bundle points: for every word $u$ the reference-chain word gives a forms map & over true columns; R\\
8 & twisted point & at bundle points: the twisted point of $O$ is integral & over true columns; R\\
9 & period at zero & at bundle points: the period class-sum identity (integrality of the reference period at $0$) & over true columns; R\\
10 & twisted period & at bundle points: the twisted periods of $O$ are admissible & over true columns; R\\
11 & real tie & at bundle points: the real-side tie holds and there are symmetric column forms with finite character at step zero & over true columns; R\\
12 & bounded rows & the indefinite realisation of every selected bundle has bounded rows & rapid decay, discrete series, Casselman--Wallach; U; R\\
13 & twist & every twisting place of the realisation of $b$ lies in the support $L.\mathrm{supp}$ of the family & from the tie (open statement); P\\
14 & primitivity $\mathrm{hPA}$ & the realised rows of $b$ are primitive and type separated & open; O\\
15 & $\Lambda'$ & the $\Lambda'$-positivity rows at $b$ & open; O, from (O4), (O8)\\
16 & divisor & the divisor condition at $b$ & open; O, from (O3), (O8)\\
17 & decomposition point & the decomposition-point condition at $b$ & open; O, from (O8)\\
18 & selection jump & the $J$-selection condition at $b$ & open; O, from (O9)\\
19 & $L^p$ & the $L^p$ bound on $D_3(3,\mathcal A)$ & Siegel--Weil; U; R\\
20 & order & the order bound (off-diagonal carrier gap) on $D_3(3,\mathcal A)$ & Siegel--Weil; U; R\\
21 & service & the service condition on $D_3(3,\mathcal A)$ for the kernels $b.R.T.\mathrm{ker}$ & Artin reciprocity; R\\
22 & zero locus & for every $k<m(D)$ the lifted projection at $(b.P.A_3,P_0,\psi_0,Z_s,k,\mu_0)$ is surjective & bundle fields; U; R\\
23 & torsion frame & the second Siegel--Weil row on $D_3(3,\mathcal A)$ & Siegel--Weil; U; R\\
24 & counting & counting constants $c_{30}$ (with $C\ge0$, $A\ge2(C_{YL}+1)$), $c_{Ic}>0$, and the irreducibility row & from the label row; R\\
25 & definite rows & the definite-side condition $\mathrm{DR}$ assembled from point summands & from rows 14, 16 and the open rows; R\\
26 & norm & $NN(D_3(3,\mathcal A))$ & hypothesis; H, see (A3)\\
\bottomrule
\end{longtable}
}

The five rows 14--18 are among the twelve open parameter statements of
Theorem~\ref{i4:thm-assembly}(c), which are proved in
\S\ref{i4:ssec-proof-assembly} from (A2). The other open statements (tie,
labels, goodness, geometry, multiplicity, $\kappa$, definite forms) are
not rows of the list but are inputs to rows 13, 24 and 25.

\begin{theorem}[all rows]\label{i4:A-thm-allrows}
Assume global Artin reciprocity, Theorem~\ref{i4:thm-sw2}, the cover and
the classification of discrete series, Harish-Chandra modules, rapid decay
of cusp forms and Casselman--Wallach, surjectivity onto $X$ and the lower
bound $\lambda_{\rm low}$. Let $p=p(N,S,\Gamma)$ be a parameter system and
$NN$ a property of family classes. Assume the open parameter statements at
$p$, and that the covolume bound, the functional condition and the norm
condition for $NN$ are linked at $p$. (At the system $p$ of
Theorem~\ref{i4:thm-assembly} these three linkage hypotheses are
discharged by (A3) and (A4) of
Proposition~\ref{i4:prop-assembly-leaves}; surjectivity onto $X$ and
$\lambda_{\rm low}$ are Proposition~\ref{i4:A-prop-funcov}(d).) Then the 26
rows of Table~\ref{i4:A-tab-rows} are row statements at $p$, modulo
Propositions~\ref{i4:A-prop-rowleaves}.
\end{theorem}

\begin{proof}
We treat the rows in the order of the table. ``Universal'' rows are
obtained from Lemma~\ref{i4:A-lem-calc}(d) or Lemma~\ref{i4:A-lem-row0}
followed by Lemma~\ref{i4:A-lem-calc}(a).

\emph{Row 1.} Part (ii) is proved below; part (i) is
Lemma~\ref{i4:A-lem-classnum}(b). Both hold for all pairs, so
Lemma~\ref{i4:A-lem-calc}(d) applies. For (ii): $L_a$ is the logarithm of a
product of a scale factor, which is at least $1$ for every size datum, and
a size factor, which is at least $1$ by the size conditions of the family
at $x>x_1(F)$.

\emph{Row 2.} By Lemma~\ref{i4:A-lem-calc}(b) applied to row 14: from
primitivity and type separation of the realised rows of $b$ one forms the
realisation datum of $b$ with unit constant and the column condition of
$b$, and the row of realisations gives surjectivity of $\iota$ on
$V_0(\mu_0)$ (Proposition~\ref{i4:A-prop-rowleaves}(R2)).

\emph{Row 3.} Proposition~\ref{i4:A-prop-rowleaves}(R4) gives the volume
ratio at every column datum. Let $W,W_0$ be hermitian spaces of equal rank
with constant signatures. Unfolding the volume ratio, each adelic unitary
volume equals the mass of the standard lattice of that rank, for every
signature ((R5)); hence the two volumes agree. The class items hold at
every rank ((R6)). This is pointwise, so Lemma~\ref{i4:A-lem-row0} applies.

\emph{Row 4.} For $0<\varepsilon\le1/100$ we have $\varepsilon<1/52$, and
for such $\varepsilon$ the second-complement inputs follow from the
$T_u$-size and $D_\Gamma$-size conditions ((R7)). The working class has
both, so Lemma~\ref{i4:A-lem-row0} applies with $\varepsilon_1=1/100$.

\emph{Rows 5, 7--11.} By Lemma~\ref{i4:A-lem-calc}(a) it suffices to prove
each for the slot $\mathrm{Good}\wedge\mathrm{True}$, which is over true
columns by Lemma~\ref{i4:A-lem-overcols}. Each row is
$\mathrm{Row}(\cdot,\mathrm{Lk},H_{\mathcal R_i})$ for the pointwise
property $\mathcal R_i$ displayed in the table. By
Lemma~\ref{i4:A-lem-row0} it suffices that eventually every built point
in the slot satisfies $\mathcal R_i$. Such a point has a column view, and
transporting the twist along the identification of the view reduces
$\mathcal R_i(\beta.D,\beta.O)$ to the same statement at the datum
$w.D(Z)$ and the object read off the true columns, which is
Proposition~\ref{i4:A-prop-rowleaves}(P1)--(P6).

\emph{Row 6.} As for row 5, with $\mathcal R$ the arithmetic row at bundle
points: it holds for every slot whose cores have a part with the row pin
for all $a$ ((R8)); a good core is built, a built record yields a part, and
the row pin follows from the reading component of the built record ((R9)).

\emph{Row 12.} For every selected bundle $b$ at a point of $F$, the theta
equality at indefinite scales ((R10)) applied to $b.B$ and the lower
bound on $n_V$ of $F$ give, by the right theta bound ((R11)), the bounded
rows of the indefinite realisation. This holds for all pairs;
Lemma~\ref{i4:A-lem-calc}(d).

\emph{Row 13.} Lemma~\ref{i4:A-lem-twist} below, via
Lemma~\ref{i4:A-lem-calc}(b) applied to the tie (every selected bundle has
a comparison datum), which is an open statement.

\emph{Rows 14--18.} Hypotheses (open parameter statements).

\emph{Rows 19, 20, 23.} For every selection restricted to
$D_3(3,\mathcal A)$: Theorem~\ref{i4:thm-sw2} gives both Siegel--Weil rows
((R12)); the $L^p$ bound follows from finite dimensionality of row spaces
and the volume ratio ((R13)); the order bound from the off-diagonal
carrier-gap estimate fed with the first Siegel--Weil row and the $L^p$
bound ((R14)); the torsion-frame row is the second Siegel--Weil row.
Lemma~\ref{i4:A-lem-calc}(d),(a).

\emph{Row 21.} The bundle linkage implies the reciprocity linkage
(Lemma~\ref{i4:A-lem-linkR}). Artin reciprocity gives the
reciprocity-kernel row for the reciprocity linkage ((R15)). A selection
for $\mathrm{Lk}_p$ induces, along $b\mapsto b.R$, a selection for the
reciprocity linkage; so the kernel row holds for $\mathrm{Lk}_p$, for the
kernels $b.R.T.\mathrm{ker}$. The closedness datum $(a,\varepsilon_1)$ and
the kernel bound are the two components of the service condition,
restricted from $\mathcal A$ to $D_3(3,\mathcal A)$.

\emph{Row 22.} For each bundle and each $k<m(D)$, the block hypothesis
(conjugation, irreducibility, lifts) follows from the family condition of
$b$, the forms tie is a field of $b$, and the block lemma ((R3)) gives
surjectivity. This holds for all pairs; Lemma~\ref{i4:A-lem-calc}(d).

\emph{Row 24.} Siegel--Weil gives the counting constants $c_{30}$ and
$c_{Ic}$ ((R16)), and with finite dimensionality the irreducibility row
from concrete automorphic data on classes with $T_u$-size ((R17)), fed
with the label ties of the label row (open statement); keep the slot $R$
of the label row; Lemma~\ref{i4:A-lem-calc}(b),(c).

\emph{Row 25.} By Lemma~\ref{i4:A-lem-calc}(c), combine the rows
primitivity, divisor, goodness, geometry, covolume, functional,
multiplicity, $\kappa$ and differential into one row; then apply
Lemma~\ref{i4:A-lem-calc}(b) with the assembly of definite rows ((R18)):
covolume, functional, multiplicity, $\kappa$ and the definite-side datum
are conjuncts; the point summands are chosen from the geometry row; and
smallness, the norm bound and the $L_6$ bound are derived from
primitivity, the divisor row, good reduction, surjectivity onto $X$ and
$\lambda_{\rm low}$ ((R19)--(R25)).

\emph{Row 26.} Hypothesis.
\end{proof}

\begin{lemma}[double cosets, part (a) of the class number formula]\label{i4:A-lem-classnum}
\begin{enumerate}
\item[(a)] Let $\Gamma,U\le K$ be subgroups of a group $G$ with $[K:U]$
finite and nonzero. Suppose $\Gamma\backslash G/K$ is finite, all
$\Gamma\cap yKy^{-1}$ are finite, and put
$m=\sum_{c\in\Gamma\backslash G/K}1/\#\mathrm{Stab}(c)$. If no
$\gamma\ne1$ in $\Gamma$ satisfies $x^{-1}\gamma x\in U$ for some $x\in G$,
then $\#(\Gamma\backslash G/U)=[K:U]\cdot m$.
\item[(b)] For every column datum $D$, level datum $L_v$ and word $u$
satisfying the conditions of row 1,
$\#\mathrm{cls}_0(u)=\mathrm{mass}_I(D,k)\cdot\mathrm{idx}(g_0,g_1)\cdot
(T_u^k-1)^f$.
\end{enumerate}
\end{lemma}
\begin{proof}
(a) The double coset $\Gamma yK$ is the union of the $U$-double cosets
$\Gamma ykU$, $k\in K$, and these correspond to the orbits of the finite
group $\mathrm{Stab}(y)=\Gamma\cap yKy^{-1}$ (acting through
$y^{-1}\cdot y$) on $K/U$. The freeness hypothesis says that these actions
are free, so $\Gamma yK$ contains $[K:U]/\#\mathrm{Stab}(y)$ double cosets
of $U$. Summing over $\Gamma\backslash G/K$ gives $[K:U]\,m$.

(b) The self-dual-over-base statement ((R1a) of
Proposition~\ref{i4:A-prop-rowleaves}) gives a self-dual
$\mathcal O_E$-lattice $N$ in $E^k$ whose adelic stabiliser contains the
level group $U$ of $L_v(u)$. The stabiliser index formula ((R1b)) gives
$[\widehat{\mathrm{Stab}}(N):U]=\mathrm{idx}(g_0,g_1)\,(T_u^k-1)^f$,
positive because each local factor is a ratio of positive cardinalities.
Classes are double cosets with coset mass $\mathrm{mass}(N)$, and
$\mathrm{mass}(N)=\mathrm{mass}_I(D,k)$ by Theorem~\ref{i4:thm-genusmass}.
Column stabilisers are trivial ((R1c)), which is the freeness hypothesis of
(a) for $\Gamma=U(V_0)(F)$. Part (a) and the bijection of
$\mathrm{cls}_0(u)$ with $\Gamma\backslash G/U$ stored in $L_v$ give the
formula.
\end{proof}

\begin{lemma}[twist row]\label{i4:A-lem-twist}
If a selected bundle $b$ has a comparison datum, then every twisting place
of its realisation lies in the support $L.\mathrm{supp}$ of its family.
\end{lemma}
\begin{proof}
Let $y$ be a twisting place of the realisation. The twisting places of the
realisation are those of the real datum, so $y$ is twisting for the real
recursion of $b.B.Q$. The recursion component of the comparison datum
gives $\mu(y)\ne\mu_0$. The support condition of the family says that
$\mu(w)\ne\mu_0$ implies $w\in L.\mathrm{supp}$. Hence $y\in
L.\mathrm{supp}$.
\end{proof}

\begin{proposition}[pointwise leaves of the rows; formally verified, \cite{CharacterFamily}]\label{i4:A-prop-rowleaves}
Under the hypotheses of Theorem~\ref{i4:A-thm-allrows}:
\begin{enumerate}
\item[(P0)] a built record yields a column view;
\item[(P1)--(P6)] the pointwise properties $\mathcal R_1,\dots,\mathcal R_6$
of rows 7, 8, 9, 10, 11, 5 hold at the datum $w.D(Z)$ of every weighted
column family with its true-column object;
\item[(R1)] (a) over every column base lattice there is a self-dual
lattice whose adelic stabiliser contains the level group; (b) the
stabiliser index formula $[\widehat{\mathrm{Stab}}(N):U]=\mathrm{idx}(g_0,g_1)
(T_u^k-1)^f$; (c) column stabilisers are trivial;
\item[(R2)] the row of realisations: primitivity and type separation of the
realised rows of $b$ imply $\iota(V_0(\mu_0))$ is everything;
\item[(R3)] the block lemma: the block hypothesis and the forms tie imply
surjectivity of the lifted projection;
\item[(R4)--(R6)] the volume ratio at every column datum (from
Theorem~\ref{i4:thm-sw2}), the unfolding to volume $=$ mass for every
signature, and the class items at every rank;
\item[(R7)] for $\varepsilon<1/52$ the second-complement inputs follow from
$T_u$-size and $D_\Gamma$-size;
\item[(R8), (R9)] delivery of the arithmetic row over true columns, and the
row pin from the reading component;
\item[(R10), (R11)] the theta equality at indefinite scales and the right
theta bound;
\item[(R12)--(R14)] the two Siegel--Weil rows, the $L^p$ bound and the
off-diagonal order bound;
\item[(R15)] the reciprocity-kernel row under Artin reciprocity;
\item[(R16), (R17)] the counting constants and the irreducibility row;
\item[(R18)--(R25)] the assembly of definite rows from point summands and
the derivation of smallness, the norm bound and the $L_6$ bound.
\end{enumerate}
\end{proposition}
\noindent The structure of its formal proof is described in \S\ref{i4:Ast-rowleaves}.
\noindent What remains in each item is a pointwise statement about one
column datum or one bundle; they are proved in the development in the
part that treats the third-place series, which
is shared with Propositions~\ref{i4:A-prop-stages} and~\ref{i4:A-prop-open}.

\subsubsection{Proof of (A1)}\label{i4:A-sss-A1}

\begin{proposition}[member inputs; (a) proved modulo Proposition~\ref{i4:A-prop-inputs-leaves}, (b)--(d) formally verified, \cite{CharacterFamily}]\label{i4:A-prop-inputs}
\begin{enumerate}
\item[(a)] Assume the eleven published statements: (P1) global Artin
reciprocity; (P2) existence of CM elliptic curves with good reduction
outside a prime; (P3) Weber's ray class field theorem; (P4) the theory of
Harish-Chandra modules; (P5) rapid decay of cusp forms; (P6) the
Casselman--Wallach globalisation; (P7) Harish-Chandra's finiteness theorem;
(P8) the classification of discrete series; (P9) existence of
discrete-series covers; (P10) the formal degree formula via the compact
dual; (P11) Theorem~\ref{i4:thm-sw2}. Then the \emph{member hypotheses}
hold. These are the conjunction of: the ray class field statement; the
principal ideal theorem; (P2); (P4); (P5); (P6); existence of the Weil
representation for unitary dual pairs; existence of discrete series;
(P7); the infinitesimal-character statement for discrete series;
uniqueness of the discrete series of $U(p,q)$ with given Harish-Chandra
parameter; finite dimensionality of rows $\mathrm{RowsFinDim}$; (P11);
existence of discrete series in the volume normalisation; the formal
degree formula in the volume normalisation; the Siegel mass ratio for CM
fields displayed below; for every Harish-Chandra parameter, existence of a
discrete series with that parameter and its uniqueness up to unitary
equivalence, and existence for every group datum; (P1); (P3). The Siegel
mass ratio reads: for every CM field $K$ (totally complex, with complex
conjugation $\bar{\cdot}$) in which every maximal ideal $P$ admits
$\theta\in\mathcal O_K$ with $\theta-\bar\theta\notin P$, and every $M\ge3$,
\[
\frac{\mathrm{mass}(\mathrm{std}_M)}{\mathrm{mass}(\mathrm{std}_{M-1})}
=\frac{((M-1)!\,(v_K/\pi)^M)^{f_K}}{\sqrt{(2v_K)^{f_K}}}\,L_\eta(K,M),
\]
where $f_K$ is the number of infinite places, $v_K=|d_K|^{1/(2f_K)}/2$, and
$L_\eta(K,M)=\prod_Pe_P(M)$ with $e_P(M)=(1-(-1)^MN(P)^{-M/2})^{-1}$ if
$\bar P=P$ and $(1-N(P)^{-M})^{-1/2}$ otherwise, the product being the
unconditional product over maximal ideals, and $\mathrm{std}_M$ the lattice
$\mathcal O_K^M$ with the standard hermitian form.

For (b)--(d) assume moreover the published inputs of
Remark~\ref{i4:A-rem-pub} and the rationality of rows $\mathrm{RowRat}$:
\item[(b)] the multiplicity-cover hypotheses hold;
\item[(c)] an unlocalised saturated state exists;
\item[(d)] for every local Langlands datum and every system
$p(N,S,\Gamma)$, eventually, at every witness point the cuspidal
multiplicity of every good irreducible admissible $\sigma$ in the
discrete-series model of the first theta space equals that of its
transport in the corresponding model of the real datum.
\end{enumerate}
\end{proposition}
\noindent The structure of the formal proof of (b)--(d) is described in \S\ref{i4:Ast-inputs}.
\noindent Each part is reduced to fifteen leaves
(the model mass ratio, consequences of the classification of discrete
series, the lowest $K$-type, the Weil datum, finite dimensionality,
unramified Howe duality, the twisted Siegel--Weil identity, removed Euler
factors, local multiplicity at pinned data, extension of adelic data, the
saturated bare-form theorem, the size ledger, the bare-stage lemma,
transport along copies, twisted multiplicity one); they occupy the
local-type-dictionary part of the development.

\noindent\emph{Remark.} The conclusion of (a) is the 21-term conjunction listed. Part (a) needs
exactly (P1)--(P11) and not $\mathrm{RowRat}$. The proof of (a) is in
\S\ref{i4:ssec-proof-r4c}.

\begin{proposition}[leaves of the member hypotheses; formally verified, \cite{CharacterFamily}]\label{i4:A-prop-inputs-leaves}
\begin{enumerate}
\item[(M1)] (P1) implies the ray class field statement and the principal
ideal theorem.
\item[(M2)] The Weil representation for unitary dual pairs exists
(no hypothesis).
\item[(M3)] (P8) implies: multiplicity one of the lowest $K$-type, the
infinitesimal-character statement, existence of discrete series, existence
in the volume normalisation, and uniqueness of the discrete series of
$U(p,q)$ with given parameter.
\item[(M4)] (P9), (P8) and the consequences of (M3) imply the three
existence and uniqueness statements for every Harish-Chandra parameter.
\item[(M5)] $\mathrm{RowsFinDim}$ holds (no hypothesis).
\item[(M6)] The Siegel mass ratio for CM fields of
Proposition~\ref{i4:A-prop-inputs}(a) holds
(no hypothesis). This is a
consequence of the Siegel--Shimura mass formula for unimodular hermitian
lattices (Gan--Hanke--Yu; reference to be confirmed); since we have not compared the exact published
normalisation, we take it from the formal verification.
\end{enumerate}
\end{proposition}
\noindent The structure of its formal proof is described in \S\ref{i4:Ast-inputs-leaves}.

\begin{proposition}[stages of the member statement; formally verified, \cite{CharacterFamily}]\label{i4:A-prop-stages}
Assume the member hypotheses, the elliptic hypotheses, the
multiplicity-cover hypotheses, an unlocalised saturated state, a local
Langlands datum $\mathcal L$ with its local Langlands hypotheses, and
$c_{\rm Dep}$ with the multiplicity condition for $p_0$. Then the member
statement holds for the slots of $p_0$, and $\mathrm{ExBL}$ holds for the
built and good slots of $p_0$ and its bundle linkage. Moreover, for an
evaluation property $A$ with $\mathrm{EvLgG}(A)$ the member statement
transfers from $p_0$ to $p_A$, and $\mathrm{ExBL}$ holds for $p_A$ whenever
the position and multiplicity conditions hold for $p_A$.
\end{proposition}
\noindent The structure of its formal proof is described in \S\ref{i4:Ast-stages}.
\noindent The proof has three steps (stage statements from the
multiplicity-cover hypotheses; a three-step datum at every start point from
$\mathcal L$ and the saturated state; the closing pair lemma), whose leaves
are four stage statements. They lie in the same
third-place part of the development as
Proposition~\ref{i4:A-prop-rowleaves}.

\begin{proof}[Proof of (A1)]
\emph{(a) of Theorem~\ref{i4:thm-assembly}.} Proposition~\ref{i4:A-prop-inputs}(a)
gives the member hypotheses, hence the Siegel mass ratio for CM fields,
(b), (c) give the multiplicity-cover hypotheses and the saturated state, and
(d) the multiplicity condition for $p_0$ and for every $p(N,S,\Gamma)$, in
particular for $p_A$. The elliptic and local Langlands hypotheses are part
of the published inputs. Proposition~\ref{i4:A-prop-stages} gives the
member statement for $p_0$, and its transfer to $p_A$: $\mathrm{EvLgG}(A)$
says that $A$ depends only on $(K,B,f,T)$ and holds eventually for large
bases, so it survives the enlargement of the base in the second pass. For
$\mathrm{ExBL}$ at $p_A$, apply the last clause of
Proposition~\ref{i4:A-prop-stages} directly to $p_A$ (not by transport from
$p_0$): every core of the extended system is a core of the original one,
so it carries the extension law; the position condition passes to the
extension by $A$; and the multiplicity condition is
Proposition~\ref{i4:A-prop-inputs}(d). The four retwisting conditions are
part of the member statement.

\emph{(d).} Wu decay, finite dimensionality of rows, the positivity
statements and the volume ratio are (R4), (R12), (R13) and the finite
dimensionality leaf of Proposition~\ref{i4:A-prop-inputs}; each is an
application of Theorems~\ref{i4:thm-sw} and~\ref{i4:thm-sw2} in the ranges
checked in
Lemma~\ref{i4:A-lem-swrange} and Remark~\ref{i4:A-rem-sw}.
\end{proof}

\subsubsection{Proof of (A2)}\label{i4:A-sss-A2}

The twelve open statements were reduced in \S\ref{i4:ssec-proof-assembly}
to the bundle-level statements listed in (A2). Two of them are formal and
proved here; the rest are stated as one smaller proposition.

\begin{definition}[comparison datum]\label{i4:A-def-tie}
A \emph{comparison datum} for a bundle $b$ consists of a recursion of the
indefinite realisation $b.B.Q$, the $\psi$-tie of the division pair of $b$,
the degree-three normalisation of $b.B.Q$, and normalisation pins of
$b.B.Q$ compatible with the recursion.
\end{definition}

\begin{lemma}[the tie from a bundle witness]\label{i4:A-lem-tie}
Let $q$ be a witness point and $b$ a bundle over $q$ with a bundle witness
$w$, real input $y$ and recursive input $r$. Assume:
(T1) the zero section of the scale datum of $w$ is the local zero section;
(T5) the local level-character pin holds for the row datum of $w$; and
normalisation pins of the realisation of $y$ exist. Then $b$ has a
comparison datum.
\end{lemma}
\begin{proof}
The recursion local row condition of $w$ gives (i) that, under (T5), every
system of normalisation pins of the realisation of $y$ is compatible with
the recursion of $r$, and (ii) the $\psi$-tie of the division pair. By (T1)
the realisation of $y$ has the degree-three normalisation
(Proposition~\ref{i4:A-prop-open}(O1)). So the real package of $y$ carries
a recursion, the $\psi$-tie, the degree-three normalisation and compatible
pins. The package of $b$ is the package of $y$ by definition of a bundle
witness, and its realisation and division pair are those of $b$. This is a
comparison datum for $b$.
\end{proof}

\begin{proposition}[bundle-level leaves; formally verified, \cite{CharacterFamily}]\label{i4:A-prop-open}
Let $p=p_0$. Eventually, for every witness point $q$ of the context whose
object has the scale of the context, every $c$ and every bundle $b$ over
$(q.D,q.O,c)$ with a bundle witness $w$:
\begin{enumerate}
\item[(O1)] (T1)--(T5) hold: the zero section of $w$ is the local zero
section; the column datum of $b$ is the coefficient-field datum of the real
input; the group of the real input is the image of the place map; the
three theta-level conjugate-character identities hold; the local
level-character pin holds; and (T1) implies the degree-three normalisation
of the realisation of the real input;
\item[(O2)] the host hypothesis holds at the stage of $q$, and with (T2)--(T4)
it gives normalisation pins of the realisation of the real input;
\item[(O3)] the realisation of $b$ satisfies $\mathrm{MI}^*_2$, and given
$\mathrm{MI}^*_2$ the pairings on the definite base and the division-row
pairings vanish;
\item[(O4)] at points with a parameter choice the $\Lambda$-rows and the
$\kappa$-rows hold;
\item[(O5)] the goodness and label rows hold (local Langlands for
$\mathrm{GL}_n$ in depth zero, Lemma~\ref{i4:A-lem-greenlab});
\item[(O6)] weak geometric compatibility holds (ray class field
statement);
\item[(O7)] the defect row on the point-evaluation side holds;
\item[(O8)] (\emph{row conditions at the companion}; it holds for
every twisted input, not only eventually) for the division pair
of the companion of every such bundle: the $\Lambda$-rows imply the
level-three $\Lambda'$-rows; the decomposition-point rows hold; the
decomposition-point condition holds; and the level-three divisor rows imply
the divisor condition.
\end{enumerate}
Moreover:
\begin{enumerate}
\item[(O9)] (\emph{join of the linked realisation}; the proof
uses only the cover of discrete series, Harish-Chandra modules,
Casselman--Wallach, Harish-Chandra finiteness, rapid decay and the
linked tie, so the hypotheses below are sufficient) assume the existence of
the Weil representation for unitary dual pairs, the existence and the cover
of discrete series, Harish-Chandra modules, the Casselman--Wallach
globalisation, Harish-Chandra's finiteness theorem and rapid decay of cusp
forms. For every parameter system $p$ of the form $p(N,S,\Gamma)$: if
primitivity, the decomposition-point condition and the tie are linked at
$p$, then the selection-jump condition is linked at $p$.
\end{enumerate}
\end{proposition}
\noindent The structure of its formal proof is described in \S\ref{i4:Ast-open}.
\noindent What is not proved here are the bundle-level verifications of the
open-origin rows (multiplicity, divisor,
$\Lambda'$, $\kappa$, goodness, labels, geometry, definite forms); (O8),
which is a direct verification on the companion; and (O9), which builds
the join of the linked realisation with the kernel from the comparison
datum, the transfer of the row placement of the spectral stage, and the
lattice of the join subspace. They are
the largest single block of the third-place part of the development (see
the size remark after Proposition~\ref{i4:A-prop-rowleaves}).

\begin{proof}[Proof of (A2)]
\emph{The tie.} Fix a bundle $b$ linked at $p$: a witness point $q$, a
bundle $b'$ over $q$ with witness $w$, and an identification
$e:(\beta.D,\beta.O,b)=(q.D,q.O,b')$. Since the scale of $\beta$ is that of
the context, so is the scale of $q.O$; hence (O1) and (O2) apply to
$(q,b',w)$, and give normalisation pins of the realisation of the real
input. Lemma~\ref{i4:A-lem-tie} gives a comparison datum for $b'$, which
is transported to $b$ along $e$. This uses no published input and no
property of the extension, so the tie is linked at every parameter
system.

\emph{The other bundle-level inputs} are (O3)--(O9).
\end{proof}

\subsubsection{Proof of (A3)}\label{i4:A-sss-A3}

The norm row (row 26) has the form
$\mathrm{Row}(Q,\mathrm{Lk},(\mathcal A,\sigma)\mapsto NN(D_3(3,\mathcal A)))$:
it does not involve the selection. Two norm properties occur: the
\emph{bound at $2,3$}, $NN_{23}(\mathcal C)=$ ``$\mathcal C$ is closed for
the property: the family $F$ has the local bound at $2$ and $3$'', and the
property $NN_{\rm LN}(\mathcal C)=$ ``$\mathcal C$ is closed for
$\mathrm{LinkN}$'', which is the norm condition $(\mathrm{LN})$ of field
(viii).

\begin{definition}[local bound at $2$ and $3$]\label{i4:A-def-at23}
A family $F$ has the \emph{local bound at $2$ and $3$} if there are $C,x_0$
such that for all $x\ge x_0$ with $x>x_1(F)$, every index $i$ and all
local data at $(F.D(x,i))$ satisfying the norm-choice setup (norm-choice
facts, the canonical hub, admissible and polynomially bounded parameters,
$L.\mathrm{sc}.L\ge8$, split bound $\ge60$, the killing condition and
finiteness of the Shafarevich groups), the bound
$\mathrm{BoundAtTwoThree}(F.L(x,i),\rho,C)$ at the primes $2$ and $3$
holds for the representation $\rho$ of the hub. Here
$\epsilon(l)=l^{-2}$, and the \emph{killing condition} for $\rho$ says: at
every twisted place $k$, with prime $l=l_k$, (K1) every $\pi\in Z^{\rm
good}_k$ has finite $s$-length and $s_\pi(\rho_k)\le\nu(l)+\epsilon(l)$,
$\nu(l)$ the cyclotomic correction at $l$; and (K2)
$t\sum_{v\in Q_k}\ell(\mathcal O_k/(\zeta_v-1))\le\epsilon(l)$, with
$Q_k$, $\zeta_v$ and $\ell$ as in \S\ref{i4:ssec-proof-at23rows}. The
\emph{finiteness of the Shafarevich groups} says that
$\mathrm{sha}_p(e)<\infty$ for every prime $p$ (including $p=2$) and
$0<e<t$.
\end{definition}

\begin{proposition}[the local bound at $2$ and $3$]\label{i4:A-prop-at23}
Assume $\mathrm{Hypotheses}_2$. For every family $F$ there are $C,x_0$
such that for all $x\ge x_0$ with $x>x_1(F)$, all $i\in I(x)$, every local
datum $(HW,F_{\rm on},\mathrm{Cmp})$ at $F.D(x,i)$, every hub datum $X$ with
representation constants $X.\rho$, and every $C_2$ satisfying the setup
condition $\mathrm{Setup}(F.L(x,i),\mathrm{nch},HW,X,C_2)$, the bound
$\mathrm{BoundAtTwoThree}(F.L(x,i),X.\rho,C)$ holds: at every twisted place
$k$ whose prime $l_k$ is $2$, or is $3$ and not in the small set of the
host, and every $\pi\in Z^{\rm good}_k$, $\mathrm{Bd}(X.\rho,k,\pi,\mathrm{mid}_F(L,C))$. Here
$\mathrm{Bd}(\rho,k,\pi,\beta)$ means
$v_{cc}(k,\pi)\le s_\pi(\rho_k)+\beta+VV(l_k)+v_{hE}(l_k)+\epsilon(l_k)$ (the
valuation of the conductor coefficient against the slope of $\rho_k$ at
$\pi$, with the three fixed local corrections), and
$\mathrm{mid}_F(L,C)=C\,fmb\log x+C\,x^C\,\ell^2$, where $\ell=x^\varepsilon$ is the scale of the datum $L$
(the quantity $L.\mathrm{sc}.L$ of the setup).
\end{proposition}
\noindent (The bound is not claimed for every $\rho$: the statement
quantifies over setups as above, and this is what
Lemma~\ref{i4:A-lem-at23} uses.)

\begin{definition}[the cost of a small prime]\label{i4:A-def-cost}
Let $L$ be a size datum at a datum $D$, with scale $x$, $b=e^{\ell}$,
$\ell=L.\mathrm{sc}.L$, column size $S_L$ and $\eta$-size $\eta_L$, and
let $N=Q_Sqd$ be the conductor of the host. Put
$h_B(L)=b+S_L+\eta_L+t+2$; for a prime $l$ put $s_l=\max(1,v_l(N))$,
$\Lambda_l=s_l+\log_l(2h_B)+1$, and
\[
\mathrm{cost}(L,A,B,l)=A\,f\,m\,h_B\,s_l+B\,f^2t^2n_V\,\Lambda_l .
\]
\end{definition}

\begin{proposition}[the row bounds at $2$ and $3$]\label{i4:A-prop-at23rows}
Assume $\mathrm{Hypotheses}_2$. (a) For every family there is $x_0$ such
that for $x\ge x_0$, $x>x_1$, every $i$ and every setup as in
Proposition~\ref{i4:A-prop-at23} at $(F.D(x,i),F.L(x,i))$, every twisted
place $k$ with $l_k=2$ and every $\pi\in Z^{\rm good}_k$ satisfy
$\mathrm{Bd}(X.\rho,k,\pi,\mathrm{cost}(L,6000,30000,l_k))$ (one may take
$x_0=0$).
(b) The same holds, with the same constants $(6000,30000)$, at the twisted
places with $l_k=3\notin$ small.
(c) (\emph{cost estimate}) Let $A,B\ge0$, $C_0\ge1$, let $L$ be a size
datum carrying a parameter choice, and let $l$ be a prime with
$v_l(N)=v_l(Q_S)$. If $1\le\log x$, $n\le x^{C_0}$ and $t\le x^{C_0}$, then
$\mathrm{cost}(L,A,B,l)\le\mathrm{mid}_F(L,C)$ for every
$C\ge18AC_0+2B(3C_0+9)+5C_0+1$.
\end{proposition}
\noindent The proof is in \S\ref{i4:ssec-proof-at23rows}; it is complete
modulo Proposition~\ref{i4:A-prop-split}.

\begin{proof}[Proof of Proposition~\ref{i4:A-prop-at23} from Proposition~\ref{i4:A-prop-at23rows}]
Let $F$ be a family. The sign conditions of (c) hold: $A,B\ge0$ by (a),
(b), and $C_0\ge1$ for the constant of the admissible parameters of the
family. The valuation hypothesis of (c) holds at $l=2$ (as
$2$ is in the small set of every host) and at $l=3$ (because $3\ne q$ by the definition of a host, and $3\ne d$ because
$d\equiv7\bmod16$). By
(a) and (c), for $x$ beyond the threshold of (a), the threshold of the
admissible parameters and $3$, at every setup the places with $l_k=2$ satisfy
$\mathrm{Bd}(X.\rho,k,\pi,\mathrm{mid}_F(L,C))$ with $C=18AC_0+2B(3C_0+9)+5C_0+1\ge0$;
by (b) and (c) the places with $l_k=3\notin$ small satisfy it with a constant
$C'\ge0$. (Every point beyond $x_1$ carries a parameter choice, a field of
the family, so (c) applies.) Since $\mathrm{mid}_F(L,\cdot)$ is nondecreasing on $[0,\infty)$ and
$\mathrm{Bd}(\rho,k,\pi,\cdot)$ is an upper bound (monotone in its last
argument), both hold with $\max(C,C')$ beyond the larger threshold.
\end{proof}

\begin{lemma}[the bound at $2,3$ is linked]\label{i4:A-lem-at23}
If every family has the local bound at $2$ and $3$, then for every
parameter system $p$ the norm row for $NN_{23}$ is linked at $p$.
\end{lemma}
\begin{proof}
A class all of whose members have a property is closed for it (with
threshold datum $0$). So $NN_{23}(D_3(3,\mathcal A))$ holds for every
$\mathcal A$, independently of $\sigma$. By Lemma~\ref{i4:A-lem-calc}(d)
the row holds for every slot and linkage, and by
Lemma~\ref{i4:A-lem-calc}(a) it is linked at $p$.
\end{proof}

\begin{proposition}[the Galois upgrade of the norm row; formally verified, \cite{CharacterFamily}]\label{i4:A-prop-galois}
Assume $\mathrm{Hypotheses}_2$, the published inputs of the third layer
(including the Chebotarev input), the compatibility of geometric Galois
representations, the series input $\mathrm{SerN}$ (below), depth-zero
support, local finiteness, local Tate duality and the local Euler
characteristic. Let $\mathcal L$ be a local Langlands datum, $c_{\rm Dep}\in\R$,
and $p_0$ the parameter system of dependence $c_{\rm Dep}$ attached to
$\mathcal L$. If the norm row for $NN_{23}$ is linked at $p_0$, then the norm
row for $NN_{\rm LN}$ is linked at $p_0$.
\end{proposition}
\noindent The structure of its formal proof is described in \S\ref{i4:Gst-galois}.
\noindent The statement has exactly eight
hypotheses. $\mathrm{SerN}$ is the conjunction of six statements:
rationality of the twisted series, their integrality, Hecke integrality,
the Hecke transform, the Hecke read tie, and the ratio identity.
It is not Section~\ref{i5:sec} and
not an extra hypothesis of the paper: in the application
it is supplied
from the
non-endoscopic third-layer hypotheses (global Artin reciprocity for
rationality and integrality, toric transport for the ratio identity) and
the compatibility of geometric Galois representations; the Hecke clauses
are proved outright.
$\mathrm{Hypotheses}_2$ consists of Fontaine's $B_{\rm cris}$, global Artin
reciprocity, local Bloch--Kato, local Tate duality, the local Euler
characteristic, local finiteness and the Bhatt--Scholze integral
comparison. The proof runs as follows: a good cut exists at every datum
satisfying the eighteen Galois preconditions and the fourfold tie
$\mathrm{Tie4}$ (Chebotarev input, local finiteness, Tate duality, Euler
characteristic); the Galois preconditions and $\mathrm{Tie4}$ hold at
members; they transfer from members to witness points along the position
condition; and the Galois linking theorem turns the linked bound at $2,3$
into $\mathrm{LinkN}$, after intersecting the two slots by
Lemma~\ref{i4:A-lem-calc}(c). What remains is that chain of Galois-side
statements.

\begin{proposition}[the series input; formally verified, \cite{CharacterFamily}]\label{i4:A-prop-sern}
The non-endoscopic third-layer hypotheses and the compatibility of
geometric Galois representations imply the published inputs of the third
layer and $\mathrm{SerN}$ (and, in addition, the elliptic hypothesis and
a local Langlands datum).
\end{proposition}
\noindent The structure of its formal proof is described in \S\ref{i4:Gst-sern}.
\noindent In $\mathrm{SerN}$, rationality and integrality of the twisted
series come from global Artin reciprocity, the ratio identity from toric
transport, and the three Hecke clauses are proved outright.

\begin{proof}[Proof of (A3)]
By Proposition~\ref{i4:A-prop-sern} the hypotheses of
Proposition~\ref{i4:A-prop-galois} that are not in
Remark~\ref{i4:rem-trust}(ii), namely the published inputs of the third
layer and $\mathrm{SerN}$, follow from the non-endoscopic third-layer
hypotheses and the compatibility of geometric Galois representations.
By Proposition~\ref{i4:A-prop-at23} every family has the local bound at
$2$ and $3$. By Lemma~\ref{i4:A-lem-at23} the norm row for $NN_{23}$ is
linked at $p_0$. By Proposition~\ref{i4:A-prop-galois} the norm row for
$NN_{\rm LN}$ is linked at $p_0$. This is (A3). At an extension $p_A$ the
same holds by the monotonicity proved in \S\ref{i4:ssec-proof-assembly}.
\end{proof}

\subsubsection{Proof of (A4)}\label{i4:A-sss-A4}

\begin{proposition}[the functional and covolume rows; formally verified, \cite{CharacterFamily}]\label{i4:A-prop-funcov}
Assume the published inputs of the third place (Remark~\ref{i4:rem-trust}(ii)),
the local theta inputs, the agreement of the theta integrals in the Weil
tensor models, the local matching of the twisted local factors, the
covolume bound for $X$ and the existence of a local Langlands datum. Then:
\begin{enumerate}
\item[(a)] there are constants $C_{F0},C_{F1}$ such that, with
$c_{\rm Dep}=\max(C_{F0},C_{F1})$ and $A$ the evaluation-grid property at
$c_{\rm Dep}$, $\mathrm{EvLgG}(A)$ holds and the functional condition is
linked at the extension $p$ of $p_0$ by $A$;
\item[(b)] the covolume bound is linked at the same $p$;
\item[(c)] the rationality of rows $\mathrm{RowRat}$ holds;
\item[(d)] surjectivity onto $X$ and the lower bound $\lambda_{\rm low}$
hold.
\end{enumerate}
\end{proposition}
\noindent The structure of its formal proof is described in \S\ref{i4:Gst-funcov}.
\noindent\emph{Remark.} The hypothesis of (d) is
the bundle of third-place hypotheses of the unitary group (about forty
published statements, together with the join statement and the Hecke--Tate
functional equation); the proof uses exactly the following. For
surjectivity onto $X$: Green's cuspidality criterion, Clozel's criterion
for cohomological representations, the $\ell$-adic characters attached to
algebraic Hecke characters (two forms), Chebotarev for degree-one primes,
Galois representations attached to polarised cuspidal representations,
base change for definite unitary groups, Bhatt--Scholze crystalline
lattices, rapid decay, the local Langlands correspondence for
$\mathrm{GL}_n$ with its local inputs (convolution, depth-zero
inertia, Brauer, Chevalley, exhaustion, multiplicities, Zelevinsky
classification, geometric lemma, Bernstein--Zelevinsky), global Artin
reciprocity, existence of discrete-series covers and Casselman--Wallach.
For $\lambda_{\rm low}$: rapid decay, Theorems~\ref{i4:thm-sw}
and~\ref{i4:thm-sw2}, the doubling method for local zeta integrals,
unramified doubling for $\mathrm{GL}$, Bernstein--Zelevinsky finiteness
(Theorem~\ref{i4:thm-bzfin}) and the formal degree formula. Several of
these (Clozel, Galois representations of polarised cuspidal
representations, base change, Bhatt--Scholze, Chebotarev, the doubling
method) are not among the published inputs of
Remark~\ref{i4:A-rem-pub}; the hypothesis ``the published inputs
of the third place'' used above is to be read as this list. We compare (a)--(c)
with the literature only at the level of the top statements.
\noindent In the development, (a) comes from the fit of the constant
function obtained from the join and the separation of functionals (by
unramified Howe duality and Siegel--Weil), together with discrete series,
Harish-Chandra modules, rapid decay, Casselman--Wallach and the ray class
field statement, applied to the open statements at $p$; (b) comes from the
same published inputs and the covolume bound for $X$, again from the open
statements at $p$; (c) is the assembly of test vectors; (d) is derived from
the published inputs of the third place.

\begin{proof}[Proof of (A4)]
This is Proposition~\ref{i4:A-prop-funcov}(a), (b). The open statements at
$p$ that it uses are Theorem~\ref{i4:thm-assembly}(c), proved in
\S\ref{i4:ssec-proof-assembly} from (A2) alone.
\end{proof}

\begin{proof}[Proof of Proposition~\ref{i4:prop-assembly-leaves}]
(A1), (A2), (A3), (A4) are proved in \S\ref{i4:A-sss-A1},
\S\ref{i4:A-sss-A2}, \S\ref{i4:A-sss-A3}, \S\ref{i4:A-sss-A4}, modulo
Propositions~\ref{i4:A-prop-rowleaves}, \ref{i4:A-prop-inputs-leaves},
\ref{i4:A-prop-inputs}(b)--(d), \ref{i4:A-prop-stages},
\ref{i4:A-prop-open}, \ref{i4:A-prop-split}, \ref{i4:A-prop-galois},
\ref{i4:A-prop-sern} and~\ref{i4:A-prop-funcov}, and the published
theorems of Remark~\ref{i4:rem-trust}(ii). Proposition~\ref{i4:A-prop-inputs}
uses $\mathrm{RowRat}$, which is Proposition~\ref{i4:A-prop-funcov}(c).
\end{proof}

\subsection{Norm primes in short intervals}\label{i4:ssec-grp-norm}

\subsubsection{The short-interval bound for norm primes}\label{i4:ssec-proof-shortnorm}

The development does not deduce Proposition~\ref{i4:thm-shortnorm} from
Theorem~\ref{i4:thm-classzeros} (which concerns Dirichlet characters, while
here the relevant $L$-functions are Hecke $L$-functions of ray class
characters of $\Q(\sqrt{-d})$). It works with an axiomatic family of
Dirichlet series, proves the analytic statements for every such family, and
then constructs a family from the ray class group. We follow that route.

\begin{definition}[$L$-families]\label{i4:def-lfam}
An \emph{$L$-family} consists of a finite abelian group $G$, reals $B$ and
$\kappa$, and functions $a,\lambda:G\times\N\to\R$ such that:
$4\le B$; $|G|\le B$; $1/B\le\kappa\le4$; $a\ge0$, $a(g,0)=0$,
$a(0,1)=1$, $a(g,1)=0$ for $g\ne0$;
$\bigl|\sum_{1\le n\le x}a(g,n)-\kappa x\bigr|\le10\sqrt x$ for $x\ge1$;
$\bigl|\sum_n a(g,n)e^{-tn}-\kappa/t\bigr|\le B^{10}t$ for $0<t\le1$;
$\lambda\ge0$ and $\sum_g\lambda(g,n)\le2\Lambda(n)$;
$a(g,n)\log n=\sum_{h\in G}\sum_{n_1n_2=n}\lambda(h,n_1)a(g-h,n_2)$; and, for
every character $\chi$ of $G$ with $\chi^2=1$ and every $n$, writing
$a_\chi(n)=\sum_g\Rel\chi(g)a(g,n)$ and $a_1=a_{\chi_0}$,
$0\le\sum_{n_1n_2=n}a_1(n_1)a_\chi(n_2)$ and
$a_1(n)\le\sum_{n_1n_2=n^2}a_1(n_1)a_\chi(n_2)$. (Since $\chi^2=1$, $\chi$
is real-valued, so this $a_\chi$ agrees with the one defined next.)

For a character $\chi$ of $G$ put $a_\chi(n)=\sum_g\chi(g)a(g,n)$,
$\lambda_\chi(n)=\sum_g\chi(g)\lambda(g,n)$, $\psi_\chi(u)=\sum_{1\le n\le
u}\lambda_\chi(n)$, $\mathrm{res}_\chi=\kappa|G|$ if $\chi=1$ and $0$
otherwise, $\mathrm{mel}_\chi(s)=\int_1^\infty\bigl(\sum_{n\le
x}a_\chi(n)-\mathrm{res}_\chi x\bigr)x^{-s-1}dx$ (holomorphic for
$\Rel s>1/2$), and $F_\chi(s)=\mathrm{res}_\chi s+(s-1)s\,\mathrm{mel}_\chi(s)$
if $\chi=1$, $F_\chi(s)=s\,\mathrm{mel}_\chi(s)$ otherwise. (By partial
summation $F_\chi(s)=(s-1)^{\delta_{\chi=1}}\sum_na_\chi(n)n^{-s}$ for
$\Rel s>1$, so $F_\chi$ is the continuation of $L(s,\chi)$ with the pole
removed.) Put $\ell(t)=\log(B(|t|+2))$. For $c>0$ and a finite set $Z\subset\R$,
$\mathrm{Exc}_\chi(c,Z)$ means: $|Z|\le1$; every $\beta\in Z$ is a simple zero
of $F_\chi$ with $1-c/\ell(0)\le\beta<1$; and every zero $s$ of $F_\chi$
with $\Rel s\ge1-c/\ell(\Iml s)$ lies in $Z$.
\end{definition}


\begin{proposition}[zeros of $L$-families]\label{i4:prop-lfam-zeros}
There are absolute constants with the following properties, for every
$L$-family and all characters $\chi,\chi'$ of $G$.
\begin{enumerate}
\item[(F1)] (\emph{zero-free region}) There is $c_1>0$ such that if
$\Rel s\ge3/4$, $F_\chi(s)=0$ and ($\chi^2\ne1$ or $\Iml s\ne0$), then
$\Rel s<1-c_1/\ell(\Iml s)$.
\item[(F2)] (\emph{one exceptional zero}) There is $c_2>0$ such that if
$\chi^2=\chi'^2=1$, $\beta,\beta'\ge3/4$, $\beta,\beta'\ge1-c_2/\ell(0)$,
$F_\chi(\beta)=F_{\chi'}(\beta')=0$, then $\chi=\chi'$, $\beta=\beta'$ and
$F'_\chi(\beta)\ne0$.
\item[(F3)] (\emph{distance from $1$}) There are $A,c_3>0$ such that if
$\chi^2=1$, $\beta\ge3/4$ and $F_\chi(\beta)=0$, then $1-\beta\ge c_3B^{-A}$.
\item[(F4)] (\emph{prime number theorem with exceptional zero}) For every
$c_0\in(0,1/20]$ there are $c,C>0$ such that if $\mathrm{Exc}_\chi(c_0,Z)$,
$u\ge2$ and $\log B\le\sqrt{\log u}$, then
\[
\Bigl|\psi_\chi(u)-\delta_{\chi=1}u+\sum_{\beta\in Z}\frac{u^\beta}\beta\Bigr|\le
Cu\,e^{-c\sqrt{\log u}} .
\]
\end{enumerate}
\end{proposition}
\noindent These are the classical theorems of de la
Vall\'ee Poussin, Gronwall--Landau--Titchmarsh (one exceptional zero) and
prime-number-theorem type, for the axiomatic family; (F3) with a power of
$B$ (one may take $A=46$) is an effective (non-Siegel) bound, available here
because of the theta axiom and the first positivity axiom. The proof, in
\S\ref{i4:ssec-proof-lfam-zeros}, uses only the axioms of
Definition~\ref{i4:def-lfam}, a lemma of Landau on the logarithmic
derivative, Jensen's formula, the truncated Perron formula and, for (F3),
an elementary integral kernel.

\begin{proposition}[the ray class family]\label{i4:prop-lfam-ray}
Let $d\equiv3\pmod4$ be a prime with $d\ge100$ and $M\ge3$. Let $K=\Q(\sqrt{-d})$,
$G$ the ray class group of $K$ modulo $M$ (nonzero ideals prime to $M$
modulo the relation $I\sim J$ iff $(\alpha)I=(\beta)J$ for some nonzero
$\alpha,\beta\in\mathcal O_K$ with $\alpha\equiv\beta\equiv1\bmod M$), written
additively; for $g\in G$ and $n\in\N$ let $a(g,n)$ be the number of nonzero
ideals of norm $n$, prime to $M$, in the class $g$, and $\lambda(g,n)$ the sum,
over those of them that are of the form $\mathfrak p^j$ ($\mathfrak p$ prime,
$j\ge1$), of $\log N\mathfrak p$. Then $(G,B,\kappa,a,\lambda)$ with
\[
B=4dM^2,\qquad\kappa=\frac{2\pi}{\sqrt d\,M^2}
\]
is an $L$-family, and every prime $p$ with $\lambda(0,p)\ne0$ can be
written $4p=A_0^2+dB_0^2$ with $A_0,B_0\in\Z$, $M\mid B_0$ and
$2M\mid A_0-2-B_0$.
\end{proposition}
\noindent The proof is in \S\ref{i4:ssec-proof-lfamray}, with the leaves
of Proposition~\ref{i4:prop-ray-leaves} proved in \S\ref{i4:Rst-ray}.

\begin{proposition}[analytic and positivity leaves of the ray class family]\label{i4:prop-ray-leaves}
In the situation of Proposition~\ref{i4:prop-lfam-ray}:
\begin{enumerate}
\item[(R1)] (\emph{count and theta}) for every $g$:
$\bigl|\sum_{1\le n\le y}a(g,n)-\kappa y\bigr|\le10\sqrt y$ for $y\ge1$, and
$\bigl|\sum_na(g,n)e^{-tn}-\kappa/t\bigr|\le(4dM^2)^{10}t$ for $0<t\le1$;
\item[(R2)] (\emph{positivity}) for every $\varepsilon:G\to\{\pm1\}$ with
$\varepsilon(0)=1$ and $\varepsilon(g+h)=\varepsilon(g)\varepsilon(h)$, and every $n$, writing
$a_1(n)=\sum_ga(g,n)$ and $a_\varepsilon(n)=\sum_g\varepsilon(g)a(g,n)$:
$0\le\sum_{n_1n_2=n}a_1(n_1)a_\varepsilon(n_2)$ and
$a_1(n)\le\sum_{n_1n_2=n^2}a_1(n_1)a_\varepsilon(n_2)$;
\item[(R3)] (\emph{class numbers}) $|G|\le h_K\,M^2$ and $h_K\le4d$.
\end{enumerate}
\end{proposition}
\noindent The proof is in \S\ref{i4:Rst-ray}.

\begin{lemma}[the trivial class in short intervals]\label{i4:lem-trivshort}
There are $A,c,c'>0$ such that for every $L$-family and all reals $X\ge1$
and $h$ with $\log B\le c\sqrt{\log X}$ and $e^{-c\sqrt{\log X}}\le h\le1$,
\[
\sum_{\lfloor X\rfloor<n\le\lfloor X(1+h)\rfloor}\lambda(0,n)\ge c'\,\frac{hX}{B^A}.
\]
\end{lemma}
\begin{proof}
\emph{Exceptional sets.} Let $c_0=\min(1/20,c_1,c_2)$ and, for each
$\chi$, let $Z_\chi$ be the set of real $\beta\ge1-c_0/\ell(0)$ with
$F_\chi(\beta)=0$. Since $\ell\ge2$, the region $\Rel s\ge1-c_0/\ell(\Iml s)$
lies in $\Rel s\ge 3/4$. By (F1) every zero there is real and belongs to a
character with $\chi^2=1$; by (F3) it is $<1$; by (F2) there is at most one
such pair $(\chi,\beta)$ over all $\chi$, and the zero is simple. Hence
$\mathrm{Exc}_\chi(c_0,Z_\chi)$ holds for every $\chi$, at most one $Z_\chi$
is nonempty, and its element $\beta_1$ satisfies $\beta_1\ge3/4$ and
$1-\beta_1\ge\delta:=\min(c_3,1)B^{-A}$.

\emph{Orthogonality.} Since $\sum_\chi\chi(g)=|G|\,\delta_{g=0}$,
$S(u):=\sum_{n\le u}\lambda(0,n)=|G|^{-1}\sum_\chi\psi_\chi(u)$. By (F4)
(with $c_0$, and constants $c_4,C_4$), for $u\ge2$ with $\log B\le\sqrt{\log u}$,
\[
\Bigl|S(u)-\frac1{|G|}\Bigl(u-\sum_\chi\sum_{\beta\in Z_\chi}\frac{u^\beta}\beta\Bigr)\Bigr|
\le C_4u\,e^{-c_4\sqrt{\log u}},
\]
the average of $|G|$ errors being bounded by the largest.

\emph{The short interval.} Let $Y=X(1+h)$, so $X\le Y\le2X$. Assume
$c\le1$; then $\log 4\le\log B\le c\sqrt{\log X}$ gives
$\sqrt{\log X}\ge\log4/c$, in particular $X\ge e$ and both applications of
(F4) are legitimate. Subtracting,
$\sum_{\lfloor X\rfloor<n\le\lfloor Y\rfloor}\lambda(0,n)=S(Y)-S(X)$ is at least
\[
\frac1{|G|}\Bigl(hX-\frac{Y^{\beta_1}-X^{\beta_1}}{\beta_1}\Bigr)-4C_4X e^{-c_4\sqrt{\log X}},
\]
with the $\beta_1$ term absent if there is no exceptional zero. For
$3/4\le\beta\le1-\delta$ and $X\ge e$ one has
$(Y^\beta-X^\beta)/\beta\le X^\beta h\le hX(1-\delta/2)$: the first
inequality is $(1+h)^\beta\le1+\beta h$, the second is
$X^{\beta-1}=e^{-(1-\beta)\log X}\le e^{-\delta}\le1-\delta/2$ for
$\delta\le1$. So the main term is at least
$hX\delta/(2|G|)\ge hX\min(c_3,1)/(2B^{A+1})$, using $|G|\le B$.

\emph{The error.} Put $\gamma=\min(c_3,1)/4$. We need
$4C_4e^{-c_4\sqrt{\log X}}\le\gamma hB^{-A-1}$. By the hypotheses
$hB^{-A-1}\ge e^{-(A+2)c\sqrt{\log X}}$. Choose
$c\le\min\bigl(1,\,c_4/(2(A+2))\bigr)$; then it suffices that
$e^{-(c_4/2)\sqrt{\log X}}\le\gamma/(4C_4)$, and since $\sqrt{\log
X}\ge\log4/c$ this holds once also $c\le c_4\log4/(2\log^+(4C_4/\gamma))$.
With these choices the error is at most half of the main term, and the lemma
holds with $A+1$ in place of $A$ and $c'=\gamma$.
\end{proof}

\begin{proof}[Proof of Proposition~\ref{i4:thm-shortnorm}]
Let $A,c_5,c'_5$ be the constants of Lemma~\ref{i4:lem-trivshort} and let
$\mathcal F$ be the family of Proposition~\ref{i4:prop-lfam-ray}, so
$B=4dM^2$. Let $N$ be the number of primes counted in the proposition, and
$Y=X(1+h)\le2X$.

\emph{Upper bound for the sum by $N$.} Since $\lambda(0,n)\le2\Lambda(n)$,
splitting the sum over $(\lfloor X\rfloor,\lfloor Y\rfloor]$ into primes
and prime powers gives
\[
\sum_{\lfloor X\rfloor<n\le\lfloor Y\rfloor}\lambda(0,n)\le2\log Y\cdot N+2(\psi(Y)-\theta(Y)),
\]
because a prime $p$ with $\lambda(0,p)\ne0$ is one of the counted primes and
then $\lambda(0,p)\le2\log p\le2\log Y$. By Chebyshev's bound
$\theta(y)\le y\log4$, $\psi(Y)-\theta(Y)=\sum_{k\ge2}\theta(Y^{1/k})\le
C_\theta\sqrt Y$ with an absolute $C_\theta$; so $\psi(Y)-\theta(Y)\le2C_\theta\sqrt X$.
Also $\log Y\le\log X+h\le2\log X$ as $\log X\ge1$.

\emph{Comparison.} Let $K=8(C_\theta+1)/c'_5$ and
$c=\min(c_5,1/(2(1+A)+2+K))$. Put $s=\sqrt{\log X}$. From $\log B\le cs$ and
$B\ge4$, $s\ge(2(1+A)+2+K)\log4\ge 2(1+A)+2+K$, so $\log X\ge1$, and
\[
K\sqrt X\,B^A\le K e^{s^2/2+Acs}\le e^{s^2/2-(1+A)s}e^{s^2/2}e^{As}=e^{-s}X\le hX,
\]
using $K\le e^{s^2/2-(1+A)s}$ (which follows from $1+y\le e^y$ and
$s/2-(1+A)\ge1+K/2$) and $h\ge e^{-cs}\ge e^{-s}$. Hence
$4C_\theta\sqrt X\le\frac{c'_5}2\,hX/B^A$, and the lemma gives
\[
4\log X\cdot N\ge c'_5\frac{hX}{B^A}-4C_\theta\sqrt X\ge\frac{c'_5}2\frac{hX}{B^A}.
\]
So $N\ge(c'_5/8)\,hX/((4dM^2)^A\log X)$, which is the claim with
$c'=c'_5/8$.
\end{proof}
\noindent Proposition~\ref{i4:prop-lfam-zeros} is proved next;
Proposition~\ref{i4:prop-lfam-ray} is proved in
\S\ref{i4:ssec-proof-lfamray} and \S\ref{i4:Rst-ray}.

\medskip\noindent\label{i4:ssec-proof-lfam-zeros}\emph{Proof of Proposition~\ref{i4:prop-lfam-zeros}.}

Throughout, an $L$-family is fixed; all constants called $C,C_Z,C_\Gamma,\dots$
or $c,\eta,\dots$ are absolute (independent of the family, of $\chi$ and of
all variables) unless a dependence is indicated. We write $s=\sigma+it$,
$\Lambda$ for von Mangoldt's function and $\zeta$ for the Riemann zeta
function. Recall $\ell(t)=\log(B(|t|+2))>2$, since $B\ge4$.

We use two published lemmas.

\begin{lemma}[Landau; Titchmarsh, \emph{The theory of the Riemann
zeta-function}, \S3.9, Lemma~$\alpha$]\label{i4:lem-landau}
There is an absolute constant $A_T$ with the following property. Let $f$
be holomorphic on the closed disc $|s-s_0|\le r$, with $f(s_0)\ne0$ and
$|f(s)/f(s_0)|<e^{M}$ there, where $M>1$. Then for $|s-s_0|\le r/4$,
\[
\Bigl|\frac{f'(s)}{f(s)}-\sum_\rho\frac1{s-\rho}\Bigr|<\frac{A_TM}{r},
\]
where $\rho$ runs through the zeros of $f$ (with multiplicity) in
$|\rho-s_0|\le r/2$.
\end{lemma}

\begin{lemma}[truncated Perron formula; Montgomery--Vaughan,
\emph{Multiplicative number theory I}, Corollary~5.3]\label{i4:lem-perron}
Let $\alpha(s)=\sum_{n\ge1}b_nn^{-s}$ converge absolutely for $\Rel s>1$,
let $\sigma_0>1$, $u\ge2$ and $T>0$. Then
\[
\sum_{n\le u}{}'\,b_n=\frac1{2\pi i}\int_{\sigma_0-iT}^{\sigma_0+iT}\alpha(s)\frac{u^s}s\,ds+R,
\]
with $|R|\le C\Bigl(\sum_{u/2<n<2u,\,n\ne u}|b_n|\min\bigl(1,\tfrac{u}{T|u-n|}\bigr)
+\frac{4^{\sigma_0}+u^{\sigma_0}}T\sum_n|b_n|n^{-\sigma_0}\Bigr)$; the
prime on the sum means that a term $n=u$ is counted with weight $1/2$.
\end{lemma}

Jensen's formula gives, in the situation of Lemma~\ref{i4:lem-landau},
\begin{equation}\label{i4:eq-jensen}
\#\{\rho:|\rho-s_0|\le r/2\}\le M/\log2 ,
\end{equation}
since $\sum_{|\rho-s_0|<r}\log(r/|\rho-s_0|)=\frac1{2\pi}\int_0^{2\pi}\log|f(s_0+re^{i\theta})/f(s_0)|\,d\theta\le M$
and each zero with $|\rho-s_0|\le r/2$ contributes at least $\log2$.

\medskip\noindent\emph{Step 0: elementary properties.} Let $A_g(x)=\sum_{1\le n\le x}a(g,n)$
and $A_\chi(x)=\sum_{1\le n\le x}a_\chi(n)=\sum_g\chi(g)A_g(x)$.

\emph{(B1) Partial sums.} For $x\ge1$, $|A_\chi(x)-\mathrm{res}_\chi x|\le10|G|\sqrt x$.
Indeed $\sum_g\chi(g)=|G|\delta_{\chi=1}$, so
$A_\chi(x)-\mathrm{res}_\chi x=\sum_g\chi(g)(A_g(x)-\kappa x)$, and each
term has modulus at most $10\sqrt x$. Hence the integral defining
$\mathrm{mel}_\chi$ converges absolutely and locally uniformly for
$\sigma>1/2$, $\mathrm{mel}_\chi$ and $F_\chi$ are holomorphic there, and
$|\mathrm{mel}_\chi(s)|\le10|G|/(\sigma-1/2)$. With $\kappa|G|\le4B$ and
$|G|\le B$ this gives $|F_1(s)|\le4B|s|+10B|s||s-1|/(\sigma-1/2)$ and, for
$\chi\ne1$, $|F_\chi(s)|\le10B|s|/(\sigma-1/2)$. On $5/8\le\sigma\le2$ we have
$1/(\sigma-1/2)\le8$ and $|s-1|\le|s|\le|t|+2$, with $|t|+2\ge2$; so
$|F_1(s)|\le4B(|t|+2)+80B(|t|+2)^2\le82B(|t|+2)^2$,
$|F_\chi(s)|\le80B(|t|+2)\le40B(|t|+2)^2$ for $\chi\ne1$, and in both cases
\begin{equation}\label{i4:eq-Fbound}
|F_\chi(s)|\le90B(|t|+2)^2\qquad(5/8\le\sigma\le2).
\end{equation}

\emph{(B2) Dirichlet series.} Since $a\ge0$ and $A_g(x)\le\kappa x+10\sqrt x$,
$L_\chi(s):=\sum_na_\chi(n)n^{-s}$ converges absolutely for $\sigma>1$, and
partial summation gives
$L_\chi(s)=s\int_1^\infty A_\chi(x)x^{-s-1}dx=s\,\mathrm{mel}_\chi(s)+\mathrm{res}_\chi s/(s-1)$,
that is $F_\chi(s)=(s-1)^{\delta_{\chi=1}}L_\chi(s)$ for $\sigma>1$. In
particular $F_1(1)=\kappa|G|>0$. If $\chi$ is real-valued ($\chi^2=1$),
$a_\chi$ is real, so $F_\chi(\bar s)=\overline{F_\chi(s)}$ and $F_\chi$ is
real on $(1/2,\infty)$.

\emph{(B3) Logarithmic derivative.} Since $\chi(g)=\chi(h)\chi(g-h)$,
multiplying the convolution axiom by $\chi(g)$ and summing over $g$ gives
$a_\chi(n)\log n=\sum_{n_1n_2=n}\lambda_\chi(n_1)a_\chi(n_2)$. Also
$|\lambda_\chi(n)|\le\sum_g\lambda(g,n)\le2\Lambda(n)$; in particular
$\lambda(g,1)=0$. Put $\Lambda_\chi(s)=\sum_n\lambda_\chi(n)n^{-s}$ and
$P_\chi(s)=\sum_{n\ge2}\lambda_\chi(n)(\log n)^{-1}n^{-s}$, absolutely
convergent for $\sigma>1$. The identity says $-L_\chi'=\Lambda_\chi L_\chi$,
and $P_\chi'=-\Lambda_\chi$; so $L_\chi e^{-P_\chi}$ has derivative $0$ and
is constant on $\sigma>1$, and letting $\sigma\to\infty$ ($L_\chi\to
a_\chi(1)=\chi(0)=1$, $P_\chi\to0$) the constant is $1$. Hence
$L_\chi=e^{P_\chi}$ has no zero in $\sigma>1$,
$|P_\chi(s)|\le2\sum_n\Lambda(n)(\log n)^{-1}n^{-\sigma}=2\log\zeta(\sigma)$,
and, for $\sigma>1$,
\begin{equation}\label{i4:eq-LambdaF}
\Lambda_\chi(s)=-\frac{L_\chi'}{L_\chi}(s)=-\frac{F_\chi'}{F_\chi}(s)+\frac{\delta_{\chi=1}}{s-1}.
\end{equation}
Consequently every zero of $F_\chi$ in $\sigma>1/2$ has $\Rel\rho\le1$.

\emph{(B4) Positivity.} If $w:G\to\C$ satisfies $\Rel w(g)\ge0$ for all
$g$, then $\Rel\sum_gw(g)\lambda(g,n)\ge0$ for every $n$, since
$\lambda\ge0$.

\medskip\noindent\emph{Step 1: the local lemma.} For $t\in\R$ and $1<\sigma_0\le3/2$
put $s_0=\sigma_0+it$ and $W=\ell(t)+\log\frac1{\sigma_0-1}$. Let
$\mathcal Z(s_0)$ be the multiset of zeros $\rho$ of $F_\chi$ with
$|\rho-s_0|\le1/8$.

\begin{lemma}\label{i4:lem-lfam-local}
There is an absolute $C_Z\ge1$ such that, for every $L$-family, every
$\chi$ and $s_0$ as above: $\#\mathcal Z(s_0)\le C_ZW$;
for $|s-s_0|\le1/16$,
$\bigl|F_\chi'/F_\chi(s)-\sum_{\rho\in\mathcal Z(s_0)}(s-\rho)^{-1}\bigr|\le C_ZW$;
and for every sub-multiset $S\subseteq\mathcal Z(s_0)$,
\[
-\Rel\frac{F_\chi'}{F_\chi}(s_0)\le C_ZW-\sum_{\rho\in S}\Rel\frac1{s_0-\rho}.
\]
\end{lemma}
\begin{proof}
On $|s-s_0|\le1/4$ we have $3/4\le\sigma\le7/4$ and $|\Iml s|\le|t|+1/4$,
so by \eqref{i4:eq-Fbound} $|F_\chi(s)|\le114B(|t|+2)^2$ and
$\log|F_\chi(s)|\le2\ell(t)+5$. At $s_0$, by (B2), (B3) and
$\zeta(\sigma)\le\sigma/(\sigma-1)$,
$|L_\chi(s_0)|\ge\zeta(\sigma_0)^{-2}\ge\frac49(\sigma_0-1)^2$, and
$|F_\chi(s_0)|\ge\frac49(\sigma_0-1)^3$ in both cases ($|s_0-1|\ge\sigma_0-1$
and $\sigma_0-1<1$). So $|F_\chi(s)/F_\chi(s_0)|\le e^M$ with
$M=2\ell(t)+6+3\log\frac1{\sigma_0-1}\le5W$ (as $\ell>2$). Lemma~\ref{i4:lem-landau}
with $r=1/4$ and \eqref{i4:eq-jensen} give the first two claims with
$C_Z=\max(20A_T,8)$. For the third take $s=s_0$: every
$\rho\in\mathcal Z(s_0)$ has $\Rel\rho\le1<\sigma_0$ by (B3), so
$\Rel(s_0-\rho)^{-1}=(\sigma_0-\Rel\rho)/|s_0-\rho|^2\ge0$, and the terms not
in $S$ may be dropped.
\end{proof}

By \eqref{i4:eq-LambdaF}, Lemma~\ref{i4:lem-lfam-local} bounds
$\Rel\Lambda_\psi(s_0)=-\Rel F_\psi'/F_\psi(s_0)+\delta_{\psi=1}\Rel(s_0-1)^{-1}$
for every character $\psi$.

\emph{Choice of $\eta$.} For $0<\eta\le1$ put $K(\eta)=3C_Z(1+\log(1/\eta))$.
Since $\eta K(\eta)\to0$ as $\eta\to0$, fix $\eta\in(0,1/8]$ with
$\eta K(\eta)\le1/100$, and put $\eta'=\eta/8$. Then
$K(\eta')\le4K(\eta)$ (as $1+\log8+y\le4(1+y)$ for $y\ge0$), so
$\eta'K(\eta')\le1/200$. If $\sigma_0=1+\vartheta/\ell(\gamma)$ with
$\vartheta\in\{\eta,\eta'\}$ and $t\in\{0,\gamma,2\gamma\}$, then, using
$\ell(0)\le\ell(\gamma)$, $\ell(2\gamma)\le\ell(\gamma)+\log2\le2\ell(\gamma)$
and $\log\ell\le\ell$, we get $W\le3\ell(\gamma)(1+\log(1/\vartheta))$ and so
\begin{equation}\label{i4:eq-CZW}
C_ZW\le K(\vartheta)\,\ell(\gamma).
\end{equation}
Also $\sigma_0\le3/2$ because $\vartheta\le1/8$ and $\ell>2$.

\medskip\noindent\emph{Step 2: proof of (F1).} Let $\rho=\beta+i\gamma$ be a zero of
$F_\chi$ with $\beta\ge3/4$ and ($\chi^2\ne1$ or $\gamma\ne0$). By (B3)
$\beta\le1$. We show $1-\beta>c_1/\ell(\gamma)$ with $c_1=\eta'/8=\eta/64$.
Suppose not.

\emph{Case 1: $\chi^2\ne1$.} Take $\sigma=1+\eta/\ell(\gamma)$. For
$w(g)=3+4\chi(g)n^{-i\gamma}+\chi(g)^2n^{-2i\gamma}$ we have $\Rel
w=3+4\cos\theta+\cos2\theta=2(1+\cos\theta)^2\ge0$, so by (B4)
$3\Lambda_1(\sigma)+4\Rel\Lambda_\chi(\sigma+i\gamma)+\Rel\Lambda_{\chi^2}(\sigma+2i\gamma)\ge0$.
Here $\chi\ne1\ne\chi^2$. By Lemma~\ref{i4:lem-lfam-local} and
\eqref{i4:eq-CZW} (with $S=\emptyset$, except $S=\{\rho\}$ at
$s_0=\sigma+i\gamma$, where $|\rho-s_0|=\sigma-\beta\le\frac98\eta/\ell(\gamma)<1/8$):
\[
0\le3\Bigl(\frac{\ell(\gamma)}\eta+K\ell(\gamma)\Bigr)+4\Bigl(K\ell(\gamma)-\frac1{\sigma-\beta}\Bigr)+K\ell(\gamma),\qquad K=K(\eta).
\]
As $\sigma-\beta\le\frac98\eta/\ell(\gamma)$, this gives
$\frac{32}{9\eta}\le\frac3\eta+8K$, i.e.\ $\frac59\le8\eta K\le\frac8{100}$,
a contradiction.

\emph{Case 2: $\chi^2=1$, $\gamma\ne0$, $|\gamma|\le\eta/(2\ell(0))$.} Take
$\sigma=1+\eta/\ell(0)$. By (B2) $\bar\rho=\beta-i\gamma\ne\rho$ is also a
zero, and $|\rho-\sigma|,|\bar\rho-\sigma|\le(\frac98+\frac12)\eta/\ell(0)<1/8$.
With $w=1+\chi\ge0$, (B4) gives $\Lambda_1(\sigma)+\Lambda_\chi(\sigma)\ge0$
(if $\chi=1$ this is $2\Lambda_1(\sigma)\ge0$). Using $S=\{\rho,\bar\rho\}$
for $F_\chi$ and $\Rel(\sigma-\rho)^{-1}+\Rel(\sigma-\bar\rho)^{-1}=2(\sigma-\beta)/((\sigma-\beta)^2+\gamma^2)$,
in either case (dividing by $2$ if $\chi=1$)
\[
0\le\frac1{\sigma-1}+2K\ell(0)-\frac{2(\sigma-\beta)}{(\sigma-\beta)^2+\gamma^2}.
\]
Since $|\gamma|\le\eta/(2\ell(0))\le(\sigma-\beta)/2$, the last term is at
least $\frac8{5(\sigma-\beta)}\ge\frac{64}{45}\frac{\ell(0)}\eta$ (as
$1-\beta\le c_1/\ell(\gamma)\le\frac{\eta}{8\ell(0)}$). Hence
$\frac{64}{45\eta}\le\frac1\eta+2K$, i.e.\ $\frac{19}{45}\le2\eta K$, a
contradiction.

\emph{Case 3: $\chi^2=1$, $|\gamma|>\eta/(2\ell(0))$.} Take
$\sigma=1+\eta'/\ell(\gamma)$ and $K'=K(\eta')$. The $3$-$4$-$1$ inequality of
Case 1 holds, now with $\chi^2=1$. Bound the three terms:
$\Lambda_1(\sigma)\le\ell(\gamma)/\eta'+K'\ell(\gamma)$;
$\Rel\Lambda_\chi(\sigma+i\gamma)\le\delta_{\chi=1}(\sigma-1)/\gamma^2+K'\ell(\gamma)-1/(\sigma-\beta)$;
$\Rel\Lambda_1(\sigma+2i\gamma)\le(\sigma-1)/(4\gamma^2)+K'\ell(\gamma)$.
Now $4\gamma^2>\eta^2/\ell(0)^2\ge\eta^2/\ell(\gamma)^2$, so
$(\sigma-1)/(4\gamma^2)\le\eta'\ell(\gamma)/\eta^2=\ell(\gamma)/(64\eta')$
and $(\sigma-1)/\gamma^2\le\ell(\gamma)/(16\eta')$. Therefore
\[
\frac4{\sigma-\beta}\le\ell(\gamma)\Bigl(\frac3{\eta'}+\frac4{16\eta'}+\frac1{64\eta'}+8K'\Bigr),
\]
and with $\sigma-\beta\le\frac98\eta'/\ell(\gamma)$ we get
$\frac{32}9-\frac{209}{64}\le8\eta'K'\le\frac8{200}$; the left side exceeds
$0.28$, a contradiction.

\medskip\noindent\emph{Step 3: proof of (F2).} Put $c_2=\eta/8$ and
$\sigma=1+\eta/\ell(0)$, $K=K(\eta)$ (for $\gamma=0$, \eqref{i4:eq-CZW}
holds with $\ell(0)$). Let $\chi,\chi',\beta,\beta'$ be as in (F2); then
$\beta,\beta'\le1$ and $|\sigma-\beta|,|\sigma-\beta'|\le\frac98\eta/\ell(0)<1/8$.

\emph{Case $\chi\ne\chi'$.} The characters are $\pm1$-valued, so
$w=(1+\chi)(1+\chi')\ge0$ and (B4) gives
$\Lambda_1+\Lambda_\chi+\Lambda_{\chi'}+\Lambda_{\chi\chi'}\ge0$ at $\sigma$.
Here $\chi\chi'\ne1$; let $m\in\{1,2\}$ be the number of entries equal to
$1$ in the list $(1,\chi,\chi',\chi\chi')$. Then $\chi$ and $\chi'$ each
occur exactly $m$ times in the list ($m=2$ exactly when one of them is
trivial, and then the list is $(1,1,\psi,\psi)$). Applying
Lemma~\ref{i4:lem-lfam-local} with $S=\{\beta\}$ at each entry equal to $\chi$,
$S=\{\beta'\}$ at each entry equal to $\chi'$, we get
$m/(\sigma-1)+4K\ell(0)-m/(\sigma-\beta)-m/(\sigma-\beta')\ge0$.

\emph{Case $\chi=\chi'$ and ($\beta\ne\beta'$ or $F_\chi'(\beta)=0$).} Then
$\{\beta,\beta'\}$ (with $\beta$ counted twice in the second alternative)
is a sub-multiset of $\mathcal Z(\sigma)$ for $F_\chi$, and $w=1+\chi\ge0$ gives
$\Lambda_1(\sigma)+\Lambda_\chi(\sigma)\ge0$, hence, as in Case 2 of Step 2,
$1/(\sigma-1)+2K\ell(0)-1/(\sigma-\beta)-1/(\sigma-\beta')\ge0$.

In both cases $\frac1{\sigma-\beta}+\frac1{\sigma-\beta'}\le\frac{\ell(0)}\eta+4K\ell(0)$,
while the left side is at least $\frac{16}9\frac{\ell(0)}\eta$; so
$\frac79\le4\eta K\le\frac4{100}$, a contradiction. Hence $\chi=\chi'$,
$\beta=\beta'$ and $F_\chi'(\beta)\ne0$.

\medskip\noindent\emph{Step 4: proof of (F3), with $A=46$.} Let $\chi^2=1$, $\beta\ge3/4$
real, $F_\chi(\beta)=0$; then $\beta\le1$ by (B3).

\emph{$\chi=1$.} By (B1), $\mathrm{mel}_1(\beta)$ is real with
$|\mathrm{mel}_1(\beta)|\le40|G|$, so
$0=F_1(\beta)\ge\beta|G|(\kappa-40(1-\beta))$, whence
$1-\beta\ge\kappa/40\ge1/(40B)$.

\emph{$\chi\ne1$.} The argument has two parts: a lower bound
$F_\chi(1)\ge c'B^{-45}$, obtained from the theta axiom and the first
positivity axiom through an integral kernel, and the mean value theorem.
Since $\chi$ is real, $a_\chi$ is real-valued. For $t>0$ put
$\Theta_\chi(t)=\sum_na_\chi(n)e^{-tn}$.

\emph{(i) The theta functions.} By the count axiom,
$a(g,n)\le A_g(n)\le\kappa n+10\sqrt n\le14n$ for $n\ge1$, so
$|a_\chi(n)|\le14|G|n\le14Bn$ for every character $\chi$, and
$\Theta_\chi(t)$ converges absolutely for $t>0$. Let
$S_0=14\sum_{n\ge1}ne^{-n/2}$ and $C_1=2S_0+4$. We claim, for every
character $\chi$ of $G$,
\begin{equation}\label{i4:eq-theta}
|\Theta_\chi(t)-\mathrm{res}_\chi/t|\le C_1B^{11}\min(t,1/t)\quad(t>0),\qquad
|\Theta_\chi(t)|\le S_0Be^{-t/2}\quad(t\ge1,\ \chi\ne1).
\end{equation}
For $0<t\le1$: $\Theta_\chi(t)-\mathrm{res}_\chi/t=\sum_g\chi(g)\bigl(\sum_na(g,n)e^{-tn}-\kappa/t\bigr)$,
and the theta axiom bounds each of the $|G|\le B$ terms by $B^{10}t$.
For $t\ge1$ and $n\ge1$, $e^{-tn}\le e^{-t/2}e^{-n/2}$, so
$|\Theta_\chi(t)|\le14B\sum_nne^{-n/2}e^{-t/2}=S_0Be^{-t/2}$; this is the
second bound, and with $e^{-t/2}\le2/t$ and $\mathrm{res}_\chi\le\kappa|G|\le4B$
it gives $|\Theta_\chi(t)-\mathrm{res}_\chi/t|\le(2S_0+4)B/t\le C_1B^{11}/t$.

\emph{(ii) The integral of $\Theta_\chi$.} Let $\chi\ne1$. Then $\Theta_\chi$ is
integrable on $(0,\infty)$ by \eqref{i4:eq-theta}, and
\[
\int_0^\infty\Theta_\chi(t)\,dt=F_\chi(1).
\]
Indeed, by partial summation ($a_\chi(0)=0$),
$\Theta_\chi(t)=\int_1^\infty A_\chi(x)\,te^{-tx}dx$; by (B1) with
$\mathrm{res}_\chi=0$, $|A_\chi(x)|\le10B\sqrt x$, and
$\int_0^\infty te^{-tx}dt=x^{-2}$, so Fubini's theorem applies
($\int_1^\infty10B\sqrt x\,x^{-2}dx<\infty$) and gives
$\int_0^\infty\Theta_\chi(t)dt=\int_1^\infty A_\chi(x)x^{-2}dx=\mathrm{mel}_\chi(1)=F_\chi(1)$.

\emph{(iii) The kernel.} For $y>0$ put
$W(y)=\int_0^\infty e^{-y/w}e^{-w}\,\frac{dw}w$ (this is $2K_0(2\sqrt y)$,
with $K_0$ the Macdonald function; we use only the following properties).
$W(y)>0$; $W$ is decreasing, so $W(v)\ge w_0:=W(1)>0$ for $0<v\le1$; from
$e^{-y/w}\le6w^3/y^3$, $W(y)\le6\Gamma(3)y^{-3}=12y^{-3}$; and the
substitution $w\mapsto wn$ gives, for $y>0$ and $n\ge1$,
$\int_0^\infty e^{-y/w}e^{-wn}\frac{dw}w=W(yn)$.

\emph{(iv) Lower bound by positivity.} Let $f(N)=\sum_{n_1n_2=N}a_1(n_1)a_\chi(n_2)$.
By the first positivity axiom, $f\ge0$; and $f(1)=a_1(1)a_\chi(1)=1$. Let
$0<v\le1$. Expanding the product of the two theta series and integrating
term by term,
\[
\int_0^\infty\Theta_1(v/w)\Theta_\chi(w)\frac{dw}w
=\sum_{m,n\ge1}a_1(m)a_\chi(n)W(vmn)=\sum_{N\ge1}f(N)W(vN);
\]
the interchange is justified by
$\sum_{m,n}|a_1(m)||a_\chi(n)|W(vmn)\le\sum_{m,n}(14B)^2mn\cdot12(vmn)^{-3}<\infty$.
All terms of the last sum are $\ge0$, so it is at least the term $N=1$:
\[
\int_0^\infty\Theta_1(v/w)\Theta_\chi(w)\frac{dw}w\ge W(v)\ge w_0 .
\]

\emph{(v) Upper bound.} Write $\Theta_1(u)=\kappa|G|/u+E(u)$; by
\eqref{i4:eq-theta} and $\min(u,1/u)\le\sqrt u$,
$|E(v/w)|\le C_1B^{11}\sqrt v/\sqrt w$, and $|\Theta_\chi(w)|\le
C_1B^{11}\min(w,1/w)$. The main term contributes
$\frac{\kappa|G|}v\int_0^\infty\Theta_\chi(w)dw=\frac{\kappa|G|}vF_\chi(1)$ by
(ii), so
\[
\int_0^\infty\Theta_1(v/w)\Theta_\chi(w)\frac{dw}w\le\frac{\kappa|G|}vF_\chi(1)+(C_1B^{11})^2\sqrt v\,K_W,
\qquad K_W=\int_0^\infty\min(w,1/w)\,w^{-3/2}dw=\tfrac83 .
\]

\emph{(vi) The value at $1$.} Let $\delta=\min\bigl(1,w_0/(2(C_1^2K_W+1))\bigr)$
and $v=(\delta/B^{22})^2\le1$. Then $(C_1B^{11})^2\sqrt vK_W=C_1^2K_W\delta\le w_0/2$,
so (iv) and (v) give $w_0/2\le\kappa|G|F_\chi(1)/v$. With $\kappa|G|\le4B$,
\[
F_\chi(1)\ge\frac{w_0v}{8B}=\frac{c'}{B^{45}},\qquad c'=\frac{\delta^2w_0}8 ,
\]
an absolute constant.

\emph{(vii) Mean value theorem.} $F_\chi$ is real on $(1/2,\infty)$ by (B2), so
$F_\chi(1)=F_\chi(1)-F_\chi(\beta)=(1-\beta)F_\chi'(\xi)$ for some
$\xi\in[\beta,1]$. Cauchy's estimate on the circle $|s-\xi|=1/8$ (on which
$5/8\le\sigma\le9/8$ and $|t|\le1/8$) with \eqref{i4:eq-Fbound} gives
$|F_\chi'(\xi)|\le8\cdot90B(17/8)^2\le3300B$. Hence
$1-\beta\ge F_\chi(1)/(3300B)\ge c'/(3300B^{46})$.

\emph{Conclusion.} Together with the case $\chi=1$
($1/(40B)\ge1/(40B^{46})$), (F3) holds with $A=46$ and
$c_3=\min(1/40,c'/3300)$.

\medskip\noindent\emph{Step 5: proof of (F4).} Let $c_0\in(0,1/20]$, let
$\mathrm{Exc}_\chi(c_0,Z)$ hold, $u\ge2$, $\log B\le\lambda:=\sqrt{\log u}$,
and put $T=e^\lambda$ ($\ge2$), $\sigma_0=1+1/\log u$. Then
$\ell(T)\le\lambda+\log(T+2)\le2\lambda+1\le4\lambda$ (as
$\lambda\ge\sqrt{\log2}$).

\emph{Perron.} Apply Lemma~\ref{i4:lem-perron} to $\Lambda_\chi(s)$, whose
coefficients satisfy $|\lambda_\chi(n)|\le2\Lambda(n)$. On $u/2<n<2u$,
$|\lambda_\chi(n)|\le2\log(2u)$, and $\sum_{u/2<n<2u,\,n\ne u}\min(1,u/(T|u-n|))\le
3+\frac{4u}T(1+\log u)$; also $\sum_n\Lambda(n)n^{-\sigma_0}\le\log u+C$.
The half-weighted term differs from $\psi_\chi(u)$ by at most $2\log u$. So
\[
\psi_\chi(u)=\frac1{2\pi i}\int_{\sigma_0-iT}^{\sigma_0+iT}\Lambda_\chi(s)\frac{u^s}s\,ds+R_1,\qquad|R_1|\le C\Bigl(\frac{u\log^2u}T+\log u\Bigr).
\]

\emph{The contour.} Of the two numbers $1-c_0/(4\ell(T))$ and
$1-c_0/(8\ell(T))$, let $\sigma_1$ be one at distance at least
$c_0/(16\ell(T))$ from the element of $Z$ (if $Z$ is empty, either). Let
$\mathcal R$ be the rectangle $\sigma_1\le\sigma\le\sigma_0$, $|t|\le T$.

\emph{Zeros in $\mathcal R$.} If $F_\chi(\rho)=0$, $\rho\in\mathcal R$, then
$\Rel\rho\ge\sigma_1\ge1-c_0/\ell(T)\ge1-c_0/\ell(\Iml\rho)$, so $\rho\in Z$
by $\mathrm{Exc}_\chi(c_0,Z)$: it is the real simple zero $\beta$, with
$\beta>\sigma_1$ (and $\beta<1$).

\emph{Bound on the three other sides.} Let $s=\sigma+it$ with $|t|\le T$ and
$\sigma_1\le\sigma\le\sigma_0$, lying on the line $\sigma=\sigma_1$ or on
$|t|=T$. If $\sigma\ge1+\eta_0/\ell(t)$ with $\eta_0=1/20$, then
$|\Lambda_\chi(s)|\le2\sum\Lambda(n)n^{-\sigma}\le2\ell(t)/\eta_0+C$. Otherwise
take $s_0=1+\eta_0/\ell(t)+it$; then $|s-s_0|\le(\eta_0+c_0/4)/\ell(t)<1/16$
and $W=\ell(t)+\log(\ell(t)/\eta_0)\le4\ell(T)$. Each $\rho\in\mathcal
Z(s_0)$ has $|\Iml\rho|\le T+1/8$, so $\ell(\Iml\rho)\le2\ell(T)$; if
$\rho\notin Z$ then by $\mathrm{Exc}_\chi$,
$\Rel\rho<1-c_0/(2\ell(T))$, so $\Rel(s-\rho)>c_0/(4\ell(T))$ and
$|s-\rho|^{-1}<4\ell(T)/c_0$; if $\rho=\beta\in Z$ then $|s-\beta|\ge
c_0/(16\ell(T))$ (on $\sigma=\sigma_1$ by the choice of $\sigma_1$, on
$|t|=T$ because $T\ge2$). By Lemma~\ref{i4:lem-lfam-local},
$|F_\chi'/F_\chi(s)|\le C_ZW+C_ZW\cdot16\ell(T)/c_0\le C\ell(T)^2/c_0$.
Finally $|s-1|^{-1}\le8\ell(T)/c_0$ on these sides. By \eqref{i4:eq-LambdaF},
$|\Lambda_\chi(s)|\le C_4\ell(T)^2/c_0$ on the three sides.

\emph{Residues.} Inside $\mathcal R$ the function
$\Lambda_\chi(s)u^s/s=\bigl(-F_\chi'/F_\chi(s)+\delta_{\chi=1}/(s-1)\bigr)u^s/s$
has a pole at $s=1$ with residue $\delta_{\chi=1}u$ (recall
$F_1(1)\ne0$) and, if $\beta\in Z$ with $\beta>\sigma_1$, a pole at $\beta$
with residue $-u^\beta/\beta$ (simple zero); $s=0$ is outside. If
$\beta\in Z$ and $\beta<\sigma_1$, then $|u^\beta/\beta|\le\frac43u^{\sigma_1}$.
By the residue theorem and the bounds on the sides,
\[
\Bigl|\psi_\chi(u)-\delta_{\chi=1}u+\sum_{\beta\in Z}\frac{u^\beta}\beta\Bigr|
\le|R_1|+\frac43u^{\sigma_1}+\frac{C_4\ell(T)^2}{c_0}\Bigl(u^{\sigma_1}\bigl(\tfrac83+2\log T\bigr)+\frac{2e\,u}T\Bigr).
\]
Here the term $2eu/T$ bounds the two horizontal sides $|t|=T$: there
$|s|\ge T$, $u^\sigma\le u^{\sigma_0}=eu$, and their length is
$\sigma_0-\sigma_1\le1/\log u+c_0/4<1$, since
$\log u\ge(\log B)^2\ge(\log4)^2>1.9$.

\emph{Conclusion.} $u^{\sigma_1}\le u\exp(-c_0\log u/(8\ell(T)))\le ue^{-c_0\lambda/32}$,
$\ell(T)^2\le16\lambda^2$, $\log T=\lambda$, $u\log^2u/T=u\lambda^4e^{-\lambda}$,
and $\log u=\lambda^2\le C'u e^{-\lambda}$. Every term is at most
$C(c_0)\,\lambda^4ue^{-c_0\lambda/32}\le C(c_0)'ue^{-c_0\lambda/64}$. This is
(F4) with $c=c_0/64$. \qed



\subsubsection{The ray class family}\label{i4:ssec-proof-lfamray}

\begin{proof}[Proof of Proposition~\ref{i4:prop-lfam-ray}]
\emph{Constants.} $B=4dM^2\ge3600\ge4$. By (R3), $|G|\le h_KM^2\le4dM^2=B$.
$\sqrt d\ge10$ and $\sqrt d\le d$. Then $\kappa\ge1/B$ means
$2\pi\cdot4dM^2\ge\sqrt dM^2$, true; $\kappa\le4$ means $2\pi\le4\sqrt dM^2$, true.

\emph{Values at $n=0,1$.} No nonzero ideal has norm $0$, so $a(g,0)=0$. The
only ideal of norm $1$ is $\mathcal O_K$, which is prime to $M$ and in the
trivial class; so $a(0,1)=1$ and $a(g,1)=0$ for $g\ne0$. Clearly $a\ge0$.

\emph{Count and theta, positivity.} These are (R1) and (R2) (for a
character $\chi$ with $\chi^2=1$, $\varepsilon=\Rel\chi=\chi$ takes values $\pm1$
and is multiplicative).

\emph{$\lambda$.} $\lambda\ge0$ as $\log N\mathfrak p\ge0$. Norms of prime
powers $\mathfrak p^j$ are prime powers, so $\sum_g\lambda(g,n)=0=\Lambda(n)$
unless $n=p^i$. For $n=p^i$: at most two primes lie over $p$; for each,
at most one power has norm $n$. If $p$ splits, the two primes have norm
$p$ and contribute at most $2\log p$; if $p$ ramifies, at most $\log p$;
if $p$ is inert, $N\mathfrak p=p^2$ and the contribution is at most $2\log p$.
So $\sum_g\lambda(g,n)\le2\Lambda(n)$.

\emph{Convolution.} For a nonzero ideal $I$, $\log NI=\sum_{\mathfrak p}v_{\mathfrak p}(I)\log N\mathfrak p=\sum_{J\mid I}\lambda_0(J)$,
where $\lambda_0(\mathfrak p^j)=\log N\mathfrak p$ ($j\ge1$) and $\lambda_0=0$ on other ideals. Summing
over the ideals $I$ prime to $M$ of norm $n$ in class $g$, and writing
$I=JI'$ (both factors prime to $M$, classes $h$ and $g-h$, norms
$n_1n_2=n$), gives
$a(g,n)\log n=\sum_h\sum_{n_1n_2=n}\lambda(h,n_1)a(g-h,n_2)$.

\emph{Norm form.} Let $\lambda(0,p)\ne0$. Then some $\mathfrak p^j$ of norm
$p$ lies in the trivial class; so $j=1$, $N\mathfrak p=p$, and
$(\alpha)\mathfrak p=(\beta)$ with $\alpha\equiv\beta\equiv1\bmod M$. Then
$\mathfrak p=(\pi)$ with $\pi=\beta/\alpha\in\mathcal O_K$, and
$\alpha(\pi-1)=\beta-\alpha\in M\mathcal O_K$; since $\alpha$ is prime to $M$,
$\pi-1\in M\mathcal O_K$. With $\omega=(1+\sqrt{-d})/2$,
$\mathcal O_K=\Z+\Z\omega$, write $\pi=1+M(u+v\omega)$, and
$\pi=(A_0+B_0\sqrt{-d})/2$ with $A_0=2+2Mu+Mv$, $B_0=Mv$. Then $M\mid B_0$,
$A_0-2-B_0=2Mu$, and $4p=4N\pi=A_0^2+dB_0^2$.
\end{proof}

\subsubsection{Proof of Proposition~\ref{i4:prop-ray-leaves}}\label{i4:Rst-ray}\label{i4:Rst-ssec}

All three items are classical facts about an imaginary quadratic field,
and the formal proof follows the classical route: (R1) is lattice-point
counting in an ellipse and Poisson summation for a Gaussian, (R2) is the
positivity of the coefficients of $\zeta_K(s)L(s,\varepsilon)$ prime by
prime, and (R3) is the exact sequence of the ray class group together with
Minkowski's bound. We give complete proofs; the only inputs are the
published theorems stated below.

\emph{Structure.}
\begin{enumerate}
\item[(R1.a)] (L) arithmetic of $K=\Q(\sqrt{-d})$: ring of integers,
units, norms, covolumes (Theorem~\ref{i4:Rst-thm-quad}).
\item[(R1.b)] (R) ideals of norm $n$ in the class $g$ correspond bijectively
to the points of a translated lattice $x_0+M\mathfrak b$ of norm
$nN\mathfrak b$ (Lemma~\ref{i4:Rst-lem-bij}).
\item[(R1.c)] (L) a reduced basis of a plane lattice
(Lagrange--Gauss; proved in Lemma~\ref{i4:Rst-lem-reduced}), giving
$a(g,n)=\#\{v\in\Z^2:Q(v)=n\}$ for a shifted binary form $Q$ with
$\alpha,\gamma\ge1$, $\alpha\gamma=dM^4/4$ (Lemma~\ref{i4:Rst-lem-form}).
\item[(R1.d)] (R) counting points of an ellipse row by row
(Lemma~\ref{i4:Rst-lem-count}).
\item[(R1.e)] (L) Poisson summation for the Gaussian
(Theorem~\ref{i4:Rst-thm-poisson}), applied in each variable
(Lemma~\ref{i4:Rst-lem-theta}).
\item[(R2)] (R) multiplicativity of $\sum_{J\mid C}\varepsilon(J)$ and the
geometric sums $1+x+\dots+x^k$, $x\in\{0,\pm1\}$.
\item[(R3)] (L) Minkowski's bound (Theorem~\ref{i4:Rst-thm-minkideal}) and
(R) the map from the ray class group to the class group, and the divisor
bound for the number of ideals of given norm.
\end{enumerate}
Nothing in this proposition is new.

\begin{theorem}[arithmetic of imaginary quadratic fields; Neukirch I (3.6), (5.2)]\label{i4:Rst-thm-quad}
Let $d\equiv3\pmod4$ be squarefree, $d>3$, $K=\Q(\sqrt{-d})\subset\C$ and
$\omega=(1+\sqrt{-d})/2$. Then $\mathcal O_K=\Z+\Z\omega$, $d_K=-d$, and
$\mathcal O_K^\times=\{\pm1\}$. For $\xi\in\mathcal O_K\smallsetminus\{0\}$,
$N((\xi))=|\xi|^2$; the absolute norm of ideals is multiplicative, and
$\mathfrak a\supseteq\mathfrak b$ implies $N\mathfrak a\mid N\mathfrak b$. A nonzero ideal $\mathfrak a$ is a
lattice in $\C=\R^2$ of covolume $N\mathfrak a\cdot\sqrt d/2$. If
$\mathfrak a+\mathfrak b=\mathcal O_K$ then $\mathfrak a\cap\mathfrak b=\mathfrak a\mathfrak b$ and (Chinese remainder theorem,
Neukirch I (3.6)) for every $y\in\mathcal O_K$ there is $x\in\mathfrak a$ with
$x-y\in\mathfrak b$.
\end{theorem}
\noindent Here $d\ge100$ is prime, so the hypotheses hold.

\begin{lemma}[ideals in a ray class]\label{i4:Rst-lem-bij}
Let $M\ge3$, $g\in G$, and let $\mathfrak b$ be an ideal prime to $M$ in the
class $-g$. Choose $x_0\in\mathfrak b$ with $x_0-1\in M\mathcal O_K$. Then
$\xi\mapsto I=(\xi)\mathfrak b^{-1}$ is a bijection from $x_0+M\mathfrak b$ onto the set of
ideals prime to $M$ in the class $g$, and $NI=|\xi|^2/N\mathfrak b$.
\end{lemma}

\begin{proof}
Since $\mathfrak b+M\mathcal O_K=\mathcal O_K$, $x_0$ exists and
$\{\xi\in\mathfrak b:\xi-1\in M\mathcal O_K\}=x_0+(\mathfrak b\cap M\mathcal O_K)=x_0+M\mathfrak b$
(Theorem~\ref{i4:Rst-thm-quad}). Such $\xi$ is nonzero ($M\nmid1$), $(\xi)$ is
prime to $M$ and divisible by $\mathfrak b$, so $I=(\xi)\mathfrak b^{-1}$ is an ideal prime to
$M$; as $I\mathfrak b=(\xi)$ with $\xi\equiv1\bmod M$, $I$ lies in the class $g$.
Conversely, if $I$ is in the class $g$, then $I\mathfrak b$ is in the trivial
class: $(\alpha)I\mathfrak b=(\beta)$ with $\alpha\equiv\beta\equiv1\bmod M$; then
$\xi=\beta/\alpha$ generates the integral ideal $I\mathfrak b$, so $\xi\in\mathfrak b$, and
$\alpha(\xi-1)=\beta-\alpha\in M\mathcal O_K$ with $\alpha$ prime to $M$ gives
$\xi-1\in M\mathcal O_K$. If $(\xi)=(\xi')$ for two such elements, $\xi'=u\xi$
with $u=\pm1$ and $(u-1)\xi\in M\mathcal O_K$; as $\xi$ is prime to $M$,
$M\mid u-1$, so $u=1$ ($M\ge3$). Finally $NI=N(\xi)/N\mathfrak b=|\xi|^2/N\mathfrak b$.
\end{proof}

\begin{lemma}[reduced basis]\label{i4:Rst-lem-reduced}
Every lattice $\Lambda\subset\C$ of covolume $\Delta$ has a basis $f_1,f_2$
with $|f_1|^2\le(2/\sqrt3)\Delta\le2\Delta$.
\end{lemma}

\begin{proof}
Let $f_1$ be a shortest nonzero vector; it is primitive, so it extends to a
basis $f_1,f_2$, and replacing $f_2$ by $f_2-nf_1$ we may assume
$|\Rel(\bar f_1f_2)|\le|f_1|^2/2$. Also $|f_2|\ge|f_1|$. Hence
$\Delta^2=|f_1|^2|f_2|^2-(\Rel\bar f_1f_2)^2\ge\frac34|f_1|^4$.
\end{proof}

\begin{lemma}[the shifted binary form]\label{i4:Rst-lem-form}
In the situation of Lemma~\ref{i4:Rst-lem-bij} there are reals
$\alpha,\gamma\ge1$ with $\alpha\gamma=dM^4/4$ and $l,\theta_1,\theta_2$ such
that, with $Q(v)=\alpha(v_1+\theta_1+lv_2)^2+\gamma(v_2+\theta_2)^2$,
$Q(\Z^2)\subset\{1,2,3,\dots\}$ and $a(g,n)=\#\{v\in\Z^2:Q(v)=n\}$ for all $n$.
\end{lemma}

\begin{proof}
The lattice $M\mathfrak b$ has covolume $\Delta=M^2N\mathfrak b\sqrt d/2$; take a basis
$f_1,f_2$ as in Lemma~\ref{i4:Rst-lem-reduced}. Write $f_2/f_1=l+ih$
($|h||f_1|^2=\Delta$) and $x_0/f_1=u_1+iu_2$. For $\xi=x_0+v_1f_1+v_2f_2$,
\[
|\xi|^2=|f_1|^2\bigl((v_1+u_1+lv_2)^2+(u_2+hv_2)^2\bigr)
=|f_1|^2(v_1+\theta_1+lv_2)^2+\frac{\Delta^2}{|f_1|^2}(v_2+\theta_2)^2,
\]
$\theta_1=u_1$, $\theta_2=u_2/h$. Put $\alpha=|f_1|^2/N\mathfrak b$ and
$\gamma=\Delta^2/(|f_1|^2N\mathfrak b)$; then $Q(v)=|\xi|^2/N\mathfrak b=NI$ by
Lemma~\ref{i4:Rst-lem-bij}, which is a positive integer, and
$a(g,n)=\#\{v:Q(v)=n\}$. Next $\alpha\gamma=\Delta^2/N\mathfrak b^2=dM^4/4$. As
$(f_1)\subseteq M\mathfrak b$, $M^2N\mathfrak b\mid|f_1|^2$, so $\alpha\ge M^2\ge1$; and
$|f_1|^2\le2\Delta$ gives $\gamma\ge\Delta/(2N\mathfrak b)=M^2\sqrt d/4\ge1$.
\end{proof}

\begin{lemma}[points of an ellipse]\label{i4:Rst-lem-count}
For $\alpha,\gamma\ge1$, any $l,\theta_1,\theta_2$ and $y\ge1$,
$\bigl|\#\{v\in\Z^2:Q(v)\le y\}-\pi y/\sqrt{\alpha\gamma}\bigr|\le5\sqrt y$.
\end{lemma}

\begin{proof}
For fixed $v_2=k$ the condition is
$\alpha(v_1+c_k)^2\le y_k:=y-\gamma(k+\theta_2)^2$; the solutions form an
interval of integers whose number $R_k$ satisfies $|R_k-w(k)|\le1$, where
$w(t)=2\sqrt{\max(0,y-\gamma(t+\theta_2)^2)/\alpha}$, and $R_k=w(k)=0$
when $y_k<0$. The $k$ with $y_k\ge0$ lie in an interval of length
$2\sqrt{y/\gamma}$, so $\sum_k|R_k-w(k)|\le2\sqrt{y/\gamma}+1$. The function
$w\ge0$ is nondecreasing, then nonincreasing; for such a function
$|\sum_{k\in\Z}w(k)-\int_\R w|\le\max w$ (compare $w(k)$ with the integral
over $[k-1,k]$ to the left of the maximum and over $[k,k+1]$ to the right;
the two integers next to the maximum cost at most $\max w$). Here
$\max w=2\sqrt{y/\alpha}$ and $\int_\R w=\pi y/\sqrt{\alpha\gamma}$ (area of an
ellipse). The total error is at most $2\sqrt y+1+2\sqrt y\le5\sqrt y$.
\end{proof}

\begin{theorem}[Poisson summation for the Gaussian; Neukirch VII (3.2)]\label{i4:Rst-thm-poisson}
For $s>0$ and $\theta\in\R$,
$\sum_{m\in\Z}e^{-s(m+\theta)^2}=\sqrt{\pi/s}\,\sum_{n\in\Z}e^{-\pi^2n^2/s}e^{2\pi in\theta}$.
\end{theorem}

\begin{lemma}[the theta bound]\label{i4:Rst-lem-theta}
For $\alpha,\gamma\ge1$ and $0<t\le1$,
$\bigl|\sum_{v\in\Z^2}e^{-tQ(v)}-\pi/(t\sqrt{\alpha\gamma})\bigr|\le3\pi(\alpha\gamma)^2t$.
\end{lemma}

\begin{proof}
By Theorem~\ref{i4:Rst-thm-poisson} and $e^{-u}\le2/u^2$,
$\bigl|\sum_me^{-s(m+\theta)^2}-\sqrt{\pi/s}\bigr|\le\sqrt{\pi/s}\,\frac{2s^2}{\pi^4}\sum_{n\ne0}\frac1{n^2}\le\sqrt{\pi/s}\,s^2$.
Write $\sum_ve^{-tQ(v)}=\sum_{v_2}w_{v_2}\sum_{v_1}u_{v_1,v_2}$ with
$u=e^{-\alpha t(v_1+\theta_1+lv_2)^2}$, $w=e^{-\gamma t(v_2+\theta_2)^2}$; the inner
sums are $A(1+O^*(a))$ and $\sum w=C(1+O^*(c))$ with $A=\sqrt{\pi/(\alpha t)}$,
$a=(\alpha t)^2$, $C=\sqrt{\pi/(\gamma t)}$, $c=(\gamma t)^2$, where $O^*(z)$ has
absolute value $\le z$. All terms are positive, so the double sum is
$AC(1+O^*(a+c+ac))$, and $AC=\pi/(t\sqrt{\alpha\gamma})$. Finally
$AC(a+c+ac)=\frac{\pi t}{\sqrt{\alpha\gamma}}(\alpha^2+\gamma^2+\alpha^2\gamma^2t^2)\le3\pi(\alpha\gamma)^2t$.
\end{proof}

\begin{proof}[Proof of (R1)]
By Lemma~\ref{i4:Rst-lem-form}, $\sum_{1\le n\le y}a(g,n)=\#\{v:Q(v)\le y\}$
and $\sum_na(g,n)e^{-tn}=\sum_ve^{-tQ(v)}$, with
$\pi/\sqrt{\alpha\gamma}=2\pi/(\sqrt dM^2)=\kappa$. Lemma~\ref{i4:Rst-lem-count}
gives the count with error $5\sqrt y\le10\sqrt y$. Lemma~\ref{i4:Rst-lem-theta}
gives the theta bound with error $3\pi(dM^4/4)^2t\le(4dM^2)^{10}t$.
\end{proof}

\begin{proof}[Proof of (R2)]
Extend $\varepsilon$ to ideals: $\chi(J)=\varepsilon(\text{class of }J)$ if $J$ is
prime to $M$, $\chi(J)=0$ otherwise; $\chi$ is completely multiplicative with
values in $\{0,\pm1\}$, and $\chi(\mathcal O_K)=1$. The map $(I,J)\mapsto(IJ,J)$
identifies pairs of ideals prime to $M$ with $NI=n_1$, $NJ=n_2$, $n_1n_2=n$,
with pairs $(C,J)$, $C$ prime to $M$ of norm $n$, $J\mid C$. Hence, for
$n\ge1$,
\[
\sum_{n_1n_2=n}a_1(n_1)a_\varepsilon(n_2)=\sum_{NC=n,\ (C,M)=1}r(C),\qquad r(C)=\sum_{J\mid C}\chi(J).
\]
By unique factorisation $r(C)=\prod_{\mathfrak p^k\parallel C}(1+\chi(\mathfrak p)+\dots+\chi(\mathfrak p)^k)$, and
for $x\in\{0,\pm1\}$ the sum $1+x+\dots+x^k$ is $\ge0$, and $\ge1$ when $k$ is
even. So every $r(C)\ge0$, which is the first claim, and $r(I^2)\ge1$. Since
$I\mapsto I^2$ is injective and maps ideals of norm $n$ prime to $M$ to
ideals of norm $n^2$ prime to $M$,
$a_1(n)\le\sum_{NI=n}r(I^2)\le\sum_{NC=n^2}r(C)$, the second claim. For $n=0$
both sides vanish ($a(g,0)=0$).
\end{proof}

\begin{theorem}[Minkowski's bound for ideal classes; Neukirch I \S5, Exercise~3, with the proof of I (6.3)]\label{i4:Rst-thm-minkideal}
Every ideal class of a number field $K$ of degree $n$ with $r_2$ complex places
contains an integral ideal of norm at most
$\frac{n!}{n^n}\bigl(\frac4\pi\bigr)^{r_2}\sqrt{|d_K|}$.
\end{theorem}

\begin{proof}[Proof of (R3)]
\emph{$|G|\le h_KM^2$.} Every ideal class $c$ contains an ideal prime to $M$:
take an integral $\mathfrak b$ in $c^{-1}$ and $\alpha\in\mathfrak b$ with
$v_\mathfrak p(\alpha)=v_\mathfrak p(\mathfrak b)$ for $\mathfrak p\mid M$ (Chinese remainder theorem); then
$(\alpha)\mathfrak b^{-1}$ is integral, prime to $M$, and in $c$. So the map $G\to\mathrm{Cl}_K$
is onto. Its kernel consists of classes of principal ideals $(\alpha)$,
$\alpha\in\mathcal O_K$ prime to $M$, and the class of $(\alpha)$ depends only on
$\alpha\bmod M$: if $\alpha'\equiv\alpha$, choose $\beta$ with $\alpha\beta\equiv1$; then
$(\alpha'\beta)(\alpha)=(\alpha\beta)(\alpha')$ with $\alpha'\beta\equiv\alpha\beta\equiv1\bmod M$. So the kernel has at
most $|(\mathcal O_K/M)^\times|\le N(M\mathcal O_K)=M^2$ elements, and $|G|\le h_KM^2$.

\emph{$h_K\le4d$.} By Theorem~\ref{i4:Rst-thm-minkideal} ($n=2$, $r_2=1$) every
class contains an integral ideal of norm $\le(2/\pi)\sqrt d\le X:=\sqrt d$.
The number of ideals of norm $m$ is at most the number $\tau(m)$ of divisors of
$m$: it is multiplicative in $m$, and for $m=p^e$ it is $e+1$, at most $1$, or
$1$ according as $p$ splits, is inert or ramifies. Hence
$h_K\le\sum_{m\le X}\tau(m)=\sum_{j\le X}\lfloor X/j\rfloor\le X(1+\log X)\le X^2=d\le4d$.
\end{proof}

\subsection{Outline of the formally verified steps}\label{i4:ssec-outline}

We outline, for each formally verified statement of this section, how its proof in
\cite{CharacterFamily} is organised.  The outline is not a proof.  Each step is classified as (L), (N) or
(R) as explained in the introduction to this annex; for an (L) step we give the published statement and
say what the formal statement needs beyond it, and for an (N) step we state it and give the idea of its
proof.  When we say that we do not describe a part of a formal proof, that part enters only through its
statement.  We treat in turn the hosts and the window data (the field of the character, the host, and
Propositions~\ref{i4:Rprop-kappaI} and~\ref{i4:Rprop-num}), the division side, the norm-row and functional
leaves, and the automorphic leaves.

\subsubsection{Hosts and window data}\label{i4:Sst-outline}

\medskip\noindent\label{i4:Sst-field}\emph{The field of the character.}

The construction behind Proposition~\ref{i4:Sprop-field} is the
following. Let $e\ge2$, and $q\equiv3\bmod8$, $d\equiv7\bmod16$ primes
with $(\frac{-d}q)=1$.

\begin{enumerate}
\item[(F1)] (R) \emph{The group and the frame.} Put $Q_S=2^{e+2}$ and
$A=(\Z/2^{e+2})^\times/\{\pm1\}$, a cyclic group of order $2^e$, with a
generator $a_0$ (for instance the class of $5$). So $n=\mathrm{ord}(\varphi)=2^e$ (the constant $c$
of the proposition is $1$), $A_0=A$ and $o=1$. In the formal proof, for a
general character $\psi$ of conductor $N_\psi$, $A$ is the image of
$(\Z/Q_S)^\times$ in $(\Z/2^{e+2})^\times/\{\pm1\}\times B$, where $B$ is the
group dual to the cyclic group generated by $\psi$, and $Q_S$ is the least
admissible divisor of $\mathrm{lcm}(2^{e+2},N_\psi)$; for $\psi=1$ this
reduces to the above.

\item[(F2)] (R) \emph{The Artin map.} Put $N=Q_Sqd$, $U=\{\pm1\}^2\times A$ and
$\mathrm{Art}(a)=((\frac aq),(\frac ad),[a])$ for $a\in(\Z/N)^\times$, where
$[a]$ is the class of $a$ modulo $2^{e+2}$ and $\pm1$. It is a surjective
homomorphism: by the Chinese remainder theorem one prescribes $a$ modulo
$2^{e+2}$, $q$ and $d$ independently. Its kernel is the set of $a$ with
$a\equiv\pm1\bmod2^{e+2}$ and $(\frac aq)=(\frac ad)=1$; this is
$\mathrm{KerChi}$. A character of $U$ trivial on $A$ is
$((s_1,s_2),a)\mapsto s_1^is_2^j$, and its composite with $\mathrm{Art}$ is
$\chi_q^i\chi_d^j$, with value $(-1)^j$ at $\iota_d$; this is
$\mathrm{TwistForm}$.

\item[(F3)] (L) \emph{The field.} Let $E\subset\Q(\zeta_N)$ be the fixed
field of $\{\sigma_a:a\in\ker\mathrm{Art}\}$, where $\sigma_a(\zeta_N)=\zeta_N^a$.
By the irreducibility of the cyclotomic polynomial,
$\mathrm{Gal}(\Q(\zeta_N)/\Q)\cong(\Z/N)^\times$ (Neukirch I \S10;
reference to be confirmed), and by Galois theory
$\mathrm{Gal}(E/\Q)\cong(\Z/N)^\times/\ker\mathrm{Art}\cong U$. Complex
conjugation is $\sigma_{-1}$, with image $((-1,-1),1)$ since $q\equiv d\equiv3\bmod4$;
it is a nontrivial central element, so $E$ is a CM field, of degree
$4\cdot2^e$, and $E\supset\Q(\sqrt{-q},\sqrt{-d})$ (Gauss sums;
$\sqrt{-q}\in\Q(\zeta_q)$, $\sqrt{-d}\in\Q(\zeta_d)$).

\item[(F4)] (L, with (N) parts) \emph{The host axioms.} The
host axioms are a list of about thirty arithmetic properties of $E$ in terms
of $\mathrm{Art}$: decomposition, inertia and Frobenius at every prime
(given by the residue of $p$ modulo the prime-to-$p$ part of $N$), the
inertia at $d$ ($\{1,\iota_d\}$), the CM type
$\mathrm{Gal}(E/K)$, the factorisation of the relative different
$\mathfrak D_{E/K}=\mathfrak b\bar{\mathfrak b}$, discriminant and
root-discriminant bounds, and norms of roots of unity. The first group is
the decomposition law in cyclotomic fields (Neukirch I \S10 and I (9.1),
(9.6)) transported to the subfield $E$. The
root-discriminant bounds are Proposition~\ref{i4:prop-rd}, and the norms of
roots of unity are Lemma~\ref{i4:lem-normmu}, both proved in this section.
We do not describe the different factorisation and the remaining clauses
in detail.

\item[(F5)] (R) \emph{The field rows.} $\mathrm{FieldRows}(H,b)$ consists of:
the $q$-sign of $\mathrm{Art}(a)$ is $(\frac aq)$; the square root
$\sqrt{-q}$ of (F3), on which $\mathrm{Gal}(E/\Q)$ acts by the $q$-sign;
exactness of $Q_S$ (for each prime $l\mid Q_S$ some $a\equiv1\bmod Q_S/l$
has nontrivial image, which holds since $Q_S=2^{e+2}$ is minimal);
$Q_S\le n^2$ ($2^{e+2}\le2^{2e}$ for $e\ge2$); $o^{64}\le2^e$; $q\ne3$; the
splitting $(\frac{-d}p)=1$ for odd $p\le b$ (a hypothesis); and the
condition $(\frac{-d}l)=1$ on odd prime divisors $l$ of $Q_S$
(Definition~\ref{i4:def-host}; vacuous here).

\item[(F6)] (R) \emph{Sizes.} $\nu=e$, $\mathrm{ord}=2^e$, $Q_S=2^{e+2}\le4\,\mathrm{ord}$,
$Q_S\mid2^{e+2}$. The formal proof takes $e_0$ at least $64$ (to cover
general $\psi$); for $\psi=1$ every $e\ge2$ works.

\item[(F7)] (R) \emph{The kept and exceptional clauses.} The kept set is
$\{1\}$. The trivial character of $U$ lies in $X^+$ and induces the trivial
Dirichlet character, so the kept set lies in $X^+$. There is no
exceptional character, and the clause then requires that the $q$-sign
character $((s_1,s_2),a)\mapsto s_1$ is not kept: its composite with
$\mathrm{Art}$ is $\chi_q\ne1$ (take $a$ with $(\frac aq)=-1$, $a\equiv1$
modulo $Q_Sd$).
\end{enumerate}

\noindent\emph{Classification.} (F1), (F2), (F5)--(F7) are (R); (F3) is
cyclotomic Galois theory (L); (F4) is the decomposition law in cyclotomic
fields (L) together with computations proved in this section and some
clauses whose proofs we do not describe. Nothing in this construction is new
mathematics; it is taken from the formal verification only because the host
axioms are not listed exactly in this paper.

\medskip\noindent\label{i4:Sst-host}\emph{The host.}

The formal proof of Proposition~\ref{i4:Sprop-host} produces the clauses of
``good'' (Definition~\ref{i4:def-host}) one by one. Let $N=N_E=Q_Sqd$; by the
hypotheses $\log N\le2x^{\varepsilon/5}(\log x)^2$ for large $x$
(Lemma~\ref{i4:lem-kits-shared}(d)).

\begin{enumerate}
\item[(H1)] (L) \emph{Norm primes.} Proposition~\ref{i4:thm-shortnorm}
(with $M=N$, $X$ with $\log X\ge8x^\varepsilon$, $h=1$) gives at least
$X^{9/10}$ primes $g\in[X,2X]$ of the form $g=\pi\bar\pi$, $\pi\in\mathcal O_K$,
$\pi\equiv1\bmod N$, for large $x$; here $X=(kb)^8$. This is proved in
this section (\S\ref{i4:ssec-proof-shortnorm}) modulo
Proposition~\ref{i4:prop-ray-leaves}.

\item[(H2)] (N, short) \emph{Two level primes with disjoint exceptional
sets.} Here $\mathrm{BB}(r,g)=\{g\}\cup\{\ell:\ \ell$ prime,
$g^\delta<\ell^{4v_\ell(g^\delta-1)}$ for some $1\le\delta\le r\}$, the set
of primes at which some $g^\delta-1$, $\delta\le r$, is unusually
divisible. It has at most $3r+1$ elements, and for fixed $g_1\in[X,2X]$ at
most $(3r+1)(1+2r^2(X^{3/4}+1))$ primes $g\in[X,2X]$ have
$\mathrm{BB}(r,g)\cap\mathrm{BB}(r,g_1)\ne\emptyset$ (for a given $\ell$ the
condition forces $g^\delta\equiv1$ modulo a large power of $\ell$, which has
few solutions in $[X,2X]$). Since $20r^3<X^{3/20}$, this is less than
$X^{9/10}$, so among the candidates of (H1) there are $g_0\ne g_1$ with
disjoint sets. The two counting facts are proved in \cite{CharacterFamily};
we use them only through their statements.

\item[(H3)] (L) \emph{Small level primes.} By Theorem~\ref{i4:thm-linnik},
applied three times to nested moduli, there are distinct primes
$p_{de},p'_{de}\equiv1\bmod15N$ and $q_0\equiv1\bmod N$ distinct from them,
all at most $N^{C_8}$. In \cite{CharacterFamily} the smallest-prime bound is
derived from a statement on primes in progressions with relative error
(Linnik's method); we do not describe its proof. Since $N^{C_8}<b\le(kb)^8\le g_i$ for large
$x$ (because $\log N\le2x^{\varepsilon/5}(\log x)^2$), these primes differ
from $g_0,g_1$.

\item[(H4)] (L) \emph{Splitting.} A prime $p\equiv1\bmod N$ splits
completely in $E\subset\Q(\zeta_N)$ (Frobenius is $\sigma_p=1$; Neukirch
I \S10), so the level primes split completely and $p_{de}$,
$p'_{de}$ have $[E:\Q]$ places above them; the element $\pi$ of (H1) gives
the norm clause of ``good''.

\item[(H5)] (L, adapted) \emph{The thick set.} There is an absolute $C$
such that for all $n_0$ and $L\ge2$ there is a set $T\subset[0,n_0/(4L))$
with $|T|\ge n_0e^{-C\sqrt{\log n_0\log L}}/L$ and \emph{no relation}: there
are no distinct $s_1,s_2,s_3\in T$ and integers $i\ne j$ with
$0<|i|,|j|\le L$ and $j(s_2-s_1)=i(s_3-s_1)$. This is Behrend's
construction (integer points on a sphere, written in a base large compared
with $L$ so that the relation holds digitwise; three distinct points of a
sphere are never collinear). Behrend, Proc.\ Nat.\ Acad.\ Sci.\ USA 32
(1946), 331--332. Extra in the formal version: all
relations with coefficients up to $L$ at once, and the explicit dependence
on $L$. Applied with $L=L_w$ and $n_0=n$ it gives the second torus set with
its lower bound and $4L_w|T_2|\le n_0$; here $L_w=O_\varepsilon(\log x)$, so
$|T_2|\ge x^{-\varepsilon}n$.

\item[(H6)] (N) \emph{The top set.} There is $C_b$ such that a set
$T_T\subset A$ with $\delta_1n/2\le|T_T|\le2\delta_1n$ exists whose shares
are within $x^{-1+\varepsilon}$ of their targets and which satisfies the
bad-prime bound: for primes $p<X'$ of inert class whose Frobenius has
order at least $x^{\varepsilon/4}$ in $A$, the weighted number of $t\in T_T$
at which the local condition fails is at most
$4C_bx^{-2}\delta_1n$ times the weight $\sum_p\log p/p$. Proof idea: choose
$T_T$ at random, each element independently with probability $\delta_1$;
compute the expectations of $|T_T|$, of the share deviations and of the
bad-prime sum (the latter by summing over Frobenius orbits, using that the
relevant primes exceed $x^{\varepsilon/(5C_8)}$, a consequence of the window
pair); a moment inequality shows that one choice satisfies all bounds at
once. For $o=1$ the share condition is vacuous (Convention (4)).

\item[(H7)] (L, with explicit uniformity) \emph{The thinness bound.} For
every primitive character $\chi$ attached to a character of $U$ and
$y\ge e^{x^{\varepsilon/3}}/3$, $|G_\chi(y)|\le y\exp(-c_6x^{\varepsilon/8})$.
This is the prime number theorem for $\chi$ with the classical zero-free
region (Z2) and explicit formula (Z4) of Theorem~\ref{i4:thm-classzeros},
once exceptional zeros are excluded: for real $\chi$ whose conductor is
divisible by $q$ or $d$ this is the box clause (iv) of $\mathrm{PairGood}$
(no zero in $(1-1/400,1)$), and for the finitely many real characters of
conductor dividing $8$ it is a fixed constant. The bound uses
$\log N\le2x^{\varepsilon/5}(\log x)^2$ against $\log y\ge x^{\varepsilon/3}/2$.
Extra in the formal version: the dependence of all constants is explicit.
The formal proof goes through $G_\chi(y)=-y\int_y^\infty\psi_0(u,\chi)u^{-2}du$
and the standard bound for $\psi_0(u,\chi)$; we describe it at the level of
statements only.

\item[(H8)] (R) \emph{The rows of the pair.} From $\mathrm{PairGood}$ and the
window $x^{1/J}\le Y\le x$: the $(q,d)$-bound with constant $C_5=4$, the box
bound, $q,d>4tr+1$ (as $\log q>Y^{\varepsilon/5}/2\ge x^{\varepsilon/(5C_0)}/2$
while $\log(4tr+1)=O(\log x)$), and $|L'/L(1,\chi_d)|\le x^\varepsilon$.

\item[(H9)] (R) \emph{Assembly.} The host is $H_0$ together with the choices
of (H2), (H3), (H5), (H6), which exists because $R$ is admitted
(\S\ref{i4:Sst-ranks}); validity, arithmeticity and the constants follow
from the above, with $C_8\ge C_0$ and the thick-set constant $C_{T_2}=2C$
from (H5).
\end{enumerate}

\noindent\emph{Classification.} New: (H6) (a probabilistic selection with
a moment computation over Frobenius orbits) and the short counting step
(H2). Published, re-derived or cited: (H1) (proved in this paper modulo
Proposition~\ref{i4:prop-ray-leaves}), (H3) (Linnik), (H4) (cyclotomic
reciprocity), (H5) (Behrend, adapted), (H7) (classical zero-free region and
explicit formula, with explicit constants). Routine: (H8), (H9). The
proposition is formally verified; the statements that
rest on \cite{CharacterFamily} and are not in the literature are the top
set (H6) and the uniform version of Behrend's construction (H5).


\medskip\noindent\label{i4:Rst-kappaI}\emph{Structure of the proof of Proposition~\ref{i4:Rprop-kappaI}.}

The proposition contains no new analysis: the depth is in hypothesis (i),
the first window theorem at the family, which is assumed. What the formal
proof does is to transport the conclusion of that theorem from its own
normalisation to the type-I excess $\kappa_I$ of the point. One point must be
made explicit, because the statement of
Definition~\ref{i4:Rdef-weight} treats $\kappa_I$ as a datum:

\emph{The excess is the least admissible excess.} In the formal
development every family carries, beyond its first threshold, the following
identity. For a prime $p$, an orbit length $o$, a set $T'$ of orbit positions
and weight vectors $a',b'$, call $\kappa\ge0$ \emph{admissible} if the local
lattice property of excess $\kappa$ holds, with the parameter
$\lceil t_I^{\rm f}(p,o,T',a',b')\rceil$ and the denominator bound
$D_I(p,o)$ of the point, in every unramified quadratic setting at $p$ and for
every dominant base weight attached to $(a',b')$. (The local lattice property
says that for every lattice of the local representation space there is a
family of edges into it whose sources carry vectors of the prescribed kind,
spanning modulo $p$ a submodule of rank at least $(1-\kappa)$ times the
dimension of the representation, with denominators bounded by $D_I$; it is
defined in \cite{CharacterFamily} and not restated here.) Then
\[
\kappa_I(p,o,T',a',b')=\inf\{\kappa:\kappa\text{ admissible}\}.
\]
The value $\kappa=1$ is admissible (the rank condition is empty), so
$0\le\kappa_I\le1$, and $\kappa_I\le\kappa$ for every admissible $\kappa$.

\emph{Steps.}
\begin{enumerate}
\item[(K1)] (R) Nonnegativity: $\kappa_I\ge0$ by the identity above.
\item[(K2)] (R) Orbit segments. If $t_0$ is good, $3G_p\le o_p$ and
$L_2<\mathrm{good}_U(\delta_1)$, then $T_{t_0}$ contains a top of type V, has at
least two segments, and every segment has gap at most $G_p$ and at most $N_1$
tops. This is an unpacking of the tiling condition that defines a good
orbit.
\item[(K3)] (R) Translation of the hypotheses. For a prime $p$ with
$p_0(k)\le p<X'$, not in the narrow band ($\rho$ in place of $\rho_B$), with
$\mathrm{HypI}(p,o_p,|T_{t_0}|)$ and (K2), the hypothesis of the first window
theorem holds at $(p,o_p,k-1,T_{t_0},G_p,N_1)$ with $Y=x^{A/2}$,
$Y_b=x^\varepsilon/20$, the scale quantities $a,b'$ and the windows
$W_A$, $W_B$. Each of its clauses is one of the fourteen clauses of
$\mathrm{HypI}$ or a direct consequence of it and of the definitions of $\rho$
and $\Lambda$.
\item[(K4)] (hypothesis (i)) The first window theorem then gives a function
$\kappa\ge0$ on $\mathrm{box}_{T_{t_0}}(W_A)\times\mathrm{box}_{T_{t_0}}(W_B)$ whose mean
is at most $C|T_{t_0}|N_1/x^{A/2}$, and such that the local lattice property
of excess $\kappa(a',b')$ holds for every parameter $t\ge t_I(a',b')$, with
$t_I$ the parameter of that theorem.
\item[(K5)] (R) Comparison of parameters: $t_I(a',b')\le t_I^{\rm f}(p,o_p,T_{t_0},a',b')$,
an inequality between two explicit expressions, using $C\ge0$,
$2C\le\kappa_1$, $C\le\kappa_2$, $p<a/b'$ and $|T_{t_0}|\le N^+(o_p)$; and the
integer windows of the theorem coincide with $W_A$ and $W_B$.
\item[(K6)] (R) By (K4) and (K5) each $\kappa(a',b')$ is admissible, so
$\kappa_I\le\kappa$ pointwise, and the mean of $\kappa_I$ is at most that of
$\kappa$. With $C_m=C$ this is the mean bound.
\end{enumerate}
\emph{What is new.} Nothing beyond hypothesis (i). Steps (K2), (K3), (K5)
are explicit verifications in
\cite{CharacterFamily}; the proposition is taken from the formal verification only because
they are not written out here. The prime $p$ ranges over primes with
$p_0(k)\le p<X'$ outside the narrow band, as in the statement.

\medskip\noindent\label{i4:Rst-num}\emph{Structure of the proof of Proposition~\ref{i4:Rprop-num}.}

All clauses are eventual inequalities between explicit functions of $x$,
derived from size information that the family carries or that
hypotheses (ii) and (iii) supply. No published theorem beyond the elementary
estimates below is used, and nothing is new; the proposition is taken from the formal
verification because the verification consists of some thirty separate
comparisons that are not written out here.

\emph{Inputs.}
\begin{enumerate}
\item[(U1)] (R) The size conditions of the family: at every point beyond its
first threshold, $k\le x^6$, $m\le x^{1/2+\varepsilon/100}$,
$|T|\le2x^{2+8\varepsilon+\varepsilon/100}$, the scale relation
$4S\le\Lambda b'/k$ and $N_m+1\le w''$. These are fields of the family; the two
size bounds of the proposition are among them.
\item[(U2)] (R) The orbit parameters: $r_G$ depends only on $\varepsilon$,
$G_p\le c_G(\varepsilon)\,x^{\varepsilon/8}\log x$ and $N_1\le c_N(\varepsilon)\log x$
for $x\ge e$, with explicit $c_G$, $c_N$; these follow from the definitions of
$G_p$ and $N_1$ through the segment-count bounds for good orbits (the latter
is not described here in detail).
\item[(U3)] (uses (ii)) Sizes of $\Lambda$ and $Q_p$: for every $\delta>0$,
eventually $x^{\varepsilon-\delta}\le\Lambda\le x^{\varepsilon+\delta}$,
$\log Q_p\le x^{\varepsilon/8+\delta}$ for $p<p_0(k)$ and
$\log Q_p\le x^{11\varepsilon/24+\delta}$ for $p<X'$. Here
$\Lambda=K_\Lambda N_1^2\log a$ and $Q_p=p^{G_p}$; the bounds follow from (U2),
$k\le x^6$ and $(\log x)^3\le x^\delta$ eventually; $K_\Lambda>0$ is hypothesis
(ii). It is used with $\delta=\varepsilon/48$.
\item[(U4)] (uses (iii)) The band: $\rho\le\rho_B\le x^{-2}$, and the sums of
$\log p/p$ over the primes $p_0(k)\le p<X'$ in the narrow band and in the band are
at most $x^{-1}$. The comparison $\rho\le\rho_B$ is where the logarithmic
window bound (iii) enters ($\rho_B$ has $K_\rho C_N\log x$ where $\rho$ has
$K_\rho N_1$). The sums are bounded by an elementary lemma: if
$0\le\rho\le1/3$, $X\ge1$, $P_0\ge1$ and $S$ is a set of primes $p\in[P_0,X)$
with $X(1-\rho)\le p^e\le X(1+\rho)$ for a fixed $e\ge1$, then
$\sum_{p\in S}\log p/p\le\log X\,(3\rho+1/P_0)$; summing over $e\le1+2\log X$
and using $\log X'=x^{\varepsilon/3}$, $\rho_B\le x^{-2}$ and $P_0=p_0(k)$ gives
$x^{-1}$ eventually.
\item[(U5)] (R) The uniform orbit constant: there is $C_{31}>0$, depending only
on the family, such that eventually, for every good representative of an
orbit of length $o\ge x^{\varepsilon/4}$ with $3G_p\le o$, the number of tops on
the orbit lies between $o\delta_1/(C_{31}\log x)$ and $40o\delta_1$, and the
gap and spacing conditions of a good orbit hold. This holds for every family
(no hypothesis).
\end{enumerate}

\emph{Steps.}
\begin{enumerate}
\item[(N1)] (R) The seven numerical conditions. With $A=26$, so
$x^{A/2}=x^{13}$, and $b'=x\exp(x^\varepsilon-x^{\varepsilon/3})$:
$K_{II}r_G^2kx^{-A}\le x^{-12}$ since $k\le x^{9/2+10\varepsilon+\varepsilon/100}$
and $r_G$ is constant; $C_mN_1/x^{13}\le x^{-12}$ by (U2);
$L_2<\mathrm{good}_U(\delta_1)$ from $514\log x\le x^{\varepsilon/8}$ and
$k\le x^6$; $3G_p\le x^{\varepsilon/4}$ by (U2); the two window widths are
$\ge1$ because $a=x\exp(x^\varepsilon)\ge x^{19}\ge kx^{13}$ eventually, and $N_*x^\varepsilon\Lambda b'/k\ge1$ (by (U3) and
$N_*\ge1$); and the total bound $(x^\varepsilon+1)m|T|x^{-12}u(2+u)\le x^{-9}/4$,
$u=x^{\varepsilon/3}$, by (U1) and $\varepsilon<1/52$. Each is of the form
$c\,x^{\alpha}\le x^{\beta}$ with $\alpha<\beta$, or $c\log x\le x^\beta$ with
$\beta>0$, for $x$ beyond a threshold depending only on $\varepsilon$, the
family constants and $C_m$.
\item[(N2)] (R) $\rho\le\rho_B$: (U4).
\item[(N3)] (R) $\mathrm{HypI}(p,o_p,|T_{t_0}|)$ at every candidate
$p\ge p_0(k)$ and good $t_0$: each of the fourteen clauses is obtained from
(U1), (U3), (U4), (U5) and (N1) by a comparison of the same kind, with
thresholds depending only on $\varepsilon$, the family constants, $C_m$ and
$C_{31}$ (we do not describe the clauses one by one).
\item[(N4)] (R) The size bounds on $m$ and $|T|$: (U1).
\end{enumerate}
The hypotheses used are therefore (ii) (in (U3), hence in (N1) and (N3)),
(iii) (in (U4), hence in (N2) and (N3)), and (i) only through the constant
$C_m$ of Proposition~\ref{i4:Rprop-kappaI}.


\subsubsection{The division side}\label{i4:Lst-ssec}

Here (R) steps are mostly thresholds, maxima of constants and transport of
structure.  The constructions of the division side share three devices, which we state
once. (i) \emph{Closedness bookkeeping.} A closed property comes with an
$\varepsilon_1$ and, for $\varepsilon\le\varepsilon_1$, a threshold
datum; conjunctions are closed with the minimum of the $\varepsilon_1$
and the entrywise maximum of the threshold data, and an implication
``$P\Rightarrow Q$'' is closed as soon as $Q$ is closed on the subclass
on which $P$ holds (R). (ii) \emph{Eventual statements.} An
``eventually'' clause is proved by taking the maximum of finitely many
thresholds, each produced by one of the steps (R). (iii) \emph{Absorption.}
Bounds of the shape $C\,P^{j}b$ with $P\le x^{M_t}$ are turned into
$x^{E}b$ for $x\ge\max(9,C)$ by Lemma~\ref{i4:Llem-bpol} (R).

\medskip\noindent\label{i4:Lst-gap}\emph{Structure of the proof of Proposition~\ref{i4:Lprop-gapleaves}.}

We first make the gap hypothesis and the two analytic properties of
(G3) explicit.

\begin{definition}[gap hypothesis; Wu decay]\label{i4:Lst-def-gap}
Let $D$ be a datum with ranks $k,m,f$ and $r=k+m-1$, and archimedean
weights $b_1(\tau)<\dots<b_m(\tau)$ at the infinite places $\tau$ of the
totally real field. $D$ satisfies the \emph{gap hypothesis} if $r\ge1801$
and $b_{j+1}(\tau)-b_j(\tau)\ge16\,r^{19}$ for every $\tau$ and every
$1\le j<m$. Let $X$ be a Hida setting over $D$: an adelic unitary group
$G$ with archimedean part $G_\infty=\prod_v U(\epsilon_v)$, for each word
$u$ an archimedean Bessel coefficient $W_u:G_\infty\to V_{\pi_u}$, a
subgroup $S_u\subseteq G_\infty$ (the Bessel subgroup) and the archimedean
test function $\varphi_\infty$. Put $\nu_u(x)=\inf_{t\in S_u}\prod_v\|(tx)_v\|_{\rm op}$
and $N_1=4n^2\sqrt{2q_{\max}}$, where $n$ is the size of the matrices and
$q_{\max}$ the largest $q$ of the signatures. $X$ has \emph{Wu decay} if
for every $u$: (a) there are a compact $\Omega\ni1$ and a bounded
measurable $\alpha$ supported in $\Omega$ with
$\int\alpha(z)W_u(xz)\,dz=W_u(x)$ for all $x$ (reproduction), and (b)
$\int_{S_u\backslash G_\infty}\|W_u(x)\|\,\nu_u(x)^{N_1}\,dx<\infty$
(weighted integrability). The archimedean coefficient is
\emph{integrable} if $\varphi_\infty\in L^1(G_\infty)$.
\end{definition}

\noindent The four leaves have independent proofs.

\medskip\noindent\textbf{(G1) The gap row.} The chain is:
\begin{enumerate}
\item[(G1.1)] (R) Given the closedness data $(\varepsilon_1,\kappa_0)$ of
$\mathrm{ResidEv}$ on $\mathcal C$, enlarge the last block of $\kappa_0$
to $\max(\kappa_0^{\rm last},29,8\sum|c|)$, the sum running over the outer
and middle constants. A family that is large for the new datum is large
for $\kappa_0$, so $\mathrm{ResidEv}$ holds, and every last constant is
at least $29$ and at least $8C_\Gamma$.
\item[(G1.2)] (N) \emph{Standing hypotheses at the family constants.}
If $8C_\Gamma\le C_{k,\rm low}$, $\mathrm{TuSize}$ and $\mathrm{ResidEv}$
hold, then eventually, at every point, the standing hypotheses
\eqref{i3:eq-standing} hold with constants $C'=C_{k,\rm low}$ and
\[
C_l=\min(C_{fL},C_w,C_\Delta,C_{\rm count},C_{h0r},C_{\rm gap}).
\]
The standing hypotheses are a list of about twenty-five explicit
inequalities among $k,m,f,n,r,L,b$, the auxiliary primes $g_0,g_1$, the
discriminant of the host and the archimedean weights; each is checked
from the size fields of the family (the parameter record and
$\mathrm{SizesAt}$), from $\mathrm{TuSize}$ for the window
$b^{m/2}\le T_u\le b^{C_{T_u}m}$, and from $\mathrm{ResidEv}$ for
$\mathrm{rd}_E\ge2$ and the splitting conditions, once
$x^{\varepsilon/100}\ge\max(2C_l,2,500/\varepsilon)$. The idea is that the
size record of a family was designed so that every standing inequality
is one of its fields or a consequence of comparing powers of $x$, $L=x^\varepsilon$
and $b=e^L$.
\item[(G1.3)] (R) \emph{Gap from the standing hypotheses.} From
$C'm^5fL\le k$ with $C'\ge29$, $m\ge2$, $f\ge m^4\ge16$, $L>1$ one gets
$k\ge29\cdot32\cdot16>1801$, hence $r\ge k\ge1801$; from
$C_lr^{19}f^5L\le b_{j+1}-b_j$ with $C_l\ge16$, $f\ge16$, $L>1$ one gets
$b_{j+1}-b_j\ge16r^{19}$.
\end{enumerate}

\medskip\noindent\textbf{(G2) The residual eventuality.} $\mathrm{ResidEv}(\mathcal F)$
says: eventually, at every point, $\mathrm{rd}_E\ge2$, every level prime
$T_u$ and the two auxiliary primes $g_0,g_1$ split completely in $E$, and
$\mathrm{disc}(E)=\mathrm{disc}(F)^2$. It holds for every family with
threshold $0$, which gives closedness of every class trivially (R).
\begin{enumerate}
\item[(G2.1)] (R) $\mathrm{rd}_E\ge2$: the host satisfies
$\mathrm{rd}_E^2\ge4\,\mathrm{ord}$ with $\mathrm{ord}\ge1$
(Proposition~\ref{i4:prop-rd}).
\item[(G2.2)] (L) \emph{Splitting of $T_u$.}

\begin{theorem}[decomposition of primes in cyclotomic fields; Neukirch, \emph{Algebraic Number Theory}, I (10.3)]\label{i4:Lst-thm-cyc}
Let $N\ge1$ and $p$ a prime with $p\nmid N$. Then $p$ is unramified in
$\Q(\zeta_N)$ and its residue degree is the order of $p$ modulo $N$. In
particular, if $p\equiv1\pmod N$ then $p$ splits completely in
$\Q(\zeta_N)$, hence in every subfield.
\end{theorem}
\noindent The level prime $T_u$ is the norm of an element $\pi\equiv1$
modulo the conductor $N_E$ of the abelian field $E$ (a field of the
datum), so $T_u\equiv1\pmod{N_E}$ and $T_u\nmid N_E$; with
$E\subseteq\Q(\zeta_{N_E})$ the theorem gives complete splitting. The
passage from $\Q(\zeta_N)$ to a subfield is multiplicativity of $e$ and
$f$ in towers and $\sum e_if_i=[E:\Q]$ (Neukirch I (8.2)).
\item[(G2.3)] (R) The splitting of $g_0,g_1$ and
$\mathrm{disc}(E)=\mathrm{disc}(F)^2$ are fields of the datum
(ramification of $E/F$ only at the places recorded in the host).
\end{enumerate}

\medskip\noindent\textbf{(G3) Wu decay.} The formal proof derives Wu decay
from a Bessel multiplicity-one statement under a gap and five further
inputs. The chain is:
\begin{enumerate}
\item[(G3.1)] (R) \emph{From the gap hypothesis to the local gap
conditions.} With $t=r+6$ (odd, since $r$ is odd) one has $t\ge r+5$,
$r+t\ge1801$, and
$b_{j+1}(\tau)-b_j(\tau)>2(r+(2r_S+2)t)^3+t+3$, a numerical consequence
of $b_{j+1}-b_j\ge16r^{19}$ and $r=2r_S+k+1$. This is the form of the gap
used by the local statements below; for the multiplicity-one statement it
gives a parameter of gap at least $120n^3$ in the sense of
Definition~\ref{i3:def-gap}, as in Lemma~\ref{i3:lem-gap}.
\item[(G3.2)] (L, proved in Section~\ref{i3:sec}) \emph{Multiplicity one.} Any two
smooth Bessel maps of a discrete series with a parameter of gap at least
$120n^3$ are proportional: Proposition~\ref{i3:prop-bessel}. The
published multiplicity-one theorems for Bessel models of unitary groups
(Jiang--Sun--Zhu; Sun--Zhu; reference to be confirmed) hold for all
irreducible Casselman--Wallach representations and contain this case;
the formal development proves the gap case directly (through boundedness
of Bessel maps on orbits and the Iwasawa, Levi and Cartan decompositions),
and Section~\ref{i3:sec} gives a written proof of it.
\item[(G3.3)] (N) \emph{Non-vanishing of the Bessel integral.} For a
discrete series $\sigma$ with parameter $\lambda$ of gap at least
$120n^3$, a lowest-$K$-type unit vector $w$, and $\pi$ of
Harish-Chandra type $u$ interlacing $\lambda$ with half-integral
differences, there are $x,x_0$ in the group with
\[
\int_{U_z\times N}\langle\sigma(x_0)w,\sigma(hn)\sigma(x)w\rangle\,
\overline{\psi_{z,c}(n)}\,\pi(h)^{-1}\,dh\,dn\neq0 .
\]
This is the archimedean existence half of the Gan--Gross--Prasad
branching for discrete series (for $U(p,q)$ it is a theorem of H.~He;
reference to be confirmed), in the form of a non-vanishing local period
integral. The formal statement asks for the absolutely convergent
integral against the lowest $K$-type, which is more than the existence of
a nonzero Bessel functional; we did not compare the two, so we count it
as (N). We do not describe its formal proof.
\item[(G3.4)] (L) \emph{Decay of matrix coefficients.} Under the same
gap, $\int_{U}|\langle w,\sigma(g)w\rangle|\,\|g\|_{\rm op}^{N_1'}\,dg<\infty$
with $N_1'=4n^2\sqrt{2(r+1)}$. This is the integrability of discrete
series coefficients far from the walls (Trombi--Varadarajan,
Hecht--Schmid; reference to be confirmed), with an explicit weight. Section~\ref{i3:sec}
proves the needed pointwise bound, Lemma~\ref{i3:lem-decay}: for gap
larger than $\mathrm{thr}^*(\epsilon)+m$ the coefficient is
$O(\|g\|^{-m})$, and integrability against $\|g\|^{N_1'}$ follows for $m$
larger than $N_1'$ plus the exponent of growth of the Haar measure. The
explicit exponent and the explicit gap are the extra part.
\item[(G3.5)] (L) \emph{Reproduction.} For a discrete series and a
lowest-$K$-type vector $w$, and every neighbourhood $U$ of $1$, there is
a continuous $\alpha$ with compact support in $U$ such that
$\int\alpha(g)\sigma(g)w\,dg=w$. This is Harish-Chandra's lemma that a
$K$-finite, $Z(\mathfrak g)$-finite vector is fixed by some
$\sigma(\alpha)$ with $\alpha\in C_c^\infty$ supported near $1$
(Harish-Chandra; reference to be confirmed).
\item[(G3.6)] (R) Two measure-theoretic facts: the unipotent group of the
Bessel datum is unimodular, and Haar measure on a finite product of real
unitary groups disintegrates along a closed subgroup $N\rtimes H$ with $H$
compact (Weil's quotient formula; reference to be confirmed).
\item[(G3.7)] (N) \emph{Local scalar reduction and assembly.} The
archimedean Bessel coefficient $W_u$ of the Hida setting is constructed
from the local Bessel functionals; by (G3.2) it is, up to a scalar, the
integral in (G3.3), which is non-zero. Reproduction (a) follows from
(G3.5) by moving $\alpha$ through the integral; weighted integrability
(b) follows from (G3.4) after the change of variables of (G3.6) and the
comparison $\nu_u(x)\le\|x\|$. The formal proof of this assembly is a
reduction to a scalar identity at each place, which we do not describe in
detail.
\end{enumerate}
\noindent Integrability of $\varphi_\infty$ is (G3.4) with weight
exponent $0$ (L).

\medskip\noindent\textbf{(G4) Finiteness from decay.} $\mathrm{FinX}$ of
a full datum $Y$ over $X$ is the conjunction of ten properties: for every
word the ``sharp'' map from rational double cosets to classes is
bijective; the Poincar\'e series attached to each word lies in the
isotypic space $HH$ cut out by the convolution operator of $X$, and the
Poincar\'e map is linear; $HH$ is finite dimensional; the fibres of the
row maps are finite; the pairs of terms are independent of the choice of
representatives; the scalars $t.G$ of the row terms are real; the Haar
measures scale as prescribed; and the level volumes are positive. The
chain:
\begin{enumerate}
\item[(G4.1)] (L) \emph{Finite dimensionality of $HH$.} This is
$\mathrm{RowsFinDim}$, an instance of Harish-Chandra's finiteness theorem.

\begin{theorem}[Harish-Chandra, \emph{Automorphic forms on semisimple Lie groups}, LNM 62 (1968), Thm.~1; reference to be confirmed]\label{i4:Lst-thm-hc}
Let $G$ be a reductive group over a number field, $K_f$ an open compact
subgroup of $G(\mathbb A_f)$, $\tau$ a finite set of $K_\infty$-types and
$J$ an ideal of finite codimension in $Z(\mathfrak g)$. The space of
automorphic forms on $G(F)\backslash G(\mathbb A)$ that are right
$K_f$-invariant, of $K_\infty$-types in $\tau$, and annihilated by $J$, is
finite dimensional.
\end{theorem}
\noindent The formal statement is about the image of a convolution
operator $T=R(\varphi_\infty\otimes e)$ on $L^2(G(F)\backslash G(\mathbb A))$,
with $e$ locally constant of compact support and $\varphi_\infty$ an
integrable discrete-series coefficient. Its formal proof passes through
reduction theory (fundamental domains, cusps, the parabolic
decomposition), shows that functions in the image are cuspidal at
infinity and of moderate growth, and concludes by a finite-dimensionality
argument from linear forms; this is the classical route. The passage
from the $L^2$ image of $T$ to automorphic forms of fixed $K$-type and
infinitesimal character is the only part not literally in the theorem
(R). In this paper the statement is
Proposition~\ref{i4:A-prop-inputs-leaves}(M5).
\item[(G4.2)] (N) \emph{Poincar\'e series in $HH$.} For a level
subgroup normalised as in the setting, and Wu decay, the Poincar\'e
series $P_u=\sum_{\gamma\in S_u(F)\backslash G(F)}\phi_u(\gamma g)$ built
from $W_u$ and the finite part converges absolutely, is square
integrable (the fold condition: a fundamental domain for $S_u(F)$, and the
unfolded integral of $|P_u|^2$ is dominated by the weighted integral of
Definition~\ref{i4:Lst-def-gap}(b)), and lies in $HH$ (the reproduction (a)
shows $P_u=R(\alpha)P_u$). The idea is the usual unfolding; the
explicit weight $\nu_u^{N_1}$ is what makes the counting of lattice
points in the unfolded domain converge.
\item[(G4.3)] (N) \emph{Neatness and bijectivity of the sharp map.} The
level subgroup at each word is neat; the sharp map is then bijective.
We do not describe the proof below the statement.
\item[(G4.4)] (R) Finiteness of fibres, independence of pairs (from the
injectivity of the embedding of the kernel), reality of $t.G$ (the flag
Gram scalar is fixed by complex conjugation), Haar scaling and positivity
of volumes: short verifications.
\end{enumerate}

\noindent\emph{What rests on \cite{CharacterFamily}.} (G1.2), (G3.3),
(G3.7), (G4.2), (G4.3); the extra uniformity in (G3.4); and the formal
proofs of the cited theorems in (G2.2), (G3.5), (G4.1), which are
re-derived in the formal development inside longer arguments rather than
isolated as separate propositions.

\medskip\noindent\label{i4:Lst-sel}\emph{Structure of the proof of Proposition~\ref{i4:Lprop-sel}(a).}

The four clauses are of different nature.

\begin{enumerate}
\item[(S1)] \emph{adm2: the period valuation bound.} The clause says:
at every word $u$, every twisted place $w$ and every $\phi$ in the
division lattice $\nu_u$, and every $F$ in $\tilde M_{\mu(w)}\cap V_0(\mu(w))$,
the period column $\mathrm{per}_u(\mu(w),j(\phi),F)$ is the image of some
$y\in LL$ with $y=0$ or $v_w(y)\ge\lfloor\mathrm{ur}_w(s_g(u,w))\rfloor$.
Its proof in \cite{CharacterFamily} has three steps.
  \begin{enumerate}
  \item[(S1.1)] (R) Eight inputs are combined by closedness bookkeeping:
  the rows ``reference chain'', ``twisted point'', ``period at zero'',
  ``twisted period'', ``level-prime valuations'' and ``arithmetic'' of the
  hypothesis structure (rows 5--10 and 6 of Table~\ref{i4:A-tab-rows}), the
  service row (row 21, turned into a family statement using
  $\mathrm{DGamSize}$), and the coprimality of the level primes to the
  discriminant of the stabiliser.
  \item[(S1.2)] (N) \emph{Coprimality of level primes and stabiliser.}
  Eventually the level primes do not divide the discriminant of the
  stabiliser order; proved from the size of $T_u$ against the order of the
  finite stabiliser. We do not describe the proof below the statement.
  \item[(S1.3)] (N) \emph{Valuation bound from the inclusion.} From the
  eight inputs: the period at a twisted place is a sum over the
  reference chain of twisted point values times admissible periods; the
  integrality of the points and of the periods at $0$, the admissibility
  of the twisted periods and the sharp column bound at $T_u$ give the
  lower bound on the valuation, rounded down. The same argument gives a
  second clause (adm12), which is not used.
  \end{enumerate}
\item[(S2)] \emph{cut.} (R, from a hypothesis) If every lattice
$\tilde M_\mu$ is stable under $\mathcal O_{LL}$, which is the lattice
stability field (ii) of the hypothesis structure, then the cut condition
holds for every used tuple: a one-line verification.
\item[(S3)] \emph{c2: the second complement.} (R, from a hypothesis) The
second-complement row (row 4) gives the family second-complement inputs;
the complement bound follows from them with a constant and a threshold
(one inequality chain).
\item[(S4)] \emph{c3.} (N) The third complement bound holds for every
family, with no hypothesis; it is an explicit estimate of the
complement against the size polynomial. We do not describe the proof beyond the
statement.
\end{enumerate}

\noindent\emph{Remark on the statement.} In the formal proof, (a) depends
on rows 4--10 and 21 and on the lattice stability field of the hypothesis
structure. This is consistent with the statement, which is made for a
class $\mathcal A$ carrying a hypothesis structure
(Proposition~\ref{i4:prop-links-leaves}).

\noindent\emph{What rests on \cite{CharacterFamily}.} (S1.2), (S1.3),
(S4). Everything else is bookkeeping around hypotheses of the structure.
No published theorem is used.

\medskip\noindent\label{i4:Lst-D}\emph{Structure of the proof of Proposition~\ref{i4:Lprop-D}.}

\noindent\textbf{(a) Maximal slope.} The chain:
\begin{enumerate}
\item[(a1)] (R) A twist selection gives, eventually, column values at
every word: for each class $c$ at level $0$ and each infinite place
$\tau$, a hermitian $\mathcal O_{LL}$-lattice $Q_{c,\tau}$, generated by
the images of the class lattice under a semilinear map
$b_{c,\tau}$ compatible with the majorant, with $Q_{c,\tau}$ good, and
an injective linear map from the level-zero space into polynomial
functions on the classes (the ``column value'').
\item[(a2)] (N) \emph{Exact slope of the column lattices.} The slope of
the column lattice equals the slope of the auxiliary lattice $Q$; a
direct computation with the polynomial model. We do not describe this step in detail.
\item[(a3)] (N) \emph{Lower slope bound for the standard lattice.} For a
decreasing $\nu\in\N^N$, let $W_\nu$ be the $\Z$-lattice of the Weyl
module of highest weight $\nu$ in $\Z[x_{ij}]_{1\le i,j\le N}$ with the
Fock (Fischer) inner product, normalised by the norm of the highest
weight vector $h_\nu$. Then
\[
\mu(W_\nu)\ge-\Bigl(\sum_l\nu_l\Bigr)\log N-\log c_\nu,\qquad
1\le c_\nu=\frac{\langle h_\nu,h_\nu\rangle}{\prod_l\nu_l!}\le(\nu_1+N)^{N(N-1)/2}.
\]
The upper bound on $c_\nu$ is a hook-length type product; the formal
proof derives it from a Capelli determinant identity. The norm of the
highest weight vector as a product over hooks is classical (it goes back
to the Capelli identity; reference to be confirmed); the explicit bound and
the slope bound are what is used, and we count them as (N).
\item[(a4)] (N) \emph{Upper slope bound with constant $4$.} An upper bound
for $\mu_{\max}$ of the column lattice by the open part of the slopes,
with constant $C_5=4$. We do not describe the proof below the statement.
\item[(a5)] (R) (a2)--(a4) give the Frobenius expansion congruence with a
constant $C$: at every word, a two-sided bound for $\mu(\Lambda^0_u)$ and
$\mu_{\max}(\Lambda^0_u)$ by $-\sum_t(a_t/k)\log T_u$ plus error terms in
$n_{TT}a\log k$, $fk(b'+S_c)\log k$, $fS_c\log T_u$, $fk^2L$.
\item[(a6)] (R) With the sizes, $\mu_{\max}(\Lambda^0_u)\le|C|(12+2|C_{T_u}|)P^4b$,
and absorption gives $x^{E_b}b$ with $E_b=4M_t+1$.
\end{enumerate}

\noindent\textbf{(b) Slope lemma.} The chain:
\begin{enumerate}
\item[(b1)] (N) \emph{Level lines.} From the level conditions, the level
prime is good in the sense of the level-$T$ family (``Lemma 5.1 from
Theorem 2'' of the formal development); it is explicit, and we do not describe it in detail.
\item[(b2)] (R) Commensurability of $\nu_u$ with the reference lattice
follows from the level conditions.
\item[(b3)] (L) \emph{Mertens.} The sums over tops
$\sum_\tau\log p_\tau/p_\tau$ are bounded through Mertens' first theorem
in the explicit form $\sum_{p<y}\log p/p\le2\log y$ ($y\ge2$; an explicit
form of Mertens' theorem, e.g.\ Rosser--Schoenfeld, Illinois J.\ Math.\ 6 (1962), 64--94).
\item[(b4)] (N) \emph{Places bound.} The slope of $\nu_u$ is compared
with that of $\Lambda^0_u$ place by place: at each place $w$ in the
support, the local length of $\nu_u$ against the reference lattice is
bounded by the gauge and frame terms and the excess $E_e(u)$, using
$(\mathrm{LG})$ for the gauge, $(\mathrm{LMu})$ for the slope count and
$\mathrm{MuRank}$ for the ranks; at places above $T_u$ the local length
is $0$; at places outside the support the lattices agree. Summing over
places gives the two-sided bound. This is the main new estimate of (b).
\end{enumerate}

\noindent\textbf{(c) Rows.} (N, conditional on rows 14, 16, 17, 25.)
At a selected point: primitivity (row 14) gives the separated types;
the definite rows (row 25) give the realisation data; the divisor rows
(row 16) give the division rows; the decomposition-point rows (row 17)
give the decomposition points. \emph{Column facts} of the indefinite
realisation then say that the row lattice of the realisation is a full
lattice and that the column maps are bijective. $\mathrm{RRFull}$ is
transport of fullness along the row identification (R).
$\mathrm{BuRational}$ is \emph{period descent}: the Bessel coefficient
of a division vector is, by the column identity, a period of the
realisation at the base point, which lies in $LL$ because the level
operator fixes the image of the forms and the column scalar is $1$
(N, short).

\noindent\textbf{(d) Sharp functional bound.} The chain:
\begin{enumerate}
\item[(d1)] (R) Rows 12, 14, 16, 25 give, at selected points, the
covolume, onto, low-weight and isometry conditions of the realisation
and the arithmetic clause (L1) with $d_g$.
\item[(d2)] (N) \emph{Pointwise functional bound.} Under these
conditions,
$\mu_{\max}(RR)+\mathrm{mean}_e\log\kappa_R(e)\le x^{C_H}b+14\lambda_0$,
$\lambda_0=\log|G_p|$. The row lattice is compared with the image of the
base lattice under the realisation; the covolume and isometry conditions
control the determinant, and the low-weight condition bounds the
archimedean constants. We do not describe this step in detail.
\item[(d3)] (R) $\lambda_0\le4b$ and absorption give
$x^{\max(C_H,1)+1}b$; the threshold bookkeeping yields $C_{F0}$.
\end{enumerate}

\noindent\textbf{(e) Pointwise size bounds.} All three bounds are
explicit term-by-term estimates from $\mathrm{SizesAt}$, the fills and the
frame definition (R), with one published input: the explicit Mertens
bound of (b3), used for $\mathrm{Ga}$ and $\mathrm{Fo}$ (L). These estimates admit a complete written version
(the gauge term through the gains and the hook products, the frame term
through the lift weights, and $h_6$ through the logarithms of the hook
lengths); they are routine and are not repeated here.

\noindent\textbf{(f) Slope count.} (R) $(\mathrm{LMu})$ is assembled from
five pointwise facts: vanishing above the levels
(Lemma~\ref{i4:lem-zeroabove}), fullness of $\nu_u$ (hypothesis),
$\mathrm{MuRank}$ (hypothesis), the $T_u$-bound
$C_{T_u}mL\le2|C_{T_u}|X''$ from $\mathrm{TuSize}$ and $m\le2n$,
$L\le x$, and the $D_\Gamma$-bound of Proposition~\ref{i4:prop-mu}(b). The
combination of these into the slope count is a short computation in the
formal development.

\noindent\emph{What rests on \cite{CharacterFamily}.} (a2)--(a4), (b1),
(b4), (c) period descent and column facts, (d2). The rest is bookkeeping
and one explicit form of Mertens' theorem.

\medskip\noindent\label{i4:Lst-pkg}\emph{Structure of the proof of Proposition~\ref{i4:Lprop-pkg}.}

The statement is a description of the input of the package lemma; its
proof in \cite{CharacterFamily} is a construction from named data. The
steps are:
\begin{enumerate}
\item[(P1)] (R) \emph{Transport.} The tied division datum $I$ compatible
with $P_k$ is transported to the datum of the selection; the column
bridge built from $I$ is tied because the tuple map, the evaluation
maps and the word map of $I$ agree with those of the selection after
transport.
\item[(P2)] (R, proved here) \emph{The $\eta$-sum.} The bound of
Lemma~\ref{i4:Llem-etatw} at the twisted places, from the division
inequality $\mathrm{DIVinN}$, the twist record and $D_0$.
\item[(P3)] (N, short) \emph{Lift exponents at the twisted places.} At
each twisted place $y$ and each index $j$, the lift exponent is
realised by a lattice weight: the lifted lattices exist (block
hypothesis of the family) and the exponent of the weight equals the
spectral exponent. One lemma on lattices over a local field.
\item[(P4)] (from hypotheses) \emph{Row data at every word.} The
definite rows (row 25) supply, at every word, chain ties, a lattice
$X_0$, the numerical record and the first columns.
\item[(P5)] (R) \emph{Assembly.} The division inputs without norm are
built from (P2)--(P4) and the arithmetic clause (L1) with $d_g$; the
norm clause is added from (P2); the result is converted to the
division-input record of the package. This is a record with some twenty
fields, each filled by one of the inputs.
\end{enumerate}

\noindent\emph{What rests on \cite{CharacterFamily}.} (P3), and the
correctness of the assembly (P5). Nothing in this construction is a
published theorem; nothing is a new estimate. As already said after the
statement, the statement itself is not made exact in this paper.

\medskip\noindent\label{i4:Lst-forms}\emph{Structure of the proof of Proposition~\ref{i4:Lprop-forms}.}

All four items start from the concrete forms at level zero at a word
$u$ in the Hida setting: a family of hermitian lattices at level $0$ with
an action of the level operators $U_v$ (``up'') and $D_v$ (``down'') at
the places $v\in VV_w$, values and Frobenius data, and a character of a
finite group with finite image (``character finite''). The items:
\begin{enumerate}
\item[(1)] \emph{Slope line.}
  \begin{enumerate}
  \item[(1a)] (N) \emph{Diamond isometry at infinity.} Character
  finiteness and admissible stabilisers imply that the diamond operators
  act isometrically for the archimedean metrics.
  \item[(1b)] (N) \emph{Orthogonality of disjoint supports.} Forms with
  disjoint supports on the classes are orthogonal for the metric.
  \item[(1c)] (N) \emph{Splitting at $T$.} The level-zero lattice splits,
  at $T$, as the sum of the ramified part and its complement, compatibly
  with the metric.
  \item[(1d)] (R) From (1a)--(1c), $\Lambda^{\rm ram}_u$ is the image of
  an orthogonal projector of determinant ratio $((T-1)/(T-2))^{-f}$ on
  the relevant factor, which gives
  $\mu(\Lambda^{\rm ram}_u)=\mu(\Lambda^0_u)-f\log(T-1)/(T-2)$.
  \end{enumerate}
\item[(2)] \emph{Second-branch inclusion.} (N) The inclusion
$T^{n_b}S_{0,w}\subseteq B_{0,w,v}$ comes from the weight-lattice
inclusions of the level-zero forms; its length is the seen-place lattice
length formula. Both are statements about lattices of polynomial
functions at one place; we do not describe it in detail.
\item[(3)] \emph{Integral Hecke condition.} (N) Integrality of $U_v$ and
$D_v$ at level $0$ is reduced to three facts: rationality of the Hecke
operators at level zero, and integrality of the ``up'' and ``down''
increments. Each increment is a sum over cosets of values of the forms,
and integrality is checked coset by coset.
\item[(4)] \emph{Ramified Hecke condition.} (N) The relations
$U_vD_v=T^{k-1}$ on the ramified subspace come from matchings of the
diagonal double cosets (``up'', ``down'' and ``exchange'' matchings at
the same place) and a formula for $D_vU_v$ at level zero. This is the
analogue, for the level operators of the family, of the classical
relation between $U_p$ and its adjoint on the $p$-new part (for $GL_2$,
the Atkin--Lehner--Li theory); we do not know a published statement for
the unitary level operators used here, so it counts as new.
\end{enumerate}

\noindent\emph{What rests on \cite{CharacterFamily}.} All four items:
(1a)--(1c), (2), (3), (4). No published theorem is used in their
formal proofs, as far as we can tell from their top levels.

\medskip\noindent\label{i4:Lst-fam}\emph{Structure of the proof of Proposition~\ref{i4:Lprop-fam}.}

\noindent\textbf{(a) Setting.} (R, with one (N) step.) The formal proof
uses $\mathrm{TuSize}$ only (so the statement holds without the minimal
global record):
\begin{enumerate}
\item[(a1)] (R) The size record gives, eventually, the number and weight
bounds of the level prime: $m\ge4$ (from $m\ge x^{2/5}$ eventually) and
the bounds $2fN_m+1<T_u$, $4a_{\max}(u)<T_u$ from $T_u\ge b^{m/2}$.
\item[(a2)] (N, short) The alcove condition and $b_{\rm top}(\tau)\ge0$
at every top, from the weight bounds of the size record.
\end{enumerate}

\noindent\textbf{(b) Items.} The proof is a list; each item is a short
consequence of earlier ones or of a hypothesis.
\begin{enumerate}
\item[(b1)] (R) Setting from (a); level lines $\mathrm{LvlF}$ from the
setting and the forms leaves of Proposition~\ref{i4:Lprop-forms};
commensurability and fullness of $\nu_u$ (Lemma~\ref{i4:Llem-lvl}); rank
condition $\mathrm{MuRank}$ from the chosen proviso.
\item[(b2)] (N) \emph{Fifth family part and count.} From the mass ratio,
``base volume equals mass'' and the two-sided $\Sigma$-degree bound
(hypotheses of the minimal link statement), the level-prime division
bounds and then the count. The mass ratio itself is an instance of the
mass formula for hermitian lattices (Proposition~\ref{i4:A-prop-inputs}(a));
here it is a hypothesis.
\item[(b3)] (R) Rank condition of the ramified-rank family from
$(\mathrm{LS})$, the proviso, $(\mathrm{LIrr})$ and $(\mathrm{LDiv})$.
\item[(b4)] (N) \emph{Slopes.} $\mathrm{SlopesM}$ from the two slope
conclusions, the count, the logarithmic-derivative bound and the column
values (the argument of Proposition~\ref{i4:Lprop-D}(a)); $\mathrm{SlopesN}$
from the places bound (Proposition~\ref{i4:Lprop-D}(b), steps (b1)--(b4)).
\item[(b5)] (R) Lattice condition from $(\mathrm{LDiv})$, fullness and
the ranks; real frame and $\mathrm{FrameF}(C_H)$ from the lattice
condition, the proviso and the matrix constant; nodes, character bound
and windows from the size record and the thick density of the chosen
data.
\end{enumerate}

\noindent\textbf{(c) Core link.} (N) Given the items, the core link
constructs the chain-tied data at every $x\ge x_0$ in the window. The
construction: by the losses $\mathrm{LossFS}(\cdot,0)$ choose a point
$i_0$ and the weight vectors $a,b$ with admissible losses; by the covering
property of the family choose, for every residue $u_0$, a point carrying
the same top data, which gives a model (a finite set of points indexed by
the residues); the frame bound gives the frame constants $(\eta,e)$ at
each of these points; the character bound gives the logarithmic
characters; the windows give the window constants. One then checks the
five properties of chain-tied data with
$K=\max(C_d,0)+C_c+C_s+C_1+C_l+6$, $c_{HH}=C_H$, $e_K$ from the windows,
$C_L=C_5$ and $x_0$ the sum of all thresholds: each property is an
inequality between the quantities just chosen and one of the items
(count, slopes, lattice, nodes), and the nodes bound the cut ratios
$a_t/w\le\lceil2a\rceil$. The idea is entirely combinatorial; the
analytic content is in the items.

\noindent\emph{What rests on \cite{CharacterFamily}.} (a2), (b2), (b4),
(c). The mass ratio is used as a hypothesis here; its proof is part of
Proposition~\ref{i4:A-prop-inputs}.


\subsubsection{The norm-row and functional leaves}\label{i4:Gst-ssec}

The propositions outlined here are formally verified, \cite{CharacterFamily}.

\medskip\noindent\label{i4:Gst-split}\emph{Structure of the proof of Proposition~\ref{i4:A-prop-split}.}

Proposition~\ref{i4:A-prop-split} is formally verified. Fix a setup, a twisted
place $k$ with prime $l=l_k$ satisfying (H1)--(H5), and $\pi\in Z^{\rm good}_k$.
The quantity $v_{cc}(k,\pi)$ is the valuation of the congruence coefficient
of the Hecke eigensystem $\pi$ in the Hecke module at $k$; the family of
Galois representations $\rho=(\rho_x)_{x\in X_k}$ with values in
$\mathrm{GL}_r(\mathcal O_k)$ and the auxiliary family $r$ with values in
$\mathrm{GL}_{n_V}(\mathcal O_k)$ are those of the hub datum. The argument
is Ribet's method: congruences between $\pi$ and the other points of the
family produce Galois cohomology classes, and the size of the congruence is
bounded by the size of a Selmer group plus local error terms. The steps
are as follows.

\medskip\noindent\textbf{Step 1 (R): the declaration.} By (H5), $l=2$ or
$(\frac{-d}l)=1$, so the places of $E$ over $l$ come in pairs exchanged by
complex conjugation. For the sign $\epsilon$ attached to the Hodge--Tate
jumps of the setup, a lemma on the place types at $l$ (used here only
through its statement) gives a declaration $(\mathrm{col},\mathrm{row})$
whose rows are conjugate to its columns and such that no place conjugate to
a column place has an embedding of the excluded type. Nothing else is
chosen; the declaration does not depend on $\pi$.

\medskip\noindent\textbf{Step 2 (L)+(N): three local statements at $l$
with explicit exponents.} Let $K/\Q_l$ be a finite extension inside the
cyclotomic level $N$ of the setup (that is, the extension is cut out by
$\mu_{l^\infty}$ at level determined by $v_l(N)$), $L$ a coefficient field
with ring of integers $\mathcal O_L$, and let $T\subset V$ be
$G_K$-stable $\mathcal O_L$-lattices in crystalline representations whose
Hodge--Tate weights lie in $[-h,h]$. The proof uses:
\begin{enumerate}
\item[(i)] (lifting) every surjection $T_1\to T_2$ of such lattices
satisfies the lifting property with exponent $e_L(l,N,h)$;
\item[(ii)] (invariants) if the weights of $V$ form a set of at most $n$
integers in $[-h,h]$ not containing $0$, then $l^{c_0(l,N,h,n)}$ kills
$(V/T)^{G_K}$;
\item[(iii)] (block condition) the $T$-block condition at level $t\le h$
holds with constant $c_B(l,N,h,n,t)$.
\end{enumerate}
These are derived from Fontaine's comparison theorem for $B_{\rm cris}$,
local Tate duality and the local Euler characteristic, the local
Bloch--Kato formalism (all in $\mathrm{Hypotheses}_2$), a lifting theorem
for crystalline lattices with the constants $(23,32,32)$ proved in the
development, and the existence of Wach modules for crystalline lattices,
also proved in the development. The published ingredients are:

\begin{theorem}[Wach, Berger; Berger, Compositio Math.\ 140 (2004), Prop.~II.1.1, Def.~III.4.1 and Prop.~III.4.2]\label{i4:Gst-thm-wach}
Let $K/\Q_l$ be finite unramified, $V$ a crystalline representation of
$G_{K}$ with Hodge--Tate weights in $[0,h]$ (convention: $\Q_l(1)$ has weight $-1$; Berger's ``positive'') and $T\subset V$ a
$G_K$-stable lattice. Then $T$ has a Wach module: a free module of rank
$\dim V$ over $\mathcal O_K[[\pi]]$ stable under $\varphi$ and $\Gamma$, of
height at most $h$, from which $T$ is recovered.
\end{theorem}

\noindent Local Tate duality and the Euler characteristic formula are
Neukirch--Schmidt--Wingberg, Thms.~7.2.6 and~7.3.1.
What is new (N) is the explicit, uniform form of (i)--(iii): the
exponents are explicit functions of $(l,N,h,n,t)$, valid at every
cyclotomic level $N$ (so for ramified $K$, outside Fontaine--Laffaille
range) and for all weights in $[-h,h]$. We do not know (i)--(iii) with
explicit exponents in the literature; we do not describe the proofs of the lifting theorem
and of the Wach-module existence in the development.

\medskip\noindent\textbf{Step 3 (N), with (L) inputs: the conditions at a
good point.} At $\pi$ one checks a list of about thirty conditions (the
``admitted'' package): continuity and semisimplicity of the family; the
decomposition of the restriction to decomposition groups at places away
from $l$ and above $l$; crystallinity and de Rham property of the pieces
$W'$ (on the columns) and $T_2$ (on the rows), with vanishing $H^0$; the
duality of Hodge--Tate multiplicities between $T_2$ and $W'$
($m_{T_2}(\tau,a)=m_{W'}(\tau,-a-1)$); counting of filtration jumps; a
finite-rank torsion bound for $H^1_f$ of $T_2$; triviality of unramified
$H^1_f$ at column places; the tame generator conditions at $Q_k$; the
lift exponents (Step~2(i)) on columns and rows; the invariant length bound
(Step~2(ii)) with constant $c_0$; the block condition (Step~2(iii)) with
constant $c_B$; and an admissibility condition at good points, which uses
the Chebotarev density theorem (Neukirch VII (13.4)) and residual
absolute irreducibility. Most items come from the fields of the setup
(norm-choice facts, canonical hub, admissible parameters) through
bookkeeping (R); the Chebotarev step and the Hodge--Tate jump window
are (N) statements of the development whose proofs we do not describe.

\medskip\noindent\textbf{Step 4 (L): local Galois cohomology.} From global
Artin reciprocity, local Tate duality and the local Euler characteristic
one obtains the finiteness of $H^1(G_{K_w},-)$ for the finite modules
that occur, the duality between $H^1$ of a module and of its Cartier dual
and the Euler characteristic formula; and the annihilator form of the
local Bloch--Kato conditions ($H^1_f$ of $V$ and of $V^*(1)$ are exact
annihilators under the Tate pairing; Bloch--Kato, Prop.~3.8;
reference to be confirmed). These are routine specialisations of the
published statements.

\medskip\noindent\textbf{Step 5 (N): the lattice-index theorem.} This is
the core. Under the conditions of Steps~3 and~4 it gives
\begin{multline*}
v_{cc}(k,\pi)\le\sum_{w\in\mathrm{row}\cup\mathrm{col}}\mathrm{len}_w
+s_\pi+\sum_{0<e<t}\mathrm{sha}_l(e)+v_l(h_E)
+t\sum_{w\in\mathrm{row}}(n_w+2)c_B\\
+4n_V(|\mathrm{col}|+1)\delta+t\,|\mathrm{col}|\,h_0
+t\sum_{v\in Q_k}\ell\bigl(\mathcal O_k/(\zeta_v-1)\bigr),
\end{multline*}
provided $l\ne g_0,g_1$, $\pi$ has finite $s$-length and
$\mathrm{sha}_l(e)<\infty$ for $0<e<t$; here $\mathrm{len}_w$ are the
local length terms at the places over $l$ (column terms from the lift
exponent $e'$ and the row length bound). The proof has the following
shape, read from the top two levels of the formal proof.
\begin{enumerate}
\item[(5a)] \emph{Normalisation.} The family $r$ is replaced by a
normalised family, block upper triangular with respect to the residual
decomposition into $t$ blocks twisted by powers of the cyclotomic
character, with the same semisimplification; all conditions of Step~3 are
transported to it.
\item[(5b)] \emph{Separating element.} The admissibility at good points
gives $g_0$ in the Galois group whose eigenvalues separate the blocks, and
idempotents $e_B$, $e_1$ cutting out the blocks.
\item[(5c)] \emph{Generalised matrix algebra.} One forms the ring datum: the
reduced Hecke algebra $T'\subset\prod_{x}\mathcal O_k$ of the points, the
$T'$-algebra $R'$ generated by the image of the Galois group in
$\prod_x M_{n_V}(\mathrm{Frac}\,\mathcal O_k)$, the block idempotents, the
element $\delta=\prod_{0<i<t}(r_0-r_i)\in T'$, and the specialisation
$R'\to M_{n_V}(\mathcal O_k)$ at $\pi$. After a rank reduction the block of
index $0$ has rank one.
\item[(5d)] \emph{Ribet's construction.} The off-diagonal parts
$e_iR'e_j$ give, after specialisation at $\pi$, extensions of the twisted
residual characters, hence classes in $H^1$ of the twists
$\mathcal O_k(j)$, $0<j<t$, and of $K_k/\mathcal O_k(j)$. The congruence
valuation is bounded by the length of the module these classes generate.
\item[(5e)] \emph{Local conditions.} At places over $l$ the classes are
crystalline up to the lift exponent and the block constant (Step~2); at
$v\in Q_k$ their failure to be unramified is bounded by
$\ell(\mathcal O_k/(\zeta_v-1))$; at column places the invariants
contribute $h_0$ and $h^W_u$. The classes therefore lie in a Selmer group
up to these error terms, and the Selmer group has length at most
$\sum_{0<e<t}\mathrm{sha}_l(e)+v_l(h_E)$ (the class-number term comes
from the unramified class field bound).
\end{enumerate}
The method is that of Ribet (Invent.\ Math.\ 34 (1976)), in the form of
generalised matrix algebras and reducibility ideals of Bella\"iche and
Chenevier (Ast\'erisque 324 (2009), Ch.~1) and of Urban (Duke Math.\ J.\
106 (2001)); all three reference to be confirmed. What is new (N) is the
quantitative statement: an explicit bound for an arbitrary number $t$ of
blocks, with explicit local error terms at places over $l$ (through
Step~2) and at the tame places $Q_k$. We know no published statement with these
terms. We do not describe the formal proof below the two levels described.

\medskip\noindent\textbf{Step 6 (R): specialisation to the split bound.}
One applies Step~5 with $e'=e_L(l,N,h_N)$,
$c_0=c_0(l,N,h_{\rm hub},4nn_V)$ and $c_B=c_B(l,N,\lfloor2b\rfloor,4nn_V,t)$.
Because the rows are conjugate to the columns, the row length terms are
bounded by the corresponding column terms, and
$\sum_{w\in\mathrm{row}\cup\mathrm{col}}\mathrm{len}_w$ becomes the column sum
in the definition of $\mathrm{bs}$; this gives (S). (The passage was
compared only at the level of statements.)

\medskip\noindent\textbf{Step 7 (R): (W) and (Z).} (W) is Step~2(ii)
applied to $W'$ at each column place, whose Hodge--Tate jumps are nonzero
by the decomposition above $l$ (Step~3); the constant is
$r\,c_0(l,h_{\rm hub},4nn_V)$. (Z) is one of the tame generator conditions
of the good point.

\medskip\noindent\textbf{Summary.} Published: Fontaine's comparison,
local Tate duality and Euler characteristic, Bloch--Kato's local
formalism, Wach modules, Chebotarev. New and resting on
\cite{CharacterFamily}: the explicit local exponents of Step~2, the
conditions at good points of Step~3, and the quantitative lattice-index
theorem of Step~5.

\medskip\noindent\label{i4:Gst-galois}\emph{Structure of the proof of Proposition~\ref{i4:A-prop-galois}.}

Proposition~\ref{i4:A-prop-galois} is formally verified. Recall that
$\mathrm{LinkN}$ for a family $F$ means: there are $C,x_0$ such that for
$x\ge x_0$, every index $i$ with the odd Shafarevich groups finite, and
every lift index $j$ of the datum $D=F.D(x,i)$, $L=F.L(x,i)$,
\[
\sum_{w}\mathrm{lw}(w)\,e_w\,\eta_j(w)\le d_{Z_0}(j)\Bigl(C\bigl(fmb\ell^2+(\log x)^2N_c\bigr)
+\sum_{l}\Bigl(\sum_{1\le e<t}\mathrm{sha}_l(e)\Bigr)\log l\Bigr),
\]
the sum on the left over the places $w$ in the support of $L$ where the
local type differs from the base type, with weight $\mathrm{lw}(w)$,
ramification index $e_w$ and $\eta$-length $\eta_j(w)$ (whenever the series
over $l$ converges). The proposition upgrades the bound at $2$ and $3$
to this global bound. The steps are as follows.

\medskip\noindent\textbf{Step 1 (R): transport to the witness points.}
The parameter system $p_0$ has witnesses that pass along the set
division; the Galois preconditions and the fourfold tie, once proved at
members (Steps~2 and~3), transfer to the witness points through the
position condition. The tie of characters at members comes from
depth-zero support and the first component of the compatibility of
geometric Galois representations.

\medskip\noindent\textbf{Step 2 (N), with (L) inputs: the Galois
preconditions at members.} At every member datum eighteen conditions
hold: the theta point conditions; the admissible datum; the first
provision; the multiplicity statement; residual irreducibility in rank
$r$; admissible residual tame eigenvalues; the weight and block-character
parameters; unramifiedness of $\mu$; depth zero of the $G$-part; the
spherical inertia datum; the $\lambda$-condition; uniqueness of
localisation; largeness of the coefficients; the norm-choice tie; good
triples; the same-twist condition; good finite types. They are derived
from the published inputs of the third layer, the compatibility of
geometric Galois representations, $\mathrm{SerN}$, depth-zero support,
local finiteness, local Tate duality and the local Euler characteristic.
Two ingredients are classical:

\begin{theorem}[Brauer--Nesbitt; Curtis--Reiner, \emph{Representation Theory of Finite Groups and Associative Algebras}, Interscience 1962, (30.16); reference to be confirmed]\label{i4:Gst-thm-bn}
Let $G$ be a group and $k$ a field. Two semisimple finite-dimensional
representations of $G$ over $k$ which have the same characteristic
polynomial at every element of $G$ are isomorphic.
\end{theorem}

\noindent Residual irreducibility is obtained from it together with a
three-prime argument of the development. The Chebotarev
input is the one recorded for Step~3 of \S\ref{i4:Gst-split}. The
remaining conditions are (N) statements whose proofs we do not describe.

\medskip\noindent\textbf{Step 3 (N): a good cut.} At a datum $D$ with the
Galois preconditions and the fourfold tie, a central-idempotent argument
gives a subset of the points (a \emph{cut}) which is good, such that the
restricted datum still has the Galois preconditions and the fourfold tie,
and such that it satisfies the norm-choice facts. The norm-choice facts
are deduced from the Galois preconditions on the cut using the Chebotarev
input, local finiteness, local Tate duality and the local Euler
characteristic (the second conclusion, including the provisions $G_5$ and
$G_6$, and the de Rham weight jumps). The cut does not change the
quantities in $\mathrm{LinkN}$ except through the good points, which
are kept.

\medskip\noindent\textbf{Step 4 (N): uniform constants.} From
$\mathrm{Hypotheses}_2$ there are $C\ge0$ and $C_T\in\N$ such that every
datum with the norm-choice facts satisfies the commuting-endomorphism
condition and the paired-character condition with these constants.

\medskip\noindent\textbf{Step 5 (N): the local bound table.} For a family
with twisting-good data, the facts of Steps~3--4 and the bound at $2$ and
$3$, one constructs, for $x$ large and every index with finite odd
Shafarevich groups, a hub datum $X$ and a constant $C_s$ such that at
every twisted place a local bound holds, by type of the place:
\begin{itemize}
\item $l\in\{2,3\}$: the hypothesis $NN_{23}$ (the bound at $2$ and $3$,
that is, Proposition~\ref{i4:A-prop-at23});
\item $l=q$ or $l=d$ (ramified places): a ``Theorem B'' bound from the
embedding-jump estimates of the setup;
\item large primes: the local bound is $0$, from a congruence bound for
the family and the killing condition;
\item medium and small inert primes and band places: length bounds read
off the hub data;
\item all places: the killing condition, the closing constant from local
Bloch--Kato, Tate duality and the Euler characteristic, the joint
condition at $B$, $C$, $D$ from Fontaine's comparison, polynomially
bounded parameters, and few and small excluded and special primes.
\end{itemize}
The constant is
\[C_s=\max(3C_2+1,\,22365,\,C_G,\,2000(C+1)+5C_1,\,C_1,\,C_{23},\,C_T),\qquad C_2=\max(C_0,1).\] The assembly is (R); the individual bounds are
(N) statements whose proofs we do not describe.

\medskip\noindent\textbf{Step 6 (N): summation.} The local bounds of
Step~5 are summed over the twisted places with the place weights; the
contribution of the Shafarevich groups is the term
$\sum_l(\sum_{1\le e<t}\mathrm{sha}_l(e))\log l$. This gives $\mathrm{LinkN}$
for the family. Inspected only at the level of its two top statements.

\medskip\noindent\textbf{Step 7 (R): slots.} The two slots (the one
carrying the Galois preconditions and the one carrying $NN_{23}$) are
intersected by Lemma~\ref{i4:A-lem-calc}(c), and closedness is inherited.

\medskip\noindent\textbf{Summary.} Published: Chebotarev,
Brauer--Nesbitt, local Tate duality and Euler characteristic, Bloch--Kato
and Fontaine (all through $\mathrm{Hypotheses}_2$ or the third-layer
inputs). New and resting on \cite{CharacterFamily}: the Galois
preconditions at members (Step~2), the good cut (Step~3), the constants
(Step~4), the local bound table (Step~5) and its summation (Step~6). No
step of this chain was identified with a published theorem.

\medskip\noindent\label{i4:Gst-sern}\emph{Structure of the proof of Proposition~\ref{i4:A-prop-sern}.}

Proposition~\ref{i4:A-prop-sern} is formally verified. Its hypotheses
are two bundles. The \emph{non-endoscopic third-layer hypotheses} consist
of: local representation theory of $\mathrm{GL}_n$ (Bernstein--Zelevinsky
finiteness, irreducibility of unitary products, the geometric lemma,
Zelevinsky's polynomial ring, totally disjoint products, Moy--Prasad and
Jacquet modules, tempered and square-integrable segments, generic
unlinked products, Langlands multisegments and exhaustion, standard
cosocles, a local Langlands correspondence, the local converse theorem,
depth-zero induction, Brauer's degree-zero theorem, Chevalley units,
Green's regular characters); global automorphic and Galois inputs (rapid
decay of cusp forms, Harish-Chandra finiteness, finiteness of automorphic
forms on $\mathrm{GL}$, Jacquet--Shalika, the cuspidal spectrum of
$\mathrm{GL}$, isobaric base change, Gelfand--Kazhdan, Galois
representations attached to polarised regular algebraic cuspidal pairs,
local-global compatibility at $l$, Clozel's purity and cohomological
criteria, global Artin reciprocity, ray class fields, $l$-adic avatars of
algebraic Hecke characters, Weil--Deligne representations of algebraic
characters, de Rham implies potentially semistable); Galois cohomology
(normalised Poitou--Tate, local and global finiteness and Euler
characteristic, unramified orthogonality, Greenberg--Wiles, Bloch--Kato's
local formalism); a join statement, toric transport and the Hecke--Tate
functional equation. The second bundle is the compatibility of geometric
Galois representations (Weil representation of unitary pairs, existence,
analyticity and classification of discrete series, trace class of cusp
forms, Dixmier--Malliavin, Flath, the Siegel--Weil formula, the doubling
method, unramified doubling for $\mathrm{GL}$, continuation of partial
Hecke $L$-functions, a height count). The steps are:

\medskip\noindent\textbf{Step 1 (L)+(R): local Langlands datum.} The
local Langlands correspondence for $\mathrm{GL}_n$ is a field of the
local hypotheses; the development checks from the listed local inputs the
normalisation properties it needs (compatibility with Artin reciprocity
in rank one, depth zero, exhaustion). The published theorem is:

\begin{theorem}[Harris--Taylor, Ann.\ of Math.\ Stud.\ 151 (2001); Henniart, Invent.\ Math.\ 139 (2000), 439--455; reference to be confirmed]\label{i4:Gst-thm-llc}
For a $p$-adic field $F$ and every $n$ there is a unique family of
bijections $\mathrm{rec}_F$ from irreducible smooth representations of
$\mathrm{GL}_n(F)$ to $n$-dimensional Frobenius-semisimple Weil--Deligne
representations of $W_F$, compatible with local class field theory for
$n=1$, with central characters and determinants, with twists and
contragredients, and preserving $L$- and $\varepsilon$-factors of pairs.
\end{theorem}

\medskip\noindent\textbf{Step 2 (R): the published inputs of the third
layer.} They are a repackaging of fields of the two bundles (theta local
inputs, rapid decay, Harish-Chandra finiteness, Green's criterion,
Zelevinsky, ray class field and Artin characters, base change for definite
unitary groups, the generic unitary dual and standard cosocles of
$\mathrm{GL}$, Chebotarev, unramified Howe duality), with Green's
criterion obtained from Green's regular characters and Chebotarev for
degree-one primes proved in the development.

\medskip\noindent\textbf{Step 3 (N): rationality of the twisted theta
series.} The twisted theta series attached to the data of the family have
Fourier coefficients in the expected field on good covers, after a phase
of absolute value $1$. Proof: rationality of the untwisted theta series
on good covers (proved in the development, no hypothesis) and the
action of the Galois group on ray class characters, which is global Artin
reciprocity. The method is that of Shimura's rationality results for
theta series (reference to be confirmed); the statement, with its specific
models and covers, is not known to us in the literature.

\medskip\noindent\textbf{Step 4 (N): integrality.} The twisted theta
series are integral at every prime power, uniformly. Proved in the
development with no hypothesis; we do not describe its proof.

\medskip\noindent\textbf{Step 5 (N): the three Hecke clauses.} Hecke
integrality, the Hecke transform (from pointwise Hecke laws, which follow
from the theta transport law) and the Hecke read tie (from a split linear
model and a fixed vacuum vector) are proved in the development with no
hypothesis. We do not describe the proof below this level.

\medskip\noindent\textbf{Step 6 (N), from the transport input: the ratio
identity.} Let $T_0\ne0$ be a toric sum and $T(q)$ the class terms of a
Hecke character $\chi$ over the Picard group of conductor
$\mathfrak f\prod_{v\in Q}v$. The ratio identity asserts that for every
character $\psi$ of this Picard group the ratio
$\sum_q\psi(q)^{-1}T(q)/T_0$ is the image of an algebraic number
$\rho(\psi)$, and that an automorphism $g$ fixing the reflex data and the
values of $\chi$ acts by $g(\rho(\psi))=g(\psi(t))^{-1}\rho(g\circ\psi)$,
where $t$ is an idele with $\mathrm{art}(t)=g^{-1}$ on the abelianisation.
It is deduced from the toric transport hypothesis (a field of the
non-endoscopic hypotheses): the transport gives algebraicity of each
twisted sum and its conjugation law, and the ratio is formed. This is a
Shimura-reciprocity statement for CM values (Shimura; reference to
be confirmed); we have not compared the transport hypothesis itself with Shimura's
statement, so the deduction is (N) given an (L)-type input, which
the formal development takes as an argument here and proves elsewhere.

\medskip\noindent\textbf{Step 7 (R): the elliptic hypothesis.} Its
twenty-five items are fields of the hypotheses or of Steps~3--6, plus:
Chebotarev for degree-one primes, Poitou--Tate from its normalised form,
local-global compatibility above $l$ from the Galois representations of
cuspidal representations, and the local Bloch--Kato condition at all
primes from Bloch--Kato, Tate duality and the Euler characteristic. One
auxiliary lemma (the ``twist front inputs'') is not described here.

\medskip\noindent\textbf{Step 8 (N): local Langlands hypothesis at
members.} Archimedean local Langlands compatibility, stability of
conjugate-dual parameters at primes not dividing the level, the point
statement from the inputs and geometric compatibility, and the cusp
parameter at members. We use it only through its statement.

\medskip\noindent\textbf{Summary.} Published, re-proved in the formal development as separate
propositions: the two
bundles above, in particular the local Langlands correspondence, the
Galois representations of polarised regular algebraic cuspidal
representations, local-global compatibility, and Poitou--Tate. New and
resting on \cite{CharacterFamily}: Steps~3--6 and~8; Step~6 depends on
the toric transport hypothesis.

\medskip\noindent\label{i4:Gst-funcov}\emph{Structure of the proof of Proposition~\ref{i4:A-prop-funcov}.}

Proposition~\ref{i4:A-prop-funcov} is formally verified. Its
hypotheses are exactly: the unitary-group third-place hypotheses (the
list of the remark after the proposition, together with the join
statement, the Hecke--Tate functional equation and the local
representation-theoretic inputs of the non-endoscopic hypotheses), the
local theta inputs, the agreement of theta integrals in the Weil tensor
models, the local matching, the covolume bound for $X$ and a local
Langlands datum. The four parts are proved in the order (d), (c), (a),
(b), since (a) uses (d).

\medskip\noindent\textbf{(d), surjectivity onto $X$.} The chain is:
\begin{enumerate}
\item[(d1)] (R) Global Artin reciprocity and the ray class field statement
follow from the existence of $l$-adic avatars of algebraic Hecke
characters; rationality of theta series is Step~3 of
\S\ref{i4:Gst-sern}.
\item[(d2)] (N) An ``exported statement'' from reciprocity and
Casselman--Wallach; page, geometric, rank, model and share statements,
proved in the development with no hypothesis; a twist statement from
these.
\item[(d3)] (N) A bad-places statement from Green's cuspidality criterion
and the model statement.
\item[(d4)] (N), with (L) inputs: the exponent statement. It uses
Clozel's criterion for cohomological representations, $l$-adic characters
attached to algebraic Hecke characters, Chebotarev for degree-one primes,
Galois representations attached to polarised cuspidal representations,
base change for definite unitary groups, Bhatt--Scholze crystalline
lattices, $l$-adic Hecke parameters above $l$ and the local Langlands
correspondence. The published theorems are:
\end{enumerate}

\begin{theorem}[Galois representations of polarised cuspidal representations; Barnet-Lamb--Gee--Geraghty--Taylor, Ann.\ of Math.\ (2) 179 (2014), Thm.~2.1.1, with Chenevier--Harris, Shin, Caraiani; at $w\mid l$ see below]\label{i4:Gst-thm-galrep}
Let $E$ be a CM field, $\pi$ a regular algebraic cuspidal automorphic
representation of $\mathrm{GL}_a(\mathbb A_E)$ and $\chi$ an algebraic
character of $\mathbb A_{E^+}^\times/E^{+\times}$ such that $(\pi,\chi)$ is
polarised. For every $l$ and $\iota:\overline\Q_l\cong\C$ there is a
continuous semisimple $r_{l,\iota}(\pi):G_E\to\mathrm{GL}_a(\overline\Q_l)$,
de Rham at places over $l$ with Hodge--Tate weights determined by the
infinity type of $\pi$, such that $\mathrm{WD}(r_{l,\iota}(\pi)|_{G_{E_w}})^{\rm F\text{-}ss}
\cong\mathrm{rec}(\pi_w\otimes|\det|^{(1-a)/2})$ at every finite place
$w$ (for $w\mid l$ with $\pi_w$ Iwahori-spherical: BLGGT Thm.~2.1.1(4); in general: Barnet-Lamb--Gee--Geraghty--Taylor, \emph{Local-global compatibility for $l=p$} I, II, and Caraiani 2014).
\end{theorem}
\noindent At $w\mid l$ the sources supply the following. BLGGT Thm.~2.1.1(3)
gives de Rham with the stated Hodge--Tate weights at every $w\mid l$, and
Thm.~2.1.1(4) the displayed isomorphism when $\pi_w$ has an Iwahori-fixed
vector. Barnet-Lamb--Gee--Geraghty--Taylor, \emph{Local-global
compatibility for $l=p$}, I (Ann.\ Fac.\ Sci.\ Toulouse Math.\ 21 (2012),
57--92) and II (Ann.\ Sci.\ ENS 47 (2014), 165--179), give it for every
$\pi_w$ up to semisimplification. Caraiani, Algebra Number Theory 8 (2014),
1597--1646, gives the monodromy operator for $a$ even and the
Frobenius-semisimplified form. The full displayed isomorphism at $w\mid l$
for $a$ odd (agreement of monodromy) is not covered by these
statements; it is used here only to classify steps of the formal development.

\begin{theorem}[base change for definite unitary groups; Labesse, \emph{Changement de base CM et s\'eries discr\`etes}, in \emph{Stabilization of the trace formula, Shimura varieties, and arithmetic applications}, Vol.~1, International Press 2011; Clozel--Harris--Labesse; reference to be confirmed]\label{i4:Gst-thm-bc}
Let $G$ be a unitary group in $n$ variables attached to $E/E^+$, definite
at all infinite places. Every automorphic representation $\sigma$ of $G$
has a base change to $\mathrm{GL}_n(\mathbb A_E)$: an isobaric sum of
conjugate self-dual discrete automorphic representations whose local
components at places split in $E$ and at places where $\sigma$ is
unramified are the local base changes of those of $\sigma$.
\end{theorem}

\begin{theorem}[Bhatt--Scholze, Cambridge J.\ Math.\ 11 (2023), Thm.~1.2]\label{i4:Gst-thm-bs}
For a complete discretely valued extension $K$ of $\Q_p$ with perfect
residue field, lattices in crystalline representations of $G_K$ are
equivalent to prismatic $F$-crystals on $(\mathcal O_K)_\Delta$.
\end{theorem}

\noindent The extra content of (d4) is the uniform bound in the family
(exponents uniform in $x$ and the index), which is not in these theorems.
\begin{enumerate}
\item[(d5)] (R) Surjectivity onto $X$ is assembled from the exported,
twist, share, bad-places and exponent statements.
\end{enumerate}

\medskip\noindent\textbf{(d), the lower bound $\lambda_{\rm low}$.}
\begin{enumerate}
\item[(d6)] (N) Eventually along the family the theta lifts have the
weight sizes $\mathrm{Sizes}(D,j,b)$; proved in the development.
\item[(d7)] (N), with (L) inputs: a lower bound for the relevant real
datum, from rapid decay, the Siegel--Weil formula in Weil's range and in
the second range (Theorems~\ref{i4:thm-sw} and~\ref{i4:thm-sw2}), the
doubling method for local zeta integrals (Piatetski-Shapiro--Rallis;
reference to be confirmed), unramified doubling for $\mathrm{GL}$,
Bernstein--Zelevinsky finiteness (Theorem~\ref{i4:thm-bzfin}) and the
formal degree in the compact-dual normalisation, combined with an
elementary lower bound for a hook-length log sum. The published inputs
enter as identities of integrals; the explicit lower bound is new.
\end{enumerate}

\medskip\noindent\textbf{(c), rationality of rows.} (N) Test vectors are
assembled from a Siegel set of the unitary group, real Howe duality,
Green's cuspidality criterion, the theta local inputs, the join, the
functional equation, the agreement of theta integrals, the local matching
and the separation of functionals (below); rationality of rows follows
from the row tests. We do not describe the proof below these two levels.

\medskip\noindent\textbf{(a), the functional condition.}
\begin{enumerate}
\item[(a1)] (N), with (L) input: the local functionals separate. This
follows from unramified Howe duality (proved in the development) and the
Siegel--Weil formula in the convergent range (Theorem~\ref{i4:thm-sw2}).
\item[(a2)] (N) An evaluation-grid constant $C_{F1}$: from real Howe
duality, the local matching, the theta local inputs, the join and the
functional equation, the agreement of theta integrals, convergence of the
twisted Weil representation (from the Siegel set), matching of the
twisted local factors and equality of twisted square theta integrals.
\item[(a3)] (N) Constant-function fit: the Fourier coefficients obtained
from the join fit a constant function, using (a1) and the lower bound
$\lambda_{\rm low}$ of (d).
\item[(a4)] (N), with (L) inputs: at the grid of size $C_{F0}$, the
geometric condition on the large grid (from the ray class field), and the
row point summands from the existence of discrete-series covers,
Harish-Chandra modules, rapid decay and the Casselman--Wallach
globalisation; with (a3) this gives the functional condition at every
point of the class.
\item[(a5)] (R) With $c_{\rm Dep}=\max(C_{F0},C_{F1})$ the open
statements at $p_0$ hold and pass to the extension $p$ by monotonicity.
\end{enumerate}

\medskip\noindent\textbf{(b), the covolume row.} (R) given the hypothesis:
the class-sum divisor condition (from the ray class field) and the row
point summands of (a4) reduce the covolume row to the covolume bound for
$X$, which is a hypothesis of the proposition.

\medskip\noindent\textbf{Summary.} Published, re-proved in the formal development as separate
propositions: the list of
the remark after Proposition~\ref{i4:A-prop-funcov}; in particular the
Galois representations of polarised cuspidal representations, base change
for definite unitary groups, Bhatt--Scholze, Siegel--Weil, the doubling
method, Casselman--Wallach and Harish-Chandra's finiteness. New and
resting on \cite{CharacterFamily}: (d2)--(d4), (d6), (d7), (c), (a1)--(a4).
(b) is bookkeeping from its hypothesis.


\subsubsection{The automorphic leaves}\label{i4:Ast-ssec}

\noindent\label{i4:Ast-key}The four Propositions~\ref{i4:A-prop-rowleaves}, \ref{i4:A-prop-inputs}(b)--(d)
with \ref{i4:A-prop-inputs-leaves}, \ref{i4:A-prop-stages} and
\ref{i4:A-prop-open} are formally verified, \cite{CharacterFamily}. We describe
the chain of steps of each formal proof, read from its top three levels; here
(R) steps are transport of data along identifications, conjunction of eventual
statements and choice of thresholds.  A large part of the published input enters the development as separate
propositions: the published statements of Remark~\ref{i4:A-rem-pub} are
isolated, taken as arguments by intermediate formal lemmas, and proved
inside the development, whose final theorem has no hypotheses. When a step
is marked (L) and the development re-derives the published theorem inside a
longer argument instead of isolating it, we say so.

\medskip\noindent\label{i4:Ast-rowleaves}\emph{Structure of the proof of Proposition~\ref{i4:A-prop-rowleaves}.}

The items fall into six groups according to the method.

\medskip\noindent\emph{Group I: integrality over true columns, (P0)--(P6).}
No published input is used in this group; every item is unconditional.
Recall that a built record of a core carries a weighted column family $w$
over $(B,f,T)$ and a lattice stage $Z$, and that the datum of the core is the
datum $w.D(Z)$ read off the true columns of $w$.
\begin{enumerate}
\item[(P0)] (R) \emph{Column view.} A built record supplies four components:
the jets condition of the column input, the row integrality (every reference
row value at the reference tuple $\mu_0$ and every twisted row value at every
used tuple is integral), the identity of the twist of the core with the
twist read off the true columns, and a tuple tied to the norm choice. These
four form the column view. They are transported along the set division of
the lattice stage by the corresponding identities of the construction, so
that the property at the built point reduces to the same property at
$w.D(Z)$.
\item[(Q1)] (R) \emph{Base signs.} For ranks $(k,m,t)$ with $m\ge2$, $m+3\le
k\le t$ and a Harish-Chandra parameter $P=(u,v,b,e)$ with $-b_0<u_{k-1}$ and
$e_j-b_j\ge(t+3)/2$: if $u_{k-1}\le-(k-1)/2$ then $e_0\ge
b_0+(t+3)/2>-u_{k-1}+(t+3)/2\ge k/2$, and if $u_{k-1}\le0$ then
$e_0>(t+3)/2>k/2$. At the top places of a built point the size data give
$u_{k-1}=-b_T-(k-1)/2$ with $b_T>0$, at the other places
$u_{k-1}=-1/2$. Hence every built base satisfies the sign law
($u_{k-1}\le-(k-1)/2$ at top places, $e_0\ge k/2$ everywhere).
\item[(Q6)] (N) \emph{Standard column comparisons.} Over true columns there
is a system of standard column comparisons: for every word $u$, a
comparison of the column forms of $w.D(Z)$ with the standard forms of the
model, compatible with the pins of the column input and with level data.
Its proof unwinds the validity condition of the column input.
\item[(P1)] (N) For every word $u$, the reference-chain word defines a map
of forms (the forms-map statement at every true column member). It follows
from the validity condition and the pins of the column input alone.
\item[(P2)] (N) \emph{Twisted point integrality.} For every word $u$ and
every used tuple $\mu\ne\mu_0$ the twisted form at $(u,\mu)$ is rational
with integral numerator. The chain is: twisted evaluation branching holds
for all tuples over true columns; it implies lower twisted compatibility;
the lower tie is integral.
\item[(P3)] (N) \emph{Period at zero.} Under row integrality: for every word
$u$, every lattice point $x$ and every form $F$ of the reference component
lying in $V_0(\mu_0)$, the reference period of $(x,F)$ is the image of an
element $y$ of the reflex field with $y=0$ or
$v_w(y)\ge-\mathrm{ur}(D,w,v_{D_\Gamma}(u,w))$ at every place $w$. Proof:
the period decomposes, at the reference tuple, along the standard column
comparisons (Q6), with integral coefficients by row integrality; the
class-by-class decomposition gives the valuation bound.
\item[(P4)] (N) \emph{Twisted periods.} Under row integrality and a tied
tuple, the twisted periods are admissible for every division datum. One
chooses a true column choice with standard chain data, decomposes the
twisted period by classes, and reduces admissibility to place tuples. At a
class $c$: restricting a vector of the class-indexed period space to $c$
commutes with the word embedding of the twisted point; row integrality gives
a lattice vector at the class representative with the given row value; the
integrality of the class pairing on integral row forms gives an integer $b$
of the reflex field whose image is the restricted period; hence $b=0$ or
$-\log|b|_w\ge0$ at every $w$.
\item[(P5)] (N) \emph{Real-side tie.} The coefficient fields over true
columns contain the required primitive roots of unity; with the standard
column comparisons this gives symmetric tied column forms whose characters
at step zero are of finite order.
\item[(P6)] (N) \emph{Level-prime column valuations.} For every word $u$
and tuple $a$ the five-point laws of the true chains give a lattice-column
chain for $(\mathrm{tw}(Z),u,a)$; the size condition $2f N_m+1<T(u)$ holds
($[E:\mathbb Q]=2f$ and the two size inequalities of the pre-datum); the
weight signs of this chain are those of the base, which satisfy the sign
law by (Q1). Stability of sharp bounds along lattice-column chains then
gives the sharp fifth conclusion at $(w.D(Z),u)$ (case $a=\mu_0$) and the
column bound $\mathrm{col}\ge-m\,E_T(u,w)$ at the places of residue
characteristic $T(u)$ (case $a\ne\mu_0$, using compatibility of the twisted
chain with the initial lattice).
\end{enumerate}
The inner statements (twisted evaluation branching, the period
decompositions, sharp-bound stability) are specific to the construction of
the true columns; we know of no published counterpart. We do not describe the proof below
the level just described.

\medskip\noindent\emph{Group II: lattices, classes and masses, (R1), (R4)--(R6).}
\begin{enumerate}
\item[(R1a)] (L) Over every column base lattice there is a self-dual lattice
whose adelic stabiliser contains the level group. The local input is that a
unimodular hermitian lattice over an unramified quadratic extension of a
$p$-adic field has an orthonormal basis (Jacobowitz's classification;
reference to be confirmed; the argument is the one recalled in the proof of
Theorem~\ref{i4:thm-genusmass}). The development proves it; we do not describe the proof
in detail.
\item[(R1b)] (N) The stabiliser index formula
$[\widehat{\mathrm{Stab}}(N):U]=\mathrm{idx}(g_0,g_1)\,(T_u^k-1)^f$: a
computation of the index of the level group in the stabiliser of the
self-dual lattice, place by place, as a ratio of orders of finite groups.
We do not describe the proof below the statement.
\item[(R1c)] (N) Column stabilisers are trivial (the freeness hypothesis of
Lemma~\ref{i4:A-lem-classnum}(a)). We do not describe the proof below the statement.
\item[(R6)] (L) Genus mass, finiteness of the class number and the class
mass at every rank: for a definite hermitian space, the set of classes in a
genus is finite and the mass is $\sum_{c}1/\#\mathrm{Aut}(c)$, independent of
the representative lattice of the genus up to isometry. Finiteness is
Borel's theorem (Borel, Publ.\ Math.\ IHES 16 (1963), 5--30, Thm.~5.1), the rest is the definition of the mass together with the
action of the rational points on the adelic lattices. The development
proves these statements without the published hypothesis; we do not describe the proofs in
detail.
\item[(R4)] (L)+(N) \emph{Volume ratio from Siegel--Weil.} Method: Siegel's
inductive computation of volumes. (i) (L) Theorem~\ref{i4:thm-sw2}, in its
form for one-vector Gaussian means of compactly supported data, gives
measurability and the mean of the one-vector theta function. (ii) (N) The
volume of the adelic quotient of $U(W)$ is compared with that of $U(W')$,
$W'=x^\perp$ for a unit vector $x$, by unfolding the one-vector theta
integral along the orbit of $x$. The six inputs of the unfolding are: the
density of self-dual lattices, the measures at the real places, the Euler
product of local densities, independence of the volume from the choice of
lattice in the isometry class, existence of a unit vector with self-dual
orthogonal complement, and the unfolding identity. Each is a classical
step of Siegel's method (Siegel, Ann.\ of Math.\ 36 (1935), 37 (1936); Weil,
\emph{Ad\`eles and algebraic groups}; reference to be confirmed); the development
proves all six with explicit normalisations, which is the new part.
\item[(R5)] (N) \emph{Volume equals mass at every signature.} For definite
$W$ the volume is the mass of its genus (class decomposition and the double
coset mass); for $Q>0$ and $P+Q\ge2$ the self-dual comparison of (R4) changes
the signature without changing the volume; for $(P,Q)=(0,1)$ the space $-W$
is definite of rank one with the same volume. Hence the volume depends only
on the rank and equals $\mathrm{mass}_I(D,P+Q)$.
\end{enumerate}

\medskip\noindent\emph{Group III: rows obtained from the Siegel--Weil formula, (R12)--(R14), (R16), (R17).}
\begin{enumerate}
\item[(R12)] (L)+(R) The rank row and the torsion-frame row on every class
with a twist selection: by (R4) the volume ratio holds for every datum, and
a ranks-and-frame statement on tied classes turns it into both closed rows.
The only published input is Theorem~\ref{i4:thm-sw2}.
\item[(R13)] (N) \emph{Loop rows.} Given finite dimensionality of rows
(Proposition~\ref{i4:A-prop-inputs-leaves}(M5)) and the volume ratio, the
loop row at level three holds on every class with $T_u$-size. The proof uses
four internal statements: multiplicity one of Bessel models at large gap
(under a uniform gap hypothesis on the Harish-Chandra parameter), a
comparison of Whittaker and Bessel functionals, an unconditional pointwise
bound for matrix coefficients of discrete series, and decay of the relevant
coefficients from the gap. Multiplicity one of Bessel models and bounds for
discrete-series matrix coefficients are published in other forms
(Harish-Chandra; Trombi--Varadarajan; Gan--Gross--Prasad for Bessel
models); the forms used here carry explicit constants and a gap
hypothesis, and are proved in the development. We count them as (N).
\item[(R14)] (N) \emph{Off-diagonal (order) bound.} From the first
Siegel--Weil row, row non-growth, finite dimensionality, the local tie, the
rank row and the loop row, the selected Theorem~S holds closed. Its proof is
the carrier-gap argument: Theorem~S for distinct words from the leaves of
Theorem~S (carrier, non-growth, uniform gap, cross-section statement), then
the off-diagonal bound from the level-three rows. This is the analytic core
of the counting and is new.
\item[(R16)] (L)+(R) Counting constants: from the volume step of
Siegel--Weil, for every $A$ there are $n_0(A)$, $C_0(A)\ge0$ with the counting
statement; one takes $A=2(C_{YL}+1)$, $C=C_0(A)$, $c=c_b=c_{Ic}=1/16$,
$C_{\mathrm{res}}=C_b=0$.
\item[(R17)] (R) Irreducibility row: closedness of the count clauses and of
the gap statement on classes with $T_u$-size, combined with the label ties
by closedness, and the all-rows-at-a-record statement at a point beyond the
maximum of the thresholds.
\end{enumerate}

\medskip\noindent\emph{Group IV: realisations, (R2), (R3), (R10), (R11), (R24).}
Published inputs: rapid decay of cusp forms, the Casselman--Wallach
globalisation, Harish-Chandra modules, the discrete-series cover, and (for
(R24)) the classification of discrete series; all are hypotheses
(Remark~\ref{i4:A-rem-pub}).
\begin{enumerate}
\item[(R11)] (L)+(N) Bounded rows. (L) Rapid decay of cusp forms on
$[U(V)]$ (every cusp form is bounded by every negative power of the height
on a Siegel domain). (N) With $n_V\ge8fr+4$ this gives a uniform bound for
the rows of the indefinite realisation; a weak characterisation of
orthonormal frames gives the extension of $\iota$; the bounded-rows clauses
follow.
\item[(R10)] (L)+(N) Firing at realisations: existence of intertwiners and
lifts at the choices (from the four published inputs), the
conjugate-reading condition, the geometric action and injectivity of
evaluation (from the level, model and frame data of the realisation).
\item[(R3)] (N) Block lemma: for a family with the condition
$\mathrm{Ok}(1)$ (conjugation, irreducibility, lifts), the lifted projection
matrix defines a surjective linear map. Linear algebra over the coefficient
field; no published input.
\item[(R2)] (R) Primitive and type-separated rows of the realisation give
surjectivity of the row map, hence $\iota(V_0(\mu_0))$ is the whole target;
$L$ is totally complex.
\item[(R24)] (L)+(N) Theta nonvanishing at indefinite scales. (N) The
idempotents of the datum are fixed by complex conjugation: the coefficients
of the Eisenstein series attached to a monic $f$ with $f(0)\ne0$ are
conjugate to those of the normalised reverse $f(0)^{-1}\mathrm{rev}(f)$
(equivariance of the label map of Lemma~\ref{i4:A-lem-greenlab}). (L) The
classification of discrete series then gives a row placement whose finite
element matches the levels, and the extension follows.
\end{enumerate}

\medskip\noindent\emph{Group V: class field theory, (R15).} (L) Global Artin
reciprocity (hypothesis). (R) The bundle linkage implies the reciprocity
linkage (Lemma~\ref{i4:A-lem-linkR}); good cores are built and have a part
((R9)(b)); the service condition follows from the standard kernel.

\medskip\noindent\emph{Group VI: bookkeeping and numerical items, (R7)--(R9), (R18)--(R23), (R25).}
\begin{enumerate}
\item[(R7)] (N) For $0<\varepsilon<1/52$, $T_u$-size and $D_\Gamma$-size,
six bounds hold along the family, each with a constant and a threshold: the
degree count, the type-prime sizes within budget, the admissible off-type
bound, the stabiliser-log bound, the off-diagonal ramified bound and the
inert-prime loss bound. Explicit estimates of the same kind as those of
\S\ref{i4:ssec-proof-Lleaves}; we do not describe it in detail.
\item[(R8)] (N) Delivery of the arithmetic row at bundle points from the
row pins, through a statement on jets of split ambient spaces and a bound
on the jet table at split primes. We do not describe the proof below these two
statements.
\item[(R9)] (R) The reading component of a rows datum identifies the column
space with the ambient space; generators and the lattice description
transfer, which is the row pin. A built core has a part.
\item[(R18), (R19)] (R) Conjunction of eight closed properties, restriction
to $D_3(3,\mathcal A)\subseteq W_3(\mathcal A)$, maximum of constants and
thresholds; the seven duties are components of membership in the class.
\item[(R20), (R21)] (N) Divisor point summands from the third divisor rows
through the rows map; row point summands from firing, the real bound, the
frames, primitivity, the divisor summands, the class-sum pin, good
reduction and compatibility. We do not describe the proof below the statement.
\item[(R22), (R23)] (R) A general transfer lemma along selections, applied
with surjectivity onto $X$ (for the norm bound) and with
$\lambda_{\mathrm{low}}$ (for the $L_6$ bound).
\item[(R25)] (N) Smallness, with constant $0$, from the visible example
(R24) and primitivity: $4fr\le\min g$ along the family, the theta limit and
the theta degree at indefinite scale.
\end{enumerate}

\medskip\noindent\emph{Summary.} Published inputs: Siegel--Weil
(Theorem~\ref{i4:thm-sw2}), Artin reciprocity, rapid decay,
Casselman--Wallach, Harish-Chandra modules, the discrete-series cover and
classification, all re-proved in the formal development as separate
propositions; Siegel's method of local densities and
Borel's finiteness theorem, re-derived. New and resting on
\cite{CharacterFamily}: the integrality statements of Group~I, the index
and freeness statements (R1b), (R1c), the explicit unfolding and
signature-independence (R4), (R5), the loop rows and the carrier-gap bound
(R13), (R14), the bounded-rows and firing statements (R10), (R11), the block
lemma (R3), real-fixedness of the idempotents (R24), and the numerical
items (R7), (R8), (R20), (R21), (R25).

\medskip\noindent\label{i4:Ast-inputs}\emph{Structure of the proof of Proposition~\ref{i4:A-prop-inputs}(b)--(d).}

Part (a) is proved in the text modulo
Proposition~\ref{i4:A-prop-inputs-leaves}. Parts (b)--(d) rest on one core
statement, the \emph{twisted multiplicity equality}, which is proved by a
comparison of simple trace formulas with pseudo-coefficients.

\medskip\noindent\emph{The core statement.} Let $Y_0$ and $W_i$ be the two hermitian
spaces of a real datum (the first theta space and the space of index $i$),
and $\Phi:U(Y_0)(\mathbb A_f)\to U(W_i)(\mathbb A_f)$ an isomorphism of the
finite-adelic groups compatible with the lattices and level data of the
datum (a \emph{finite isometry}). For every irreducible admissible
representation $\sigma$ of $U(Y_0)(\mathbb A_f)$ that is good (its local
types at the level primes are those prescribed by the datum),
\begin{equation}\label{i4:Ast-eq-multeq}
m_{\mathrm{cusp}}(\pi_\infty^{(0)}\otimes\sigma)
=m_{\mathrm{cusp}}(\pi_\infty^{(i)}\otimes(\sigma\circ\Phi^{-1})),
\end{equation}
where $\pi_\infty^{(0)}$, $\pi_\infty^{(i)}$ are the discrete series of
$U(Y_0)(\mathbb R)$, $U(W_i)(\mathbb R)$ with the Harish-Chandra parameters
of the datum, and $m_{\mathrm{cusp}}$ is the multiplicity in the space of
cusp forms. This is the statement ``multiplicity of $\sigma$ in the
discrete-series model of the first theta space equals that of its
transport'' of (d).

\medskip\noindent\emph{Steps of the proof of \eqref{i4:Ast-eq-multeq}.}
\begin{enumerate}
\item[(M-a)] (L), re-derived. \emph{Archimedean pseudo-coefficients.} For an
odd rank $r$, signs $\varepsilon$ and an integral Harish-Chandra parameter
$\lambda$ satisfying a gap condition, and a discrete series $D$ of
$U(J_\varepsilon)$ with parameter $\lambda$, there is a compactly
supported smooth function $f$ on $U(J_\varepsilon)$ which is a
pseudo-coefficient of $D$ in the sense of the development (an isolating
property of its traces relative to $D$; we have not compared the formal definition
with the published one). This is the analogue, for these
groups, of the following theorem.
\begin{theorem}[Clozel--Delorme, Invent.\ Math.\ 77 (1984), 427--453; II, Ann.\ Sci.\ ENS (4) 23 (1990), 193--228; reference to be confirmed]\label{i4:Ast-thm-pc}
Let $G$ be a connected reductive group over $\mathbb R$ with $G(\mathbb R)$
having a compact Cartan subgroup and compact centre, and $\pi$ a discrete
series representation of $G(\mathbb R)$. There is a $K$-finite
$f\in C_c^\infty(G(\mathbb R))$ with $\mathrm{tr}\,\pi'(f)=1$ if
$\pi'\cong\pi$ and $\mathrm{tr}\,\pi'(f)=0$ for every other irreducible
tempered representation $\pi'$.
\end{theorem}
\noindent What the development proves is more explicit (the function is
constructed from orbital integrals of a profile family, with a homotopy and
a limit argument, and orbital integrals on non-elliptic classes are
treated separately) and assumes the gap condition and $r$ odd. Whether the
formal pseudo-coefficient property is implied by the published theorem we
have not determined; the step therefore counts as re-derived, not cited.
\item[(M-b)] (L) \emph{Finiteness of the cuspidal spectrum of a given type},
from rapid decay of cusp forms, Harish-Chandra's finiteness theorem and the
Dixmier--Malliavin factorisation (all hypotheses).
\item[(M-c)] (N) \emph{Type vectors.} Green's cuspidality criterion gives,
for a good $\sigma$, a vector of the prescribed type; local isometries and
an inner automorphism give an adelic frame isometry matching the two
sides; the parameters of the two discrete series match and satisfy the gap
condition.
\item[(M-d)] (L)+(N) \emph{Spectral side.} Under the hypothesis that
convolution operators by smooth compactly supported functions are of trace
class on the cuspidal spectrum (hypothesis), the trace of the test function
$f_\infty\otimes e_\sigma$ ($e_\sigma$ the type idempotent) on the cuspidal
spectrum is expressed through the cuspidal multiplicities, and the
constituents other than the one in question are separated by their traces.
The separation uses the trace on Iwahori double cosets and the separation of
exponents of twisted discrete series, which rests on Zelevinsky's
polynomial ring structure and the geometric lemma of Bernstein--Zelevinsky
(hypotheses; Theorem~\ref{i4:thm-bz}).
\item[(M-e)] (N) \emph{Geometric side.} The trace equals a geometric
expansion in which only regular semisimple elliptic classes contribute:
regular classes in a compact set are finite (from a Siegel-type finiteness
statement proved in the development), the cuspidal operators have
vanishing integrals along unipotent radicals, the parabolic and non-field
terms vanish, and an integrality statement for the integrals over regular
classes is used. These are read from the top level of the proof; we do not
describe the proofs of the individual vanishing statements. The
expansion is a simple trace formula in the sense of Arthur and of
Deligne--Kazhdan (reference to be confirmed), derived in the development for
this specific test function rather than quoted.
\item[(M-f)] (N) \emph{Comparison.} The traces of the matched test
functions on the two sides are shown to be equal, using the geometric
expansion (M-e), the finite isometry $\Phi$ and the matching of the two
archimedean parameters (M-c); equality of traces and (M-d) give
\eqref{i4:Ast-eq-multeq}. How the archimedean orbital integrals of the two
pseudo-coefficients are matched we do not describe.
\end{enumerate}
The method is the classical comparison of multiplicities between inner
forms by the simple trace formula (as in Clozel's and Harris--Taylor's
proofs of multiplicity statements for unitary groups; references to
be confirmed). The development proves (M-a), (M-e) and (M-f) from first
principles; (M-b), (M-d) use the published inputs as separate propositions, which
the development also proves.

\medskip\noindent\emph{Part (d).} (R) Eventually, at every witness point, take a size
datum from the size ledger and host frames from the host hypothesis; the
bare-stage lemma (from Green's criterion) gives a second lattice stage
$Z'$ with the same group, and the multiplicity statement transfers from
$Z'$ to $Z$ along the same-group transfer and the copy of the datum. At
$Z'$ apply \eqref{i4:Ast-eq-multeq}. The published inputs are those of
(M-b), (M-d) and Green's criterion; the elliptic and Galois inputs are not
used.

\medskip\noindent\emph{Part (b).} The multiplicity-cover hypotheses are the conjunction
of five clauses.
\begin{enumerate}
\item[(b1)] (L) the existence of discrete-series covers (hypothesis);
\item[(b2)] (L) a package of published statements on theta lifts: unramified
Howe duality, rapid decay, Harish-Chandra finiteness, Harish-Chandra modules,
the cover, Casselman--Wallach, the Weil--Siegel formula, the convergent
Siegel--Weil formula (Theorems~\ref{i4:thm-sw}, \ref{i4:thm-sw2}), the
local zeta integrals of the doubling method, unramified doubling for
$\mathrm{GL}$, Bernstein--Zelevinsky finiteness (Theorem~\ref{i4:thm-bzfin}),
Green's cuspidality criterion, existence of the Weil representation for
unitary dual pairs, and Schur orthogonality for square integrable
representations. All are hypotheses except three, which the development
proves: unramified Howe duality (the unramified clause of
Theorem~\ref{i4:thm-howe}, re-derived; Remark~\ref{i4:A-rem-howe}), the
Weil representation ((M2) below), and the doubling local zeta integrals,
which are deduced from meromorphic continuation of Siegel Eisenstein
series, a second Eisenstein input and rapid decay
(Piatetski-Shapiro--Rallis, Part~A of S.~Gelbart, I.~Piatetski-Shapiro, S.~Rallis, \emph{Explicit constructions of automorphic $L$-functions}, LNM 1254 (1987)). From the
package one derives the nonvanishing of the Hecke form through the doubling
method (the Rallis inner product formula; reference to be confirmed);
\item[(b3)] (L), re-derived: for a unitary Hecke character $\chi$ of a
number field and a finite set $T$ of places, the partial $L$-function
$L^T(s,\chi)$ is, on $\mathrm{Re}\,s>1$, the product of an entire function
(for $\chi$ not a power of the norm; for $\chi=1$, of $(s-1)^{-1}$ times an
entire function) and the finite product of removed Euler factors, and it
does not vanish on $\mathrm{Re}\,s\ge1$ (Theorem~\ref{i4:thm-hecke} and the
nonvanishing on $\mathrm{Re}\,s=1$);
\item[(b4)] (N) the multiplicity equality at pinned data, which is
\eqref{i4:Ast-eq-multeq} through the transfer of part (d);
\item[(b5)] (R) the weighted version: (b4) together with the rationality
of rows $\mathrm{RowRat}$ (hypothesis) gives the multiplicity equality and
the row test at start data.
\end{enumerate}

\medskip\noindent\emph{Part (c).} (N) An unlocalised saturated state exists, eventually
for large bases, at every start datum. Two steps: (c1) fine adelic data
extend to larger coefficient fields (from the cover, Harish-Chandra modules,
Casselman--Wallach, the Weil representation, existence of discrete series,
Harish-Chandra finiteness); (c2) the saturated bare-form theorem: from the
member hypotheses and the multiplicity-cover hypotheses, a bare form with
empty localisation set and vanishing $\eta$ exists. We do not describe the proof below these
two statements.

\medskip\noindent\emph{Summary.} Published, re-proved in the formal development as separate
propositions: the inputs listed in (b2)
except the three re-derived ones, Dixmier--Malliavin, trace class of cusp
forms, Zelevinsky's ring and the geometric lemma. Published but re-derived:
pseudo-coefficients (with a gap hypothesis), unramified Howe duality, the
Weil representation, Hecke's continuation and nonvanishing, the doubling
zeta integrals. New: the simple trace formula for the specific test
functions and the comparison (M-c), (M-e), (M-f), and the saturated state
(c).

\medskip\noindent\label{i4:Ast-inputs-leaves}\emph{Structure of the proof of Proposition~\ref{i4:A-prop-inputs-leaves}.}

\begin{enumerate}
\item[(M1)] (L), re-derived from (P1). \emph{Ray class fields.} From an Artin
map for $K$ whose kernel is the closure of $K^\times$ times the connected
component, one obtains, for every nonzero modulus $\mathfrak m$, the abelian
extension $L/K$ with $\mathrm{Gal}(L/K)$ isomorphic to the ray class group
mod $\mathfrak m$, primes $\mathfrak p\nmid\mathfrak m$ unramified and
mapped to their Frobenius: the Artin map factors through the idele class
group modulo the (open) ray subgroup. \emph{Principal ideals.}
\begin{theorem}[principal ideal theorem, Furtw\"angler, Abh.\ Math.\ Sem.\ Hamburg 7 (1929), 14--36]\label{i4:Ast-thm-pit}
Let $H$ be the Hilbert class field of a number field $K$. Every ideal of
$\mathcal O_K$ becomes principal in $\mathcal O_H$.
\end{theorem}
\noindent The development proves the form it needs (for totally complex
$L$ and the ray class field of conductor one of $L$) by the classical
route: the Galois group of $H_2/K$, $H_2$ the Hilbert class field of $H$,
has commutator subgroup $\mathrm{Gal}(H_2/H)$; the extension of ideals
corresponds under Artin reciprocity to the transfer into the commutator
subgroup; and the transfer of a group into its commutator subgroup is
trivial (the group-theoretic theorem of Artin--Furtwangler, re-derived).
\item[(M2)] (L), re-derived. \emph{Weil representation.} For a unitary dual
pair over a CM extension, an additive character $\psi$ and splitting
characters $\chi_S,\chi_T$ with the restriction condition, a Weil
representation of $U(S)\times U(T)$ on the Schwartz space exists
(Kudla, Israel J.\ Math.\ 87 (1994); Harris--Kudla--Sweet, JAMS 9 (1996);
reference to be confirmed). The development constructs it from the oriented
coordinates of the Schrodinger model, first for powers of an unramified
character $\tilde\eta$ and then for arbitrary splitting characters by
twisting.
\item[(M3)] (L)+(R) From the classification of discrete series
(hypothesis): existence, exhaustion, uniqueness, the infinitesimal
character, multiplicity one of the lowest $K$-type, the $K$-types above the
lowest one, and strictness of the lowest $K$-type. Harish-Chandra's
classification and the lowest-$K$-type theorems of Schmid and
Hecht--Schmid (reference to be confirmed). The
deduction is a case analysis of the hypothesis.
\item[(M4)] (R) The three statements for every Harish-Chandra parameter are
the components of (P9), the existence in (M3) and uniqueness.
\item[(M5)] (L), re-derived. Harish-Chandra's finiteness theorem
(Harish-Chandra, LNM 62 (1968), Thm.~1; reference to be confirmed): the space of automorphic forms of given $K$-type
and $Z(\mathfrak g)$-eigencharacter is finite dimensional. The formal
statement is about the isotypic piece of $L^2$ of the adelic quotient cut
out by a convolution operator; its proof passes from adelic to classical
quotients by congruence subgroups, uses a fundamental domain with cusps and
the reduction to $K$-finite vectors, and bounds dimensions by linear forms
evaluated on a discrete-series bound. This passage, not the finiteness
theorem itself, is what the development adds.
\item[(M6)] (L)+(N). The Siegel mass ratio for CM fields. The development
derives it from a mass ratio for model lattices, obtained by Siegel's
inductive method: a unimodular binary hermitian lattice represents the
required values, and a local-index statement computes the ratio of masses
in ranks $M$ and $M-1$ as a product of local densities. The published
counterpart is the mass formula for unimodular hermitian lattices
(Shimura; Gan--Hanke--Yu; reference to be confirmed);
we have not compared the normalisation, so the explicit product formula
stated in Proposition~\ref{i4:A-prop-inputs}(a) is counted as new.
\end{enumerate}

\medskip\noindent\label{i4:Ast-stages}\emph{Structure of the proof of Proposition~\ref{i4:A-prop-stages}.}

The member statement is proved in three steps; the published inputs enter
in two groups, the automorphic group of Remark~\ref{i4:A-rem-pub} and the
\emph{elliptic hypotheses}, which are Galois-cohomological: Poitou--Tate
duality for totally complex fields, cohomological dimension two, tame
local Galois groups, local finiteness, local Tate duality, the local and
global Euler characteristic formulas, global finiteness, orthogonality of
unramified classes, the Greenberg--Wiles formula, Artin's theorem on ray
class characters, Chebotarev's theorem for degree-one primes, Fontaine's
theorem that de Rham representations are potentially semistable (Berger,
Invent.\ Math.\ 148 (2002), 219--284), finiteness of automorphic
forms on $\mathrm{GL}$, together with internal statements on Hecke
integrality, rational ray levels, twisted series and the Hecke reading. All
of these are hypotheses of the proposition.

\medskip\noindent\emph{Step 1: the stage statements.} A \emph{state} is a point of a
start datum; its levels form a chain (same-Galois level, intermediate level,
three successive intermediate levels, strictly henselian acyclic level).
\begin{enumerate}
\item[(s$_{Mt}$)] (N) Every state has a multiplicity datum: the Gram matrix
of every part is the standard one and the standard lattice is diagonal, so
a standard rational frame exists. No hypothesis.
\item[(isot)] (N) Internal isotypy at states, from the Weil representation,
existence of discrete series, the multiplicity-cover hypotheses (through
\eqref{i4:Ast-eq-multeq}) and the third-layer inputs: the idempotent of the
Green character is stable under inversion, and the multiplicity equality
gives the isotypic-rank condition.
\item[(s$_{Rat}$)] (N) Rationality data: coefficient fields and finite
generation of forms, a nonzero character of the host groups, the Hecke data
(from rapid decay and the Hecke hypotheses), integrality of the Hecke
$L$-values, and eigenvalues with the eigenvalue law from isotypy.
\item[(par), (s$_{Cusp}$)] (L)+(N) The parameter statement: the row of a
same-Galois level is a good cuspidal parameter, from the lift and
local-component statements and the base-change identities. Cuspidality of
rows: the Galois polynomial lifts almost everywhere by Frobenius splitting
and, after twisting by the local Langlands datum, everywhere (Chebotarev,
Galois and finiteness inputs). The published content is local-global
compatibility at unramified places (Satake parameters versus Frobenius
eigenvalues), which here is part of the hypotheses; the bookkeeping of
twists is new.
\item[(s$_A$)] (L)+(N) Same-Galois replacement: the types-true theorem for
ray classes (from the local Langlands datum and the Berger, Green, Artin,
ray class and Chebotarev inputs) gives a layer with a Galois
compatibility; continuity of Weil--Deligne representations (part of the
archimedean comparison of the local Langlands datum) gives the same-Galois
level.
\item[(levels), (s$_C$), (fin)] (N) The three intermediate levels with the
list statement; the height-counting character datum, the theta data and
the $\eta$-statement (Casselman--Wallach and Harish-Chandra modules enter
only here); the finiteness statements for $N_{\mathrm{law}}$ and the
finiteness list statement. We do not describe the proof below the statement.
\end{enumerate}

\medskip\noindent\emph{Step 2: three-step data.} (N) The fine adelic datum extends to
larger fields (as in part (c) of \S\ref{i4:Ast-inputs}); the weight-case
condition follows from the multiplicity-cover hypotheses; the finiteness
statements give the second pass; with the stage statements this gives a
three-step datum at every start point, the level-three conditions and the
position condition.

\medskip\noindent\emph{Step 3: export and closing.}
\begin{enumerate}
\item[(E)] (L)+(N) \emph{Export from automorphy.} For every rows datum with
a column pin and a norm-choice core with the extension law there is a twist
norm choice; the series realisation from the published inputs, applied to
it and to the cuspidal Galois datum built from it, gives exportability.
This is the only step that uses the elliptic hypotheses (Poitou--Tate
duality, the Euler characteristic formulas, the Greenberg--Wiles formula
and the others listed above; reference to be confirmed). How the cuspidal
Galois datum is built from them we describe only at the level of this statement.
\item[(C)] (N)+(R) \emph{Closing.} From the ray class field, principal
ideal and CM-curve statements and the Siegel mass ratio: the base
condition, the passage from three-step to four-step data, the ambient,
division and theta existence statements, export, the second pass and the
goodness of cores built for the second system give the member statement.
\item[(ExBL)] (N) Uniformity of the geometric parameter: from Green's
criterion, the cover, Harish-Chandra modules, Casselman--Wallach, the Weil
representation, discrete series, Harish-Chandra finiteness, rapid decay,
Artin reciprocity and Weber's ray class field theorem, together with the
multiplicity condition for $p_0$ (part (d) of
Proposition~\ref{i4:A-prop-inputs}), $\mathrm{ExBL}$ holds for the built
and good slots.
\item[(A)] (R) Transfer along an evaluation property $A$ with
$\mathrm{EvLgG}(A)$: $A$ depends only on $(K,B,f,T)$ and holds eventually for
large bases, so it survives the enlargement of the base.
\end{enumerate}

\medskip\noindent\emph{Summary.} Every published input is a hypothesis; none is
re-derived in this part. The new content is the construction of states,
levels and three-step data (Steps 1--2) and the export and closing
(Step 3); the twisted multiplicity equality enters through the
multiplicity-cover hypotheses. We do not describe the proofs below the level described.

\medskip\noindent\label{i4:Ast-open}\emph{Structure of the proof of Proposition~\ref{i4:A-prop-open}.}

The bundle-level leaves are the twelve open parameter statements of
Theorem~\ref{i4:thm-assembly}(c), proved at $p=p_0$ from the third-place
hypotheses. The formal proof derives them in the following order.
\begin{enumerate}
\item[(O1), (O2)] (R)+(N) \emph{Tie.} The identities (T1)--(T5) are read off
the bundle witness (zero section, coefficient field datum, place map,
conjugate-character identities at the theta levels, level-character pin);
the host hypothesis holds at the stage of the witness point and with
(T2)--(T4) gives normalisation pins. Although the statement lists seven
published inputs, the proof of the tie uses none of them. We do not describe the proof
below the witness identities.
\item[(prim), (dec)] (N) Primitivity and the decomposition-point condition
are linked at every parameter system; no published input. They are proved
from the companion row conditions (O8).
\item[(O8)] (N) \emph{Row conditions at the companion}, for every twisted
input: the $\Lambda$-rows imply the level-three $\Lambda'$-rows; the
decomposition-point rows and condition hold; the level-three divisor rows
imply the divisor condition. A direct verification on the division pair
of the companion.
\item[(O3)] (L)+(N) \emph{Multiplicity row $\mathrm{MI}^*_2$, divisor and
definite base.} Published inputs: Harish-Chandra finiteness, rapid decay,
the cover, Harish-Chandra modules, Casselman--Wallach, the Weil
representation, the Weil--Siegel formula, the convergent Siegel--Weil
formula in both forms, the doubling method, unramified doubling for
$\mathrm{GL}$, the classification of discrete series, and strong
approximation:
\begin{theorem}[strong approximation; Kneser, Platonov--Rapinchuk, \emph{Algebraic Groups and Number Theory} (1994), \S7.4, Thm.~7.12; reference to be confirmed]\label{i4:Ast-thm-sa}
Let $G$ be a connected, simply connected, absolutely almost simple group
over a number field $F$ and $S$ a finite set of places with
$\prod_{v\in S}G(F_v)$ noncompact. Then $G(F)$ is dense in
$G(\mathbb A_F^S)$.
\end{theorem}
\noindent It enters for special unitary groups; we have not compared the exact form of the
hypothesis in the development with this statement. Unramified Howe duality is used in the re-derived form.
(N): the second form of Siegel--Weil gives the diagonal row; at every
witness point the row slot is populated; this gives $\mathrm{MI}^*_2$. Given
$\mathrm{MI}^*_2$ and the boundary law, the divisor rows follow, and the
pairings on the definite base and the division-row pairings vanish.
\item[(O4)] (L)+(N) \emph{$\Lambda$- and $\kappa$-rows.} At points with a
parameter choice: the isotypic rank condition at the witness point; finite
types through the full Hecke eigensystems; the first part of (O8) turns
finite types into the $\Lambda'$-rows. The $\kappa$-rows follow from the
$\Lambda'$-rows by \eqref{i4:Ast-eq-multeq} (the twisted multiplicity
equality, \S\ref{i4:Ast-inputs}) and the identification of the divisor
embedding constant with the rows constant at the infinite place. Published
inputs: those of (M-b), (M-d) and Green's criterion.
\item[(O5)] (L)+(N) \emph{Goodness and labels.} Green's theorem on regular
characters (Theorem~\ref{i4:A-thm-green}) gives
the bijection of Lemma~\ref{i4:A-lem-greenlab}; with Zelevinsky's
polynomial ring (hypothesis) every built core has good finite types
(eventually, for large bases). The labels: the generator pin of the
extension law and the label tie of the rows give a Hida setting with good
operator, isometry and the three ties.
\item[(O6)] (L)+(N) \emph{Weak geometric compatibility}, from the ray class
field statement (Proposition~\ref{i4:A-prop-inputs-leaves}(M1)), the
extension law, the divisor row and the tie: for every comparison datum the
point summands of the corresponding part satisfy the point summand
condition, the class-sum pin, compatibility and the condition
$\mathrm{More}$.
\item[(O7)] (N) The defect row on the point-evaluation side: real tie with
symmetric tied column forms of finite step-zero character, compatible
column characters, and kernel weights with the kernel ties and sizes. No
published input.
\item[(O9)] (L)+(N) \emph{Join of the linked realisation.} From the Weil
representation, discrete series and their cover, Harish-Chandra modules,
Casselman--Wallach, Harish-Chandra finiteness and rapid decay: if
primitivity, the decomposition-point condition and the tie are linked, so
is the selection-jump condition. The proof builds the join of the linked
realisation with the kernel of the comparison datum, transfers the row
placement of the spectral stage, and constructs the lattice of the join
subspace. We do not describe the proof below these three constructions.
\end{enumerate}

\medskip\noindent\emph{Summary.} Published, re-proved in the formal development as separate
propositions: the automorphic inputs of
Remark~\ref{i4:A-rem-pub}, the Siegel--Weil formulas, the doubling method,
strong approximation, the ray class field statement, Green's theorem.
New and resting on \cite{CharacterFamily}: the witness identities of the tie,
the companion row conditions (O8), the deduction of $\mathrm{MI}^*_2$ and the
divisor rows, the $\Lambda'$-rows and the passage to the $\kappa$-rows, the
geometric and defect rows, and the join (O9).

\medskip\noindent\textbf{Status.}
\emph{Proved here:} Theorems~\ref{i4:thm-top} and~\ref{i4:thm-final}, hence Theorem~\ref{M:thm:alpoge-families}, from the four
inputs (Propositions~\ref{i4:hyp-loss}, \ref{i4:hyp-div}, \ref{i4:hyp-S}, \ref{i4:hyp-selmer}) and the statements listed in
Remark~\ref{i4:rem-trust}; the logical part of the existence argument (thresholds, closed properties, connectedness from
design existence, design existence); the root-discriminant estimate, the kit parameters and the host supply modulo the
formally verified host constructions; the twist bounds and the window data; the closed estimates on the division side
modulo Propositions~\ref{i4:Lprop-sel}--\ref{i4:Lprop-fam}; realisation and its size leaves; the reduction of the automorphic
construction to pointwise and bundle-level statements, with the row calculus and the assembly of (A1)--(A4); and the
short-interval bound for norm primes (Proposition~\ref{i4:thm-shortnorm}) with the ray class family and its leaves.
\emph{Published theorems, quoted:} those listed in Remark~\ref{i4:rem-trust}(ii).
\emph{Formally verified, \cite{CharacterFamily}:} those listed in Remark~\ref{i4:rem-trust}(i); their formally verified steps
are outlined in \S\ref{i4:ssec-outline}.  The formal development proves, as separate propositions, the published statements
on which these rest, and its final theorem has no hypotheses.











\begin{thebibliography}{9}
\bibitem{Main} D.~A.~Trejo Pizzo, \emph{The sign of the radial intersection form on the idele class group}.
\bibitem{CharacterFamily} L.~Alpoge, character-family argument for
Dirichlet $L$-functions, author's source archive accompanying the
November 2026 remark, \url{https://alpo.ge/lean.tar.gz.part01}
(first of fourteen parts).
\end{thebibliography}
\end{document}
